% Auto-generated from data/segments.json + layout.json (+ transforms) — do not edit by hand.
% volume: 32Abhi04
% mode: reading
% segments: 3526
% transform_rules: 0
% toc_marks: 520

% Volume layout defaults (layout.json ``layout``).
\csromanlayoutapply{1}{1.5}{21.6pt}{5pt}{6.3pt}{65pt}{2.5em}%

% Reading physical-page layout (layout.json ``page_layout_reading_mode``).
\csromanreadingpagelayoutvolume{1}{1.5}{21.6pt}{5pt}{6.3pt}{65pt}{2.5em}%
\csromanreadingpagelayoutenable

\pitaka{\textbf{อภิธมฺมปิฏก}}
\csromanmatikaanchor{mtk.0}
\csromantocmarknopage{chapter}{กถาวตฺถุปาฬิ}{mtk.0}
\bookmark[dest={mtk.0},level=0]{กถาวตฺถุปาฬิ}
\gambhira{\textbf{กถาวตฺถุปาฬิ}}
\namakkaram{นโม ตสฺส ภควโต อรหโต สมฺมาสมฺพุทฺธสฺส.}
\csromanmatikaanchor{mtk.1}
\csromanmatikaanchor{mtk.2}
\csromantocmarknopage{chapter}{1. ปุคฺคลกถา}{mtk.1}
\bookmark[dest={mtk.1},level=0]{1. ปุคฺคลกถา}
\csromantocmarkref{section}{1. สุทฺธสจฺจิกฏฺฐ}{mtk.2}
\bookmark[dest={mtk.2},level=1]{1. สุทฺธสจฺจิกฏฺฐ}
\csromanheader{1}{1. สุทฺธสจฺจิกฏฺฐ}
\proseitem{1}{\csromanfolio{1}\csromansymbolfootnote{*}{อิมินา ลกฺขเณน สกวาทีปุจฺฉา ทสฺสิตา.}ปุคฺคโล อุปลพฺภติ สจฺจิกฏฺฐปรมตฺเถนาติ\footnote{สจฺจิกฏฺฐปรมฏฺเฐนาติ (สฺยา, อิ, ก-สี)}, อามนฺตา.\csromanspacer{} โย สจฺจิกฏฺโฐ ปรมตฺโถ, ตโต โส ปุคฺคโล อุปลพฺภติ สจฺจิกฏฺฐปรมตฺเถนาติ, น เหวํ วตฺตพฺเพ.}

\prose{อาชานาหิ นิคฺคหํ\csromandash{}หญฺจิ ปุคฺคโล อุปลพฺภติ สจฺจิกฏฺฐปรมตฺเถน, เตน วต เร วตฺตพฺเพ “โย สจฺจิกฏฺโฐ ปรมตฺโถ, ตโต โส ปุคฺคโล อุปลพฺภติ สจฺจิกฏฺฐปรมตฺเถนา”ติ.\csromanspacer{} ยํ ตตฺถ วเทสิ “วตฺตพฺเพ โข ‘ปุคฺคโล อุปลพฺภติ สจฺจิกฏฺฐปรมตฺเถน’, โน จ วตฺตพฺเพ ‘โย สจฺจิกฏฺโฐ ปรมตฺโถ, ตโต โส ปุคฺคโล อุปลพฺภติ}

\prose{โน เจ ปน วตฺตพฺเพ “โย สจฺจิกฏฺโฐ ปรมตฺโถ, ตโต โส ปุคฺคโล อุปลพฺภติ สจฺจิกฏฺฐปรมตฺเถนา”ติ, โน จ วต เร\footnote{โน วต เร (สฺยา, อิ)} วตฺตพฺเพ \csromanfolio{2}“ปุคฺคโล อุปลพฺภติ สจฺจิกฏฺฐปรมตฺเถนา”ติ.\csromanspacer{} ยํ ตตฺถ วเทสิ “วตฺตพฺเพ โข ‘ปุคฺคโล อุปลพฺภติ สจฺจิกฏฺฐปรมตฺเถน’, โน จ วตฺตพฺเพ ‘โย สจฺจิกฏฺโฐ ปรมตฺโถ, ตโต โส ปุคฺคโล อุปลพฺภติ อนุโลมปญฺจกํ.}

\proseitem{2}{\csromansymbolfootnote{+}{อิมินา ลกฺขเณน ปรวาทีปุจฺฉา ทสฺสิตา.}ปุคฺคโล นุปลพฺภติ สจฺจิกฏฺฐปรมตฺเถนาติ, อามนฺตา.\csromanspacer{} โย สจฺจิกฏฺโฐ ปรมตฺโถ, ตโต โส ปุคฺคโล นุปลพฺภติ สจฺจิกฏฺฐปรมตฺเถนาติ, น เหวํ วตฺตพฺเพ.}

\prose{อาชานาหิ ปฏิกมฺมํ\csromandash{}หญฺจิ ปุคฺคโล นุปลพฺภติ สจฺจิกฏฺฐปรมตฺเถน, เตน วต เร วตฺตพฺเพ “โย สจฺจิกฏฺโฐ ปรมตฺโถ, ตโต โส ปุคฺคโล นุปลพฺภติ สจฺจิกฏฺฐปรมตฺเถนา”ติ.\csromanspacer{} ยํ ตตฺถ วเทสิ “วตฺตพฺเพ โข ‘ปุคฺคโล นุปลพฺภติ สจฺจิกฏฺฐปรมตฺเถน’, โน จ วตฺตพฺเพ ‘โย สจฺจิกฏฺโฐ ปรมตฺโถ, ตโต โส ปุคฺคโล นุปลพฺภติ}

\prose{โน เจ ปน วตฺตพฺเพ “โย สจฺจิกฏฺโฐ ปรมตฺโถ, ตโต โส ปุคฺคโล นุปลพฺภติ สจฺจิกฏฺฐปรมตฺเถนา”ติ, โน จ วต เร วตฺตพฺเพ “ปุคฺคโล นุปลพฺภติ สจฺจิกฏฺฐปรมตฺเถนา”ติ.\csromanspacer{} ยํ ตตฺถ วเทสิ “วตฺตพฺเพ โข ‘ปุคฺคโล นุปลพฺภติ สจฺจิกฏฺฐปรมตฺเถน’, โน จ วตฺตพฺเพ ‘โย สจฺจิกฏฺโฐ ปรมตฺโถ, ตโต โส ปุคฺคโล นุปลพฺภติ ปฏิกมฺมจตุกฺกํ.}

\proseitem{3}{ตฺวํ เจ ปน มญฺญสิ “วตฺตพฺเพ โข ‘ปุคฺคโล นุปลพฺภติ สจฺจิกฏฺฐปรมตฺเถน’, โน จ วตฺตพฺเพ ‘โย สจฺจิกฏฺโฐ ปรมตฺโถ, ตโต โส ปุคฺคโล นุปลพฺภติ สจฺจิกฏฺฐปรมตฺเถนา’ติ”.\csromanspacer{} เตน ตว\footnote{ตฺวํ (สฺยา) ฏีกา โอโลเกตพฺพา.} ตตฺถ เหตาย ปฏิญฺญาย เหวํ ปฏิชานนฺตํ\footnote{ฏีกา โอโลเกตพฺพา.} เหวํ นิคฺคเหตพฺเพ, อถ ตํ นิคฺคณฺหาม, สุนิคฺคหิโต จ\footnote{สุนิคฺคหิโตว (สฺยา)} โหสิ\csromandash{}}

\prose{\csromanfolio{3}หญฺจิ ปุคฺคโล นุปลพฺภติ สจฺจิกฏฺฐปรมตฺเถน, เตน วต เร วตฺตพฺเพ “โย สจฺจิกฏฺโฐ ปรมตฺโถ, ตโต โส ปุคฺคโล นุปลพฺภติ สจฺจิกฏฺฐปรมตฺเถนา”ติ.\csromanspacer{} ยํ ตตฺถ วเทสิ “วตฺตพฺเพ โข ‘ปุคฺคโล นุปลพฺภติ สจฺจิกฏฺฐปรมตฺเถน’, โน จ วตฺตพฺเพ ‘โย สจฺจิกฏฺโฐ ปรมตฺโถ, ตโต โส ปุคฺคโล นุปลพฺภติ สจฺจิกฏฺฐปรมตฺเถนา’ติ”,}

\prose{โน เจ ปน วตฺตพฺเพ “โย สจฺจิกฏฺโฐ ปรมตฺโถ, ตโต โส ปุคฺคโล นุปลพฺภติ สจฺจิกฏฺฐปรมตฺเถนา”ติ, โน จ วต เร วตฺตพฺเพ “ปุคฺคโล นุปลพฺภติ สจฺจิกฏฺฐปรมตฺเถนา”ติ.\csromanspacer{} ยํ ตตฺถ วเทสิ “วตฺตพฺเพ โข ‘ปุคฺคโล นุปลพฺภติ สจฺจิกฏฺฐปรมตฺเถน’, โน จ วตฺตพฺเพ ‘โย สจฺจิกฏฺโฐ ปรมตฺโถ, ตโต โส ปุคฺคโล นุปลพฺภติ สจฺจิกฏฺฐปรมตฺเถนา’ติ”, อิทํ เต มิจฺฉา.\csromanspacer{} นิคฺคหจตุกฺกํ.}

\proseitem{4}{เอเส เจ ทุนฺนิคฺคหิเต เหวเมวํ\footnote{เหวเมว (สฺยา)} ตตฺถ ทกฺข.\csromanspacer{} วตฺตพฺเพ โข “ปุคฺคโล อุปลพฺภติ สจฺจิกฏฺฐปรมตฺเถน”, โน จ วตฺตพฺเพ “โย สจฺจิกฏฺโฐ ปรมตฺโถ, ตโต โส ปุคฺคโล อุปลพฺภติ สจฺจิกฏฺฐปรมตฺเถนา”ติ’.\csromanspacer{} โน จ มยํ ตยา ตตฺถ เหตาย ปฏิญฺญาย เหวํ ปฏิชานนฺตา เหวํ นิคฺคเหตพฺพา, อถ มํ นิคฺคณฺหาสิ, ทุนฺนิคฺคหิตา จ\footnote{ทุนฺนิคฺคหิตาว (สฺยา)} โหม\csromandash{}}

\prose{หญฺจิ ปุคฺคโล อุปลพฺภติ สจฺจิกฏฺฐปรมตฺเถน, เตน วต เร วตฺตพฺเพ “โย สจฺจิกฏฺโฐ ปรมตฺโถ, ตโต โส ปุคฺคโล อุปลพฺภติ สจฺจิกฏฺฐปรมตฺเถนา”ติ.\csromanspacer{} ยํ ตตฺถ วเทสิ “วตฺตพฺเพ โข ‘ปุคฺคโล อุปลพฺภติ สจฺจิกฏฺฐปรมตฺเถน’, โน จ วตฺตพฺเพ ‘โย สจฺจิกฏฺโฐ ปรมตฺโถ, ตโต โส ปุคฺคโล อุปลพฺภติ สจฺจิกฏฺฐปรมตฺเถนา’ติ”,}

\prose{โน เจ ปน วตฺตพฺเพ “โย สจฺจิกฏฺโฐ ปรมตฺโถ, ตโต โส ปุคฺคโล อุปลพฺภติ สจฺจิกฏฺฐปรมตฺเถนา”ติ, โน จ วต เร วตฺตพฺเพ “ปุคฺคโล อุปลพฺภติ สจฺจิกฏฺฐปรมตฺเถนา”ติ.\csromanspacer{} ยํ ตตฺถ วเทสิ “วตฺตพฺเพ โข ‘ปุคฺคโล อุปลพฺภติ สจฺจิกฏฺฐปรมตฺเถน’, โน จ วตฺตพฺเพ ‘โย \csromanfolio{4}สจฺจิกฏฺโฐ ปรมตฺโถ, ตโต โส ปุคฺคโล อุปลพฺภติ สจฺจิกฏฺฐปรมตฺเถนา’ติ”, อิทํ เต มิจฺฉา.\csromanspacer{} อุปนยนจตุกฺกํ.}

\proseitem{5}{น เหวํ นิคฺคเหตพฺเพ, เตน หิ ยํ นิคฺคณฺหาสิ\csromandash{}หญฺจิ ปุคฺคโล อุปลพฺภติ สจฺจิกฏฺฐปรมตฺเถน, เตน วต เร วตฺตพฺเพ “โย สจฺจิกฏฺโฐ ปรมตฺโถ, ตโต โส ปุคฺคโล อุปลพฺภติ สจฺจิกฏฺฐปรมตฺเถนา”ติ.\csromanspacer{} ยํ ตตฺถ วเทสิ “วตฺตพฺเพ โข ‘ปุคฺคโล อุปลพฺภติ สจฺจิกฏฺฐปรมตฺเถน’, โน จ วตฺตพฺเพ ‘โย สจฺจิกฏฺโฐ ปรมตฺโถ, ตโต โส ปุคฺคโล อุปลพฺภติ}

\prose{โน เจ ปน วตฺตพฺเพ “โย สจฺจิกฏฺโฐ ปรมตฺโถ, ตโต โส ปุคฺคโล อุปลพฺภติ สจฺจิกฏฺฐปรมตฺเถนา”ติ, โน จ วต เร วตฺตพฺเพ “ปุคฺคโล อุปลพฺภติ สจฺจิกฏฺฐปรมตฺเถนา”ติ.\csromanspacer{} ยํ ตตฺถ วเทสิ “วตฺตพฺเพ โข ‘ปุคฺคโล อุปลพฺภติ สจฺจิกฏฺฐปรมตฺเถน, โน จ วตฺตพฺเพ ‘โย สจฺจิกฏฺโฐ ปรมตฺโถ, ตโต โส ปุคฺคโล อุปลพฺภติ สจฺจิกฏฺฐปรมตฺเถนา’ติ”, อิทํ เต มิจฺฉา.\csromanspacer{} เตน หิ เย กเต นิคฺคเห, เส นิคฺคเห ทุกฺกเฏ.\csromanspacer{} สุกเต ปฏิกมฺเม.\csromanspacer{} สุกตา ปฏิปาทนาติ.}

\vspace{\tipitakaparskip}
\csromancenter{นิคฺคมนจตุกฺกํ.}
\csromancenter{ปฐโม นิคฺคโห.}
\csromanheader{2}{1. สุทฺธสจฺจิกฏฺฐ}
\proseitem{6}{\csromansymbolmark{+}ปุคฺคโล นุปลพฺภติ สจฺจิกฏฺฐปรมตฺเถนาติ, อามนฺตา.\csromanspacer{} โย สจฺจิกฏฺโฐ ปรมตฺโถ, ตโต โส ปุคฺคโล นุปลพฺภติ สจฺจิกฏฺฐปรมตฺเถนาติ, น เหวํ วตฺตพฺเพ.}

\prose{อาชานาหิ นิคฺคหํ\csromandash{}หญฺจิ ปุคฺคโล นุปลพฺภติ สจฺจิกฏฺฐปรมตฺเถน, เตน วต เร วตฺตพฺเพ “โย สจฺจิกฏฺโฐ ปรมตฺโถ, ตโต โส ปุคฺคโล นุปลพฺภติ สจฺจิกฏฺฐปรมตฺเถนา”ติ.\csromanspacer{} ยํ ตตฺถ วเทสิ “วตฺตพฺเพ โข \csromanfolio{5}‘ปุคฺคโล นุปลพฺภติ สจฺจิกฏฺฐปรมตฺเถน’, โน จ วตฺตพฺเพ ‘โย สจฺจิกฏฺโฐ ปรมตฺโถ, ตโต โส ปุคฺคโล นุปลพฺภติ}

\prose{โน เจ ปน วตฺตพฺเพ “โย สจฺจิกฏฺโฐ ปรมตฺโถ, ตโต โส ปุคฺคโล นุปลพฺภติ สจฺจิกฏฺฐปรมตฺเถนา”ติ, โน จ วต เร วตฺตพฺเพ “ปุคฺคโล นุปลพฺภติ สจฺจิกฏฺฐปรมตฺเถนา”ติ.\csromanspacer{} ยํ ตตฺถ วเทสิ “วตฺตพฺเพ โข ‘ปุคฺคโล นุปลพฺภติ สจฺจิกฏฺฐปรมตฺเถน’, โน จ วตฺตพฺเพ ‘โย สจฺจิกฏฺโฐ ปรมตฺโถ, ตโต โส ปุคฺคโล นุปลพฺภติ ปจฺจนีกปญฺจกํ.}

\proseitem{7}{\csromansymbolmark{*}ปุคฺคโล อุปลพฺภติ สจฺจิกฏฺฐปรมตฺเถนาติ, อามนฺตา.\csromanspacer{} โย สจฺจิกฏฺโฐ ปรมตฺโถ, ตโต โส ปุคฺคโล อุปลพฺภติ สจฺจิกฏฺฐปรมตฺเถนาติ, น เหวํ วตฺตพฺเพ.}

\prose{อาชานาหิ ปฏิกมฺมํ\csromandash{}หญฺจิ ปุคฺคโล อุปลพฺภติ สจฺจิกฏฺฐปรมตฺเถน, เตน วต เร วตฺตพฺเพ “โย สจฺจิกฏฺโฐ ปรมตฺโถ, ตโต โส ปุคฺคโล อุปลพฺภติ สจฺจิกฏฺฐปรมตฺเถนา”ติ.\csromanspacer{} ยํ ตตฺถ วเทสิ “วตฺตพฺเพ โข ‘ปุคฺคโล อุปลพฺภติ สจฺจิกฏฺฐปรมตฺเถน’, โน จ วตฺตพฺเพ ‘โย สจฺจิกฏฺโฐ ปรมตฺโถ, ตโต โส ปุคฺคโล อุปลพฺภติ}

\prose{โน เจ ปน วตฺตพฺเพ “โย สจฺจิกฏฺโฐ ปรมตฺโถ, ตโต โส ปุคฺคโล อุปลพฺภติ สจฺจิกฏฺฐปรมตฺเถนา”ติ, โน จ วต เร วตฺตพฺเพ “ปุคฺคโล อุปลพฺภติ สจฺจิกฏฺฐปรมตฺเถนา”ติ.\csromanspacer{} ยํ ตตฺถ วเทสิ “วตฺตพฺเพ โข ‘ปุคฺคโล อุปลพฺภติ สจฺจิกฏฺฐปรมตฺเถน’, โน จ วตฺตพฺเพ ‘โย สจฺจิกฏฺโฐ ปรมตฺโถ, ตโต โส ปุคฺคโล อุปลพฺภติ ปฏิกมฺมจตุกฺกํ.}

\proseitem{8}{ตฺวํ เจ ปน มญฺญสิ “วตฺตพฺเพ โข ‘ปุคฺคโล อุปลพฺภติ สจฺจิกฏฺฐปรมตฺเถน’, โน จ วตฺตพฺเพ ‘โย สจฺจิกฏฺโฐ ปรมตฺโถ, ตโต โส \csromanfolio{6}ปุคฺคโล อุปลพฺภติ สจฺจิกฏฺฐปรมตฺเถนา’ติ”.\csromanspacer{} เตน ตว ตตฺถ เหตาย ปฏิญฺญาย เหวํ ปฏิชานนฺตํ เหวํ นิคฺคเหตพฺเพ, อถ ตํ นิคฺคณฺหาม, สุนิคฺคหิโต จ โหสิ\csromandash{}}

\prose{หญฺจิ ปุคฺคโล อุปลพฺภติ สจฺจิกฏฺฐปรมตฺเถน, เตน วต เร วตฺตพฺเพ “โย สจฺจิกฏฺโฐ ปรมตฺโถ, ตโต โส ปุคฺคโล อุปลพฺภติ สจฺจิกฏฺฐปรมตฺเถนา”ติ.\csromanspacer{} ยํ ตตฺถ วเทสิ “วตฺตพฺเพ โข ‘ปุคฺคโล อุปลพฺภติ สจฺจิกฏฺฐปรมตฺเถน’, โน จ วตฺตพฺเพ ‘โย สจฺจิกฏฺโฐ ปรมตฺโถ, ตโต โส ปุคฺคโล อุปลพฺภติ สจฺจิกฏฺฐปรมตฺเถนา’ติ”,}

\prose{โน เจ ปน วตฺตพฺเพ “โย สจฺจิกฏฺโฐ ปรมตฺโถ, ตโต โส ปุคฺคโล อุปลพฺภติ สจฺจิกฏฺฐปรมตฺเถนา”ติ, โน จ วต เร วตฺตพฺเพ “ปุคฺคโล อุปลพฺภติ สจฺจิกฏฺฐปรมตฺเถนา”ติ.\csromanspacer{} ยํ ตตฺถ วเทสิ “วตฺตพฺเพ โข ‘ปุคฺคโล อุปลพฺภติ สจฺจิกฏฺฐปรมตฺเถน’, โน จ วตฺตพฺเพ ‘โย สจฺจิกฏฺโฐ ปรมตฺโถ, ตโต โส ปุคฺคโล อุปลพฺภติ สจฺจิกฏฺฐปรมตฺเถนา’ติ”, อิทํ เต มิจฺฉา.\csromanspacer{} นิคฺคหจตุกฺกํ.}

\proseitem{9}{เอเส เจ ทุนฺนิคฺคหิเต เหวเมวํ ตตฺถ ทกฺข.\csromanspacer{} วตฺตพฺเพ โข “ปุคฺคโล นุปลพฺภติ สจฺจิกฏฺฐปรมตฺเถน”, โน จ วตฺตพฺเพ “โย สจฺจิกฏฺโฐ ปรมตฺโถ, ตโต โส ปุคฺคโล นุปลพฺภติ สจฺจิกฏฺฐปรมตฺเถนา”ติ.\csromanspacer{} โน จ มยํ ตยา ตตฺถ เหตาย ปฏิญฺญาย เหวํ ปฏิชานนฺตา เหวํ นิคฺคเหตพฺพา, อถ มํ นิคฺคณฺหาสิ, ทุนฺนิคฺคหิตา จ โหม\csromandash{}}

\prose{หญฺจิ ปุคฺคโล นุปลพฺภติ สจฺจิกฏฺฐปรมตฺเถน, เตน วต เร วตฺตพฺเพ “โย สจฺจิกฏฺโฐ ปรมตฺโถ, ตโต โส ปุคฺคโล นุปลพฺภติ สจฺจิกฏฺฐปรมตฺเถนา”ติ.\csromanspacer{} ยํ ตตฺถ วเทสิ “วตฺตพฺเพ โข ‘ปุคฺคโล นุปลพฺภติ สจฺจิกฏฺฐปรมตฺเถน’, โน จ วตฺตพฺเพ ‘โย สจฺจิกฏฺโฐ ปรมตฺโถ, ตโต โส ปุคฺคโล นุปลพฺภติ สจฺจิกฏฺฐปรมตฺเถนา’ติ”,}

\prose{โน เจ ปน วตฺตพฺเพ “โย สจฺจิกฏฺโฐ ปรมตฺโถ, ตโต โส ปุคฺคโล นุปลพฺภติ สจฺจิกฏฺฐปรมตฺเถนา”ติ, โน จ วต เร วตฺตพฺเพ “ปุคฺคโล นุปลพฺภติ สจฺจิกฏฺฐปรมตฺเถนา”ติ.\csromanspacer{} ยํ ตตฺถ วเทสิ “วตฺตพฺเพ โข ‘ปุคฺคโล นุปลพฺภติ สจฺจิกฏฺฐปรมตฺเถน’, โน จ วตฺตพฺเพ ‘โย \csromanfolio{7}สจฺจิกฏฺโฐ ปรมตฺโถ, ตโต โส ปุคฺคโล นุปลพฺภติ สจฺจิกฏฺฐปรมตฺเถนา’ติ”, อิทํ เต มิจฺฉา.\csromanspacer{} อุปนยนจตุกฺกํ.}

\proseitem{10}{น เหวํ นิคฺคเหตพฺเพ, เตน หิ ยํ นิคฺคณฺหาสิ\csromandash{}หญฺจิ ปุคฺคโล นุปลพฺภติ สจฺจิกฏฺฐปรมตฺเถน, เตน วต เร วตฺตพฺเพ “โย สจฺจิกฏฺโฐ ปรมตฺโถ, ตโต โส ปุคฺคโล นุปลพฺภติ สจฺจิกฏฺฐปรมตฺเถนา”ติ.\csromanspacer{} ยํ ตตฺถ วเทสิ “วตฺตพฺเพ โข ‘ปุคฺคโล นุปลพฺภติ สจฺจิกฏฺฐปรมตฺเถน’, โน จ วตฺตพฺเพ ‘โย สจฺจิกฏฺโฐ ปรมตฺโถ, ตโต โส ปุคฺคโล นุปลพฺภติ}

\prose{โน เจ ปน วตฺตพฺเพ “โย สจฺจิกฏฺโฐ ปรมตฺโถ, ตโต โส ปุคฺคโล นุปลพฺภติ สจฺจิกฏฺฐปรมตฺเถนา”ติ, โน จ วต เร วตฺตพฺเพ “ปุคฺคโล นุปลพฺภติ สจฺจิกฏฺฐปรมตฺเถนา”ติ.\csromanspacer{} ยํ ตตฺถ วเทสิ “วตฺตพฺเพ โข ‘ปุคฺคโล นุปลพฺภติ สจฺจิกฏฺฐปรมตฺเถน’, โน จ วตฺตพฺเพ ‘โย สจฺจิกฏฺโฐ ปรมตฺโถ, ตโต โส ปุคฺคโล นุปลพฺภติ, สจฺจิกฏฺฐปรมตฺเถนา’ติ”, อิทํ เต มิจฺฉา.\csromanspacer{} เตน หิ เย กเต นิคฺคเห เส นิคฺคเห ทุกฺกเฏ.\csromanspacer{} สุกเต ปฏิกมฺเม, สุกตา ปฏิปาทนาติ.}

\vspace{\tipitakaparskip}
\csromancenter{นิคฺคมนจตุกฺกํ.}
\csromancenter{ทุติโย นิคฺคโห.}
\csromanmatikaanchor{mtk.3}
\csromantocmarkref{section}{2. โอกาสสจฺจิกฏฺฐ}{mtk.3}
\bookmark[dest={mtk.3},level=1]{2. โอกาสสจฺจิกฏฺฐ}
\csromanheader{1}{2. โอกาสสจฺจิกฏฺฐ}
\proseitem{11}{\csromansymbolmark{*}ปุคฺคโล อุปลพฺภติ สจฺจิกฏฺฐปรมตฺเถนาติ, อามนฺตา.\csromanspacer{} สพฺพตฺถ ปุคฺคโล อุปลพฺภติ สจฺจิกฏฺฐปรมตฺเถนาติ, น เหวํ วตฺตพฺเพ.}

\prose{อาชานาหิ นิคฺคหํ\csromandash{}หญฺจิ ปุคฺคโล อุปลพฺภติ สจฺจิกฏฺฐปรมตฺเถน, เตน วต เร วตฺตพฺเพ “สพฺพตฺถ ปุคฺคโล อุปลพฺภติ สจฺจิกฏฺฐปรมตฺเถนา”ติ.\csromanspacer{} ยํ ตตฺถ วเทสิ “วตฺตพฺเพ โข ‘ปุคฺคโล อุปลพฺภติ สจฺจิกฏฺฐปรมตฺเถน’, โน จ วตฺตพฺเพ ‘สพฺพตฺถ ปุคฺคโล อุปลพฺภติ สจฺจิกฏฺฐปรมตฺเถนา’ติ”, มิจฺฉา.}

\prose{\csromanfolio{8}โน เจ ปน วตฺตพฺเพ “สพฺพตฺถ ปุคฺคโล อุปลพฺภติ สจฺจิกฏฺฐปรมตฺเถนา”ติ, โน จ วต เร วตฺตพฺเพ “ปุคฺคโล อุปลพฺภติ สจฺจิกฏฺฐปรมตฺเถนา”ติ.\csromanspacer{} ยํ ตตฺถ วเทสิ “วตฺตพฺเพ โข ‘ปุคฺคโล อุปลพฺภติ สจฺจิกฏฺฐปรมตฺเถน’, โน จ วตฺตพฺเพ ‘สพฺพตฺถ ปุคฺคโล อุปลพฺภติ สจฺจิกฏฺฐปรมตฺเถนา’ติ”, มิจฺฉา.\csromanspacer{}\csromanspacer{}(เปยฺยาล.)}

\vspace{\tipitakaparskip}
\csromancenter{ตติโย นิคฺคโห.}
\csromanmatikaanchor{mtk.4}
\csromantocmarkref{section}{3. กาลสจฺจิกฏฺฐ}{mtk.4}
\bookmark[dest={mtk.4},level=1]{3. กาลสจฺจิกฏฺฐ}
\csromanheader{1}{3. กาลสจฺจิกฏฺฐ}
\proseitem{12}{\csromansymbolmark{*}ปุคฺคโล อุปลพฺภติ สจฺจิกฏฺฐปรมตฺเถนาติ, อามนฺตา.\csromanspacer{} สพฺพทา ปุคฺคโล อุปลพฺภติ สจฺจิกฏฺฐปรมตฺเถนาติ, น เหวํ วตฺตพฺเพ.}

\prose{อาชานาหิ นิคฺคหํ\csromandash{}หญฺจิ ปุคฺคโล อุปลพฺภติ สจฺจิกฏฺฐปรมตฺเถน, เตน วต เร วตฺตพฺเพ “สพฺพทา ปุคฺคโล อุปลพฺภติ สจฺจิกฏฺฐปรมตฺเถนา”ติ.\csromanspacer{} ยํ ตตฺถ วเทสิ “วตฺตพฺเพ โข ‘ปุคฺคโล อุปลพฺภติ สจฺจิกฏฺฐปรมตฺเถน’, โน จ วตฺตพฺเพ ‘สพฺพทา ปุคฺคโล อุปลพฺภติ สจฺจิกฏฺฐปรมตฺเถนา’ติ”, มิจฺฉา.}

\prose{โน เจ ปน วตฺตพฺเพ “สพฺพทา ปุคฺคโล อุปลพฺภติ สจฺจิกฏฺฐปรมตฺเถนา”ติ, โน จ วต เร วตฺตพฺเพ “ปุคฺคโล อุปลพฺภติ สจฺจิกฏฺฐปรมตฺเถนา”ติ.\csromanspacer{} ยํ ตตฺถ วเทสิ “วตฺตพฺเพ โข ‘ปุคฺคโล อุปลพฺภติ สจฺจิกฏฺฐปรมตฺเถน’, โน จ วตฺตพฺเพ ‘สพฺพทา ปุคฺคโล อุปลพฺภติ สจฺจิกฏฺฐปรมตฺเถนา’ติ”, มิจฺฉา ฯเปฯ.}

\vspace{\tipitakaparskip}
\csromancenter{จตุตฺโถ นิคฺคโห.}
\csromanmatikaanchor{mtk.5}
\csromantocmarkref{section}{4. อวยวสจฺจิกฏฺฐ}{mtk.5}
\bookmark[dest={mtk.5},level=1]{4. อวยวสจฺจิกฏฺฐ}
\csromanheader{1}{4. อวยวสจฺจิกฏฺฐ}
\proseitem{13}{\csromansymbolmark{*}ปุคฺคโล อุปลพฺภติ สจฺจิกฏฺฐปรมตฺเถนาติ, อามนฺตา.\csromanspacer{} สพฺเพสุ ปุคฺคโล อุปลพฺภติ สจฺจิกฏฺฐปรมตฺเถนาติ, น เหวํ วตฺตพฺเพ.}

\prose{\csromanfolio{9}อาชานาหิ นิคฺคหํ\csromandash{}หญฺจิ ปุคฺคโล อุปลพฺภติ สจฺจิกฏฺฐปรมตฺเถน, เตน วต เร วตฺตพฺเพ “สพฺเพสุ ปุคฺคโล อุปลพฺภติ สจฺจิกฏฺฐปรมตฺเถนา”ติ.\csromanspacer{} ยํ ตตฺถ วเทสิ “วตฺตพฺเพ โข ‘ปุคฺคโล อุปลพฺภติ สจฺจิกฏฺฐปรมตฺเถน’, โน จ วตฺตพฺเพ ‘สพฺเพสุ ปุคฺคโล อุปลพฺภติ สจฺจิกฏฺฐปรมตฺเถนา’ติ”, มิจฺฉา}

\prose{โน เจ ปน วตฺตพฺเพ “สพฺเพสุ ปุคฺคโล อุปลพฺภติ สจฺจิกฏฺฐปรมตฺเถนา”ติ, โน จ วต เร วตฺตพฺเพ “ปุคฺคโล อุปลพฺภติ สจฺจิกฏฺฐปรมตฺเถนา”ติ.\csromanspacer{} ยํ ตตฺถ วเทสิ “วตฺตพฺเพ โข ‘ปุคฺคโล อุปลพฺภติ สจฺจิกฏฺฐปรมตฺเถน’, โน จ วตฺตพฺเพ ‘สพฺเพสุ ปุคฺคโล อุปลพฺภติ สจฺจิกฏฺฐปรมตฺเถนา’ติ”, มิจฺฉา ฯเปฯ.}

\vspace{\tipitakaparskip}
\csromancenter{ปญฺจโม นิคฺคโห.}
\csromanheader{2}{2. โอกาสสจฺจิกฏฺฐ}
\proseitem{14}{\csromansymbolmark{+}ปุคฺคโล นุปลพฺภติ สจฺจิกฏฺฐปรมตฺเถนาติ, อามนฺตา.\csromanspacer{} สพฺพตฺถ ปุคฺคโล นุปลพฺภติ สจฺจิกฏฺฐปรมตฺเถนาติ, น เหวํ วตฺตพฺเพ.}

\prose{อาชานาหิ นิคฺคหํ\csromandash{}หญฺจิ ปุคฺคโล นุปลพฺภติ สจฺจิกฏฺฐปรมตฺเถน, เตน วต เร วตฺตพฺเพ “สพฺพตฺถ ปุคฺคโล นุปลพฺภติ สจฺจิกฏฺฐปรมตฺเถนา”ติ.\csromanspacer{} ยํ ตตฺถ วเทสิ “วตฺตพฺเพ โข ‘ปุคฺคโล นุปลพฺภติ สจฺจิกฏฺฐปรมตฺเถน’, โน จ วตฺตพฺเพ ‘สพฺพตฺถ ปุคฺคโล นุปลพฺภติ สจฺจิกฏฺฐปรมตฺเถนา’ติ”, มิจฺฉา.}

\prose{โน เจ ปน วตฺตพฺเพ “สพฺพตฺถ ปุคฺคโล นุปลพฺภติ สจฺจิกฏฺฐปรมตฺเถนา”ติ, โน จ วต เร วตฺตพฺเพ “ปุคฺคโล นุปลพฺภติ สจฺจิกฏฺฐปรมตฺเถนา”ติ.\csromanspacer{} ยํ ตตฺถ วเทสิ “วตฺตพฺเพ โข ‘ปุคฺคโล นุปลพฺภติ สจฺจิกฏฺฐปรมตฺเถน’, โน จ วตฺตพฺเพ ‘สพฺพตฺถ ปุคฺคโล นุปลพฺภติ สจฺจิกฏฺฐปรมตฺเถนา’ติ”, มิจฺฉา ฯเปฯ.}

\vspace{\tipitakaparskip}
\csromancenter{ฉฏฺโฐ นิคฺคโห.}
\csromanheader{2}{3. กาลสจฺจิกฏฺฐ}
\proseitem{15}{\csromanfolio{10}\csromansymbolmark{+}ปุคฺคโล นุปลพฺภติ สจฺจิกฏฺฐปรมตฺเถนาติ, อามนฺตา.\csromanspacer{} สพฺพทา ปุคฺคโล นุปลพฺภติ สจฺจิกฏฺฐปรมตฺเถนาติ, น เหวํ วตฺตพฺเพ.}

\prose{อาชานาหิ นิคฺคหํ\csromandash{}หญฺจิ ปุคฺคโล นุปลพฺภติ สจฺจิกฏฺฐปรมตฺเถน, เตน วต เร วตฺตพฺเพ “สพฺพทา ปุคฺคโล นุปลพฺภติ สจฺจิกฏฺฐปรมตฺเถนา”ติ.\csromanspacer{} ยํ ตตฺถ วเทสิ “วตฺตพฺเพ โข ‘ปุคฺคโล นุปลพฺภติ สจฺจิกฏฺฐปรมตฺเถน’, โน จ วตฺตพฺเพ ‘สพฺพทา ปุคฺคโล นุปลพฺภติ สจฺจิกฏฺฐปรมตฺเถนา’ติ”, มิจฺฉา.}

\prose{โน เจ ปน วตฺตพฺเพ “สพฺพทา ปุคฺคโล นุปลพฺภติ สจฺจิกฏฺฐปรมตฺเถนา”ติ, โน จ วต เร วตฺตพฺเพ “ปุคฺคโล นุปลพฺภติ สจฺจิกฏฺฐปรมตฺเถนา”ติ.\csromanspacer{} ยํ ตตฺถ วเทสิ “วตฺตพฺเพ โข ‘ปุคฺคโล นุปลพฺภติ สจฺจิกฏฺฐปรมตฺเถน’, โน จ วตฺตพฺเพ ‘สพฺพทา ปุคฺคโล นุปลพฺภติ สจฺจิกฏฺฐปรมตฺเถนา’ติ”, มิจฺฉา ฯเปฯ.}

\vspace{\tipitakaparskip}
\csromancenter{สตฺตโม นิคฺคโห.}
\csromanheader{2}{4. อวยวสจฺจิกฏฺฐ}
\proseitem{16}{\csromansymbolmark{+}ปุคฺคโล นุปลพฺภติ สจฺจิกฏฺฐปรมตฺเถนาติ, อามนฺตา.\csromanspacer{} สพฺเพสุ ปุคฺคโล นุปลพฺภติ สจฺจิกฏฺฐปรมตฺเถนาติ, น เหวํ วตฺตพฺเพ.}

\prose{อาชานาหิ นิคฺคหํ\csromandash{}หญฺจิ ปุคฺคโล นุปลพฺภติ สจฺจิกฏฺฐปรมตฺเถน, เตน วต เร วตฺตพฺเพ “สพฺเพสุ ปุคฺคโล นุปลพฺภติ สจฺจิกฏฺฐปรมตฺเถนา”ติ.\csromanspacer{} ยํ ตตฺถ วเทสิ “วตฺตพฺเพ โข ‘ปุคฺคโล นุปลพฺภติ สจฺจิกฏฺฐปรมตฺเถน’, โน จ วตฺตพฺเพ ‘สพฺเพสุ ปุคฺคโล นุปลพฺภติ สจฺจิกฏฺฐปรมตฺเถนา’ติ”, มิจฺฉา.}

\prose{โน เจ ปน วตฺตพฺเพ “สพฺเพสุ ปุคฺคโล นุปลพฺภติ สจฺจิกฏฺฐปรมตฺเถนา”ติ, โน จ วต เร วตฺตพฺเพ “ปุคฺคโล นุปลพฺภติ สจฺจิกฏฺฐปรมตฺเถนา”ติ.\csromanspacer{} ยํ ตตฺถ วเทสิ “วตฺตพฺเพ โข ‘ปุคฺคโล นุปลพฺภติ \csromanfolio{11}สจฺจิกฏฺฐปรมตฺเถน’, โน จ วตฺตพฺเพ ‘สพฺเพสุ ปุคฺคโล นุปลพฺภติ สจฺจฺจิกฏฺฐปรมตฺเถนา’ติ”, มิจฺฉา ฯเปฯ.}

\vspace{\tipitakaparskip}
\csromancenter{อฏฺฐกนิคฺคโห.}
\csromanmatikaanchor{mtk.6}
\csromantocmarkref{section}{5. สุทฺธิกสํสนฺทน}{mtk.6}
\bookmark[dest={mtk.6},level=1]{5. สุทฺธิกสํสนฺทน}
\csromanheader{1}{5. สุทฺธิกสํสนฺทน}
\proseitem{17}{\csromansymbolmark{*}ปุคฺคโล อุปลพฺภติ สจฺจิกฏฺฐปรมตฺเถน, รูปญฺจ อุปลพฺภติ สจฺจิกฏฺฐปรมตฺเถนาติ, อามนฺตา.\csromanspacer{} อญฺญํ รูปํ อญฺโญ ปุคฺคโลติ, น เหวํ วตฺตพฺเพ.}

\prose{อาชานาหิ นิคฺคหํ\csromandash{}หญฺจิ ปุคฺคโล อุปลพฺภติ สจฺจิกฏฺฐปรมตฺเถน, รูปญฺจ อุปลพฺภติ สจฺจิกฏฺฐปรมตฺเถน, เตน วต เร วตฺตพฺเพ “อญฺญํ รูปํ อญฺโญ ปุคฺคโล”ติ.\csromanspacer{} ยํ ตตฺถ วเทสิ “วตฺตพฺเพ โข ‘ปุคฺคโล อุปลพฺภติ สจฺจิกฏฺฐปรมตฺเถน, รูปญฺจ อุปลพฺภติ สจฺจิกฏฺฐปรมตฺเถน’.\csromanspacer{} โน จ วตฺตพฺเพ ‘อญฺญํ รูปํ อญฺโญ ปุคฺคโล’ติ”, มิจฺฉา.}

\prose{โน เจ ปน วตฺตพฺเพ “อญฺญํ รูปํ อญฺโญ ปุคฺคโล”ติ, โน จ วต เร วตฺตพฺเพ “ปุคฺคโล อุปลพฺภติ สจฺจิกฏฺฐปรมตฺเถน, รูปญฺจ อุปลพฺภติ สจฺจิกฏฺฐปรมตฺเถนา”ติ.\csromanspacer{} ยํ ตตฺถ วเทสิ “วตฺตพฺเพ โข ‘ปุคฺคโล อุปลพฺภติ สจฺจิกฏฺฐปรมตฺเถน, รูปญฺจ อุปลพฺภติ สจฺจิกฏฺฐปรมตฺเถน’, โน จ วตฺตพฺเพ ‘อญฺญํ รูปํ อญฺโญ ปุคฺคโล”ติ”, มิจฺฉา ฯเปฯ.}

\proseitem{18}{\csromansymbolmark{*}ปุคฺคโล อุปลพฺภติ สจฺจิกฏฺฐปรมตฺเถน, เวทนา จ อุปลพฺภติ สจฺจิกฏฺฐปรมตฺเถน ฯเปฯ สญฺญา จ อุปลพฺภติ ฯเปฯ สงฺขารา จ อุปลพฺภนฺติ ฯเปฯ วิญฺญาณญฺจ อุปลพฺภติ สจฺจิกฏฺฐปรมตฺเถนาติ, อามนฺตา.\csromanspacer{} อญฺญํ วิญฺญาณํ อญฺโญ ปุคฺคโลติ, น เหวํ วตฺตพฺเพ.}

\prose{อาชานาหิ นิคฺคหํ\csromandash{}หญฺจิ ปุคฺคโล อุปลพฺภติ สจฺจิกฏฺฐปรมตฺเถน, วิญฺญาณญฺจ อุปลพฺภติ สจฺจิกฏฺฐปรมตฺเถน, เตน วต เร วตฺตพฺเพ “อญฺญํ วิญฺญาณํ อญฺโญ ปุคฺคโล”ติ.\csromanspacer{} ยํ ตตฺถ วเทสิ “วตฺตพฺเพ โข ‘ปุคฺคโล อุปลพฺภติ สจฺจิกฏฺฐปรมตฺเถน, วิญฺญาณญฺจ อุปลพฺภติ สจฺจิกฏฺฐปรมตฺเถน’, โน จ วตฺตพฺเพ ‘อญฺญํ วิญฺญาณํ อญฺโญ ปุคฺคโล’ติ”, มิจฺฉา.}

\prose{\csromanfolio{12}โน เจ ปน วตฺตพฺเพ “อญฺญํ วิญฺญาณํ อญฺโญ ปุคฺคโล”ติ, โน จ วต เร วตฺตพฺเพ “ปุคฺคโล อุปลพฺภติ สจฺจิกฏฺฐปรมตฺเถน, วิญฺญาณญฺจ อุปลพฺภติ สจฺจิกฏฺฐปรมตฺเถนา”ติ.\csromanspacer{} ยํ ตตฺถ วเทสิ “วตฺตพฺเพ โข ‘ปุคฺคโล อุปลพฺภติ สจฺจิกฏฺฐปรมตฺเถน, วิญฺญาณญฺจ อุปลพฺภติ สจฺจิกฏฺฐปรมตฺเถน’, โน จ วตฺตพฺเพ ‘อญฺญํ วิญฺญาณํ อญฺโญ ปุคฺคโล’ติ”, มิจฺฉา ฯเปฯ.}

\proseitem{19}{\csromansymbolmark{*}ปุคฺคโล อุปลพฺภติ สจฺจิกฏฺฐปรมตฺเถน, จกฺขายตนญฺจ อุปลพฺภติ สจฺจิกฏฺฐปรมตฺเถน ฯเปฯ โสตายตนญฺจ อุปลพฺภติ.\csromanspacer{} ฆานายตนญฺจ อุปลพฺภติ.\csromanspacer{} ชิวฺหายตนญฺจ อุปลพฺภติ.\csromanspacer{} กายายตนญฺจ อุปลพฺภติ.\csromanspacer{} รูปายตนญฺจ อุปลพฺภติ.\csromanspacer{} สทฺทายตนญฺจ อุปลพฺภติ.\csromanspacer{} คนฺธายตนญฺจ อุปลพฺภติ.\csromanspacer{} รสายตนญฺจ อุปลพฺภติ.\csromanspacer{} โผฏฺฐพฺพายตนญฺจ อุปลพฺภติ.\csromanspacer{} มนายตนญฺจ อุปลพฺภติ.\csromanspacer{} ธมฺมายตนญฺจ อุปลพฺภติ สจฺจิกฏฺฐปรมตฺเถน ฯเปฯ.}

\proseitem{20}{\csromansymbolmark{*}จกฺขุธาตุ จ อุปลพฺภติ สจฺจิกฏฺฐปรมตฺเถน ฯเปฯ.\csromanspacer{} โสตธาตุ จ อุปลพฺภติ.\csromanspacer{} ฆานธาตุ จ อุปลพฺภติ.\csromanspacer{} ชิวฺหาธาตุ จ อุปลพฺภติ.\csromanspacer{} กายธาตุ จ อุปลพฺภติ.\csromanspacer{} รูปธาตุ จ อุปลพฺภติ.\csromanspacer{} สทฺทธาตุ จ อุปลพฺภติ.\csromanspacer{} คนฺธธาตุ จ อุปลพฺภติ.\csromanspacer{} รสธาตุ จ อุปลพฺภติ.\csromanspacer{} โผฏฺฐพฺพธาตุ จ อุปลพฺภติ.\csromanspacer{} จกฺขุวิญฺญาณธาตุ จ อุปลพฺภติ.\csromanspacer{} โสตวิญฺญาณธาตุ จ อุปลพฺภติ.\csromanspacer{} ฆานวิญฺญาณธาตุ จ อุปลพฺภติ.\csromanspacer{} ชิวฺหาวิญฺญาณธาตุ จ อุปลพฺภติ.\csromanspacer{} กายวิญฺญาณธาตุ จ อุปลพฺภติ.\csromanspacer{} มโนธาตุ จ อุปลพฺภติ.\csromanspacer{} มโนวิญฺญาณธาตุ จ อุปลพฺภติ.\csromanspacer{} ธมฺมธาตุ จ อุปลพฺภติ สจฺจิกฏฺฐปรมตฺเถน ฯเปฯ.}

\proseitem{21}{\csromansymbolmark{*}จกฺขุนฺทฺริยญฺจ อุปลพฺภติ สจฺจิกฏฺฐปรมตฺเถน ฯเปฯ.\csromanspacer{} โสตินฺทฺริยญฺจ อุปลพฺภติ.\csromanspacer{} ฆานินฺทฺริยญฺจ อุปลพฺภติ.\csromanspacer{} ชิวฺหินฺทฺริยญฺจ อุปลพฺภติ.\csromanspacer{} กายินฺทฺริยญฺจ อุปลพฺภติ.\csromanspacer{} มนินฺทฺริยญฺจ อุปลพฺภติ.\csromanspacer{} ชีวิตินฺทฺริยญฺจ อุปลพฺภติ.\csromanspacer{} อิตฺถินฺทฺริยญฺจ อุปลพฺภติ.\csromanspacer{} ปุริสินฺทฺริยญฺจ อุปลพฺภติ.\csromanspacer{} สุขินฺทฺริยญฺจ อุปลพฺภติ.\csromanspacer{} ทุกฺขินฺทฺริยญฺจ อุปลพฺภติ.\csromanspacer{} โสมนสฺสินฺทฺริยญฺจ อุปลพฺภติ.\csromanspacer{} โทมนสฺสินฺทฺริยญฺจ อุปลพฺภติ.\csromanspacer{} อุเปกฺขินฺทฺริยญฺจ อุปลพฺภติ.\csromanspacer{} สทฺธินฺทฺริยญฺจ อุปลพฺภติ.\csromanspacer{} วีริยินฺทฺริยญฺจ อุปลพฺภติ.\csromanspacer{} สตินฺทฺริยญฺจ อุปลพฺภติ.\csromanspacer{} สมาธินฺทฺริยญฺจ อุปลพฺภติ.\csromanspacer{} ปญฺญินฺทฺริยญฺจ อุปลพฺภติ.\csromanspacer{} อนญฺญาตญฺญสฺสามีตินฺทฺริยญฺจ อุปลพฺภติ.\csromanspacer{} อญฺญินฺทฺริยญฺจ อุปลพฺภติ.\csromanspacer{} อญฺญาตาวินฺทฺริยญฺจ อุปลพฺภติ สจฺจิกฏฺฐปรมตฺเถนาติ, อามนฺตา.\csromanspacer{} อญฺญํ อญฺญาตาวินฺทฺริยํ อญฺโญ ปุคฺคโลติ, น เหวํ วตฺตพฺเพ.}

\prose{\csromanfolio{13}อาชานาหิ นิคฺคหํ\csromandash{}หญฺจิ ปุคฺคโล อุปลพฺภติ สจฺจิกฏฺฐปรมตฺเถน, อญฺญาตาวินฺทฺริยญฺจ อุปลพฺภติ สจฺจิกฏฺฐปรมตฺเถน, เตน วต เร วตฺตพฺเพ “อญฺญํ อญฺญาตาวินฺทฺริยํ อญฺโญ ปุคฺคโล”ติ.\csromanspacer{} ยํ ตตฺถ วเทสิ “วตฺตพฺเพ โข ‘ปุคฺคโล อุปลพฺภติ สจฺจิกฏฺฐปรมตฺเถน, อญฺญาตาวินฺทฺริยญฺจ อุปลพฺภติ สจฺจิกฏฺฐปรมตฺเถน’, โน จ วตฺตพฺเพ ‘อญฺญํ อญฺญาตาวินฺทฺริยํ อญฺโญ ปุคฺคโล’ติ”, มิจฺฉา.}

\prose{โน เจ ปน วตฺตพฺเพ “อญฺญํ อญฺญาตาวินฺทฺริยํ อญฺโญ ปุคฺคโล”ติ, โน จ วต เร วตฺตพฺเพ “ปุคฺคโล อุปลพฺภติ สจฺจิกฏฺฐปรมตฺเถน, อญฺญาตาวินฺทฺริยญฺจ อุปลพฺภติ สจฺจิกฏฺฐปรมตฺเถนา”ติ.\csromanspacer{} ยํ ตตฺถ วเทสิ “วตฺตพฺเพ โข ‘ปุคฺคโล อุปลพฺภติ สจฺจิกฏฺฐปรมตฺเถน, อญฺญาตาวินฺทฺริยญฺจ อุปลพฺภติ สจฺจิกฏฺฐปรมตฺเถน’, โน จ วตฺตพฺเพ ‘อญฺญํ อญฺญาตาวินฺทฺริยํ อญฺโญ ปุคฺคโล’ติ”, มิจฺฉา ฯเปฯ.}

\proseitem{22}{\csromansymbolmark{+}ปุคฺคโล นุปลพฺภติ สจฺจิกฏฺฐปรมตฺเถนาติ, อามนฺตา.\csromanspacer{} วุตฺตํ ภควตา “อตฺถิ ปุคฺคโล อตฺตหิตาย ปฏิปนฺโน”\footnote{ปุคฺคลปญฺญตฺติ 109; อํ 1. 408 ปิฏฺเฐสุ.}, รูปญฺจ อุปลพฺภติ สจฺจิกฏฺฐปรมตฺเถนาติ, อามนฺตา.\csromanspacer{} อญฺญํ รูปํ อญฺโญ ปุคฺคโลติ, น เหวํ วตฺตพฺเพ.}

\prose{อาชานาหิ ปฏิกมฺมํ\csromandash{}หญฺจิ วุตฺตํ ภควตา “อตฺถิ ปุคฺคโล อตฺตหิตาย ปฏิปนฺโน”, รูปญฺจ อุปลพฺภติ สจฺจิกฏฺฐปรมตฺเถน, เตน วต เร วตฺตพฺเพ “อญฺญํ รูปํ อญฺโญ ปุคฺคโล”ติ.\csromanspacer{} ยํ ตตฺถ วเทสิ “วตฺตพฺเพ โข ‘วุตฺตํ ภควตา อตฺถิ ปุคฺคโล อตฺตหิตาย ปฏิปนฺโน, รูปญฺจ อุปลพฺภติ สจฺจิกฏฺฐปรมตฺเถน’, โน จ วตฺตพฺเพ ‘อญฺญํ รูปํ อญฺโญ ปุคฺคโล’ติ”, มิจฺฉา.}

\prose{โน เจ ปน วตฺตพฺเพ “อญฺญํ รูปํ อญฺโญ ปุคฺคโล”ติ, โน จ วต เร วตฺตพฺเพ “วุตฺตํ ภควตา อตฺถิ ปุคฺคโล อตฺตหิตาย ปฏิปนฺโน, รูปญฺจ อุปลพฺภติ สจฺจิกฏฺฐปรมตฺเถนา”ติ.\csromanspacer{} ยํ ตตฺถ วเทสิ “วตฺตพฺเพ โข ‘วุตฺตํ ภควตา อตฺถิ ปุคฺคโล อตฺตหิหาย ปฏิปนฺโน, รูปญฺจ อุปลพฺภติ สจฺจิกฏฺฐปรมตฺเถน’, โน จ วตฺตพฺเพ ‘อญฺญํ รูปํ ‘อญฺโญ ปุคฺคโล’ติ”, มิจฺฉา ฯเปฯ.}

\proseitem{23}{\csromanfolio{14}\csromansymbolmark{+}ปุคฺคโล นุปลพฺภติ สจฺจิกฏฺฐปรมตฺเถนาติ, อามนฺตา.\csromanspacer{} วุตฺตํ ภควตา “อตฺถิ ปุคฺคโล อตฺตหิตาย ปฏิปนฺโน”, เวทนา จ อุปลพฺภติ ฯเปฯ สญฺญา จ อุปลพฺภติ ฯเปฯ สงฺขารา จ อุปลพฺภนฺติ ฯเปฯ วิญฺญาณญฺจ อุปลพฺภติ สจฺจิกฏฺฐปรมตฺเถนาติ, อามนฺตา.\csromanspacer{} อญฺญํ วิญฺญาณํ อญฺโญ ปุคฺคโลติ, น เหวํ วตฺตพฺเพ.}

\prose{อาชานาหิ ปฏิกมฺมํ\csromandash{}หญฺจิ วุตฺตํ ภควตา ‘อตฺถิ ปุคฺคโล อตฺตหิตาย ปฏิปนฺโน”, วิญฺญาณญฺจ อุปลพฺภติ สจฺจิกฏฺฐปรมตฺเถน, เตน วต เร วตฺตพฺเพ “อญฺญํ วิญฺญาณํ อญฺโญ ปุคฺคโล”ติ.\csromanspacer{} ยํ ตตฺถ วเทสิ “วตฺตพฺเพ โข ‘วุตฺตํ ภควตา อตฺถิ ปุคฺคโล อตฺตหิตาย ปฏิปนฺโน, วิญฺญาณญฺจ อุปลพฺภติ สจฺจิกฏฺฐปรมตฺเถน’, โน จ วตฺตพฺเพ อญฺญํ วิญฺญาณํ อญฺโญ ปุคฺคโล’ติ”, มิจฺฉา.}

\prose{โน เจ ปน วตฺตพฺเพ “อญฺญํ วิญฺญาณํ อญฺโญ ปุคฺคโล”ติ, โน จ วต เร วตฺตพฺเพ “วุตฺตํ ภควตา ‘อตฺถิ ปุคฺคโล อตฺตหิตาย ปฏิปนฺโน’ วิญฺญาณญฺจ อุปลพฺภติ สจฺจิกฏฺฐปรมตฺเถนา”ติ.\csromanspacer{} ยํ ตตฺถ วเทสิ “วตฺตพฺเพ โข ‘วุตฺตํ ภควตา อตฺถิ ปุคฺคโล อตฺตหิตาย ปฏิปนฺโน, วิญฺญาณญฺจ อุปลพฺภติ สจฺจิกฏฺฐปรมตฺเถน’, โน จ วตฺตพฺเพ ‘อญฺญํ วิญฺญาณํ อญฺโญ ปุคฺคโล’ติ”, มิจฺฉา ฯเปฯ.}

\proseitem{24}{\csromansymbolmark{+}ปุคฺคโล นุปลพฺภติ สจฺจิกฏฺฐปรมตฺเถนาติ, อามนฺตา.\csromanspacer{} วุตฺตํ ภควตา “อตฺถิ ปุคฺคโล อตฺตหิตาย ปฏิปนฺโน”, จกฺขายตนญฺจ อุปลพฺภติ สจฺจิกฏฺฐปรมตฺเถน ฯเปฯ โสตายตนญฺจ อุปลพฺภติ ฯเปฯ ธมฺมายตนญฺจ อุปลพฺภติ สจฺจิกฏฺฐปรมตฺเถน ฯเปฯ.}

\proseitem{25}{\csromansymbolmark{+}จกฺขุธาตุ จ อุปลพฺภติ สจฺจิกฏฺฐปรมตฺเถน ฯเปฯ.\csromanspacer{} กายธาตุ จ อุปลพฺภติ ฯเปฯ.\csromanspacer{} รูปธาตุ จ อุปลพฺภติ ฯเปฯ.\csromanspacer{} โผฏฺฐพฺพธาตุ จ อุปลพฺภติ ฯเปฯ.\csromanspacer{} จกฺขุวิญฺญาณธาตุ จ อุปลพฺภติ ฯเปฯ.\csromanspacer{} มโนวิญฺญาณธาตุ จ อุปลพฺภติ ฯเปฯ.\csromanspacer{} ธมฺมธาตุ จ อุปลพฺภติ สจฺจิกฏฺฐปรมตฺเถน ฯเปฯ.}

\proseitem{26}{\csromansymbolmark{+}จกฺขุนฺทฺริยญฺจ อุปลพฺภติ สจฺจิกฏฺฐปรมตฺเถน ฯเปฯ.\csromanspacer{} โสตินฺทฺริยญฺจ อุปลพฺภติ สจฺจิกฏฺฐปรมตฺเถน ฯเปฯ.\csromanspacer{} อญฺญินฺทฺริยญฺจ\footnote{อญฺญาตาวินฺทฺริยญฺจ (พหูสุ)} อุปลพฺภติ สจฺจิกฏฺฐปรมตฺเถน ฯเปฯ.}

\proseitem{27}{\csromanfolio{15}\csromansymbolmark{+}ปุคฺคโล นุปลพฺภติ สจฺจิกฏฺฐปรมตฺเถนาติ, อามนฺตา.\csromanspacer{} วุตฺตํ ภควตา “อตฺถิ ปุคฺคโล อตฺตหิตาย ปฏิปนฺโน”, อญฺญาตาวินฺทฺริยญฺจ อุปลพฺภติ สจฺจิกฏฺฐปรมตฺเถนาติ, อามนฺตา.\csromanspacer{} อญฺญํ อญฺญาตาวินฺทฺริยํ อญฺโญ ปุคฺคโลติ, น เหวํ วตฺตพฺเพ.}

\prose{อาชานาหิ ปฏิกมฺมํ\csromandash{}หญฺจิ วุตฺตํ ภควตา “อตฺถิ ปุคฺคโล อตฺตหิตาย ปฏิปนฺโน”, อญฺญาตาวินฺทฺริยญฺจ อุปลพฺภติ สจฺจิกฏฺฐปรมตฺเถน, เตน วต เร วตฺตพฺเพ “อญฺญํ อญฺญาตาวินฺทฺริยํ อญฺโญ ปุคฺคโล”ติ.\csromanspacer{} ยํ ตตฺถ วเทสิ “วตฺตพฺเพ โข ‘วุตฺตํ ภควตา อตฺถิ ปุคฺคโล อตฺตหิตาย ปฏิปนฺโน, อญฺญาตาวินฺทฺริยญฺจ อุปลพฺภติ สจฺจิกฏฺฐปรมตฺเถน’, โน จ วตฺตพฺเพ ‘อญฺญํ อญฺญาตาวินฺทฺริยํ อญฺโญ ปุคฺคโล’ติ”, มิจฺฉา.}

\prose{โน เจ ปน วตฺตพฺเพ “อญฺญํ อญฺญาตาวินฺทฺริยํ อญฺโญ ปุคฺคโล”ติ, โน จ วต เร วตฺตพฺเพ “วุตฺตํ ภควตา ‘อตฺถิ ปุคฺคโล อตฺตหิตาย ปฏิปนฺโน’, อญฺญาตาวินฺทฺริยญฺจ อุปลพฺภติ สจฺจิกฏฺฐปรมตฺเถนา”ติ.\csromanspacer{} ยํ ตตฺถ วเทสิ “วตฺตพฺเพ โข ‘วุตฺตํ ภควตา อตฺถิ ปุคฺคโล อตฺตหิตาย ปฏิปนฺโน, อญฺญาตาวินฺทฺริยญฺจ อุปลพฺภติ สจฺจิกฏฺฐปรมตฺเถน’, โน จ วตฺตพฺเพ ‘อญฺญํ อญฺญาตาวินฺทฺริยํ อญฺโญ ปุคฺคโล’ติ”, มิจฺฉา ฯเปฯ.}

\vspace{\tipitakaparskip}
\csromancenter{สุทฺธิกสํสนฺทนา.}
\csromanmatikaanchor{mtk.7}
\csromantocmarkref{section}{6. โอปมฺมสํสนฺทน}{mtk.7}
\bookmark[dest={mtk.7},level=1]{6. โอปมฺมสํสนฺทน}
\csromanheader{1}{6. โอปมฺมสํสนฺทน}
\proseitem{28}{\csromansymbolmark{*}รูปํ อุปลพฺภติ สจฺจิกฏฺฐปรมตฺเถน, เวทนา จ อุปลพฺภติ สจฺจิกฏฺฐปรมตฺเถน, อญฺญํ รูปํ อญฺญา เวทนาติ, อามนฺตา.\csromanspacer{} ปุคฺคโล อุปลพฺภติ สจฺจิกฏฺฐปรมตฺเถน, รูปญฺจ อุปลพฺภติ สจฺจิกฏฺฐปรมตฺเถนาติ, อามนฺตา.\csromanspacer{} อญฺญํ รูปํ อญฺโญ ปุคฺคโลติ, น เหวํ วตฺตพฺเพ.}

\prose{อาชานาหิ นิคฺคหํ\csromandash{}หญฺจิ รูปํ อุปลพฺภติ สจฺจิกฏฺฐปรมตฺเถน, เวทนา จ อุปลพฺภติ สจฺจิกฏฺฐปรมตฺเถน, อญฺญํ รูปํ อญฺญา เวทนา.\csromanspacer{} ปุคฺคโล อุปลพฺภติ สจฺจิกฏฺฐปรมตฺเถน, รูปญฺจ อุปลพฺภติ สจฺจิกฏฺฐปรมตฺเถน.\csromanspacer{} เตน วต เร วตฺตพฺเพ “อญฺญํ รูปํ อญฺโญ ปุคฺคโล”ติ.\csromanspacer{} ยํ ตตฺถ วเทสิ “วตฺตพฺเพ โข ‘รูปํ อุปลพฺภติ สจฺจิกฏฺฐปรมตฺเถน, เวทนา จ อุปลพฺภติ สจฺจิกฏฺฐปรมตฺเถน, อญฺญํ รูปํ อญฺญา เวทนา.\csromanspacer{} ปุคฺคโล อุปลพฺภติ สจฺจิกฏฺฐปรมตฺเถน, \csromanfolio{16}รูปญฺจ อุปลพฺภติ สจฺจิกฏฺฐปรมตฺเถน’, โน จ วตฺตพฺเพ ‘อญฺญํ รูปํ อญฺโญ ปุคฺคโล’ติ”, มิจฺฉา.}

\prose{โน เจ ปน วตฺตพฺเพ “อญฺญํ รูปํ อญฺโญ ปุคฺคโล”ติ, โน จ วต เร วตฺตพฺเพ “รูปํ อุปลพฺภติ สจฺจิกฏฺฐปรมตฺเถน, เวทนา จ อุปลพฺภติ สจฺจิกฏฺฐปรมตฺเถน, อญฺญํ รูปํ อญฺญา เวทนา.\csromanspacer{} ปุคฺคโล อุปลพฺภติ สจฺจิกฏฺฐปรมตฺเถน, รูปญฺจ อุปลพฺภติ สจฺจิกฏฺฐปรมตฺเถนา”ติ.\csromanspacer{} ยํ ตตฺถ วเทสิ “วตฺตพฺเพ โข ‘รูปํ อุปลพฺภติ สจฺจิกฏฺฐปรมตฺเถน, เวทนา จ อุปลพฺภติ สจฺจิกฏฺฐปรมตฺเถน, อญฺญํ รูปํ อญฺญา เวทนา.\csromanspacer{} ปุคฺคโล อุปลพฺภติ สจฺจิกฏฺฐปรมตฺเถน, รูปญฺจ อุปลพฺภติ สจฺจิกฏฺฐปรมตฺเถน’, โน จ วตฺตพฺเพ ‘อญฺญํ รูปํ อญฺโญ ปุคฺคโล’ติ”, มิจฺฉา ฯเปฯ.}

\proseitem{29}{\csromansymbolmark{*}รูปํ อุปลพฺภติ สจฺจิกฏฺฐปรมตฺเถน, สญฺญา จ อุปลพฺภติ ฯเปฯ สงฺขารา จ อุปลพฺภนฺติ ฯเปฯ วิญฺญาณญฺจ อุปลพฺภติ สจฺจิกฏฺฐปรมตฺเถน, อญฺญํ รูปํ อญฺญํ วิญฺญาณนฺติ, อามนฺตา.\csromanspacer{} ปุคฺคโล อุปลพฺภติ สจฺจิกฏฺฐปรมตฺเถน, รูปญฺจ อุปลพฺภติ สจฺจิกฏฺฐปรมตฺเถนาติ, อามนฺตา.\csromanspacer{} อญฺญํ รูปํ อญฺโญ ปุคฺคโลติ, น เหวํ วตฺตพฺเพ.}

\prose{อาชานาหิ นิคฺคหํ\csromandash{}หญฺจิ รูปํ อุปลพฺภติ สจฺจิกฏฺฐปรมตฺเถน, วิญฺญาณญฺจ อุปลพฺภติ สจฺจิกฏฺฐปรมตฺเถน, อญฺญํ รูปํ อญฺญํ วิญฺญาณํ.\csromanspacer{} ปุคฺคโล อุปลพฺภติ สจฺจิกฏฺฐปรมตฺเถน, รูปญฺจ อุปลพฺภติ สจฺจิกฏฺฐปรมตฺเถน.\csromanspacer{} เตน วต เร วตฺตพฺเพ “อญฺญํ รูปํ อญฺโญ ปุคฺคโล”ติ.\csromanspacer{} ยํ ตตฺถ วเทสิ “วตฺตพฺเพ โข ‘รูปํ อุปลพฺภติ สจฺจิกฏฺฐปรมตฺเถน, วิญฺญาณญฺจ อุปลพฺภติ สจฺจิกฏฺฐปรมตฺเถน, อญฺญํ รูปํ อญฺญํ วิญฺญาณํ.\csromanspacer{} ปุคฺคโล อุปลพฺภติ สจฺจิกฏฺฐปรมตฺเถน, รูปญฺจ อุปลพฺภติ สจฺจิกฏฺฐปรมตฺเถน’, โน จ วตฺตพฺเพ ‘อญฺญํ รูปํ อญฺโญ ปุคฺคโล’ติ”, มิจฺฉา.}

\prose{โน เจ ปน วตฺตพฺเพ “อญฺญํ รูปํ อญฺโญ ปุคฺคโล”ติ, โน จ วต เร วตฺตพฺเพ “รูปํ อุปลพฺภติ สจฺจิกฏฺฐปรมตฺเถน, วิญฺญาณญฺจ อุปลพฺภติ สจฺจิกฏฺฐปรมตฺเถน, อญฺญํ รูปํ อญฺญํ วิญฺญาณํ.\csromanspacer{} ปุคฺคโล อุปลพฺภติ สจฺจิกฏฺฐปรมตฺเถน, รูปญฺจ อุปลพฺภติ สจฺจิกฏฺฐปรมตฺเถนาติ.\csromanspacer{} ยํ ตตฺถ วเทสิ “วตฺตพฺเพ โข ‘รูปํ อุปลพฺภติ สจฺจิกฏฺฐปรมตฺเถน, วิญฺญาณญฺจ อุปลพฺภติ สจฺจิกฏฺฐปรมตฺเถน, อญฺญํ รูปํ อญฺญํ วิญฺญาณํ.\csromanspacer{} ปุคฺคโล อุปลพฺภติ สจฺจิกฏฺฐปรมตฺเถน, รูปญฺจ \csromanfolio{17}อุปลพฺภติ สจฺจิกฏฺฐปรมตฺเถน’, โน จ วตฺตพฺเพ ‘อญฺญํ รูปํ อญฺโญ ปุคฺคโล’ติ”, มิจฺฉา ฯเปฯ.}

\proseitem{30}{\csromansymbolmark{*}เวทนา อุปลพฺภติ สจฺจิกฏฺฐปรมตฺเถน, สญฺญา จ อุปลพฺภติ ฯเปฯ สงฺขารา จ อุปลพฺภนฺติ ฯเปฯ วิญฺญาณญฺจ อุปลพฺภติ ฯเปฯ รูปญฺจ อุปลพฺภติ สจฺจิกฏฺฐปรมตฺเถน ฯเปฯ.}

\proseitem{31}{\csromansymbolmark{*}สญฺญา อุปลพฺภติ สจฺจิกฏฺฐปรมตฺเถน, สงฺขารา จ อุปลพฺภนฺติ ฯเปฯ วิญฺญาณญฺจ อุปลพฺภติ ฯเปฯ รูปญฺจ อุปลพฺภติ ฯเปฯ เวทนา จ อุปลพฺภติ สจฺจิกฏฺฐปรมตฺเถน ฯเปฯ.}

\proseitem{32}{\csromansymbolmark{*}สงฺขารา อุปลพฺภนฺติ สจฺจิกฏฺฐปรมตฺเถน, วิญฺญาณญฺจ อุปลพฺภติ ฯเปฯ รูปญฺจ อุปลพฺภติ ฯเปฯ เวทนา จ อุปลพฺภติ ฯเปฯ สญฺญา จ อุปลพฺภติ สจฺจิกฏฺฐปรมตฺเถน ฯเปฯ.}

\proseitem{33}{\csromansymbolmark{*}วิญฺญาณํ อุปลพฺภติ สจฺจิกฏฺฐปรมตฺเถน, รูปญฺจ อุปลพฺภติ ฯเปฯ เวทนา จ อุปลพฺภติ ฯเปฯ สญฺญา จ อุปลพฺภติ ฯเปฯ สงฺขาราจ อุปลพฺภนฺติ สจฺจิกฏฺฐปรมตฺเถน, อญฺญํ วิญฺญาณํ อญฺเญ สงฺขาราติ, อามนฺตา.\csromanspacer{} ปุคฺคโล อุปลพฺภติ สจฺจิกฏฺฐปรมตฺเถน, วิญฺญาณญฺจ อุปลพฺภติ สจฺจิกฏฺฐปรมตฺเถนาติ, อามนฺตา.\csromanspacer{} อญฺญํ วิญฺญาณํ อญฺโญ ปุคฺคโลติ, น เหวํ วตฺตพฺเพ.}

\prose{อาชานาหิ นิคฺคหํ\csromandash{}หญฺจิ วิญฺญาณํ อุปลพฺภติ สจฺจิกฏฺฐปรมตฺเถน, สงฺขารา จ อุปลพฺภนฺติ สจฺจิกฏฺฐปรมตฺเถน, อญฺญํ วิญฺญาณํ อญฺเญ สงฺขารา.\csromanspacer{} ปุคฺคโล อุปลพฺภติ สจฺจิกฏฺฐปรมตฺเถน, วิญฺญาณญฺจ อุปลพฺภติ สจฺจิกฏฺฐปรมตฺเถน.\csromanspacer{} เตน วต เร วตฺตพฺเพ “อญฺญํ วิญฺญาณํ อญฺโญ ปุคฺคโล”ติ.\csromanspacer{} ยํ ตตฺถ วเทสิ “วตฺตพฺเพ โข ‘วิญฺญาณํ อุปลพฺภติ สจฺจิกฏฺฐปรมตฺเถน, สงฺขารา จ อุปลพฺภนฺติ สจฺจิกฏฺฐปรมตฺเถน, อญฺญํ วิญฺญาณํ อญฺเญ สงฺขารา.\csromanspacer{} ปุคฺคโล อุปลพฺภติ สจฺจิกฏฺฐปรมตฺเถน, วิญฺญาณญฺจ อุปลพฺภติ สจฺจิกฏฺฐปรมตฺเถน’, โน จ วตฺตพฺเพ ‘อญฺญํ วิญฺญานํ อญฺโญ ปุคฺคโล’ติ”, มิจฺฉา.}

\prose{โน เจ ปน วตฺตพฺเพ “อญฺญํ วิญฺญาณํ อญฺโญ ปุคฺคโล”ติ, โน จ วต เร วตฺตพฺเพ “วิญฺญาณํ อุปลพฺภติ สจฺจิกฏฺฐปรมตฺเถน, สงฺขารา จ อุปลพฺภนฺติ สจฺจิกฏฺฐปรมตฺเถน, อญฺญํ วิญฺญาณํ อญฺเญ สงฺขารา.\csromanspacer{} ปุคฺคโล อุปลพฺภติ สจฺจิกฏฺฐปรมตฺเถน, วิญฺญาณญฺจ อุปลพฺภติ สจฺจิกฏฺฐปรมตฺเถนา”ติ.\csromanspacer{} ยํ ตตฺถ \csromanfolio{18}วเทสิ “วตฺตพฺเพ โข ‘วิญฺญาณํ อุปลพฺภติ สจฺจิกฏฺฐปรมตฺเถน, สงฺขารา จ อุปลพฺภนฺติ สจฺจิกฏฺฐปรมตฺเถน, อญฺญํ วิญฺญาณํ อญฺเญ สงฺขารา.\csromanspacer{} ปุคฺคโล อุปลพฺภติ สจฺจิกฏฺฐปรมตฺเถน, วิญฺญาณญฺจ อุปลพฺภติ สจฺจิกฏฺฐปรมตฺเถน’, โน จ วตฺตพฺเพ ‘อญฺญํ วิญฺญาณํ อญฺโญ ปุคฺคโล’ติ”, มิจฺฉา ฯเปฯ.}

\proseitem{34}{\csromansymbolmark{*}จกฺขายตนํ อุปลพฺภติ สจฺจิกฏฺฐปรมตฺเถน, โสตายตนญฺจ อุปลพฺภติ ฯเปฯ ธมฺมายตนญฺจ อุปลพฺภติ สจฺจิกฏฺฐปรมตฺเถน ฯเปฯ.\csromanspacer{} โสตายตนํ อุปลพฺภติ ฯเปฯ.\csromanspacer{} ธมฺมายตนํ อุปลพฺภติ สจฺจิกฏฺฐปรมตฺเถน, จกฺขายตนญฺจ อุปลพฺภติ ฯเปฯ มนายตนญฺจ อุปลพฺภติ สจฺจิกฏฺฐปรมตฺเถน ฯเปฯ.}

\proseitem{35}{\csromansymbolmark{*}จกฺขุธาตุ อุปลพฺภติ สจฺจิกฏฺฐปรมตฺเถน, โสตธาตุ จ อุปลพฺภติ ฯเปฯ ธมฺมธาตุ จ อุปลพฺภติ สจฺจิกฏฺฐปรมตฺเถน ฯเปฯ.\csromanspacer{} โสตธาตุ อุปลพฺภติ ฯเปฯ.\csromanspacer{} ธมฺมธาตุ อุปลพฺภติ สจฺจิกฏฺฐปรมตฺเถน, จกฺขุธาตุ จ อุปลพฺภติ ฯเปฯ มโนวิญฺญาณธาตุ จ อุปลพฺภติ สจฺจิกฏฺฐปรมตฺเถน ฯเปฯ.}

\proseitem{36}{\csromansymbolmark{*}จกฺขุนฺทฺริยํ อุปลพฺภติ สจฺจิกฏฺฐปรมตฺเถน, โสตินฺทฺริยญฺจ อุปลพฺภติ ฯเปฯ อญฺญาตาวินฺทฺริยญฺจ อุปลพฺภติ สจฺจิกฏฺฐปรมตฺเถน ฯเปฯ.\csromanspacer{} โสตินฺทฺริยํ อุปลพฺภติ ฯเปฯ.\csromanspacer{} อญฺญาตาวินฺทฺริยํ อุปลพฺภติ สจฺจิกฏฺฐปรมตฺเถน ฯเปฯ จกฺขุนฺทฺริยญฺจ อุปลพฺภติ ฯเปฯ อญฺญินฺทฺริยญฺจ อุปลพฺภติ สจฺจิกฏฺฐปรมตฺเถน, อญฺญํ อญฺญาตาวินฺทฺริยํ อญฺญํ อญฺญินฺทฺริยนฺติ, อามนฺตา.\csromanspacer{} ปุคฺคโล อุปลพฺภติ สจฺจิกฏฺฐปรมตฺเถน, อญฺญาตาวินฺทฺริยญฺจ อุปลพฺภติ สจฺจิกฏฺฐปรมตฺเถนาติ, อามนฺตา.\csromanspacer{} อญฺญํ อญฺญาตาวินฺทฺริยํ อญฺโญ ปุคฺคโลติ, น เหวํ วตฺตพฺเพ.}

\prose{อาชานาหิ นิคฺคหํ\csromandash{}หญฺจิ อญฺญาตาวินฺทฺริยํ อุปลพฺภติ สจฺจิกฏฺฐปรมตฺเถน, อญฺญินฺทฺริยญฺจ อุปลพฺภติ สจฺจิกฏฺฐปรมตฺเถน, อญฺญํ อญฺญาตาวินฺทฺริยํ อญฺญํ อญฺญินฺทฺริยํ.\csromanspacer{} ปุคฺคโล อุปลพฺภติ สจฺจิกฏฺฐปรมตฺเถน, อญฺญาตาวินฺทฺริยญฺจ อุปลพฺภติ สจฺจิกฏฺฐปรมตฺเถน.\csromanspacer{} เตน วต เร วตฺตพฺเพ “อญฺญํ อญฺญาตาวินฺทฺริยํ อญฺโญ ปุคฺคโล”ติ.\csromanspacer{} ยํ ตตฺถ วเทสิ “วตฺตพฺเพ โข ‘อญฺญาตาวินฺทฺริยํ อุปลพฺภติ สจฺจิกฏฺฐปรมตฺเถน, อญฺญินฺทฺริยญฺจ อุปลพฺภติ สจฺจิกฏฺฐปรมตฺเถน, อญฺญํ อญฺญาตาวินฺทฺริยํ อญฺญํ อญฺญินฺทฺริยํ.\csromanspacer{} ปุคฺคโล อุปลพฺภติ สจฺจิกฏฺฐปรมตฺเถน, อญฺญาตาวินฺทฺริยญฺจ อุปลพฺภติ สจฺจิกฏฺฐปรมตฺเถน’, โน จ วตฺตพฺเพ ‘อญฺญํ อญฺญาตาวินฺทฺริยํ อญฺโญ ปุคฺคโล’ติ”, มิจฺฉา.}

\prose{\csromanfolio{19}โน เจ ปน วตฺตพฺเพ “อญฺญํ อญฺญาตาวินฺทฺริยํ อญฺโญ ปุคฺคโล”ติ, โน จ วต เร วตฺตพฺเพ “อญฺญาตาวินฺทฺริยํ อุปลพฺภติ สจฺจิกฏฺฐปรมตฺเถน, อญฺญินฺทฺริยญฺจ อุปลพฺภติ สจฺจิกฏฺฐปรมตฺเถน, อญฺญํ อญฺญาตาวินฺทฺริยํ อญฺญํ อญฺญินฺทฺริยํ.\csromanspacer{} ปุคฺคโล อุปลพฺภติ สจฺจิกฏฺฐปรมตฺเถน, อญฺญาตาวินฺทฺริยญฺจ อุปลพฺภติ สจฺจิกฏฺฐปรมตฺเถนา”ติ.\csromanspacer{} ยํ ตตฺถ วเทสิ “วตฺตพฺเพ โข ‘อญฺญาตาวินฺทฺริยํ อุปลพฺภติ สจฺจิกฏฺฐปรมตฺเถน, อญฺญินฺทฺริยญฺจ อุปลพฺภติ สจฺจิกฏฺฐปรมตฺเถน, อญฺญํ อญฺญาตาวินฺทฺริยํ อญฺญํ อญฺญินฺทฺริยํ.\csromanspacer{} ปุคฺคโล อุปลพฺภติ สจฺจิกฏฺฐปรมตฺเถน, อญฺญาตาวินฺทฺริยญฺจ อุปลพฺภติ สจฺจิกฏฺฐปรมตฺเถน’, โน จ วตฺตพฺเพ ‘อญฺญํ อญฺญาตาวินฺทฺริยํ อญฺโญ ปุคฺคโล’ติ”, มิจฺฉา ฯเปฯ.}

\proseitem{37}{\csromansymbolmark{+}รูปํ อุปลพฺภติ สจฺจิกฏฺฐปรมตฺเถน, เวทนา จ อุปลพฺภติ สจฺจิกฏฺฐปรมตฺเถน, อญฺญํ รูปํ อญฺญา เวทนาติ, อามนฺตา.\csromanspacer{} วุตฺตํ ภควตา “อตฺถิ ปุคฺคโล อตฺตหิตาย ปฏิปนฺโน”, รูปญฺจ อุปลพฺภติ สจฺจิกฏฺฐปรมตฺเถนาติ, อามนฺตา.\csromanspacer{} อญฺญํ รูปํ อญฺโญ ปุคฺคโลติ, น เหวํ วตฺตพฺเพ.}

\prose{อาชานาหิ ปฏิกมฺมํ\csromandash{}หญฺจิ รูปํ อุปลพฺภติ สจฺจิกฏฺฐปรมตฺเถน, เวทนา จ อุปลพฺภติ สจฺจิกฏฺฐปรมตฺเถน, อญฺญํ รูปํ อญฺญา เวทนา.\csromanspacer{} วุตฺตํ ภควตา\csromandash{}อตฺถิ ปุคฺคโล อตฺตหิตาย ปฏิปนฺโน, รูปญฺจ อุปลพฺภติ สจฺจิกฏฺฐปรมตฺเถน, เตน วต เร วตฺตพฺเพ “อญฺญํ รูปํ อญฺโญ ปุคฺคโล”ติ.\csromanspacer{} ยํ ตตฺถ วเทสิ “วตฺตพฺเพ โข ‘รูปํ อุปลพฺภติ สจฺจิกฏฺฐปรมตฺเถน, เวทนา จ อุปลพฺภติ สจฺจิกฏฺฐปรมตฺเถน, อญฺญํ รูปํ อญฺญา เวทนา.\csromanspacer{} วุตฺตํ ภควตา\csromandash{} อตฺถิ ปุคฺคโล อตฺตหิตาย ปฏิปนฺโน, รูปญฺจ อุปลพฺภติ สจฺจิกฏฺฐปรมตฺเถน’, โน จ วตฺตพฺเพ ‘อญฺญํ รูปํ อญฺโญ ปุคฺคโล’ติ”,}

\prose{โน เจ ปน วตฺตพฺเพ “อญฺญํ รูปํ อญฺโญ ปุคฺคโล”ติ, โน จ วต เร วตฺตพฺเพ “รูปํ อุปลพฺภติ สจฺจิกฏฺฐปรมตฺเถน, เวทนา จ อุปลพฺภติ สจฺจิกฏฺฐปรมตฺเถน, อญฺญํ รูปํ อญฺญา เวทนา.\csromanspacer{} วุตฺตํ ภควตา\csromandash{}อตฺถิ ปุคฺคโล อตฺตหิตาย ปฏิปนฺโน, รูปญฺจ อุปลพฺภติ สจฺจิกฏฺฐปรมตฺเถนา”ติ.\csromanspacer{} ยํ ตตฺถ วเทสิ “วตฺตพฺเพ โข ‘รูปํ อุปลพฺภติ สจฺฉิกฏฺฐปรมตฺเถน, เวทนา จ อุปลพฺภติ สจฺจิกฏฺฐปรมตฺเถน, อญฺญํ รูปํ อญฺญา เวทนา.\csromanspacer{} วุตฺตํ ภควตา\csromandash{} อตฺถิ ปุคฺคโล อตฺตหิตาย ปฏิปนฺโน, รูปญฺจ อุปลพฺภติ สจฺจิกฏฺฐปรมตฺเถน’, โน จ วตฺตพฺเพ ‘อญฺญํ รูปํ อญฺโญ ปุคฺคโล’ติ”, มิจฺฉา ฯเปฯ.}

\proseitem{38}{\csromanfolio{20}\csromansymbolmark{+}รูปํ อุปลพฺภติ สจฺจิกฏฺฐปรมตฺเถน, สญฺญา จ อุปลพฺภติ.\csromanspacer{} สงฺขารา จ อุปลพฺภนฺติ.\csromanspacer{} วิญฺญาณญฺจ อุปลพฺภติ สจฺจิกฏฺฐปรมตฺเถน ฯเปฯ.}

\proseitem{39}{\csromansymbolmark{+}เวทนา อุปลพฺภติ สจฺจิกฏฺฐปรมตฺเถน, สญฺญา จ อุปลพฺภติ.\csromanspacer{} สงฺขารา จ อุปลพฺภนฺติ.\csromanspacer{} วิญฺญาณญฺจ อุปลพฺภติ.\csromanspacer{} รูปญฺจ อุปลพฺภติ สจฺจิกฏฺฐปรมตฺเถน ฯเปฯ.}

\proseitem{40}{\csromansymbolmark{+}สญฺญา อุปลพฺภติ สจฺจิกฏฺฐปรมตฺเถน, สงฺขารา จ อุปลพฺภนฺติ.\csromanspacer{} วิญฺญาณญฺจ อุปลพฺภติ.\csromanspacer{} รูปญฺจ อุปลพฺภติ.\csromanspacer{} เวทนา จ อุปลพฺภติ สจฺจิกฏฺฐปรมตฺเถน ฯเปฯ.}

\proseitem{41}{\csromansymbolmark{+}สงฺขารา อุปลพฺภนฺติ สจฺจิกฏฺฐปรมตฺเถน, วิญฺญาณญฺจ อุปลพฺภติ.\csromanspacer{} รูปญฺจ อุปลพฺภติ.\csromanspacer{} เวทนา จ อุปลพฺภติ.\csromanspacer{} สญฺญา จ อุปลพฺภติ สจฺจิกฏฺฐปรมตฺเถน ฯเปฯ.}

\proseitem{42}{\csromansymbolmark{+}วิญฺญาณํ อุปลพฺภติ สจฺจิกฏฺฐปรมตฺเถน, รูปญฺจ อุปลพฺภติ.\csromanspacer{} เวทนา จ อุปลพฺภติ.\csromanspacer{} สญฺญา จ อุปลพฺภติ.\csromanspacer{} สงฺขารา จ อุปลพฺภนฺติ สจฺจิกฏฺฐปรมตฺเถน ฯเปฯ.}

\proseitem{43}{\csromansymbolmark{+}จกฺขายตนํ อุปลพฺภติ สจฺจิกฏฺฐปรมตฺเถน, โสตายตนญฺจ อุปลพฺภติ ฯเปฯ ธมฺมายตนญฺจ อุปลพฺภติ สจฺจิกฏฺฐปรมตฺเถน ฯเปฯ.\csromanspacer{} โสตายตนํ อุปลพฺภติ ฯเปฯ.\csromanspacer{} ธมฺมายตนํ อุปลพฺภติ สจฺจิกฏฺฐปรมตฺเถน, จกฺขายตนญฺจ อุปลพฺภติ ฯเปฯ มนายตนญฺจ อุปลพฺภติ สจฺจิกฏฺฐปรมตฺเถน ฯเปฯ.}

\proseitem{44}{\csromansymbolmark{+}จกฺขุธาตุ อุปลพฺภติ สจฺจิกฏฺฐปรมตฺเถน, โสตธาตุ จ อุปลพฺภติ ฯเปฯ ธมฺมธาตุ จ อุปลพฺภติ สจฺจิกฏฺฐปรมตฺเถน ฯเปฯ.\csromanspacer{} โสตธาตุ อุปลพฺภติ สจฺจิกฏฺฐปรมตฺเถน ฯเปฯ.\csromanspacer{} ธมฺมธาตุ อุปลพฺภติ สจฺจิกฏฺฐปรมตฺเถน, จกฺขุธาตุ จ อุปลพฺภติ ฯเปฯ มโนวิญฺญาณธาตุ จ อุปลพฺภติ สจฺจิกฏฺฐปรมตฺเถน ฯเปฯ.}

\proseitem{45}{\csromansymbolmark{+}จกฺขุนฺทฺริยํ อุปลพฺภติ สจฺจิกฏฺฐปรมตฺเถน, โสตินฺทฺริยญฺจ อุปลพฺภติ ฯเปฯ อญฺญาตาวินฺทฺริยญฺจ อุปลพฺภติ สจฺจิกฏฺฐปรมตฺเถน ฯเปฯ.\csromanspacer{} โสตินฺทฺริยํ อุปลพฺภติ ฯเปฯ.\csromanspacer{} อญฺญาตาวินฺทฺริยํ อุปลพฺภติ สจฺจิกฏฺฐปรมตฺเถน, จกฺขุนฺทฺริยญฺจ อุปลพฺภติ ฯเปฯ อญฺญินฺทฺริยญฺจ อุปลพฺภติ สจฺจิกฏฺฐปรมตฺเถน, อญฺญํ \csromanfolio{21}อญฺญาตาวินฺทฺริยํ อญฺญํ อญฺญินฺทฺริยนฺติ, อามนฺตา.\csromanspacer{} วุตฺตํ ภควตา “อตฺถิ ปุคฺคโล อตฺตหิตาย ปฏิปนฺโน”, อญฺญาตาวินฺทฺริยญฺจ อุปลพฺภติ สจฺจิกฏฺฐปรมตฺเถนาติ, อามนฺตา.\csromanspacer{} อญฺญํ อญฺญาตาวินฺทฺริยํ อญฺโญ ปุคฺคโลติ, น เหวํ วตฺตพฺเพ.}

\prose{อาชานาหิ ปฏิกมฺมํ\csromandash{}หญฺจิ อญฺญาตาวินฺทฺริยํ อุปลพฺภติ สจฺจิกฏฺฐปรมตฺเถน, อญฺญินฺทฺริยญฺจ อุปลพฺภติ สจฺจิกฏฺฐปรมตฺเถน, อญฺญํ อญฺญาตาวินฺทฺริยํ อญฺญํ อญฺญินฺทฺริยํ.\csromanspacer{} วุตฺตํ ภควตา\csromandash{}อตฺถิ ปุคฺคโล อตฺตหิตาย ปฏิปนฺโน, อญฺญาตาวินฺทฺริยญฺจ อุปลพฺภติ สจฺจิกฏฺฐปรมตฺเถน, เตน วต เร วตฺตพฺเพ “อญฺญํ อญฺญาตาวินฺทฺริยํ อญฺโญ ปุคฺคโล”ติ.\csromanspacer{} ยํ ตตฺถ วเทสิ “วตฺตพฺเพ โข ‘อญฺญาตาวินฺทฺริยํ อุปลพฺภติ สจฺจิกฏฺฐปรมตฺเถน, อญฺญินฺทฺริยญฺจ อุปลพฺภติ สจฺจิกฏฺฐปรมตฺเถน, อญฺญํ อญฺญาตาวินฺทฺริยํ อญฺญํ อญฺญินฺทฺริยํ.\csromanspacer{} วุตฺตํ ภควตา\csromandash{}อตฺถิ ปุคฺคโล อตฺตหิตาย ปฏิปนฺโน, อญฺญาตาวินฺทฺริยญฺจ อุปลพฺภติ สจฺจิกฏฺฐปรมตฺเถน’, โน จ วตฺตพฺเพ ‘อญฺญํ อญฺญาตาวินฺทฺริยํ อญฺโญ ปุคฺคโล’ติ”, มิจฺฉา.}

\prose{โน เจ ปน วตฺตพฺเพ “อญฺญํ อญฺญาตาวินฺทฺริยํ อญฺโญ ปุคฺคโล”ติ, โน จ วต เร วตฺตพฺเพ “อญฺญาตาวินฺทฺริยํ อุปลพฺภติ สจฺจิกฏฺฐปรมตฺเถน, อญฺญินฺทฺริยญฺจ อุปลพฺภติ สจฺจิกฏฺฐปรมตฺเถน, อญฺญํ อญฺญาตาวินฺทฺริยํ อญฺญํ อญฺญินฺทฺริยํ.\csromanspacer{} วุตฺตํ ภควตา\csromandash{}อตฺถิ ปุคฺคโล อตฺตหิตาย ปฏิปนฺโน, อญฺญาตาวินฺทฺริยญฺจ อุปลพฺภติ สจฺจิกฏฺฐปรมตฺเถนา”ติ.\csromanspacer{} ยํ ตตฺถ วเทสิ “วตฺตพฺเพ โข ‘อญฺญาตาวินฺทฺริยํ อุปลพฺภติ สจฺจิกฏฺฐปรมตฺเถน, อญฺญินฺทฺริยญฺจ อุปลพฺภติ สจฺจิกฏฺฐปรมตฺเถน, อญฺญํ อญฺญาตาวินฺทฺริยํ อญฺญํ อญฺญินฺทฺริยํ.\csromanspacer{} วุตฺตํ ภควตา\csromandash{}อตฺถิ ปุคฺคโล อตฺตหิตาย ปฏิปนฺโน, อญฺญาตาวินฺทฺริยญฺจ อุปลพฺภติ สจฺจิกฏฺฐปรมตฺเถน’, โน จ วตฺตพฺเพ ‘อญฺญํ อญฺญาตาวินฺทฺริยํ อญฺโญ ปุคฺคโล’ติ”, มิจฺฉา ฯเปฯ.}

\vspace{\tipitakaparskip}
\csromancenter{โอปมฺมสํสนฺทนํ.}
\csromanmatikaanchor{mtk.8}
\csromantocmarkref{section}{7. จตุกฺกนยสํสนฺทน}{mtk.8}
\bookmark[dest={mtk.8},level=1]{7. จตุกฺกนยสํสนฺทน}
\csromanheader{1}{7. จตุกฺกนยสํสนฺทน}
\proseitem{46}{\csromansymbolmark{*}ปุคฺคโล อุปลพฺภติ สจฺจิกฏฺฐปรมตฺเถนาติ, อามนฺตา.\csromanspacer{} รูปํ ปุคฺคโลติ, น เหวํ วตฺตพฺเพ.}

\prose{\csromanfolio{22}อาชานาหิ นิคฺคหํ\csromandash{}หญฺจิ ปุคฺคโล อุปลพฺภติ สจฺจิกฏฺฐปรมตฺเถน, เตน วต เร วตฺตพฺเพ “รูปํ ปุคฺคโล”ติ.\csromanspacer{} ยํ ตตฺถ วเทสิ “วตฺตพฺเพ โข ‘ปุคฺคโล อุปลพฺภติ สจฺจิกฏฺฐปรมตฺเถน’, โน จ วตฺตพฺเพ ‘รูปํ ปุคฺคโล’ติ”, มิจฺฉา.}

\prose{โน เจ ปน วตฺตพฺเพ “รูปํ ปุคฺคโล”ติ, โน จ วต เร วตฺตพฺเพ “ปุคฺคโล อุปลพฺภติ สจฺจิกฏฺฐปรมตฺเถนา”ติ.\csromanspacer{} ยํ ตตฺถ วเทสิ “วตฺตพฺเพ โข ‘ปุคฺคโล อุปลพฺภติ สจฺจิกฏฺฐปรมตฺเถน’, โน จ วตฺตพฺเพ ‘รูปํ ปุคฺคโล’ติ”, มิจฺฉา ฯเปฯ.}

\proseitem{47}{\csromansymbolmark{*}ปุคฺคโล อุปลพฺภติ สจฺจิกฏฺฐปรมตฺเถนาติ, อามนฺตา.\csromanspacer{} รูปสฺมิํ ปุคฺคโล ฯเปฯ.\csromanspacer{} อญฺญตฺร รูปา ปุคฺคโล ฯเปฯ.\csromanspacer{} ปุคฺคลสฺมิํ รูปนฺติ, น เหวํ วตฺตพฺเพ.}

\prose{อาชานาหิ นิคฺคหํ\csromandash{}หญฺจิ ปุคฺคโล อุปลพฺภติ สจฺจิกฏฺฐปรมตฺเถน, เตน วต เร วตฺตพฺเพ “ปุคฺคลสฺมิํ รูปนฺ”ติ.\csromanspacer{} ยํ ตตฺถ วเทสิ “วตฺตพฺเพ โข ‘ปุคฺคโล อุปลพฺภติ สจฺจิกฏฺฐปรมตฺเถน’, โน จ วตฺตพฺเพ ‘ปุคฺคลสฺมิํ รูปนฺ’ติ”, มิจฺฉา.}

\prose{โน เจ ปน วตฺตพฺเพ “ปุคฺคลสฺมิํ รูปนฺ”ติ, โน จ วต เร วตฺตพฺเพ “ปุคฺคโล อุปลพฺภติ สจฺจิกฏฺฐปรมตฺเถนา”ติ.\csromanspacer{} ยํ ตตฺถ วเทสิ “วตฺตพฺเพ โข ‘ปุคฺคโล อุปลพฺภติ สจฺจิกฏฺฐปรมตฺเถน’, โน จ วตฺตพฺเพ ‘ปุคฺคลสฺมิํ รูปนฺ’ติ”, มิจฺฉา ฯเปฯ.}

\proseitem{48}{\csromansymbolmark{*}ปุคฺคโล อุปลพฺภติ สจฺจิกฏฺฐปรมตฺเถนาติ, อามนฺตา.\csromanspacer{} เวทนา ปุคฺคโล ฯเปฯ.\csromanspacer{} เวทนาย ปุคฺคโล ฯเปฯ.\csromanspacer{} อญฺญตฺร เวทนาย ปุคฺคโล ฯเปฯ.\csromanspacer{} ปุคฺคลสฺมิํ เวทนา ฯเปฯ.}

\prose{สญฺญา ปุคฺคโล ฯเปฯ.\csromanspacer{} สญฺญาย ปุคฺคโล ฯเปฯ.\csromanspacer{} อญฺญตฺร สญฺญาย ปุคฺคโล ฯเปฯ ปุคฺคลสฺมิํ สญฺญา ฯเปฯ.}

\prose{สงฺขารา ปุคฺคโล ฯเปฯ.\csromanspacer{} สงฺขาเรสุ ปุคฺคโล ฯเปฯ.\csromanspacer{} อญฺญตฺร สงฺขาเรหิ ปุคฺคโล ฯเปฯ ปุคฺคลสฺมิํ สงฺขารา ฯเปฯ.}

\prose{วิญฺญาณํ ปุคฺคโล ฯเปฯ.\csromanspacer{} วิญฺญาณสฺมิํ ปุคฺคโล ฯเปฯ.\csromanspacer{} อญฺญตฺร วิญฺญาณา ปุคฺคโล ฯเปฯ ปุคฺคลสฺมิํ วิญฺญาณนฺติ, น เหวํ วตฺตพฺเพ.}

\prose{อาชานาหิ นิคฺคหํ\csromandash{}หญฺจิ ปุคฺคโล อุปลพฺภติ สจฺจิกฏฺฐปรมตฺเถน, เตน วต เร วตฺตพฺเพ “ปุคฺคลสฺมิํ วิญฺญาณนฺ”ติ.\csromanspacer{} ยํ ตตฺถ วเทสิ “วตฺตพฺเพ \csromanfolio{23}โข ‘ปุคฺคโล อุปลพฺภติ สจฺจิกฏฺฐปรมตฺเถน’, โน จ วตฺตพฺเพ ‘ปุคฺคลสฺมิํ วิญฺญาณนฺ’ติ”, มิจฺฉา.}

\prose{โน เจ ปน วตฺตพฺเพ “ปุคฺคลสฺมิํ วิญฺญาณนฺ”ติ, โน จ วต เร วตฺตพฺเพ “ปุคฺคโล อุปลพฺภติ สจฺจิกฏฺฐปรมตฺเถนา”ติ.\csromanspacer{} ยํ ตตฺถ วเทสิ “วตฺตพฺเพ โข ‘ปุคฺคโล อุปลพฺภติ สจฺจิกฏฺฐปรมตฺเถน’, โน จ วตฺตพฺเพ ‘ปุคฺคลสฺมิํ วิญฺญาณนฺ’ติ”, มิจฺฉา ฯเปฯ.}

\proseitem{49}{\csromansymbolmark{*}ปุคฺคโล อุปลพฺภติ สจฺจิกฏฺฐปรมตฺเถนาติ, อามนฺตา.\csromanspacer{} จกฺขายตนํ ปุคฺคโล ฯเปฯ.\csromanspacer{} จกฺขายตนสฺมิํ ปุคฺคโล ฯเปฯ.\csromanspacer{} อญฺญตฺร จกฺขายตนา ปุคฺคโล ฯเปฯ ปุคฺคลสฺมิํ จกฺขายตนํ ฯเปฯ.\csromanspacer{} ธมฺมายตนํ ปุคฺคโล ฯเปฯ.ธมฺมายตนสฺมิํ ปุคฺคโล ฯเปฯ.\csromanspacer{} อญฺญตฺร ธมฺมายตนา ปุคฺคโล ฯเปฯ.\csromanspacer{} ปุคฺคลสฺมิํ ธมฺมายตนํ ฯเปฯ.}

\prose{จกฺขุธาตุ ปุคฺคโล ฯเปฯ.\csromanspacer{} จกฺขุธาตุยา ปุคฺคโล ฯเปฯ.\csromanspacer{} อญฺญตฺร จกฺขุธาตุยา ปุคฺคโล ฯเปฯ.\csromanspacer{} ปุคฺคลสฺมิํ จกฺขุธาตุ ฯเปฯ.\csromanspacer{} ธมฺมธาตุ ปุคฺคโล ฯเปฯ.\csromanspacer{} ธมฺมธาตุยา ปุคฺคโล ฯเปฯ.\csromanspacer{} อญฺญตฺร ธมฺมธาตุยา ปุคฺคโล ฯเปฯ.\csromanspacer{} ปุคฺคลสฺมิํ ธมฺมธาตุ ฯเปฯ.}

\prose{จกฺขุนฺทฺริยํ ปุคฺคโล ฯเปฯ.\csromanspacer{} จกฺขุนฺทฺริยสฺมิํ ปุคฺคโล ฯเปฯ.\csromanspacer{} อญฺญตฺร จกฺขุนฺทฺริยา ปุคฺคโล ฯเปฯ.\csromanspacer{} ปุคฺคลสฺมิํ จกฺขุนฺทฺริยํ ฯเปฯ.\csromanspacer{} อญฺญาตาวินฺทฺริยํ ปุคฺคโล ฯเปฯ.\csromanspacer{} อญฺญาตาวินฺทฺริยสฺมิํ ปุคฺคโล ฯเปฯ.\csromanspacer{} อญฺญตฺร อญฺญาตาวินฺทฺริยา ปุคฺคโล ฯเปฯ.\csromanspacer{} ปุคฺคลสฺมิํ อญฺญาตาวินฺทฺริยนฺติ, น เหวํ วตฺตพฺเพ.}

\prose{อาชานาหิ นิคฺคหํ\csromandash{}หญฺจิ ปุคฺคโล อุปลพฺภติ สจฺจิกฏฺฐปรมตฺเถน, เตน วต เร วตฺตพฺเพ “ปุคฺคลสฺมิํ อญฺญาตาวินฺทฺริยนฺ”ติ.\csromanspacer{} ยํ ตตฺถ วเทสิ “วตฺตพฺเพ โข ‘ปุคฺคโล อุปลพฺภติ สจฺจิกฏฺฐปรมตฺเถน’, โน จ วตฺตพฺเพ ‘ปุคฺคลสฺมิํ อญฺญาตาวินฺทฺริยนฺ’ติ”,}

\prose{โน เจ ปน วตฺตพฺเพ “ปุคฺคลสฺมิํ อญฺญาตาวินฺทฺริยนฺ”ติ, โน จ วต เร วตฺตพฺเพ “ปุคฺคโล อุปลพฺภติ สจฺจิกฏฺฐปรมตฺเถนา”ติ.\csromanspacer{} ยํ ตตฺถ วเทสิ “วตฺตพฺเพ โข ‘ปุคฺคโล อุปลพฺภติ สจฺจิกฏฺฐปรมตฺเถน’, โน จ วตฺตพฺเพ ‘ปุคฺคลสฺมิํ อญฺญาตาวินฺทฺริยนฺ’ติ “, มิจฺฉา ฯเปฯ.}

\proseitem{50}{\csromansymbolmark{+}ปุคฺคโล นุปลพฺภติ สจฺจิกฏฺฐปรมตฺเถนาติ, อามนฺตา.\csromanspacer{} วุตฺตํ ภควตา “อตฺถิ ปุคฺคโล อตฺตหิตาย ปฏิปนฺโน”ติ, อามนฺตา.\csromanspacer{} รูปํ ปุคฺคโลติ, น เหวํ วตฺตพฺเพ.}

\prose{\csromanfolio{24}อาชานาหิ ปฏิกมฺมํ\csromandash{}หญฺจิ วุตฺตํ ภควตา\csromandash{}อตฺถิ ปุคฺคโล อตฺตหิตาย ปฏิปนฺโน, เตน วต เร วตฺตพฺเพ “รูปํ ปุคฺคโล”ติ.\csromanspacer{} ยํ ตตฺถ วเทสิ “วตฺตพฺเพ โข ‘วุตฺตํ ภควตา\csromandash{}อตฺถิ ปุคฺคโล อตฺตหิตาย ปฏิปนฺโน’, โน จ วตฺตพฺเพ ‘รูปํ ปุคฺคโล’ติ”, มิจฺฉา.}

\prose{โน เจ ปน วตฺตพฺเพ “รูปํ ปุคฺคโล”ติ, โน จ วต เร วตฺตพฺเพ “วุตฺตํ ภควตา\csromandash{}อตฺถิ ปุคฺคโล อตฺตหิตาย ปฏิปนฺโน”ติ.\csromanspacer{} ยํ ตตฺถ วเทสิ ‘วตฺตพฺเพ โข ‘วุตฺตํ ภควตา\csromandash{}อตฺถิ ปุคฺคโล อตฺตหิตาย ปฏิปนฺโน’, โน จ วตฺตพฺเพ ‘รูปํ ปุคฺคโล’ติ”, มิจฺฉา ฯเปฯ.}

\proseitem{51}{\csromansymbolmark{+}ปุคฺคโล นุปลพฺภติ สจฺจิกฏฺฐปรมตฺเถนาติ, อามนฺตา.\csromanspacer{} วุตฺตํ ภควตา “อตฺถิ ปุคฺคโล อตฺตหิตาย ปฏิปนฺโน”ติ, อามนฺตา.\csromanspacer{} รูปสฺมิํ ปุคฺคโล ฯเปฯ.\csromanspacer{} อญฺญตฺร รูปา ปุคฺคโล ฯเปฯ.\csromanspacer{} ปุคฺคลสฺมิํ รูปํ ฯเปฯ.}

\prose{เวทนา ปุคฺคโล ฯเปฯ.\csromanspacer{} เวทนาย ปุคฺคโล ฯเปฯ.\csromanspacer{} อญฺญตฺร เวทนาย ปุคฺคโล ฯเปฯ.ปุคฺคลสฺมิํ เวทนา ฯเปฯ.}

\prose{สญฺญา ปุคฺคโล ฯเปฯ.\csromanspacer{} สญฺญาย ปุคฺคโล ฯเปฯ.\csromanspacer{} อญฺญตฺร สญฺญาย ปุคฺคโล ฯเปฯ.\csromanspacer{} ปุคฺคลสฺมิํ สญฺญา ฯเปฯ.}

\prose{สงฺขารา ปุคฺคโล ฯเปฯ.\csromanspacer{} สงฺขาเรสุ ปุคฺคโล ฯเปฯ.\csromanspacer{} อญฺญตฺร สงฺขาเรหิ ปุคฺคโล ฯเปฯ.\csromanspacer{} ปุคฺคลสฺมิํ สงฺขารา ฯเปฯ.}

\prose{วิญฺญาณํ ปุคฺคโล ฯเปฯ.\csromanspacer{} วิญฺญาณสฺมิํ ปุคฺคโล ฯเปฯ.\csromanspacer{} อญฺญตฺร วิญฺญาณา ปุคฺคโล ฯเปฯ.\csromanspacer{} ปุคฺคลสฺมิํ วิญฺญาณนฺติ, น เหวํ วตฺตพฺเพ.}

\prose{อาชานาหิ ปฏิกมฺมํ\csromandash{}หญฺจิ วุตฺตํ ภควตา\csromandash{}อตฺถิ ปุคฺคโล อตฺตหิตาย ปฏิปนฺโน, เตน วต เร วตฺตพฺเพ “ปุคฺคลสฺมิํ วิญฺญาณนฺ”ติ.\csromanspacer{} ยํ ตตฺถ วเทสิ “วตฺตพฺเพ โข ‘วุตฺตํ ภควตา\csromandash{}อตฺถิ ปุคฺคโล อตฺตหิตาย ปฏิปนฺโน’, โน จ วตฺตพฺเพ ‘ปุคฺคลสฺมิํ วิญฺญาณนฺ’ติ”, มิจฺฉา.}

\prose{โน เจ ปน วตฺตพฺเพ “ปุคฺคลสฺมิํ วิญฺญาณนฺ”ติ, โน จ วต เร วตฺตพฺเพ “วุตฺตํ ภควตา\csromandash{}อตฺถิ ปุคฺคโล อตฺตหิตาย ปฏิปนฺโน”ติ.\csromanspacer{} ยํ ตตฺถ วเทสิ “วตฺตพฺเพ โข ‘วุตฺตํ ภควตา\csromandash{}อตฺถิ ปุคฺคโล อตฺตหิตาย ปฏิปนฺโน’, โน จ วตฺตพฺเพ ‘ปุคฺคลสฺมิํ วิญฺญาณนฺ’ติ”, มิจฺฉา ฯเปฯ.}

\proseitem{52}{\csromanfolio{25}\csromansymbolmark{+}ปุคฺคโล นุปลพฺภติ สจฺจิกฏฺฐปรมตฺเถนาติ, อามนฺตา.\csromanspacer{} วุตฺตํ ภควตา “อตฺถิ ปุคฺคโล อตฺตหิตาย ปฏิปนฺโน”ติ, อามนฺตา.\csromanspacer{} จกฺขายตนํ ปุคฺคโล ฯเปฯ.\csromanspacer{} จกฺขายตนสฺมิํ ปุคฺคโล ฯเปฯ.\csromanspacer{} อญฺญตฺร จกฺขายตนา ปุคฺคโล ฯเปฯ.\csromanspacer{} ปุคฺคลสฺมิํ จกฺขายตนํ ฯเปฯ.\csromanspacer{} ธมฺมายตนํ ปุคฺคโล ฯเปฯ.\csromanspacer{} ธมฺมายตนสฺมิํ ปุคฺคโล ฯเปฯ.\csromanspacer{} อญฺญตฺร ธมฺมายตนา ปุคฺคโล ฯเปฯ.\csromanspacer{} ปุคฺคลสฺมิํ ธมฺมายตนํ ฯเปฯ.}

\prose{จกฺขุธาตุ ปุคฺคโล ฯเปฯ.\csromanspacer{} จกฺขุธาตุยา ปุคฺคโล ฯเปฯ.\csromanspacer{} อญฺญตฺร จกฺขุธาตุยา ปุคฺคโล ฯเปฯ.\csromanspacer{} ปุคฺคลสฺมิํ จกฺขุธาตุ ฯเปฯ.\csromanspacer{} ธมฺมธาตุ ปุคฺคโล ฯเปฯ.\csromanspacer{} ธมฺมธาตุยา ปุคฺคโล ฯเปฯ.\csromanspacer{} อญฺญตฺร ธมฺมธาตุยา ปุคฺคโล ฯเปฯ.\csromanspacer{} ปุคฺคลสฺมิํ ธมฺมธาตุ ฯเปฯ.}

\prose{จกฺขุนฺทฺริยํ ปุคฺคโล ฯเปฯ.\csromanspacer{} จกฺขุนฺทฺริยสฺมิํ ปุคฺคโล ฯเปฯ.\csromanspacer{} อญฺญตฺร จกฺขุนฺทฺริยา ปุคฺคโล ฯเปฯ.\csromanspacer{} ปุคฺคลสฺมิํ จกฺขุนฺทฺริยํ ฯเปฯ.\csromanspacer{} อญฺญาตาวินฺทฺริยํ ปุคฺคโล ฯเปฯ.\csromanspacer{} อญฺญาตาวินฺทฺริยสฺมิํ ปุคฺคโล ฯเปฯ.\csromanspacer{} อญฺญตฺร อญฺญาตาวินฺทฺริยา ปุคฺคโล ฯเปฯ.\csromanspacer{} ปุคฺคลสฺมิํ อญฺญาตาวินฺทฺริยนฺติ, น เหวํ วตฺตพฺเพ.}

\prose{อาชานาหิ ปฏิกมฺมํ\csromandash{}หญฺจิ วุตฺตํ ภควตา\csromandash{}อตฺถิ ปุคฺคโล อตฺตหิตาย ปฏิปนฺโน, เตน วต เร วตฺตพฺเพ “ปุคฺคลสฺมิํ อญฺญาตาวินฺทฺริยนฺ”ติ.\csromanspacer{} ยํ ตตฺถ วเทสิ “วตฺตพฺเพ โข ‘วุตฺตํ ภควตา\csromandash{} อตฺถิ ปุคฺคโล อตฺตหิตาย ปฏิปนฺโน’, โน จ วตฺตพฺเพ ‘ปุคฺคลสฺมิํ อญฺญาตาวินฺทฺริยนฺ’ติ, มิจฺฉา.}

\prose{โน เจ ปน วตฺตพฺเพ “ปุคฺคลสฺมิํ อญฺญาตาวินฺทฺริยนฺ”ติ, โน จ วต เร วตฺตพฺเพ “วุตฺตํ ภควตา\csromandash{}อตฺถิ ปุคฺคโล อตฺตหิตาย ปฏิปนฺโน”ติ.\csromanspacer{} ยํ ตตฺถ วเทสิ “วตฺตพฺเพ โข ‘วุตฺตํ ภควตา\csromandash{}อตฺถิ ปุคฺคโล อตฺตหิตาย ปฏิปนฺโน’, โน จ วตฺตพฺเพ ‘ปุคฺคลสฺมิํ อญฺญาตาวินฺทฺริยนฺ’ติ”, มิจฺฉา ฯเปฯ.}

\vspace{\tipitakaparskip}
\csromancenter{จตุกฺกนยสํสนฺทนํ.}
\csromanmatikaanchor{mtk.9}
\csromantocmarkref{section}{8. ลกฺขณยุตฺติ}{mtk.9}
\bookmark[dest={mtk.9},level=1]{8. ลกฺขณยุตฺติ}
\csromanheader{1}{8. ลกฺขณยุตฺติ}
\proseitem{53}{\csromansymbolmark{*}ปุคฺคโล อุปลพฺภติ สจฺจิกฏฺฐปรมตฺเถนาติ, อามนฺตา.\csromanspacer{} ปุคฺคโล สปฺปจฺจโย ฯเปฯ.\csromanspacer{} ปุคฺคโล อปฺปจฺจโย.\csromanspacer{} ปุคฺคโล สงฺขโต.\csromanspacer{} ปุคฺคโล อสงฺขโต.\csromanspacer{} ปุคฺคโล สสฺสโต.\csromanspacer{} ปุคฺคโล อสสฺสโต.\csromanspacer{} ปุคฺคโล สนิมิตฺโต.\csromanspacer{} ปุคฺคโล อนิมิตฺโตติ, น เหวํ วตฺตพฺเพ.\csromanspacer{}\csromanspacer{}(สํขิตฺตํ.)}

\proseitem{54}{\csromanfolio{26}\csromansymbolmark{+}ปุคฺคโล นุปลพฺภติ สจฺจิกฏฺฐปรมตฺเถนาติ, อามนฺตา.\csromanspacer{} วุตฺตํ ภควตา “อตฺถิ ปุคฺคโล อตฺตหิตาย ปฏิปนฺโน”ติ, อามนฺตา.\csromanspacer{} ปุคฺคโล สปฺปจฺจโย ฯเปฯ.\csromanspacer{} ปุคฺคโล อปฺปจฺจโย.\csromanspacer{} ปุคฺคโล สงฺขโต.\csromanspacer{} ปุคฺคโล อสงฺขโต.\csromanspacer{} ปุคฺคโล สสฺสโต.\csromanspacer{} ปุคฺคโล อสสฺสโต.\csromanspacer{} ปุคฺคโล สนิมิตฺโต.\csromanspacer{} ปุคฺคโล อนิมิตฺโตติ, น เหวํ วตฺตพฺเพ.\csromanspacer{}\csromanspacer{}(สํขิตฺตํ.)}

\vspace{\tipitakaparskip}
\csromancenter{ลกฺขณยุตฺติกถา.}
\csromanmatikaanchor{mtk.10}
\csromantocmarkref{section}{9. วจนโสธน}{mtk.10}
\bookmark[dest={mtk.10},level=1]{9. วจนโสธน}
\csromanheader{1}{9. วจนโสธน}
\proseitem{55}{\csromansymbolmark{*}ปุคฺคโล อุปลพฺภติ, อุปลพฺภติ ปุคฺคโลติ? ปุคฺคโล อุปลพฺภติ, อุปลพฺภติ เก หิ จิ ปุคฺคโล เก หิ จิ น ปุคฺคโลติ.\csromanspacer{} ปุคฺคโล เก หิ จิ อุปลพฺภติ เก หิ จิ น อุปลพฺภตีติ, น เหวํ วตฺตพฺเพ ฯเปฯ.}

\proseitem{56}{\csromansymbolmark{*}ปุคฺคโล สจฺจิกฏฺโฐ, สจฺจิกฏฺโฐ ปุคฺคโลติ? ปุคฺคโล สจฺจิกฏฺโฐ, สจฺจิกฏฺโฐ เก หิ จิ ปุคฺคโล เก หิ จิ น ปุคฺคโลติ.\csromanspacer{} ปุคฺคโล เก หิ จิ สจฺจิกฏฺโฐ เก หิ จิ น สจฺจิกฏฺโฐติ, น เหวํ วตฺตพฺเพ ฯเปฯ.}

\proseitem{57}{\csromansymbolmark{*}ปุคฺคโล วิชฺชมาโน, วิชฺชมาโน ปุคฺคโลติ? ปุคฺคโล วิชฺชมาโน, วิชฺชมาโน เก หิ จิ ปุคฺคโล เก หิ จิ น ปุคฺคโลติ.\csromanspacer{} ปุคฺคโล เก หิ จิ วิชฺชมาโน เก หิ จิ น วิชฺชมาโนติ, น เหวํ วตฺตพฺเพ ฯเปฯ.}

\proseitem{58}{\csromansymbolmark{*}ปุคฺคโล สํวิชฺชมาโน, สํวิชฺชมาโน ปุคฺคโลติ? ปุคฺคโล สํวิชฺชมาโน, สํวิชฺชมาโน เก หิ จิ ปุคฺคโล เก หิ จิ น ปุคฺคโลติ.\csromanspacer{} ปุคฺคโล เก หิ จิ สํวิชฺชมาโน เก หิ จิ น สํวิชฺชมาโนติ, น เหวํ วตฺตพฺเพ ฯเปฯ.}

\proseitem{59}{\csromansymbolmark{*}ปุคฺคโล อตฺถิ, อตฺถิ ปุคฺคโลติ? ปุคฺคโล อตฺถิ, อตฺถิ เก หิ จิ ปุคฺคโล เก หิ จิ น ปุคฺคโลติ.\csromanspacer{} ปุคฺคโล เก หิ จิ อตฺถิ เก หิ จิ นตฺถีติ, น เหวํ วตฺตพฺเพ ฯเปฯ.}

\proseitem{60}{\csromansymbolmark{*}ปุคฺคโล อตฺถิ, อตฺถิ น สพฺโพ ปุคฺคโลติ, อามนฺตา ฯเปฯ.\csromanspacer{} ปุคฺคโล นตฺถิ, นตฺถิ น สพฺโพ ปุคฺคโลติ, น เหวํ วตฺตพฺเพ.\csromanspacer{}\csromanspacer{}(สํขิตฺตํ.)}

\vspace{\tipitakaparskip}
\csromancenter{วจนโสธนํ.}
\csromanmatikaanchor{mtk.11}
\csromantocmarkref{section}{10. ปญฺญตฺตานุโยค}{mtk.11}
\bookmark[dest={mtk.11},level=1]{10. ปญฺญตฺตานุโยค}
\csromanheader{1}{10. ปญฺญตฺตานุโยค}
\proseitem{61}{\csromanfolio{27}\csromansymbolmark{*}รูปธาตุยา รูปี ปุคฺคโลติ, อามนฺตา.\csromanspacer{} กามธาตุยา กามี ปุคฺคโลติ, น เหวํ วตฺตพฺเพ ฯเปฯ.}

\proseitem{62}{\csromansymbolmark{*}รูปธาตุยา รูปิโน สตฺตาติ, อามนฺตา.\csromanspacer{} กามธาตุยา กามิโน สตฺตาติ, น เหวํ วตฺตพฺเพ ฯเปฯ.}

\proseitem{63}{\csromansymbolmark{*}อรูปธาตุยา อรูปี ปุคฺคโลติ, อามนฺตา.\csromanspacer{} กามธาตุยา กามี ปุคฺคโลติ, น เหวํ วตฺตพฺเพ ฯเปฯ.}

\proseitem{64}{\csromansymbolmark{*}อรูปธาตุยา อรูปิโน สตฺตาติ, อามนฺตา.\csromanspacer{} กามธาตุยา กามิโน สตฺตาติ, น เหวํ วตฺตพฺเพ ฯเปฯ.}

\proseitem{65}{\csromansymbolmark{*}รูปธาตุยา รูปี ปุคฺคโล อรูปธาตุยา อรูปี ปุคฺคโล, อตฺถิ จ โกจิ รูปธาตุยา จุโต อรูปธาตุํ อุปปชฺชตีติ, อามนฺตา.\csromanspacer{} รูปี ปุคฺคโล อุปจฺฉินฺโน, อรูปี ปุคฺคโล ชาโตติ, น เหวํ วตฺตพฺเพ ฯเปฯ.}

\proseitem{66}{\csromansymbolmark{*}รูปธาตุยา รูปิโน สตฺตา อรูปธาตุยา อรูปิโน สตฺตา, อตฺถิ จ โกจิ รูปธาตุยา จุโต อรูปธาตุํ อุปปชฺชตีติ, อามนฺตา.\csromanspacer{} รูปี สตฺโต อุปจฺฉินฺโน, อรูปี สตฺโต ชาโตติ, น เหวํ วตฺตพฺเพ ฯเปฯ.}

\proseitem{67}{\csromansymbolmark{*}“กาโย”ติ วา “สรีรนฺ”ติ วา, “สรีรนฺ”ติ วา “กาโย”ติ วา กายํ อปฺปิยํ กริตฺวา เอเสเส เอกฏฺเฐ สเม สมภาเค ตชฺชาเตติ, อามนฺตา. “ปุคฺคโล”ติ วา “ชีโว”ติ วา, “ชีโว”ติ วา “ปุคฺคโล”ติ วา ปุคฺคลํ อปฺปิยํ กริตฺวา เอเสเส เอกฏฺเฐ สเม สมภาเค ตชฺชาเตติ, อามนฺตา.\csromanspacer{} อญฺโญ กาโย อญฺโญ ปุคฺคโลติ, อามนฺตา.\csromanspacer{} อญฺญํ ชีวํ อญฺญํ สรีรนฺติ, น เหวํ วตฺตพฺเพ.}

\prose{อาชานาหิ นิคฺคหํ\csromandash{}หญฺจิ “กาโย’ติ วา “สรีรนฺ”ติ วา, “สรีรนฺ”ติ วา “กาโย”ติ วา กายํ อปฺปิยํ กริตฺวา เอเสเส เอกฏฺเฐ สเม สมภาเค ตชฺชาเต, “ปุคฺคโล”ติ วา “ชีโว”ติ วา, “ชีโว”ติ วา “ปุคฺคโล”ติ วา ปุคฺคลํ อปฺปิยํ กริตฺวา เอเสเส เอกฏฺเฐ สเม สมภาเค ตชฺชาเต, อญฺโญ กาโย อญฺโญ ปุคฺคโล.\csromanspacer{} เตน วต เร วตฺตพฺเพ “อญฺญํ ชีวํ อญฺญํ สรีรนฺ”ติ.\csromanspacer{} ยํ ตตฺถ วเทสิ “วตฺตพฺเพ โข \csromanfolio{28}‘กาโยติ วา สรีรนฺติ วา, สรีรนฺติ วา กาโยติ วา กายํ อปฺปิยํ กริตฺวา เอเสเส เอกฏฺเฐ สเม สมภาเค ตชฺชาเต, ปุคฺคโลติ วา ชีโวติ วา, ชีโวติ วา ปุคฺคโลติ วา ปุคฺคลํ อปฺปิยํ กริตฺวา เอเสเส เอกฏฺเฐ สเม สมภาเค ตชฺชาเต, อญฺโญ กาโย อญฺโญ ปุคฺคโล’, โน จ วตฺตพฺเพ ‘อญฺญํ ชีวํ อญฺญํ สรีรนฺ’ติ”, มิจฺฉา.}

\prose{โน เจ ปน วตฺตพฺเพ “อญฺญํ ชีวํ อญฺญํ สรีรนฺ”ติ, โน จ วต เร วตฺตพฺเพ “กาโยติ วา สรีรนฺติ วา, สรีรนฺติ วา กาโยติ วา กายํ อปฺปิยํ กริตฺวา เอเสเส เอกฏฺเฐ สเม สมภาเค ตชฺชาเต, ปุคฺคโลติ วา ชีโวติ วา, ชีโวติ วา ปุคฺคโลติ วา ปุคฺคลํ อปฺปิยํ กริตฺวา เอเสเส เอกฏฺเฐ สเม สมภาเค ตชฺชาเต, อญฺโญ กาโย อญฺโญ ปุคฺคโล”ติ.\csromanspacer{} ยํ ตตฺถ วเทสิ “วตฺตพฺเพ โข ‘กาโยติ วา สรีรนฺติ วา, สรีรนฺติ วา กาโยติ วา กายํ อปฺปิยํ กริตฺวา เอเสเส เอกฏฺเฐ สเม สมภาเค ตชฺชาเต, ปุคฺคโลติ วา ชีโวติ วา, ชีโวติ วา ปุคฺคโลติ วา ปุคฺคลํ อปฺปิยํ กริตฺวา เอเสเส เอกฏฺเฐ สเม สมภาเค ตชฺชาเต, อญฺโญ กาโย อญฺโญ ปุคฺคโล’, โน จ วตฺตพฺเพ ‘อญฺญํ ชีวํ อญฺญํ สรีรนฺ’ติ”, มิจฺฉา ฯเปฯ.}

\proseitem{68}{\csromansymbolmark{+}“กาโย”ติ วา “สรีรนฺ”ติ วา, “สรีรนฺ”ติ วา “กาโย”ติ วา กายํ อปฺปิยํ กริตฺวา เอเสเส เอกฏฺเฐ สเม สมภาเค ตชฺชาเตติ, อามนฺตา.\csromanspacer{} วุตฺตํ ภควตา “อตฺถิ ปุคฺคโล อตฺตหิตาย ปติปนฺโน”ติ, อามนฺตา.\csromanspacer{} อญฺโญ กาโย อญฺโญ ปุคฺคโลติ, น เหวํ วตฺตพฺเพ.}

\prose{อาชานาหิ ปฏิกมฺมํ\csromandash{}หญฺจิ “กาโย”ติ วา “สรีรนฺ”ติ วา, “สรีรนฺ”ติ วา “กาโย”ติ วา กายํ อปฺปิยํ กริตฺวา เอเสเส เอกฏฺเฐ สเม สมภาเค ตชฺชาเต, วุตฺตํ ภควตา\csromandash{}อตฺถิ ปุคฺคโล อตฺตหิตาย ปฏิปนฺโน, เตน วต เร วตฺตพฺเพ “อญฺโญ กาโย อญฺโญ ปุคฺคโล”ติ.\csromanspacer{} ยํ ตตฺถ วเทสิ “วตฺตพฺเพ โข ‘กาโยติ วา สรีรนฺติ วา, สรีรนฺติ วา กาโยติ วา กายํ อปฺปิยํ กริตฺวา เอเสเส เอกฏฺเฐ สเม สมภาเค ตชฺชาเต, วุตฺตํ ภควตา\csromandash{}อตฺถิ ปุคฺคโล อตฺตหิตาย ปฏิปนฺโน’, โน จ วตฺตพฺเพ ‘อญฺโญ กาโย อญฺโญ ปุคฺคโล’ติ”, มิจฺฉา.}

\prose{\csromanfolio{29}โน เจ ปน วตฺตพฺเพ “อญฺโญ กาโย อญฺโญ ปุคฺคโล”ติ, โน จ วต เร วตฺตพฺเพ “กาโยติ วา สรีรนฺติ วา, สรีรนฺติ วา กาโยติ วา กายํ อปฺปิยํ กริตฺวา เอเสเส เอกฏฺเฐ สเม สมภาเค ตชฺชาเต, วุตฺตํ ภควตา\csromandash{}อตฺถิ ปุคฺคโล อตฺตหิตาย ปฏิปนฺโน”ติ.\csromanspacer{} ยํ ตตฺถ วเทสิ “วตฺตพฺเพ โข ‘กาโยติ วา สรีรนฺติ วา, สรีรนฺติ วา กาโยติ วา กายํ อปฺปิยํ กริตฺวา เอเสเส เอกฏฺเฐ สเม สมภาเค ตชฺชาเต, วุตฺตํ ภควตา\csromandash{}อตฺถิ ปุคฺคโล อตฺตหิตาย ปฏิปนฺโน’, โน จ วตฺตพฺเพ ‘อญฺโญ กาโย อญฺโญ ปุคฺคโล’ติ”, มิจฺฉา.\csromanspacer{}\csromanspacer{}(สํขิตฺตํ.)}

\vspace{\tipitakaparskip}
\csromancenter{ปญฺญตฺตานุโยโค.}
\csromanmatikaanchor{mtk.12}
\csromantocmarkref{section}{11. คติ-อนุโยค}{mtk.12}
\bookmark[dest={mtk.12},level=1]{11. คติ-อนุโยค}
\csromanheader{1}{11. คติ-อนุโยค}
\proseitem{69}{\csromansymbolmark{*}ปุคฺคโล สนฺธาวติ อสฺมา โลกา ปรํ โลกํ, ปรสฺมา โลกา อิมํ โลกนฺติ, อามนฺตา.\csromanspacer{} โส ปุคฺคโล สนฺธาวติ อสฺมา โลกา ปรํ โลกํ, ปรสฺมา โลกา อิมํ โลกนฺติ, น เหวํ วตฺตพฺเพ ฯเปฯ.}

\proseitem{70}{\csromansymbolmark{*}ปุคฺคโล สนฺธาวติ อสฺมา โลกา ปรํ โลกํ, ปรสฺมา โลกา อิมํ โลกนฺติ, อามนฺตา.\csromanspacer{} อญฺโญ ปุคฺคโล สนฺธาวติ อสฺมา โลกา ปรํ โลกํ, ปรสฺมา โลกา อิมํ โลกนฺติ, น เหวํ วตฺตพฺเพ ฯเปฯ.}

\proseitem{71}{\csromansymbolmark{*}ปุคฺคโล สนฺธาวติ อสฺมา โลกา ปรํ โลกํ, ปรสฺมา โลกา อิมํ โลกนฺติ, อามนฺตา.\csromanspacer{} โส จ อญฺโญ จ สนฺธาวติ อสฺมา โลกา ปรํ โลกํ, ปรสฺมา โลกา อิมํ โลกนฺติ, น เหวํ วตฺตพฺเพ ฯเปฯ.}

\proseitem{72}{\csromansymbolmark{*}ปุคฺคโล สนฺธาวติ อสฺมา โลกา ปรํ โลกํ, ปรสฺมา โลกา อิมํ โลกนฺติ, อามนฺตา.\csromanspacer{} เนว โส สนฺธาวติ, น อญฺโญ สนฺธาวติ อสฺมา โลกา ปรํ โลกํ, ปรสฺมา โลกา อิมํ โลกนฺติ, น เหวํ วตฺตพฺเพ ฯเปฯ.}

\proseitem{73}{\csromanfolio{30}\csromansymbolmark{*}ปุคฺคโล สนฺธาวติ อสฺมา โลกา ปรํ โลกํ, ปรสฺมา โลกา อิมํ โลกนฺติ, อามนฺตา.\csromanspacer{} โส ปุคฺคโล สนฺธาวติ, อญฺโญ ปุคฺคโล สนฺธาวติ, โส จ อญฺโญ จ สนฺธาวติ, เนว โส สนฺธาวติ, น อญฺโญ สนฺธาวติ อสฺมา โลกา ปรํ โลกํ, ปรสฺมา โลกา อิมํ โลกนฺติ, น เหวํ วตฺตพฺเพ ฯเปฯ.}

\proseitem{74}{\csromansymbolmark{+}น วตฺตพฺพํ “ปุคฺคโล สนฺธาวติ อสฺมา โลกา ปรํ โลกํ, ปรสฺมา โลกา อิมํ โลกนฺ”ติ, อามนฺตา.\csromanspacer{} นนุ วุตฺตํ ภควตา}

\csromangathameasure{“ส สตฺตกฺขตฺตุปรมํ, สนฺธาวิตฺวาน ปุคฺคโล. \\ ทุกฺขสฺสนฺตกโร โหติ, สพฺพสํโยชนกฺขยา”ติ.}
\csromangathagroup{\csromangathastanza{“ส สตฺตกฺขตฺตุปรมํ, สนฺธาวิตฺวาน ปุคฺคโล. \\ ทุกฺขสฺสนฺตกโร โหติ, สพฺพสํโยชนกฺขยา”ติ\footnote{สํ 1. 393 ปิฏฺเฐ; ขุ 1. 207 อิติวุตฺตเกปิ.}.}}

\prose{อตฺเถว สุตฺตนฺโตติ, อามนฺตา.\csromanspacer{} เตน หิ ปุคฺคโล สนฺธาวติ อสฺมา โลกา ปรํ โลกํ, ปรสฺมา โลกา อิมํ โลกนฺติ.}

\proseitem{75}{\csromansymbolmark{+}น วตฺตพฺพํ “ปุคฺคโล สนฺธาวติ อสฺมา โลกา ปรํ โลกํ, ปรสฺมา โลกา อิมํ โลกนฺ”ติ, อามนฺตา.\csromanspacer{} นนุ วุตฺตํ ภควตา “อนมตคฺโคยํ\footnote{อนมตคฺคายํ (ก)} ภิกฺขเว สํสาโร ปุพฺพโกฏิ น ปญฺญายติ, อวิชฺชานีวรณานํ สตฺตานํ ตณฺหาสํโยชนานํ สนฺธาวตํ สํสรตนฺ”ติ\footnote{สํ 1. 387 ปิฏฺเฐ.} อตฺเถว สุตฺตนฺโตติ, อามนฺตา.\csromanspacer{} เตน หิ ปุคฺคโล สนฺธาวติ อสฺมา โลกา ปรํ โลกํ, ปรสฺมา โลกา อิมํ โลกนฺติ.}

\proseitem{76}{\csromansymbolmark{*}ปุคฺคโล สนฺธาวติ อสฺมา โลกา ปรํ โลกํ, ปรสฺมา โลกา อิมํ โลกนฺติ, อามนฺตา.\csromanspacer{} เสฺวว ปุคฺคโล สนฺธาวติ อสฺมา โลกา ปรํ โลกํ, ปรสฺมา โลกา อิมํ โลกนฺติ, น เหวํ วตฺตพฺเพ ฯเปฯ.}

\proseitem{77}{\csromansymbolmark{*}เสฺวว ปุคฺคโล สนฺธาวติ อสฺมา โลกา ปรํ โลกํ, ปรสฺมา โลกา อิมํ โลกนฺติ, อามนฺตา.\csromanspacer{} อตฺถิ โกจิ มนุสฺโส หุตฺวา เทโว โหตีติ, อามนฺตา.\csromanspacer{} เสฺวว มนุสฺโส โส เทโวติ, น เหวํวตฺตพฺเพ ฯเปฯ.}

\proseitem{78}{\csromansymbolmark{*}เสฺวว มนุสฺโส โส เทโวติ, อามนฺตา.\csromanspacer{} มนุสฺโส หุตฺวา เทโว โหติ, เทโว หุตฺวา มนุสฺโส โหติ, มนุสฺสภูโต อญฺโญ, เทโว อญฺโญ, มนุสฺสภูโต เสฺววายํ สนฺธาวตีติ, มิจฺฉา ฯเปฯ.}

\prose{\csromanfolio{31}สเจ หิ สนฺธาวติ เสฺวว ปุคฺคโล อิโต จุโต ปรํ โลกํ, อนญฺโญ, เหวํ มรณํ น เหหิติ, ปาณาติปาโตปิ นุปลพฺภติ.\csromanspacer{} กมฺมํ อตฺถิ, กมฺมวิปาโก อตฺถิ, กตานํ กมฺมานํ วิปาโก อตฺถิ, กุสลากุสเล วิปจฺจมาเน เสฺววายํ สนฺธาวตีติ, มิจฺฉา.}

\proseitem{79}{\csromansymbolmark{*}เสฺวว ปุคฺคโล สนฺธาวติ อสฺมา โลกา ปรํ โลกํ, ปรสฺมา โลกา อิมํ โลกนฺติ, อามนฺตา.\csromanspacer{} อตฺถิ โกจิ มนุสฺโส หุตฺวา ยกฺโข โหติ.\csromanspacer{} เปโต โหติ.\csromanspacer{} เนรยิโก โหติ.\csromanspacer{} ติรจฺฉานคโต โหติ.\csromanspacer{} โอฏฺโฐ โหติ.\csromanspacer{} โคโณ โหติ.\csromanspacer{} คทฺรโภ โหติ.\csromanspacer{} สูกโร โหติ.\csromanspacer{} มหิํโส\footnote{มหิโส (สี, สฺยา, กํ, อิ)} โหตีติ, อามนฺตา.\csromanspacer{} เสฺวว มนุสฺโส โส มหิํโสติ, น เหวํ วตฺตพฺเพ ฯเปฯ.}

\proseitem{80}{\csromansymbolmark{*}เสฺวว มนุสฺโส โส มหิํโสติ, อามนฺตา.\csromanspacer{} มนุสฺโส หุตฺวา มหิํโส โหติ, มหิํโส หุตฺวา มนุสฺโส โหติ, มนุสฺสภูโต อญฺโญ, มหิํโส อญฺโญ, มนุสฺสภูโต เสฺววายํ สนฺธาวตีติ, มิจฺฉา ฯเปฯ.}

\prose{สเจ หิ สนฺธาวติ เสฺวว ปุคฺคโล อิโต จุโต ปรํ โลกํ, อนญฺโญ, เหวํ มรณํ น เหหิติ, ปาณาติปาโตปิ นุปลพฺภติ.\csromanspacer{} กมฺมํ อตฺถิ, กมฺมวิปาโก อตฺถิ, กตานํ กมฺมานํ วิปาโก อตฺถิ, กุสลากุสเล วิปจฺจมาเน เสฺววายํ สนฺธาวตีติ, มิจฺฉา.}

\proseitem{81}{\csromansymbolmark{*}เสฺวว ปุคฺคโล สนฺธาวติ อสฺมา โลกา ปรํ โลกํ, ปรสฺมา โลกา อิมํ โลกนฺติ, อามนฺตา.\csromanspacer{} อตฺถิ โกจิ ขตฺติโย หุตฺวา พฺราหฺมโณ โหตีติ, อามนฺตา.\csromanspacer{} เสฺวว ขตฺติโย โส พฺราหฺมโณติ, น เหวํ วตฺตพฺเพ ฯเปฯ.}

\proseitem{82}{\csromansymbolmark{*}อตฺถิ โกจิ ขตฺติโย หุตฺวา เวสฺโส โหติ.\csromanspacer{} สุทฺโท โหตีติ, อามนฺตา.\csromanspacer{} เสฺวว ขตฺติโย โส สุทฺโทติ, น เหวํ วตฺตพฺเพ ฯเปฯ.}

\proseitem{83}{\csromansymbolmark{*}อตฺถิ โกจิ พฺราหฺมโณ หุตฺวา เวสฺโส โหติ.\csromanspacer{} สุทฺโท โหติ.\csromanspacer{} ขตฺติโย โหตีติ, อามนฺตา.\csromanspacer{} เสฺวว พฺราหฺมโณ โส ขตฺติโยติ, น เหวํ วตฺตพฺเพ ฯเปฯ.}

\proseitem{84}{\csromanfolio{32}\csromansymbolmark{*}อตฺถิ โกจิ เวสฺโส หุตฺวา สุทฺโท โหติ.\csromanspacer{} ขตฺติโย โหติ.\csromanspacer{} พฺราหฺมโณ โหตีติ, อามนฺตา.\csromanspacer{} เสฺวว เวสฺโส โส พฺราหฺมโณติ, น เหวํ วตฺตพฺเพ ฯเปฯ.}

\proseitem{85}{\csromansymbolmark{*}อตฺถิ โกจิ สุทฺโท หุตฺวา ขตฺติโย โหติ.\csromanspacer{} พฺราหฺมโณ โหติ.\csromanspacer{} เวสฺโส โหตีติ, อามนฺตา.\csromanspacer{} เสฺวว สุทฺโท โส เวสฺโสติ, น เหวํ วตฺตพฺเพ ฯเปฯ.}

\proseitem{86}{\csromansymbolmark{*}เสฺวว ปุคฺคโล สนฺธาวติ อสฺมา โลกา ปรํ โลกํ, ปรสฺมา โลกา อิมํ โลกนฺติ, อามนฺตา.\csromanspacer{} หตฺถจฺฉินฺโน หตฺถจฺฉินฺโนว โหติ.\csromanspacer{} ปาทจฺฉินฺโน ปาทจฺฉินฺโนว โหติ.\csromanspacer{} หตฺถปาทจฺฉินฺโน หตฺถปาทจฺฉินฺโนว โหติ.\csromanspacer{} กณฺณจฺฉินฺโน.\csromanspacer{} นาสจฺฉินฺโน.\csromanspacer{} กณฺณนาสจฺฉินฺโน.\csromanspacer{} องฺคุลิจฺฉินฺโน.\csromanspacer{} อฬจฺฉินฺโน.\csromanspacer{} กณฺฑรจฺฉินฺโน.\csromanspacer{} กุณิหตฺถโก.\csromanspacer{} ผณหตฺถโก.\csromanspacer{} กุฏฺฐิโย.\csromanspacer{} คณฺฑิโย.\csromanspacer{} กิลาสิโย.\csromanspacer{} โสสิโย.\csromanspacer{} อปมาริโย.\csromanspacer{} โอฏฺโฐ.\csromanspacer{} โคโณ.\csromanspacer{} คทฺรโภ.\csromanspacer{} สูกโร.\csromanspacer{} มหิํโส มหิํโสว โหตีติ, น เหวํ วตฺตพฺเพ ฯเปฯ.}

\proseitem{87}{\csromansymbolmark{+}น วตฺตพฺพํ “เสฺวว ปุคฺคโล สนฺธาวติ อสฺมา โลกา ปรํ โลกํ, ปรสฺมา โลกา อิมํ โลกนฺ”ติ, อามนฺตา.\csromanspacer{} นนุ โสตาปนฺโน ปุคฺคโล มนุสฺสโลกา จุโต เทวโลกํ อุปปนฺโน, ตตฺถปิ โสตาปนฺโนว โหตีติ, อามนฺตา.}

\prose{หญฺจิ โสตาปนฺโน ปุคฺคโล มนุสฺสโลกา จุโต เทวโลกํ อุปปนฺโน, ตตฺถปิ โสตาปนฺโนว โหติ, เตน วต เร วตฺตพฺเพ “เสฺวว ปุคฺคโล สนฺธาวติ อสฺมา โลกา ปรํ โลกํ, ปรสฺมา โลกา อิมํ โลกนฺ”ติ.}

\proseitem{88}{\csromansymbolmark{*}โสตาปนฺโน ปุคฺคโล มนุสฺสโลกา จุโต เทวโลกํ อุปปนฺโน, ตตฺถปิ โสตาปนฺโนว โหตีติ กตฺวา เตน จ การเณน เสฺวว ปุคฺคโล สนฺธาวติ อสฺมา โลกา ปรํ โลกํ, ปรสฺมา โลกา อิมํ โลกนฺติ, อามนฺตา.\csromanspacer{} โสตาปนฺโน ปุคฺคโล มนุสฺสโลกา จุโต เทวโลกํ อุปปนฺโน, ตตฺถปิ มนุสฺโส โหตีติ กตฺวา, น เหวํ วตฺตพฺเพ ฯเปฯ.}

\proseitem{89}{\csromanfolio{33}\csromansymbolmark{*}เสฺวว ปุคฺคโล สนฺธาวติ อสฺมา โลกา ปรํ โลกํ, ปรสฺมา โลกา อิมํ โลกนฺติ, อามนฺตา.\csromanspacer{} อนญฺโญ อวิคโต สนฺธาวตีติ, น เหวํ วตฺตพฺเพ ฯเปฯ.}

\proseitem{90}{\csromansymbolmark{*}อนญฺโญ อวิคโต สนฺธาวตีติ, อามนฺตา.\csromanspacer{} หตฺถจฺฉินฺโน หตฺถจฺฉินฺโนว โหติ.\csromanspacer{} ปาทจฺฉินฺโน ปาทจฺฉินฺโนว โหติ.\csromanspacer{} หตฺถปาทจฺฉินฺโน หตฺถปาทจฺฉินฺโนว โหติ.\csromanspacer{} กณฺณจฺฉินฺโน.\csromanspacer{} นาสจฺฉินฺโน.\csromanspacer{} กณฺณนาสจฺฉินฺโน.\csromanspacer{} องฺคุลิจฺฉินฺโน.\csromanspacer{} อฬจฺฉินฺโน.\csromanspacer{} กณฺฑรจฺฉินฺโน.\csromanspacer{} กุณิหตฺถโก.\csromanspacer{} ผณหตฺถโก.\csromanspacer{} กุฏฺฐิโย.\csromanspacer{} คณฺฑิโย.\csromanspacer{} กิลาสิโย.\csromanspacer{} โสสิโย.\csromanspacer{} อปมาริโย.\csromanspacer{} โอฏฺโฐ.\csromanspacer{} โคโณ.\csromanspacer{} คทฺรโภ.\csromanspacer{} สูกโร.\csromanspacer{} มหิํโส มหิํโสว โหตีติ, น เหวํ วตฺตพฺเพ ฯเปฯ.}

\proseitem{91}{\csromansymbolmark{*}เสฺวว ปุคฺคโล สนฺธาวติ อสฺมา โลกา ปรํ โลกํ, ปรสฺมา โลกา อิมํ โลกนฺติ, อามนฺตา.\csromanspacer{} สรูโป สนฺธาวตีติ, น เหวํ วตฺตพฺเพ ฯเปฯ.\csromanspacer{} สรูโป สนฺธาวตีติ, อามนฺตา.\csromanspacer{} ตํ ชีวํ ตํ สรีรนฺติ, น เหวํ วตฺตพฺเพ ฯเปฯ.}

\prose{สเวทโน ฯเปฯ.\csromanspacer{} สสญฺโญ ฯเปฯ.\csromanspacer{} สสงฺขาโร ฯเปฯ.\csromanspacer{} สวิญฺญาโณ สนฺธาวตีติ, น เหวํ วตฺตพฺเพ ฯเปฯ.\csromanspacer{} สวิญฺญาโณ สนฺธาวตีติ, อามนฺตา.\csromanspacer{} ตํ ชีวํ ตํ สรีรนฺติ, น เหวํ วตฺตพฺเพ ฯเปฯ.}

\proseitem{92}{\csromansymbolmark{*}เสฺวว ปุคฺคโล สนฺธาวติ อสฺมา โลกา ปรํ โลกํ, ปรสฺมา โลกา อิมํ โลกนฺติ, อามนฺตา.\csromanspacer{} อรูโป สนฺธาวตีติ, น เหวํ วตฺตพฺเพ ฯเปฯ.\csromanspacer{} อรูโป สนฺธาวตีติ, อามนฺตา.\csromanspacer{} อญฺญํ ชีวํ อญฺญํ สรีรนฺติ, น เหวํ วตฺตพฺเพ ฯเปฯ.}

\prose{อเวทโน ฯเปฯ.\csromanspacer{} อสญฺโญ ฯเปฯ.\csromanspacer{} อสงฺขาโร ฯเปฯ.\csromanspacer{} อวิญฺโญโณ สนฺธาวตีติ, น เหวํ วตฺตพฺเพ ฯเปฯ.\csromanspacer{} อวิญฺญาโณ สนฺธาวตีติ, อามนฺตา.\csromanspacer{} อญฺญํ ชีวํ อญฺญํ สรีรนฺติ, น เหวํ วตฺตพฺเพ ฯเปฯ.}

\proseitem{93}{\csromansymbolmark{*}เสฺวว ปุคฺคโล สนฺธาวติ อสฺมา โลกา ปรํ โลกํ, ปรสฺมา โลกา อิมํ โลกนฺติ, อามนฺตา.\csromanspacer{} รูปํ สนฺธาวตีติ, น เหวํ วตฺตพฺเพ ฯเปฯ.\csromanspacer{} รูปํ สนฺธาวตีติ, อามนฺตา.\csromanspacer{} ตํ ชีวํ ตํ สรีรนฺติ, น เหวํ วตฺตพฺเพ ฯเปฯ.}

\prose{\csromanfolio{34}เวทนา ฯเปฯ.\csromanspacer{} สญฺญา ฯเปฯ.\csromanspacer{} สงฺขารา ฯเปฯ.\csromanspacer{} วิญฺญาณํ สนฺธาวตีติ, น เหวํ วตฺตพฺเพ ฯเปฯ.\csromanspacer{} วิญฺญาณํ สนฺธาวตีติ, อามนฺตา.\csromanspacer{} ตํ ชีวํ ตํ สรีรนฺติ, น เหวํ วตฺตพฺเพ ฯเปฯ.}

\proseitem{94}{\csromansymbolmark{*}เสฺวว ปุคฺคโล สนฺธาวติ อสฺมา โลกา ปรํ โลกํ, ปรสฺมา โลกา อิมํ โลกนฺติ, อามนฺตา.\csromanspacer{} รูปํ น สนฺธาวตีติ, น เหวํ วตฺตพฺเพ ฯเปฯ.\csromanspacer{} รูปํ น สนฺธาวตีติ, อามนฺตา.\csromanspacer{} อญฺญํ ชีวํ อญฺญํ สรีรนฺติ, น เหวํ วตฺตพฺเพ ฯเปฯ.}

\prose{เวทนา ฯเปฯ.\csromanspacer{} สญฺญา ฯเปฯ.\csromanspacer{} สงฺขารา ฯเปฯ.\csromanspacer{} วิญฺญาณํ น สนฺธาวตีติ, น เหวํ วตฺตพฺเพ ฯเปฯ.\csromanspacer{} วิญฺญาณํ น สนฺธาวตีติ, อามนฺตา.\csromanspacer{} อญฺญํ ชีวํ อญฺญํ สรีรนฺติ, น เหวํ วตฺตพฺเพ.\csromanspacer{}\csromanspacer{}(สํขิตฺตํ.)}

\csromangathameasure{ขนฺเธสุ ภิชฺชมาเนสุ, โส เจ ภิชฺชติ ปุคฺคโล. \\ อุจฺเฉทา ภวติ ทิฏฺฐิ, ยา พุทฺเธน วิวชฺชิตา. \\ ขนฺเธสุ ภิชฺชมาเนสุ, โน เจ ภิชฺชติ ปุคฺคโล. \\ ปุคฺคโล สสฺสโต โหติ, นิพฺพาเนน สมสโมติ.}
\csromangathagroup{\csromangathastanza{ขนฺเธสุ ภิชฺชมาเนสุ, โส เจ ภิชฺชติ ปุคฺคโล. \\ อุจฺเฉทา ภวติ ทิฏฺฐิ, ยา พุทฺเธน วิวชฺชิตา.}\csromangathastanza{ขนฺเธสุ ภิชฺชมาเนสุ, โน เจ ภิชฺชติ ปุคฺคโล. \\ ปุคฺคโล สสฺสโต โหติ, นิพฺพาเนน สมสโมติ.}}

\vspace{\tipitakaparskip}
\csromancenter{คติ-อนุโยโค.}
\csromanmatikaanchor{mtk.13}
\csromantocmarkref{section}{12. อุปาทาปญฺญตฺตานุโยค}{mtk.13}
\bookmark[dest={mtk.13},level=1]{12. อุปาทาปญฺญตฺตานุโยค}
\csromanheader{1}{12. อุปาทาปญฺญตฺตานุโยค}
\proseitem{95}{\csromansymbolmark{*}รูปํ อุปาทาย ปุคฺคลสฺส ปญฺญตฺตีติ, อามนฺตา.\csromanspacer{} รูปํ อนิจฺจํ สงฺขตํ ปฏิจฺจสมุปฺปนฺนํ ขยธมฺมํ วยธมฺมํ วิราคธมฺมํ นิโรธธมฺมํ วิปริณามธมฺมนฺติ, อามนฺตา.\csromanspacer{} ปุคฺคโลปิ อนิจฺโจ สงฺขโต ปฏิจฺจสมุปฺปนฺโน ขยธมฺโม วยธมฺโม วิราคธมฺโม นิโรธธมฺโม วิปริณามธมฺโมติ, น เหวํ วตฺตพฺเพ ฯเปฯ.}

\proseitem{96}{\csromansymbolmark{*}เวทนํ อุปาทาย.\csromanspacer{} สญฺญํ อุปาทาย.\csromanspacer{} สงฺขาเร อุปาทาย.\csromanspacer{} วิญฺญาณํ อุปาทาย ปุคฺคลสฺส ปญฺญตฺตีติ, อามนฺตา.\csromanspacer{} วิญฺญาณํ อนิจฺจํ สงฺขตํ ปฏิจฺจสมุปฺปนฺนํ ขยธมฺมํ วยธมฺมํ วิราคธมฺมํ นิโรธธมฺมํ วิปริณามธมฺมนฺติ, อามนฺตา.\csromanspacer{} ปุคฺคโลปิ อนิจฺโจ สงฺขโต ปฏิจฺจสมุปฺปนฺโน ขยธมฺโม วยธมฺโม วิราคธมฺโม นิโรธธมฺโม วิปริณามธมฺโมติ, น เหวํ วตฺตพฺเพ ฯเปฯ.}

\proseitem{97}{\csromansymbolmark{*}รูปํ อุปาทาย ปุคฺคลสฺส ปญฺญตฺตีติ, อามนฺตา.\csromanspacer{} นีลํ รูปํ อุปาทาย นีลกสฺส ปุคฺคลสฺส ปญฺญตฺตีติ, น เหวํ วตฺตพฺเพ ฯเปฯ.\csromanspacer{} ปีตํ รูปํ อุปาทาย. \csromanfolio{35}โลหิตํ รูปํ อุปาทาย.\csromanspacer{} โอทาตํ รูปํ อุปาทาย.\csromanspacer{} สนิทสฺสนํ รูปํ อุปาทาย.\csromanspacer{} อนิทสฺสนํ รูปํ อุปาทาย.\csromanspacer{} สปฺปฏิฆํ รูปํ อุปาทาย.\csromanspacer{} อปฺปฏิฆํ รูปํ อุปาทาย อปฺปฏิฆสฺส ปุคฺคลสฺส ปญฺญตฺตีติ, น เหวํ วตฺตพฺเพ ฯเปฯ.}

\proseitem{98}{\csromansymbolmark{*}เวทนํ อุปาทาย ปุคฺคลสฺส ปญฺญตฺตีติ, อามนฺตา.\csromanspacer{} กุสลํ เวทนํ อุปาทาย กุสลสฺส ปุคฺคลสฺส ปญฺญตฺตีติ, น เหวํ วตฺตพฺเพ ฯเปฯ.\csromanspacer{} กุสลํ เวทนํ อุปาทาย กุสลสฺส ปุคฺคลสฺส ปญฺญตฺตีติ, อามนฺตา.\csromanspacer{} กุสลา เวทนา สผลา สวิปากา อิฏฺฐผลา กนฺตผลา มนุญฺญผลา อเสจนกผลา สุขุทฺรยา สุขวิปากาติ, อามนฺตา.\csromanspacer{} กุสโลปิ ปุคฺคโล สผโล สวิปาโก อิฏฺฐผโล กนฺตผโล มนุญฺญผโล อเสจนกผโล สุขุทฺรโย สุขวิปาโกติ, น เหวํ วตฺตพฺเพ ฯเปฯ.}

\proseitem{99}{\csromansymbolmark{*}เวทนํ อุปาทาย ปุคฺคลสฺส ปญฺญตฺตีติ, อามนฺตา.\csromanspacer{} อกุสลํ เวทนํ อุปาทาย อกุสลสฺส ปุคฺคลสฺส ปญฺญตฺตีติ, น เหวํ วตฺตพฺเพ ฯเปฯ.\csromanspacer{} อกุสลํ เวทนํ อุปาทาย อกุสลสฺส ปุคฺคลสฺส ปญฺญตฺตีติ, อามนฺตา.\csromanspacer{} อกุสลา เวทนา สผลา สวิปากา อนิฏฺฐผลา อกนฺตผลา อมนุญฺญผลา เสจนกผลา ทุกฺขุทฺรยา ทุกฺขวิปากาติ, อามนฺตา.\csromanspacer{} อกุสโลปิ ปุคฺคโล สผโล สวิปาโก อนิฏฺฐผโล อกนฺตผโล อมนุญฺญผโล เสจนกผโล ทุกฺขุทฺรโย ทุกฺขวิปาโกติ, น เหวํ วตฺตพฺเพ ฯเปฯ.}

\proseitem{100}{\csromansymbolmark{*}เวทนํ อุปาทาย ปุคฺคลสฺส ปญฺญตฺตีติ, อามนฺตา.\csromanspacer{} อพฺยากตํ เวทนํ อุปาทาย อพฺยากตสฺส ปุคฺคลสฺส ปญฺญตฺตีติ, น เหวํ วตฺตพฺเพ ฯเปฯ.\csromanspacer{} อพฺยากตํ เวทนํ อุปาทาย อพฺยากตสฺส ปุคฺคลสฺส ปญฺญตฺตีติ, อามนฺตา.\csromanspacer{} อพฺยากตา เวทนา อนิจฺจา สงฺขตา ปฏิจฺจสมุปฺปนฺนา ขยธมฺมา วยธมฺมา วิราคธมฺมา นิโรธธมฺมา วิปริณามธมฺมาติ, อามนฺตา.\csromanspacer{} อพฺยากโตปิ ปุคฺคโล อนิจฺโจ สงฺขโต ปฏิจฺจสมุปฺปนฺโน ขยธมฺโม วยธมฺโม วิราคธมฺโม นิโรธธมฺโม วิปริณามธมฺโมติ, น เหวํ วตฺตพฺเพ ฯเปฯ.}

\proseitem{101}{\csromansymbolmark{*}สญฺญํ อุปาทาย.\csromanspacer{} สงฺขาเร อุปาทาย.\csromanspacer{} วิญฺญาณํ อุปาทาย ปุคฺคลสฺส ปญฺญตฺตีติ, อามนฺตา.\csromanspacer{} กุสลํ วิญฺญาณํ อุปาทาย กุสลสฺส ปุคฺคลสฺส ปญฺญตฺตีติ, น เหวํ วตฺตพฺเพ ฯเปฯ.\csromanspacer{} กุสลํ วิญฺญาณํ อุปาทาย กุสลสฺส ปุคฺคลสฺส ปญฺญตฺตีติ, อามนฺตา.\csromanspacer{} กุสลํ วิญฺญาณํ สผลํ สวิปากํ อิฏฺฐผลํ กนฺตผลํ มนุญฺญผลํ อเสจนกผลํ สุขุทฺรยํ \csromanfolio{36}สุขวิปากนฺติ, อามนฺตา.\csromanspacer{} กุสโลปิ ปุคฺคโล สผโล สวิปาโก อิฏฺฐผโล กนฺตผโล มนุญฺญผโล อเสจนกผโล สุขุทฺรโย สุขวิปาโกติ, น เหวํ วตฺตพฺเพ ฯเปฯ.}

\proseitem{102}{\csromansymbolmark{*}วิญฺญาณํ อุปาทาย ปุคฺคลสฺส ปญฺญตฺตีติ, อามนฺตา.\csromanspacer{} อกุสลํ วิญฺญาณํ อุปาทาย อกุสลสฺส ปุคฺคลสฺส ปญฺญตฺตีติ, น เหวํ วตฺตพฺเพ ฯเปฯ.\csromanspacer{} อกุสลํ วิญฺญาณํ อุปาทาย อกุสลสฺส ปุคฺคลสฺส ปญฺญตฺตีติ, อามนฺตา.\csromanspacer{} อกุสลํ วิญฺญาณํ สผลํ สวิปากํ อนิฏฺฐผลํ อกนฺตผลํ อมนุญฺญผลํ เสจนกผลํ ทุกฺขุทฺรยํ ทุกฺขวิปากนฺติ, อามนฺตา.\csromanspacer{} อกุสโลปิ ปุคฺคโล สผโล สวิปาโก อนิฏฺฐผโล อกนฺตผโล อมนุญฺญผโล เสจนกผโล ทุกฺขุทฺรโย ทุกฺขวิปาโกติ, น เหวํ วตฺตพฺเพ ฯเปฯ.}

\proseitem{103}{\csromansymbolmark{*}วิญฺญาณํ อุปาทาย ปุคฺคลสฺส ปญฺญตฺตีติ, อามนฺตา.\csromanspacer{} อพฺยากตํ วิญฺญาณํ อุปาทาย อพฺยากตสฺส ปุคฺคลสฺส ปญฺญตฺตีติ, น เหวํ วตฺตพฺเพ ฯเปฯ.\csromanspacer{} อพฺยากตํ วิญฺญาณํ อุปาทาย อพฺยากตสฺส ปุคฺคลสฺส ปญฺญตฺตีติ, อามนฺตา.\csromanspacer{} อพฺยากตํ วิญฺญาณํ อนิจฺจํ สงฺขตํ ปฏิจฺจสมุปฺปนฺนํ ขยธมฺมํ วยธมฺมํ วิราคธมฺมํ นิโรธธมฺมํ วิปริณามธมฺมนฺติ, อามนฺตา.\csromanspacer{} อพฺยากโตปิ ปุคฺคโล อนิจฺโจ สงฺขโต ปฏิจฺจสมุปฺปนฺโน ขยธมฺโม วยธมฺโม วิราคธมฺโม นิโรธธมฺโม วิปริณามธมฺโมติ, น เหวํ วตฺตพฺเพ ฯเปฯ.}

\proseitem{104}{\csromansymbolmark{*}จกฺขุํ อุปาทาย “จกฺขุมา ปุคฺคโล”ติ วตฺตพฺโพติ, อามนฺตา.\csromanspacer{} จกฺขุมฺหิ นิรุทฺเธ “จกฺขุมา ปุคฺคโล นิรุทฺโธ”ติ วตฺตพฺโพติ, น เหวํ วตฺตพฺเพ ฯเปฯ.\csromanspacer{} โสตํ อุปาทาย.\csromanspacer{} ฆานํ อุปาทาย.\csromanspacer{} ชิวฺหํ อุปาทาย.\csromanspacer{} กายํ อุปาทาย.\csromanspacer{} มนํ อุปาทาย “มนวา ปุคฺคโล”ติ วตฺตพฺโพติ, อามนฺตา.\csromanspacer{} มนมฺหิ นิรุทฺเธ “มนวา ปุคฺคโล นิรุทฺโธ”ติ วตฺตพฺโพติ, น เหวํ วตฺตพฺเพ.}

\proseitem{105}{\csromansymbolmark{*}มิจฺฉาทิฏฺฐิํ อุปาทาย “มิจฺฉาทิฏฺฐิโย ปุคฺคโล”ติ วตฺตพฺโพติ, อามนฺตา.\csromanspacer{} มิจฺฉาทิฏฺฐิยา นิรุทฺธาย “มิจฺฉาทิฏฺฐิโย ปุคฺคโล นิรุทฺโธ”ติ วตฺตพฺโพติ, น เหวํ วตฺตพฺเพ.\csromanspacer{} มิจฺฉาสงฺกปฺปํ อุปาทาย.\csromanspacer{} มิจฺฉาวาจํ อุปาทาย.\csromanspacer{} มิจฺฉากมฺมนฺตํ อุปาทาย.\csromanspacer{} มิจฺฉา-อาชีวํ อุปาทาย.\csromanspacer{} มิจฺฉาวายามํ อุปาทาย.\csromanspacer{} มิจฺฉาสติํ อุปาทาย.\csromanspacer{} มิจฺฉาสมาธิํ อุปาทาย “มิจฺฉาสมาธิโย ปุคฺคโล”ติ วตฺตพฺโพติ, อามนฺตา.\csromanspacer{} มิจฺฉาสมาธิมฺหิ นิรุทฺเธ “มิจฺฉาสมาธิโย ปุคฺคโล นิรุทฺโธ”ติ วตฺตพฺโพติ, น เหวํ วตฺตพฺเพ.}

\proseitem{106}{\csromanfolio{37}\csromansymbolmark{*}สมฺมาทิฏฺฐิํ อุปาทาย “สมฺมาทิฏฺฐิโย ปุคฺคโล”ติ วตฺตพฺโพติ, อามนฺตา.\csromanspacer{} สมฺมาทิฏฺฐิยา นิรุทฺธาย “สมฺมาทิฏฺฐิโย ปุคฺคโล นิรุทฺโธ”ติ วตฺตพฺโพติ, น เหวํ วตฺตพฺเพ ฯเปฯ.\csromanspacer{} สมฺมาสงฺกปฺปํ อุปาทาย.\csromanspacer{} สมฺมาวาจํ อุปาทาย.\csromanspacer{} สมฺมากมฺมนฺตํ อุปาทาย.\csromanspacer{} สมฺมา-อาชีวํ อุปาทาย.\csromanspacer{} สมฺมาวายามํ อุปาทาย.\csromanspacer{} สมฺมาสติํ อุปาทาย.\csromanspacer{} สมฺมาสมาธิํ อุปาทาย “สมฺมาสมาธิโย ปุคฺคโล”ติ วตฺตพฺโพติ, อามนฺตา.\csromanspacer{} สมฺมาสมาธิมฺหิ นิรุทฺเธ “สมฺมาสมาธิโย ปุคฺคโล นิรุทฺโธ”ติ วตฺตพฺโพติ, น เหวํ วตฺตพฺเพ ฯเปฯ.}

\proseitem{107}{\csromansymbolmark{*}รูปํ อุปาทาย, เวทนํ อุปาทาย ปุคฺคลสฺส ปญฺญตฺตีติ, อามนฺตา.\csromanspacer{} ทฺวินฺนํ ขนฺธานํ อุปาทาย ทฺวินฺนํ ปุคฺคลานํ ปญฺญตฺตีติ, น เหวํ วตฺตพฺเพ ฯเปฯ.\csromanspacer{} รูปํ อุปาทาย, เวทนํ อุปาทาย, สญฺญํ อุปาทาย, สงฺขาเร อุปาทาย, วิญฺญาณํ อุปาทาย ปุคฺคลสฺส ปญฺญตฺตีติ, อามนฺตา.\csromanspacer{} ปญฺจนฺนํ ขนฺธานํ อุปาทาย ปญฺจนฺนํ ปุคฺคลานํ ปญฺญตฺตีติ, น เหวํ วตฺตพฺเพ ฯเปฯ.}

\proseitem{108}{\csromansymbolmark{*}จกฺขายตนํ อุปาทาย โสตายตนํ อุปาทาย, ปุคฺคลสฺส ปญฺญตฺตีติ, อามนฺตา.\csromanspacer{} ทฺวินฺนํ อายตนานํ อุปาทาย ทฺวินฺนํ ปุคฺคลานํ ปญฺญตฺตีติ, น เหวํ วตฺตพฺเพ ฯเปฯ.\csromanspacer{} จกฺขายตนํ อุปาทาย, โสตายตนํ อุปาทาย ฯเปฯ.\csromanspacer{} ธมฺมายตนํ อุปาทาย ปุคฺคลสฺส ปญฺญตฺตีติ, อามนฺตา.\csromanspacer{} ทฺวาทสนฺนํ อายตนานํ อุปาทาย ทฺวาทสนฺนํ ปุคฺคลานํ ปญฺญตฺตีติ, น เหวํ วตฺตพฺเพ ฯเปฯ.}

\proseitem{109}{\csromansymbolmark{*}จกฺขุธาตุํ อุปาทาย โสตธาตุํ อุปาทาย, ปุคฺคลสฺส ปญฺญตฺตีติ, อามนฺตา.\csromanspacer{} ทฺวินฺนํ ธาตูนํ อุปาทาย ทฺวินฺนํ ปุคฺคลานํ ปญฺญตฺตีติ, น เหวํ วตฺตพฺเพ ฯเปฯ.\csromanspacer{} จกฺขุธาตุํ อุปาทาย, โสตธาตุํ อุปาทาย ฯเปฯ.\csromanspacer{} ธมฺมธาตุํ อุปาทาย ปุคฺคลสฺส ปญฺญตฺตีติ, อามนฺตา.\csromanspacer{} อฏฺฐารสนฺนํ ธาตูนํ อุปาทาย อฏฺฐารสนฺนํ ปุคฺคลานํ ปญฺญตฺตีติ, น เหวํ วตฺตพฺเพ ฯเปฯ.}

\proseitem{110}{\csromansymbolmark{*}จกฺขุนฺทฺริยํ อุปาทาย, โสตินฺทฺริยํ อุปาทาย ปุคฺคลสฺส ปญฺญตฺตีติ, อามนฺตา.\csromanspacer{} ทฺวินฺนํ อินฺทฺริยานํ อุปาทาย ทฺวินฺนํ ปุคฺคลานํ ปญฺญตฺตีติ, น เหวํ วตฺตพฺเพ ฯเปฯ.\csromanspacer{} จกฺขุนฺทฺริยํ อุปาทาย, โสตินฺทฺริยํ อุปาทาย ฯเปฯ.\csromanspacer{} อญฺญาตาวินฺทฺริยํ อุปาทาย ปุคฺคลสฺส ปญฺญตฺตีติ, อามนฺตา.\csromanspacer{} พาวีสตีนํ\footnote{พาวีสติยา (?)} อินฺทฺริยานํ อุปาทาย พาวีสตีนํ1 ปุคฺคลานํ ปญฺญตฺตีติ, น เหวํ วตฺตพฺเพ ฯเปฯ.}

\proseitem{111}{\csromanfolio{38}\csromansymbolmark{*}เอกโวการภวํ อุปาทาย เอกสฺส ปุคฺคลสฺส ปญฺญตฺตีติ, อามนฺตา.\csromanspacer{} จตุโวการภวํ อุปาทาย จตุนฺนํ ปุคฺคลานํ ปญฺญตฺตีติ, น เหวํ วตฺตพฺเพ ฯเปฯ.\csromanspacer{} เอกโวการภวํ อุปาทาย เอกสฺส ปุคฺคลสฺส ปญฺญตฺตีติ, อามนฺตา.\csromanspacer{} ปญฺจโวการภวํ อุปาทาย ปญฺจนฺนํ ปุคฺคลานํ ปญฺญตฺตีติ, น เหวํ วตฺตพฺเพ ฯเปฯ.\csromanspacer{} เอกโวการภเว เอโกว ปุคฺคโลติ, อามนฺตา.\csromanspacer{} จตุโวการภเว จตฺตาโรว\footnote{จตฺตาโร (?)} ปุคฺคลาติ, น เหวํ วตฺตพฺเพ ฯเปฯ.\csromanspacer{} เอกโวการภเว เอโกว ปุคฺคโลติ, อามนฺตา.\csromanspacer{} ปญฺจโวการภเว ปญฺเจว ปุคฺคลาติ, น เหวํ วตฺตพฺเพ ฯเปฯ.}

\proseitem{112}{\csromansymbolmark{*}ยถา รุกฺขํ อุปาทาย ฉายาย ปญฺญตฺติ, เอวเมวํ รูปํ อุปาทาย ปุคฺคลสฺส ปญฺญตฺตีติ. ( )\footnote{(อามนฺตา) (?) เอวํ อนนฺตรวารตฺตเยปิ.} ยถา รุกฺขํ อุปาทาย ฉายาย ปญฺญตฺติ รุกฺโขปิ อนิจฺโจ ฉายาปิ อนิจฺจา, เอวเมวํ รูปํ อุปาทาย ปุคฺคลสฺส ปญฺญตฺติ รูปมฺปิ อนิจฺจํ ปุคฺคโลปิ อนิจฺโจติ, น เหวํ วตฺตพฺเพ ฯเปฯ.\csromanspacer{} ยถา รุกฺขํ อุปาทาย ฉายาย ปญฺญตฺติ อญฺโญ รุกฺโข, อญฺญา ฉายา, เอวเมวํ รูปํ อุปาทาย ปุคฺคลสฺส ปญฺญตฺติ อญฺญํ รูปํ, อญฺโญ ปุคฺคโลติ, น เหวํ วตฺตพฺเพ ฯเปฯ.}

\proseitem{113}{\csromansymbolmark{*}ยถา คามํ อุปาทาย คามิกสฺส ปญฺญตฺติ, เอวเมวํ รูปํ อุปาทาย ปุคฺคลสฺส ปญฺญตฺตีติ.\csromanspacer{} ยถา คามํ อุปาทาย คามิกสฺส ปญฺญตฺติ อญฺโญ คาโม, อญฺโญ คามิโก, เอวเมวํ รูปํ อุปาทาย ปุคฺคลสฺส ปญฺญตฺติ อญฺญํ รูปํ, อญฺโญ ปุคฺคโลติ, น เหวํ วตฺตพฺเพ ฯเปฯ.}

\proseitem{114}{\csromansymbolmark{*}ยถา รฏฺฐํ อุปาทาย รญฺโญ ปญฺญตฺติ, เอวเมวํ รูปํ อุปาทาย ปุคฺคลสฺส ปญฺญตฺตีติ.\csromanspacer{} ยถา รฏฺฐํ อุปาทาย รญฺโญ ปญฺญตฺติ อญฺญํ รฏฺฐํ, อญฺโญ ราชา, เอวเมวํ รูปํ อุปาทาย ปุคฺคลสฺส ปญฺญตฺติ อญฺญํ รูปํ, อญฺโญ ปุคฺคโลติ, น เหวํ วตฺตพฺเพ ฯเปฯ.}

\proseitem{115}{\csromansymbolmark{*}ยถา น นิคโฬ เนคฬิโก, ยสฺส นิคโฬ โส เนคฬิโก, เอวเมวํ น รูปํ รูปวา, ยสฺส รูปํ โส รูปวาติ.\csromanspacer{} ยถา น นิคโฬ เนคฬิโก, ยสฺส นิคโฬ โส เนคฬิโก, อญฺโญ นิคโฬ, อญฺโญ เนคฬิโท, เอวเมวํ น รูปํ รูปวา, ยสฺส รูปํ โส รูปวา, อญฺญํ รูปํ, อญฺโญ รูปวาติ, น เหวํ วตฺตพฺเพ ฯเปฯ.}

\proseitem{116}{\csromanfolio{39}\csromansymbolmark{*}จิตฺเต จิตฺเต ปุคฺคลสฺส ปญฺญตฺตีติ, อามนฺตา.\csromanspacer{} จิตฺเต จิตฺเต ปุคฺคโล ชายติ ชียติ มียติ จวติ อุปปชฺชตีติ, น เหวํ วตฺตพฺเพ ฯเปฯ.\csromanspacer{} ทุติเย จิตฺเต อุปฺปนฺเน น วตฺตพฺพํ “โส”ติ วา “อญฺโญ”ติ วาติ, อามนฺตา.\csromanspacer{} ทุติเย จิตฺเต อุปฺปนฺเน น วตฺตพฺพํ “กุมารโก”ติ วา “กุมาริกา”ติ วาติ, วตฺตพฺพํ.}

\prose{อาชานาหิ นิคฺคหํ\csromandash{}หญฺจิ ทุติเย จิตฺเต อุปฺปนฺเน น วตฺตพฺพํ “โส”ติ วา “อญฺโญ”ติ วา, เตน วต เร วตฺตพฺเพ “ทุติเย จิตฺเต อุปฺปนฺเน น วตฺตพฺพํ ‘กุมารโก’ติ วา ‘กุมาริกา’ติ วา”ติ.\csromanspacer{} ยํ ตตฺถ วเทสิ “วตฺตพฺเพ โข ‘ทุติเย จิตฺเต อุปฺปนฺเน น วตฺตพฺพํ ‘โส’ติ วา ‘อญฺโญ’ติ วา, ทุติเย จิตฺเต อุปฺปนฺเน วตฺตพฺพํ ‘กุมารโก’ติ วา ‘กุมาริกา’ติ วา’ติ”, มิจฺฉา.}

\prose{หญฺจิ วา ปน ทุติเย จิตฺเต อุปฺปนฺเน วตฺตพฺพํ “กุมารโก”ติ วา “กุมาริกา”ติ วา, เตน วต เร วตฺตพฺเพ “ทุติเย จิตฺเต อุปฺปนฺเน วตฺตพฺพํ ‘โส’ติ วา ‘อญฺโญ’ติ วา”ติ.\csromanspacer{} ยํ ตตฺถ วเทสิ “วตฺตพฺเพ โข ‘ทุติเย จิตฺเต อุปฺปนฺเน น วตฺตพฺพํ ‘โส’ติ วา ‘อญฺโญ’ติ วา ทุติเย จิตฺเต อุปฺปนฺเน วตฺตพฺพํ ‘กุมารโก’ติ วา ‘กุมาริกา’ติ วา’ติ”, มิจฺฉา.}

\proseitem{117}{\csromansymbolmark{*}ทุติเย จิตฺเต อุปฺปนฺเน น วตฺตพฺพํ “โส”ติ วา “อญฺโญ”ติ วาติ, อามนฺตา.\csromanspacer{} ทุติเย จิตฺเต อุปฺปนฺเน น วตฺตพฺพํ “อิตฺถี”ติ วา “ปุริโส”ติ วา “คหฏฺโฐ”ติ วา “ปพฺพชิโต”ติ วา “เทโว”ติ วา “มนุสฺโส”ติ วาติ, วตฺตพฺพํ.}

\prose{อาชานาหิ นิคฺคหํ\csromandash{}หญฺจิ ทุติเย จิตฺเต อุปฺปนฺเน น วตฺตพฺพํ “โส”ติ วา “อญฺโญ”ติ วา, เตน วต เร วตฺตพฺเพ “ทุติเย จิตฺเต อุปฺปนฺเน น วตฺตพฺพํ ‘เทโว’ติ วา ‘มนุสฺโส’ติ วา”ติ.\csromanspacer{} ยํ ตตฺถ วเทสิ “วตฺตพฺเพ โข ‘ทุติเย จิตฺเต อุปฺปนฺเน น วตฺตพฺพํ ‘โส’ติ วา ‘อญฺโญ’ติ วา, ทุติเย จิตฺเต อุปฺปนฺเน วตฺตพฺพํ ‘เทโว’ติ วา ‘มนุสฺโส’ติ วา’ติ”, มิจฺฉา.}

\prose{หญฺจิ วา ปน ทุติเย จิตฺเต อุปฺปนฺเน วตฺตพฺพํ “เทโว”ติ วา “มนุสฺโส”ติ วา, เตน วต เร วตฺตพฺเพ “ทุติเย จิตฺเต อุปฺปนฺเน วตฺตพฺพํ ‘โส’ติ วา ‘อญฺโญ’ติ วา”ติ.\csromanspacer{} ยํ ตตฺถ วเทสิ “วตฺตพฺเพ โข ‘ทุติเย จิตฺเต อุปฺปนฺเน น วตฺตพฺพํ ‘โส’ติ วา ‘อญฺโญ’ติ วา, ทุติเย จิตฺเต อุปฺปนฺเน วตฺตพฺพํ ‘เทโว’ติ วา ‘มนุสฺโส’ติ วา’ติ”, มิจฺฉา ฯเปฯ.}

\proseitem{118}{\csromanfolio{40}\csromansymbolmark{+}น วตฺตพฺพํ “ปุคฺคโล อุปลพฺภติ สจฺจิกฏฺฐปรมตฺเถนา”ติ, อามนฺตา.\csromanspacer{} นนุ โย ปสฺสติ, ยํ ปสฺสติ, เยน ปสฺสติ.\csromanspacer{} โส ปสฺสติ, ตํ ปสฺสติ, เตน ปสฺสตีติ, อามนฺตา.\csromanspacer{} หญฺจิ โย ปสฺสติ, ยํ ปสฺสติ, เยน ปสฺสติ.\csromanspacer{} โส ปสฺสติ, ตํ ปสฺสติ, เตน ปสฺสติ.\csromanspacer{} เตน วต เร วตฺตพฺเพ “ปุคฺคโล อุปลพฺภติ สจฺจิกฏฺฐปรมตฺเถนา”ติ.}

\proseitem{119}{\csromansymbolmark{+}น วตฺตพฺพํ “ปุคฺคโล อุปลพฺภติ สจฺจิกฏฺฐปรมตฺเถนา”ติ, อามนฺตา.\csromanspacer{} นนุ โย สุณาติ ฯเปฯ.\csromanspacer{} โย ฆายติ.\csromanspacer{} โย สายติ.\csromanspacer{} โย ผุสติ.\csromanspacer{} โย วิชานาติ, ยํ วิชานาติ, เยน วิชานาติ.\csromanspacer{} โส วิชานาติ, ตํ วิชานาติ, เตน วิชานาตีติ, อามนฺตา.\csromanspacer{} หญฺจิ โย วิชานาติ, ยํ วิชานาติ, เยน วิชานาติ.\csromanspacer{} โส วิชานาติ, ตํ วิชานาติ, เตน วิชานาติ.\csromanspacer{} เตน วต เร วตฺตพฺเพ “ปุคฺคโล อุปลพฺภติ สจฺจิกฏฺฐปรมตฺเถนา”ติ.}

\proseitem{120}{\csromansymbolmark{*}ปุคฺคโล อุปลพฺภติ สจฺจิกฏฺฐปรมตฺเถนาติ, อามนฺตา.\csromanspacer{} นนุ โย น ปสฺสติ, ยํ น ปสฺสติ, เยน น ปสฺสติ.\csromanspacer{} โส น ปสฺสติ, ตํ น ปสฺสติ, เตน น ปสฺสตีติ, อามนฺตา.\csromanspacer{} หญฺจิ โย น ปสฺสติ, ยํ น ปสฺสติ, เยน น ปสฺสติ.\csromanspacer{} โส น ปสฺสติ, ตํ น ปสฺสติ, เตน น ปสฺสติ.\csromanspacer{} โน จ วต เร วตฺตพฺเพ “ปุคฺคโล อุปลพฺภติ สจฺจิกฏฺฐปรมตฺเถนา”ติ.}

\prose{ปุคฺคโล อุปลพฺภติ สจฺจิกฏฺฐปรมตฺเถนาติ, อามนฺตา.\csromanspacer{} นนุ โย น สุณาติ ฯเปฯ.\csromanspacer{} โย น ฆายติ.\csromanspacer{} โย น สายติ.\csromanspacer{} โย น ผุสติ.\csromanspacer{} โย น วิชานาติ, ยํ น วิชานาติ, เยน น วิชานาติ.\csromanspacer{} โส น วิชานาติ, ตํ น วิชานาติ, เตน น วิชานาตีติ, อามนฺตา.\csromanspacer{} หญฺจิ โย น วิชานาติ, ยํ น วิชานาติ, เยน น วิชานาติ.\csromanspacer{} โส น วิชานาติ, ตํ น วิชานาติ, เตน น วิชานาติ.\csromanspacer{} โน จ วต เร วตฺตพฺเพ “ปุคฺคโล อุปลพฺภติ สจฺจิกฏฺฐปรมตฺเถนา”ติ.}

\proseitem{121}{\csromansymbolmark{+}น วตฺตพฺพํ “ปุคฺคโล อุปลพฺภติ สจฺจิกฏฺฐปรมตฺเถนา”ติ, อามนฺตา.\csromanspacer{} นนุ วุตฺตํ ภควตา “ปสฺสามหํ ภิกฺขเว ทิพฺเพน จกฺขุนา วิสุทฺเธน อติกฺกนฺตมานุสเกน สตฺเต จวมาเน อุปปชฺชมาเน หีเน ปณีเต สุวณฺเณ ทุพฺพณฺเณ สุคเต ทุคฺคเต ยถากมฺมูปเค สตฺเต ปชานามี”ติ\footnote{ม 1. 165 ปิฏฺเฐ โถกํ วิสทิสํ.} อตฺเถว สุตฺตนฺโตติ, อามนฺตา.\csromanspacer{} เตน หิ ปุคฺคโล อุปลพฺภติ สจฺจิกฏฺฐปรมตฺเถนาติ.}

\proseitem{122}{\csromanfolio{41}\csromansymbolmark{*}วุตฺตํ ภควตา “ปสฺสามหํ ภิกฺขเว ทิพฺเพน จกฺขุนา วิสุทฺเธน อติกฺกนฺตมานุสเกน สตฺเต จวมาเน อุปปชฺชมาเน หีเน ปณีเต สุวณฺเณ ทุพฺพณฺเณ สุคเต ทุคฺคเต ยถากมฺมูปเค สตฺเต ปชานามี”ติ กตฺวา เตเนว การเณน ปุคฺคโล อุปลพฺภติ สจฺจิกฏฺฐปรมตฺเถนาติ, อามนฺตา.\csromanspacer{} ภควา ทิพฺเพน จกฺขุนา วิสุทฺเธน อติกฺกนฺตมานุสเกน รูปํ ปสฺสติ ปุคฺคลํ ปสฺสตีติ, รูปํ ปสฺสติ.\csromanspacer{} รูปํ ปุคฺคโล, รูปํ จวติ, รูปํ อุปปชฺชติ, รูปํ ยถากมฺมูปคนฺติ, น เหวํ วตฺตพฺเพ.}

\prose{ภควา ทิพฺเพน จกฺขุนา วิสุทฺเธน อติกฺกนฺตมานุสเกน รูปํ ปสฺสติ ปุคฺคลํ ปสฺสตีติ, ปุคฺคลํ ปสฺสติ.\csromanspacer{} ปุคฺคโล รูปํ รูปายตนํ รูปธาตุ นีลํ ปีตกํ โลหิตกํ โอทาตํ จกฺขุวิญฺเญยฺยํ จกฺขุสฺมิํ ปฏิหญฺญติ, จกฺขุสฺส อาปาถํ อาคจฺฉตีติ, น เหวํ วตฺตพฺเพ.}

\prose{ภควา ทิพฺเพน จกฺขุนา วิสุทฺเธน อติกฺกนฺตมานุสเกน รูปํ ปสฺสติ ปุคฺคลํ ปสฺสตีติ, อุโภ ปสฺสติ.\csromanspacer{} อุโภ รูปํ รูปายตนํ รูปธาตุ, อุโภ นีลา, อุโภ ปีตกา, อุโภ โลหิตกา, อุโภ โอทาตา, อุโภ จกฺขุวิญฺเญยฺยา, อุโภ จกฺขุสฺมิํ ปฏิหญฺญนฺติ, อุโภ จกฺขุสฺส อาปาถํ อาคจฺฉนฺติ, อุโภ จวนฺติ, อุโภ อุปปชฺชนฺติ, อุโภ ยถากมฺมูปคาติ, น เหวํ วตฺตพฺเพ.}

\vspace{\tipitakaparskip}
\csromancenter{อุปาทาปญฺญตฺตานุโยโค.}
\csromanmatikaanchor{mtk.14}
\csromantocmarkref{section}{13. ปุริสการานุโยค}{mtk.14}
\bookmark[dest={mtk.14},level=1]{13. ปุริสการานุโยค}
\csromanheader{1}{13. ปุริสการานุโยค}
\proseitem{123}{\csromansymbolmark{+}กลฺยาณปาปกานิ กมฺมานิ อุปลพฺภนฺตีติ, อามนฺตา.\csromanspacer{} กลฺยาณปาปกานํ กมฺมานํ กตฺตา กาเรตา อุปลพฺภตีติ, น เหวํ วตฺตพฺเพ ฯเปฯ.}

\proseitem{124}{\csromansymbolmark{*}กลฺยาณปาปกานิ กมฺมานิ อุปลพฺภนฺตีติ กลฺยาณปาปกานํ กมฺมานํ กตฺตา กาเรตา อุปลพฺภตีติ, อามนฺตา.\csromanspacer{} ตสฺส กตฺตา กาเรตา อุปลพฺภตีติ, น เหวํ วตฺตพฺเพ ฯเปฯ.}

\proseitem{125}{\csromansymbolmark{*}ตสฺส กตฺตา กาเรตา อุปลพฺภตีติ, อามนฺตา.\csromanspacer{} ตสฺส ตสฺเสว นตฺถิ ทุกฺขสฺส อนฺตกิริยา, นตฺถิ วฏฺฏุปจฺเฉโท, นตฺถิ อนุปาทาปรินิพฺพานนฺติ, น เหวํ วตฺตพฺเพ ฯเปฯ.}

\proseitem{126}{\csromanfolio{42}\csromansymbolmark{*}กลฺยาณปาปกานิ กมฺมานิ อุปลพฺภนฺตีติ กลฺยาณปาปกานํ กมฺมานํ กตฺตา กาเรตา อุปลพฺภตีติ, อามนฺตา.\csromanspacer{} ปุคฺคโล อุปลพฺภตีติ ปุคฺคลสฺส กตฺตา กาเรตา อุปลพฺภตีติ, น เหวํ วตฺตพฺเพ ฯเปฯ.}

\proseitem{127}{\csromansymbolmark{*}กลฺยาณปาปกานิ กมฺมานิ อุปลพฺภนฺตีติ กลฺยาณปาปกานํ กมฺมานํ กตฺตา กาเรตา อุปลพฺภตีติ, อามนฺตา.\csromanspacer{} นิพฺพานํ อุปลพฺภตีติ นิพฺพานสฺส กตฺตา กาเรตา อุปลพฺภตีติ, น เหวํ วตฺตพฺเพ ฯเปฯ.}

\proseitem{128}{\csromansymbolmark{*}กลฺยาณปาปกานิ กมฺมานิ อุปลพฺภนฺตีติ กลฺยาณปาปกนํ กมฺมานํ กตฺตา กาเรตา อุปลพฺภตีติ, อามนฺตา.\csromanspacer{} มหาปถวี อุปลพฺภตีติ มหาปถวิยา กตฺตา กาเรตา อุปลพฺภตีติ, น เหวํ วตฺตพฺเพ ฯเปฯ.}

\proseitem{129}{\csromansymbolmark{*}กลฺยาณปาปกานิ กมฺมานิ อุปลพฺภนฺตีติ กลฺยาณปาปกานํ กมฺมานํ กตฺตา กาเรตา อุปลพฺภตีติ, อามนฺตา.\csromanspacer{} มหาสมุทฺโท อุปลพฺภตีติ มหาสมุทฺทสฺส กตฺตา กาเรตา อุปลพฺภตีติ, น เหวํ วตฺตพฺเพ ฯเปฯ.}

\proseitem{130}{\csromansymbolmark{*}กลฺยาณปาปกานิ กมฺมานิ อุปลพฺภนฺตีติ กลฺยาณปาปกานํ กมฺมานํ กตฺตา กาเรตา อุปลพฺภตีติ, อามนฺตา.\csromanspacer{} สิเนรุปพฺพตราชา อุปลพฺภตีติ สิเนรุสฺส ปพฺพตราชสฺส กตฺตา กาเรตา อุปลพฺภตีติ, น เหวํ วตฺตพฺเพ ฯเปฯ.}

\proseitem{131}{\csromansymbolmark{*}กลฺยาณปาปกานิ กมฺมานิ อุปลพฺภนฺตีติ กลฺยาณปาปกานํ กมฺมานํ กตฺตา กาเรตา อุปลพฺภตีติ, อามนฺตา.\csromanspacer{} อาโป อุปลพฺภตีติ อาปสฺส กตฺตา กาเรตา อุปลพฺภตีติ น เหวํ วตฺตพฺเพ ฯเปฯ.}

\proseitem{132}{\csromansymbolmark{*}กลฺยาณปาปกานิ กมฺมานิ อุปลพฺภนฺตีติ กลฺยาณปาปกานํ กมฺมานํ กตฺตา กาเรตา อุปลพฺภตีติ, อามนฺตา.\csromanspacer{} เตโช อุปลพฺภตีติ เตชสฺส กตฺตา กาเรตา อุปลพฺภตีติ, น เหวํ วตฺตพฺเพ ฯเปฯ.}

\proseitem{133}{\csromansymbolmark{*}กลฺยาณปาปกานิ กมฺมานิ อุปลพฺภนฺตีติ กลฺยาณปาปกานํ กมฺมานํ กตฺตา กาเรตา อุปลพฺภตีติ, อามนฺตา.\csromanspacer{} วาโย อุปลพฺภตีติ วายสฺส กตฺตา กาเรตา อุปลพฺภตีติ, น เหวํ วตฺตพฺเพ ฯเปฯ.}

\proseitem{134}{\csromansymbolmark{*}กลฺยาณปาปกานิ กมฺมานิ อุปลพฺภนฺตีติ กลฺยาณปาปกานํ กมฺมานํ กตฺตา กาเรตา อุปลพฺภตีติ, อามนฺตา.\csromanspacer{} ติณกฏฺฐวนปฺปตโย \csromanfolio{43}อุปลพฺภนฺตีติ ติณกฏฺฐวนปฺปตีนํ กตฺตา กาเรตา อุปลพฺภตีติ, น เหวํ วตฺตพฺเพ ฯเปฯ.}

\proseitem{135}{\csromansymbolmark{*}กลฺยาณปาปกานิ กมฺมานิ อุปลพฺภนฺตีติ กลฺยาณปาปกานํ กมฺมานํ กตฺตา กาเรกา อุปลพฺภตีติ, อามนฺตา.\csromanspacer{} อญฺญานิ กลฺยาณปาปกานิ กมฺมานิ, อญฺโญ กลฺยาณปาปกานํ กมฺมานํ กตฺตา กาเรตาติ, น เหวํ วตฺตพฺเพ ฯเปฯ.}

\proseitem{136}{\csromansymbolmark{+}กลฺยาณปาปกานํ กมฺมานํ วิปาโก อุปลพฺภตีติ, อามนฺตา.\csromanspacer{} กลฺยาณปาปกานํ กมฺมานํ วิปากปฏิสํเวที อุปลพฺภตีติ, น เหวํ วตฺตพฺเพ ฯเปฯ.}

\proseitem{137}{\csromansymbolmark{*}กลฺยาณปาปกานํ กมฺมานํ วิปาโก อุปลพฺภตีติ กลฺยาณปาปกานํ กมฺมานํ วิปากปฏิสํเวที อุปลพฺภตีติ, อามนฺตา.\csromanspacer{} ตสฺส ปฏิสํเวที อุปลพฺภตีติ, น เหวํ วตฺตพฺเพ ฯเปฯ.}

\proseitem{138}{\csromansymbolmark{*}ตสฺส ปฏิสํเวที อุปลพฺภตีติ, อามนฺตา.\csromanspacer{} ตสฺส ตสฺเสว นตฺถิ ทุกฺขสฺส อนฺตกิริยา, นตฺถิ วฏฺฏุปจฺเฉโท, นตฺถิ อนุปาทาปรินิพฺพานนฺติ, น เหวํ วตฺตพฺเพ ฯเปฯ.}

\proseitem{139}{\csromansymbolmark{*}กลฺยาณปาปกานํ กมฺมานํ วิปาโก อุปลพฺภตีติ กลฺยาณปาปกานํ กมฺมานํ วิปากปฏิสํเวที อุปลพฺภตีติ, อามนฺตา.\csromanspacer{} ปุคฺคโล อุปลพฺภตีติ ปุคฺคลสฺส ปฏิสํเวที อุปลพฺภตีติ, น เหวํ วตฺตพฺเพ ฯเปฯ.}

\proseitem{140}{\csromansymbolmark{*}กลฺยาณปาปกานํ กมฺมานํ วิปาโก อุปลพฺภตีติ กลฺยาณปาปกานํ กมฺมานํ วิปากปฏิสํเวที อุปลพฺภตีติ, อามนฺตา.\csromanspacer{} นิพฺพานํ อุปลพฺภตีติ นิพฺพานสฺส ปฏิสํเวที อุปลพฺภตีติ, น เหวํ วตฺตพฺเพ ฯเปฯ.}

\proseitem{141}{\csromansymbolmark{*}กลฺยาณปาปกานํ กมฺมานํ วิปาโก อุปลพฺภตีติ กลฺยาณปาปกานํ กมฺมานํ วิปากปฏิสํเวที อุปลพฺภตีติ, อามนฺตา.\csromanspacer{} มหาปถวี อุปลพฺภตีติ ฯเปฯ.\csromanspacer{} มหาสมุทฺโท อุปลพฺภตีติ.\csromanspacer{} สิเนรุปพฺพตราชา อุปลพฺภตีติ.\csromanspacer{} อาโป อุปลพฺภตีติ.\csromanspacer{} เตโช อุปลพฺภตีติ.\csromanspacer{} วาโย อุปลพฺภตีติ ฯเปฯ.\csromanspacer{} ติณกฏฺฐวนปฺปตโย อุปลพฺภนฺตีติ ติณกฏฺฐวนปฺปตีนํ ปฏิสํเวที อุปลพฺภตีติ, น เหวํ วตฺตพฺเพ ฯเปฯ.}

\proseitem{142}{\csromanfolio{44}\csromansymbolmark{*}กลฺยาณปาปกานํ กมฺมานํ วิปาโก อุปลพฺภตีติ กลฺยาณปาปกานํ กมฺมานํ วิปากปฏิสํเวที อุปลพฺภตีติ, อามนฺตา.\csromanspacer{} อญฺโญ กลฺยาณปาปกานํ กมฺมานํ วิปาโก, อญฺโญ กลฺยาณปาปกานํ กมฺมานํ วิปากปฏิสํเวทีติ, น เหวํ วตฺตพฺเพ ฯเปฯ.}

\proseitem{143}{\csromansymbolmark{+}ทิพฺพํ สุขํ อุปลพฺภตีติ, อามนฺตา.\csromanspacer{} ทิพฺพสฺส สุขสฺส ปฏิสํเวที อุปลพฺภตีติ, น เหวํ วตฺตพฺเพ ฯเปฯ.}

\proseitem{144}{\csromansymbolmark{*}ทิพฺพํ สุขํ อุปลพฺภตีติ ทิพฺพสฺส สุขสฺส ปฏิสํเวที อุปลพฺภตีติ, อามนฺตา.\csromanspacer{} ตสฺส ปฏิสํเวที อุปลพฺภตีติ, น เหวํ วตฺตพฺเพ ฯเปฯ.}

\proseitem{145}{\csromansymbolmark{*}ตสฺส ปฏิสํเวที อุปลพฺภตีติ, อามนฺตา.\csromanspacer{} ตสฺส ตสฺเสว นตฺถิ ทุกฺขสฺส อนฺตกิริยา, นตฺถิ วฏฺฏุปจฺเฉโท, นตฺถิ อนุปาทาปรินิพฺพานนฺติ, น เหวํ วตฺตพฺเพ ฯเปฯ.}

\proseitem{146}{\csromansymbolmark{*}ทิพฺพํ สุขํ อุปลพฺภตีติ ทิพฺพสฺส สุขสฺส ปฏิสํเวที อุปลพฺภตีติ, อามนฺตา.\csromanspacer{} ปุคฺคโล อุปลพฺภตีติ, ปุคฺคลสฺส ปฏิสํเวที อุปลพฺภตีติ, น เหวํ วตฺตพฺเพ ฯเปฯ.}

\proseitem{147}{\csromansymbolmark{*}ทิพฺพํ สุขํ อุปลพฺภตีติ ทิพฺพสฺส สุขสฺส ปฏิสํเวที อุปลพฺภตีติ, อามนฺตา.\csromanspacer{} นิพฺพานํ อุปลพฺภตีติ นิพฺพานสฺส ปฏิสํเวที อุปลพฺภตีติ, น เหวํ วตฺตพฺเพ ฯเปฯ.}

\proseitem{148}{\csromansymbolmark{*}ทิพฺพํ สุขํ อุปลพฺภตีติ ทิพฺพสฺส สุขสฺส ปฏิสํเวที อุปลพฺภตีติ, อามนฺตา.\csromanspacer{} มหาปถวี อุปลพฺภตีติ.\csromanspacer{} มหาสมุทฺโท อุปลพฺภตีติ.\csromanspacer{} สิเนรุปพฺพตราชา อุปลพฺภตีติ.\csromanspacer{} อาโป อุปลพฺภตีติ.\csromanspacer{} เตโช อุปลพฺภตีติ.\csromanspacer{} วาโย อุปลพฺภตีติ ฯเปฯ.\csromanspacer{} ติณกฏฺฐวนปฺปตโย อุปลพฺภนฺตีติ ติณกฏฺฐวนปฺปตีนํ ปฏิสํเวที อุปลพฺภตีติ, น เหวํ วตฺตพฺเพ ฯเปฯ.}

\proseitem{149}{\csromansymbolmark{*}ทิพฺพํ สุขํ อุปลพฺภตีติ ทิพฺพสฺส สุขสฺส ปฏิสํเวที อุปลพฺภตีติ, อามนฺตา.\csromanspacer{} อญฺญํ ทิพฺพํ สุขํ, อญฺโญ ทิพฺพสฺส สุขสฺส ปฏิสํเวทีติ, น เหวํ วตฺตพฺเพ ฯเปฯ.}

\proseitem{150}{\csromansymbolmark{+}มานุสกํ สุขํ อุปลพฺภตีติ, อามนฺตา.\csromanspacer{} มานุสกสฺส สุขสฺส ปฏิสํเวที อุปลพฺภตีติ, น เหวํ วตฺตพฺเพ ฯเปฯ.}

\proseitem{151}{\csromanfolio{45}\csromansymbolmark{*}มานุสกํ สุขํ อุปลพฺภตีติ มานุสกสฺส สุขสฺส ปฏิสํเวที อุปลพฺภตีติ, อามนฺตา.\csromanspacer{} ตสฺส ปฏิสํเวที อุปลพฺภตีติ, น เหวํ วตฺตพฺเพ ฯเปฯ.}

\proseitem{152}{\csromansymbolmark{*}ตสฺส ปฏิสํเวที อุปลพฺภตีติ, อามนฺตา.\csromanspacer{} ตสฺส ตสฺเสว นตฺถิ ทุกฺขสฺส อนฺตกิริยา, นตฺถิ วฏฺฏุปจฺเฉโท, นตฺถิ อนุปาทาปรินิพฺพานนฺติ, น เหวํ วตฺตพฺเพ ฯเปฯ.}

\proseitem{153}{\csromansymbolmark{*}มานุสกํ สุขํ อุปลพฺภตีติ มานุสกสฺส สุขสฺส ปฏิสํเวที อุปลพฺภตีติ, อามนฺตา.\csromanspacer{} ปุคฺคโล อุปลพฺภตีติ ปุคฺคลสฺส ปฏิสํเวที อุปลพฺภตีติ, น เหวํ วตฺตพฺเพ ฯเปฯ.}

\proseitem{154}{\csromansymbolmark{*}มานุสกํ สุขํ อุปลพฺภตีติ มานุสกสฺส สุขสฺส ปฏิสํเวที อุปลพฺภตีติ, อามนฺตา.\csromanspacer{} นิพฺพานํ อุปลพฺภตีติ นิพฺพานสฺส ปฏิสํเวที อุปลพฺภตีติ, น เหวํ วตฺตพฺเพ ฯเปฯ.}

\proseitem{155}{\csromansymbolmark{*}มานุสกํ สุขํ อุปลพฺภตีติ มานุสกสฺส สุขสฺส ปฏิสํเวที อุปลพฺภตีติ, อามนฺตา.\csromanspacer{} มหาปถวี อุปลพฺภตีติ ฯเปฯ.\csromanspacer{} มหาสมุทฺโท อุปลพฺภตีติ.\csromanspacer{} สิเนรุปพฺพตราชา อุปลพฺภตีติ.\csromanspacer{} อาโป อุปลพฺภตีติ.\csromanspacer{} เตโช อุปลพฺภตีติ.\csromanspacer{} วาโย อุปลพฺภตีติ.\csromanspacer{} ติณกฏฺฐวนปฺปตโย อุปลพฺภนฺตีติ ติณกฏฺฐวนปฺปตีนํ ปฏิสํเวที อุปลพฺภตีติ, น เหวํ วตฺตพฺเพ ฯเปฯ.}

\proseitem{156}{\csromansymbolmark{*}มานุสกํ สุขํ อุปลพฺภตีติ มานุสกสฺส สุขสฺส ปฏิสํเวที อุปลพฺภตีติ, อามนฺตา.\csromanspacer{} อญฺญํ มานุสกํ สุขํ, อญฺโญ มานุสกสฺส สุขสฺส ปฏิสํเวทีติ, น เหวํ วตฺตพฺเพ ฯเปฯ.}

\proseitem{157}{\csromansymbolmark{+}อาปายิกํ ทุกฺขํ อุปลพฺภตีติ, อามนฺตา.\csromanspacer{} อาปายิกสฺส ทุกฺขสฺส ปฏิสํเวที อุปลพฺภตีติ, น เหวํ วตฺตพฺเพ ฯเปฯ.}

\proseitem{158}{\csromansymbolmark{*}อาปายิกํ ทุกฺขํ อุปลพฺภตีติ อาปายิกสฺส ทุกฺขสฺส ปฏิสํเวที อุปลพฺภตีติ, อามนฺตา.\csromanspacer{} ตสฺส ปฏิสํเวที อุปลพฺภตีติ, น เหวํ วตฺตพฺเพ ฯเปฯ.}

\proseitem{159}{\csromansymbolmark{*}ตสฺส ปฏิสํเวที อุปลพฺภตีติ, อามนฺตา.\csromanspacer{} ตสฺส ตสฺเสว นตฺถิ ทุกฺขสฺส อนฺตกิริยา, นตฺถิ วฏฺฏุปจฺเฉโท, นตฺถิ อนุปาทาปรินิพฺพานนฺติ, น เหวํ วตฺตพฺเพ ฯเปฯ.}

\proseitem{160}{\csromanfolio{46}\csromansymbolmark{*}อาปายิกํ ทุกฺขํ อุปลพฺภตีติ อาปายิกสฺส ทุกฺขสฺส ปฏิสํเวที อุปลพฺภตีติ, อามนฺตา.\csromanspacer{} ปุคฺคโล อุปลพฺภตีติ ปุคฺคลสฺส ปฏิสํเวที อุปลพฺภตีติ, น เหวํ วตฺตพฺเพ ฯเปฯ.}

\proseitem{161}{\csromansymbolmark{*}อาปายิกํ ทุกฺขํ อุปลพฺภตีติ อาปายิกสฺส ทุกฺขสฺส ปฏิสํเวที อุปลพฺภตีติ, อามนฺตา.\csromanspacer{} นิพฺพานํ อุปลพฺภตีติ นิพฺพานสฺส ปฏิสํเวที อุปลพฺภตีติ, น เหวํ วตฺตพฺเพ ฯเปฯ.}

\proseitem{162}{\csromansymbolmark{*}อาปายิกํ ทุกฺขํ อุปลพฺภตีติ อาปายิกสฺส ทุกฺขสฺส ปฏิสํเวที อุปลพฺภตีติ, อามนฺตา.\csromanspacer{} มหาปถวี อุปลพฺภตีติ ฯเปฯ.\csromanspacer{} มหาสมุทฺโท อุปลพฺภตีติ.\csromanspacer{} สิเนรุปพฺพตราชา อุปลพฺภตีติ.\csromanspacer{} อาโป อุปลพฺภตีติ.\csromanspacer{} เตโช อุปลพฺภตีติ.\csromanspacer{} วาโย อุปลพฺภตีติ.\csromanspacer{} ติณกฏฺฐวนปฺปตโย อุปลพฺภนฺตีติ ติณกฏฺฐวนปฺปตีนํ ปฏิสํเวที อุปลพฺภตีติ, น เหวํ วตฺตพฺเพ ฯเปฯ.}

\proseitem{163}{\csromansymbolmark{*}อาปายิกํ ทุกฺขํ อุปลพฺภตีติ อาปายิกสฺส ทุกฺขสฺส ปฏิสํเวที อุปลพฺภตีติ, อามนฺตา.\csromanspacer{} อญฺญํ อาปายิกํ ทุกฺขํ, อญฺโญ อาปายิกสฺส ทุกฺขสฺส ปฏิสํเวทีติ, น เหวํ วตฺตพฺเพ ฯเปฯ.}

\proseitem{164}{\csromansymbolmark{+}เนรยิกํ ทุกฺขํ อุปลพฺภตีติ, อามนฺตา.\csromanspacer{} เนรยิกสฺส ทุกฺขสฺส ปฏิสํเวที อุปลพฺภตีติ, น เหวํ วตฺตพฺเพ.}

\prose{เนรยิกํ ทุกฺขํ อุปลพฺภตีติ เนรยิกสฺส ทุกฺขสฺส ปฏิสํเวที อุปลพฺภตีติ, อามนฺตา.\csromanspacer{} ตสฺส ปฏิสํเวที อุปลพฺภตีติ, น เหวํ วตฺตพฺเพ ฯเปฯ.}

\proseitem{165}{\csromansymbolmark{*}ตสฺส ปฏิสํเวที อุปลพฺภตีติ, อามนฺตา.\csromanspacer{} ตสฺส ตสฺเสว นตฺถิ ทุกฺขสฺส อนฺตกิริยา, นตฺถิ วฏฺฏุปจฺเฉโท, นตฺถิ อนุปาทาปรินิพฺพานนฺติ, น เหวํ วตฺตพฺเพ ฯเปฯ.}

\proseitem{166}{\csromansymbolmark{*}เนรยิกํ ทุกฺขํ อุปลพฺภตีติ เนรยิกสฺส ทุกฺขสฺส ปฏิสํเวที อุปลพฺภตีติ, อามนฺตา.\csromanspacer{} ปุคฺคโล อุปลพฺภตีติ ปุคฺคลสฺส ปฏิสํเวที อุปลพฺภตีติ, น เหวํ วตฺตพฺเพ ฯเปฯ.}

\proseitem{167}{\csromansymbolmark{*}เนรยิกํ ทุกฺขํ อุปลพฺภตีติ เนรยิกสฺส ทุกฺขสฺส ปฏิสํเวที อุปลพฺภตีติ, อามนฺตา.\csromanspacer{} นิพฺพานํ อุปลพฺภตีติ นิพฺพานสฺส ปฏิสํเวที อุปลพฺภตีติ, น เหวํ วตฺตพฺเพ ฯเปฯ.}

\proseitem{168}{\csromanfolio{47}\csromansymbolmark{*}เนรยิกํ ทุกฺขํ อุปลพฺภตีติ เนรยิกสฺส ทุกฺขสฺส ปฏิสํเวที อุปลพฺภตีติ, อามนฺตา.\csromanspacer{} มหาปถวี อุปลพฺภตีติ ฯเปฯ.\csromanspacer{} มหาสมุทฺโท อุปลพฺภตีติ.\csromanspacer{} สิเนรุปพฺพตราชา อุปลพฺภตีติ.\csromanspacer{} อาโป อุปลพฺภตีติ.\csromanspacer{} เตโช อุปลพฺภตีติ.\csromanspacer{} วาโย อุปลพฺภตีติ.\csromanspacer{} ติณกฏฺฐวนปฺปตโย อุปลพฺภนฺตีติ ติณกฏฺฐวนปฺปตีนํ ปฏิสํเวที อุปลพฺภตีติ, น เหวํ วตฺตพฺเพ ฯเปฯ.}

\proseitem{169}{\csromansymbolmark{*}เนรยิกํ ทุกฺขํ อุปลพฺภตีติ เนรยิกสฺส ทุกฺขสฺส ปฏิสํเวที อุปลพฺภตีติ, อามนฺตา.\csromanspacer{} อญฺญํ เนรยิกํ ทุกฺขํ, อญฺโญ เนรยิกสฺส ทุกฺขสฺส ปฏิสํเวทีติ, น เหวํ วตฺตพฺเพ ฯเปฯ.}

\proseitem{170}{\csromansymbolmark{*}กลฺยาณปาปกานิ กมฺมานิ อุปลพฺภนฺตีติ กลฺยาณปาปกานํ กมฺมานํ กตฺตา กาเรตา วิปากปฏิสํเวที อุปลพฺภตีติ, อามนฺตา.\csromanspacer{} โส กโรติ, โส ปฏิสํเวเทตีติ, น เหวํ วตฺตพฺเพ ฯเปฯ.}

\proseitem{171}{\csromansymbolmark{*}โส กโรติ, โส ปฏิสํเวเทตีติ, อามนฺตา.\csromanspacer{} สยํกตํ สุขทุกฺขนฺติ, น เหวํ วตฺตพฺเพ ฯเปฯ.}

\proseitem{172}{\csromansymbolmark{*}กลฺยาณปาปกานิ กมฺมานิ อุปลพฺภนฺตีติ กลฺยาณปาปกานํ กมฺมานํ กตฺตา กาเรตา วิปากปฏิสํเวที อุปลพฺภตีติ, อามนฺตา.\csromanspacer{} อญฺโญ กโรติ, อญฺโญ ปฏิสํเวเทตีติ, น เหวํ วตฺตพฺเพ ฯเปฯ.}

\proseitem{173}{\csromansymbolmark{*}อญฺโญ กโรติ, อญฺโญ ปฏิสํเวเทตีติ, อามนฺตา.\csromanspacer{} ปรํกตํ สุขทุกฺขนฺติ, น เหวํ วตฺตพฺเพ ฯเปฯ.}

\proseitem{174}{\csromansymbolmark{*}กลฺยาณปาปกานิ กมฺมานิ อุปลพฺภนฺตีติ กลฺยาณปาปกานํ กมฺมานํ กตฺตา กาเรตา วิปากปฏิสํเวที อุปลพฺภตีติ, อามนฺตา.\csromanspacer{} โส จ อญฺโญ จ กโรนฺติ, โส จ อญฺโญ จ ปฏิสํเวเทนฺตีติ, น เหวํ วตฺตพฺเพ ฯเปฯ.}

\proseitem{175}{\csromansymbolmark{*}โส จ อญฺโญ จ กโรนฺติ, โส จ อญฺโญ จ ปฏิสํเวเทนฺตีติ, อามนฺตา.\csromanspacer{} สยํกตญฺจ ปรํกตญฺจ สุขทุกฺขนฺติ, น เหวํ วตฺตพฺเพ ฯเปฯ.}

\proseitem{176}{\csromansymbolmark{*}กลฺยาณปาปกานิ กมฺมานิ อุปลพฺภนฺตีติ กลฺยาณปาปกานํ กมฺมานํ กตฺตา กาเรตา วิปากปฏิสํเวที อุปลพฺภตีติ, อามนฺตา.\csromanspacer{} เนว \csromanfolio{48}โส กโรติ น โส ปฏิสํเวเทติ, น อญฺโญ กโรติ น อญฺโญ ปฏิสํเวเทตีติ, น เหวํ วตฺตพฺเพ ฯเปฯ.}

\proseitem{177}{\csromansymbolmark{*}เนวโส กโรติ น โส ปฏิสํเวเทติ, น อญฺโญ กโรติ น อญฺโญ ปฏิสํเวเทตีติ, อามนฺตา.\csromanspacer{} อสยํการํ อปรํการํ อธิจฺจสมุปฺปนฺนํ สุขทุกฺขนฺติ, น เหวํ วตฺตพฺเพ ฯเปฯ.}

\proseitem{178}{\csromansymbolmark{*}กลฺยาณปาปกานิ กมฺมานิ อุปลพฺภนฺตีติ กลฺยาณปาปกานํ กมฺมานํ กตฺตา กาเรตา วิปากปฏิสํเวที อุปลพฺภตีติ, อามนฺตา.\csromanspacer{} โส กโรติ, โส ปฏิสํเวเทติ, อญฺโญ กโรติ อญฺโญ ปฏิสํเวเทติ, โส จ อญฺโญ จ กโรนฺติ โส จ อญฺโญ จ ปฏิสํเวเทนฺติ, เนว โส กโรติ น โส ปฏิสํเวเทติ, น อญฺโญ กโรติ น อญฺโญ ปฏิสํเวเทตีติ, น เหวํ วตฺตพฺเพ ฯเปฯ.}

\proseitem{179}{\csromansymbolmark{*}โส กโรติ, โส ปฏิสํเวเทติ.\csromanspacer{} อญฺโญ กโรติ, อญฺโญ ปฏิสํเวเทติ.\csromanspacer{} โส จ อญฺโญ จ กโรนฺติ, โส จ อญฺโญ จ ปฏิสํเวเทนฺติ.\csromanspacer{} เนว โส กโรติ, น โส ปฏิสํเวเทติ.\csromanspacer{} น อญฺโญ กโรติ, น อญฺโญ ปฏิสํเวเทตีติ, อามนฺตา.\csromanspacer{} สยํกตํ สุขทุกฺขํ, ปรํกตํ สุขทุกฺขํ, สยํกตญฺจ ปรํกตญฺจ สุขทุกฺขํ, อสยํการํ อปรํการํ อธิจฺจสมุปฺปนฺนํ สุขทุกฺขนฺติ, น เหวํ วตฺตพฺเพ ฯเปฯ.}

\proseitem{180}{\csromansymbolmark{+}กมฺมํ อตฺถีติ, อามนฺตา.\csromanspacer{} กมฺมการโก อตฺถีติ, น เหวํ วตฺตพฺเพ ฯเปฯ.}

\proseitem{181}{\csromansymbolmark{*}กมฺมํ อตฺถีติ กมฺมการโก อตฺถีติ, อามนฺตา.\csromanspacer{} ตสฺส การโก อตฺถีติ, น เหวํ วตฺตพฺเพ ฯเปฯ.}

\proseitem{182}{\csromansymbolmark{*}ตสฺส การโก อตฺถีติ, อามนฺตา.\csromanspacer{} ตสฺส ตสฺเสว นตฺถิ ทุกฺขสฺส อนฺตกิริยา, นตฺถิ วฏฺฏุปจฺเฉโท, นตฺถิ อนุปาทาปรินิพฺพานนฺติ, น เหวํ วตฺตพฺเพ ฯเปฯ.}

\proseitem{183}{\csromansymbolmark{*}กมฺมํ อตฺถีติ กมฺมการโก อตฺถีติ, อามนฺตา.\csromanspacer{} ปุคฺคโล อตฺถีติ ปุคฺคลสฺส การโก อตฺถีติ, น เหวํ วตฺตพฺเพ ฯเปฯ.}

\proseitem{184}{\csromanfolio{49}\csromansymbolmark{*}กมฺมํ อตฺถีติ กมฺมการโก อตฺถีติ, อามนฺตา.\csromanspacer{} นิพฺพานํ อตฺถีติ นิพฺพานสฺส การโก อตฺถีติ, น เหวํ วตฺตพฺเพ ฯเปฯ.}

\proseitem{185}{\csromansymbolmark{*}กมฺมํ อตฺถีติ กมฺมการโก อตฺถีติ, อามนฺตา.\csromanspacer{} มหาปถวี อตฺถีติ ฯเปฯ.\csromanspacer{} มหาสมุทฺโท อตฺถีติ.\csromanspacer{} สิเนรุปพฺพตราชา อตฺถีติ.\csromanspacer{} อาโป อตฺถีติ.\csromanspacer{} เตโช อตฺถีติ.\csromanspacer{} วาโย อตฺถีติ.\csromanspacer{} ติณกฏฺฐวนปฺปตโย อตฺถีติ ติณกฏฺฐวนปฺปตีนํ การโก อตฺถีติ, น เหวํ วตฺตพฺเพ ฯเปฯ.}

\proseitem{186}{\csromansymbolmark{*}กมฺมํ อตฺถีติ กมฺมการโก อตฺถีติ, อามนฺตา.\csromanspacer{} อญฺญํ กมฺมํ, อญฺโญ กมฺมการโกติ, น เหวํ วตฺตพฺเพ ฯเปฯ.}

\proseitem{187}{\csromansymbolmark{+}วิปาโก อตฺถีติ, อามนฺตา.\csromanspacer{} วิปากปฏิสํเวที อตฺถีติ, น เหวํ วตฺตพฺเพ ฯเปฯ.}

\proseitem{188}{\csromansymbolmark{*}วิปาโก อตฺถีติ วิปากปฏิสํเวที อตฺถีติ, อามนฺตา.\csromanspacer{} ตสฺส ปฏิสํเวที อตฺถีติ, น เหวํ วตฺตพฺเพ ฯเปฯ.}

\proseitem{189}{\csromansymbolmark{*}ตสฺส ปฏิสํเวที อตฺถีติ, อามนฺตา.\csromanspacer{} ตสฺส ตสฺเสว นตฺถิ ทุกฺขสฺส อนฺตกิริยา, นตฺถิ วฏฺฏุปจฺเฉโท, นตฺถิ อนุปาทาปรินิพฺพานนฺติ, น เหวํ วตฺตพฺเพ ฯเปฯ.\csromanspacer{} วิปาโก อตฺถีติ วิปากปฏิสํเวที อตฺถีติ, อามนฺตา.\csromanspacer{} ปุคฺคโล อตฺถีติ ปุคฺคลสฺส ปฏิสํเวที อตฺถีติ, น เหวํ วตฺตพฺเพ ฯเปฯ.}

\proseitem{190}{\csromansymbolmark{*}วิปาโก อตฺถีติ วิปากปฏิสํเวที อตฺถีติ, อามนฺตา.\csromanspacer{} นิพฺพานํ อตฺถีติ นิพฺพานสฺส ปฏิสํเวที อตฺถีติ, น เหวํ วตฺตพฺเพ ฯเปฯ.}

\proseitem{191}{\csromansymbolmark{*}วิปาโก อตฺถีติ วิปากปฏิสํเวที อตฺถีติ, อามนฺตา.\csromanspacer{} มหาปถวี อตฺถีติ ฯเปฯ.\csromanspacer{} มหาสมุทฺโท อตฺถีติ.\csromanspacer{} สิเนรุปพฺพตราชา อตฺถีติ.\csromanspacer{} อาโป อตฺถีติ.\csromanspacer{} เตโช อตฺถีติ.\csromanspacer{} วาโย อตฺถีติ.\csromanspacer{} ติณกฏฺฐวนปฺปตโย อตฺถีติ ติณกฏฺฐวนปฺปตีนํ ปฏิสํเวที อตฺถีติ, น เหวํ วตฺตพฺเพ ฯเปฯ.}

\proseitem{192}{\csromansymbolmark{*}วิปาโก อตฺถีติ วิปากปฏิสํเวที อตฺถีติ, อามนฺตา.\csromanspacer{} อญฺโญ วิปาโก, อญฺโญ วิปากปฏิสํเวทีติ, น เหวํ วตฺตพฺเพ.\csromanspacer{}\csromanspacer{}(สํขิตฺตํ.)}

\vspace{\tipitakaparskip}
\csromancenter{กลฺยาณวคฺโค ปฐโม.}
\csromanmatikaanchor{mtk.15}
\csromantocmarkref{section}{14. อภิญฺญานุโยค}{mtk.15}
\bookmark[dest={mtk.15},level=1]{14. อภิญฺญานุโยค}
\csromanheader{1}{14. อภิญฺญานุโยค}
\proseitem{193}{\csromanfolio{50}\csromansymbolmark{+}น วตฺตพฺพํ “ปุคฺคโล อุปลพฺภติ สจฺจิกฏฺฐปรมตฺเถนา”ติ, อามนฺตา.\csromanspacer{} นนุ อตฺถิ โกจิ อิทฺธิํ วิกุพฺพตีติ, อามนฺตา.\csromanspacer{} หญฺจิ อตฺถิ โกจิ อิทฺธิํ วิกุพฺพติ, เตน วต เร วตฺตพฺเพ “ปุคฺคโล อุปลพฺภติ สจฺจิกฏฺฐปรมตฺเถนา”ติ.}

\proseitem{194}{\csromansymbolmark{+}น วตฺตพฺพํ “ปุคฺคโล อุปลพฺภติ สจฺจิกฏฺฐปรมตฺเถนา”ติ, อามนฺตา.\csromanspacer{} นนุ อตฺถิ โกจิ ทิพฺพาย โสตธาตุยา สทฺทํ สุณาติ ฯเปฯ ปรจิตฺตํ วิชานาติ.\csromanspacer{} ปุพฺเพนิวาสํ อนุสฺสรติ.\csromanspacer{} ทิพฺเพน จกฺขุนา รูปํ ปสฺสติ.\csromanspacer{} อาสวานํ ขยํ สจฺฉิกโรตีติ, อามนฺตา.\csromanspacer{} หญฺจิ อตฺถิ โกจิ อาสวานํ ขยํ สจฺฉิกโรติ, เตน วต เร วตฺตพฺเพ “ปุคฺคโล อุปลพฺภติ สจฺจิกฏฺฐปรมตฺเถนา”ติ.}

\proseitem{195}{\csromansymbolmark{*}อตฺถิ โกจิ อิทฺธิํ วิกุพฺพตีติ กตฺวา เตน จ การเณน ปุคฺคโล อุปลพฺภติ สจฺจิกฏฺฐปรมตฺเถนาติ, อามนฺตา.\csromanspacer{} โย อิทฺธิํ วิกุพฺพติ, เสฺวว ปุคฺคโล.\csromanspacer{} โย อิทฺธิํ น วิกุพฺพติ, น โส ปุคฺคโลติ, น เหวํ วตฺตพฺเพ ฯเปฯ.}

\proseitem{196}{\csromansymbolmark{*}โย ทิพฺพาย โสตธาตุยา สทฺทํ สุณาติ ฯเปฯ.\csromanspacer{} โย ปรจิตฺตํ วิชานาติ.\csromanspacer{} โย ปุพฺเพนิวาสํ อนุสฺสรติ.\csromanspacer{} โย ทิพฺเพน จกฺขุนา รูปํ ปสฺสติ.\csromanspacer{} โย อาสวานํ ขยํ สจฺฉิกโรติ, เสฺวว ปุคฺคโล.\csromanspacer{} โย อาสวานํ ขยํ น สจฺฉิกโรติ, น โส ปุคฺคโลติ, น เหวํ วตฺตพฺเพ ฯเปฯ.}

\vspace{\tipitakaparskip}
\csromancenter{อภิญฺญานุโยโค.}
\csromanmatikaanchor{mtk.16}
\csromantocmarkref{subsection}{15-18. ญาตกานุโยคาทิ}{mtk.16}
\bookmark[dest={mtk.16},level=2]{15-18. ญาตกานุโยคาทิ}
\csromanheader{2}{15-18. ญาตกานุโยคาทิ}
\proseitem{197}{\csromansymbolmark{+}น วตฺตพฺพํ “ปุคฺคโล อุปลพฺภติ สจฺจิกฏฺฐปรมตฺเถนา”ติ, อามนฺตา.\csromanspacer{} นนุ มาตา อตฺถีติ, อามนฺตา.\csromanspacer{} หญฺจิ มาตา อตฺถิ, เตน วต เร วตฺตพฺเพ “ปุคฺคโล อุปลพฺภติ สจฺจิกฏฺฐปรมตฺเถนา”ติ.}

\proseitem{198}{\csromansymbolmark{+}น วตฺตพฺพํ “ปุคฺคโล อุปลพฺภติ สจฺจิกฏฺฐปรมตฺเถนา”ติ, อามนฺตา.\csromanspacer{} นนุ ปิตา อตฺถิ ฯเปฯ ภาตา อตฺถิ.\csromanspacer{} ภคินี อตฺถิ.\csromanspacer{} ขตฺติโย \csromanfolio{51}อตฺถิ.\csromanspacer{} พฺราหฺมโณ อตฺถิ.\csromanspacer{} เวสฺโส อตฺถิ.\csromanspacer{} สุทฺโท อตฺถิ.\csromanspacer{} คหฏฺโฐ อตฺถิ.\csromanspacer{} ปพฺพชิโต อตฺถิ.\csromanspacer{} เทโว อตฺถิ.\csromanspacer{} มนุสฺโส อตฺถีติ, อามนฺตา.\csromanspacer{} หญฺจิ มนุสฺโส อตฺถิ, เตน วต เร วตฺตพฺเพ “ปุคฺคโล อุปลพฺภติ สจฺจิกฏฺฐปรมตฺเถนา”ติ.}

\proseitem{199}{\csromansymbolmark{*}มาตา-อตฺถีติ กตฺวา เตน จ การเณน ปุคฺคโล อุปลพฺภติ สจฺจิกฏฺฐปรมตฺเถนาติ, อามนฺตา.\csromanspacer{} อตฺถิ โกจิ น มาตา หุตฺวา มาตา โหตีติ, อามนฺตา.\csromanspacer{} อตฺถิ โกจิ น ปุคฺคโล หุตฺวา ปุคฺคโล โหตีติ, น เหวํ วตฺตพฺเพ ฯเปฯ.\csromanspacer{} อตฺถิ โกจิ น ปิตา หุตฺวา ฯเปฯ น ภาตา หุตฺวา.\csromanspacer{} น ภคินี หุตฺวา.\csromanspacer{} น ขตฺติโย หุตฺวา.\csromanspacer{} น พฺราหฺมโณ หุตฺวา.\csromanspacer{} น เวสฺโส หุตฺวา.\csromanspacer{} น สุทฺโท หุตฺวา.\csromanspacer{} น คหฏฺโฐ หุตฺวา.\csromanspacer{} น ปพฺพชิโต หุตฺวา.\csromanspacer{} น เทโว หุตฺวา.\csromanspacer{} น มนุสฺโส หุตฺวา มนุสฺโส โหตีติ, อามนฺตา.\csromanspacer{} อตฺถิ โกจิ น ปุคฺคโล หุตฺวา ปุคฺคโล โหตีติ, น เหวํ วตฺตพฺเพ ฯเปฯ.}

\proseitem{200}{\csromansymbolmark{*}มาตา อตฺถีติ กตฺวา เตน จ การเณน ปุคฺคโล อุปลพฺภติ สจฺจิกฏฺฐปรมตฺเถนาติ, อามนฺตา.\csromanspacer{} อตฺถิ โกจิ มาตา หุตฺวา น มาตา โหตีติ, อามนฺตา.\csromanspacer{} อตฺถิ โกจิ ปุคฺคโล หุตฺวา น ปุคฺคโล โหตีติ, น เหวํ วตฺตพฺเพ ฯเปฯ.}

\prose{อตฺถิ โกจิ ปิตา หุตฺวา.\csromanspacer{} ภาตา หุตฺวา.\csromanspacer{} ภคินี หุตฺวา.\csromanspacer{} ขตฺติโย หุตฺวา.\csromanspacer{} พฺราหฺมโณ หุตฺวา.\csromanspacer{} เวสฺโส หุตฺวา.\csromanspacer{} สุทฺโท หุตฺวา.\csromanspacer{} คหฏฺโฐ หุตฺวา.\csromanspacer{} ปพฺพชิโต หุตฺวา.\csromanspacer{} เทโว หุตฺวา.\csromanspacer{} มนุสฺโส หุตฺวา น มนุสฺโส โหตีติ, อามนฺตา.\csromanspacer{} อตฺถิ โกจิ ปุคฺคโล หุตฺวา น ปุคฺคโล โหตีติ, น เหวํ วตฺตพฺเพ ฯเปฯ.}

\csromanmatikaanchor{mtk.17}
\csromantocmarkref{subsection}{19. ปฏิเวธานุโยค}{mtk.17}
\bookmark[dest={mtk.17},level=2]{19. ปฏิเวธานุโยค}
\csromanheader{2}{19. ปฏิเวธานุโยค}
\proseitem{201}{\csromansymbolmark{+}น วตฺตพฺพํ “ปุคฺคโล อุปลพฺภติ สจฺจิกฏฺฐปรมตฺเถนา”ติ, อามนฺตา.\csromanspacer{} นนุ โสตาปนฺโน อตฺถีติ, อามนฺตา.\csromanspacer{} หญฺจิ โสตาปนฺโน อตฺถิ, เตน วต เร วตฺตพฺเพ “ปุคฺคโล อุปลพฺภติ สจฺจิกฏฺฐปรมตฺเถนา”ติ.}

\proseitem{202}{\csromansymbolmark{+}น วตฺตพฺพํ “ปุคฺคโล อุปลพฺภติ สจฺจิกฏฺฐปรมตฺเถนา”ติ, อามนฺตา.\csromanspacer{} นนุ สกทาคามี อตฺถิ ฯเปฯ.\csromanspacer{} อนาคามี อตฺถิ.\csromanspacer{} อรหา อตฺถิ.\csromanspacer{} อุภโตภาควิมุตฺโต อตฺถิ.\csromanspacer{} ปญฺญาวิมุตฺโต อตฺถิ.\csromanspacer{} กฺèยสกฺขิ\footnote{กายสกฺขี (สฺยา)} อตฺถิ. \csromanfolio{52}ทิฏฺฐิปฺปตฺโต อตฺถิ.\csromanspacer{} สทฺธาวิมุตฺโต อตฺถิ.\csromanspacer{} ธมฺมานุสารี อตฺถิ.\csromanspacer{} สทฺธานุสารี อตฺถีติ, อามนฺตา.\csromanspacer{} หญฺจิ สทฺธานุสารี อตฺถิ, เตน วต เร วตฺตพฺเพ “ปุคฺคโล อุปลพฺภติ สจฺจิกฏฺฐปรมตฺเถนา”ติ.}

\proseitem{203}{\csromansymbolmark{*}โสตาปนฺโน อตฺถีติ กตฺวา เตน จ การเณน ปุคฺคเล อุปลพฺภติ สจฺจิกฏฺฐปรมตฺเถนาติ, อามนฺตา.\csromanspacer{} อตฺถิ โกจิ น โสตาปนฺโน หุตฺวา โสตาปนฺโน โหตีติ, อามนฺตา.\csromanspacer{} อตฺถิ โกจิ น ปุคฺคโล หุตฺวา ปุคฺคโล โหตีติ, น เหวํ วตฺตพฺเพ ฯเปฯ.}

\proseitem{204}{\csromansymbolmark{*}อตฺถิ โกจิ น สกทาคามี หุตฺวา.\csromanspacer{} น อนาคามี หุตฺวา.\csromanspacer{} น อรหา หุตฺวา.\csromanspacer{} น อุภโตภาควิมุตฺโต หุตฺวา.\csromanspacer{} น ปญฺญาวิมุตฺโต หุตฺวา.\csromanspacer{} น กายสกฺขิ หุตฺวา.\csromanspacer{} น ทิฏฺฐิปฺปตฺโต หุตฺวา.\csromanspacer{} น สทฺธาวิมุตฺโต หุตฺวา.\csromanspacer{} น ธมฺมานุสารี หุตฺวา.\csromanspacer{} น สทฺธานุสารี หุตฺวา สทฺธานุสารี โหตีติ, อามนฺตา.\csromanspacer{} อตฺถิ โกจิ น ปุคฺคโล หุตฺวา ปุคฺคโล โหตีติ, น เหวํ วตฺตพฺเพ ฯเปฯ.}

\proseitem{205}{\csromansymbolmark{*}โสตาปนฺโน อตฺถีติ กตฺวา เตน จ การเณน ปุคฺคโล อุปลพฺภติ สจฺจิกฏฺฐปรมตฺเถนาติ, อามนฺตา.\csromanspacer{} อตฺถิ โกจิ โสตาปนฺโน หุตฺวา น โสตาปนฺโน โหตีติ, อามนฺตา.\csromanspacer{} อตฺถิ โกจิ ปุคฺคโล หุตฺวา น ปุคฺคโล โหตีติ, น เหวํ วตฺตพฺเพ ฯเปฯ.}

\prose{อตฺถิ โกจิ สกทาคามี หุตฺวา.\csromanspacer{} อนาคามี หุตฺวา น อนาคามี โหตีติ, อามนฺตา.\csromanspacer{} อตฺถิ โกจิ ปุคฺคโล หุตฺวา น ปุคฺคโล โหตีติ, น เหวํ วตฺตพฺเพ ฯเปฯ.}

\csromanmatikaanchor{mtk.18}
\csromantocmarkref{subsection}{20. สํฆานุโยค}{mtk.18}
\bookmark[dest={mtk.18},level=2]{20. สํฆานุโยค}
\csromanheader{2}{20. สํฆานุโยค}
\proseitem{206}{\csromansymbolmark{+}น วตฺตพฺพํ “ปุคฺคโล อุปลพฺภติ สจฺจิกฏฺฐปรมตฺเถนา”ติ, อามนฺตา.\csromanspacer{} นนุ จตฺตาโร ปุริสยุคา อฏฺฐ ปุริสปุคฺคลา อตฺถีติ, อามนฺตา.\csromanspacer{} หญฺจิ จตฺตาโร ปุริสยุคา อฏฺฐ ปุริสปุคฺคลา อตฺถิ, เตน วต เร วตฺตพฺเพ “ปุคฺคโล อุปลพฺภติ สจฺจิกฏฺฐปรมตฺเถนา”ติ.}

\proseitem{207}{\csromansymbolmark{*}จตฺตาโร ปุริสยุคา อฏฺฐ ปุริสปุคฺคลา อตฺถีติ กตฺวา เตน จ การเณน ปุคฺคโล อุปลพฺภติ สจฺจิกฏฺฐปรมตฺเถนาติ, อามนฺตา. \csromanfolio{53}จตฺตาโร ปุริสยุคา อฏฺฐ ปุริสปุคฺคลา พุทฺธปาตุภาวา ปาตุภวนฺตีติ, อามนฺตา.\csromanspacer{} ปุคฺคโล พุทฺธปาตุภาวา ปาตุภวตีติ, น เหวํ วตฺตพฺเพ ฯเปฯ. ปุคฺคโล พุทฺธปาตุภาวา ปาตุภวตีติ, อามนฺตา.\csromanspacer{} พุทฺธสฺส ภควโต ปรินิพฺพุเต อุจฺฉินฺโน ปุคฺคโล, นตฺถิ ปุคฺคโลติ, น เหวํ วตฺตพฺเพ ฯเปฯ.}

\csromanmatikaanchor{mtk.19}
\csromantocmarkref{subsection}{21. สจฺจิกฏฺฐสภาคานุโยค}{mtk.19}
\bookmark[dest={mtk.19},level=2]{21. สจฺจิกฏฺฐสภาคานุโยค}
\csromanheader{2}{21. สจฺจิกฏฺฐสภาคานุโยค}
\proseitem{208}{\csromansymbolmark{*}ปุคฺคโล อุปลพฺภติ สจฺจิกฏฺฐปรมตฺเถนาติ, อามนฺตา.\csromanspacer{} ปุคฺคโล สงฺขโตติ, น เหวํ วตฺตพฺเพ ฯเปฯ.\csromanspacer{} ปุคฺคโล อสงฺขโตติ, น เหวํ วตฺตพฺเพ ฯเปฯ.\csromanspacer{} ปุคฺคโล เนว สงฺขโต นาสงฺขโตติ, น เหวํ วตฺตพฺเพ.}

\proseitem{209}{\csromansymbolmark{*}ปุคฺคโล เนว สงฺขโต นาสงฺขโตติ, อามนฺตา.\csromanspacer{} สงฺขตญฺจ อสงฺขตญฺจ ฐเปตฺวา อตฺถญฺญา ตติยา โกฏีติ, น เหวํ วตฺตพฺเพ ฯเปฯ.}

\proseitem{210}{\csromansymbolmark{*}สงฺขตญฺจ อสงฺขตญฺจ ฐเปตฺวา อตฺถญฺญา ตติยา โกฏีติ, อามนฺตา.\csromanspacer{} นนุ วุตฺตํ ภควตา “เทฺวมา ภิกฺขเว ธาตุโย.\csromanspacer{} กตมา เทฺว, สงฺขตา จ ธาตุ อสงฺขตา จ ธาตุ.\csromanspacer{} อิมา โข ภิกฺขเว เทฺว ธาตุโย”ติ\footnote{ม 3. 108 ปิฏฺเฐ. อาลปนมตฺตเมว นานํ.} อตฺเถว สุตฺตนฺโตติ, อามนฺตา.\csromanspacer{} เตน หิ น วตฺตพฺพํ “สงฺขตญฺจ อสงฺขตญฺจ ฐเปตฺวา อตฺถญฺญา ตติยา โกฏี”ติ.}

\proseitem{211}{\csromansymbolmark{*}ปุคฺคโล เนว สงฺขโต นาสงฺขโตติ, อามนฺตา.\csromanspacer{} อญฺญํ สงฺขตํ, อญฺญํ อสงฺขตํ, อญฺโญ ปุคฺคโลติ, น เหวํ วตฺตพฺเพ ฯเปฯ.}

\proseitem{212}{\csromansymbolmark{*}ขนฺธา สงฺขตา, นิพฺพานํ อสงฺขตํ, ปุคฺคโล เนว สงฺขโต นาสงฺขโตติ, อามนฺตา.\csromanspacer{} อญฺเญ ขนฺธา, อญฺญํ นิพฺพานํ, อญฺโญ ปุคฺคโลติ, น เหวํ วตฺตพฺเพ ฯเปฯ.}

\proseitem{213}{\csromansymbolmark{*}รูปํ สงฺขตํ, นิพฺพานํ อสงฺขตํ, ปุคฺคโล เนว สงฺขโต นาสงฺขโตติ, อามนฺตา.\csromanspacer{} อญฺญํ รูปํ, อญฺญํ นิพฺพานํ, อญฺโญ ปุคฺคโลติ, น เหวํ วตฺตพฺเพ.\csromanspacer{} เวทนา.\csromanspacer{} สญฺญา.\csromanspacer{} สงฺขารา.\csromanspacer{} วิญฺญาณํ สงฺขตํ, นิพฺพานํ อสงฺขตํ ปุคฺคโล เนว สงฺขโต นาสงฺขโตติ, อามนฺตา.\csromanspacer{} อญฺญํ วิญฺญาณํ, อญฺญํ นิพฺพานํ, อญฺโญ ปุคฺคโลติ, น เหวํ วตฺตพฺเพ ฯเปฯ.}

\proseitem{214}{\csromanfolio{54}\csromansymbolmark{*}ปุคฺคลสฺส อุปฺปาโท ปญฺญายติ, วโย ปญฺญายติ, ฐิตสฺส อญฺญถตฺตํ ปญฺญายตีติ, อามนฺตา.\csromanspacer{} ปุคฺคโล สงฺขโตติ, น เหวํ วตฺตพฺเพ ฯเปฯ.\csromanspacer{} วุตฺตํ ภควตา “ตีณิมานิ ภิกฺขเว สงฺขตสฺส สงฺขตลกฺขณานิ, สงฺขตานํ ภิกฺขเว ธมฺมานํ\footnote{กตมานิ ตีณิ (อํ 1. 150)} อุปฺปาโท ปญฺญายติ, วโย ปญฺญายติ, ฐิตานํ อญฺญถตฺตํ ปญฺญายตี”ติ\footnote{อํ 1. 150 ปิฏฺเฐ.}, ปุคฺคลสฺส อุปฺปาโท ปญฺญายติ, วโย ปญฺญายติ, ฐิตสฺส อญฺญถตฺตํ ปญฺญายติ, เตน หิ ปุคฺคโล สงฺขโตติ.}

\proseitem{215}{\csromansymbolmark{*}ปุคฺคลสฺส น อุปฺปาโท ปญฺญายติ, น วโย ปญฺญายติ, น ฐิตสฺส อญฺญถตฺตํ ปญฺญายตีติ, อามนฺตา.\csromanspacer{} ปุคฺคโล อสงฺขโตติ, น เหวํ วตฺตพฺเพ ฯเปฯ.\csromanspacer{} วุตฺตํ ภควตา “ตีณิมานิ ภิกฺขเว อสงฺขตสฺส อสงฺขตลกฺขณานิ, อสงฺขตานํ ภิกฺขเว ธมฺมานํ น อุปฺปาโท ปญฺญายติ, น วโย ปญฺญายติ, น ฐิตานํ อญฺญถตฺตํ ปญฺญายตี”ติ\footnote{อํ 1. 150 ปิฏฺเฐ.}, ปุคฺคลสฺส น อุปฺปาโท ปญฺญายติ, น วโย ปญฺญายติ, น ฐิตสฺส อญฺญถตฺตํ ปญฺญายติ, เตน หิ ปุคฺคโล อสงฺขโตติ.}

\proseitem{216}{\csromansymbolmark{*}ปรินิพฺพุโต ปุคฺคโล อตฺถ’ตฺถมฺหิ, นตฺถ’ตฺถมฺหีติ, อตฺถ’ตฺถมฺหีติ.\csromanspacer{} ปรินิพฺพุโต ปุคฺคโล สสฺสโตติ, น เหวํ วตฺตพฺเพ ฯเปฯ นตฺถ’ตฺถมฺหีติ.\csromanspacer{} ปรินิพฺพุโต ปุคฺคโล อุจฺฉินฺโนติ, น เหวํ วตฺตพฺเพ ฯเปฯ.}

\proseitem{217}{\csromansymbolmark{*}ปุคฺคโล กิํ นิสฺสาย ติฏฺฐตีติ, ภวํ นิสฺสาย ติฏฺฐตีติ.\csromanspacer{} ภโว อนิจฺโจ สงฺขโต ปฏิจฺจสมุปฺปนฺโน ขยธมฺโม วยธมฺโม วิราคธมฺโม นิโรธธมฺโม วิปริณามธมฺโมติ, อามนฺตา.\csromanspacer{} ปุคฺคโลปิ อนิจฺโจ สงฺขโต ปฏิจฺจสมุปฺปนฺโน ขยธมฺโม วยธมฺโม วิราคธมฺโม นิโรธธมฺโม วิปริณามธมฺโมติ, น เหวํ วตฺตพฺเพ ฯเปฯ.}

\proseitem{218}{\csromansymbolmark{+}น วตฺตพฺพํ “ปุคฺคโล อุปลพฺภติ สจฺจิกฏฺฐปรมตฺเถนา”ติ, อามนฺตา.\csromanspacer{} นนุ อตฺถิ โกจิ สุขํ เวทนํ เวทิยมาโน “สุขํ เวทนํ เวทิยามี”ติ ปชานาตีติ, อามนฺตา.\csromanspacer{} หญฺจิ อตฺถิ โกจิ สุขํ เวทนํ เวทิยมาโน “สุขํ เวทนํ เวทิยามี”ติ ปชานาติ, เตน วต เร วตฺตพฺเพ “ปุคฺคโล อุปลพฺภติ สจฺจิกฏฺฐปรมตฺเถนา”ติ.}

\proseitem{219}{\csromanfolio{55}\csromansymbolmark{+}น วตฺตพฺพํ “ปุคฺคโล อุปลพฺภติ สจฺจิกฏฺฐปรมตฺเถนา”ติ, อามนฺตา.\csromanspacer{} นนุ อตฺถิ โกจิ ทุกฺขํ เวทนํ เวทิยมาโน ฯเปฯ อทุกฺขมสุขํ เวทนํ เวทิยมาโน “อทุกฺขมสุขํ เวทนํ เวทิยามี”ติ ปชานาตีติ, อามนฺตา.\csromanspacer{} หญฺจิ อตฺถิ โกจิ อทุกฺขมสุขํ เวทนํ เวทิยมาโน “อทุกฺขมสุขํ เวทนํ เวทิยามี”ติ ปชานาติ, เตน วต เร วตฺตพฺเพ “ปุคฺคโล อุปลพฺภติ สจฺจิกฏฺฐปรมตฺเถนา”ติ.}

\proseitem{220}{\csromansymbolmark{*}อตฺถิ โกจิ สุขํ เวทนํ เวทิยมาโน “สุขํ เวทนํ เวทิยามี”ติ ปชานาตีติ กตฺวา เตน จ การเณน ปุคฺคโล อุปลพฺภติ สจฺจิกฏฺฐปรมตฺเถนาติ, อามนฺตา.\csromanspacer{} โย สุขํ เวทนํ เวทิยมาโน “สุขํ เวทนํ เวทิยามี”ติ ปชานาติ, เสฺวว ปุคฺคโล.\csromanspacer{} โย สุขํ เวทนํ เวทิยมาโน “สุขํ เวทนํ เวทิยามี”ติ น ปชานาติ, น โส ปุคฺคโลติ, น เหวํ วตฺตพฺเพ ฯเปฯ.}

\prose{โย ทุกฺขํ เวทนํ เวทิยมาโน ฯเปฯ.\csromanspacer{} โย อทุกฺขมสุขํ เวทนํ เวทิยมาโน “อทุกฺขมสุขํ เวทนํ เวทิยามี”ติ ปชานาติ, เสฺวว ปุคฺคโล.\csromanspacer{} โย อทุกฺขมสุขํ เวทนํ เวทิยมาโน “อทุกฺขมสุขํ เวทนํ เวทิยามี”ติ น ปชานาติ, น โส ปุคฺคโลติ, น เหวํ วตฺตพฺเพ ฯเปฯ.}

\proseitem{221}{\csromansymbolmark{*}อตฺถิ โกจิ สุขํ เวทนํ เวทิยมาโน “สุขํ เวทนํ เวทิยามี”ติ ปชานาตีติ กตฺวา เตน จ การเณน ปุคฺคโล อุปลพฺภติ สจฺจิกฏฺฐปรมตฺเถนาติ, อามนฺตา.\csromanspacer{} อญฺญา สุขา เวทนา, อญฺโญ สุขํ เวทนํ เวทิยมาโน “สุขํ เวทนํ เวทิยามี”ติ ปชานาตีติ, น เหวํ วตฺตพฺเพ ฯเปฯ.\csromanspacer{} อญฺญา ทุกฺขา เวทนา ฯเปฯ.\csromanspacer{} อญฺญา อทุกฺขมสุขา เวทนา, อญฺโญ อทุกฺขมสุขํ เวทนํ เวทิยมาโน “อทุกฺขมสุขํ เวทนํ เวทิยามี”ติ ปชานาตีติ, น เหวํ วตฺตพฺเพ ฯเปฯ.}

\proseitem{222}{\csromansymbolmark{+}น วตฺตพฺพํ “ปุคฺคโล อุปลพฺภติ สจฺจิกฏฺฐปรมตฺเถนา”ติ, อามนฺตา.\csromanspacer{} นนุ อตฺถิ โกจิ กาเย กายานุปสฺสี วิหรตีติ, อามนฺตา.\csromanspacer{} หญฺจิ อตฺถิ โกจิ กาเย กายานุปสฺสี วิหรติ, เตน วต เร วตฺตพฺเพ “ปุคฺคโล อุปลพฺภติ สจฺจิกฏฺฐปรมตฺเถนา”ติ.}

\proseitem{223}{\csromansymbolmark{+}น วตฺตพฺพํ “ปุคฺคโล อุปลพฺภติ สจฺจิกฏฺฐปรมตฺเถนา”ติ, อามนฺตา.\csromanspacer{} นนุ อตฺถิ โกจิ เวทนาสุ ฯเปฯ.\csromanspacer{} จิตฺเต.\csromanspacer{} ธมฺเมสุ ธมฺมานุปสฺสี \csromanfolio{56}วิหรตีติ, อามนฺตา.\csromanspacer{} หญฺจิ อตฺถิ โกจิ ธมฺเมสุ ธมฺมานุปสฺสี วิหรติ, เตน วต เร วตฺตพฺเพ “ปุคฺคโล อุปลพฺภติ สจฺจิกฏฺฐปรมตฺเถนา”ติ.}

\proseitem{224}{\csromansymbolmark{*}อตฺถิ โกจิ กาเย กายานุปสฺสี วิหรตีติ กตฺวา เตน จ การเณน ปุคฺคโล อุปลพฺภติ สจฺจิกฏฺฐปรมตฺเถนาติ, อามนฺตา.\csromanspacer{} โย กาเย กายานุปสฺสี วิหรติ, เสฺวว ปุคฺคโล.\csromanspacer{} โย น กาเย กายานุปสฺสี วิหรติ, น โส ปุคฺคโลติ, น เหวํ วตฺตพฺเพ ฯเปฯ.}

\prose{โย เวทนาสุ ฯเปฯ.\csromanspacer{} จิตฺเต.\csromanspacer{} ธมฺเมสุ ธมฺมานุปสฺสี วิหรติ, เสฺวว ปุคฺคโล.\csromanspacer{} โย น ธมฺเมสุ ธมฺมานุปสฺสี วิหรติ, น โส ปุคฺคโลติ, น เหวํ วตฺตพฺเพ ฯเปฯ.}

\proseitem{225}{\csromansymbolmark{*}อตฺถิ โกจิ กาเย กายานุปสฺสี วิหรตีติ กตฺวา เตน จ การเณน ปุคฺคโล อุปลพฺภติ สจฺจิกฏฺฐปรมตฺเถนาติ, อามนฺตา.\csromanspacer{} อญฺโญ กาโย, อญฺโญ กาเย กายานุปสฺสี วิหรตีติ, น เหวํ วตฺตพฺเพ ฯเปฯ.\csromanspacer{} อญฺญา เวทนา.\csromanspacer{} อญฺญํ จิตฺตํ.\csromanspacer{} อญฺเญ ธมฺมา, อญฺโญ ธมฺเมสุ ธมฺมานุปสฺสี วิหรตีติ, น เหวํ วตฺตพฺเพ ฯเปฯ.}

\proseitem{226}{\csromansymbolmark{*}ปุคฺคโล อุปลพฺภติ สจฺจิกฏฺฐปรมตฺเถนาติ, อามนฺตา.\csromanspacer{} นนุ วุตฺตํ ภควตา\csromandash{}}

\csromangathameasure{“สุญฺญโต โลกํ อเวกฺขสฺสุ, โมฆราช สทา สโต. \\ อตฺตานุทิฏฺฐิํ อูหจฺจ, เอวํ มจฺจุตโร สิยา. \\ เอวํ โลกํ อเวกฺขนฺตํ, มจฺจุราชา น ปสฺสตี”ติ.}
\csromangathagroup{\csromangathastanza{“สุญฺญโต โลกํ อเวกฺขสฺสุ, โมฆราช สทา สโต. \\ อตฺตานุทิฏฺฐิํ อูหจฺจ\footnote{โอหจฺจ (สฺยา), อุหจฺจ (ก)}, เอวํ มจฺจุตโร สิยา. \\ เอวํ โลกํ อเวกฺขนฺตํ, มจฺจุราชา น ปสฺสตี”ติ\footnote{โอหจฺจ (สฺยา), อุหจฺจ (ก)}.}}

\prose{อตฺเถว สุตฺตนฺโตติ, อามนฺตา.\csromanspacer{} เตน หิ น วตฺตพฺพํ “ปุคฺคโล อุปลพฺภติ สจฺจิกฏฺฐปรมตฺเถนา”ติ.}

\proseitem{227}{\csromansymbolmark{*}ปุคฺคโล อเวกฺขตีติ, อามนฺตา.\csromanspacer{} สห รูเปน อเวกฺขติ, วินา รูเปน อเวกฺขตีติ, สห รูเปน อเวกฺขตีติ.\csromanspacer{} ตํ ชีวํ ตํ สรีรนฺติ, น เหวํ วตฺตพฺเพ ฯเปฯ.\csromanspacer{} วินา รูเปน อเวกฺขตีติ.\csromanspacer{} อญฺญํ ชีวํ อญฺญํ สรีรนฺติ, น เหวํ วตฺตพฺเพ ฯเปฯ.}

\proseitem{228}{\csromanfolio{57}\csromansymbolmark{*}ปุคฺคโล อเวกฺขตีติ, อามนฺตา.\csromanspacer{} อพฺภนฺตรคโต อเวกฺขติ, พหิทฺธา นิกฺขมิตฺวา อเวกฺขตีติ, อพฺภนฺตรคโต อเวกฺขตีติ.\csromanspacer{} ตํ ชีวํ ตํ สรีรนฺติ, น เหวํ วตฺตพฺเพ ฯเปฯ พหิทฺธา นิกฺขมิตฺวา อเวกฺขตีติ.\csromanspacer{} อญฺญํ ชีวํ อญฺญํ สรีรนฺติ, น เหวํ วตฺตพฺเพ ฯเปฯ.}

\proseitem{229}{\csromansymbolmark{+}น วตฺตพฺพํ “ปุคฺคโล อุปลพฺภติ สจฺจิกฏฺฐปรมตฺเถนา”ติ, อามนฺตา.\csromanspacer{} นนุ ภควา สจฺจวาที กาลวาที ภูตวาที ตถวาที อวิตถวาที อนญฺญถวาทีติ, อามนฺตา.\csromanspacer{} วุตฺตํ ภควตา “อตฺถิ ปุคฺคโล อตฺตหิตาย ปฏิปนฺโน”ติ อตฺเถว สุตฺตนฺโตติ, อามนฺตา.\csromanspacer{} เตน หิ ปุคฺคโล อุปลพฺภติ สจฺจิกฏฺฐปรมตฺเถนาติ.}

\proseitem{230}{\csromansymbolmark{+}น วตฺตพฺพํ “ปุคฺคโล อุปลพฺภติ สจฺจิกฏฺฐปรมตฺเถนา”ติ, อามนฺตา.\csromanspacer{} นนุ ภควา สจฺจวาที กาลวาที ภูตวาที ตถวาที อวิตถวาที อนญฺญถวาทีติ, อามนฺตา.\csromanspacer{} วุตฺตํ ภควตา “เอกปุคฺคโล ภิกฺขเว โลเก อุปฺปชฺชมาโน อุปฺปชฺชติ พหุชนหิตาย พหุชนสุขาย โลกานุกมฺปาย อตฺถาย หิตาย สุขาย เทวมนุสฺสานนฺ”ติ\footnote{อํ 1. 21 ปิฏฺเฐ.} อตฺเถว สุตฺตนฺโตติ, อามนฺตา.\csromanspacer{} เตน หิ ปุคฺคโล อุปลพฺภติ สจฺจิกฏฺฐปรมตฺเถนาติ.}

\proseitem{231}{\csromansymbolmark{*}ปุคฺคโล อุปลพฺภติ สจฺจิกฏฺฐปรมตฺเถนาติ, อามนฺตา.\csromanspacer{} นนุ ภควา สจฺจวาที กาลวาที ภูตวาที ตถวาที อวิตถวาที อนญฺญถวาทีติ, อามนฺตา.\csromanspacer{} วุตฺตํ ภควตา “สพฺเพ ธมฺมา อนตฺตา”ติ\footnote{ม 1. 292 ปิฏฺเฐ; ขุ 1. 53 ธมฺมปเทปิ.} อตฺเถว สุตฺตนฺโตติ, อามนฺตา.\csromanspacer{} เตน หิ น วตฺตพฺพํ “ปุคฺคโล อุปลพฺภติ สจฺจิกฏฺฐปรมตฺเถนา”ติ.}

\proseitem{232}{\csromansymbolmark{*}ปุคฺคโล อุปลพฺภติ สจฺจิกฏฺฐปรมตฺเถนาติ, อามนฺตา.\csromanspacer{} นนุ ภควา สจฺจวาที กาลวาที ภูตวาที ตถวาที อวิตถวาที อนญฺญถวาทีติ, อามนฺตา.\csromanspacer{} วุตฺตํ ภควตา “ทุกฺขเมว อุปฺปชฺชมานํ อุปฺปชฺชติ, ทุกฺขเมว\footnote{ทุกฺขํ (สํ 1. 258 ปิฏฺเฐ.)} นิรุชฺฌมานํ นุรุชฺฌตีติ น กงฺขติ น วิจิกิจฺฉติ, อปรปฺปจฺจยญฺญาณเมวสฺส เอตฺถ โหติ.\csromanspacer{} เอตฺตาวตา โข กจฺจาน สมฺมาทิฏฺฐิ โหตี”ติ\footnote{สํ 1. 257-8 ปิฏฺเฐ.} อตฺเถว สุตฺตนฺโตติ, อามนฺตา.\csromanspacer{} เตน หิ น วตฺตพฺพํ “ปุคฺคโล อุปลพฺภติ สจฺจิกฏฺฐปรมตฺเถนา”ติ.}

\proseitem{233}{\csromanfolio{58}\csromansymbolmark{*}ปุคฺคโล อุปลพฺภติ สจฺจิกฏฺฐปรมตฺเถนาติ, อามนฺตา.\csromanspacer{} นนุ วชิรา ภิกฺขุนี มารํ ปาปิมนฺตํ เอตทโวจ\csromandash{}}

\csromangathameasure{“กินฺนุ สตฺโตติ ปจฺเจสิ, มาร ทิฏฺฐิคตํ นุ เต. \\ สุทฺธสงฺขารปุญฺโชยํ, นยิธ สตฺตุปลพฺภติ. \\ ยถา หิ องฺคสมฺภารา, โหติ สทฺโท ‘รโถ’อิติ. \\ เอวํ ขนฺเธสุ สนฺเตสุ, โหติ ‘สตฺโต’ติ สมฺมุติ. \\ ทุกฺขเมว หิ สมฺโภติ, ทุกฺขํ ติฏฺฐติ เวติ จ. \\ นาญฺญตฺร ทุกฺขา สมฺโภติ, นาญฺญํ ทุกฺขา นิรุชฺฌตี”ติ.}
\csromangathagroup{\csromangathastanza{“กินฺนุ สตฺโตติ ปจฺเจสิ, มาร ทิฏฺฐิคตํ นุ เต. \\ สุทฺธสงฺขารปุญฺโชยํ, นยิธ สตฺตุปลพฺภติ.}\csromangathastanza{ยถา หิ\footnote{ยถาปิ (พหูสุ)} องฺคสมฺภารา, โหติ สทฺโท ‘รโถ’อิติ. \\ เอวํ ขนฺเธสุ สนฺเตสุ, โหติ ‘สตฺโต’ติ สมฺมุติ\footnote{ยถาปิ (พหูสุ)}.}\csromangathastanza{ทุกฺขเมว หิ สมฺโภติ, ทุกฺขํ ติฏฺฐติ เวติ จ. \\ นาญฺญตฺร ทุกฺขา สมฺโภติ, นาญฺญํ ทุกฺขา นิรุชฺฌตี”ติ\footnote{สํ 1. 137 ปิฏฺเฐ.}.}}

\prose{อตฺเถว สุตฺตนฺโตติ, อามนฺตา.\csromanspacer{} เตน หิ น วตฺตพฺพํ “ปุคฺคโล อุปลพฺภติ สจฺจิกฏฺฐปรมตฺเถนา”ติ.}

\proseitem{234}{\csromansymbolmark{*}ปุคฺคโล อุปลพฺภติ สจฺจิกฏฺฐปรมตฺเถนาติ, อามนฺตา.\csromanspacer{} นนุ อายสฺมา อานนฺโท ภควนฺตํ เอตทโวจ\csromandash{}สุญฺโญ โลโก สุญฺโญ โลโกติ ภนฺเต วุจฺจติ, กิตฺตาวตา นุ โข ภนฺเต “สุญฺโญ โลโก”ติ วุจฺจตีติ.\csromanspacer{} ยสฺมา โข อานนฺท สุญฺญํ อตฺเตน วา อตฺตนิเยน วา, ตสฺมา “สุญฺโญ โลโก”ติ วุจฺจติ.\csromanspacer{} กิญฺจานนฺท สุญฺญํ อตฺเตน วา อตฺตนิเยน วา, จกฺขุํ โข อานนฺท สุญฺญํ อตฺเตน วา อตฺตนิเยน วา, รูปา สุญฺญา ฯเปฯ จกฺขุวิญฺญาณํ สุญฺญํ.\csromanspacer{} จกฺขุสมฺผสฺโส สุญฺโญ.\csromanspacer{} ยมฺปิทํ จกฺขุสมฺผสฺสปจฺจยา อุปฺปชฺชติ เวทยิตํ สุขํ วา ทุกฺขํ วา อทุกฺขมสุขํ วา, ตมฺปิ สุญฺญํ อตฺเตน วา อตฺตนิเยน วา.\csromanspacer{} โสตํ สุญฺญํ ฯเปฯ.\csromanspacer{} สทฺทา สุญฺญา.\csromanspacer{} ฆานํ สุญฺญํ.\csromanspacer{} คนฺธา สุญฺญา.\csromanspacer{} ชิวฺหา สุญฺญา.\csromanspacer{} รสา สุญฺญา.\csromanspacer{} กาโย สุญฺโญ.\csromanspacer{} โผฏฺฐพฺพา สุญฺญา.\csromanspacer{} มโน สุญฺโญ.\csromanspacer{} ธมฺมา สุญฺญา.\csromanspacer{} มโนวิญฺญาณํ สุญฺญํ.\csromanspacer{} มโนสมฺผสฺโส สุญฺโญ.\csromanspacer{} ยมฺปิทํ มโนสมฺผสฺสปจฺจยา อุปฺปชฺชติ เวทยิตํ สุขํ วา ทุกฺขํ วา อทุกฺขมสุขํ วา, ตมฺปิ สุญฺญํ อตฺเตน วา อตฺตนิเยน วา.\csromanspacer{} ยสฺมา โข อานนฺท สุญฺญํ อตฺเตน วา อตฺตนิเยน วา, ตสฺมา “สุญฺโญ โลโก”ติ วุจฺจตีติ\footnote{สํ 2. 280 ปิฏฺเฐ.} อตฺเถว สุตฺตนฺโตติ, อามนฺตา.\csromanspacer{} เตน หิ น วตฺตพฺพํ “ปุคฺคโล อุปลพฺภติ สจฺจิกฏฺฐปรมตฺเถนา”ติ.}

\proseitem{235}{\csromanfolio{59}\csromansymbolmark{*}ปุคฺคโล อุปลพฺภติ สจฺจิกฏฺฐปรมตฺเถนาติ, อามนฺตา.\csromanspacer{} นนุ ภควา สจฺจวาที กาลวาที ภูตวาที ตถวาที อวิตถวาที อนญฺญถวาทีติ, อามนฺตา.\csromanspacer{} วุตฺตํ ภควตา “อตฺตนิ วา ภิกฺขเว สติ ‘อตฺตนิยํ เม’ติ อสฺสาติ, เอวํ ภนฺเต.\csromanspacer{} อตฺตนิเย วา ภิกฺขเว สติ ‘อตฺตา เม’ติ อสฺสาติ, เอวํ ภนฺเต.\csromanspacer{} อตฺตนิ จ ภิกฺขเว อตฺตนิเย จ สจฺจโต เถตโต อนุปลพฺภิยมาเน\footnote{อนุปลพฺภมาเน (ม 1. 191 ปิฏฺเฐ.)} ยมฺปิทํ\footnote{ยมิทํ (สฺยา) ยมฺปิตํ (ม 1. 191 ปิฏฺเฐ.)} ทิฏฺฐิฏฺฐานํ โส โลโก โส อตฺตา โส เปจฺจ ภวิสฺสามิ นิจฺโจ ธุโว สสฺสโต อวิปริณามธมฺโม สสฺสติสมํ ตเถว ฐสฺสามีติ, นนฺวายํ ภิกฺขเว เกวโล ปริปูโร พาลธมฺโมติ.\csromanspacer{} กิญฺหิ โน สิยา ภนฺเต, เกวโล หิ ภนฺเต ปริปูโร พาลธมฺโม”ติ\footnote{ม 1. 191 ปิฏฺเฐ.} อตฺเถว สุตฺตนฺโตติ, อามนฺตา.\csromanspacer{} เตน หิ น วตฺตพฺพํ “ปุคฺคโล อุปลพฺภติ สจฺจิกฏฺฐปรมตฺเถนา”ติ.}

\proseitem{236}{\csromansymbolmark{*}ปุคฺคโล อุปลพฺภติ สจฺจิกฏฺฐปรมตฺเถนาติ, อามนฺตา.\csromanspacer{} นนุ ภควา สจฺจวาที กาลวาที ภูตวาที ตถวาที อวิตถวาที อนญฺญถวาทีติ, อามนฺตา.\csromanspacer{} วุตฺตํ ภควตา “ตโยเม เสนิย สตฺถาโร สนฺโต สํวิชฺชมานา โลกสฺมิํ.\csromanspacer{} กตเม ตโย, อิธ เสนิย เอกจฺโจ สตฺถา ทิฏฺเฐว ธมฺเม อตฺตานํ สจฺจโต เถตโต ปญฺญาเปติ, อภิสมฺปรายญฺจ อตฺตานํ สจฺจโต เถตโต ปญฺญาเปติ.}

\prose{อิธ ปน เสนิย เอกจฺโจ สตฺถา ทิฏฺเฐว หิ โข ธมฺเม อตฺตานํ สจฺจโต เถตโต ปญฺญาเปติ, โน จ โข อภิสมฺปรายํ อตฺตานํ สจฺจโต เถตโต ปญฺญาเปติ.}

\prose{อิธ ปน เสนิย เอกจฺโจ สตฺถา ทิฏฺเฐ เจว ธมฺเม อตฺตานํ สจฺจโต เถตโต น ปญฺญาเปติ, อภิสมฺปรายญฺจ อตฺตานํ สจฺจโต เถตโต น ปญฺญาเปติ.}

\prose{ตตฺร เสนิย ยฺวายํ สตฺถา ทิฏฺเฐ เจว ธมฺเม อตฺตานํ สจฺจโต เถตโต ปญฺญาเปติ, อภิสมฺปรายญฺจ อตฺตานํ สจฺจโต เถตโต ปญฺญาเปติ, อยํ วุจฺจติ เสนิย สตฺถา สสฺสตวาโท.}

\prose{\csromanfolio{60}ตตฺร เสนิย ยฺวายํ สตฺถา ทิฏฺเฐว หิ โข ธมฺเม อตฺตานํ สจฺจโต เถตโต ปญฺญาเปติ, โน จ โข อภิสมฺปรายํ อตฺตานํ สจฺจโต เถตโต ปญฺญาเปติ, อยํ วุจฺจติ เสนิย สตฺถา อุจฺเฉทวาโท.}

\prose{ตตฺร เสนิย ยฺวายํ สตฺถา ทิฏฺเฐ เจว ธมฺเม อตฺตานํ สจฺจโต เถตโต น ปญฺญาเปติ, อภิสมฺปรายญฺจ อตฺตานํ สจฺจโต เถตโต น ปญฺญาเปติ, อยํ วุจฺจติ เสนิย สตฺถา สมฺมาสมฺพุทฺโธ.\csromanspacer{} อิเม โข เสนิย ตโย สตฺถาโร สนฺโต สํวิชฺชมานา โลกสฺมินฺติ\footnote{อภิ 3. 143 ปุคฺคลปญฺญตฺติยํ, อตฺถโต เอกํ.} อตฺเถว สุตฺตนฺโตติ, อามนฺตา.\csromanspacer{} เตน หิ น วตฺตพฺพํ “ปุคฺคโล อุปลพฺภติ สจฺจิกฏฺฐปรมตฺเถนา”ติ.}

\proseitem{237}{\csromansymbolmark{*}ปุคฺคโล อุปลพฺภติ สจฺจิกฏฺฐปรมตฺเถนาติ, อามนฺตา.\csromanspacer{} นนุ ภควา สจฺจวาที กาลวาที ภูตวาที ตถวาที อวิตถวาที อนญฺญถวาทีติ, อามนฺตา.\csromanspacer{} วุตฺตํ ภควตา “สปฺปิกุมฺโภ”ติ, อามนฺตา.\csromanspacer{} อตฺถิ โกจิ สปฺปิสฺส กุมฺภํ กโรตีติ, น เหวํ วตฺตพฺเพ ฯเปฯ.\csromanspacer{} เตน หิ น วตฺตพฺพํ “ปุคฺคโล อุปลพฺภติ สจฺจิกฏฺฐปรมตฺเถนา”ติ.}

\proseitem{238}{\csromansymbolmark{*}ปุคฺคโล อุปลพฺภติ สจฺจิกฏฺฐปรมตฺเถนาติ, อามนฺตา.\csromanspacer{} นนุ ภควา สจฺจวาที กาลวาที ภูตวาที ตถวาที อวิตถวาที อนญฺญถวาทีติ, อามนฺตา.\csromanspacer{} วุตฺตํ ภควตา “เตลกุมฺโภ.\csromanspacer{} มธุกุมฺโภ.\csromanspacer{} ผาณิตกุมฺโภ.\csromanspacer{} ขีรกุมฺโภ.\csromanspacer{} อุทกกุมฺโภ.\csromanspacer{} ปานียถาลกํ.\csromanspacer{} ปานียโกสกํ.\csromanspacer{} ปานียสราวกํ.\csromanspacer{} นิจฺจภตฺตํ.\csromanspacer{} ธุวยาคู”ติ, อามนฺตา.\csromanspacer{} อตฺถิ กาจิ ยาคุ นิจฺจา ธุวา สสฺสตา อวิปริณามธมฺมาติ, น เหวํ วตฺตพฺเพ ฯเปฯ.\csromanspacer{} เตน หิ น วตฺตพฺพํ “ปุคฺคโล อุปลพฺภติ สจฺจิกฏฺฐปรมตฺเถนา”ติ.\csromanspacer{}\csromanspacer{}(สํขิตฺตํ.)}

\csromangathameasure{อฏฺฐกนิคฺคหเปยฺยาลา, สนฺธาวนิยา อุปาทาย. \\ จิตฺเตน ปญฺจมํ กลฺยาณํ, อิทฺธิสุตฺตาหรเณน อฏฺฐมํ.}
\csromangathagroup{\csromangathastanza{อฏฺฐกนิคฺคหเปยฺยาลา, สนฺธาวนิยา อุปาทาย. \\ จิตฺเตน ปญฺจมํ กลฺยาณํ, อิทฺธิสุตฺตาหรเณน อฏฺฐมํ.}}

\nitthitam{\textbf{ปุคฺคลกถา นิฏฺฐิตา.}}
\csromanmatikaanchor{mtk.20}
\csromantocmarkref{chapter}{2. ปริหานิกถา}{mtk.20}
\bookmark[dest={mtk.20},level=0]{2. ปริหานิกถา}
\chapterheadpageread{2. ปริหานิกถา}
\csromanheader{2}{1. วาทยุตฺติปริหานิ}
\proseitem{239}{\csromanfolio{61}\csromansymbolmark{*}ปริหายติ อรหา อรหตฺตาติ, อามนฺตา.\csromanspacer{} สพฺพตฺถ อรหา อรหตฺตา ปริหายตีติ, น เหวํ วตฺตพฺเพ ฯเปฯ.\csromanspacer{} สพฺพตฺถ อรหา อรหตฺตา ปริหายตีติ, อามนฺตา.\csromanspacer{} สพฺพตฺถ อรหโต ปริหานีติ, น เหวํ วตฺตพฺเพ ฯเปฯ.}

\prose{\csromansymbolmark{*}ปริหายติ อรหา อรหตฺตาติ, อามนฺตา.\csromanspacer{} สพฺพทา อรหา อรหตฺตา ปริหายตีติ, น เหวํ วตฺตพฺเพ ฯเปฯ.\csromanspacer{} สพฺพทา อรหา อรหตฺตา ปริหายตีติ, อามนฺตา.\csromanspacer{} สพฺพทา อรหโต ปริหานีติ, น เหวํ วตฺตพฺเพ ฯเปฯ.}

\prose{\csromansymbolmark{*}ปริหายติ อรหา อรหตฺตาติ, อามนฺตา.\csromanspacer{} สพฺเพว อรหนฺโต อรหตฺตา ปริหายนฺตีติ, น เหวํ วตฺตพฺเพ ฯเปฯ.\csromanspacer{} สพฺเพว อรหนฺโต อรหตฺตา ปริหายนฺตีติ, อามนฺตา.\csromanspacer{} สพฺเพสํเยว อรหนฺตานํ ปริหานีติ, น เหวํ วตฺตพฺเพ ฯเปฯ.}

\prose{\csromansymbolmark{*}ปริหายติ อรหา อรหตฺตาติ, อามนฺตา.\csromanspacer{} อรหา อรหตฺตา ปริหายมาโน จตูหิ ผเลหิ ปริหายตีติ, น เหวํ วตฺตพฺเพ ฯเปฯ.}

\prose{\csromansymbolmark{+}จตูหิ สตสหสฺเสหิ เสฏฺฐี เสฏฺฐิตฺตํ กาเรนฺโต สตสหสฺเส ปริหีเน เสฏฺฐิตฺตา ปริหีโน โหตีติ, อามนฺตา.\csromanspacer{} สพฺพสาปเตยฺยา ปริหีโน โหตีติ, น เหวํ วตฺตพฺเพ ฯเปฯ.}

\prose{\csromansymbolmark{*}จตูหิ สตสหสฺเสหิ เสฏฺฐี เสฏฺฐิตฺตํ กาเรนฺโต สตสหสฺเส ปริหีเน ภพฺโพ สพฺพสาปเตยฺยา ปริหายิตุนฺติ, อามนฺตา.\csromanspacer{} อรหา อรหตฺตา ปริหายมาโน ภพฺโพ จตูหิ ผเลหิ ปริหายิตุนฺติ, น เหวํ วตฺตพฺเพ ฯเปฯ.}

\csromanheader{2}{2. อริยปุคฺคลสํสนฺทนปริหานิ}
\proseitem{240}{\csromansymbolmark{*}ปริหายติ อรหา อรหตฺตาติ, อามนฺตา.\csromanspacer{} ปริหายติ โสตาปนฺโน โสตาปตฺติผลาติ, น เหวํ วตฺตพฺเพ ฯเปฯ.}

\prose{\csromanfolio{62}\csromansymbolmark{*}ปริหายติ อรหา อรหตฺตาติ, อามนฺตา.\csromanspacer{} ปริหายติ สกทาคามี สกทาคามิผลาติ, น เหวํ วตฺตพฺเพ ฯเปฯ.}

\prose{\csromansymbolmark{*}ปริหายติ อรหา อรหตฺตาติ, อามนฺตา.\csromanspacer{} ปริหายติ อนาคามี อนาคามิผลาติ, น เหวํ วตฺตพฺเพ ฯเปฯ.}

\prose{\csromansymbolmark{*}ปริหายติ อนาคามี อนาคามิผลาติ, อามนฺตา.\csromanspacer{} ปริหายติ โสตาปนฺโน โสตาปตฺติผลาติ, น เหวํ วตฺตพฺเพ ฯเปฯ.}

\prose{\csromansymbolmark{*}ปริหายติ อนาคามี อนาคามิผลาติ, อามนฺตา.\csromanspacer{} ปริหายติ สกทาคามี สกทาคามิผลาติ, น เหวํ วตฺตพฺเพ ฯเปฯ.}

\prose{\csromansymbolmark{*}ปริหายติ สกทาคามี สกทาคามิผลาติ, อามนฺตา.\csromanspacer{} ปริหายติ โสตาปนฺโน โสตาปตฺติผลาติ, น เหวํ วตฺตพฺเพ ฯเปฯ.}

\prose{\csromansymbolmark{*}น ปริหายติ โสตาปนฺโน โสตาปตฺติผลาติ, อามนฺตา.\csromanspacer{} น ปริหายติ อรหา อรหตฺตาติ, น เหวํ วตฺตพฺเพ ฯเปฯ.}

\prose{\csromansymbolmark{*}น ปริหายติ สกทาคามี สกทาคามิผลาติ, อามนฺตา.\csromanspacer{} น ปริหายติ อรหา อรหตฺตาติ, น เหวํ วตฺตพฺเพ ฯเปฯ.}

\prose{\csromansymbolmark{*}น ปริหายติ อนาคามี อนาคามิผลาติ, อามนฺตา.\csromanspacer{} น ปริหายติ อรหา อรหตฺตาติ, น เหวํ วตฺตพฺเพ ฯเปฯ.}

\prose{\csromansymbolmark{*}น ปริหายติ โสตาปนฺโน โสตาปตฺติผลาติ, อามนฺตา.\csromanspacer{} น ปริหายติ อนาคามี อนาคามิผลาติ, น เหวํ วตฺตพฺเพ ฯเปฯ.}

\prose{\csromansymbolmark{*}น ปริหายติ สกทาคามี สกทาคามิผลาติ, อามนฺตา.\csromanspacer{} น ปริหายติ อนาคามี อนาคามิผลาติ, น เหวํ วตฺตพฺเพ ฯเปฯ.}

\prose{\csromansymbolmark{*}น ปริหายติ โสตาปนฺโน โสตาปตฺติผลาติ, อามนฺตา.\csromanspacer{} น ปริหายติ สกทาคามี สกทาคามิผลาติ, น เหวํ วตฺตพฺเพ ฯเปฯ.}

\proseitem{241}{\csromansymbolmark{*}ปริหายติ อรหา อรหตฺตาติ, อามนฺตา.\csromanspacer{} อรหา อรหตฺตา ปริหายมาโน กตฺถ สณฺฐาตีติ, อนาคามิผเลติ.\csromanspacer{} อนาคามี อนาคามิผลา ปริหายมาโน กตฺถ สณฺฐาตีติ, สกทาคามิผเลติ.\csromanspacer{} สกทาคามี สกทาคามิผลา ปริหายมาโน กตฺถ สณฺฐาตีติ, โสตาปตฺติผเลติ.\csromanspacer{} โสตาปนฺโน โสตาปตฺติผลา ปริหายมาโน ปุถุชฺชนภูมิยํ สณฺฐาตีติ, น เหวํ วตฺตพฺเพ.}

\csromanheader{2}{2. ปริหานิกถา}
\prose{\csromanfolio{63}อาชานาหิ นิคฺคหํ\csromandash{}หญฺจิ อรหา อรหตฺตา ปริหายมาโน อนาคามิผเล สณฺฐาติ, อนาคามี อนาคามิผลา ปริหายมาโน สกทาคามิผเล สณฺฐาติ, สกทาคามี สกทาคามิผลา ปริหายมาโน โสตาปตฺติผเล สณฺฐาติ, เตน วต เร วตฺตพฺเพ “โสตาปนฺโน โสตาปตฺติผลา ปริหายมาโน ปุถุชฺชนภูมิยํ สณฺฐาตี”ติ.}

\prose{\csromansymbolmark{*}อรหา อรหตฺตา ปริหายมาโน โสตาปตฺติผเล สณฺฐาตีติ, อามนฺตา.\csromanspacer{} โสตาปตฺติผลสฺส อนนฺตรา อรหตฺตํเยว สจฺฉิกโรตีติ, น เหวํ วตฺตพฺเพ ฯเปฯ.}

\proseitem{242}{\csromansymbolmark{*}ปริหายติ อรหา อรหตฺตาติ, อามนฺตา.\csromanspacer{} ปริหายติ โสตาปนฺโน โสตาปตฺติผลาติ, น เหวํ วตฺตพฺเพ ฯเปฯ.\csromanspacer{} กสฺส พหุตรา กิเลสา ปหีนา อรหโต วา โสตาปนฺนสฺส วาติ, อรหโต.\csromanspacer{} หญฺจิ อรหโต พหุตรา กิเลสา ปหีนา, ปริหายติ อรหา อรหตฺตา, เตน วต เร วตฺตพฺเพ “ปริหายติ โสตาปนฺโน โสตาปตฺติผลา”ติ.}

\prose{\csromansymbolmark{*}ปริหายติ อรหา อรหตฺตาติ, อามนฺตา.\csromanspacer{} ปริหายติ สกทาคามี สกทาคามิผลาติ, น เหวํ วตฺตพฺเพ ฯเปฯ.\csromanspacer{} กสฺส พหุตรา กิเลสา ปหีนา อรหโต วา สกทาคามิสฺส วาติ, อรหโต.\csromanspacer{} หญฺจิ อรหโต พหุตรา กิเลสา ปหีนา, ปริหายติ อรหา อรหตฺตา, เตน วต เร วตฺตพฺเพ “ปริหายติ สกทาคามี สกทาคามิผลา”ติ.}

\prose{\csromansymbolmark{*}ปริหายติ อรหา อรหตฺตาติ, อามนฺตา.\csromanspacer{} ปริหายติ อนาคามี อนาคามิผลาติ, น เหวํ วตฺตพฺเพ ฯเปฯ.\csromanspacer{} กสฺส พหุตรา กิเลสา ปหีนา อรหโต วา อนาคามิสฺส วาติ, อรหโต.\csromanspacer{} หญฺจิ อรหโต พหุตรา กิเลสา ปหีนา, ปริหายติ อรหา อรหตฺตา, เตน วต เร วตฺตพฺเพ “ปริหายติ อนาคามี อนาคามิผลา”ติ.}

\prose{\csromansymbolmark{*}ปริหายติ อนาคามี อนาคามิผลาติ, อามนฺตา.\csromanspacer{} ปริหายติ โสตาปนฺโน โสตาปตฺติผลาติ, น เหวํ วตฺตพฺเพ ฯเปฯ.\csromanspacer{} กสฺส พหุตรา กิเลสา ปหีนา อนาคามิสฺส วา โสตาปนฺนสฺส วาติ, อนาคามิสฺส.\csromanspacer{} หญฺจิ อนาคามิสฺส พหุตรา กิเลสา ปหีนา, ปริหายติ อนาคามี อนาคามิผลา, เตน วต เร วตฺตพฺเพ “ปริหายติ โสตาปนฺโน โสตาปตฺติผลา”ติ.}

\proseitem{243}{\csromanfolio{64}\csromansymbolmark{*}ปริหายติ อนาคามี อนาคามิผลาติ, อามนฺตา.\csromanspacer{} ปริหายติ สกทาคามี สกทาคามิผลาติ, น เหวํ วตฺตพฺเพ ฯเปฯ.\csromanspacer{} กสฺส พหุตรา กิเลสา ปหีนา อนาคามิสฺส วา สกทาคามิสฺส วาติ, อนาคามิสฺส.\csromanspacer{} หญฺจิ อนาคามิสฺส พหุตรา กิเลสา ปหีนา, ปริหายติ อนาคามี อนาคามิผลา, เตน วต เร วตฺตพฺเพ “ปริหายติ สกทาคามี สกทาคามิผลา”ติ.}

\prose{\csromansymbolmark{*}ปริหายติ สกทาคามี สกทาคามิผลาติ, อามนฺตา.\csromanspacer{} ปริหายติ โสตาปนฺโน โสตาปตฺติผลาติ, น เหวํ วตฺตพฺเพ ฯเปฯ.\csromanspacer{} กสฺส พหุตรา กิเลสา ปหีนา สกทาคามิสฺส วา โสตาปนฺนสฺส วาติ, สกทาคามิสฺส.\csromanspacer{} หญฺจิ สกทาคามิสฺส พหุตรา กิเลสา ปหีนา, ปริหายติ สกทาคามี สกทาคามิผลา, เตน วต เร วตฺตพฺเพ “ปริหายติ โสตาปนฺโน โสตาปตฺติผลา”ติ.}

\proseitem{244}{\csromansymbolmark{*}ปริหายติ อรหา อรหตฺตาติ, อามนฺตา.\csromanspacer{} ปริหายติ โสตาปนฺโน โสตาปตฺติผลาติ, น เหวํ วตฺตพฺเพ ฯเปฯ.\csromanspacer{} กสฺส อธิมตฺตา มคฺคภาวนา อรหโต วา โสตาปนฺนสฺส วาติ, อรหโต.\csromanspacer{} หญฺจิ อรหโต อธิมตฺตา มคฺคภาวนา, ปริหายติ อรหา อรหตฺตา, เตน วต เร วตฺตพฺเพ “ปริหายติ โสตาปนฺโน โสตาปตฺติผลา”ติ.}

\prose{\csromansymbolmark{*}ปริหายติ อรหา อรหตฺตาติ, อามนฺตา.\csromanspacer{} ปริหายติ โสตาปนฺโน โสตาปตฺติผลาติ, น เหวํ วตฺตพฺเพ ฯเปฯ.\csromanspacer{} กสฺส อธิมตฺตา สติปฏฺฐานภาวนา ฯเปฯ สมฺมปฺปธานภาวนา.\csromanspacer{} อิทฺธิปาทภาวนา.\csromanspacer{} อินฺทฺริยภาวนา.\csromanspacer{} พลภาวนา.\csromanspacer{} โพชฺฌงฺคภาวนา อรหโต วา โสตาปนฺนสฺส วาติ, อรหโต.\csromanspacer{} หญฺจิ อรหโต อธิมตฺตา โพชฺฌงฺคภาวนา, ปริหายติ อรหา อรหตฺตา, เตน วต เร วตฺตพฺเพ “ปริหายติ โสตาปนฺโน โสตาปตฺติผลา”ติ.}

\prose{\csromansymbolmark{*}ปริหายติ อรหา อรหตฺตาติ, อามนฺตา.\csromanspacer{} ปริหายติ สกทาคามี สกทาคามิผลาติ, น เหวํ วตฺตพฺเพ ฯเปฯ.\csromanspacer{} กสฺส อธิมตฺตา มคฺคภาวนา ฯเปฯ โพชฺฌงฺคภาวนา อรหโต วา สกทาคามิสฺส วาติ, อรหโต.\csromanspacer{} หญฺจิ อรหโต อธิมตฺตา โพชฺฌงฺคภาวนา, ปริหายติ อรหา อรหตฺตา, เตน วต เร วตฺตพฺเพ “ปริหายติ สกทาคามี สกทาคามิผลา”ติ.}

\csromanheader{2}{2. ปริหานิกถา}
\prose{\csromanfolio{65}\csromansymbolmark{*}ปริหายติ อรหา อรหตฺตาติ, อามนฺตา.\csromanspacer{} ปริหายติ อนาคามี อนาคามิผลาติ, น เหวํ วตฺตพฺเพ ฯเปฯ.\csromanspacer{} กสฺส อธิมตฺตา มคฺคภาวนา ฯเปฯ โพชฺฌงฺคภาวนา อรหโต วา อนาคามิสฺส วาติ, อรหโต.\csromanspacer{} หญฺจิ อรหโต อธิมตฺตา โพชฺฌงฺคภาวนา, ปริหายติ อรหา อรหตฺตา, เตน วต เร วตฺตพฺเพ “ปริหายติ อนาคามี อนาคามิผลา”ติ.}

\proseitem{245}{\csromansymbolmark{*}ปริหายติ อนาคามี อนาคามิผลาติ, อามนฺตา.\csromanspacer{} ปริหายติ โสตาปนฺโน โสตาปตฺติผลาติ, น เหวํ วตฺตพฺเพ ฯเปฯ.\csromanspacer{} กสฺส อธิมตฺตา มคฺคภาวนา ฯเปฯ โพชฺฌงฺคภาวนา อนาคามิสฺส วา โสตาปนฺนสฺส วาติ, อนาคามิสฺส.\csromanspacer{} หญฺจิ อนาคามิสฺส อธิมตฺตา โพชฺฌงฺคภาวนา, ปริหายติ อนาคามี อนาคามิผลา, เตน วต เร วตฺตพฺเพ “ปริหายติ โสตาปนฺโน โสตาปตฺติผลา”ติ.}

\prose{\csromansymbolmark{*}ปริหายติ อนาคามี อนาคามิผลาติ, อามนฺตา.\csromanspacer{} ปริหายติ สกทาคามี สกทาคามิผลาติ, น เหวํ วตฺตพฺเพ ฯเปฯ.\csromanspacer{} กสฺส อธิมตฺตา มคฺคภาวนา ฯเปฯ โพชฺฌงฺคภาวนา อนาคามิสฺส วา สกทาคามิสฺส วาติ, อนาคามิสฺส.\csromanspacer{} หญฺจิ อนาคามิสฺส อธิมตฺตา โพชฺฌงฺคภาวนา, ปริหายติ อนาคามี อนาคามิผลา, เตน วต เร วตฺตพฺเพ “ปริหายติ สกทาคามี สกทาคามิผลา”ติ.}

\proseitem{246}{\csromansymbolmark{*}ปริหายติ สกทาคามี สกทาคามิผลาติ, อามนฺตา.\csromanspacer{} ปริหายติ โสตาปนฺโน โสตาปตฺติผลาติ.\csromanspacer{} น เหวํ วตฺตพฺเพ ฯเปฯ.\csromanspacer{} กสฺส อธิมตฺตา มคฺคภาวนา ฯเปฯ โพชฺฌงฺคภาวนา สกทาคามิสฺส วา โสตาปนฺนสฺส วาติ, สกทาคามิสฺส.\csromanspacer{} หญฺจิ สกทาคามิสฺส อธิมตฺตา โหชฺฌงฺคภาวนา, ปริหายติ สกทาคามี สกทาคามิผลา, เตน เวต เร วตฺตพฺเพ “ปริหายติ โสตาปนฺโน โสตาปตฺติผลา”ติ ฯเปฯ.}

\proseitem{247}{\csromansymbolmark{*}อรหตา ทุกฺขํ ทิฏฺฐํ, ปริหายติ อรหา อรหตฺตาติ, อามนฺตา.\csromanspacer{} โสตาปนฺเนน ทุกฺขํ ทิฏฺฐํ, ปริหายติ โสตาปนฺโน โสตาปตฺติผลาติ, น เหวํ วตฺตพฺเพ ฯเปฯ.}

\prose{\csromansymbolmark{*}อรหตา สมุทโย ทิฏฺโฐ, ปริหายติ อรหา อรหตฺตาติ, อามนฺตา.\csromanspacer{} โสตาปนฺเนน สมุทโย ทิฏฺโฐ, ปริหายติ โสตาปนฺโน โสตาปตฺติผลาติ, น เหวํ วตฺตพฺเพ ฯเปฯ.}

\prose{\csromanfolio{66}\csromansymbolmark{*}อรหตา นิโรโธ ทิฏฺโฐ, ปริหายติ อรหา อรหตฺตาติ, อามนฺตา.\csromanspacer{} โสตาปนฺเนน นิโรโธ ทิฏฺโฐ, ปริหายติ โสตาปนฺโน โสตาปตฺติผลาติ, น เหวํ วตฺตพฺเพ ฯเปฯ.}

\prose{\csromansymbolmark{*}อรหตา มคฺโค ทิฏฺโฐ, ปริหายติ อรหา อรหตฺตาติ, อามนฺตา.\csromanspacer{} โสตาปนฺเนน มคฺโค ทิฏฺโฐ, ปริหายติ โสตาปนฺโน โสตาปตฺติผลาติ, น เหวํ วตฺตพฺเพ ฯเปฯ.}

\prose{\csromansymbolmark{*}อรหตา จตฺตาริ สจฺจานิ ทิฏฺฐานิ, ปริหายติ อรหา อรหตฺตาติ, อามนฺตา.\csromanspacer{} โสตาปนฺเนน จตฺตาริ สจฺจานิ ทิฏฺฐานิ, ปริหายติ โสตาปนฺโน โสตาปตฺติผลาติ, น เหวํ วตฺตพฺเพ ฯเปฯ.}

\prose{\csromansymbolmark{*}อรหตา ทุกฺขํ ทิฏฺฐํ, ปริหายติ อรหา อรหตฺตาติ, อามนฺตา.\csromanspacer{} สกทาคามินา ทุกฺขํ ทิฏฺฐํ, ปริหายติ สกทาคามี สกทาคามิผลาติ, น เหวํ วตฺตพฺเพ ฯเปฯ.}

\prose{\csromansymbolmark{*}อรหตา สมุทโย ทิฏฺโฐ ฯเปฯ นิโรโธ ทิฏฺโฐ ฯเปฯ มคฺโค ทิฏฺโฐ ฯเปฯ จตฺตาริ สจฺจานิ ทิฏฺฐานิ, ปริหายติ อรหา อรหตฺตาติ, อามนฺตา.\csromanspacer{} สกทาคามินา จตฺตาริ สจฺจานิ ทิฏฺฐานิ, ปริหายติ สกทาคามี สกทาคามิผลาติ, น เหวํ วตฺตพฺเพ ฯเปฯ.}

\prose{\csromansymbolmark{*}อรหตา ทุกฺขํ ทิฏฺฐํ, ปริหายติ อรหา อรหตฺตาติ, อามนฺตา.\csromanspacer{} อนาคามินา ทุกฺขํ ทิฏฺฐํ, ปริหายติ อนาคามี อนาคามิผลาติ, น เหวํ วตฺตพฺเพ ฯเปฯ.}

\prose{\csromansymbolmark{*}อรหตา สมุทโย ทิฏฺโฐ ฯเปฯ นิโรโธ ทิฏฺโฐ ฯเปฯ มคฺโค ทิฏฺโฐ ฯเปฯ จตฺตาริ สจฺจานิ ทิฏฺฐานิ, ปริหายติ อรหา อรหตฺตาติ, อามนฺตา.\csromanspacer{} อนาคามินา จตฺตาริ สจฺจานิ ทิฏฺฐานิ, ปริหายติ อนาคามี อนาคามิผลาติ, น เหวํ วตฺตพฺเพ ฯเปฯ.}

\proseitem{248}{\csromansymbolmark{*}อนาคามินา ทุกฺขํ ทิฏฺฐํ, ปริหายติ อนาคามี อนาคามิผลาติ, อามนฺตา.\csromanspacer{} โสตาปนฺเนน ทุกฺขํ ทิฏฺฐํ, ปริหายติ โสตาปนฺโน โสตาปตฺติผลาติ, น เหวํ วตฺตพฺเพ ฯเปฯ.}

\prose{\csromansymbolmark{*}อนาคามินา สมุทโย ทิฏฺโฐ ฯเปฯ นิโรโธ ทิฏฺโฐ ฯเปฯ มคฺโค ทิฏฺโฐ ฯเปฯ จตฺตาริ สจฺจานิ ทิฏฺฐานิ, ปริหายติ อนาคามี อนาคามิผลาติ,}

\csromanheader{2}{2. ปริหานิกถา}
\prose{\csromanfolio{67}อามนฺตา.\csromanspacer{} โสตาปนฺเนน จตฺตาริ สจฺจานิ ทิฏฺฐานิ, ปริหายติ โสตาปนฺโน โสตาปตฺติผลาติ, น เหวํ วตฺตพฺเพ ฯเปฯ.}

\prose{\csromansymbolmark{*}อนาคามินา ทุกฺขํ ทิฏฺฐํ, ปริหายติ อนาคามี อนาคามิผลาติ, อามนฺตา.\csromanspacer{} สกทาคามินา ทุกฺขํ ทิฏฺฐํ, ปริหายติ สกทาคามี สกทาคามิผลาติ, น เหวํ วตฺตพฺเพ ฯเปฯ.}

\prose{\csromansymbolmark{*}อนาคามินา สมุทโย ทิฏฺโฐ ฯเปฯ นิโรโธ ทิฏฺโฐ ฯเปฯ มคฺโค ทิฏฺโฐ ฯเปฯ จตฺตาริ สจฺจานิ ทิฏฺฐานิ, ปริหายติ อนาคามี อนาคามิผลาติ, อามนฺตา.\csromanspacer{} สกทาคามินา จตฺตาริ สจฺจานิ ทิฏฺฐานิ, ปริหายติ สกทาคามี สกทาคามิผลาติ, น เหวํ วตฺตพฺเพ ฯเปฯ.}

\proseitem{249}{\csromansymbolmark{*}สกทาคามินา ทุกฺขํ ทิฏฺฐํ, ปริหายติ สกทาคามี สกทาคามิผลาติ, อามนฺตา.\csromanspacer{} โสตาปนฺเนน ทุกฺขํ ทิฏฺฐํ, ปริหายติ โสตาปนฺโน โสตาปตฺติผลาติ, น เหวํ วตฺตพฺเพ ฯเปฯ.}

\prose{\csromansymbolmark{*}สกทาคามินา สมุทโย ทิฏฺโฐ ฯเปฯ นิโรโธ ทิฏฺโฐ ฯเปฯ มคฺโค ทิฏฺโฐ ฯเปฯจตฺตาริ สจฺจานิ ทิฏฺฐานิ, ปริหายติ สกทาคามี สกทาคามิผลาติ, อามนฺตา.\csromanspacer{} โสตาปนฺเนน จตฺตริ สจฺจานิ ทิฏฺฐานิ, ปริหายติ โสตาปนฺโน โสตาปตฺติผลาติ.\csromanspacer{} น เหวํ วตฺตพฺเพ ฯเปฯ.}

\proseitem{250}{\csromansymbolmark{*}โสตาปนฺเนน ทุกฺขํ ทิฏฺฐํ, น ปริหายติ โสตาปนฺโน โสตาปตฺติผลาติ, อามนฺตา.\csromanspacer{} อรหตา ทุกฺขํ ทิฏฺฐํ, น ปริหายติ อรหา อรหตฺตาติ, น เหวํ วตฺตพฺเพ ฯเปฯ.}

\prose{\csromansymbolmark{*}โสตาปนฺเนน สมุทโย ทิฏฺโฐ ฯเปฯ นิโรโธ ทิฏฺโฐ ฯเปฯ มคฺโค ทิฏฺโฐ ฯเปฯ จตฺตาริ สจฺจานิ ทิฏฺฐานิ, น ปริหายติ โสตาปนฺโน โสตาปตฺติผลาติ, อามนฺตา.\csromanspacer{} อรหตา จตฺตาริ สจฺจานิ ทิฏฺฐานิ, น ปริหายติ อรหา อรหตฺตาติ, น เหวํ วตฺตพฺเพ ฯเปฯ.}

\prose{\csromansymbolmark{*}สกทาคามินา ทุกฺขํ ทิฏฺฐํ ฯเปฯ จตฺตาริ สจฺจานิ ทิฏฺฐานิ, น ปริหายติ สกทาคามี สกทาคามิผลาติ, อามนฺตา.\csromanspacer{} อรหตา จตฺตาริ สจฺจานิ ทิฏฺฐานิ, น ปริหายติ อรหา อรหตฺตาติ, น เหวํ วตฺตพฺเพ ฯเปฯ.}

\prose{\csromansymbolmark{*}อนาคามินา ทุกฺขํ ทิฏฺฐํ ฯเปฯ จตฺตาริ สจฺจานิ ทิฏฺฐานิ, น ปริหายติ อนาคามี อนาคามิผลาติ, อามนฺตา.\csromanspacer{} อรหตา จตฺตาริ สจฺจานิ ทิฏฺฐานิ, น ปริหายติ อรหา อรหตฺตาติ, น เหวํ วตฺตพฺเพ ฯเปฯ.}

\prose{\csromanfolio{68}\csromansymbolmark{*}โสตาปนฺเนน ทุกฺขํ ทิฏฺฐํ ฯเปฯ จตฺตาริ สจฺจานิ ทิฏฺฐานิ, น ปริหายติ โสตาปนฺโน โสตาปตฺติผลาติ, อามนฺตา.\csromanspacer{} อนาคามินา ทุกฺขํ ทิฏฺฐํ ฯเปฯ จตฺตาริ สจฺจานิ ทิฏฺฐานิ, น ปริหายติ อนาคามี อนาคามิผลาติ, น เหวํ วตฺตพฺเพ ฯเปฯ.}

\prose{\csromansymbolmark{*}สกทาคามินา ทุกฺขํ ทิฏฺฐํ ฯเปฯ จตฺตาริ สจฺจานิ ทิฏฺฐานิ, น ปริหายติ สกทาคามี สกทาคามิผลาติ, อามนฺตา.\csromanspacer{} อนาคามินา ทุกฺขํ ทิฏฺฐํ ฯเปฯ จตฺตาริ สจฺจานิ ทิฏฺฐานิ, น ปริหายติ อนาคามี อนาคามิผลาติ, น เหวํ วตฺตพฺเพ ฯเปฯ.}

\prose{\csromansymbolmark{*}โสตาปนฺเนน ทุกฺขํ ทิฏฺฐํ ฯเปฯ จตฺตาริ สจฺจานิ ทิฏฺฐานิ, น ปริหายติ, โสตาปนฺโน โสตาปตฺติผลาติ, อามนฺตา.\csromanspacer{} สกทาคามินา ทุกฺขํ ทิฏฺฐํ ฯเปฯ จตฺตาริ สจฺจานิ ทิฏฺฐานิ, น ปริหายติ สกทาคามี สกทาคามิผลาติ, น เหวํ วตฺตพฺเพ ฯเปฯ.}

\proseitem{251}{\csromansymbolmark{*}อรหโต ราโค ปหีโน, ปริหายติ อรหา อรหตฺตาติ, อามนฺตา.\csromanspacer{} โสตาปนฺนสฺส สกฺกายทิฏฺฐิ ปหีนา, ปริหายติ โสตาปนฺโน โสตาปตฺติผลาติ, น เหวํ วตฺตพฺเพ ฯเปฯ.}

\prose{\csromansymbolmark{*}อรหโต ราโค ปหีโน ปริหายติ อรหา อรหตฺตาติ, อามนฺตา.\csromanspacer{} โสตาปนฺนสฺส วิจิกิจฺฉา ปหีนา ฯเปฯ สีลพฺพตปรามาโส ปหีโน ฯเปฯ อปายคมนีโย ราโค ปหีโน ฯเปฯ อปายคมนีโย โทโส ปหีโน ฯเปฯ อปายคมนีโย โมโห ปหีโน, ปริหายติ โสตาปนฺโน โสตาปตฺติผลาติ, น เหวํ วตฺตพฺเพ ฯเปฯ.}

\prose{\csromansymbolmark{*}อรหโต โทโส ปหีโน ฯเปฯ โมโห ปหีโน.\csromanspacer{} มาโน ปหีโน.\csromanspacer{} ทิฏฺฐิ ปหีนา.\csromanspacer{} วิจิกิจฺฉา ปหีนา.\csromanspacer{} ถินํ\footnote{ถีนํ (สี, สฺยา, กํ, อิ)} ปหีนํ.\csromanspacer{} อุทฺธจฺจํ ปหีนํ.\csromanspacer{} อหิริกํ ปหีนํ ฯเปฯ.\csromanspacer{} อโนตฺตปฺปํ ปหีนํ, ปริหายติ อรหา อรหตฺตาติ, อามนฺตา.\csromanspacer{} โสตาปนฺนสฺส สกฺกายทิฏฺฐิ ปหีนา, ปริหายติ โสตาปนฺโน โสตาปตฺติผลาติ, น เหวํ วตฺตพฺเพ ฯเปฯ.}

\prose{\csromansymbolmark{*}อรหโต อโนตฺตปฺปํ ปหีนํ, ปริหายติ อรหา อรหตฺตาติ, อามนฺตา.\csromanspacer{} โสตาปนฺนสฺส วิจิกิจฺฉา ปหีนา ฯเปฯ สีลพฺพตปรามาโส ปหีโน ฯเปฯ อปายคมนีโย ราโค ปหีโน ฯเปฯ อปายคมนีโย}

\csromanheader{2}{2. ปริหานิกถา}
\prose{\csromanfolio{69}โทโส ปหีโน ฯเปฯ อปายคมนีโย โมโห ปหีโน ปริหายติ โสตาปนฺโน โสตาปตฺติผลาติ, น เหวํ วตฺตพฺเพ ฯเปฯ.}

\prose{\csromansymbolmark{*}อรหโต ราโค ปหีโน ปริหายติ, อรหา อรหตฺตาติ, อามนฺตา.\csromanspacer{} สกทาคามิสฺส สกฺกายทิฏฺฐิ ปหีนา, ปริหายติ สกทาคามี สกทาคามิผลาติ, น เหวํ วตฺตพฺเพ.}

\prose{\csromansymbolmark{*}อรหโต ราโค ปหีโน, ปริหายติ อรหา อรหตฺตาติ, อามนฺตา.\csromanspacer{} สกทาคามิสฺส วิจิกิจฺฉา ปหีนา ฯเปฯ สีลพฺพตปรามาโส ปหีโน ฯเปฯ โอฬาริโก กามราโค ปหีโน.\csromanspacer{} โอฬาริโก พฺยาปาโท ปหีโน, ปริหายติ สกทาคามี สกทาคามิผลาติ, น เหวํ วตฺตพฺเพ ฯเปฯ.}

\prose{\csromansymbolmark{*}อรหโต โทโส ปหีโน ฯเปฯ อโนตฺตปฺปํ ปหีนํ, ปริหายติ อรหา อรหตฺตาติ, อามนฺตา.\csromanspacer{} สกทาคามิสฺส สกฺกายทิฏฺฐิ ปหีนา ฯเปฯ โอฬาริโก พฺยาปาโท ปหีโน, ปริหายติ สกทาคามี สกทาคามิผลาติ, น เหวํ วตฺตพฺเพ.}

\prose{\csromansymbolmark{*}อรหโต ราโค ปหีโน, ปริหายติ อรหา อรหตฺตาติ, อามนฺตา.\csromanspacer{} อนาคามิสฺส สกฺกายทิฏฺฐิ ปหีนา, ปริหายติ อนาคามี อนาคามิผลาติ, น เหวํ วตฺตพฺเพ ฯเปฯ.}

\prose{\csromansymbolmark{*}อรหโต ราโค ปหีโน, ปริหายติ อรหา อรหตฺตาติ, อามนฺตา.\csromanspacer{} อนาคามิสฺส วิจิกิจฺฉา ปหีนา ฯเปฯ สีลพฺพตปรามาโส ปหีโน.\csromanspacer{} อณุสหคโต กามราโค ปหีโน.\csromanspacer{} อณุสหคโต พฺยาปาโท ปหีโน, ปริหายติ อนาคามี อนาคามิผลาติ, น เหวํ วตฺตพฺเพ.}

\prose{\csromansymbolmark{*}อรหโต โทโส ปหีโน ฯเปฯ อโนตฺตปฺปํ ปหีนํ, ปริหายติ อรหา อรหตฺตาติ, อามนฺตา.\csromanspacer{} อนาคามิสฺส สกฺกายทิฏฺฐิ ปหีนา ฯเปฯ อณุสหคโต พฺยาปาโท ปหีโน, ปริหายติ อนาคามี อนาคามิผลาติ, น เหวํ วตฺตพฺเพ.}

\proseitem{252}{\csromansymbolmark{*}อนาคามิสฺส สกฺกายทิฏฺฐิ ปหีนา, ปริหายติ อนาคามี อนาคามิผลาติ, อามนฺตา.\csromanspacer{} โสตาปนฺนสฺส สกฺกายทิฏฺฐิ ปหีนา, ปริหายติ โสตาปนฺโน โสตาปตฺติผลาติ, น เหวํ วตฺตพฺเพ.}

\prose{\csromanfolio{70}\csromansymbolmark{*}อนาคามิสฺส สกฺกยทิฏฺฐิ ปหีนา, ปริหายติ อนาคามี อนาคามิผลาติ, อามนฺตา.\csromanspacer{} โสตาปนฺนสฺส วิจิกิจฺฉา ปหีนา ฯเปฯ อปายคมนีโย โมโห ปหีโน, ปริหายติ โสตาปนฺโน โสตาปตฺติผลาติ, น เหวํ วตฺตพฺเพ.}

\prose{\csromansymbolmark{*}อนาคามิสฺส วิจิกิจฺฉา ปหีนา ฯเปฯ อณุสหคโต พฺยาปาโท ปหีโน, ปริหายติ อนาคามี อนาคามิผลาติ, อามนฺตา.\csromanspacer{} โสตาปนฺนสฺส สกฺกายทิฏฺฐิ ปหีนา ฯเปฯ.\csromanspacer{} อปายคมนีโย โมโห ปหีโน, ปริหายติ โสตาปนฺโน โสตาปตฺติผลาติ, น เหวํ วตฺตพฺเพ.}

\prose{\csromansymbolmark{*}อนาคามิสฺส สกฺกายทิฏฺฐิ ปหีนา, ปริหายติ อนาคามี อนาคามิผลาติ, อามนฺตา.\csromanspacer{} สกทาคามิสฺส สกฺกายทิฏฺฐิ ปหีนา, ปริหายติ สกทาคามี สกทาคามิผลาติ, น เหวํ วตฺตพฺเพ.}

\prose{\csromansymbolmark{*}อนาคามิสฺส สกฺกายทิฏฺฐิ ปหีนา, ปริหายติ อนาคามี อนาคามิผลาติ, อามนฺตา.\csromanspacer{} สกทาคามิสฺส วิจิกิจฺฉา ปหีนา ฯเปฯ.\csromanspacer{} สีลพฺพตปรามาโส ปหีโน.\csromanspacer{} โอฬาริโก กามราโค ปหีโน.\csromanspacer{} โอฬาริโก พฺยาปาโท ปหีโน, ปริหายติ สกทาคามี สกทาคามิผลาติ, น เหวํ วตฺตพฺเพ.}

\prose{\csromansymbolmark{*}อนาคามิสฺส วิจิกิจฺฉา ปหีนา ฯเปฯ.\csromanspacer{} อณุสหคโต พฺยาปาโท ปหีโน, ปริหายติ อนาคามี อนาคามิผลาติ, อามนฺตา.\csromanspacer{} สกทาคามิสฺส สกฺกายทิฏฺฐิ ปหีนา ฯเปฯ.\csromanspacer{} โอฬาริโก พฺยาปาโท ปหีโน, ปริหายติ สกทาคามี สกทาคามิผลาติ, น เหวํ วตฺตพฺเพ.}

\proseitem{253}{\csromansymbolmark{*}สกทาคามิสฺส สกฺกายทิฏฺฐิ ปหีนา, ปริหายติ สกทาคามี สกทาคามิผลาติ, อามนฺตา.\csromanspacer{} โสตาปนฺนสฺส สกฺกายทิฏฺฐิ ปหีนา, ปริหายติ โสตาปนฺโน โสตาปตฺติผลาติ, น เหวํ วตฺตพฺเพ.}

\prose{\csromansymbolmark{*}สกทาคามิสฺส สกฺกายทิฏฺฐิ ปหีนา, ปริหายติ สกทาคามี สกทาคามิผลาติ, อามนฺตา.\csromanspacer{} โสตาปนฺนสฺส วิจิกิจฺฉา ปหีนา ฯเปฯ.\csromanspacer{} อปายคมนีโย โมโห ปหีโน, ปริหายติ โสตาปนฺโน โสตาปตฺติผลาติ, น เหวํ วตฺตพฺเพ.}

\prose{\csromansymbolmark{*}สกทาคามิสฺส วิจิกิจฺฉา ปหีนา ฯเปฯ.\csromanspacer{} โอฬาริโก กามราโค ปหีโน.\csromanspacer{} โอฬาริโก พฺยาปาโท ปหีโน, ปริหายติ สกทาคามี}

\csromanheader{2}{2. ปริหานิกถา}
\prose{\csromanfolio{71}สกทาคามิผลาติ, อามนฺตา.\csromanspacer{} โสตาปนฺนสฺส สกฺกายทิฏฺฐิ ปหีนา ฯเปฯ.\csromanspacer{} อปายคมนีโย โมโห ปหีโน, ปริหายติ โสตาปนฺโน โสตาปตฺติผลาติ, น เหวํ วตฺตพฺเพ.}

\proseitem{254}{\csromansymbolmark{*}โสตาปนฺนสฺส สกฺกายทิฏฺฐิ ปหีนา, น ปริหายติ โสตาปนฺโน โสตาปตฺติผลาติ, อามนฺตา.\csromanspacer{} อรหโต ราโค ปหีโน, น ปริหายติ อรหา อรหตฺตาติ, น เหวํ วตฺตพฺเพ.}

\prose{\csromansymbolmark{*}โสตาปนฺนสฺส สกฺกายทิฏฺฐิ ปหีนา, น ปริหายติ โสตาปนฺโน โสตาปตฺติผลาติ, อามนฺตา.\csromanspacer{} อรหโต โทโส ปหีโน ฯเปฯ.\csromanspacer{} อโนตฺตปฺปํ ปหีนํ, น ปริหายติ อรหา อรหตฺตาติ, น เหวํ วตฺตพฺเพ.}

\prose{\csromansymbolmark{*}โสตาปนฺนสฺส วิจิกิจฺฉา ปหีนา ฯเปฯ.\csromanspacer{} อปายคมนีโย โมโห ปหีโน, น ปริหายติ โสตาปนฺโน โสตาปตฺติผลาติ, อามนฺตา.\csromanspacer{} อรหโต ราโค ปหีโน ฯเปฯ.\csromanspacer{} อโนตฺตปฺปํ ปหีนํ, น ปริหายติ อรหา อรหตฺตาติ, น เหวํ วตฺตพฺเพ.}

\proseitem{255}{\csromansymbolmark{*}สกทาคามิสฺส สกฺกายทิฏฺฐิ ปหีนา, น ปริหายติ สกทาคามี สกทาคามิผลาติ, อามนฺตา.\csromanspacer{} อรหโต ราโค ปหีโน ฯเปฯ.\csromanspacer{} อโนตฺตปฺปํ ปหีนํ, น ปริหายติ อรหา อรหตฺตาติ, น เหวํ วตฺตพฺเพ.}

\prose{\csromansymbolmark{*}สกทาคามิสฺส วิจิกิจฺฉา ปหีนา ฯเปฯ.\csromanspacer{} โอฬาริโก พฺยาปาโท ปหีโน, น ปริหายติ สกทาคามี สกทาคามิผลาติ, อามนฺตา.\csromanspacer{} อรหโต ราโค ปหีโน ฯเปฯ.\csromanspacer{} อโนตฺตปฺปํ ปหีนํ, น ปริหายติ อรหา อรหตฺตาติ, น เหวํ วตฺตพฺเพ.}

\proseitem{256}{\csromansymbolmark{*}อนาคามิสฺส สกฺกายทิฏฺฐิ ปหีนา, น ปริหายติ อนาคามี อนาคามิผลาติ, อามนฺตา.\csromanspacer{} อรหโต ราโค ปหีโน ฯเปฯ อโนตฺตปฺปํ ปหีนํ, น ปริหายติ อรหา อรหตฺตาติ, น เหวํ วตฺตพฺเพ.}

\prose{\csromansymbolmark{*}อนาคามิสฺส วิจิกิจฺฉา ปหีนา ฯเปฯ อณุสหคโต พฺยาปาโท ปหีโน, น ปริหายติ อนาคามี อนาคามิผลาติ, อามนฺตา.\csromanspacer{} อรหโต ราโค ปหีโน ฯเปฯ อโนตฺตปฺปํ ปหีนํ, น ปริหายติ อรหา อรหตฺตาติ, น เหวํ วตฺตพฺเพ.}

\proseitem{257}{\csromanfolio{72}\csromansymbolmark{*}โสตาปนฺนสฺส สกฺกายทิฏฺฐิ ปหีนา, น ปริหายติ โสตาปนฺโน โสตาปตฺติผลาติ, อามนฺตา.\csromanspacer{} อนาคามิสฺส สกฺกายทิฏฺฐิ ปหีนา.\csromanspacer{} อณุสหคโต พฺยาปาโท ปหีโน, น ปริหายติ อนาคามี อนาคามิผลาติ, น เหวํ วตฺตพฺเพ.}

\prose{\csromansymbolmark{*}โสตาปนฺนสฺส วิจิกิจฺฉา ปหีนา ฯเปฯ อปายคมนีโย โมโห ปหีโน, น ปริหายติ โสตาปนฺโน โสตาปตฺติผลาติ, อามนฺตา.\csromanspacer{} อนาคามิสฺส สกฺกายทิฏฺฐิ ปหีนา ฯเปฯ อณุสหคโต พฺยาปาโท ปหีโน, น ปริหายติ อนาคามี อนาคามิผลาติ, น เหวํ วตฺตพฺเพ.}

\proseitem{258}{\csromansymbolmark{*}สกทาคามิสฺส สกฺกายทิฏฺฐิ ปหีนา, น ปริหายติ สกทาคามี สกทาคามิผลาติ, อามนฺตา.\csromanspacer{} อนาคามิสฺส สกฺกายทิฏฺฐิ ปหีนา ฯเปฯ อณุสหคโต พฺยาปาโท ปหีโน, น ปริหายติ อนาคามี อนาคามิผลาติ, น เหวํ วตฺตพฺเพ.}

\prose{\csromansymbolmark{*}สกทาคามิสฺส วิจิกิจฺฉา ปหีนา ฯเปฯ โอฬาริโก พฺยาปาโท ปหีโน, น ปริหายติ สกทาคามี สกทาคามิผลาติ, อามนฺตา.\csromanspacer{} อนาคามิสฺส สกฺกายทิฏฺฐิ ปหีนา ฯเปฯ อณุสหคโต พฺยาปาโท ปหีโน, น ปริหายติ อนาคามี อนาคามิผลาติ, น เหวํ วตฺตพฺเพ.}

\proseitem{259}{\csromansymbolmark{*}โสตาปนฺนสฺส สกฺกายทิฏฺฐิ ปหีนา, น ปริหายติ โสตาปนฺโน โสตาปตฺติผลาติ, อามนฺตา.\csromanspacer{} สกทาคามิสฺส สกฺกายทิฏฺฐิ ปหีนา ฯเปฯ โอฬาริโก พฺยาปาโท ปหีโน, น ปริหายติ สกทาคามี สกทาคามิผลาติ, น เหวํ วตฺตพฺเพ.}

\prose{\csromansymbolmark{*}โสตาปนฺนสฺส วิจิกิจฺฉา ปหีนา ฯเปฯ อปายคมนีโย โมโห ปหีโน, น ปริหายติ โสตาปนฺโน โสตาปตฺติผลาติ, อามนฺตา.\csromanspacer{} สกทาคามิสฺส สกฺกายทิฏฺฐิ ปหีนา ฯเปฯ โอฬาริโก พฺยาปาโท ปหีโน, น ปริหายติ สกทาคามี สกทาคามิผลาติ, น เหวํ วตฺตพฺเพ.}

\proseitem{260}{\csromansymbolmark{*}ปริหายติ อรหา อรหตฺตาติ, อามนฺตา.\csromanspacer{} นนุ อรหโต ราโค ปหีโน อุจฺฉินฺนมูโล ตาลาวตฺถุกโต อนภาวํกโต\footnote{อนภาวกโต (สี)} อายติํ อนุปฺปาทธมฺโมติ, อามนฺตา.\csromanspacer{} หญฺจิ อรหโต ราโค ปหีโน}

\csromanheader{2}{2. ปริหานิกถา}
\prose{\csromanfolio{73}อุจฺฉินฺนมูโล ตาลาวตฺถุกโต อนภาวํกโต อายติํ อนุปฺปาทธมฺโม, โน จ วต เร\footnote{โน วต เร (สฺยา, อิ) เอวมุปริปิ.} วตฺตพฺเพ “ปริหายติ อรหา อรหตฺตา”ติ.}

\prose{\csromansymbolmark{*}ปริหายติ อรหา อรหตฺตาติ, อามนฺตา.\csromanspacer{} นนุ อรหโต โทโส ปหีโน ฯเปฯ โมโห ปหีโน.\csromanspacer{} มาโน ปหีโน.\csromanspacer{} ทิฏฺฐิ ปหีนา.\csromanspacer{} วิจิกิจฺฉา ปหีนา.\csromanspacer{} ถินํ ปหีนํ.\csromanspacer{} อุทฺธจฺจํ ปหีนํ.\csromanspacer{} อหิริกํ ปหีนํ.\csromanspacer{} อโนตฺตปฺปํ ปหีนํ อุจฺฉินฺนมูลํ ตาลาวตฺถุกตํ อนภาวํกตํ อายติํ อนุปฺปาทธมฺมนฺติ, อามนฺตา.\csromanspacer{} หญฺจิ อรหโต อโนตฺตปฺปํ ปหีนํ อุจฺฉินฺนมูลํ ตาลาวตฺถุกตํ อนภาวํกตํ อายติํ อนุปฺปาทธมฺมํ, โน จ วต เร วตฺตพฺเพ “ปริหายติ อรหา อรหตฺตา”ติ.}

\prose{\csromansymbolmark{*}ปริหายติ อรหา อรหตฺตาติ, อามนฺตา.\csromanspacer{} นนุ อรหโต ราคปฺปหานาย มคฺโค ภาวิโตติ, อามนฺตา.\csromanspacer{} หญฺจิ อรหโต ราคปฺปหานาย มคฺโค ภาวิโต, โน จ วต เร วตฺตพฺเพ “ปริหายติ อรหา อรหตฺตา”ติ.}

\prose{\csromansymbolmark{*}ปริหายติ อรหา อรหตฺตาติ, อามนฺตา.\csromanspacer{} นนุ อรหโต ราคปฺปหานาย สติปฏฺฐานา ภาวิตา ฯเปฯ สมฺมปฺปธานา ภาวิตา.\csromanspacer{} อิทฺธิปาทา ภาวิตา.\csromanspacer{} อินฺทฺริยา ภาวิตา.\csromanspacer{} พลา ภาวิตา.\csromanspacer{} โพชฺฌงฺคา ภาวิตาติ, อามนฺตา.\csromanspacer{} หญฺจิ อรหโต ราคปฺปหานาย โพชฺฌงฺคา ภาวิตา, โน จ วต เร วตฺตพฺเพ “ปริหายติ อรหา อรหตฺตา”ติ.}

\prose{\csromansymbolmark{*}ปริหายติ อรหา อรหตฺตาติ, อามนฺตา.\csromanspacer{} นนุ อรหโต โทสปฺปหานาย ฯเปฯ อโนตฺตปฺปปหานาย มคฺโค ภาวิโต ฯเปฯ โพชฺฌงฺคา ภาวิตาติ, อามนฺตา.\csromanspacer{} หญฺจิ อรหโต อโนตฺตปฺปปหานาย โพชฺฌงฺคา ภาวิตา, โน จ วต เร วตฺตพฺเพ “ปริหายติ อรหา อรหตฺตา”ติ.}

\proseitem{261}{\csromansymbolmark{*}ปริหายติ อรหา อรหตฺตาติ, อามนฺตา.\csromanspacer{} นนุ อรหา วีตราโค วีตโทโส วีตโมโห กตกรณีโย โอหิตภาโร อนุปฺปตฺตสทตฺโถ ปริกฺขีณภวสํโยชโน สมฺมทญฺญาวิมุตฺโต อุกฺขิตฺตปลิโฆ สํกิณฺณปริโข อพฺพูเฬฺหสิโก นิรคฺคโฬ อริโย ปนฺนทฺธโช ปนฺนภาโร วิสญฺญุตฺโต สุวิชิตวิชโย, ทุกฺขํ ตสฺส ปริญฺญาตํ, สมุทโย ปหีโน, นิโรโธ สจฺฉิกโต, มคฺโค ภาวิโต, อภิญฺเญยฺยํ \csromanfolio{74}อภิญฺญาตํ, ปริญฺเญยฺยํ ปริญฺญาตํ, ปหาตพฺพํ ปหีนํ, ภาเวตพฺพํ ภาวิตํ, สจฺฉิกาตพฺพํ สจฺฉิกตนฺติ, อามนฺตา.\csromanspacer{} หญฺจิ อรหา วีตราโค วีตโทโส วีตโมโห ฯเปฯ สจฺฉิกาตพฺพํ สจฺฉิกตํ, โน จ วต เร วตฺตพฺเพ “ปริหายติ อรหา อรหตฺตา”ติ.}

\proseitem{262}{\csromansymbolmark{*}ปริหายติ อรหา อรหตฺตาติ, ( )\footnote{(อามนฺตา.) (ก)} สมยวิมุตฺโต อรหา อรหตฺตา ปริหายติ, อสมยวิมุตฺโต อรหา อรหตฺตา น ปริหายตีติ.\csromanspacer{} สมยวิมุตฺโต อรหา อรหตฺตา ปริหายตีติ, อามนฺตา.\csromanspacer{} อสมยวิมุตฺโต อรหา อรหตฺตา ปริหายตีติ\footnote{น ปริหายตีติ (ก)}, น เหวํ วตฺตพฺเพ.}

\prose{\csromansymbolmark{*}อสมยวิมุตฺโต อรหา อรหตฺตา น ปริหายตีติ, อามนฺตา.\csromanspacer{} สมยวิมุตฺโต อรหา อรหตฺตา น ปริหายตีติ, น เหวํ วตฺตพฺเพ.}

\prose{\csromansymbolmark{*}สมยวิมุตฺตสฺส อรหโต ราโค ปหีโน, ปริหายติ สมยวิมุตฺโต อรหา อรหตฺตาติ, อามนฺตา.\csromanspacer{} อสมยวิมุตฺตสฺส อรหโต ราโค ปหีโน, ปริหายติ อสมยวิมุตฺโต อรหา อรหตฺตาติ, น เหวํ วตฺตพฺเพ.}

\prose{\csromansymbolmark{*}สมยวิมุตฺตสฺส อรหโต โทโส ปหีโน ฯเปฯ อโนตฺตปฺปํ ปหีนํ, ปริหายติ สมยวิมุตฺโต อรหา อรหตฺตาติ, อามนฺตา.\csromanspacer{} อสมยวิมุตฺตสฺส อรหโต อโนตฺตปฺปํ ปหีนํ, ปริหายติ อสมยวิมุตฺโต อรหา อรหตฺตาติ, น เหวํ วตฺตพฺเพ.}

\prose{\csromansymbolmark{*}สมยวิมุตฺตสฺส อรหโต ราคปฺปหานาย มคฺโค ภาวิโต, ปริหายติ สมยวิมุตฺโต อรหา อรหตฺตาติ, อามนฺตา.\csromanspacer{} อสมยวิมุตฺตสฺส อรหโต ราคปฺปหานาย มคฺโค ภาวิโต, ปริหายติ อสมยวิมุตฺโต อรหา อรหตฺตาติ, น เหวํ วตฺตพฺเพ.}

\prose{\csromansymbolmark{*}สมยวิมุตฺตสฺส อรหโต ราคปฺปหานาย สติปฏฺฐานา ภาวิตา ฯเปฯ สมฺมปฺปธานา ภาวิตา.\csromanspacer{} อิทฺธิปาทา ภาวิตา.\csromanspacer{} อินฺทฺริยา ภาวิตา.\csromanspacer{} พลา ภาวิตา.\csromanspacer{} โพชฺฌงฺคา ภาวิตา, ปริหายติ สมยวิมุตฺโต อรหา อรหตฺตาติ, อามนฺตา.\csromanspacer{} อสมยวิมุตฺตสฺส อรหโต ราคปฺปหานาย}

\csromanheader{2}{2. ปริหานิกถา}
\prose{\csromanfolio{75}สติปฏฺฐานา ภาวิตา ฯเปฯ โพชฺฌงฺคา ภาวิตา, ปริหายติ อสมยวิมุตฺโต อรหา อรหตฺตาติ, น เหวํ วตฺตพฺเพ.}

\prose{\csromansymbolmark{*}สมยวิมุตฺตสฺส อรหโต โทสปฺปหานาย ฯเปฯอโนตฺตปฺปปหานาย มคฺโค ภาวิโต ฯเปฯ โพชฺฌงฺคา ภาวิตา, ปริหายติ สมยวิมุตฺโต อรหา อรหตฺตาติ, อามนฺตา.\csromanspacer{} อสมยวิมุตฺตสฺส อรหโต อโนตฺตปฺปปหานาย มคฺโค ภาวิโต ฯเปฯ โพชฺฌงฺคา ภาวิตา, ปริหายติ อสมยวิมุตฺโต อรหา อรหตฺตาติ, น เหวํ วตฺตพฺเพ ฯเปฯ.}

\prose{\csromansymbolmark{*}สมยวิมุตฺโต อรหา วีตราโค วีตโทโส วีตโมโห กตกรณีโย โอหิตภาโร อนุปฺปตฺตสทตฺโถ ปริกฺขีณภวสํโยชโน สมฺมทญฺญาวิมุตฺโต อุกฺขิตฺตปลิโฆ สํกิณฺณปริโข อพฺพูเฬฺหสิโก นิรคฺคโฬ อริโย ปนฺนทฺธโช ปนฺนภาโร วิสญฺญุตฺโต สุวิชิตวิชโย, ทุกฺขํ ตสฺส ปริญฺญาตํ, สมุทโย ปหีโน, นิโรโธ สจฺฉิกโต, มคฺโค ภาวิโต, อภิญฺเญยฺยํ อภิญฺญาตํ, ปริญฺเญยฺยํ ปริญฺญาตํ, ปหาตพฺพํ ปหีนํ, ภาเวตพฺพํ ภาวิตํ, สจฺฉิกาตพฺพํ สจฺฉิกตํ, ปริหายติ สมยวิมุตฺโต อรหา อรหตฺตาติ, อามนฺตา.\csromanspacer{} อสมยวิมุตฺโต อรหา วีตราโค วีตโทโส วีตโมโห ฯเปฯ สจฺฉิกาตพฺพํ สจฺฉิกตํ, ปริหายติ อสมยวิมุตฺโต อรหา อรหตฺตาติ, น เหวํ วตฺตพฺเพ.}

\proseitem{263}{\csromansymbolmark{*}อสมยวิมุตฺตสฺส อรหโต ราโค ปหีโน, น ปริหายติ อสมยวิมุตฺโต อรหา อรหตฺตาติ, อามนฺตา.\csromanspacer{} สมยวิมุตฺตสฺส อรหโต ราโค ปหีโน, น ปริหายติ สมยวิมุตฺโต อรหา อรหตฺตาติ, น เหวํ วตฺตพฺเพ.}

\prose{\csromansymbolmark{*}อสมยวิมุตฺตสฺส อรหโต โทโส ปหีโน ฯเปฯ อโนตฺตปฺปํ ปหีนํ, น ปริหายติ อสมยวิมุตฺโต อรหา อรหตฺตาติ, อามนฺตา.\csromanspacer{} สมยวิมุตฺตสฺส อรหโต อโนตฺตปฺปํ ปหีนํ, น ปริหายติ สมยวิมุตฺโต อรหา อรหตฺตาติ, น เหวํ วตฺตพฺเพ.}

\prose{\csromansymbolmark{*}อสมยวิมุตฺตสฺส อรหโต ราคปฺปหานาย มคฺโค ภาวิโต ฯเปฯ โพชฺฌงฺคา ภาวิตา, น ปริหายติ อสมยวิมุตฺโต อรหา อรหตฺตาติ, อามนฺตา.\csromanspacer{} สมยวิมุตฺตสฺส อรหโต ราคปฺปหานาย \csromanfolio{76}มคฺโค ภาวิโต ฯเปฯ โพชฺฌงฺคา ภาวิตา, น ปริหายติ สมยวิมุตฺโต อรหา อรหตฺตาติ, น เหวํ วตฺตพฺเพ.}

\prose{\csromansymbolmark{*}อสมยวิมุตฺตสฺส อรหโต โทสปฺปหานาย ฯเปฯ อโนตฺตปฺปปหานาย มคฺโค ภาวิโต ฯเปฯ โพชฺฌงฺคา ภาวิตา, น ปริหายติ อสมยวิมุตฺโต อรหา อรหตฺตาติ, อามนฺตา.\csromanspacer{} สมยวิมุตฺตสฺส อรหโต อโนตฺตปฺปปหานาย มคฺโค ภาวิโต ฯเปฯ โพชฺฌงฺคา ภาวิตา, น ปริหายติ สมยวิมุตฺโต อรหา อรหตฺตาติ, น เหวํ วตฺตพฺเพ.}

\prose{\csromansymbolmark{*}อสมยวิมุตฺโต อรหา วีตราโค วีตโทโส วีตโมโห ฯเปฯ สจฺฉิกาตพฺพํ สจฺฉิกตํ, น ปริหายติ อสมยวิมุตฺโต อรหา อรหตฺตาติ, อามนฺตา.\csromanspacer{} สมยวิมุตฺโต อรหา วีตราโค วีตโทโส วีตโมโห ฯเปฯ สจฺฉิกาตพฺพํ สจฺฉิกตํ, น ปริหายติ สมยวิมุตฺโต อรหา อรหตฺตาติ, น เหวํ วตฺตพฺเพ ฯเปฯ.}

\proseitem{264}{\csromansymbolmark{*}ปริหายติ อรหา อรหตฺตาติ, อามนฺตา.\csromanspacer{} สาริปุตฺโต เถโร ปริหายิตฺถ อรหตฺตาติ, น เหวํ วตฺตพฺเพ.\csromanspacer{} มหาโมคฺคลฺลาโน\footnote{มหาโมคฺคลาโน (ก)} เถโร.\csromanspacer{} มหากสฺสโป เถโร.\csromanspacer{} มหากจฺจายโน\footnote{มหากจฺจาโน (ม 1. 156 ปิฏฺฐาทีสุ)} เถโร.\csromanspacer{} มหาโกฏฺฐิโก\footnote{มหาโกฏฺฐิโต (สี, สฺยา, กํ, อิ)} เถโร.\csromanspacer{} มหาปนฺถโก เถโร ปริหายิตฺถ อรหตฺตาติ, น เหวํ วตฺตพฺเพ.}

\prose{\csromansymbolmark{*}สาริปุตฺโต เถโร น ปริหายิตฺถ อรหตฺตาติ, อามนฺตา.\csromanspacer{} หญฺจิ สาริปุตฺโต เถโร น ปริหายิตฺถ อรหตฺตา, โน จ วต เร วตฺตพฺเพ “ปริหายติ อรหา อรหตฺตา”ติ.}

\prose{\csromansymbolmark{*}มหาโมคฺคลฺลาโน เถโร.\csromanspacer{} มหากสฺสโป เถโร.\csromanspacer{} มหากจฺจายโน เถโร.\csromanspacer{} มหาโกฏฺฐิโก เถโร.\csromanspacer{} มหาปนฺถโก เถโร น ปริหายิตฺถ อรหตฺตาติ, อามนฺตา.\csromanspacer{} หญฺจิ มหาปนฺถโก เถโร น ปริหายิตฺถ อรหตฺตา, โน จ วต เร วตฺตพฺเพ “ปริหายติ อรหา อรหตฺตา”ติ.}

\vspace{\tipitakaparskip}
\csromancenter{อริยปุคฺคลสํสนฺทนํ.}
\csromanheader{2}{2. ปริหานิกถา}
\csromanheader{2}{3. สุตฺตสาธนปริหานิ}
\proseitem{265}{\csromanfolio{77}\csromansymbolmark{*}ปริหายติ อรหา อรหตฺตาติ, อามนฺตา.\csromanspacer{} นนุ วุตฺตํ ภควตา\csromandash{}}

\csromangathameasure{“อุจฺจาวจา หิ ปฏิปทา, สมเณน ปกาสิตา. \\ น ปารํ ทิคุณํ ยนฺติ, นยิทํ เอกคุณํ มุตนฺ”ติ.}
\csromangathagroup{\csromangathastanza{“อุจฺจาวจา หิ ปฏิปทา\footnote{ปฏิปาทา (ฏฺฐ)}, สมเณน ปกาสิตา. \\ น ปารํ ทิคุณํ ยนฺติ, นยิทํ เอกคุณํ มุตนฺ”ติ\footnote{ปฏิปาทา (ฏฺฐ)}.}}

\prose{อตฺเถว สุตฺตนฺโตติ, อามนฺตา.\csromanspacer{} เตน หิ น วตฺตพฺพํ “ปริหายติ อรหา อรหตฺตา”ติ.}

\prose{\csromansymbolmark{*}ปริหายติ อรหา อรหตฺตาติ, อามนฺตา.\csromanspacer{} อตฺถิ ฉินฺนสฺส เฉทิยนฺติ, น เหวํ วตฺตพฺเพ.}

\prose{\csromansymbolmark{*}อตฺถิ ฉินฺนสฺส เฉทิยนฺติ, อามนฺตา.\csromanspacer{} นนุ วุตฺตํ ภควตา\csromandash{}}

\csromangathameasure{“วีตตโณฺห อนาทาโน, กิจฺจํ ยสฺส น วิชฺชติ. \\ ฉินฺนสฺส เฉทิยํ นตฺถิ, โอฆปาโส สมูหโต”ติ.}
\csromangathagroup{\csromangathastanza{“วีตตโณฺห อนาทาโน, กิจฺจํ ยสฺส น วิชฺชติ. \\ ฉินฺนสฺส เฉทิยํ นตฺถิ, โอฆปาโส สมูหโต”ติ.}}

\prose{อตฺเถว สุตฺตนฺโตติ, อามนฺตา.\csromanspacer{} เตน หิ น วตฺตพฺพํ “อตฺถิ ฉินฺนสฺส เฉทิยนฺ”ติ.}

\proseitem{266}{\csromansymbolmark{*}ปริหายติ อรหา อรหตฺตาติ, อามนฺตา.\csromanspacer{} อตฺถิ กตสฺส ปติจโยติ, น เหวํ วตฺตพฺเพ.}

\prose{\csromansymbolmark{*}อตฺถิ กตสฺส ปติจโยติ, อามนฺตา.\csromanspacer{} นนุ วุตฺตํ ภควตา\csromandash{}}

\csromangathameasure{“ตสฺส สมฺมา วิมุตฺตสฺส, สนฺตจิตฺตสฺส ภิกฺขุโน. \\ กตสฺส ปติจโย นตฺถิ, กรณียํ น วิชฺชติ. \\ เสโล ยถา เอกคฺฆโน, วาเตน น สมีรติ. \\ เอวํ รูปา รสา สทฺทา, คนฺธา ผสฺสา จ เกวลา. \\ อิฏฺฐา ธมฺมา อนิฏฺฐา จ, นปฺปเวเธนฺติ ตาทิโน. \\ ฐิตํ จิตฺตํ วิปฺปมุตฺตํ, วยํ จสฺสานุปสฺสตี”ติ.}
\csromangathagroup{\csromangathastanza{“ตสฺส สมฺมา วิมุตฺตสฺส, สนฺตจิตฺตสฺส ภิกฺขุโน. \\ กตสฺส ปติจโย นตฺถิ, กรณียํ น วิชฺชติ.}\csromangathastanza{เสโล ยถา เอกคฺฆโน, วาเตน น สมีรติ. \\ เอวํ รูปา รสา สทฺทา, คนฺธา ผสฺสา จ เกวลา.}\csromangathastanza{อิฏฺฐา ธมฺมา อนิฏฺฐา จ, นปฺปเวเธนฺติ ตาทิโน. \\ ฐิตํ จิตฺตํ วิปฺปมุตฺตํ, วยํ จสฺสานุปสฺสตี”ติ\footnote{วิ 3. 272 ปิฏฺเฐ; อํ 2. 333 ปิฏฺเฐปิ.}.}}

\prose{อตฺเถว สุตฺตนฺโตติ, อามนฺตา.\csromanspacer{} เตน หิ น วตฺตพฺพํ “อตฺถิ กตสฺส ปติจโย”ติ.}

\proseitem{267}{\csromanfolio{78}\csromansymbolmark{+}น วตฺตพฺพํ “ปริหายติ อรหา อรหตฺตา”ติ, อามนฺตา.\csromanspacer{} นนุ วุตฺตํ ภควตา\csromandash{}ปญฺจิเม ภิกฺขเว ธมฺมา สมยวิมุตฺตสฺส ภิกฺขุโน ปริหานาย สํวตฺตนฺติ.\csromanspacer{} กตเม ปญฺจ, กมฺมารามตา ภสฺสารามตา นิทฺทารามตา สงฺคณิการามตา ยถาวิมุตฺตํ จิตฺตํ น ปจฺจเวกฺขติ.\csromanspacer{} อิเม โข ภิกฺขเว ปญฺจ ธมฺมา สมยวิมุตฺตสฺส ภิกฺขุโน ปริหานาย สํวตฺตนฺตีติ\footnote{อํ 2. 153 ปิฏฺเฐ.} อตฺเถว สุตฺตนฺโตติ, อามนฺตา.\csromanspacer{} เตน หิ ปริหายติ “อรหา อรหตฺตา”ติ.}

\prose{\csromansymbolmark{*}อตฺถิ อรหโต กมฺมารามตาติ, น เหวํ วตฺตพฺเพ.}

\prose{\csromansymbolmark{*}อตฺถิ อรหโต กมฺมารามตาติ, อามนฺตา.\csromanspacer{} อตฺถิ อรหโต ราโค กามราโค กามราคปริยุฏฺฐานํ กามราคสํโยชนํ กาโมโฆ กามโยโค กามจฺฉนฺทนีวรณนฺติ, น เหวํ วตฺตพฺเพ.}

\prose{\csromansymbolmark{*}อตฺถิ อรหโต ภสฺสารามตา.\csromanspacer{} อตฺถิ อรหโต นิทฺทารามตา.\csromanspacer{} อตฺถิ อรหโต สงฺคณิการามตาติ, น เหวํ วตฺตพฺเพ.}

\prose{\csromansymbolmark{*}อตฺถิ อรหโต สงฺคณิการามตาติ, อามนฺตา.\csromanspacer{} อตฺถิ อรหโต ราโค กามราโค กามราคปริยุฏฺฐานํ กามราคสํโยชนํ กาโมโฆ กามโยโค กามจฺฉนฺทนีวรณนฺติ, น เหวํ วตฺตพฺเพ.}

\proseitem{268}{\csromansymbolmark{*}ปริหายติ อรหา อรหตฺตาติ, อามนฺตา.\csromanspacer{} อรหา อรหตฺตา ปริหายมาโน กิํ ปริยุฏฺฐิโต ปริหายตีติ, ราคปริยุฏฺฐิโต ปริหายตีติ.\csromanspacer{} ปริยุฏฺฐานํ กิํ ปฏิจฺจ อุปฺปชฺชตีติ, อนุสยํ ปฏิจฺจ อุปฺปชฺชตีติ.\csromanspacer{} อตฺถิ อรหโต อนุสยาติ, น เหวํ วตฺตพฺเพ.}

\prose{\csromansymbolmark{*}อตฺถิ อรหโต อนุสยาติ, อามนฺตา.\csromanspacer{} อตฺถิ อรหโต กามราคานุสโย ปฏิฆานุสโย มานานุสโย ทิฏฺฐานุสโย วิจิกิจฺฉานุสโย ภวราคานุสโย อวิชฺชานุสโยติ, น เหวํ วตฺตพฺเพ.}

\prose{โทสปริยุฏฺฐิโต ปริหายตีติ ฯเปฯ.\csromanspacer{} โมหปริยุฏฺฐิโต ปริหายตีติ.\csromanspacer{} ปริยุฏฺฐานํ กิํ ปฏิจฺจ อุปฺปชฺชตีติ, อนุสยํ ปฏิจฺจ อุปฺปชฺชตีติ.\csromanspacer{} อตฺถิ อรหโต อนุสยาติ, น เหวํ วตฺตพฺเพ ฯเปฯ.}

\prose{\csromansymbolmark{*}อตฺถิ อรหโต อนุสยาติ, อามนฺตา.\csromanspacer{} อตฺถิ อรหโต กามราคานุสโย ฯเปฯ อวิชฺชานุสโยติ, น เหวํ วตฺตพฺเพ.}

\csromanheader{2}{2. ปริหานิกถา}
\prose{\csromanfolio{79}\csromansymbolmark{*}ปริหายติ อรหา อรหตฺตาติ, อามนฺตา.\csromanspacer{} อรหโต อรหตฺตา ปริหายมานสฺส กิํ อุปจยํ คจฺฉตีติ, ราโค อุปจยํ คจฺฉตีติ.\csromanspacer{} สกฺกายทิฏฺฐิ อุปจยํ คจฺฉตีติ.\csromanspacer{} วิจิกิจฺฉา อุปจยํ คจฺฉตีติ.\csromanspacer{} สีลพฺพตปรามาโส อุปจยํ คจฺฉตีติ, น เหวํ วตฺตพฺเพ.\csromanspacer{} โทโส อุปจยํ คจฺฉตีติ ฯเปฯ.\csromanspacer{} โมโห อุปจยํ คจฺฉตีติ.\csromanspacer{} สกฺกายทิฏฺฐิ อุปจยํ คจฺฉตีติ.\csromanspacer{} วิจิกิจฺฉา อุปจยํ คจฺฉตีติ.\csromanspacer{} สีลพฺพตปรามาโส อุปจยํ คจฺฉตีติ, น เหวํ วตฺตพฺเพ.}

\prose{\csromansymbolmark{*}ปริหายติ อรหา อรหตฺตาติ, อามนฺตา.\csromanspacer{} อรหา อาจินตีติ, น เหวํ วตฺตพฺเพ.\csromanspacer{} อรหา อปจินตีติ, น เหวํ วตฺตพฺเพ.\csromanspacer{} อรหา ปชหตีติ, น เหวํ วตฺตพฺเพ.\csromanspacer{} อรหา อุปาทิยตีติ, น เหวํ วตฺตพฺเพ.\csromanspacer{} อรหา วิสิเนตีติ, น เหวํ วตฺตพฺเพ.\csromanspacer{} อรหา อุสฺสิเนตีติ, น เหวํ วตฺตพฺเพ.\csromanspacer{} อรหา วิธูเปตีติ, น เหวํ วตฺตพฺเพ.\csromanspacer{} อรหา สนฺธูเปตีติ, น เหวํ วตฺตพฺเพ.}

\prose{นนุ อรหา เนวาจินติ น อปจินติ อปจินิตฺวา ฐิโตติ, อามนฺตา.\csromanspacer{} หญฺจิ อรหา เนวาจินติ น อปจินติ อปจินิตฺวา ฐิโต, โน จ วต เร วตฺตพฺเพ “ปริหายติ อรหา อรหตฺตา”ติ.}

\prose{นนุ อรหา เนว ปชหติ น อุปาทิยติ ปชหิตฺวา ฐิโตติ, อามนฺตา.\csromanspacer{} หญฺจิ อรหา เนว ปชหติ น อุปาทิยติ ปชหิตฺวา ฐิโต, โน จ วต เร วตฺตพฺเพ “ปริหายติ อรหา อรหตฺตา”ติ.}

\prose{นนุ อรหา เนว วิสิเนติ น อุสฺสิเนติ วิสินิตฺวา ฐิโตติ, อามนฺตา.\csromanspacer{} หญฺจิ อรหา เนว วิสิเนติ น อุสฺสิเนติ วิสินิตฺวา ฐิโต, โน จ วต เร วตฺตพฺเพ “ปริหายติ อรหา อรหตฺตา”ติ.}

\prose{นนุ อรหา เนว วิธูเปติ น สนฺธูเปติ วิธูเปตฺวา ฐิโตติ, อามนฺตา.\csromanspacer{} หญฺจิ อรหา เนว วิธูเปติ น สนฺธูเปติ วิธูเปตฺวา ฐิโต, โน จ วต เร วตฺตพฺเพ “ปริหายติ อรหา อรหตฺตา”ติ.}

\nitthitam{\textbf{ปริหานิกถา นิฏฺฐิตา.}}
\csromanmatikaanchor{mtk.21}
\csromantocmarkref{chapter}{3. พฺรหฺมจริยกถา (ก)}{mtk.21}
\bookmark[dest={mtk.21},level=0]{3. พฺรหฺมจริยกถา (ก)}
\chapterheadpageread{3. พฺรหฺมจริยกถา}
\csromanheader{2}{1. สุทฺธพฺรหฺมจริยกถา}
\proseitem{269}{\csromanfolio{80}\csromansymbolmark{*}นตฺถิ เทเวสุ พฺรหฺมจริยวาโสติ, อามนฺตา.\csromanspacer{} สพฺเพ เทวา ชฬา เอลมูคา\footnote{เอฬมูคา (สฺยา)} อวิญฺญู หตฺถสํวาจิกา นปฺปฏิพลา สุภาสิตทุพฺภาสิตานํ อตฺถมญฺญาตุํ, สพฺเพ เทวา น พุทฺเธ ปสนฺนา น ธมฺเม ปสนฺนา น สํเฆ ปสนฺนา, น พุทฺธํ ภควนฺตํ ปยิรุปาสนฺติ, น พุทฺธํ ภควนฺตํ ปญฺหํ ปุจฺฉนฺติ, น พุทฺเธน ภควตา ปเญฺห วิสฺสชฺชิเต อตฺตมนา, สพฺเพ เทวา กมฺมาวรเณน สมนฺนาคตา กิเลสาวรเณน สมนฺนาคตา วิปากาวรเณน สมนฺนาคตา อสฺสทฺธา อจฺฉนฺทิกา ทุปฺปญฺญา อภพฺพา นิยามํ โอกฺกมิตุํ กุสเลสุ ธมฺเมสุ สมฺมตฺตํ, สพฺเพ เทวา มาตุฆาตกา ปิตุฆาตกา อรหนฺตฆาตกา รุหิรุปฺปาทกา สํฆเภทกา, สพฺเพ เทวา ปาณาติปาติโน อทินฺนาทายิโน กาเมสุมิจฺฉาจาริโน มุสาวาทิโน ปิสุณาวาจา ผรุสาวาจา สมฺผปฺปลาปิโน อภิชฺฌาลุโน พฺยาปนฺนจิตฺตา มิจฺฉาทิฏฺฐิกาติ, น เหวํ วตฺตพฺเพ ฯเปฯ.}

\prose{นนุ อตฺถิ เทวา อชฬา อเนลมูคา วิญฺญู น หตฺถสํวาจิกา ปฏิพลา สุภาสิตทุพฺภาสิตานํ อตฺถมญฺญาตุํ, อตฺถิ เทวา พุทฺเธ ปสนฺนา ธมฺเม ปสนฺนา สํเฆ ปสนฺนา, พุทฺธํ ภควนฺตํ ปยิรุปาสนฺติ, พุทฺธํ ภควนฺตํ ปญฺหํ ปุจฺฉนฺติ, พุทฺเธน ภควตา ปเญฺห วิสฺสชฺชิเต อตฺตมนา โหนฺติ, อตฺถิ เทวา น กมฺมาวรเณน สมนฺนาคตา น กิเลสาวรเณน สมนฺนาคตา น วิปากาวรเณน สมนฺนาคตา สทฺธา ฉนฺทิกา ปญฺญวนฺโต ภพฺพา นิยามํ โอกฺกมิตุํ กุสเลสุ ธมฺเมสุ สมฺมตฺตํ, อตฺถิ เทวา น มาตุฆาตกา น ปิตุฆาตกา น อรหนฺตฆาตกา น รุหิรุปฺปาทกา น สํฆเภทกา, อตฺถิ เทวา น ปาณาติปาติโน น อทินฺนาทายิโน น กาเมสุมิจฺฉาจาริโน น มุสาวาทิโน น ปิสุณาวาจา น ผรุสาวาจา น สมฺผปฺปลาปิโน น อภิชฺฌาลุโน อพฺยาปนฺนจิตฺตา สมฺมาทิฏฺฐิกาติ, อามนฺตา.}

\prose{หญฺจิ อตฺถิ เทวา อชฬา อเนลมูคา วิญฺญู น หตฺถสํวาจิกา ปฏิพลา สุภาสิตทุพฺภาสิตานํ อตฺถมญฺญาตุํ ฯเปฯ อตฺถิ เทวา พุทฺเธ ปสนฺนา ฯเปฯ สมฺมาทิฏฺฐิกา, โน จ วต เร วตฺตพฺเพ “นตฺถิ เทเวสุ พฺรหฺมจริยวาโส”ติ.}

\csromanheader{2}{3. พฺรหฺมจริยกถา}
\proseitem{270}{\csromanfolio{81}\csromansymbolmark{+}อตฺถิ เทเวสุ พฺรหฺมจริยวาโสติ, อามนฺตา.\csromanspacer{} อตฺถิ ตตฺถ ปพฺพชฺชา มุณฺฑิยํ กาสาวธารณา ปตฺตธารณา, เทเวสุ สมฺมาสมฺพุทฺธา อุปฺปชฺชนฺติ, ปจฺเจกสมฺพุทฺธา อุปฺปชฺชนฺติ, สาวกยุคํ อุปฺปชฺชตีติ, น เหวํ วตฺตพฺเพ ฯเปฯ.}

\prose{\csromansymbolmark{*}เทเวสุ ปพฺพชฺชา นตฺถีติ นตฺถิ เทเวสุ พฺรหฺมจริยวาโสติ, อามนฺตา.\csromanspacer{} ยตฺถ อตฺถิ ปพฺพชฺชา, ตตฺเถว พฺรหฺมจริยวาโส.\csromanspacer{} ยตฺถ นตฺถิ ปพฺพชฺชา, นตฺถิ ตตฺถ พฺรหฺมจริยวาโสติ, น เหวํ วตฺตพฺเพ ฯเปฯ.\csromanspacer{} ยตฺถ อตฺถิ ปพฺพชฺชา, ตตฺเถว พฺรหฺมจริยวาโส.\csromanspacer{} ยตฺถ นตฺถิ ปพฺพชฺชา, นตฺถิ ตตฺถ พฺรหฺมจริยวาโสติ, อามนฺตา.\csromanspacer{} โย ปพฺพชติ, ตสฺเสว พฺรหฺมจริยวาโส.\csromanspacer{} โย น ปพฺพชติ, นตฺถิ ตสฺส พฺรหฺมจริยวาโสติ, น เหวํ วตฺตพฺเพ ฯเปฯ.}

\prose{\csromansymbolmark{*}เทเวสุ มุณฺฑิยํ นตฺถีติ นตฺถิ เทเวสุ พฺรหฺมจริยวาโสติ, อามนฺตา.\csromanspacer{} ยตฺถ อตฺถิ มุณฺฑิยํ, ตตฺเถว พฺรหฺมจริยวาโส.\csromanspacer{} ยตฺถ นตฺถิ มุณฺฑิยํ, นตฺถิ ตตฺถ พฺรหฺมจริยวาโสติ, น เหวํ วตฺตพฺเพ ฯเปฯ.\csromanspacer{} ยตฺถ อตฺถิ มุณฺฑิยํ, ตตฺเถว พฺรหฺมจริยวาโส.\csromanspacer{} ยตฺถ นตฺถิ มุณฺฑิยํ, นตฺถิ ตตฺถ พฺรหฺมจริยวาโสติ, อามนฺตา.\csromanspacer{} โย มุณฺโฑ โหติ, ตสฺเสว พฺรหฺมจริยวาโส.\csromanspacer{} โย มุณฺโฑ น โหติ, นตฺถิ ตสฺส พฺรหฺมจริยวาโสติ, น เหวํ วตฺตพฺเพ ฯเปฯ.}

\prose{\csromansymbolmark{*}เทเวสุ กาสาวธารณา นตฺถีติ นตฺถิ เทเวสุ พฺรหฺมจริยวาโสติ, อามนฺตา.\csromanspacer{} ยตฺถ อตฺถิ กาสาวธารณา, ตตฺเถว พฺรหฺมจริยวาโส.\csromanspacer{} ยตฺถ นตฺถิ กาสาวธารณา, นตฺถิ ตตฺถ พฺรหฺมจริยวาโสติ, น เหวํ วตฺตพฺเพ ฯเปฯ.\csromanspacer{} ยตฺถ อตฺถิ กาสาวธารณา, ตตฺเถว พฺรหฺมจริยวาโส.\csromanspacer{} ยตฺถ นตฺถิ กาสาวธารณา, นตฺถิ ตตฺถ พฺรหฺมจริยวาโสติ, อามนฺตา.\csromanspacer{} โย กาสาวํ ธาเรติ, ตสฺเสว พฺรหฺมจริยวาโส.\csromanspacer{} โย กาสาวํ น ธาเรติ, นตฺถิ ตสฺส พฺรหฺมจริยวาโสติ, น เหวํ วตฺตพฺเพ ฯเปฯ.}

\prose{\csromansymbolmark{*}เทเวสุ ปตฺตธารณา นตฺถีติ นตฺถิ เทเวสุ พฺรหฺมจริยวาโสติ, อามนฺตา.\csromanspacer{} ยตฺถ อตฺถิ ปตฺตธารณา, ตตฺเถว พฺรหฺมจริยวาโส.\csromanspacer{} ยตฺถ นตฺถิ ปตฺตธารณา, นตฺถิ ตตฺถ พฺรหฺมจริยวาโสติ, น เหวํ วตฺตพฺเพ ฯเปฯ.\csromanspacer{} ยตฺถ อตฺถิ ปตฺตธารณา, ตตฺเถว พฺรหฺมจริยวาโส.\csromanspacer{} ยตฺถ นตฺถิ ปตฺตธารณา, นตฺถิ ตตฺถ พฺรหฺมจริยวาโสติ, อามนฺตา.\csromanspacer{} โย ปตฺตํ ธาเรติ, ตสฺเสว พฺรหฺมจริยวาโส.\csromanspacer{} โย ปตฺตํ น ธาเรติ, นตฺถิ ตสฺส พฺรหฺมจริยวาโสติ, น เหวํ วตฺตพฺเพ ฯเปฯ.}

\prose{\csromanfolio{82}\csromansymbolmark{*}เทเวสุ สมฺมาสมฺพุทฺธา นุปฺปชฺชนฺตีติ นตฺถิ เทเวสุ พฺรหฺมจริยวาโสติ, อามนฺตา.\csromanspacer{} ยตฺถ สมฺมาสมฺพุทฺธา อุปฺปชฺชนฺติ, ตตฺเถว พฺรหฺมจริยวาโส.\csromanspacer{} ยตฺถ สมฺมาสมฺพุทฺธา นุปฺปชฺชนฺติ, นตฺถิ ตตฺถ พฺรหฺมจริยวาโสติ, น เหวํ วตฺตพฺเพ ฯเปฯ.\csromanspacer{} ยตฺถ สมฺมาสมฺพุทฺธา อุปฺปชฺชนฺติ, ตตฺเถว พฺรหฺมจริยวาโส.\csromanspacer{} ยตฺถ สมฺมาสมฺพุทฺธา นุปฺปชฺชนฺติ, นตฺถิ ตตฺถ พฺรหฺมจริยวาโสติ, อามนฺตา.\csromanspacer{} ลุมฺพินิยา ภควา ชาโต, โพธิยา มูเล อภิสมฺพุทฺโธ, พาราณสิยํ ภควตา ธมฺมจกฺกํ ปวตฺติตํ, ตตฺเถว พฺรหฺมจริยวาโส.\csromanspacer{} นตฺถญฺญตฺร พฺรหฺมจริยวาโสติ, น เหวํ วตฺตพฺเพ ฯเปฯ.}

\prose{\csromansymbolmark{*}เทเวสุ ปจฺเจกสมฺพุทฺธา นุปฺปชฺชนฺตีติ นตฺถิ เทเวสุ พฺรหฺมจริยวาโสติ, อามนฺตา.\csromanspacer{} ยตฺถ ปจฺเจกสมฺพุทฺธา อุปฺปชฺชนฺติ, ตตฺเถว พฺรหฺมจริยวาโส.\csromanspacer{} ยตฺถ ปจฺเจกสมฺพุทฺธา นุปฺปชฺชนฺติ, นตฺถิ ตตฺถ พฺรหฺมจริยวาโสติ, น เหวํ วตฺตพฺเพ ฯเปฯ.\csromanspacer{} ยตฺถ ปจฺเจกสมฺพุทฺธา อุปฺปชฺชนฺติ, ตตฺเถว พฺรหฺมจริยวาโส.\csromanspacer{} ยตฺถ ปจฺเจกสมฺพุทฺธา นุปฺปชฺชนฺติ, นตฺถิ ตตฺถ พฺรหฺมจริยวาโสติ, อามนฺตา.\csromanspacer{} มชฺฌิเมสุ ชนปเทสุ ปจฺเจกสมฺพุทฺธา อุปฺปชฺชนฺติ, ตตฺเถว พฺรหฺมจริยวาโส.\csromanspacer{} นตฺถญฺญตฺร พฺรหฺมจริยวาโสติ, น เหวํ วตฺตพฺเพ ฯเปฯ.}

\prose{\csromansymbolmark{*}เทเวสุ สาวกยุคํ นุปฺปชฺชตีติ นตฺถิ เทเวสุ พฺรหฺมจริยวาโสติ, อามนฺตา.\csromanspacer{} ยตฺถ สาวกยุคํ อุปฺปชฺชติ, ตตฺเถว พฺรหฺมจริยวาโส.\csromanspacer{} ยตฺถ สาวกยุคํ นุปฺปชฺชติ, นตฺถิ ตตฺถ พฺรหฺมจริยวาโสติ, น เหวํ วตฺตพฺเพ ฯเปฯ.\csromanspacer{} ยตฺถ สาวกยุคํ อุปฺปชฺชติ, ตตฺเถว พฺรหฺมจริยวาโส.\csromanspacer{} ยตฺถ สาวกยุคํ นุปฺปชฺชติ, นตฺถิ ตตฺถ พฺรหฺมจริยวาโสติ, อามนฺตา.\csromanspacer{} มคเธสุ สาวกยุคํ อุปฺปนฺนํ, ตตฺเถว พฺรหฺมจริยวาโส.\csromanspacer{} นตฺถญฺญตฺร พฺรหฺมจริยวาโสติ, น เหวํ วตฺตพฺเพ ฯเปฯ.}

\proseitem{271}{\csromansymbolmark{+}อตฺถิ เทเวสุ พฺรหฺมจริยวาโสติ, อามนฺตา.\csromanspacer{} สพฺพเทเวสุ อตฺถิ พฺรหฺมจริยวาโสติ, น เหวํ วตฺตพฺเพ ฯเปฯ.}

\prose{\csromansymbolmark{*}อตฺถิ มนุสฺเสสุ พฺรหฺมจริยวาโสติ, อามนฺตา.\csromanspacer{} สพฺพมนุสฺเสสุ อตฺถิ พฺรหฺมจริยวาโสติ, น เหวํ วตฺตพฺเพ ฯเปฯ.}

\prose{\csromansymbolmark{+}อตฺถิ เทเวสุ พฺรหฺมจริยวาโสติ, อามนฺตา.\csromanspacer{} อสญฺญสตฺเตสุ เทเวสุ อตฺถิ พฺรหฺมจริยวาโสติ, น เหวํ วตฺตพฺเพ ฯเปฯ.}

\csromanheader{2}{3. พฺรหฺมจริยกถา}
\prose{\csromanfolio{83}\csromansymbolmark{*}อตฺถิ มนุสฺเสสุ พฺรหฺมจริยวาโสติ, อามนฺตา.\csromanspacer{} ปจฺจนฺติเมสุ ชนปเทสุ อตฺถิ พฺรหฺมจริยวาโส มิลกฺเขสุ\footnote{มิลกฺขูสุ (สฺยา, ก)} อวิญฺญาตาเรสุ, ยตฺถ นตฺถิ คติ ภิกฺขูนํ ภิกฺขุนีนํ อุปาสกานํ อุปาสิกานนฺติ, น เหวํ วตฺตพฺเพ.}

\prose{\csromansymbolmark{+}อตฺถิ เทเวสุ พฺรหฺมจริยวาโสติ? อตฺถิ ยตฺถ อตฺถิ, อตฺถิ ยตฺถ นตฺถีติ.\csromanspacer{} อสญฺญสตฺเตสุ เทเวสุ อตฺถิ ยตฺถ อตฺถิ, อตฺถิ ยตฺถ นตฺถิ พฺรหฺมจริยวาโส, สญฺญสตฺเตสุ\footnote{อสญฺญสตฺเตสุ (ก)} เทเวสุ อตฺถิ ยตฺถ อตฺถิ, อตฺถิ ยตฺถ นตฺถิ พฺรหฺมจริยวาโสติ, น เหวํ วตฺตพฺเพ.}

\prose{\csromansymbolmark{+}เทเวสุ อตฺถิ ยตฺถ อตฺถิ, อตฺถิ ยตฺถ นตฺถิ พฺรหฺมจริยวาโสติ, อามนฺตา.\csromanspacer{} กตฺถ อตฺถิ, กตฺถ นตฺถีติ, อสญฺญสตฺเตสุ เทเวสุ นตฺถิ พฺรหฺมจริยวาโส, สญฺญสตฺเตสุ2 เทเวสุ อตฺถิ พฺรหฺมจริยวาโสติ.\csromanspacer{} อสญฺญสตฺเตสุ เทเวสุ นตฺถิ พฺรหฺมจริยวาโสติ, อามนฺตา.\csromanspacer{} สญฺญสตฺเตสุ2 เทเวสุ นตฺถิ พฺรหฺมจริยวาโสติ, น เหวํ วตฺตพฺเพ.}

\prose{\csromansymbolmark{+}สญฺญสตฺเตสุ2 เทเวสุ อตฺถิ พฺรหฺมจริยวาโสติ, อามนฺตา.\csromanspacer{} อสญฺญสตฺเตสุ เทเวสุ อตฺถิ พฺรหฺมจริยวาโสติ, น เหวํ วตฺตพฺเพ.}

\prose{\csromansymbolmark{*}อตฺถิ มนุสฺเสสุ พฺรหฺมจริยวาโสติ? อตฺถิ ยตฺถ อตฺถิ, อตฺถิ ยตฺถ นตฺถีติ.\csromanspacer{} ปจฺจนฺติเมสุ ชนปเทสุ อตฺถิ ยตฺถ อตฺถิ, อตฺถิ ยตฺถ นตฺถิ พฺรหฺมจริยวาโส มิลกฺเขสุ อวิญฺญาตาเรสุ ยตฺถ นตฺถิ คติ ภิกฺขูนํ ภิกฺขุนีนํ อุปาสกานํ อุปาสิกานํ, มชฺฌิเมสุ ชนปเทสุ อตฺถิ ยตฺถ อตฺถิ, อตฺถิ ยตฺถ นตฺถิ พฺรหฺมจริยวาโสติ, น เหวํ วตฺตพฺเพ.}

\prose{\csromansymbolmark{*}มนุสฺเสสุ อตฺถิ ยตฺถ อตฺถิ, อตฺถิ ยตฺถ นตฺถิ พฺรหฺมจริยวาโสติ, อามนฺตา.\csromanspacer{} กตฺถ อตฺถิ, กตฺถ นตฺถีติ, ปจฺจนฺติเมสุ ชนปเทสุ นตฺถิ พฺรหฺมจริยวาโส มิลกฺเขสุ อวิญฺญาตาเรสุ ยตฺถ นตฺถิ คติ ภิกฺขูนํ ภิกฺขุนีนํ อุปาสกานํ อุปาสิกานํ, มชฺฌิเมสุ ชนปเทสุ อตฺถิ พฺรหฺมจริยวาโสติ.\csromanspacer{} ปจฺจนฺติเมสุ ชนปเทสุ นตฺถิ พฺรหฺมจริยวาโส มิลกฺเขสุ อวิญฺญาตาเรสุ ยตฺถ นตฺถิ คติ ภิกฺขูนํ ภิกฺขุนีนํ อุปาสกานํ อุปาสิกานนฺติ, อามนฺตา.\csromanspacer{} มชฺฌิเมสุ ชนปเทสุ นตฺถิ พฺรหฺมจริยวาโสติ, น เหวํ วตฺตพฺเพ.}

\prose{\csromanfolio{84}\csromansymbolmark{*}มชฺฌิเมสุ ชนปเทสุ อตฺถิ พฺรหฺมจริยวาโสติ, อามนฺตา.\csromanspacer{} ปจฺจนฺติเมสุ ชนปเทสุ อตฺถิ พฺรหฺมจริยวาโส มิลกฺเขสุ อวิญฺญาตาเรสุ ยตฺถ นตฺถิ คติ ภิกฺขูนํ ภิกฺขุนีนํ อุปาสกานํ อุปาสิกานนฺติ, น เหวํ วตฺตพฺเพ.}

\prose{\csromansymbolmark{+}อตฺถิ เทเวสุ พฺรหฺมจริยวาโสติ, อามนฺตา.\csromanspacer{} นนุ วุตฺตํ ภควตา “ตีหิ ภิกฺขเว ฐาเนหิ ชมฺพุทีปกา มนุสฺสา อุตฺตรกุรุเก จ มนุสฺเส อธิคฺคณฺหนฺติ เทเว จ ตาวติํเส.\csromanspacer{} กตเมหิ ตีหิ, สูรา สติมนฺโต อิธ พฺรหฺมจริยวาโส”ติ\footnote{อํ 3. 198 ปิฏฺเฐ.} อตฺเถว สุตฺตนฺโตติ, อามนฺตา.\csromanspacer{} เตน หิ นตฺถิ เทเวสุ พฺรหฺมจริยวาโสติ.}

\prose{\csromansymbolmark{*}สาวตฺถิยํ วุตฺตํ ภควตา “อิธ พฺรหฺมจริยวาโส”ติ, อามนฺตา.\csromanspacer{} สาวตฺถิยํเยว พฺรหฺมจริยวาโส, นตฺถิ อญฺญตฺร พฺรหฺมจริยวาโสติ, น เหวํ วตฺตพฺเพ.}

\proseitem{272}{\csromansymbolmark{*}อนาคามิสฺส ปุคฺคลสฺส ปญฺโจรมฺภาคิยานิ สํโยชนานิ ปหีนานิ, ปญฺจุทฺธมฺภาคิยานิ สํโยชนานิ อปฺปหีนานิ, อิโต จุตสฺส ตตฺถ อุปปนฺนสฺส กุหิํ ผลุปฺปตฺตีติ, ตตฺเถว.\csromanspacer{} หญฺจิ อนาคามิสฺส ปุคฺคลสฺส ปญฺโจรมฺภาคิยานิ สํโยชนานิ ปหีนานิ ปญฺจุทฺธมฺภาคิยานิ สํโยชนานิ อปฺปหีนานิ อิโต จุตสฺส ตตฺถ อุปปนฺนสฺส ตหิํ ผลุปฺปตฺติ, โน จ วต เร วตฺตพฺเพ “นตฺถิ เทเวสุ พฺรหฺมจริยวาโส”ติ.}

\prose{\csromansymbolmark{*}อนาคามิสฺส ปุคฺคลสฺส ปญฺโจรมฺภาคิยานิ สํโยชนานิ ปหีนานิ, ปญฺจุทฺธมฺภาคิยานิ สํโยชนานิ อปฺปหีนานิ, อิโต จุตสฺส ตตฺถ อุปปนฺนสฺส กุหิํ ภาโรหรณํ, กุหิํ ทุกฺขปริญฺญาตํ, กุหิํ กิเลสปฺปหานํ, กุหิํ นิโรธสจฺฉิกิริยา, กุหิํ อกุปฺปปฏิเวโธติ, ตตฺเถว.\csromanspacer{} หญฺจิ อนาคามิสฺส ปุคฺคลสฺส ปญฺโจรมฺภาคิยานิ สํโยชนานิ ปหีนานิ, ปญฺจุทฺธมฺภาคิยานิ สํโยชนานิ อปฺปหีนานิ.\csromanspacer{} อิโต จุตสฺส ตตฺถ อุปปนฺนสฺส ตหิํ อกุปฺปปฏิเวโธ, โน จ วต เร วตฺตพฺเพ “นตฺถิ เทเวสุ พฺรหฺมจริยวาโส”ติ.}

\prose{\csromansymbolmark{*}อนาคามิสฺส ปุคฺคลสฺส ปญฺโจรมฺภาคิยานิ สํโยชนานิ ปหีนานิ, ปญฺจุทฺธมฺภาคิยานิ สํโยชนานิ อปฺปหีนานิ, อิโต จุตสฺส ตตฺถ อุปปนฺนสฺส}

\vspace{\tipitakaparskip}
\csromancenter{พฺรหฺมจริยกถา ตหิํ ผลุปฺปตฺติ ตหิํ ภาโรหรณํ ตหิํ ทุกฺขปริญฺญาตํ ตหิํ กิเลสปฺปหานํ ตหิํ นิโรธสจฺฉิกิริยา ตหิํ อกุปฺปปฏิเวโธ, เกนฏฺเฐน วเทสิ “นตฺถิ เทเวสุ พฺรหฺมจริยวาโส”ติ? หนฺท หิ อนาคามี ปุคฺคโล อิธ ภาวิเตน มคฺเคน ตตฺถ ผลํ สจฺฉิกโรตีติ\footnote{สจฺฉิกโรติ (พหูสุ)}.}
\csromanheader{2}{2. สํสนฺทนพฺรหฺมจริยกถา}
\proseitem{273}{\csromanfolio{85}อนาคามี ปุคฺคโล อิธ ภาวิเตน มคฺเคน ตตฺถ ผลํ สจฺฉิกโรตีติ, อามนฺตา.\csromanspacer{} โสตาปนฺโน ปุคฺคโล ตตฺถ ภาวิเตน มคฺเคน อิธ ผลํ สจฺฉิกโรตีติ, น เหวํ วตฺตพฺเพ.}

\prose{\csromansymbolmark{*}อนาคามี ปุคฺคโล อิธ ภาวิเตน มคฺเคน ตตฺถ ผลํ สจฺฉิกโรตีติ, อามนฺตา.\csromanspacer{} สกทาคามี ปุคฺคโล อิธ ปรินิพฺพายิปุคฺคโล\footnote{อิธปรินิพฺพายี (?)} ตตฺถ ภาวิเตน มคฺเคน อิธ ผลํ สจฺฉิกโรตีติ, น เหวํ วตฺตพฺเพ.}

\prose{\csromansymbolmark{*}โสตาปนฺโน ปุคฺคโล อิธ ภาวิเตน มคฺเคน อิธ ผลํ สจฺฉิกโรตีติ, อามนฺตา.\csromanspacer{} อนาคามี ปุคฺคโล ตตฺถ ภาวิเตน มคฺเคน ตตฺถ ผลํ สจฺฉิกโรตีติ, น เหวํ วตฺตพฺเพ.}

\prose{\csromansymbolmark{*}สกทาคามี ปุคฺคโล อิธ ปรินิพฺพายิปุคฺคโล อิธ ภาวิเตน มคฺเคน อิธ ผลํ สจฺฉิกโรตีติ, อามนฺตา.\csromanspacer{} อนาคามี ปุคฺคโล ตตฺถ ภาวิเตน มคฺเคน ตตฺถ ผลํ สจฺฉิกโรตีติ, น เหวํ วตฺตพฺเพ ฯเปฯ.}

\prose{\csromansymbolmark{*}อิธ วิหาย นิฏฺฐสฺส ปุคฺคลสฺส มคฺโค จ ภาวียติ, น จ กิเลสา ปหียนฺตีติ, อามนฺตา.\csromanspacer{} โสตาปตฺติผลสจฺฉิกิริยาย ปฏิปนฺนสฺส ปุคฺคลสฺส มคฺโค จ ภาวียติ น จ กิเลสา ปหียนฺตีติ, น เหวํ วตฺตพฺเพ ฯเปฯ.}

\prose{\csromansymbolmark{*}อิธ วิหาย นิฏฺฐสฺส ปุคฺคลสฺส มคฺโค จ ภาวียติ, น จ กิเลสา ปหียนฺตีติ, อามนฺตา.\csromanspacer{} สกทาคามิผลสจฺฉิกิริยาย ปฏิปนฺนสฺส ปุคฺคลสฺส ฯเปฯ.\csromanspacer{} อรหตฺตสจฺฉิกิริยาย ปฏิปนฺนสฺส ปุคฺคลสฺส มคฺโค จ ภาวียติ, น จ กิเลสา ปหียนฺตีติ, น เหวํ วตฺตพฺเพ ฯเปฯ.}

\prose{\csromansymbolmark{*}โสตาปตฺติผลสจฺฉิกิริยาย ปฏิปนฺนสฺส ปุคฺคลสฺส อปุพฺพํ อจริมํ มคฺโค จ ภาวียติ, กิเลสา จ ปหียนฺตีติ, อามนฺตา.\csromanspacer{} อิธ วิหาย \csromanfolio{86}นิฏฺฐสฺส ปุคฺคลสฺส อปุพฺพํ อจริมํ มคฺโค จ ภาวียติ, กิเลสา จ ปหียนฺตีติ, น เหวํ วตฺตพฺเพ.}

\prose{\csromansymbolmark{*}สกทาคามิผลสจฺฉิกิริยาย ปฏิปนฺนสฺส ปุคฺคลสฺส ฯเปฯ.\csromanspacer{} อรหตฺตสจฺฉิกิริยาย ปฏิปนฺนสฺส ปุคฺคลสฺส อปุพฺพํ อจริมํ มคฺโค จ ภาวียติ, กิเลสา จ ปหียนฺตีติ, อามนฺตา.\csromanspacer{} อิธ วิหาย นิฏฺฐสฺส ปุคฺคลสฺส อปุพฺพํ อจริมํ มคฺโค จ ภาวียติ, กิเลสา จ ปหียนฺตีติ, น เหวํ วตฺตพฺเพ.}

\prose{\csromansymbolmark{*}อนาคามี ปุคฺคโล กตกรณีโย ภาวิตภาวโน ตตฺถ อุปปชฺชตีติ, อามนฺตา.\csromanspacer{} อรหา อุปปชฺชตีติ, น เหวํ วตฺตพฺเพ.}

\prose{\csromansymbolmark{*}อรหา อุปปชฺชตีติ, อามนฺตา.\csromanspacer{} อตฺถิ อรหโต ปุนพฺภโวติ, น เหวํ วตฺตพฺเพ.}

\prose{\csromansymbolmark{*}อตฺถิ อรหโต ปุนพฺภโวติ, อามนฺตา.\csromanspacer{} อรหา ภเวน ภวํ คจฺฉติ, คติยา คติํ คจฺฉติ, สํสาเรน สํสารํ คจฺฉติ, อุปปตฺติยา อุปปตฺติํ คจฺฉตีติ, น เหวํ วตฺตพฺเพ.}

\prose{\csromansymbolmark{*}อนาคามี ปุคฺคโล กตกรณีโย ภาวิตภาวโน อโนหฏภาโร ตตฺถ อุปปชฺชตีติ, อามนฺตา.\csromanspacer{} ภาโรหรณาย ปุน มคฺคํ ภาเวตีติ, น เหวํ วตฺตพฺเพ.}

\prose{\csromansymbolmark{*}อนาคามี ปุคฺคโล กตกรณีโย ภาวิตภาวโน อปริญฺญาตทุกฺโข อปฺปหีนกิเลโส อสจฺฉิกตนิโรโธ อปฺปฏิวิทฺธากุปฺโป ตตฺถ อุปปชฺชตีติ, อามนฺตา.\csromanspacer{} อกุปฺปปฏิเวธาย ปุน มคฺคํ ภาเวตีติ, น เหวํ วตฺตพฺเพ.}

\prose{\csromansymbolmark{*}อนาคามี ปุคฺคโล กตกรณีโย ภาวิตภาวโน อโนหฏภาโร ตตฺถ อุปปชฺชติ, น จ ภาโรหรณาย ปุน มคฺคํ ภาเวตีติ, อามนฺตา.\csromanspacer{} อโนหฏภาโร จ ตตฺถ ปรินิพฺพายตีติ, น เหวํ วตฺตพฺเพ.}

\prose{\csromansymbolmark{*}อนาคามี ปุคฺคโล กตกรณีโย ภาวิตภาวโน อปริญฺญาตทุกฺโข อปฺปหีนกิเลโส อสจฺฉิกตนิโรโธ อปฺปฏิวิทฺธากุปฺโป ตตฺถ อุปปชฺชติ, น จ อกุปฺปปฏิเวธาย ปุน มคฺคํ ภาเวตีติ, อามนฺตา.\csromanspacer{} อปฺปฏิวิทฺธากุปฺโป จ ตตฺถ ปรินิพฺพายตีติ, น เหวํ วตฺตพฺเพ, ยถา มิโค สลฺเลน วิทฺโธ ทูรมฺปิ คนฺตฺวา กาลํ กโรติ, เอวเมวํ อนาคามี ปุคฺคโล อิธ ภาวิเตน มคฺเคน ตตฺถ ผลํ สจฺฉิกโรตีติ.}

\csromanheader{2}{3. พฺรหฺมจริยกถา}
\prose{\csromanfolio{87}\csromansymbolmark{*}ยถา มิโค สลฺเลน วิทฺโธ ทูรมฺปิ คนฺตฺวา สสลฺโลว กาลํ กโรติ, เอวเมวํ อนาคามี ปุคฺคโล อิธ ภาวิเตน มคฺเคน ตตฺถ สสลฺโลว ปรินิพฺพายตีติ, น เหวํ วตฺตพฺเพ ฯเปฯ.}

\nitthitam{พฺรหฺมจริยกถา นิฏฺฐิตา.}
\csromanmatikaanchor{mtk.22}
\csromantocmarkref{chapter}{3. โอธิโสกถา (ข)}{mtk.22}
\bookmark[dest={mtk.22},level=0]{3. โอธิโสกถา (ข)}
\chapterhead{3. โอธิโสกถา}
\proseitem{274}{\csromansymbolmark{*}โอธิโสธิโส กิเลเส ชหตีติ, อามนฺตา.\csromanspacer{} โสตาปตฺติผลสจฺฉิกิริยาย ปฏิปนฺโน ปุคฺคโล ทุกฺขทสฺสเนน กิํ ชหตีติ, สกฺกายทิฏฺฐิํ วิจิกิจฺฉํ สีลพฺพตปรามาสํ ตเทกฏฺเฐ จ กิเลเส เอกเทเส ชหตีติ.\csromanspacer{} เอกเทสํ โสตาปนฺโน, เอกเทสํ น โสตาปนฺโน, เอกเทสํ โสตาปตฺติผลปฺปตฺโต\footnote{โสตาปตฺติผลํ ปตฺโต (?)} ปฏิลทฺโธ อธิคโต สจฺฉิกโต อุปสมฺปชฺช วิหรติ, กาเยน ผุสิตฺวา วิหรติ, เอกเทสํ น กาเยน ผุสิตฺวา วิหรติ, เอกเทสํ สตฺตกฺขตฺตุปรโม, โกลํโกโล, เอกพีชี พุทฺเธ อเวจฺจปฺปสาเทน สมนฺนาคโต, ธมฺเม ฯเปฯ สํเฆ ฯเปฯ อริยกนฺเตหิ สีเลหิ สมนฺนาคโต, เอกเทสํ อริยกนฺเตหิ สีเลหิ น สมนฺนาคโตติ, น เหวํ วตฺตพฺเพ ฯเปฯ.}

\prose{\csromansymbolmark{*}สมุทยทสฺสเนน กิํ ชหตีติ, สกฺกายทิฏฺฐิํ ชหติ วิจิกิจฺฉํ สีลพฺพตปรามาสํ ตเทกฏฺเฐ จ กิเลเส เอกเทเส ชหตีติ.\csromanspacer{} เอกเทสํ โสตาปนฺโน, เอกเทสํ น โสตาปนฺโน ฯเปฯ.\csromanspacer{} เอกเทสํ อริยกนฺเตหิ สีเลหิ สมนฺนาคโต, เอกเทสํ อริยกนฺเตหิ สีเลหิ น สมนฺนาคโตติ, น เหวํ วตฺตพฺเพ ฯเปฯ.}

\prose{\csromansymbolmark{*}นิโรธทสฺสเนน กิํ ชหตีติ, วิจิกิจฺฉํ สีลพฺพตปรามาสํ ตเทกฏฺเฐ จ กิเลเส เอกเทเส ชหตีติ.\csromanspacer{} เอกเทสํ โสตาปนฺโน, เอกเทสํ น โสตาปนฺโน ฯเปฯ.\csromanspacer{} เอกเทสํ อริยกนฺเตหิ สีเลหิ สมนฺนาคโต, เอกเทสํ อริยกนฺเตหิ สีเลหิ น สมนฺนาคโตติ, น เหวํ วตฺตพฺเพ ฯเปฯ.}

\prose{\csromansymbolmark{*}มคฺคทสฺสเนน กิํ ชหตีติ, สีลพฺพตปรามาสํ ตเทกฏฺเฐ จ กิเลเส ชหตีติ.\csromanspacer{} เอกเทสํ โสตาปนฺโน, เอกเทสํ น โสตาปนฺโน, \csromanfolio{88}เอกเทสํ โสตาปตฺติผลปฺปตฺโต ปฏิลทฺโธ อธิคโต สจฺฉิกโต อุปสมฺปชฺช วิหรติ, กาเยน ผุสิตฺวา วิหรติ, เอกเทสํ น กาเยน ผุสิตฺวา วิหรติ, เอกเทสํ สตฺตกฺขตฺตุปรโม, โกลํโกโล, เอกพีชี พุทฺเธ อเวจฺจปฺปสาเทน สมนฺนาคโต, ธมฺเม ฯเปฯ สํเฆ ฯเปฯ อริยกนฺเตหิ สีเลหิ สมนฺนาคโต, เอกเทสํ อริยกนฺเตหิ สีเลหิ น สมนฺนาคโตติ, น เหวํ วตฺตพฺเพ ฯเปฯ.}

\proseitem{275}{\csromansymbolmark{*}สกทาคามิผลสจฺฉิกิริยาย ปฏิปนฺโน ปุคฺคโล ทุกฺขทสฺสเนน กิํ ชหตีติ, โอฬาริกํ กามราคํ ชหติ, โอฬาริกํ พฺยาปาทํ ตเทกฏฺเฐ จ กิเลเส เอกเทเส ชหตีติ.\csromanspacer{} เอกเทสํ สกทาคามี, เอกเทสํ น สกทาคามี, เอกเทสํ สกทาคามิผลปฺปตฺโต\footnote{สกทาคามิผลํ ปตฺโต (?)} ปฏิลทฺโธ อธิคโต สจฺฉิกโต อุปสมฺปชฺช วิหรติ, กาเยน ผุสิตฺวา วิหรติ, เอกเทสํ น กาเยน ผุสิตฺวา วิหรตีติ, น เหวํ วตฺตพฺเพ ฯเปฯ.}

\prose{\csromansymbolmark{*}สมุทยทสฺสเนน กิํ ชหตีติ, โอฬาริกํ กามราคํ ชหติ, โอฬาริกํ พฺยาปาทํ ตเทกฏฺเฐ จ กิเลเส เอกเทเส ชหตีติ.\csromanspacer{} เอกเทสํ สกทาคามี, เอกเทสํ น สกทาคามี, เอกเทสํ สกทาคามิผลปฺปตฺโต ปฏิลทฺโธ อธิคโต สจฺฉิกโต อุปสมฺปชฺช วิหรติ, กาเยน ผุสิตฺวา วิหรติ, เอกเทสํ น กาเยน ผุสิตฺวา วิหรตีติ, น เหวํ วตฺตพฺเพ ฯเปฯ.}

\prose{\csromansymbolmark{*}นิโรธทสฺสเนน กิํ ชหตีติ, โอฬาริกํ พฺยาปาทํ ชหติ, ตเทกฏฺเฐ จ กิเลเส เอกเทเส ชหตีติ.\csromanspacer{} เอกเทสํ สกทาคามี, เอกเทสํ น สกทาคามี, เอกเทสํ สกทาคามิผลปฺปตฺโต ปฏิลทฺโธ อธิคโต สจฺฉิกโต อุปสมฺปชฺช วิหรติ, กาเยน ผุสิตฺวา วิหรติ, เอกเทสํ น กาเยน ผุสิตฺวา วิหรตีติ, น เหวํ วตฺตพฺเพ ฯเปฯ.}

\prose{\csromansymbolmark{*}มคฺคทสฺสเนน กิํ ชหตีติ, โอฬาริกํ พฺยาปาทํ ชหติ, ตเทกฏฺเฐ จ กิเลเส ชหตีติ.\csromanspacer{} เอกเทสํ สกทาคามี, เอกเทสํ น สกทาคามี, เอกเทสํ สกทาคามิผลปฺปตฺโต ปฏิลทฺโธ อธิคโต สจฺฉิกโต อุปสมฺปชฺช วิหรติ, กาเยน ผุสิตฺวา วิหรติ, เอกเทสํ น กาเยน ผุสิตฺวา วิหรตีติ, น เหวํ วตฺตพฺเพ ฯเปฯ.}

\csromanheader{2}{3. พฺรหฺมจริยกถา}
\proseitem{276}{\csromanfolio{89}\csromansymbolmark{*}อนาคามิผลสจฺฉิกิริยาย ปฏิปนฺโน ปุคฺคโล ทุกฺขทสฺสเนน กิํ ชหตีติ, อณุสหคตํ กามราคํ ชหติ, อณุสหคตํ พฺยาปาทํ ตเทกฏฺเฐ จ กิเลเส เอกเทเส ชหตีติ.\csromanspacer{} เอกเทสํ อนาคามี, เอกเทสํ น อนาคามี, เอกเทสํ อนาคามิผลปฺปตฺโต\footnote{อนาคามิผลํ ปตฺโต (?)} ปฏิลทฺโธ อธิคโต สจฺฉิกโต อุปสมฺปชฺช วิหรติ, กาเยน ผุสิตฺวา วิหรติ, เอกเทสํ น กาเยน ผุสิตฺวา วิหรติ, เอกเทสํ อนฺตราปรินิพฺพายี, อุปหจฺจปรินิพฺพายี, อสงฺขารปรินิพฺพายี, สสงฺขารปรินิพฺพายี, อุทฺธํโสโต อกนิฏฺฐคามี, เอกเทสํ น อุทฺธํโสโต อกนิฏฺฐคามีติ, น เหวํ วตฺตพฺเพ ฯเปฯ.}

\prose{\csromansymbolmark{*}สมุทยทสฺสเนน กิํ ชหตีติ, อณุสหคตํ กามราคํ ชหติ, อณุสหคตํ พฺยาปาทํ ตเทกฏฺเฐ จ กิเลเส เอกเทเส ชหตีติ.\csromanspacer{} เอกเทสํ อนาคามี, เอกเทสํ น อนาคามี ฯเปฯ.\csromanspacer{} เอกเทสํ อุทฺธํโสโต อกนิฏฺฐคามี, เอกเทสํ น อุทฺธํโสโต อกนิฏฺฐคามีติ, น เหวํ วตฺตพฺเพ ฯเปฯ.}

\prose{\csromansymbolmark{*}นิโรธทสฺสเนน กิํ ชหตีติ, อณุสหคตํ พฺยาปาทํ ชหติ, ตเทกฏฺเฐ จ กิเลเส เอกเทเส ชหตีติ.\csromanspacer{} เอกเทสํ อนาคามี, เอกเทสํ น อนาคามี ฯเปฯ.\csromanspacer{} เอกเทสํ อุทฺธํโสโต อกนิฏฺฐคามี เอกเทสํ น อุทฺธํโสโต อกนิฏฺฐคามีติ, น เหวํ วตฺตพฺเพ ฯเปฯ.}

\prose{\csromansymbolmark{*}มคฺคทสฺสเนน กิํ ชหตีติ, ตเทกฏฺเฐ จ กิเลเส ชหตีติ.\csromanspacer{} เอกเทสํ อนาคามี, เอกเทสํ น อนาคามี, เอกเทสํ อนาคามิผลปฺปตฺโต ปฏิลทฺโธ อธิคโต สจฺฉิกโต อุปสมฺปชฺช วิหรติ, กาเยน ผุสิตฺวา วิหรติ, เอกเทสํ น กาเยน ผุสิตฺวา วิหรติ, เอกเทสํ อนฺตราปรินิพฺพายี, อุปหจฺจปรินิพฺพายี, อสงฺขารปรินิพฺพายี, สสงฺขารปรินิพฺพายี, อุทฺธํโสโต อกนิฏฺฐคามี, เอกเทสํ น อุทฺธํโสโต อกนิฏฺฐคามีติ, น เหวํ วตฺตพฺเพ ฯเปฯ.}

\proseitem{277}{\csromansymbolmark{*}อรหตฺตสจฺฉิกิริยาย ปฏิปนฺโน ปุคฺคโล ทุกฺขทสฺสเนน กิํ ชหตีติ, รูปราคํ อรูปราคํ มานํ อุทฺธจฺจํ อวิชฺชํ ตเทกฏฺเฐ จ กิเลเส เอกเทเส ชหตีติ.\csromanspacer{} เอกเทสํ อรหา, เอกเทสํ น อรหา, เอกเทสํ \csromanfolio{90}อรหตฺตปฺปตฺโต\footnote{อรหตฺตํ ปตฺโต (?)} ปฏิลทฺโธ อธิคโต สจฺฉิกโต อุปสมฺปชฺช วิหรติ, กาเยน ผุสิตฺวา วิหรติ, เอกเทสํ น กาเยน ผุสิตฺวา วิหรติ, เอกเทสํ วีตราโค วีตโทโส วีตโมโห กตกรณีโย โอหิตภาโร อนุปฺปตฺตสทตฺโถ ปริกฺขีณภวสํโยชโน สมฺมทญฺญาวิมุตฺโต อุกฺขิตฺตปลิโฆ สํกิณฺณปริโข อพฺพูเฬฺหสิโก นิรคฺคโฬ อริโย ปนฺนทฺธโช ปนฺนภาโร วิสญฺญุตฺโต สุวิชิตวิชโย, ทุกฺขํ ตสฺส ปริญฺญาตํ, สมุทโย ปหีโน, นิโรโธ สจฺฉิกโต, มคฺโค ภาวิโต, อภิญฺเญยฺยํ อภิญฺญาตํ, ปริญฺเญยฺยํ ปริญฺญาตํ, ปหาตพฺพํ ปหีนํ, ภาเวตพฺพํ ภาวิตํ, สจฺฉิกาตพฺพํ สจฺฉิกตํ, (เอกเทสํ สจฺฉิกาตพฺพํ สจฺฉิกตํ)\footnote{( ) (?)} เอกเทสํ สจฺฉิกาตพฺพํ น สจฺฉิกตนฺติ, น เหวํ วตฺตพฺเพ ฯเปฯ.}

\prose{\csromansymbolmark{*}สมุทยทสฺสเนน กิํ ชหตีติ, รูปราคํ อรูปราคํ ชหติ, มานํ อุทฺธจฺจํ อวิชฺชํ ตเทกฏฺเฐ จ กิเลเส เอกเทเส ชหตีติ.\csromanspacer{} เอกเทสํ อรหา, เอกเทสํ น อรหา ฯเปฯ.\csromanspacer{} เอกเทสํ สจฺฉิกาตพฺพํ สจฺฉิกตํ, เอกเทสํ สจฺฉิกาตพฺพํ น สจฺฉิกตนฺติ, น เหวํ วตฺตพฺเพ ฯเปฯ.}

\prose{\csromansymbolmark{*}นิโรธทสฺสเนน กิํ ชหตีติ, มานํ ชหติ, อุทฺธจฺจํ อวิชฺชํ ตเทกฏฺเฐ จ กิเลเส เอกเทเส ชหตีติ.\csromanspacer{} เอกเทสํ อรหา, เอกเทสํ น อรหา ฯเปฯ.\csromanspacer{} เอกเทสํ สจฺฉิกาตพฺพํ สจฺฉิกตํ, เอกเทสํ สจฺฉิกาตพฺพํ น สจฺฉิกตนฺติ, น เหวํ วตฺตพฺเพ ฯเปฯ.}

\prose{\csromansymbolmark{*}มคฺคทสฺสเนน กิํ ชหตีติ, อุทฺธจฺจํ อวิชฺชํ ตเทกฏฺเฐ จ กิเลเส ชหตีติ.\csromanspacer{} เอกเทสํ อรหา, เอกเทสํ น อรหา, เอกเทสํ อรหตฺตปฺปตฺโต ปฏิลทฺโธ อธิคโต สจฺฉิกโต อุปสมฺปชฺช วิหรติ, กาเยน ผุสิตฺวา วิหรติ, เอกเทสํ น กาเยน ผุสิตฺวา วิหรติ, เอกเทสํ วีตราโค วีตโทโส วีตโมโห กตกรณีโย โอหิตภาโร อนุปฺปตฺตสทตฺโถ ปริกฺขีณภวสํโยชโน สมฺมทญฺญาวิมุตฺโต อุกฺขิตฺตปลิโฆ สํกิณฺณปริโข อพฺพูเฬฺหสิโก นิรคฺคโฬ อริโย ปนฺนทฺธโช ปนฺนภาโร วิสญฺญุตฺโต สุวิชิตวิชโย, ทุกฺขํ ตสฺส ปริญฺญาตํ, สมุทโย ปหีโน, นิโรโธ สจฺฉิกโต, มคฺโค ภาวิโต, อภิญฺเญยฺยํ อภิญฺญาตํ, ปริญฺเญยฺยํ ปริญฺญาตํ, ปหาตพฺพํ ปหีนํ, ภาเวตพฺพํ}

\vspace{\tipitakaparskip}
\csromancenter{พฺรหฺมจริยกถา ภาวิตํ, สจฺฉิกาตพฺพํ สจฺฉิกตํ, เอกเทสํ สจฺฉิกาตพฺพํ สจฺฉิกตํ, เอกเทสํ สจฺฉิกาตพฺพํ น สจฺฉิกตนฺติ, น เหวํ วตฺตพฺเพ ฯเปฯ.}
\proseitem{278}{\csromanfolio{91}\csromansymbolmark{+}น วตฺตพฺพํ “โอธิโสธิโส กิเลเส ชหตี”ติ, อามนฺตา.\csromanspacer{} นนุ วุตฺตํ ภควตา\csromandash{}}

\csromangathameasure{“อนุปุพฺเพน เมธาวี, โถกํ โถกํ ขเณ ขเณ. \\ กมฺมาโร รชตสฺเสว, นิทฺธเม มลมตฺตโน”ติ.}
\csromangathagroup{\csromangathastanza{“อนุปุพฺเพน เมธาวี, โถกํ โถกํ ขเณ ขเณ. \\ กมฺมาโร รชตสฺเสว, นิทฺธเม มลมตฺตโน”ติ\footnote{ขุ 1. 48 ธมฺมปเท.}.}}

\prose{อตฺเถว สุตฺตนฺโตติ, อามนฺตา.\csromanspacer{} เตน หิ วตฺตพฺพํ “โอธิโสธิโส กิเลเส ชหตี”ติ.}

\prose{\csromansymbolmark{*}โอธิโสธิโส กิเลเส ชหตีติ, อามนฺตา.\csromanspacer{} นนุ วุตฺตํ ภควตา\csromandash{}}

\csromangathameasure{“สหาวสฺส ทสฺสนสมฺปทาย, \\ ตยสฺสุ ธมฺมา ชหิตา ภวนฺติ. \\ สกฺกายทิฏฺฐี วิจิกิจฺฉิตญฺจ, \\ สีลพฺพตํ วาปิ ยทตฺถิ กิญฺจิ. \\ จตูหปาเยหิ จ วิปฺปมุตฺโต, \\ ฉจฺจาภิฐานานิ อภพฺพ กาตุนฺ”ติ.}
\csromangathagroup{\csromangathastanza{“สหาวสฺส ทสฺสนสมฺปทาย, \\ ตยสฺสุ ธมฺมา ชหิตา ภวนฺติ. \\ สกฺกายทิฏฺฐี วิจิกิจฺฉิตญฺจ, \\ สีลพฺพตํ วาปิ ยทตฺถิ กิญฺจิ. \\ จตูหปาเยหิ จ วิปฺปมุตฺโต, \\ ฉจฺจาภิฐานานิ อภพฺพ\footnote{อภพฺโพ (สี, สฺยา)} กาตุนฺ”ติ\footnote{ขุ 1. 6, 313 ปิฏฺเฐสุ.}.}}

\prose{อตฺเถว สุตฺตนฺโตติ, อามนฺตา.\csromanspacer{} เตน หิ น วตฺตพฺพํ “โอธิโสธิโส กิเลเส ชหตี”ติ.}

\prose{\csromansymbolmark{*}โอธิโสธิโส กิเลเส ชหตีติ, อามนฺตา.\csromanspacer{} นนุ วุตฺตํ ภควตา “ยสฺมิํ ภิกฺขเว สมเย อริยสาวกสฺส วิรชํ วีตมลํ ธมฺมจกฺขุํ อุทปาทิ ‘ยํ กิญฺจิ สมุทยธมฺมํ, สพฺพํ ตํ นิโรธธมฺมนฺ’ติ, สห ทสฺสนุปฺปาทา ภิกฺขเว อริยสาวกสฺส ตีณิ สํโยชนานิ ปหียนฺติ, สกฺกายทิฏฺฐิ วิจิกิจฺฉา สีลพฺพตปรามาโส”ติ\footnote{อํ 1. 244 ปิฏฺเฐ.} อตฺเถว สุตฺตนฺโตติ, อามนฺตา.\csromanspacer{} เตน หิ น วตฺตพฺพํ”โอธิโสธิโส กิเลเส ชหตี”ติ.}

\nitthitam{โอธิโสกถา นิฏฺฐิตา.}
\csromanmatikaanchor{mtk.23}
\csromantocmarkref{chapter}{4. ชหติกถา}{mtk.23}
\bookmark[dest={mtk.23},level=0]{4. ชหติกถา}
\chapterheadpageread{4. ชหติกถา}
\csromanheader{2}{1. นสุตฺตาหรณกถา}
\proseitem{279}{\csromanfolio{92}\csromansymbolmark{*}ชหติ ปุถุชฺชโน กามราคพฺยาปาทนฺติ, อามนฺตา.\csromanspacer{} อจฺจนฺตํ ชหติ.\csromanspacer{} อนวเสสํ ชหติ.\csromanspacer{} อปฺปฏิสนฺธิยํ ชหติ.\csromanspacer{} สมูลํ ชหติ.\csromanspacer{} สตณฺหํ ชหติ.\csromanspacer{} สานุสยํ ชหติ.\csromanspacer{} อริเยน ญาเณน ชหติ.\csromanspacer{} อริเยน มคฺเคน ชหติ.\csromanspacer{} อกุปฺปํ ปฏิวิชฺฌนฺโต ชหติ.\csromanspacer{} อนาคามิผลํ สจฺฉิกโรนฺโต ชหตีติ, น เหวํ วตฺตพฺเพ ฯเปฯ.}

\prose{\csromansymbolmark{*}วิกฺขมฺเภติ ปุถุชฺชโน กามราคพฺยาปาทนฺติ, อามนฺตา.\csromanspacer{} อจฺจนฺตํ วิกฺขมฺเภติ.\csromanspacer{} อนวเสสํ วิกฺขมฺเภติ.\csromanspacer{} อปฺปฏิสนฺธิยํ วิกฺขมฺเภติ.\csromanspacer{} สมูลํ วิกฺขมฺเภติ.\csromanspacer{} สตณฺหํ วิกฺขมฺเภติ.\csromanspacer{} สานุสยํ วิกฺขมฺเภติ.\csromanspacer{} อริเยน ญาเณน วิกฺขมฺเภติ.\csromanspacer{} อริเยน มคฺเคน วิกฺขมฺเภติ.\csromanspacer{} อกุปฺปํ ปฏิวิชฺฌนฺโต วิกฺขมฺเภติ.\csromanspacer{} อนาคามิผลํ สจฺฉิกโรนฺโต วิกฺขมฺเภตีติ, น เหวํ วตฺตพฺเพ ฯเปฯ.}

\prose{\csromansymbolmark{*}ชหติ อนาคามิผลสจฺฉิกิริยาย ปฏิปนฺโน ปุคฺคโล กามราคพฺยาปาทํ, โส จ อจฺจนฺตํ ชหติ.\csromanspacer{} อนวเสสํ ชหติ ฯเปฯ.\csromanspacer{} อนาคามิผลํ สจฺฉิกโรนฺโต ชหตีติ, อามนฺตา.\csromanspacer{} ชหติ ปุถุชฺชโน กามราคพฺยาปาทํ, โส จ อจฺจนฺตํ ชหติ, อนวเสสํ ชหติ ฯเปฯ.\csromanspacer{} อนาคามิผลํ สจฺฉิกโรนฺโต ชหตีติ.\csromanspacer{} น เหวํ วตฺตพฺเพ ฯเปฯ.}

\prose{\csromansymbolmark{*}วิกฺขมฺเภติ อนาคามิผลสจฺฉิกิริยาย ปฏิปนฺโน ปุคฺคโล กามราคพฺยาปาทํ, โส จ อจฺจนฺตํ วิกฺขมฺเภติ.\csromanspacer{} อนวเสสํ วิกฺขมฺเภติ ฯเปฯ.\csromanspacer{} อนาคามิผลํ สจฺฉิกโรนฺโต วิกฺขมฺเภตีติ, อามนฺตา.\csromanspacer{} วิกฺขมฺเภติ ปุถุชฺชโน กามราคพฺยาปาทํ, โส จ อจฺจนฺตํ วิกฺขมฺเภติ.\csromanspacer{} อนวเสสํ วิกฺขมฺเภติ ฯเปฯ.\csromanspacer{} อนาคามิผลํ สจฺฉิกโรนฺโต วิกฺขมฺเภตีติ, น เหวํ วตฺตพฺเพ ฯเปฯ.}

\prose{\csromansymbolmark{*}ชหติ ปุถุชฺชโน กามราคพฺยาปาทํ, โส จ น อจฺจนฺตํ ชหติ.\csromanspacer{} น อนวเสสํ ชหติ ฯเปฯ.\csromanspacer{} น อนาคามิผลํ สจฺฉิกโรนฺโต ชหตีติ, อามนฺตา.\csromanspacer{} ชหติ อนาคามิผลสจฺฉิกิริยาย ปฏิปนฺโน ปุคฺคโล กามราคพฺยาปาทํ, โส จ น อจฺจนฺตํ ชหติ ฯเปฯ.\csromanspacer{} น อนาคามิผลํ สจฺฉิกโรนฺโต ชหตีติ, น เหวํ วตฺตพฺเพ ฯเปฯ.}

\csromanheader{2}{4. ชหติกถา}
\prose{\csromanfolio{93}\csromansymbolmark{*}วิกฺขมฺเภติ ปุถุชฺชโน กามราคพฺยาปาทํ, โส จ น อจฺจนฺตํ วิกฺขมฺเภติ.\csromanspacer{} น อนวเสสํ วิกฺขมฺเภติ ฯเปฯ.\csromanspacer{} น อนาคามิผลํ สจฺฉิกโรนฺโต วิกฺขมฺเภตีติ, อามนฺตา.\csromanspacer{} วิกฺขมฺเภติ อนาคามิผลสจฺฉิกิริยาย ปฏิปนฺโน ปุคฺคโล กามราคพฺยาปาทํ, โส จ น อจฺจนฺตํ วิกฺขมฺเภติ.\csromanspacer{} น อนวเสสํ วิกฺขมฺเภติ ฯเปฯ.\csromanspacer{} น อนาคามิผลํ สจฺฉิกโรนฺโต วิกฺขมฺเภตีติ, น เหวํ วตฺตพฺเพ ฯเปฯ.}

\prose{\csromansymbolmark{*}ชหติ ปุถุชฺชโน กามราคพฺยาปาทนฺติ, อามนฺตา.\csromanspacer{} กตเมน มคฺเคนาติ, รูปาวจเรน มคฺเคนาติ.\csromanspacer{} รูปาวจโร มคฺโค นิยฺยานิโก ขยคามี โพธคามี อปจยคามี อนาสโว อสํโยชนิโย อคนฺถนิโย อโนฆนิโย อโยคนิโย อนีวรณิโย อปรามฏฺโฐ อนุปาทานิโย อสํกิเลสิโยติ, น เหวํ วตฺตพฺเพ.\csromanspacer{} นนุ รูปาวจโร มคฺโค อนิยฺยานิโก น ขยคามี น โพธคามี น อปจยคามี สาสโว สํโยชนิโย ฯเปฯ สํกิเลสิโยติ, อามนฺตา.\csromanspacer{} หญฺจิ รูปาวจโร มคฺโค อนิยฺยานิโก น ขยคามี ฯเปฯ สํกิเลสิโย, โน จ วต เร วตฺตพฺเพ “ชหติ ปุถุชฺชโน รูปาวจเรน มคฺเคน กามราคพฺยาปาทนฺ”ติ.}

\prose{\csromansymbolmark{*}ชหติ อนาคามิผลสจฺฉิกิริยาย ปฏิปนฺโน ปุคฺคโล อนาคามิมคฺเคน กามราคพฺยาปาทํ, โส จ มคฺโค นิยฺยานิโก ขยคามี โพธคามี อปจยคามี อนาสโว ฯเปฯ อสํกิเลสิโยติ, อามนฺตา.\csromanspacer{} ชหติ ปุถุชฺชโน รูปาวจเรน มคฺเคน กามราคพฺยาปาทํ, โส จ มคฺโค นิยฺยานิโก ขยคามี โพธคามี อปจยคามี อนาสโว ฯเปฯ อสํกิเลสิโยติ, น เหวํ วตฺตพฺเพ ฯเปฯ.}

\prose{\csromansymbolmark{*}ชหติ ปุถุชฺชโน รูปาวจเรน มคฺเคน กามราคพฺยาปาทํ, โส จ มคฺโค อนิยฺยานิโก น ขยคามี น โพธคามี น อปจยคามี สาสโว ฯเปฯ สํกิเลสิโยติ, อามนฺตา.\csromanspacer{} ชหติ อนาคามิผลสจฺฉิกิริยาย ปฏิปนฺโน ปุคฺคโล อนาคามิมคฺเคน กามราคพฺยาปาทํ, โส จ มคฺโค อนิยฺยานิโก น ขยคามี น โพธคามี น อปจยคามี สาสโว ฯเปฯ สํกิเลสิโยติ, น เหวํ วตฺตพฺเพ ฯเปฯ.}

\proseitem{280}{\csromansymbolmark{*}ปุถุชฺชโน กาเมสุ วีตราโค สห ธมฺมาภิสมยา อนาคามิผเล สณฺฐาตีติ, อามนฺตา.\csromanspacer{} อรหตฺเต สณฺฐาตีติ, น เหวํ วตฺตพฺเพ ฯเปฯ.}

\prose{\csromanfolio{94}\csromansymbolmark{*}ปุถุชฺชโน กาเมสุ วีตราโค สห ธมฺมาภิสมยา อนาคามิผเล สณฺฐาตีติ, อามนฺตา.\csromanspacer{} อปุพฺพํ อจริมํ ตโย มคฺเค ภาเวตีติ, น เหวํ วตฺตพฺเพ ฯเปฯ. \csromansymbolmark{*}อปุพฺพํ อจริมํ ตโย มคฺเค ภาเวตีติ, อามนฺตา.\csromanspacer{} อปุพฺพํ อจริมํ ตีณิ สามญฺญผลานิ สจฺฉิกโรตีติ, น เหวํ วตฺตพฺเพ ฯเปฯ.}

\prose{\csromansymbolmark{*}อปุพฺพํ อจริมํ ตีณิ สามญฺญผลานิ สจฺฉิกโรตีติ, อามนฺตา.\csromanspacer{} ติณฺณํ ผสฺสานํ ติสฺสนฺนํ เวทนานํ ติสฺสนฺนํ สญฺญานํ ติสฺสนฺนํ เจตนานํ ติณฺณํ จิตฺตานํ ติสฺสนฺนํ สทฺธานํ ติณฺณํ วีริยานํ ติสฺสนฺนํ สตีนํ ติณฺณํ สมาธีนํ ติสฺสนฺนํ ปญฺญานํ สโมธานํ โหตีติ, น เหวํ วตฺตพฺเพ ฯเปฯ.}

\prose{\csromansymbolmark{*}ปุถุชฺชโน กาเมสุ วีตราโค สห ธมฺมาภิสมยา อนาคามิผเล สณฺฐาตีติ, อามนฺตา.\csromanspacer{} โสตาปตฺติมคฺเคนาติ, น เหวํ วตฺตพฺเพ ฯเปฯ.}

\prose{สกทาคามิมคฺเคนาติ, น เหวํ วตฺตพฺเพ.\csromanspacer{} กตเมน มคฺเคนาติ, อนาคามิมคฺเคนาติ.\csromanspacer{} อนาคามิมคฺเคน สกฺกายทิฏฺฐิํ วิจิกิจฺฉํ สีลพฺพตปรามาสํ ชหตีติ, น เหวํ วตฺตพฺเพ ฯเปฯ.}

\csromanheader{2}{2. สุตฺตาหรณกถา}
\proseitem{281}{\csromansymbolmark{*}อนาคามิมคฺเคน สกฺกายทิฏฺฐิํ วิจิกิจฺฉํ สีลพฺพตปรามาสํ ชหตีติ, อามนฺตา.\csromanspacer{} นนุ ติณฺณํ สํโยชนานํ ปหานา โสตาปตฺติผลํ วุตฺตํ ภควตาติ, อามนฺตา.\csromanspacer{} หญฺจิ ติณฺณํ สํโยชนานํ ปหานา โสตาปตฺติผลํ วุตฺตํ ภควตา, โน จ วต เร วตฺตพฺเพ “อนาคามิมคฺเคน สกฺกายทิฏฺฐิํ วิจิกิจฺฉํ สีลพฺพตปรามาสํ ชหตี”ติ.\csromanspacer{} อนาคามิมคฺเคน โอฬาริกํ กามราคํ โอฬาริกํ พฺยาปาทํ ชหตีติ, น เหวํ วตฺตพฺเพ ฯเปฯ.}

\prose{\csromansymbolmark{*}อนาคามิมคฺเคน โอฬาริกํ กามราคํ โอฬาริกํ พฺยาปาทํ ชหตีติ, อามนฺตา.\csromanspacer{} นนุ กามราคพฺยาปาทานํ ตนุภาวา สกทาคามิผลํ วุตฺตํ ภควตาติ, อามนฺตา.\csromanspacer{} หญฺจิ กามราคพฺยาปาทานํ ตนุภาวา สกทาคามิผลํ วุตฺตํ ภควตา, โน จ วต เร วตฺตพฺเพ “อนาคามิมคฺเคน โอฬาริกํ กามราคํ โอฬาริกํ พฺยาปาทํ ชหตี”ติ.}

\prose{\csromansymbolmark{*}ปุถุชฺชโน กาเมสุ วีตราโค สห ธมฺมาภิสมยา อนาคามิผเล สณฺฐาตีติ, อามนฺตา.\csromanspacer{} เย เกจิ ธมฺมํ อภิสเมนฺติ, สพฺเพ เต สห ธมฺมาภิสมยา อนาคามิผเล สณฺฐหนฺตีติ, น เหวํ วตฺตพฺเพ ฯเปฯ.}

\csromanheader{2}{4. ชหติกถา}
\prose{\csromanfolio{95}\csromansymbolmark{+}น วตฺตพฺพํ “ชหติ ปุถุชฺชโน กามราคพฺยาปาทนฺ”ติ, อามนฺตา.\csromanspacer{} นนุ วุตฺตํ ภควตา\csromandash{}}

\csromangathameasure{“อเหสุํ เต อตีตํเส, ฉ สตฺถาโร ยสสฺสิโน. \\ นิรามคนฺธา กรุเณธิมุตฺตา, กามสํโยชนาติคา. \\ กามราคํ วิราเชตฺวา, พฺรหฺมโลกูปคา อหุ. \\ อเหสุํ สาวกา เตสํ, อเนกานิ สตานิปิ. \\ นิรามคนฺธา กรุเณธิมุตฺตา, กามสํโยชนาติคา. \\ กามราคํ วิราเชตฺวา, พฺรหฺมโลกูปคา อหู”ติ.}
\csromangathagroup{\csromangathastanza{“อเหสุํ เต\footnote{อหิํสกา (อํ 2. 327 ปิฏฺเฐ)} อตีตํเส, ฉ สตฺถาโร ยสสฺสิโน. \\ นิรามคนฺธา กรุเณธิมุตฺตา\footnote{อหิํสกา (อํ 2. 327 ปิฏฺเฐ)}, กามสํโยชนาติคา.}\csromangathastanza{กามราคํ วิราเชตฺวา, พฺรหฺมโลกูปคา อหุ. \\ อเหสุํ สาวกา เตสํ, อเนกานิ สตานิปิ.}\csromangathastanza{นิรามคนฺธา กรุเณธิมุตฺตา, กามสํโยชนาติคา. \\ กามราคํ วิราเชตฺวา, พฺรหฺมโลกูปคา อหู”ติ\footnote{อํ 2. 327 ปิฏฺเฐ.}.}}

\prose{อตฺเถว สุตฺตนฺโตติ, อมนฺตา.\csromanspacer{} เตน หิ ชหติ ปุถุชฺชโน กามราคพฺยาปาทนฺติ.}

\prose{\csromansymbolmark{*}ชหติ ปุถุชฺชโน กามราคพฺยาปาทนฺติ, อามนฺตา.\csromanspacer{} นนุ วุตฺตํ ภควตา “โส หิ นาม ภิกฺขเว สุเนตฺโต สตฺถา เอวํ ทีฆายุโก สมาโน เอวํ จิรฏฺฐิติโก อปริมุตฺโต อโหสิ ชาติยา ชราย มรเณน โสเกหิ ปริเทเวหิ ทุกฺเขหิ โทมนสฺเสหิ อุปายาเสหิ อปริมุตฺโต ทุกฺขสฺมาติ วทามิ.\csromanspacer{} ตํ กิสฺส เหตุ, จตุนฺนํ ธมฺมานํ อนนุโพธา อปฺปฏิเวธา.\csromanspacer{} กตเมสํ จตุนฺนํ, อริยสฺส สีลสฺส อนนุโพธา อปฺปฏิเวธา.\csromanspacer{} อริยสฺส สมาธิสฺส.\csromanspacer{} อริยาย ปญฺญาย.\csromanspacer{} อริยาย วิมุตฺติยา อนนุโพธา อปฺปฏิเวธา.\csromanspacer{} ตยิทํ ภิกฺขเว อริยํ สีลํ อนุพุทฺธํ ปฏิวิทฺธํ, อริโย สมาธิ อนุพุทฺโธ ปฏิวิทฺโธ, อริยา ปญฺญา อนุพุทฺธา ปฏิวิทฺธา, อริยา วิมุตฺติ อนุพุทฺธา ปฏิวิทฺธา, อุจฺฉินฺนา ภวตณฺหา, ขีณา ภวเนตฺติ, นตฺถิ ทานิ ปุนพฺภโวติ.}

\csromangathameasure{สีลํ สมาธิ ปญฺญา จ, วิมุตฺติ จ อนุตฺตรา. \\ อนุพุทฺธา อิเม ธมฺมา, โคตเมน ยสสฺสินา. \\ อิติ พุทฺโธ อภิญฺญาย, ธมฺมมกฺขาสิ ภิกฺขุนํ. \\ ทุกฺขสฺสนฺตกโร สตฺถา, จกฺขุมา ปรินิพฺพุโต”ติ.}
\csromangathagroup{\csromangathastanza{สีลํ สมาธิ ปญฺญา จ, วิมุตฺติ จ อนุตฺตรา. \\ อนุพุทฺธา อิเม ธมฺมา, โคตเมน ยสสฺสินา.}\csromangathastanza{อิติ พุทฺโธ อภิญฺญาย, ธมฺมมกฺขาสิ ภิกฺขุนํ. \\ ทุกฺขสฺสนฺตกโร สตฺถา, จกฺขุมา ปรินิพฺพุโต”ติ\footnote{อํ 2. 477 ปิฏฺเฐ.}.}}

\prose{อตฺเถว สุตฺตนฺโตติ, อามนฺตา.\csromanspacer{} เตน หิ น วตฺตพฺพํ “ชหติ ปุถุชฺชโน กามราคพฺยาปาทนฺ”ติ.}

\nitthitam{ชหติกถา นิฏฺฐิตา.}
\csromanmatikaanchor{mtk.24}
\csromantocmarkref{chapter}{5. สพฺพมตฺถีติกถา}{mtk.24}
\bookmark[dest={mtk.24},level=0]{5. สพฺพมตฺถีติกถา}
\chapterheadpageread{5. สพฺพมตฺถีติกถา}
\csromanheader{2}{1. วาทยุตฺติ}
\proseitem{282}{\csromanfolio{96}\csromansymbolmark{*}สพฺพมตฺถีติ, อามนฺตา.\csromanspacer{} สพฺพตฺถ สพฺพมตฺถีติ, น เหวํ วตฺตพฺเพ.\csromanspacer{} สพฺพมตฺถีติ, อามนฺตา.\csromanspacer{} สพฺพทา สพฺพมตฺถีติ, น เหวํ วตฺตพฺเพ.\csromanspacer{} สพฺพมตฺถีติ, อามนฺตา.\csromanspacer{} สพฺเพน สพฺพํ สพฺพมตฺถีติ, น เหวํ วตฺตพฺเพ.\csromanspacer{} สพฺพมตฺถีติ, อามนฺตา.\csromanspacer{} สพฺเพสุ สพฺพมตฺถีติ, น เหวํ วตฺตพฺเพ.\csromanspacer{} สพฺพมตฺถีติ, อามนฺตา.\csromanspacer{} อโยคนฺติ กตฺวา สพฺพมตฺถีติ, น เหวํ วตฺตพฺเพ.\csromanspacer{} สพฺพมตฺถีติ, อามนฺตา.\csromanspacer{} ยมฺปิ นตฺถิ, ตมฺปตฺถีติ, น เหวํ วตฺตพฺเพ.\csromanspacer{} สพฺพมตฺถีติ, อามนฺตา.\csromanspacer{} สพฺพมตฺถีติ ยา ทิฏฺฐิ สา ทิฏฺฐิ มิจฺฉาทิฏฺฐีติ, ยา ทิฏฺฐิ สา ทิฏฺฐิ สมฺมาทิฏฺฐีติ เหวมตฺถีติ, น เหวํ วตฺตพฺเพ.\csromanspacer{}\csromanspacer{}(สํขิตฺตํ).\csromanspacer{}\csromanspacer{}วาทยุตฺติ.}

\csromanheader{2}{2. กาลสํสนฺทน}
\proseitem{283}{\csromansymbolmark{*}อตีตํ อตฺถีติ, อามนฺตา.\csromanspacer{} นนุ อตีตํ นิรุทฺธํ วิคตํ วิปริณตํ อตฺถงฺคตํ อพฺภตฺถงฺคตนฺติ, อามนฺตา.\csromanspacer{} หญฺจิ อตีตํ นิรุทฺธํ วิคตํ วิปริณตํ อตฺถงฺคตํ อพฺภตฺถงฺคตํ, โน จ วต เร วตฺตพฺเพ “อตีตํ อตฺถี”ติ.}

\prose{\csromansymbolmark{*}อนาคตํ อตฺถีติ, อามนฺตา.\csromanspacer{} นนุ อนาคตํ อชาตํ อภูตํ อสญฺชาตํ อนิพฺพตฺตํ อนภินิพฺพตฺตํ อปาตุภูตนฺติ, อามนฺตา.\csromanspacer{} หญฺจิ อนาคตํ อชาตํ อภูตํ อสญฺชาตํ อนิพฺพตฺตํ อนภินิพฺพตฺตํ อปาตุภูตํ, โน จ วต เร วตฺตพฺเพ “อนาคตํ อตฺถี”ติ.}

\prose{\csromansymbolmark{*}ปจฺจุปฺปนฺนํ อตฺถิ ปจฺจุปฺปนฺนํ อนิรุทฺธํ อวิคตํ อวิปริณตํ น อตฺถงฺคตํ อพฺภตฺถงฺคตนฺติ, อามนฺตา.\csromanspacer{} อตีตํ อตฺถิ อตีตํ อนิรุทฺธํ อวิคตํ อวิปริณตํ น อตฺถงฺคตํ น อพฺภตฺถงฺคตนฺติ, น เหวํ วตฺตพฺเพ ฯเปฯ.\csromanspacer{} ปจฺจุปฺปนฺนํ อตฺถิ ปจฺจุปฺปนฺนํ ชาตํ ภูตํ สญฺชาตํ นิพฺพตฺตํ อภินิพฺพตฺตํ ปาตุภูตนฺติ, อามนฺตา.\csromanspacer{} อนาคตํ อตฺถิ อนาคตํ ชาตํ ภูตํ สญฺชาตํ นิพฺพตฺตํ อภินิพฺพตฺตํ ปาตุภูตนฺติ, น เหวํ วตฺตพฺเพ ฯเปฯ.}

\prose{\csromansymbolmark{*}อตีตํ อตฺถิ อตีตํ นิรุทฺธํ วิคตํ วิปริณตํ อตฺถงฺคตํ อพฺภตฺถงฺคตนฺติ, อามนฺตา.\csromanspacer{} ปจฺจุปฺปนฺนํ อตฺถิ ปจฺจุปฺปนฺนํ นิรุทฺธํ วิคตํ วิปริณตํ อตฺถงฺคตํ อพฺภตฺถงฺคตนฺติ, น เหวํ วตฺตพฺเพ ฯเปฯ.\csromanspacer{} อนาคตํ อตฺถิ อนาคตํ อชาตํ}

\vspace{\tipitakaparskip}
\csromancenter{สพฺพมตฺถีติกถา อภูตํ อสญฺชาตํ อนิพฺพตฺตํ อนภินิพฺพตฺตํ อปาตุภูตนฺติ, อามนฺตา.\csromanspacer{} ปจฺจุปฺปนฺนํ อตฺถิ ปจฺจุปฺปนฺนํ อชาตํ อภูตํ อสญฺชาตํ อนิพฺพตฺตํ อนภินิพฺพตฺตํ อปาตุภูตนฺติ, น เหวํ วตฺตพฺเพ.}
\proseitem{284}{\csromanfolio{97}\csromansymbolmark{*}อตีตํ รูปํ อตฺถีติ, อามนฺตา.\csromanspacer{} นนุ อตีตํ รูปํ นิรุทฺธํ วิคตํ วิปริณตํ อตฺถงฺคตํ อพฺภตฺถงฺคตนฺติ, อามนฺตา.\csromanspacer{} หญฺจิ อตีตํ รูปํ นิรุทฺธํ ฯเปฯ อพฺภตฺถงฺคตํ, โน จ วต เร วตฺตพฺเพ “อตีตํ รูปํ อตฺถี”ติ.}

\prose{\csromansymbolmark{*}อนาคตํ รูปํ อตฺถีติ, อามนฺตา.\csromanspacer{} นนุ อนาคตํ รูปํ อชาตํ อภูตํ อสญฺชาตํ อนิพฺพตฺตํ อนภินิพฺพตฺตํ อปาตุภูตนฺติ, อามนฺตา.\csromanspacer{} หญฺจิ อนาคตํ รูปํ อชาตํ ฯเปฯ อปาตุภูตํ, โน จ วต เร วตฺตพฺเพ “อนาคตํ รูปํ อตฺถี”ติ.}

\prose{\csromansymbolmark{*}ปจฺจุปฺปนฺนํ รูปํ อตฺถิ ปจฺจุปฺปนฺนํ รูปํ อนิรุทฺธํ อวิคตํ อวิปริณตํ น อตฺถงฺคตํ น อพฺภตฺถงฺคตนฺติ, อามนฺตา.\csromanspacer{} อตีตํ รูปํ อตฺถิ อตีตํ รูปํ อนิรุทฺธํ อวิคตํ อวิปริณตํ น อตฺถงฺคตํ น อพฺภตฺถงฺคตนฺติ, น เหวํ วตฺตพฺเพ.}

\prose{\csromansymbolmark{*}ปจฺจุปฺปนฺนํ รูปํ อตฺถิ ปจฺจุปฺปนฺนํ รูปํ ชาตํ ภูตํ สญฺชาตํ นิพฺพตฺตํ อภินิพฺพตฺตํ ปาตุภูตนฺติ, อามนฺตา.\csromanspacer{} อนาคตํ รูปํ อตฺถิ อนาคตํ รูปํ ชาตํ ภูตํ สญฺชาตํ นิพฺพตฺตํ อภินิพฺพตฺตํ ปาตุภูตนฺติ, น เหวํ วตฺตพฺเพ.}

\prose{\csromansymbolmark{*}อตีตํ รูปํ อตฺถิ อตีตํ รูปํ นิรุทฺธํ วิคตํ วิปริณตํ อตฺถงฺคตํ อพฺภตฺถงฺคตนฺติ, อามนฺตา.\csromanspacer{} ปจฺจุปฺปนฺนํ รูปํ อตฺถิ ปจฺจุปฺปนฺนํ รูปํ นิรุทฺธํ วิคตํ วิปริณตํ อตฺถงฺคตํ อพฺภตฺถงฺคตนฺติ, น เหวํ วตฺตพฺเพ.}

\prose{\csromansymbolmark{*}อนาคตํ รูปํ อตฺถิ อนาคตํ รูปํ อชาตํ อภูตํ อสญฺชาตํ อนิพฺพตฺตํ อนภินิพฺพตฺตํ อปาตุภูตนฺติ, อามนฺตา.\csromanspacer{} ปจฺจุปฺปนฺนํ รูปํ อตฺถิ ปจฺจุปฺปนฺนํ รูปํ อชาตํ อภูตํ อสญฺชาตํ อนิพฺพตฺตํ อนภินิพฺพตฺตํ อปาตุภูตนฺติ, น เหวํ วตฺตพฺเพ.}

\prose{\csromansymbolmark{*}อตีตา เวทนา อตฺถิ ฯเปฯ สญฺญา อตฺถิ.\csromanspacer{} สงฺขารา อตฺถิ.\csromanspacer{} วิญฺญาณํ อตฺถีติ, อามนฺตา.\csromanspacer{} นนุ อตีตํ วิญฺญาณํ นิรุทฺธํ วิคตํ วิปริณตํ อตฺถงฺคตํ อพฺภตฺถงฺคตนฺติ, อามนฺตา.\csromanspacer{} หญฺจิ อตีตํ วิญฺญาณํ นิรุทฺธํ ฯเปฯ อพฺภตฺถงฺคตํ, โน จ วต เร วตฺตพฺเพ “อตีตํ วิญฺญาณํ อตฺถี”ติ.}

\prose{\csromanfolio{98}\csromansymbolmark{*}อนาคตํ วิญฺญาณํ อตฺถีติ, อามนฺตา.\csromanspacer{} นนุ อนาคตํ วิญฺญาณํ อชาตํ อภูตํ อสญฺชาตํ อนิพฺพตฺตํ อนภินิพฺพตฺตํ อปาตุภูตนฺติ, อามนฺตา.\csromanspacer{} หญฺจิ อนาคตํ วิญฺญาณํ อชาตํ ฯเปฯ อปาตุภูตํ, โน จ วต เร วตฺตพฺเพ “อนาคตํ วิญฺญาณํ อตฺถี”ติ.}

\prose{\csromansymbolmark{*}ปจฺจุปฺปนฺนํ วิญฺญาณํ อตฺถิ ปจฺจุปฺปนฺนํ วิญฺญาณํ อนิรุทฺธํ ฯเปฯ น อพฺภตฺถงฺคตนฺติ, อามนฺตา.\csromanspacer{} อตีตํ วิญฺญาณํ อตฺถิ อตีตํ วิญฺญาณํ อนิรุทฺธํ ฯเปฯ น อพฺภตฺถงฺคตนฺติ, น เหวํ วตฺตพฺเพ.}

\prose{\csromansymbolmark{*}ปจฺจุปฺปนฺนํ วิญฺญาณํ อตฺถิ ปจฺจุปฺปนฺนํ วิญฺญาณํ ชาตํ ฯเปฯ ปาตุภูตนฺติ, อามนฺตา.\csromanspacer{} อนาคตํ วิญฺญาณํ อตฺถิ อนาคตํ วิญฺญาณํ ชาตํ ฯเปฯ ปาตุภูตนฺติ, น เหวํ วตฺตพฺเพ.}

\prose{\csromansymbolmark{*}อตีตํ วิญฺญาณํ อตฺถิ อตีตํ วิญฺญาณํ นิรุทฺธํ ฯเปฯ อพฺภตฺถงฺคตนฺติ, อามนฺตา.\csromanspacer{} ปจฺจุปฺปนฺนํ วิญฺญาณํ อตฺถิ ปจฺจุปฺปนฺนํ วิญฺญาณํ นิรุทฺธํ ฯเปฯ อพฺภตฺถงฺคตนฺติ, น เหวํ วตฺตพฺเพ.\csromanspacer{} อนาคตํ วิญฺญาณํ อตฺถิ อนาคตํ วิญฺญาณํ อชาตํ ฯเปฯ อปาตุภูตนฺติ, อามนฺตา.}

\prose{\csromansymbolmark{*}ปจฺจุปฺปนฺนํ วิญฺญาณํ อตฺถิ ปจฺจุปฺปนฺนํ วิญฺญาณํ อชาตํ ฯเปฯ อปาตุภูตนฺติ, น เหวํ วตฺตพฺเพ.}

\proseitem{285}{\csromansymbolmark{*}“ปจฺจุปฺปนฺนนฺ”ติ วา “รูปนฺ”ติ วา, “รูปนฺ”ติ วา “ปจฺจุปฺปนฺนนฺ”ติ วา ปจฺจุปฺปนฺนํ รูปํ อปฺปิยํ กริตฺวา เอเสเส เอกฏฺเฐ สเม สมภาเค ตชฺชาเตติ, อามนฺตา.\csromanspacer{} ปจฺจุปฺปนฺนํ รูปํ นิรุชฺฌมานํ ปจฺจุปฺปนฺนภาวํ ชหตีติ, อามนฺตา.\csromanspacer{} รูปภาวํ ชหตีติ, น เหวํ วตฺตพฺเพ ฯเปฯ.}

\prose{\csromansymbolmark{*}“ปจฺจุปฺปนฺนนฺ”ติ วา “รูปนฺ”ติ วา, “รูปนฺ”ติ วา “ปจฺจุปฺปนฺนนฺ”ติ วา ปจฺจุปฺปนฺนํ รูปํ อปฺปิยํ กริตฺวา เอเสเส เอกฏฺเฐ สเม สมภาเค ตชฺชาเตติ, อามนฺตา.\csromanspacer{} ปจฺจุปฺปนฺนํ รูปํ นิรุชฺฌมานํ รูปภาวํ น ชหตีติ, อามนฺตา.\csromanspacer{} ปจฺจุปฺปนฺนภาวํ น ชหตีติ, น เหวํ วตฺตพฺเพ ฯเปฯ.}

\prose{\csromansymbolmark{+}“โอทาตนฺ”ติ วา “วตฺถนฺ”ติ วา, “วตฺถนฺ”ติ วา “โอทาตนฺ”ติ วา โอทาตํ วตฺถํ อปฺปิยํ กริตฺวา เอเสเส เอกฏฺเฐ สเม สมภาเค ตชฺชาเตติ, อามนฺตา.\csromanspacer{} โอทาตํ วตฺถํ รชฺชมานํ โอทาตภาวํ ชหตีติ, อามนฺตา.\csromanspacer{} วตฺถภาวํ ชหตีติ, น เหวํ วตฺตพฺเพ.}

\csromanheader{2}{5. สพฺพมตฺถีติกถา}
\prose{\csromanfolio{99}\csromansymbolmark{+}“โอทาตนฺ”ติ วา “วตฺถนฺ”ติ วา, “วตฺถนฺ”ติ วา “โอทาตนฺ”ติ วา โอทาตํ วตฺถํ อปฺปิยํ กริตฺวา เอเสเส เอกฏฺเฐ สเม สมภาเค ตชฺชาเตติ, อามนฺตา.\csromanspacer{} โอทาตํ วตฺถํ รชฺชมานํ วตฺถภาวํ น ชหตีติ, อามนฺตา.\csromanspacer{} โอทาตภาวํ น ชหตีติ, น เหวํ วตฺตพฺเพ ฯเปฯ.}

\proseitem{286}{\csromansymbolmark{*}รูปํ รูปภาวํ น ชหตีติ, อามนฺตา.\csromanspacer{} รูปํ นิจฺจํ ธุวํ สสฺสตํ อวิปริณามธมฺมนฺติ, น เหวํ วตฺตพฺเพ.\csromanspacer{} นนุ รูปํ รูปภาวํ น ชหตีติ รูปํ อนิจฺจํ อธุวํ อสสฺสตํ วิปริณามธมฺมนฺติ, อามนฺตา.\csromanspacer{} หญฺจิ รูปํ อนิจฺจํ อธุวํ อสสฺสตํ วิปริณามธมฺมํ, โน จ วต เร วตฺตพฺเพ “รูปํ รูปภาวํ น ชหตี”ติ.}

\prose{\csromansymbolmark{*}นิพฺพานํ นิพฺพานภาวํ น ชหตีติ นิพฺพานํ นิจฺจํ ธุวํ สสฺสตํ อวิปริณามธมฺมนฺติ, อามนฺตา.\csromanspacer{} รูปํ รูปภาวํ น ชหตีติ\footnote{น ชหติ (สี, ก)} รูปํ นิจฺจํ ธุวํ สสฺสตํ อวิปริณามธมฺมนฺติ, น เหวํ วตฺตพฺเพ ฯเปฯ.\csromanspacer{} รูปํ รูปภาวํ น ชหติ รูปํ อนิจฺจํ อธุวํ อสสฺสตํ วิปริณามธมฺมนฺติ, อามนฺตา.\csromanspacer{} นิพฺพานํ นิพฺพานภาวํ น ชหติ นิพฺพานํ อนิจฺจํ อธุวํ อสสฺสตํ วิปริณามธมฺมนฺติ, น เหวํ วตฺตพฺเพ ฯเปฯ.}

\prose{\csromansymbolmark{*}อตีตํ อตฺถิ อตีตํ อตีตภาวํ น ชหตีติ, อามนฺตา.\csromanspacer{} อนาคตํ อตฺถิ อนาคตํ อนาคตภาวํ น ชหตีติ, น เหวํ วตฺตพฺเพ.\csromanspacer{} อตีตํ อตฺถิ อตีตํ อตีตภาวํ น ชหตีติ, อามนฺตา.\csromanspacer{} ปจฺจุปฺปนฺนํ อตฺถิ ปจฺจุปฺปนฺนํ ปจฺจุปฺปนฺนภาวํ น ชหตีติ, น เหวํ วตฺตพฺเพ ฯเปฯ.}

\prose{\csromansymbolmark{*}อนาคตํ อตฺถิ อนาคตํ อนาคตภาวํ ชหตีติ\footnote{น ชหตีติ (พหูสุ) อฏฺฐกถา โอโลเกตพฺพา.}, อามนฺตา.\csromanspacer{} อตีตํ อตฺถิ อตีตํ อตีตภาวํ ชหตีติ2, น เหวํ วตฺตพฺเพ ฯเปฯ.}

\prose{\csromansymbolmark{*}ปจฺจุปฺปนฺนํ อตฺถิ ปจฺจุปฺปนฺนํ ปจฺจุปฺปนฺนภาวํ ชหตีติ2, อามนฺตา.\csromanspacer{} อตีตํ อตฺถิ อตีตํ อตีตภาวํ ชหตีติ2, น เหวํ วตฺตพฺเพ.}

\prose{\csromansymbolmark{*}อตีตํ อตฺถิ อตีตํ อตีตภาวํ น ชหตีติ, อามนฺตา.\csromanspacer{} อตีตํ นิจฺจํ ธุวํ สสฺสตํ อวิปริณามธมฺมนฺติ, น เหวํ วตฺตพฺเพ ฯเปฯ.\csromanspacer{} นนุ อตีตํ อนิจฺจํ อธุวํ อสสฺสตํ วิปริณามธมฺมนฺติ, อามนฺตา.\csromanspacer{} หญฺจิ อตีตํ อนิจฺจํ อธุวํ อสสฺสตํ วิปริณามธมฺมํ, โน จ วต เร วตฺตพฺเพ “อตีตํ อตฺถิ อตีตํ อตีตภาวํ น ชหตี”ติ.}

\prose{\csromanfolio{100}\csromansymbolmark{*}นิพฺพานํ อตฺถิ นิพฺพานํ นิพฺพานภาวํ น ชหตีติ นิพฺพานํ นิจฺจํ ธุวํ สสฺสตํ อวิปริณามธมฺมนฺติ, อามนฺตา.\csromanspacer{} อตีตํ อตฺถิ อตีตํ อตีตภาวํ น ชหตีติ อตีตํ นิจฺจํ ธุวํ สสฺสตํ อวิปริณามธมฺมนฺติ, น เหวํ วตฺตพฺเพ ฯเปฯ.}

\prose{\csromansymbolmark{*}อตีตํ อตฺถิ อตีตํ อตีตภาวํ น ชหติ อตีตํ อนิจฺจํ อธุวํ อสสฺสตํ วิปริณามธมฺมนฺติ, อามนฺตา.\csromanspacer{} นิพฺพานํ อตฺถิ นิพฺพานํ นิพฺพานภาวํ น ชหติ นิพฺพานํ อนิจฺจํ อธุวํ อสสฺสตํ วิปริณามธมฺมนฺติ, น เหวํ วตฺตพฺเพ ฯเปฯ.}

\proseitem{287}{\csromansymbolmark{*}อตีตํ รูปํ อตฺถิ อตีตํ รูปํ อตีตภาวํ น ชหตีติ, อามนฺตา.\csromanspacer{} อนาคตํ รูปํ อตฺถิ อนาคตํ รูปํ อนาคตภาวํ น ชหตีติ, น เหวํ วตฺตพฺเพ ฯเปฯ.\csromanspacer{} อตีตํ รูปํ อตฺถิ อตีตํ รูปํ อตีตภาวํ น ชหตีติ, อามนฺตา.\csromanspacer{} ปจฺจุปฺปนฺนํ รูปํ อตฺถิ ปจฺจุปฺปนฺนํ รูปํ ปจฺจุปฺปนฺนภาวํ น ชหตีติ, น เหวํ วตฺตพฺเพ ฯเปฯ.}

\prose{\csromansymbolmark{*}อนาคตํ รูปํ อตฺถิ อนาคตํ รูปํ อนาคตภาวํ ชหตีติ\footnote{น ชหตีติ (พหูสุ) อนุโลมปโญฺหเยว.}, อามนฺตา.\csromanspacer{} อตีตํ รูปํ อตฺถิ อตีตํ รูปํ อตีตภาวํ ชหตีติ1, น เหวํ วตฺตพฺเพ ฯเปฯ.}

\prose{\csromansymbolmark{*}ปจฺจุปฺปนฺนํ รูปํ อตฺถิ ปจฺจุปฺปนฺนํ รูปํ ปจฺจุปฺปนฺนภาวํ ชหตีติ1, อามนฺตา.\csromanspacer{} อตีตํ รูปํ อตฺถิ อตีตํ รูปํ อตีตภาวํ ชหตีติ1, น เหวํ วตฺตพฺเพ ฯเปฯ.}

\prose{\csromansymbolmark{*}อตีตํ รูปํ อตฺถิ อตีตํ รูปํ อตีตภาวํ น ชหตีติ อามนฺตา.\csromanspacer{} อตีตํ รูปํ นิจฺจํ ธุวํ สสฺสตํ อวิปริณามธมฺมนฺติ, น เหวํ วตฺตพฺเพ ฯเปฯ.\csromanspacer{} นนุ อตีตํ รูปํ อนิจฺจํ อธุวํ อสสฺสตํ วิปริณามธมฺมนฺติ, อามนฺตา.\csromanspacer{} หญฺจิ อตีตํ รูปํ อนิจฺจํ ฯเปฯ วิปริณามธมฺมํ, โน จ วต เร วตฺตพฺเพ “อตีตํ รูปํ อตฺถิ อตีตํ รูปํ อตีตภาวํ น ชหติ”ติ.}

\prose{\csromansymbolmark{*}นิพฺพานํ อตฺถิ นิพฺพานํ นิพฺพานภาวํ น ชหติ\footnote{น ชหตีติ (?) ปุริมปเญฺหหิ สํสนฺเทตพฺพํ.} นิพฺพานํ นิจฺจํ ธุวํ สสฺสตํ อวิปริณามธมฺมนฺติ, อามนฺตา.\csromanspacer{} อตีตํ รูปํ อตฺถิ อตีตํ รูปํ อตีตภาวํ น ชหติ2 อตีตํ รูปํ นิจฺจํ ธุวํ สสฺสตํ อวิปริณามธมฺมนฺติ, น เหวํ วตฺตพฺเพ ฯเปฯ.}

\csromanheader{2}{5. สพฺพมตฺถีติกถา}
\prose{\csromanfolio{101}\csromansymbolmark{*}อตีตํ รูปํ อตฺถิ อตีตํ รูปํ อตีตภาวํ น ชหติ อตีตํ รูปํ อนิจฺจํ อธุวํ อสสฺสตํ วิปริณามธมฺมนฺติ, อามนฺตา.\csromanspacer{} นิพฺพานํ อตฺถิ นิพฺพานํ นิพฺพานภาวํ น ชหติ นิพฺพานํ อนิจฺจํ อธุวํ อสสฺสตํ วิปริณามธมฺมนฺติ, น เหวํ วตฺตพฺเพ ฯเปฯ.}

\prose{\csromansymbolmark{*}อตีตา เวทนา อตฺถิ.\csromanspacer{} อตีตา สญฺญา อตฺถิ.\csromanspacer{} อตีตา สงฺขารา อตฺถิ.\csromanspacer{} อตีตํ วิญฺญาณํ อตฺถิ อตีตํ วิญฺญาณํ อตีตภาวํ น ชหตีติ, อามนฺตา.\csromanspacer{} อนาคตํ วิญฺญาณํ อตฺถิ อนาคตํ วิญฺญาณํ อนาคตภาวํ น ชหตีติ, น เหวํ วตฺตพฺเพ ฯเปฯ.\csromanspacer{} อตีตํ วิญฺญาณํ อตฺถิ อตีตํ วิญฺญาณํ อตีตภาวํ น ชหตีติ, อามนฺตา.\csromanspacer{} ปจฺจุปฺปนฺนํ วิญฺญาณํ อตฺถิ ปจฺจุปฺปนฺนํ วิญฺญาณํ ปจฺจุปฺปนฺนภาวํ น ชหตีติ, น เหวํ วตฺตพฺเพ ฯเปฯ.}

\prose{\csromansymbolmark{*}อนาคตํ วิญฺญาณํ อตฺถิ อนาคตํ วิญฺญาณํ อนาคตภาวํ ชหตีติ,\footnote{น ชหตีติ (พหูสุ) อนุโลมปโญฺหเยว.} อามนฺตา.\csromanspacer{} อตีตํ วิญฺญาณํ อตฺถิ อตีตํ วิญฺญาณํ อตีตภาวํ ชหตีติ1, น เหวํ วตฺตพฺเพ ฯเปฯ.}

\prose{\csromansymbolmark{*}ปจฺจุปฺปนฺนํ วิญฺญาณํ อตฺถิ ปจฺจุปฺปนฺนํ วิญฺญาณํ ปจฺจุปฺปนฺนภาวํ ชหตีติ1, อามนฺตา.\csromanspacer{} อตีตํ วิญฺญาณํ อตฺถิ อตีตํ วิญฺญาณํ อตีตภาวํ ชหตีติ1, น เหวํ วตฺตพฺเพ ฯเปฯ.}

\prose{\csromansymbolmark{*}อตีตํ วิญฺญาณํ อตฺถิ อตีตํ วิญฺญาณํ อตีตภาวํ น ชหตีติ, อามนฺตา.\csromanspacer{} อตีตํ วิญฺญาณํ นิจฺจํ ธุวํ สสฺสตํ อวิปริณามธมฺมนฺติ, น เหวํ วตฺตพฺเพ ฯเปฯ.\csromanspacer{} นนุ อตีตํ วิญฺญาณํ อนิจฺจํ อธุวํ อสสฺสตํ วิปริณามธมฺมนฺติ, อามนฺตา.\csromanspacer{} หญฺจิ อตีตํ วิญฺญาณํ อนิจฺจํ อธุวํ อสสฺสตํ วิปริณามธมฺมํ, โน จ วต เร วตฺตพฺเพ “อตีตํ วิญฺญาณํ อตฺถิ อตีตํ วิญฺญาณํ อตีตภาวํ น ชหตี”ติ.}

\prose{\csromansymbolmark{*}นิพฺพานํ อตฺถิ นิพฺพานํ นิพฺพานภาวํ น ชหติ\footnote{น ชหตีติ (?) ปุริมปเญฺหหิ สํสนฺเทตพฺพํ.} นิพฺพานํ นิจฺจํ ธุวํ สสฺสตํ อวิปริณามธมฺมนฺติ, อามนฺตา.\csromanspacer{} อตีตํ วิญฺญาณํ อตฺถิ อตีตํ วิญฺญาณํ อตีตภาวํ น ชหติ2 อตีตํ วิญฺญาณํ นิจฺจํ ธุวํ สสฺสตํ อวิปริณามธมฺมนฺติ, น เหวํ วตฺตพฺเพ ฯเปฯ.}

\prose{\csromansymbolmark{*}อตีตํ วิญฺญาณํ อตฺถิ อตีตํ วิญฺญาณํ อตีตภาวํ น ชหติ อตีตํ วิญฺญาณํ อนิจฺจํ อธุวํ อสสฺสตํ วิปริณามธมฺมนฺติ, อามนฺตา.\csromanspacer{} นิพฺพานํ อตฺถิ}

\prose{\csromanfolio{102}นิพฺพานํ นิพฺพานภาวํ น ชหติ นิพฺพานํ อนิจฺจํ อธุวํ อสสฺสตํ วิปริณามธมฺมนฺติ, น เหวํ วตฺตพฺเพ ฯเปฯ.}

\csromanheader{2}{\textbf{วจนโสธนา}}
\proseitem{288}{\csromansymbolmark{*}อตีตํ นฺวตฺถีติ, อามนฺตา.\csromanspacer{} หญฺจิ อตีตํ นฺวตฺถิ, อตีตํ อตฺถีติ มิจฺฉา.\csromanspacer{} หญฺจิ วา ปน อตฺถิ นฺวาตีตํ, อตฺถิ อตีตนฺติ มิจฺฉา.\csromanspacer{} อนาคตํ นฺวตฺถีติ, อามนฺตา.\csromanspacer{} หญฺจิ อนาคตํ นฺวตฺถิ, อนาคตํ อตฺถีติ มิจฺฉา.\csromanspacer{} หญฺจิ วา ปน อตฺถิ นฺวานาคตํ, อตฺถิ อนาคตนฺติ มิจฺฉา.}

\prose{\csromansymbolmark{*}อนาคตํ หุตฺวา ปจฺจุปฺปนฺนํ โหตีติ, อามนฺตา.\csromanspacer{} ตญฺเญว อนาคตํ ตํ ปจฺจุปฺปนฺนนฺติ, น เหวํ วตฺตพฺเพ ฯเปฯ.\csromanspacer{} ตญฺเญว อนาคตํ ตํ ปจฺจุปฺปนฺนนฺติ, อามนฺตา.\csromanspacer{} หุตฺวา โหติ หุตฺวา โหตีติ, น เหวํ วตฺตพฺเพ ฯเปฯ.\csromanspacer{} หุตฺวา โหติ หุตฺวา โหตีติ, อามนฺตา.\csromanspacer{} น หุตฺวา น โหติ น หุตฺวา น โหตีติ, น เหวํ วตฺตพฺเพ ฯเปฯ.}

\prose{\csromansymbolmark{*}ปจฺจุปฺปนฺนํ หุตฺวา อตีตํ โหตีติ, อามนฺตา.\csromanspacer{} ตญฺเญว ปจฺจุปฺปนฺนํ ตํ อตีตนฺติ, น เหวํ วตฺตพฺเพ ฯเปฯ.\csromanspacer{} ตญฺเญว ปจฺจุปฺปนฺนํ ตํ อตีตนฺติ, อามนฺตา.\csromanspacer{} หุตฺวา โหติ หุตฺวา โหตีติ, น เหวํ วตฺตพฺเพ ฯเปฯ.\csromanspacer{} หุตฺวา โหติ, หุตฺวา โหตีติ, อามนฺตา.\csromanspacer{} น หุตฺวา น โหติ น หุตฺวา น โหตีติ, น เหวํ วตฺตพฺเพ ฯเปฯ.}

\prose{\csromansymbolmark{*}อนาคตํ หุตฺวา ปจฺจุปฺปนฺนํ โหติ, ปจฺจุปฺปนฺนํ หุตฺวา อตีตํ โหตีติ, อามนฺตา.\csromanspacer{} ตญฺเญว อนาคตํ ตํ ปจฺจุปฺปนฺนํ ตํ อตีตนฺติ, น เหวํ วตฺตพฺเพ ฯเปฯ.\csromanspacer{} ตญฺเญว อนาคตํ ตํ ปจฺจุปฺปนฺนํ ตํ อตีตนฺติ, อามนฺตา.\csromanspacer{} หุตฺวา โหติ, หุตฺวา โหตีติ, น เหวํ วตฺตพฺเพ ฯเปฯ.\csromanspacer{} หุตฺวา โหติ, หุตฺวา โหตีติ, อามนฺตา.\csromanspacer{} น หุตฺวา น โหติ, น หุตฺวา น โหตีติ, น เหวํ วตฺตพฺเพ.}

\csromanheader{2}{\textbf{อตีตจกฺขุรูปาทิกถา}}
\proseitem{289}{\csromansymbolmark{*}อตีตํ จกฺขุํ อตฺถิ รูปา อตฺถิ จกฺขุวิญฺญาณํ อตฺถิ อาโลโก อตฺถิ มนสิกาโร อตฺถีติ, อามนฺตา.\csromanspacer{} อตีเตน จกฺขุนา อตีตํ รูปํ ปสฺสตีติ, น เหวํ วตฺตพฺเพ ฯเปฯ.\csromanspacer{} อตีตํ โสตํ อตฺถิ สทฺทา อตฺถิ โสตวิญฺญาณํ อตฺถิ อากาโส อตฺถิ มนสิกาโร อตฺถีติ, อามนฺตา.}

\vspace{\tipitakaparskip}
\csromancenter{สพฺพมตฺถีติกถา อตีเตน โสเตน อตีตํ สทฺทํ สุณาตีติ, น เหวํ วตฺตพฺเพ ฯเปฯ.\csromanspacer{} อตีตํ ฆานํ อตฺถิ คนฺธา อตฺถิ ฆานวิญฺญาณํ อตฺถิ วาโย อตฺถิ มนสิกาโร อตฺถีติ, อามนฺตา.\csromanspacer{} อตีเตน ฆาเนน อตีตํ คนฺธํ ฆายตีติ, น เหวํ วตฺตพฺเพ ฯเปฯ.\csromanspacer{} อตีตา ชิวฺหา อตฺถิ รสา อตฺถิ ชิวฺหาวิญฺญาณํ อตฺถิ อาโป อตฺถิ มนสิกาโร อตฺถีติ, อามนฺตา.\csromanspacer{} อตีตาย ชิวฺหาย อตีตํ รสํ สายตีติ, น เหวํ วตฺตพฺเพ ฯเปฯ.\csromanspacer{} อตีโต กาโย อตฺถิ โผฏฺฐพฺพา อตฺถิ กายวิญฺญาณํ อตฺถิ ปถวี อตฺถิ มนสิกาโร อตฺถีติ, อามนฺตา.\csromanspacer{} อตีเตน กาเยน อตีตํ โผฏฺฐพฺพํ ผุสตีติ, น เหวํ วตฺตพฺเพ ฯเปฯ.\csromanspacer{} อตีโต มโน อตฺถิ ธมฺมา อตฺถิ มโนวิญฺญาณํ อตฺถิ วตฺถุํ อตฺถิ มนสิกาโร อตฺถีติ, อามนฺตา.\csromanspacer{} อตีเตน มเนน อตีตํ ธมฺมํ วิชานาตีติ, น เหวํ วตฺตพฺเพ ฯเปฯ.}
\prose{\csromanfolio{103}\csromansymbolmark{*}อนาคตํ จกฺขุํ อตฺถิ รูปา อตฺถิ จกฺขุวิญฺญานํ อตฺถิ อาโลโก อตฺถิ มนสิกาโร อตฺถีติ, อามนฺตา.\csromanspacer{} อนาคเตน จกฺขุนา อนาคตํ รูปํ ปสฺสตีติ, น เหวํ วตฺตพฺเพ ฯเปฯ.\csromanspacer{} อนาคตํ โสตํ อตฺถิ.\csromanspacer{} ฆานํ อตฺถิ.\csromanspacer{} ชิวฺหา อตฺถิ.\csromanspacer{} กาโย อตฺถิ.\csromanspacer{} มโน อตฺถิ.\csromanspacer{} ธมฺมา อตฺถิ.\csromanspacer{} มโนวิญฺญาณํ อตฺถิ.\csromanspacer{} วตฺถุํ อตฺถิ.\csromanspacer{} มนสิกาโร อตฺถีติ, อามนฺตา.\csromanspacer{} อนาคเตน มเนน อนาคตํ ธมฺมํ วิชานาตีติ, น เหวํ วตฺตพฺเพ ฯเปฯ.}

\prose{\csromansymbolmark{*}ปจฺจุปฺปนฺนํ จกฺขุํ อตฺถิ รูปา อตฺถิ จกฺขุวิญฺญาณํ อตฺถิ อาโลโก อตฺถิ มนสิกาโร อตฺถิ, ปจฺจุปฺปนฺเนน จกฺขุนา ปจฺจุปฺปนฺนํ รูปํ ปสฺสตีติ, อามนฺตา.\csromanspacer{} อตีตํ จกฺขุํ อตฺถิ รูปา อตฺถิ จกฺขุวิญฺญาณํ อตฺถิ อาโลโก อตฺถิ มนสิกาโร อตฺถิ, อตีเตน จกฺขุนา อตีตํ รูปํ ปสฺสตีติ, น เหวํ วตฺตพฺเพ ฯเปฯ.\csromanspacer{} ปจฺจุปฺปนฺนํ โสตํ อตฺถิ.\csromanspacer{} ฆานํ อตฺถิ.\csromanspacer{} ชิวฺหา อตฺถิ.\csromanspacer{} กาโย อตฺถิ.\csromanspacer{} มโน อตฺถิ.\csromanspacer{} ธมฺมา อตฺถิ.\csromanspacer{} มโนวิญฺญาณํ อตฺถิ.\csromanspacer{} วตฺถุํ อตฺถิ.\csromanspacer{} มนสิกาโร อตฺถิ, ปจฺจุปฺปนฺเนน มเนน ปจฺจุปฺปนฺนํ ธมฺมํ วิชานาตีติ, อามนฺตา.\csromanspacer{} อตีโต มโน อตฺถิ ธมฺมา อตฺถิ มโนวิญฺญาณํ อตฺถิ วตฺถุํ อตฺถิ มนสิกาโร อตฺถิ, อตีเตน มเนน อตีตํ ธมฺมํ วิชานาตีติ, น เหวํ วตฺตพฺเพ ฯเปฯ.}

\prose{\csromansymbolmark{*}ปจฺจุปฺปนฺนํ จกฺขุํ อตฺถิ รูปา อตฺถิ จกฺขุวิญฺญาณํ อตฺถิ อาโลโก อตฺถิ มนสิกาโร อตฺถิ, ปจฺจุปฺปนฺเนน จกฺขุนา ปจฺจุปฺปนฺนํ รูปํ ปสฺสตีติ, อามนฺตา.\csromanspacer{} อนาคตํ จกฺขุํ อตฺถิ รูปา อตฺถิ จกฺขุวิญฺญาณํ อตฺถิ อาโลโก อตฺถิ \csromanfolio{104}มนโสกาโร อตฺถิ, อนาคเตน จกฺขุนา อนาคตํ รูปํ ปสฺสตีติ, น เหวํ วตฺตพฺเพ ฯเปฯ.\csromanspacer{} ปจฺจุปฺปนฺนํ โสตํ อตฺถิ.\csromanspacer{} ฆานํ อตฺถิ.\csromanspacer{} ชิวฺหา อตฺถิ.\csromanspacer{} กาโย อตฺถิ.\csromanspacer{} มโน อตฺถิ ธมฺมา อตฺถิ มโนวิญฺญาณํ อตฺถิ วตฺถุํ อตฺถิ มนสิกาโร อตฺถิ, ปจฺจุปฺปนฺเนน มเนน ปจฺจุปฺปนฺนํ ธมฺมํ วิชานาตีติ, อามนฺตา.\csromanspacer{} อนาคโต มโน อตฺถิ ธมฺมา อตฺถิ มโนวิญฺญาณํ อตฺถิ วตฺถุํ อตฺถิ มนสิกาโร อตฺถิ, อนาคเตน มเนน อนาคตํ ธมฺมํ วิชานาตีติ, น เหวํ วตฺตพฺเพ ฯเปฯ.}

\prose{\csromansymbolmark{*}อตีตํ จกฺขุํ อตฺถิ รูปา อตฺถิ จกฺขุวิญฺญาณํ อตฺถิ อาโลโก อตฺถิ มนสิกาโร อตฺถิ, น จ อตีเตน จกฺขุนา อตีตํ รูปํ ปสฺสตีติ, อามนฺตา.\csromanspacer{} ปจฺจุปฺปนฺนํ จกฺขุํ อตฺถิ รูปา อตฺถิ จกฺขุวิญฺญาณํ อตฺถิ อาโลโก อตฺถิ มนสิกาโร อตฺถิ, น จ ปจฺจุปฺปนฺเนน จกฺขุนา ปจฺจุปฺปนฺนํ รูปํ ปสฺสตีติ, น เหวํ วตฺตพฺเพ ฯเปฯ.\csromanspacer{} อตีตํ โสตํ อตฺถิ.\csromanspacer{} ฆานํ อตฺถิ.\csromanspacer{} ชิวฺหา อตฺถิ.\csromanspacer{} กาโย อตฺถิ.\csromanspacer{} มโน อตฺถิ ธมฺมา อตฺถิ มโนวิญฺญาณํ อตฺถิ วตฺถุํ อตฺถิ มนสิกาโร อตฺถิ, น จ อตีเตน มเนน อตีตํ ธมฺมํ วิชานาตีติ, อามนฺตา.\csromanspacer{} ปจฺจุปฺปนฺโน มโน อตฺถิ ธมฺมา อตฺถิ มโนวิญฺญาณํ อตฺถิ วตฺถุํ อตฺถิ มนสิกาโร อตฺถิ, น จ ปจฺจุปฺปนฺเนน มเนน ปจฺจุปฺปนฺนํ ธมฺมํ วิชานาตีติ, น เหวํ วตฺตพฺเพ ฯเปฯ.}

\prose{\csromansymbolmark{*}อนาคตํ จกฺขุํ อตฺถิ รูปา อตฺถิ จกฺขุวิญฺญาณํ อตฺถิ อาโลโก อตฺถิ มนสิกาโร อตฺถิ, น จ อนาคเตน จกฺขุนา อนาคตํ รูปํ ปสฺสตีติ, อามนฺตา.\csromanspacer{} ปจฺจุปฺปนฺนํ จกฺขุํ อตฺถิ รูปา อตฺถิ จกฺขุวิญฺญาณํ อตฺถิ อาโลโก อตฺถิ มนสิกาโร อตฺถิ, น จ ปจฺจุปฺปนฺเนน จกฺขุนา ปจฺจุปฺปนฺนํ รูปํ ปสฺสตีติ, น เหวํ วตฺตพฺเพ ฯเปฯ.\csromanspacer{} อนาคตํ โสตํ อตฺถิ.\csromanspacer{} ฆานํ อตฺถิ.\csromanspacer{} ชิวฺหา อตฺถิ.\csromanspacer{} กาโย อตฺถิ.\csromanspacer{} มโน อตฺถิ ธมฺมา อตฺถิ มโนวิญฺญาณํ อตฺถิ วตฺถุํ อตฺถิ มนสิกาโร อตฺถิ, น จ อนาคเตน มเนน อนาคตํ ธมฺมํ วิชานาตีติ, อามนฺตา.\csromanspacer{} ปจฺจุปฺปนฺโน มโน อตฺถิ ธมฺมา อตฺถิ มโนวิญฺญาณํ อตฺถิ วตฺถุํ อตฺถิ มนสิกาโร อตฺถิ, น จ ปจฺจุปฺปนฺเนน มเนน ปจฺจุปฺปนฺนํ ธมฺมํ วิชานาตีติ, น เหวํ วตฺตพฺเพ ฯเปฯ.}

\csromanheader{2}{\textbf{อตีตญาณาทิกถา}}
\proseitem{290}{\csromansymbolmark{*}อตีตํ ญาณํ อตฺถีติ, อามนฺตา.\csromanspacer{} เตน ญาเณน ญาณกรณียํ กโรตีติ, น เหวํ วตฺตพฺเพ ฯเปฯ.\csromanspacer{} เตน ญาเณน ญาณกรณียํ}

\vspace{\tipitakaparskip}
\csromancenter{สพฺพมตฺถีติกถา กโรตีติ, อามนฺตา.\csromanspacer{} เตน ญาเณน ทุกฺขํ ปริชานาติ, สมุทยํ ปชหติ, นิโรธํ สจฺฉิกโรติ, มคฺคํ ภาเวตีติ, น เหวํ วตฺตพฺเพ ฯเปฯ.}
\prose{\csromanfolio{105}\csromansymbolmark{*}อนาคตํ ญาณํ อตฺถีติ, อามนฺตา.\csromanspacer{} เตน ญาเณน ญาณกรณียํ กโรตีติ, น เหวํ วตฺตพฺเพ ฯเปฯ.\csromanspacer{} เตน ญาเณน ญาณกรณียํ กโรตีติ, อามนฺตา.\csromanspacer{} เตน ญาเณน ทุกฺขํ ปริชานาติ, สมุทยํ ปชหติ, นิโรธํ สจฺฉิกโรติ, มคฺคํ ภาเวตีติ, น เหวํ วตฺตพฺเพ ฯเปฯ.}

\prose{\csromansymbolmark{*}ปจฺจุปฺปนฺนํ ญาณํ อตฺถิ เตน ญาเณน ญาณกรณียํ กโรตีติ, อามนฺตา.\csromanspacer{} อตีตํ ญาณํ อตฺถิ เตน ญาเณน ญาณกรณียํ กโรตีติ, น เหวํ วตฺตพฺเพ ฯเปฯ.\csromanspacer{} ปจฺจุปฺปนฺนํ ญาณํ อตฺถิ เตน ญาเณน ทุกฺขํ ปริชานาติ, สมุทยํ ปชหติ, นิโรธํ สจฺฉิกโรติ, มคฺคํ ภาเวตีติ, อามนฺตา.\csromanspacer{} อตีตํ ญาณํ อตฺถิ เตน ญาเณน ทุกฺขํ ปริชานาติ, สมุทยํ ปชหติ, นิโรธํ สจฺฉิกโรติ, มคฺคํ ภาเวตีติ น เหวํ วตฺตพฺเพ ฯเปฯ.\csromanspacer{} ปจฺจุปฺปนฺนํ ญาณํ อตฺถิ เตน ญาเณน ญาณกรณียํ กโรตีติ, อามนฺตา.\csromanspacer{} อนาคตํ ญาณํ อตฺถิ เตน ญาเณน ญาณกรณียํ กโรตีติ, น เหวํ วตฺตพฺเพ ฯเปฯ.\csromanspacer{} ปจฺจุปฺปนฺนํ ญาณํ อตฺถิ เตน ญาเณน ทุกฺขํ ปริชานาติ, สมุทยํ ปชหติ, นิโรธํ สจฺฉิกโรติ, มคฺคํ ภาเวตีติ, อามนฺตา.\csromanspacer{} อนาคตํ ญาณํ อตฺถิ เตน ญาเณน ทุกฺขํ ปริชานาติ, สมุทยํ ปชหติ, นิโรธํ สจฺฉิกโรติ, มคฺคํ ภาเวตีติ, น เหวํ วตฺตพฺเพ ฯเปฯ.}

\prose{\csromansymbolmark{*}อตีตํ ญาณํ อตฺถิ, น จ เตน ญาเณน ญาณกรณียํ กโรตีติ, อามนฺตา.\csromanspacer{} ปจฺจุปฺปนฺนํ ญาณํ อตฺถิ, น จ เตน ญาเณน ญาณกรณียํ กโรตีติ, น เหวํ วตฺตพฺเพ ฯเปฯ.\csromanspacer{} อตีตํ ญาณํ อตฺถิ, น จ เตน ญาเณน ทุกฺขํ ปริชานาติ, สมุทยํ ปชหติ, นิโรธํ สจฺฉิกโรติ, มคฺคํ ภาเวตีติ, อามนฺตา.\csromanspacer{} ปจฺจุปฺปนฺนํ ญาณํ อตฺถิ, น จ เตน ญาเณน ทุกฺขํ ปริชานาติ, สมุทยํ ปชหติ, นิโรธํ สจฺฉิกโรติ, มคฺคํ ภาเวตีติ, น เหวํ วตฺตพฺเพ ฯเปฯ.}

\prose{\csromansymbolmark{*}อนาคตํ ญาณํ อตฺถิ, น จ เตน ญาเณน ญาณกรณียํ กโรตีติ, อามนฺตา.\csromanspacer{} ปจฺจุปฺปนฺนํ ญาณํ อตฺถิ, น จ เตน ญาเณน ญาณกรณียํ กโรตีติ, น เหวํ วตฺตพฺเพ ฯเปฯ.\csromanspacer{} อนาคตํ ญาณํ อตฺถิ, น จ เตน ญาเณน ทุกฺขํ ปริชานาติ, สมุทยํ ปชหติ, นิโรธํ สจฺฉิกโรติ, มคฺคํ ภาเวตีติ, \csromanfolio{106}อามนฺตา.\csromanspacer{} ปจฺจุปฺปนฺนํ ญาณํ อตฺถิ, น จ เตน ญาเณน ทุกฺขํ ปริชานาติ, สมุทยํ ปชหติ, นิโรธํ สจฺฉิกโรติ, มคฺคํ ภาเวตีติ, น เหวํ วตฺตพฺเพ ฯเปฯ.}

\csromanheader{2}{\textbf{อรหนฺตาทิกถา}}
\proseitem{291}{\csromansymbolmark{*}อรหโต อตีโต ราโค อตฺถีติ, อามนฺตา.\csromanspacer{} อรหา เตน ราเคน สราโคติ, น เหวํ วตฺตพฺเพ ฯเปฯ.\csromanspacer{} อรหโต อตีโต โทโส อตฺถีติ, อามนฺตา.\csromanspacer{} อรหา เตน โทเสน สโทโสติ, น เหวํ วตฺตพฺเพ ฯเปฯ.\csromanspacer{} อรหโต อตีโต โมโห อตฺถีติ, อามนฺตา.\csromanspacer{} อรหา เตน โมเหน สโมโหติ, น เหวํ วตฺตพฺเพ ฯเปฯ.\csromanspacer{} อรหโต อตีโต มาโน อตฺถีติ, อามนฺตา.\csromanspacer{} อรหา เตน มาเนน สมาโนติ, น เหวํ วตฺตพฺเพ ฯเปฯ.\csromanspacer{} อรหโต อตีตา ทิฏฺฐิ อตฺถีติ, อามนฺตา.\csromanspacer{} อรหา ตาย ทิฏฺฐิยา สทิฏฺฐิโกติ, น เหวํ วตฺตพฺเพ ฯเปฯ.\csromanspacer{} อรหโต อตีตา วิจิกิจฺฉา อตฺถีติ, อามนฺตา.\csromanspacer{} อรหา ตาย วิจิกิจฺฉาย สวิจิกิจฺโฉติ, น เหวํ วตฺตพฺเพ ฯเปฯ.\csromanspacer{} อรหโต อตีตํ ถินํ อตฺถีติ, อามนฺตา.\csromanspacer{} อรหา เตน ถิเนน สถิโนติ, น เหวํ วตฺตพฺเพ ฯเปฯ.\csromanspacer{} อรหโต อตีตํ อุทฺธจฺจํ อตฺถีติ, อามนฺตา.\csromanspacer{} อรหา เตน อุทฺธจฺเจน ส-อุทฺธจฺโจติ, น เหวํ วตฺตพฺเพ ฯเปฯ.\csromanspacer{} อรหโต อตีตํ อหิริกํ อตฺถีติ, อามนฺตา.\csromanspacer{} อรหา เตน อหิริเกน ส-อหิริโกติ, น เหวํ วตฺตพฺเพ ฯเปฯ.\csromanspacer{} อรหโต อตีตํ อโนตฺตปฺปํ อตฺถีติ, อามนฺตา.\csromanspacer{} อรหา เตน อโนตฺตปฺเปน ส-อโนตฺตปฺปีติ, น เหวํ วตฺตพฺเพ ฯเปฯ.}

\prose{\csromansymbolmark{*}อนาคามิสฺส อตีตา สกฺกายทิฏฺฐิ อตฺถีติ, อามนฺตา.\csromanspacer{} อนาคามี ตาย ทิฏฺฐิยา สทิฏฺฐิโกติ, น เหวํ วตฺตพฺเพ ฯเปฯ.\csromanspacer{} อนาคามิสฺส อตีตา วิจิกิจฺฉา อตฺถิ.\csromanspacer{} อตีโต สีลพฺพตปรามาโส อตฺถิ.\csromanspacer{} อตีโต อณุสหคโต กามราโค อตฺถิ.\csromanspacer{} อตีโต อณุสหคโต พฺยาปาโท อตฺถีติ, อามนฺตา.\csromanspacer{} อนาคามี เตน พฺยาปาเทน พฺยาปนฺนจิตฺโตติ, น เหวํ วตฺตพฺเพ ฯเปฯ.}

\prose{\csromansymbolmark{*}สกทาคามิสฺส อตีตา สกฺกายทิฏฺฐิ อตฺถีติ, อามนฺตา.\csromanspacer{} สกทาคามี ตาย ทิฏฺฐิยา สทิฏฺฐิโกติ, น เหวํ วตฺตพฺเพ ฯเปฯ.\csromanspacer{} สกทาคามิสฺส อตีตา วิจิกิจฺฉา อตฺถิ.\csromanspacer{} อตีโต สีลพฺพตปรามาโส อตฺถิ.}

\vspace{\tipitakaparskip}
\csromancenter{สพฺพมตฺถีติกถา อตีโต โอฬาริโก กามราโค อตฺถิ.\csromanspacer{} อตีโต โอฬาริโก พฺยาปาโท อตฺถีติ, อามนฺตา.\csromanspacer{} สกทาคามี เตน พฺยาปาเทน พฺยาปนฺนจิตฺโตติ, น เหวํ วตฺตพฺเพ ฯเปฯ.}
\prose{\csromanfolio{107}\csromansymbolmark{*}โสตาปนฺนสฺส อตีตา สกฺกายทิฏฺฐิ อตฺถีติ, อามนฺตา.\csromanspacer{} โสตาปนฺโน ตาย ทิฏฺฐิยา สทิฏฺฐิโกติ, น เหวํ วตฺตพฺเพ ฯเปฯ.\csromanspacer{} โสตาปนฺนสฺส อตีตา วิจิกิจฺฉา อตฺถิ.\csromanspacer{} อตีโต สีลพฺพตปรามาโส อตฺถิ.\csromanspacer{} อตีโต อปายคมนีโย ราโค อตฺถิ.\csromanspacer{} อตีโต อปายคมนีโย โทโส อตฺถิ.\csromanspacer{} อตีโต อปายคมนีโย โมโห อตฺถีติ, อามนฺตา.\csromanspacer{} โสตาปนฺโน เตน โมเหน สโมโหติ, น เหวํ วตฺตพฺเพ ฯเปฯ.}

\proseitem{292}{\csromansymbolmark{*}ปุถุชฺชนสฺส อตีโต ราโค อตฺถิ, ปุถุชฺชโน เตน ราเคน สราโคติ, อามนฺตา.\csromanspacer{} อรหโต อตีโต ราโค อตฺถิ, อรหา เตน ราเคน สราโคติ, น เหวํ วตฺตพฺเพ ฯเปฯ.\csromanspacer{} ปุถุชฺชนสฺส อตีโต โทโส อตฺถิ ฯเปฯ.\csromanspacer{} อตีตํ อโนตฺตปฺปํ อตฺถิ, ปุถุชฺชโน เตน อโนตฺตปฺเปน อโนตฺตปฺปีติ, อามนฺตา.\csromanspacer{} อรหโต อตีตํ อโนตฺตปฺปํ อตฺถิ, อรหา เตน อโนตฺตปฺเปน อโนตฺตปฺปีติ, น เหวํ วตฺตพฺเพ ฯเปฯ.}

\prose{\csromansymbolmark{*}ปุถุชฺชนสฺส อตีตา สกฺกายทิฏฺฐิ อตฺถิ, ปุถุชฺชโน ตาย ทิฏฺฐิยา สทิฏฺฐิโกติ, อามนฺตา.\csromanspacer{} อนาคามิสฺส อตีตา สกฺกายทิฏฺฐิ อตฺถิ, อนาคามี ตาย ทิฏฺฐิยา สทิฏฺฐิโกติ, น เหวํ วตฺตพฺเพ ฯเปฯ.\csromanspacer{} ปุถุชฺชนสฺส อตีตา วิจิกิจฺฉา อตฺถิ ฯเปฯ.\csromanspacer{} อตีโต อณุสหคโต พฺยาปาโท อตฺถิ, ปุถุชฺชโน เตน พฺยาปาเทน พฺยาปนฺนจิตฺโตติ, อามนฺตา.\csromanspacer{} อนาคามิสฺส อตีโต อณุสหคโต พฺยาปาโท อตฺถิ, อนาคามี เตน พฺยาปาเทน พฺยาปนฺนจิตฺโตติ, น เหวํ วตฺตพฺเพ ฯเปฯ.}

\prose{\csromansymbolmark{*}ปุถุชฺชนสฺส อตีตา สกฺกายทิฏฺฐิ อตฺถิ, ปุถุชฺชโน ตาย ทิฏฺฐิยา สทิฏฺฐิโกติ, อามนฺตา.\csromanspacer{} สกทาคามิสฺส อตีตา สกฺกายทิฏฺฐิ อตฺถิ, สกทาคามี ตาย ทิฏฺฐิยา สทิฏฺฐิโกติ, น เหวํ วตฺตพฺเพ ฯเปฯ.\csromanspacer{} ปุถุชฺชนสฺส อตีตา วิจิกิจฺฉา อตฺถิ.\csromanspacer{} อตีโต โอฬาริโก พฺยาปาโท อตฺถิ, ปุถุชฺชโน เตน พฺยาปาเทน พฺยาปนฺนจิตฺโตติ, อามนฺตา.\csromanspacer{} สกทาคามิสฺส อตีโต โอฬาริโก พฺยาปาโท อตฺถิ, สกทาคามี เตน พฺยาปาเทน พฺยาปนฺนจิตฺโตติ, น เหวํ วตฺตพฺเพ ฯเปฯ.}

\prose{\csromanfolio{108}\csromansymbolmark{*}ปุถุชฺชนสฺส อตีตา สกฺกายทิฏฺฐิ อตฺถิ, ปุถุชฺชโน ตาย ทิฏฺฐิยา สทิฏฺฐิโกติ, อามนฺตา.\csromanspacer{} โสตาปนฺนสฺส อตีตา สกฺกายทิฏฺฐิ อตฺถิ, โสตาปนฺโน ตาย ทิฏฺฐิยา สทิฏฺฐิโกติ, น เหวํ วตฺตพฺเพ ฯเปฯ.\csromanspacer{} ปุถุชฺชนสฺส อตีตา วิจิกิจฺฉา อตฺถิ ฯเปฯ.\csromanspacer{} อตีโต อปายคมนีโย โมโห อตฺถิ, ปุถุชฺชโน เตน โมเหน สโมโหติ, อามนฺตา.\csromanspacer{} โสตาปนฺนสฺส อตีโต อปายคมนีโย โมโห อตฺถิ, โสตาปนฺโน เตน โมเหน สโมโหติ, น เหวํ วตฺตพฺเพ ฯเปฯ.}

\prose{\csromansymbolmark{*}อรหโต อตีโต ราโค อตฺถิ, น จ อรหา เตน ราเคน สราโคติ, อามนฺตา.\csromanspacer{} ปุถุชฺชนสฺส อตีโต ราโค อตฺถิ, น จ ปุถุชฺชโน เตน ราเคน สราโคติ, น เหวํ วตฺตพฺเพ ฯเปฯ.\csromanspacer{} อรหโต อตีโต โทโส อตฺถิ ฯเปฯ.\csromanspacer{} อตีตํ อโนตฺตปฺปํ อตฺถิ, น จ อรหา เตน อโนตฺตปฺเปน อโนตฺตปฺปีติ, อามนฺตา.\csromanspacer{} ปุถุชฺชนสฺส อตีตํ อโนตฺตปฺปํ อตฺถิ, น จ ปุถุชฺชโน เตน อโนตฺตปฺเปน อโนตฺตปฺปีติ, น เหวํ วตฺตพฺเพ ฯเปฯ.}

\prose{\csromansymbolmark{*}อนาคามิสฺส อตีตา สกฺกายทิฏฺฐิ อตฺถิ, น จ อนาคามี ตาย ทิฏฺฐิยา สทิฏฺฐิโกติ, อามนฺตา.\csromanspacer{} ปุถุชฺชนสฺส อตีตา สกฺกายทิฏฺฐิ อตฺถิ, น จ ปุถุชฺชโน ตาย ทิฏฺฐิยา สทิฏฺฐิโกติ, น เหวํ วตฺตพฺเพ ฯเปฯ.\csromanspacer{} อนาคามิสฺส อตีตา วิจิกิจฺฉา อตฺถิ ฯเปฯ.\csromanspacer{} อตีโต อณุสหคโต พฺยาปาโท อตฺถิ, น จ อนาคามี เตน พฺยาปาเทน พฺยาปนฺนจิตฺโตติ, อามนฺตา.\csromanspacer{} ปุถุชฺชนสฺส อตีโต อณุสหคโต พฺยาปาโท อตฺถิ, น จ ปุถุชฺชโน เตน พฺยาปาเทน พฺยาปนฺนจิตฺโตติ, น เหวํ วตฺตพฺเพ ฯเปฯ.}

\prose{\csromansymbolmark{*}สกทาคามิสฺส อตีตา สกฺกายทิฏฺฐิ อตฺถิ, น จ สกทาคามี ตาย ทิฏฺฐิยา สทิฏฺฐิโกติ, อามนฺตา.\csromanspacer{} ปุถุชฺชนสฺส อตีตา สกฺกายทิฏฺฐิ อตฺถิ, น จ ปุถุชฺชโน ตาย ทิฏฺฐิยา สทิฏฺฐิโกติ, น เหวํ วตฺตพฺเพ ฯเปฯ.\csromanspacer{} สกทาคามิสฺส อตีตา วิจิกิจฺฉา อตฺถิ ฯเปฯ.\csromanspacer{} อตีโต โอฬาริโก พฺยาปาโท อตฺถิ, น จ สกทาคามี เตน พฺยาปาเทน พฺยาปนฺนจิตฺโตติ, อามนฺตา.\csromanspacer{} ปุถุชฺชนสฺส อตีโต โอฬาริโก พฺยาปาโท อตฺถิ, น จ ปุถุชฺชโน เตน พฺยาปาเทน พฺยาปนฺนจิตฺโตติ, น เหวํ วตฺตพฺเพ ฯเปฯ. \csromansymbolmark{*}โสตาปนฺนสฺส อตีตา สกฺกายทิฏฺฐิ อตฺถิ, น จ โสตาปนฺโน ตาย ทิฏฺฐิยา สทิฏฺฐิโกติ, อามนฺตา.\csromanspacer{} ปุถุชฺชนสฺส อตีตา}

\vspace{\tipitakaparskip}
\csromancenter{สพฺพมตฺถีติกถา สกฺกายทิฏฺฐิ อตฺถิ, น จ ปุถุชฺชโน ตาย ทิฏฺฐิยา สทิฏฺฐิโกติ, น เหวํ วตฺตพฺเพ ฯเปฯ.\csromanspacer{} โสตาปนฺนสฺส อตีตา วิจิกิจฺฉา อตฺถิ ฯเปฯ.\csromanspacer{} อตีโต อปายคมนีโย โมโห อตฺถิ, น จ โสตาปนฺโน เตน โมเหน สโมโหติ, อามนฺตา.\csromanspacer{} ปุถุชฺชนสฺส อตีโต อปายคมนีโย โมโห อตฺถิ, น จ ปุถุชฺชโน เตน โมเหน สโมโหติ, น เหวํ วตฺตพฺเพ ฯเปฯ.}
\csromanheader{2}{\textbf{อตีตหตฺถาทิกถา}}
\proseitem{293}{\csromanfolio{109}\csromansymbolmark{*}อตีตา หตฺถา อตฺถีติ, อามนฺตา.\csromanspacer{} อตีเตสุ หตฺเถสุ สติ อาทานนิกฺเขปนํ ปญฺญายตีติ, น เหวํ วตฺตพฺเพ ฯเปฯ.\csromanspacer{} อตีตา ปาทา อตฺถีติ, อามนฺตา.\csromanspacer{} อตีเตสุ ปาเทสุ สติ อภิกฺกมปฏิกฺกโม ปญฺญายตีติ, น เหวํ วตฺตพฺเพ ฯเปฯ.\csromanspacer{} อตีตา ปพฺพา อตฺถีติ, อามนฺตา.\csromanspacer{} อตีเตสุ ปพฺเพสุ สติ สมิญฺชนปสารณํ ปญฺญายตีติ, น เหวํ วตฺตพฺเพ ฯเปฯ.\csromanspacer{} อตีโต กุจฺฉิ อตฺถีติ, อามนฺตา.\csromanspacer{} อตีตสฺมิํ กุจฺฉิสฺมิํ สติ ชิฆจฺฉา ปิปาสา ปญฺญายตีติ, น เหวํ วตฺตพฺเพ ฯเปฯ.}

\prose{\csromansymbolmark{*}อตีโต กาโย อตฺถีติ, อามนฺตา.\csromanspacer{} อตีโต กาโย ปคฺคหนิคฺคหุปโค เฉทนเภทนุปโค กาเกหิ คิชฺเฌหิ กุลเลหิ สาธารโณติ, น เหวํ วตฺตพฺเพ ฯเปฯ.\csromanspacer{} อตีเต กาเย วิสํ กเมยฺย.\csromanspacer{} สตฺถํ กเมยฺย.\csromanspacer{} อคฺคิ กเมยฺยาติ, น เหวํ วตฺตพฺเพ ฯเปฯ.\csromanspacer{} ลพฺภา อตีโต กาโย อทฺทุพนฺธเนน พนฺธิตุํ, รชฺชุพนฺธเนน พนฺธิตุํ, สงฺขลิกพนฺธเนน พนฺธิตุํ, คามพนฺธเนน พนฺธิตุํ, นิคมพนฺธเนน พนฺธิตุํ, นครพนฺธเนน พนฺธิตุํ, ชนปทพนฺธเนน พนฺธิตุํ, กณฺฐปญฺจเมหิ พนฺธเนหิ พนฺธิตุนฺติ, น เหวํ วตฺตพฺเพ ฯเปฯ.}

\prose{\csromansymbolmark{*}อตีโต อาโป อตฺถีติ, อามนฺตา.\csromanspacer{} เตน อาเปน อาปกรณียํ กโรตีติ, น เหวํ วตฺตพฺเพ ฯเปฯ.\csromanspacer{} อตีโต เตโช อตฺถีติ, อามนฺตา.\csromanspacer{} เตน เตเชน เตชกรณียํ กโรตีติ, น เหวํ วตฺตพฺเพ ฯเปฯ.\csromanspacer{} อตีโต วาโย อตฺถีติ, อามนฺตา.\csromanspacer{} เตน วาเยน วายกรณียํ กโรตีติ, น เหวํ วตฺตพฺเพ ฯเปฯ.}

\csromanheader{2}{\textbf{อตีตกฺขนฺธาทิสโมธานกถา}}
\proseitem{294}{\csromansymbolmark{*}อตีโต รูปกฺขนฺโธ อตฺถิ.\csromanspacer{} อนาคโต รูปกฺขนฺโธ อตฺถิ.\csromanspacer{} ปจฺจุปฺปนฺโน รูปกฺขนฺโธ อตฺถีติ, อามนฺตา.\csromanspacer{} ตโย รูปกฺขนฺธาติ, น เหวํ \csromanfolio{110}วตฺตพฺเพ ฯเปฯ.\csromanspacer{} อตีตา ปญฺจกฺขนฺธา อตฺถิ.\csromanspacer{} อนาคตา ปญฺจกฺขนฺธา อตฺถิ.\csromanspacer{} ปจฺจุปฺปนฺนา ปญฺจกฺขนฺธา อตฺถีติ, อามนฺตา.\csromanspacer{} ปนฺนรสกฺขนฺธาติ, น เหวํ วตฺตพฺเพ ฯเปฯ.}

\prose{\csromansymbolmark{*}อตีตํ จกฺขายตนํ อตฺถิ.\csromanspacer{} อนาคตํ จกฺขายตนํ อตฺถิ.\csromanspacer{} ปจฺจุปฺปนฺนํ จกฺขายตนํ อตฺถีติ, อามนฺตา.\csromanspacer{} ตีณิ จกฺขายตนานีติ, น เหวํ วตฺตพฺเพ ฯเปฯ.\csromanspacer{} อตีตานิ ทฺวาทสายตนานิ อตฺถิ.\csromanspacer{} อนาคตานิ ทฺวาทสายตนานิ อตฺถิ.\csromanspacer{} ปจฺจุปฺปนฺนานิ ทฺวาทสายตนานิ อตฺถีติ, อามนฺตา.\csromanspacer{} ฉตฺติํสายตนานีติ, น เหวํ วตฺตพฺเพ ฯเปฯ.}

\prose{\csromansymbolmark{*}อตีตา จกฺขุธาตุ อตฺถิ.\csromanspacer{} อนาคตา จกฺขุธาตุ อตฺถิ.\csromanspacer{} ปจฺจุปฺปนฺนา จกฺขุธาตุ อตฺถีติ, อามนฺตา.\csromanspacer{} ติสฺโส จกฺขุธาตุโยติ, น เหวํ วตฺตพฺเพ ฯเปฯ.\csromanspacer{} อตีตา อฏฺฐารส ธาตุโย อตฺถิ.\csromanspacer{} อนาคตา อฏฺฐารส ธาตุโย อตฺถิ.\csromanspacer{} ปจฺจุปฺปนฺนา อฏฺฐารส ธาตุโย อตฺถีติ, อามนฺตา.\csromanspacer{} จตุปญฺญาส ธาตุโยติ, น เหวํ วตฺตพฺเพ ฯเปฯ.}

\prose{\csromansymbolmark{*}อตีตํ จกฺขุนฺทฺริยํ อตฺถิ.\csromanspacer{} อนาคตํ จกฺขุนฺทฺริยํ อตฺถิ.\csromanspacer{} ปจฺจุปฺปนฺนํ จกฺขุนฺทฺริยํ อตฺถีติ, อามนฺตา.\csromanspacer{} ตีณิ จกฺขุนฺทฺริยานีติ, น เหวํ วตฺตพฺเพ ฯเปฯ.\csromanspacer{} อตีตานิ พาวีสตินฺทฺริยานิ อตฺถิ.\csromanspacer{} อนาคตานิ พาวีสตินฺทฺริยานิ อตฺถิ.\csromanspacer{} ปจฺจุปฺปนฺนานิ พาวีสตินฺทฺริยานิ อตฺถีติ, อามนฺตา.\csromanspacer{} ฉสฏฺฐินฺทฺริยานีติ, น เหวํ วตฺตพฺเพ ฯเปฯ.}

\prose{\csromansymbolmark{*}อตีโต ราชา จกฺกวตฺตี อตฺถิ.\csromanspacer{} อนาคโต ราชา จกฺกวตฺตี อตฺถิ.\csromanspacer{} ปจฺจุปฺปนฺโน ราชา จกฺกวตฺตี อตฺถีติ, อามนฺตา.\csromanspacer{} ติณฺณนฺนํ ราชูนํ จกฺกวตฺตีนํ สมฺมุขีภาโว โหตีติ, น เหวํ วตฺตพฺเพ ฯเปฯ.}

\prose{\csromansymbolmark{*}อตีโต สมฺมาสมฺพุทฺโธ อตฺถิ.\csromanspacer{} อนาคโต สมฺมาสมฺพุทฺโธ อตฺถิ.\csromanspacer{} ปจฺจุปฺปนฺโน สมฺมาสมฺพุทฺโธ อตฺถีติ, อามนฺตา.\csromanspacer{} ติณฺณนฺนํ สมฺมาสมฺพุทฺธานํ สมฺมุขีภาโว โหตีติ, น เหวํ วตฺตพฺเพ ฯเปฯ.}

\csromanheader{2}{\textbf{ปทโสธนกถา}}
\proseitem{295}{\csromansymbolmark{*}อตีตํ อตฺถีติ, อามนฺตา.\csromanspacer{} อตฺถิ อตีตนฺติ, อตฺถิ สิยา อตีตํ, สิยา นฺวาตีตนฺติ.}

\prose{อาชานาหิ นิคฺคหํ\csromandash{}หญฺจิ อตีตํ อตฺถิ อตฺถิ สิยา อตีตํ, สิยา นฺวาตีตํ, เตนาตีตํ นฺวาตีตํ, นฺวาตีตํ อตีตนฺติ.\csromanspacer{} ยํ ตตฺถ}

\vspace{\tipitakaparskip}
\csromancenter{สพฺพมตฺถีติกถา วเทสิ “วตฺตพฺเพ โข ‘อตีตํ อตฺถิ อตฺถิ สิยา อตีตํ, สิยา นฺวาตีตํ, เตนาตีตํ นฺวาตีตํ, นฺวาตีตํ อตีตนฺ’ติ”, มิจฺฉา.}
\prose{\csromanfolio{111}โน เจ ปน อตีตํ นฺวาตีตํ นฺวาตีตํ, อตีตนฺติ, โน จ วต เร วตฺตพฺเพ “อตีตํ อตฺถิ อตฺถิ สิยา อตีตํ, สิยา นฺวาตีตนฺ”ติ.\csromanspacer{} ยํ ตตฺถ วเทสิ “วตฺตพฺเพ โข ‘อตีตํ อตฺถิ อตฺถิ สิยา อตีตํ, สิยา นฺวาตีตํ, เตนาตีตํ นฺวาตีตํ, นฺวาตีตํ อตีตนฺ’ติ”, มิจฺฉา.}

\prose{\csromansymbolmark{*}อนาคตํ อตฺถีติ, อามนฺตา.\csromanspacer{} อตฺถิ อนาคตนฺติ, อตฺถิ สิยา อนาคตํ, สิยา นฺวานาคตนฺติ.}

\prose{อาชานาหิ นิคฺคหํ\csromandash{}หญฺจิ อนาคตํ อตฺถิ อตฺถิ สิยา อนาคตํ สิยา นฺวานาคตํ, เตนานาคตํ นฺวานาคตํ นฺวานาคตํ, อนาคตนฺติ.\csromanspacer{} ยํ ตตฺถ วเทสิ “วตฺตพฺเพ โข ‘อนาคตํ อตฺถิ อตฺถิ สิยา อนาคตํ, สิยา นฺวานาคตํ, เตนานาคตํ นฺวานาคตํ, นฺวานาคตํ อนาคตนฺ’ติ”, มิจฺฉา.}

\prose{โน เจ ปนานาคตํ นฺวานาคตํ, นฺวานาคตํ อนาคตนฺติ, โน จ วต เร วตฺตพฺเพ “อนาคตํ อตฺถิ อตฺถิ สิยา อนาคตํ, สิยา นฺวานาคตนฺ”ติ.\csromanspacer{} ยํ ตตฺถ วเทสิ “วตฺตพฺเพ โข ‘อนาคตํ อตฺถิ อตฺถิ สิยา อนาคตํ, สิยา นฺวานาคตํ, เตนานาคตํ นฺวานาคตํ, นฺวานาคตํ อนาคตนฺ’ติ”, มิจฺฉา.}

\prose{\csromansymbolmark{+}ปจฺจุปฺปนฺนํ อตฺถีติ, อามนฺตา.\csromanspacer{} อตฺถิ ปจฺจุปฺปนฺนนฺติ, อตฺถิ สิยา ปจฺจุปฺปนฺนํ, สิยา โน ปจฺจุปฺปนฺนนฺติ.}

\prose{อาชานาหิ นิคฺคหํ\csromandash{}หญฺจิ ปจฺจุปฺปนฺนํ อตฺถิ อตฺถิ สิยา ปจฺจุปฺปนฺนํ, สิยา โน ปจฺจุปฺปนฺนํ, เตน ปจฺจุปฺปนฺนํ โน ปจฺจุปฺปนฺนํ, โน ปจฺจุปฺปนฺนํ ปจฺจุปฺปนฺนนฺติ.\csromanspacer{} ยํ ตตฺถ วเทสิ “วตฺตพฺเพ โข ‘ปจฺจุปฺปนฺนํ อตฺถิ อตฺถิ สิยา ปจฺจุปฺปนฺนํ, สิยา โน ปจฺจุปฺปนฺนํ, เตน ปจฺจุปฺปนฺนํ โน ปจฺจุปฺปนฺนํ, โน ปจฺจุปฺปนฺนํ ปจฺจุปฺปนฺนนฺ’ติ”, มิจฺฉา.}

\prose{โน เจ ปน ปจฺจุปฺปนฺนํ โน ปจฺจุปฺปนฺนํ, โน ปจฺจุปฺปนฺนํ ปจฺจุปฺปนฺนนฺติ, โน จ วต เร วตฺตพฺเพ “ปจฺจุปฺปนฺนํ อตฺถิ อตฺถิ สิยา ปจฺจุปฺปนฺนํ, สิยา โน ปจฺจุปฺปนฺนนฺ”ติ.\csromanspacer{} ยํ ตตฺถ วเทสิ “วตฺตพฺเพ โข ‘ปจฺจุปฺปนฺนํ อตฺถิ อตฺถิ \csromanfolio{112}สิยา ปจฺจุปฺปนฺนํ, สิยา โน ปจฺจุปฺปนฺนํ, เตน ปจฺจุปฺปนฺนํ โน ปจฺจุปฺปนฺนํ, โน ปจฺจุปฺปนฺนํ ปจฺจุปฺปนฺนนฺ’ติ”, มิจฺฉา.}

\prose{\csromansymbolmark{+}นิพฺพานํ อตฺถีติ, อามนฺตา.\csromanspacer{} อตฺถิ นิพฺพานนฺติ, อตฺถิ สิยา นิพฺพานํ, สิยา โน นิพฺพานนฺติ.}

\prose{อาชานาหิ นิคฺคหํ\csromandash{}หญฺจิ นิพฺพานํ อตฺถิ อตฺถิ สิยา นิพฺพานํ, สิยา โน นิพฺพานํ, เตน นิพฺพานํ โน นิพฺพานํ, โน นิพฺพานํ นิพฺพานนฺติ.\csromanspacer{} ยํ ตตฺถ วเทสิ “วตฺตพฺเพ โข ‘นิพฺพานํ อตฺถิ อตฺถิ สิยา นิพฺพานํ, สิยา โน นิพฺพานํ, เตน นิพฺพานํ โน นิพฺพานํ, โน นิพฺพานํ นิพฺพานนฺ’ติ”, มิจฺฉา.}

\prose{โน เจ ปน นิพฺพานํ โน นิพฺพานํ, โน นิพฺพานํ นิพฺพานนฺติ, โน จ วต เร วตฺตพฺเพ “นิพฺพานํ อตฺถิ อตฺถิ สิยา นิพฺพานํ, สิยา โน นิพฺพานนฺ”ติ.\csromanspacer{} ยํ ตตฺถ วเทสิ “วตฺตพฺเพ โข ‘นิพฺพานํ อตฺถิ อตฺถิ สิยา นิพฺพานํ, สิยา โน นิพฺพานํ, เตน นิพฺพานํ โน นิพฺพานํ, โน นิพฺพานํ นิพฺพานนฺ’ติ”, มิจฺฉา.}

\csromanheader{2}{\textbf{สุตฺตสาธน}}
\proseitem{296}{\csromansymbolmark{+}น วตฺตพฺพํ “อตีตํ อตฺถิ, อนาคตํ อตฺถี”ติ, อามนฺตา.\csromanspacer{} นนุ วุตฺตํ ภควตา “ยํ กิญฺจิ ภิกฺขเว รูปํ อตีตานาคตปจฺจุปฺปนฺนํ อชฺฌตฺตํ วา พหิทฺธา วา โอฬาริกํ วา สุขุมํ วา หีนํ วา ปณีตํ วา ยํ ทูเร สนฺติเก วา, อยํ วุจฺจติ รูปกฺขนฺโธ.\csromanspacer{} ยา กาจิ เวทนา.\csromanspacer{} ยา กาจิ สญฺญา.\csromanspacer{} เย เกจิ สงฺขารา.\csromanspacer{} ยํ กิญฺจิ วิญฺญาณํ อตีตานาคตปจฺจุปฺปนฺนํ อชฺฌตฺตํ วา พหิทฺธา วา โอฬาริกํ วา สุขุมํ วา หีนํ วา ปณีตํ วา ยํ ทูเร สนฺติเก วา, อยํ วุจฺจติ วิญฺญาณกฺขนฺโธ”ติ\footnote{สํ 2. 39 ปิฏฺเฐ.} อตฺเถว สุตฺตนฺโตติ, อามนฺตา.\csromanspacer{} เตน หิ อตีตํ อตฺถิ, อนาคตํ อตฺถีติ.}

\prose{\csromansymbolmark{*}อตีตํ อตฺถิ, อนาคตํ อตฺถีติ, อามนฺตา.\csromanspacer{} นนุ วุตฺตํ ภควตา\csromandash{} ตโยเม ภิกฺขเว นิรุตฺติปถา อธิวจนปถา ปญฺญตฺติปถา อสํกิณฺณา อสํกิณฺณปุพฺพา น สํกิยนฺติ น สํกิยิสฺสนฺติ อปฺปฏิกุฏฺฐา สมเณหิ พฺราหฺมเณหิ วิญฺญูหิ.\csromanspacer{} กตเม ตโย, ยํ ภิกฺขเว รูปํ อตีตํ นิรุทฺธํ วิคตํ วิปริณตํ “อโหสี”ติ ตสฺส สงฺขา “อโหสี”ติ ตสฺส สมญฺญา “อโหสี”ติ ตสฺส ปญฺญตฺติ, น ตสฺส สงฺขา “อตฺถี”ติ น ตสฺส สงฺขา “ภวิสฺสตี”ติ.\csromanspacer{} ยา เวทนา ฯเปฯ.\csromanspacer{} ยา สญฺญา.\csromanspacer{} เย สงฺขารา.\csromanspacer{} ยํ}

\vspace{\tipitakaparskip}
\csromancenter{สพฺพมตฺถีติกถา วิญฺญาณํ อตีตํ นิรุทฺธํ วิคตํ วิปริณตํ “อโหสี”ติ ตสฺส สงฺขา “อโหสี”ติ ตสฺส สมญฺญา “อโหสี”ติ ตสฺส ปญฺญตฺติ, น ตสฺส สงฺขา “อตฺถี”ติ น ตสฺส สงฺขา “ภวิสฺสตี”ติ.}
\prose{\csromanfolio{113}ยํ ภิกฺขเว รูปํ อชาตํ อปาตุภูตํ “ภวิสฺสตี”ติ ตสฺส สงฺขา “ภวิสฺสตี”ติ ตสฺส สมญฺญา “ภวิสฺสตี”ติ ตสฺส ปญฺญตฺติ, น ตสฺส สงฺขา “อตฺถี”ติ น ตสฺส สงฺขา “อโหสี”ติ.\csromanspacer{} ยา เวทนา ฯเปฯ.\csromanspacer{} ยา สญฺญา.\csromanspacer{} เย สงฺขารา.\csromanspacer{} ยํ วิญฺญาณํ อชาตํ อปาตุภูตํ “ภวิสฺสตี”ติ ตสฺส สงฺขา “ภวิสฺสตี”ติ ตสฺส สมญฺญา “ภวิสฺสตี”ติ ตสฺส ปญฺญตฺติ, น ตสฺส สงฺขา “อตฺถี”ติ น ตสฺส สงฺขา “อโหสี”ติ.}

\prose{ยํ ภิกฺขเว รูปํ ชาตํ ปาตุภูตํ “อตฺถี”ติ ตสฺส สงฺขา “อตฺถี”ติ ตสฺส สมญฺญา “อตฺถี”ติ ตสฺส ปญฺญตฺติ, น ตสฺส สงฺขา “อโหสี”ติ น ตสฺส สงฺขา “ภวิสฺสตี”ติ.\csromanspacer{} ยา เวทนา ฯเปฯ.\csromanspacer{} ยา สญฺญา.\csromanspacer{} เย สงฺขารา.\csromanspacer{} ยํ วิญฺญาณํ ชาตํ ปาตุภูตํ “อตฺถี”ติ ตสฺส สงฺขา “อตฺถี”ติ ตสฺส สมญฺญา “อตฺถี”ติ ตสฺส ปญฺญตฺติ, น ตสฺส สงฺขา “อโหสี”ติ น ตสฺส สงฺขา “ภวิสฺสตี”ติ.\csromanspacer{} อิเม โข ภิกฺขเว ตโย นิรุตฺติปถา อธิวจนปถา ปญฺญตฺติปถา อสํกิณฺณา อสํกิณฺณปุพฺพา น สํกิยนฺติ น สํกิยิสฺสนฺติ อปฺปฏิกุฏฺฐา สมเณหิ พฺรหฺมเณหิ วิญฺญูหิ.}

\prose{เยปิ เต ภิกฺขเว อเหสุํ อุกฺกลา วสฺสภญฺญา อเหตุกวาทา อกิริยวาทา นตฺถิกวาทา, เตปิเม ตโย นิรุตฺติปเถ อธิวจนปเถ ปญฺญตฺติปเถ น ครหิตพฺพํ น ปฏิกฺโกสิตพฺพํ อมญฺญิํสุ.\csromanspacer{} ตํ กิสฺส เหตุ, นินฺทาพฺยาโรส-อุปารมฺภภยาติ\footnote{สํ 2. 59 ปิฏฺเฐ.} อตฺเถว สุตฺตนฺโตติ, อามนฺตา.\csromanspacer{} เตน หิ น วตฺตพฺพํ “อตีตํ อตฺถิ อนาคตํ อตฺถี”ติ.}

\prose{\csromansymbolmark{*}อตีตํ อตฺถีติ, อามนฺตา.\csromanspacer{} นนุ อายสฺมา ผคฺคุโน ภควนฺตํ เอตทโวจ\csromandash{}อตฺถิ นุ โข ตํ ภนฺเต จกฺขุํ เยน จกฺขุนา อตีเต พุทฺเธ ปรินิพฺพุเต ฉินฺนปปญฺเจ ฉินฺนวฏุเม ปริยาทินฺนวฏฺเฏ สพฺพทุกฺขวีติวตฺเต ปญฺญาปยมาโน ปญฺญาเปยฺยาติ.\csromanspacer{} อตฺถิ นุ โข สา ภนฺเต ชิวฺหา ฯเปฯ.\csromanspacer{} อตฺถิ นุ โข โส ภนฺเต มโน เยน มเนน อตีเต พุทฺเธ ปรินิพฺพุเต ฉินฺนปปญฺเจ ฉินฺนวฏุเม ปริยาทินฺนวฏฺเฏ สพฺพทุกฺขปีติวตฺเต ปญฺญาปยมาโน ปญฺญาเปยฺยาติ.}

\prose{\csromanfolio{114}นตฺถิ โข ตํ ผคฺคุน จกฺขุํ เยน จกฺขุนา อตีเต พุทฺเธ ปรินิพฺพุเต ฉินฺนปปญฺเจ ฉินฺนวฏุเม ปริยาทินฺนวฏฺเฏ สพฺพทุกฺขวีติวตฺเต ปญฺญาปยมาโน ปญฺญาเปยฺย.\csromanspacer{} นตฺถิ โข สา ผคฺคุน ชิวฺหา ฯเปฯ.\csromanspacer{} นตฺถิ นุ โข โส ผคฺคุน มโน เยน มเนน อตีเต พุทฺเธ ปรินิพฺพุเต ฉินฺนปปญฺเจ ฉินฺนวฏุเม ปริยาทินฺนวฏฺเฏ สพฺพทุกฺขวีติวตฺเต ปญฺญาปยมาโน ปญฺญาเปยฺยาติ\footnote{สํ 2. 278 ปิฏฺเฐ.} อตฺเถว สุตฺตนฺโตติ, อามนฺตา.\csromanspacer{} เตน หิ น วตฺตพฺพํ “อตีตํ อตฺถี”ติ.}

\prose{\csromansymbolmark{*}อตีตํ อตฺถีติ, อามนฺตา.\csromanspacer{} นนุ อายสฺมา นนฺทโก เอตทโวจ\csromandash{}อหุ ปุพฺเพ โลโภ, ตทหุ อกุสลํ, โส เอตรหิ นตฺถิ, อิจฺเจตํ กุสลํ.\csromanspacer{} อหุ ปุพฺเพ โทโส.\csromanspacer{} อหุ ปุพฺเพ โมโห, ตทหุ อกุสลํ, โส เอตรหิ นตฺถิ, อิจฺเจตํ กุสลนฺติ\footnote{อํ 1. 197 ปิฏฺเฐ สาฬฺหสุตฺเต.} อตฺเถว สุตฺตนฺโตติ, อามนฺตา.\csromanspacer{} เตน หิ น วตฺตพฺพํ “อตีตํ อตฺถี”ติ.}

\prose{\csromansymbolmark{+}น “วตฺตพฺพํ อนาคตํ อตฺถี”ติ, อามนฺตา.\csromanspacer{} นนุ วุตฺตํ ภควตา\csromandash{} กพฬีกาเร เจ ภิกฺขเว อาหาเร อตฺถิ ราโค, อตฺถิ นนฺที, อตฺถิ ตณฺหา, ปติฏฺฐิตํ ตตฺถ วิญฺญาณํ วิรูฬฺหํ.\csromanspacer{} ยตฺถ ปติฏฺฐิตํ วิญฺญาณํ วิรูฬฺหํ, อตฺถิ ตตฺถ นามรูปสฺส อวกฺกนฺติ.\csromanspacer{} ยตฺถ อตฺถิ นามรูปสฺส อวกฺกนฺติ, อตฺถิ ตตฺถ สงฺขารานํ วุทฺธิ.\csromanspacer{} ยตฺถ อตฺถิ สงฺขารานํ วุทฺธิ, อตฺถิ ตตฺถ อายติํ ปุนพฺภวาภินิพฺพตฺติ.\csromanspacer{} ยตฺถ อตฺถิ อายติํ ปุนพฺภวาภินิพฺพตฺติ, อตฺถิ ตตฺถ อายติํ ชาติชรามรณํ.\csromanspacer{} ยตฺถ อตฺถิ อายติํ ชาติชรามรณํ สโสกํ ตํ ภิกฺขเว สรชํ\footnote{สทรํ (สํ 1. 325 ปิฏฺเฐ) ตเทว ยุตฺตตรํ.} ส-อุปายาสนฺติ วทามิ.}

\prose{ผสฺเส เจ ภิกฺขเว อาหาเร.\csromanspacer{} มโนสญฺเจตนาย เจ ภิกฺขเว อาหาเร.\csromanspacer{} วิญฺญาเณ เจ ภิกฺขเว อาหาเร อตฺถิ ราโค, อตฺถิ นนฺที ฯเปฯ สรชํ ส- อุปายาสนฺติ วทามีติ\footnote{สํ 1. 325 ปิฏฺเฐ.} อตฺเถว สุตฺตนฺโตติ, อามนฺตา.\csromanspacer{} เตน หิ อนาคตํ อตฺถีติ.}

\prose{\csromansymbolmark{*}อนาคตํ อตฺถีติ, อามนฺตา.\csromanspacer{} นนุ วุตฺตํ ภควตา\csromandash{}กพฬีกาเร เจ ภิกฺขเว อาหาเร นตฺถิ ราโค, นตฺถิ นนฺที, นตฺถิ ตณฺหา, อปฺปติฏฺฐิตํ ตตฺถ วิญฺญาณํ อวิรูฬฺหํ.\csromanspacer{} ยตฺถ วิญฺญาณํ อปฺปติฏฺฐิตํ อวิรูฬฺหํ, นตฺถิ ตตฺถ นามรูปสฺส อวกฺกนฺติ.\csromanspacer{} ยตฺถ นตฺถิ นามรูปสฺส อวกฺกนฺติ, นตฺถิ ตตฺถ สงฺขารานํ วุทฺธิ.\csromanspacer{} ยตฺถ นตฺถิ สงฺขารานํ วุทฺธิ, นตฺถิ ตตฺถ อายติํ ปุนพฺภวาภินิพฺพตฺติ.}

\vspace{\tipitakaparskip}
\csromancenter{อตีตกฺขนฺธาทิกถา ยตฺถ นตฺถิ อายติํ ปุนพฺภวาภินิพฺพตฺติ, นตฺถิ ตตฺถ อายติํ ชาติชรามรณํ.\csromanspacer{} ยตฺถ นตฺถิ อายติํ ชาติชรามรณํ, อโสกํ ตํ ภิกฺขเว อรชํ อนุปายาสนฺติ วทามิ.}
\prose{\csromanfolio{115}ผสฺเส เจ ภิกฺขเว อาหาเร.\csromanspacer{} มโนสญฺเจตนาย เจ ภิกฺขเว อาหาเร.\csromanspacer{} วิญฺญาเณ เจ ภิกฺขเว อาหาเร นตฺถิ ราโค, นตฺถิ นนฺที ฯเปฯ อรชํ อนุปายาสนฺติ วทามีติ\footnote{สํ 1. 326 ปิฏฺเฐ.} อตฺเถว สุตฺตนฺโตติ, อามนฺตา.\csromanspacer{} เตน หิ น วตฺตพฺพํ “อนาคตํ อตฺถี”ติ.}

\nitthitam{สพฺพมตฺถีติกถา นิฏฺฐิตา.}
\csromanmatikaanchor{mtk.25}
\csromantocmarkref{chapter}{6. อตีตกฺขนฺธาทิกถา}{mtk.25}
\bookmark[dest={mtk.25},level=0]{6. อตีตกฺขนฺธาทิกถา}
\chapterhead{6. อตีตกฺขนฺธาทิกถา}
\csromanheader{2}{1. นสุตฺตสาธน}
\proseitem{297}{\csromansymbolmark{+}อตีตํ ขนฺธาติ, อามนฺตา.\csromanspacer{} อตีตํ อตฺถีติ, น เหวํ วตฺตพฺเพ ฯเปฯ.\csromanspacer{} อตีตํ อายตนนฺติ, อามนฺตา.\csromanspacer{} อตีตํ อตฺถีติ, น เหวํ วตฺตพฺเพ ฯเปฯ.\csromanspacer{} อตีตํ ธาตูติ, อามนฺตา.\csromanspacer{} อตีตํ อตฺถีติ, น เหวํ วตฺตพฺเพ ฯเปฯ.\csromanspacer{} อตีตํ ขนฺธา ธาตุ อายตนนฺติ\footnote{ขนฺธธาตุ-อายตนนฺติ (สฺยา)}, อามนฺตา.\csromanspacer{} อตีตํ อตฺถีติ, น เหวํ วตฺตพฺเพ ฯเปฯ.}

\prose{\csromansymbolmark{+}อนาคตํ ขนฺธาติ, อามนฺตา.\csromanspacer{} อนาคตํ อตฺถีติ, น เหวํ วตฺตพฺเพ ฯเปฯ.\csromanspacer{} อนาคตํ อายตนนฺติ, อามนฺตา.\csromanspacer{} อนาคตํ อตฺถีติ, น เหวํ วตฺตพฺเพ ฯเปฯ.\csromanspacer{} อนาคตํ ธาตูติ, อามนฺตา.\csromanspacer{} อนาคตํ อตฺถีติ, น เหวํ วตฺตพฺเพ ฯเปฯ.\csromanspacer{} อนาคตํ ขนฺธา ธาตุ อายตนนฺติ, อามนฺตา.\csromanspacer{} อนาคตํ อตฺถีติ, น เหวํ วตฺตพฺเพ ฯเปฯ.}

\prose{\csromansymbolmark{+}ปจฺจุปฺปนฺนํ ขนฺธา ปจฺจุปฺปนฺนํ อตฺถีติ, อามนฺตา.\csromanspacer{} อตีตํ ขนฺธา อตีตํ อตฺถีติ, น เหวํ วตฺตพฺเพ ฯเปฯ.\csromanspacer{} ปจฺจุปฺปนฺนํ อายตนํ ปจฺจุปฺปนฺนํ อตฺถีติ, อามนฺตา.\csromanspacer{} อตีตํ อายตนํ อตีตํ อตฺถีติ, น เหวํ วตฺตพฺเพ ฯเปฯ.\csromanspacer{} ปจฺจุปฺปนฺนํ ธาตุ ปจฺจุปฺปนฺนํ อตฺถีติ, อามนฺตา.\csromanspacer{} อตีตํ ธาตุ อตีตํ อตฺถีติ, น เหวํ วตฺตพฺเพ ฯเปฯ.\csromanspacer{} ปจฺจุปฺปนฺนํ ขนฺธา ธาตุ อายตนํ ปจฺจุปฺปนฺนํ อตฺถีติ, อามนฺตา.\csromanspacer{} อตีตํ ขนฺธา ธาตุ อายตนํ อตีตํ อตฺถีติ, น เหวํ วตฺตพฺเพ ฯเปฯ.}

\prose{\csromanfolio{116}\csromansymbolmark{+}ปจฺจุปฺปนฺนํ ขนฺธา ปจฺจุปฺปนฺนํ อตฺถีติ, อามนฺตา.\csromanspacer{} อนาคตํ ขนฺธา อนาคตํ อตฺถีติ, น เหวํ วตฺตพฺเพ ฯเปฯ.\csromanspacer{} ปจฺจุปฺปนฺนํ อายตนํ ปจฺจุปฺปนฺนํ อตฺถีติ, อามนฺตา.\csromanspacer{} อนาคตํ อายตนํ อนาคตํ อตฺถีติ, น เหวํ วตฺตพฺเพ ฯเปฯ.\csromanspacer{} ปจฺจุปฺปนฺนํ ธาตุ ปจฺจุปฺปนฺนํ อตฺถีติ, อามนฺตา.\csromanspacer{} อนาคตํ ธาตุ อนาคตํ อตฺถีติ, น เหวํ วตฺตพฺเพ ฯเปฯ.\csromanspacer{} ปจฺจุปฺปนฺนํ ขนฺธา ธาตุ อายตนํ ปจฺจุปฺปนฺนํ อตฺถีติ, อามนฺตา, อนาคตํ ขนฺธา ธาตุ อายตนํ อนาคตํ อตฺถีติ, น เหวํ วตฺตพฺเพ ฯเปฯ.}

\prose{\csromansymbolmark{+}อตีตํ ขนฺธา อตีตํ นตฺถีติ, อามนฺตา.\csromanspacer{} ปจฺจุปฺปนฺนํ ขนฺธา ปจฺจุปฺปนฺนํ นตฺถีติ, น เหวํ วตฺตพฺเพ ฯเปฯ.\csromanspacer{} อตีตํ อายตนํ อตีตํ นตฺถีติ, อามนฺตา.\csromanspacer{} ปจฺจุปฺปนฺนํ อายตนํ ปจฺจุปฺปนฺนํ นตฺถีติ, น เหวํ วตฺตพฺเพ ฯเปฯ.\csromanspacer{} อตีตํ ธาตุ อตีตํ นตฺถีติ, อามนฺตา.\csromanspacer{} ปจฺจุปฺปนฺนํ ธาตุ ปจฺจุปฺปนฺนํ นตฺถีติ, น เหวํ วตฺตพฺเพ ฯเปฯ.\csromanspacer{} อตีตํ ขนฺธา ธาตุ อายตนํ อตีตํ นตฺถีติ, อามนฺตา.\csromanspacer{} ปจฺจุปฺปนฺนํ ขนฺธา ธาตุ อายตนํ ปจฺจุปฺปนฺนํ นตฺถีติ, น เหวํ วตฺตพฺเพ ฯเปฯ.}

\prose{\csromansymbolmark{+}อนาคตํ ขนฺธา อนาคตํ นตฺถีติ, อามนฺตา.\csromanspacer{} ปจฺจุปฺปนฺนํ ขนฺธา ปจฺจุปฺปนฺนํ นตฺถีติ, น เหวํ วตฺตพฺเพ ฯเปฯ.\csromanspacer{} อนาคตํ อายตนํ ฯเปฯ.\csromanspacer{} อนาคตํ ธาตุ ฯเปฯ.\csromanspacer{} อนาคตํ ขนฺธา ธาตุ อายตนํ อนาคตํ นตฺถีติ, อามนฺตา.\csromanspacer{} ปจฺจุปฺปนฺนํ ขนฺธา ธาตุ อายตนํ ปจฺจุปฺปนฺนํ นตฺถีติ, น เหวํ วตฺตพฺเพ ฯเปฯ.}

\prose{\csromansymbolmark{+}อตีตํ รูปํ ขนฺโธติ, อามนฺตา.\csromanspacer{} อตีตํ รูปํ อตฺถีติ, น เหวํ วตฺตพฺเพ ฯเปฯ.\csromanspacer{} อตีตํ รูปํ อายตนนฺติ, อามนฺตา.\csromanspacer{} อตีตํ รูปํ อตฺถีติ, น เหวํ วตฺตพฺเพ ฯเปฯ.\csromanspacer{} อตีตํ รูปํ ธาตูติ, อามนฺตา.\csromanspacer{} อตีตํ รูปํ อตฺถีติ, น เหวํ วตฺตพฺเพ ฯเปฯ.\csromanspacer{} อตีตํ รูปํ ขนฺธา ธาตุ อายตนนฺติ, อามนฺตา.\csromanspacer{} อตีตํ รูปํ อตฺถีติ, น เหวํ วตฺตพฺเพ ฯเปฯ.}

\prose{\csromansymbolmark{+}อนาคตํ รูปํ ขนฺโธติ, อามนฺตา.\csromanspacer{} อนาคตํ รูปํ อตฺถีติ, น เหวํ วตฺตพฺเพ ฯเปฯ.\csromanspacer{} อนาคตํ รูปํ อายตนํ ฯเปฯ.\csromanspacer{} อนาคตํ รูปํ ธาตุ ฯเปฯ.\csromanspacer{} อนาคตํ รูปํ ขนฺธา ธาตุ อายตนนฺติ, อามนฺตา.\csromanspacer{} อนาคตํ รูปํ อตฺถีติ, น เหวํ วตฺตพฺเพ ฯเปฯ.}

\prose{\csromansymbolmark{+}ปจฺจุปฺปนฺนํ รูปํ ขนฺโธ ปจฺจุปฺปนฺนํ รูปํ อตฺถีติ, อามนฺตา.\csromanspacer{} อตีตํ รูปํ ขนฺโธ อตีตํ รูปํ อตฺถีติ, น เหวํ วตฺตพฺเพ ฯเปฯ.\csromanspacer{} ปจฺจุปฺปนฺนํ รูปํ อายตนํ ฯเปฯ.\csromanspacer{} ปจฺจุปฺปนฺนํ รูปํ ธาตุ ฯเปฯ.\csromanspacer{} ปจฺจุปฺปนฺนํ รูปํ ขนฺธา ธาตุ อายตนํ ปจฺจุปฺปนฺนํ รูปํ อตฺถีติ, อามนฺตา.\csromanspacer{} อตีตํ รูปํ ขนฺธา ธาตุ อายตนํ อตีตํ รูปํ อตฺถีติ, น เหวํ วตฺตพฺเพ ฯเปฯ.}

\csromanheader{2}{6. อตีตกฺขนฺธาทิกถา}
\prose{\csromanfolio{117}\csromansymbolmark{+}ปจฺจุปฺปนฺนํ รูปํ ขนฺโธ ปจฺจุปฺปนฺนํ รูปํ อตฺถีติ, อามนฺตา.\csromanspacer{} อนาคตํ รูปํ ขนฺโธ อนาคตํ รูปํ อตฺถีติ, น เหวํ วตฺตพฺเพ ฯเปฯ.\csromanspacer{} ปจฺจุปฺปนฺนํ อายตนํ ฯเปฯ.\csromanspacer{} ปจฺจุปฺปนฺนํ รูปํ ธาตุ ฯเปฯ.\csromanspacer{} ปจฺจุปฺปนฺนํ รูปํ ขนฺธา ธาตุ อายตนํ ปจฺจุปฺปนฺนํ รูปํ อตฺถีติ, อามนฺตา.\csromanspacer{} อนาคตํ รูปํ ขนฺธา ธาตุ อายตนํ อนาคตํ รูปํ อตฺถีติ, น เหวํ วตฺตพฺเพ ฯเปฯ.}

\prose{\csromansymbolmark{+}อตีตํ รูปํ ขนฺโธ อตีตํ รูปํ นตฺถีติ, อามนฺตา.\csromanspacer{} ปจฺจุปฺปนฺนํ รูปํ ขนฺโธ ปจฺจุปฺปนฺนํ รูปํ นตฺถีติ, น เหวํ วตฺตพฺเพ ฯเปฯ.\csromanspacer{} อตีตํ รูปํ อายตนํ ฯเปฯ.\csromanspacer{} อตีตํ รูปํ ธาตุ ฯเปฯ.\csromanspacer{} อตีตํ รูปํ ขนฺธา ธาตุ อายตนํ อตีตํ รูปํ นตฺถีติ, อามนฺตา.\csromanspacer{} ปจฺจุปฺปนฺนํ รูปํ ขนฺธา ธาตุ อายตนํ ปจฺจุปฺปนฺนํ รูปํ นตฺถีติ, น เหวํ วตฺตพฺเพ ฯเปฯ.}

\prose{\csromansymbolmark{+}อนาคตํ รูปํ ขนฺโธ อนาคตํ รูปํ นตฺถีติ, อามนฺตา.\csromanspacer{} ปจฺจุปฺปนฺนํ รูปํ ขนฺโธ ปจฺจุปฺปนฺนํ รูปํ นตฺถีติ, น เหวํ วตฺตพฺเพ ฯเปฯ.\csromanspacer{} อนาคตํ รูปํ อายตนํ ฯเปฯ.\csromanspacer{} อนาคตํ รูปํ ธาตุ ฯเปฯ.\csromanspacer{} อนาคตํ รูปํ ขนฺธา ธาตุ อายตนํ อนาคตํ รูปํ นตฺถีติ, อามนฺตา.\csromanspacer{} ปจฺจุปฺปนฺนํ รูปํ ขนฺธา ธาตุ อายตนํ ปจฺจุปฺปนฺนํ รูปํ นตฺถีติ, น เหวํ วตฺตพฺเพ ฯเปฯ.}

\prose{\csromansymbolmark{+}อตีตา เวทนา.\csromanspacer{} อตีตา สญฺญา.\csromanspacer{} อตีตา สงฺขารา.\csromanspacer{} อตีตํ วิญฺญาณํ ขนฺโธติ, อามนฺตา.\csromanspacer{} อตีตํ วิญฺญาณํ อตฺถีติ, น เหวํ วตฺตพฺเพ ฯเปฯ.\csromanspacer{} อตีตํ วิญฺญานฺ−ํ อายตนํ ฯเปฯ.\csromanspacer{} อตีตํ วิญฺญาณํ ธาตุ ฯเปฯ.\csromanspacer{} อตีตํ วิญฺญาณํ ขนฺธา ธาตุ อายตนนฺติ, อามนฺตา.\csromanspacer{} อตีตํ วิญฺญาณํ อตฺถีติ, น เหวํ วตฺตพฺเพ ฯเปฯ.}

\prose{\csromansymbolmark{+}อนาคตํ วิญฺญาณํ ขนฺโธติ, อามนฺตา.\csromanspacer{} อนาคตํ วิญฺญาณํ อตฺถีติ, น เหวํ วตฺตพฺเพ ฯเปฯ.\csromanspacer{} อนาคตํ วิญฺญาณํ อายตนํ ฯเปฯ.\csromanspacer{} อนาคตํ วิญฺญาณํ ธาตุ ฯเปฯ.\csromanspacer{} อนาคตํ วิญฺญาณํ ขนฺธา ธาตุ อายตนนฺติ, อามนฺตา.\csromanspacer{} อนาคตํ วิญฺญาณํ อตฺถีติ, น เหวํ วตฺตพฺเพ ฯเปฯ.}

\prose{\csromansymbolmark{+}ปจฺจุปฺปนฺนํ วิญฺญาณํ ขนฺโธ ปจฺจุปฺปนฺนํ วิญฺญาณํ อตฺถีติ, อามนฺตา.\csromanspacer{} อตีตํ วิญฺญาณํ ขนฺโธ อตีตํ วิญฺญาณํ อตฺถีติ, น เหวํ วตฺตพฺเพ ฯเปฯ.\csromanspacer{} ปจฺจุปฺปนฺนํ วิญฺญาณํ อายตนํ ฯเปฯ.\csromanspacer{} ปจฺจุปฺปนฺนํ วิญฺญาณํ ธาตุ ฯเปฯ.\csromanspacer{} ปจฺจุปฺปนฺนํ วิญฺญาณํ ขนฺธา ธาตุ อายตนํ ปจฺจุปฺปนฺนํ วิญฺญาณํ อตฺถีติ, อามนฺตา.\csromanspacer{} อตีตํ วิญฺญาณํ ขนฺธา ธาตุ อายตนํ อตีตํ วิญฺญาณํ อตฺถีติ, น เหวํ วตฺตพฺเพ ฯเปฯ.}

\prose{\csromanfolio{118}\csromansymbolmark{+}ปจฺจุปฺปนฺนํ วิญฺญาณํ ขนฺโธ ปจฺจุปฺปนฺนํ วิญฺญานฺ−ํ อตฺถีติ, อามนฺตา.\csromanspacer{} อนาคตํ วิญฺญาณํ ขนฺโธ อนาคตํ วิญฺญาณํ อตฺถีติ, น เหวํ วตฺตพฺเพ ฯเปฯ.\csromanspacer{} ปจฺจุปฺปนฺนํ วิญฺญาณํ อายตนํ ฯเปฯ.\csromanspacer{} ปจฺจุปฺปนฺนํ วิญฺญาณํ ธาตุ ฯเปฯ.\csromanspacer{} ปจฺจุปฺปนฺนํ วิญฺญาณํ ขนฺธา ธาตุ อายตนํ ปจฺจุปฺปนฺนํ วิญฺญาณํ อตฺถีติ, อามนฺตา.\csromanspacer{} อนาคตํ วิญฺญาณํ ขนฺธา ธาตุ อายตนํ อนาคตํ วิญฺญาณํ อตฺถีติ, น เหวํ วตฺตพฺเพ ฯเปฯ.}

\prose{\csromansymbolmark{+}อตีตํ วิญฺญาณํ ขนฺโธ อตีตํ วิญฺญาณํ นตฺถีติ, อามนฺตา.\csromanspacer{} ปจฺจุปฺปนฺนํ วิญฺญาณํ ขนฺโธ ปจฺจุปฺปนฺนํ วิญฺญาณํ นตฺถีติ, น เหวํ วตฺตพฺเพ ฯเปฯ.\csromanspacer{} อตีตํ วิญฺญาณํ อายตนํ ฯเปฯ.\csromanspacer{} อตีตํ วิญฺญาณํ ธาตุ ฯเปฯ.\csromanspacer{} อตีตํ วิญฺญาณํ ขนฺธา ธาตุ อายตนํ อตีตํ วิญฺญาณํ นตฺถีติ, อามนฺตา.\csromanspacer{} ปจฺจุปฺปนฺนํ วิญฺญาณํ ขนฺธา ธาตุ อายตนํ ปจฺจุปฺปนฺนํ วิญฺญาณํ นตฺถีติ, น เหวํ วตฺตพฺเพ ฯเปฯ.\csromanspacer{} อนาคตํ วิญฺญาณํ ขนฺโธ อนาคตํ วิญฺญาณํ นตฺถีติ, อามนฺตา.\csromanspacer{} ปจฺจุปฺปนฺนํ วิญฺญาณํ ขนฺโธ ปจฺจุปฺปนฺนํ วิญฺญาณํ นตฺถีติ, น เหวํ วตฺตพฺเพ ฯเปฯ.\csromanspacer{} อนาคตํ วิญฺญาณํ อายตนํ ฯเปฯ.\csromanspacer{} อนาคตํ วิญฺญาณํ ธาตุ ฯเปฯ.\csromanspacer{} อนาคตํ วิญฺญาณํ ขนฺธา ธาตุ อายตนํ อนาคตํ วิญฺญาณํ นตฺถีติ, อามนฺตา.\csromanspacer{} ปจฺจุปฺปนฺนํ วิญฺญาณํ ขนฺธา ธาตุ อายตนํ ปจฺจุปฺปนฺนํ วิญฺญาณํ นตฺถีติ, น เหวํ วตฺตพฺเพ ฯเปฯ.}

\csromanheader{2}{2. สุตฺตสาธน}
\proseitem{298}{\csromansymbolmark{*}น วตฺตพฺพํ “อตีตานาคตา ขนฺธา ธาตุ อายตนํ นตฺถิ เจ’เต”ติ, อามนฺตา.\csromanspacer{} นนุ วุตฺตํ ภควตา “ตโยเม ภิกฺขเว นิรุตฺติปถา อธิวจนปถา ปญฺญตฺติ ฯเปฯ วิญฺญูหีติ ฯเปฯ” อตฺเถว สุตฺตนฺโตติ, อามนฺตา.\csromanspacer{} เตน หิ น วตฺตพฺพํ “อตีตานาคตา ขนฺธา ธาตุ อายตนํ นตฺถิ เจ’เต”ติ.}

\prose{\csromansymbolmark{+}อตีตานาคตา ขนฺธา ธาตุ อายตนํ นตฺถิ เจเตติ, อามนฺตา.\csromanspacer{} นนุ วุตฺตํ ภควตา “ยํ กิญฺจิ ภิกฺขเว รูปํ อตีตานาคตปจฺจุปฺปนฺนํ อชฺฌตฺตํ วา พหิทฺธา วา โอฬาริกํ วา สุขุมํ วา หีนํ วา ปณีตํ วา ยํ ทูเร สนฺติเก วา, อยํ วุจฺจติ รูปกฺขนฺโธ.\csromanspacer{} ยา กาจิ เวทนา.\csromanspacer{} ยา กาจิ สญฺญา.\csromanspacer{} เย เกจิ สงฺขารา.\csromanspacer{} ยํ กิญฺจิ วิญฺญาณํ อตีตานาคตปจฺจุปฺปนฺนํ ฯเปฯ อยํ วุจฺจติ วิญฺญาณกฺขนฺโธ”ติ อตฺเถว สุตฺตนฺโตติ, อามนฺตา.\csromanspacer{} เตน หิ น วตฺตพฺพํ “อตีตานาคตา ขนฺธา ธาตุ อายตนํ นตฺถิ เจ’เต”ติ.}

\nitthitam{อตีตกฺขนฺธาทิกถา นิฏฺฐิตา.}
\csromanmatikaanchor{mtk.26}
\csromantocmarkref{chapter}{7. เอกจฺจํอตฺถีติกถา}{mtk.26}
\bookmark[dest={mtk.26},level=0]{7. เอกจฺจํอตฺถีติกถา}
\chapterheadpageread{7. เอกจฺจํอตฺถีติกถา \textbf{7. เอกจฺจํอตฺถีติกถา}}
\csromanheader{2}{1. อตีตาทิ-เอกจฺจกถา}
\proseitem{299}{\csromanfolio{119}\csromansymbolmark{*}อตีตํ อตฺถีติ? เอกจฺจํ อตฺถิ, เอกจฺจํ นตฺถีติ.\csromanspacer{} เอกจฺจํ นิรุทฺธํ, เอกจฺจํ น นิรุทฺธํ.\csromanspacer{} เอกจฺจํ วิคตํ, เอกจฺจํ อวิคตํ.\csromanspacer{} เอกจฺจํ อตฺถงฺคตํ, เอกจฺจํ น อตฺถงฺคตํ.\csromanspacer{} เอกจฺจํ อพฺภตฺถงฺคตํ, เอกจฺจํ น อพฺภตฺถงฺคตนฺติ, น เหวํ วตฺตพฺเพ ฯเปฯ.}

\prose{\csromansymbolmark{*}อตีตํ เอกจฺจํ อตฺถิ, เอกจฺจํ นตฺถีติ, อามนฺตา.\csromanspacer{} อตีตา อวิปกฺกวิปากา ธมฺมา เอกจฺเจ อตฺถิ, เอกจฺเจ นตฺถีติ, น เหวํ วตฺตพฺเพ ฯเปฯ.\csromanspacer{} อตีตํ เอกจฺจํ อตฺถิ, เอกจฺจํ นตฺถีติ, อามนฺตา.\csromanspacer{} อตีตา วิปกฺกวิปากา ธมฺมา เอกจฺเจ อตฺถิ, เอกจฺเจ นตฺถีติ, น เหวํ วตฺตพฺเพ ฯเปฯ.\csromanspacer{} อตีตํ เอกจฺจํ อตฺถิ, เอกจฺจํ นตฺถีติ, อามนฺตา.\csromanspacer{} อตีตา อวิปากา ธมฺมา เอกจฺเจ อตฺถิ, เอกจฺเจ นตฺถีติ, น เหวํ วตฺตพฺเพ ฯเปฯ.}

\prose{\csromansymbolmark{*}อตีตํ เอกจฺจํ อตฺถิ เอกจฺจํ นตฺถีติ, อามนฺตา.\csromanspacer{} กิํ อตฺถิ, กิํ นตฺถีติ, อตีตา อวิปกฺกวิปากา ธมฺมา เต อตฺถิ, อตีตา วิปกฺกวิปากา ธมฺมา เต นตฺถีติ.\csromanspacer{} อตีตา อวิปกฺกวิปากา ธมฺมา เต อตฺถีติ, อามนฺตา.\csromanspacer{} อตีตา วิปกฺกวิปากา ธมฺมา เต อตฺถีติ, น เหวํ วตฺตพฺเพ ฯเปฯ.\csromanspacer{} อตีตา อวิปกฺกวิปากา ธมฺมา เต อตฺถีติ, อามนฺตา.\csromanspacer{} อตีตา อวิปากา ธมฺมา เต อตฺถีติ\footnote{อตีตา อวิปกฺกวิปากา ธมฺมา เต นตฺถีติ (ก)}, น เหวํ วตฺตพฺเพ ฯเปฯ.}

\prose{\csromansymbolmark{*}อตีตา วิปกฺกวิปากา ธมฺมา เต นตฺถีติ, อามนฺตา.\csromanspacer{} อตีตา อวิปกฺกวิปากา ธมฺมา เต นตฺถีติ, น เหวํ วตฺตพฺเพ ฯเปฯ.}

\prose{\csromansymbolmark{*}อตีตา อวิปากา\footnote{วิปกฺกวิปากา (สฺยา), อวิปกฺกวิปากา (ก)} ธมฺมา เต นตฺถีติ, อามนฺตา.\csromanspacer{} อตีตา อวิปกฺกวิปากา\footnote{อวิปากา (สฺยา)} ธมฺมา เต นตฺถีติ.\csromanspacer{} น เหวํ วตฺตพฺเพ ฯเปฯ.}

\prose{\csromansymbolmark{*}อตีตา อวิปกฺกวิปากา ธมฺมา เต อตฺถีติ, อามนฺตา.\csromanspacer{} นนุ อตีตา อวิปกฺกวิปากา ธมฺมา นิรุทฺธาติ, อามนฺตา.\csromanspacer{} หญฺจิ อตีตา อวิปกฺกวิปากา ธมฺมา นิรุทฺธา, โน จ วต เร วตฺตพฺเพ “อตีตา อวิปกฺกวิปากา ธมฺมา เต\footnote{ธมฺมา นิรุทฺธา เต (สฺยา, ก)} อตฺถี”ติ.}

\prose{\csromanfolio{120}\csromansymbolmark{*}อตีตา อวิปกฺกวิปากา ธมฺมา นิรุทฺธา เต อตฺถีติ, อามนฺตา.\csromanspacer{} อตีตา วิปกฺกวิปากา ธมฺมา นิรุทฺธา เต อตฺถีติ, น เหวํ วตฺตพฺเพ ฯเปฯ.\csromanspacer{} อตีตา อวิปกฺกวิปากา ธมฺมา นิรุทฺธา เต อตฺถีติ, อามนฺตา.\csromanspacer{} อตีตา อวิปากา\footnote{อวิปกฺกวิปากา (สี, ก)} ธมฺมา นิรุทฺธา เต อตฺถีติ, น เหวํ วตฺตพฺเพ ฯเปฯ.}

\prose{\csromansymbolmark{*}อตีตา วิปกฺกวิปากา ธมฺมา นิรุทฺธา เต นตฺถีติ, อามนฺตา.\csromanspacer{} อตีตา อวิปกฺกวิปากา ธมฺมา นิรุทฺธา เต นตฺถีติ, น เหวํ วตฺตพฺเพ ฯเปฯ.}

\prose{\csromansymbolmark{*}อตีตา อวิปากา\footnote{วิปกฺกวิปากา (สฺยา), อวิปกฺกวิปากา (ก)} ธมฺมา นิรุทฺธา เต นตฺถีติ, อามนฺตา.\csromanspacer{} อตีตา อวิปกฺกวิปากา\footnote{อวิปากา (สฺยา)} ธมฺมา นิรุทฺธา เต นตฺถีติ, น เหวํ วตฺตพฺเพ ฯเปฯ.}

\prose{\csromansymbolmark{*}อตีตา อวิปกฺกวิปากา ธมฺมา นิรุทฺธา เต อตฺถีติ, อามนฺตา.\csromanspacer{} อตีตา วิปกฺกวิปากา ธมฺมา นิรุทฺธา เต นตฺถีติ, อามนฺตา.\csromanspacer{} อตีตา เอกเทสํ วิปกฺกวิปากา ธมฺมา เอกเทสํ อวิปกฺกวิปากา ธมฺมา นิรุทฺธา เต เอกจฺเจ อตฺถิ, เอกจฺเจ นตฺถีติ, น เหวํ วตฺตพฺเพ ฯเปฯ.}

\prose{\csromansymbolmark{+}น วตฺตพฺพํ “อตีตา อวิปกฺกวิปากา ธมฺมา เต อตฺถี”ติ, อามนฺตา.\csromanspacer{} นนุ อตีตา อวิปกฺกวิปากา ธมฺมา วิปจฺจิสฺสนฺตีติ, อามนฺตา.\csromanspacer{} หญฺจิ อตีตา อวิปกฺกวิปากา ธมฺมา วิปจฺจิสฺสนฺติ, เตน วต เร วตฺตพฺเพ “อตีตา อวิปกฺกวิปากา ธมฺมา เต อตฺถี”ติ}

\prose{\csromansymbolmark{*}อตีตา อวิปกฺกวิปากา ธมฺมา วิปจฺจิสฺสนฺตีติ กตฺวา เต อตฺถีติ, อามนฺตา.\csromanspacer{} วิปจฺจิสฺสนฺตีติ กตฺวา ปจฺจุปฺปนฺนาติ, น เหวํ วตฺตพฺเพ ฯเปฯ.\csromanspacer{} วิปจฺจิสฺสนฺตีติ กตฺวา ปจฺจุปฺปนฺนาติ, อามนฺตา.\csromanspacer{} ปจฺจุปฺปนฺนา ธมฺมา นิรุชฺฌิสฺสนฺตีติ กตฺวา เต นตฺถีติ, น เหวํ วตฺตพฺเพ ฯเปฯ.}

\csromanheader{2}{2. อนาคตาทิ-เอกจฺจกถา}
\proseitem{300}{อนาคตํ อตฺถีติ? เอกจฺจํ อตฺถิ เอกจฺจํ นตฺถีติ.\csromanspacer{} เอกจฺจํ ชาตํ, เอกจฺจํ อชาตํ.\csromanspacer{} เอกจฺจํ สญฺชาตํ, เอกจฺจํ อสญฺชาตํ.\csromanspacer{} เอกจฺจํ นิพฺพตฺตํ, เอกจฺจํ อนิพฺพตฺตํ.\csromanspacer{} เอกจฺจํ ปาตุภูตํ, เอกจฺจํ อปาตุภูตนฺติ, น เหวํ วตฺตพฺเพ ฯเปฯ.}

\prose{อนาคตํ เอกจฺจํ อตฺถิ, เอกจฺจํ นตฺถีติ, อามนฺตา.\csromanspacer{} อนาคตา อุปฺปาทิโน ธมฺมา เอกจฺเจ อตฺถิ, เอกจฺเจ นตฺถีติ, น เหวํ วตฺตพฺเพ ฯเปฯ.}

\vspace{\tipitakaparskip}
\csromancenter{สติปฏฺฐานกถา อนาคตํ เอกจฺจํ อตฺถิ, เอกจฺจํ นตฺถีติ, อามนฺตา.\csromanspacer{} อนาคตา อนุปฺปาทิโน ธมฺมา เอกจฺเจ อตฺถิ, เอกจฺเจ นตฺถีติ, น เหวํ วตฺตพฺเพ ฯเปฯ.}
\prosecont{\csromanfolio{121}\csromansymbolmark{*}อนาคตํ เอกจฺจํ อตฺถิ, เอกจฺจํ นตฺถีติ, อามนฺตา.\csromanspacer{} กิํ อตฺถิ.\csromanspacer{} กิํ นตฺถีติ, อนาคตา อุปฺปาทิโน ธมฺมา เต อตฺถิ, อนาคตา อนุปฺปาทิโน ธมฺมา เต นตฺถีติ.\csromanspacer{} อนาคตา อุปฺปาทิโน ธมฺมา เต อตฺถีติ, อามนฺตา.\csromanspacer{} อนาคตา อนุปฺปาทิโน ธมฺมา เต อตฺถีติ, น เหวํ วตฺตพฺเพ ฯเปฯ.\csromanspacer{} อนาคตา อนุปฺปาทิโน ธมฺมา เต นตฺถีติ, อามนฺตา.\csromanspacer{} อนาคตา อุปฺปาทิโน ธมฺมา เต นตฺถีติ, น เหวํ วตฺตพฺเพ ฯเปฯ. \csromansymbolmark{*}อนาคตา อุปฺปาทิโน ธมฺมา เต อตฺถีติ, อามนฺตา.\csromanspacer{} นนุ อนาคตา อุปฺปาทิโน ธมฺมา อชาตาติ, อามนฺตา.\csromanspacer{} หญฺจิ อนาคตา อุปฺปาทิโน ธมฺมา อชาตา, โน จ วต เร วตฺตพฺเพ “อนาคตา อุปฺปาทิโน ธมฺมา เต อตฺถี”ติ.}

\prose{\csromansymbolmark{*}อนาคตา อุปฺปาทิโน ธมฺมา อชาตา เต อตฺถีติ, อามนฺตา.\csromanspacer{} อนาคตา อนุปฺปาทิโน ธมฺมา อชาตา เต อตฺถีติ, น เหวํ วตฺตพฺเพ ฯเปฯ.\csromanspacer{} อนาคตา อนุปฺปาทิโน ธมฺมา อชาตา เต นตฺถีติ, อามนฺตา.\csromanspacer{} อนาคตา อุปฺปาทิโน ธมฺมา อชาตา เต นตฺถีติ, น เหวํ วตฺตพฺเพ ฯเปฯ.}

\prose{\csromansymbolmark{+}น วตฺตพฺพํ “อนาคตา อุปฺปาทิโน ธมฺมา เต อตฺถี”ติ, อามนฺตา.\csromanspacer{} นนุ อนาคตา อุปฺปาทิโน ธมฺมา อุปฺปชฺชิสฺสนฺตีติ, อามนฺตา.\csromanspacer{} หญฺจิ อนาคตา อุปฺปาทิโน ธมฺมา อุปฺปชฺชิสฺสนฺติ, เตน วต เร วตฺตพฺเพ “อนาคตา อุปฺปาทิโน ธมฺมา เต อตฺถี”ติ.}

\prose{\csromansymbolmark{*}อนาคตา อุปฺปาทิโน ธมฺมา อุปฺปชฺชิสฺสนฺตีติ กตฺวา เต อตฺถีติ, อามนฺตา.\csromanspacer{} อุปฺปชฺชิสฺสนฺตีติ กตฺวา ปจฺจุปฺปนฺนาติ, น เหวํ วตฺตพฺเพ ฯเปฯ.\csromanspacer{} อุปฺปชฺชิสฺสนฺตีติ กตฺวา ปจฺจุปฺปนฺนาติ, อามนฺตา.\csromanspacer{} ปจฺจุปฺปนฺนา ธมฺมา นิรุชฺฌิสฺสนฺตีติ กตฺวา เต นตฺถีติ, น เหวํ วตฺตพฺเพ ฯเปฯ.}

\nitthitam{เอกจฺจํ อตฺถีติกถา นิฏฺฐิตา.}
\csromanmatikaanchor{mtk.27}
\csromantocmarkref{chapter}{8. สติปฏฺฐานกถา}{mtk.27}
\bookmark[dest={mtk.27},level=0]{8. สติปฏฺฐานกถา}
\chapterhead{8. สติปฏฺฐานกถา}
\proseitem{301}{\csromansymbolmark{*}สพฺเพ ธมฺมา สติปฏฺฐานาติ, อามนฺตา.\csromanspacer{} สพฺเพ ธมฺมา สติ, สตินฺทฺริยํ, สติพลํ, สมฺมาสติ, สติสมฺโพชฺฌงฺโค, เอกายนมคฺโค, ขยคามี \csromanfolio{122}โพธคามี.\csromanspacer{} อปจยคามี, อนาสวา, อสํโยชนิยา, อคนฺถนิยา, อโนฆนิยา, อโยคนิยา, อนีวรณิยา, อปรามฏฺฐา, อนุปาทานิยา, อสํกิเลสิกา, สพฺเพ ธมฺมา พุทฺธานุสฺสติ, ธมฺมานุสฺสติ, สํฆานุสฺสติ, สีลานุสฺสติ, จาคานุสฺสติ, เทวตานุสฺสติ, อานาปานสฺสติ, มรณานุสฺสติ, กายคตาสติ, อุปสมานุสฺสตีติ, น เหวํ วตฺตพฺเพ ฯเปฯ.}

\prose{\csromansymbolmark{*}สพฺเพ ธมฺมา สติปฏฺฐานาติ, อามนฺตา.\csromanspacer{} จกฺขายตนํ สติปฏฺฐานนฺติ, น เหวํ วตฺตพฺเพ ฯเปฯ.\csromanspacer{} จกฺขายตนํ สติปฏฺฐานนฺติ, อามนฺตา.\csromanspacer{} จกฺขายตนํ สติ, สตินฺทฺริยํ, สติพลํ, สมฺมาสติ, สติสมฺโพชฺฌงฺโค, เอกายนมคฺโค, ขยคามี, โพธคามี, อปจยคามี, อนาสวํ, อสํโยชนิยํ ฯเปฯ อสํกิเลสิกํ, จกฺขายตนํ พุทฺธานุสฺสติ, ธมฺมานุสฺสติ, สํฆานุสฺสติ, สีลานุสฺสติ, จาคานุสฺสติ, เทวตานุสฺสติ, อานาปานสฺสติ, มรณานุสฺสติ, กายคตาสติ, อุปสมานุสฺสตีติ, น เหวํ วตฺตพฺเพ ฯเปฯ.\csromanspacer{} โสตายตนํ.\csromanspacer{} ฆานายตนํ.\csromanspacer{} ชิวฺหายตนํ.\csromanspacer{} กายายตนํ.\csromanspacer{} รูปายตนํ.\csromanspacer{} สทฺทายตนํ.\csromanspacer{} คนฺธายตนํ.\csromanspacer{} รสายตนํ.\csromanspacer{} โผฏฺฐพฺพายตนํ.\csromanspacer{} ราโค.\csromanspacer{} โทโส.\csromanspacer{} โมโห.\csromanspacer{} มาโน.\csromanspacer{} ทิฏฺฐิ.\csromanspacer{} วิจิกิจฺฉา.\csromanspacer{} ถินํ.\csromanspacer{} อุทฺธจฺจํ.\csromanspacer{} อหิริกํ.\csromanspacer{} อโนตฺตปฺปํ.\csromanspacer{} สติปฏฺฐานนฺติ, น เหวํ วตฺตพฺเพ ฯเปฯ.\csromanspacer{} อโนตฺตปฺปํ สติปฏฺฐานนฺติ, อามนฺตา.\csromanspacer{} อโนตฺตปฺปํ, สติ, สตินฺทฺริยํ, สติพลํ, สมฺมาสติ ฯเปฯ กายคตาสติ, อุปสมานุสฺสตีติ, น เหวํ วตฺตพฺเพ ฯเปฯ.}

\prose{\csromansymbolmark{*}สติ สติปฏฺฐานา, สา จ สตีติ, อามนฺตา.\csromanspacer{} จกฺขายตนํ สติปฏฺฐานํ, ตญฺจ สตีติ, น เหวํ วตฺตพฺเพ ฯเปฯ.\csromanspacer{} สติ สติปฏฺฐานา, สา จ สตีติ, อามนฺตา.\csromanspacer{} โสตายตนํ ฯเปฯ.\csromanspacer{} กายายตนํ.\csromanspacer{} รูปายตนํ ฯเปฯ.\csromanspacer{} โผฏฺฐพฺพายตนํ.\csromanspacer{} ราโค.\csromanspacer{} โทโส.\csromanspacer{} โมโห ฯเปฯ.\csromanspacer{} อโนตฺตปฺปํ สติปฏฺฐานํ, ตญฺจ สตีติ, น เหวํ วตฺตพฺเพ ฯเปฯ.}

\prose{\csromansymbolmark{*}จกฺขายตนํ สติปฏฺฐานํ, ตญฺจ น สตีติ, อามนฺตา.\csromanspacer{} สติ สติปฏฺฐานา, สา จ น สตีติ, น เหวํ วตฺตพฺเพ ฯเปฯ.\csromanspacer{} โสตายตนํ ฯเปฯ.\csromanspacer{} กายายตนํ.\csromanspacer{} รูปายตนํ ฯเปฯ.\csromanspacer{} โผฏฺฐพฺพายตนํ.\csromanspacer{} ราโค.\csromanspacer{} โทโส.\csromanspacer{} โมโห ฯเปฯ.\csromanspacer{} อโนตฺตปฺปํ สติปฏฺฐานํ, ตญฺจ น สตีติ, อามนฺตา.\csromanspacer{} สติ สติปฏฺฐานา, สา จ น สตีติ, น เหวํ วตฺตพฺเพ ฯเปฯ.}

\csromanheader{2}{8. สติปฏฺฐานกถา}
\proseitem{302}{\csromanfolio{123}\csromansymbolmark{+}น วตฺตพฺพํ “สพฺเพ ธมฺมา สติปฏฺฐานา”ติ, อามนฺตา.\csromanspacer{} นนุ สพฺเพ ธมฺเม อารพฺภ สติ สนฺติฏฺฐตีติ, อามนฺตา.\csromanspacer{} หญฺจิ สพฺเพ ธมฺเม อารพฺภ สติ สนฺติฏฺฐตีติ, เตน วต เร วตฺตพฺเพ “สพฺเพ ธมฺมา สติปฏฺฐานา”ติ.}

\prose{\csromansymbolmark{*}สพฺพํ ธมฺมํ อารพฺภ สติ สนฺติฏฺฐตีติ สพฺเพ ธมฺมา สติปฏฺฐานาติ, อามนฺตา.\csromanspacer{} สพฺพํ ธมฺมํ อารพฺภ ผสฺโส สนฺติฏฺฐตีติ สพฺเพ ธมฺมา ผสฺสปฏฺฐานาติ, น เหวํ วตฺตพฺเพ ฯเปฯ.}

\prose{\csromansymbolmark{*}สพฺพํ ธมฺมํ อารพฺภ สติ สนฺติฏฺฐตีติ สพฺเพ ธมฺมา สติปฏฺฐานาติ, อามนฺตา.\csromanspacer{} สพฺพํ ธมฺมํ อารพฺภ เวทนา สนฺติฏฺฐติ.\csromanspacer{} สญฺญา สนฺติฏฺฐติ.\csromanspacer{} เจตนา สนฺติฏฺฐติ.\csromanspacer{} จิตฺตํ สนฺติฏฺฐตีติ สพฺเพ ธมฺมา จิตฺตปฏฺฐานาติ, น เหวํ วตฺตพฺเพ ฯเปฯ.}

\prose{\csromansymbolmark{*}สพฺเพ ธมฺมา สติปฏฺฐานาติ, อามนฺตา.\csromanspacer{} สพฺเพ สตฺตา อุปฏฺฐิตสติโน สติยา สมนฺนาคตา สติยา สโมหิตา สพฺเพสํ สตฺตานํ สติ ปจฺจุปฏฺฐิตาติ, น เหวํ วตฺตพฺเพ ฯเปฯ.}

\proseitem{303}{\csromansymbolmark{*}สพฺเพ ธมฺมา สติปฏฺฐานาติ, อามนฺตา.\csromanspacer{} นนุ วุตฺตํ ภควตา “อมตํ เต ภิกฺขเว น ปริภุญฺชนฺติ เย กายคตาสติํ น ปริภุญฺชนฺติ, อมตํ เต ภิกฺขเว ปริภุญฺชนฺติ เย กายคตาสติํ ปริภุญฺชนฺตี”ติ\footnote{อํ 1. 47 ปิฏฺเฐ.} อตฺเถว สุตฺตนฺโตติ, อามนฺตา.\csromanspacer{} สพฺเพ สตฺตา กายคตาสติํ ปริภุญฺชนฺติ ปฏิลภนฺติ อาเสวนฺติ ภาเวนฺติ พหุลีกโรนฺตีติ, น เหวํ วตฺตพฺเพ ฯเปฯ.}

\prose{\csromansymbolmark{*}สพฺเพ ธมฺมา สติปฏฺฐานาติ, อามนฺตา.\csromanspacer{} นนุ วุตฺตํ ภควตา “เอกายโน อยํ ภิกฺขเว มคฺโค สตฺตานํ วิสุทฺธิยา โสกปริเทวานํ สมติกฺกมาย ทุกฺขโทมนสฺสานํ อตฺถงฺคมาย ญายสฺส อธิคมาย นิพฺพานสฺส สจฺฉิกิริยาย ยทิทํ จตฺตาโร สติปฏฺฐานา”ติ\footnote{ที 2. 231; ม 1. 70; สํ 3. 123 ปิฏฺเฐสุ.} อตฺเถว สุตฺตนฺโตติ, อามนฺตา.\csromanspacer{} สพฺเพ ธมฺมา เอกายนมคฺโคติ, น เหวํ วตฺตพฺเพ ฯเปฯ.}

\prose{\csromansymbolmark{*}สพฺเพ ธมฺมา สติปฏฺฐานาติ, อามนฺตา.\csromanspacer{} นนุ วุตฺตํ ภควตา\csromandash{} “รญฺโญ ภิกฺขเว จกฺกวตฺติสฺส ปาตุภาวา สตฺตนฺนํ รตนานํ ปาตุภาโว โหติ.\csromanspacer{} กตเมสํ สตฺตนฺนํ, จกฺกรตนสฺส ปาตุภาโว โหติ, หตฺถิรตนสฺส ปาตุภาโว โหติ, อสฺสรตนสฺส.\csromanspacer{} มณิรตนสฺส.\csromanspacer{} อิตฺถิรตนสฺส.\csromanspacer{} คหปติรตนสฺส.\csromanspacer{} ปริณายกรตนสฺส ปาตุภาโว \csromanfolio{124}โหติ.\csromanspacer{} รญฺโญ ภิกฺขเว จกฺกวตฺติสฺส ปาตุภาวา อิเมสํ สตฺตนฺนํ รตนานํ ปาตุภาโว โหติ.}

\prose{ตถาคตสฺส ภิกฺขเว ปาตุภาวา อรหโต สมฺมาสมฺพุทฺธสฺส สตฺตนฺนํ โพชฺฌงฺครตนานํ ปาตุภาโว โหติ.\csromanspacer{} กตเมสํ สตฺตนฺนํ, สติสมฺโพชฺฌงฺครตนสฺส ปาตุภาโว โหติ, ธมฺมวิจยสมฺโพชฺฌงฺครตนสฺส ปาตุภาโว โหติ, วีริยสมฺโพชฺฌงฺครตนสฺส ปาตุภาโว โหติ, ปีติสมฺโพชฺฌงฺครตนสฺส ปาตุภาโว โหติ, ปสฺสทฺธิสมฺโพชฺฌงฺครตนสฺส ปาตุภาโว โหติ, สมาธิสมฺโพชฺฌงฺครตนสฺส ปาตุภาโว โหติ, อุเปกฺขาสมฺโพชฺฌงฺครตนสฺส ปาตุภาโว โหติ.\csromanspacer{} ตถาคตสฺส ภิกฺขเว ปาตุภาวา อรหโต สมฺมาสมฺพุทฺธสฺส อิเมสํ สตฺตนฺนํ โพชฺฌงฺครตนานํ ปาตุภาโว โหตี”ติ\footnote{สํ 3. 87 ปิฏฺเฐ.} อตฺเถว สุตฺตนฺโตติ, อามนฺตา.\csromanspacer{} ตถาคตสฺส ปาตุภาวา อรหโต สมฺมาสมฺพุทฺธสฺส สพฺเพ ธมฺมา สติสมฺโพชฺฌงฺครตนาว โหนฺตีติ, น เหวํ วตฺตพฺเพ ฯเปฯ.\csromanspacer{} สพฺเพ ธมฺมา สติปฏฺฐานาติ, อามนฺตา.\csromanspacer{} สพฺเพ ธมฺมา สมฺมปฺปธานา.\csromanspacer{} อิทฺธิปาทา.\csromanspacer{} อินฺทฺริยา.\csromanspacer{} พลา.\csromanspacer{} โพชฺฌงฺคาติ, น เหวํ วตฺตพฺเพ ฯเปฯ.}

\nitthitam{สติปฏฺฐานกถา นิฏฺฐิตา.}
\csromanmatikaanchor{mtk.28}
\csromantocmarkref{chapter}{9. เหวตฺถิกถา}{mtk.28}
\bookmark[dest={mtk.28},level=0]{9. เหวตฺถิกถา}
\chapterhead{9. เหวตฺถิกถา}
\proseitem{304}{\csromansymbolmark{*}อตีตํ อตฺถีติ? เหวตฺถิ, เหว นตฺถีติ.\csromanspacer{} เสวตฺถิ, เสว นตฺถีติ, น เหวํ วตฺตพฺเพ ฯเปฯ.\csromanspacer{} เสวตฺถิ, เสว นตฺถีติ, อามนฺตา.\csromanspacer{} อตฺถฏฺโฐ นตฺถฏฺโฐ, นตฺถฏฺโฐ อตฺถฏฺโฐ, อตฺถิภาโว นตฺถิภาโว, นตฺถิภาโว อตฺถิภาโว, อตฺถีติ วา นตฺถีติ วา, นตฺถีติ วา อตฺถีติ วา เอเสเส เอกฏฺเฐ สเม สมภาเค ตชฺชาเตติ, น เหวํ วตฺตพฺเพ ฯเปฯ.}

\prose{\csromansymbolmark{*}อนาคตํ อตฺถีติ, เหวตฺถิ, เหว นตฺถีติ.\csromanspacer{} เสวตฺถิ, เสว นตฺถีติ, น เหวํ วตฺตพฺเพ ฯเปฯ.\csromanspacer{} เสวตฺถิ, เสวนตฺถีติ, อามนฺตา.\csromanspacer{} อตฺถฏฺโฐ นตฺถฏฺโฐ, นตฺถฏฺโฐ อตฺถฏฺโฐ, อตฺถิภาโว นตฺถิภาโว, นตฺถิภาโว อตฺถิภาโว, อตฺถีติ วา นตฺถีติ วา, นตฺถีติ วา อตฺถีติ วา เอเสเส เอกฏฺเฐ สเม สมภาเค ตชฺชาเตติ, น เหวํ วตฺตพฺเพ ฯเปฯ.}

\csromanheader{2}{9. เหวตฺถิกถา}
\prose{\csromanfolio{125}\csromansymbolmark{*}ปจฺจุปฺปนฺนํ อตฺถีติ? เหวตฺถิ, เหว นตฺถีติ.\csromanspacer{} เสวตฺถิ, เสว นตฺถีติ, น เหวํ วตฺตพฺเพ ฯเปฯ.\csromanspacer{} เสวตฺถิ, เสว นตฺถีติ, อามนฺตา.\csromanspacer{} อตฺถฏฺโฐ นตฺถฏฺโฐ ฯเปฯ.\csromanspacer{} สเม สมภาเค ตชฺชาเตติ, น เหวํ วตฺตพฺเพ ฯเปฯ.}

\proseitem{305}{\csromansymbolmark{*}อตีตํ เหวตฺถิ, เหว นตฺถีติ, อามนฺตา.\csromanspacer{} กินฺตตฺถิ, กินฺติ นตฺถีติ, อตีตํ “อตีตนฺ”ติ เหวตฺถิ, อตีตํ “อนาคตนฺ”ติ เหว นตฺถิ, อตีตํ “ปจฺจุปฺปนฺนนฺ”ติ เหว นตฺถีติ.\csromanspacer{} เสวตฺถิ, เสว นตฺถีติ, น เหวํ วตฺตพฺเพ ฯเปฯ.\csromanspacer{} เสวตฺถิ, เสว นตฺถีติ, อามนฺตา.\csromanspacer{} อตฺถฏฺโฐ นตฺถฏฺโฐ, นตฺถฏฺโฐ อตฺถฏฺโฐ, อตฺถิภาโว นตฺถิภาโว, นตฺถิภาโว อตฺถิภาโว, อตฺถีติ วา นตฺถีติ วา, นตฺถีติ วา อตฺถีติ วา เอเสเส เอกฏฺเฐ สเม สมภาเค ตชฺชาเตติ, น เหวํ วตฺตพฺเพ ฯเปฯ.}

\prose{\csromansymbolmark{*}อนาคตํ เหวตฺถิ, เหว นตฺถีติ, อามนฺตา.\csromanspacer{} กินฺตตฺถิ, กินฺติ นตฺถีติ, อนาคตํ “อนาคตนฺ”ติ เหวตฺถิ, อนาคตํ “อตีตนฺ”ติ เหว นตฺถิ, อนาคตํ “ปจฺจุปฺปนฺนนฺ”ติ เหว นตฺถีติ.\csromanspacer{} เสวตฺถิ, เสว นตฺถีติ, น เหวํ วตฺตพฺเพ ฯเปฯ.\csromanspacer{} เสวตฺถิ, เสว นตฺถีติ, อามนฺตา.\csromanspacer{} อตฺถฏฺโฐ นตฺถฏฺโฐ, นตฺถฏฺโฐ อตฺถฏฺโฐ ฯเปฯ สเม สมภาเค ตชฺชาเตติ, น เหวํ วตฺตพฺเพ ฯเปฯ.}

\prose{\csromansymbolmark{*}ปจฺจุปฺปนฺนํ เหวตฺถิ, เหว นตฺถีติ, อามนฺตา.\csromanspacer{} กินฺตตฺถิ, กินฺติ นตฺถีติ, ปจฺจุปฺปนฺนํ “ปจฺจุปฺปนฺนนฺ”ติ เหวตฺถิ, ปจฺจุปฺปนฺนํ “อตีตนฺ”ติ เหว นตฺถิ, ปจฺจุปฺปนฺนํ “อนาคตนฺ”ติ เหว นตฺถีติ.\csromanspacer{} เสวตฺถิ, เสว นตฺถีติ, น เหวํ วตฺตพฺเพ ฯเปฯ.\csromanspacer{} เสวตฺถิ, เสว นตฺถีติ, อามนฺตา.\csromanspacer{} อตฺถฏฺโฐ นตฺถฏฺโฐ ฯเปฯ สเม สมภาเค ตชฺชาเตติ, น เหวํ วตฺตพฺเพ ฯเปฯ.}

\prose{\csromansymbolmark{+}น วตฺตพฺพํ “อตีตํ เหวตฺถิ, เหว นตฺถิ.\csromanspacer{} อนาคตํ เหวตฺถิ, เหว นตฺถิ.\csromanspacer{} ปจฺจุปฺปนฺนํ เหวตฺถิ, เหว นตฺถี”ติ, อามนฺตา.\csromanspacer{} อตีตํ “อนาคตนฺ”ติ เหวตฺถิ.\csromanspacer{} อตีตํ “ปจฺจุปฺปนฺนนฺ”ติ เหวตฺถิ.\csromanspacer{} อนาคตํ “อตีตนฺ”ติ เหวตฺถิ.\csromanspacer{} อนาคตํ “ปจฺจุปฺปนฺนนฺ”ติ เหวตฺถิ.\csromanspacer{} ปจฺจุปฺปนฺนํ “อตีตนฺ”ติ เหวตฺถิ.\csromanspacer{} ปจฺจุปฺปนฺนํ “อนาคตนฺ”ติ เหวตฺถีติ, น เหวํ วตฺตพฺเพ ฯเปฯ.\csromanspacer{} เตน หิ อตีตํ เหวตฺถิ, เหว นตฺถิ.\csromanspacer{} อนาคตํ เหวตฺถิ, เหว นตฺถิ.\csromanspacer{} ปจฺจุปฺปนฺนํ เหวตฺถิ, เหว นตฺถีติ.}

\proseitem{306}{\csromansymbolmark{*}รูปํ อตฺถีติ? เหวตฺถิ, เหว นตฺถีติ.\csromanspacer{} เสวตฺถิ, เสว นตฺถีติ, น เหวํ วตฺตพฺเพ ฯเปฯ.\csromanspacer{} เสวตฺถิ, เสว นตฺถีติ, อามนฺตา.\csromanspacer{} อตฺถฏฺโฐ นตฺถฏฺโฐ, \csromanfolio{126}นตฺถฏฺโฐ อตฺถฏฺโฐ, อตฺถิภาโว นตฺถิภาโว, นตฺถิภาโว อตฺถิภาโว, อตฺถีติ วา นตฺถีติ วา, นตฺถีติ วา อตฺถีติ วา เอเสเส เอกฏฺเฐ สเม สมภาเค ตชฺชาเตติ, น เหวํ วตฺตพฺเพ ฯเปฯ.}

\prose{\csromansymbolmark{*}เวทนา.\csromanspacer{} สญฺญา.\csromanspacer{} สงฺขารา.\csromanspacer{} วิญฺญาณํ อตฺถีติ? เหวตฺถิ, เหว นตฺถีติ.\csromanspacer{} เสวตฺถิ, เสว นตฺถีติ, น เหวํ วตฺตพฺเพ ฯเปฯ.\csromanspacer{} เสวตฺถิ, เสว นตฺถีติ, อามนฺตา.\csromanspacer{} อตฺถฏฺโฐ นตฺถฏฺโฐ ฯเปฯ สเม สมภาเค ตชฺชาเตติ, น เหวํ วตฺตพฺเพ ฯเปฯ.}

\prose{\csromansymbolmark{*}รูปํ เหวตฺถิ, เหว นตฺถีติ, อามนฺตา.\csromanspacer{} กินฺตตฺถิ, กินฺติ นตฺถีติ.\csromanspacer{} รูปํ “รูปนฺ”ติ เหวตฺถิ.\csromanspacer{} รูปํ “เวทนา”ติ เหว นตฺถิ ฯเปฯ.\csromanspacer{} รูปํ “สญฺญา”ติ เหว นตฺถิ ฯเปฯ.\csromanspacer{} รูปํ “สงฺขารา”ติ เหว นตฺถิ ฯเปฯ.\csromanspacer{} รูปํ “วิญฺญาณนฺ”ติ เหว นตฺถีติ.\csromanspacer{} เสวตฺถิ, เสว นตฺถีติ, น เหวํ วตฺตพฺเพ ฯเปฯ.\csromanspacer{} เสวตฺถิ, เสว นตฺถีติ, อามนฺตา.\csromanspacer{} อตฺถฏฺโฐ นตฺถฏฺโฐ ฯเปฯ สเม สมภาเค ตชฺชาเตติ, น เหวํ วตฺตพฺเพ ฯเปฯ.}

\prose{\csromansymbolmark{*}เวทนา.\csromanspacer{} สญฺญา.\csromanspacer{} สงฺขารา.\csromanspacer{} วิญฺญาณํ เหวตฺถิ, เหว นตฺถีติ, อามนฺตา.\csromanspacer{} กินฺตตฺถิ, กินฺติ นตฺถีติ.\csromanspacer{} วิญฺญาณํ “วิญฺญาณนฺ”ติ เหวตฺถิ.\csromanspacer{} วิญฺญาณํ “รูปนฺ”ติ เหว นตฺถิ ฯเปฯ.\csromanspacer{} วิญฺญาณํ “เวทนา”ติ เหว นตฺถิ ฯเปฯ.\csromanspacer{} วิญฺญาณํ “สญฺญา”ติ เหว นตฺถิ ฯเปฯ.\csromanspacer{} วิญฺญาณํ “สงฺขารา”ติ เหว นตฺถีติ.\csromanspacer{} เสวตฺถิ, เสว นตฺถีติ, น เหวํ วตฺตพฺเพ ฯเปฯ.\csromanspacer{} เสวตฺถิ, เสว นตฺถีติ, อามนฺตา.\csromanspacer{} อตฺถฏฺโฐ นตฺถฏฺโฐ ฯเปฯ สเม สมภาเค ตชฺชาเตติ, น เหวํ วตฺตพฺเพ ฯเปฯ.}

\prose{\csromansymbolmark{+}น วตฺตพฺพํ “รูปํ เหวตฺถิ, เหว นตฺถี”ติ. “เวทนา.\csromanspacer{} สญฺญา.\csromanspacer{} สงฺขารา.\csromanspacer{} วิญฺญาณํ เหวตฺถิ, เหว นตฺถี”ติ, อามนฺตา.\csromanspacer{} รูปํ “เวทนา”ติ เหวตฺถิ ฯเปฯ.\csromanspacer{} รูปํ “สญฺญา”ติ เหวตฺถิ ฯเปฯ.\csromanspacer{} รูปํ “สงฺขารา”ติ เหวตฺถิ ฯเปฯ.\csromanspacer{} รูปํ “วิญฺญาณนฺ”ติ เหวตฺถิ.\csromanspacer{} เวทนา.\csromanspacer{} สญฺญา.\csromanspacer{} สงฺขารา.\csromanspacer{} วิญฺญาณํ “รูปนฺ”ติ เหวตฺถิ.\csromanspacer{} วิญฺญาณํ “เวทนา”ติ เหวตฺถิ.\csromanspacer{} วิญฺญาณํ “สญฺญา”ติ เหวตฺถิ.\csromanspacer{} วิญฺญาณํ “สงฺขารา”ติ เหวตฺถีติ, น เหวํ วตฺตพฺเพ ฯเปฯ.\csromanspacer{} เตน หิ รูปํ เหวตฺถิ, เหว นตฺถิ.\csromanspacer{} เวทนา.\csromanspacer{} สญฺญา.\csromanspacer{} สงฺขารา.\csromanspacer{} วิญฺญาณํ เหวตฺถิ, เหว นตฺถีติ.}

\nitthitam{เหวตฺถิกถา นิฏฺฐิตา.}
\chapterheadpageread{2. ทุติยวคฺค \textbf{ตสฺสุทฺทานํ}}
\csromanmatikaanchor{mtk.29}
\csromanmatikaanchor{mtk.30}
\csromantocmarknopage{chapter}{2. ทุติยวคฺค}{mtk.29}
\bookmark[dest={mtk.29},level=0]{2. ทุติยวคฺค}
\csromantocmarkref{subsection}{(10) 1. ปรูปหารกถา}{mtk.30}
\bookmark[dest={mtk.30},level=2]{(10) 1. ปรูปหารกถา}
\csromangathameasure{อุปลพฺโภ ปริหานิ, พฺรหฺมจริยวาโส โอธิโส. \\ ปริญฺญา กามราคปฺปหานํ, สพฺพตฺถิวาโท อายตนํ. \\ อตีตานาคโต สุภงฺโค, สพฺเพ ธมฺมา สติปฏฺฐานา.}
\csromangathagroup{\csromangathastanza{\csromanfolio{127}อุปลพฺโภ ปริหานิ, พฺรหฺมจริยวาโส โอธิโส. \\ ปริญฺญา กามราคปฺปหานํ, สพฺพตฺถิวาโท อายตนํ. \\ อตีตานาคโต สุภงฺโค\footnote{อตีตานาคเตสุ ภาโค (สฺยา)}, สพฺเพ ธมฺมา สติปฏฺฐานา.}}

\vspace{\tipitakaparskip}
\csromancenter{เหวตฺถิ เหว นตฺถีติ.}
\csromancenter{มหาวคฺโค.}
\chapterhead{2. ทุติยวคฺค}
\vspace{\tipitakaparskip}
\csromancenter{\textbf{(10) 1.}\csromanspacer{}\textbf{ ปรูปหารกถา}}
\proseitem{307}{\csromansymbolmark{*}อตฺถิ อรหโต อสุจิ สุกฺกวิสฺสฏฺฐีติ, อามนฺตา.\csromanspacer{} อตฺถิ อรหโต ราโค กามราโค กามราคปริยุฏฺฐานํ กามราคสํโยชนํ กาโมโฆ กามโยโค กามจฺฉนฺทนีวรณนฺติ, น เหวํ วตฺตพฺเพ ฯเปฯ.}

\prose{\csromansymbolmark{*}นตฺถิ อรหโต ราโค กามราโค กามราคปริยุฏฺฐานํ กามราคสํโยชนํ กาโมโฆ กามโยโค กามจฺฉนฺทนีวรณนฺติ, อามนฺตา.\csromanspacer{} หญฺจิ นตฺถิ อรหโต ราโค กามราโค กามราคปริยุฏฺฐานํ กามราคสํโยชนํ กาโมโฆ กามโยโค กามจฺฉนฺทนีวรณํ, โน จ วต เร วตฺตพฺเพ “อตฺถิ อรหโต อสุจิ สุกฺกวิสฺสฏฺฐี”ติ.}

\prose{\csromansymbolmark{*}อตฺถิ ปุถุชฺชนสฺส อสุจิ สุกฺกวิสฺสฏฺฐิ, อตฺถิ ตสฺส ราโค กามราโค กามราคปริยุฏฺฐานํ กามราคสํโยชนํ กาโมโฆ กามโยโค กามจฺฉนฺทนีวรณนฺติ, อามนฺตา.\csromanspacer{} อตฺถิ อรหโต อสุจิ สุกฺกวิสฺสฏฺฐิ, อตฺถิ ตสฺส ราโค กามราโค กามราคปริยุฏฺฐานํ ฯเปฯ กามจฺฉนฺทนีวรณนฺติ, น เหวํ วตฺตพฺเพ ฯเปฯ.}

\prose{\csromansymbolmark{*}อตฺถิ อรหโต อสุจิ สุกฺกวิสฺสฏฺฐิ, นตฺถิ ตสฺส ราโค กามราโค กามราคปริยุฏฺฐานํ ฯเปฯ กามจฺฉนฺทนีวรณนฺติ, อามนฺตา.\csromanspacer{} อตฺถิ ปุถุชฺชนสฺส อสุจิ สุกฺกวิสฺสฏฺฐิ, นตฺถิ ตสฺส ราโค กามราโค กามราคปริยุฏฺฐานํ ฯเปฯ กามจฺฉนฺทนีวรณนฺติ, น เหวํ วตฺตพฺเพ ฯเปฯ.}

\prose{\csromanfolio{128}\csromansymbolmark{*}อตฺถิ อรหโต อสุจิ สุกฺกวิสฺสฏฺฐีติ, อามนฺตา.\csromanspacer{} เกนฏฺเฐนาติ? หนฺท หิ มารกายิกา เทวตา อรหโต อสุจิํ สุกฺกวิสฺสฏฺฐิํ อุปสํหรนฺตีติ.}

\prose{\csromansymbolmark{*}มารกายิกา เทวตา อรหโต อสุจิํ สุกฺกวิสฺสฏฺฐิํ อุปสํหรนฺตีติ, อามนฺตา.\csromanspacer{} อตฺถิ มารกายิกานํ เทวตานํ อสุจิ สุกฺกวิสฺสฏฺฐีติ, น เหวํ วตฺตพฺเพ ฯเปฯ.}

\prose{\csromansymbolmark{*}นตฺถิ มารกายิกานํ เทวตานํ อสุจิ สุกฺกวิสฺสฏฺฐีติ, อามนฺตา.\csromanspacer{} หญฺจิ นตฺถิ มารกายิกานํ เทวตานํ อสุจิ สุกฺกวิสฺสฏฺฐิ, โน จ วต เร วตฺตพฺเพ “มารกายิกา เทวตา อรหโต อสุจิํ สุกฺกวิสฺสฏฺฐิํ อุปสํหรนฺตี”ติ.}

\prose{\csromansymbolmark{*}มารกายิกา เทวตา อรหโต อสุจิํ สุกฺกวิสฺสฏฺฐิํ อุปสํหรนฺตีติ, อามนฺตา.\csromanspacer{} มารกายิกา เทวตา อตฺตโน อสุจิํ สุกฺกวิสฺสฏฺฐิํ อุปสํหรนฺติ, อญฺเญสํ อสุจิํ สุกฺกวิสฺสฏฺฐิํ อุปสํหรนฺติ, ตสฺส อสุจิํ สุกฺกวิสฺสฏฺฐิํ อุปสํหรนฺตีติ, น เหวํ วตฺตพฺเพ ฯเปฯ.}

\prose{\csromansymbolmark{*}มารกายิกา เทวตา เนว อตฺตโน น อญฺเญสํ น ตสฺส อสุจิํ สุกฺกวิสฺสฏฺฐิํ อุปสํหรนฺตีติ, อามนฺตา.\csromanspacer{} หญฺจิ มารกายิกา เทวตา เนว อตฺตโน น อญฺเญสํ น ตสฺส อสุจิํ สุกฺกวิสฺสฏฺฐิํ อุปสํหรนฺติ, โน จ วต เร วตฺตพฺเพ “มารกายิกา เทวตา อรหโต อสุจิํ สุกฺกวิสฺสฏฺฐิํ อุปสํหรนฺตี”ติ.}

\prose{\csromansymbolmark{*}มารกายิกา เทวตา อรหโต อสุจิํ สุกฺกวิสฺสฏฺฐิํ อุปสํหรนฺตีติ, อามนฺตา.\csromanspacer{} โลมกูเปหิ อุปสํหรนฺตีติ, น เหวํ วตฺตพฺเพ ฯเปฯ.}

\proseitem{308}{\csromansymbolmark{*}มารกายิกา เทวตา อรหโต อสุจิํ สุกฺกวิสฺสฏฺฐิํ อุปสํหรนฺตีติ, อามนฺตา.\csromanspacer{} กิํ การณาติ, หนฺท หิ วิมติํ คาหยิสฺสามาติ.\csromanspacer{} อตฺถิ อรหโต วิมตีติ, น เหวํ วตฺตพฺเพ ฯเปฯ.}

\prose{\csromansymbolmark{*}อตฺถิ อรหโต วิมตีติ, อามนฺตา.\csromanspacer{} อตฺถิ อรหโต สตฺถริ วิมติ, ธมฺเม วิมติ, สํเฆ วิมติ, สิกฺขาย วิมติ, ปุพฺพนฺเต วิมติ, อปรนฺเต วิมติ, ปุพฺพนฺตาปรนฺเต วิมติ, อิทปฺปจฺจยตาปฏิจฺจสมุปฺปนฺเนสุ ธมฺเมสุ วิมตีติ, น เหวํ วตฺตพฺเพ ฯเปฯ.}

\chapterheadpageread{2. ทุติยวคฺค}
\prose{\csromanfolio{129}\csromansymbolmark{*}นตฺถิ อรหโต สตฺถริ วิมติ, ธมฺเม วิมติ, สํเฆ วิมติ, สิกฺขาย วิมติ, ปุพฺพนฺเต วิมติ, อปรนฺเต วิมติ, ปุพฺพนฺตาปรนฺเต วิมติ, อิทปฺปจฺจยตาปฏิจฺจสมุปฺปนฺเนสุ ธมฺเมสุ วิมตีติ, อามนฺตา.\csromanspacer{} หญฺจิ นตฺถิ อรหโต สตฺถริ วิมติ ฯเปฯ อิทปฺปจฺจยตาปฏิจฺจสมุปฺปนฺเนสุ ธมฺเมสุ วิมติ, โน จ วต เร วตฺตพฺเพ “อตฺถิ อรหโต วิมตี”ติ.}

\prose{\csromansymbolmark{*}อตฺถิ ปุถุชฺชนสฺส วิมติ, อตฺถิ ตสฺส สตฺถริ วิมติ ฯเปฯ อิทปฺปจฺจยตาปฏิจฺจสมุปฺปนฺเนสุ ธมฺเมสุ วิมตีติ, อามนฺตา.\csromanspacer{} อตฺถิ อรหโต วิมติ อตฺถิ ตสฺส สตฺถริ วิมติ ฯเปฯ อิทปฺปจฺจยตาปฏิจฺจสมุปฺปนฺเนสุ ธมฺเมสุ วิมตีติ, น เหวํ วตฺตพฺเพ ฯเปฯ.}

\prose{\csromansymbolmark{*}อตฺถิ อรหโต วิมติ, นตฺถิ ตสฺส สตฺถริ วิมติ ฯเปฯ อิทปฺปจฺจยตาปฏิจฺจสมุปฺปนฺเนสุ ธมฺเมสุ วิมตีติ, อามนฺตา.\csromanspacer{} อตฺถิ ปุถุชฺชนสฺส วิมติ, นตฺถิ ตสฺส สตฺถริ วิมติ ฯเปฯ อิทปฺปจฺจยตาปฏิจฺจสมุปฺปนฺเนสุ ธมฺเมสุ วิมตีติ, น เหวํ วตฺตพฺเพ ฯเปฯ.}

\prose{\csromansymbolmark{*}อตฺถิ อรหโต อสุจิ สุกฺกวิสฺสฏฺฐีติ, อามนฺตา.\csromanspacer{} อรหโต อสุจิ สุกฺกวิสฺสฏฺฐิ กิสฺส นิสฺสนฺโทติ, อสิตปีตขายิตสายิตสฺส นิสฺสนฺโทติ.\csromanspacer{} อรหโต อสุจิ สุกฺกวิสฺสฏฺฐิ อสิตปีตขายิตสายิตสฺส นิสฺสนฺโทติ, อามนฺตา.\csromanspacer{} เย เกจิ อสนฺติ ปิวนฺติ ขาทนฺติ สายนฺติ, สพฺเพสํเยว อตฺถิ อสุจิ สุกฺกวิสฺสฏฺฐีติ, น เหวํ วตฺตพฺเพ ฯเปฯ.}

\prose{\csromansymbolmark{*}เย เกจิ อสนฺติ ปิวนฺติ ขาทนฺติ สายนฺติ, สพฺเพสํเยว อตฺถิ อสุจิ สุกฺกวิสฺสฏฺฐีติ, อามนฺตา.\csromanspacer{} ทารกา อสนฺติ ปิวนฺติ ขาทนฺติ สายนฺติ, อตฺถิ ทารกานํ อสุจิ สุกฺกวิสฺสฏฺฐีติ, น เหวํ วตฺตพฺเพ.}

\prose{\csromansymbolmark{*}ปณฺฑกา อสนฺติ ปิวนฺติ ขาทนฺติ สายนฺติ, อตฺถิ ปณฺฑกานํ อสุจิ สุกฺกวิสฺสฏฺฐีติ, น เหวํ วตฺตพฺเพ ฯเปฯ.}

\prose{\csromansymbolmark{*}เทวา อสนฺติ ปิวนฺติ ขาทนฺติ สายนฺติ, อตฺถิ เทวตานํ อสุจิ สุกฺกวิสฺสฏฺฐีติ, น เหวํ วตฺตพฺเพ ฯเปฯ.}

\proseitem{309}{\csromansymbolmark{*}อรหโต อสุจิ สุกฺกวิสฺสฏฺฐิ อสิตปีตขายิตสายิตสฺส นิสฺสนฺโทติ, อามนฺตา.\csromanspacer{} อตฺถิ ตสฺส อาสโยติ, น เหวํ วตฺตพฺเพ ฯเปฯ.}

\prose{\csromanfolio{130}\csromansymbolmark{*}อรหโต อุจฺจารปสฺสาโว อสิตปีตขายิตสายิตสฺส นิสฺสนฺโท, อตฺถิ ตสฺส อาสโยติ, อามนฺตา.\csromanspacer{} อรหโต อสุจิ สุกฺกวิสฺสฏฺฐิ อสิตปีตขายิตสายิตสฺส นิสฺสนฺโท, อตฺถิ ตสฺส อาสโยติ, น เหวํ วตฺตพฺเพ. \csromansymbolmark{*}อรหโต อสุจิ สุกฺกวิสฺสฏฺฐิ อสิตปีตขายิตสายิตสฺส นิสฺสนฺโท, นตฺถิ ตสฺส อาสโยติ, อามนฺตา.\csromanspacer{} อรหโต อุจฺจารปสฺสาโว อสิตปีตขายิตสายิตสฺส นิสฺสนฺโท, นตฺถิ ตสฺส อาสโยติ, น เหวํ วตฺตพฺเพ ฯเปฯ.}

\proseitem{310}{\csromansymbolmark{*}อตฺถิ อรหโต อสุจิ สุกฺกวิสฺสฏฺฐีติ, อามนฺตา.\csromanspacer{} อรหา เมถุนํ ธมฺมํ ปฏิเสเวยฺย, เมถุนํ อุปฺปาเทยฺย\footnote{เมถุนํ ธมฺมํ อุปฺปาเทยฺย (สี, สฺยา)}, ปุตฺตสมฺพาธสยนํ อชฺฌาวเสยฺย, กาสิกจนฺทนํ ปจฺจนุภเวยฺย, มาลาคนฺธวิเลปนํ ธาเรยฺย, ชาตรูปรชตํ สาทิเยยฺยาติ, น เหวํ วตฺตพฺเพ.}

\prose{\csromansymbolmark{*}อตฺถิ ปุถุชฺชนสฺส อสุจิ สุกฺกวิสฺสฏฺฐิ, ปุถุชฺชโน เมถุนํ ธมฺมํ ปฏิเสเวยฺย, เมถุนํ อุปฺปาเทยฺย, ปุตฺตสมฺพาธสยนํ อชฺฌาวเสยฺย, กาสิกจนฺทนํ ปจฺจนุภเวยฺย, มาลาคนฺธวิเลปนํ ธาเรยฺย, ชาตรูปรชตํ สาทิเยยฺยาติ, อามนฺตา.\csromanspacer{} อตฺถิ อรหโต อสุจิ สุกฺกวิสฺสฏฺฐิ อรหา เมถุนํ ธมฺมํ ปฏิเสเวยฺย, เมถุนํ อุปฺปาเทยฺย, ปุตฺตสมฺพาธสยนํ อชฺฌาวเสยฺย, กาสิกจนฺทนํ ปจฺจนุภเวยฺย, มาลาคนฺธวิเลปนํ ธาเรยฺย, ชาตรูปรชตํ สาทิเยยฺยาติ, น เหวํ วตฺตพฺเพ ฯเปฯ.}

\prose{\csromansymbolmark{*}อตฺถิ อรหโต อสุจิ สุกฺกวิสฺสฏฺฐิ, น จ อรหา เมถุนํ ธมฺมํ ปฏิเสเวยฺย, เมถุนํ อุปฺปาเทยฺย, ปุตฺตสมฺพาธสยนํ อชฺฌาวเสยฺย, กาสิกจนฺทนํ ปจฺจนุภเวยฺย, มาลาคนฺธวิเลปนํ ธาเรยฺย, ชาตรูปรชตํ สาทิเยยฺยาติ, อามนฺตา.\csromanspacer{} อตฺถิ ปุถุชฺชนสฺส อสุจิ สุกฺกวิสฺสฏฺฐิ, น จ ปุถุชฺชโน เมถุนํ ธมฺมํ ปฏิเสเวยฺย, เมถุนํ อุปฺปาเทยฺย, ปุตฺตสมฺพาธสยนํ อชฺฌาวเสยฺย, กาสิกจนฺทนํ ปจฺจนุภเวยฺย, มาลาคนฺธวิเลปนํ ธาเรยฺย, ชาตรูปรชตํ สาทิเยยฺยาติ, น เหวํ วตฺตพฺเพ.}

\prose{\csromansymbolmark{*}อตฺถิ อรหโต อสุจิ สุกฺกวิสฺสฏฺฐีติ, อามนฺตา.\csromanspacer{} นนุ อรหโต ราโค ปหีโน อุจฺฉินฺนมูโล ตาลาวตฺถุกโต อนภาวํกโต}

\vspace{\tipitakaparskip}
\csromancenter{ทุติยวคฺค อายติํ อนุปฺปาทธมฺโมติ, อามนฺตา.\csromanspacer{} หญฺจิ อรหโต ราโค ปหีโน อุจฺฉินฺนมูโล ตาลาวตฺถุกโต อนภาวํกโต อายติํ อนุปฺปาทธมฺโม, โน จ วต เร วตฺตพฺเพ “อตฺถิ อรหโต อสุจิ สุกฺกวิสฺสฏฺฐี”ติ.}
\prose{\csromanfolio{131}\csromansymbolmark{*}อตฺถิ อรหโต อสุจิ สุกฺกวิสฺสฏฺฐีติ, อามนฺตา.\csromanspacer{} นนุ อรหโต โทโส ปหีโน ฯเปฯ โมโห ปหีโน.\csromanspacer{} มาโน ปหีโน.\csromanspacer{} ทิฏฺฐิ ปหีนา.\csromanspacer{} วิจิกิจฺฉา ปหีนา.\csromanspacer{} ถินํ ปหีนํ.\csromanspacer{} อุทฺธจฺจํ ปหีนํ.\csromanspacer{} อหิริกํ ปหีนํ ฯเปฯ.\csromanspacer{} อโนตฺตปฺปํ ปหีนํ อุจฺฉินฺนมูลํ ตาลาวตฺถุกตํ อนภาวํกตํ อายติํ อนุปฺปาทธมฺมนฺติ, อามนฺตา.\csromanspacer{} หญฺจิ อรหโต อโนตฺตปฺปํ ปหีนํ อุจฺฉินฺนมูลํ ตาลาวตฺถุกตํ อนภาวํกตํ อายติํ อนุปฺปาทธมฺมํ, โน จ วต เร วตฺตพฺเพ “อตฺถิ อรหโต อสุจิ สุกฺกวิสฺสฏฺฐี”ติ.}

\proseitem{311}{\csromansymbolmark{*}อตฺถิ อรหโต อสุจิ สุกฺกวิสฺสฏฺฐีติ, อามนฺตา.\csromanspacer{} นนุ อรหโต ราคปฺปหานาย มคฺโค ภาวิโตติ, อามนฺตา.\csromanspacer{} หญฺจิ อรหโต ราคปฺปหานาย มคฺโค ภาวิโต, โน จ วต เร วตฺตพฺเพ “อตฺถิ อรหโต อสุจิ สุกฺกวิสฺสฏฺฐี”ติ.}

\prose{\csromansymbolmark{*}อตฺถิ อรหโต อสุจิ สุกฺกวิสฺสฏฺฐีติ, อามนฺตา.\csromanspacer{} นนุ อรหโต ราคปฺปหานาย สติปฏฺฐานา ภาวิตา ฯเปฯ สมฺมปฺปธานา ภาวิตา.\csromanspacer{} อิทฺธิปาทา ภาวิตา.\csromanspacer{} อินฺทฺริยา ภาวิตา.\csromanspacer{} พลา ภาวิตา ฯเปฯ โพชฺฌงฺคา ภาวิตาติ, อามนฺตา.\csromanspacer{} หญฺจิ อรหโต ราคปฺปหานาย โพชฺฌงฺคา ภาวิตา, โน จ วต เร วตฺตพฺเพ “อตฺถิ อรหโต อสุจิ สุกฺกวิสฺสฏฺฐี”ติ.}

\prose{\csromansymbolmark{*}อตฺถิ อรหโต อสุจิ สุกฺกวิสฺสฏฺฐีติ, อามนฺตา.\csromanspacer{} นนุ อรหโต โทสปฺปหานาย ฯเปฯ โมหปฺปหานาย ฯเปฯ อโนตฺตปฺปปหานาย มคฺโค ภาวิโต ฯเปฯ โพชฺฌงฺคา ภาวิตาติ, อามนฺตา.\csromanspacer{} หญฺจิ อรหโต อโนตฺตปฺปปหานาย โพชฺฌงฺคา ภาวิตา, โน จ วต เร วตฺตพฺเพ “อตฺถิ อรหโต อสุจิ สุกฺกวิสฺสฏฺฐี”ติ.}

\prose{\csromansymbolmark{*}อตฺถิ อรหโต อสุจิ สุกฺกวิสฺสฏฺฐีติ, อามนฺตา.\csromanspacer{} นนุ อรหา วีตราโค วีตโทโส วีตโมโห กตกรณีโย โอหิตภาโร อนุปฺปตฺตสทตฺโถ ปริกฺขีณภวสํโยชโน สมฺมทญฺญาวิมุตฺโต อุกฺขิตฺตปลิโฆ สํกิณฺณปริโข อพฺพูเฬฺหสิโก นิรคฺคโฬ อริโย \csromanfolio{132}ปนฺนทฺธโช ปนฺนภาโร วิสญฺญุตฺโต สุวิชิตวิชโย, ทุกฺขํ ตสฺส ปริญฺญาตํ, สมุทโย ปหีโน, นิโรโธ สจฺฉิกโต, มคฺโค ภาวิโต, อภิญฺเญยฺยํ อภิญฺญาตํ, ปริญฺเญยฺยํ ปริญฺญาตํ, ปหาตพฺพํ ปหีนํ, ภาเวตพฺพํ ภาวิตํ, สจฺฉิกาตพฺพํ สจฺฉิกตนฺติ, อามนฺตา.\csromanspacer{} หญฺจิ อรหา วีตราโค วีตโทโส วีตโมโห กตกรณีโย ฯเปฯ สจฺฉิกาตพฺพํ สจฺฉิกตํ, โน จ วต เร วตฺตพฺเพ “อตฺถิ อรหโต อสุจิ สุกฺกวิสฺสฏฺฐี”ติ.}

\proseitem{312}{\csromansymbolmark{*}อตฺถิ อรหโต อสุจิ สุกฺกวิสฺสฏฺฐีติ? สธมฺมกุสลสฺส อรหโต อตฺถิ อสุจิ สุกฺกวิสฺสฏฺฐิ, ปรธมฺมกุสลสฺส อรหโต นตฺถิ อสุจิ สุกฺกวิสฺสฏฺฐีติ.\csromanspacer{} สธมฺมกุสลสฺส อรหโต อตฺถิ อสุจิ สุกฺกวิสฺสฏฺฐีติ, อามนฺตา.\csromanspacer{} ปรธมฺมกุสลสฺส อรหโต อตฺถิ อสุจิ สุกฺกวิสฺสฏฺฐีติ, น เหวํ วตฺตพฺเพ ฯเปฯ.}

\prose{\csromansymbolmark{*}ปรธมฺมกุสลสฺส อรหโต นตฺถิ อสุจิ สุกฺกวิสฺสฏฺฐีติ, อามนฺตา.\csromanspacer{} สธมฺมกุสลสฺส อรหโต นตฺถิ อสุจิ สุกฺกวิสฺสฏฺฐีติ, น เหวํ วตฺตพฺเพ ฯเปฯ.}

\prose{\csromansymbolmark{*}สธมฺมกุสลสฺส อรหโต ราโค ปหีโน, อตฺถิ ตสฺส อสุจิ สุกฺกวิสฺสฏฺฐีติ, อามนฺตา.\csromanspacer{} ปรธมฺมกุสลสฺส อรหโต ราโค ปหีโน, อตฺถิ ตสฺส อสุจิ สุกฺกวิสฺสฏฺฐีติ, น เหวํ วตฺตพฺเพ ฯเปฯ.}

\prose{\csromansymbolmark{*}สธมฺมกุสลสฺส อรหโต โทโส ปหีโน ฯเปฯ โมโห ปหีโน ฯเปฯ อโนตฺตปฺปํ ปหีนํ, อตฺถิ ตสฺส อสุจิ สุกฺกวิสฺสฏฺฐีติ, อามนฺตา.\csromanspacer{} ปรธมฺมกุสลสฺส อรหโต อโนตฺตปฺปํ ปหีนํ, อตฺถิ ตสฺส อสุจิ สุกฺกวิสฺสฏฺฐีติ, น เหวํ วตฺตพฺเพ ฯเปฯ.}

\prose{\csromansymbolmark{*}สธมฺมกุสลสฺส อรหโต ราคปฺปหานาย มคฺโค ภาวิโต ฯเปฯ โพชฺฌงฺคา ภาวิตา ฯเปฯ โทสปฺปหานาย ฯเปฯ โมหปฺปหานาย ฯเปฯ อโนตฺตปฺปปหานาย มคฺโค ภาวิโต ฯเปฯ โพชฺฌงฺคา ภาวิตา, อตฺถิ ตสฺส อสุจิ สุกฺกวิสฺสฏฺฐีติ, อามนฺตา.\csromanspacer{} ปรธมฺมกุสลสฺส อรหโต อโนตฺตปฺปปหานาย โพชฺฌงฺคา ภาวิตา, อตฺถิ ตสฺส อสุจิ สุกฺกวิสฺสฏฺฐีติ, น เหวํ วตฺตพฺเพ ฯเปฯ.}

\prose{\csromansymbolmark{*}สธมฺมกุสโล อรหา วีตราโค วีตโทโส วีตโมโห ฯเปฯ สจฺฉิกาตพฺพํ สจฺฉิกตํ, อตฺถิ ตสฺส อสุจิ สุกฺกวิสฺสฏฺฐีติ, อามนฺตา.}

\vspace{\tipitakaparskip}
\csromancenter{ทุติยวคฺค ปรธมฺมกุสโล อรหา วีตราโค วีตโทโส วีตโมโห ฯเปฯ สจฺฉิกาตพฺพํ สจฺฉิกตํ, อตฺถิ ตสฺส อสุจิ สุกฺกวิสฺสฏฺฐีติ, น เหวํ วตฺตพฺเพ ฯเปฯ.}
\prose{\csromanfolio{133}\csromansymbolmark{*}ปรธมฺมกุสลสฺส อรหโต ราโค ปหีโน, นตฺถิ ตสฺส อสุจิ สุกฺกวิสฺสฏฺฐีติ, อามนฺตา.\csromanspacer{} สธมฺมกุสลสฺส อรหโต ราโค ปหีโน, นตฺถิ ตสฺส อสุจิ สุกฺกวิสฺสฏฺฐีติ, น เหวํ วตฺตพฺเพ ฯเปฯ.}

\prose{\csromansymbolmark{*}ปรธมฺมกุสลสฺส อรหโต โทโส ปหีโน ฯเปฯ โมโห ปหีโน ฯเปฯ อโนตฺตปฺปํ ปหีนํ, นตฺถิ ตสฺส อสุจิ สุกฺกวิสฺสฏฺฐีติ, อามนฺตา.\csromanspacer{} สธมฺมกุสลสฺส อรหโต อโนตฺตปฺปํ ปหีนํ, นตฺถิ ตสฺส อสุจิ สุกฺกวิสฺสฏฺฐีติ, น เหวํ วตฺตพฺเพ ฯเปฯ.}

\prose{\csromansymbolmark{*}ปรธมฺมกุสลสฺส อรหโต ราคปฺปหานาย มคฺโค ภาวิโต ฯเปฯ โพชฺฌงฺคา ภาวิตา ฯเปฯ โทสปฺปหานาย ฯเปฯ โมหปฺปหานาย ฯเปฯ อโนตฺตปฺปปหานาย มคฺโค ภาวิโต ฯเปฯ โพชฺฌงฺคา ภาวิตา, นตฺถิ ตสฺส อสุจิ สุกฺกวิสฺสฏฺฐีติ, อามนฺตา.\csromanspacer{} สธมฺมกุสลสฺส อรหโต อโนตฺตปฺปปหานาย โพชฺฌงฺคา ภาวิตา, นตฺถิ ตสฺส อสุจิ สุกฺกวิสฺสฏฺฐีติ, น เหวํ วตฺตพฺเพ ฯเปฯ. \csromansymbolmark{*}ปรธมฺมกุสโล อรหา วีตราโค วีตโทโส วีตโมโห ฯเปฯ สจฺฉิกาตพฺพํ สจฺฉิกตํ, นตฺถิ ตสฺส อสุจิ สุกฺกวิสฺสฏฺฐีติ, อามนฺตา.\csromanspacer{} สธมฺมกุสโล อรหา วีตราโค วีตโทโส วีตโมโห ฯเปฯ สจฺฉิกาตพฺพํ สจฺฉิกตํ, นตฺถิ ตสฺส อสุจิ สุกฺกวิสฺสฏฺฐีติ, น เหวํ วตฺตพฺเพ ฯเปฯ.}

\proseitem{313}{\csromansymbolmark{*}อตฺถิ อรหโต อสุจิ สุกฺกวิสฺสฏฺฐีติ, อามนฺตา.\csromanspacer{} นนุ วุตฺตํ ภควตา “เย เต\footnote{เย เกจิ (สี)} ภิกฺขเว ภิกฺขู ปุถุชฺชนา สีลสมฺปนฺนา สตา สมฺปชานา นิทฺทํ โอกฺกมนฺติ, เตสํ อสุจิ น มุจฺจติ.\csromanspacer{} เยปิ เต ภิกฺขเว พาหิรกา อิสโย กาเมสุ วีตราคา, เตสมฺปิ อสุจิ น มุจฺจติ.\csromanspacer{} อฏฺฐานเมตํ ภิกฺขเว อนวกาโส ยํ อรหโต อสุจิ มุจฺเจยฺยา”ติ\footnote{วิ 3. 409 ปิฏฺเฐ อตฺถโต สมานํ; อํ 2. 219 ปิฏฺเฐปิ ปสฺสิตพฺพํ.} อตฺเถว สุตฺตนฺโตติ, อามนฺตา.\csromanspacer{} เตน หิ น วตฺตพฺพํ “อตฺถิ อรหโต อสุจิ สุกฺกวิสฺสฏฺฐี”ติ.}

\csromanmatikaanchor{mtk.31}
\csromantocmarkref{subsection}{(11) 2. อญฺญาณกถา}{mtk.31}
\bookmark[dest={mtk.31},level=2]{(11) 2. อญฺญาณกถา}
\prose{\csromanfolio{134}\csromansymbolmark{+}น วตฺตพฺพํ “อตฺถิ อรหโต ปรูปหาโร”ติ, อามนฺตา.\csromanspacer{} นนุ อรหโต จีวรปิณฺฑปาตเสนาสนคิลานปจฺจยเภสชฺชปริกฺขารํ ปเร อุปสํหเรยฺยุนฺติ, อามนฺตา.\csromanspacer{} หญฺจิ อรหโต จีวรปิณฺฑปาตเสนาสนคิลานปจฺจยเภสชฺชปริกฺขารํ ปเร อุปสํหเรยฺยุํ, เตน วต เร วตฺตพฺเพ “อตฺถิ อรหโต ปรูปหาโร”ติ.}

\prose{\csromansymbolmark{*}อรหโต จีวรปิณฺฑปาตเสนาสนคิลานปจฺจยเภสชฺชปริกฺขารํ ปเร อุปสํหเรยฺยุนฺติ อตฺถิ อรหโต ปรูปหาโรติ, อามนฺตา.\csromanspacer{} อรหโต โสตาปตฺติผลํ วา สกทาคามิผลํ วา อนาคามิผลํ วา อรหตฺตํ วา ปเร อุปสํหเรยฺยุนฺติ, น เหวํ วตฺตพฺเพ ฯเปฯ.}

\nitthitam{ปรูปหารกถา นิฏฺฐิตา.}
\chapterhead{2. ทุติยวคฺค}
\vspace{\tipitakaparskip}
\csromancenter{\textbf{(11) 2.}\csromanspacer{}\textbf{ อญฺญาณกถา}}
\proseitem{314}{\csromansymbolmark{*}อตฺถิ อรหโต อญฺญาณนฺติ, อามนฺตา.\csromanspacer{} อตฺถิ อรหโต อวิชฺชา อวิชฺโชโฆ อวิชฺชาโยโค อวิชฺชานุสโย อวิชฺชาปริยุฏฺฐานํ อวิชฺชาสํโยชนํ อวิชฺชานีวรณนฺติ, น เหวํ วตฺตพฺเพ.\csromanspacer{} นตฺถิ อรหโต อวิชฺชา อวิชฺโชโฆ อวิชฺชาโยโค อวิชฺชานุสโย อวิชฺชาปริยุฏฺฐานํ อวิชฺชาสํโยชนํ อวิชฺชานีวรณนฺติ, อามนฺตา.\csromanspacer{} หญฺจิ นตฺถิ อรหโต อวิชฺชา อวิชฺโชโฆ อวิชฺชาโยโค อวิชฺชานุสโย อวิชฺชาปริยุฏฺฐานํ อวิชฺชาสํโยชนํ อวิชฺชานีวรณํ, โน จ วต เร วตฺตพฺเพ “อตฺถิ อรหโต อญฺญาณนฺ”ติ.}

\prose{\csromansymbolmark{*}อตฺถิ ปุถุชฺชนสฺส อญฺญาณํ, อตฺถิ ตสฺส อวิชฺชา อวิชฺโชโฆ อวิชฺชาโยโค อวิชฺชานุสโย อวิชฺชาปริยุฏฺฐานํ อวิชฺชาสํโยชนํ อวิชฺชานีวรณนฺติ, อามนฺตา.\csromanspacer{} อตฺถิ อรหโต อญฺญาณํ, อตฺถิ ตสฺส อวิชฺชา อวิชฺโชโฆ อวิชฺชาโยโค อวิชฺชานุสโย อวิชฺชาปริยุฏฺฐานํ อวิชฺชาสํโยชนํ อวิชฺชานีวรณนฺติ, น เหวํ วตฺตพฺเพ.}

\prose{\csromansymbolmark{*}อตฺถิ อรหโต อญฺญาณํ, นตฺถิ ตสฺส อวิชฺชา อวิชฺโชโฆ อวิชฺชาโยโค อวิชฺชานุสโย อวิชฺชาปริยุฏฺฐานํ อวิชฺชาสํโยชนํ}

\vspace{\tipitakaparskip}
\csromancenter{ทุติยวคฺค อวิชฺชานีวรณนฺติ, อามนฺตา.\csromanspacer{} อตฺถิ ปุถุชฺชนสฺส อญฺญาณํ, นตฺถิ ตสฺส อวิชฺชา อวิชฺโชโฆ อวิชฺชาโยโค อวิชฺชานุสโย อวิชฺชาปริยุฏฺฐานํ อวิชฺชาสํโยชนํ อวิชฺชานีวรณนฺติ, น เหวํ วตฺตพฺเพ.}
\prose{\csromanfolio{135}\csromansymbolmark{*}อตฺถิ อรหโต อญฺญาณนฺติ, อามนฺตา.\csromanspacer{} อรหา อญฺญาณปกโต ปาณํ หเนยฺย.\csromanspacer{} อทินฺนํ อาทิเยยฺย.\csromanspacer{} มุสา ภเณยฺย.\csromanspacer{} ปิสุณํ ภเณยฺย.\csromanspacer{} ผรุสํ ภเณยฺย.\csromanspacer{} สมฺผํ ปลเปยฺย.\csromanspacer{} สนฺธิํ ฉินฺเทยฺย.\csromanspacer{} นิลฺโลปํ หเรยฺย.\csromanspacer{} เอกาคาริยํ กเรยฺย.\csromanspacer{} ปริปนฺเถ ติฏฺเฐยฺย.\csromanspacer{} ปรทารํ คจฺเฉยฺย.\csromanspacer{} คามฆาตํ กเรยฺย.\csromanspacer{} นิคมฆาตํ กเรยฺยาติ, น เหวํ วตฺตพฺเพ.}

\prose{\csromansymbolmark{*}อตฺถิ ปุถุชฺชนสฺส อญฺญาณํ ปุถุชฺชโน อญฺญาณปกโต ปาณํ หเนยฺย.\csromanspacer{} อทินฺนํ อาทิเยยฺย.\csromanspacer{} มุสา ภเณยฺย ฯเปฯ คามฆาตํ กเรยฺย.\csromanspacer{} นิคมฆาตํ กเรยฺยาติ, อามนฺตา.\csromanspacer{} อตฺถิ อรหโต อญฺญาณํ อรหา อญฺญาณปกโต ปาณํ หเนยฺย.\csromanspacer{} อทินฺนํ อาทิเยยฺย ฯเปฯ คามฆาตํ กเรยฺย.\csromanspacer{} นิคมฆาตํ กเรยฺยาติ, น เหวํ วตฺตพฺเพ.}

\prose{\csromansymbolmark{*}อตฺถิ อรหโต อญฺญาณํ น จ อรหา อญฺญาณปกโต ปาณํ หเนยฺย.\csromanspacer{} อทินฺนํ อาทิเยยฺย ฯเปฯ คามฆาตํ กเรยฺย.\csromanspacer{} นิคมฆาตํ กเรยฺยาติ, อามนฺตา.\csromanspacer{} อตฺถิ ปุถุชฺชนสฺส อญฺญาณํ, น จ ปุถุชฺชโน อญฺญาณปกโต ปาณํ หเนยฺย.\csromanspacer{} อทินฺนํ อาทิเยยฺย ฯเปฯ คามฆาตํ กเรยฺย.\csromanspacer{} นิคมฆาตํ กเรยฺยาติ, น เหวํ วตฺตพฺเพ.}

\prose{\csromansymbolmark{*}อตฺถิ อรหโต อญฺญาณนฺติ, อามนฺตา.\csromanspacer{} อตฺถิ อรหโต สตฺถริ อญฺญานํ, ธมฺเม อญฺญาณํ, สํเฆ อญฺญาณํ สิกฺขาย อญฺญาณํ, ปุพฺพนฺเต อญฺญาณํ, อปรนฺเต อญฺญาณํ, ปุพฺพนฺตาปรนฺเต อญฺญาณํ, อิทปฺปจฺจยตาปฏิจฺจสมุปฺปนฺเนสุ ธมฺเมสุ อญฺญาณนฺติ, น เหวํ วตฺตพฺเพ.}

\prose{\csromansymbolmark{*}นตฺถิ อรหโต สตฺถริ อญฺญาณํ, ธมฺเม อญฺญาณํ, สํเฆ อญฺญาณํ, สิกฺขาย อญฺญาณํ, ปุพฺพนฺเต อญฺญาณํ, อปรนฺเต อญฺญาณํ, ปุพฺพนฺตาปรนฺเต อญฺญาณํ, อิทปฺปจฺจยตาปฏิจฺจสมุปฺปนฺเนสุ ธมฺเมสุ อญฺญาณนฺติ, อามนฺตา.\csromanspacer{} หญฺจิ นตฺถิ อรหโต สตฺถริ อญฺญาณํ, ธมฺเม อญฺญาณํ, สํเฆ อญฺญาณํ ฯเปฯ อิทปฺปจฺจยตาปฏิจฺจ สมุปฺปนฺเนสุ ธมฺเมสุ อญฺญาณํ, โน จ วต เร วตฺตพฺเพ “อตฺถิ อรหโต อญฺญาณนฺ”ติ. \csromansymbolmark{*}อตฺถิ ปุถุชฺชนสฺส อญฺญาณํ, อตฺถิ ตสฺส สตฺถริ อญฺญาณํ, ธมฺเม อญฺญาณํ, สํเฆ อญฺญาณํ ฯเปฯ อิทปฺปจฺจยตาปฏิจฺจสมุปฺปนฺเนสุ ธมฺเมสุ \csromanfolio{136}อญฺญาณนฺติ, อามนฺตา.\csromanspacer{} อตฺถิ อรหโต อญฺญาณํ, อตฺถิ ตสฺส สตฺถริ อญฺญาณํ,, ธมฺเม อญฺญาณํ, สํเฆ อญฺญาณํ ฯเปฯ อิทปฺปจฺจยตาปฏิจฺจสมุปฺปนฺเนสุ ธมฺเมสุ อญฺญาณนฺติ, น เหวํ วตฺตพฺเพ.}

\prose{\csromansymbolmark{*}อตฺถิ อรหโต อญฺญาณํ, นตฺถิ ตสฺส สตฺถริ อญฺญาณํ, ธมฺเม อญฺญาณํมฺ สํเฆ อญฺญาณํ ฯเปฯ อิทปฺปจฺจยตาปฏิจฺจสมุปฺปนฺเนสุ ธมฺเมสุ อญฺญาณนฺติ, อามนฺตา.\csromanspacer{} อตฺถิ ปุถุชฺชนสฺส อญฺญาณํ, นตฺถิ ตสฺส สตฺถริ อญฺญาณํ, ธมฺเม อญฺญาณํ, สํเฆ อญฺญาณํ ฯเปฯ อิทปฺปจฺจยตาปฏิจฺจสมุปฺปนฺเนสุ ธมฺเมสุ อญฺญาณนฺติ, น เหวํ วตฺตพฺเพ.}

\proseitem{315}{\csromansymbolmark{*}อตฺถิ อรหโต อญฺญาณนฺติ, อามนฺตา.\csromanspacer{} นนุ อรหโต ราโค ปหีโน อุจฺฉินฺนมูโล ตาลาวตฺถุกโต อนภาวํกโต อายติํ อนุปฺปาทธมฺโมติ, อามนฺตา.\csromanspacer{} หญฺจิ อรหโต ราโค ปหีโน อุจฺฉินฺนมูโล ตาลาวตฺถุกโต อนภาวํกโต อายติํ อนุปฺปาทธมฺโม, โน จ วต เร วตฺตพฺเพ “อตฺถิ อรหโต อญฺญาณนฺ”ติ.}

\prose{\csromansymbolmark{*}อตฺถิ อรหโต อญฺญาณนฺติ, อามนฺตา.\csromanspacer{} นนุ อรหโต โทโส ปหีโน ฯเปฯ โมโห ปหีโน ฯเปฯ อโนตฺตปฺปํ ปหีนํ อุจฺฉินฺนมูลํ ตาลาวตฺถุกตํ อนภาวํกตํ อายติํ อนุปฺปาทธมฺมนฺติ, อามนฺตา.\csromanspacer{} หญฺจิ อรหโต อโนตฺตปฺปํ ปหีนํ อุจฺฉินฺนมูลํ ตาลาวตฺถุกตํ อนภาวํกตํ อายติํ อนุปฺปาทธมฺมํ, โน จ วต เร วตฺตพฺเพ “อตฺถิ อรหโต อญฺญาณนฺ”ติ.}

\prose{\csromansymbolmark{*}อตฺถิ อรหโต อญฺญาณนฺติ, อามนฺตา.\csromanspacer{} นนุ อรหโต ราคปฺปหานาย มคฺโค ภาวิโต ฯเปฯ โพชฺฌงฺคา ภาวิตาติ, อามนฺตา.\csromanspacer{} หญฺจิ อรหโต ราคปฺปหานาย โพชฺฌงฺคา ภาวิตา, โน จ วต เร วตฺตพฺเพ “อตฺถิ อรหโต อญฺญาณนฺ”ติ.}

\prose{\csromansymbolmark{*}อตฺถิ อรหโต อญฺญาณนฺติ, อามนฺตา.\csromanspacer{} นนุ อรหโต โทสปฺปหานาย ฯเปฯ โมหปฺปหานาย ฯเปฯ.\csromanspacer{} อโนตฺตปฺปปหานาย มคฺโค ภาวิโต ฯเปฯ โพชฺฌงฺคา ภาวิตาติ, อามนฺตา.\csromanspacer{} หญฺจิ อรหโต อโนตฺตปฺปปหานาย โพชฺฌงฺคา ภาวิตา, โน จ วต เร วตฺตพฺเพ “อตฺถิ อรหโต อญฺญาณนฺ”ติ.}

\chapterheadpageread{2. ทุติยวคฺค}
\prose{\csromanfolio{137}\csromansymbolmark{*}อตฺถิ อรหโต อญฺญาณนฺติ, อามนฺตา.\csromanspacer{} นนุ อรหา วีตราโค วีตโทโส วีตโมโห ฯเปฯ สจฺฉิกาตพฺพํ สจฺฉิกตนฺติ, อามนฺตา.\csromanspacer{} หญฺจิ อรหา วีตราโค ฯเปฯ สจฺฉิกาตพฺพํ สจฺฉิกตํ, โน จ วต เร วตฺตพฺเพ “อตฺถิ อรหโต อญฺญาณนฺ”ติ.}

\proseitem{316}{\csromansymbolmark{*}อตฺถิ อรหโต อญฺญาณนฺติ? สธมฺมกุสลสฺส อรหโต อตฺถิ อญฺญาณํ, ปรธมฺมกุสลสฺส อรหโต นตฺถิ อญฺญาณนฺติ.\csromanspacer{} สธมฺมกุสลสฺส อรหโต อตฺถิ อญฺญาณนฺติ, อามนฺตา.\csromanspacer{} ปรธมฺมกุสลสฺส อรหโต อตฺถิ อญฺญาณนฺติ, น เหวํ วตฺตพฺเพ ฯเปฯ.}

\prose{\csromansymbolmark{*}ปรธมฺมกุสลสฺส อรหโต นตฺถิ อญฺญาณนฺติ, อามนฺตา.\csromanspacer{} สธมฺมกุสลสฺส อรหโต นตฺถิ อญฺญาณนฺติ, น เหวํ วตฺตพฺเพ ฯเปฯ.}

\prose{\csromansymbolmark{*}สธมฺมกุสลสฺส อรหโต ราโค ปหีโน, อตฺถิ ตสฺส อญฺญาณนฺติ, อามนฺตา.\csromanspacer{} ปรธมฺมกุสลสฺส อรหโต ราโค ปหีโน, อตฺถิ ตสฺส อญฺญาณนฺติ, น เหวํ วตฺตพฺเพ ฯเปฯ.}

\prose{\csromansymbolmark{*}สธมฺมกุสลสฺส อรหโต โทโส ปหีโน ฯเปฯ โมโห ปหีโน ฯเปฯ อโนตฺตปฺปํ ปหีนํ, อตฺถิ ตสฺส อญฺญาณนฺติ, อามนฺตา.\csromanspacer{} ปรธมฺมกุสลสฺส อรหโต อโนตฺตปฺปํ ปหีนํ, อตฺถิ ตสฺส อญฺญาณนฺติ, น เหวํ วตฺตพฺเพ ฯเปฯ.}

\prose{\csromansymbolmark{*}สธมฺมกุสลสฺส อรหโต ราคปฺปหานาย มคฺโค ภาวิโต ฯเปฯ โพชฺฌงฺคา ภาวิตา, อตฺถิ ตสฺส อญฺญาณนฺติ, อามนฺตา.\csromanspacer{} ปรธมฺมกุสลสฺส อรหโต ราคปฺปหานาย โพชฺฌงฺคา ภาวิตา, อตฺถิ ตสฺส อญฺญาณนฺติ, น เหวํ วตฺตพฺเพ ฯเปฯ.}

\prose{\csromansymbolmark{*}สธมฺมกุสลสฺส อรหโต โทสปฺปหานาย ฯเปฯ โมหปฺปหานาย ฯเปฯ อโนตฺตปฺปปหานาย มคฺโค ภาวิโต ฯเปฯ โพชฺฌงฺคา ภาวิตา, อตฺถิ ตสฺส อญฺญาณนฺติ, อามนฺตา.\csromanspacer{} ปรธมฺมกุสลสฺส อรหโต อโนตฺตปฺปปหานาย โพชฺฌงฺคา ภาวิตา, อตฺถิ ตสฺส อญฺญาณนฺติ, น เหวํ วตฺตพฺเพ ฯเปฯ.}

\prose{\csromansymbolmark{*}สธมฺมกุสโล อรหา วีตราโค วีตโทโส วีตโมโห ฯเปฯ สจฺฉิกาตพฺพํ สจฺฉิกตํ, อตฺถิ ตสฺส อญฺญาณนฺติ, อามนฺตา.\csromanspacer{} ปรธมฺมกุสโล อรหา วีตราโค วีตโทโส วีตโมโห ฯเปฯ สจฺฉิกาตพฺพํ สจฺฉิกตํ, อตฺถิ ตสฺส อญฺญาณนฺติ, น เหวํ วตฺตพฺเพ ฯเปฯ.}

\prose{\csromanfolio{138}\csromansymbolmark{*}ปรธมฺมกุสลสฺส อรหโต ราโค ปหีโน, นตฺถิ ตสฺส อญฺญาณนฺติ, อามนฺตา.\csromanspacer{} สธมฺมกุสลสฺส อรหโต ราโค ปหีโน, นตฺถิ ตสฺส อญฺญาณนฺติ, น เหวํ วตฺตพฺเพ ฯเปฯ.}

\prose{\csromansymbolmark{*}ปรธมฺมกุสลสฺส อรหโต โทโส ปหีโน ฯเปฯ โมโห ปหีโน ฯเปฯ อโนตฺตปฺปํ ปหีนํ, นตฺถิ ตสฺส อญฺญาณนฺติ, อามนฺตา.\csromanspacer{} สธมฺมกุสลสฺส อรหโต อโนตฺตปฺปํ ปหีนํมฺ นตฺถิ ตสฺส อญฺญาณนฺติ, น เหวํ วตฺตพฺเพ ฯเปฯ.}

\prose{\csromansymbolmark{*}ปรธมฺมกุสลสฺส อรหโต ราคปฺปหานาย มคฺโค ภาวิโต ฯเปฯ โพชฺฌงฺคา ภาวิตา ฯเปฯ โทสปฺปหานาย.\csromanspacer{} โมหปฺปหานาย ฯเปฯ อโนตฺตปฺปปหานาย มคฺโค ภาวิโต ฯเปฯ โพชฺฌงฺคา ภาวิตา, นตฺถิ ตสฺส อญฺญาณนฺติ, อามนฺตา.\csromanspacer{} สธมฺมกุสลสฺส อรหโต อโนตฺตปฺปปหานาย โพชฺฌงฺคา ภาวิตา, นตฺถิ ตสฺส อญฺญาณนฺติ, น เหวํ วตฺตพฺเพ ฯเปฯ.}

\prose{\csromansymbolmark{*}ปรธมฺมกุสโล อรหา วีตราโค วีตโทโส วีตโมโห ฯเปฯ สจฺฉิกาตพฺพํ สจฺฉิกตํ, นตฺถิ ตสฺส อญฺญาณนฺติ, อามนฺตา.\csromanspacer{} สธมฺมกุสโล อรหา วีตราโค วีตโทโส วีตโมโห ฯเปฯ สจฺฉิกาตพฺพํ สจฺฉิกตํ, นตฺถิ ตสฺส อญฺญาณนฺติ, น เหวํ วตฺตพฺเพ ฯเปฯ.}

\proseitem{317}{\csromansymbolmark{*}อตฺถิ อรหโต อญฺญาณนฺติ, อามนฺตา.\csromanspacer{} นนุ วุตฺตํ ภควตา “ชานโตหํ\footnote{ชานตาหํ (สี), ชานตฺวาหํ (สฺยา, อิ, ก)} ภิกฺขเว ปสฺสโต อาสวานํ ขยํ วทามิ, โน อชานโต โน อปสฺสโต.\csromanspacer{} กิญฺจิ ภิกฺขเว ชานโต กิํ ปสฺสโต อาสวานํ ขโย โหติ, ‘อิติ รูปํ, อิติ รูปสฺส สมุทโย, อิติ รูปสฺส อตฺถงฺคโม.\csromanspacer{} อิติ เวทนา ฯเปฯ.\csromanspacer{} อิติ สญฺญา.\csromanspacer{} อิติ สงฺขารา.\csromanspacer{} อิติ วิญฺญาณํ, อิติ วิญฺญาณสฺส สมุทโย, อิติ วิญฺญาณสฺส อตฺถงฺคโมติ.\csromanspacer{} เอวํ โข ภิกฺขเว ชานโต เอวํ ปสฺสโต อาสวานํ ขโย โหตี”ติ\footnote{สํ 2. 124; ขุ 1. 265; สํ 3. 380; ม 1. 9 ปิฏฺเฐสุ.} อตฺเถว สุตฺตนฺโตติ, อามนฺตา.\csromanspacer{} เตน หิ น วตฺตพฺพํ “อตฺถิ อรหโต อญฺญาณนฺ”ติ.}

\prose{\csromansymbolmark{*}อตฺถิ อรหโต อญฺญาณนฺติ, อามนฺตา.\csromanspacer{} นนุ วุตฺตํ ภควตา “ชานโตหํ ภิกฺขเว ปสฺสโต อาสวานํ ขยํ วทามิ, โน อชานโต โน อปสฺสโต.\csromanspacer{} กิญฺจ ภิกฺขเว ชานโต กิํ ปสฺสโต อาสวานํ}

\vspace{\tipitakaparskip}
\csromancenter{ทุติยวคฺค ขโย โหติ, ‘อิทํ ทุกฺขนฺ’ติ ภิกฺขเว ชานโต ปสฺสโต อาสวานํ ขโย โหติ, ‘อยํ ทุกฺขสมุทโย’ติ ชานโต ปสฺสโต อาสวานํ ขโย โหติ, ‘อยํ ทุกฺขนิโรโธ’ติ ชานโต ปสฺสโต อาสวานํ ขโย โหติ, ‘อยํ ทุกฺขนิโรธคามินี ปฏิปทา’ติ ชานโต ปสฺสโต อาสวานํ ขโย โหติ.\csromanspacer{} เอวํ โข ภิกฺขเว ชานโต เอวํ ปสฺสโต อาสวานํ ขโย โหตี”ติ\footnote{สํ 3. 380 ปิฏฺเฐ.} อตฺเถว สุตฺตนฺโตติ, อามนฺตา.\csromanspacer{} เตน หิ น วตฺตพฺพํ “อตฺถิ อรหโต อญฺญาณนฺ”ติ.}
\prose{\csromanfolio{139}\csromansymbolmark{*}อตฺถิ อรหโต อญฺญาณนฺติ, อามนฺตา.\csromanspacer{} นนุ วุตฺตํ ภควตา “สพฺพํ ภิกฺขเว อนภิชานํ อปริชานํ อวิราชยํ อปฺปชหํ อภพฺโพ ทุกฺขกฺขยาย, สพฺพญฺจ โข ภิกฺขเว อภิชานํ ปริชานํ วิราชยํ ปชหํ ภพฺโพ ทุกฺขกฺขยายา”ติ\footnote{สํ 2. 249 ปิฏฺเฐ; ขุ 1. 197 อิติวุตฺตเกปิ.} อตฺเถว สุตฺตนฺโตติ, อามนฺตา.\csromanspacer{} เตน หิ น วตฺตพฺพํ “อตฺถิ อรหโต อญฺญาณนฺ”ติ.}

\prose{\csromansymbolmark{*}อตฺถิ อรหโต อญฺญาณนฺติ, อามนฺตา.\csromanspacer{} นนุ วุตฺตํ ภควตา\csromandash{}}

\csromangathameasure{“สหาวสฺส ทสฺสนสมฺปทาย, \\ ตยสฺสุ ธมฺมา ชหิตา ภวนฺติ. \\ สกฺกายทิฏฺฐี วิจิกิจฺฉิตญฺจ, \\ สีลพฺพตํ วาปิ ยทตฺถิ กิญฺจิ. \\ จตูหปาเยหิ จ วิปฺปมุตฺโต, \\ ฉจฺจาภิฐานานิ อภพฺพ กาตุนฺ”ติ.}
\csromangathagroup{\csromangathastanza{“สหาวสฺส ทสฺสนสมฺปทาย, \\ ตยสฺสุ ธมฺมา ชหิตา ภวนฺติ. \\ สกฺกายทิฏฺฐี วิจิกิจฺฉิตญฺจ, \\ สีลพฺพตํ วาปิ ยทตฺถิ กิญฺจิ. \\ จตูหปาเยหิ จ วิปฺปมุตฺโต, \\ ฉจฺจาภิฐานานิ อภพฺพ กาตุนฺ”ติ.}}

\prose{อตฺเถว สุตฺตนฺโตติ, อามนฺตา.\csromanspacer{} เตน หิ น วตฺตพฺพํ “อตฺถิ อรหโต อญฺญาณนฺ”ติ.}

\prose{\csromansymbolmark{*}อตฺถิ อรหโต อญฺญาณนฺติ, อามนฺตา.\csromanspacer{} นนุ วุตฺตํ ภควตา “ยสฺมิํ ภิกฺขเว สมเย อริยสาวกสฺส วิรชํ วีตมลํ ธมฺมจกฺขุํ อุทปาทิ ‘ยํ กิญฺจิ สมุทยธมฺมํ, สพฺพํ ตํ นิโรธธมฺมนฺ’ติ.\csromanspacer{} สห ทสฺสนุปฺปาทา ภิกฺขเว อริยสาวกสฺส ตีณิ สํโยชนานิ ปหียนฺติ, สกฺกายทิฏฺฐิ วิจิกิจฺฉา สีลพฺพตปรามาโส”ติ อตฺเถว สุตฺตนฺโตติ, อามนฺตา.\csromanspacer{} เตน หิ น วตฺตพฺพํ “อตฺถิ อรหโต อญฺญาณนฺ”ติ.}

\csromanmatikaanchor{mtk.32}
\csromantocmarkref{subsection}{(12) 3. กงฺขากถา}{mtk.32}
\bookmark[dest={mtk.32},level=2]{(12) 3. กงฺขากถา}
\prose{\csromanfolio{140}\csromansymbolmark{+}น วตฺตพฺพํ “อตฺถิ อรหโต อญฺญาณนฺ”ติ, อามนฺตา.\csromanspacer{} นนุ อรหา อิตฺถิปุริสานํ นามโคตฺตํ น ชาเนยฺย, มคฺคามคฺคํ น ชาเนยฺย, ติณกฏฺฐวนปฺปตีนํ นามํ น ชาเนยฺยาติ, อามนฺตา.\csromanspacer{} หญฺจิ อรหา อิตฺถิปุริสานํ นามคฺโคตฺตํ น ชาเนยฺย, มคฺคามคฺคํ น ชาเนยฺย, ติณกฏฺฐวนปฺปตีนํ นามํ น ชาเนยฺย, เตน วต เร วตฺตพฺเพ “อตฺถิ อรหโต อญฺญาณนฺ”ติ.}

\prose{\csromansymbolmark{*}อรหา อิตฺถิปุริสานํ นามโคตฺตํ น ชาเนยฺย, มคฺคามคฺคํ น ชาเนยฺย, ติณกฏฺฐวนปฺปตีนํ นามํ น ชาเนยฺยาติ, อตฺถิ อรหโต อญฺญาณนฺติ, อามนฺตา.\csromanspacer{} อรหา โสตาปตฺติผลํ วา สกทคามิผลํ วา อนาคามิผลํ วา อรหตฺตํ วา น ชาเนยฺยาติ, น เหวํ วตฺตพฺเพ ฯเปฯ.}

\nitthitam{อญฺญาณกถา นิฏฺฐิตา.}
\chapterhead{2. ทุติยวคฺค}
\vspace{\tipitakaparskip}
\csromancenter{\textbf{(12) 3.}\csromanspacer{}\textbf{ กงฺขากถา}}
\proseitem{318}{\csromansymbolmark{*}อตฺถิ อรหโต กงฺขาติ, อามนฺตา.\csromanspacer{} อตฺถิ อรหโต วิจิกิจฺฉา วิจิกิจฺฉาปริยุฏฺฐานํ วิจิกิจฺฉาสํโยชนํ วิจิกิจฺฉานีวรณนฺติ, น เหวํ วตฺตพฺเพ ฯเปฯ.}

\prose{\csromansymbolmark{*}นตฺถิ อรหโต วิจิกิจฺฉา วิจิกิจฺฉาปริยุฏฺฐานํ วิจิกิจฺฉาสํโยชนํ วิจิกิจฺฉานีวรณนฺติ, อามนฺตา.\csromanspacer{} หญฺจิ นตฺถิ อรหโต วิจิกิจฺฉา วิจิกิจฺฉาปริยุฏฺฐานํ วิจิกิจฺฉาสํโยชนํ วิจิกิจฺฉานีวรณํ, โน จ วต เร วตฺตพฺเพ “อตฺถิ อรหโต กงฺขา”ติ.}

\prose{\csromansymbolmark{*}อตฺถิ ปุถุชฺชนสฺส กงฺขา, อตฺถิ ตสฺส วิจิกิจฺฉา วิจิกิจฺฉาปริยุฏฺฐานํ วิจิกิจฺฉาสํโยชนํ วิจิกิจฺฉานีวรณนฺติ, อามนฺตา.\csromanspacer{} อตฺถิ อรหโต กงฺขา, อตฺถิ ตสฺส วิจิกิจฺฉา วิจิกิจฺฉาปริยุฏฺฐานํ วิจิกิจฺฉาสํโยชนํ วิจิกิจฺฉานีวรณนฺติ, น เหวํ วตฺตพฺเพ ฯเปฯ.}

\prose{\csromansymbolmark{*}อตฺถิ อรหโต กงฺขา, นตฺถิ ตสฺส วิจิกิจฺฉา วิจิกิจฺฉาปริยุฏฺฐานํ วิจิกิจฺฉาสํโยชนํ วิจิกิจฺฉานีวรณนฺติ, อามนฺตา.\csromanspacer{} อตฺถิ ปุถุชฺชนสฺส กงฺขา,}

\vspace{\tipitakaparskip}
\csromancenter{ทุติยวคฺค นตฺถิ ตสฺส วิจิกิจฺฉา วิจิกิจฺฉาปริยุฏฺฐานํ วิจิกิจฺฉาสํโยชนํ วิจิกิจฺฉานีวรณนฺติ, น เหวํ วตฺตพฺเพ ฯเปฯ.}
\prose{\csromanfolio{141}\csromansymbolmark{*}อตฺถิ อรหโต กงฺขาติ, อามนฺตา.\csromanspacer{} อตฺถิ อรหโต สตฺถริ กงฺขา, ธมฺเม กงฺขา, สํเฆ กงฺขา, สิกฺขาย กงฺขา, ปุพฺพนฺเต กงฺขา, อปรนฺเต กงฺขา, ปุพฺพนฺตาปรนฺเต กงฺขา, อิทปฺปจฺจยตาปฏิจฺจสมุปฺปนฺเนสุ ธมฺเมสุ กงฺขาติ, น เหวํ วตฺตพฺเพ.}

\prose{\csromansymbolmark{*}นตฺถิ อรหโต สตฺถริ กงฺขา ธมฺเม กงฺขา ฯเปฯ อิทปฺปจฺจยตาปฏิจฺจ สมุปฺปนฺเนสุ ธมฺเมสุ กงฺขาติ, อามนฺตา.\csromanspacer{} หญฺจิ นตฺถิ อรหโต สตฺถริ กงฺขา, ธมฺเม กงฺขา ฯเปฯ อิทปฺปจฺจยตาปฏิจฺจสมุปฺปนฺเนสุ ธมฺเมสุ กงฺขา, โน จ วต เร วตฺตพฺเพ “อตฺถิ อรหโต กงฺขา”ติ.}

\prose{\csromansymbolmark{*}อตฺถิ ปุถุชฺชนสฺส กงฺขา, อตฺถิ ตสฺส สตฺถริ กงฺขา, ธมฺเม กงฺขา ฯเปฯ อิทปฺปจฺจยตาปฏิจฺจสมุปฺปนฺเนสุ ธมฺเมสุ กงฺขาติ, อามนฺตา.\csromanspacer{} อตฺถิ อรหโต กงฺขา, อตฺถิ ตสฺส สตฺถริ กงฺขา, ธมฺเม กงฺขา ฯเปฯ อิทปฺปจฺจยตาปฏิจฺจสมุปฺปนฺเนสุ ธมฺเมสุ กงฺขาติ, น เหวํ วตฺตพฺเพ.}

\prose{\csromansymbolmark{*}อตฺถิ อรหโต กงฺขา, นตฺถิ ตสฺส สตฺถริ กงฺขา, ธมฺเม กงฺขา ฯเปฯ อิทปฺปจฺจยตาปฏิจฺจสมุปฺปนฺเนสุ ธมฺเมสุ กงฺขาติ, อามนฺตา.\csromanspacer{} อตฺถิ ปุถุชฺชนสฺส กงฺขา, นตฺถิ ตสฺส สตฺถริ กงฺขา, ธมฺเม กงฺขา ฯเปฯ อิทปฺปจฺจยตาปฏิจฺจสมุปฺปนฺเนสุ ธมฺเมสุ กงฺขาติ, น เหวํ วตฺตพฺเพ.}

\proseitem{319}{\csromansymbolmark{*}อตฺถิ อรหโต กงฺขาติ, อามนฺตา.\csromanspacer{} นนุ อรหโต ราโค ปหีโน อุจฺฉินฺนมูโล ตาลาวตฺถุกโต อนภาวํกโต อายติํ อนุปฺปาทธมฺโมติ, อามนฺตา.\csromanspacer{} หญฺจิ อรหโต ราโค ปหีโน อุจฺฉินฺนมูโล ตาลาวตฺถุกโต อนภาวํกโต อายติํ อนุปฺปาทธมฺโม, โน จ วต เร วตฺตพฺเพ “อตฺถิ อรหโต กงฺขา”ติ.}

\prose{\csromansymbolmark{*}อตฺถิ อรหโต กงฺขาติ, อามนฺตา.\csromanspacer{} นนุ อรหโต โทโส ปหีโน ฯเปฯ โมโห ปหีโน ฯเปฯ อโนตฺตปฺปํ ปหีนํ ฯเปฯ ราคปฺปหานาย มคฺโค ภาวิโต ฯเปฯ โพชฺฌงฺคา ภาวิตา ฯเปฯ โทสปฺปหานาย ฯเปฯ อโนตฺตปฺปปหานาย มคฺโค ภาวิโต ฯเปฯ โพชฺฌงฺคา ภาวิตา.\csromanspacer{} นนุ อรหา วีตราโค วีตโทโส วีตโมโห ฯเปฯ สจฺฉิกาตพฺพํ สจฺฉิกตนฺติ, อามนฺตา.\csromanspacer{} หญฺจิ อรหา วีตราโค วีตโทโส วีตโมโห ฯเปฯ สจฺฉิกาตพฺพํ สจฺฉิกตํ, โน จ วต เร วตฺตพฺเพ “อตฺถิ อรหโต กงฺขา”ติ.}

\proseitem{320}{\csromanfolio{142}\csromansymbolmark{*}อตฺถิ อรหโต กงฺขาติ? สธมฺมกุสลสฺส อรหโต อตฺถิ กงฺขา, ปรธมฺมกุสลสฺส อรหโต นตฺถิ กงฺขาติ.\csromanspacer{} สธมฺมกุสลสฺส อรหโต อตฺถิ กงฺขาติ, อามนฺตา.\csromanspacer{} ปรธมฺมกุสลสฺส อรหโต อตฺถิ กงฺขาติ, น เหวํ วตฺตพฺเพ ฯเปฯ.}

\prose{\csromansymbolmark{*}ปรธมฺมกุสลสฺส อรหโต นตฺถิ กงฺขาติ, อามนฺตา.\csromanspacer{} สธมฺมกุสลสฺส อรหโต นตฺถิ กงฺขาติ, น เหวํ วตฺตพฺเพ ฯเปฯ.}

\prose{\csromansymbolmark{*}สธมฺมกุสลสฺส อรหโต ราโค ปหีโน, อตฺถิ ตสฺส กงฺขาติ, อามนฺตา.\csromanspacer{} ปรธมฺมกุสลสฺส อรหโต ราโค ปหีโน, อตฺถิ ตสฺส กงฺขาติ, น เหวํ วตฺตพฺเพ ฯเปฯ.}

\prose{\csromansymbolmark{*}สธมฺมกุสลสฺส อรหโต โทโส ปหีโน ฯเปฯ โมโห ปหีโน ฯเปฯ อโนตฺตปฺปํ ปหีนํ ฯเปฯ ราคปฺปหานาย มคฺโค ภาวิโต ฯเปฯ โพชฺฌงฺคา ภาวิตา ฯเปฯ โทสปฺปหานาย ฯเปฯ โมหปฺปหานาย ฯเปฯ อโนตฺตปฺปปหานาย มคฺโค ภาวิโต ฯเปฯ โพชฺฌงฺคา ภาวิตา ฯเปฯ.\csromanspacer{} สธมฺมกุสโล อรหา วีตราโค วีตโทโส วีตโมโห ฯเปฯ สจฺฉิกาตพฺพํ สจฺฉิกตํ, อตฺถิ ตสฺส กงฺขาติ, อามนฺตา.\csromanspacer{} ปรธมฺมกุสโล อรหา วีตราโค วีตโทโส วีตโมโห ฯเปฯ สจฺฉิกาตพฺพํ สจฺฉิกตํ, อตฺถิ ตสฺส กงฺขาติ, น เหวํ วตฺตพฺเพ ฯเปฯ.\csromanspacer{} ปรธมฺมกุสลสฺส อรหโต ราโค ปหีโน, นตฺถิ ตสฺส กงฺขาติ, อามนฺตา.\csromanspacer{} สธมฺมกุสลสฺส อรหโต ราโค ปหีโน, นตฺถิ ตสฺส กงฺขาติ, น เหวํ วตฺตพฺเพ ฯเปฯ.}

\prose{\csromansymbolmark{*}ปรธมฺมกุสลสฺส อรหโต โทโส ปหีโน ฯเปฯ โมโห ปหีโน ฯเปฯ อโนตฺตปฺปํ ปหีนํ ฯเปฯ ราคปฺปหานาย มคฺโค ภาวิโต ฯเปฯ โพชฺฌงฺคา ภาวิตา ฯเปฯ โทสปฺปหานาย ฯเปฯ โมหปฺปหานาย ฯเปฯ อโนตฺตปฺปปหานาย มคฺโค ภาวิโต ฯเปฯ โพชฺฌงฺคา ภาวิตา ฯเปฯ.\csromanspacer{} ปรธมฺมกุสโล อรหา วีตราโค วีตโทโส วีตโมโห ฯเปฯ สจฺฉิกาตพฺพํ สจฺฉิกตํ, นตฺถิ ตสฺส กงฺขาติ, อามนฺตา.\csromanspacer{} สธมฺมกุสโล อรหา วีตราโค วีตโทโส วีตโมโห ฯเปฯ สจฺฉิกาตพฺพํ สจฺฉิกตํ, นตฺถิ ตสฺส กงฺขาติ, น เหวํ วตฺตพฺเพ ฯเปฯ.}

\proseitem{321}{\csromansymbolmark{*}อตฺถิ อรหโต กงฺขาติ, อามนฺตา.\csromanspacer{} นนุ วุตฺตํ ภควตา “ชานโตหํ ภิกฺขเว ปสฺสโต อาสวานํ ขยํ วทามิ, โน อชานโต}

\vspace{\tipitakaparskip}
\csromancenter{ทุติยวคฺค โน อปสฺสโต.\csromanspacer{} กิญฺจ ภิกฺขเว ชานโต กิํ ปสฺสโต อาสวานํ ขโย โหติ, ‘อิติ รูปํ ฯเปฯ อิติ วิญฺญาณสฺส อตฺถงฺคโม’ติ.\csromanspacer{} เอวํ โข ภิกฺขเว ชานโต เอวํ ปสฺสโต อาสวานํ ขโย โหตี”ติ, อตฺเถว สุตฺตนฺโตติ, อามนฺตา.\csromanspacer{} เตน หิ น วตฺตพฺพํ “อตฺถิ อรหโต กงฺขา”ติ.}
\prose{\csromanfolio{143}\csromansymbolmark{*}อตฺถิ อรหโต กงฺขาติ, อามนฺตา.\csromanspacer{} นนุ วุตฺตํ ภควตา “ชานโตหํ ภิกฺขเว ปสฺสโต อาสวานํ ขยํ วทามิ, โน อชานโต โน อปสฺสโต.\csromanspacer{} กิญฺจ ภิกฺขเว ชานโต กิํ ปสฺสโต อาสวานํ ขโย โหติ, ‘อิทํ ทุกฺขนฺ’ติ ภิกฺขเว ฯเปฯ ‘อยํ ทุกฺขนิโรธคามินี ปฏิปทา’ติ ชานโต ปสฺสโต อาสวานํ ขโย โหติ.\csromanspacer{} เอวํ โข ภิกฺขเว ชานโต เอวํ ปสฺสโต อาสวานํ ขโย โหตี”ติ อตฺเถว สุตฺตนฺโตติ, อามนฺตา.\csromanspacer{} เตน หิ น วตฺตพฺพํ “อตฺถิ อรหโต กงฺขา”ติ.}

\prose{\csromansymbolmark{*}อตฺถิ อรหโต กงฺขาติ, อามนฺตา.\csromanspacer{} นนุ วุตฺตํ ภควตา “สพฺพํ ภิกฺขเว อนภิชานํ อปริชานํ อวิราชยํ อปฺปชหํ อภพฺโพ ทุกฺขกฺขยาย, สพฺพญฺจ โข ภิกฺขเว อภิชานํ ปริชานํ วิราชยํ ปชหํ ภพฺโพ ทุกฺขกฺขยายา”ติ อตฺเถว สุตฺตนฺโตติ, อามนฺตา.\csromanspacer{} เตน หิ น วตฺตพฺพํ “อตฺถิ อรหโต กงฺขา”ติ.}

\prose{\csromansymbolmark{*}อตฺถิ อรหโต กงฺขาติ, อามนฺตา.\csromanspacer{} นนุ วุตฺตํ ภควตา “สหาวสฺส ทสฺสนสมฺปทาย ฯเปฯ ฉจฺจาภิฐานานิ อภพฺพ กาตุนฺ”ติ อตฺเถว สุตฺตนฺโตติ, อามนฺตา.\csromanspacer{} เตน หิ น วตฺตพฺพํ “อตฺถิ อรหโต กงฺขา”ติ.}

\prose{\csromansymbolmark{*}อตฺถิ อรหโต กงฺขาติ, อามนฺตา.\csromanspacer{} นนุ วุตฺตํ ภควตา “ยสฺมิํ ภิกฺขเว สมเย อริยสาวกสฺส วิรชํ วีตมลํ ธมฺมจกฺขุํ อุทปาทิ ‘ยํ กิญฺจิ สมุทยธมฺมํ, สพฺพํ ตํ นิโรธธมฺมนฺ’ติ.\csromanspacer{} สห ทสฺสนุปฺปาทา ภิกฺขเว อริยสาวกสฺส ตีณิ สํโยชนานิ ปหียนฺติ สกฺกายทิฏฺฐิ วิจิกิจฺฉา สีลพฺพตปรามาโส”ติ อตฺเถว สุตฺตนฺโตติ, อามนฺตา.\csromanspacer{} เตน หิ น วตฺตพฺพํ “อตฺถิ อรหโต กงฺขา”ติ.}

\prose{\csromansymbolmark{*}อตฺถิ อรหโต กงฺขาติ, อามนฺตา.\csromanspacer{} นนุ วุตฺตํ ภควตา\csromandash{}}

\csromangathameasure{“ยทา หเว ปาตุภวนฺติ ธมฺมา, \\ อาตาปิโน ฌายโต พฺราหฺมณสฺส. \\ อถสฺส กงฺขา วปยนฺติ สพฺพา, \\ ยโต ปชานาติ สเหตุธมฺมนฺ”ติ. \\ “ยทา หเว ปาตุภวนฺติ ธมฺมา, \\ อาตาปิโน ฌายโต พฺราหฺมณสฺส. \\ อถสฺส กงฺขา วปยนฺติ สพฺพา, \\ ยโต ขยํ ปจฺจยานํ อเวที”ติ. \\ “ยทา หเว ปาตุภวนฺติ ธมฺมา, \\ อาตาปิโน ฌายโต พฺราหฺมณสฺส. \\ วิธูปยํ ติฏฺฐติ มารเสนํ, \\ สูริโยว โอภาสยมนฺตลิกฺขนฺ”ติ. \\ “ยา กาจิ กงฺขา อิธ วา หุรํ วา, \\ สกเวทิยา วา ปรเวทิยา วา. \\ ฌายิโน ตา ปชหนฺติ สพฺพา, \\ อาตาปิโน พฺรหฺมจริยํ จรนฺตา”ติ. \\ “เย กงฺขาสมติกฺกนฺตา, กงฺขาภูเตสุ ปาณิสุ. \\ อสํสยา วิสํยุตฺตา, เตสุ ทินฺนํ มหปฺผลนฺ”ติ. \\ “เอตาทิสี ธมฺมปกาสเนตฺถ, \\ น ตตฺถ กิํ กงฺขติ โกจิ สาวโก. \\ นิตฺถิณฺณ-โอฆํ วิจิกิจฺฉฉินฺนํ, \\ พุทฺธํ นมสฺสาม ชินํ ชนินฺทา”ติ.}
\csromangathagroup{\csromangathastanza{“ยทา หเว ปาตุภวนฺติ ธมฺมา, \\ อาตาปิโน ฌายโต พฺราหฺมณสฺส. \\ อถสฺส กงฺขา วปยนฺติ สพฺพา, \\ ยโต ปชานาติ สเหตุธมฺมนฺ”ติ.}\csromangathastanza{\csromanfolio{144}“ยทา หเว ปาตุภวนฺติ ธมฺมา, \\ อาตาปิโน ฌายโต พฺราหฺมณสฺส. \\ อถสฺส กงฺขา วปยนฺติ สพฺพา, \\ ยโต ขยํ ปจฺจยานํ อเวที”ติ.}\csromangathastanza{“ยทา หเว ปาตุภวนฺติ ธมฺมา, \\ อาตาปิโน ฌายโต พฺราหฺมณสฺส. \\ วิธูปยํ ติฏฺฐติ มารเสนํ, \\ สูริโยว\footnote{สุริโยว (สี, สฺยา, กํ, อิ)} โอภาสยมนฺตลิกฺขนฺ”ติ\footnote{วิ 3. 2 ปิฏฺเฐ; ขุ 1. 78-79 ปิฏฺเฐสุปิ.}.}\csromangathastanza{“ยา กาจิ กงฺขา อิธ วา หุรํ วา, \\ สกเวทิยา วา ปรเวทิยา วา. \\ ฌายิโน\footnote{เย ฌายิโน (อุทาเน)} ตา ปชหนฺติ สพฺพา, \\ อาตาปิโน พฺรหฺมจริยํ จรนฺตา”ติ\footnote{เย ฌายิโน (อุทาเน)}. \\ “เย กงฺขาสมติกฺกนฺตา, กงฺขาภูเตสุ ปาณิสุ. \\ อสํสยา วิสํยุตฺตา, เตสุ ทินฺนํ มหปฺผลนฺ”ติ.}\csromangathastanza{“เอตาทิสี ธมฺมปกาสเนตฺถ, \\ น ตตฺถ กิํ กงฺขติ โกจิ สาวโก. \\ นิตฺถิณฺณ-โอฆํ\footnote{นิติณฺณ-โอฆํ (สี, ก), นิตฺติณฺณ-โอฆํ (สฺยา, กํ)} วิจิกิจฺฉฉินฺนํ, \\ พุทฺธํ นมสฺสาม ชินํ ชนินฺทา”ติ\footnote{นิติณฺณ-โอฆํ (สี, ก), นิตฺติณฺณ-โอฆํ (สฺยา, กํ)}.}}

\prose{อตฺเถว สุตฺตนฺโตติ, อามนฺตา.\csromanspacer{} เตน หิ น วตฺตพฺพํ “อตฺถิ อรหโต กงฺขา”ติ.}

\prose{\csromansymbolmark{+}น วตฺตพฺพํ “อตฺถิ อรหโต กงฺขา”ติ, อามนฺตา.\csromanspacer{} นนุ อรหา อิตฺถิปุริสานํ นามโคตฺเต กงฺเขยฺย, มคฺคามคฺเค กงฺเขยฺย, ติณกฏฺฐวนปฺปตีนํ นาเม กงฺเขยฺยาติ, อามนฺตา.\csromanspacer{} หญฺจิ อรหา อิตฺถิปุริสานํ นามโคตฺเต กงฺเขยฺย, มคฺคามคฺเค กงฺเขยฺย, ติณกฏฺฐวนปฺปตีนํ นาเม กงฺเขยฺย, เตน วต เร วตฺตพฺเพ “อตฺถิ อรหโต กงฺขา”ติ.}

\chapterheadpageread{2. ทุติยวคฺค}
\csromanmatikaanchor{mtk.33}
\csromantocmarkref{subsection}{(13) 4. ปรวิตารณกถา}{mtk.33}
\bookmark[dest={mtk.33},level=2]{(13) 4. ปรวิตารณกถา}
\prose{\csromanfolio{145}\csromansymbolmark{*}อรหา อิตฺถิปุริสานํ นามโคตฺเต กงฺเขยฺย, มคฺคามคฺเค กงฺเขยฺย, ติณกฏฺฐวนปฺปตีนํ นาเม กงฺเขยฺยาติ อตฺถิ อรหโต กงฺขาติ, อามนฺตา.\csromanspacer{} อรหา โสตาปตฺติผเล วา สกทาคามิผเล วา อนาคามิผเล วา อรหตฺเต วา กงฺเขยฺยาติ, น เหวํ วตฺตพฺเพ ฯเปฯ.}

\nitthitam{กงฺขากถา นิฏฺฐิตา.}
\chapterhead{2. ทุติยวคฺค}
\vspace{\tipitakaparskip}
\csromancenter{\textbf{(13) 4.}\csromanspacer{}\textbf{ ปรวิตารณกถา}}
\proseitem{322}{\csromansymbolmark{*}อตฺถิ อรหโต ปรวิตารณาติ, อามนฺตา.\csromanspacer{} อรหา ปรเนยฺโย ปรปตฺติโย ปรปจฺจโย ปรปฏิพทฺธภู, น ชานาติ น ปสฺสติ สมฺมูโฬฺห อสมฺปชาโนติ, น เหวํ วตฺตพฺเพ ฯเปฯ.}

\prose{\csromansymbolmark{*}นนุ อรหา น ปรเนยฺโย น ปรปตฺติโย น ปรปจฺจโย น ปรปฏิพทฺธภู, ชานาติ ปสฺสติ อสมฺมูโฬฺห สมฺปชาโนติ, อามนฺตา.\csromanspacer{} หญฺจิ อรหา น ปรเนยฺโย น ปรปตฺติโย น ปรปจฺจโย น ปรปฏิพทฺธภู, ชานาติ ปสฺสติ อสมฺมูโฬฺห สมฺปชาโน, โน จ วต เร วตฺตพฺเพ “อตฺถิ อรหโต ปรวิตารณา”ติ.}

\prose{\csromansymbolmark{*}อตฺถิ ปุถุชฺชนสฺส ปรวิตารณา, โส จ ปรเนยฺโย ปรปตฺติโย ปรปจฺจโย ปรปฏิพทฺธภู, น ชานาติ น ปสฺสติ สมฺมูโฬฺห อสมฺปชาโนติ, อามนฺตา.\csromanspacer{} อตฺถิ อรหโต ปรวิตารณา, โส จ ปรเนยฺโย ปรปตฺติโย ปรปจฺจโย ปรปฏิพทฺธภู, น ชานาติ น ปสฺสติ สมฺมูโฬฺห อสมฺปชาโนติ, น เหวํ วตฺตพฺเพ ฯเปฯ.}

\prose{\csromansymbolmark{*}อตฺถิ อรหโต ปรวิตารณา, โส จ น ปรเนยฺโย น ปรปตฺติโย น ปรปจฺจโย น ปรปฏิพทฺธภู, ชานาติ ปสฺสติ อสมฺมูโฬฺห สมฺปชาโนติ, อามนฺตา.\csromanspacer{} อตฺถิ ปุถุชฺชนสฺส ปรวิตารณา, โส จ น ปรเนยฺโย น ปรปตฺติโย น ปรปจฺจโย น ปรปฏิพทฺธภู, ชานาติ ปสฺสติ อสมฺมูโฬฺห สมฺปชาโนติ, น เหวํ วตฺตพฺเพ ฯเปฯ.}

\prose{\csromansymbolmark{*}อตฺถิ อรหโต ปรวิตารณาติ, อามนฺตา.\csromanspacer{} อตฺถิ อรหโต สตฺถริ ปรวิตารณา, ธมฺเม ปรวิตารณา, สํเฆ ปรวิตารณา, \csromanfolio{146}สิกฺขาย ปรวิตารณา, ปุพฺพนฺเต ปรวิตารณา, อปรนฺเต ปรวิตารณามฺ ปุพฺพนฺตาปรนฺเต ปรวิตารณา, อิทปฺปจฺจยตาปฏิจฺจสมุปฺปนฺเนสุ ธมฺเมสุ ปรวิตารณาติ, น เหวํ วตฺตพฺเพ ฯเปฯ.}

\prose{\csromansymbolmark{*}นตฺถิ อรหโต สตฺถริ ปรวิตารณา, ธมฺเม ปรวิตารณา ฯเปฯ อิทปฺปจฺจยตาปฏิจฺจสมุปฺปนฺเนสุ ธมฺเมสุ ปรวิตารณาติ, อามนฺตา.\csromanspacer{} หญฺจิ นตฺถิ อรหโต สตฺถริ ปรวิตารณา, ธมฺเม ปรวิตารณา ฯเปฯ อิทปฺปจฺจยตาปฏิจฺจ สมุปฺปนฺเนสุ ธมฺเมสุ ปรวิตารณา, โน จ วต เร วตฺตพฺเพ “อตฺถิ อรหโต ปรวิตารณา”ติ.}

\prose{\csromansymbolmark{*}อตฺถิ ปุถุชฺชนสฺส ปรวิตารณา, อตฺถิ ตสฺส สตฺถริ ปรวิตารณา, ธมฺเม ปรวิตารณา ฯเปฯ อิทปฺปจฺจยตาปฏิจฺจสมุปฺปนฺเนสุ ธมฺเมสุ ปรวิตารณาติ, อามนฺตา.\csromanspacer{} อตฺถิ อรหโต ปรวิตารณา, อตฺถิ ตสฺส สตฺถริ ปรวิตารณา, ธมฺเม ปรวิตารณา ฯเปฯ อิทปฺปจฺจยตาปฏิจฺจสมุปฺปนฺเนสุ ธมฺเมสุ ปรวิตารณาติ, น เหวํ วตฺตพฺเพ ฯเปฯ. \csromansymbolmark{*}อตฺถิ อรหโต ปรวิตารณา, นตฺถิ ตสฺส สตฺถริ ปรวิตารณา ธมฺเม ปรวิตารณา ฯเปฯ อิทปฺปจฺจยตาปฏิจฺจสมุปฺปนฺเนสุ ธมฺเมสุ ปรวิตารณาติ, อามนฺตา.\csromanspacer{} อตฺถิ ปุถุชฺชนสฺส ปรวิตารณา, นตฺถิ ตสฺส สตฺถริ ปรวิตารณา, ธมฺเม ปรวิตารณา ฯเปฯ อิทปฺปจฺจยตาปฏิจฺจสมุปฺปนฺเนสุ ธมฺเมสุ ปรวิตารณาติ, น เหวํ วตฺตพฺเพ ฯเปฯ.}

\proseitem{323}{\csromansymbolmark{*}อตฺถิ อรหโต ปรวิตารณาติ, อามนฺตา.\csromanspacer{} นนุ อรหโต ราโค ปหีโน อุจฺฉินฺนมูโล ตาลาวตฺถุกโต อนภาวํกโต อายติํ อนุปฺปาทธมฺโมติ, อามนฺตา.\csromanspacer{} หญฺจิ อรหโต ราโค ปหีโน อุจฺฉินฺนมูโล ตาลาวตฺถุกโต อนภาวํกโต อายติํ อนุปฺปาทธมฺโม, โน จ วต เร วตฺตพฺเพ “อตฺถิ อรหโต ปรวิตารณา”ติ.}

\prose{\csromansymbolmark{*}อตฺถิ อรหโต ปรวิตารณาติ, อามนฺตา.\csromanspacer{} นนุ อรหโต โทโส ปหีโน ฯเปฯ โมโห ปหีโน ฯเปฯ อโนตฺตปฺปํ ปหีนํ อุจฺฉินฺนมูลํ ตาลาวตฺถุกตํ อนภาวํกตํ อายติํ อนุปฺปาทธมฺมํ ฯเปฯ ราคปฺปหานาย มคฺโค ภาวิโต ฯเปฯ โพชฺฌงฺคา ภาวิตา ฯเปฯ โทสปฺปหานาย ฯเปฯ อโนตฺตปฺปปหานาย มคฺโค ภาวิโต ฯเปฯ โพชฺฌงฺคา ภาวิตา ฯเปฯ.}

\vspace{\tipitakaparskip}
\csromancenter{ทุติยวคฺค นนุ อรหา วีตราโค วีตโทโส วีตโมโห ฯเปฯ สจฺฉิกาตพฺพํ สจฺฉิกตนฺติ, อามนฺตา.\csromanspacer{} หญฺจิ อรหา วีตราโค วีตโทโส วีตโมโห สจฺฉิกาตพฺพํ สจฺฉิกตํ, โน จ วต เร วตฺตพฺเพ “อตฺถิ อรหโต ปรวิตารณา”ติ.}
\proseitem{324}{\csromanfolio{147}\csromansymbolmark{*}อตฺถิ อรหโต ปรวิตารณาติ? สธมฺมกุสลสฺส อรหโต อตฺถิ ปรวิตารณา, ปรธมฺมกุสลสฺส อรหโต นตฺถิ ปรวิตารณาติ.\csromanspacer{} สธมฺมกุสลสฺส อรหโต อตฺถิ ปรวิตารณาติ, อามนฺตา.\csromanspacer{} ปรธมฺมกุสลสฺส อรหโต อตฺถิ ปรวิตารณาติ, น เหวํ วตฺตพฺเพ ฯเปฯ.}

\prose{\csromansymbolmark{*}ปรธมฺมกุสลสฺส อรหโต นตฺถิ ปรวิตารณาติ, อามนฺตา.\csromanspacer{} สธมฺมกุสลสฺส อรหโต นตฺถิ ปรวิตารณาติ, น เหวํ วตฺตพฺเพ ฯเปฯ.}

\prose{\csromansymbolmark{*}สธมฺมกุสลสฺส อรหโต ราโค ปหีโน, อตฺถิ ตสฺส ปรวิตารณาติ, อามนฺตา.\csromanspacer{} ปรธมฺมกุสลสฺส อรหโต ราโค ปหีโน, อตฺถิ ตสฺส ปรวิตารณาติ, น เหวํ วตฺตพฺเพ.}

\prose{\csromansymbolmark{*}สธมฺมกุสลสฺส อรหโต โทโส ปหีโน ฯเปฯ โมโห ปหีโน ฯเปฯ อโนตฺตปฺปํ ปหีนํ ฯเปฯ ราคปฺปหานาย มคฺโค ภาวิโต ฯเปฯ โพชฺฌงฺคา ภาวิตา ฯเปฯ โทสปฺปหานาย ฯเปฯ อโนตฺตปฺปปหานาย มคฺโค ภาวิโต ฯเปฯ โพชฺฌงฺคา ภาวิตา ฯเปฯ.\csromanspacer{} สธมฺมกุสโล อรหา วีตราโค วีตโทโส วีตโมโห ฯเปฯ สจฺฉิกาตพฺพํ สจฺฉิกตํ, อตฺถิ ตสฺส ปรวิตารณาติ, อามนฺตา.\csromanspacer{} ปรธมฺมกุสโล อรหา วีตราโค วีตโทโส วีตโมโห ฯเปฯ สจฺฉิกาตพฺพํ สจฺฉิกตํ, อตฺถิ ตสฺส ปรวิตารณาติ, น เหวํ วตฺตพฺเพ ฯเปฯ.}

\prose{\csromansymbolmark{*}ปรธมฺมกุสลสฺส อรหโต ราโค ปหีโน, นตฺถิ ตสฺส ปรวิตารณาติ, อามนฺตา.\csromanspacer{} สธมฺมกุสลสฺส อรหโต ราโค ปหีโน, นตฺถิ ตสฺส ปรวิตารณาติ, น เหวํ วตฺตพฺเพ.}

\prose{\csromansymbolmark{*}ปรธมฺมกุสลสฺส อรหโต โทโส ปหีโน.\csromanspacer{} โมโห ปหีโน ฯเปฯ อโนตฺตปฺปํ ปหีนํ ฯเปฯ ราคปฺปหานาย มคฺโค ภาวิโต ฯเปฯ โพชฺฌงฺคา ภาวิตา ฯเปฯ โทสปฺปหานาย ฯเปฯ อโนตฺตปฺปปหานาย มคฺโค ภาวิโต ฯเปฯ โพชฺฌงฺคา ภาวิตา ฯเปฯ.\csromanspacer{} ปรธมฺมกุสโล อรหา \csromanfolio{148}วิตราโค วีตโทโส วีตโมโห ฯเปฯ สจฺฉิกาตพฺพํ สจฺฉิกตํ, นตฺถิ ตสฺส ปรวิตารณาติ, อามนฺตา.\csromanspacer{} สธมฺมกุสโล อรหา วีตราโค วีตโทโส วีตโมโห ฯเปฯ สจฺฉิกาตพฺพํ สจฺฉิกตํ.\csromanspacer{} นตฺถิ ตสฺส ปรวิตารณาติ, น เหวํ วตฺตพฺเพ ฯเปฯ.}

\proseitem{325}{\csromansymbolmark{*}อตฺถิ อรหโต ปรวิตารณาติ, อามนฺตา.\csromanspacer{} นนุ วุตฺตํ ภควตา “ชานโตหํ ภิกฺขเว ปสฺสโต อาสวานํ ขยํ วทามิ, โน อชานโต โน อปสฺสโต.\csromanspacer{} กิญฺจิ ภิกฺขเว ชานโต กิํ ปสฺสโต อาสวานํ ขโย โหติ, ‘อิติ รูปํ ฯเปฯ อิติ วิญฺญาณสฺส อตฺถงฺคโม’ติ.\csromanspacer{} เอวํ โข ภิกฺขเว ชานโต เอวํ ปสฺสโต อาสวานํ ขโย โหตี”ติ อตฺเถว สุตฺตนฺโตติ, อามนฺตา.\csromanspacer{} เตน หิ น วตฺตพฺพํ “อตฺถิ อรหโต ปรวิตารณา”ติ.}

\prose{\csromansymbolmark{*}อตฺถิ อรหโต ปรวิตารณาติ, อามนฺตา.\csromanspacer{} นนุ วุตฺตํ ภควตา “ชานโตหํ ภิกฺขเว ปสฺสโต อาสวานํ ขยํ วทามิ, โน อชานโต โน อปสฺสโต, กิญฺจ ภิกฺขเว ชานโต กิํ ปสฺสโต อาสวนํ ขโย โหติ, ‘อิทํ ทุกฺขนฺ’ติ ภิกฺขเว ชานโต ปสฺสโต อาสวานํ ขโย โหติ, ‘อยํ ทุกฺขสมุทโย’ติ ฯเปฯ ‘อยํ ทุกฺขนิโรโธ’ติ ฯเปฯ ‘อยํ ทุกฺขนิโรธคามินี ปฏิปทา’ติ ชานโต ปสฺสโต อาสวานํ ขโย โหติ.\csromanspacer{} เอวํ โข ภิกฺขเว ชานโต เอวํ ปสฺสโต อาสวานํ ขโย โหตี”ติ อตฺเถว สุตฺตนฺโตติ, อามนฺตา.\csromanspacer{} เตน หิ น วตฺตพฺพํ “อตฺถิ อรหโต ปรวิตารณา”ติ.}

\prose{\csromansymbolmark{*}อตฺถิ อรหโต ปรวิตารณาติ, อามนฺตา.\csromanspacer{} นนุ วุตฺตํ ภควตา “สพฺพํ ภิกฺขเว อนภิชานํ อปริชานํ อวิราชยํ อปฺปชหํ อภพฺโพ ทุกฺขกฺขยาย.\csromanspacer{} สพฺพญฺจ โข ภิกฺขเว อภิชานํ ปริชานํ วิราชยํ ปชหํ ภพฺโพ ทุกฺขกฺขยายา”ติ อตฺเถว สุตฺตนฺโตติ, อามนฺตา.\csromanspacer{} เตน หิ น วตฺตพฺพํ “อตฺถิ อรหโต ปรวิตารณา”ติ.}

\prose{\csromansymbolmark{*}อตฺถิ อรหโต ปรวิตารณาติ, อามนฺตา.\csromanspacer{} นนุ วุตฺตํ ภควตา “สหาวสฺส ทสฺสนสมฺปทาย ฯเปฯ ฉจฺจาภิฐานานิ อภพฺพ กาตุนฺ”ติ อตฺเถว สุตฺตนฺโตติ, อามนฺตา.\csromanspacer{} เตน หิ น วตฺตพฺพํ “อตฺถิ อรหโต ปรวิตารณา”ติ.}

\chapterheadpageread{2. ทุติยวคฺค}
\prose{\csromanfolio{149}\csromansymbolmark{*}อตฺถิ อรหโต ปรวิตารณาติ, อามนฺตา.\csromanspacer{} นนุ วุตฺตํ ภควตา “ยสฺมิํ ภิกฺขเว สมเย อริยสาวกสฺส วิรชํ วีตมลํ ธมฺมจกฺขุํ อุทปาทิ ‘ยํ กิญฺจิ สมุทยธมฺมํ, สพฺพํ ตํ นิโรธธมฺมนฺ’ติ.\csromanspacer{} สห ทสฺสนุปฺปาทา ภิกฺขเว อริยสาวกสฺส ตีณิ สํโยชนานิ ปหียนฺติ สกฺกายทิฏฺฐิ วิจิกิจฺฉา สีลพฺพตปรามาโส”ติ อตฺเถว สุตฺตนฺโตติ, อามนฺตา.\csromanspacer{} เตน หิ น วตฺตพฺพํ “อตฺถิ อรหโต ปรวิตารณา”ติ.}

\prose{\csromansymbolmark{*}อตฺถิ อรหโต ปรวิตารณาติ, อามนฺตา.\csromanspacer{} นนุ วุตฺตํ ภควตา\csromandash{}}

\csromangathameasure{“นาหํ สหิสฺสามิ ปโมจนาย, \\ กถํกถิํ โธตก กญฺจิ โลเก. \\ ธมฺมญฺจ เสฏฺฐํ อภิชานมาโน, \\ เอวํ ตุวํ โอฆมิมํ ตเรสี”ติ.}
\csromangathagroup{\csromangathastanza{“นาหํ สหิสฺสามิ\footnote{คมิสฺสามิ (สี, สฺยา, ก), สมีหามิ (อิ)} ปโมจนาย, \\ กถํกถิํ โธตก กญฺจิ\footnote{คมิสฺสามิ (สี, สฺยา, ก), สมีหามิ (อิ)} โลเก. \\ ธมฺมญฺจ เสฏฺฐํ อภิชานมาโน, \\ เอวํ ตุวํ โอฆมิมํ ตเรสี”ติ\footnote{คมิสฺสามิ (สี, สฺยา, ก), สมีหามิ (อิ)}.}}

\prose{อตฺเถว สุตฺตนฺโตติ, อามนฺตา.\csromanspacer{} เตน หิ น วตฺตพฺพํ “อตฺถิ อรหโต ปรวิตารณา”ติ.}

\prose{\csromansymbolmark{+}น วตฺตพฺพํ “อตฺถิ อรหโต ปรวิตารณา”ติ, อามนฺตา.\csromanspacer{} นนุ อรหโต อิตฺถิปุริสานํ นามโคตฺตํ ปเร วิตาเรยฺยุํ, มคฺคามคฺคํ ปเร วิตาเรยฺยุํ, ติณกฏฺฐวนปฺปตีนํ นามํ ปเร วิตาเรยฺยุนฺติ, อามนฺตา.\csromanspacer{} หญฺจิ อรหโต อิตฺถิปุริสานํ นามโคตฺตํ ปเร วิตาเรยฺยุํ, มคฺคามคฺคํ ปเร วิตาเรยฺยุํ, ติณกฏฺฐวนปฺปตีนํ นามํ ปเร วิตาเรยฺยุํ, เตน วต เร วตฺตพฺเพ “อตฺถิ อรหโต ปรวิตารณา”ติ.}

\prose{\csromansymbolmark{*}อรหโต อิตฺถิปุริสานํ นามโคตฺตํ ปเร วิตาเรยฺยุํ, มคฺคามคฺคํ ปเร วิตาเรยฺยุํ, ติณกฏฺฐวนปฺปตีนํ นามํ ปเร วิตาเรยฺยุนฺติ อตฺถิ อรหโต ปรวิตารณาติ, อามนฺตา.\csromanspacer{} อรหโต โสตาปตฺติผลํ วา สกทาคามิผลํ วา อนาคามิผลํ วา อรหตฺตํ วา ปเร วิตาเรยฺยุนฺติ, น เหวํ วตฺตพฺเพ ฯเปฯ.}

\nitthitam{ปรวิตารณกถา นิฏฺฐิตา.}
\chapterheadpageread{2. ทุติยวคฺค}
\vspace{\tipitakaparskip}
\csromancenter{\textbf{(14) 5.}\csromanspacer{}\textbf{ วจีเภทกถา}}
\csromanmatikaanchor{mtk.34}
\csromantocmarkref{subsection}{(14) 5. วจีเภทกถา}{mtk.34}
\bookmark[dest={mtk.34},level=2]{(14) 5. วจีเภทกถา}
\proseitem{326}{\csromanfolio{150}\csromansymbolmark{*}สมาปนฺนสฺส อตฺถิ วจีเภโทติ, อามนฺตา.\csromanspacer{} สพฺพตฺถ สมาปนฺนานํ อตฺถิ วจีเภโทติ, น เหวํ วตฺตพฺเพ ฯเปฯ.}

\prose{\csromansymbolmark{*}สมาปนฺนสฺส อตฺถิ วจีเภโทติ, อามนฺตา.\csromanspacer{} สพฺพทา สมาปนฺนานํ อตฺถิ วจีเภโทติ, น เหวํ วตฺตพฺเพ ฯเปฯ.}

\prose{\csromansymbolmark{*}สมาปนฺนสฺส อตฺถิ วจีเภโทติ, อามนฺตา.\csromanspacer{} สพฺเพสํ สมาปนฺนานํ อตฺถิ วจีเภโทติ, น เหวํ วตฺตพฺเพ ฯเปฯ.}

\prose{\csromansymbolmark{*}สมาปนฺนสฺส อตฺถิ วจีเภโทติ, อามนฺตา.\csromanspacer{} สพฺพสมาปตฺตีสุ อตฺถิ วจีเภโทติ, น เหวํ วตฺตพฺเพ ฯเปฯ.}

\prose{\csromansymbolmark{*}สมาปนฺนสฺส อตฺถิ วจีเภโทติ, อามนฺตา.\csromanspacer{} สมาปนฺนสฺส อตฺถิ กายเภโทติ, น เหวํ วตฺตพฺเพ ฯเปฯ.}

\prose{\csromansymbolmark{*}สมาปนฺนสฺส นตฺถิ กายเภโทติ, อามนฺตา.\csromanspacer{} สมาปนฺนสฺส นตฺถิ วจีเภโทติ, น เหวํ วตฺตพฺเพ ฯเปฯ.}

\prose{\csromansymbolmark{*}สมาปนฺนสฺส อตฺถิ วาจา, อตฺถิ วจีเภโทติ, อามนฺตา.\csromanspacer{} สมาปนฺนสฺส อตฺถิ กาโย, อตฺถิ กายเภโทติ, น เหวํ วตฺตพฺเพ ฯเปฯ.}

\prose{\csromansymbolmark{*}สมาปนฺนสฺส อตฺถิ กาโย, นตฺถิ กายเภโทติ, อามนฺตา.\csromanspacer{} สมาปนฺนสฺส อตฺถิ วาจา, นตฺถิ วจีเภโทติ, น เหวํ วตฺตพฺเพ ฯเปฯ.}

\proseitem{327}{\csromansymbolmark{*}“ทุกฺขนฺ”ติ ชานนฺโต “ทุกฺขนฺ”ติ วาจํ ภาสตีติ, อามนฺตา. “สมุทโย”ติ ชานนฺโต “สมุทโย”ติ วาจํ ภาสตีติ, น เหวํ วตฺตพฺเพ ฯเปฯ.}

\prose{\csromansymbolmark{*}“ทุกฺขนฺ”ติ ชานนฺโต “ทุกฺขนฺ”ติ วาจํ ภาสตีติ, อามนฺตา. “นิโรโธ”ติ ชานนฺโต “นิโรโธ”ติ วาจํ ภาสตีติ, น เหวํ วตฺตพฺเพ ฯเปฯ.}

\prose{\csromansymbolmark{*}“ทุกฺขนฺ”ติ ชานนฺโต “ทุกฺขนฺ”ติ วาจํ ภาสตีติ, อามนฺตา. “มคฺโค”ติ ชานนฺโต “มคฺโค”ติ วาจํ ภาสตีติ, น เหวํ วตฺตพฺเพ ฯเปฯ.}

\chapterheadpageread{2. ทุติยวคฺค}
\prose{\csromanfolio{151}\csromansymbolmark{*}“สมุทโย”ติ ชานนฺโต น จ “สมุทโย”ติ วาจํ ภาสตีติ, อามนฺตา. “ทุกฺขนฺ”ติ ชานนฺโต น จ “ทุกฺขนฺ”ติ วาจํ ภาสตีติ, น เหวํ วตฺตพฺเพ ฯเปฯ.}

\prose{\csromansymbolmark{*}“นิโรโธ”ติ ชานนฺโต น จ “นิโรโธ”ติ วาจํ ภาสตีติ, อามนฺตา. “ทุกฺขนฺ”ติ ชานนฺโต น จ “ทุกฺขนฺ”ติ วาจํ ภาสตีติ, น เหวํ วตฺตพฺเพ ฯเปฯ.}

\prose{\csromansymbolmark{*}“มคฺโค”ติ ชานนฺโต น จ “มคฺโค”ติ วาจํ ภาสตีติ, อามนฺตา. “ทุกฺขนฺ”ติ ชานนฺโต น จ “ทุกฺขนฺ”ติ วาจํ ภาสตีติ, น เหวํ วตฺตพฺเพ ฯเปฯ.}

\proseitem{328}{\csromansymbolmark{*}สมาปนฺนสฺส อตฺถิ วจีเภโทติ, อามนฺตา.\csromanspacer{} ญาณํ กิํโคจรนฺติ, ญาณํ สจฺจโคจรนฺติ.\csromanspacer{} โสตํ สจฺจโคจรนฺติ, น เหวํ วตฺตพฺเพ ฯเปฯ.}

\prose{\csromansymbolmark{*}สมาปนฺนสฺส อตฺถิ วจีเภโทติ, อามนฺตา, โสตํ กิํโคจรนฺติ.\csromanspacer{} โสตํ สทฺทโคจรนฺติ.\csromanspacer{} ญาณํ สทฺทโคจรนฺติ, น เหวํ วตฺตพฺเพ.}

\prose{\csromansymbolmark{*}สมาปนฺนสฺส อตฺถิ วจีเภโท ญาณํ สจฺจโคจรํ, โสตํ สทฺทโคจรนฺติ, อามนฺตา.\csromanspacer{} หญฺจิ ญาณํ สจฺจโคจรํ, โสตํ สทฺทโคจรํ, โน จ วต เร วตฺตพฺเพ “สมาปนฺนสฺส อตฺถิ วจีเภโท”ติ.}

\prose{\csromansymbolmark{*}สมาปนฺนสฺส อตฺถิ วจีเภโท ญาณํ สจฺจโคจรํ, โสตํ สทฺทโคจรนฺติ, อามนฺตา.\csromanspacer{} ทฺวินฺนํ ผสฺสานํ, ทฺวินฺนํ เวทนานํ, ทฺวินฺนํ สญฺญานํ, ทฺวินฺนํ เจตนานํ, ทฺวินฺนํ จิตฺตานํ สโมธานํ โหตีติ, น เหวํ วตฺตพฺเพ.}

\proseitem{329}{\csromansymbolmark{*}สมาปนฺนสฺส อตฺถิ วจีเภโทติ, อามนฺตา.\csromanspacer{} ปถวีกสิณํ สมาปตฺติํ สมาปนฺนสฺส อตฺถิ วจีเภโทติ, น เหวํ วตฺตพฺเพ.}

\prose{\csromansymbolmark{*}สมาปนฺนสฺส อตฺถิ วจีเภโทติ, อามนฺตา.\csromanspacer{} อาโปกสิณํ ฯเปฯ.\csromanspacer{} เตโชกสิณํ.\csromanspacer{} วาโยกสิณํ.\csromanspacer{} นีลกสิณํ.\csromanspacer{} ปีตกสิณํ.\csromanspacer{} โลหิตกสิณํ.\csromanspacer{} โอทาตกสิณํ.\csromanspacer{} อากาสานญฺจายตนํ.\csromanspacer{} วิญฺญาณญฺจายตนํ.\csromanspacer{} อากิญฺจญฺญายตนํ ฯเปฯ.\csromanspacer{} เนวสญฺญานาสญฺญายตนํ สมาปนฺนสฺส อตฺถิ วจีเภโทติ, น เหวํ วตฺตพฺเพ ฯเปฯ.}

\prose{\csromansymbolmark{*}ปถวีกสิณํ สมาปตฺติํ สมาปนฺนสฺส นตฺถิ วจีเภโทติ, อามนฺตา.\csromanspacer{} หญฺจิ ปถวีกสิณํ สมาปตฺติํ สมาปนฺนสฺส นตฺถิ วจีเภโท, โน จ วต เร วตฺตพฺเพ “สมาปนฺนสฺส อตฺถิ วจีเภโท”ติ.}

\prose{\csromanfolio{152}อาโปกสิณํ.\csromanspacer{} เตโชกสิณํ.\csromanspacer{} วาโยกสิณํ.\csromanspacer{} นีลกสิณํ.\csromanspacer{} ปีตกสิณํ.\csromanspacer{} โลหิตกสิณํ.\csromanspacer{} โอทาตกสิณํ.\csromanspacer{} อากาสานญฺจายตนํ.\csromanspacer{} วิญฺญาณญฺจายตนํ.\csromanspacer{} อากิญฺจญฺญายตนํ.\csromanspacer{} เนวสญฺญานาสญฺญายตนํ สมาปนฺนสฺส นตฺถิ วจีเภโทติ, อามนฺตา.\csromanspacer{} หญฺจิ เนวสญฺญานาสญฺญายตนํ สมาปนฺนสฺส นตฺถิ วจีเภโท, โน จ วต เร วตฺตพฺเพ “สมาปนฺนสฺส อตฺถิ วจีเภโท”ติ.}

\prose{\csromansymbolmark{*}สมาปนฺนสฺส อตฺถิ วจีเภโทติ, อามนฺตา.\csromanspacer{} โลกิยสมาปตฺติํ สมาปนฺนสฺส อตฺถิ วจีเภโทติ, น เหวํ วตฺตพฺเพ ฯเปฯ.}

\prose{\csromansymbolmark{*}สมาปนฺนสฺส อตฺถิ วจีเภโทติ, อามนฺตา.\csromanspacer{} โลกิยํ ปฐมํ ฌานํ สมาปนฺนสฺส อตฺถิ วจีเภโทติ, น เหวํ วตฺตพฺเพ ฯเปฯ.\csromanspacer{} สมาปนฺนสฺส อตฺถิ วจีเภโทติ, อามนฺตา.\csromanspacer{} โลกิยํ ทุติยํ ฌานํ.\csromanspacer{} ตติยํ ฌานํ.\csromanspacer{} จตุตฺถํ ฌานํ สมาปนฺนสฺส อตฺถิ วจีเภโทติ, น เหวํ วตฺตพฺเพ ฯเปฯ.}

\prose{\csromansymbolmark{*}โลกิยํ สมาปตฺติํ สมาปนฺนสฺส นตฺถิ วจีเภโทติ, อามนฺตา.\csromanspacer{} หญฺจิ โลกิยํ สมาปตฺติํ สมาปนฺนสฺส นตฺถิ วจีเภโท, โน จ วต เร วตฺตพฺเพ “สมาปนฺนสฺส อตฺถิ วจีเภโท”ติ.}

\prose{\csromansymbolmark{*}โลกิยํ ปฐมํ ฌานํ สมาปนฺนสฺส นตฺถิ วจีเภโทติ, อามนฺตา.\csromanspacer{} หญฺจิ โลกิยํ ปฐมํ ฌานํ สมาปนฺนสฺส นตฺถิ วจีเภโท, โน จ วต เร วตฺตพฺเพ “สมาปนฺนสฺส อตฺถิ วจีเภโท”ติ. \csromansymbolmark{*}โลกิยํ ทุติยํ ฌานํ.\csromanspacer{} ตติยํ ฌานํ.\csromanspacer{} จตุตฺถํ ฌานํ สมาปนฺนสฺส นตฺถิ วจีเภโทติ, อามนฺตา.\csromanspacer{} หญฺจิ โลกิยํ จตุตฺถํ ฌานํ สมาปนฺนสฺส นตฺถิ วจีเภโท, โน จ วต เร วตฺตพฺเพ “สมาปนฺนสฺส อตฺถิ วจีเภโท”ติ.}

\proseitem{330}{\csromansymbolmark{*}โลกุตฺตรํ ปฐมํ ฌานํ สมาปนฺนสฺส อตฺถิ วจีเภโทติ, อามนฺตา.\csromanspacer{} โลกิยํ ปฐมํ ฌานํ สมาปนฺนสฺส อตฺถิ วจีเภโทติ, น เหวํ วตฺตพฺเพ ฯเปฯ.}

\prose{\csromansymbolmark{*}โลกุตฺตรํ ปฐมํ ฌานํ สมาปนฺนสฺส อตฺถิ วจีเภโทติ, อามนฺตา.\csromanspacer{} โลกิยํ ทุติยํ ฌานํ.\csromanspacer{} ตติยํ ฌานํ.\csromanspacer{} จตุตฺถํ ฌานํ สมาปนฺนสฺส อตฺถิ วจีเภโทติ, น เหวํ วตฺตพฺเพ ฯเปฯ.}

\chapterheadpageread{2. ทุติยวคฺค}
\prose{\csromanfolio{153}\csromansymbolmark{*}โลกิยํ ปฐมํ ฌานํ สมาปนฺนสฺส นตฺถิ วจีเภโทติ, อามนฺตา.\csromanspacer{} โลกุตฺตรํ ปฐมํ ฌานํ สมาปนฺนสฺส นตฺถิ วจีเภโทติ, น เหวํ วตฺตพฺเพ ฯเปฯ.}

\prose{\csromansymbolmark{*}โลกิยํ ทุติยํ ฌานํ.\csromanspacer{} ตติยํ ฌานํ.\csromanspacer{} จตุตฺถํ ฌานํ สมาปนฺนสฺส นตฺถิ วจีเภโทติ, อามนฺตา.\csromanspacer{} โลกุตฺตรํ ปฐมํ ฌานํ สมาปนฺนสฺส นตฺถิ วจีเภโทติ, น เหวํ วตฺตพฺเพ ฯเปฯ.}

\proseitem{331}{\csromansymbolmark{*}โลกุตฺตรํ ปฐมํ ฌานํ สมาปนฺนสฺส อตฺถิ วจีเภโทติ, อามนฺตา.\csromanspacer{} โลกุตฺตรํ ทุติยํ ฌานํ สมาปนฺนสฺส อตฺถิ วจีเภโทติ, น เหวํ วตฺตพฺเพ ฯเปฯ.}

\prose{\csromansymbolmark{*}โลกุตฺตรํ ปฐมํ ฌานํ สมาปนฺนสฺส อตฺถิ วจีเภโทติ, อามนฺตา.\csromanspacer{} โลกุตฺตรํ ตติยํ ฌานํ.\csromanspacer{} จตุตฺถํ ฌานํ สมาปนฺนสฺส อตฺถิ วจีเภโทติ, น เหวํ วตฺตพฺเพ ฯเปฯ.}

\prose{\csromansymbolmark{*}โลกุตฺตรํ ทุติยํ ฌานํ สมาปนฺนสฺส นตฺถิ วจีเภโทติ, อามนฺตา.\csromanspacer{} โลกุตฺตรํ ปฐมํ ฌานํ สมาปนฺนสฺส นตฺถิ วจีเภโทติ, น เหวํ วตฺตพฺเพ ฯเปฯ.}

\prose{\csromansymbolmark{*}โลกุตฺตรํ ทุติยํ ฌานํ.\csromanspacer{} ตติยํ ฌานํ.\csromanspacer{} จตุตฺถํ ฌานํ สมาปนฺนสฺส นตฺถิ วจีเภโทติ, อามนฺตา.\csromanspacer{} โลกุตฺตรํ ปฐมํ ฌานํ สมาปนฺนสฺส นตฺถิ วจีเภโทติ, น เหวํ วตฺตพฺเพ ฯเปฯ.}

\proseitem{332}{\csromansymbolmark{+}น วตฺตพฺพํ “สมาปนฺนสฺส อตฺถิ วจีเภโท”ติ, อามนฺตา.\csromanspacer{} นนุ วิตกฺกวิจารา วจีสงฺขารา วุตฺตา ภควตา, ปฐมํ ฌานํ สมาปนฺนสฺส อตฺถิ วิตกฺกวิจาราติ, อามนฺตา.\csromanspacer{} หญฺจิ วิตกฺกวิจารา วจีสงฺขารา วุตฺตา ภควตา, ปฐมํ ฌานํ สมาปนฺนสฺส อตฺถิ วิตกฺกวิจารา, เตน วต เร วตฺตพฺเพ “สมาปนฺนสฺส อตฺถิ วจีเภโท”ติ.}

\prose{\csromansymbolmark{*}วิตกฺกวิจารา วจีสงฺขารา วุตฺตา ภควตา, ปฐมํ ฌานํ สมาปนฺนสฺส อตฺถิ วิตกฺกวิจาราติ\footnote{อตฺถิ วิตกฺกวิจารา (สฺยา) เอวมุปริปิ.} อตฺถิ ตสฺส วจีเภโทติ, อามนฺตา.\csromanspacer{} ปถวีกสิณํ ปฐมํ ฌานํ สมาปนฺนสฺส อตฺถิ วิตกฺกวิจารา, อตฺถิ ตสฺส วจีเภโทติ, น เหวํ วตฺตพฺเพ ฯเปฯ.}

\prose{\csromanfolio{154}\csromansymbolmark{*}วิตกฺกวิจารา วจีสงฺขารา วุตฺตา ภควตา, ปฐมํ ฌานํ สมาปนฺนสฺส อตฺถิ วิตกฺกวิจาราติ อตฺถิ ตสฺส วจีเภโทติ, อามนฺตา.\csromanspacer{} อาโปกสิณํ เตโชกสิณํ.\csromanspacer{} วาโยกสิณํ.\csromanspacer{} นีลกสิณํ.\csromanspacer{} ปีตกสิณํ.\csromanspacer{} โลหิตกสิณํ.\csromanspacer{} โอทาตกสิณํ ปฐมํ ฌานํ สมาปนฺนสฺส อตฺถิ วิตกฺกวิจารา, อตฺถิ ตสฺส วจีเภโทติ, น เหวํ วตฺตพฺเพ ฯเปฯ.}

\prose{\csromansymbolmark{+}น วตฺตพฺพํ “สมาปนฺนสฺส อตฺถิ วจีเภโท”ติ, อามนฺตา.\csromanspacer{} นนุ วิตกฺกสมุฏฺฐานา วาจา วุตฺตา ภควตา, ปฐมํ ฌานํ สมาปนฺนสฺส อตฺถิ วิตกฺกวิจาราติ, อามนฺตา.\csromanspacer{} หญฺจิ วิตกฺกสมุฏฺฐานา วาจา วุตฺตา ภควตา, ปฐมํ ฌานํ สมาปนฺนสฺส อตฺถิ วิตกฺกวิจารา, เตน วเต เร วตฺตพฺเพ “สมาปนฺนสฺส อตฺถิ วจีเภโท”ติ.}

\prose{\csromansymbolmark{*}วิตกฺกสมุฏฺฐานา วาจา วุตฺตา ภควตา, ปฐมํ ฌานํ สมาปนฺนสฺส อตฺถิ วิตกฺกวิจาราติ อตฺถิ ตสฺส วจีเภโทติ, อามนฺตา.\csromanspacer{} สญฺญาสมุฏฺฐานา วาจา วุตฺตา ภควตา, ทุติยํ ฌานํ สมาปนฺนสฺส อตฺถิ สญฺญา, อตฺถิ ตสฺส วิตกฺกวิจาราติ, น เหวํ วตฺตพฺเพ ฯเปฯ.}

\prose{\csromansymbolmark{*}วิตกฺกสมุฏฺฐานา วาจา วุตฺตา ภควตา, ปฐมํ ฌานํ สมาปนฺนสฺส อตฺถิ วิตกฺกวิจาราติ อตฺถิ ตสฺส วจีเภโทติ, อามนฺตา.\csromanspacer{} สญฺญาสมุฏฺฐานา วาจา วุตฺตา ภควตา, ตติยํ ฌานํ ฯเปฯ.\csromanspacer{} จตุตฺถํ ฌานํ.\csromanspacer{} อากาสานญฺจายตนํ.\csromanspacer{} วิญฺญาณญฺจายตนํ.\csromanspacer{} อากิญฺจญฺญายตนํ สมาปนฺนสฺส อตฺถิ สญฺญา, อตฺถิ ตสฺส วิตกฺกวิจาราติ, น เหวํ วตฺตพฺเพ ฯเปฯ.}

\proseitem{333}{\csromansymbolmark{*}สมาปนฺนสฺส อตฺถิ วจีเภโทติ, อามนฺตา.\csromanspacer{} นนุ “ปฐมํ ฌานํ สมาปนฺนสฺส วาจา นิรุทฺธา โหตี”ติ\footnote{สํ 2. 418, 420 ปิฏฺเฐสุ.} อตฺเถว สุตฺตนฺโตติ, อามนฺตา.\csromanspacer{} หญฺจิ “ปฐมํ ฌานํ สมาปนฺนสฺสวาจา นิรุทฺธา โหตี”ติ อตฺเถว สุตฺตนฺโตติ, โน จ วต เร วตฺตพฺเพ “สมาปนฺนสฺส อตฺถิ วจีเภโท”ติ.}

\prose{\csromansymbolmark{*}“ปฐมํ ฌานํ สมาปนฺนสฺส วาจา นิรุทฺธา โหตี”ติ1 อตฺเถว สุตฺตนฺโต”ติ อตฺถิ ตสฺส วจีเภโทติ, อามนฺตา. “ทุติยํ ฌานํ สมาปนฺนสฺส วิตกฺกวิจารา นิรุทฺธา โหนฺตีติ1 อตฺเถว สุตฺตนฺโต”ติ อตฺถิ ตสฺส วิตกฺกวิจาราติ, น เหวํ วตฺตพฺเพ ฯเปฯ.}

\prose{\csromansymbolmark{*}“ปฐมํ ฌานํ สมาปนฺนสฺส วาจา นิรุทฺธา โหตีติ1 อตฺเถว สุตฺตนฺโต”ติ อตฺถิ ตสฺส วจีเภโทติ, อามนฺตา. “ตติยํ ฌานํ}

\vspace{\tipitakaparskip}
\csromancenter{ทุติยวคฺค สมาปนฺนสฺส ปีติ นิรุทฺธา โหติ.\csromanspacer{} จตุตฺถํ ฌานํ สมาปนฺนสฺส อสฺสาสปสฺสาสา นิรุทฺธา โหนฺติ.\csromanspacer{} อากาสานญฺจายตนํ สมาปนฺนสฺส รูปสญฺญา นิรุทฺธา โหติ.\csromanspacer{} วิญฺญาณญฺจายตนํ สมาปนฺนสฺส อากาสานญฺจายตนสญฺญา นิรุทฺธา โหติ.\csromanspacer{} อากิญฺจญฺญายตนํ สมาปนฺนสฺส วิญฺญาณญฺจายตนสญฺญา นิรุทฺธา โหติ.\csromanspacer{} เนวสญฺญานาสญฺญายตนํ สมาปนฺนสฺส อากิญฺจญฺญายตนสญฺญา นิรุทฺธา โหติ.\csromanspacer{} สญฺญาเวทยิตนิโรธํ สมาปนฺนสฺส สญฺญา จ เวทนา จ นิรุทฺธา โหนฺตีติ\footnote{สํ 2. 418, 420; อํ 3. 208; ที 3. 221, 255 ปิฏฺเฐสุ.} อตฺเถว สุตฺตนฺโต”ติ อตฺถิ ตสฺส สญฺญา จ เวทนา จาติ, น เหวํ วตฺตพฺเพ ฯเปฯ.}
\prose{\csromanfolio{155}\csromansymbolmark{+}น วตฺตพฺพํ “สมาปนฺนสฺส อตฺถิ วจีเภโท”ติ, อามนฺตา.\csromanspacer{} นนุ ปฐมสฺส ฌานสฺส สทฺโท กณฺฑโก\footnote{กณฺฏโก (สฺยา) อํ 3. 363 ปิฏฺเฐปิ.} วุตฺโต ภควตาติ\footnote{อํ 3. 363 ปิฏฺเฐ กณฺฏกสุตฺตํ นิสฺสาย ปุจฺฉติ.}, อามนฺตา.\csromanspacer{} หญฺจิ ปฐมสฺส ฌานสฺส สทฺโท กณฺฑโก วุตฺโต ภควตา, เตน วต เร วตฺตพฺเพ “สมาปนฺนสฺส อตฺถิ วจีเภโท”ติ.}

\prose{\csromansymbolmark{*}ปฐมสฺส ฌานสฺส สทฺโท กณฺฑโก วุตฺโต ภควตาติ สมาปนฺนสฺส อตฺถิ วจีเภโทติ, อามนฺตา.\csromanspacer{} ทุติยสฺส ฌานสฺส วิตกฺกวิจารา กณฺฑโก วุตฺโต ภควตา.\csromanspacer{} ตติยสฺส ฌานสฺส ปีติ กณฺฑโก วุตฺโต ภควตา.\csromanspacer{} จตุตฺถสฺส ฌานสฺส อสฺสาสปสฺสาสา กณฺฑโก วุตฺโต ภควตา\footnote{อํ 3. 363 ปิฏฺเฐ.}.\csromanspacer{} อากาสานญฺจายตนํ สมาปนฺนสฺส รูปสญฺญา กณฺฑโก วุตฺโต ภควตา.\csromanspacer{} วิญฺญาณญฺจายตนํ สมาปนฺนสฺส อากาสานญฺจายตนสญฺญา กณฺฑโก วุตฺโต ภควตา.\csromanspacer{} อากิญฺจญฺญายตนํ สมาปนฺนสฺส วิญฺญาณญฺจายตนสญฺญา กณฺฑโก วุตฺโต ภควตา.\csromanspacer{} เนวสญฺญานาสญฺญายตนํ สมาปนฺนสฺส อากิญฺจญฺญายตนสญฺญา กณฺฑโก วุตฺโต ภควตา.\csromanspacer{} สญฺญาเวทยิตนิโรธํ สมาปนฺนสฺส สญฺญา จ เวทนา จ กณฺฑโก วุตฺโต ภควตา4, อตฺถิ ตสฺส สญฺญา จ เวทนา จาติ, น เหวํ วตฺตพฺเพ ฯเปฯ.}

\prose{\csromansymbolmark{+}น วตฺตพฺพํ “สมาปนฺนสฺส อตฺถิ วจีเภโท”ติ, อามนฺตา.\csromanspacer{} นนุ วุตฺตํ ภควตา “สิขิสฺส อานนฺท ภควโต อรหโต สมฺมาสมฺพุทฺธสฺส อภิภู นาม สาวโก พฺรหฺมโลเก ฐิโต ทสสหสฺสิโลกธาตุํ สเรน วิญฺญาเปสิ\csromandash{}}

\csromanmatikaanchor{mtk.35}
\csromanmatikaanchor{mtk.36}
\csromantocmarkref{subsection}{(15) 6. ทุกฺขาหารกถา}{mtk.35}
\bookmark[dest={mtk.35},level=2]{(15) 6. ทุกฺขาหารกถา}
\csromantocmarkref{subsection}{(16) 7. จิตฺติฏฺฐิติกถา}{mtk.36}
\bookmark[dest={mtk.36},level=2]{(16) 7. จิตฺติฏฺฐิติกถา}
\csromangathameasure{อารพฺภถ นิกฺกมถ, ยุญฺชถ พุทฺธสาสเน. \\ ธุนาถ มจฺจุโน เสนํ, นฬาคารํว กุญฺชโร. \\ โย อิมสฺมิํ ธมฺมวินเย, อปฺปมตฺโต วิหสฺสติ. \\ ปหาย ชาติสํสารํ, ทุกฺขสฺสนฺตํ กริสฺสตี”ติ.}
\csromangathagroup{\csromangathastanza{\csromanfolio{156}อารพฺภถ นิกฺกมถ, ยุญฺชถ พุทฺธสาสเน. \\ ธุนาถ มจฺจุโน เสนํ, นฬาคารํว กุญฺชโร.}\csromangathastanza{โย อิมสฺมิํ ธมฺมวินเย, อปฺปมตฺโต วิหสฺสติ\footnote{วิเหสฺสติ (สี, สฺยา, กํ), วิหริสฺสติ (ก)}. \\ ปหาย ชาติสํสารํ, ทุกฺขสฺสนฺตํ กริสฺสตี”ติ\footnote{วิเหสฺสติ (สี, สฺยา, กํ), วิหริสฺสติ (ก)}.}}

\prose{อตฺเถว สุตฺตนฺโตติ, อามนฺตา.\csromanspacer{} เตน หิ สมาปนฺนสฺส อตฺถิ วจีเภโทติ.}

\nitthitam{วจีเภทกถา นิฏฺฐิตา.}
\chapterhead{2. ทุติยวคฺค}
\vspace{\tipitakaparskip}
\csromancenter{\textbf{(15) 6.}\csromanspacer{}\textbf{ ทุกฺขาหารกถา}}
\proseitem{334}{\csromansymbolmark{*}ทุกฺขาหาโร มคฺคงฺคํ มคฺคปริยาปนฺนนฺติ, อามนฺตา.\csromanspacer{} เย เกจิ “ทุกฺขนฺ”ติ วาจํ ภาสนฺติ, สพฺเพ เต มคฺคํ ภาเวนฺตีติ, น เหวํ วตฺตพฺเพ.}

\prose{\csromansymbolmark{*}เย เกจิ “ทุกฺขนฺ”ติ วาจํ ภาสนฺติ, สพฺเพ เต มคฺคํ ภาเวนฺตีติ, อามนฺตา.\csromanspacer{} พาลปุถุชฺชนา “ทุกฺขนฺ”ติ วาจํ ภาสนฺติ, พาลปุถุชฺชนา มคฺคํ ภาเวนฺตีติ, น เหวํ วตฺตพฺเพ.\csromanspacer{} มาตุฆาตโก.\csromanspacer{} ปิตุฆาตโก.\csromanspacer{} อรหนฺตฆาตโก.\csromanspacer{} รุหิรุปฺปาทโก\footnote{โลหิตุปฺปาทโก (สี, ก) อญฺญฏฺฐาเนสุ ปน รุหิรุปฺปาทโกเตฺวว ทิสฺสติ.}.\csromanspacer{} สํฆเภทโก “ทุกฺขนฺ”ติ วาจํ ภาสติ, สํฆเภทโก มคฺคํ ภาเวตีติ, น เหวํ วตฺตพฺเพ ฯเปฯ.}

\nitthitam{ทุกฺขาหารกถา นิฏฺฐิตา.}
\chapterhead{2. ทุติยวคฺค}
\vspace{\tipitakaparskip}
\csromancenter{\textbf{(16) 7.}\csromanspacer{}\textbf{ จิตฺตฏฺฐิติกถา}}
\proseitem{335}{\csromansymbolmark{*}เอกํ จิตฺตํ ทิวสํ ติฏฺฐตีติ, อามนฺตา.\csromanspacer{} อุปฑฺฒทิวโส อุปฺปาทกฺขโณ, อุปฑฺฒทิวโส วยกฺขโณติ, น เหวํ วตฺตพฺเพ.}

\chapterheadpageread{2. ทุติยวคฺค}
\prose{\csromanfolio{157}\csromansymbolmark{*}เอกํ จิตฺตํ เทฺว ทิวเส ติฏฺฐตีติ, อามนฺตา.\csromanspacer{} ทิวโส อุปฺปาทกฺขโณ, ทิวโส วยกฺขโณติ, น เหวํ วตฺตพฺเพ.}

\prose{\csromansymbolmark{*}เอกํ จิตฺตํ จตฺตาโร ทิวเส ติฏฺฐติ.\csromanspacer{} อฏฺฐ ทิวเส ติฏฺฐติ.\csromanspacer{} ทส ทิวเส ติฏฺฐติ.\csromanspacer{} วีสติ ทิวเส ติฏฺฐติ.\csromanspacer{} มาสํ ติฏฺฐติ.\csromanspacer{} เทฺว มาเส ติฏฺฐติ.\csromanspacer{} จตฺตาโร มาเส ติฏฺฐติ.\csromanspacer{} อฏฺฐ มาเส ติฏฺฐติ.\csromanspacer{} ทส มาเส ติฏฺฐติ.\csromanspacer{} สํวจฺฉรํ ติฏฺฐติ.\csromanspacer{} เทฺว วสฺสานิ ติฏฺฐติ.\csromanspacer{} จตฺตาริ วสฺสานิ ติฏฺฐติ.\csromanspacer{} อฏฺฐ วสฺสานิ ติฏฺฐติ.\csromanspacer{} ทส วสฺสานิ ติฏฺฐติ.\csromanspacer{} วีสติ วสฺสานิ ติฏฺฐติ.\csromanspacer{} ติํส วสฺสานิ ติฏฺฐติ.\csromanspacer{} จตฺตารีส วสฺสานิ ติฏฺฐติ.\csromanspacer{} ปญฺญาส วสฺสานิ ติฏฺฐิติ.\csromanspacer{} วสฺสสตํ ติฏฺฐติ.\csromanspacer{} เทฺว วสฺสสตานิ ติฏฺฐติ.\csromanspacer{} จตฺตาริ วสฺสสตานิ ติฏฺฐติ.\csromanspacer{} ปญฺจ วสฺสสตานิ ติฏฺฐติ.\csromanspacer{} วสฺสสหสฺสํ ติฏฺฐติ.\csromanspacer{} เทฺว วสฺสสหสฺสานิ ตีฏฺฐติ.\csromanspacer{} จตฺตาริ วสฺสสหสฺสานิ ติฏฺฐติ.\csromanspacer{} อฏฺฐ วสฺสสหสฺสานิ ติฏฺฐติ.\csromanspacer{} โสฬส วสฺสสหสฺสานิ ติฏฺฐติ.\csromanspacer{} กปฺปํ ติฏฺฐติ.\csromanspacer{} เทฺว กปฺเป ติฏฺฐติ.\csromanspacer{} จตฺตาโร กปฺเป ตีฏฺฐติ.\csromanspacer{} อฏฺฐ กปฺเป ติฏฺฐติ.\csromanspacer{} โสฬส กปฺเป ติฏฺฐติ.\csromanspacer{} พาตฺติํส กปฺเป ติฏฺฐติ.\csromanspacer{} จตุสฏฺฐิ กปฺเป ติฏฺฐติ.\csromanspacer{} ปญฺจ กปฺปสตานิ ติฏฺฐติ.\csromanspacer{} กปฺปสหสฺสานิ ติฏฺฐติ.\csromanspacer{} เทฺว กปฺปสหสฺสานิ ติฏฺฐติ.\csromanspacer{} จตฺตาริ กปฺปสหสฺสานิ ติฏฺฐติ.\csromanspacer{} อฏฺฐ กปฺปสหสฺสานิ ติฏฺฐติ.\csromanspacer{} โสฬส กปฺปสหสฺสานิ ติฏฺฐติ.\csromanspacer{} วีสติ กปฺปสหสฺสานิ ติฏฺฐติ.\csromanspacer{} จตฺตารีส กปฺปสหสฺสานิ ตีฏฺฐติ.\csromanspacer{} สฏฺฐิ กปฺปสหสฺสานิ ติฏฺฐติ.\csromanspacer{} จตุราสีติ กปฺปสหสฺสานิ ติฏฺฐตีติ, อามนฺตา.\csromanspacer{} เทฺวจตฺตารีส กปฺปสหสฺสานิ อุปฺปาทกฺขโณ, เทฺวจตฺตารีส กปฺปสหสฺสานิ วยกฺขโณติ, น เหวํ วตฺตพฺเพ.}

\prose{\csromansymbolmark{*}เอกํ จิตฺตํ ทิวสํ ติฏฺฐตีติ, อามนฺตา.\csromanspacer{} อตฺถญฺเญ ธมฺมา เอกาหํ พหุมฺปิ อุปฺปชฺชิตฺวา นิรุชฺฌนฺตีติ, อามนฺตา.\csromanspacer{} เต ธมฺมา จิตฺเตน ลหุปริวตฺตาติ, น เหวํ วตฺตพฺเพ.}

\prose{\csromansymbolmark{*}เต ธมฺมา จิตฺเตน ลหุปริวตฺตาติ, อามนฺตา.\csromanspacer{} นนุ วุตฺตํ ภควตา “นาหํ ภิกฺขเว อญฺญํ เอกธมฺมมฺปิ สมนุปสฺสามิ เอวํ ลหุปริวตฺตํ ยถยิทํ จิตฺตํ, ยาวญฺจิทํ ภิกฺขเว อุปมาปิ น สุกรา ยาว ลหุปริวตฺตํ จิตฺตนฺ”ติ\footnote{อํ 1. 9 ปิฏฺเฐ.} อตฺเถว สุตฺตนฺโตติ, อามนฺตา.\csromanspacer{} เตน หิ น วตฺตพฺพํ “เต ธมฺมา จิตฺเตน ลหุปริวตฺตา”ติ.}

\prose{\csromanfolio{158}\csromansymbolmark{*}เต ธมฺมา จิตฺเตน ลหุปริวตฺตาติ, อามนฺตา.\csromanspacer{} นนุ วุตฺตํ ภควตา “เสยฺยถาปิ ภิกฺขเว มกฺกโฏ อรญฺเญ ปวเน จรมาโน สาขํ คณฺหติ, ตํ มุญฺจิตฺวา อญฺญํ คณฺหติ, ตํ มุญฺจิตฺวา อญฺญํ คณฺหติ.\csromanspacer{} เอวเมว โข ภิกฺขเว ยมิทํ\footnote{ยทิทํ (พหูสุ)} วุจฺจติ จิตฺตํ อิติปิ มโน อิติปิ วิญฺญาณํ อิติปิ, ตํ รตฺติยา จ ทิวสสฺส จ อญฺญเทว อุปฺปชฺชติ อญฺญํ นิรุชฺฌตี”ติ\footnote{สํ 1. 320 ปิฏฺเฐ.} อตฺเถว สุตฺตนฺโตติ, อามนฺตา.\csromanspacer{} เตน หิ น วตฺตพฺพํ “เต ธมฺมา จิตฺเตน ลหุปริวตฺตา”ติ.}

\proseitem{336}{\csromansymbolmark{*}เอกํ จิตฺตํ ทิวสํ ติฏฺฐตีติ, อามนฺตา.\csromanspacer{} จกฺขุวิญฺญาณํ ทิวสํ ติฏฺฐตีติ, น เหวํ วตฺตพฺเพ.\csromanspacer{} โสตวิญฺญาณํ ฯเปฯ.\csromanspacer{} ฆานวิญฺญาณํ.\csromanspacer{} ชิวฺหาวิญฺญาณํ.\csromanspacer{} กายวิญฺญาณํ.\csromanspacer{} อกุสลํ จิตฺตํ.\csromanspacer{} ราคสหคตํ.\csromanspacer{} โทสสหคตํ.\csromanspacer{} โมหสหคตํ.\csromanspacer{} มานสหคตํ.\csromanspacer{} ทิฏฺฐิสหคตํ.\csromanspacer{} วิจิกิจฺฉาสหคตํ.\csromanspacer{} ถินสหคตํ.\csromanspacer{} อุทฺธจฺจสหคตํ.\csromanspacer{} อหิริกสหคตํ.\csromanspacer{} อโนตฺตปฺปสหคตํ จิตฺตํ ทิวสํ ติฏฺฐตีติ, น เหวํ วตฺตพฺเพ ฯเปฯ.}

\prose{\csromansymbolmark{*}เอกํ จิตฺตํ ทิวสํ ติฏฺฐตีติ, อามนฺตา.\csromanspacer{} เยเนว จิตฺเตน จกฺขุนา รูปํ ปสฺสติ, เตเนว จิตฺเตน โสเตน สทฺทํ สุณาติ ฯเปฯ ฆาเนน คนฺธํ ฆายติ.\csromanspacer{} ชิวฺหาย รสํ สายติ.\csromanspacer{} กาเยน โผฏฺฐพฺพํ ผุสติ.\csromanspacer{} มนสา ธมฺมํ วิชานาติ ฯเปฯ.\csromanspacer{} เยเนว จิตฺเตน มนสา ธมฺมํ วิชานาติ, เตเนว จิตฺเตน จกฺขุนา รูปํ ปสฺสติ ฯเปฯ โสเตน สทฺทํ สุณาติ.\csromanspacer{} ฆาเนน คนฺธํ ฆายติ.\csromanspacer{} ชิวฺหาย รสํ สายติ ฯเปฯ กาเยน โผฏฺฐพฺพํ ผุสตีติ, น เหวํ วตฺตพฺเพ ฯเปฯ.}

\prose{\csromansymbolmark{*}เอกํ จิตฺตํ ทิวสํ ติฏฺฐตีติ, อามนฺตา.\csromanspacer{} เยเนว จิตฺเตน อภิกฺกมติ, เตเนว จิตฺเตน ปฏิกฺกมติ.\csromanspacer{} เยเนว จิตฺเตน ปฏิกฺกมติ, เตเนว จิตฺเตน อภิกฺกมติ.\csromanspacer{} เยเนว จิตฺเตน อาโลเกติ, เตเนว จิตฺเตน วิโลเกติ.\csromanspacer{} เยเนว จิตฺเตน วิโลเกติ, เตเนว จิตฺเตน อาโลเกติ.\csromanspacer{} เยเนว จิตฺเตน สมิญฺเชติ, เตเนว จิตฺเตน ปสาเรติ.\csromanspacer{} เยเนว จิตฺเตน ปสาเรติ, เตเนว จิตฺเตน สมิญฺเชตีติ, น เหวํ วตฺตพฺเพ ฯเปฯ.}

\chapterheadpageread{2. ทุติยวคฺค}
\proseitem{337}{\csromanfolio{159}\csromansymbolmark{*}อากาสานญฺจายตนูปคานํ เทวานํ เอกํ จิตฺตํ ยาวตายุกํ ติฏฺฐตีติ, อามนฺตา.\csromanspacer{} มนุสฺสานํ เอกํ จิตฺตํ ยาวตายุกํ ติฏฺฐตีติ, น เหวํ วตฺตพฺเพ.}

\prose{\csromansymbolmark{*}อากาสานญฺจายตนูปคานํ เทวานํ เอกํ จิตฺตํ ยาวตายุกํ ติฏฺฐตีติ, อามนฺตา.\csromanspacer{} จาตุมหาราชิกานํ เทวานํ ฯเปฯ.\csromanspacer{} ตาวติํสานํ เทวานํ.\csromanspacer{} ยามานํ เทวานํ.\csromanspacer{} ตุสิตานํ เทวานํ.\csromanspacer{} นิมฺมานรตีนํ เทวานํ.\csromanspacer{} ปรนิมฺมิตวสวตฺตีนํ เทวานํ.\csromanspacer{} พฺรหฺมปาริสชฺชานํ เทวานํ.\csromanspacer{} พฺรหฺมปุโรหิตานํ เทวานํ.\csromanspacer{} มหาพฺรหฺมานํ เทวานํ.\csromanspacer{} ปริตฺตาภานํ เทวานํ.\csromanspacer{} อปฺปมาณาภานํ เทวานํ.\csromanspacer{} อาภสฺสรานํ เทวานํ.\csromanspacer{} ปริตฺตสุภานํ เทวานํ.\csromanspacer{} อปฺปมาณสุภานํ เทวานํ.\csromanspacer{} สุภกิณฺหานํ เทวานํ.\csromanspacer{} เวหปฺผลานํ เทวานํ.\csromanspacer{} อวิหานํ เทวานํ.\csromanspacer{} อตปฺปานํ เทวานํ.\csromanspacer{} สุทสฺสานํ เทวานํ.\csromanspacer{} สุทสฺสีนํ เทวานํ.\csromanspacer{} อกนิฏฺฐานํ เทวานํ เอกํ จิตฺตํ ยาวตายุกํ ติฏฺฐตีติ, น เหวํ วตฺตพฺเพ.}

\prose{\csromansymbolmark{*}อากาสานญฺจายตนูปคานํ เทวานํ วีสติ กปฺปสหสฺสานิ อายุปฺปมาณํ, อากาสานญฺจายตนูปคานํ เทวานํ เอกํ จิตฺตํ วีสติ กปฺปสหสฺสานิ ติฏฺฐตีติ, อามนฺตา.\csromanspacer{} มนุสฺสานํ วสฺสสตํ อายุปฺปมาณํ มนุสฺสานํ เอกํ จิตฺตํ วสฺสสตํ ติฏฺฐตีติ, น เหวํ วตฺตพฺเพ.}

\prose{\csromansymbolmark{*}อากาสานญฺจายตนูปคานํ เทวานํ วีสติ กปฺปสหสฺสานิ อายุปฺปมาณํ, อากาสานญฺจายตนูปคานํ เทวานํ เอกํ จิตฺตํ วีสติ กปฺปสหสฺสานิ ติฏฺฐตีติ, อามนฺตา.\csromanspacer{} จาตุมหาราชิกานํ เทวานํ ปญฺจ วสฺสสตานิ อายุปฺปมาณํ, จาตุมหาราชิกานํ เทวานํ เอกํ จิตฺตํ ปญฺจ วสฺสสตานิ ติฏฺฐติ.\csromanspacer{} วสฺสสหสฺสํ ติฏฺฐติ.\csromanspacer{} เทฺว วสฺสสหสฺสานิ ติฏฺฐติ.\csromanspacer{} จตฺตาริ วสฺสสหสฺสานิ ติฏฺฐติ.\csromanspacer{} อฏฺฐ วสฺสสหสฺสานิ ติฏฺฐติ.\csromanspacer{} โสฬส วสฺสสหสฺสานิ ติฏฺฐติ.\csromanspacer{} กปฺปสฺส ตติยภาคํ ติฏฺฐติ.\csromanspacer{} อุปฑฺฒกปฺปํ ติฏฺฐติ.\csromanspacer{} เอกํ กปฺปํ ติฏฺฐติ.\csromanspacer{} เทฺว กปฺเป ติฏฺฐติ.\csromanspacer{} จตฺตาโร กปฺเป ติฏฺฐติ.\csromanspacer{} อฏฺฐ กปฺเป ติฏฺฐติ.\csromanspacer{} โสฬส กปฺเป ติฏฺฐติ.\csromanspacer{} พาตฺติํส กปฺเป ติฏฺฐติ.\csromanspacer{} จตุสฏฺฐิ กปฺเป ติฏฺฐติ.\csromanspacer{} ปญฺจ กปฺปสตานิ ติฏฺฐติ.\csromanspacer{} กปฺปสหสฺสํ ติฏฺฐติ.\csromanspacer{} เทฺว กปฺปสหสฺสานิ ติฏฺฐติ.\csromanspacer{} จตฺตาริ กปฺปสหสฺสานิ ติฏฺฐติ.\csromanspacer{} อฏฺฐ กปฺปสหสฺสานิ ติฏฺฐติ.\csromanspacer{} อกนิฏฺฐานํ เทวานํ โสฬส กปฺปสหสฺสานิ อายุปฺปมาณํ, อกนิฏฺฐานํ เทวานํ เอกํ จิตฺตํ โสฬส กปฺปสหสฺสานิ ติฏฺฐตีติ, น เหวํ วตฺตพฺเพ.}

\csromanmatikaanchor{mtk.37}
\csromantocmarkref{subsection}{(17) 8. กุกฺกุฬกถา}{mtk.37}
\bookmark[dest={mtk.37},level=2]{(17) 8. กุกฺกุฬกถา}
\prose{\csromanfolio{160}\csromansymbolmark{+}อากาสานญฺจายตนูปคานํ เทวานํ จิตฺตํ มุหุตฺตํ มุหุตฺตํ อุปฺปชฺชติ, มุหุตฺตํ มุหุตฺตํ นิรุชฺฌตีติ, อามนฺตา.\csromanspacer{} อากาสานญฺจายตนูปคา เทวา มุหุตฺตํ มุหุตฺตํ จวนฺติ, มุหุตฺตํ มุหุตฺตํ อุปฺปชฺชนฺตีติ, น เหวํ วตฺตพฺเพ.}

\prose{\csromansymbolmark{*}อากาสานญฺจายตนูปคานํ เทวานํ เอกํ จิตฺตํ ยาวตายุกํ ติฏฺฐตีติ, อามนฺตา.\csromanspacer{} อากาสานญฺจายตนูปคา เทวา เยเนว จิตฺเตน อุปฺปชฺชนฺติ, เตเนว จิตฺเตน จวนฺตีติ, น เหวํ วตฺตพฺเพ ฯเปฯ.}

\nitthitam{จิตฺตฏฺฐิติกถา นิฏฺฐิตา.}
\chapterhead{2. ทุติยวคฺค}
\vspace{\tipitakaparskip}
\csromancenter{\textbf{(17) 8.}\csromanspacer{}\textbf{ กุกฺกุฬกถา}}
\proseitem{338}{\csromansymbolmark{*}สพฺเพ สงฺขารา อโนธิํ กตฺวา กุกฺกุฬาติ, อามนฺตา.\csromanspacer{} นนุ อตฺถิ สุขา เวทนา กายิกํ สุขํ เจตสิกํ สุขํ ทิพฺพํ สุขํ มานุสกํ สุขํ ลาภสุขํ สกฺการสุขํ ยานสุขํ สยนสุขํ อิสฺสริยสุขํ อาธิปจฺจสุขํ คิหิสุขํ สามญฺญสุขํ สาสวํ สุขํ อนาสวํ สุขํ อุปธิสุขํ นิรูปธิสุขํ สามิสํ สุขํ นิรามิสํ สุขํ สปฺปีติกํ สุขํ นิปฺปีติกํ สุขํ ฌานสุขํ วิมุตฺติสุขํ กามสุขํ เนกฺขมฺมสุขํ วิเวกสุขํ อุปสมสุขํ สมฺโพธสุขนฺติ, อามนฺตา.\csromanspacer{} หญฺจิ อตฺถิ สุขา เวทนา ฯเปฯ สมฺโพธสุขํ, โน จ วต เร วตฺตพฺเพ “สพฺเพ สงฺขารา อโนธิํ กตฺวา กุกฺกุฬา”ติ.}

\prose{\csromansymbolmark{*}สพฺเพ สงฺขารา อโนธิํ กตฺวา กุกฺกุฬาติ, อามนฺตา.\csromanspacer{} สพฺเพ สงฺขารา ทุกฺขา เวทนา กายิกํ ทุกฺขํ เจตสิกํ ทุกฺขํ โสกปริเทวทุกฺขโทมนสฺส-อุปายาสาติ, น เหวํ วตฺตพฺเพ.}

\prose{\csromansymbolmark{+}น วตฺตพฺพํ “สพฺเพ สงฺขารา อโนธิํ กตฺวา กุกฺกุฬา”ติ, อามนฺตา.\csromanspacer{} นนุ วุตฺตํ ภควตา “สพฺพํ ภิกฺขเว อาทิตฺตํ, กิญฺจ ภิกฺขเว สพฺพํ อาทิตฺตํ, จกฺขุํ ภิกฺขเว อาทิตฺตํ รูปา อาทิตฺตา จกฺขุวิญฺญาณํ อาทิตฺตํ จกฺขุสมฺผสฺโส อาทิตฺโต.\csromanspacer{} ยมิทํ\footnote{ยมฺปิทํ (สี, สฺยา, สํ 2. 251 ปิฏฺเฐปิ.)} จกฺขุสมฺผสฺสปจฺจยา อุปฺปชฺชติ เวทยิตํ สุขํ วา ทุกฺขํ วา อทุกฺขมสุขํ วา, ตมฺปิ อาทิตฺตํ.\csromanspacer{} เกน อาทิตฺตํ, ราคคฺคินา โทสคฺคินา โมหคฺคินา อาทิตฺตํ, ชาติยา ชราย มรเณน โสเกหิ ปริเทเวหิ}

\vspace{\tipitakaparskip}
\csromancenter{ทุติยวคฺค ทุกฺเขหิ โทมนสฺเสหิ อุปายาเสหิ อาทิตฺตนฺติ วทามิ.\csromanspacer{} โสตํ อาทิตฺตํ, สทฺทา อาทิตฺตา ฯเปฯ.\csromanspacer{} ฆานํ อาทิตฺตํ, คนฺธา อาทิตฺตา ฯเปฯ.\csromanspacer{} ชิวฺหา อาทิตฺตา, รสา อาทิตฺตา ฯเปฯ.\csromanspacer{} กาโย อาทิตฺโต, โผฏฺฐพฺพา อาทิตฺตา ฯเปฯ.\csromanspacer{} มโน อาทิตฺโต, ธมฺมา อาทิตฺตา.\csromanspacer{} มโนวิญฺญาณํ อาทิตฺตํ, มโนสมฺผสฺโส อาทิตฺโต.\csromanspacer{} ยมิทํ มโนสมฺผสฺสปจฺจยา อุปฺปชฺชติ เวทยิตํ สุขํ วา ทุกฺขํ วา อทุกฺขมสุขํ วา, ตมฺปิ อาทิตฺตํ.\csromanspacer{} เกน อาทิตฺตํ, ราคคฺคินา โทสคฺคินา โมหคฺคินา อาทิตฺตํ, ชาติยา ชราย มรเณน โสเกหิ ปริเทเวหิ ทุกฺเขหิ โทมนสฺเสหิ อุปายาเสหิ อาทิตฺตนฺติ วทามี”ติ\footnote{วิ 3. 44; สํ 2. 251 ปิฏฺเฐสุ.} อตฺเถว สุตฺตนฺโตติ, อามนฺตา.\csromanspacer{} เตน หิ วตฺตพฺพํ\footnote{เตน หิ (สี, สฺยา), เตน หิ น วตฺตพฺพํ (ก)} “สพฺเพ สงฺขารา อโนธิํ กตฺวา กุกฺกุฬา”ติ.}
\prose{\csromanfolio{161}\csromansymbolmark{*}สพฺเพ สงฺขารา อโนธิํ กตฺวา กุกฺกุฬาติ, อามนฺตา.\csromanspacer{} นนุ วุตฺตํ ภควตา “ปญฺจิเม ภิกฺขเว กามคุณา.\csromanspacer{} กตเม ปญฺจ, จกฺขุวิญฺเญยฺยา รูปา อิฏฺฐา กนฺตา มนาปา ปิยรูปา กามูปสํหิตา รชนียา, โสตวิญฺเญยฺยา สทฺทา ฯเปฯ ฆานวิญฺเญยฺยา คนฺธา.\csromanspacer{} ชิวฺหาวิญฺเญยฺยา รสา.\csromanspacer{} กายวิญฺเญยฺยา โผฏฺฐพฺพา อิฏฺฐา กนฺตา มนาปา ปิยรูปา กามูปสํหิตา รชนียา.\csromanspacer{} อิเม โข ภิกฺขเว ปญฺจ กามคุณา”ติ\footnote{ม 1. 119; สํ 2. 427 ปิฏฺเฐสุ.} อตฺเถว สุตฺตนฺโตติ, อามนฺตา.\csromanspacer{} เตน หิ น วตฺตพฺพํ “สพฺเพ สงฺขารา อโนธิํ กตฺวา กุกฺกุฬา”ติ.}

\prose{\csromansymbolmark{+}น วตฺตพฺพํ “สพฺเพ สงฺขารา อโนธิํ กตฺวา กุกฺกุฬา”ติ, อามนฺตา.\csromanspacer{} นนุ วุตฺตํ ภควตา “ลาภา โว ภิกฺขเว สุลทฺธํ โว ภิกฺขเว ขโณ โว ปฏิลทฺโธ\footnote{ปฏิวิทฺโธ (พหูสุ)} พฺรหฺมจริยวาสาย.\csromanspacer{} ทิฏฺฐา มยา ภิกฺขเว ฉ ผสฺสายตนิกา นาม นิรยา.\csromanspacer{} ตตฺถ ยํ กิญฺจิ จกฺขุนา รูปํ ปสฺสติ, อนิฏฺฐรูปญฺเญว ปสฺสติ โน อิฏฺฐรูปํ, อกนฺตรูปญฺเญว ปสฺสติ โน กนฺตรูปํ, อมนาปรูปญฺเญว ปสฺสติ โน มนาปรูปํ.\csromanspacer{} ยํ กิญฺจิ โสเตน สทฺทํ สุณาติ ฯเปฯ ฆาเนน คนฺธํ ฆายติ.\csromanspacer{} ชิวฺหาย รสํ สายติ.\csromanspacer{} กาเยน โผฏฺฐพฺพํ ผุสติ.\csromanspacer{} มนสา ธมฺมํ วิชานาติ.\csromanspacer{} อนิฏฺฐรูปญฺเญว วิชานาติ โน อิฏฺฐรูปํ, อกนฺตรูปญฺเญว วิชานาติ โน กนฺตรูปํ, อมนาปรูปญฺเญว วิชานาติ โน มนาปรูปนฺ”ติ\footnote{สํ 2. 341 ปิฏฺเฐ.} อตฺเถว สุตฺตนฺโตติ, อามนฺตา.\csromanspacer{} เตน หิ วตฺตพฺพํ\footnote{เตน หิ น วตฺตพฺพํ (สฺยา, ก)} “สพฺเพ สงฺขารา อโนธิํ กตฺวา กุกฺกุฬา”ติ.}

\prose{\csromanfolio{162}\csromansymbolmark{*}สพฺเพ สงฺขารา อโนธิํ กตฺวา กุกฺกุฬาติ, อามนฺตา.\csromanspacer{} นนุ วุตฺตํ ภควตา “ลาภา โว ภิกฺขเว สุลทฺธํ โว ภิกฺขเว ขโณ โว ปฏิลทฺโธ พฺรหฺมจริยวาสาย.\csromanspacer{} ทิฏฺฐา มยา ภิกฺขเว ฉ ผสฺสายตนิกา นาม สคฺคา.\csromanspacer{} ตตฺถ ยํ กิญฺจิ จกฺขุนา รูปํ ปสฺสติ, อิฏฺฐรูปญฺเญว ปสฺสติ โน อนิฏฺฐรูปํ, กนฺตรูปญฺเญว ปสฺสติ โน อกนฺตรูปํ, มนาปรูปญฺเญว ปสฺสติ โน อมนาปรูปํ.\csromanspacer{} ยํ กิญฺจิ โสเตน สทฺทํ สุณาติ ฯเปฯ ฆาเนน คนฺธํ ฆายติ.\csromanspacer{} ชิวฺหาย รสํ สายติ.\csromanspacer{} กาเยน โผฏฺฐพฺพํ ผุสติ.\csromanspacer{} มนสา ธมฺมํ วิชานาติ.\csromanspacer{} อิฏฺฐรูปญฺเญว วิชานาติ โน อนิฏฺฐรูปํ, กนฺตรูปญฺเญว วิชานาติ โน อกนฺตรูปํ, มนาปรูปญฺเญว วิชานาติ โน อมนาปรูปนฺ”ติ\footnote{สํ 2. 341 ปิฏฺเฐ.} อตฺเถว สุตฺตนฺโตติ, อามนฺตา.\csromanspacer{} เตน หิ น วตฺตพฺพํ “สพฺเพสงฺขารา อโนธิํ กตฺวา กุกฺกุฬา”ติ.}

\prose{\csromansymbolmark{+}น วตฺตพฺพํ “สพฺเพ สงฺขารา อโนธิํ กตฺวา กุกฺกุฬา”ติ, อามนฺตา.\csromanspacer{} นนุ ‘ยทนิจฺจํ ตํ ทุกฺขํ’\footnote{สํ 2. 19 ปิฏฺเฐ.} วุตฺตํ ภควตา, ‘สพฺเพ สงฺขารา อนิจฺจา’ติ\footnote{ขุ 1. 53; ม 1. 292 ปิฏฺเฐสุ.}, อามนฺตา.\csromanspacer{} หญฺจิ ‘ยทนิจฺจํ ตํ ทุกฺขํ’ วุตฺตํ ภควตา, ‘สพฺเพ สงฺขารา อนิจฺจา’, เตน วต เร วตฺตพฺเพ “สพฺเพ สงฺขารา อโนธิํ กตฺวา กุกฺกุฬา”ติ.}

\prose{\csromansymbolmark{*}สพฺเพ สงฺขารา อโนธิํ กตฺวา กุกฺกุฬาติ, อามนฺตา.\csromanspacer{} ทานํ อนิฏฺฐผลํ อกนฺตผลํ อมนุญฺญผลํ เสจนกผลํ ทุกฺขุทฺรยํ ทุกฺขวิปากนฺติ, น เหวํ วตฺตพฺเพ.}

\prose{สีลํ ฯเปฯ.\csromanspacer{} อุโปสโถ ฯเปฯ.\csromanspacer{} ภาวนา ฯเปฯ.\csromanspacer{} พฺรหฺมจริยํ อนิฏฺฐผลํ อกนฺตผลํ อมนุญฺญผลํ เสจนกผลํ ทุกฺขุทฺรยํ ทุกฺขวิปากนฺติ, น เหวํ วตฺตพฺเพ.}

\prose{นนุ ทานํ อิฏฺฐผลํ กนฺตผลํ มนุญฺญผลํ อเสจนกผลํ สุขุทฺรยํ สุขวิปากนฺติ, อามนฺตา.\csromanspacer{} หญฺจิ ทานํ อิฏฺฐผลํ กนฺตผลํ มนุญฺญผลํ อเสจนกผลํ สุขุทฺรยํ สุขวิปากํ, โน จ วต เร วตฺตพฺเพ “สพฺเพ สงฺขารา อโนธิํ กตฺวา กุกฺกุฬา”ติ.}

\prose{นนุ สีลํ.\csromanspacer{} อุโปสโถ.\csromanspacer{} ภาวนา.\csromanspacer{} พฺรหฺมจริยํ อิฏฺฐผลํ กนฺตผลํ มนุญฺญผลํ อเสจนกผลํ สุขุทฺรยํ สุขวิปากนฺติ, อามนฺตา.\csromanspacer{} หญฺจิ พฺรหฺมจริยํ อิฏฺฐผลํ กนฺตผลํ มนุญฺญผลํ อเสจนกผลํ สุขุทฺรยํ สุขวิปากนฺติ, โน จ วต เร วตฺตพฺเพ “สพฺเพ สงฺขารา อโนธิํ กตฺวา กุกฺกุฬา”ติ.}

\chapterheadpageread{2. ทุติยวคฺค}
\csromanmatikaanchor{mtk.38}
\csromantocmarkref{subsection}{(18) 9. อนุปุพฺพาภิสมยกถา}{mtk.38}
\bookmark[dest={mtk.38},level=2]{(18) 9. อนุปุพฺพาภิสมยกถา}
\prose{\csromanfolio{163}\csromansymbolmark{*}สพฺเพ สงฺขารา อโนธิํ กตฺวา กุกฺกุฬาติ, อามนฺตา.\csromanspacer{} นนุ วุตฺตํ ภควตา\csromandash{}}

\csromangathameasure{“สุโข วิเวโก ตุฏฺฐสฺส, สุตธมฺมสฺส ปสฺสโต. \\ อพฺยาปชฺชํ สุขํ โลเก, ปาณภูเตสุ สํยโม. \\ สุขา วิราคตา โลเก, กามานํ สมติกฺกโม. \\ อสฺมิมานสฺส โย วินโย, เอตํ เว ปรมํ สุขํ. \\ ตํ สุเขน สุขํ ปตฺตํ, อจฺจนฺตสุขเมว ตํ. \\ ติสฺโส วิชฺชา อนุปฺปตฺตา, เอตํ เว ปรมํ สุขนฺ”ติ.}
\csromangathagroup{\csromangathastanza{“สุโข วิเวโก ตุฏฺฐสฺส, สุตธมฺมสฺส ปสฺสโต. \\ อพฺยาปชฺชํ สุขํ โลเก, ปาณภูเตสุ สํยโม.}\csromangathastanza{สุขา วิราคตา โลเก, กามานํ สมติกฺกโม. \\ อสฺมิมานสฺส โย วินโย, เอตํ เว ปรมํ สุขํ\footnote{วิ 3. 4 ปิฏฺเฐ; ขุ 1. 88 อุทาเน จ.}.}\csromangathastanza{ตํ สุเขน สุขํ ปตฺตํ, อจฺจนฺตสุขเมว ตํ. \\ ติสฺโส วิชฺชา อนุปฺปตฺตา, เอตํ เว ปรมํ สุขนฺ”ติ.}}

\prose{อตฺเถว สุตฺตนฺโตติ, อามนฺตา.\csromanspacer{} เตน หิ น วตฺตพฺพํ “สพฺเพ สงฺขารา อโนธิํ กตฺวา กุกฺกุฬา”ติ.}

\nitthitam{กุกฺกุฬกถา นิฏฺฐิตา.}
\chapterhead{2. ทุติยวคฺค}
\vspace{\tipitakaparskip}
\csromancenter{\textbf{(18) 9.}\csromanspacer{}\textbf{ อนุปุพฺพาภิสมยกถา}}
\proseitem{339}{\csromansymbolmark{*}อนุปุพฺพาภิสมโยติ, อามนฺตา.\csromanspacer{} อนุปุพฺเพน โสตาปตฺติมคฺคํ ภาเวตีติ, น เหวํ วตฺตพฺเพ.\csromanspacer{} อนุปุพฺเพน โสตาปตฺติมคฺคํ ภาเวตีติ, อามนฺตา.\csromanspacer{} อนุปุพฺเพน โสตาปตฺติผลํ สจฺฉิกโรตีติ, น เหวํ วตฺตพฺเพ.}

\prose{\csromansymbolmark{*}อนุปุพฺพาภิสมโยติ, อามนฺตา.\csromanspacer{} อนุปุพฺเพน สกทาคามิมคฺคํ ภาเวตีติ, น เหวํ วตฺตพฺเพ.\csromanspacer{} อนุปุพฺเพน สกทาคามิมคฺคํ ภาเวตีติ, อามนฺตา.\csromanspacer{} อนุปุพฺเพน สกทาคามิผลํ สจฺฉิกโรตีติ, น เหวํ วตฺตพฺเพ.}

\prose{\csromansymbolmark{*}อนุปุพฺพาภิสมโยติ, อามนฺตา.\csromanspacer{} อนุปุพฺเพน อนาคามิมคฺคํ ภาเวตีติ, น เหวํ วตฺตพฺเพ.\csromanspacer{} อนุปุพฺเพน อนาคามิมคฺคํ ภาเวตีติ, อามนฺตา.\csromanspacer{} อนุปุพฺเพน อนาคามิผลํ สจฺฉิกโรตีติ, น เหวํ วตฺตพฺเพ.}

\prose{\csromansymbolmark{*}อนุปุพฺพาภิสมโยติ, อามนฺตา.\csromanspacer{} อนุปุพฺเพน อรหตฺตมคฺคํ ภาเวตีติ, น เหวํ วตฺตพฺเพ.\csromanspacer{} อนุปุพฺเพน อรหตฺตมคฺคํ ภาเวตีติ, อามนฺตา.\csromanspacer{} อนุปุพฺเพน อรหตฺตผลํ สจฺฉิกโรตีติ, น เหวํ วตฺตพฺเพ.}

\proseitem{340}{\csromanfolio{164}\csromansymbolmark{*}โสตาปตฺติผลสจฺฉิกิริยาย ปฏิปนฺโน ปุคฺคโล ทุกฺขทสฺสเนน กิํ ชหตีติ? สกฺกายทิฏฺฐิํ วิจิกิจฺฉํ สีลพฺพตปรามาสํ ตเทกฏฺเฐ จ กิเลเส จตุภาคํ ชหตีติ.\csromanspacer{} จตุภาคํ โสตาปนฺโน จตุภาคํ น โสตาปนฺโน จตุภาคํ โสตาปตฺติผลปฺปตฺโต ปฏิลทฺโธ อธิคโต สจฺฉิกโต อุปสมฺปชฺช วิหรติ, กาเยน ผุสิตฺวา วิหรติ, จตุภาคํ น กาเยน ผุสิตฺวา วิหรติ, จตุภาคํ สตฺตกฺขตฺตุปรโม โกลํโกโล เอกพีชี พุทฺเธ อเวจฺจปฺปสาเทน สมนฺนาคโต.\csromanspacer{} ธมฺเม ฯเปฯ สํเฆ ฯเปฯ อริยกนฺเตหิ สีเลหิ สมนฺนาคโต จตุภาคํ น อริยกนฺเตหิ สีเลหิ สมนฺนาคโตติ, น เหวํ วตฺตพฺเพ.}

\prose{\csromansymbolmark{*}สมุทยทสฺสเนน ฯเปฯ.\csromanspacer{} นิโรธทสฺสเนน ฯเปฯ.\csromanspacer{} มคฺคทสฺสเนน กิํ ชหตีติ, สกฺกายทิฏฺฐิํ วิจิกิจฺฉํ สีลพฺพตปรามาสํ ตเทกฏฺเฐ จ กิเลเส จตุภาคํ ชหตีติ.\csromanspacer{} จตุภาคํ โสตาปนฺโน จตุภาคํ น โสตาปนฺโน จตุภาคํ โสตาปตฺติผลปฺปโต ปฏิลทฺโธ อธิคโต สจฺฉิกโต อุปสมฺปชฺช วิหรติ กาเยน ผุสิตฺวา วิหรติ.\csromanspacer{} จตุภาคํ น กาเยน ผุสิตฺวา วิหรติ จตุภาคํ สตฺตกฺขตฺตุปรโม โกลํโกโล เอกพีชี พุทฺเธ อเวจฺจปฺปสาเทน สมนฺนาคโต.\csromanspacer{} ธมฺเม ฯเปฯ สํเฆ ฯเปฯ อริยกนฺเตหิ สีเลหิ สมนฺนาคโต จตุภาคํ น อริยกนฺเตหิ สีเลหิ สมนฺนาคโตติ, น เหวํ วตฺตพฺเพ ฯเปฯ.}

\proseitem{341}{\csromansymbolmark{*}สกทาคามิผลสจฺฉิกิริยาย ปฏิปนฺโน ปุคฺคโล ทุกฺขทสฺสเนน กิํ ชหตีติ, โอฬาริกํ กามราคํ โอฬาริกํ พฺยาปาทํ ตเทกฏฺเฐ จ กิเลเส จตุภาคํ ชหตีติ.\csromanspacer{} จตุภาคํ สกทาคามี จตุภาคํ น สกทาคามี.\csromanspacer{} จตุภาคํ สกทาคามิผลปฺปตฺโต ปฏิลทฺโท อธิคโต สจฺฉิกโต อุปสมฺปชฺช วิหรติ.\csromanspacer{} กาเยน ผุสิตฺวา วิหรติ.\csromanspacer{} จตุภาคํ น กาเยน ผุสิตฺวา วิหรตีติ, น เหวํ วตฺตพฺเพ ฯเปฯ.\csromanspacer{} สมุทยทสฺสเนน ฯเปฯ.\csromanspacer{} นิโรธทสฺสเนน ฯเปฯ.\csromanspacer{} มคฺคทสฺสเนน กิํ ชหตีติ, โอฬาริกํ กามราคํ โอฬาริกํ พฺยาปาทํ ตเทกฏฺเฐ จ กิเลเส จตุภาคํ ชหตีติ.\csromanspacer{} จตุภาคํ สกทาคามี จตุภาคํ น สกทาคามี จตุภาคํ สกทาคามิผลปฺปตฺโต ปฏิลทฺโธ อธิคโต สจฺฉิกโต อุปสมฺปชฺช วิหรติ.\csromanspacer{} กาเยน ผุสิตฺวา วิหรติ.\csromanspacer{} จตุภาคํ น กาเยน ผุสิตฺวา วิหรตีติ, น เหวํ วตฺตพฺเพ ฯเปฯ.}

\chapterheadpageread{2. ทุติยวคฺค}
\proseitem{342}{\csromanfolio{165}\csromansymbolmark{*}อนาคามิผลสจฺฉิกิริยาย ปฏิปนฺโน ปุคฺคโล ทุกฺขทสฺสเนน กิํ ชหตีติ, อณุสหคตํ กามราคํ อณุสหคตํ พฺยาปาทํ ตเทกฏฺเฐ จ กิเลเส จตุภาคํ ชหตีติ.\csromanspacer{} จตุภาคํ อนาคามี จตุภาคํ น อนาคามี จตุภาคํ อนาคามิผลปฺปตฺโต ปฏิลทฺโธ อธิคโต สจฺฉิกโต อุปสมฺปชฺช วิหรติ.\csromanspacer{} กาเยน ผุสิตฺวา วิหรติ.\csromanspacer{} จตุภาคํ น กาเยน ผุสิตฺวา วิหรติ.\csromanspacer{} จตุภาคํ อนฺตราปรินิพฺพายี ฯเปฯ.\csromanspacer{} อุปหจฺจปรินิพฺพายี.\csromanspacer{} อสงฺขารปรินิพฺพายี.\csromanspacer{} สสงฺขารปรินิพฺพายี.\csromanspacer{} อุทฺธํโสโต อกนิฏฺฐคามี จตุภาคํ น อุทฺธํโสโต น อกนิฏฺฐคามีติ, น เหวํ วตฺตพฺเพ ฯเปฯ.}

\prose{\csromansymbolmark{*}สมุทยทสฺสเนน ฯเปฯ.\csromanspacer{} นิโรธทสฺสเนน ฯเปฯ.\csromanspacer{} มคฺคทสฺสเนน กิํ ชหตีติ, อณุสหคตํ กามราคํ อณุสหคตํ พฺยาปาทํ ตเทกฏฺเฐ จ กิเลเส จตุภาคํ ชหตีติ.\csromanspacer{} จตุภาคํ อนาคามี จตุภาคํ น อนาคามี จตุภาคํ อนาคามิผลปฺปตฺโต ปฏิลทฺโธ อธิคโต สจฺฉิกโต อุปสมฺปชฺช วิหรติ.\csromanspacer{} กาเยน ผุสิตฺวา วิหรติ.\csromanspacer{} จตุภาคํ น กาเยน ผุสิตฺวา วิหรติ.\csromanspacer{} จตุภาคํ อนฺตราปรินิพฺพายี ฯเปฯ.\csromanspacer{} อุปหจฺจปรินิพฺพายี.\csromanspacer{} อสงฺขารปรินิพฺพายี.\csromanspacer{} สสงฺขารปรินิพฺพายี.\csromanspacer{} อุทฺธํโสโต อกนิฏฺฐคามี จตุภาคํ น อุทฺธํโสโต อกนิฏฺฐคามีติ, น เหวํ วตฺตพฺเพ ฯเปฯ.}

\proseitem{343}{\csromansymbolmark{*}อรหตฺตสจฺฉิกิริยาย ปฏิปนฺโน ปุคฺคโล ทุกฺขทสฺสเนน กิํ ชหตีติ, รูปราคํ อรูปราคํ มานํ อุทฺธจฺจํ อวิชฺชํ ตเทกฏฺเฐ จ กิเลเส จตุภาคํ ชหตีติ.\csromanspacer{} จตุภาคํ อรหา จตุภาคํ น อรหา จตุภาคํ อรหตฺตปฺปตฺโต ปฏิลทฺโธ อธิคโต สจฺฉิกโต อุปสมฺปชฺช วิหรติ, กาเยน ผุสิตฺวา วิหรติ.\csromanspacer{} จตุภาคํ น กาเยน ผุสิตฺวา วิหรติ.\csromanspacer{} จตุภาคํ วีตราโค ฯเปฯ.\csromanspacer{} วีตโทโส.\csromanspacer{} วีตโมโห ฯเปฯ.\csromanspacer{} กตกรณีโย โอหิตภาโร อนุปฺปตฺตสทตฺโถ ปริกฺขีณภวสํโยชโน สมฺมทญฺญาวิมุตฺโต อุกฺขิตฺตปลิโฆ สํกิณฺณปริโข อพฺพูเฬฺหสิโก นิรคฺคโฬ อริโย ปนฺนทฺธโช ปนฺนภาโร วิสญฺญุตฺโต สุวิชิตวิชโย, ทุกฺขํ ตสฺส ปริญฺญาตํ, สมุทโย ปหีโน, นิโรโธ สจฺฉิกโต, มคฺโค ภาวิโต, อภิญฺเญยฺยํ อภิญฺญาตํ, ปริญฺเญยฺยํ ปริญฺญาตํ, ปหาตพฺพํ ปหีนํ, ภาเวตพฺพํ ภาวิตํ ฯเปฯ สจฺฉิกาตพฺพํ สจฺฉิกตํ.\csromanspacer{} จตุภาคํ สจฺฉิกาตพฺพํ น สจฺฉิกตนฺติ, น เหวํ วตฺตพฺเพ ฯเปฯ.}

\prose{\csromanfolio{166}\csromansymbolmark{*}สมุทยทสฺสเนน.\csromanspacer{} นิโรธทสฺสเนน.\csromanspacer{} มคฺคทสฺสเนน กิํ ชหตีติ, รูปราคํ อรูปราคํ มานํ อุทฺธจฺจํ อวิชฺชํ ตเทกฏฺเฐ จ กิเลเส จตุภาคํ ชหตีติ.\csromanspacer{} จตุภาคํ อรหา จตุภาคํ น อรหา จตุภาคํ อรหตฺตปฺปตฺโต ปฏิลทฺโธ อธิคโต สจฺฉิกโต อุปสมฺปชฺช วิหรติ.\csromanspacer{} กาเยน ผุสิตฺวา วิหรติ.\csromanspacer{} จตุภาคํ น กาเยน ผุสิตฺวา วิหรติ.\csromanspacer{} จตุภาคํ วีตราโค ฯเปฯ วีตโทโส วีตโมโห กตกรณีโย โอหิตภาโร อนุปฺปตฺตสทตฺโถ ปริกฺขีณภวสํโยชโน สมฺมทญฺญาวิมุตฺโต อุกฺขิตฺตปลิโฆ สํกิณฺณปริโข อพฺพูเฬฺหสิโก นิรคฺคโฬ อริโย ปนฺนทฺธโช ปนฺนภาโร วิสญฺญุตฺโต สุวิชิตวิชโย, ทุกฺขํ ตสฺส ปริญฺญาตํ, สมุทโย ปหีโน, นิโรโธ สจฺฉิกโต, มคฺโค ภาวิโต, อภิญฺเญยฺยํ อภิญฺญาตํ, ปริญฺเญยฺยํ ปริญฺญาตํ, ปหาตพฺพํ ปหีนํ, ภาวิตพฺพํ ภาวิตํ ฯเปฯ สจฺฉิกาตพฺพํ สจฺฉิกตํ จตุภาคํ สจฺฉิกาตพฺพํ น สจฺฉิกตนฺติ, น เหวํ วตฺตพฺเพ ฯเปฯ.}

\proseitem{344}{\csromansymbolmark{*}โสตาปตฺติผลสจฺฉิกิริยาย ปฏิปนฺโน ปุคฺคโล ทุกฺขํ ทกฺขนฺโต “ปฏิปนฺนโก”ติ วตฺตพฺโพติ, อามนฺตา.\csromanspacer{} ทุกฺเข ทิฏฺเฐ “ผเล ฐิโต”ติ วตฺตพฺโพติ, น เหวํ วตฺตพฺเพ.\csromanspacer{} สมุทยํ ทกฺขนฺโต ฯเปฯ.\csromanspacer{} นิโรธํ ทกฺขนฺโต “ปฏิปนฺนโก”ติ วตฺตพฺโพติ, อามนฺตา.\csromanspacer{} นิโรเธ ทิฏฺเฐ “ผเล ฐิโก”ติ วตฺตพฺโพติ, น เหวํ วตฺตพฺเพ.}

\prose{\csromansymbolmark{*}โสตาปตฺติผลสจฺฉิกิริยาย ปฏิปนฺโน ปุคฺคโล มคฺคํ ทกฺขนฺโต “ปฏิปนฺนโก”ติ วตฺตพฺโพ, มคฺเค ทิฏฺเฐ “ผเล ฐิโต”ติ วตฺตพฺโพติ, อามนฺตา.\csromanspacer{} ทุกฺขํ ทกฺขนฺโต “ปฏิปนฺนโก”ติ วตฺตพฺโพ, ทุกฺเข ทิฏฺเฐ “ผเล ฐิโต”ติ วตฺตพฺโพติ, น เหวํ วตฺตพฺเพ ฯเปฯ.\csromanspacer{} มคฺคํ ทกฺขนฺโต “ปฏิปนฺนโก”ติ วตฺตพฺโพ, มคฺเค ทิฏฺเฐ “ผเล ฐิโต”ติ วตฺตพฺโพติ, อามนฺตา.\csromanspacer{} สมุทยํ ทกฺขนฺโต ฯเปฯ.\csromanspacer{} นิโรธํ ทกฺขนฺโต “ปฏิปนฺนโก”ติ วตฺตพฺโพ, นิโรเธ ทิฏฺเฐ “ผเล ฐิโต”ติ วตฺตพฺโพติ, น เหวํ วตฺตพฺเพ ฯเปฯ.}

\prose{\csromansymbolmark{*}โสตาปตฺติผลสจฺฉิกิริยาย ปฏิปนฺโน ปุคฺคโล ทุกฺขํ ทกฺขนฺโต “ปฏิปนฺนโก”ติ วตฺตพฺโพ, ทุกฺเข ทิฏฺเฐ น วตฺตพฺพํ “ผเล ฐิโตติ วตฺตพฺโพ”ติ, อามนฺตา.\csromanspacer{} มคฺคํ ทกฺขนฺโต “ปฏิปนฺนโก”ติ วตฺตพฺโพ, มคฺเค ทิฏฺเฐ น วตฺตพฺพํ “ผเล ฐิโตติ วตฺตพฺโพ”ติ, น เหวํ วตฺตพฺเพ ฯเปฯ.\csromanspacer{} สมุทยํ ทกฺขนฺโต.\csromanspacer{} นิโรธํ ทกฺขนฺโต “ปฏิปนฺนโก”ติ วตฺตพฺโพ, นิโรเธ ทิฏฺเฐ น วตฺตพฺพํ “ผเล ฐิโตติ วตฺตพฺโพ”ติ, อามนฺตา.\csromanspacer{} มคฺคํ ทกฺขนฺโต}

\vspace{\tipitakaparskip}
\csromancenter{ทุติยวคฺค “ปฏิปนฺนโก”ติ วตฺตพฺโพ, มคฺเค ทิฏฺเฐ น วตฺตพฺพํ “ผเล ฐิโตติ วตฺตพฺโพ”ติ, น เหวํ วตฺตพฺเพ ฯเปฯ.}
\prose{\csromanfolio{167}\csromansymbolfootnote{*}{สกวาทีปุจฺฉาลกฺขณํ.}โสตาปตฺติผลสจฺฉิกิริยาย ปฏิปนฺโน ปุคฺคโล ทุกฺขํ ทกฺขนฺโต “ปฏิปนฺนโก”ติ วตฺตพฺโพ, ทุกฺเข ทิฏฺเฐ น วตฺตพฺพํ “ผเล ฐิโตติ วตฺตพฺโพ”ติ,อามนฺตา.\csromanspacer{} นิรตฺถิยํ ทุกฺขทสฺสนนฺติ, น เหวํ วตฺตพฺเพ ฯเปฯ.\csromanspacer{} สมุทยํ ทกฺขนฺโต ฯเปฯ.\csromanspacer{} นิโรธํ ทกฺขนฺโต “ปฏิปนฺนโก”ติ วตฺตพฺโพ.\csromanspacer{} นิโรเธ ทิฏฺเฐ น วตฺตพฺพํ “ผเล ฐิโตติ วตฺตพฺโพ”ติ, อามนฺตา.\csromanspacer{} นิรตฺถิยํ นิโรธทสฺสนนฺติ, น เหวํ วตฺตพฺเพ ฯเปฯ.}

\proseitem{345}{\csromansymbolfootnote{+}{ปรวาทีปุจฺฉาลกฺขณํ.}ทุกฺเข ทิฏฺเฐ จตฺตาริ สจฺจานิ ทิฏฺฐานิ โหนฺตีติ, อามนฺตา.\csromanspacer{} ทุกฺขสจฺจํ จตฺตาริ สจฺจานีติ, น เหวํ วตฺตพฺเพ ฯเปฯ.}

\prose{\csromansymbolmark{*}รูปกฺขนฺเธ อนิจฺจโต ทิฏฺเฐ ปญฺจกฺขนฺธา อนิจฺจโต ทิฏฺฐา โหนฺตีติ, อามนฺตา.\csromanspacer{} รูปกฺขนฺโธ ปญฺจกฺขนฺธาติ, น เหวํ วตฺตพฺเพ ฯเปฯ.}

\prose{\csromansymbolmark{*}จกฺขายตเน อนิจฺจโต ทิฏฺเฐ ทฺวาทสายตนานิ อนิจฺจโต ทิฏฺฐานิ โหนฺตีติ, อามนฺตา.\csromanspacer{} จกฺขายตนํ ทฺวาทสายตนานีติ, น เหวํ วตฺตพฺเพ ฯเปฯ.}

\prose{\csromansymbolmark{*}จกฺขุธาตุยา อนิจฺจโต ทิฏฺฐาย อฏฺฐารส ธาตุโย อนิจฺจโต ทิฏฺฐา โหนฺตีติ, อามนฺตา.\csromanspacer{} จกฺขุธาตุ อฏฺฐารส ธาตุโยติ, น เหวํ วตฺตพฺเพ ฯเปฯ.}

\prose{\csromansymbolmark{*}จกฺขุนฺทฺริเย อนิจฺจโต ทิฏฺเฐ พาวีสตินฺทฺริยานิ อนิจฺจโต ทิฏฺฐานิ โหนฺตีติ, อามนฺตา.\csromanspacer{} จกฺขุนฺทฺริยํ พาวีสตินฺทฺริยานีติ, น เหวํ วตฺตพฺเพ ฯเปฯ.}

\prose{\csromansymbolmark{*}จตูหิ ญาเณหิ โสตาปตฺติผลํ สจฺฉิกโรตีติ, อามนฺตา.\csromanspacer{} จตฺตาริ โสตาปตฺติผลานีติ, น เหวํ วตฺตพฺเพ ฯเปฯ.\csromanspacer{} อฏฺฐหิ ญาเณหิ โสตาปตฺติผลํ สจฺฉิกโรตีติ, อามนฺตา.\csromanspacer{} อฏฺฐ โสตาปตฺติผลานีติ, น เหวํ วตฺตพฺเพ ฯเปฯ.\csromanspacer{} ทฺวาทสหิ ญาเณหิ โสตาปตฺติผลํ สจฺฉิกโรตีติ, อามนฺตา.\csromanspacer{} ทฺวาทส โสตาปตฺติผลานีติ, น เหวํ วตฺตพฺเพ ฯเปฯ.\csromanspacer{} จตุจตฺตารีสาย ญาเณหิ โสตาปตฺติผลํ สจฺฉิกโรตีติ, อามนฺตา.\csromanspacer{} จตุจตฺตารีสํ โสตาปตฺติผลานีติ, น เหวํ วตฺตพฺเพ ฯเปฯ.\csromanspacer{} สตฺตสตฺตติยา ญาเณหิ โสตาปตฺติผลํ สจฺฉิกโรตีติ, อามนฺตา.\csromanspacer{} สตฺตสตฺตติ โสตาปตฺติผลานีติ, น เหวํ วตฺตพฺเพ ฯเปฯ.}

\proseitem{346}{\csromanfolio{168}\csromansymbolmark{+}น วตฺตพฺพํ อนุปุพฺพาภิสมโยติ, อามนฺตา.\csromanspacer{} นนุ วุตฺตํ ภควตา\csromandash{}“เสยฺยถาปิ ภิกฺขเว มหาสมุทฺโท อนุปุพฺพนินฺโน อนุปุพฺพโปโณ อนุปุพฺพปพฺภาโร น อายตเกเนว ปปาโต.\csromanspacer{} เอวเมว โข ภิกฺขเว อิมสฺมิํ ธมฺมวินเย อนุปุพฺพสิกฺขา อนุปุพฺพกิริยา อนุปุพฺพปฏิปทา น อายตเกเนว อญฺญาปฏิเวโธ”ติ\footnote{วิ 4. 421; อํ 3. 46 ปิฏฺเฐ; ขุ 1. 141 อุทาเน จ.} อตฺเถว สุตฺตนฺโตติ, อามนฺตา.\csromanspacer{} เตน หิ อนุปุพฺพาภิสมโยติ.}

\prose{\csromansymbolmark{+}น วตฺตพฺพํ อนุปุพฺพาภิสมโยติ, อามนฺตา.\csromanspacer{} นนุ วุตฺตํ ภควตา\csromandash{}}

\csromangathameasure{“อนุปุพฺเพน เมธาวี, โถกํ โถกํ ขเณ ขเณ. \\ กมฺมาโร รชตสฺเสว, นิทฺธเม มลมตฺตโน”ติ.}
\csromangathagroup{\csromangathastanza{“อนุปุพฺเพน เมธาวี, โถกํ โถกํ ขเณ ขเณ. \\ กมฺมาโร รชตสฺเสว, นิทฺธเม มลมตฺตโน”ติ\footnote{ขุ 1. 48 ธมฺมปเท.}.}}

\prose{อตฺเถว สุตฺตนฺโตติ, อามนฺตา.\csromanspacer{} เตน หิ อนุปุพฺพาภิสมโยติ.}

\prose{\csromansymbolmark{*}อนุปุพฺพาภิสมโยติ, อามนฺตา.\csromanspacer{} นนฺวายสฺมา ควมฺปติ เถโร ภิกฺขู เอตทโวจ\csromandash{}สมฺมุขา เมตํ อาวุโส ภควโต สุตํ สมฺมุขา ปฏิคฺคหิตํ “โย ภิกฺขเว ทุกฺขํ ปสฺสติ, ทุกฺขสมุทยมฺปิ โส ปสฺสติ, ทุกฺขนิโรธมฺปิ ปสฺสติ, ทุกฺขนิโรธคามินิํ ปฏิปทมฺปิ ปสฺสติ.\csromanspacer{} โย ทุกฺขสมุทยํ ปสฺสติ, ทุกฺขมฺปิ โส ปสฺสติ, ทุกฺขนิโรธมฺปิ ปสฺสติ, ทุกฺขนิโรธคามินิํ ปฏิปทมฺปิ ปสฺสติ.\csromanspacer{} โย ทุกฺขนิโรธํ ปสฺสติ, ทุกฺขมฺปิ โส ปสฺสติ, ทุกฺขสมุทยมฺปิ ปสฺสติ, ทุกฺขนิโรธคามินิํ ปฏิปทมฺปิ ปสฺสติ.\csromanspacer{} โย ทุกฺขนิโรธคามินิํ ปฏิปทํ ปสฺสติ, ทุกฺขมฺปิ โส ปสฺสติ, ทุกฺขสมุทยมฺปิ ปสฺสติ, ทุกฺขนิโรธมฺปิ ปสฺสตี”ติ\footnote{สํ 3. 382 ปิฏฺเฐ.} อตฺเถว สุตฺตนฺโตติ, อามนฺตา.\csromanspacer{} เตน หิ น วตฺตพฺพํ อนุปุพฺพาภิสมโยติ.}

\prose{\csromansymbolmark{*}อนุปุพฺพาภิสมโยติ, อามนฺตา.\csromanspacer{} นนุ วุตฺตํ ภควตา\csromandash{}}

\csromangathameasure{“สหาวสฺส ทสฺสนสมฺปทาย, \\ ตยสฺสุ ธมฺมา ชหิตา ภวนฺติ. \\ สกฺกายทิฏฺฐี วิจิกิจฺฉิตญฺจ, \\ สีลพฺพตํ วาปิ ยทตฺถิ กิญฺจิ. \\ จตูหปาเยหิ จ วิปฺปมุตฺโต, \\ ฉจฺจาภิฐานานิ อภพฺพ กาตุนฺ”ติ.}
\csromangathagroup{\csromangathastanza{“สหาวสฺส ทสฺสนสมฺปทาย, \\ ตยสฺสุ ธมฺมา ชหิตา ภวนฺติ. \\ สกฺกายทิฏฺฐี วิจิกิจฺฉิตญฺจ, \\ สีลพฺพตํ วาปิ ยทตฺถิ กิญฺจิ. \\ จตูหปาเยหิ จ วิปฺปมุตฺโต, \\ ฉจฺจาภิฐานานิ อภพฺพ กาตุนฺ”ติ.}}

\chapterheadpageread{2. ทุติยวคฺค}
\csromanmatikaanchor{mtk.39}
\csromantocmarkref{subsection}{(19) 10. โวหารกถา}{mtk.39}
\bookmark[dest={mtk.39},level=2]{(19) 10. โวหารกถา}
\prose{\csromanfolio{169}อตฺเถว สุตฺตนฺโตติ, อามนฺตา.\csromanspacer{} เตน หิ น วตฺตพฺพํ “อนุปุพฺพาภิสมโย”ติ.}

\prose{\csromansymbolmark{*}อนุปุพฺพาภิสมโยติ, อามนฺตา.\csromanspacer{} นนุ วุตฺตํ ภควตา “ยสฺมิํ ภิกฺขเว สมเย อริยสาวกสฺส วิรชํ วีตมลํ ธมฺมจกฺขุํ อุทปาทิ ‘ยํ กิญฺจิ สมุทยธมฺมํ, สพฺพํ ตํ นิโรธธมฺมนฺ’ติ.\csromanspacer{} สห ทสฺสนุปฺปาทา ภิกฺขเว อริยสาวกสฺส ตีณิ สํโยชนานิ ปหียนฺติ สกฺกายทิฏฺฐิ วิจิกิจฺฉา สีลพฺพตปรามาโส”ติ อตฺเถว สุตฺตนฺโตติ, อามนฺตา.\csromanspacer{} เตน หิ น วตฺตพฺพํ อนุปุพฺพาภิสมโยติ.}

\nitthitam{อนุปุพฺพาภิสมยกถา นิฏฺฐิต.}
\chapterhead{2. ทุติยวคฺค}
\vspace{\tipitakaparskip}
\csromancenter{\textbf{(19) 10.}\csromanspacer{}\textbf{ โวหารกถา}}
\proseitem{347}{\csromansymbolmark{*}พุทฺธสฺส ภควโต โวหาโร โลกุตฺตโรติ, อามนฺตา.\csromanspacer{} โลกุตฺตเร โสเต ปฏิหญฺญติ โน โลกิเย.\csromanspacer{} โลกุตฺตเรน วิญฺญาเณน ปฏิวิชานนฺติ โน โลกิเยน.\csromanspacer{} สาวกา ปฏิวิชานนฺติ โน ปุถุชฺชนาติ, น เหวํ วตฺตพฺเพ ฯเปฯ.}

\prose{\csromansymbolmark{*}นนุ พุทฺธสฺส ภควโต โวหาโร โลกิเย โสเต ปฏิหญฺญตีติ, อามนฺตา.\csromanspacer{} หญฺจิ พุทฺธสฺส ภควโต โวหาโร โลกิเย โสเต ปฏิหญฺญติ, โน จ วต เร วตฺตพฺเพ “พุทฺธสฺส ภควโต โวหาโร โลกุตฺตโร”ติ.}

\prose{\csromansymbolmark{*}นนุ พุทฺธสฺส ภควโต โวหารํ โลกิเยน วิญฺญาเณน ปฏิวิชานนฺตีติ, อามนฺตา.\csromanspacer{} หญฺจิ พุทฺธสฺส ภควโต โวหารํ โลกิเยน วิญฺญาเณน ปฏิวิชานนฺติ โน จ วต เร วตฺตพฺเพ “พุทฺธสฺส ภควโต โวหาโร โลกุตฺตโร”ติ.}

\prose{\csromansymbolmark{*}นนุ พุทฺธสฺส ภควโต โวหารํ ปุถุชฺชนา ปฏิวิชานนฺตีติ, อามนฺตา.\csromanspacer{} หญฺจิ พุทฺธสฺส ภควโต โวหารํ ปุถุชฺชนา ปฏิวิชานนฺติ, โน จ วต เร วตฺตพฺเพ “พุทฺธสฺส ภควโต โวหาโร โลกุตฺตโร”ติ.}

\proseitem{348}{\csromanfolio{170}\csromansymbolmark{*}พุทฺธสฺส ภควโต โวหาโร โลกุตฺตโรติ, อามนฺตา.\csromanspacer{} มคฺโค ผลํ นิพฺพานํ โสตาปตฺติมคฺโค โสตาปตฺติผลํ สกทาคามิมคฺโค สกทาคามิผลํ อนาคามิมคฺโค อนาคามิผลํ อรหตฺตมคฺโค อรหตฺตผลํ.\csromanspacer{} สติปฏฺฐานํ สมฺมปฺปธานํ อิทฺธิปาโท อินฺทฺริยํ พลํ โพชฺฌงฺโคติ, น เหวํ วตฺตพฺเพ ฯเปฯ.}

\prose{\csromansymbolmark{*}พุทฺธสฺส ภควโต โวหาโร โลกุตฺตโรติ, อามนฺตา.\csromanspacer{} อตฺถิ เกจิ พุทฺธสฺส ภควโต โวหารํ สุณนฺตีติ, อามนฺตา.\csromanspacer{} โลกุตฺตโร ธมฺโม โสตวิญฺเญยฺโย, โสตสฺมิํ ปฏิหญฺญติ, โสตสฺส อาปาถํ อาคจฺฉตีติ, น เหวํ วตฺตพฺเพ ฯเปฯ.}

\prose{\csromansymbolmark{*}นนุ โลกุตฺตโร ธมฺโม น โสตวิญฺเญยฺโย, น โสตสฺมิํ ปฏิหญฺญติ, น โสตสฺส อาปาถํ อาคจฺฉตีติ, อามนฺตา.\csromanspacer{} หญฺจิ โลกุตฺตโร ธมฺโม น โสตวิญฺเญยฺโย, น โสตสฺมิํ ปฏิหญฺญติ, น โสตสฺส อาปาถํ อาคจฺฉติ, โน จ วต เร วตฺตพฺเพ “พุทฺธสฺส ภควโต โวหาโร โลกุตฺตโร”ติ.}

\proseitem{349}{\csromansymbolmark{*}พุทฺธสฺส ภควโต โวหาโร โลกุตฺตโรติ, อามนฺตา.\csromanspacer{} อตฺถิ เกจิ พุทฺธสฺส ภควโต โวหาเร รชฺเชยฺยุนฺติ, อามนฺตา.\csromanspacer{} โลกุตฺตโร ธมฺโม ราคฏฺฐานิโย รชนีโย กมนีโย มทนีโย พนฺธนีโย มุจฺฉนีโยติ, น เหวํ วตฺตพฺเพ ฯเปฯ.}

\prose{\csromansymbolmark{*}นนุ โลกุตฺตโร ธมฺโม น ราคฏฺฐานิโย น รชนีโย น กมนีโย น มทนีโย น พนฺธนีโย น มุจฺฉนีโยติ, อามนฺตา.\csromanspacer{} หญฺจิ โลกุตฺตโร ธมฺโม น ราคฏฺฐานิโย น รชนีโย น กมนีโย น มทนีโย น พนฺธนีโย น มุจฺฉนีโย, โน จ วต เร วตฺตพฺเพ “พุทฺธสฺส ภควโต โวหาโร โลกุตฺตโร”ติ.}

\prose{\csromansymbolmark{*}พุทฺธสฺส ภควโต โวหาโร โลกุตฺตโรติ, อามนฺตา.\csromanspacer{} อตฺถิ เกจิ พุทฺธสฺส ภควโต โวหาเร ทุสฺเสยฺยุนฺติ, อามนฺตา.\csromanspacer{} โลกุตฺตโร ธมฺโม โทสฏฺฐานิโย โกปฏฺฐานิโย ปฏิฆฏฺฐานิโยติ, น เหวํ วตฺตพฺเพ ฯเปฯ.}

\prose{\csromansymbolmark{*}นนุ โลกุตฺตโร ธมฺโม น โทสฏฺฐานิโย น โกปฏฺฐานิโย น ปฏิฆฏฺฐานิโยติ, อามนฺตา.\csromanspacer{} หญฺจิ โลกุตฺตโร ธมฺโม น โทสฏฺฐานิโย}

\vspace{\tipitakaparskip}
\csromancenter{ทุติยวคฺค น โกปฏฺฐานิโย น ปฏิฆฏฺฐานิโย, โน จ วต เร วตฺตพฺเพ “พุทฺธสฺส ภควโต โวหาโร โลกุตฺตโร”ติ.}
\prose{\csromanfolio{171}\csromansymbolmark{*}พุทฺธสฺส ภควโต โวหาโร โลกุตฺตโรติ, อามนฺตา.\csromanspacer{} อตฺถิ เกจิ พุทฺธสฺส ภควโต โวหาเร มุเยฺหยฺยุนฺติ, อามนฺตา.\csromanspacer{} โลกุตฺตโร ธมฺโม โมหฏฺฐานิโย อญฺญาณกรโณ อจกฺขุกรโณ ปญฺญานิโรธิโก วิฆาตปกฺขิโก อนิพฺพานสํวตฺตนิโกติ, น เหวํ วตฺตพฺเพ ฯเปฯ.}

\prose{\csromansymbolmark{*}นนุ โลกุตฺตโร ธมฺโม น โมหฏฺฐานิโย น อญฺญาณกรโณ น อจกฺขุกรโณ ปญฺญาวุทฺธิโก อวิฆาตปกฺขิโก นิพฺพานสํวตฺตนิโกติ, อามนฺตา.\csromanspacer{} หญฺจิ โลกุตฺตโร ธมฺโม น โมหฏฺฐานิโย น อญฺญาณกรโณ น อจกฺขุกรโณ ปญฺญาวุทฺธิโก อวิฆาตปกฺขิโก นิพฺพานสํวตฺตนิโก, โน จ วต เร วตฺตพฺเพ “พุทฺธสฺส ภควโต โวหาโร โลกุตฺตโร”ติ.}

\proseitem{350}{\csromansymbolmark{*}พุทฺธสฺส ภควโต โวหาโร โลกุตฺตโรติ, อามนฺตา.\csromanspacer{} เย เกจิ พุทฺธสฺส ภควโต โวหารํ สุณนฺติ, สพฺเพ เต มคฺคํ ภาเวนฺตีติ, น เหวํ วตฺตพฺเพ ฯเปฯ.}

\prose{\csromansymbolmark{*}เย เกจิ พุทฺธสฺส ภควโต โวหารํ สุณนฺติ, สพฺเพ เต มคฺคํ ภาเวนฺตีติ, อามนฺตา.\csromanspacer{} พาลปุถุชฺชนา พุทฺธสฺส ภควโต โวหารํ สุณนฺติ, พาลปุถุชฺชนา มคฺคํ ภาเวนฺตีติ, น เหวํ วตฺตพฺเพ ฯเปฯ.\csromanspacer{} มาตุฆาตโก มคฺคํ ภาเวติ ฯเปฯ.\csromanspacer{} ปิตุฆาตโก.\csromanspacer{} อรหนฺตฆาตโก.\csromanspacer{} รุหิรุปฺปาทโก.\csromanspacer{} สํฆเภทโก พุทฺธสฺส ภควโต โวหารํ สุณาติ, สํฆเภทโก มคฺคํ ภาเวตีติ, น เหวํ วตฺตพฺเพ ฯเปฯ.}

\proseitem{351}{\csromansymbolmark{+}ลพฺภา โสวณฺณมยาย ลฏฺฐิยา ธญฺญปุญฺโชปิ สุวณฺณปุญฺโชปิ อาจิกฺขิตุนฺติ, อามนฺตา.\csromanspacer{} เอวเมวํ ภควา โลกุตฺตเรน โวหาเรน โลกิยมฺปิ โลกุตฺตรมฺปิ ธมฺมํ โวหรตีติ.}

\prose{\csromansymbolmark{*}ลพฺภา เอลณฺฑิยาย ลฏฺฐิยา ธญฺญปุญฺโชปิ สุวณฺณปุญฺโชปิ อาจิกฺขิตุนฺติ, อามนฺตา.\csromanspacer{} เอวเมวํ ภควา โลกิเยน โวหาเรน โลกิยมฺปิ โลกุตฺตรมฺปิ ธมฺมํ โวหรตีติ.}

\proseitem{352}{\csromanfolio{172}\csromansymbolmark{*}พุทฺธสฺส ภควโต โวหาโร โลกิยํ โวหรนฺตสฺส โลกิโย โหติ, โลกุตฺตรํ โวหรนฺตสฺส โลกุตฺตโร โหตีติ, อามนฺตา.\csromanspacer{} โลกิยํ โวหรนฺตสฺส โลกิเย โสเต ปฏิหญฺญติ, โลกุตฺตรํ โวหรนฺตสฺส โลกุตฺตเร โสเต ปฏิหญฺญติ, โลกิยํ โวหรนฺตสฺส โลกิเยน วิญฺญาเณน ปฏิวิชานนฺติ, โลกุตฺตรํ โวหรนฺตสฺส โลกุตฺตเรน วิญฺญาเณน ปฏิวิชานนฺติ, โลกิยํ โวหรนฺตสฺส ปุถุชฺชนา ปฏิวิชานนฺติ, โลกุตฺตรํ โวหรนฺตสฺส สาวกา ปฏิวิชานนฺตีติ, น เหวํ วตฺตพฺเพ ฯเปฯ.}

\prose{\csromansymbolmark{+}น วตฺตพฺพํ พุทฺธสฺส ภควโต โวหาโร โลกิยํ โวหรนฺตสฺส โลกิโย โหติ, โลกุตฺตรํ โวหรนฺตสฺส โลกุตฺตโร โหตีติ, อามนฺตา.\csromanspacer{} นนุ ภควา โลกิยมฺปิ โลกุตฺตรมฺปิ ธมฺมํ โวหรตีติ, อามนฺตา.\csromanspacer{} หญฺจิ ภควา โลกิยมฺปิ โลกุตฺตรมฺปิ ธมฺมํ โวหรติ, เตน วต เร วตฺตพฺเพ “พุทฺธสฺส ภควโต โวหาโร โลกิยํ โวหรนฺตสฺส โลกิโย โหติ, โลกุตฺตรํ โวหรนฺตสฺส โลกุตฺตโร โหตี”ติ.}

\prose{\csromansymbolmark{*}พุทฺธสฺส ภควโต โวหาโร โลกิยํ โวหรนฺตสฺส โลกิโย โหติ, โลกุตฺตรํ โวหรนฺตสฺส โลกุตฺตโร โหตีติ, อามนฺตา.\csromanspacer{} มคฺคํ โวหรนฺตสฺส มคฺโค โหติ, อมคฺคํ โวหรนฺตสฺส อมคฺโค โหติ, ผลํ โวหรนฺตสฺส ผลํ โหติ, อผลํ โวหรนฺตสฺส อผลํ โหติ, นิพฺพานํ โวหรนฺตสฺส นิพฺพานํ โหติ, อนิพฺพานํ โวหรนฺตสฺส อนิพฺพานํ โหติ, สงฺขตํ โวหรนฺตสฺส สงฺขตํ โหติ, อสงฺขตํ โวหรนฺตสฺส อสงฺขตํ โหติ, รูปํ โวหรนฺตสฺส รูปํ โหติ, อรูปํ โวหรนฺตสฺส อรูปํ โหติ, เวทนํ โวหรนฺตสฺส เวทนา โหติ, อเวทนํ โวหรนฺตสฺส อเวทนา โหติ, สญฺญํ โวหรนฺตสฺส สญฺญา โหติ, อสญฺญํ โวหรนฺตสฺส อสญฺญา โหติ, สงฺขาเร โวหรนฺตสฺส สงฺขารา โหนฺติ, อสงฺขาเร โวหรนฺตสฺส อสงฺขารา โหนฺติ, วิญฺญาณํ โวหรนฺตสฺส วิญฺญาณํ โหติ, อวิญฺญาณํ โวหรนฺตสฺส อวิญฺญาณํ โหตีติ, น เหวํ วตฺตพฺเพ ฯเปฯ.}

\nitthitam{โวหารกถา นิฏฺฐิตา.}
\chapterheadpageread{2. ทุติยวคฺค \textbf{2. ทุติยวคฺค}}
\vspace{\tipitakaparskip}
\csromancenter{\textbf{(20) 11.}\csromanspacer{}\textbf{ นิโรธกถา}}
\csromanmatikaanchor{mtk.40}
\csromantocmarkref{subsection}{(20) 11. นิโรธกถา}{mtk.40}
\bookmark[dest={mtk.40},level=2]{(20) 11. นิโรธกถา}
\proseitem{353}{\csromanfolio{173}\csromansymbolmark{*}เทฺว นิโรธาติ, อามนฺตา.\csromanspacer{} เทฺว ทุกฺขนิโรธาติ, น เหวํ วตฺตพฺเพ ฯเปฯ.\csromanspacer{} เทฺว ทุกฺขนิโรธาติ, อามนฺตา.\csromanspacer{} เทฺว นิโรธสจฺจานีติ, น เหวํ วตฺตพฺเพ ฯเปฯ.\csromanspacer{} เทฺว นิโรธสจฺจานีติ, อามนฺตา.\csromanspacer{} เทฺว ทุกฺขสจฺจานีติ, น เหวํ วตฺตพฺเพ ฯเปฯ.\csromanspacer{} เทฺว นิโรธสจฺจานีติ, อามนฺตา.\csromanspacer{} เทฺว สมุทยสจฺจานีติ, น เหวํ วตฺตพฺเพ ฯเปฯ.\csromanspacer{} เทฺว นิโรธสจฺจานีติ, อามนฺตา.\csromanspacer{} เทฺว มคฺคสจฺจานีติ, น เหวํ วตฺตพฺเพ ฯเปฯ.}

\prose{\csromansymbolmark{*}เทฺว นิโรธสจฺจานีติ, อามนฺตา.\csromanspacer{} เทฺว ตาณานิ ฯเปฯ.\csromanspacer{} เทฺว เลณานิ.\csromanspacer{} เทฺว สรณานิ.\csromanspacer{} เทฺว ปรายณานิ.\csromanspacer{} เทฺว อจฺจุตานิ.\csromanspacer{} เทฺว อมตานิ.\csromanspacer{} เทฺว นิพฺพานานีติ, น เหวํ วตฺตพฺเพ ฯเปฯ.}

\prose{\csromansymbolmark{*}เทฺว นิพฺพานานีติ, อามนฺตา.\csromanspacer{} อตฺถิ ทฺวินฺนํ นิพฺพานานํ อุจฺจนีจตา หีนปณีตตา อุกฺกํสาวกํโส สีมา วา เภโท วา ราชิ วา อนฺตริกา วาติ, น เหวํ วตฺตพฺเพ ฯเปฯ.}

\prose{\csromansymbolmark{*}เทฺว นิโรธาติ, อามนฺตา.\csromanspacer{} นนุ อปฺปฏิสงฺขานิรุทฺเธ สงฺขาเร ปฏิสงฺขา นิโรเธนฺตีติ, อามนฺตา.\csromanspacer{} หญฺจิ อปฺปฏิสงฺขานิรุทฺเธ สงฺขาเร ปฏิสงฺขา นิโรเธนฺติ, โน จ วต เร วตฺตพฺเพ “เทฺว นิโรธา”ติ.}

\prose{\csromansymbolmark{+}น วตฺตพฺพํ “เทฺว นิโรธา”ติ, อามนฺตา.\csromanspacer{} นนุ อปฺปฏิสงฺขานิรุทฺธาปิ สงฺขารา อจฺจนฺตภคฺคา, ปฏิสงฺขานิรุทฺธาปิ สงฺขารา อจฺจนฺตภคฺคาติ, อามนฺตา.\csromanspacer{} หญฺจิ อปฺปฏิสงฺขานิรุทฺธาปิ สงฺขารา อจฺจนฺตภคฺคา, ปฏิสงฺขานิรุทฺธาปิ สงฺขารา อจฺจนฺตภคฺคา, เตน วต เร วตฺตพฺเพ “เทฺว นิโรธา”ติ.}

\prose{\csromansymbolmark{*}เทฺว นิโรธาติ, อามนฺตา.\csromanspacer{} ปฏิสงฺขานิรุทฺธาปิ\footnote{ปฏิสงฺขานิรุทฺธา (?)} สงฺขารา อริยมคฺคํ อาคมฺม นิรุทฺธาติ, อามนฺตา.\csromanspacer{} อปฺปฏิสงฺขานิรุทฺธา สงฺขารา อริยมคฺคํ อาคมฺม นิรุทฺธาติ, น เหวํ วตฺตพฺเพ ฯเปฯ.}

\csromanmatikaanchor{mtk.41}
\csromanmatikaanchor{mtk.42}
\csromantocmarknopage{chapter}{3. ตติยวคฺค}{mtk.41}
\bookmark[dest={mtk.41},level=0]{3. ตติยวคฺค}
\csromantocmarkref{subsection}{(21) 1. พลกถา}{mtk.42}
\bookmark[dest={mtk.42},level=2]{(21) 1. พลกถา}
\prose{\csromanfolio{174}\csromansymbolmark{*}เทฺว นิโรธาติ, อามนฺตา.\csromanspacer{} ปฏิสงฺขานิรุทฺธา สงฺขารา น ปุน อุปฺปชฺชนฺตีติ, อามนฺตา.\csromanspacer{} อปฺปฏิสงฺขานิรุทฺธา สงฺขารา น ปุน อุปฺปชฺชนฺตีติ, น เหวํ วตฺตพฺเพ ฯเปฯ.\csromanspacer{} เตน หิ น วตฺตพฺพํ “เทฺว นิโรธา”ติ.\csromanspacer{} นิโรธกถา นิฏฺฐิตา.\csromanspacer{} ทุติโย วคฺโค.}

\titlehead{\textbf{ตสฺสุทฺทานํ}}
\csromangathameasure{ปรูปหาโร อญฺญาณํ, กงฺขา ปรวิตารณา. \\ วจีเภโท ทุกฺขาหาโร, จิตฺตฏฺฐิติ จ กุกฺกุฬา. \\ อนุปุพฺพาภิสมโย, โวหาโร จ นิโรธโกติ.}
\csromangathagroup{\csromangathastanza{ปรูปหาโร อญฺญาณํ, กงฺขา ปรวิตารณา. \\ วจีเภโท ทุกฺขาหาโร, จิตฺตฏฺฐิติ จ กุกฺกุฬา. \\ อนุปุพฺพาภิสมโย, โวหาโร จ นิโรธโกติ.}}

\chapterhead{3. ตติยวคฺค}
\vspace{\tipitakaparskip}
\csromancenter{\textbf{(21) 1.}\csromanspacer{}\textbf{ พลกถา}}
\proseitem{354}{\csromansymbolmark{*}ตถาคตพลํ สาวกสาธารณนฺติ, อามนฺตา.\csromanspacer{} ตถาคตพลํ สาวกพลํ, สาวกพลํ ตถาคตพลนฺติ, น เหวํ วตฺตพฺเพ ฯเปฯ.}

\prose{\csromansymbolmark{*}ตถาคตพลํ สาวกสาธารณนฺติ, อามนฺตา.\csromanspacer{} ตญฺเญว ตถาคตพลํ, ตํ สาวกพลํ, ตญฺเญว สาวกพลํ, ตํ ตถาคตพลนฺติ, น เหวํ วตฺตพฺเพ ฯเปฯ.}

\prose{\csromansymbolmark{*}ตถาคตพลํ สาวกสาธารณนฺติ, อามนฺตา.\csromanspacer{} ยาทิสํ ตถาคตพลํ, ตาทิสํ สาวกพลํ, ยาทิสํ สาวกพลํ, ตาทิสํ ตถาคตพลนฺติ, น เหวํ วตฺตพฺเพ ฯเปฯ.}

\prose{\csromansymbolmark{*}ตถาคตพลํ สาวกสาธารณนฺติ, อามนฺตา.\csromanspacer{} ยาทิโส ตถาคตสฺส ปุพฺพโยโค ปุพฺพจริยา ธมฺมกฺขานํ ธมฺมเทสนา, ตาทิโส สาวกสฺส ปุพฺพโยโค ปุพฺพจริยา ธมฺมกฺขานํ ธมฺมเทสนาติ, น เหวํ วตฺตพฺเพ ฯเปฯ.}

\prose{\csromansymbolmark{*}ตถาคตพลํ สาวกสาธารณนฺติ, อามนฺตา.\csromanspacer{} ตถาคโต ชิโน สตฺถา สมฺมาสมฺพุทฺโธ สพฺพญฺญู สพฺพทสฺสาวี ธมฺมสฺสามี ธมฺมปฺปฏิสรโณติ,}

\vspace{\tipitakaparskip}
\csromancenter{ตติยวคฺค อามนฺตา.\csromanspacer{} สาวโก ชิโน สตฺถา สมฺมาสมฺพุทฺโธ สพฺพญฺญู สพฺพทสฺสาวี ธมฺมสฺสามี ธมฺมปฺปฏิสรโณติ, น เหวํ วตฺตพฺเพ ฯเปฯ.}
\prose{\csromanfolio{175}\csromansymbolmark{*}ตถคตพลํ สาวกสาธารณนฺติ, อามนฺตา.\csromanspacer{} ตถคโต อนุปฺปนฺนสฺส มคฺคสฺส อุปฺปาเทตา อสญฺชาตสฺส มคฺคสฺส สญฺชเนตา อนกฺขาตสฺส มคฺคสฺส อกฺขาตา มคฺคญฺญู มคฺควิทู มคฺคโกวิโทติ, อามนฺตา.\csromanspacer{} สาวโก อนุปฺปนฺนสฺส มคฺคสฺส อุปฺปาเทตา อสญฺชาตสฺส มคฺคสฺส สญฺชเนตา อนกฺขาตสฺส มคฺคสฺส อกฺขาตา มคฺคญฺญู มคฺควิทู มคฺคโกวิโทติ, น เหวํ วตฺตพฺเพ ฯเปฯ.}

\prose{\csromansymbolmark{*}อินฺทฺริยปโรปริยตฺตํ ยถาภูตํ ญาณํ ตถาคตพลํ สาวกสาธารณนฺติ, อามนฺตา.\csromanspacer{} สาวโก สพฺพญฺญู สพฺพทสฺสาวีติ, น เหวํ วตฺตพฺเพ ฯเปฯ.}

\proseitem{355}{\csromansymbolmark{+}สาวโก ฐานาฐานํ ชานาตีติ, อามนฺตา.\csromanspacer{} หญฺจิ สาวโก ฐานาฐานํ ชานาติ, เตน วต เร วตฺตพฺเพ “ฐานาฐานํ ยถาภูตํ ญาณํ ตถาคตพลํ สาวกสาธารณนฺ”ติ.}

\prose{\csromansymbolmark{+}สาวโก อตีตานาคตปจฺจุปฺปนฺนานํ กมฺมสมาทานานํ ฐานโส เหตุโส วิปากํ ชานาตีติ, อามนฺตา.\csromanspacer{} หญฺจิ สาวโก อตีตานาคตปจฺจุปฺปนฺนานํ กมฺมสมาทานานํ ฐานโส เหตุโส วิปากํ ชานาติ, เตน วต เร วตฺตพฺเพ “อตีตานาคตปจฺจุปฺปนฺนานํ กมฺมสมาทานานํ ฐานโส เหตุโส วิปากํ ยถาภูตํ ญาณํ ตถาคตพลํ สาวกสาธารณนฺ”ติ.}

\prose{\csromansymbolmark{+}สาวโก สพฺพตฺถคามินิํ ปฏิปทํ ชานาตีติ, อามนฺตา.\csromanspacer{} หญฺจิ สาวโก สพฺพตฺถคามินิํ ปฏิปทํ ชานาติ, เตน วต เร วตฺตพฺเพ “สพฺพตฺถคามินิํ ปฏิปทํ ยถาภูตํ ญาณํ ตถาคตพลํ สาวกสาธารณนฺ”ติ.}

\prose{\csromansymbolmark{+}สาวโก อเนกธาตุํ นานาธาตุํ โลกํ ชานาตีติ, อามนฺตา.\csromanspacer{} หญฺจิ สาวโก อเนกธาตุํ นานาธาตุํ โลกํ ชานาติ, เตน วต เร วตฺตพฺเพ “อเนกธาตุํ นานาธาตุํ โลกํ ยถาภูตํ ญาณํ ตถาคตพลํ สาวกสาธารณนฺ”ติ.}

\prose{\csromansymbolmark{+}สาวโก สตฺตานํ นานาธิมุตฺติกตํ ชานาตีติ, อามนฺตา.\csromanspacer{} หญฺจิ สาวกา สตฺตานํ นานาธิมุตฺติกตํ ชานาติ, เตน วต เร วตฺตพฺเพ \csromanfolio{176}“สตฺตานํ นานาธิมุตฺติกตํ ยถาภูตํ ญาณํ ตถาคตพลํ สาวกสาธารณนฺ”ติ.}

\prose{\csromansymbolmark{+}สาวโก ฌานวิโมกฺขสมาธิสมาปตฺตีนํ สํกิเลสํ โวทานํ วุฏฺฐานํ ชานาตีติ, อามนฺตา.\csromanspacer{} หญฺจิ สาวโก ฌานวิโมกฺขสมาธิสมาปตฺตีนํ สํกิเลสํ โวทานํ วุฏฺฐานํ ชานาติ, เตน วต เร วตฺตพฺเพ “ฌานวิโมกฺขสมาธิสมาปตฺตีนํ สํกิเลสํ โวทานํ วุฏฺฐานํ ยถาภูตํ ญาณํ ตถาคตพลํ สาวกสาธารณนฺติ.}

\prose{\csromansymbolmark{+}สาวโก ปุพฺเพนิวาสานุสฺสติํ ชานาตีติ, อามนฺตา.\csromanspacer{} หญฺจิ สาวโก ปุพฺเพนิวาสานุสฺสติํ ชานาติ, เตเน วต เร วตฺตพฺเพ “ปุพฺเพนิวาสานุสฺสติ ยถาภูตํ ญาณํ ตถาคตพลํ สาวกสาธารณนฺ”ติ.}

\prose{\csromansymbolmark{+}สาวโก สตฺตานํ จุตูปปาตํ ชานาตีติ, อามนฺตา.\csromanspacer{} หญฺจิ สาวโก สตฺตานํ จุตูปปาตํ ชานาติ, เตน วต เร วตฺตพฺเพ “สตฺตานํ จุตูปปาตํ ยถาภูตํ ญาณํ ตถาคตพลํ สาวกสาธารณนฺ”ติ.}

\prose{\csromansymbolmark{+}นนุ ตถาคตสฺสาปิ อาสวา ขีณา สาวกสฺสาปิ อาสวา ขีณาติ, อามนฺตา.\csromanspacer{} อตฺถิ กิญฺจิ นานากรณํ ตถาคตสฺส วา สาวกสฺส วา อาสวกฺขเยน วา อาสวกฺขยํ, วิมุตฺติยา วา วิมุตฺตีติ, นตฺถิ.\csromanspacer{} หญฺจิ นตฺถิ กิญฺจิ นานากรณํ ตถาคตสฺส วา สาวกสฺส วา อาสวกฺขเยน วา อาสวกฺขยํ, วิมุตฺติยา วา วิมุตฺติ, เตน วต เร วตฺตพฺเพ “อาสวานํ ขเย ยถาภูตํ ญาณํ ตถาคตพลํ สาวกสาธารณนฺ”ติ.}

\proseitem{356}{\csromansymbolmark{+}อาสวานํ ขเย ยถาภูตํ ญาณํ ตถาคตพลํ สาวกสาธารณนฺติ, อามนฺตา.\csromanspacer{} ฐานาฐาเน ยถาภูตํ ญาณํ ตถาคตพลํ สาวกสาธารณนฺติ, น เหวํ วตฺตพฺเพ ฯเปฯ.}

\prose{\csromansymbolmark{+}อาสวานํ ขเย ยถาภูตํ ญาณํ ตถาคตพลํ สาวกสาธารณนฺติ, อามนฺตา.\csromanspacer{} สตฺตานํ จุตูปปาเต ยถาภูตํ ญาณํ ตถาคตพลํ สาวกสาธารณนฺติ, น เหวํ วตฺตพฺเพ ฯเปฯ.}

\prose{\csromansymbolmark{+}ฐานาฐาเน ยถาภูตํ ญาณํ ตถาคตพลํ สาวก- อสาธารณนฺติ, อามนฺตา.\csromanspacer{} อาสวานํ ขเย ยถาภูตํ ญาณํ ตถาคตพลํ สาวก-อสาธารณนฺติ, น เหวํ วตฺตพฺเพ ฯเปฯ.}

\chapterheadpageread{3. ตติยวคฺค}
\csromanmatikaanchor{mtk.43}
\csromantocmarkref{subsection}{(22) 2. อริยนฺติกถา}{mtk.43}
\bookmark[dest={mtk.43},level=2]{(22) 2. อริยนฺติกถา}
\prose{\csromanfolio{177}\csromansymbolmark{+}สตฺตานํ จุตูปปาเต ยถาภูตํ ญาณํ ตถาคตพลํ สาวก- อสาธารณนฺติ, อามนฺตา.\csromanspacer{} อาสวานํ ขเย ยถาภูตํ ญาณํ ตถาคตพลํ สาวก-อสาธารณนฺติ, น เหวํ วตฺตพฺเพ ฯเปฯ.}

\prose{\csromansymbolmark{+}อินฺทฺริยปโรปริยตฺตํ ยถาภูตํ ญาณํ ตถาคตพลํ สาวก- อสาธารณนฺติ, อามนฺตา.\csromanspacer{} ฐานาฐาเน ยถาภูตํ ญาณํ ตถาคตพลํ สาวก-อสาธารณนฺติ, น เหวํ วตฺตพฺเพ\footnote{อยเมตฺถ สาธารณปกฺขํ สนฺธาย ปฏีกฺเขโป (ฏีกา โอโลเกตพฺพา)}.ป-.}

\prose{\csromansymbolmark{+}อินฺทฺริยปโรปริยตฺตํ ยถาภูตํ ญาณํ ตถาคตพลํ สาวก- อสาธารณนฺติ, อามนฺตา ฯเปฯ.\csromanspacer{} อาสวานํ ขเย ยถาภูตํ ญาณํ ตถาคตพลํ สาวก-อสาธารณนฺติ, น เหวํ วตฺตพฺเพ ฯเปฯ.}

\prose{\csromansymbolmark{+}ฐานาฐาเน ยถาภูตํ ญาณํ ตถาคตพลํ สาวกสาธารณนฺติ, อามนฺตา.\csromanspacer{} อินฺทฺริยปโรปริยตฺตํ ยถาภูตํ ญาณํ ตถาคตพลํ สาวกสาธารณนฺติ, น เหวํ วตฺตพฺเพ ฯเปฯ.}

\prose{\csromansymbolmark{+}อาสวานํ ขเย ยถาภูตํ ญาณํ ตถาคตพลํ สาวกสาธารณนฺติ, อามนฺตา.\csromanspacer{} อินฺทฺริยปโรปริยตฺตํ ยถาภูตํ ญาณํ ตถาคตพลํ สาวกสาธารณนฺติ, น เหวํ วตฺตพฺเพ ฯเปฯ.}

\nitthitam{พลกถา นิฏฺฐิตา.}
\chapterhead{3. ตติยวคฺค}
\vspace{\tipitakaparskip}
\csromancenter{\textbf{(22) 2.}\csromanspacer{}\textbf{ อริยนฺติกถา}}
\proseitem{357}{\csromansymbolmark{*}ฐานาฐาเน ยถาภูตํ ญาณํ ตถาคตพลํ อริยนฺติ, อามนฺตา.\csromanspacer{} มคฺโค ผลํ นิพฺพานํ โสตาปตฺติมคฺโค โสตาปตฺติผลํ, สกทาคามิมคฺโค สกทาคามิผลํ, อนาคามิมคฺโค อนาคามิผลํ, อรหตฺตมคฺโค อรหตฺตผลํ, สติปฏฺฐานํ สมฺมปฺปธานํ อิทฺธิปาโท อินฺทฺริยํ พลํ โพชฺฌงฺโคติ, น เหวํ วตฺตพฺเพ ฯเปฯ.}

\prose{\csromansymbolmark{*}ฐานาฐาเน ยถาภูตํ ญาณํ ตถาคตพลํ อริยนฺติ, อามนฺตา.\csromanspacer{} สุญฺญตารมฺมณนฺติ, น เหวํ วตฺตพฺเพ ฯเปฯ.\csromanspacer{} สุญฺญตารมฺมณนฺติ, อามนฺตา. \csromanfolio{178}ฐานาฐานญฺจ มนสิ กโรติ, สุญฺญตญฺจ มนสิ กโรตีติ, น เหวํ วตฺตพฺเพ ฯเปฯ. \csromansymbolmark{*}ฐานาฐานญฺจ มนสิ กโรติ, สุญฺญตญฺจ มนสิ กโรตีติ, อามนฺตา.\csromanspacer{} ทฺวินฺนํ ผสฺสานํ ทฺวินฺนํ จิตฺตานํ สโมธานํ โหตีติ, น เหวํ วตฺตพฺเพ ฯเปฯ.\csromanspacer{} ฐานาฐาเน ยถาภูตํ ญาณํ ตถาคตพลํ อริยนฺติ, อามนฺตา.\csromanspacer{} อนิมิตฺตารมฺมณํ ฯเปฯ.\csromanspacer{} อปฺปณิหิตารมฺมณนฺติ, น เหวํ วตฺตพฺเพ ฯเปฯ.\csromanspacer{} อปฺปณิหิตารมฺมณนฺติ, อามนฺตา.\csromanspacer{} ฐานาฐานญฺจ มนสิ กโรติ, อปฺปณิหิตญฺจ มนสิ กโรตีติ, น เหวํ วตฺตพฺเพ ฯเปฯ. \csromansymbolmark{*}ฐานาฐานญฺจ มนสิ กโรติ, อปฺปณิหิตญฺจ มนสิ กโรตีติ, อามนฺตา.\csromanspacer{} ทฺวินฺนํ ผสฺสานํ ทฺวินฺนํ จิตฺตานํ สโมธานํ โหตีติ, น เหวํ วตฺตพฺเพ ฯเปฯ.}

\proseitem{358}{\csromansymbolmark{*}สติปฏฺฐานา อริยา สุญฺญตารมฺมณาติ, อามนฺตา.\csromanspacer{} ฐานาฐาเน ยถาภูตํ ญาณํ ตถาคตพลํ อริยํ สุญฺญาตารมฺมณนฺติ, น เหวํ วตฺตพฺเพ ฯเปฯ.}

\prose{\csromansymbolmark{*}สติปฏฺฐานา อริยา อนิมิตฺตารมฺมณา.\csromanspacer{} อปฺปณิหิตารมฺมณาติ, อามนฺตา.\csromanspacer{} ฐานาฐาเน ยถาภูตํ ญาณํ ตถาคตพลํ อริยํ อปฺปณิหิตารมฺมณนฺติ, น เหวํ วตฺตพฺเพ ฯเปฯ.}

\prose{\csromansymbolmark{*}สมฺมปฺปธานา.\csromanspacer{} อิทฺธิปาทา.\csromanspacer{} อินฺทฺริยา.\csromanspacer{} พลา.\csromanspacer{} โพชฺฌงฺคา อริยา สุญฺญตารมฺมณาติ, อามนฺตา.\csromanspacer{} ฐานาฐาเน ยถาภูตํ ญาณํ ตถาคตพลํ อริยํ สุญฺญตารมฺมณนฺติ, น เหวํ วตฺตพฺเพ ฯเปฯ.}

\prose{\csromansymbolmark{*}โพชฺฌงฺคา อริยา อนิมิตฺตารมฺมณา อปฺปณิหิตารมฺมณาติ, อามนฺตา.\csromanspacer{} ฐานาฐาเน ยถาภูตํ ญาณํ ตถาคตพลํ อริยํ อปฺปณิหิตารมฺมณนฺติ, น เหวํ วตฺตพฺเพ ฯเปฯ.}

\proseitem{359}{\csromansymbolmark{*}ฐานาฐาเน ยถาภูตํ ญาณํ ตถาคตพลํ อริยํ น วตฺตพฺพํ “สุญฺญตารมฺมณนฺ”ติ, อามนฺตา.\csromanspacer{} สติปฏฺฐานา อริยา น วตฺตพฺพา “สุญฺญตารมฺมณา”ติ, น เหวํ วตฺตพฺเพ ฯเปฯ.}

\prose{\csromansymbolmark{*}ฐานาฐาเน ยถาภูตํ ญาณํ ตถาคตพลํ อริยํ น วตฺตพฺพํ “อนิมิตฺตารมฺมณํ” ฯเปฯ “อปฺปณิหิตารมฺมณนฺ”ติ, อามนฺตา.\csromanspacer{} สติปฏฺฐานา อริยา น วตฺตพฺพา “อปฺปณิหิตารมฺมณา”ติ, น เหวํ วตฺตพฺเพ ฯเปฯ.}

\chapterheadpageread{3. ตติยวคฺค}
\prose{\csromanfolio{179}\csromansymbolmark{*}ฐานาฐาเน ยถาภูตํ ญาณํ ตถาคตพลํ อริยํ น วตฺตพฺพํ “สุญฺญตารมฺมณํ” ฯเปฯ “อนิมิตฺตารมฺมณํ” ฯเปฯ “อปฺปณิหิตารมฺมณนฺ”ติ, อามนฺตา.\csromanspacer{} สมฺมปฺปธานํ ฯเปฯ.\csromanspacer{} โพชฺฌงฺคา อริยา น วตฺตพฺพา อปฺปณิหิตารมฺมณาติ, น เหวํ วตฺตพฺเพ ฯเปฯ.}

\proseitem{360}{\csromansymbolmark{*}สตฺตานํ จุตูปปาเต ยถาภูตํ ญาณํ ตถาคตพลํ อริยนฺติ, อามนฺตา.\csromanspacer{} มคฺโค ผลํ นิพฺพานํ โสตาปตฺติมคฺโค โสตาปตฺติผลํ ฯเปฯ โพชฺฌงฺโคติ, น เหวํ วตฺตพฺเพ ฯเปฯ.}

\prose{\csromansymbolmark{*}สตฺตานํ จุตูปปาเต ยถาภูตํ ญาณํ ตถาคตพลํ อริยนฺติ, อามนฺตา.\csromanspacer{} สุญฺญตารมฺมณนฺติ, น เหวํ วตฺตพฺเพ ฯเปฯ.\csromanspacer{} สุญฺญตารมฺมณนฺติ, อามนฺตา.\csromanspacer{} สตฺตานํ จุตูปปาตญฺจ มนสิ กโรติ, สุญฺญตญฺจ มนสิ กโรตีติ, น เหวํ วตฺตพฺเพ ฯเปฯ.}

\prose{\csromansymbolmark{*}สตฺตานํ จุตูปปาตญฺจ มนสิ กโรติ สุญฺญตญฺจ มนสิ กโรตีติ, อามนฺตา.\csromanspacer{} ทฺวินฺนํ ผสฺสานํ ทฺวินฺนํ จิตฺตานํ สโมธานํ โหตีติ, น เหวํ วตฺตพฺเพ ฯเปฯ.}

\prose{\csromansymbolmark{*}สตฺตานํ จุตูปปาเต ยถาภูตํ ญาณํ ตถาคตพลํ อริยนฺติ, อามนฺตา.\csromanspacer{} อนิมิตฺตารมฺมณํ.\csromanspacer{} อปฺปณิหิตารมฺมณนฺติ, น เหวํ วตฺตพฺเพ ฯเปฯ.\csromanspacer{} อปฺปณิหิตารมฺมณนฺติ, อามนฺตา.\csromanspacer{} สตฺตานํ จุตูปปาตญฺจ มนสิ กโรติ, อปฺปณิหิตญฺจ มนสิ กโรตีติ, น เหวํ วตฺตพฺเพ ฯเปฯ.}

\prose{\csromansymbolmark{*}สตฺตานํ จุตูปปาตญฺจ มนสิ กโรติ, อปฺปณิหิตญฺจ มนสิ กโรตีติ, อามนฺตา.\csromanspacer{} ทฺวินฺนํ ผสฺสานํ ทฺวินฺนํ จิตฺตานํ สโมธานํ โหตีติ, น เหวํ วตฺตพฺเพ ฯเปฯ.}

\proseitem{361}{\csromansymbolmark{*}สติปฏฺฐานา อริยา สุญฺญตารมฺมณา ฯเปฯ อนิมิตฺตารมฺมณา ฯเปฯ อปฺปณิหิตารมฺมณาติ, อามนฺตา.\csromanspacer{} สตฺตานํ จุตูปปาเต ยถาภูตํ ญาณํ ตถาคตพลํ อริยํ อปฺปณิหิตารมฺมณนฺติ, น เหวํ วตฺตพฺเพ ฯเปฯ.}

\prose{\csromansymbolmark{*}สมฺมปฺปธานํ ฯเปฯ โพชฺฌงฺคา อริยา สุญฺญตารมฺมณา ฯเปฯ อนิมิตฺตารมฺมณา ฯเปฯ อปฺปณิหิตารมฺมณาติ, อามนฺตา.\csromanspacer{} สตฺตานํ จุตูปปาเต ยถาภูตํ ญาณํ ตถาคตพลํ อริยํ อปฺปณิหิตารมฺมณนฺติ, น เหวํ วตฺตพฺเพ ฯเปฯ.}

\prose{\csromanfolio{180}\csromansymbolmark{*}สตฺตานํ จุตูปปาเต ยถาภูตํ ญาณํ ตถาคตพลํ อริยํ น วตฺตพฺพํ “สุญฺญตารมฺมณํ” ฯเปฯ “อนิมิตฺตารมฺมณํ” ฯเปฯ “อปฺปณิหิตารมฺมณนฺ”ติ, อามนฺตา.\csromanspacer{} สติปฏฺฐานา อริยา น วตฺตพฺพา “อปฺปณิหิตารมฺมณา”ติ, น เหวํ วตฺตพฺเพ ฯเปฯ.}

\proseitem{362}{\csromansymbolmark{*}สตฺตานํ จุตูปปาเต ยถาภูตํ ญาณํ ตถาคตพลํ อริยํ น วตฺตพฺพํ “สุญฺญตารมฺมณํ” ฯเปฯ “อนิมิตฺตารมฺมณํ ฯเปฯ “อปฺปณิหิตารมฺมณนฺ”ติ, อามนฺตา.\csromanspacer{} สมฺมปฺปธานา ฯเปฯ.\csromanspacer{} โพชฺฌงฺคา อริยา น วตฺตพฺพา “อปฺปณิหิตารมฺมณา”ติ, น เหวํ วตฺตพฺเพ ฯเปฯ.}

\prose{\csromansymbolmark{+}อาสวานํ ขเย ยถาภูตํ ญาณํ ตถาคตพลํ อริยนฺติ, อามนฺตา.\csromanspacer{} ฐานาฐาเน ยถาภูตํ ญาณํ ตถาคตพลํ อริยนฺติ, น เหวํ วตฺตพฺเพ ฯเปฯ.}

\prose{\csromansymbolmark{+}อาสวานํ ขเย ยถาภูตํ ญาณํ ตถาคตพลํ อริยนฺติ, อามนฺตา.\csromanspacer{} สตฺตานํ จุตูปปาเต ยถาภูตํ ญาณํ ตถาคตพลํ อริยนฺติ, น เหวํ วตฺตพฺเพ ฯเปฯ.}

\prose{\csromansymbolmark{+}ฐานาฐาเน ยถาภูตํ ญาณํ ตถาคตพลํ น วตฺตพฺพํ “อริยนฺ”ติ, อามนฺตา.\csromanspacer{} อาสวานํ ขเย ยถาภูตํ ญาณํ ตถาคตพลํ น วตฺตพฺพํ “อริยนฺ”ติ, น เหวํ วตฺตพฺเพ ฯเปฯ.}

\prose{\csromansymbolmark{+}สตฺตานํ จุตูปปาเต ยถาภูตํ ญาณํ ตถาคตพลํ น วตฺตพฺพํ “อริยนฺ”ติ, อามนฺตา.\csromanspacer{} อาสวานํ ขเย ยถาภูตํ ญาณํ ตถาคตพลํ น วตฺตพฺพํ “อริยนฺ”ติ, น เหวํ วตฺตพฺเพ ฯเปฯ.}

\prose{\csromansymbolmark{+}อาสวานํ ขเย ยถาภูตํ ญาณํ ตถาคตพลํ อริยํ สุญฺญตารมฺมณนฺติ, อามนฺตา.\csromanspacer{} ฐานาฐาเน ยถาภูตํ ญาณํ ตถาคตพลํ อริยํ สุญฺญตารมฺมณนฺติ, น เหวํ วตฺตพฺเพ ฯเปฯ.}

\prose{\csromansymbolmark{+}อาสวานํ ขเย ยถาภูตํ ญาณํ ตถาคตพลํ อริยํ อนิมิตฺตารมฺมณํ.\csromanspacer{} อปฺปณิหิตารมฺมณนฺติ, อามนฺตา.\csromanspacer{} ฐานาฐาเน ยถาภูตํ ญาณํ ตถาคตพลํ อริยํ อปฺปณิหิตารมฺมณนฺติ, น เหวํ วตฺตพฺเพ ฯเปฯ.}

\prose{\csromansymbolmark{+}อาสวานํ ขเย ยถาภูตํ ญาณํ ตถาคตพลํ อริยํ สุญฺญตารมฺมณํ ฯเปฯ อนิมิตฺตารมฺมณํ ฯเปฯ อปฺปณิหิตารมฺมณนฺติ, อามนฺตา ฯเปฯ.\csromanspacer{} สตฺตานํ}

\vspace{\tipitakaparskip}
\csromancenter{ตติยวคฺค จุตูปปาเต ยถาภูตํ ญาณํ ตถาคตพลํ อริยํ อปฺปณิหิตารมฺมณนฺติ, น เหวํ วตฺตพฺเพ ฯเปฯ.}
\csromanmatikaanchor{mtk.44}
\csromantocmarkref{subsection}{(23) 3. วิมุตฺติกถา}{mtk.44}
\bookmark[dest={mtk.44},level=2]{(23) 3. วิมุตฺติกถา}
\prose{\csromanfolio{181}\csromansymbolmark{+}ฐานาฐาเน ยถาภูตํ ญาณํ ตถาคตพลํ อริยํ น วตฺตพฺพํ “สุญฺญตารมฺมณนฺ”ติ, อามนฺตา.\csromanspacer{} อาสวานํ ขเย ยถาภูตํ ญาณํ ตถาคตพลํ อริยํ น วตฺตพฺพํ “สุญฺญตารมฺมณนฺ”ติ, น เหวํ วตฺตพฺเพ ฯเปฯ.}

\prose{\csromansymbolmark{+}ฐานาฐาเน ยถาภูตํ ญาณํ ตถาคตพลํ อริยํ น วตฺตพฺพํ “อนิมิตฺตารมฺมณํ.\csromanspacer{} อปฺปณิหิตารมฺมณนฺ”ติ, อามนฺตา.\csromanspacer{} อาสวานํ ขเย ยถาภูตํ ญาณํ ตถาคตพลํ อริยํ น วตฺตพฺพํ “อปฺปณิหิตารมฺมณนฺ”ติ, น เหวํ วตฺตพฺเพ ฯเปฯ.}

\prose{\csromansymbolmark{+}สตฺตานํ จุตูปปาเต ยถาภูตํ ญาณํ ตถาคตพลํ อริยํ น วตฺตพฺพํ “สุญฺญตารมฺมณํ” ฯเปฯ “อนิมิตฺตารมฺมณํ” ฯเปฯ “อปฺปณิหิตารมฺมณนฺ”ติ, อามนฺตา.\csromanspacer{} อาสวานํ ขเย ยถาภูตํ ญาณํ ตถาคตพลํ อริยํ น วตฺตพฺพํ “อปฺปณิหิตารมฺมณนฺ”ติ, น เหวํ วตฺตพฺเพ ฯเปฯ.}

\nitthitam{อริยนฺติกถา นิฏฺฐิตา.}
\chapterhead{3. ตติยวคฺค}
\vspace{\tipitakaparskip}
\csromancenter{\textbf{(23) 3.}\csromanspacer{}\textbf{ วิมุตฺติกถา}}
\proseitem{363}{\csromansymbolmark{*}สราคํ จิตฺตํ วิมุจฺจตีติ, อามนฺตา.\csromanspacer{} ราคสหคตํ ราคสหชาตํ ราคสํสฏฺฐํ ราคสมฺปยุตฺตํ ราคสหภุ ราคานุปริวตฺติ อกุสลํ โลกิยํ สาสวํ สํโยชนิยํ คนฺถนิยํ โอฆนิยํ โยคนิยํ นีวรณิยํ ปรามฏฺฐํ อุปาทานิยํ สํกิเลสิยํ จิตฺตํ วิมุจฺจตีติ, น เหวํ วตฺตพฺเพ ฯเปฯ.}

\prose{\csromansymbolmark{*}สผสฺสํ จิตฺตํ วิมุจฺจติ, ผสฺโส จ จิตฺตญฺจ อุโภ วิมุจฺจนฺตีติ, อามนฺตา.\csromanspacer{} สราคํ จิตฺตํ วิมุจฺจติ, ราโค จ จิตฺตญฺจ อุโภ วิมุจฺจนฺตีติ, น เหวํ วตฺตพฺเพ ฯเปฯ.}

\prose{\csromansymbolmark{*}สเวทนํ ฯเปฯ.\csromanspacer{} สสญฺญํ ฯเปฯ.\csromanspacer{} สเจตนํ ฯเปฯ.\csromanspacer{} สปญฺญํ\footnote{สสญฺญํ (สี, ก)} จิตฺตํ วิมุจฺจติ, ปญฺญา จ จิตฺตญฺจ อุโภ วิมุจฺจนฺตีติ, อามนฺตา.\csromanspacer{} สราคํ จิตฺตํ วิมุจฺจติ, ราโค จ จิตฺตญฺจ อุโภ วิมุจฺจนฺตีติ, น เหวํ วตฺตพฺเพ ฯเปฯ.}

\prose{\csromanfolio{182}\csromansymbolmark{*}สผสฺสํ สราคํ จิตฺตํ วิมุจฺจติ, ผสฺโส จ จิตฺตญฺจ อุโภ วิมุจฺจนฺตีติ, อามนฺตา.\csromanspacer{} ราโค จ จิตฺตญฺจ อุโภ วิมุจฺจนฺตีติ, น เหวํ วตฺตพฺเพ ฯเปฯ.}

\prose{\csromansymbolmark{*}สเวทนํ สราคํ ฯเปฯ.\csromanspacer{} สสญฺญํ สราคํ ฯเปฯ.\csromanspacer{} สเจตนํ สราคํ ฯเปฯ.\csromanspacer{} สปญฺญํ สราคํ จิตฺตํ วิมุจฺจติ, ปญฺญา จ จิตฺตญฺจ อุโภ วิมุจฺจนฺตีติ, อามนฺตา.\csromanspacer{} ราโค จ จิตฺตญฺจ อุโภ วิมุจฺจนฺตีติ, น เหวํ วตฺตพฺเพ ฯเปฯ.}

\proseitem{364}{\csromansymbolmark{*}สโทสํ จิตฺตํ วิมุจฺจตีติ, อามนฺตา.\csromanspacer{} โทสสหคตํ โทสสหชาตํ โทสสํสฏฺฐํ โทสสมฺปยุตฺตํ โทสสหภุ โทสานุปริวตฺติ อกุสลํ โลกิยํ สาสวํ ฯเปฯ สํกิเลสิยํ จิตฺตํ วิมุจฺจตีติ, น เหวํ วตฺตพฺเพ ฯเปฯ.}

\prose{\csromansymbolmark{*}สผสฺสํ จิตฺตํ วิมุจฺจติ, ผสฺโส จ จิตฺตญฺจ อุโภ วิมุจฺจนฺตีติ, อามนฺตา.\csromanspacer{} สโทสํ จิตฺตํ วิมุจฺจติ, โทโส จ จิตฺตญฺจ อุโภ วิมุจฺจนฺตีติ, น เหวํ วตฺตพฺเพ ฯเปฯ.}

\prose{\csromansymbolmark{*}สเวทนํ ฯเปฯ.\csromanspacer{} สสญฺญํ ฯเปฯ.\csromanspacer{} สเจตนํ ฯเปฯ.\csromanspacer{} สปญฺญํ จิตฺตํ วิมุจฺจติ, ปญฺญา จ จิตฺตญฺจ อุโภ วิมุจฺจนฺตีติ, อามนฺตา.\csromanspacer{} สโทสํ จิตฺตํ วิมุจฺจติ, โทโส จ จิตฺตญฺจ อุโภ วิมุจฺจนฺตีติ น เหวํ วตฺตพฺเพ ฯเปฯ.}

\prose{\csromansymbolmark{*}สผสฺสํ สโทสํ จิตฺตํ วิมุจฺจติ, ผสฺโส จ จิตฺตญฺจ อุโภ วิมุจฺจนฺตีติ, อามนฺตา.\csromanspacer{} สโทสํ จิตฺตํ วิมุจฺจติ, โทโส จ จิตฺตญฺจ อุโภ วิมุจฺจนฺตีติ, น เหวํ วตฺตพฺเพ ฯเปฯ.}

\prose{\csromansymbolmark{*}สเวทนํ สโทสํ.\csromanspacer{} สสญฺญํ สโทสํ.\csromanspacer{} สเจตนํ สโทสํ.\csromanspacer{} สปญฺญํ สโทสํ จิตฺตํ วิมุจฺจติ, ปญฺญา จ จิตฺตญฺจ อุโภ วิมุจฺจนฺตีติ, อามนฺตา.\csromanspacer{} โทโส จ จิตฺตญฺจ อุโภ วิมุจฺจนฺตีติ, น เหวํ วตฺตพฺเพ ฯเปฯ.}

\proseitem{365}{\csromansymbolmark{*}สโมหํ จิตฺตํ วิมุจฺจตีติ, อามนฺตา.\csromanspacer{} โมหสหคตํ โมหสหชาตํ โมหสํสฏฺฐํ โมหสมฺปยุตฺตํ โมหสหภุ โมหานุปริวตฺติ อกุสลํ โลกิยํ สาสวํ ฯเปฯ สํกิเลสิยํ จิตฺตํ วิมุจฺจตีติ, น เหวํ วตฺตพฺเพ ฯเปฯ.}

\prose{\csromansymbolmark{*}สผสฺสํ จิตฺตํ วิมุจฺจติ, ผสฺโส จ จิตฺตญฺจ อุโภ วิมุจฺจนฺตีติ, อามนฺตา.\csromanspacer{} สโมหํ จิตฺตํ วิมุจฺจติ, โมโห จ จิตฺตญฺจ อุโภ วิมุจฺจนฺตีติ, น เหวํ วตฺตพฺเพ ฯเปฯ.}

\chapterheadpageread{3. ตติยวคฺค}
\csromanmatikaanchor{mtk.45}
\csromantocmarkref{subsection}{(24) 4. วิมุจฺจมานกถา}{mtk.45}
\bookmark[dest={mtk.45},level=2]{(24) 4. วิมุจฺจมานกถา}
\prose{\csromanfolio{183}\csromansymbolmark{*}สเวทนํ.\csromanspacer{} สสญฺญํ.\csromanspacer{} สเจตนํ.\csromanspacer{} สปญฺญํ จิตฺตํ วิมุจฺจติ, ปญฺญา จ จิตฺตญฺจ อุโภ วิมุจฺจนฺตีติ, อามนฺตา.\csromanspacer{} สโมหํ จิตฺตํ วิมุจฺจติ โมโห จ จิตฺตญฺจ อุโภ วิมุจฺจนฺตีติ, น เหวํ วตฺตพฺเพ ฯเปฯ.}

\prose{\csromansymbolmark{*}สผสฺสํ สโมหํ จิตฺตํ วิมุจฺจติ, ผสฺโส จ จิตฺตญฺจ อุโภ วิมุจฺจนฺตีติ, อามนฺตา.\csromanspacer{} โมโห จ จิตฺตญฺจ อุโภ วิมุจฺจนฺตีติ, น เหวํ วตฺตพฺเพ ฯเปฯ.}

\prose{\csromansymbolmark{*}สเวทนํ สโมหํ.\csromanspacer{} สสญฺญํ สโมหํ.\csromanspacer{} สเจตนํ สโมหํ ฯเปฯ.\csromanspacer{} สปญฺญํ สโมหํ จิตฺตํ วิมุจฺจติ, ปญฺญา จ จิตฺตญฺจ อุโภ วิมุจฺจนฺตีติ, อามนฺตา.\csromanspacer{} โมโห จ จิตฺตญฺจ อุโภ วิมุจฺจนฺตีติ, น เหวํ วตฺตพฺเพ ฯเปฯ.}

\prose{\csromansymbolmark{*}สราคํ สโทสํ สโมหํ จิตฺตํ วิมุจฺจตีติ, อามนฺตา.\csromanspacer{} วีตราคํ วิตโทสํ วีตโมหํ นิกฺกิเลสํ จิตฺตํ วิมุจฺจตีติ, น เหวํ วตฺตพฺเพ ฯเปฯ.\csromanspacer{} เตน หิ น วตฺตพฺพํ “สราคํ สโทสํ สโมหํ จิตฺตํ วิมุจฺจตี”ติ.}

\nitthitam{วิมุตฺติกถา นิฏฺฐิตา.}
\chapterhead{3. ตติยวคฺค}
\vspace{\tipitakaparskip}
\csromancenter{\textbf{(24) 4.}\csromanspacer{}\textbf{ วิมุจฺจมานกถา}}
\proseitem{366}{\csromansymbolmark{*}วิมุตฺตํ วิมุจฺจมานนฺติ, อามนฺตา.\csromanspacer{} เอกเทสํ วิมุตฺตํ, เอกเทสํ อวิมุตฺตนฺติ, น เหวํ วตฺตพฺเพ ฯเปฯ.}

\prose{\csromansymbolmark{*}เอกเทสํ วิมุตฺตํ, เอกเทสํ อวิมุตฺตนฺติ, อามนฺตา.\csromanspacer{} เอกเทสํ โสตาปนฺโน, เอกเทสํ น โสตาปนฺโน, เอกเทสํ โสตาปตฺติผลปฺปตฺโต ปฏิลทฺโธ อธิคโต สจฺฉิกโต อุปสมฺปชฺช วิหรติ, กาเยน ผุสิตฺวา วิหรติ, เอกเทสํ น กาเยน ผุสิตฺวา วิหรติ, เอกเทสํ สตฺตกฺขตฺตุปรโม โกลํโกโล เอกพีชี พุทฺเธ อเวจฺจปฺปสาเทน สมนฺนาคโต ธมฺเม ฯเปฯ สํเฆ ฯเปฯ อริยกนฺเตหิ สีเลหิ สมนฺนาคโต, เอกเทสํ อริยกนฺเตหิ สีเลหิ น สมนฺนาคโตติ, น เหวํ วตฺตพฺเพ ฯเปฯ.}

\prose{\csromansymbolmark{*}เอกเทสํ วิมุตฺตํ, เอกเทสํ อวิมุตฺตนฺติ, อามนฺตา.\csromanspacer{} เอกเทสํ สกทาคามี, เอกเทสํ น สกทาคามี, เอกเทสํ สกทาคามิผลปฺปตฺโต ปฏิลทฺโธ อธิคโต สจฺฉิกโต อุปสมฺปชฺช วิหรติ, กาเยน ผุสิตฺวา วิหรติ, เอกเทสํ น กาเยน ผุสิตฺวา วิหรตีติ, น เหวํ วตฺตพฺเพ ฯเปฯ.}

\prose{\csromanfolio{184}\csromansymbolmark{*}เอกเทสํ วิมุตฺตํ, เอกเทสํ อวิมุตฺตนฺติ, อามนฺตา.\csromanspacer{} เอกเทสํ อนาคามี, เอกเทสํ น อนาคามี, เอกเทสํ อนาคามิผลปฺปตฺโต ปฏิลทฺโธ อธิคโต สจฺฉิกโต อุปสมฺปชฺช วิหรติ, กาเยน ผุสิตฺวา วิหรติ, เอกเทสํ น กาเยน ผุสิตฺวา วิหรติ, เอกเทสํ อนฺตราปรินิพฺพายี, อุปหจฺจปรินิพฺพายี, อสงฺขารปรินิพฺพายี, สสงฺขารปรินิพฺพายี, อุทฺธํโสโต-อกนิฏฺฐคามี, เอกเทสํ น อุทฺธํโสโต-อกนิฏฺฐคามีติ, น เหวํ วตฺตพฺเพ ฯเปฯ.}

\prose{\csromansymbolmark{*}เอกเทสํ วิมุตฺตํ, เอกเทสํ อวิมุตฺตนฺติ, อามนฺตา.\csromanspacer{} เอกเทสํ อรหา, เอกเทสํ น อรหา, เอกเทสํ อรหตฺตปฺปตฺโต ปฏิลทฺโธ อธิคโต สจฺฉิกโต อุปสมฺปชฺช วิหรติ, กาเยน ผุสิตฺวา วิหรติ, เอกเทสํ น กาเยน ผุสิตฺวา วิหรติ, เอกเทสํ วีตราโค วีตโทโส วีตโมโห ฯเปฯ เอกเทสํ สจฺฉิกาตพฺพํ สจฺฉิกตํ, เอกเทสํ สจฺฉิกาตพฺพํ น สจฺฉิกตนฺติ, น เหวํ วตฺตพฺเพ ฯเปฯ.}

\prose{\csromansymbolmark{*}วิมุตฺตํ วิมุจฺจมานนฺติ, อามนฺตา.\csromanspacer{} อุปฺปาทกฺขเณ วิมุตฺตํ, ภงฺคกฺขเณ วิมุจฺจมานนฺติ, น เหวํ วตฺตพฺเพ ฯเปฯ.}

\proseitem{367}{\csromansymbolmark{+}น วตฺตพฺพํ “วิมุตฺตํ วิมุจฺจมานนฺ”ติ, อามนฺตา.\csromanspacer{} นนุ วุตฺตํ ภควตา “ตสฺส เอวํ ชานโต เอวํ ปสฺสโต กามาสวาปิ จิตฺตํ วิมุจฺจติ, ภวาสวาปิ จิตฺตํ วิมุจฺจติ, อวิชฺชาสวาปิ จิตฺตํ วิมุจฺจตี”ติ1 อตฺเถว สุตฺตนฺโตติ, อามนฺตา.\csromanspacer{} เตน หิ วตฺตพฺพํ\footnote{เตน หิ (สี, สฺยา, กํ), เตน หิ น วตฺตพฺพํ (ก)} “วิมุตฺตํ วิมุจฺจมานนฺ”ติ.}

\prose{\csromansymbolmark{*}วิมุตฺตํ วิมุจฺจมานนฺติ, อามนฺตา.\csromanspacer{} นนุ วุตฺตํ ภควตา “โส เอวํ สมาหิเต จิตฺเต ปริสุทฺเธ ปริโยทาเต อนงฺคเณ วิคตูปกฺกิเลเส มุทุภูเต กมฺมนิเย ฐิเต อาเนญฺชปฺปตฺเต อาสวานํ ขยญาณาย จิตฺตํ อภินินฺนาเมตี”ติ3 อตฺเถว สุตฺตนฺโตติ, อามนฺตา.\csromanspacer{} เตน หิ น วตฺตพฺพํ “วิมุตฺตํ วิมุจฺจมานนฺ”ติ.}

\prose{\csromansymbolmark{*}อตฺถิ จิตฺตํ วิมุจฺจมานนฺติ, อามนฺตา.\csromanspacer{} อตฺถิ จิตฺตํ รชฺชมานํ ทุสฺสมานํ มุยฺหมานํ กิลิสฺสมานนฺติ, น เหวํ วตฺตพฺเพ ฯเปฯ.\csromanspacer{} นนุ รตฺตญฺเจว อรตฺตญฺจ, ทุฏฺฐญฺเจว อทุฏฺฐญฺจ, มูฬฺหญฺเจว อมูฬฺหญฺจ, ฉินฺนญฺเจว อฉินฺนญฺจ, ภินฺนญฺเจว อภินฺนญฺจ, กตญฺเจว}

\vspace{\tipitakaparskip}
\csromancenter{ตติยวคฺค อกตญฺจาติ, อามนฺตา.\csromanspacer{} หญฺจิ รตฺตญฺเจว อรตฺตญฺจ, ทุฏฺฐญฺเจว อทุฏฺฐญฺจ, มูฬฺหญฺเจว อมูฬฺหญฺจ, ฉินฺนญฺเจว อฉินฺนญฺจ, ภินฺนญฺเจว อภินฺนญฺจ, กตญฺเจว อกตญฺจ, โน จ วต เร วตฺตพฺเพ “อตฺถิ จิตฺตํ วิมุจฺจมานนฺ”ติ.}
\nitthitam{วิมุจฺจมานกถา นิฏฺฐิตา.}
\chapterhead{3. ตติยวคฺค}
\vspace{\tipitakaparskip}
\csromancenter{\textbf{(25) 5.}\csromanspacer{}\textbf{ อฏฺฐมกกถา}}
\csromanmatikaanchor{mtk.46}
\csromantocmarkref{subsection}{(25) 5. อฏฺฐมกกถา}{mtk.46}
\bookmark[dest={mtk.46},level=2]{(25) 5. อฏฺฐมกกถา}
\proseitem{368}{\csromanfolio{185}\csromansymbolmark{*}อฏฺฐมกสฺส ปุคฺคลสฺส ทิฏฺฐิปริยุฏฺฐานํ ปหีนนฺติ, อามนฺตา.\csromanspacer{} อฏฺฐมโก ปุคฺคโล โสตาปนฺโน โสตาปตฺติผลปฺปตฺโต ปฏิลทฺโธ อธิคโต สจฺฉิกโต อุปสมฺปชฺช วิหรติ, กาเยน ผุสิตฺวา วิหรตีติ, น เหวํ วตฺตพฺเพ ฯเปฯ.}

\prose{\csromansymbolmark{*}อฏฺฐมกสฺส ปุคฺคลสฺส วิจิกิจฺฉาปริยุฏฺฐานํ ปหีนนฺติ, อามนฺตา.\csromanspacer{} อฏฺฐมโก ปุคฺคโล โสตาปนฺโน โสตาปตฺติผลปฺปตฺโต ฯเปฯ กาเยน ผุสิตฺวา วิหรตีติ, น เหวํ วตฺตพฺเพ ฯเปฯ.}

\prose{\csromansymbolmark{*}อฏฺฐมกสฺส ปุคฺคลสฺส ทิฏฺฐิปริยุฏฺฐานํ ปหีนนฺติ, อามนฺตา.\csromanspacer{} อฏฺฐมกสฺส ปุคฺคลสฺส ทิฏฺฐานุสโย ปหีโนติ, น เหวํ วตฺตพฺเพ ฯเปฯ.\csromanspacer{} อฏฺฐมกสฺส ปุคฺคลสฺส ทิฏฺฐิปริยุฏฺฐานํ ปหีนนฺติ, อามนฺตา.\csromanspacer{} อฏฺฐมกสฺส ปุคฺคลสฺส วิจิกิจฺฉานุสโย.\csromanspacer{} สีลพฺพตปรามาโส ปหีโนติ, น เหวํ วตฺตพฺเพ ฯเปฯ.}

\prose{\csromansymbolmark{*}อฏฺฐมกสฺส ปุคฺคลสฺส วิจิกิจฺฉาปริยุฏฺฐานํ ปหีนนฺติ, อามนฺตา.\csromanspacer{} อฏฺฐมกสฺส ปุคฺคลสฺส วิจิกิจฺฉานุสโย ปหีโนติ, น เหวํ วตฺตพฺเพ ฯเปฯ.\csromanspacer{} อฏฺฐมกสฺส ปุคฺคลสฺส วิจิกิจฺฉาปริยุฏฺฐานํ ปหีนนฺติ, อามนฺตา.\csromanspacer{} อฏฺฐมกสฺส ปุคฺคลสฺส ทิฏฺฐานุสโย.\csromanspacer{} สีลพฺพตปรามาโส ปหีโนติ, น เหวํ วตฺตพฺเพ ฯเปฯ.}

\prose{\csromansymbolmark{*}อฏฺฐมกสฺส ปุคฺคลสฺส ทิฏฺฐานุสโย อปฺปหีโนติ, อามนฺตา.\csromanspacer{} อฏฺฐมกสฺส ปุคฺคลสฺส ทิฏฺฐิปริยุฏฺฐานํ อปฺปหีนนฺติ, น เหวํ วตฺตพฺเพ ฯเปฯ.\csromanspacer{} อฏฺฐมกสฺส ปุคฺคลสฺส ทิฏฺฐานุสโย อปฺปหีโนติ, อามนฺตา.\csromanspacer{} อฏฺฐมกสฺส ปุคฺคลสฺส วิจิกิจฺฉาปริยุฏฺฐานํ อปฺปหีนนฺติ, น เหวํ วตฺตพฺเพ ฯเปฯ.}

\prose{\csromanfolio{186}\csromansymbolmark{*}อฏฺฐมกสฺส ปุคฺคลสฺส วิจิกิจฺฉานุสโย.\csromanspacer{} สีลพฺพตปรามาโส อปฺปหีโนติ, อามนฺตา.\csromanspacer{} อฏฺฐมกสฺส ปุคฺคลสฺส ทิฏฺฐิปริยุฏฺฐานํ อปฺปหีนนฺติ, น เหวํ วตฺตพฺเพ ฯเปฯ.\csromanspacer{} อฏฺฐมกสฺส ปุคฺคลสฺส สีลพฺพตปรามาโส อปฺปหีโนติ, อามนฺตา.\csromanspacer{} อฏฺฐมกสฺส ปุคฺคลสฺส วิจิกิจฺฉาปริยุฏฺฐานํ อปฺปหีนนฺติ, น เหวํ วตฺตพฺเพ ฯเปฯ.}

\proseitem{369}{\csromansymbolmark{*}อฏฺฐมกสฺส ปุคฺคลสฺส ทิฏฺฐิปริยุฏฺฐานํ ปหีนนฺติ, อามนฺตา.\csromanspacer{} อฏฺฐมกสฺส ปุคฺคลสฺส ทิฏฺฐิปริยุฏฺฐานปหานาย มคฺโค ภาวิโตติ, น เหวํ วตฺตพฺเพ ฯเปฯ.\csromanspacer{} อฏฺฐมกสฺส ปุคฺคลสฺส ทิฏฺฐิปริยุฏฺฐานํ ปหีนนฺติ, อามนฺตา.\csromanspacer{} อฏฺฐมกสฺส ปุคฺคลสฺส ทิฏฺฐิปริยุฏฺฐานปหานาย สติปฏฺฐานา ภาวิตา ฯเปฯ สมฺมปฺปธานา ฯเปฯ โพชฺฌงฺคา ภาวิตาติ, น เหวํ วตฺตพฺเพ.}

\prose{\csromansymbolmark{*}อฏฺฐมกสฺส ปุคฺคลสฺส วิจิกิจฺฉาปริยุฏฺฐานํ ปหีนนฺติ, อามนฺตา.\csromanspacer{} อฏฺฐมกสฺส ปุคฺคลสฺส วิจิกิจฺฉาปริยุฏฺฐานปหานาย มคฺโค ภาวิโต ฯเปฯ โพชฺฌงฺคา ภาวิตาติ, น เหวํ วตฺตพฺเพ ฯเปฯ.}

\prose{\csromansymbolmark{*}อฏฺฐมกสฺส ปุคฺคลสฺส ทิฏฺฐิปริยุฏฺฐานปหานาย มคฺโค อภาวิโตติ, อามนฺตา.\csromanspacer{} อมคฺเคน ปหีนํ โลกิเยน สาสเวน ฯเปฯ สํกิเลสิเยนาติ, น เหวํ วตฺตพฺเพ ฯเปฯ.\csromanspacer{} อฏฺฐมกสฺส ปุคฺคลสฺส ทิฏฺฐิปริยุฏฺฐานปหานาย สติปฏฺฐานา ฯเปฯ โพชฺฌงฺคา อภาวิตาติ, อามนฺตา.\csromanspacer{} อมคฺเคน ปหีนํ โลกิเยน สาสเวน ฯเปฯ สํกิเลสิเยนาติ, น เหวํ วตฺตพฺเพ ฯเปฯ.}

\prose{\csromansymbolmark{*}อฏฺฐมกสฺส ปุคฺคลสฺส วิจิกิจฺฉาปริยุฏฺฐานปหานาย มคฺโค อภาวิโต ฯเปฯ สติปฏฺฐานา ฯเปฯ โพชฺฌงฺคา อภาวิตาติ, อามนฺตา.\csromanspacer{} อมคฺเคน ปหีนํ โลกิเยน สาสเวน ฯเปฯ สํกิเลสิเยนาติ, น เหวํ วตฺตพฺเพ ฯเปฯ.}

\proseitem{370}{\csromansymbolmark{+}น วตฺตพฺพํ “อฏฺฐมกสฺส ปุคฺคลสฺส ทิฏฺฐิปริยุฏฺฐานํ ปหีนนฺ”ติ, อามนฺตา.\csromanspacer{} อุปฺปชฺชิสฺสตีติ? นุปฺปชฺชิสฺสติ.\csromanspacer{} หญฺจิ นุปฺปชฺชิสฺสติ, เตน วต เร วตฺตพฺเพ “อฏฺฐมกสฺส ปุคฺคลสฺส ทิฏฺฐิปริยุฏฺฐานํ ปหีนนฺ”ติ.}

\prose{\csromansymbolmark{+}น วตฺตพฺพํ “อฏฺฐมกสฺส ปุคฺคลสฺส วิจิกิจฺฉาปริยุฏฺฐานํ ปหีนนฺ”ติ, อามนฺตา.\csromanspacer{} อุปฺปชฺชิสฺสตีติ? นุปฺปชฺชิสฺสติ.\csromanspacer{} หญฺจิ นุปฺปชฺชิสฺสติ, เตน วต เร วตฺตพฺเพ “อฏฺฐมกสฺส ปุคฺคลสฺส วิจิกิจฺฉาปริยุฏฺฐานํ ปหีนนฺ”ติ.}

\chapterheadpageread{3. ตติยวคฺค}
\csromanmatikaanchor{mtk.47}
\csromantocmarkref{subsection}{(26) 6. อฏฺฐมกสฺส-อินฺทฺริยกถา}{mtk.47}
\bookmark[dest={mtk.47},level=2]{(26) 6. อฏฺฐมกสฺส-อินฺทฺริยกถา}
\prose{\csromanfolio{187}\csromansymbolmark{*}อฏฺฐมกสฺส ปุคฺคลสฺส ทิฏฺฐิปริยุฏฺฐานํ นุปฺปชฺชิสฺสตีติ กตฺวา ปหีนนฺติ, อามนฺตา.\csromanspacer{} อฏฺฐมกสฺส ปุคฺคลสฺส ทิฏฺฐานุสโย นุปฺปชฺชิสฺสตีติ กตฺวา ปหีโนติ, น เหวํ วตฺตพฺเพ ฯเปฯ.}

\prose{\csromansymbolmark{*}อฏฺฐมกสฺส ปุคฺคลสฺส ทิฏฺฐิปริยุฏฺฐานํ นุปฺปชฺชิสฺสตีติ กตฺวา ปหีนนฺติ, อามนฺตา.\csromanspacer{} อฏฺฐมกสฺส ปุคฺคลสฺส วิจิกิจฺฉานุสโย.\csromanspacer{} สีลพฺพตปรามาโส นุปฺปชฺชิสฺสตีติ กตฺวา ปหีโนติ, น เหวํ วตฺตพฺเพ ฯเปฯ.}

\prose{\csromansymbolmark{*}อฏฺฐมกสฺส ปุคฺคลสฺส วิจิกิจฺฉาปริยุฏฺฐานํ นุปฺปชฺชิสฺสตีติ กตฺวา ปหีนนฺติ, อามนฺตา.\csromanspacer{} อฏฺฐมกสฺส ปุคฺคลสฺส วิจิกิจฺฉานุสโย.\csromanspacer{} สีลพฺพตปรามาโส นุปฺปชฺชิสฺสตีติ กตฺวา ปหีโนติ, น เหวํ วตฺตพฺเพ ฯเปฯ.}

\prose{\csromansymbolmark{*}อฏฺฐมกสฺส ปุคฺคลสฺส ทิฏฺฐิปริยุฏฺฐานํ นุปฺปชฺชิสฺสตีติ กตฺวา ปหีนนฺติ, อามนฺตา.\csromanspacer{} โคตฺรภุโน ปุคฺคลสฺส ทิฏฺฐิปริยุฏฺฐานํ นุปฺปชฺชิสฺสตีติ กตฺวา ปหีนนฺติ, น เหวํ วตฺตพฺเพ ฯเปฯ.\csromanspacer{} อฏฺฐมกสฺส ปุคฺคลสฺส วิจิกิจฺฉาปริยุฏฺฐานํ นุปฺปชฺชิสฺสตีติ กตฺวา ปหีนนฺติ, อามนฺตา.\csromanspacer{} โคตฺรภุโน ปุคฺคลสฺส วิจิกิจฺฉาปริยุฏฺฐานํ นุปฺปชฺชิสฺสตีติ กตฺวา ปหีนนฺติ, น เหวํ วตฺตพฺเพ ฯเปฯ.}

\nitthitam{อฏฺฐมกกถา นิฏฺฐิตา.}
\chapterhead{3. ตติยวคฺค}
\vspace{\tipitakaparskip}
\csromancenter{\textbf{(26) 6.}\csromanspacer{}\textbf{ อฏฺฐมกสฺส อินฺทฺริยกถา}}
\proseitem{371}{\csromansymbolmark{*}อฏฺฐมกสฺส ปุคฺคลสฺส นตฺถิ สทฺธินฺทฺริยนฺติ, อามนฺตา.\csromanspacer{} อฏฺฐมกสฺส ปุคฺคลสฺส นตฺถิ สทฺธาติ, น เหวํ วตฺตพฺเพ.\csromanspacer{} อฏฺฐมกสฺส ปุคฺคลสฺส นตฺถิ วีริยินฺทฺริยํ ฯเปฯ.\csromanspacer{} นตฺถิ สตินฺทฺริยํ ฯเปฯ.\csromanspacer{} นตฺถิ สมาธินฺทฺริยํ ฯเปฯ.\csromanspacer{} นตฺถิ ปญฺญินฺทฺริยนฺติ, อามนฺตา.\csromanspacer{} อฏฺฐมกสฺส ปุคฺคลสฺส นตฺถิ ปญฺญาติ, น เหวํ วตฺตพฺเพ ฯเปฯ.}

\prose{\csromansymbolmark{*}อฏฺฐมกสฺส ปุคฺคลสฺส อตฺถิ สทฺธาติ, อามนฺตา.\csromanspacer{} อฏฺฐมกสฺส ปุคฺคลสฺส อตฺถิ สทฺธินฺทฺริยนฺติ, น เหวํ วตฺตพฺเพ ฯเปฯ.\csromanspacer{} อฏฺฐมกสฺส ปุคฺคลสฺส อตฺถิ วีริยํ ฯเปฯ.\csromanspacer{} อตฺถิ สติ.\csromanspacer{} อตฺถิ สมาธิ, อตฺถิ ปญฺญาติ, อามนฺตา.\csromanspacer{} อฏฺฐมกสฺส ปุคฺคลสฺส อตฺถิ ปญฺญินฺทฺริยนฺติ, น เหวํ วตฺตพฺเพ ฯเปฯ.}

\prose{\csromansymbolmark{*}อฏฺฐมกสฺส ปุคฺคลสฺส อตฺถิ มโน, อตฺถิ มนินฺทฺริยนฺติ, อามนฺตา.\csromanspacer{} อฏฺฐมกสฺส ปุคฺคลสฺส อตฺถิ สทฺธา, อตฺถิ สทฺธินฺทฺริยนฺติ, น เหวํ วตฺตพฺเพ ฯเปฯ. \csromanfolio{188}อฏฺฐมกสฺส ปุคฺคลสฺส อตฺถิ มโน, อตฺถิ มนินฺทฺริยนฺติ, อามนฺตา.\csromanspacer{} อฏฺฐมกสฺส ปุคฺคลสฺส อตฺถิ ปญฺญา, อตฺถิ ปญฺญินฺทฺริยนฺติ, น เหวํ วตฺตพฺเพ ฯเปฯ. \csromansymbolmark{*}อฏฺฐมกสฺส ปุคฺคลสฺส อตฺถิ โสมนสฺสํ, อตฺถิ โสมนสฺสินฺทฺริยํ.\csromanspacer{} อตฺถิ ชีวิตํ, อตฺถิ ชีวิตินฺทฺริยนฺติ, อามนฺตา.\csromanspacer{} อฏฺฐมกสฺส ปุคฺคลสฺส อตฺถิ สทฺธา, อตฺถิ สทฺธินฺทฺริยนฺติ, น เหวํ วตฺตพฺเพ ฯเปฯ.\csromanspacer{} อฏฺฐมกสฺส ปุคฺคลสฺส อตฺถิ ชีวิตํ, อตฺถิ ชีวิตินฺทฺริยนฺติ, อามนฺตา.\csromanspacer{} อฏฺฐมกสฺส ปุคฺคลสฺส ฯเปฯ.\csromanspacer{} อตฺถิ ปญฺญา, อตฺถิ ปญฺญินฺทฺริยนฺติ, น เหวํ วตฺตพฺเพ ฯเปฯ. \csromansymbolmark{*}อฏฺฐมกสฺส ปุคฺคลสฺส อตฺถิ สทฺธา, นตฺถิ สทฺธินฺทฺริยนฺติ, อามนฺตา.\csromanspacer{} อฏฺฐมกสฺส ปุคฺคลสฺส อตฺถิ มโน, นตฺถิ มนินฺทฺริยนฺติ, น เหวํ วตฺตพฺเพ ฯเปฯ. \csromansymbolmark{*}อฏฺฐมกสฺส ปุคฺคลสฺส อตฺถิ สทฺธา, นตฺถิ สทฺธินฺทฺริยนฺติ, อามนฺตา.\csromanspacer{} อฏฺฐมกสฺส ปุคฺคลสฺส อตฺถิ โสมนสฺสํ, นตฺถิ โสมนสฺสินฺทฺริยนฺติ ฯเปฯ.\csromanspacer{} อตฺถิ ชีวิตํ, นตฺถิ ชีวิตินฺทฺริยนฺติ, น เหวํ วตฺตพฺเพ ฯเปฯ.\csromanspacer{} อฏฺฐมกสฺส ปุคฺคลสฺส อตฺถิ ปญฺญา, นตฺถิ ปญฺญินฺทฺริยนฺติ, อามนฺตา.\csromanspacer{} อฏฺฐมกสฺส ปุคฺคลสฺส อตฺถิ มโน, นตฺถิ มนินฺทฺริยนฺติ.\csromanspacer{} อตฺถิ โสมนสฺสํ, นตฺถิ โสมนสฺสินฺทฺริยนฺติ.\csromanspacer{} อตฺถิ ชีวิตํ, นตฺถิ ชีวิตินฺทฺริยนฺติ, น เหวํ วตฺตพฺเพ ฯเปฯ.}

\prose{\csromansymbolmark{*}อฏฺฐมกสฺส ปุคฺคลสฺส นตฺถิ สทฺธินฺทฺริยนฺติ, อามนฺตา.\csromanspacer{} อฏฺฐมโก ปุคฺคโล อสฺสทฺโธติ, น เหวํ วตฺตพฺเพ ฯเปฯ.\csromanspacer{} อฏฺฐมกสฺส ปุคฺคลสฺส นตฺถิ วีริยินฺทฺริยนฺติ, อามนฺตา.\csromanspacer{} อฏฺฐมโก ปุคฺคโล กุสีโต หีนวีริโยติ, น เหวํ วตฺตพฺเพ ฯเปฯ.\csromanspacer{} อฏฺฐมกสฺส ปุคฺคลสฺส นตฺถิ สตินฺทฺริยนฺติ, อามนฺตา.\csromanspacer{} อฏฺฐมโก ปุคฺคโล มุฏฺฐสฺสติ อสมฺปชาโนติ, น เหวํ วตฺตพฺเพ ฯเปฯ.\csromanspacer{} อฏฺฐมกสฺส ปุคฺคลสฺส นตฺถิ สมาธินฺทฺริยนฺติ, อามนฺตา.\csromanspacer{} อฏฺฐมโก ปุคฺคโล อสมาหิโต วิพฺภนฺตจิตฺโตติ, น เหวํ วตฺตพฺเพ ฯเปฯ.\csromanspacer{} อฏฺฐมกสฺส ปุคฺคลสฺส นตฺถิ ปญฺญินฺทฺริยนฺติ, อามนฺตา.\csromanspacer{} อฏฺฐมโก ปุคฺคโล ทุปฺปญฺโญ เอลมูโคติ, น เหวํ วตฺตพฺเพ ฯเปฯ. \csromansymbolmark{*}อฏฺฐมกสฺส ปุคฺคลสฺส อตฺถิ สทฺธา สา จ สทฺธา, นิยฺยานิกาติ, อามนฺตา.\csromanspacer{} หญฺจิ อฏฺฐมกสฺส ปุคฺคลสฺส อตฺถิ สทฺธา, สา จ สทฺธา นิยฺยานิกา, โน จ วต เร วตฺตพฺเพ “อฏฺฐมกสฺส ปุคฺคลสฺส นตฺถิ สทฺธินฺทฺริยนฺ”ติ.\csromanspacer{} อฏฺฐมกสฺส ปุคฺคลสฺส อตฺถิ วีริยํ, ตญฺจ วีริยํ นิยฺยานิกํ.\csromanspacer{} อตฺถิ สติ, สา จ สติ นิยฺยานิกา.\csromanspacer{} อตฺถิ สมาธิ, โส จ สมาธิ นิยฺยานิโก.\csromanspacer{} อตฺถิ ปญฺญา, สา จ ปญฺญา นิยฺยานิกาติ, อามนฺตา.\csromanspacer{} หญฺจิ อฏฺฐมกสฺส}

\vspace{\tipitakaparskip}
\csromancenter{ตติยวคฺค ปุคฺคลสฺส อตฺถิ ปญฺญา, สา จ ปญฺญา นิยฺยานิกา, โน จ วต เร วตฺตพฺเพ “อฏฺฐมกสฺส ปุคฺคลสฺส นตฺถิ ปญฺญินฺทฺริยนฺ”ติ.}
\proseitem{372}{\csromanfolio{189}\csromansymbolmark{*}สกทาคามิผลสจฺฉิกิริยาย ปฏิปนฺนสฺส ปุคฺคลสฺส อตฺถิ สทฺธา, อตฺถิ สทฺธินฺทฺริยนฺติ, อามนฺตา.\csromanspacer{} อฏฺฐมกสฺส ปุคฺคลสฺส อตฺถิ สทฺธา, อตฺถิ สทฺธินฺทฺริยนฺติ, น เหวํ วตฺตพฺเพ ฯเปฯ.\csromanspacer{} สกทาคามิผลสจฺฉิกิริยาย ปฏิปนฺนสฺส ปุคฺคลสฺส อตฺถิ ปญฺญา, อตฺถิ ปญฺญินฺทฺริยนฺติ, อามนฺตา.\csromanspacer{} อฏฺฐมกสฺส ปุคฺคลสฺส อตฺถิ ปญฺญา, อตฺถิ ปญฺญินฺทฺริยนฺติ, น เหวํ วตฺตพฺเพ ฯเปฯ.}

\prose{\csromansymbolmark{*}อนาคามิผลสจฺฉิกิริยาย ปฏิปนฺนสฺส ปุคฺคลสฺส.\csromanspacer{} อรหตฺตสจฺฉิกิริยาย ปฏิปนฺนสฺส ปุคฺคลสฺส อตฺถิ สทฺธา, อตฺถิ สทฺธินฺทฺริยํ ฯเปฯ.\csromanspacer{} อตฺถิ ปญฺญา, อตฺถิ ปญฺญินฺทฺริยนฺติ, อามนฺตา.\csromanspacer{} อฏฺฐมกสฺส ปุคฺคลสฺส อตฺถิ ปญฺญา, อตฺถิ ปญฺญินฺทฺริยนฺติ, น เหวํ วตฺตพฺเพ ฯเปฯ.}

\prose{\csromansymbolmark{*}อฏฺฐมกสฺส ปุคฺคลสฺส อตฺถิ สทฺธา, นตฺถิ สทฺธินฺทฺริยนฺติ, อามนฺตา.\csromanspacer{} สกทาคามิผลสจฺฉิกิริยาย ปฏิปนฺนสฺส ปุคฺคลสฺส อตฺถิ สทฺธา, นตฺถิ สทฺธินฺทฺริยนฺติ, น เหวํ วตฺตพฺเพ ฯเปฯ.\csromanspacer{} อฏฺฐมกสฺส ปุคฺคลสฺส อตฺถิ ปญฺญา, นตฺถิ ปญฺญินฺทฺริยนฺติ, อามนฺตา.\csromanspacer{} สกทาคามิผลสจฺฉิกิริยาย ปฏิปนฺนสฺส ปุคฺคลสฺส อตฺถิ ปญฺญา, นตฺถิ ปญฺญินฺทฺริยนฺติ, น เหวํ วตฺตพฺเพ ฯเปฯ.}

\prose{\csromansymbolmark{*}อฏฺฐมกสฺส ปุคฺคลสฺส อตฺถิ สทฺธา, นตฺถิ สทฺธินฺทฺริยนฺติ ฯเปฯ.\csromanspacer{} อตฺถิ ปญฺญา, นตฺถิ ปญฺญินฺทฺริยนฺติ, อามนฺตา.\csromanspacer{} อนาคามิผลสจฺฉิกิริยาย ปฏิปนฺนสฺส ปุคฺคลสฺส.\csromanspacer{} อรหตฺตสจฺฉิกิริยาย ปฏิปนฺนสฺส ปุคฺคลสฺส อตฺถิ ปญฺญา, นตฺถิ ปญฺญินฺทฺริยนฺติ, น เหวํ วตฺตพฺเพ ฯเปฯ.}

\csromanmatikaanchor{mtk.48}
\csromantocmarkref{subsection}{(27) 7. ทิพฺพจกฺขุกถา}{mtk.48}
\bookmark[dest={mtk.48},level=2]{(27) 7. ทิพฺพจกฺขุกถา}
\prose{\csromansymbolmark{*}อฏฺฐมกสฺส ปุคฺคลสฺส นตฺถิ ปญฺจินฺทฺริยานีติ, อามนฺตา.\csromanspacer{} นนุ วุตฺตํ ภควตา “ปญฺจิมานิ ภิกฺขเว อินฺทฺริยานิ.\csromanspacer{} กตมานิ ปญฺจ, สทฺธินฺทฺริยํ, วีริยินฺทฺริยํ, สตินฺทฺริยํ, สมาธินฺทฺริยํ, ปญฺญินฺทฺริยํ.\csromanspacer{} อิมานิ โข ภิกฺขเว ปญฺจินฺทฺริยานิ.\csromanspacer{} อิเมสํ โข ภิกฺขเว ปญฺจนฺนํ อินฺทฺริยานํ สมตฺตา ปริปูรตฺตา อรหา โหติ, ตโต มุทุตเรหิ อรหตฺตสจฺฉิกิริยาย ปฏิปนฺโน โหติ, ตโต มุทุตเรหิ อนาคามี โหติ, ตโต มุทุตเรหิ อนาคามิผลสจฺฉิกิริยาย ปฏิปนฺโน โหติ, ตโต มุทุตเรหิ สกทาคามี โหติ, ตโต มุทุตเรหิ สกทาคามิผลสจฺฉิกิริยาย ปฏิปนฺโน โหติ, ตโต มุทุตเรหิ โสตาปนฺโน โหติ, ตโต มุทุตเรหิ โสตาปตฺติผลสจฺฉิกิริยาย ปฏิปนฺโน โหติ.\csromanspacer{} ยสฺส โข ภิกฺขเว \csromanfolio{190}อิมานิ ปญฺจินฺทฺริยานิ สพฺเพน สพฺพํ สพฺพถา สพฺพํ นตฺถิ, ตมหํ ‘พาหิโร ปุถุชฺชนปกฺเข ฐิโต’ติ วทามี”ติ\footnote{สํ 3. 178 ปิฏฺเฐ.} อตฺเถว สุตฺตนฺโตติ, อามนฺตา.\csromanspacer{} อฏฺฐมโก ปุคฺคโล พาหิโร ปุถุชฺชนปกฺเข ฐิโตติ, น เหวํ วตฺตพฺเพ ฯเปฯ.\csromanspacer{} เตน หิ อฏฺฐมกสฺส ปุคฺคลสฺส อตฺถิ ปญฺจินฺทฺริยานีติ.}

\nitthitam{อฏฺฐมกสฺส อินฺทฺริยกถา นิฏฺฐิตา.}
\chapterhead{3. ตติยวคฺค}
\vspace{\tipitakaparskip}
\csromancenter{\textbf{(27) 7.}\csromanspacer{}\textbf{ ทิพฺพจกฺขุกถา}}
\proseitem{373}{\csromansymbolmark{*}มํสจกฺขุํ ธมฺมุปตฺถทฺธํ ทิพฺพจกฺขุํ โหตีติ, อามนฺตา.\csromanspacer{} มํสจกฺขุํ ทิพฺพจกฺขุํ, ทิพฺพจกฺขุํ มํสจกฺขุนฺติ, น เหวํ วตฺตพฺเพ ฯเปฯ.}

\prose{\csromansymbolmark{*}มํสจกฺขุํ ธมฺมุปตฺถทฺธํ ทิพฺพจกฺขุํ โหตีติ, อามนฺตา.\csromanspacer{} ยาทิสํ มํสจกฺขุํ, ตาทิสํ ทิพฺพจกฺขุํ, ยาทิสํ ทิพฺพจกฺขุํ, ตาทิสํ มํสจกฺขุนฺติ, น เหวํ วตฺตพฺเพ ฯเปฯ.}

\prose{\csromansymbolmark{*}มํสจกฺขุํ ธมฺมุปตฺถทฺธํ ทิพฺพจกฺขุํ โหตีติ, อามนฺตา.\csromanspacer{} ตญฺเญว มํสจกฺขุํ ตํ ทิพฺพจกฺขุํ, ตํ ทิพฺพจกฺขุํ ตํ มํสจกฺขุนฺติ, น เหวํ วตฺตพฺเพ ฯเปฯ.}

\prose{\csromansymbolmark{*}มํสจกฺขุํ ธมฺมุปตฺถทฺธํ ทิพฺพจกฺขุํ โหตีติ, อามนฺตา.\csromanspacer{} ยาทิโส มํสจกฺขุสฺส วิสโย อานุภาโว โคจโร, ตาทิโส ทิพฺพสฺส จกฺขุสฺส วิสโย อานุภาโว โคจโรติ, น เหวํ วตฺตพฺเพ ฯเปฯ.}

\prose{\csromansymbolmark{*}มํสจกฺขุํ ธมฺมุปตฺถทฺธํ ทิพฺพจกฺขุํ โหตีติ, อามนฺตา.\csromanspacer{} อุปาทิณฺณํ\footnote{อุปาทินฺนํ (สฺยา, กํ)} หุตฺวา อนุปาทิณฺณํ โหตีติ, น เหวํ วตฺตพฺเพ ฯเปฯ.}

\prose{\csromansymbolmark{*}อุปาทิณฺณํ หุตฺวา อนุปาทิณฺณํ โหตีติ, อามนฺตา.\csromanspacer{} กามาวจรํ หุตฺวา รูปาวจรํ โหตีติ, น เหวํ วตฺตพฺเพ ฯเปฯ.}

\prose{\csromansymbolmark{*}กามาวจรํ หุตฺวา รูปาวจรํ โหตีติ, อามนฺตา.\csromanspacer{} รูปาวจรํ หุตฺวา อรูปาวจรํ โหตีติ, น เหวํ วตฺตพฺเพ ฯเปฯ.}

\prose{\csromansymbolmark{*}รูปาวจรํ หุตฺวา อรูปาวจรํ โหตีติ, อามนฺตา.\csromanspacer{} ปริยาปนฺนํ หุตฺวา อปริยาปนฺนํ โหตีติ, น เหวํ วตฺตพฺเพ ฯเปฯ.}

\chapterheadpageread{3. ตติยวคฺค}
\csromanmatikaanchor{mtk.49}
\csromantocmarkref{subsection}{(28) 8. ทิพฺพโสตกถา}{mtk.49}
\bookmark[dest={mtk.49},level=2]{(28) 8. ทิพฺพโสตกถา}
\proseitem{374}{\csromanfolio{191}\csromansymbolmark{*}มํสจกฺขุํ ธมฺมุปตฺถทฺธํ ทิพฺพจกฺขุํ โหตีติ, อามนฺตา.\csromanspacer{} ทิพฺพจกฺขุํ ธมฺมุปตฺถทฺธํ มํสจกฺขุํ โหตีติ, น เหวํ วตฺตพฺเพ ฯเปฯ.}

\prose{\csromansymbolmark{*}มํสจกฺขุํ ธมฺมุปตฺถทฺธํ ทิพฺพจกฺขุํ โหตีติ, อามนฺตา.\csromanspacer{} ทิพฺพจกฺขุํ ธมฺมุปตฺถทฺธํ ปญฺญาจกฺขุํ โหตีติ, น เหวํ วตฺตพฺเพ ฯเปฯ.}

\prose{\csromansymbolmark{*}มํสจกฺขุํ ธมฺมุปตฺถทฺธํ ทิพฺพจกฺขุํ โหตีติ, อามนฺตา.\csromanspacer{} ทิพฺพจกฺขุํ ธมฺมุปตฺถทฺธํ มํสจกฺขุํ โหตีติ, น เหวํ วตฺตพฺเพ ฯเปฯ.}

\prose{\csromansymbolmark{*}มํสจกฺขุํ ธมฺมุปตฺถทฺธํ ทิพฺพจกฺขุํ โหตีติ, อามนฺตา.\csromanspacer{} เทฺวว จกฺขูนีติ, น เหวํ วตฺตพฺเพ ฯเปฯ.\csromanspacer{} เทฺวว จกฺขูนีติ, อามนฺตา.\csromanspacer{} นนุ ตีณิ จกฺขูนิ วุตฺตานิ ภควตา “มํสจกฺขุํ ทิพฺพจกฺขุํ ปญฺญาจกฺขุนฺ”ติ, อามนฺตา.\csromanspacer{} หญฺจิ ตีณิ จกฺขูนิ วุตฺตานิ ภควตา มํสจกฺขุํ ทิพฺพจกฺขุํ ปญฺญาจกฺขุํ, โน จ วต เร วตฺตพฺเพ “เทฺวว จกฺขูนี”ติ.}

\prose{\csromansymbolmark{*}เทฺวว จกฺขูนีติ, อามนฺตา.\csromanspacer{} นนุ วุตฺตํ ภควตา “ตีณิมานิ ภิกฺขเว จกฺขูนิ.\csromanspacer{} กตมานิ ตีณิ, มํสจกฺขุํ ทิพฺพจกฺขุํ ปญฺญาจกฺขุนฺติ.\csromanspacer{} อิมานิ โข ภิกฺขเว ตีณิ จกฺขูนีติ.}

\csromangathameasure{มํสจกฺขุํ ทิพฺพจกฺขุํ, ปญฺญาจกฺขุํ อนุตฺตรํ. \\ เอตานิ ตีณิ จกฺขูนิ, อกฺขาสิ ปุริสุตฺตโม. \\ มํสจกฺขุสฺส อุปฺปาโท, มคฺโค ทิพฺพสฺส จกฺขุโน. \\ ยทา จ ญาณํ อุทปาทิ, ปญฺญาจกฺขุํ อนุตฺตรํ. \\ ตสฺส จกฺขุสฺส ปฏิลาภา, สพฺพทุกฺขา ปมุจฺจตี”ติ. อตฺเถว สุตฺตนฺโตติ, อามนฺตา. เตน หิ น วตฺตพฺพํ “เทฺวว จกฺขูนี”ติ.}
\csromangathagroup{\csromangathastanza{มํสจกฺขุํ ทิพฺพจกฺขุํ, ปญฺญาจกฺขุํ อนุตฺตรํ. \\ เอตานิ ตีณิ จกฺขูนิ, อกฺขาสิ ปุริสุตฺตโม.}\csromangathastanza{มํสจกฺขุสฺส อุปฺปาโท, มคฺโค ทิพฺพสฺส จกฺขุโน. \\ ยทา จ ญาณํ อุทปาทิ, ปญฺญาจกฺขุํ อนุตฺตรํ. \\ ตสฺส จกฺขุสฺส ปฏิลาภา, สพฺพทุกฺขา ปมุจฺจตี”ติ\footnote{ขุ 1. 231 อิติวุตฺตเก.}.\csromanspacer{} อตฺเถว สุตฺตนฺโตติ, อามนฺตา.\csromanspacer{} เตน หิ น วตฺตพฺพํ “เทฺวว จกฺขูนี”ติ.}}

\nitthitam{ทิพฺพจกฺขุกถา นิฏฺฐิตา.}
\chapterhead{3. ตติยวคฺค}
\vspace{\tipitakaparskip}
\csromancenter{\textbf{(28) 8.}\csromanspacer{}\textbf{ ทิพฺพโสตกถา}}
\proseitem{375}{\csromansymbolmark{*}มํสโสตํ ธมฺมุปตฺถทฺธํ, ทิพฺพโสตํ โหตีติ, อามนฺตา.\csromanspacer{} มํสโสตํ ทิพฺพโสตํ, ทิพฺพโสตํ มํสโสตนฺติ, น เหวํ วตฺตพฺเพ ฯเปฯ.}

\prose{\csromanfolio{192}\csromansymbolmark{*}มํสโสตํ ธมฺมุปตฺถทฺธํ ทิพฺพโสตํ โหตีติ, อามนฺตา.\csromanspacer{} ยาทิสํ มํสโสตํ, ตาทิสํ ทิพฺพโสตํ, ยาทิสํ ทิพฺพโสตํ.\csromanspacer{} ตาทิสํ มํสโสตนฺติ, น เหวํ วตฺตพฺเพ ฯเปฯ.}

\prose{\csromansymbolmark{*}มํสโสตํ ธมฺมุปตฺถทฺธํ, ทิพฺพโสตํ โหตีติ, อามนฺตา.\csromanspacer{} ตญฺเญว มํสโสตํ ตํ ทิพฺพโสตํ, ตํ ทิพฺพโสตํ ตํ มํสโสตนฺติ, น เหวํ วตฺตพฺเพ ฯเปฯ.}

\prose{\csromansymbolmark{*}มํสโสตํ ธมฺมุปตฺถทฺธํ, ทิพฺพโสตํ โหตีติ, อามนฺตา.\csromanspacer{} ยาทิโส มํสโสตสฺส วิสโย อานุภาโว โคจโร, ตาทิโส ทิพฺพสฺส โสตสฺส วิสโย อานุภาโว โคจโรติ, น เหวํ วตฺตพฺเพ ฯเปฯ.}

\prose{\csromansymbolmark{*}มํสโสตํ ธมฺมุปตฺถทฺธํ, ทิพฺพโสตํ โหตีติ, อามนฺตา.\csromanspacer{} อุปาทิณฺณํ หุตฺวา อนุปาทิณฺณํ โหตีติ, น เหวํ วตฺตพฺเพ ฯเปฯ. \csromansymbolmark{*}อุปาทิณฺณํ หุตฺวา อนุปาทิณฺณํ โหตีติ, อามนฺตา.\csromanspacer{} กามาวจรํ หุตฺวา รูปาวจรํ โหตีติ, น เหวํ วตฺตพฺเพ ฯเปฯ.}

\prose{\csromansymbolmark{*}กามาวจรํ หุตฺวา รูปาวจรํ โหตีติ, อามนฺตา.\csromanspacer{} รูปาวจรํ หุตฺวา อรูปาวจรํ โหตีติ, น เหวํ วตฺตพฺเพ ฯเปฯ.}

\prose{\csromansymbolmark{*}รูปาวจรํ หุตฺวา อรูปาวจรํ โหตีติ, อามนฺตา.\csromanspacer{} ปริยาปนฺนํ หุตฺวา อปริยาปนฺนํ โหตีติ, น เหวํ วตฺตพฺเพ ฯเปฯ.}

\proseitem{376}{\csromansymbolmark{*}มํสโสตํ ธมฺมุปตฺถทฺธํ, ทิพฺพโสตํ โหตีติ, อามนฺตา.\csromanspacer{} ทิพฺพโสตํ ธมฺมุปตฺถทฺธํ, มํสโสตํ โหตีติ, น เหวํ วตฺตพฺเพ ฯเปฯ.}

\prose{\csromansymbolmark{*}มํสโสตํ ธมฺมุปตฺถทฺธํ, ทิพฺพโสตํ โหตีติ, อามนฺตา.\csromanspacer{} เอกํเยว โสตนฺติ, น เหวํ วตฺตพฺเพ ฯเปฯ.}

\prose{\csromansymbolmark{*}เอกํเยว โสตนฺติ, อามนฺตา.\csromanspacer{} นนุ เทฺว โสตานิ วุตฺตานิ ภควตา “มํสโสตํ ทิพฺพโสตนฺ”ติ, อามนฺตา.\csromanspacer{} หญฺจิ เทฺว โสตานิ วุตฺตานิ ภควตา มํสโสตํ ทิพฺพโสตํ, โน จ วต เร วตฺตพฺเพ “เอกญฺเญว โสตนฺ”ติ.}

\nitthitam{ทิพฺพโสตกถา นิฏฺฐิตา.}
\chapterheadpageread{3. ตติยวคฺค \textbf{3. ตติยวคฺค}}
\vspace{\tipitakaparskip}
\csromancenter{\textbf{(29) 9.}\csromanspacer{}\textbf{ ยถากมฺมูปคตญาณกถา}}
\csromanmatikaanchor{mtk.50}
\csromantocmarkref{subsection}{(29) 9. ยถากมฺมูปคตญาณกถา}{mtk.50}
\bookmark[dest={mtk.50},level=2]{(29) 9. ยถากมฺมูปคตญาณกถา}
\proseitem{377}{\csromanfolio{193}\csromansymbolmark{*}ยถากมฺมูปคตํ ญาณํ\footnote{ยถากมฺมูปคตญาณํ (สฺยา, กํ)} ทิพฺพจกฺขุนฺติ, อามนฺตา.\csromanspacer{} ยถากมฺมูปคตญฺจ มนสิ กโรติ, ทิพฺเพน จกฺขุนา รูปํ ปสฺสตีติ, น เหวํ วตฺตพฺเพ ฯเปฯ.}

\prose{\csromansymbolmark{*}ยถากมฺมูปคตญฺจ มนสิ กโรติ, ทิพฺเพน จกฺกุนา รูปํ ปสฺสตีติ, อามนฺตา.\csromanspacer{} ทฺวินฺนํ ผสฺสานํ ทฺวินฺนํ จิตฺตานํ สโมธานํ โหตีติ, น เหวํ วตฺตพฺเพ ฯเปฯ.}

\prose{\csromansymbolmark{*}ยถากมฺมูปคตํ ญาณํ ทิพฺพจกฺขุนฺติ, อามนฺตา. “อิเม วต โภนฺโต สตฺตา”ติ จ มนสิ กโรติ, “กายทุจฺจริเตน สมนฺนาคตา”ติ จ มนสิ กโรติ, “วจีทุจฺจริเตน สมนฺนาคตา”ติ จ มนสิ กโรติ, “มโนทุจฺจริเตน สมนฺนาคตา”ติ จ มนสิ กโรติ, “อริยานํ อุปวาทกา”ติ จ มนสิ กโรติ, “มิจฺฉาทิฏฺฐิกา”ติ จ มนสิ กโรติ, “มิจฺฉาทิฏฺฐิกมฺมสมาทานา”ติ จ มนสิ กโรติ, “เต กายสฺส เภทา ปรํ มรณา อปายํ ทุคฺคติํ วินิปาตํ นิรยํ อุปปนฺนา”ติ จ มนสิ กโรติ, “อิเม วา ปน โภนฺโต สตฺตา”ติ จ มนสิ กโรติ, “กายสุจริเตน สมนฺนาคตา”ติ จ มนสิ กโรติ, “วจีสุจริเตน สมนฺนาคตา”ติ จ มนสิ กโรติ, “มโนสุจริเตน สมนฺนาคตา”ติ จ มนสิ กโรติ, “อริยานํ อนุปวาทกา”ติ จ มนสิ กโรติ, “สมฺมาทิฏฺฐิกา”ติ จ มนสิ กโรติ, “สมฺมาทิฏฺฐิกมฺมสมาทานา”ติ จ มนสิ กโรติ, “เต กายสฺส เภทา ปรํ มรณา สุคติํ สคฺคํ โลกํ อุปปนฺนา”ติ จ มนสิ กโรติ.\csromanspacer{} ทิพฺเพน จกฺขุนา รูปํ ปสฺสตีติ, น เหวํ วตฺตพฺเพ ฯเปฯ.}

\prose{\csromansymbolmark{*}“เต กายสฺส เภทา ปรํ มรณา สุคติํ สคฺคํ โลกํ อุปปนฺนา”ติ จ มนสิ กโรติ.\csromanspacer{} ทิพฺเพน จกฺกุนา รูปํ ปสฺสตีติ, อามนฺตา.\csromanspacer{} ทฺวินฺนํ ผสฺสานํ ทฺวินฺนํ จิตฺตานํ สโมธานํ โหตีติ, น เหวํ วตฺตพฺเพ ฯเปฯ.}

\csromanmatikaanchor{mtk.51}
\csromantocmarkref{subsection}{(30) 10. สํวรกถา}{mtk.51}
\bookmark[dest={mtk.51},level=2]{(30) 10. สํวรกถา}
\proseitem{378}{\csromansymbolmark{*}ยถากมฺมูปคตํ ญาณํ ทิพฺพจกฺขุนฺติ, อามนฺตา.\csromanspacer{} อตฺถิ โกจิ อทิพฺพจกฺขุโก ทิพฺพจกฺขุํ อปฺปฏิลทฺโธ อนธิคโต อสจฺฉิกโต ยถากมฺมูปคตํ ชานาตีติ, อามนฺตา.\csromanspacer{} หญฺจิ อตฺถิ โกจิ อทิพฺพจกฺขุโก \csromanfolio{194}ทิพฺพจกฺขุํ อปฺปฏิลทฺโธ อนธิคโต อสจฺฉิกโต ยถากมฺมูปคตํ ชานาติ, โน จ วต เร วตฺตพฺเพ “ยถากมฺมูปคตํ ญาณํ ทิพฺพจกฺขุนฺ”ติ.}

\prose{\csromansymbolmark{*}ยถากมฺมูปคตํ ญาณํ ทิพฺพจกฺขุนฺติ, อามนฺตา.\csromanspacer{} อายสฺมา สาริปุตฺโต ยถากมฺมูปคตํ ญาณํ ชานาตีติ, อามนฺตา.\csromanspacer{} หญฺจิ อายสฺมา สาริปุตฺโต ยถากมฺมูปคตํ ญาณํ ชานาติ, โน จ วต เร วตฺตพฺเพ “ยถากมฺมูปคตํ ญาณํ ทิพฺพจกฺขุนฺ”ติ.}

\prose{\csromansymbolmark{*}ยถากมฺมูปคตํ ญาณํ ทิพฺพจกฺขุนฺติ, อามนฺตา.\csromanspacer{} อายสฺมา สาริปุตฺโต ยถากมฺมูปคตํ ญาณํ ชานาตีติ, อามนฺตา.\csromanspacer{} อตฺถายสฺมโต สาริปุตฺตสฺส ทิพฺพจกฺขุนฺติ, น เหวํ วตฺตพฺเพ ฯเปฯ.}

\prose{\csromansymbolmark{*}อตฺถายสฺมโต สาริปุตฺตสฺส ทิพฺพจกฺขุนฺติ, อามนฺตา.\csromanspacer{} นนุ อายสฺมา สาริปุตฺโต เอตทโวจ\csromandash{}}

\csromangathameasure{“เนว ปุพฺเพนิวาสาย, นปิ ทิพฺพสฺส จกฺขุโน. \\ เจโตปริยาย อิทฺธิยา, โสตธาตุวิสุทฺธิยา. \\ จุติยา อุปปตฺติยา, ปณิธิ เม น วิชฺชตี”ติ. อตฺเถว สุตฺตนฺโตติ, อามนฺตา. เตน หิ น วตฺตพฺพํ “ยถากมฺมูปคตํ ญาณํ ทิพฺพจกฺขุนฺ”ติ.}
\csromangathagroup{\csromangathastanza{“เนว ปุพฺเพนิวาสาย, นปิ ทิพฺพสฺส จกฺขุโน. \\ เจโตปริยาย อิทฺธิยา, โสตธาตุวิสุทฺธิยา. \\ จุติยา อุปปตฺติยา, ปณิธิ เม น วิชฺชตี”ติ\footnote{ขุ 2. 344 เถรคาถายํ.}.\csromanspacer{} อตฺเถว สุตฺตนฺโตติ, อามนฺตา.\csromanspacer{} เตน หิ น วตฺตพฺพํ “ยถากมฺมูปคตํ ญาณํ ทิพฺพจกฺขุนฺ”ติ.}}

\nitthitam{ยถากมฺมูปคตญาณกถา นิฏฺฐิตา.}
\chapterhead{3. ตติยวคฺค}
\vspace{\tipitakaparskip}
\csromancenter{\textbf{(30) 10.}\csromanspacer{}\textbf{ สํวรกถา}}
\proseitem{379}{\csromansymbolmark{*}อตฺถิ เทเวสุ สํวโรติ, อามนฺตา.\csromanspacer{} อตฺถิ เทเวสุ อสํวโรติ, น เหวํ วตฺตพฺเพ ฯเปฯ.}

\prose{\csromansymbolmark{*}นตฺถิ เทเวสุ อสํวโรติ, อามนฺตา.\csromanspacer{} นตฺถิ เทเวสุ สํวโรติ, น เหวํ วตฺตพฺเพ ฯเปฯ.}

\chapterheadpageread{3. ตติยวคฺค}
\prose{\csromanfolio{195}\csromansymbolmark{*}นนุ อสํวรา สํวโร สีลํ, อตฺถิ เทเวสุ สํวโรติ, อามนฺตา.\csromanspacer{} อตฺถิ เทเวสุ อสํวโร, ยมฺหา อสํวรา สํวโร สีลนฺติ, น เหวํ วตฺตพฺเพ ฯเปฯ.}

\prose{อาชานาหิ นิคฺคหํ\csromandash{}หญฺจิ อสํวรา สํวโร สีลํ, อตฺถิ เทเวสุ สํวโร, เตน วต เร วตฺตพฺเพ “อตฺถิ เทเวสุ อสํวโร, ยมฺหา อสํวรา สํวโร สีลนฺ”ติ.\csromanspacer{} ยํ ตตฺถ วเทสิ “วตฺตพฺเพ โข ‘อสํวรา สํวโร สีลํ, อตฺถิ เทเวสุ สํวโร’, โน จ วตฺตพฺเพ ‘อตฺถิ เทเวสุ อสํวโร, ยมฺหา อสํวรา สํวโร สีลนฺ’ติ”, มิจฺฉา.}

\prose{โน เจ ปน วตฺตพฺเพ “อตฺถิ เทเวสุ อสํวโร, ยมฺหา อสํวรา สํวโร สีลนฺ”ติ, โน จ วต เร วตฺตพฺเพ “อสํวรา สํวโร สีลํ, อตฺถิ เทเวสุ สํวโร”ติ.\csromanspacer{} ยํ ตตฺถ วเทสิ “วตฺตพฺเพ โข ‘อสํวรา สํวโร สีลํ, อตฺถิ เทเวสุ สํวโร’, โน จ วตฺตพฺเพ ‘อตฺถิ เทเวสุ อสํวโร, ยมฺหา อสํวรา สํวโร สีลนฺ’ติ, มิจฺฉา.}

\prose{\csromansymbolmark{*}อตฺถิ มนุสฺเสสุ สํวโร, อตฺถิ ตตฺถ อสํวโรติ, อามนฺตา.\csromanspacer{} อตฺถิ เทเวสุ สํวโร, อตฺถิ ตตฺถ อสํวโรติ, น เหวํ วตฺตพฺเพ ฯเปฯ.}

\prose{\csromansymbolmark{*}อตฺถิ เทเวสุ สํวโร, นตฺถิ ตตฺถ อสํวโรติ, อามนฺตา.\csromanspacer{} อตฺถิ มนุสฺเสสุ สํวโร, นตฺถิ ตตฺถ อสํวโรติ, น เหวํ วตฺตพฺเพ ฯเปฯ.}

\proseitem{380}{\csromansymbolmark{*}อตฺถิ เทเวสุ ปาณาติปาตา เวรมณีติ, อามนฺตา.\csromanspacer{} อตฺถิ เทเวสุ ปาณาติปาโตติ, น เหวํ วตฺตพฺเพ ฯเปฯ.}

\prose{\csromansymbolmark{*}อตฺถิ เทเวสุ สุราเมรยมชฺชปมาทฏฺฐานา เวรมณีติ, อามนฺตา.\csromanspacer{} อตฺถิ เทเวสุ สุราเมรยมชฺชปมาทฏฺฐานนฺติ, น เหวํ วตฺตพฺเพ ฯเปฯ.}

\prose{\csromansymbolmark{*}นตฺถิ เทเวสุ ปาณาติปาโตติ, อามนฺตา.\csromanspacer{} นตฺถิ เทเวสุ ปาณาติปาตา เวรมณีติ, น เหวํ วตฺตพฺเพ ฯเปฯ.}

\prose{\csromansymbolmark{*}นตฺถิ เทเวสุ สุราเมรยมชฺชปมาทฏฺฐานนฺติ, อามนฺตา.\csromanspacer{} นตฺถิ เทเวสุ สุราเมรยมชฺชปมาทฏฺฐานา เวรมณีติ, น เหวํ วตฺตพฺเพ ฯเปฯ.}

\prose{\csromansymbolmark{*}อตฺถิ มนุสฺเสสุ ปาณาติปาตา เวรมณิ, อตฺถิ ตตฺถ ปาณาติปาโตติ, อามนฺตา.\csromanspacer{} อตฺถิ เทเวสุ ปาณาติปาตา เวรมณิ, อตฺถิ ตตฺถ ปาณาติปาโตติ, น เหวํ วตฺตพฺเพ ฯเปฯ.}

\csromanmatikaanchor{mtk.52}
\csromantocmarkref{subsection}{(31) 11. อสญฺญกถา}{mtk.52}
\bookmark[dest={mtk.52},level=2]{(31) 11. อสญฺญกถา}
\prose{\csromanfolio{196}\csromansymbolmark{*}อตฺถิ มนุสฺเสสุ สุราเมรยมชฺชปมาทฏฺฐานา เวรมณิ, อตฺถิ ตตฺถ สุราเมรยมชฺชปมาทฏฺฐานนฺติ, อามนฺตา.\csromanspacer{} อตฺถิ เทเวสุ สุราเมรยมชฺชปมาทฏฺฐานา เวรมณิ, อตฺถิ ตตฺถ สุราเมรยมชฺชปมาทฏฺฐานนฺติ, น เหวํ วตฺตพฺเพ ฯเปฯ. \csromansymbolmark{*}อตฺถิ เทเวสุ ปาณาติปาตา เวรมณิ, นตฺถิ ตตฺถ ปาณาติปาโตติ, อามนฺตา.\csromanspacer{} อตฺถิ มนุสฺเสสุ ปาณาติปาตา เวรมณิ, นตฺถิ ตตฺถ ปาณาติปาโตติ, น เหวํ วตฺตพฺเพ ฯเปฯ. \csromansymbolmark{*}อตฺถิ เทเวสุ สุราเมรยมชฺชปมาทฏฺฐานา เวรมณิ, นตฺถิ ตตฺถ สุราเมรยมชฺชปมาทฏฺฐานนฺติ, อามนฺตา.\csromanspacer{} อตฺถิ มนุสฺเสสุ สุราเมรยมชฺชปมาทฏฺฐานา เวรมณิ, นตฺถิ ตตฺถ สุราเมรยมชฺชปมาทฏฺฐานนฺติ, น เหวํ วตฺตพฺเพ ฯเปฯ.}

\prose{\csromansymbolmark{+}นตฺถิ เทเวสุ สํวโรติ, อามนฺตา.\csromanspacer{} สพฺเพ เทวา ปาณาติปาติโน อทินฺนาทายิโน กาเมสุมิจฺฉาจาริโน มุสาวาทิโน สุราเมรยมชฺชปมาทฏฺฐายิโนติ, น เหวํ วตฺตพฺเพ ฯเปฯ.\csromanspacer{} เตน หิ อตฺถิ เทเวสุ สํวโรติ.}

\nitthitam{สํวรกถา นิฏฺฐิตา.}
\chapterhead{3. ตติยวคฺค}
\vspace{\tipitakaparskip}
\csromancenter{\textbf{(31) 11.}\csromanspacer{}\textbf{ อสญฺญกถา}}
\proseitem{381}{\csromansymbolmark{*}อสญฺญสตฺเตสุ สญฺญา อตฺถีติ, อามนฺตา.\csromanspacer{} สญฺญาภโว สญฺญาคติ สญฺญาสตฺตาวาโส สญฺญาสํสาโร สญฺญาโยนิ สญฺญตฺตภาวปฏิลาโภติ, น เหวํ วตฺตพฺเพ ฯเปฯ.}

\prose{นนุ อสญฺญภโว อสญฺญคติ อสญฺญสตฺตาวาโส อสญฺญสํสาโร อสญฺญโยนิ อสญฺญตฺตภาวปฏิลาโภติ, อามนฺตา.\csromanspacer{} หญฺจิ อสญฺญภโว อสญฺญคติ อสญฺญสตฺตาวาโส อสญฺญสํสาโร อสญฺญโยนิ อสญฺญตฺตภาวปฏิลาโภ, โน จ วต เร วตฺตพฺเพ “อสญฺญสตฺเตสุ สญฺญา อตฺถี”ติ.}

\chapterheadpageread{3. ตติยวคฺค}
\prose{\csromanfolio{197}\csromansymbolmark{*}อสญฺญสตฺเตสุ สญฺญา อตฺถีติ, อามนฺตา.\csromanspacer{} ปญฺจโวการภโว คติ สตฺตาวาโส สํสาโร โยนิ อตฺตภาวปฏิลาโภติ, น เหวํ วตฺตพฺเพ ฯเปฯ.}

\prose{นนุ เอกโวการภโว คติ สตฺตาวาโส สํสาโร โยนิ อตฺตภาวปฏิลาโภติ, อามนฺตา.\csromanspacer{} หญฺจิ เอกโวการภโว คติ สตฺตาวาโส สํสาโร โยนิ อตฺตภาวปฏิลาโภ, โน จ วต เร วตฺตพฺเพ “อสญฺญสตฺเตสุ สญฺญา อตฺถี”ติ.}

\prose{\csromansymbolmark{*}อสญฺญสตฺเตสุ สญฺญา อตฺถีติ, อามนฺตา.\csromanspacer{} ตาย สญฺญาย สญฺญากรณียํ กโรตีติ, น เหวํ วตฺตพฺเพ ฯเปฯ.}

\proseitem{382}{\csromansymbolmark{*}มนุสฺเสสุ สญฺญา อตฺถิ, โส จ สญฺญาภโว สญฺญาคติ สญฺญาสตฺตาวาโส สญฺญาสํสาโร สญฺญาโยนิ สญฺญตฺตภาวปฏิลาโภติ, อามนฺตา.\csromanspacer{} อสญฺญสตฺเตสุ สญฺญา อตฺถิ, โส จ สญฺญาภโว สญฺญาคติ สตฺตาวาโส สํสาโร โยนิ อตฺตภาวปฏิลาโภติ, น เหวํ วตฺตพฺเพ ฯเปฯ.}

\prose{\csromansymbolmark{*}มนุสฺเสสุ สญฺญา อตฺถิ, โส จ ปญฺจโวการภโว คติ สตฺตาวาโส สํสาโร โยนิ อตฺตภาวปฏิลาโภติ, อามนฺตา.\csromanspacer{} อสญฺญสตฺเตสุ สญฺญา อตฺถิ, โส จ ปญฺจโวการภโว คติ สตฺตาวาโส สํสาโร โยนิ อตฺตภาวปฏิลาโภติ, น เหวํ วตฺตพฺเพ ฯเปฯ.}

\prose{\csromansymbolmark{*}มนุสฺเสสุ สญฺญา อตฺถิ, ตาย สญฺญาย สญฺญากรณียํ กโรตีติ, อามนฺตา.\csromanspacer{} อสญฺญสตฺเตสุ สญฺญา อตฺถิ, ตาย สญฺญาย สญฺญากรณียํ กโรตีติ, น เหวํ วตฺตพฺเพ ฯเปฯ.}

\prose{\csromansymbolmark{*}อสญฺญสตฺเตสุ สญฺญา อตฺถิ, โส จ อสญฺญภโว อสญฺญคติ อสญฺญสตฺตาวาโส อสญฺญสํสาโร อสญฺญโยนิ อสญฺญตฺตภาวปฏิลาโภติ, อามนฺตา.\csromanspacer{} มนุสฺเสสุ สญฺญา อตฺถิ, โส จ อสญฺญภโว อสญฺญคติ ฯเปฯ อสญฺญตฺตภาวปฏิลาโภติ, น เหวํ วตฺตพฺเพ ฯเปฯ.}

\prose{\csromansymbolmark{*}อสญฺญสตฺเตสุ สญฺญา อตฺถิ, โส จ เอกโวการภโว คติ สตฺตาวาโส สํสาโร โยนิ อตฺตภาวปฏิลาโภติ, อามนฺตา.\csromanspacer{} มนุสฺเสสุ สญฺญา อตฺถิ, โส จ เอกโวการภโว คติ ฯเปฯ อตฺตภาวปฏิลาโภติ, น เหวํ วตฺตพฺเพ ฯเปฯ.}

\csromanmatikaanchor{mtk.53}
\csromantocmarkref{subsection}{(32) 12. เนวสญฺญานาสญฺญายตนกถา}{mtk.53}
\bookmark[dest={mtk.53},level=2]{(32) 12. เนวสญฺญานาสญฺญายตนกถา}
\prose{\csromanfolio{198}\csromansymbolmark{*}อสญฺญสตฺเตสุ สญฺญา อตฺถิ, น จ ตาย สญฺญาย สญฺญากรณียํ กโรตีติ, อามนฺตา.\csromanspacer{} มนุสฺเสสุ สญฺญา อตฺถิ, น จ ตาย สญฺญาย สญฺญากรณียํ กโรตีติ, น เหวํ วตฺตพฺเพ ฯเปฯ.}

\proseitem{383}{\csromansymbolmark{+}น วตฺตพฺพํ “อสญฺญสตฺเตสุ สญฺญา อตฺถี”ติ, อามนฺตา.\csromanspacer{} นนุ วุตฺตํ ภควตา “สนฺติ ภิกฺขเว อสญฺญสตฺตา นาม เทวา, สญฺญุปฺปาทา จ ปน เต เทวา ตมฺหา กายา จวนฺตี”ติ\footnote{ที 1. 27 ปิฏฺเฐ.} อตฺเถว สุตฺตนฺโตติ, อามนฺตา.\csromanspacer{} เตน หิ อสญฺญสตฺเตสุ สญฺญา อตฺถีติ.}

\prose{\csromansymbolmark{*}อสญฺญสตฺเตสุ สญฺญา อตฺถีติ? กิญฺจิ\footnote{กญฺจิ (สฺยา)} กาเล อตฺถิ, กิญฺจิ กาเล นตฺถีติ.\csromanspacer{} กิญฺจิ กาเล สญฺญสตฺตา กิญฺจิ กาเล อสญฺญสตฺตา กิญฺจิ กาเล สญฺญภโว กิญฺจิ กาเล อสญฺญภโว กิญฺจิ กาเล ปญฺจโวการภโว กิญฺจิ กาเล เอกโวการภโวติ, น เหวํ วตฺตพฺเพ ฯเปฯ.}

\prose{\csromansymbolmark{*}อสญฺญสตฺเตสุ สญฺญา กิญฺจิ กาเล อตฺถิ, กิญฺจิ กาเล นตฺถีติ, อามนฺตา.\csromanspacer{} กํ กาลํ อตฺถิ, กํ กาลํ นตฺถีติ, จุติกาเล อุปปตฺติกาเล อตฺถิ, ฐิติกาเล นตฺถีติ.\csromanspacer{} จุติกาเล อุปปตฺติกาเล สญฺญสตฺตา, ฐิติกาเล อสญฺญสตฺตา.\csromanspacer{} จุติกาเล อุปปตฺติกาเล สญฺญภโว, ฐิติกาเล อสญฺญภโว.\csromanspacer{} จุติกาเล อุปปตฺติกาเล ปญฺจโวการภโว, ฐิติกาเล เอกโวการภโวติ, น เหวํ วตฺตพฺเพ ฯเปฯ.}

\nitthitam{อสญฺญกถา นิฏฺฐิตา.}
\chapterhead{3. ตติยวคฺค}
\vspace{\tipitakaparskip}
\csromancenter{\textbf{(32) 12.}\csromanspacer{}\textbf{ เนวสญฺญานาสญฺญายตนกถา}}
\proseitem{384}{\csromansymbolmark{*}เนวสญฺญานาสญฺญายตเน น วตฺตพฺพํ “สญฺญา อตฺถี”ติ, อามนฺตา.\csromanspacer{} อสญฺญภโว อสญฺญคติ อสญฺญสตฺตาวาโส อสญฺญสํสาโร อสญฺญโยนิ อสญฺญตฺตภาวปฏิลาโภติ, เน เหวํ วตฺตพฺเพ ฯเปฯ.}

\chapterheadpageread{3. ตติยวคฺค}
\prose{\csromanfolio{199}นนุ สญฺญาภโว สญฺญาคติ สญฺญาสตฺตาวาโส สญฺญาสํสาโร สญฺญาโยนิ สญฺญตฺตภาวปฏิลาโภติ, อามนฺตา.\csromanspacer{} หญฺจิ สญฺญาภโว สญฺญาคติ ฯเปฯ สญฺญตฺตภาวปฏิลาโภ, โน จ วต เร วตฺตพฺเพ “เนวสญฺญานาสญฺญายตเน น วตฺตพฺพํ ‘สญฺญา อตฺถี’ติ”.}

\prose{\csromansymbolmark{*}เนวสญฺญานาสญฺญายตเน น วตฺตพฺพํ “สญฺญา อตฺถี”ติ, อามนฺตา.\csromanspacer{} เอกโวการภโว คติ ฯเปฯ อตฺตภาวปฏิลาโภติ, น เหวํ วตฺตพฺเพ ฯเปฯ.}

\prose{นนุ จตุโวการภโว คติ ฯเปฯ อตฺตภาวปฏิลาโภติ, อามนฺตา.\csromanspacer{} หญฺจิ จตุโวการภโว คติ ฯเปฯ อตฺตภาวปฏิลาโภ, โน จ วต เร วตฺตพฺเพ “เนวสญฺญานาสญฺญายตเน น วตฺตพฺพํ ‘สญฺญา อตฺถี’ติ”.}

\proseitem{385}{\csromansymbolmark{*}อสญฺญสตฺเตสุ น วตฺตพฺพํ “สญฺญา อตฺถิ”, โส จ อสญฺญภโว อสญฺญคติ อสญฺญสตฺตาวาโส อสญฺญสํสาโร อสญฺญโยนิ อสญฺญตฺตภาวปฏิลาโภติ, อามนฺตา.\csromanspacer{} เนวสญฺญานาสญฺญายตเน น วตฺตพฺพํ “สญฺญา อตฺถิ”, โส จ อสญฺญภโว อสญฺญคติ อสญฺญสตฺตาวาโส อสญฺญสํสาโร อสญฺญโยนิ อสญฺญตฺตภาวปฏิลาโภติ, น เหวํ วตฺตพฺเพ ฯเปฯ.}

\prose{\csromansymbolmark{*}อสญฺญสตฺเตสุ น วตฺตพฺพํ “สญฺญา อตฺถิ”, โส จ เอกโวการภโว คติ ฯเปฯ อตฺตภาวปฏิลาโภติ, อามนฺตา.\csromanspacer{} เนวสญฺญานาสญฺญายตเน น วตฺตพฺพํ “สญฺญา อตฺถิ”, โส จ เอกโวการภโว คติ สตฺตาวาโส สํสาโร โยนิ อตฺตภาวปฏิลาโภติ, น เหวํ วตฺตพฺเพ ฯเปฯ.}

\prose{\csromansymbolmark{*}เนวสญฺญานาสญฺญายตเน น วตฺตพฺพํ “สญฺญา อตฺถิ”, โส จ สญฺญาภโว สญฺญาคติ ฯเปฯ สญฺญตฺตภาวปฏิลาโภติ, อามนฺตา.\csromanspacer{} อสญฺญสตฺเตสุ น วตฺตพฺพํ “สญฺญา อตฺถิ”, โส จ สญฺญาภโว สญฺญาคติ ฯเปฯ สญฺญตฺตภาวปฏิลาโภติ, น เหวํ วตฺตพฺเพ ฯเปฯ.}

\prose{\csromansymbolmark{*}เนวสญฺญานาสญฺญายตเน น วตฺตพฺพํ “สญฺญา อตฺถิ”, โส จ จตุโวการภโว คติ ฯเปฯ อตฺตภาวปฏิลาโภติ, อามนฺตา.\csromanspacer{} อสญฺญสตฺเตสุ น วตฺตพฺพํ “สญฺญา อตฺถิ”, โส จ จตุโวการภโว คติ ฯเปฯ อตฺตภาวปฏิลาโภติ, น เหวํ วตฺตพฺเพ ฯเปฯ.}

\prose{\csromanfolio{200}\csromansymbolmark{*}เนวสญฺญานาสญฺญายตเน น วตฺตพฺพํ “สญฺญา อตฺถี”ติ, อามนฺตา.\csromanspacer{} นนุ เนวสญฺญานาสญฺญายตนํ จตุโวการภโวติ, อามนฺตา.\csromanspacer{} หญฺจิ เนวสญฺญานาสญฺญายตนํ จตุโวการภโว, โน จ วต เร วตฺตพฺเพ “เนวสญฺญานาสญฺญายตเน น วตฺตพฺพํ ‘สญฺญา อตฺถี’ติ.}

\proseitem{386}{\csromansymbolmark{*}เนวสญฺญานาสญฺญายตนํ จตุโวการภโว, เนวสญฺญานาสญฺญายตเน น วตฺตพฺพํ “สญฺญา อตฺถี”ติ, อามนฺตา.\csromanspacer{} อากาสานญฺจายตนํ จตุโวการภโว, อากาสานญฺจายตเน น วตฺตพฺพํ “สญฺญา อตฺถี”ติ, เน เหวํ วตฺตพฺเพ ฯเปฯ.}

\prose{\csromansymbolmark{*}เนวสญฺญานาสญฺญายตนํ จตุโวการภโว, เนวสญฺญานาสญฺญายตเน น วตฺตพฺพํ “สญฺญา อตฺถี”ติ, อามนฺตา.\csromanspacer{} วิญฺญาณญฺจายตนํ ฯเปฯ.\csromanspacer{} อากิญฺจญฺญายตนํ จตุโวการภโว, อากิญฺจญฺญายตเน น วตฺตพฺพํ “สญฺญา อตฺถี”ติ, น เหวํ วตฺตพฺเพ ฯเปฯ.}

\prose{\csromansymbolmark{*}อากาสานญฺจายตนํ จตุโวการภโว อตฺถิ ตตฺถ สญฺญาติ, อามนฺตา.\csromanspacer{} เนวสญฺญานาสญฺญายตนํ จตุโวการภโว อตฺถิ ตตฺถ สญฺญาติ, น เหวํ วตฺตพฺเพ ฯเปฯ.}

\prose{\csromansymbolmark{*}วิญฺญาณญฺจายตนํ ฯเปฯ.\csromanspacer{} อากิญฺจญฺญายตนํ จตุโวการภโว, อตฺถิ ตตฺถ สญฺญาติ, อามนฺตา.\csromanspacer{} เนวสญฺญานาสญฺญายตนํ จตุโวการภโว, อตฺถิ ตตฺถ สญฺญาติ, น เหวํ วตฺตพฺเพ ฯเปฯ.}

\prose{\csromansymbolmark{*}เนวสญฺญานาสญฺญายตเน น วตฺตพฺพํ “สญฺญา อตฺถี”ติ วา “นตฺถี”ติ วาติ, อามนฺตา.\csromanspacer{} นนุ เนวสญฺญานาสญฺญายตนํ จตุโวการภโวติ, อามนฺตา.\csromanspacer{} หญฺจิ เนวสญฺญานาสญฺญายตนํ จตุโวการภโว, โน จ วต เร วตฺตพฺเพ “เนวสญฺญานาสญฺญายตเน น วตฺตพฺพํ “สญฺญา อตฺถี”ติ วา “นตฺถี”ติ วาติ.}

\prose{\csromansymbolmark{*}เนวสญฺญานาสญฺญายตนํ จตุโวการภโว, เนวสญฺญานาสญฺญายตเน น วตฺตพฺพํ “สญฺญา อตฺถี”ติ วา “นตฺถี”ติ วาติ, อามนฺตา.\csromanspacer{} อากาสานญฺจายตนํ ฯเปฯ.\csromanspacer{} วิญฺญาณญฺจายตนํ ฯเปฯ.\csromanspacer{} อากิญฺจญฺญายตนํ จตุโวการภโว, อากิญฺจญฺญายตเน น วตฺตพฺพํ “สญฺญา อตฺถี”ติ วา “นตฺถี”ติ วาติ, น เหวํ วตฺตพฺเพ ฯเปฯ.}

\chapterheadpageread{3. ตติยวคฺค}
\prose{\csromanfolio{201}\csromansymbolmark{*}อากาสานญฺจายตนํ จตุโวการภโว, อตฺถิ ตตฺถ สญฺญาติ, อามนฺตา.\csromanspacer{} เนวสญฺญานาสญฺญายตนํ จตุโวการภโว, อตฺถิ ตตฺถ สญฺญาติ, น เหวํ วตฺตพฺเพ ฯเปฯ.}

\prose{\csromansymbolmark{*}วิญฺญาณญฺจายตนํ ฯเปฯ.\csromanspacer{} อากิญฺจญฺญายตนํ จตุโวการภโว, อตฺถิ ตตฺถ สญฺญาติ, อามนฺตา.\csromanspacer{} เนวสญฺญานาสญฺญายตนํ จตุโวการภโว, อตฺถิ ตตฺถ สญฺญาติ, น เหวํ วตฺตพฺเพ ฯเปฯ.}

\prose{\csromansymbolmark{+}เนวสญฺญานาสญฺญายตเน น วตฺตพฺพํ\footnote{เนวสญฺญานาสญฺญายตเน วตฺตพฺพํ (?)} “สญฺญา อตฺถี”ติ วา “นตฺถี”ติ วาติ, อามนฺตา.\csromanspacer{} นนุ เนวสญฺญานาสญฺญายตนนฺติ, อามนฺตา.\csromanspacer{} หญฺจิ เนวสญฺญานาสญฺญายตนํ, เตน วต เร วตฺตพฺเพ “เนวสญฺญานาสญฺญายตเน น วตฺตพฺพํ “สญฺญา อตฺถี”ติ วา “นตฺถี”ติ วาติ.}

\prose{\csromansymbolmark{*}เนวสญฺญานาสญฺญายตนนฺติ กตฺวา เนวสญฺญานาสญฺญายตเน น วตฺตพฺพํ “สญฺญา อตฺถี”ติ วา “นตฺถี”ติ วาติ, อามนฺตา.\csromanspacer{} อทุกฺขมสุขา เวทนาติ กตฺวา อทุกฺขมสุขาย เวทนาย\footnote{อทุกฺขมสุขา เวทนา (สี, ก)} น วตฺตพฺพํ “เวทนา”ติ วา “อเวทนา”ติ วาติ, น เหวํ วตฺตพฺเพ ฯเปฯ.\csromanspacer{} เนวสญฺญานาสญฺญายตนกถา นิฏฺฐิตา.\csromanspacer{} ตติโย วคฺโค.}

\titlehead{\textbf{ตสฺสุทฺทานํ}}
\csromangathameasure{พลํ สาธารณํ อริยํ, สราคํ จิตฺตํ วิมุจฺจติ. \\ วิมุตฺตํ วิมุจฺจมานํ, อตฺถิ จิตฺตํ วิมุจฺจมานํ. \\ อฏฺฐมกสฺส ปุคฺคลสฺส, ทิฏฺฐิปริยุฏฺฐานํ ปหีนํ. \\ อฏฺฐมกสฺส ปุคฺคลสฺส, นตฺถิ ปญฺจินฺทฺริยานิ จกฺขุํ. \\ โสตํ ธมฺมุปตฺถทฺธํ, ยถากมฺมูปคตํ ญาณํ. \\ เทเวสุ สํวโร อสญฺญ\csromandash{}สตฺเตสุ สญฺญา เอวเมว ภวคฺคนฺติ.}
\csromangathagroup{\csromangathastanza{พลํ สาธารณํ อริยํ, สราคํ จิตฺตํ วิมุจฺจติ. \\ วิมุตฺตํ วิมุจฺจมานํ, อตฺถิ จิตฺตํ วิมุจฺจมานํ.}\csromangathastanza{อฏฺฐมกสฺส ปุคฺคลสฺส, ทิฏฺฐิปริยุฏฺฐานํ ปหีนํ. \\ อฏฺฐมกสฺส ปุคฺคลสฺส, นตฺถิ ปญฺจินฺทฺริยานิ จกฺขุํ.}\csromangathastanza{โสตํ ธมฺมุปตฺถทฺธํ, ยถากมฺมูปคตํ ญาณํ. \\ เทเวสุ สํวโร อสญฺญ\csromandash{}สตฺเตสุ สญฺญา เอวเมว ภวคฺคนฺติ.}}

\csromanmatikaanchor{mtk.54}
\csromantocmarknopage{chapter}{4. จตุตฺถวคฺค}{mtk.54}
\bookmark[dest={mtk.54},level=0]{4. จตุตฺถวคฺค}
\chapterheadpageread{4. จตุตฺถวคฺค}
\vspace{\tipitakaparskip}
\csromancenter{\textbf{(33) 1.}\csromanspacer{}\textbf{ คิหิสฺส-อรหาติกถา}}
\csromanmatikaanchor{mtk.55}
\csromantocmarkref{subsection}{(33) 1. คิหิสฺส-อรหาติกถา}{mtk.55}
\bookmark[dest={mtk.55},level=2]{(33) 1. คิหิสฺส-อรหาติกถา}
\proseitem{387}{\csromanfolio{202}\csromansymbolmark{*}คิหิ’สฺส อรหาติ, อามนฺตา.\csromanspacer{} อตฺถิ อรหโต คิหิสํโยชนนฺติ, น เหวํ วตฺตพฺเพ ฯเปฯ.\csromanspacer{} นตฺถิ อรหโต คิหิสํโยชนนฺติ, อามนฺตา.\csromanspacer{} หญฺจิ นตฺถิ อรหโต คิหิสํโยชนํ, โน จ วต เร วตฺตพฺเพ “คิหิ’สฺส อรหา”ติ.}

\prose{\csromansymbolmark{*}คิหิ’สฺส อรหาติ, อามนฺตา.\csromanspacer{} นนุ อรหโต คิหิสํโยชนํ ปหีนํ อุจฺฉินฺนมูลํ ตาลาวตฺถุกตํ อนภาวํกตํ อายติํ อนุปฺปาทธมฺมนฺติ, อามนฺตา.\csromanspacer{} หญฺจิ อรหโต คิหิสํโยชนํ ปหีนํ อุจฺฉินฺนมูลํ ตาลาวตฺถุกตํ อนภาวํกตํ อายติํ อนุปฺปาทธมฺมํ, โน จ วต เร วตฺตพฺเพ “คิหิ’สฺส อรหา”ติ.}

\prose{\csromansymbolmark{*}คิหิ’สฺส อรหาติ, อามนฺตา.\csromanspacer{} อตฺถิ โกจิ คิหี คิหิสํโยชนํ อปฺปหาย ทิฏฺเฐว ธมฺเม ทุกฺขสฺสนฺตกโรติ, นตฺถิ.\csromanspacer{} หญฺจิ นตฺถิ โกจิ คิหี คิหิสํโยชนํ อปฺปหาย ทิฏฺเฐว ธมฺเม ทุกฺขสฺสนฺตกโร, โน จ วต เร วตฺตพฺเพ “คิหิ’สฺส อรหา”ติ.}

\prose{\csromansymbolmark{*}คิหิ’สฺส อรหาติ, อามนฺตา.\csromanspacer{} นนุ วจฺฉโคตฺโต ปริพฺพาชโก ภควนฺตํ เอตทโวจ\csromandash{}“อตฺถิ นุ โข โภ โคตม โกจิ คิหี คิหิสํโยชนํ อปฺปหาย กายสฺส เภทา ทุกฺขสฺสนฺตกโร”ติ.\csromanspacer{} นตฺถิ โข วจฺฉ โกจิ คิหี คิหิสํโยชนํ อปฺปหาย กายสฺส เภทา ทุกฺขสฺสนฺตกโรติ\footnote{ม 2. 149 ปิฏฺเฐ.} อตฺเถว สุตฺตนฺโตติ, อามนฺตา.\csromanspacer{} เตน หิ น วตฺตพฺพํ “คิหิ’สฺส อรหา”ติ.}

\prose{\csromansymbolmark{*}คิหิ’สฺส อรหาติ, อามนฺตา.\csromanspacer{} อรหา เมถุนํ ธมฺมํ ปฏิเสเวยฺย, เมถุนํ อุปฺปาเทยฺย, ปุตฺตสมฺพาธสยนํ อชฺฌาวเสยฺย, กาสิกจนฺทนํ ปจฺจนุภเวยฺย, มาลาคนฺธวิเลปนํ ธาเรยฺย, ชาตรูปรชตํ สาทิเยยฺย, อเชฬกํ ปฏิคฺคเณฺหยฺย, กุกฺกุฏสูกรํ ปฏิคฺคเณฺหยฺย, หตฺถิควสฺสวฬวํ ปฏิคฺคเณฺหยฺย, ติตฺติรวฏฺฏกโมรกปิญฺชรํ\footnote{...กปิญฺชลํ (สฺยา, กํ, อิ)} ปฏิคฺคเณฺหยฺย, จิตฺตวณฺฑวาลโมฬิํ\footnote{ปีตวณฺฏวาลโมฬิกํ (สฺยา, กํ, อิ)} ธาเรยฺย, โอทาตานิ วตฺถานิ ทีฆทสานิ ธาเรยฺย, ยาวชีวํ อคาริยภูโต อสฺสาติ, น เหวํ วตฺตพฺเพ ฯเปฯ.}

\chapterheadpageread{4. จตุตฺถวคฺค}
\csromanmatikaanchor{mtk.56}
\csromantocmarkref{subsection}{(34) 2. อุปปตฺติกถา}{mtk.56}
\bookmark[dest={mtk.56},level=2]{(34) 2. อุปปตฺติกถา}
\prose{\csromanfolio{203}\csromansymbolmark{+}น วตฺตพฺพํ “คิหิ’สฺส อรหา”ติ, อามนฺตา.\csromanspacer{} นนุ ยโส กุลปุตฺโต, อุตฺติโย คหปติ, เสตุ มาณโว, คิหิพฺยญฺชเนน อรหตฺตํ ปตฺโตติ, อามนฺตา.\csromanspacer{} หญฺจิ ยโส กุลปุตฺโต, อุตฺติโย คหปติ, เสตุ มาณโว, คิหิพฺยญฺชเนน อรหตฺตํ ปตฺโต, เตน วต เร วตฺตพฺเพ “คิหิ’สฺส อรหา”ติ.}

\nitthitam{คิหิสฺส อรหาติกถา นิฏฺฐิตา.}
\chapterhead{4. จตุตฺถวคฺค}
\vspace{\tipitakaparskip}
\csromancenter{\textbf{(34) 2.}\csromanspacer{}\textbf{ อุปปตฺติกถา}}
\proseitem{388}{\csromansymbolmark{*}สห อุปปตฺติยา อรหาติ, อามนฺตา.\csromanspacer{} สห อุปปตฺติยา โสตาปนฺโน โหตีติ, น เหวํ วตฺตพฺเพ ฯเปฯ.}

\prose{\csromansymbolmark{*}สห อุปปตฺติยา อรหาติ, อามนฺตา.\csromanspacer{} สห อุปปตฺติยา สกทาคามี โหตีติ, น เหวํ วตฺตพฺเพ ฯเปฯ.}

\prose{\csromansymbolmark{*}สห อุปปตฺติยา อรหาติ, อามนฺตา.\csromanspacer{} สห อุปปตฺติยา อนาคามี โหตีติ, น เหวํ วตฺตพฺเพ ฯเปฯ.}

\prose{\csromansymbolmark{*}สห อุปปตฺติยา โสตาปนฺโน น โหตีติ, อามนฺตา.\csromanspacer{} หญฺจิ สห อุปปตฺติยา โสตาปนฺโน น โหติ, โน จ วต เร วตฺตพฺเพ “สห อุปปตฺติยา อรหา”ติ.}

\prose{\csromansymbolmark{*}สห อุปปตฺติยา สกทาคามี น โหตีติ, อามนฺตา.\csromanspacer{} หญฺจิ สห อุปปตฺติยา สกทาคามี น โหติ, โน จ วต เร วตฺตพฺเพ “สห อุปปตฺติยา อรหา”ติ.}

\prose{\csromansymbolmark{*}สห อุปปตฺติยา อนาคามี น โหตีติ, อามนฺตา.\csromanspacer{} หญฺจิ สห อุปปตฺติยา อนาคามี น โหติ, โน จ วต เร วตฺตพฺเพ “สห อุปปตฺติยา อรหา”ติ.}

\proseitem{389}{\csromansymbolmark{*}สห อุปปตฺติยา อรหาติ, อามนฺตา.\csromanspacer{} สาริปุตฺโต เถโร สห อุปปตฺติยา อรหาติ, น เหวํ วตฺตพฺเพ.\csromanspacer{} มหาโมคฺคลฺลาโน \csromanfolio{204}เถโร ฯเปฯ.\csromanspacer{} มหากสฺสโป เถโร ฯเปฯ.\csromanspacer{} มหากจฺจาโน เถโร ฯเปฯ.\csromanspacer{} มหาโกฏฺฐิโก เถโร ฯเปฯ.\csromanspacer{} มหาปนฺถโก เถโร สห อุปปตฺถิยา อรหาติ, น เหวํ วตฺตพฺเพ ฯเปฯ.}

\prose{\csromansymbolmark{*}สาริปุตฺโต เถโร น สห อุปปตฺติยา อรหาติ, อามนฺตา.\csromanspacer{} หญฺจิ สาริปุตฺโต เถโร น สห อุปปตฺติยา อรหา, โน จ วต เร วตฺตพฺเพ “สห อุปปตฺติยา อรหา”ติ.}

\prose{\csromansymbolmark{*}มหาโมคฺคลฺลาโน เถโร ฯเปฯ.\csromanspacer{} มหากสฺสโป เถโร ฯเปฯ.\csromanspacer{} มหากจฺจาโน เถโร ฯเปฯ.\csromanspacer{} มหาโกฏฺฐิโก เถโร ฯเปฯ.\csromanspacer{} มหาปนฺถโก เถโร น สห อุปปตฺติยา อรหาติ, อามนฺตา.\csromanspacer{} หญฺจิ มหาปนฺถโก เถโร น สห อุปปตฺติยา อรหา, โน จ วต เร วตฺตพฺเพ “สห อุปปตฺติยา อรหา”ติ.}

\proseitem{390}{\csromansymbolmark{*}สห อุปปตฺติยา อรหาติ, อามนฺตา.\csromanspacer{} อุปปตฺเตสิเยน จิตฺเตน อรหตฺตํ สจฺฉิกโรติ โลกิเยน สาสเวน ฯเปฯ สํกิเลสิเยนาติ, น เหวํ วตฺตพฺเพ ฯเปฯ.}

\prose{\csromansymbolmark{*}สห อุปปตฺติยา อรหาติ, อามนฺตา.\csromanspacer{} อุปปตฺเตสิยํ จิตฺตํ นิยฺยานิกํ ขยคามี โพธคามี อปจยคามี อนาสวํ ฯเปฯ อสํกิเลสิยนฺติ, น เหวํ วตฺตพฺเพ ฯเปฯ.}

\prose{\csromansymbolmark{*}นนุ อุปปตฺเตสิยํ จิตฺตํ อนิยฺยานิกํ น ขยคามิ น โพธคามิ น อปจยคามิ สาสวํ ฯเปฯ สํกิเลสิยนฺติ, อามนฺตา.\csromanspacer{} หญฺจิ อุปปตฺเตสิยํ จิตฺตํ อนิยฺยานิกํ น ขยคามิ น โพธคามิ น อปจยคามิ สาสวํ ฯเปฯ สํกิเลสิยํ, โน จ วต เร วตฺตพฺเพ “สห อุปปตฺติยา อรหา”ติ.}

\prose{\csromansymbolmark{*}สห อุปปตฺติยา อรหาติ, อามนฺตา.\csromanspacer{} อุปปตฺเตสิเยน จิตฺเตน ราคํ ปชหติ.\csromanspacer{} โทสํ ปชหติ.\csromanspacer{} โมหํ ปชหติ.\csromanspacer{} อโนตฺตปฺปํ ปชหตีติ, น เหวํ วตฺตพฺเพ ฯเปฯ.}

\prose{\csromansymbolmark{*}สห อุปปตฺติยา อรหาติ, อามนฺตา.\csromanspacer{} อุปปตฺเตสิยํ จิตฺตํ มคฺโค.\csromanspacer{} สติปฏฺฐานํ ฯเปฯ.\csromanspacer{} สมฺมปฺปธานํ.\csromanspacer{} อิทฺธิปาโท.\csromanspacer{} อินฺทฺริยํ.\csromanspacer{} พลํ.\csromanspacer{} โพชฺฌงฺโคติ, น เหวํ วตฺตพฺเพ ฯเปฯ.}

\prose{\csromansymbolmark{*}สห อุปปตฺติยา อรหาติ, อามนฺตา.\csromanspacer{} อุปปตฺเตสิเยน จิตฺเตน ทุกฺขํ ปริชานาติ, สมุทยํ ปชหติ, นิโรธํ สจฺฉิกโรติ, มคฺคํ ภาเวตีติ,}

\vspace{\tipitakaparskip}
\csromancenter{จตุตฺถวคฺค น เหวํ วตฺตพฺเพ ฯเปฯ สห อุปปตฺติยา อรหาติ, อามนฺตา.\csromanspacer{} จุติจิตฺตํ มคฺคจิตฺตํ อุปปตฺเตสิยํ จิตฺตํ ผลจิตฺตนฺติ, น เหวํ วตฺตพฺเพ ฯเปฯ.}
\nitthitam{อุปปตฺติกถา นิฏฺฐิตา.}
\chapterhead{4. จตุตฺถวคฺค}
\vspace{\tipitakaparskip}
\csromancenter{\textbf{(35) 3.}\csromanspacer{}\textbf{ อนาสวกถา}}
\csromanmatikaanchor{mtk.57}
\csromantocmarkref{subsection}{(35) 3. อนาสวกถา}{mtk.57}
\bookmark[dest={mtk.57},level=2]{(35) 3. อนาสวกถา}
\proseitem{391}{\csromanfolio{205}\csromansymbolmark{*}อรหโต สพฺเพ ธมฺมา อนาสวาติ, อามนฺตา.\csromanspacer{} มคฺโค.\csromanspacer{} ผลํ.\csromanspacer{} นิพฺพานํ.\csromanspacer{} โสตาปตฺติมคฺโค.\csromanspacer{} โสตาปตฺติผลํ.\csromanspacer{} สกทาคามิมคฺโค.\csromanspacer{} สกทาคามิผลํ.\csromanspacer{} อนาคามิมคฺโค.\csromanspacer{} อนาคามิผลํ.\csromanspacer{} อรหตฺตมคฺโค.\csromanspacer{} อรหตฺตผลํ.\csromanspacer{} สติปฏฺฐานํ.\csromanspacer{} สมฺมปฺปธานํ.\csromanspacer{} อิทฺธิปาโท.\csromanspacer{} อินฺทฺริยํ.\csromanspacer{} พลํ.\csromanspacer{} โพชฺฌงฺโคติ, น เหวํ วตฺตพฺเพ ฯเปฯ.}

\prose{อรหโต สพฺเพ ธมฺมา อนาสวาติ, อามนฺตา.\csromanspacer{} อรหโต จกฺขุํ อนาสวนฺติ, น เหวํ วตฺตพฺเพ ฯเปฯ.\csromanspacer{} อรหโต จกฺขุํ อนาสวนฺติ, อามนฺตา.\csromanspacer{} มคฺโค.\csromanspacer{} ผลํ.\csromanspacer{} นิพฺพานํ.\csromanspacer{} โสตาปตฺติมคฺโค.\csromanspacer{} โสตาปตฺติผลํ ฯเปฯ.\csromanspacer{} โพชฺฌงฺโคติ, น เหวํ วตฺตพฺเพ ฯเปฯ.}

\prose{\csromansymbolmark{*}อรหโต โสตํ ฯเปฯ.\csromanspacer{} อรหโต ฆานํ.\csromanspacer{} อรหโต ชิวฺหา.\csromanspacer{} อรหโต กาโย อนาสโวติ, น เหวํ วตฺตพฺเพ ฯเปฯ.\csromanspacer{} อรหโต กาโย อนาสโวติ, อามนฺตา.\csromanspacer{} มคฺโค.\csromanspacer{} ผลํ.\csromanspacer{} นิพฺพานํ.\csromanspacer{} โสตาปตฺติมคฺโค.\csromanspacer{} โสตาปตฺติผลํ ฯเปฯ.\csromanspacer{} โพชฺฌงฺโคติ, น เหวํ วตฺตพฺเพ ฯเปฯ.}

\prose{\csromansymbolmark{*}อรหโต กาโย อนาสโวติ, อามนฺตา.\csromanspacer{} อรหโต กาโย ปคฺคหนิคฺคหุปโค เฉทนเภทนุปโค กาเกหิ คิชฺเฌหิ กุลเลหิ สาธารโณติ, อามนฺตา.\csromanspacer{} อนาสโว ธมฺโม ปคฺคหนิคฺคหุปโค เฉทนเภทนุปโค กาเกหิ คิชฺเฌหิ กุลเลหิ สาธารโณติ, น เหวํ วตฺตพฺเพ ฯเปฯ.}

\prose{\csromansymbolmark{*}อรหโต กาเย วิสํ กเมยฺย, สตฺถํ กเมยฺย, อคฺคิ กเมยฺยาติ, อามนฺตา.\csromanspacer{} อนาสเว ธมฺเม วิสํ กเมยฺย, สตฺถํ กเมยฺย, อคฺคิ กเมยฺยาติ, น เหวํ วตฺตพฺเพ ฯเปฯ.}

\prose{\csromanfolio{206}\csromansymbolmark{*}ลพฺภา อรหโต กาโย อทฺทุพนฺธเนน พนฺธิตุํ, รชฺชุพนฺธเนน พนฺธิตุํ, สงฺขลิกพนฺธเนน พนฺธิตุํ, คามพนฺธเนน พนฺธิตุํ, นิคมพนฺธเนน พนฺธิตุํ, นครพนฺธเนน พนฺธิตุํ, ชนปทพนฺธเนน พนฺธิตุํ, กณฺฐปญฺจเมหิ พนฺธเนหิ พนฺธิตุนฺติ, อามนฺตา.\csromanspacer{} ลพฺภา อนาสโว ธมฺโม อทฺทุพนฺธเนน พนฺธิตุํ, รชฺชุพนฺธเนน พนฺธิตุํ, สงฺขลิกพนฺธเนน พนฺธิตุํ, คาม นิคม นคร ชนปท พนฺธเนน พนฺธิตุํ, กณฺฐปญฺจเมหิ พนฺธเนหิ พนฺธิตุนฺติ, น เหวํ วตฺตพฺเพ ฯเปฯ.}

\proseitem{392}{\csromansymbolmark{*}ยทิ อรหา ปุถุชฺชนสฺส จีวรํ เทติ, อนาสวํ หุตฺวา สาสวํ โหตีติ, น เหวํ วตฺตพฺเพ ฯเปฯ อนาสวํ หุตฺวา สาสวํ โหตีติ, อามนฺตา.\csromanspacer{} ตญฺเญว อนาสวํ ตํ สาสวนฺติ, น เหวํ วตฺตพฺเพ ฯเปฯ.\csromanspacer{} ตญฺเญว อนาสวํ ตํ สาสวนฺติ, อามนฺตา.\csromanspacer{} มคฺโค อนาสโว หุตฺวา สาสโว โหตีติ, น เหวํ วตฺตพฺเพ ฯเปฯ.\csromanspacer{} ผลํ.\csromanspacer{} สติปฏฺฐานํ.\csromanspacer{} สมฺมปฺปธานํ.\csromanspacer{} อิทฺธิปาโท อินฺทฺริยํ.\csromanspacer{} พลํ.\csromanspacer{} โพชฺฌงฺโค อนาสโวหุตฺวา สาสโว โหตีติ, น เหวํ วตฺตพฺเพ ฯเปฯ.}

\prose{\csromansymbolmark{*}ยทิ อรหา ปุถุชฺชนสฺส ปิณฺฑปาตํ เทติ, เสนาสนํ เทติ, คิลานปจฺจย เภสชฺชปริกฺขารํ เทติ, อนาสโว หุตฺวา สาสโว โหตีติ, น เหวํ วตฺตพฺเพ ฯเปฯ.\csromanspacer{} อานาสโว หุตฺวา สาสโว โหตีติ, อามนฺตา.\csromanspacer{} ตญฺเญว อนาสวํ ตํ สาสวนฺติ, น เหวํ วตฺตพฺเพ ฯเปฯ.\csromanspacer{} ตญฺเญว อนาสวํ ตํ สาสวนฺติ, อามนฺตา.\csromanspacer{} มคฺโค อนาสโว หุตฺวา สาสโว โหตีติ, น เหวํ วตฺตพฺเพ ฯเปฯ.\csromanspacer{} ผลํ.\csromanspacer{} สติปฏฺฐานํ.\csromanspacer{} สมฺมปฺปธานํ.\csromanspacer{} อิทฺธิปาโท.\csromanspacer{} อินฺทฺริยํ.\csromanspacer{} พลํ.\csromanspacer{} โพชฺฌงฺโค อนาสโว หุตฺวา สาสโว โหตีติ, น เหวํ วตฺตพฺเพ ฯเปฯ.}

\prose{\csromansymbolmark{*}ยทิ ปุถุชฺชโน อรหโต จีวรํ เทติ, สาสวํ หุตฺวา อนาสวํ โหตีติ, น เหวํ วตฺตพฺเพ ฯเปฯ.\csromanspacer{} สาสวํ หุตฺวา อนาสวํ โหตีติ, อามนฺตา.\csromanspacer{} ตญฺเญว สาสวํ ตํ อนาสวนฺติ, น เหวํ วตฺตพฺเพ ฯเปฯ.\csromanspacer{} ตญฺเญว สาสวํ ตํ อนาสวนฺติ, อามนฺตา.\csromanspacer{} ราโค สาสโว หุตฺวา อนาสโว โหตีติ, น เหวํ วตฺตพฺเพ ฯเปฯ.\csromanspacer{} โทโส ฯเปฯ.\csromanspacer{} โมโห ฯเปฯ.\csromanspacer{} อโนตฺตปฺปํ สาสวํ หุตฺวา อนาสวํ โหตีติ, น เหวํ วตฺตพฺเพ ฯเปฯ.}

\prose{\csromansymbolmark{*}ยทิ ปุถุชฺชโน อรหโต ปิณฺฑปาตํ เทติ, เสนาสนํ เทติ, คิลานปจฺจยเภสชฺชปริกฺขารํ เทติ, สาสโว หุตฺวา อนาสโว โหตีติ, น เหวํ วตฺตพฺเพ ฯเปฯ.\csromanspacer{} สาสโว หุตฺวา อนาสโว โหตีติ อามนฺตา.\csromanspacer{} ตญฺเญว สาสวํ ตํ อนาสวนฺติ, น เหวํ วตฺตพฺเพ ฯเปฯ.\csromanspacer{} ตญฺเญว}

\vspace{\tipitakaparskip}
\csromancenter{จตุตฺถวคฺค สาสวํ ตํ อนาสวนฺติ, อามนฺตา.\csromanspacer{} ราโค สาสโว หุตฺวา อนาสโว โหตีติ, น เหวํ วตฺตพฺเพ ฯเปฯ.\csromanspacer{} โทโส ฯเปฯ.\csromanspacer{} โมโห ฯเปฯ.\csromanspacer{} อโนตฺตปฺปํ สาสวํ หุตฺวา อนาสวํ โหตีติ, น เหวํ วตฺตพฺเพ ฯเปฯ.}
\csromanmatikaanchor{mtk.58}
\csromantocmarkref{subsection}{(36) 4. สมนฺนาคตกถา}{mtk.58}
\bookmark[dest={mtk.58},level=2]{(36) 4. สมนฺนาคตกถา}
\prose{\csromanfolio{207}\csromansymbolmark{+}น วตฺตพฺพํ “อรหโต สพฺเพ ธมฺมา อนาสวา”ติ, อามนฺตา.\csromanspacer{} นนุ อรหา อนาสโวติ, อามนฺตา.\csromanspacer{} หญฺจิ อรหา อนาสโว, เตน วต เร วตฺตพฺเพ “อรหโต สพฺเพ ธมฺมา อนาสวา”ติ.}

\nitthitam{อนาสวกถา นิฏฺฐิตา.}
\chapterhead{4. จตุตฺถวคฺค}
\vspace{\tipitakaparskip}
\csromancenter{\textbf{(36) 4.}\csromanspacer{}\textbf{ สมนฺนาคตกถา}}
\proseitem{393}{\csromansymbolmark{*}อรหา จตูหิ ผเลหิ สมนฺนาคโตติ, อามนฺตา.\csromanspacer{} อรหา จตูหิ ผสฺเสหิ.\csromanspacer{} จตูหิ เวทนาหิ.\csromanspacer{} จตูหิ สญฺญาหิ.\csromanspacer{} จตูหิ เจตนาหิ.\csromanspacer{} จตูหิ จิตฺเตหิ.\csromanspacer{} จตูหิ สทฺธาหิ.\csromanspacer{} จตูหิ วีริเยหิ.\csromanspacer{} จตูหิ สตีติ.\csromanspacer{} จตูหิ สมาธีหิ.\csromanspacer{} จตูหิ ปญฺญาหิ สมนฺนาคโตติ, น เหวํ วตฺตพฺเพ ฯเปฯ.}

\prose{\csromansymbolmark{*}อนาคามี ตีหิ ผเลหิ สมนฺนาคโตติ, อามนฺตา.\csromanspacer{} อนาคามี ตีหิ ผสฺเสหิ ฯเปฯ.\csromanspacer{} ตีหิ ปญฺญาหิ สมนฺนาคโตติ, น เหวํ วตฺตพฺเพ ฯเปฯ.}

\prose{\csromansymbolmark{*}สกทาคามี ทฺวีหิ ผเลหิ สมนฺนาคโตติ, อามนฺตา.\csromanspacer{} สกทาคามี ทฺวีหิ ผสฺเสหิ ฯเปฯ.\csromanspacer{} ทฺวีหิ ปญฺญาหิ สมนฺนาคโตติ, น เหวํ วตฺตพฺเพ ฯเปฯ.}

\prose{\csromansymbolmark{*}อรหา โสตาปตฺติผเลน สมนฺนาคโตติ, อามนฺตา.\csromanspacer{} อรหา โสตาปนฺโน สตฺตกฺขตฺตุปรโม โกลํโกโล เอกพีชีติ, น เหวํ วตฺตพฺเพ ฯเปฯ.\csromanspacer{} อรหา สกทาคามิผเลน สมนฺนาคโตติ, อามนฺตา.\csromanspacer{} อรหา สกทาคามีติ, น เหวํ วตฺตพฺเพ ฯเปฯ.\csromanspacer{} อรหา อนาคามิผเลน สมนฺนาคโตติ, อามนฺตา.\csromanspacer{} อรหา อนาคามี อนฺตราปรินิพฺพายี, อุปหจฺจปรินิพฺพายี, อสงฺขารปรินิพฺพายี, สสงฺขารปรินิพฺพายี, อุทฺธํโสโต อกนิฏฺฐคามีติ, น เหวํ วตฺตพฺเพ ฯเปฯ.}

\prose{\csromansymbolmark{*}อนาคามี โสตาปตฺติผเลน สมนฺนาคโตติ, อามนฺตา.\csromanspacer{} อนาคามี โสตาปนฺโน สตฺตกฺขตฺตุปรโม, โกลํโกโล, เอกพีชีติ, น เหวํ \csromanfolio{208}วตฺตพฺเพ ฯเปฯ.\csromanspacer{} อนาคามี สกทาคามิผเลน สมนฺนาคโตติ, อามนฺตา.\csromanspacer{} อนาคามี สกทาคามีติ, น เหวํ วตฺตพฺเพ ฯเปฯ.}

\prose{\csromansymbolmark{*}สกทาคามี โสตาปตฺติผเลน สมนฺนาคโตติ, อามนฺตา.\csromanspacer{} สกทาคามี โสตาปนฺโน สตฺตกฺขตฺตุปรโม, โกลํโกโล, เอกพีชีติ, น เหวํ วตฺตพฺเพ ฯเปฯ.}

\proseitem{394}{\csromansymbolmark{*}โสตาปตฺติผเลน สมนฺนาคโต “โสตาปนฺโน”ติ วตฺตพฺโพติ, อามนฺตา.\csromanspacer{} อรหา โสตาปตฺติผเลน สมนฺนาคโตติ, อามนฺตา.\csromanspacer{} เสฺวว อรหา, โส โสตาปนฺโนติ, น เหวํ วตฺตพฺเพ ฯเปฯ.}

\prose{\csromansymbolmark{*}สกทาคามิผเลน สมนฺนาคโต “สกทาคามี”ติ วตฺตพฺโพติ, อามนฺตา.\csromanspacer{} อรหา สกทาคามิผเลน สมนฺนาคโตติ, อามนฺตา.\csromanspacer{} เสฺวว อรหา, โส สกทาคามีติ, น เหวํ วตฺตพฺเพ ฯเปฯ.}

\prose{\csromansymbolmark{*}อนาคามิผเลน สมนฺนาคโต “อนาคามี”ติ วตฺตพฺโพติ, อามนฺตา.\csromanspacer{} อรหา อนาคามิผเลน สมนฺนาคโตติ, อามนฺตา.\csromanspacer{} เสฺวว อรหา, โส อนาคามีติ, น เหวํ วตฺตพฺเพ ฯเปฯ.}

\prose{\csromansymbolmark{*}โสตาปตฺติผเลน สมนฺนาคโต “โสตาปนฺโน”ติ วตฺตพฺโพติ, อามนฺตา.\csromanspacer{} อนาคามี โสตาปตฺติผเลน สมนฺนาคโตติ, อามนฺตา.\csromanspacer{} เสฺวว อนาคามี, โส โสตาปนฺโนติ, น เหวํ วตฺตพฺเพ ฯเปฯ.}

\prose{\csromansymbolmark{*}สกทาคามิผเลน สมนฺนาคโต “สกทาคามี”ติ วตฺตพฺโพติ, อามนฺตา.\csromanspacer{} อนาคามี สกทาคามิผเลน สมนฺนาคโตติ, อามนฺตา.\csromanspacer{} เสฺวว อนาคามี, โส สกทาคามีติ, น เหวํ วตฺตพฺเพ ฯเปฯ.}

\prose{\csromansymbolmark{*}โสตาปตฺติผเลน สมนฺนาคโต “โสตาปนฺโน”ติ วตฺตพฺโพติ, อามนฺตา.\csromanspacer{} สกทาคามี โสตาปตฺติผเลน สมนฺนาคโตติ, อามนฺตา.\csromanspacer{} เสฺวว สกทาคามี, โส โสตาปนฺโนติ, น เหวํ วตฺตพฺเพ ฯเปฯ.}

\proseitem{395}{\csromansymbolmark{*}อรหา โสตาปตฺติผเลน สมนฺนาคโตติ, อามนฺตา.\csromanspacer{} นนุ อรหา โสตาปตฺติผลํ วีติวตฺโตติ, อามนฺตา.\csromanspacer{} หญฺจิ อรหา โสตาปตฺติผลํ วีติวตฺโต, โน จ วต เร วตฺตพฺเพ “อรหา โสตาปตฺติผเลน สมนฺนาคโต”ติ.}

\chapterheadpageread{4. จตุตฺถวคฺค}
\prose{\csromanfolio{209}\csromansymbolmark{*}อรหา โสตาปตฺติผลํ วีติวตฺโต เตน สมนฺนาคโตติ, อามนฺตา.\csromanspacer{} อรหา โสตาปตฺติมคฺคํ วีติวตฺโต, สกฺกายทิฏฺฐิํ, วิจิกิจฺฉํ, สีลพฺพตปรามาสํ, อปายคมนียํ ราคํ, อปายคมนียํ โทสํ, อปายคมนียํ โมหํ วีติวตฺโต เตน สมนฺนาคโตติ, น เหวํ วตฺตพฺเพ ฯเปฯ.}

\prose{\csromansymbolmark{*}อรหา สกทาคามิผเลน สมนฺนาคโตติ, อามนฺตา.\csromanspacer{} นนุ อรหา สกทาคามิผลํ วีติวตฺโตติ, อามนฺตา.\csromanspacer{} หญฺจิ อรหา สกทาคามิผลํ วีติวตฺโต, โน จ วต เร วตฺตพฺเพ “อรหา สกทาคามิผเลน สมนฺนาคโต”ติ.}

\prose{\csromansymbolmark{*}อรหา สกทาคามิผลํ วีติวตฺโต เตน สมนฺนาคโตติ, อามนฺตา.\csromanspacer{} อรหา สกทาคามิมคฺคํ วีติวตฺโต, โอฬาริกํ กามราคํ, โอฬาริกํ พฺยาปาทํ วีติวตฺโต เตน สมนฺนาคโตติ, น เหวํ วตฺตพฺเพ ฯเปฯ.}

\prose{\csromansymbolmark{*}อรหา อนาคามิผเลน สมนฺนาคโตติ, อามนฺตา.\csromanspacer{} นนุ อรหา อนาคามิผลํ วีติวตฺโตติ, อามนฺตา.\csromanspacer{} หญฺจิ อรหา อนาคามิผลํ วีติวตฺโต, โน จ วต เร วตฺตพฺเพ “อรหา อนาคามิผเลน สมนฺนาคโต”ติ.}

\prose{\csromansymbolmark{*}อรหา อนาคามิผลํ วีติวตฺโต เตน สมนฺนาคโตติ, อามนฺตา.\csromanspacer{} อรหา อนาคามิมคฺคํ วีติวตฺโต, อณุสหคตํ กามราคํ, อณุสหคตํ พฺยาปาทํ วีติวตฺโต เตน สมนฺนาคโตติ, น หิวํ วตฺตพฺเพ ฯเปฯ.}

\prose{\csromansymbolmark{*}อนาคามี โสตาปตฺติผเลน สมนฺนาคโตติ, อามนฺตา.\csromanspacer{} นนุ อนาคามี โสตาปตฺติผลํ วีติวตฺโตติ, อามนฺตา.\csromanspacer{} หญฺจิ อนาคามี โสตาปตฺติผลํ วีติวตฺโต, โน จ วต เร วตฺตพฺเพ “อนาคามี โสตาปตฺติผเลน สมนฺนาคโต”ติ.}

\prose{\csromansymbolmark{*}อนาคามี โสตาปตฺติผลํ วีติวตฺโต เตน สมนฺนาคโตติ, อามนฺตา.\csromanspacer{} อนาคามี โสตาปตฺติมคฺคํ วีติวตฺโต, สกฺกายทิฏฺฐิํ ฯเปฯ อปายคมนียํ โมหํ วีติวตฺโต เตน สมนฺนาคโตติ, น เหวํ วตฺตพฺเพ ฯเปฯ.}

\prose{\csromansymbolmark{*}อนาคามี สกทาคามิผเลน สมนฺนาคโตติ, อามนฺตา.\csromanspacer{} นนุ อนาคามี สกทาคามิผลํ วีติวตฺโตติ, อามนฺตา.\csromanspacer{} หญฺจิ อนาคามี สกทาคามิผลํ วีติวตฺโต, โน จ วต เร วตฺตพฺเพ “อนาคามี สกทาคามิผเลน สมนฺนาคโตติ.}

\prose{\csromanfolio{210}\csromansymbolmark{*}อนาคามี สกทาคามิผลํ วีติวตฺโต เตน สมนฺนาคโตติ, อามนฺตา.\csromanspacer{} อนาคามี สกทาคามิมคฺคํ วีติวตฺโต, โอฬาริกํ กามราคํ, โอฬาริกํ พฺยาปาทํ วีติวตฺโต เตน สมนฺนาคโตติ, น เหวํ วตฺตพฺเพ ฯเปฯ.}

\prose{\csromansymbolmark{*}สกทาคามี โสตาปตฺติผเลน สมนฺนาคโตติ, อามนฺตา.\csromanspacer{} นนุ สกทาคามี โสตาปตฺติผลํ วีติวตฺโตติ, อามนฺตา.\csromanspacer{} หญฺจิ สกทาคามี โสตาปตฺติผลํ วีติวตฺโต, โน จ วต เร วตฺตพฺเพ “สกทาคามี โสตาปตฺติผเลน สมนฺนาคโตติ.}

\prose{\csromansymbolmark{*}สกทาคามี โสตาปตฺติผลํ วีติวตฺโต เตน สมนฺนาคโตติ, อามนฺตา.\csromanspacer{} สกทาคามี โสตาปตฺติมคฺคํ วีติวตฺโต, สกฺกายทิฏฺฐิํ ฯเปฯ อปายคมนียํ โมหํ วีติวตฺโต เตน สมนฺนาคโตติ, น เหวํ วตฺตพฺเพ ฯเปฯ.}

\proseitem{396}{\csromansymbolmark{+}น วตฺตพฺพํ “อรหา จตูหิ ผเลหิ สมนฺนาคโต”ติ, อามนฺตา.\csromanspacer{} นนุ อรหตา จตฺตาริ ผลานิ ปฏิลทฺธานิ, เตหิ จ อปริหีโนติ, อามนฺตา.\csromanspacer{} หญฺจิ อรหตา จตฺตาริ ผลานิ ปฏิลทฺธานิ เตหิ จ อปริหีโน, เตน วต เร วตฺตพฺเพ “อรหา จตูหิ ผเลหิ สมนฺนาคโต”ติ.}

\prose{\csromansymbolmark{+}น วตฺตพฺพํ “อนาคามี ตีหิ ผเลหิ สมนฺนาคโต’ติ, อามนฺตา.\csromanspacer{} นนุ อนาคามินา ตีณิ ผลานิ ปฏิลทฺธานิ เตหิ จ อปริหีโนติ, อามนฺตา.\csromanspacer{} หญฺจิ อนาคามินา ตีณิ ผลานิ ปฏิลทฺธานิ, เตหิ จ อปริหีโน, เตน วต เร วตฺตพฺเพ “อนาคามี ตีหิ ผเลหิ สมนฺนาคโต”ติ.}

\prose{\csromansymbolmark{+}น วตฺตพฺพํ “สกทาคามี ทฺวีหิ ผเลหิ สมนฺนาคโต”ติ, อามนฺตา.\csromanspacer{} นนุ สกทาคามินา เทฺว ผลานิ ปฏิลทฺธานิ, เตหิ จ อปริหีโนติ, อามนฺตา.\csromanspacer{} หญฺจิ สกทาคามินา เทฺว ผลานิ ปฏิลทฺธานิ เตหิ จ อปริหีโน, เตน วต เร วตฺตพฺเพ “สกทาคามี ทฺวีหิ ผเลหิ สมนฺนาคโต”ติ.}

\prose{\csromansymbolmark{*}อรหตา จตฺตาริ ผลานิ ปฏิลทฺธานิ, เตหิ จ อปริหีโนติ อรหา จตูหิ ผเลหิ สมนฺนาคโตติ, อามนฺตา.\csromanspacer{} อรหตา จตฺตาโร มคฺคา ปฏิลทฺธา, เตหิ จ อปริหีโนติ อรหา จตูหิ มคฺเคหิ สมนฺนาคโตติ, น เหวํ วตฺตพฺเพ ฯเปฯ.}

\prose{\csromansymbolmark{*}อนาคามินา ตีณิ ผลานิ ปฏิลทฺธานิ, เตหิ จ อปริหีโนติ อนาคามี ตีหิ ผเลหิ สมนฺนาคโตติ, อามนฺตา.\csromanspacer{} อนาคามินา}

\vspace{\tipitakaparskip}
\csromancenter{จตุตฺถวคฺค ตโย มคฺคา ปฏิลทฺธา, เตหิ จ อปริหีโนติ อนาคามี ตีหิ มคฺเคหิ สมนฺนาคโตติ, น เหวํ วตฺตพฺเพ ฯเปฯ.}
\csromanmatikaanchor{mtk.59}
\csromantocmarkref{subsection}{(37) 5. อุเปกฺขาสมนฺนาคตกถา}{mtk.59}
\bookmark[dest={mtk.59},level=2]{(37) 5. อุเปกฺขาสมนฺนาคตกถา}
\prose{\csromanfolio{211}\csromansymbolmark{*}สกทาคามินา เทฺว ผลานิ ปฏิลทฺธานิ, เตหิ จ อปริหีโนติ สกทาคามี ทฺวีหิ ผเลหิ สมนฺนาคโตติ, อามนฺตา.\csromanspacer{} สกทาคามินา เทฺว มคฺคา ปฏิลทฺธา, เตหิ จ อปริหีโนติ สกทาคามี ทฺวีหิ มคฺเคหิ สมนฺนาคโตติ, น เหวํ วตฺตพฺเพ ฯเปฯ.}

\nitthitam{สมนฺนาคตกถา นิฏฺฐิตา.}
\chapterhead{4. จตุตฺถวคฺค}
\vspace{\tipitakaparskip}
\csromancenter{\textbf{(37) 5.}\csromanspacer{}\textbf{ อุเปกฺขาสมนฺนาคตกถา}}
\proseitem{397}{\csromansymbolmark{*}อรหา ฉหิ อุเปกฺขาหิ สมนฺนาคโตติ, อามนฺตา.\csromanspacer{} อรหา ฉหิ ผสฺเสหิ.\csromanspacer{} ฉหิ เวทนาหิ.\csromanspacer{} ฉหิ สญฺญาหิ ฯเปฯ.\csromanspacer{} ฉหิ ปญฺญาหิ สมนฺนาคโตติ, น เหวํ วตฺตพฺเพ ฯเปฯ.}

\prose{\csromansymbolmark{*}อรหา ฉหิ อุเปกฺขาหิ สมนฺนาคโตติ, อามนฺตา.\csromanspacer{} อรหา จกฺขุนา รูปํ ปสฺสนฺโต โสเตน สทฺทํ สุณาติ, ฆาเนน คนฺธํ ฆายติ, ชิวฺหาย รสํ สายติ, กาเยน โผฏฺฐพฺพํ ผุสติ, มนสา ธมฺมํ วิชานาติ ฯเปฯ มนสา ธมฺมํ วิชานนฺโต จกฺขุนา รูปํ ปสฺสติ, โสเตน สทฺทํ สุณาติ, ฆาเนน คนฺธํ ฆายติ, ชิวฺหาย รสํ สายติ, กาเยน โผฏฺฐพฺพํ ผุสตีติ, น เหวํ วตฺตพฺเพ ฯเปฯ.}

\prose{\csromansymbolmark{*}อรหา ฉหิ อุเปกฺขาหิ สมนฺนาคโตติ, อามนฺตา.\csromanspacer{} สตตํ สมิตํ อพฺโพกิณฺณํ ฉหิ อุเปกฺขาหิ สมนฺนาคโต สโมหิโต\footnote{สมาหิโต (สฺยา)}, ฉ อุเปกฺขาโย ปจฺจุปฏฺฐิตาติ, น เหวํ วตฺตพฺเพ ฯเปฯ.}

\prose{\csromansymbolmark{+}น วตฺตพฺพํ “อรหา ฉหิ อุเปกฺขาหิ สมนฺนาคโต”ติ, อามนฺตา.\csromanspacer{} นนุ อรหา ฉฬงฺคุเปกฺโขติ, อามนฺตา.\csromanspacer{} หญฺจิ อรหา ฉฬงฺคุเปกฺโข, เตน วต เร วตฺตพฺเพ “อรหา ฉหิ อุเปกฺขาหิ สมนฺนาคโต”ติ ฯเปฯ.}

\nitthitam{อุเปกฺขาสมนฺนาคตกถา นิฏฺฐิตา.}
\chapterheadpageread{4. จตุตฺถวคฺค}
\vspace{\tipitakaparskip}
\csromancenter{\textbf{(38) 6.}\csromanspacer{}\textbf{ โพธิยาพุทฺโธติกถา}}
\csromanmatikaanchor{mtk.60}
\csromantocmarkref{subsection}{(38) 6. โพธิยาพุทฺโธติกถา}{mtk.60}
\bookmark[dest={mtk.60},level=2]{(38) 6. โพธิยาพุทฺโธติกถา}
\proseitem{398}{\csromanfolio{212}\csromansymbolmark{*}โพธิยา พุทฺโธติ, อามนฺตา.\csromanspacer{} โพธิยา นิรุทฺธาย วิคตาย ปฏิปสฺสทฺธาย อพุทฺโธ โหตีติ, น เหวํ วตฺตพฺเพ ฯเปฯ.}

\prose{\csromansymbolmark{*}โพธิยา พุทฺโธติ, อามนฺตา.\csromanspacer{} อตีตาย โพธิยา พุทฺโธติ, น เหวํ วตฺตพฺเพ ฯเปฯ.\csromanspacer{} อตีตาย โพธิยา พุทฺโธติ, อามนฺตา.\csromanspacer{} ตาย โพธิยา โพธิกรณียํ กโรตีติ, น เหวํ วตฺตพฺเพ ฯเปฯ.}

\prose{\csromansymbolmark{*}ตาย โพธิยา โพธิกรณียํ กโรตีติ, อามนฺตา.\csromanspacer{} ตาย โพธิยา ทุกฺขํ ปริชานาติ, สมุทยํ ปชหติ, นิโรธํ สจฺฉิกโรติ, มคฺคํ ภาเวตีติ, น เหวํ วตฺตพฺเพ ฯเปฯ.}

\prose{\csromansymbolmark{*}โพธิยา พุทฺโธติ, อามนฺตา.\csromanspacer{} อนาคตาย โพธิยา พุทฺโธติ, น เหวํ วตฺตพฺเพ ฯเปฯ.}

\prose{\csromansymbolmark{*}อนาคตาย โพธิยา พุทฺโธติ, อามนฺตา.\csromanspacer{} ตาย โพธิยา โพธิกรณียํ กโรตีติ, น เหวํ วตฺตพฺเพ ฯเปฯ.}

\prose{\csromansymbolmark{*}ตาย โพธิยา โพธิกรณียํ กโรตีติ, อามนฺตา.\csromanspacer{} ตาย โพธิยา ทุกฺขํ ปริชานาติ ฯเปฯ มคฺคํ ภาเวตีติ, น เหวํ วตฺตพฺเพ ฯเปฯ.}

\prose{\csromansymbolmark{*}ปจฺจุปฺปนฺนาย โพธิยา พุทฺโธ, ตาย โพธิยา โพธิกรณียํ กโรตีติ, อามนฺตา.\csromanspacer{} อตีตาย โพธิยา พุทฺโธ, ตาย โพธิยา โพธิกรณียํ กโรตีติ, น เหวํ วตฺตพฺเพ ฯเปฯ.}

\prose{\csromansymbolmark{*}ปจฺจุปฺปนฺนาย โพธิยา พุทฺโธ, ตาย โพธิยา ทุกฺขํ ปริชานาติ, สมุทยํ ปชหติ, นิโรธํ สจฺฉิกโรติ, มคฺคํ ภาเวตีติ, อามนฺตา.\csromanspacer{} อตีตาย โพธิยา พุทฺโธ, ตาย โพธิยา ทุกฺขํ ปริชานาติ ฯเปฯ มคฺคํ ภาเวตีติ, น เหวํ วตฺตพฺเพ ฯเปฯ.}

\prose{\csromansymbolmark{*}ปจฺจุปฺปนฺนาย โพธิยา พุทฺโธ, ตาย โพธิยา โพธิกรณียํ กโรตีติ, อามนฺตา.\csromanspacer{} อนาคตาย โพธิยา พุทฺโธ, ตาย โพธิยา โพธิกรณียํ กโรตีติ, น เหวํ วตฺตพฺเพ ฯเปฯ.}

\chapterheadpageread{4. จตุตฺถวคฺค}
\prose{\csromanfolio{213}\csromansymbolmark{*}ปจฺจุปฺปนฺนาย โพธิยา พุทฺโธ, ตาย โพธิยา ทุกฺขํ ปริชานาติ ฯเปฯ มคฺคํ ภาเวตีติ, อามนฺตา.\csromanspacer{} อนาคตาย โพธิยา พุทฺโธ, ตาย โพธิยา ทุกฺขํ ปริชานาติ ฯเปฯ มคฺคํ ภาเวตีติ, น เหวํ วตฺตพฺเพ ฯเปฯ.}

\prose{\csromansymbolmark{*}อตีตาย โพธิยา พุทฺโธ, น จ ตาย โพธิยา โพธิกรณียํ กโรตีติ, อามนฺตา.\csromanspacer{} ปจฺจุปฺปนฺนาย โพธิยา พุทฺโธ, น จ ตาย โพธิยา โพธิกรณียํ กโรตีติ, น เหวํ วตฺตพฺเพ ฯเปฯ.}

\prose{\csromansymbolmark{*}อตีตาย โพธิยา พุทฺโธ, น จ ตาย โพธิยา ทุกฺขํ ปริชานาติ ฯเปฯ มคฺคํ ภาเวตีติ, อามนฺตา.\csromanspacer{} ปจฺจุปฺปนฺนาย โพธิยา พุทฺโธ, น จ ตาย โพธิยา ทุกฺขํ ปริชานาติ ฯเปฯ มคฺคํ ภาเวตีติ, น เหวํ วตฺตพฺเพ ฯเปฯ.}

\prose{\csromansymbolmark{*}อนาคตาย โพธิยา พุทฺโธ, น จ ตาย โพธิยา โพธิกรณียํ กโรตีติ, อามนฺตา.\csromanspacer{} ปจฺจุปฺปนฺนาย โพธิยา พุทฺโธ, น จ ตาย โพธิยา โพธิกรณียํ กโรตีติ, น เหวํ วตฺตพฺเพ ฯเปฯ.}

\prose{\csromansymbolmark{*}อนาคตาย โพธิยา พุทฺโธ, น จ ตาย โพธิยา ทุกฺขํ ปริชานาติ ฯเปฯ.\csromanspacer{} มคฺคํ ภาเวตีติ, อามนฺตา.\csromanspacer{} ปจฺจุปฺปนฺนาย โพธิยา พุทฺโธ, น จ ตาย โพธิยา ทุกฺขํ ปริชานาติ ฯเปฯ มคฺคํ ภาเวตีติ, น เหวํ วตฺตพฺเพ ฯเปฯ.}

\proseitem{399}{\csromansymbolmark{*}อตีตาย โพธิยา พุทฺโธ, อนาคตาย โพธิยา พุทฺโธ, ปจฺจุปฺปนฺนาย โพธิยา พุทฺโธติ, อามนฺตา.\csromanspacer{} ตีหิ โพธีหิ พุทฺโธติ, น เหวํ วตฺตพฺเพ ฯเปฯ.}

\prose{\csromansymbolmark{*}ตีหิ โพธีหิ พุทฺโธติ, อามนฺตา.\csromanspacer{} สตตํ สมิตํ อพฺโพกิณฺณํ ตีหิ โพธีหิ สมนฺนาคโต สโมหิโต, ติสฺโส โพธิโย ปจฺจุปฏฺฐิตาติ, น เหวํ วตฺตพฺเพ ฯเปฯ.}

\prose{\csromansymbolmark{+}น วตฺตพฺพํ “โพธิยา พุทฺโธ”ติ, อามนฺตา.\csromanspacer{} นนุ โพธิปฏิลาภา พุทฺโธติ, อามนฺตา.\csromanspacer{} หญฺจิ โพธิปฏิลาภา พุทฺโธ, เตน วต เร วตฺตพฺเพ “โพธิยา พุทฺโธ”ติ.}

\prose{\csromansymbolmark{*}โพธิปฏิลาภา พุทฺโธติ โพธิยา พุทฺโธติ, อามนฺตา.\csromanspacer{} โพธิปฏิลาภา โพธีติ, น เหวํ วตฺตพฺเพ ฯเปฯ.}

\nitthitam{โพธิยา พุทฺโธติกถา นิฏฺฐิตา.}
\chapterheadpageread{4. จตุตฺถวคฺค}
\vspace{\tipitakaparskip}
\csromancenter{\textbf{(39) 7.}\csromanspacer{}\textbf{ ลกฺขณกถา}}
\csromanmatikaanchor{mtk.61}
\csromantocmarkref{subsection}{(39) 7. ลกฺขณกถา}{mtk.61}
\bookmark[dest={mtk.61},level=2]{(39) 7. ลกฺขณกถา}
\proseitem{400}{\csromanfolio{214}\csromansymbolmark{*}ลกฺขณสมนฺนาคโต โพธิสตฺโตติ, อามนฺตา.\csromanspacer{} ปเทสลกฺขเณหิ สมนฺนาคโต ปเทสโพธิสตฺโตติ, น เหวํ วตฺตพฺเพ ฯเปฯ.}

\prose{\csromansymbolmark{*}ลกฺขณสมนฺนาคโต โพธิสตฺโตติ, อามนฺตา.\csromanspacer{} ติภาคลกฺขเณหิ สมนฺนาคโต ติภาคโพธิสตฺโตติ, น เหวํ วตฺตพฺเพ ฯเปฯ.}

\prose{\csromansymbolmark{*}ลกฺขณสมนฺนาคโต โพธิสตฺโตติ, อามนฺตา.\csromanspacer{} อุปฑฺฒลกฺขเณหิ สมนฺนาคโต อุปฑฺฒโพธิสตฺโตติ, น เหวํ วตฺตพฺเพ ฯเปฯ.}

\prose{\csromansymbolmark{*}ลกฺขณสมนฺนาคโต โพธิสตฺโตติ, อามนฺตา.\csromanspacer{} จกฺกวตฺติสตฺโต ลกฺขณสมนฺนาคโต, จกฺกวตฺติสตฺโต โพธิสตฺโตติ, น เหวํ วตฺตพฺเพ ฯเปฯ.}

\prose{\csromansymbolmark{*}จกฺกวตฺติสตฺโต ลกฺขณสมนฺนาคโต, จกฺกวตฺติสตฺโต โพธิสตฺโตติ, อามนฺตา.\csromanspacer{} ยาทิโส โพธิสตฺตสฺส ปุพฺพโยโค ปุพฺพจริยา ธมฺมกฺขานํ ธมฺมเทสนา, ตาทิโส จกฺกวตฺติสตฺตสฺส ปุพฺพโยโค ปุพฺพจริยา ธมฺมกฺขานํ ธมฺมเทสนาติ, น เหวํ วตฺตพฺเพ ฯเปฯ.}

\proseitem{401}{\csromansymbolmark{*}ยถา โพธิสตฺตสฺส ชายมานสฺส เทวา ปฐมํ ปฏิคฺคณฺหนฺติ ปจฺฉา มนุสฺสา\footnote{ที 2. 12; ม 3. 161 ปิฏฺฐาทีสุ วุตฺตํ นิสฺสาย ปุจฺฉติ.}, เอวเมวํ จกฺกวตฺติสตฺตสฺส ชายมานสฺส เทวา ปฐมํ ปฏิคฺคณฺหนฺติ ปจฺฉา มนุสฺสาติ, น เหวํ วตฺตพฺเพ ฯเปฯ.}

\prose{\csromansymbolmark{*}ยถา โพธิสตฺตสฺส ชายมานสฺส จตฺตาโร นํ เทวปุตฺตา ปฏิคฺคเหตฺวา มาตุ ปุรโต ฐเปนฺติ “อตฺตมนา เทวิ โหหิ, มเหสกฺโข เต ปุตฺโต อุปฺปนฺโน”ติ.\csromanspacer{} เอวเมวํ จกฺกวตฺติสตฺตสฺส ชายมานสฺส จตฺตาโร นํ เทวปุตฺตา ปฏิคฺคเหตฺวา มาตุ ปุรโต ฐเปนฺติ “อตฺตมนา เทวิ โหหิ, มเหสกฺโข เต ปุตฺโต อุปฺปนฺโน”ติ, น เหวํ วตฺตพฺเพ ฯเปฯ.}

\prose{\csromansymbolmark{*}ยถา โพธิสตฺตสฺส ชายมานสฺส เทฺว อุทกสฺส ธารา อนฺตลิกฺขา ปาตุภวนฺติ เอกา สีตสฺส เอกา อุณฺหสฺส, เยน โพธิสตฺตสฺส อุทกกิจฺจํ กโรนฺติ มาตุ จ.\csromanspacer{} เอวเมวํ จกฺกวตฺติสตฺตสฺส ชายมานสฺส เทฺว อุทกสฺส ธารา อนฺตลิกฺขา ปาตุภวนฺติ เอกา สีตสฺส เอกา อุณฺหสฺส,}

\vspace{\tipitakaparskip}
\csromancenter{จตุตฺถวคฺค เยน จกฺกวตฺติสตฺตสฺส อุทกกิจฺจํ กโรนฺติ มาตุ จาติ, น เหวํ วตฺตพฺเพ ฯเปฯ.}
\prose{\csromanfolio{215}\csromansymbolmark{*}ยถา สมฺปติชาโต โพธิสตฺโต สเมหิ ปาเทหิ ปติฏฺฐหิตฺวา อุตฺตเรน อภิมุโข สตฺตปทวีติหาเรน คจฺฉติ เสตมฺหิ ฉตฺเต อนุธาริยมาเน, สพฺพา จ ทิสา วิโลเกติ, อาสภิญฺจ วาจํ ภาสติ “อคฺโคหมสฺมิ โลกสฺส, เชฏฺโฐหมสฺมิ โลกสฺส, เสฏฺโฐหมสฺมิ โลกสฺส, อยมนฺติมา ชาติ, นตฺถิ ทานิ ปุนพฺภโว”ติ.\csromanspacer{} เอวเมวํ สมฺปติชาโต จกฺกวตฺติสตฺโต สเมหิ ปาเทหิ ปติฏฺฐหิตฺวา อุตฺตเรน อภิมุโข สตฺตปทวีติหาเรน คจฺฉติ เสตมฺหิ ฉตฺเต อนุธาริยมาเน, สพฺพา จ ทิสา วิโลเกติ, อาสภิญฺจ วาจํ ภาสติ “อคฺโคหมสฺมิ โลกสฺส, เชฏฺโฐหมสฺมิ โลกสฺส, เสฏฺโฐหมสฺมิ โลกสฺส, อยมนฺติมา ชาติ, นตฺถิ ทานิ ปุนพฺภโว”ติ, น เหวํ วตฺตพฺเพ ฯเปฯ.}

\prose{\csromansymbolmark{*}ยถา โพธิสตฺตสฺส ชายมานสฺส มหโต อาโลกสฺส มหโต โอภาสสฺส มหโต ภูมิจาลสฺส ปาตุภาโว โหติ, เอวเมวํ จกฺกวตฺติสตฺตสฺส ชายมานสฺส มหโต อาโลกสฺส มหโต โอภาสสฺส มหโต ภูมิจาลสฺส ปาตุภาโว โหตีติ, น เหวํ วตฺตพฺเพ ฯเปฯ.}

\prose{\csromansymbolmark{*}ยถา โพธิสตฺตสฺส ปกติกาโย สมนฺตา พฺยามํ โอภาสติ, เอวเมวํ จกฺกวตฺติสตฺตสฺส ปกติกาโย สมนฺตา พฺยามํ โอภาสตีติ, น เหวํ วตฺตพฺเพ ฯเปฯ.}

\prose{\csromansymbolmark{*}ยถา โพธิสตฺโต มหาสุปินํ ปสฺสติ, เอวเมวํ จกฺกวตฺติสตฺโต มหาสุปินํ ปสฺสตีติ, น เหวํ วตฺตพฺเพ ฯเปฯ.}

\csromanmatikaanchor{mtk.62}
\csromantocmarkref{subsection}{(40) 8. นิยาโมกฺกนฺติกถา}{mtk.62}
\bookmark[dest={mtk.62},level=2]{(40) 8. นิยาโมกฺกนฺติกถา}
\proseitem{402}{\csromansymbolmark{+}น วตฺตพฺพํ “ลกฺขณสมนฺนาคโต โพธิสตฺโต”ติ, อามนฺตา.\csromanspacer{} นนุ วุตฺตํ ภควตา “ทฺวตฺติํสิมานิ ภิกฺขเว มหาปุริสสฺส มหาปุริสลกฺขณานิ, เยหิ สมนฺนาคตสฺส มหาปุริสสฺส เทฺวว คติโย ภวนฺติ อนญฺญา\footnote{น อญฺญา (ก)}.\csromanspacer{} สเจ อคารํ อชฺฌาวสติ, ราชา โหติ จกฺกวตฺตี ธมฺมิโก ธมฺมราชา จาตุรนฺโต วิชิตาวี ชนปทตฺถาวริยปฺปตฺโต สตฺตรตนสมนฺนาคโต, ตสฺสิมานิ สตฺต รตนานิ ภวนฺติ.\csromanspacer{} เสยฺยถิทํ, จกฺกรตนํ หตฺถิรตนํ \csromanfolio{216}อสฺสรตนํ มณิรตนํ อิตฺถิรตนํ คหปติรตนํ ปริณายกรตนเมว สตฺตมํ.\csromanspacer{} ปโรสหสฺสํ โข ปนสฺส ปุตฺตา ภวนฺติ สูรา วีรงฺครูปา ปรเสนปฺปมทฺทนา.\csromanspacer{} โส อิมํ ปถวิํ สาครปริยนฺตํ อทณฺเฑน อสตฺเถน ธมฺเมน อภิวิชิย อชฺฌาวสติ, สเจ โข ปน อคารสฺมา อนคาริยํ ปพฺพชติ, อรหํ โหติ สมฺมาสมฺพุทฺโธ โลเก วิวฏฺฏจฺฉโท”ติ\footnote{ที 3. 117 ปิฏฺเฐ.} อตฺเถว สุตฺตนฺโตติ, อามนฺตา.\csromanspacer{} เตน หิ ลกฺขณสมนฺนาคโต โพธิสตฺโตติ.}

\nitthitam{ลกฺขณกถา นิฏฺฐิตา.}
\chapterhead{4. จตุตฺถวคฺค}
\vspace{\tipitakaparskip}
\csromancenter{\textbf{(40) 8.}\csromanspacer{}\textbf{ นิยาโมกฺกนฺติกถา}}
\proseitem{403}{\csromansymbolmark{*}โพธิสตฺโต กสฺสปสฺส ภควโต ปาวจเน โอกฺกนฺตนิยาโม จริตพฺรหฺมจริโยติ\footnote{ม 2. 236 ฆฏิการสุตฺตํ นิสฺสาย ปุจฺฉติ.}, อามนฺตา.\csromanspacer{} โพธิสตฺโต กสฺสปสฺส ภควโต สาวโกติ, น เหวํ วตฺตพฺเพ ฯเปฯ.}

\prose{\csromansymbolmark{*}โพธิสตฺโต กสฺสปสฺส ภควโต สาวโกติ, อามนฺตา.\csromanspacer{} สาวโก หุตฺวา พุทฺโธ โหตีติ, น เหวํ วตฺตพฺเพ ฯเปฯ.}

\prose{\csromansymbolmark{*}สาวโก หุตฺวา พุทฺโธ โหตีติ, อามนฺตา.\csromanspacer{} อนุสฺสวิโยติ, น เหวํ วตฺตพฺเพ ฯเปฯ.}

\prose{\csromansymbolmark{*}อนุสฺสวิโยติ, อามนฺตา.\csromanspacer{} นนุ ภควา สยมฺภูติ, อามนฺตา.\csromanspacer{} หญฺจิ ภควา สยมฺภู, โน จ วต เร วตฺตพฺเพ “อนุสฺสวิโย”ติ.}

\prose{\csromansymbolmark{*}โพธิสตฺโต กสฺสปสฺส ภควโต ปาวจเน โอกฺกนฺตนิยาโม จริตพฺรหฺมจริโยติ, อามนฺตา.\csromanspacer{} ภควตา โพธิยา มูเล ตีเณว สามญฺญผลานิ อภิสมฺพุทฺธานีติ, น เหวํ วตฺตพฺเพ ฯเปฯ.}

\prose{\csromansymbolmark{*}นนุ ภควตา โพธิยา มูเล จตฺตาริ สามญฺญผลานิ อภิสมฺพุทฺธานีติ, อามนฺตา.\csromanspacer{} หญฺจิ ภควตา โพธิยา มูเล จตฺตาริ สมญฺญผลานิ อภิสมฺพุทฺธานิ, โน จ วต เร วตฺตพฺเพ “โพธิสตฺโต กสฺสปสฺส ภควโต ปาวจเน โอกฺกนฺตนิยาโม จริตพฺรหฺมจริโย”ติ.}

\chapterheadpageread{4. จตุตฺถวคฺค}
\proseitem{404}{\csromanfolio{217}\csromansymbolmark{*}โพธิสตฺโต กสฺสปสฺส ภควโต ปาวจเน โอกฺกนฺตนิยาโม จริตพฺรหฺมจริโยติ, อามนฺตา.\csromanspacer{} โพธิสตฺโต ทุกฺกรการิยํ อกาสีติ, อามนฺตา.\csromanspacer{} ทสฺสนสมฺปนฺโน ปุคฺคโล ทุกฺกรการิยํ กเรยฺยาติ, น เหวํ วตฺตพฺเพ ฯเปฯ.}

\prose{\csromansymbolmark{*}โพธิสตฺโต อปรนฺตปํ\footnote{อมรํ ตปํ (สํ 1. 104 ปิฏฺเฐ)} อกาสิ, อญฺญํ สตฺถารํ อุทฺทิสีติ, อามนฺตา.\csromanspacer{} ทสฺสนสมฺปนฺโน ปุคฺคโล อญฺญํ สตฺถารํ อุทฺทิเสยฺยาติ, น เหวํ วตฺตพฺเพ ฯเปฯ.}

\prose{\csromansymbolmark{*}อายสฺมา อานนฺโท ภควโต ปาวจเน โอกฺกนฺตนิยาโม จริตพฺรหฺมจริโย, อายสฺมา อานนฺโท ภควโต สาวโกติ, อามนฺตา.\csromanspacer{} โพธิสตฺโต กสฺสปสฺส ภควโต ปาวจเน โอกฺกนฺตนิยาโม จริตพฺรหฺมจริโย, โพธิสตฺโต กสฺสปสฺส ภควโต สาวโกติ, น เหวํ วตฺตพฺเพ ฯเปฯ.}

\prose{\csromansymbolmark{*}จิตฺโต คหปติ, หตฺถโก อาฬวโก ภควโต ปาวจเน โอกฺกนฺตนิยาโม จริตพฺรหฺมจริโย, จิตฺโต คหปติ, หตฺถโก อาฬวโก ภควโต สาวโกติ, อามนฺตา.\csromanspacer{} โพธิสตฺโต กสฺสปสฺส ภควโต ปาวจเน โอกฺกนฺตนิยาโม จริตพฺรหฺมจริโย โพธิสตฺโต กสฺสปสฺส ภควโต สาวโกติ, น เหวํ วตฺตพฺเพ ฯเปฯ.}

\prose{\csromansymbolmark{*}โพธิสตฺโต กสฺสปสฺส ภควโต ปาวจเน โอกฺกนฺตนิยาโม จริตพฺรหฺมจริโย, น จ กสฺสปสฺส ภควโต สาวโกติ, อามนฺตา.\csromanspacer{} อายสฺมา อานนฺโท ภควโต ปาวจเน โอกฺกนฺตนิยาโม จริตพฺรหฺมจริโย, น จ ภควโต สาวโกติ, น เหวํ วตฺตพฺเพ ฯเปฯ.}

\prose{\csromansymbolmark{*}โพธิสตฺโต กสฺสปสฺส ภควโต ปาวจเน โอกฺกนฺตนิยาโม จริตพฺรหฺมจริโย, น จ กสฺสปสฺส ภควโต สาวโกติ, อามนฺตา.\csromanspacer{} จิตฺโต คหปติ หตฺถโก อาฬวโก ภควโต ปาวจเน โอกฺกนฺตนิยาโม จริตพฺรหฺมจริโย, น จ ภควโต สาวโกติ, น เหวํ วตฺตพฺเพ ฯเปฯ.}

\prose{\csromansymbolmark{*}โพธิสตฺโต กสฺสปสฺส ภควโต ปาวจเน โอกฺกนฺตนิยาโม จริตพฺรหฺมจริโย, น จ กสฺสปสฺส ภควโต สาวโกติ, อามนฺตา.\csromanspacer{} สาวโก ชาติํ วีติวตฺโต อสาวโก โหตีติ, น เหวํ วตฺตพฺเพ ฯเปฯ.}

\proseitem{405}{\csromanfolio{218}\csromansymbolmark{+}น วตฺตพฺพํ “โพธิสตฺโต กสฺสปสฺส ภควโต ปาวจเน โอกฺกนฺตนิยาโม จริตพฺรหฺมจริโย”ติ, อามนฺตา.\csromanspacer{} นนุ วุตฺตํ ภควตา “กสฺสเป อหํ อานนฺท ภควติ พฺรหฺมจริยํ อจริํ อายติํ สมฺโพธายา”ติ อตฺเถว สุตฺตนฺโตติ, อามนฺตา.\csromanspacer{} เตน หิ น วตฺตพฺพํ “โพธิสตฺโต กสฺสปสฺส ภควโต ปาวจเน โอกฺกนฺตนิยาโม จริตพฺรหฺมจริโย”ติ.}

\prose{\csromansymbolmark{*}โพธิสตฺโต กสฺสปสฺส ภควโต ปาวจเน โอกฺกนฺตนิยาโม จริตพฺรหฺมจริโยติ, อามนฺตา.\csromanspacer{} นนุ วุตฺตํ ภควตา\csromandash{}}

\csromangathameasure{“สพฺพาภิภู สพฺพวิทูหมสฺมิ, \\ สพฺเพสุ ธมฺเมสุ อนุปลิตฺโต. \\ สพฺพญฺชโห ตณฺหากฺขเย วิมุตฺโต, \\ สยํ อภิญฺญาย กมุทฺทิเสยฺยํ. \\ น เม อาจริโย อตฺถิ, สทิโส เม น วิชฺชติ. \\ สเทวกสฺมิํ โลกสฺมิํ, นตฺถิ เม ปฏิปุคฺคโล. \\ อหํ หิ อรหา โลเก, อหํ สตฺถา อนุตฺตโร. \\ เอโกมฺหิ สมฺมาสมฺพุทฺโธ, สีติภูโตสฺมิ นิพฺพุโต. \\ ธมฺมจกฺกํ ปวตฺเตตุํ, คจฺฉามิ กาสินํ ปุรํ. \\ อนฺธีภูตสฺมิํ โลกสฺมิํ, อาหญฺฉํ อมตทุนฺทุภินฺ”ติ.}
\csromangathagroup{\csromangathastanza{“สพฺพาภิภู สพฺพวิทูหมสฺมิ, \\ สพฺเพสุ ธมฺเมสุ อนุปลิตฺโต. \\ สพฺพญฺชโห ตณฺหากฺขเย วิมุตฺโต, \\ สยํ อภิญฺญาย กมุทฺทิเสยฺยํ. \\ น เม อาจริโย อตฺถิ, สทิโส เม น วิชฺชติ. \\ สเทวกสฺมิํ โลกสฺมิํ, นตฺถิ เม ปฏิปุคฺคโล.}\csromangathastanza{อหํ หิ อรหา โลเก, อหํ สตฺถา อนุตฺตโร. \\ เอโกมฺหิ สมฺมาสมฺพุทฺโธ, สีติภูโตสฺมิ นิพฺพุโต.}\csromangathastanza{ธมฺมจกฺกํ ปวตฺเตตุํ, คจฺฉามิ กาสินํ ปุรํ. \\ อนฺธีภูตสฺมิํ\footnote{อนฺธภูตสฺมิํ (สี, สฺยา, กํ, อิ)} โลกสฺมิํ, อาหญฺฉํ\footnote{อาหญฺญิํ (ก)} อมตทุนฺทุภินฺ”ติ\footnote{ทุทฺรภินฺติ (ก)}.}}

\csromanhangingbegin
\hangingheaditem{405}{ยถา โข ตฺวํ อาวุโส ปฏิชานาสิ, อรหสิ อนนฺตชิโนติ.}
\hangingbody{“มาทิสา เว ชินา โหนฺติ, เย ปตฺตา อาสวกฺขยํ.}
\hangingbody{ชิตา เม ปาปกา ธมฺมา, ตสฺมาหํ อุปก ชิโน”ติ4.}
\csromanhangingend

\prose{อตฺเถว สุตฺตนฺโตติ, อามนฺตา.\csromanspacer{} เตน หิ น วตฺตพฺพํ “โพธิสตฺโต กสฺสปสฺส ภควโต ปาวจเน โอกฺกนฺตนิยาโม จริตพฺรหฺมจริโย”ติ.}

\prose{\csromansymbolmark{*}โพธิสตฺโต กสฺสปสฺส ภควโต ปาวจเน โอกฺกนฺตนิยาโม จริตพฺรหฺมจริโยติ, อามนฺตา.\csromanspacer{} นนุ วุตฺตํ ภควตา\csromandash{}“อิทํ ทุกฺขํ อริยสจฺจนฺ”ติ เม ภิกฺขเว ปุพฺเพ อนนุสฺสุเตสุ ธมฺเมสุ จกฺขุํ อุทปาทิ, ญานํ อุทปาทิ,}

\vspace{\tipitakaparskip}
\csromancenter{จตุตฺถวคฺค ปญฺญา อุทปาทิ, วิชฺชา อุทปาทิ, อาโลโก อุทปาทิ. “ตํ โข ปนิทํ ทุกฺขํ อริยสจฺจํ ปริญฺเญยฺยนฺ”ติ เม ภิกฺขเว ฯเปฯ ปริญฺญาตนฺ”ติ เม ภิกฺขเว ปุพฺเพ อนนุสฺสุเตสุ ธมฺเมสุ จกฺขุํ อุทปาทิ ฯเปฯ อาโลโก อุทปาทิ. “อิทํ ทุกฺขสมุทยํ อริยสจฺจนฺ”ติ เม ภิกฺขเว ฯเปฯ “ตํ โข ปนิทํ ทุกฺขสมุทยํ\footnote{ทุกฺขสมุทโย (สฺยา, กํ)} อริยสจฺจํ ปหาตพฺพนฺ”ติ เม ภิกฺขเว ฯเปฯ ปหีนนฺ”ติ เม ภิกฺขเว ฯเปฯ “อิทํ ทุกฺขนิโรธํ\footnote{ทุกฺขนิโรโธ (สฺยา, กํ)} อริยสจฺจนฺ”ติ เม ภิกฺขเว ฯเปฯ “ตํ โข ปนิทํ ทุกฺขนิโรธํ อริยสจฺจํ สจฺฉิกาตพฺพนฺ”ติ เม ภิกฺขเว ฯเปฯ สจฺฉิกตนฺ”ติ เม ภิกฺขเว ฯเปฯ “อิทํ ทุกฺขนิโรธคามินี ปฏิปทา อริยสจฺจนฺ”ติ เม ภิกฺขเว ฯเปฯ “ตํ โข ปนิทํ ทุกฺขนิโรธคามินี ปฏิปทา อริยสจฺจํ ภาเวตพฺพนฺ”ติ เม ภิกฺขเว ฯเปฯ ภาวิตนฺ”ติ เม ภิกฺขเว ปุพฺเพ อนนุสฺสุเตสุ ธมฺเมสุ จกฺขุํ อุทปาทิ, ญาณํ อุทปาทิ, ปญฺญา อุทปาทิ, วิชฺชา อุทปาทิ, อาโลโก อุทปาทีติ\footnote{วิ 3. 15; สํ 3. 369; ขุ 9. 330 ปิฏฺเฐ จ.} อตฺเถว สุตฺตนฺโตติ, อามนฺตา.\csromanspacer{} เตน หิ น วตฺตพฺพํ “โพธิสตฺโต กสฺสปสฺส ภควโต ปาวจเน โอกฺกนฺตนิยาโม จริตพฺรหฺมจริโย”ติ.}
\nitthitam{นิยาโมกฺกนฺติกถา นิฏฺฐิตา.}
\chapterhead{4. จตุตฺถวคฺค}
\vspace{\tipitakaparskip}
\csromancenter{\textbf{(41) 9.}\csromanspacer{}\textbf{ อปราปิสมนฺนาคตกถา}}
\csromanmatikaanchor{mtk.63}
\csromantocmarkref{subsection}{(41) 9. อปราปิสมนฺนาคตกถา}{mtk.63}
\bookmark[dest={mtk.63},level=2]{(41) 9. อปราปิสมนฺนาคตกถา}
\proseitem{406}{\csromanfolio{219}\csromansymbolmark{*}อรหตฺตสจฺฉิกิริยาย ปฏิปนฺโน ปุคฺคโล ตีหิ ผเลหิ สมนฺนาคโตติ, อามนฺตา.\csromanspacer{} อรหตฺตสจฺฉิกิริยาย ปฏิปนฺโน ปุคฺคโล จตูหิ ผสฺเสหิ.\csromanspacer{} จตูหิ เวทนาหิ.\csromanspacer{} จตูหิ สญฺญาหิ.\csromanspacer{} จตูหิ เจตนาหิ.\csromanspacer{} จตูหิ จิตฺเตหิ.\csromanspacer{} จตูหิ สทฺธาหิ.\csromanspacer{} จตูหิ วีริเยหิ.\csromanspacer{} จตูหิ สตีหิ.\csromanspacer{} จตูหิ สมาธีหิ.\csromanspacer{} จตูหิ ปญฺญาหิ สมนฺนาคโตติ, น เหวํ วตฺตพฺเพ ฯเปฯ.}

\prose{\csromansymbolmark{*}อนาคามิผลสจฺฉิ กิริยาย ปฏิปนฺโน ปุคฺคโล ทฺวีหิ ผเลหิ สมนฺนาคโตติ, อามนฺตา.\csromanspacer{} อนาคามิผลสจฺฉิกิริยาย ปฏิปนฺโน ปุคฺคโล ตีหิ ผสฺเสหิ.\csromanspacer{} ตีหิ เวทนาหิ ฯเปฯ.\csromanspacer{} ตีหิ ปญฺญาหิ สมนฺนาคโตติ, น เหวํ วตฺตพฺเพ ฯเปฯ.}

\prose{\csromanfolio{220}\csromansymbolmark{*}สกทาคามิผลสจฺฉิกิริยาย ปฏิปนฺโน ปุคฺคโล โสตาปตฺติผเลน สมนฺนาคโตติ, อามนฺตา.\csromanspacer{} สกทาคามิผลสจฺฉิกิริยาย ปฏิปนฺโน ปุคฺคโล ทฺวีหิ ผสฺเสหิ.\csromanspacer{} ทฺวีหิ เวทนาหิ ฯเปฯ.\csromanspacer{} ทฺวีหิ ปญฺญาหิ สมนฺนาคโตติ, น เหวํ วตฺตพฺเพ ฯเปฯ.}

\prose{\csromansymbolmark{*}อรหตฺตสจฺฉิกิริยาย ปฏิปนฺโน ปุคฺคโล โสตาปตฺติผเลน สมนฺนาคโตติ, อามนฺตา.\csromanspacer{} อรหตฺตสจฺฉิกิริยาต ปฏิปนฺโน ปุคฺคโล โสตาปนฺโน สตฺตกฺขตฺตุปรโม, โกลํโกโล, เอกพีชีติ, น เหวํ วตฺตพฺเพ ฯเปฯ.}

\prose{\csromansymbolmark{*}อรหตฺตสจฺฉิกิริยาย ปฏิปนฺโน ปุคฺคโล สกทาคามิผเลน สมนฺนาคโตติ, อามนฺตา.\csromanspacer{} อรหตฺตสจฺฉิกิริยาย ปฏิปนฺโน ปุคฺคโล สกทาคามีติ, น เหวํ วตฺตพฺเพ ฯเปฯ.}

\prose{\csromansymbolmark{*}อรหตฺตสจฺฉิกิริยาย ปฏิปนฺโน ปุคฺคโล อนาคามิผเลน สมนฺนาคโตติ, อามนฺตา.\csromanspacer{} อรหตฺตสจฺฉิกิริยาย ปฏิปนฺโน ปุคฺคโล อนาคามี อนฺตราปรินิพฺพายี, อุปหจฺจปรินิพฺพายี, อสงฺขารปรินิพฺพายี, สสงฺขารปรินิพฺพายี, อุทฺธํโสโต อกนิฏฺฐคามีติ, น เหวํ วตฺตพฺเพ ฯเปฯ.}

\prose{\csromansymbolmark{*}อนาคามิผลสจฺฉิกิริยาย ปฏิปนฺโน ปุคฺคโล โสตาปตฺติผเลน สมนฺนาคโตติ, อามนฺตา.\csromanspacer{} อนาคามิผลสจฺฉิกิริยาย ปฏิปนฺโน ปุคฺคโล โสตาปนฺโน สตฺตกฺขตฺตุปรโม, โกลํโกโล, เอกพีชีติ, น เหวํ วตฺตพฺเพ ฯเปฯ.}

\prose{\csromansymbolmark{*}อนาคามิผลสจฺฉิ กิริยาย ปฏิปนฺโน ปุคฺคโล สกทาคามิผเลน สมนฺนาคโตติ, อามนฺตา.\csromanspacer{} อนาคามิผลสจฺฉิกิริยาย ปฏิปนฺโน ปุคฺคโล สกทาคามีติ, น เหวํ วตฺตพฺเพ ฯเปฯ.}

\prose{\csromansymbolmark{*}สกทาคามิผลสจฺฉิกิริยาย ปฏิปนฺโน ปุคฺคโล โสตาปตฺติผเลน สมนฺนาคโตติ, อามนฺตา.\csromanspacer{} สกทาคามิผลสจฺฉิกิริยาย ปฏิปนฺโน ปุคฺคโล โสตาปนฺโน สตฺตกฺขตฺตุปรโม, โกลํโกโล, เอกวีชีติ, น เหวํ วตฺตพฺเพ ฯเปฯ.}

\proseitem{407}{\csromansymbolmark{*}โสตาปตฺติผเลน สมนฺนาคโต “โสตาปนฺโน”ติ วตฺตพฺโพติ, อามนฺตา.\csromanspacer{} อรหตฺตสจฺฉิกิริยาย ปฏิปนฺโน ปุคฺคโล}

\vspace{\tipitakaparskip}
\csromancenter{จตุตฺถวคฺค “โสตาปตฺติผเลน สมนฺนาคโต”ติ, อามนฺตา.\csromanspacer{} เสฺวว อรหตฺตสจฺฉิกิริยาย ปฏิปนฺโน ปุคฺคโล, โส โสตาปนฺโนติ, น เหวํ วตฺตพฺเพ ฯเปฯ.}
\prose{\csromanfolio{221}\csromansymbolmark{*}สกทาคามิผเลน สมนฺนาคโต “สกทาคามี”ติ วตฺตพฺโพติ, อามนฺตา.\csromanspacer{} อรหตฺตสจฺฉิกิริยาย ปฏิปนฺโน ปุคฺคโล สกทาคามิผเลน สมนฺนาคโตติ, อามนฺตา.\csromanspacer{} เสฺวว อรหตฺตสจฺฉิกิริยาย ปฏิปนฺโน ปุคฺคโล, โส สกทาคามีติ, น เหวํ วตฺตพฺเพ ฯเปฯ.}

\prose{\csromansymbolmark{*}อนาคามิผเลน สมนฺนาคโต “อนาคามี”ติ วตฺตพฺโพติ, อามนฺตา.\csromanspacer{} อรหตฺตสจฺฉิกิริยาย ปฏิปนฺโน ปุคฺคโล อนาคามิผเลน สมนฺนาคโตติ, อามนฺตา.\csromanspacer{} เสฺวว อรหตฺตสจฺฉิกิริยาย ปฏิปนฺโน ปุคฺคโล, โส อนาคามีติ, น เหวํ วตฺตพฺเพ ฯเปฯ.}

\prose{\csromansymbolmark{*}โสตาปตฺติผเลน สมนฺนาคโต “โสตาปนฺโน”ติ วตฺตพฺโพติ, อามนฺตา.\csromanspacer{} อนาคามิผลสจฺฉิกิริยาย ปฏิปนฺโน ปุคฺคโล โสตาปตฺติผเลน สมนฺนาคโตติ, อามนฺตา.\csromanspacer{} เสฺวว อนาคามิผลสจฺฉิกิริยาย ปฏิปนฺโน ปุคฺคโล, โส โสตาปนฺโนติ, น เหวํ วตฺตพฺเพ ฯเปฯ.}

\prose{\csromansymbolmark{*}สกทาคามิผเลน สมนฺนาคโต “สกทาคามี”ติ วตฺตพฺโพติ, อามนฺตา.\csromanspacer{} อนาคามิผลสจฺฉิกิริยาย ปฏิปนฺโน ปุคฺคโล สกทาคามิผเลน สมนฺนาคโตติ, อามนฺตา.\csromanspacer{} เสฺวว อนาคามิผลสจฺฉิกิริยาย ปฏิปนฺโน ปุคฺคโล, โส สกทาคามีติ, น เหวํ วตฺตพฺเพ ฯเปฯ.}

\prose{\csromansymbolmark{*}โสตาปตฺติผเลน สมนฺนาคโต “โสตาปนฺโน”ติ วตฺตพฺโพติ, อามนฺตา.\csromanspacer{} สกทาคามิผลสจฺฉิกิริยาย ปฏิปนฺโน ปุคฺคโล โสตาปตฺติผเลน สมนฺนาคโตติ, อามนฺตา.\csromanspacer{} เสฺวว สกทาคามิผลสจฺฉิกิริยาย ปฏิปนฺโน ปุคฺคโล, โส โสตาปนฺโนติ, น เหวํ วตฺตพฺเพ ฯเปฯ.}

\proseitem{408}{\csromansymbolmark{*}อรหตฺตสจฺฉิกิริยาย ปฏิปนฺโน ปุคฺคโล โสตาปตฺติผเลน สมนฺนาคโตติ, อามนฺตา.\csromanspacer{} นนุ อรหตฺตสจฺฉิกิริยาย ปฏิปนฺโน ปุคฺคโล โสตาปตฺติผลํ วีติวตฺโตติ, อามนฺตา.\csromanspacer{} หญฺจิ อรหตฺตสจฺฉิกิริยาย ปฏิปนฺโน ปุคฺคโล โสตาปตฺติผลํ วีติวตฺโต, โน จ วต เร วตฺตพฺเพ “อรหตฺตสจฺฉิกิริยาย ปฏิปนฺโน ปุคฺคโล โสตาปตฺติผเลน สมนฺนาคโต”ติ.}

\prose{\csromanfolio{222}\csromansymbolmark{*}อรหตฺตสจฺฉิกิริยาย ปฏิปนฺโน ปุคฺคโล โสตาปตฺติผลํ วีติวตฺโต เตน สมนฺนาคโตติ, อามนฺตา.\csromanspacer{} อรหตฺตสจฺฉิกิริยาย ปฏิปนฺโน ปุคฺคโล โสตาปตฺติมคฺคํ วีติวตฺโต, สกฺกายทิฏฺฐิํ วิจิกิจฺฉํ สีลพฺพตปรามาสํ อปายคมนียํ ราคํ อปายคมนียํ โทสํ อปายคมนียํ โมหํ วีติวตฺโต เตน สมนฺนาคโตติ, น เหวํ วตฺตพฺเพ ฯเปฯ.}

\prose{\csromansymbolmark{*}อรหตฺตสจฺฉิกิริยาย ปฏิปนฺโน ปุคฺคโล สกทาคามิผเลน สมนฺนาคโตติ, อามนฺตา.\csromanspacer{} นนุ อรหตฺตสจฺฉิกิริยาย ปฏิปนฺโน ปุคฺคโล สกทาคามิผลํ วีติวตฺโตติ, อามนฺตา.\csromanspacer{} หญฺจิ อรหตฺตสจฺฉิกิริยาย ปฏิปนฺโน ปุคฺคโล สกทาคามิผลํ วีติวตฺโต, โน จ วต เร วตฺตพฺเพ “อรหตฺตสจฺฉิกิริยาย ปฏิปนฺโน ปุคฺคโล สกทาคามิผเลน สมนฺนาคโต”ติ.}

\prose{\csromansymbolmark{*}อรหตฺตสจฺฉิกิริยาย ปฏิปนฺโน ปุคฺคโล สกทาคามิผลํ วีติวตฺโต เตน สมนฺนาคโตติ, อามนฺตา.\csromanspacer{} อรหตฺตสจฺฉิกิริยาย ปฏิปนฺโน ปุคฺคโล สกทาคามิมคฺคํ วีติวตฺโต, โอฬาริกํ กามราคํ โอฬาริกํ พฺยาปาทํ วีติวตฺโต เตน สมนฺนาคโตติ, น เหวํ วตฺตพฺเพ ฯเปฯ.}

\prose{\csromansymbolmark{*}อรหตฺตสจฺฉิกิริยาย ปฏิปนฺโน ปุคฺคโล อนาคามิผเลน สมนฺนาคโตติ, อามนฺตา.\csromanspacer{} นนุ อรหตฺตสจฺฉิกิริยาย ปฏิปนฺโน ปุคฺคโล อนาคามิผลํ วีติวตฺโตติ, อามนฺตา.\csromanspacer{} หญฺจิ อรหตฺตสจฺฉิกิริยาย ปฏิปนฺโน ปุคฺคโล อนาคามิผลํ วีติวตฺโต, โน จ วต เร วตฺตพฺเพ “อรหตฺตสจฺฉิกิริยาย ปฏิปนฺโน ปุคฺคโล อนาคามิผเลน สมนฺนาคโต”ติ.}

\prose{\csromansymbolmark{*}อรหตฺตสจฺฉิกิริยาย ปฏิปนฺโน ปุคฺคโล อนาคามิผลํ วีติวตฺโต เตน สมนฺนาคโตติ, อามนฺตา.\csromanspacer{} อรหตฺตสจฺฉิกิริยาย ปฏิปนฺโน ปุคฺคโล อนาคามิมคฺคํ วีติวตฺโต, อณุสหคตํ กามราคํ อณุสหคตํ พฺยาปาทํ วีติวตฺโต เตน สมนฺนาคโตติ, น เหวํ วตฺตพฺเพ ฯเปฯ.}

\proseitem{409}{\csromansymbolmark{*}อนาคามิผลสจฺฉิกิริยาย ปฏิปนฺโน ปุคฺคโล โสตาปตฺติผเลน สมนฺนาคโตติ, อามนฺตา.\csromanspacer{} นนุ อนาคามิผลสจฺฉิกิริยาย ปฏิปนฺโน ปุคฺคโล โสตาปตฺติผลํ วีติวตฺโตติ, อามนฺตา.\csromanspacer{} หญฺจิ อนาคามิผลสจฺฉิกิริยาย ปฏิปนฺโน ปุคฺคโล โสตาปตฺติผลํ วีติวตฺโต, โน จ วต เร วตฺตพฺเพ “อนาคามิผลสจฺฉิกิริยาย ปฏิปนฺโน ปุคฺคโล โสตาปตฺติผเลน สมนฺนาคโต”ติ.}

\chapterheadpageread{4. จตุตฺถวคฺค}
\prose{\csromanfolio{223}\csromansymbolmark{*}อนาคามิผลสจฺฉิกิริยาย ปฏิปนฺโน ปุคฺคโล โสตาปตฺติผลํ วีติวตฺโต เตน สมนฺนาคโตติ, อามนฺตา.\csromanspacer{} อนาคามิผลสจฺฉิกิริยาย ปฏิปนฺโน ปุคฺคโล โสตาปตฺติมคฺคํ วีติวตฺโต, สกฺกายทิฏฺฐิํ ฯเปฯ อปายคมนียํ โมหํ วีติวตฺโต เตน สมนฺนาคโตติ, น เหวํ วตฺตพฺเพ ฯเปฯ.}

\prose{\csromansymbolmark{*}อนาคามิผลสจฺฉิกิริยาย ปฏิปนฺโน ปุคฺคโล สกทาคามิผเลน สมนฺนาคโตติ, อามนฺตา.\csromanspacer{} นนุ อนาคามิผลสจฺฉิกิริยาย ปฏิปนฺโน ปุคฺคโล สกทาคามิผลํ วีติวตฺโตติ, อามนฺตา.\csromanspacer{} หญฺจิ อนาคามิผลสจฺฉิกิริยาย ปฏิปนฺโน ปุคฺคโล สกทาคามิผลํ วีติวตฺโต, โน จ วต เร วตฺตพฺเพ “อนาคามิผลสจฺฉิกิริยาย ปฏิปนฺโน ปุคฺคโล สกทาคามิผเลน สมนฺนาคโต”ติ.}

\prose{\csromansymbolmark{*}อนาคามิผลสจฺฉิกิริยาย ปฏิปนฺโน ปุคฺคโล สกทาคามิผลํ วีติวตฺโต เตน สมนฺนาคโตติ, อามนฺตา.\csromanspacer{} อนาคามิผลสจฺฉิกิริยาย ปฏิปนฺโน ปุคฺคโล สกทาคามิมคฺคํ วีติวตฺโต, โอฬาริกํ กามราคํ โอฬาริกํ พฺยาปาทํ วีติวตฺโต เตน สมนฺนาคโตติ, น เหวํ วตฺตพฺเพ ฯเปฯ.}

\proseitem{410}{\csromansymbolmark{*}สกทาคามิผลสจฺฉิกิริยาย ปฏิปนฺโน ปุคฺคโล โสตาปตฺติผเลน สมนฺนาคโตติ, อามนฺตา.\csromanspacer{} นนุ สกทาคามิผลสจฺฉิกิริยาย ปฏิปนฺโน ปุคฺคโล โสตาปตฺติผลํ วีติวตฺโตติ, อามนฺตา.\csromanspacer{} หญฺจิ สกทาคามิผลสจฺฉิกิริยาย ปฏิปนฺโน ปุคฺคโล โสตาปตฺติผลํ วีติวตฺโต, โน จ วต เร วตฺตพฺเพ “สกทาคามิผลสจฺฉิกิริยาย ปฏิปนฺโน ปุคฺคโล โสตาปตฺติผเลน สมนฺนาคโต”ติ.}

\prose{\csromansymbolmark{*}สกทาคามิผลสจฺฉิกิริยาย ปฏิปนฺโน ปุคฺคโล โสตาปตฺติผลํ วีติวตฺโต เตน สมนฺนาคโตติ, อามนฺตา.\csromanspacer{} สกทาคามิผลสจฺฉิกิริยาย ปฏิปนฺโน ปุคฺคโล โสตาปตฺติมคฺคํ วีติวตฺโต, สกฺกายทิฏฺฐิํ ฯเปฯ อปายคมนียํ โมหํ วีติวตฺโต เตน สมนฺนาคโตติ, น เหวํ วตฺตพฺเพ ฯเปฯ.}

\proseitem{411}{\csromansymbolmark{+}น วตฺตพฺพํ “อรหตฺตสจฺฉิกิริยาย ปฏิปนฺโน ปุคฺคโล ตีหิ ผเลหิ สมนฺนาคโต”ติ, อามนฺตา.\csromanspacer{} นนุ อรหตฺตสจฺฉิกิริยาย ปฏิปนฺเนน ปุคฺคเลน ตีณิ ผลานิ ปฏิลทฺธานิ, เตหิ จ อปริหีโนติ, อามนฺตา.\csromanspacer{} หญฺจิ อรหตฺตสจฺฉิกิริยาย ปฏิปนฺเนน ปุคฺคเลน ตีณิ ผลานิ ปฏิลทฺธานิ, \csromanfolio{224}เตหิ จ อปริหีโน, เตน วต เร วตฺตพฺเพ “อรหตฺตสจฺฉิกิริยาย ปฏิปนฺโน ปุคฺคโล ตีหิ ผเลหิ สมนฺนาคโต”ติ.}

\prose{\csromansymbolmark{+}น วตฺตพฺพํ “อนาคามิผลสจฺฉิกิริยาย ปฏิปนฺโน ปุคฺคโล ทฺวีหิ ผเลหิ สมนฺนาคโต”ติ, อามนฺตา.\csromanspacer{} นนุ อนาคามิผลสจฺฉิกิริยาย ปฏิปนฺเนน ปุคฺคเลน เทฺว ผลานิ ปฏิลทฺธานิ, เตหิ จ อปริหีโนติ, อามนฺตา.\csromanspacer{} หญฺจิ อนาคามิผลสจฺฉิกิริยาย ปฏิปนฺเนน ปุคฺคเลน เทฺว ผลานิ ปฏิลทฺธานิ, เตหิ จ อปริหีโน, เตน วต เร วตฺตพฺเพ “อนาคามิผลสจฺฉิกิริยาย ปฏิปนฺโน ปุคฺคโล ทฺวีหิ ผเลหิ สมนฺนาคโต”ติ.}

\prose{\csromansymbolmark{+}น วตฺตพฺพํ “สกทาคามิผลสจฺฉิกิริยาย ปฏิปนฺโน ปุคฺคโล โสตาปตฺติผเลน สมนฺนาคโต”ติ, อามนฺตา.\csromanspacer{} นนุ สกทาคามิผลสจฺฉิกิริยาย ปฏิปนฺเนน ปุคฺคเลน โสตาปตฺติผลํ ปฏิลทฺธํ, เตน จ อปริหีโนติ, อามนฺตา.\csromanspacer{} หญฺจิ สกทาคามิผลสจฺฉิกิริยาย ปฏิปนฺเนน ปุคฺคเลน โสตาปตฺติผลํ ปฏิลทฺธํ, เตน จ อปริหีโน, เตน วต เร วตฺตพฺเพ “สกทาคามิผลสจฺฉิกิริยาย ปฏิปนฺโน ปุคฺคโล โสตาปตฺติผเลน สมนฺนาคโต”ติ.}

\proseitem{412}{\csromansymbolmark{*}อรหตฺตสจฺฉิกิริยาย ปฏิปนฺเนน ปุคฺคเลน ตีณิ ผลานิ ปฏิลทฺธานิ, เตหิ จ อปริหีโนติ อรหตฺตสจฺฉิกิริยาย ปฏิปนฺโน ปุคฺคโล ตีหิ ผเลหิ สมนฺนาคโตติ, อามนฺตา.\csromanspacer{} อรหตฺตสจฺฉิกิริยาย ปฏิปนฺเนน ปุคฺคเลน จตฺตาโร มคฺคา ปฏิลทฺธา, เตหิ จ อปริหีโนติ อรหตฺตสจฺฉิกิริยาย ปฏิปนฺโน ปุคฺคโล จตูหิ มคฺเคหิ สมนฺนาคโตติ, น เหวํ วตฺตพฺเพ ฯเปฯ.}

\prose{\csromansymbolmark{*}อนาคามิผลสจฺฉิกิริยาย ปฏิปนฺเนน ปุคฺคเลน เทฺว ผลานิ ปฏิลทฺธานิ, เตหิ จ อปริหีโนติ อนาคามิผลสจฺฉิกิริยาย ปฏิปนฺโน ปุคฺคโล ทฺวีหิ ผเลหิ สมนฺนาคโตติ, อามนฺตา.\csromanspacer{} อนาคามิผลสจฺฉิกิริยาย ปฏิปนฺเนน ปุคฺคเลน ตโย มคฺคา ปฏิลทฺธา, เตหิ จ อปริหีโนติ อนาคามิผลสจฺฉิกิริยาย ปฏิปนฺโน ปุคฺคโล ตีหิ มคฺเคหิ สมนฺนาคโตติ, น เหวํ วตฺตพฺเพ ฯเปฯ.}

\prose{\csromansymbolmark{*}สกทาคามิผลสจฺฉิกิริยาย ปฏิปนฺเนน ปุคฺคเลน โสตาปตฺติผลํ ปฏิลทฺธํ, เตน จ อปริหีโนติ สกทาคามิผลสจฺฉิกิริยาย ปฏิปนฺโน}

\vspace{\tipitakaparskip}
\csromancenter{จตุตฺถวคฺค ปุคฺคโล โสตาปตฺติผเลน สมนฺนาคโตติ, อามนฺตา.\csromanspacer{} สกทาคามิผลสจฺฉิกิริยาย ปฏิปนฺเนน ปุคฺคเลน เทฺว มคฺคา ปฏิลทฺธา, เตหิ จ อปริหีโนติ, สกทาคามิผลสจฺฉิกิริยาย ปฏิปนฺโน ปุคฺคโล ทฺวีหิ มคฺเคหิ สมนฺนาคโตติ, น เหวํ วตฺตพฺเพ ฯเปฯ.}
\nitthitam{อปราปิ สมนฺนาคตกถา นิฏฺฐิตา.}
\chapterhead{4. จตุตฺถวคฺค}
\vspace{\tipitakaparskip}
\csromancenter{\textbf{(42) 10.}\csromanspacer{}\textbf{ สพฺพสํโยชนปฺปหานกถา}}
\csromanmatikaanchor{mtk.64}
\csromantocmarkref{subsection}{(42) 10. สพฺพสํโยชนปฺปหานกถา}{mtk.64}
\bookmark[dest={mtk.64},level=2]{(42) 10. สพฺพสํโยชนปฺปหานกถา}
\proseitem{413}{\csromanfolio{225}\csromansymbolmark{*}สพฺพสํโยชนานํ ปหานํ อรหตฺตนฺติ, อามนฺตา.\csromanspacer{} อรหตฺตมคฺเคน สพฺเพ สํโยชนา ปหียนฺตีติ, น เหวํ วตฺตพฺเพ ฯเปฯ.}

\prose{\csromansymbolmark{*}อรหตฺตมคฺเคน สพฺเพ สํโยชนา ปหียนฺตีติ, อามนฺตา.\csromanspacer{} อรหตฺตมคฺเคน สกฺกายทิฏฺฐิํ วิจิกิจฺฉํ สีลพฺพตปรามาสํ ปชหตีติ, น เหวํ วตฺตพฺเพ ฯเปฯ.}

\prose{\csromansymbolmark{*}อรหตฺตมคฺเคน สกฺกายทิฏฺฐิํ วิจิกิจฺฉํ สีลพฺพตปรามาสํ ปชหตีติ, อามนฺตา.\csromanspacer{} นนุ ติณฺณํ สํโยชนานํ ปหานํ โสตาปตฺติผลํ วุตฺตํ ภควตาติ, อามนฺตา.\csromanspacer{} หญฺจิ ติณฺณํ สํโยชนานํ ปหานํ โสตาปตฺติผลํ วุตฺตํ ภควตา, โน จ วต เร วตฺตพฺเพ “อรหตฺตมคฺเคน สพฺเพ สํโยชนา ปหียนฺตี”ติ.}

\proseitem{414}{\csromansymbolmark{*}อรหตฺตมคฺเคน สพฺเพ สํโยชนา ปหียนฺตีติ, อามนฺตา.\csromanspacer{} อรหตฺตมคฺเคน โอฬาริกํ กามราคํ โอฬาริกํ พฺยาปาทํ ปชหตีติ, น เหวํ วตฺตพฺเพ ฯเปฯ.}

\proseitem{415}{\csromansymbolmark{*}อรหตฺตมคฺเคน โอฬาริกํ กามราคํ โอฬาริกํ พฺยาปาทํ ปชหตีติ, อามนฺตา.\csromanspacer{} นนุ กามราคพฺยาปาทานํ ตนุภาวํ สกทาคามิผลํ วุตฺตํ ภควตาติ, อามนฺตา.\csromanspacer{} หญฺจิ กามราคพฺยาปาทานํ ตนุภาวํ สกทาคามิผลํ วุตฺตํ ภควตา, โน จ วต เร วตฺตพฺเพ “อรหตฺตมคฺเคน สพฺเพ สํโยชนา ปหียนฺตี”ติ.}

\prose{\csromanfolio{226}\csromansymbolmark{*}อรหตฺตมคฺเคน สพฺเพ สํโยชนา ปหียนฺตีติ, อามนฺตา.\csromanspacer{} อรหตฺตมคฺเคน อณุสหคตํ กามราคํ อณุสหคตํ พฺยาปาทํ ปชหตีติ, น เหวํ วตฺตพฺเพ ฯเปฯ.}

\proseitem{416}{\csromansymbolmark{*}อรหตฺตมคฺเคน อณุสหคตํ กามราคํ อณุสหคตํ พฺยาปาทํ ปชหตีติ, อามนฺตา.\csromanspacer{} นนุ กามราคพฺยาปาทานํ อนวเสสปฺปหานํ อนาคามิผลํ วุตฺตํ ภควตาติ, อามนฺตา.\csromanspacer{} หญฺจิ กามราคพฺยาปาทานํ อนวเสสปฺปหานํ อนาคามิผลํ วุตฺตํ ภควตา, โน จ วต เร วตฺตพฺเพ “อรหตฺตมคฺเคน สพฺเพ สํโยชนา ปหียนฺตี”ติ.}

\prose{\csromansymbolmark{*}อรหตฺตมคฺเคน สพฺเพ สํโยชนา ปหียนฺตีติ, อามนฺตา.\csromanspacer{} นนุ รูปราค อรูปราค มาน อุทฺธจฺจ อวิชฺชาย อนวเสสปฺปหานํ อรหตฺตํ วุตฺตํ ภควตาติ, อามนฺตา.\csromanspacer{} หญฺจิ รูปราค อรูปราค มาน อุทฺธจฺจ อวิชฺชาย อนวเสสปฺปหานํ อรหตฺตํ วุตฺตํ ภควตา, โน จ วต เร วตฺตพฺเพ “อรหตฺตมคฺเคน สพฺเพ สํโยชนา ปหียนฺตี”ติ.}

\proseitem{417}{\csromansymbolmark{+}น วตฺตพฺพํ “สพฺพสํโยชนานํ ปหานํ อรหตฺตนฺ”ติ, อามนฺตา.\csromanspacer{} นนุ อรหโต สพฺเพ สํโยชนา ปหีนาติ, อามนฺตา.\csromanspacer{} หญฺจิ อรหโต สพฺเพ สํโยชนา ปหีนา, เตน วต เร วตฺตพฺเพ “สพฺพสํโยชนานํ ปหานํ อรหตฺตนฺ”ติ.\csromanspacer{} สพฺพสํโยชนปฺปหานกถา นิฏฺฐิตา.\csromanspacer{} จตุตฺโถ วคฺโค.}

\titlehead{\textbf{ตสฺสุทฺทานํ}}
\prose{คิหิสฺส อรหา, สห อุปปตฺติยา อรหา, อรหโต สพฺเพ ธมฺมา อนาสวา, อรหา จตูหิ ผเลหิ สมนฺนาคโต, เอวเมวํ ฉหิ อุเปกฺขาหิ, โพธิยา พุทฺโธ, สลกฺขณสมนฺนาคโต, โพธิสตฺโต โอกฺกนฺตนิยาโม จริตพฺรหฺมจริโย, ปฏิปนฺนโก ผเลน สมนฺนาคโต, สพฺพสํโยชนานํ ปหานํ อรหตฺตนฺติ.}

\csromanmatikaanchor{mtk.65}
\csromantocmarknopage{chapter}{5. ปญฺจมวคฺค}{mtk.65}
\bookmark[dest={mtk.65},level=0]{5. ปญฺจมวคฺค}
\chapterheadpageread{5. ปญฺจมวคฺค \textbf{5. ปญฺจมวคฺค}}
\vspace{\tipitakaparskip}
\csromancenter{\textbf{(43) 1.}\csromanspacer{}\textbf{ วิมุตฺติกถา}}
\csromanmatikaanchor{mtk.66}
\csromantocmarkref{subsection}{(43) 1. วิมุตฺติกถา}{mtk.66}
\bookmark[dest={mtk.66},level=2]{(43) 1. วิมุตฺติกถา}
\proseitem{418}{\csromanfolio{227}\csromansymbolmark{*}วิมุตฺติญาณํ วิมุตฺตนฺติ, อามนฺตา.\csromanspacer{} ยํ กิญฺจิ วิมุตฺติญาณํ, สพฺพํ ตํ วิมุตฺตนฺติ, น เหวํ วตฺตพฺเพ ฯเปฯ.\csromanspacer{} วิมุตฺติญาณํ วิมุตฺตนฺติ, อามนฺตา.\csromanspacer{} ปจฺจเวกฺขณญาณํ วิมุตฺตนฺติ, น เหวํ วตฺตพฺเพ ฯเปฯ.}

\prose{\csromansymbolmark{*}วิมุตฺติญาณํ วิมุตฺตนฺติ, อามนฺตา.\csromanspacer{} โคตฺรภุโน ปุคฺคลสฺส วิมุตฺติญาณํ วิมุตฺตนฺติ, น เหวํ วตฺตพฺเพ ฯเปฯ.\csromanspacer{} โสตาปตฺติผลสจฺฉิกิริยาย ปฏิปนฺนสฺส ปุคฺคลสฺส วิมุตฺติญาณํ วิมุตฺตนฺติ, อามนฺตา.\csromanspacer{} โสตาปนฺนสฺส ญาณํ, โสตาปตฺติผลํ ปตฺตสฺส ปฏิลทฺธสฺส อธิคตสฺส สจฺฉิกตสฺส ญาณนฺติ, น เหวํ วตฺตพฺเพ ฯเปฯ.\csromanspacer{} สกทาคามิผลสจฺฉิกิริยาย ปฏิปนฺนสฺส ปุคฺคลสฺส วิมุตฺติญาณํ วิมุตฺตนฺติ, อามนฺตา.\csromanspacer{} สกทาคามิสฺส ญาณํ, สกทาคามิผลํ ปตฺตสฺส ปฏิลทฺธสฺส อธิคตสฺส สจฺฉิกตสฺส ญาณนฺติ, น เหวํ วตฺตพฺเพ ฯเปฯ.\csromanspacer{} อนาคามิผลสจฺฉิกิริยาย ปฏิปนฺนสฺส ปุคฺคลสฺส วิมุตฺติญาณํ วิมุตฺตนฺติ, อามนฺตา.\csromanspacer{} อนาคามิสฺส ญาณํ, อนาคามิผลํ ปตฺตสฺส ปฏิลทฺธสฺส อธิคตสฺส สจฺฉิกตสฺส ญาณนฺติ, น เหวํ วตฺตพฺเพ ฯเปฯ.\csromanspacer{} อรหตฺตสจฺฉิกิริยาย ปฏิปนฺนสฺส ปุคฺคลสฺส วิมุตฺติญาณํ วิมุตฺตนฺติ, อามนฺตา.\csromanspacer{} อรหโต ญาณํ, อรหตฺตํ ปตฺตสฺส ปฏิลทฺธสฺส อธิคตสฺส สจฺฉิกตสฺส ญาณนฺติ, น เหวํ วตฺตพฺเพ ฯเปฯ.}

\proseitem{419}{\csromansymbolmark{*}โสตาปตฺติผลสมงฺคิสฺส ปุคฺคลสฺส วิมุตฺติญาณํ วิมุตฺตนฺติ, อามนฺตา.\csromanspacer{} โสตาปตฺติผลสจฺฉิกิริยาย ปฏิปนฺนสฺส ปุคฺคลสฺส วิมุตฺติญาณํ วิมุตฺตนฺติ, น เหวํ วตฺตพฺเพ ฯเปฯ.\csromanspacer{} สกทาคามิผลสมงฺคิสฺส ปุคฺคลสฺส วิมุตฺติญาณํ วิมุตฺตนฺติ, อามนฺตา.\csromanspacer{} สกทาคามิผลสจฺฉิกิริยาย ปฏิปนฺนสฺส ปุคฺคลสฺส วิมุตฺติญาณํ วิมุตฺตนฺติ, น เหวํ วตฺตพฺเพ ฯเปฯ.\csromanspacer{} อนาคามิผลสมงฺคิสฺส ปุคฺคลสฺส วิมุตฺติญาณํ วิมุตฺตนฺติ, อามนฺตา.\csromanspacer{} อนาคามิผลสจฺฉิกิริยาย ปฏิปนฺนสฺส ปุคฺคลสฺส วิมุตฺติญาณํ วิมุตฺตนฺติ, น เหวํ วตฺตพฺเพ ฯเปฯ.\csromanspacer{} อรหตฺตผลสมงฺคิสฺส ปุคฺคลสฺส วิมุตฺติญาณํ วิมุตฺตนฺติ, อามนฺตา.\csromanspacer{} อรหตฺตสจฺฉิกิริยาย ปฏิปนฺนสฺส ปุคฺคลสฺส วิมุตฺติญาณํ วิมุตฺตนฺติ, น เหวํ วตฺตพฺเพ ฯเปฯ.}

\csromanmatikaanchor{mtk.67}
\csromantocmarkref{subsection}{(44) 2. อเสขญาณกถา}{mtk.67}
\bookmark[dest={mtk.67},level=2]{(44) 2. อเสขญาณกถา}
\proseitem{420}{\csromansymbolmark{*}โสตาปตฺติผลสมงฺคิสฺส ปุคฺคลสฺส วิมุตฺติญาณํ วิมุตฺตนฺตํ จ ผลํ ปตฺตสฺส ญาณนฺติ, อามนฺตา.\csromanspacer{} โสตาปตฺติผลสจฺฉิกิริยาย ปฏิปนฺนสฺส \csromanfolio{228}ปุคฺคลสฺส วิมุตฺติญาณํ วิมุตฺตนฺตํ จ ผลํ ปตฺตสฺส ญาณนฺติ, น เหวํ วตฺตพฺเพ ฯเปฯ.\csromanspacer{} สกทาคามิผลสมงฺคิสฺส ปุคฺคลสฺส วิมุตฺติญาณํ วิมุตฺตนฺตํ จ ผลํ ปตฺตสฺส ญาณนฺติ, อามนฺตา.\csromanspacer{} สกทาคามิผลสจฺฉิกิริยาย ปฏิปนฺนสฺส ปุคฺคลสฺส วิมุตฺติญาณํ วิมุตฺตนฺตํ จ ผลํ ปตฺตสฺส ญาณนฺติ, น เหวํ วตฺตพฺเพ ฯเปฯ.\csromanspacer{} อนาคามิผลสมงฺคิสฺส ปุคฺคลสฺส วิมุตฺติญาณํ วิมุตฺตนฺตํ จ ผลํ ปตฺตสฺส ญาณนฺติ, อามนฺตา.\csromanspacer{} อนาคามิผลสจฺฉิกิริยาย ปฏิปนฺนสฺส ปุคฺคลสฺส วิมุตฺติญาณํ วิมุตฺตนฺตํ จ ผลํ ปตฺตสฺส ญาณนฺติ, น เหวํ วตฺตพฺเพ ฯเปฯ.\csromanspacer{} อรหตฺตผลสมงฺคิสฺส ปุคฺคลสฺส วิมุตฺติญาณํ วิมุตฺตนฺตํ จ ผลํ ปตฺตสฺส ญาณนฺติ, อามนฺตา.\csromanspacer{} อรหตฺตสจฺฉิกิริยาย ปฏิปนฺนสฺส ปุคฺคลสฺส วิมุตฺติญาณํ วิมุตฺตนฺตํ จ ผลํ ปตฺตสฺส ญาณนฺติ, น เหวํ วตฺตพฺเพ ฯเปฯ.}

\nitthitam{วิมุตฺติกถา นิฏฺฐิตา.}
\chapterhead{5. ปญฺจมวคฺค}
\vspace{\tipitakaparskip}
\csromancenter{\textbf{(44) 2.}\csromanspacer{}\textbf{ อเสขญาณกถา}}
\proseitem{421}{\csromansymbolmark{*}เสขสฺส อเสขํ ญาณํ อตฺถีติ, อามนฺตา.\csromanspacer{} เสโข อเสขํ ธมฺมํ ชานาติ ปสฺสติ, ทิฏฺฐํ วิทิตํ สจฺฉิกตํ อุปสมฺปชฺช วิหรติ.\csromanspacer{} กาเยน ผุสิตฺวา วิหรตีติ, น เหวํ วตฺตพฺเพ ฯเปฯ.\csromanspacer{} นนุ เสโข อเสขํ ธมฺมํ น ชานาติ น ปสฺสติ อทิฏฺฐํ อวิทิตํ อสจฺฉิกตํ น อุปสมฺปชฺช วิหรติ.\csromanspacer{} น กาเยน ผุสิตฺวา วิหรตีติ, อามนฺตา.\csromanspacer{} หญฺจิ เสโข อเสขํ ธมฺมํ น ชานาติ น ปสฺสติ อทิฏฺฐํ อวิทิตํ อสจฺฉิกตํ น อุปสมฺปชฺช วิหรติ, น กาเยน ผุสิตฺวา วิหรติ.\csromanspacer{} โน จ วต เร วตฺตพฺเพ “เสขสฺส อเสขํ ญาณํ อตฺถี”ติ.}

\prose{\csromansymbolmark{*}อเสขสฺส อเสขํ ญาณํ อตฺถิ อเสโข อเสขํ ธมฺมํ ชานาติ, ปสฺสติ, ทิฏฺฐํ วิทิตํ สจฺฉิกตํ อุปสมฺปชฺช วิหรติ, กาเยน ผุสิตฺวา วิหรตีติ, อามนฺตา.\csromanspacer{} เสขสฺส อเสขํ ญาณํ อตฺถิ เสโข อเสขํ ธมฺมํ ชานาติ ปสฺสติ, ทิฏฺฐํ วิทิตํ สจฺฉิกตํ อุปสมฺปชฺช วิหรติ กาเยน ผุสิตฺวา วิหรตีติ, น เหวํ วตฺตพฺเพ ฯเปฯ.}

\chapterheadpageread{5. ปญฺจมวคฺค}
\csromanmatikaanchor{mtk.68}
\csromantocmarkref{subsection}{(45) 3. วิปรีตกถา}{mtk.68}
\bookmark[dest={mtk.68},level=2]{(45) 3. วิปรีตกถา}
\proseitem{422}{\csromanfolio{229}\csromansymbolmark{*}เสขสฺส อเสขํ ญาณํ อตฺถิ เสโข อเสขํ ธมฺมํ น ชานาติ น ปสฺสติ, อทิฏฺฐํ อวิทิตํ อสจฺฉิกตํ น อุปสมฺปชฺช วิหรติ, น กาเยน ผุสิตฺวา วิหรตีติ, อามนฺตา.\csromanspacer{} อเสขสฺส อเสขํ ญาณํ อตฺถิ อเสโข อเสขํ ธมฺมํ น ชานาติ น ปสฺสติ อทิฏฺฐํ อวิทิตํ อสจฺฉิกตํ น อุปสมฺปชฺช วิหรติ น กาเยน ผุสิตฺวา วิหรตีติ, น เหวํ วตฺตพฺเพ ฯเปฯ.}

\prose{\csromansymbolmark{*}เสขสฺส อเสขํ ญาณํ อตฺถีติ, อามนฺตา.\csromanspacer{} โคตฺรภุโน ปุคฺคลสฺส โสตาปตฺติมคฺเค ญาณํ อตฺถีติ, น เหวํ วตฺตพฺเพ ฯเปฯ.\csromanspacer{} โสตาปตฺติผลสจฺฉิกิริยาย ปฏิปนฺนสฺส ปุคฺคลสฺส โสตาปตฺติผเล ญาณํ อตฺถีติ, น เหวํ วตฺตพฺเพ ฯเปฯ.\csromanspacer{} สกทาคามิผล.\csromanspacer{} อนาคามิผล.\csromanspacer{} อรหตฺตสจฺฉิกิริยาย ปฏิปนฺนสฺส ปุคฺคลสฺส อรหตฺเต ญาณํ อตฺถีติ, น เหวํ วตฺตพฺเพ ฯเปฯ.}

\proseitem{423}{\csromansymbolmark{+}น วตฺตพฺพํ “เสขสฺส อเสขํ ญาณํ อตฺถี”ติ, อามนฺตา.\csromanspacer{} นนุ อายสฺมา อานนฺโท เสโข “ภควา อุฬาโร”ติ ชานาติ, “สาริปุตฺโต เถโร, มหาโมคฺคลฺลาโน เถโร อุฬาโร”ติ ชานาตีติ, อามนฺตา.\csromanspacer{} หญฺจิ อายสฺมา อานนฺโท เสโข “ภควา อุฬาโร”ติ ชานาติ, “สาริปุตฺโต เถโร, มหาโมคฺคลฺลาโน เถโร อุฬาโร”ติ ชานาติ, เตน วต เร วตฺตพฺเพ “เสขสฺส อเสขํ ญาณํ อตฺถี”ติ.}

\nitthitam{อเสขญาณกถา นิฏฺฐิตา.}
\chapterhead{5. ปญฺจมวคฺค}
\vspace{\tipitakaparskip}
\csromancenter{\textbf{(45) 3.}\csromanspacer{}\textbf{ วิปรีตกถา}}
\proseitem{424}{\csromansymbolmark{*}ปถวีกสิณํ สมาปตฺติํ\footnote{ปถวีกสิณสมาปตฺติํ (สี, ก)} สมาปนฺนสฺส วิปรีเต ญาณนฺติ, อามนฺตา.\csromanspacer{} อนิจฺเจ “นิจฺจนฺ”ติ วิปริเยโสติ, น เหวํ วตฺตพฺเพ ฯเปฯ.\csromanspacer{} ทุกฺเข “สุขนฺ”ติ ฯเปฯ.\csromanspacer{} อนตฺตนิ “อตฺตา”ติ ฯเปฯ.\csromanspacer{} อสุเภ “สุภนฺ”ติ วิปริเยโสติ, น เหวํ วตฺตพฺเพ ฯเปฯ.}

\prose{\csromanfolio{230}\csromansymbolmark{*}ปถวีกสิณํ สมาปตฺติํ สมาปนฺนสฺส วิปรีเต ญาณนฺติ, อามนฺตา.\csromanspacer{} อกุสลนฺติ, น เหวํ วตฺตพฺเพ ฯเปฯ.\csromanspacer{} นนุ กุสลนฺติ, อามนฺตา.\csromanspacer{} หญฺจิ กุสลํ, โน จ วต เร วตฺตพฺเพ “ปถวีกสิณํ สมาปตฺติํ สมาปนฺนสฺส วิปรีเต ญาณนฺ”ติ.}

\prose{\csromansymbolmark{*}อนิจฺเจ “นิจฺจนฺ”ติ วิปริเยโส, โส จ อกุสโลติ, อามนฺตา.\csromanspacer{} ปถวีกสิณํ สมาปตฺติํ สมาปนฺนสฺส วิปรีเต ญาณํ ตญฺจ อกุสลนฺติ, น เหวํ วตฺตพฺเพ ฯเปฯ.\csromanspacer{} ทุกฺเข “สุขนฺ”ติ ฯเปฯ.\csromanspacer{} อนตฺตนิ “อตฺตา”ติ ฯเปฯ.\csromanspacer{} อสุเภ “สุภนฺติ” วิปริเยโส, โส จ อกุสโลติ, อามนฺตา.\csromanspacer{} ปถวีกสิณํ สมาปตฺติํ สมาปนฺนสฺส วิปรีเต ญาณํ ตญฺจ อกุสลนฺติ, น เหวํ วตฺตพฺเพ ฯเปฯ.}

\proseitem{425}{\csromansymbolmark{*}ปถวีกสิณํ สมาปตฺติํ สมาปนฺนสฺส วิปรีเต ญาณํ, ตญฺจ อกุสลนฺติ, อามนฺตา.\csromanspacer{} อนิจฺเจ “นิจฺจนฺ”ติ วิปริเยโส, โส จ กุสโลติ, น เหวํ วตฺตพฺเพ ฯเปฯ.\csromanspacer{} ปถวีกสิณํ สมาปตฺติํ สมาปนฺนสฺส วิปรีเต ญาณํ, ตญฺจ อกุสลนฺติ, อามนฺตา.\csromanspacer{} ทุกฺเข “สุขนฺ”ติ ฯเปฯ.\csromanspacer{} อนตฺตนิ “อตฺตา”ติ ฯเปฯ.\csromanspacer{} อสุเภ “สุภนฺ”ติ วิปริเยโส, โส จ กุสโลติ, น เหวํ วตฺตพฺเพ ฯเปฯ.}

\proseitem{426}{\csromansymbolmark{*}ปถวีกสิณํ สมาปตฺติํ สมาปนฺนสฺส วิปรีเต ญาณนฺติ, อามนฺตา.\csromanspacer{} อรหา ปถวีกสิณํ สมาปตฺติํ สมาปชฺเชยฺยาติ, อามนฺตา.\csromanspacer{} หญฺจิ อรหา ปถวีกสิณํ สมาปตฺติํ สมาปชฺเชยฺย, โน จ วต เร วตฺตพฺเพ “ปถวีกสิณํ สมาปตฺติํ สมาปนฺนสฺส วิปรีเต ญาณนฺ”ติ.}

\prose{\csromansymbolmark{*}ปถวีกสิณํ สมาปตฺติํ สมาปนฺนสฺส วิปรีเต ญาณํ, อรหา ปถวีกสิณํ สมาปตฺติํ สมาปชฺเชยฺยาติ, อามนฺตา.\csromanspacer{} อตฺถิ อรหโต วิปริเยโสติ, น เหวํ วตฺตพฺเพ ฯเปฯ.}

\prose{\csromansymbolmark{*}อตฺถิ อรหโต วิปริเยโสติ, อามนฺตา.\csromanspacer{} อตฺถิ อรหโต สญฺญาวิปริเยโส จิตฺตวิปริเยโส ทิฏฺฐิวิปริเยโสติ, น เหวํ วตฺตพฺเพ ฯเปฯ.}

\prose{\csromansymbolmark{*}นตฺถิ อรหโต สญฺญาวิปริเยโส จิตฺตวิปริเยโส ทิฏฺฐิวิปริเยโสติ, อามนฺตา.\csromanspacer{} หญฺจิ นตฺถิ อรหโต สญฺญาวิปริเยโส}

\vspace{\tipitakaparskip}
\csromancenter{ปญฺจมวคฺค จิตฺตวิปริเยโส ทิฏฺฐิวิปริเยโส, โน จ วต เร วตฺตพฺเพ “อตฺถิ อรหโต วิปริเยโส”ติ.}
\csromanmatikaanchor{mtk.69}
\csromantocmarkref{subsection}{(46) 4. นิยามกถา}{mtk.69}
\bookmark[dest={mtk.69},level=2]{(46) 4. นิยามกถา}
\proseitem{427}{\csromanfolio{231}\csromansymbolmark{+}น วตฺตพฺพํ “ปถวีกสิณํ สมาปตฺติํ สมาปนฺนสฺส วิปรีเต ญาณนฺ”ติ, อามนฺตา.\csromanspacer{} ปถวีกสิณํ สมาปตฺติํ สมาปชฺชนฺตสฺส สพฺเพว ปถวีติ, น เหวํ วตฺตพฺเพ.\csromanspacer{} เตน หิ ปถวีกสิณํ สมาปตฺติํ สมาปนฺนสฺส วิปรีเต ญาณนฺติ.}

\prose{\csromansymbolmark{*}ปถวีกสิณํ สมาปตฺติํ สมาปนฺนสฺส วิปรีเต ญาณนฺติ, อามนฺตา.\csromanspacer{} นนุ ปถวี อตฺถิ, อตฺถิ จ โกจิ ปถวิํ ปถวิโต สมาปชฺชตีติ, อามนฺตา.\csromanspacer{} หญฺจิ ปถวี อตฺถิ, อตฺถิ จ โกจิ ปถวิํ ปถวิโต สมาปชฺชติ, โน จ วต เร วตฺตพฺเพ “ปถวีกสิณํ สมาปตฺติํ สมาปนฺนสฺส วิปรีเต ญาณนฺ”ติ.}

\prose{\csromansymbolmark{*}ปถวี อตฺถิ, ปถวิํ ปถวิโต สมาปชฺชนฺตสฺส วิปรีตํ โหตีติ, อามนฺตา.\csromanspacer{} นิพฺพานํ อตฺถิ, นิพฺพานํ นิพฺพานโต สมาปชฺชนฺตสฺส วิปรีตํ โหตีติ, น เหวํ วตฺตพฺเพ ฯเปฯ.\csromanspacer{} เตน หิ น วตฺตพฺพํ “ปถวีกสิณํ สมาปตฺติํ สมาปนฺนสฺส วิปรีเต ญาณนฺ”ติ.}

\nitthitam{วิปรีตกถา นิฏฺฐิตา.}
\chapterhead{5. ปญฺจมวคฺค}
\vspace{\tipitakaparskip}
\csromancenter{\textbf{(46) 4.}\csromanspacer{}\textbf{ นิยามกถา}}
\proseitem{428}{\csromansymbolmark{*}อนิยตสฺส นิยามคมนาย อตฺถิ ญาณนฺติ, อามนฺตา.\csromanspacer{} นิยตสฺส อนิยามคมนาย อตฺถิ ญาณนฺติ, น เหวํ วตฺตพฺเพ ฯเปฯ.}

\prose{\csromansymbolmark{*}นิยตสฺส อนิยามคมนาย นตฺถิ ญาณนฺติ, อามนฺตา.\csromanspacer{} อนิยตสฺส นิยามคมนาย นตฺถิ ญาณนฺติ, น เหวํ วตฺตพฺเพ ฯเปฯ.}

\prose{\csromansymbolmark{*}อนิยตสฺส นิยามคมนาย อตฺถิ ญาณนฺติ, อามนฺตา.\csromanspacer{} นิยตสฺส นิยามคมนาย อตฺถิ ญาณนฺติ, น เหวํ วตฺตพฺเพ ฯเปฯ.}

\prose{\csromansymbolmark{*}นิยตสฺส นิยามคมนาย นตฺถิ ญาณนฺติ, อามนฺตา.\csromanspacer{} อนิยตสฺส นิยามคมนาย นตฺถิ ญาณนฺติ, น เหวํ วตฺตพฺเพ ฯเปฯ.}

\prose{\csromanfolio{232}\csromansymbolmark{*}อนิยตสฺส นิยามคมนาย อตฺถิ ญาณนฺติ, อามนฺตา.\csromanspacer{} อนิยตสฺส อนิยามคมนาย อตฺถิ ญาณนฺติ, น เหวํ วตฺตพฺเพ ฯเปฯ.}

\prose{\csromansymbolmark{*}อนิยตสฺส อนิยามคมนาย นตฺถิ ญาณนฺติ, อามนฺตา.\csromanspacer{} อนิยตสฺส นิยามคมนาย นตฺถิ ญาณนฺติ, น เหวํ วตฺตพฺเพ ฯเปฯ.}

\proseitem{429}{\csromansymbolmark{*}อนิยตสฺส นิยามคมนาย อตฺถิ ญาณนฺติ, อามนฺตา.\csromanspacer{} อนิยตสฺส นิยามคมนาย อตฺถิ นิยาโมติ, น เหวํ วตฺตพฺเพ ฯเปฯ.}

\prose{\csromansymbolmark{*}อนิยตสฺส นิยามคมนาย อตฺถิ ญาณนฺติ, อามนฺตา.\csromanspacer{} อนิยตสฺส นิยามคมนาย อตฺถิ สติปฏฺฐานา.\csromanspacer{} สมฺมปฺปธานา.\csromanspacer{} อิทฺธิปาทา.\csromanspacer{} อินฺทฺริยา.\csromanspacer{} พลา.\csromanspacer{} โพชฺฌงฺคาติ, น เหวํ วตฺตพฺเพ ฯเปฯ.}

\prose{\csromansymbolmark{*}อนิยตสฺส นิยามคมนาย นตฺถิ นิยาโมติ, อามนฺตา.\csromanspacer{} หญฺจิ อนิยตสฺส นิยามคมนาย นตฺถิ นิยาโม, โน วต เร วตฺตพฺเพ “อนิยตสฺส นิยามคมนาย อตฺถิ ญาณนฺ”ติ.}

\prose{\csromansymbolmark{*}อนิยตสฺส นิยามคมนาย นตฺถิ สติปฏฺฐานา ฯเปฯ โพชฺฌงฺคาติ, อามนฺตา.\csromanspacer{} หญฺจิ อนิยตสฺส นิยามคมนาย นตฺถิ โพชฺฌงฺคา, โน วต เร วตฺตพฺเพ “อนิยตสฺส นิยามคมนาย อตฺถิ ญาณนฺ”ติ.}

\proseitem{430}{\csromansymbolmark{*}อนิยตสฺส นิยามคมนาย อตฺถิ ญาณนฺติ, อามนฺตา.\csromanspacer{} โคตฺรภุโน ปุคฺคลสฺส โสตาปตฺติมคฺเค ญาณํ อตฺถีติ, น เหวํ วตฺตพฺเพ ฯเปฯ.\csromanspacer{} โสตาปตฺติผลสจฺฉิกิริยาย ปฏิปนฺนสฺส ปุคฺคลสฺส โสตาปตฺติผเล ญาณํ อตฺถีติ, น เหวํ วตฺตพฺเพ ฯเปฯ.\csromanspacer{} อรหตฺตสจฺฉิกิริยาย ปฏิปนฺนสฺส ปุคฺคลสฺส อรหตฺเต ญาณํ อตฺถีติ, น เหวํ วตฺตพฺเพ ฯเปฯ.}

\proseitem{431}{\csromansymbolmark{+}น วตฺตพฺพํ “อนิยตสฺส นิยามคมนาย อตฺถิ ญาณนฺ”ติ, อามนฺตา.\csromanspacer{} นนุ ภควา ชานาติ”อยํ ปุคฺคโล สมฺมตฺตนิยามํ โอกฺกมิสฺสติ, ภพฺโพ อยํ ปุคฺคโล ธมฺมํ อภิสเมตุนฺ”ติ, อามนฺตา.\csromanspacer{} หญฺจิ ภควา ชานาติ “อยํ ปุคฺคโล สมฺมตฺตนิยามํ โอกฺกมิสฺสติ, ภพฺโพ อยํ ปุคฺคโล ธมฺมํ อภิสเมตุํ”, เตน วต เร วตฺตพฺเพ “อนิยตสฺส นิยามคมนาย อตฺถิ ญาณนฺ”ติ.}

\nitthitam{นิยามกถา นิฏฺฐิตา.}
\chapterheadpageread{5. ปญฺจมวคฺค \textbf{5. ปญฺจมวคฺค}}
\vspace{\tipitakaparskip}
\csromancenter{\textbf{(47) 5.}\csromanspacer{}\textbf{ ปฏิสมฺภิทากถา}}
\csromanmatikaanchor{mtk.70}
\csromanmatikaanchor{mtk.71}
\csromantocmarkref{subsection}{(47) 5. ปฏิสมฺภิทากถา}{mtk.70}
\bookmark[dest={mtk.70},level=2]{(47) 5. ปฏิสมฺภิทากถา}
\csromantocmarkref{subsection}{(48) 6. สมฺมุติญาณกถา}{mtk.71}
\bookmark[dest={mtk.71},level=2]{(48) 6. สมฺมุติญาณกถา}
\proseitem{423}{\csromanfolio{233}\csromansymbolmark{*}สพฺพํ ญาณํ ปฏิสมฺภิทาติ, อามนฺตา.\csromanspacer{} สมฺมุติญาณํ ปฏิสมฺภิทาติ, น เหวํ วตฺตพฺเพ ฯเปฯ.\csromanspacer{} สมฺมุติญาณํ ปฏิสมฺภิทาติ, อามนฺตา.\csromanspacer{} เย เกจิ สมฺมุติํ ชานนฺติ, สพฺเพ เต ปฏิสมฺภิทาปตฺตาติ, น เหวํ วตฺตพฺเพ ฯเปฯ.\csromanspacer{} สพฺพํ ญาณํ ปฏิสมฺภิทาติ, อามนฺตา.\csromanspacer{} เจโตปริยาเย ญาณํ ปฏิสมฺภิทาติ, น เหวํ วตฺตพฺเพ ฯเปฯ.\csromanspacer{} เจโตปริยาเย ญาณํ ปฏิสมฺภิทาติ, อามนฺตา.\csromanspacer{} เย เกจิ ปรจิตฺตํ ชานนฺติ สพฺเพ เต ปฏิสมฺภิทาปตฺตาติ, น เหวํ วตฺตพฺเพ ฯเปฯ.}

\prose{\csromansymbolmark{*}สพฺพํ ญาณํ ปฏิสมฺภิทาติ, อามนฺตา.\csromanspacer{} สพฺพา ปญฺญา ปฏิสมฺภิทาติ, น เหวํ วตฺตพฺเพ ฯเปฯ.\csromanspacer{} สพฺพา ปญฺญา ปฏิสมฺภิทาติ, อามนฺตา.\csromanspacer{} ปถวีกสิณํ สมาปตฺติํ สมาปนฺนสฺส อตฺถิ ปญฺญา, สา ปญฺญา ปฏิสมฺภิทาติ, น เหวํ วตฺตพฺเพ ฯเปฯ.\csromanspacer{} อาโปกสิณํ ฯเปฯ.\csromanspacer{} เตโชกสิณํ ฯเปฯ.\csromanspacer{} วาโยกสิณํ ฯเปฯ.\csromanspacer{} นีลกสิณํ ฯเปฯ.\csromanspacer{} ปีตกสิณํ ฯเปฯ.\csromanspacer{} โลหิตกสิณํ ฯเปฯ.\csromanspacer{} โอทาตกสิณํ ฯเปฯ.\csromanspacer{} อากาสานญฺจายตนํ ฯเปฯ.\csromanspacer{} วิญฺญาณญฺจายตนํ ฯเปฯ.\csromanspacer{} อากิญฺจญฺญายตนํ ฯเปฯ.\csromanspacer{} เนวสญฺญานาสญฺญายตนํ สมาปนฺนสฺส ฯเปฯ.\csromanspacer{} ทานํ ททนฺตสฺส ฯเปฯ.\csromanspacer{} จีวรํ ททนฺตสฺส ฯเปฯ.\csromanspacer{} ปิณฺฑปาตํ ททนฺตสฺส ฯเปฯ.\csromanspacer{} เสนาสนํ ททนฺตสฺส ฯเปฯ.\csromanspacer{} คิลานปจฺจยเภสชฺชปริกฺขารํ ททนฺตสฺส อตฺถิ ปญฺญา, สา ปญฺญา ปฏิสมฺภิทาติ, น เหวํ วตฺตพฺเพ ฯเปฯ.}

\proseitem{433}{\csromansymbolmark{+}น วตฺตพฺพํ “สพฺพํ ญาณํ ปฏิสมฺภิทา”ติ, อามนฺตา.\csromanspacer{} อตฺถิ โลกุตฺตรา ปญฺญา, สา ปญฺญา น ปฏิสมฺภิทาติ, น เหวํ วตฺตพฺเพ ฯเปฯ.\csromanspacer{} เตน หิ สพฺพํ ญาณํ ปฏิสมฺภิทาติ.}

\nitthitam{ปฏิสมฺภิทากถา นิฏฺฐิตา.}
\chapterhead{5. ปญฺจมวคฺค}
\vspace{\tipitakaparskip}
\csromancenter{\textbf{(48) 6.}\csromanspacer{}\textbf{ สมฺมุติญาณกถา}}
\csromanmatikaanchor{mtk.72}
\csromantocmarkref{subsection}{(49) 7. จิตฺตารมฺมณกถา}{mtk.72}
\bookmark[dest={mtk.72},level=2]{(49) 7. จิตฺตารมฺมณกถา}
\proseitem{434}{\csromansymbolmark{+}น วตฺตพฺพํ “สมฺมุติญาณํ สจฺจารมฺมณญฺเญว น อญฺญารมฺมณนฺ”ติ, อามนฺตา.\csromanspacer{} นนุ ปถวีกสิณํ สมาปตฺติํ สมาปนฺนสฺส อตฺถิ ญาณํ, \csromanfolio{234}ปถวีกสิณญฺจ สมฺมุติสจฺจมฺหีติ, อามนฺตา.\csromanspacer{} หญฺจิ ปถวีกสิณํ สมาปตฺติํ สมาปนฺนสฺส อตฺถิ ญาณํ, ปถวีกสิณญฺจ สมฺมุติสจฺจมฺหิ, เตน วต เร วตฺตพฺเพ “สมฺมุติญาณํ สจฺจารมฺมณญฺเญว น อญฺญารมฺมณนฺ”ติ.}

\prose{\csromansymbolmark{+}น วตฺตพฺพํ “สมฺมุติญาณํ สจฺจารมฺมณญฺเญว น อญฺญารมฺมณนฺ”ติ, อามนฺตา ฯเปฯ.\csromanspacer{} นนุ อาโปกสิณํ ฯเปฯ.\csromanspacer{} เตโชกสิณํ ฯเปฯ.\csromanspacer{} คิลานปจฺจยเภสชฺชปริกฺขารํ ททนฺตสฺส อตฺถิ ญาณํ.\csromanspacer{} คิลานปจฺจยเภสชฺชปริกฺขาโร จ สมฺมุติสจฺจมฺหีติ, อามนฺตา.\csromanspacer{} หญฺจิ คิลานปจฺจยเภสชฺชปริกฺขารํ ททนฺตสฺส อตฺถิ ญาณํ, คิลานปจฺจยเภสชฺชปริกฺขาโร จ สมฺมุติสจฺจมฺหิ, เตน วต เร วตฺตพฺเพ “สมฺมุติญาณํ สจฺจารมฺมณญฺเญว น อญฺญารมฺมณนฺ”ติ.}

\proseitem{435}{\csromansymbolmark{*}สมฺมุติญาณํ สจฺจารมฺมณญฺเญว น อญฺญารมฺมณนฺติ, อามนฺตา.\csromanspacer{} เตน ญาเณน ทุกฺขํ ปริชานาติ.\csromanspacer{} สมุทยํ ปชหติ.\csromanspacer{} นิโรธํ สจฺฉิกโรติ.\csromanspacer{} มคฺคํ ภาเวตีติ.\csromanspacer{} น เหวํ วตฺตพฺเพ ฯเปฯ.}

\nitthitam{สมฺมุติญาณกถา นิฏฺฐิตา.}
\chapterhead{5. ปญฺจมวคฺค}
\vspace{\tipitakaparskip}
\csromancenter{\textbf{(49) 7.}\csromanspacer{}\textbf{ จิตฺตารมฺมณกถา}}
\proseitem{436}{\csromansymbolmark{*}เจโตปริยาเย ญาณํ จิตฺตารมฺมณญฺเญว น อญฺญารมฺมณนฺติ, อามนฺตา.\csromanspacer{} นนุ อตฺถิ โกจิ “สราคํ จิตฺตํ สราคํ จิตฺตนฺ”ติ ปชานาตีติ, อามนฺตา.\csromanspacer{} หญฺจิ อตฺถิ โกจิ “สราคํ จิตฺตํ สราคํ จิตฺตนฺ”ติ ปชานาติ, โน จ วต เร วตฺตพฺเพ “เจโตปริยาเย ญาณํ จิตฺตารมฺมณญฺเญว น อญฺญารมฺมณนฺ”ติ.}

\prose{\csromansymbolmark{*}นนุ อตฺถิ โกจิ วีตราคํ จิตฺตํ ฯเปฯ สโทสํ จิตฺตํ, วีตโทสํ จิตฺตํ สโมหํ จิตฺตํ, วีตโมหํ จิตฺตํ, สํขิตฺตํ จิตฺตํ, วิกฺขิตฺตํ จิตฺตํ, มหคฺคตํ จิตฺตํ, อมหคฺคตํ จิตฺตํ, ส-อุตฺตรํ จิตฺตํ, อนุตฺตรํ จิตฺตํ, สมาหิตํ จิตฺตํ, อสมาหิตํ จิตฺตํ, วิมุตฺตํ จิตฺตํ ฯเปฯ อวิมุตฺตํ จิตฺตํ “อวิมุตฺตํ จิตฺตนฺ”ติ ปชานาตีติ, อามนฺตา.\csromanspacer{} หญฺจิ อตฺถิ โกจิ “อวิมุตฺตํ จิตฺตํ อวิมุตฺตํ จิตฺตนฺ”ติ ปชานาติ, โน จ วต เร วตฺตพฺเพ “เจโตปริยาเย ญาณํ จิตฺตารมฺมณญฺเญว น อญฺญารมฺมณนฺ”ติ.}

\chapterheadpageread{5. ปญฺจมวคฺค}
\csromanmatikaanchor{mtk.73}
\csromantocmarkref{subsection}{(50) 8. อนาคตญาณกถา}{mtk.73}
\bookmark[dest={mtk.73},level=2]{(50) 8. อนาคตญาณกถา}
\proseitem{437}{\csromanfolio{235}\csromansymbolmark{*}ผสฺสารมฺมเณ ญาณํ วตฺตพฺพํ “เจโตปริยาเย ญาณนฺ”ติ, อามนฺตา.\csromanspacer{} หญฺจิ ผสฺสารมฺมเณ ญาณํ วตฺตพฺพํ เจโตปริยาเย ญาณํ, โน จ วต เร วตฺตพฺเพ “เจโตปริยาเย ญาณํ จิตฺตารมฺมณญฺเญว น อญฺญารมฺมณนฺ”ติ, เวทนารมฺมเณ ญาณํ ฯเปฯ สญฺญารมฺมเณ ญาณํ, เจตนารมฺมเณ ญาณํ, จิตฺตารมฺมเณ ญาณํ, สทฺธารมฺมเณ ญาณํ, วีริยารมฺมเณ ญาณํ, สตารมฺมเณ ญาณํ, สมาธารมฺมเณ ญาณํ, ปญฺญารมฺมเณ ญาณํ, ราคารมฺมเณ ญาณํ, โทสารมฺมเณ ญาณํ, โมหารมฺมเณ ญาณํ ฯเปฯ อโนตฺตปฺปารมฺมเณ ญาณํ วตฺตพฺพํ “เจโตปริยาเย ญาณนฺ”ติ, อามนฺตา.\csromanspacer{} หญฺจิ อโนตฺตปฺปารมฺมเณ ญาณํ วตฺตพฺพํ เจโตปริยาเย ญาณํ, โน จ วต เร วตฺตพฺเพ “เจโตปริยาเย ญาณํ จิตฺตารมฺมณญฺเญว น อญฺญารมฺมณนฺ”ติ.}

\prose{\csromansymbolmark{*}ผสฺสารมฺมเณ ญาณํ น วตฺตพฺพํ “เจโตปริยาเย ญาณนฺ”ติ, อามนฺตา.\csromanspacer{} ผสฺสปริยาเย ญาณนฺติ, น เหวํ วตฺตพฺเพ ฯเปฯ.\csromanspacer{} เวทนารมฺมเณ ญาณํ ฯเปฯ.\csromanspacer{} สญฺญารมฺมเณ ญาณํ ฯเปฯ.\csromanspacer{} อโนตฺตปฺปารมฺมเณ ญาณํ น วตฺตพฺพํ “เจโตปริยาเย ญาณนฺ”ติ, อามนฺตา.\csromanspacer{} อโนตฺตปฺปปริยาเย ญาณนฺติ, น เหวํ วตฺตพฺเพ ฯเปฯ.}

\proseitem{438}{\csromansymbolmark{+}น วตฺตพฺพํ “เจโตปริยาเย ญาณํ จิตฺตารมฺมณญฺเญว น อญฺญารมฺมณนฺ”ติ, อามนฺตา.\csromanspacer{} นนุ เจโตปริยาเย ญาณนฺติ, อามนฺตา.\csromanspacer{} หญฺจิ เจโตปริยาเย ญาณํ, เตน วต เร วตฺตพฺเพ “เจโตปริยาเย ญาณํ จิตฺตารมฺมณญฺเญว น อญฺญารมฺมณนฺ”ติ.}

\nitthitam{จิตฺตารมฺมณกถา นิฏฺฐิตา.}
\chapterhead{5. ปญฺจมวคฺค}
\vspace{\tipitakaparskip}
\csromancenter{\textbf{(50) 8.}\csromanspacer{}\textbf{ อนาคตญาณกถา}}
\csromanmatikaanchor{mtk.74}
\csromantocmarkref{subsection}{(51) 9. ปฏุปฺปนฺนกถา}{mtk.74}
\bookmark[dest={mtk.74},level=2]{(51) 9. ปฏุปฺปนฺนกถา}
\proseitem{439}{\csromansymbolmark{*}อนาคเต ญาณํ อตฺถีติ, อามนฺตา.\csromanspacer{} อนาคตํ มูลโต ชานาติ, เหตุโต ชานาติ, นิทานโต ชานาติ, สมฺภวโต ชานาติ, ปภวโต ชานาติ, สมุฏฺฐานโต ชานาติ, อาหารโต \csromanfolio{236}ชานาติ, อารมฺมณโต ชานาติ, ปจฺจยโต ชานาติ, สมุทยโต ชานาตีติ, น เหวํ วตฺตพฺเพ ฯเปฯ.}

\prose{\csromansymbolmark{*}อนาคเต ญาณํ อตฺถีติ, อามนฺตา.\csromanspacer{} อนาคตํ เหตุปจฺจยตํ ชานาติ, อารมฺมณปจฺจยตํ ชานาติ, อธิปติปจฺจยตํ ชานาติ, อนนฺตรปจฺจยตํ ชานาติ, สมนนฺตรปจฺจยตํ ชานาตีติ, น เหวํ วตฺตพฺเพ ฯเปฯ.}

\prose{\csromansymbolmark{*}อนาคเต ญาณํ อตฺถีติ, อามนฺตา.\csromanspacer{} โคตฺรภุโน ปุคฺคลสฺส โสตาปตฺติมคฺเค ญาณํ อตฺถีติ, น เหวํ วตฺตพฺเพ ฯเปฯ.\csromanspacer{} โสตาปตฺติผลสจฺฉิกิริยาย ปฏิปนฺนสฺส ปุคฺคลสฺส โสตาปตฺติผเล ญาณํ อตฺถีติ, น เหวํ วตฺตพฺเพ ฯเปฯ.}

\prose{\csromansymbolmark{*}สกทาคามิ ฯเปฯ.\csromanspacer{} อนาคามิ ฯเปฯ.\csromanspacer{} อรหตฺตสจฺฉิกิริยาย ปฏิปนฺนสฺส ปุคฺคลสฺส อรหตฺเต ญาณํ อตฺถีติ, น เหวํ วตฺตพฺเพ ฯเปฯ.}

\proseitem{440}{\csromansymbolmark{+}น วตฺตพฺพํ “อนาคเต ญาณํ อตฺถิ”ติ, อามนฺตา.\csromanspacer{} นนุ วุตฺตํ ภควตา “ปาฏลิปุตฺตสฺส โข อานนฺท ตโย อนฺตรายา ภวิสฺสนฺติ อคฺคิโต วา อุทกโต วา มิถุเภทา วา”ติ\footnote{วิ 3. 323; ที 2. 74; ขุ 1. 187 ปิฏฺเฐสุ.} อตฺเถว สุตฺตนฺโตติ, อามนฺตา.\csromanspacer{} เตน หิ อนาคเต ญาณํ อตฺถีติ.}

\nitthitam{อนาคตญาณกถา นิฏฺฐิตา.}
\chapterhead{5. ปญฺจมวคฺค}
\vspace{\tipitakaparskip}
\csromancenter{\textbf{(51) 9.}\csromanspacer{}\textbf{ ปฏุปฺปนฺนกถา}}
\proseitem{441}{\csromansymbolmark{*}ปฏุปฺปนฺเน\footnote{ปจฺจุปฺปนฺเน (สี, สฺยา)} ญาณํ อตฺถีติ, อามนฺตา.\csromanspacer{} เตน ญาเณน ตํ ญาณํ, ชานาตีติ, น เหวํ วตฺตพฺเพ ฯเปฯ.\csromanspacer{} เตน ญาเณน ตํ ญาณํ ชานาตีติ อามนฺตา.\csromanspacer{} เตน ญาเณน ตํ ญาณํ “ญาณนฺ”ติ ชานาตีติ, น เหวํ วตฺตพฺเพ ฯเปฯ.\csromanspacer{} เตน ญาเณน ตํ ญาณํ “ญาณนฺ”ติ ชานาตีติ, อามนฺตา.\csromanspacer{} ตํ ญาณํ ตสฺส ญาณสฺส อารมฺมณนฺติ, น เหวํ วตฺตพฺเพ ฯเปฯ.}

\chapterheadpageread{5. ปญฺจมวคฺค}
\csromanmatikaanchor{mtk.75}
\csromantocmarkref{subsection}{(52) 10. ผลญาณกถา}{mtk.75}
\bookmark[dest={mtk.75},level=2]{(52) 10. ผลญาณกถา}
\prose{\csromanfolio{237}\csromansymbolmark{*}ตํ ญาณํ ตสฺส ญาณสฺส อารมฺมณนฺติ, อามนฺตา.\csromanspacer{} เตน ผสฺเสน ตํ ผสฺสํ ผุสติ.\csromanspacer{} ตาย เวทนาย ตํ เวทนํ เวเทติ.\csromanspacer{} ตาย สญฺญาย ตํ สญฺญํ สญฺชานาติ.\csromanspacer{} ตาย เจตนาย ตํ เจตนํ เจเตติ.\csromanspacer{} เตน จิตฺเตน ตํ จิตฺตํ จินฺเตติ.\csromanspacer{} เตน วิตกฺเกน ตํ วิตกฺกํ วิตกฺเกติ.\csromanspacer{} เตน วิจาเรน ตํ วิจารํ วิจาเรติ.\csromanspacer{} ตาย ปีติยา ตํ ปีติํ ปิยายติ.\csromanspacer{} ตาย สติยา ตํ สติํ สรติ.\csromanspacer{} ตาย ปญฺญาย ตํ ปญฺญํ ปชานาติ.\csromanspacer{} เตน ขคฺเคน ตํ ขคฺคํ ฉินฺทติ.\csromanspacer{} เตน ผรสุนา ตํ ผรสุํ ตจฺฉติ.\csromanspacer{} ตาย กุธาริยา ตํ กุธาริํ ตจฺฉติ.\csromanspacer{} ตาย วาสิยา ตํ วาสิํ ตจฺฉติ.\csromanspacer{} ตาย สูจิยา ตํ สูจิํ สิพฺเพติ.\csromanspacer{} เตน องฺคุลคฺเคน ตํ องฺคุลคฺคํ ปรามสติ.\csromanspacer{} เตน นาสิกคฺเคน ตํ นาสิกคฺคํ ปรามสติ.\csromanspacer{} เตน มตฺถเกน ตํ มตฺถกํ ปรามสติ.\csromanspacer{} เตน คูเถน ตํ คูถํ โธวติ.\csromanspacer{} เตน มุตฺเตน ตํ มุตฺตํ โธวติ.\csromanspacer{} เตน เขเฬน ตํ เขฬํ โธวติ.\csromanspacer{} เตน ปุพฺเพน ตํ ปุพฺพํ โธวติ.\csromanspacer{} เตน โลหิเตน ตํ โลหิตํ โธวตีติ.\csromanspacer{} น เหวํ วตฺตพฺเพ ฯเปฯ.}

\proseitem{442}{\csromansymbolmark{+}น วตฺตพฺพํ “ปฏุปฺปนฺเน ญาณํ อตฺถี”ติ, อามนฺตา.\csromanspacer{} นนุ สพฺพสงฺขาเร อนิจฺจโต ทิฏฺเฐ ตมฺปิ ญาณํ อนิจฺจโต ทิฏฺฐํ โหตีติ, อามนฺตา.\csromanspacer{} หญฺจิ สพฺพสงฺขาเร อนิจฺจโต ทิฏฺเฐ ตมฺปิ ญาณํ อนิจฺจโต ทิฏฺฐํ โหติ, เตน วต เร วตฺตพฺเพ “ปฏุปฺปนฺเน ญาณํ อตฺถี”ติ.}

\nitthitam{ปฏุปฺปนฺนญาณกถา\footnote{ปจฺจุปฺปนฺนกถา (สี)} นิฏฺฐิตา.}
\chapterhead{5. ปญฺจมวคฺค}
\vspace{\tipitakaparskip}
\csromancenter{\textbf{(52) 10.}\csromanspacer{}\textbf{ ผลญาณกถา}}
\proseitem{443}{\csromansymbolmark{*}สาวกสฺส ผเล ญาณํ อตฺถีติ, อามนฺตา.\csromanspacer{} สาวโก ผลสฺส กตํ ปญฺญาเปตีติ.\csromanspacer{} น เหวํ วตฺตพฺเพ ฯเปฯ.}

\prose{\csromansymbolmark{*}สาวกสฺส ผเล ญาณํ อตฺถีติ, อามนฺตา.\csromanspacer{} อตฺถิ สาวกสฺส ผลปโรปริยตฺติ อินฺทฺริยปโรปริยตฺติ ปุคฺคลปโรปริยตฺตีติ, น เหวํ วตฺตพฺเพ ฯเปฯ.}

\prose{\csromanfolio{238}\csromansymbolmark{*}สาวกสฺส ผเล ญาณํ อตฺถีติ, อามนฺตา.\csromanspacer{} อตฺถิ สาวกสฺส ขนฺธปญฺญตฺติ อายตนปญฺญตฺติ ธาตุปญฺญตฺติ สจฺจปญฺญตฺติ อินฺทฺริยปญฺญตฺติ ปุคฺคลปญฺญตฺตีติ, น เหวํ วตฺตพฺเพ ฯเปฯ.}

\prose{\csromansymbolmark{*}สาวกสฺส ผเล ญาณํ อตฺถีติ, อามนฺตา.\csromanspacer{} สาวโก ชิโน สตฺถา สมฺมาสมฺพุทฺโธ สพฺพญฺญู สพฺพทสฺสาวี ธมฺมสฺสามี ธมฺมปฺปฏิสรโณติ, น เหวํ วตฺตพฺเพ ฯเปฯ.}

\prose{\csromansymbolmark{*}สาวกสฺส ผเล ญาณํ อตฺถีติ, อามนฺตา.\csromanspacer{} สาวโก อนุปฺปนฺนสฺส มคฺคสฺส อุปฺปาเทตา อสญฺชาตสฺส มคฺคสฺส สญฺชเนตา อนกฺขาตสฺส มคฺคสฺส อกฺขาตา มคฺคญฺญู มคฺควิทู มคฺคโกวิโทติ, น เหวํ วตฺตพฺเพ ฯเปฯ.}

\proseitem{444}{\csromansymbolmark{+}น วตฺตพฺพํ “สาวกสฺส ผเล ญาณํ อตฺถี”ติ, อามนฺตา.\csromanspacer{} สาวโก อญฺญาณีติ, น เหวํ วตฺตพฺเพ ฯเปฯ.\csromanspacer{} เตน หิ สาวกสฺส ผเล ญาณํ อตฺถีติ ฯเปฯ.\csromanspacer{} ผลญาณกถา นิฏฺฐิตา.\csromanspacer{} ปญฺจโม วคฺโค.}

\titlehead{\textbf{ตสฺสุทฺทานํ}}
\prose{วิมุตฺติญาณํ วิมุตฺตํ, เสขสฺส อเสขํ ญาณํ, วิปรีเต ญาณํ, อนิยตสฺส นิยามคมนาย อตฺถิ ญาณํ, สพฺพํ ญาณํ ปฏิสมฺภิทาติ, สมฺมุติญาณํ, เจโตปริยาเย ญาณํ.\csromanspacer{} อนาคเต ญาณํ, ปฏุปฺปนฺเน ญาณํ, สาวกสฺส ผเล ญาณนฺติ.\csromanspacer{} มหาปณฺณาสโก.}

\csromanheader{2}{\textbf{ตสฺสาปิ อุทฺทานํ}}
\csromangathameasure{สตฺตุปลทฺธิํ, อุปหรโต, พลํ, คิหิสฺส อรหา จ, วิมุตฺติปญฺจมนฺติ.}
\csromangathagroup{\csromangathastanza{สตฺตุปลทฺธิํ, อุปหรโต, พลํ, คิหิสฺส อรหา จ, วิมุตฺติปญฺจมนฺติ.}}

\csromanmatikaanchor{mtk.76}
\csromantocmarknopage{chapter}{6. ฉฏฺฐวคฺค}{mtk.76}
\bookmark[dest={mtk.76},level=0]{6. ฉฏฺฐวคฺค}
\chapterheadpageread{6. ฉฏฺฐวคฺค \textbf{6. ฉฏฺฐวคฺค}}
\vspace{\tipitakaparskip}
\csromancenter{\textbf{(53) 1.}\csromanspacer{}\textbf{ นิยามกถา}}
\csromanmatikaanchor{mtk.77}
\csromantocmarkref{subsection}{(53) 1. นิยามกถา}{mtk.77}
\bookmark[dest={mtk.77},level=2]{(53) 1. นิยามกถา}
\proseitem{445}{\csromanfolio{239}\csromansymbolmark{*}นิยาโม อสงฺขโตติ, อามนฺตา.\csromanspacer{} นิพฺพานํ ตาณํ เลณํ สรณํ ปรายนํ อจฺจุตํ อมตนฺติ, น เหวํ วตฺตพฺเพ ฯเปฯ.}

\prose{\csromansymbolmark{*}นิยาโม อสงฺขโต, นิพฺพานํ อสงฺขตนฺติ, อามนฺตา.\csromanspacer{} เทฺว อสงฺขตานีติ, น เหวํ วตฺตพฺเพ ฯเปฯ.}

\prose{\csromansymbolmark{*}เทฺว อสงฺขตานีติ, อามนฺตา.\csromanspacer{} เทฺว ตาณานิ เทฺว เลณานิ เทฺว สรณานิ เทฺว ปรายนานิ เทฺว อจฺจุตานิ เทฺว อมตานิ เทฺว นิพฺพานานีติ, น เหวํ วตฺตพฺเพ ฯเปฯ.}

\prose{\csromansymbolmark{*}เทฺว นิพฺพานานีติ, อามนฺตา.\csromanspacer{} อตฺถิ ทฺวินฺนํ นิพฺพานานํ อุจฺจนีจตา หีนปณีตตา อุกฺกํสาวกํโส สีมา วา เภโท วา ราชิ วา อนฺตริกา วาติ, น เหวํ วตฺตพฺเพ ฯเปฯ.}

\prose{\csromansymbolmark{*}นิยาโม อสงฺขโตติ, อามนฺตา.\csromanspacer{} อตฺถิ เกจิ นิยามํ โอกฺกมนฺติ ปฏิลภนฺติ อุปฺปาเทนฺติ สมุปฺปาเทนฺติ อุฏฺฐาเปนฺติ สมุฏฺฐาเปนฺติ นิพฺพตฺเตนฺติ อภินิพฺพตฺเตนฺติ ชเนนฺติ สญฺชเนนฺตีติ, อามนฺตา.\csromanspacer{} อตฺถิ เกจิ อสงฺขตํ โอกฺกมนฺติ ปฏิลภนฺติ อุปฺปาเทนฺติ สมุปฺปาเทนฺติ อุฏฺฐาเปนฺติ สมุฏฺฐาเปนฺติ นิพฺพตฺเตนฺติ อภินิพฺพตฺเตนฺติ ชเนนฺติ สญฺชเนนฺตีติ, น เหวํ วตฺตพฺเพ ฯเปฯ.}

\proseitem{446}{\csromansymbolmark{*}นิยาโม อสงฺขโตติ, อามนฺตา.\csromanspacer{} มคฺโค อสงฺขโตติ, น เหวํ วตฺตพฺเพ ฯเปฯ.}

\prose{\csromansymbolmark{*}มคฺโค สงฺขโตติ, อามนฺตา.\csromanspacer{} นิยาโม สงฺขโตติ, น เหวํ วตฺตพฺเพ ฯเปฯ.}

\prose{\csromansymbolmark{*}โสตาปตฺตินิยาโม อสงฺขโตติ, อามนฺตา.\csromanspacer{} โสตาปตฺติมคฺโค อสงฺขโตติ, น เหวํ วตฺตพฺเพ ฯเปฯ.}

\prose{\csromansymbolmark{*}โสตาปตฺติมคฺโค สงฺขโตติ, อามนฺตา.\csromanspacer{} โสตาปตฺตินิยาโม สงฺขโตติ, น เหวํ วตฺตพฺเพ ฯเปฯ.}

\csromanmatikaanchor{mtk.78}
\csromantocmarkref{subsection}{(54) 2. ปฏิจฺจสมุปฺปาทกถา}{mtk.78}
\bookmark[dest={mtk.78},level=2]{(54) 2. ปฏิจฺจสมุปฺปาทกถา}
\prose{\csromansymbolmark{*}สกทาคามินิยาโม ฯเปฯ.\csromanspacer{} อนาคามินิยาโม ฯเปฯ.\csromanspacer{} อรหตฺตนิยาโม อสงฺขโตติ, อามนฺตา.\csromanspacer{} อรหตฺตมคฺโค อสงฺขโตติ, น เหวํ \csromanfolio{240}วตฺตพฺเพ ฯเปฯ.\csromanspacer{} อรหตฺตมคฺโค สงฺขโตติ, อามนฺตา.\csromanspacer{} อรหตฺตนิยาโม สงฺขโตติ, น เหวํ วตฺตพฺเพ ฯเปฯ.}

\prose{\csromansymbolmark{*}โสตาปตฺตินิยาโม อสงฺขโต ฯเปฯ.\csromanspacer{} อรหตฺตนิยาโม อสงฺขโต นิพฺพานํ อสงฺขตนฺติ, อามนฺตา.\csromanspacer{} ปญฺจ อสงฺขตานีติ, น เหวํ วตฺตพฺเพ ฯเปฯ.\csromanspacer{} ปญฺจ อสงฺขตานีติ, อามนฺตา.\csromanspacer{} ปญฺจ ตาณานิ ฯเปฯ.\csromanspacer{} อนฺตริกา วาติ, น เหวํ วตฺตพฺเพ ฯเปฯ.}

\prose{\csromansymbolmark{*}นิยาโม อสงฺขโตติ, อามนฺตา.\csromanspacer{} มิจฺฉตฺตนิยาโม อสงฺขโตติ, น เหวํ วตฺตพฺเพ ฯเปฯ.\csromanspacer{} มิจฺฉตฺตนิยาโม สงฺขโตติ, อามนฺตา.\csromanspacer{} สมฺมตฺตนิยาโม สงฺขโตติ, น เหวํ วตฺตพฺเพ ฯเปฯ.}

\proseitem{447}{\csromansymbolmark{+}น วตฺตพฺพํ “นิยาโม อสงฺขโต”ติ, อามนฺตา.\csromanspacer{} นิยาเม อุปฺปชฺช นิรุทฺเธ อนิยโต โหตีติ, น เหวํ วตฺตพฺเพ ฯเปฯ.\csromanspacer{} เตน หิ นิยาโม อสงฺขโตติ.\csromanspacer{} มิจฺฉตฺตนิยาเม อุปฺปชฺช นิรุทฺเธ อนิยโต โหตีติ, น เหวํ วตฺตพฺเพ ฯเปฯ.\csromanspacer{} เตน หิ มิจฺฉตฺตนิยาโม อสงฺขโตติ.}

\nitthitam{นิยามกถา นิฏฺฐิตา.}
\chapterhead{6. ฉฏฺฐวคฺค}
\vspace{\tipitakaparskip}
\csromancenter{\textbf{(54) 2.}\csromanspacer{}\textbf{ ปฏิจฺจสมุปฺปาทกถา}}
\proseitem{448}{\csromansymbolmark{*}ปฏิจฺจสมุปฺปาโท อสงฺขโตติ, อามนฺตา.\csromanspacer{} นิพฺพานํ ตาณํ เลณํ สรณํ ปรายนํ อจฺจุตํ อมตนฺติ, น เหวํ วตฺตพฺเพ ฯเปฯ.\csromanspacer{} ปฏิจฺจสมุปฺปาโท อสงฺขโต, นิพฺพานํ อสงฺขตนฺติ, อามนฺตา.\csromanspacer{} เทฺว อสงฺขตานีติ, น เหวํ วตฺตพฺเพ ฯเปฯ.\csromanspacer{} เทฺว อสงฺขตานีติ, อามนฺตา.\csromanspacer{} เทฺว ตาณานิ เทฺว เลณานิ เทฺว สรณานิ เทฺว ปรายนานิ เทฺว อจฺจุตานิ เทฺว อมตานิ เทฺว นิพฺพานานีติ, น เหวํ วตฺตพฺเพ ฯเปฯ.\csromanspacer{} เทฺว นิพฺพานานีติ, อามนฺตา.\csromanspacer{} อตฺถิ ทฺวินฺนํ นิพฺพานานํ อุจฺจนีจตา หีนปณีตตา อุกฺกํสาวกํโส สีมา วา เภโท วา ราชิ วา อนฺตริกา วาติ, น เหวํ วตฺตพฺเพ ฯเปฯ.}

\proseitem{449}{\csromansymbolmark{*}ปฏิจฺจสมุปฺปาโท อสงฺขโตติ, อามนฺตา.\csromanspacer{} อวิชฺชา อสงฺขตาติ, น เหวํ วตฺตพฺเพ ฯเปฯ.\csromanspacer{} อวิชฺชา สงฺขตาติ, อามนฺตา.\csromanspacer{} ปฏิจฺจสมุปฺปาโท}

\vspace{\tipitakaparskip}
\csromancenter{ฉฏฺฐวคฺค สงฺขโตติ, น เหวํ วตฺตพฺเพ ฯเปฯ.\csromanspacer{} ปฏิจฺจสมุปฺปาโท อสงฺขโตติ, อามนฺตา.\csromanspacer{} อวิชฺชาปจฺจยา สงฺขารา อสงฺขตาติ, น เหวํ วตฺตพฺเพ ฯเปฯ.\csromanspacer{} อวิชฺชาปจฺจยา สงฺขารา สงฺขตาติ, อามนฺตา.\csromanspacer{} ปฏิจฺจสมุปฺปาโท สงฺขโตติ, น เหวํ วตฺตพฺเพ ฯเปฯ.\csromanspacer{} ปฏิจฺจสมุปฺปาโท อสงฺขโตติ, อามนฺตา.\csromanspacer{} สงฺขารปจฺจยา วิญฺญาณํ อสงฺขตนฺติ, น เหวํ วตฺตพฺเพ ฯเปฯ.\csromanspacer{} สงฺขารปจฺจยา วิญฺญาณํ สงฺขตนฺติ, อามนฺตา.\csromanspacer{} ปฏิจฺจสมุปฺปาโท สงฺขโตติ, น เหวํ วตฺตพฺเพ ฯเปฯ.\csromanspacer{} ปฏิจฺจสมุปฺปาโท อสงฺขโตติ, อามนฺตา.\csromanspacer{} วิญฺญาณปจฺจยา นามรูปํ อสงฺขตนฺติ, น เหวํ วตฺตพฺเพ ฯเปฯ.\csromanspacer{} วิญฺญาณปจฺจยา นามรูปํ สงฺขตนฺติ, อามนฺตา.\csromanspacer{} ปฏิจฺจสมุปฺปาโท สงฺขโตติ, น เหวํ วตฺตพฺเพ ฯเปฯ.\csromanspacer{} ปฏิจฺจสมุปฺปาโท อสงฺขโตติ, อามนฺตา.\csromanspacer{} ชาติปจฺจยา ชรามรณํ อสงฺขตนฺติ, น เหวํ วตฺตพฺเพ ฯเปฯ.\csromanspacer{} ชาติปจฺจยา ชรามรณํ สงฺขตนฺติ, อามนฺตา.\csromanspacer{} ปฏิจฺจสมุปฺปาโท สงฺขโตติ, น เหวํ วตฺตพฺเพ ฯเปฯ.}
\proseitem{450}{\csromanfolio{241}\csromansymbolmark{+}น วตฺตพฺพํ “ปฏิจฺจสมุปฺปาโท อสงฺขโต”ติ, อามนฺตา.\csromanspacer{} นนุ วุตฺตํ ภควตา “ชาติปจฺจยา ภิกฺขเว ชรามรณํ อุปฺปาทา วา ตถาคตานํ อนุปฺปาทา วา ตถาคตานํ ฐิตาว สา ธาตุ ธมฺมฏฺฐิตตา ธมฺมนิยามตา อิทปฺปจฺจยตา, ตํ ตถาคโต อภิสมฺพุชฺฌติ อภิสเมติ, อภิสมฺพุชฺฌิตฺวา อภิสเมตฺวา อาจิกฺขติ เทเสติ ปญฺญเปติ ปฏฺฐเปติ วิวรติ วิภชติ อุตฺตานิํ กโรติ.\csromanspacer{} ปสฺสถาติ จาห.\csromanspacer{} ชาติปจฺจยา ภิกฺขเว ชารามรณํ.\csromanspacer{} ภวปจฺจยา ภิกฺขเว ชาติ ฯเปฯ.\csromanspacer{} อวิชฺชาปจฺจยา ภิกฺขเว สงฺขารา อุปฺปาทา วา ตถาคตานํ อนุปฺปาทา วา ตถาคตานํ ฐิตาว สา ธาตุ ฯเปฯ.\csromanspacer{} ปสฺสถาติ จาห.\csromanspacer{} อวิชฺชาปจฺจยา ภิกฺขเว สงฺขารา, อิติ โข ภิกฺขเว ยา ตตฺร ตถตา อวิตถตา อนญฺญถตา อิทปฺปจฺจยตา, อยํ วุจฺจติ ภิกฺขเว ปฏิจฺจสมุปฺปาโท”ติ\footnote{สํ 1. 264 ปิฏฺเฐ.} อตฺเถว สุตฺตนฺโตติ, อามนฺตา.\csromanspacer{} เตน หิ ปฏิจฺจสมุปฺปาโท อสงฺขโตติ.}

\proseitem{451}{\csromansymbolmark{*}อวิชฺชาปจฺจยา สงฺขาราติ ยา ตตฺถ ธมฺมฏฺฐิตตา ธมฺมนิยามตา อสงฺขตา นิพฺพานํ อสงฺขตนฺติ, อามนฺตา.\csromanspacer{} เทฺว อสงฺขตานีติ, น เหวํ วตฺตพฺเพ ฯเปฯ.\csromanspacer{} เทฺว อสงฺขตานีติ, อามนฺตา.\csromanspacer{} เทฺว ตาณานิ ฯเปฯ.\csromanspacer{} อนฺตริกา วาติ, น เหวํ วตฺตพฺเพ ฯเปฯ.}

\csromanmatikaanchor{mtk.79}
\csromantocmarkref{subsection}{(55) 3. สจฺจกถา}{mtk.79}
\bookmark[dest={mtk.79},level=2]{(55) 3. สจฺจกถา}
\prose{\csromanfolio{242}\csromansymbolmark{*}อวิชฺชาปจฺจยา สงฺขาราติ ยา ตตฺถ ธมฺมฏฺฐิตตา ธมฺมนิยามตา อสงฺขตา, สงฺขารปจฺจยา วิญฺญาณนฺติ ยา ตตฺถ ธมฺมฏฺฐิตตา ธมฺมนิยามตา อสงฺขตา, นิพฺพานํ อสงฺขตนฺติ, อามนฺตา.\csromanspacer{} ตีณิ อสงฺขตานีติ, น เหวํ วตฺตพฺเพ ฯเปฯ.\csromanspacer{} ตีณิ อสงฺขตานีติ, อามนฺตา.\csromanspacer{} ตีณิ ตาณานิ ฯเปฯ.\csromanspacer{} อนฺตริกา วาติ, น เหวํ วตฺตพฺเพ ฯเปฯ.}

\prose{\csromansymbolmark{*}อวิชฺชาปจฺจยา สงฺขาราติ ยา ตตฺถ ธมฺมฏฺฐิตตา ธมฺมนิยามตา อสงฺขตา, สงฺขารปจฺจยา วิญฺญาณนฺติ ยา ตตฺถ ธมฺมฏฺฐิตตา ธมฺมนิยามตา อสงฺขตา ฯเปฯ.\csromanspacer{} ชาติปจฺจยา ชรามรณนฺติ ยา ตตฺถ ธมฺมฏฺฐิตตา ธมฺมนิยามตา อสงฺขตา, นิพฺพานํ อสงฺขตนฺติ, อามนฺตา.\csromanspacer{} ทฺวาทส อสงฺขตานีติ, น เหวํ วตฺตพฺเพ ฯเปฯ.\csromanspacer{} ทฺวาทส อสงฺขตานีติ, อามนฺตา.\csromanspacer{} ทฺวาทส ตาณานิ, ทฺวาทส เลณานิ ฯเปฯ.\csromanspacer{} อนฺตริกา วาติ, น เหวํ วตฺตพฺเพ ฯเปฯ.}

\nitthitam{ปฏิจฺจสมุปฺปาทกถา นิฏฺฐิตา.}
\chapterhead{6. ฉฏฺฐวคฺค}
\vspace{\tipitakaparskip}
\csromancenter{\textbf{(55) 3.}\csromanspacer{}\textbf{ สจฺจกถา}}
\proseitem{452}{\csromansymbolmark{*}จตฺตาริ สจฺจานิ อสงฺขตานีติ, อามนฺตา.\csromanspacer{} จตฺตาริ ตาณานิ จตฺตาริ เลณานิ จตฺตาริ สรณานิ จตฺตาริ ปรายนานิ จตฺตาริ อจฺจุตานิ จตฺตาริ อมตานิ จตฺตาริ นิพฺพานานีติ, น เหวํ วตฺตพฺเพ ฯเปฯ.\csromanspacer{} จตฺตาริ นิพฺพานานีติ, อามนฺตา.\csromanspacer{} อตฺถิ จตุนฺนํ นิพฺพานานํ อุจฺจนีจตา หีนปณีตตา อุกฺกํสาวกํโส สีมา วา เภโท วา ราชิ วา อนฺตริกา วาติ, น เหวํ วตฺตพฺเพ ฯเปฯ.}

\prose{\csromansymbolmark{*}ทุกฺขสจฺจํ อสงฺขตนฺติ, อามนฺตา.\csromanspacer{} ทุกฺขํ อสงฺขตนฺติ, น เหวํ วตฺตพฺเพ ฯเปฯ.\csromanspacer{} ทุกฺขสจฺจํ อสงฺขตนฺติ, อามนฺตา.\csromanspacer{} กายิกํ ทุกฺขํ เจตสิกํ ทุกฺขํ โสกปริเทวทุกฺขโทมนสฺส-อุปายาสา อสงฺขตาติ, น เหวํ วตฺตพฺเพ ฯเปฯ.\csromanspacer{} สมุทยสจฺจํ อสงฺขตนฺติ, อามนฺตา.\csromanspacer{} สมุทโย อสงฺขโตติ, น เหวํ วตฺตพฺเพ ฯเปฯ.\csromanspacer{} สมุทยสจฺจ อสงฺขตนฺติ, อามนฺตา.\csromanspacer{} กามตณฺหา ภวตณฺหา วิภวตณฺหา อสงฺขตาติ, น เหวํ วตฺตพฺเพ ฯเปฯ.\csromanspacer{} มคฺคสจฺจํ อสงฺขตนฺติ, อามนฺตา.\csromanspacer{} มคฺโค อสงฺขโตติ, น เหวํ วตฺตพฺเพ ฯเปฯ.\csromanspacer{} มคฺคสจฺจํ อสงฺขตนฺติ, อามนฺตา.\csromanspacer{} สมฺมาทิฏฺฐิ ฯเปฯ สมฺมาสมาธิ อสงฺขโตติ, น เหวํ วตฺตพฺเพ ฯเปฯ.}

\chapterheadpageread{6. ฉฏฺฐวคฺค}
\prose{\csromanfolio{243}\csromansymbolmark{*}ทุกฺขํ สงฺขตนฺติ, อามนฺตา.\csromanspacer{} ทุกฺขสจฺจํ สงฺขตนฺติ, น เหวํ วตฺตพฺเพ ฯเปฯ.\csromanspacer{} กายิกํ ทุกฺขํ เจตสิกํ ทุกฺขํ โสกปริเทวทุกฺขโทมนสฺส-อุปายาสา สงฺขตาติ, อามนฺตา.\csromanspacer{} ทุกฺขสจฺจํ สงฺขตนฺติ, น เหวํ วตฺตพฺเพ ฯเปฯ.\csromanspacer{} สมุทโย สงฺขโตติ, อามนฺตา.\csromanspacer{} สมุทยสจฺจํ สงฺขตนฺติ, น เหวํ วตฺตพฺเพ ฯเปฯ.\csromanspacer{} กามตณฺหา ภวตณฺหา วิภวตณฺหา สงฺขตาติ, อามนฺตา.\csromanspacer{} สมุทยสจฺจํ สงฺขตนฺติ, น เหวํ วตฺตพฺเพ ฯเปฯ.\csromanspacer{} มคฺโค สงฺขโตติ, อามนฺตา.\csromanspacer{} มคฺคสจฺจํ สงฺขตนฺติ, น เหวํ วตฺตพฺเพ ฯเปฯ.\csromanspacer{} สมฺมาทิฏฺฐิ ฯเปฯ สมฺมาสมาธิ สงฺขโตติ, อามนฺตา.\csromanspacer{} มคฺคสจฺจํ สงฺขตนฺติ, น เหวํ วตฺตพฺเพ ฯเปฯ.}

\proseitem{453}{\csromansymbolmark{*}นิโรธสจฺจํ อสงฺขตํ, นิโรโธ อสงฺขโตติ, อามนฺตา.\csromanspacer{} ทุกฺขสจฺจํ อสงฺขตํ, ทุกฺขํ อสงฺขตนฺติ, น เหวํ วตฺตพฺเพ ฯเปฯ.\csromanspacer{} นิโรธสจฺจํ อสงฺขตํ, นิโรโธ อสงฺขโตติ, อามนฺตา.\csromanspacer{} สมุทยสจฺจํ อสงฺขตํ, สมุทโย อสงฺขโตติ, น เหวํ วตฺตพฺเพ ฯเปฯ.\csromanspacer{} นิโรธสจฺจํ อสงฺขตํ, นิโรโธ อสงฺขโตติ, อามนฺตา.\csromanspacer{} มคฺคสจฺจํ อสงฺขตํ, มคฺโค อสงฺขโตติ, น เหวํ วตฺตพฺเพ ฯเปฯ.}

\prose{\csromansymbolmark{*}ทุกฺขสจฺจํ อสงฺขตํ, ทุกฺขํ สงฺขตนฺติ, อามนฺตา.\csromanspacer{} นิโรธสจฺจํ อสงฺขตํ, นิโรโธ สงฺขโตติ, น เหวํ วตฺตพฺเพ ฯเปฯ.\csromanspacer{} สมุทยสจฺจํ อสงฺขตํ, สมุทโย สงฺขโตติ, อามนฺตา.\csromanspacer{} นิโรธสจฺจํ อสงฺขตํ, นิโรโธ สงฺขโตติ, น เหวํ วตฺตพฺเพ ฯเปฯ.\csromanspacer{} มคฺคสจฺจํ อสงฺขตํ, มคฺโค สงฺขโตติ, อามนฺตา.\csromanspacer{} นิโรธสจฺจํ อสงฺขตํ.\csromanspacer{} นิโรโธ สงฺขโตติ, น เหวํ วตฺตพฺเพ ฯเปฯ.}

\proseitem{454}{\csromansymbolmark{+}น วตฺตพฺพํ “จตฺตาริ สจฺจานิ อสงฺขตานี”ติ, อามนฺตา.\csromanspacer{} นนุ วุตฺตํ ภควตา “จตฺตาริมานิ ภิกฺขเว ตถานิ อวิตถานิ อนญฺญถานิ.\csromanspacer{} กตมานิ จตฺตาริ, ‘อิทํ ทุกฺขนฺ’ติ ภิกฺขเว ตถเมตํ อวิตถเมตํ อนญฺญถเมตํ ฯเปฯ ‘อยํ ทุกฺขสมุทโย’ติ ฯเปฯ ‘อยํ ทุกฺขนิโรโธ’ติ ฯเปฯ ‘อยํ ทุกฺขนิโรธคามินี ปฏิปทา’ติ ตถเมตํ อวิตถเมตํ อนญฺญถเมตํ, อิมานิ โข ภิกฺขเว จตฺตาริ ตถานิ อวิตถานิ อนญฺญถานี”ติ\footnote{สํ 3. 377 ปิฏฺเฐ.}, อตฺเถว สุตฺตนฺโตติ, อามนฺตา.\csromanspacer{} เตน หิ จตฺตาริ สจฺจานิ อสงฺขตานีติ.}

\nitthitam{สจฺจกถา นิฏฺฐิตา.}
\chapterheadpageread{6. ฉฏฺฐวคฺค}
\vspace{\tipitakaparskip}
\csromancenter{\textbf{(56) 4.}\csromanspacer{}\textbf{ อารุปฺปกถา}}
\csromanmatikaanchor{mtk.80}
\csromantocmarkref{subsection}{(56) 4. อารุปฺปกถา}{mtk.80}
\bookmark[dest={mtk.80},level=2]{(56) 4. อารุปฺปกถา}
\proseitem{455}{\csromanfolio{244}\csromansymbolmark{*}อากาสานญฺจายตนํ อสงฺขตนฺติ, อามนฺตา.\csromanspacer{} นิพฺพานํ ตาณํ เลณํ สรณํ ปรายนํ อจฺจุตํ อมตนฺติ, น เหวํ วตฺตพฺเพ ฯเปฯ.\csromanspacer{} อากาสานญฺจายตนํ อสงฺขตํ, นิพฺพานํ อสงฺขตนฺติ, อามนฺตา.\csromanspacer{} เทฺว อสงฺขตานีติ, น เหวํ วตฺตพฺเพ ฯเปฯ.\csromanspacer{} เทฺว อสงฺขตานีติ, อามนฺตา.\csromanspacer{} เทฺว ตาณานิ ฯเปฯ อนฺตริกา วาติ, น เหวํ วตฺตพฺเพ ฯเปฯ.}

\prose{\csromansymbolmark{*}อากาสานญฺจายตนํ อสงฺขตนฺติ, อามนฺตา.\csromanspacer{} อากาสานญฺจายตนํ ภโว คติ สตฺตาวาโส สํสาโร โยนิ วิญฺญาณฏฺฐิติ อตฺตภาวปฏิลาโภติ, อามนฺตา.\csromanspacer{} อสงฺขตํ ภโว คติ สตฺตาวาโส สํสาโร โยนิ วิญฺญาณฏฺฐิติ อตฺตภาวปฏิลาโภติ, น เหวํ วตฺตพฺเพ ฯเปฯ.}

\prose{\csromansymbolmark{*}อตฺถิ อากาสานญฺจายตนูปคํ กมฺมนฺติ, อามนฺตา.\csromanspacer{} อตฺถิ อสงฺขตูปคํ กมฺมนฺติ, น เหวํ วตฺตพฺเพ ฯเปฯ.\csromanspacer{} อตฺถิ อากาสานญฺจายตนูปคา สตฺตาติ, อามนฺตา.\csromanspacer{} อตฺถิ อสงฺขตูปคา สตฺตาติ, น เหวํ วตฺตพฺเพ ฯเปฯ.}

\prose{\csromansymbolmark{*}อากาสานญฺจายตเน สตฺตา ชายนฺติ ชียนฺติ มียนฺติ จวนฺติ อุปปชฺชนฺตีติ, อามนฺตา.\csromanspacer{} อสงฺขเต สตฺตา ชายนฺติ ชียนฺติ มียนฺติ จวนฺติ อุปปชฺชนฺตีติ, น เหวํ วตฺตพฺเพ ฯเปฯ.\csromanspacer{} อากาสานญฺจายตเน อตฺถิ เวทนา.\csromanspacer{} สญฺญา, สงฺขารา, วิญฺญาณนฺติ, อามนฺตา.\csromanspacer{} อสงฺขเต อตฺถิ เวทนา, สญฺญา, สงฺขารา, วิญฺญาณนฺติ, น เหวํ วตฺตพฺเพ ฯเปฯ.\csromanspacer{} อากาสานญฺจายตนํ จตุโวการภโวติ, อามนฺตา.\csromanspacer{} อสงฺขตํ จตุโวการภโวติ, น เหวํ วตฺตพฺเพ ฯเปฯ.}

\proseitem{456}{\csromansymbolmark{+}น วตฺตพฺพํ “จตฺตาโร อารุปฺปา อสงฺขตา”ติ, อามนฺตา.\csromanspacer{} นนุ จตฺตาโร อารุปฺปา อเนชา วุตฺตา ภควตาติ, อามนฺตา.\csromanspacer{} หญฺจิ จตฺตาโร อารุปฺปา อเนชา วุตฺตา ภควตา, เตน วต เร วตฺตพฺเพ “จตฺตาโร อารุปฺปา อสงฺขตา”ติ.}

\nitthitam{อารุปฺปกถา นิฏฺฐิตา.}
\chapterheadpageread{6. ฉฏฺฐวคฺค \textbf{6. ฉฏฺฐวคฺค}}
\vspace{\tipitakaparskip}
\csromancenter{\textbf{(57) 5.}\csromanspacer{}\textbf{ นิโรธสมาปตฺติกถา}}
\csromanmatikaanchor{mtk.81}
\csromantocmarkref{subsection}{(57) 5. นิโรธสมาปตฺติกถา}{mtk.81}
\bookmark[dest={mtk.81},level=2]{(57) 5. นิโรธสมาปตฺติกถา}
\proseitem{457}{\csromanfolio{245}\csromansymbolmark{*}นิโรธสมาปตฺติ อสงฺขตาติ, อามนฺตา.\csromanspacer{} นิพฺพานํ ตาณํ เลณํ สรณํ ปรายนํ อจฺจุตํ อมตนฺติ, น เหวํ วตฺตพฺเพ ฯเปฯ.\csromanspacer{} นิโรธสมาปตฺติ อสงฺขตา, นิพฺพานํ อสงฺขตนฺติ, อามนฺตา.\csromanspacer{} เทฺว อสงฺขตานีติ, น เหวํ วตฺตพฺเพ ฯเปฯ.\csromanspacer{} เทฺว อสงฺขตานีติ, อามนฺตา.\csromanspacer{} เทฺว ตาณานิ ฯเปฯ อนฺตริกา วาติ.\csromanspacer{} น เหวํ วตฺตพฺเพ ฯเปฯ.}

\prose{\csromansymbolmark{*}นิโรธสมาปตฺติ อสงฺขตาติ, อามนฺตา.\csromanspacer{} อตฺถิ เกจิ นิโรธํ สมาปชฺชนฺติ ปฏิลภนฺติ อุปฺปาเทนฺติ สมุปฺปาเทนฺติ อุฏฺฐเปนฺติ สมุฏฺฐเปนฺติ นิพฺพตฺเตนฺติ อภินิพฺพตฺเตนฺติ ชเนนฺติ สญฺชเนนฺตีติ, อามนฺตา.\csromanspacer{} อตฺถิ เกจิ อสงฺขตํ สมาปชฺชนฺติ ปฏิลภนฺติ อุปฺปาเทนฺติ สมุปฺปาเทนฺติ อุฏฺฐเปนฺติ สมุฏฺฐเปนฺติ นิพฺพตฺเตนฺติ อภินิพฺพตฺเตนฺติ ชเนนฺติ สญฺชเนนฺตีติ, น เหวํ วตฺตพฺเพ ฯเปฯ.}

\proseitem{458}{\csromansymbolmark{*}นิโรธา โวทานํ วุฏฺฐานํ ปญฺญายตีติ, อามนฺตา.\csromanspacer{} อสงฺขตา โวทานํ วุฏฺฐานํ ปญฺญายตีติ, น เหวํ วตฺตพฺเพ ฯเปฯ.\csromanspacer{} นิโรธํ สมาปชฺชนฺตสฺส ปฐมํ นิรุชฺฌติ วจีสงฺขาโร, ตโต กายสงฺขาโร, ตโต จิตฺตสงฺขาโรติ, อามนฺตา.\csromanspacer{} อสงฺขตํ สมาปชฺชนฺตสฺส ปฐมํ นิรุชฺฌติ วจีสงฺขาโร, ตโต กายสงฺขาโร, ตโต จิตฺตสงฺขาโรติ, น เหวํ วตฺตพฺเพ ฯเปฯ.\csromanspacer{} นิโรธา วุฏฺฐหนฺตสฺส ปฐมํ อุปฺปชฺชติ จิตฺตสงฺขาโร, ตโต กายสงฺขาโร, ตโต วจีสงฺขาโรติ, อามนฺตา.\csromanspacer{} อสงฺขตา วุฏฺฐหนฺตสฺส ปฐมํ อุปฺปชฺชติ จิตฺตสงฺขาโร, ตโต กายสงฺขาโร, ตโต วจีสงฺขาโรติ, น เหวํ วตฺตพฺเพ ฯเปฯ.}

\prose{\csromansymbolmark{*}นิโรธา วุฏฺฐิตํ ตโย ผสฺสา ผุสนฺติ, สุญฺญโต ผสฺโส, อนิมิตฺโต ผสฺโส, อปฺปณิหิโต ผสฺโสติ, อามนฺตา.\csromanspacer{} อสงฺขตา วุฏฺฐิตํ ตโย ผสฺสา ผุสนฺติ, สุญฺญโต ผสฺโส, อนิมิตฺโต ผสฺโส, อปฺปณิหิโต ผสฺโสติ, น เหวํ วตฺตพฺเพ ฯเปฯ.}

\prose{\csromansymbolmark{*}นิโรธา วุฏฺฐิตสฺส วิเวกนินฺนํ จิตฺตํ โหติ วิเวกโปณํ วิเวกปพฺภารนฺติ, อามนฺตา.\csromanspacer{} อสงฺขตา วุฏฺฐิตสฺส วิเวกนินฺนํ จิตฺตํ โหติ วิเวกโปณํ วิเวกปพฺภารนฺติ, น เหวํ วตฺตพฺเพ ฯเปฯ.}

\csromanmatikaanchor{mtk.82}
\csromantocmarkref{subsection}{(58) 6. อากาสกถา}{mtk.82}
\bookmark[dest={mtk.82},level=2]{(58) 6. อากาสกถา}
\proseitem{459}{\csromanfolio{246}\csromansymbolmark{+}น วตฺตพฺพํ “นิโรธสมาปตฺติ อสงฺขตา”ติ, อามนฺตา.\csromanspacer{} สงฺขตาติ, น เหวํ วตฺตพฺเพ ฯเปฯ.\csromanspacer{} เตน หิ นิโรธสมาปตฺติ อสงฺขตาติ.}

\nitthitam{นิโรธสมาปตฺติกถา นิฏฺฐิตา.}
\chapterhead{6. ฉฏฺฐวคฺค}
\vspace{\tipitakaparskip}
\csromancenter{\textbf{(58) 6.}\csromanspacer{}\textbf{ อากาสกถา}}
\proseitem{460}{\csromansymbolmark{*}อากาโส อสงฺขโตติ, อามนฺตา.\csromanspacer{} นิพฺพานํ ตาณํ เลณํ สรณํ ปรายนํ อจฺจุตํ อมตนฺติ, น เหวํ วตฺตพฺเพ ฯเปฯ.\csromanspacer{} อากาโส อสงฺขโต, นิพฺพานํ อสงฺขตนฺติ, อามนฺตา.\csromanspacer{} เทฺว อสงฺขตานีติ, น เหวํ วตฺตพฺเพ ฯเปฯ.\csromanspacer{} เทฺว อสงฺขตานีติ, อามนฺตา.\csromanspacer{} เทฺว ตาณานิ ฯเปฯ อนฺตริกา วาติ, น เหวํ วตฺตพฺเพ ฯเปฯ.}

\prose{\csromansymbolmark{*}อากาโส อสงฺขโตติ, อามนฺตา.\csromanspacer{} อตฺถิ เกจิ อนากาสํ อากาสํ กโรนฺตีติ, อามนฺตา.\csromanspacer{} อตฺถิ เกจิ สงฺขตํ อสงฺขตํ กโรนฺตีติ, น เหวํ วตฺตพฺเพ ฯเปฯ.\csromanspacer{} อตฺถิ เกจิ อากาสํ อนากาสํ กโรนฺตีติ, อามนฺตา.\csromanspacer{} อตฺถิ เกจิ อสงฺขตํ สงฺขตํ กโรนฺตีติ, น เหวํ วตฺตพฺเพ ฯเปฯ.}

\prose{\csromansymbolmark{*}อากาเส ปกฺขิโน คจฺฉนฺติ, จนฺทิมสูริยา คจฺฉนฺติ, ตารกรูปานิ คจฺฉนฺติ, อิทฺธิํ วิกุพฺพนฺติ, พาหุํ จาเลนฺติ, ปาณิํ จาเลนฺติ, เลฑฺฑุํ ขิปนฺติ, ลคุฬํ ขิปนฺติ, อิทฺธิํ\footnote{สตฺติํ (?)} ขิปนฺติ, อุสุํ ขิปนฺตีติ, อามนฺตา.\csromanspacer{} อสงฺขเต ปกฺขิโน คจฺฉนฺติ, จนฺทิมสูริยา คจฺฉนฺติ, ตารกรูปานิ คจฺฉนฺติ, อิทฺธิํ วิกุพฺพนฺติ, พาหุํ จาเลนฺติ, ปาณิํ จาเลนฺติ, เลฑฺฑุํ ขิปนฺติ, ลคุฬํ ขิปนฺติ, อิทฺธิํ ขิปนฺติ, อุสุํ ขิปนฺตีติ, น เหวํ วตฺตพฺเพ ฯเปฯ.}

\proseitem{461}{\csromansymbolmark{*}อากาสํ ปริวาเรตฺวา ฆรานิ กโรนฺติ โกฏฺฐานิ กโรนฺตีติ, อามนฺตา.\csromanspacer{} อสงฺขตํ ปริวาเรตฺวา ฆรานิ กโรนฺติ โกฏฺฐานิ กโรนฺตีติ, น เหวํ วตฺตพฺเพ ฯเปฯ.}

\chapterheadpageread{6. ฉฏฺฐวคฺค}
\csromanmatikaanchor{mtk.83}
\csromantocmarkref{subsection}{(59) 7. อากาโสสนิทสฺสโนติกถา}{mtk.83}
\bookmark[dest={mtk.83},level=2]{(59) 7. อากาโสสนิทสฺสโนติกถา}
\prose{\csromanfolio{247}\csromansymbolmark{*}อุทปาเน ขญฺญมาเน อนากาโส อากาโส โหตีติ, อามนฺตา.\csromanspacer{} สงฺขตํ อสงฺขตํ โหตีติ, น เหวํ วตฺตพฺเพ ฯเปฯ.}

\prose{\csromansymbolmark{*}ตุจฺฉ-อุทปาเน ปูริยมาเน ตุจฺฉโกฏฺเฐ ปูริยมาเน ตุจฺฉกุมฺภิยา ปูริยมานาย อากาโส อนฺตรธายตีติ, อามนฺตา.\csromanspacer{} อสงฺขตํ อนฺตรธายตีติ, น เหวํ วตฺตพฺเพ ฯเปฯ.}

\proseitem{462}{\csromansymbolmark{+}น วตฺตพฺพํ “อากาโส อสงฺขโต”ติ, อามนฺตา.\csromanspacer{} อากาโส สงฺขโตติ, น เหวํ วตฺตพฺเพ ฯเปฯ.\csromanspacer{} เตน หิ อากาโส อสงฺขโตติ.}

\nitthitam{อากาสกถา นิฏฺฐิตา.}
\chapterhead{6. ฉฏฺฐวคฺค}
\vspace{\tipitakaparskip}
\csromancenter{\textbf{(59) 7.}\csromanspacer{}\textbf{ อากาโสสนิทสฺสโนติกถา}}
\proseitem{463}{\csromansymbolmark{*}อากาโส สนิทสฺสโนติ, อามนฺตา.\csromanspacer{} รูปํ รูปายตนํ รูปธาตุ นีลํ ปีตกํ โลหิตกํ โอทาตํ จกฺขุวิญฺเญยฺยํ จกฺขุสฺมิํ ปฏิหญฺญติ จกฺขุสฺส อาปาถํ อาคจฺฉตีติ, น เหวํ วตฺตพฺเพ ฯเปฯ.}

\prose{\csromansymbolmark{*}อากาโส สนิทสฺสโนติ, อามนฺตา.\csromanspacer{} จกฺขุญฺจ ปฏิจฺจ อากาสญฺจ อุปฺปชฺชติ จกฺขุวิญฺญาณนฺติ, น เหวํ วตฺตพฺเพ ฯเปฯ.}

\prose{\csromansymbolmark{*}จกฺขุญฺจ ปฏิจฺจ อากาสญฺจ อุปฺปชฺชติ จกฺขุวิญฺญาณนฺติ, อามนฺตา. “จกฺขุญฺจ ปฏิจฺจ อากาสญฺจ อุปฺปชฺชติ จกฺขุวิญฺญาณนฺ”ติ อตฺเถว\footnote{อตฺถิ(?)} สุตฺตนฺโตติ, นตฺถิ. “จกฺขุญฺจ ปฏิจฺจ รูเป จ อุปฺปชฺชติ จกฺขุวิญฺญาณนฺ”ติ\footnote{ม 1. 326; ม 3. 328; สํ 2. 261 ปิฏฺฐาทีสุ.} อตฺเถว สุตฺตนฺโตติ, อามนฺตา.\csromanspacer{} หญฺจิ “จกฺขุญฺจ ปฏิจฺจ รูเป จ อุปฺปชฺชติ จกฺขุวิญฺญาณนฺ”ติ อตฺเถว สุตฺตนฺโต, โน จ วต เร วตฺตพฺเพ “จกฺขุญฺจ ปฏิจฺจ อากาสญฺจ อุปฺปชฺชติ จกฺขุวิญฺญาณนฺ”ติ.}

\csromanmatikaanchor{mtk.84}
\csromantocmarkref{subsection}{(60) 8. ปถวีธาตุสนิทสฺสนาติ-อาทิกถา}{mtk.84}
\bookmark[dest={mtk.84},level=2]{(60) 8. ปถวีธาตุสนิทสฺสนาติ-อาทิกถา}
\proseitem{464}{\csromansymbolmark{+}น วตฺตพฺพํ “อากาโส สนิทสฺสโน”ติ, อามนฺตา.\csromanspacer{} นนุ ปสฺสติ ทฺวินฺนํ รุกฺขานํ อนฺตรํ ทฺวินฺนํ ถมฺภานํ อนฺตรํ ตาฬจฺฉิทฺทํ วาตปานจฺฉิทฺทนฺติ, \csromanfolio{248}อามนฺตา.\csromanspacer{} หญฺจิ ปสฺสติ ทฺวินฺนํ รูกฺขานํ อนฺตรํ ทฺวินฺนํ ถมฺภานํ อนฺตรํ ตาฬจฺฉิทฺทํ วาตปานจฺฉิทฺทํ, เตน วต เร วตฺตพฺเพ “อากาโส สนิทสฺสโน”ติ.}

\nitthitam{อากาโส สนิทสฺสโนติกถา นิฏฺฐิตา.}
\chapterhead{6. ฉฏฺฐวคฺค}
\vspace{\tipitakaparskip}
\csromancenter{\textbf{(60) 8.}\csromanspacer{}\textbf{ ปถวีธาตุสนิทสฺสนาติ-อาทิกถา}}
\proseitem{465}{\csromansymbolmark{*}ปถวีธาตุ สนิทสฺสนาติ, อามนฺตา.\csromanspacer{} รูปํ รูปายตนํ รูปธาตุ นีลํ ปีตกํ โลหิตกํ โอทาตํ จกฺขุวิญฺเญยฺยํ จกฺขุสฺมิํ ปฏิหญฺญติ จกฺขุสฺส อาปาถํ อาคจฺฉตีติ, น เหวํ วตฺตพฺเพ ฯเปฯ.}

\prose{\csromansymbolmark{*}ปถวีธาตุ สนิทสฺสนาติ, อามนฺตา.\csromanspacer{} จกฺขุญฺจ ปฏิจฺจ ปถวีธาตุญฺจ อุปฺปชฺชติ จกฺขุวิญฺญาณนฺติ, น เหวํ วตฺตพฺเพ ฯเปฯ.}

\prose{\csromansymbolmark{*}จกฺขุญฺจ ปฏิจฺจ ปถวีธาตุญฺจ อุปฺปชฺชติ จกฺขุวิญฺญาณนฺติ, อามนฺตา. “จกฺขุญฺจ ปฏิจฺจ ปถวีธาตุญฺจ อุปฺปชฺชติ จกฺขุวิญฺญาณนฺ”ติ อตฺเถว สุตฺตนฺโตติ, นตฺถิ. “จกฺขุญฺจ ปฏิจฺจ รูเป จ อุปฺปชฺชติ จกฺขุวิญฺญาณนฺ”ติ อตฺเถว สุตฺตนฺโตติ, อามนฺตา.\csromanspacer{} หญฺจิ “จกฺขุญฺจ ปฏิจฺจ รูเป จ อุปฺปชฺชติ จกฺขุวิญฺญาณนฺ”ติ อตฺเถว สุตฺตนฺโต, โน จ วต เร วตฺตพฺเพ “จกฺขุญฺจ ปฏิจฺจ ปถวีธาตุญฺจ อุปฺปชฺชติ จกฺขุวิญฺญาณนฺ”ติ.}

\proseitem{466}{\csromansymbolmark{+}น วตฺตพฺพํ “ปถวีธาตุ สนิทสฺสนา”ติ, อามนฺตา.\csromanspacer{} นนุ ปสฺสติ\footnote{ปสฺสสิ (สี, สฺยา, ก) เอวมุปริปิ.} ภูมิํ ปาสาณํ ปพฺพตนฺติ, อามนฺตา.\csromanspacer{} หญฺจิ ปสฺสติ1 ภูมิํ ปาสาณํ ปพฺพตํ, เตน วต เร วตฺตพฺเพ “ปถวีธาตุ สนิทสฺสนา”ติ ฯเปฯ.}

\prose{\csromansymbolmark{+}น วตฺตพฺพํ “อาโปธาตุ สนิทสฺสนา”ติ, อามนฺตา.\csromanspacer{} นนุ ปสฺสติ อุทกนฺติ, อามนฺตา.\csromanspacer{} หญฺจิ ปสฺสติ อุทกํ, เตน วต เร วตฺตพฺเพ “อาโปธาตุ สนิทสฺสนา”ติ ฯเปฯ.}

\prose{\csromansymbolmark{+}น วตฺตพฺพํ “เตโชธาตุ สนิทสฺสนา”ติ, อามนฺตา.\csromanspacer{} นนุ ปสฺสติ อคฺคิํ ชลนฺตนฺติ, อามนฺตา.\csromanspacer{} หญฺจิ ปสฺสติ อคฺคิํ ชลนฺตํ, เตน วต เร วตฺตพฺเพ “เตโชธาตุ สนิทสฺสนา”ติ ฯเปฯ.}

\chapterheadpageread{6. ฉฏฺฐวคฺค}
\csromanmatikaanchor{mtk.85}
\csromanmatikaanchor{mtk.86}
\csromantocmarkref{subsection}{(61) 9. จกฺขุนฺทฺริยํสนิทสฺสนนฺติ-อาทิกถา}{mtk.85}
\bookmark[dest={mtk.85},level=2]{(61) 9. จกฺขุนฺทฺริยํสนิทสฺสนนฺติ-อาทิกถา}
\csromantocmarkref{subsection}{(62) 10. กายกมฺมํสนิทสฺสนนฺติกถา}{mtk.86}
\bookmark[dest={mtk.86},level=2]{(62) 10. กายกมฺมํสนิทสฺสนนฺติกถา}
\prose{\csromanfolio{249}\csromansymbolmark{+}น วตฺตพฺพํ “วาโยธาตุ สนิทสฺสนา”ติ, อามนฺตา.\csromanspacer{} นนุ ปสฺสติ วาเตน รุกฺเข สญฺจาลิยมาเนติ, อามนฺตา.\csromanspacer{} หญฺจิ ปสฺสติ วาเตน รุกฺเข สญฺจาลิยมาเน, เตน วต เร วตฺตพฺเพ “วาโยธาตุ สนิทสฺสนา”ติ ฯเปฯ.}

\nitthitam{ปถวีธาตุ สนิทสฺสนาติ-อาทิกถา นิฏฺฐิตา.}
\chapterhead{6. ฉฏฺฐวคฺค}
\vspace{\tipitakaparskip}
\csromancenter{\textbf{(61) 9.}\csromanspacer{}\textbf{ จกฺขุนฺทฺริยํสนิทสฺสนนฺติ-อาทิกถา}}
\proseitem{467}{\csromansymbolmark{*}จกฺขุนฺทฺริยํ สนิทสฺสนนฺติ, อามนฺตา.\csromanspacer{} รูปํ รูปายตนํ รูปธาตุ ฯเปฯ.\csromanspacer{} จกฺขุสฺส อาปาถํ อาคจฺฉตีติ, น เหวํ วตฺตพฺเพ ฯเปฯ.}

\prose{\csromansymbolmark{*}จกฺขุนฺทฺริยํ สนิทสฺสนนฺติ, อามนฺตา.\csromanspacer{} จกฺขุญฺจ ปฏิจฺจ จกฺขุนฺทฺริยญฺจ อุปฺปชฺชติ จกฺขุวิญฺญาณนฺติ, น เหวํ วตฺตพฺเพ ฯเปฯ.}

\prose{\csromansymbolmark{*}จกฺขุญฺจ ปฏิจฺจ จกฺขุนฺทฺริยญฺจ อุปฺปชฺชติ จกฺขุวิญฺญาณนฺติ, อามนฺตา. “จกฺขุญฺจ ปฏิจฺจ จกฺขุนฺทฺริยญฺจ อุปฺปชฺชติ จกฺขุวิญฺญาณนฺ”ติ อตฺเถว สุตฺตนฺโตติ, นตฺถิ. “จกฺขุญฺจ ปฏิจฺจ รูเป จ อุปฺปชฺชติ จกฺขุวิญฺญาณนฺ”ติ อตฺเถว สุตฺตนฺโตติ, อามนฺตา.\csromanspacer{} หญฺจิ “จกฺขุญฺจ ปฏิจฺจ รูเป จ อุปฺปชฺชติ จกฺขุวิญฺญาณนฺ”ติ อตฺเถว สุตฺตนฺโต, โน จ วต เร วตฺตพฺเพ “จกฺขุญฺจ ปฏิจฺจ จกฺขุนฺทฺริยญฺจ อุปฺปชฺชติ จกฺขุวิญฺญาณนฺ”ติ.}

\proseitem{468}{\csromansymbolmark{+}น วตฺตพฺพํ “ปญฺจินฺทฺริยานิ สนิทสฺสนานี”ติ, อามนฺตา.\csromanspacer{} นนุ ปสฺสติ จกฺขุํ, โสตํ, ฆานํ, ชิวฺหํ, กายนฺติ, อามนฺตา.\csromanspacer{} หญฺจิ ปสฺสติ จกฺขุํ, โสตํ, ฆานํ, ชิวฺหํ, กายํ, เตน วต เร วตฺตพฺเพ “ปญฺจินฺทฺริยานิ สนิทสฺสนานี”ติ ฯเปฯ.}

\nitthitam{จกฺขุนฺทฺริยํ สนิทสฺสนนฺติ-อาทิกถา นิฏฺฐิตา.}
\chapterhead{6. ฉฏฺฐวคฺค}
\vspace{\tipitakaparskip}
\csromancenter{\textbf{(62) 10.}\csromanspacer{}\textbf{ กายกมฺมํสนิทสฺสนนฺติกถา}}
\proseitem{469}{\csromansymbolmark{*}กายกมฺมํ สนิทสฺสนนฺติ, อามนฺตา.\csromanspacer{} รูปํ รูปายตนํ รูปธาตุ นีลํ ปีตกํ โลหิตกํ โอทาตํ จกฺขุวิญฺเญยฺยํ จกฺขุสฺมิํ ปฏิหญฺญติ จกฺขุสฺส อาปาถํ อาคจฺฉตีติ, น เหวํ วตฺตพฺเพ ฯเปฯ.}

\csromanmatikaanchor{mtk.87}
\csromanmatikaanchor{mtk.88}
\csromantocmarknopage{chapter}{7. สตฺตมวคฺค}{mtk.87}
\bookmark[dest={mtk.87},level=0]{7. สตฺตมวคฺค}
\csromantocmarkref{subsection}{(63) 1. สงฺคหิตกถา}{mtk.88}
\bookmark[dest={mtk.88},level=2]{(63) 1. สงฺคหิตกถา}
\prose{\csromanfolio{250}\csromansymbolmark{*}กายกมฺมํ สนิทสฺสนนฺติ, อามนฺตา.\csromanspacer{} จกฺขุญฺจ ปฏิจฺจ กายกมฺมญฺจ อุปฺปชฺชติ จกฺขุวิญฺญาณนฺติ, น เหวํ วตฺตพฺเพ ฯเปฯ.}

\prose{\csromansymbolmark{*}จกฺขุญฺจ ปฏิจฺจ กายกมฺมญฺจ อุปฺปชฺชติ จกฺขุวิญฺญณนฺติ, อามนฺตา. “จกฺขุญฺจ ปฏิจฺจ กายกมฺมญฺจ อุปฺปชฺชติ จกฺขุวิญฺญาณนฺ”ติ อตฺเถว สุตฺตนฺโตติ, นตฺถิ. “จกฺขุญฺจ ปฏิจฺจ รูเป จ อุปฺปชฺชติ จกฺขุวิญฺญาณนฺ”ติ อตฺเถว สุตฺตนฺโตติ, อามนฺตา.\csromanspacer{} หญฺจิ “จกฺขุญฺจ ปฏิจฺจ รูเป จ อุปฺปชฺชติ จกฺขุวิญฺญาณนฺ”ติ อตฺเถว สุตฺตนฺโต, โน จ วต เร วตฺตพฺเพ “จกฺขุญฺจ ปฏิจฺจ กายกมฺมญฺจ อุปฺปชฺชติ จกฺขุวิญฺญาณนฺ”ติ.}

\proseitem{470}{\csromansymbolmark{+}น วตฺตพฺพํ “กายกมฺมํ สนิทสฺสนนฺ”ติ, อามนฺตา.\csromanspacer{} นนุ ปสฺสติ อภิกฺกมนฺตํ ปฏิกฺกมนฺตํ อาโลเกนฺตํ วิโลเกนฺตํ สมิญฺเชนฺตํ ปสาเรนฺตนฺติ, อามนฺตา.\csromanspacer{} หญฺจิ ปสฺสติ อภิกฺกมนฺตํ ปฏิกฺกมนฺตํ อาโลเกนฺตํ วิโลเกนฺตํ สมิญฺเชนฺตํ ปสาเรนฺตํ, เตน วต เร วตฺตพฺเพ “กายกมฺมํ สนิทสฺสนนฺ”ติ.\csromanspacer{} กายกมฺมํ สนิทสฺสนนฺติกถา นิฏฺฐิตา.\csromanspacer{} ฉฏฺโฐ วคฺโค.}

\titlehead{\textbf{ตสฺสุทฺทานํ}}
\prose{นิยโม อสงฺขโต, ปฏิจฺจสมุปฺปาโท อสงฺขโต, จตฺตาริ สจฺจานิ อสงฺขตานิ, จตฺตาโร อารุปฺปา อสงฺขตา, นิโรธสมาปตฺติ อสงฺขตา, อากาโส อสงฺขโต, อากาโส สนิทสฺสโน, จตฺตาโร มหาภูตา, ปญฺจินฺทฺริยานิ, ตเถว กายกมฺมนฺติ.}

\chapterhead{7. สตฺตมวคฺค}
\vspace{\tipitakaparskip}
\csromancenter{\textbf{(63) 1.}\csromanspacer{}\textbf{ สงฺคหิตกถา}}
\proseitem{471}{\csromansymbolmark{*}นตฺถิ เกจิ ธมฺมา เกหิจิ ธมฺเมหิ สงฺคหิตาติ\footnote{สงฺคหีตาติ (อิ)}, อามนฺตา.\csromanspacer{} นนุ อตฺถิ เกจิ ธมฺมา เกหิจิ ธมฺเมหิ คณนํ คจฺฉนฺติ อุทฺเทสํ คจฺฉนฺติ ปริยาปนฺนาติ, อามนฺตา.\csromanspacer{} หญฺจิ อตฺถิ เกจิ ธมฺมา เกหิจิ ธมฺเมหิ คณนํ}

\vspace{\tipitakaparskip}
\csromancenter{สตฺตมวคฺค คจฺฉนฺติ อุทฺเทสํ คจฺฉนฺติ ปริยาปนฺนา, โน จ วต เร วตฺตพฺเพ “นตฺถิ เกจิ ธมฺมา เกหิจิ ธมฺเมหิ สงฺคหิตา”ติ.}
\prose{\csromanfolio{251}\csromansymbolmark{*}จกฺขายตนํ กตมกฺขนฺธคณนํ\footnote{กตมํ ขนฺธคณนํ (สี, อิ, ก)} คจฺฉตีติ, รูปกฺขนฺธคณนํ คจฺฉตีติ.\csromanspacer{} หญฺจิ จกฺขายตนํ รูปกฺขนฺธคณนํ คจฺฉติ, เตน วต เร วตฺตพฺเพ “จกฺขายตนํ รูปกฺขนฺเธน สงฺคหิตนฺ”ติ.\csromanspacer{} โสตายตนํ ฯเปฯ.\csromanspacer{} ฆานายตนํ ฯเปฯ.\csromanspacer{} ชิวฺหายตนํ ฯเปฯ.\csromanspacer{} กายายตนํ กตมกฺขนฺธคณนํ คจฺฉตีติ, รูปกฺขนฺธคณนํ คจฺฉตีติ, หญฺจิ กายายตนํ รูปกฺขนฺธคณนํ คจฺฉติ, เตน วต เร วตฺตพฺเพ “กายายตนํ รูปกฺขนฺเธน สงฺคหิตนฺ”ติ.}

\prose{\csromansymbolmark{*}รูปายตนํ ฯเปฯ.\csromanspacer{} สทฺทายตนํ ฯเปฯ.\csromanspacer{} คนฺธายตนํ ฯเปฯ.\csromanspacer{} รสายตนํ ฯเปฯ.\csromanspacer{} โผฏฺฐพฺพายตนํ กตมกฺขนฺธคณนํ คจฺฉตีติ, รูปกฺขนฺธคณนํ คจฺฉตีติ.\csromanspacer{} หญฺจิ โผฏฺฐพฺพายตนํ รูปกฺขนฺธคณนํ คจฺฉติ, เตน วต เร วตฺตพฺเพ “โผฏฺฐพฺพายตนํ รูปกฺขนฺเธน สงฺคหิตนฺติ.}

\prose{\csromansymbolmark{*}สุขา เวทนา กตมกฺขนฺธคณนํ คจฺฉตีติ, เวทนากฺขนฺธคณนํ คจฺฉตีติ.\csromanspacer{} หญฺจิ สุขา เวทนา เวทนากฺขนฺธคณนํ คจฺฉติ, เตน วต เร วตฺตพฺเพ “สุขา เวทนา เวทนากฺขนฺเธน สงฺคหิตา”ติ, ทุกฺขา เวทนา ฯเปฯ อทุกฺขมสุขา เวทนา กตมกฺขนฺธคณนํ คจฺฉตีติ, เวทนากฺขนฺธคณนํ คจฺฉตีติ, หญฺจิ อทุกฺขมสุขา เวทนา เวทนากฺขนฺธคณนํ คจฺฉติ, เตน วต เร วตฺตพฺเพ “อทุกฺขมสุขา เวทนา เวทนากฺขนฺเธน สงฺคหิตา”ติ.}

\prose{\csromansymbolmark{*}จกฺขุสมฺผสฺสชา สญฺญา กตมกฺขนฺธคณนํ คจฺฉตีติ, สญฺญากฺขนฺธคณนํ คจฺฉตีติ, หญฺจิ จกฺขุสมฺผสฺสชา สญฺญา สญฺญากฺขนฺธคณนํ คจฺฉติ, เตน วต เร วตฺตพฺเพ “จกฺขุสมฺผสฺสชา สญฺญา สญฺญากฺขนฺเธน สงฺคหิตา”ติ.\csromanspacer{} โสตสมฺผสฺสชา สญฺญา ฯเปฯ.\csromanspacer{} มโนสมฺผสฺสชา สญฺญา กตมกฺขนฺธคณนํ คจฺฉตีติ, สญฺญากฺขนฺธคณนํ คจฺฉตีติ.\csromanspacer{} หญฺจิ มโนสมฺผสฺสชา สญฺญา สญฺญากฺขนฺธคณนํ คจฺฉติ, เตน วต เร วตฺตพฺเพ “มโนสมฺผสฺสชา สญฺญา สญฺญากฺขนฺเธน สงฺคหิตา”ติ.}

\prose{\csromansymbolmark{*}จกฺขุสมฺผสฺสชา เจตนา ฯเปฯ.\csromanspacer{} มโนสมฺผสฺสชา เจตนา กตมกฺขนฺธคณนํ คจฺฉตีติ, สงฺขารกฺขนฺธคณนํ คจฺฉตีติ.\csromanspacer{} หญฺจิ มโนสมฺผสฺสชา เจตนา สงฺขารกฺขนฺธคณนํ คจฺฉติ, เตน วต เร วตฺตพฺเพ “มโนสมฺผสฺสชา เจตนา สงฺขารกฺขนฺเธน สงฺคหิตาติ.}

\csromanmatikaanchor{mtk.89}
\csromantocmarkref{subsection}{(64) 2. สมฺปยุตฺตกถา}{mtk.89}
\bookmark[dest={mtk.89},level=2]{(64) 2. สมฺปยุตฺตกถา}
\prose{\csromanfolio{252}\csromansymbolmark{*}จกฺขุวิญฺญาณํ ฯเปฯ.\csromanspacer{} มโนวิญฺญาณํ กตมกฺขนฺธคณนํ คจฺฉตีติ, วิญฺญาณกฺขนฺธคณนํ คจฺฉตีติ.\csromanspacer{} หญฺจิ มโนวิญฺญาณํ วิญฺญาณกฺขนฺธคณนํ คจฺฉติ, เตน วต เร วตฺตพฺเพ “มโนวิญฺญาณํ วิญฺญาณกฺขนฺเธน สงฺคหิตนฺ”ติ.}

\proseitem{472}{\csromansymbolmark{+}ยถา ทาเมน วา โยตฺเตน วา เทฺว พลีพทฺทา สงฺคหิตา, สิกฺกาย ปิณฺฑปาโต สงฺคหิโต, สา คทฺทุเลน สงฺคหิโต.\csromanspacer{} เอวเมว เต ธมฺมา เตหิ ธมฺเมหิ สงฺคหิตาติ.1 หญฺจิ ทาเมน วา โยตฺเตน วา เทฺว พลีพทฺทา สงฺคหิตา, สิกฺกาย ปิณฺฑปาโต สงฺคหิโต, สา คทฺทุเลน สงฺคหิโต, เตน วต เร วตฺตพฺเพ “อตฺถิ เกจิ ธมฺมา เกหิจิ ธมฺเมหิ สงฺคหิตา”ติ1.}

\nitthitam{สงฺคหิตกถา นิฏฺฐิตา.}
\chapterhead{7. สตฺตมวคฺค}
\vspace{\tipitakaparskip}
\csromancenter{\textbf{(64) 2.}\csromanspacer{}\textbf{ สมฺปยุตฺตกถา}}
\proseitem{473}{\csromansymbolmark{*}นตฺถิ เกจิ ธมฺมา เกหิจิ ธมฺเมหิ สมฺปยุตฺตาติ, อามนฺตา.\csromanspacer{} นนุ อตฺถิ เกจิ ธมฺมา เกหิจิ ธมฺเมหิ สหคตา สหชาตา สํสฏฺฐา เอกุปฺปาทา เอกนิโรธา เอกวตฺถุกา เอการมฺมณาติ, อามนฺตา.\csromanspacer{} หญฺจิ อตฺถิ เกจิ ธมฺมา เกหิจิ ธมฺเมหิ สหคตา สหชาตา สํสฏฺฐา เอกุปฺปาทา เอกนิโรธา เอกวตฺถุกา เอการมฺมณา, โน จ วต เร วตฺตพฺเพ “นตฺถิ เกจิ ธมฺมา เกหิจิ ธมฺเมหิ สมฺปยุตฺตา”ติ.}

\prose{\csromansymbolmark{*}เวทนากฺขนฺโธ สญฺญากฺขนฺเธน สหชาโตติ, อามนฺตา.\csromanspacer{} หญฺจิ เวทนากฺขนฺโธ สญฺญากฺขนฺเธน สหชาโต, เตน วต เร วตฺตพฺเพ “เวทนากฺขนฺโธ สญฺญากฺขนฺเธน สมฺปยุตฺโต”ติ.}

\prose{\csromansymbolmark{*}เวทนากฺขนฺโธ สงฺขารกฺขนฺเธน.\csromanspacer{} วิญฺญาณกฺขนฺเธน สหชาโตติ, อามนฺตา.\csromanspacer{} หญฺจิ เวทนากฺขนฺโธ วิญฺญาณกฺขนฺเธน สหชาโต, เตน วต เร วตฺตพฺเพ “เวทนากฺขนฺโธ วิญฺญาณกฺขนฺเธน สมฺปยุตฺโต”ติ.}

\chapterheadpageread{7. สตฺตมวคฺค}
\csromanmatikaanchor{mtk.90}
\csromantocmarkref{subsection}{(65) 3. เจตสิกกถา}{mtk.90}
\bookmark[dest={mtk.90},level=2]{(65) 3. เจตสิกกถา}
\prose{\csromanfolio{253}\csromansymbolmark{*}สญฺญากฺขนฺโธ.\csromanspacer{} สงฺขารกฺขนฺโธ.\csromanspacer{} วิญฺญาณกฺขนฺโธ.\csromanspacer{} เวทนากฺขนฺเธน.\csromanspacer{} สญฺญากฺขนฺเธน.\csromanspacer{} สงฺขารกฺขนฺเธน สหชาโตติ, อามนฺตา.\csromanspacer{} หญฺจิ วิญฺญาณกฺขนฺโธ สงฺขารกฺขนฺเธน สหชาโต, เตน วต เร วตฺตพฺเพ “วิญฺญาณกฺขนฺโธ สงฺขารกฺขนฺเธน สมฺปยุตฺโต”ติ.}

\proseitem{474}{\csromansymbolmark{+}ยถา ติลมฺหิ เตลํ อนุคตํ อนุปวิฏฺฐํ, อุจฺฉุมฺหิ รโส อนุคโต อนุปวิฏฺโฐ.\csromanspacer{} เอวเมว เต ธมฺมา เตหิ ธมฺเมหิ อนุคตา อนุปวิฏฺฐาติ, น เหวํ วตฺตพฺเพ ฯเปฯ.}

\nitthitam{สมฺปยุตฺตกถา นิฏฺฐิตา.}
\chapterhead{7. สตฺตมวคฺค}
\vspace{\tipitakaparskip}
\csromancenter{\textbf{(65) 3.}\csromanspacer{}\textbf{ เจตสิกกถา}}
\proseitem{475}{\csromansymbolmark{*}นตฺถิ เจตสิโก ธมฺโมติ, อามนฺตา.\csromanspacer{} นนุ อตฺธิ เกจิ ธมฺมา จิตฺเตน สหคตา สหชาตา สํสฏฺฐา สมฺปยุตฺตา เอกุปฺปาทา เอกนิโรธา เอกวตฺถุกา เอการมฺมณาติ, อามนฺตา.\csromanspacer{} หญฺจิ อตฺถิ เกจิ ธมฺมา จิตฺเตน สหคตา สหชาตา สํสฏฺฐา สมฺปยุตฺตา เอกุปฺปาทา เอกนิโรธา เอกวตฺถุกา เอการมฺมณา, โน จ วต เร วตฺตพฺเพ “นตฺถิ เจตสิโก ธมฺโม”ติ.}

\prose{\csromansymbolmark{*}ผสฺโส จิตฺเตน สหชาโตติ, อามนฺตา.\csromanspacer{} หญฺจิ ผสฺโส จิตฺเตน สหชาโต, เตน วต เร วตฺตพฺเพ “ผสฺโส เจตสิโก”ติ.\csromanspacer{} เวทนา ฯเปฯ.\csromanspacer{} สญฺญา.\csromanspacer{} เจตนา.\csromanspacer{} สทฺธา.\csromanspacer{} วีริยํ.\csromanspacer{} สติ.\csromanspacer{} สมาธิ.\csromanspacer{} ปญฺญา.\csromanspacer{} ราโค.\csromanspacer{} โทโส.\csromanspacer{} โมโห ฯเปฯ.\csromanspacer{} อโนตฺตปฺปํ จิตฺเตน สหชาตนฺติ, อามนฺตา.\csromanspacer{} หญฺจิ อโนตฺตปฺปํ จิตฺเตน สหชาตํ, เตน วต เร วตฺตพฺเพ “อโนตฺตปฺปํ เจตสิกนฺ”ติ.}

\csromanmatikaanchor{mtk.91}
\csromantocmarkref{subsection}{(66) 4. ทานกถา}{mtk.91}
\bookmark[dest={mtk.91},level=2]{(66) 4. ทานกถา}
\proseitem{476}{\csromansymbolmark{+}จิตฺเตน สหชาตาติ กตฺวา เจตสิกาติ, อามนฺตา.\csromanspacer{} ผสฺเสน สหชาตาติ กตฺวา ผสฺสสิกาติ\footnote{ผสฺสิกาติ (อิ, ฏฺฐ)}, อามนฺตา.\csromanspacer{} จิตฺเตน \csromanfolio{254}สหชาตาติ กตฺวา เจตสิกาติ, อามนฺตา.\csromanspacer{} เวทนาย.\csromanspacer{} สญฺญาย.\csromanspacer{} เจตนาย.\csromanspacer{} สทฺธาย.\csromanspacer{} วีริเยน.\csromanspacer{} สติยา.\csromanspacer{} สมาธินา.\csromanspacer{} ปญฺญาย.\csromanspacer{} ราเคน.\csromanspacer{} โทเสน.\csromanspacer{} โมเหน ฯเปฯ.\csromanspacer{} อโนตฺตปฺเปน สหชาตาติ กตฺวา อโนตฺตปฺปาสิกาติ\footnote{อโนตฺตปฺปิกาติ (?)}, อามนฺตา.}

\proseitem{477}{\csromansymbolmark{*}นตฺถิ เจตสิโก ธมฺโมติ, อามนฺตา.\csromanspacer{} นนุ วุตฺตํ ภควตา\csromandash{}}

\csromangathameasure{“จิตฺตญฺหิทํ เจตสิกา จ ธมฺมา, \\ อนตฺตโต สํวิทิตสฺส โหนฺติ. \\ หีนปฺปณีตํ ตทุภเย วิทิตฺวา, \\ สมฺมทฺทโส เวทิ ปโลกธมฺมนฺ”ติ.}
\csromangathagroup{\csromangathastanza{“จิตฺตญฺหิทํ เจตสิกา จ ธมฺมา, \\ อนตฺตโต สํวิทิตสฺส โหนฺติ. \\ หีนปฺปณีตํ ตทุภเย วิทิตฺวา, \\ สมฺมทฺทโส เวทิ ปโลกธมฺมนฺ”ติ.}}

\prose{อตฺเถว สุตฺตนฺโตติ, อามนฺตา.\csromanspacer{} เตน หิ อตฺถิ เจตสิโก ธมฺโมติ.}

\prose{\csromansymbolmark{*}นตฺถิ เจตสิโก ธมฺโมติ, อามนฺตา.\csromanspacer{} นนุ วุตฺตํ ภควตา “อิธ เกวฏฺฏ ภิกฺขุ ปรสตฺตานํ ปรปุคฺคลานํ จิตฺตมฺปิ อาทิสติ, เจตสิกมฺปิ อาทิสติ, วิตกฺกิตมฺปิ อาทิสติ, วิจาริตมฺปิ อาทิสติ ‘เอวมฺปิ เต มโน, อิตฺถมฺปิ เต มโน, อิติปิ เต จิตฺตนฺ’ติ”\footnote{ที 1. 207 ปิฏฺเฐ.} อตฺเถว สุตฺตนฺโตติ, อามนฺตา.\csromanspacer{} เตน หิ อตฺถิ เจตสิโก ธมฺโมติ.}

\nitthitam{เจตสิกกถา นิฏฺฐิตา.}
\chapterhead{7. สตฺตมวคฺค}
\vspace{\tipitakaparskip}
\csromancenter{\textbf{(66) 4.}\csromanspacer{}\textbf{ ทานกถา}}
\proseitem{478}{\csromansymbolmark{*}เจตสิโก ธมฺโม ทานนฺติ, อามนฺตา.\csromanspacer{} ลพฺภา เจตสิโก ธมฺโม ปเรสํ ทาตุนฺติ, น เหวํ วตฺตพฺเพ ฯเปฯ.\csromanspacer{} ลพฺภา เจตสิโก ธมฺโม ปเรสํ ทาตุนฺติ, อามนฺตา.\csromanspacer{} ลพฺภา ผสฺโส ปเรสํ ทาตุนฺติ, น เหวํ วตฺตพฺเพ ฯเปฯ.\csromanspacer{} ลพฺภา เวทนา ฯเปฯ สญฺญา.\csromanspacer{} เจตนา.\csromanspacer{} สทฺธา.\csromanspacer{} วีริยํ.\csromanspacer{} สติ.\csromanspacer{} สมาธิ.\csromanspacer{} ปญฺญา ปเรสํ ทาตุนฺติ, น เหวํ วตฺตพฺเพ ฯเปฯ.}

\chapterheadpageread{7. สตฺตมวคฺค}
\proseitem{479}{\csromanfolio{255}\csromansymbolmark{+}น วตฺตพฺพํ “เจตสิโก ธมฺโม ทานนฺ”ติ, อามนฺตา.\csromanspacer{} ทานํ อนิฏฺฐผลํ อกนฺตผลํ อมนุญฺญผลํ เสจนกผลํ ทุกฺขุทฺรยํ ทุกฺขวิปากนฺติ, น เหวํ วตฺตพฺเพ ฯเปฯ.\csromanspacer{} นนุ ทานํ อิฏฺฐผลํ กนฺตผลํ มนุญฺญผลํ อเสจนกผลํ สุขุทฺรยํ สุขวิปากนฺติ, อามนฺตา.\csromanspacer{} หญฺจิ ทานํ อิฏฺฐผลํ กนฺตผลํ มนุญฺญผลํ อเสจนกผลํ สุขุทฺรยํ สุขวิปากํ, เตน วต เร วตฺตพฺเพ “เจตสิโก ธมฺโม ทานนฺ”ติ.}

\prose{\csromansymbolmark{+}ทานํ อิฏฺฐผลํ วุตฺตํ ภควตา “จีวรํ ทานนฺ”ติ\footnote{อุปริ วุจฺจมานาย ปรวาทีปุจฺฉาย สทิสา, อฏฺฐกถา โอโลเกตพฺพา.}, อามนฺตา.\csromanspacer{} จีวรํ อิฏฺฐผลํ กนฺตผลํ มนุญฺญผลํ อเสจนกผลํ สุขุทฺรยํ สุขวิปากนฺติ, น เหวํ วตฺตพฺเพ ฯเปฯ.\csromanspacer{} ทานํ อิฏฺฐผลํ วุตฺตํ ภควตา “ปิณฺฑปาโต เสนาสนํ คิลานปจฺจยเภสชฺชปริกฺขาโร ทานนฺ”ติ, อามนฺตา.\csromanspacer{} คิลานปจฺจยเภสชฺชปริกฺขาโร อิฏฺฐผโล กนฺตผโล มนุญฺญผโล อเสจนกผโล สุขุทฺรโย สุขวิปาโกติ, น เหวํ วตฺตพฺเพ ฯเปฯ.}

\proseitem{480}{\csromansymbolmark{+}น วตฺตพฺพํ “เจตสิโก ธมฺโม ทานนฺ”ติ, อามนฺตา.\csromanspacer{} นนุ วุตฺตํ ภควตา\csromandash{}}

\csromangathameasure{“สทฺธา หิริยํ กุสลํ จ ทานํ, \\ ธมฺมา เอเต สปฺปุริสานุยาตา. \\ เอตํ หิ มคฺคํ ทิวิยํ วทนฺติ, \\ เอเตน หิ คจฺฉติ เทวโลกนฺ”ติ.}
\csromangathagroup{\csromangathastanza{“สทฺธา หิริยํ กุสลํ จ ทานํ, \\ ธมฺมา เอเต สปฺปุริสานุยาตา. \\ เอตํ หิ มคฺคํ ทิวิยํ วทนฺติ, \\ เอเตน หิ คจฺฉติ เทวโลกนฺ”ติ\footnote{อํ 3. 68 ปิฏฺเฐ.}.}}

\prose{อตฺเถว สุตฺตนฺโตติ, อามนฺตา.\csromanspacer{} เตน หิ เจตสิโก ธมฺโม ทานนฺติ.}

\prose{\csromansymbolmark{+}น วตฺตพฺพํ “จิตสิโก ธมฺโม ทานนฺ”ติ, อามนฺตา.\csromanspacer{} นนุ วุตฺตํ ภควตา “ปญฺจิมานิ ภิกฺขเว ทานานิ มหาทานานิ อคฺคญฺญานิ รตฺตญฺญานิ วํสญฺญานิ โปราณานิ อสํกิณฺณานิ อสํกิณฺณปุพฺพานิ น สํกิยนฺติ น สํกิยิสฺสนฺติ อปฺปฏิกุฏฺฐานิ สมเณหิ พฺราหฺมเณหิ วิญฺญูหิ.\csromanspacer{} กตมานิ ปญฺจ, อิธ ภิกฺขเว อริยสาวโก ปาณาติปาตํ ปหาย ปาณาติปาตา ปฏิวิรโต โหติ.\csromanspacer{} ปาณาติปาตา ปฏิวิรโต ภิกฺขเว อริยสาวโก อปริมาณานํ สตฺตานํ อภยํ เทติ อเวรํ เทติ อพฺยาพชฺฌํ\footnote{อพฺยาปชฺฌํ (สฺยา, ก)} เทติ.\csromanspacer{} อปริมาณานํ สตฺตานํ อภยํ ทตฺวา อเวรํ ทตฺวา อพฺยาพชฺฌํ ทตฺวา อปริมาณสฺส อภยสฺส อเวรสฺส อพฺยาพชฺฌสฺส ภาคี โหติ.\csromanspacer{} อิทํ \csromanfolio{256}ภิกฺขเว ปฐมํ ทานํ มหาทานํ อคฺคญฺญํ รตฺตญฺญํ วํสญฺญํ โปราณํ อสํกิณฺณํ อสํกิณฺณปุพฺพํ น สํกิยติ น สํกิยิสฺสติ, อปฺปฏิกุฏฺฐํ สมเณหิ พฺราหฺมเณหิ วิญฺญูหิ.\csromanspacer{} ปุน จปรํ ภิกฺขเว อริยสาวโก อทินฺนาทานํ ปหาย ฯเปฯ กาเมสุมิจฺฉาจารํ ปหาย ฯเปฯ มุสาวาทํ ปหาย ฯเปฯ สุราเมรยมชฺช ปมาทฏฺฐานํ ปหาย สุราเมรยมชฺชปมาทฏฺฐานา ปฏิวิรโต โหติ, สุราเมรยมชฺชปมาทฏฺฐานา ปฏิวิรโต ภิกฺขเว อริยสาวโก อปริมาณานํ สตฺตานํ อภยํ เทติ อเวรํ เทติ อพฺยาพชฺฌํ เทติ.\csromanspacer{} อปริมาณานํ สตฺตานํ อภยํ ทตฺวา อเวรํ ทตฺวา อพฺยาพชฺฌํ ทตฺวา อปริมาณสฺส อภยสฺส อเวรสฺส อพฺยาพชฺฌสฺส ภาคี โหติ.\csromanspacer{} อิทํ ภิกฺขเว ปญฺจมํ ทานํ มหาทานํ อคฺคญฺญํ รตฺตญฺญํ วํสญฺญํ โปราณํ อสํกิณฺณํ อสํกิณฺณปุพฺพํ น สํกิยติ น สํกิยิสฺสติ, อปฺปฏิกุฏฺฐํ สมเณหิ พฺราหฺมเณหิ วิญฺญูหิ.\csromanspacer{} อิมานิ โข ภิกฺขเว ปญฺจ ทานานิ มหาทานานิ อคฺคญฺญานิ รตฺตญฺญานิ วํสญฺญานิ โปราณานิ อสํกิณฺณานิ อสํกิณฺณปุพฺพานิ น สํกิยนฺติ น สํกิยิสฺสนฺติ, อปฺปฏิกุฏฺฐานิ สมเณหิ พฺราหฺมเณหิ วิญฺญูหี”ติ\footnote{อํ 3. 76 ปิฏฺเฐ.} อตฺเถว สุตฺตนฺโตติ, อามนฺตา.\csromanspacer{} เตน หิ เจตสิโก ธมฺโม ทานนฺติ.}

\proseitem{481}{\csromansymbolmark{*}น วตฺตพฺพํ “เทยฺยธมฺโม ทานนฺ”ติ, อามนฺตา.\csromanspacer{} นนุ วุตฺตํ ภควตา “อิเธกจฺโจ อนฺนํ เทติ, ปานํ เทติ, วตฺถํ เทติ, ยานํ เทติ, มาลํ เทติ, คนฺธํ เทติ, วิเลปนํ เทติ, เสยฺยํ เทติ, อาวสถํ เทติ, ปทีเปยฺยํ เทตี”ติ\footnote{สํ 2. 206 ปิฏฺเฐ, โถกํ วิสทิสํ.} อตฺเถว สุตฺตนฺโตติ, อามนฺตา.\csromanspacer{} เตน หิ เทยฺยธมฺโม ทานนฺติ.}

\proseitem{482}{\csromansymbolmark{+}เทยฺยธมฺโม ทานนฺติ, อามนฺตา.\csromanspacer{} เทยฺยธมฺโม อิฏฺฐผโล กนฺตผโล มนุญฺญผโล อเสจนกผโล สุขุทฺรโย สุขวิปาโกติ, น เหวํ วตฺตพฺเพ ฯเปฯ.\csromanspacer{} ทานํ อิฏฺฐผลํ วุตฺตํ ภควตา “จีวรํ ทานนฺ”ติ, อามนฺตา.\csromanspacer{} จีวรํ อิฏฺฐผลํ กนฺตผลํ มนุญฺญผลํ อเสจนกผลํ สุขุทฺรยํ สุขวิปากนฺติ, น เหวํ วตฺตพฺเพ ฯเปฯ.\csromanspacer{} ทานํ อิฏฺฐผลํ วุตฺตํ ภควตา “ปิณฺฑปาโต ทานํ.\csromanspacer{} เสนาสนํ ทานํ.\csromanspacer{} คิลานปจฺจยเภสชฺชปริกฺขาโร ทานนฺ”ติ, อามนฺตา.\csromanspacer{} คิลานปจฺจยเภสชฺชปริกฺขาโร อิฏฺฐผโล กนฺตผโล มนุญฺญผโล}

\vspace{\tipitakaparskip}
\csromancenter{สตฺตมวคฺค อเสจนกผโล สุขุทฺรโย สุขวิปาโกติ, น เหวํ วตฺตพฺเพ ฯเปฯ.\csromanspacer{} เตน หิ น วตฺตพฺพํ “เทยฺยธมฺโม ทานนฺ”ติ.}
\nitthitam{ทานกถา นิฏฺฐิตา.}
\chapterhead{7. สตฺตมวคฺค}
\vspace{\tipitakaparskip}
\csromancenter{\textbf{(67) 5.}\csromanspacer{}\textbf{ ปริโภคมย ปุญฺญกถา}}
\csromanmatikaanchor{mtk.92}
\csromantocmarkref{subsection}{(67) 5. ปริโภคมยปุญฺญกถา}{mtk.92}
\bookmark[dest={mtk.92},level=2]{(67) 5. ปริโภคมยปุญฺญกถา}
\proseitem{483}{\csromanfolio{257}\csromansymbolmark{*}ปริโภคมยํ ปุญฺญํ วฑฺฒตีติ, อามนฺตา.\csromanspacer{} ปริโภคมโย ผสฺโส วฑฺฒติ, เวทนา วฑฺฒติ, สญฺญา วฑฺฒติ, เจตนา วฑฺฒติ, จิตฺตํ วฑฺฒติ, สทฺธา วฑฺฒติ, วีริยํ วฑฺฒติ, สติ วฑฺฒติ, สมาธิ วฑฺฒติ, ปญฺญา วฑฺฒตีติ, น เหวํ วตฺตพฺเพ ฯเปฯ.}

\prose{\csromansymbolmark{*}ปริโภคมยํ ปุญฺญํ วฑฺฒตีติ, อามนฺตา.\csromanspacer{} ลตา วิย วฑฺฒติ, มาลุวา วิย วฑฺฒติ, รุกฺโข วิย วฑฺฒติ, ติณํ วิย วฑฺฒติ, มุญฺชปุญฺโช วิย วฑฺฒตีติ, น เหวํ วตฺตพฺเพ ฯเปฯ.}

\proseitem{484}{\csromansymbolmark{*}ปริโภคมยํ ปุญฺญํ วฑฺฒตีติ, อามนฺตา.\csromanspacer{} ทายโก ทานํ ทตฺวา น สมนฺนาหรติ โหติ ปุญฺญนฺติ, อามนฺตา.\csromanspacer{} อนาวฏฺเฏนฺตสฺส\footnote{อนาวฏฺฏนฺตสฺส (สี, อิ, ก), อนาวชฺชนฺตสฺส (สฺยา)} โหติ.\csromanspacer{} อนาโภคสฺส โหติ.\csromanspacer{} อสมนฺนาหรนฺตสฺส โหติ.\csromanspacer{} อมนสิกโรนฺตสฺส โหติ.\csromanspacer{} อเจตยนฺตสฺส โหติ.\csromanspacer{} อปตฺถยนฺตสฺส โหติ.\csromanspacer{} อปฺปณิทหนฺตสฺส โหตีติ, น เหวํ วตฺตพฺเพ ฯเปฯ.\csromanspacer{} นนุ อาวฏฺเฏนฺตสฺส โหติ.\csromanspacer{} อาโภคสฺส โหติ.\csromanspacer{} สมนฺนาหรนฺตสฺส โหติ.\csromanspacer{} มนสิกโรนฺตสฺส โหติ.\csromanspacer{} เจตยนฺตสฺส โหติ.\csromanspacer{} ปตฺถยนฺตสฺส โหติ.\csromanspacer{} ปณิทหนฺตสฺส โหตีติ, อามนฺตา.\csromanspacer{} หญฺจิ อาวฏฺเฏนฺตสฺส โหติ.\csromanspacer{} อาโภคสฺส โหติ.\csromanspacer{} สมนฺนาหรนฺตสฺส โหติ.\csromanspacer{} มนสิกโรนฺตสฺส โหติ.\csromanspacer{} เจตยนฺตสฺส โหติ.\csromanspacer{} ปตฺถยนฺตสฺส โหติ.\csromanspacer{} ปณิทหนฺตสฺส โหติ, โน จ วต เร วตฺตพฺเพ “ปริโภคมยํ ปุญฺญํ วฑฺฒตี”ติ.}

\proseitem{485}{\csromansymbolmark{*}ปริโภคมยํ ปุญฺญํ วฑฺฒตีติ, อามนฺตา.\csromanspacer{} ทายโก ทานํ ทตฺวา กามวิตกฺกํ วิตกฺเกติ.\csromanspacer{} พฺยาปาทวิตกฺกํ วิตกฺเกติ.\csromanspacer{} วิหิํสาวิตกฺกํ \csromanfolio{258}วิตกฺเกติ โหติ ปุญฺญนฺติ, อามนฺตา.\csromanspacer{} ทฺวินฺนํ ผสฺสานํ ฯเปฯ.\csromanspacer{} ทฺวินฺนํ จิตฺตานํ สโมธานํ โหตีติ, น เหวํ วตฺตพฺเพ ฯเปฯ.\csromanspacer{} ทฺวินฺนํ ผสฺสานํ ฯเปฯ.\csromanspacer{} ทฺวินฺนํ จิตฺตานํ สโมธานํ โหตีติ, อามนฺตา.\csromanspacer{} กุสลากุสลา สาวชฺชานวชฺชา หีนปณีตา กณฺหสุกฺกสปฺปฏิภาคา ธมฺมา สมฺมุขีภาวํ อาคจฺฉนฺตีติ, น เหวํ วตฺตพฺเพ ฯเปฯ.\csromanspacer{} กุสลากุสลา สาวชฺชานวชฺชา หีนปณีตา กณฺหสุกฺกสปฺปฏิภาคา ธมฺมา สมฺมุขีภาวํ อาคจฺฉนฺตีติ, อามนฺตา.\csromanspacer{} นนุ วุตฺตํ ภควตา “จตฺตาริมานิ ภิกฺขเว สุวิทูรวิทูรานิ.\csromanspacer{} กตมานิ จตฺตาริ, นภญฺจ ภิกฺขเว ปถวี จ อิทํ ปฐมํ สุวิทูรวิทูรํ.\csromanspacer{} โอริมญฺจ ภิกฺขเว ตีรํ สมุทฺทสฺส ปาริมญฺจ ตีรํ อิทํ ทุติยํ สุวิทูรวิทูรํ.\csromanspacer{} ยโต จ ภิกฺขเว เวโรจโน อพฺภุเทติ ยตฺถ จ อตฺถเมติ อิทํ ตติยํ สุวิทูรวิทูรํ.\csromanspacer{} สตญฺจ ภิกฺขเว ธมฺโม อสตญฺจ ธมฺโม อิทํ จตุตฺถํ สุวิทูรวิทูรํ.\csromanspacer{} อิมานิ โข ภิกฺขเว จตฺตาริ สุวิทูรวิทูรานีติ.}

\csromangathameasure{นภญฺจ ทูเร ปถวี จ ทูเร, \\ ปารํ สมุทฺทสฺส ตทาหุ ทูเร. \\ ยโต จ เวโรจโน อพฺภุเทติ, \\ ปภงฺกโร ยตฺถ จ อตฺถเมติ. \\ ตโต หเว ทูรตรํ วทนฺติ, \\ สตญฺจ ธมฺมํ อสตญฺจ ธมฺมํ. \\ อพฺยายิโก โหติ สตํ สมาคโม, \\ ยาวมฺปิ ติฏฺเฐยฺย ตเถว โหติ. \\ ขิปฺปํ หิ เวติ อสตํ สมาคโม, \\ ตสฺมา สตํ ธมฺโม อสพฺภิ อารกา”ติ.}
\csromangathagroup{\csromangathastanza{นภญฺจ ทูเร ปถวี จ ทูเร, \\ ปารํ สมุทฺทสฺส ตทาหุ ทูเร. \\ ยโต จ เวโรจโน อพฺภุเทติ, \\ ปภงฺกโร ยตฺถ จ อตฺถเมติ.}\csromangathastanza{ตโต หเว ทูรตรํ วทนฺติ, \\ สตญฺจ ธมฺมํ อสตญฺจ ธมฺมํ. \\ อพฺยายิโก โหติ สตํ สมาคโม, \\ ยาวมฺปิ ติฏฺเฐยฺย ตเถว โหติ. \\ ขิปฺปํ หิ เวติ\footnote{ขิปฺปํหเวติ (พหูสุ)} อสตํ สมาคโม, \\ ตสฺมา สตํ ธมฺโม อสพฺภิ อารกา”ติ\footnote{ขิปฺปํหเวติ (พหูสุ)}.}}

\prose{อตฺเถว สุตฺตนฺโตติ, อามนฺตา.\csromanspacer{} เตน หิ น วตฺตพฺพํ “กุสลากุสลา สาวชฺชานวชฺชา หีนปณีตา กณฺหสุกฺกสปฺปฏิภาคา ธมฺมา สมฺมุขีภาวํ อาคจฺฉนฺตี”ติ.}

\proseitem{486}{\csromansymbolmark{+}น วตฺตพฺพํ “ปริโภคมยํ ปุญฺญํ วฑฺฒตี”ติ, อามนฺตา.\csromanspacer{} นนุ วุตฺตํ ภควตา\csromandash{}}

\csromangathameasure{“อารามโรปา วนโรปา, เย ชนา เสตุการกา. \\ ปปญฺจ อุทปานญฺจ, เย ททนฺติ อุปสฺสยํ.}
\csromangathagroup{\csromangathastanza{“อารามโรปา วนโรปา, เย ชนา เสตุการกา. \\ ปปญฺจ อุทปานญฺจ, เย ททนฺติ อุปสฺสยํ.}}

\chapterheadpageread{7. สตฺตมวคฺค}
\csromangathameasure{เตสํ ทิวา จ รตฺโต จ, สทา ปุญฺญํ ปวฑฺฒติ. \\ ธมฺมฏฺฐา สีลสมฺปนฺนา, เต ชนา สคฺคคามิโน”ติ.}
\csromangathagroup{\csromangathastanza{\csromanfolio{259}เตสํ ทิวา จ รตฺโต จ, สทา ปุญฺญํ ปวฑฺฒติ. \\ ธมฺมฏฺฐา สีลสมฺปนฺนา, เต ชนา สคฺคคามิโน”ติ\footnote{สํ 1. 30 ปิฏฺเฐ.}.}}

\prose{อตฺเถว สุตฺตนฺโตติ, อามนฺตา.\csromanspacer{} เตน หิ ปริโภคมยํ ปุญฺญํ วฑฺฒตีติ.}

\prose{\csromansymbolmark{+}น วตฺตพฺพํ “ปริโภคมยํ ปุญฺญํ วฑฺฒตี”ติ, อามนฺตา.\csromanspacer{} นนุ วุตฺตํ ภควตา “จตฺตาโรเม ภิกฺขเว ปุญฺญาภิสนฺทา กุสลาภิสนฺทา สุขสฺสาหารา โสวคฺคิกา สุขวิปากา สคฺคสํวตฺตนิกา อิฏฺฐาย กนฺตาย มนาปาย หิตาย สุขาย สํวตฺตนฺติ.\csromanspacer{} กตเม จตฺตาโร, ยสฺส ภิกฺขเว ภิกฺขุ จีวรํ ปริภุญฺชมาโน อปฺปมาณํ เจโตสมาธิํ อุปสมฺปชฺช วิหรติ.\csromanspacer{} อปฺปมาโณ ตสฺส ปุญฺญาภิสนฺโท กุสลาภิสนฺโท สุขสฺสาหาโร โสวคฺคิโก สุขวิปาโก สคฺคสํวตฺตนิโก อิฏฺฐาย กนฺตาย มนาปาย หิตาย สุขาย สํวตฺตติ.\csromanspacer{} ยสฺส ภิกฺขเว ภิกฺขุ ปิณฺฑปาตํ ปริภุญฺชมาโน ฯเปฯ เสนาสนํ ปริภุญฺชมาโน ฯเปฯ คิลานปจฺจยเภสชฺชปริกฺขารํ ปริภุญฺชมาโน อปฺปมาณํ เจโตสมาธิํ อุปสมฺปชฺช วิหรติ.\csromanspacer{} อปฺปมาโณ ตสฺส ปุญฺญาภิสนฺโท กุสลาภิสนฺโท สุขสฺสาหาโร โสวคฺคิโก สุขวิปาโก สคฺคสํวตฺตนิโก อิฏฺฐาย กนฺตาย มนาปาย หิตาย สุขาย สํวตฺตติ.\csromanspacer{} อิเม โข ภิกฺขเว จตฺตาโร ปุญฺญาภิสนฺทา กุสลาภิสนฺทา สุขสฺสาหารา โสวคฺคิกา สุขวิปากา สคฺคสํวตฺตนิกา อิฏฺฐาย กนฺตาย มนาปาย หิตาย สุขาย สํวตฺตนฺตี”ติ\footnote{อํ 1. 365 ปิฏฺเฐ.} อตฺเถว สุตฺตนฺโตติ, อามนฺตา.\csromanspacer{} เตน หิ ปริโภคมยํ ปุญฺญํ วฑฺฒตีติ.}

\proseitem{487}{\csromansymbolmark{*}ปริโภคมยํ ปุญฺญํ วฑฺฒตีติ, อามนฺตา.\csromanspacer{} ทายโก ทานํ เทติ, ปฏิคฺคาหโก ปฏิคฺคเหตฺวา น ปริภุญฺชติ ฉฑฺเฑติ วิสฺสชฺเชติ, โหติ ปุญฺญนฺติ, อามนฺตา.\csromanspacer{} หญฺจิ ทายโก ทานํ เทติ, ปฏิคฺคาหโก ปฏิคฺคเหตฺวา น ปริภุญฺชติ ฉฑฺเฑติ วิสฺสชฺเชติ, โหติ ปุญฺญํ, โน จ วต เร วตฺตพฺเพ “ปริโภคมยํ ปุญฺญํ วฑฺฒตี”ติ.}

\csromanmatikaanchor{mtk.93}
\csromantocmarkref{subsection}{(68) 6. อิโตทินฺนกถา}{mtk.93}
\bookmark[dest={mtk.93},level=2]{(68) 6. อิโตทินฺนกถา}
\prose{\csromansymbolmark{*}ปริโภคมยํ ปุญฺญํ วฑฺฒตีติ, อามนฺตา.\csromanspacer{} ทายโก ทานํ เทติ, ปฏิคฺคาหเก ปฏิคฺคหิเต ราชาโน วา หรนฺติ, โจรา วา หรนฺติ, อคฺคิ วา ทหติ, อุทกํ วา วหติ, อปฺปิยา วา ทายาทา หรนฺติ, โหติ ปุญฺญนฺติ, \csromanfolio{260}อามนฺตา.\csromanspacer{} หญฺจิ ทายโก ทานํ เทติ, ปฏิคฺคาหเก ปฏิคฺคหิเต ราชาโน วา หรนฺติ, โจรา วา หรนฺติ, อคฺคิ วา ทหติ, อุทกํ วา วหติ, อปฺปิยา วา ทายาทา หรนฺติ, โหติ ปุญฺญํ, โน จ วต เร วตฺตพฺเพ “ปริโภคมยํ ปุญฺญํ วฑฺฒตี”ติ.}

\nitthitam{ปริโภคมยปุญฺญกถา นิฏฺฐิตา.}
\chapterhead{7. สตฺตมวคฺค}
\vspace{\tipitakaparskip}
\csromancenter{\textbf{(68) 6.}\csromanspacer{}\textbf{ อิโตทินฺนกถา}}
\proseitem{488}{\csromansymbolmark{*}อิโต ทินฺเนน ตตฺถ ยาเปนฺตีติ, อามนฺตา.\csromanspacer{} อิโต จีวรํ เทนฺติ ตํ จีวรํ ตตฺถ ปริภุญฺชนฺตีติ, น เหวํ วตฺตพฺเพ ฯเปฯ.\csromanspacer{} อิโต ปิณฺฑปาตํ เทนฺติ, อิโต เสนาสนํ เทนฺติ, อิโต คิลานปจฺจยเภสชฺชปริกฺขารํ เทนฺติ, อิโต ขาทนียํ เทนฺติ, อิโต โภชนียํ เทนฺติ, อิโต ปานียํ เทนฺติ, ตํ ปานียํ ตตฺถ ปริภุญฺชนฺตีติ, น เหวํ วตฺตพฺเพ ฯเปฯ.}

\prose{\csromansymbolmark{*}อิโต ทินฺเนน ตตฺถ ยาเปนฺตีติ, อามนฺตา.\csromanspacer{} อญฺโญ อญฺญสฺส การโก, ปรกตํ สุขทุกฺขํ, อญฺโญ กโรติ อญฺโญ ปฏิสํเวเทตีติ, น เหวํ วตฺตพฺเพ ฯเปฯ.}

\proseitem{489}{\csromansymbolmark{+}น วตฺตพฺพํ “อิโต ทินฺเนน ตตฺถ ยาเปนฺตี”ติ, อามนฺตา.\csromanspacer{} นนุ เปตา อตฺตโน อตฺถาย ทานํ เทนฺตํ อนุโมเทนฺติ จิตฺตํ ปสาเทนฺติ ปีติํ อุปฺปาเทนฺติ โสมนสฺสํ ปฏิลภนฺตีติ, อามนฺตา.\csromanspacer{} หญฺจิ เปตา อตฺตโน อตฺถาย ทานํ เทนฺตํ อนุโมเทนฺติ จิตฺตํ ปสาเทนฺติ ปีติํ อุปฺปาเทนฺติ โสมนสฺสํ ปฏิลภนฺติ, เตน วต เร วตฺตพฺเพ “อิโต ทินฺเนน ตตฺถ ยาเปนฺตี”ติ.}

\proseitem{490}{\csromansymbolmark{+}น วตฺตพฺพํ “อิโต ทินฺเนน ตตฺถ ยาเปนฺตี”ติ, อามนฺตา.\csromanspacer{} นนุ วุตฺตํ ภควตา\csromandash{}}

\csromangathameasure{“อุนฺนเม อุทกํ วุฏฺฐํ, ยถานินฺนํ ปวตฺตติ. \\ เอวเมว อิโต ทินฺนํ, เปตานํ อุปกปฺปติ.}
\csromangathagroup{\csromangathastanza{“อุนฺนเม อุทกํ วุฏฺฐํ, ยถานินฺนํ ปวตฺตติ. \\ เอวเมว อิโต ทินฺนํ, เปตานํ อุปกปฺปติ.}}

\chapterheadpageread{7. สตฺตมวคฺค}
\csromangathameasure{ยถา วาริวหา ปูรา, ปริปูเรนฺติ สาครํ. \\ เอวเมว อิโต ทินฺนํ, เปตานํ อุปกปฺปติ. \\ น หิ ตตฺถ กสี อตฺถิ, โครกฺเขตฺถ น วิชฺชติ. \\ วณิชฺชา ตาทิสี นตฺถิ, หิรญฺเญน กยากยํ. \\ อิโต ทินฺเนน ยาเปนฺติ, เปตา กาลงฺกตา ตหินฺ”ติ.}
\csromangathagroup{\csromangathastanza{\csromanfolio{261}ยถา วาริวหา ปูรา, ปริปูเรนฺติ สาครํ. \\ เอวเมว อิโต ทินฺนํ, เปตานํ อุปกปฺปติ.}\csromangathastanza{น หิ ตตฺถ กสี อตฺถิ, โครกฺเขตฺถ น วิชฺชติ. \\ วณิชฺชา ตาทิสี นตฺถิ, หิรญฺเญน กยากยํ\footnote{กยกฺกยํ (สี, อิ)}. \\ อิโต ทินฺเนน ยาเปนฺติ, เปตา กาลงฺกตา ตหินฺ”ติ\footnote{กยกฺกยํ (สี, อิ)}.}}

\prose{อตฺเถว สุตฺตนฺโตติ, อามนฺตา.\csromanspacer{} เตน หิ อิโต ทินฺเนน ตตฺถ ยาเปนฺตีติ.}

\proseitem{491}{\csromansymbolmark{+}น วตฺตพฺพํ “อิโต ทินฺเนน ตตฺถ ยาเปนฺตี”ติ, อามนฺตา.\csromanspacer{} นนุ วุตฺตํ ภควตา “ปญฺจิมานิ ภิกฺขเว ฐานานิ สมฺปสฺสนฺตา มาตาปิตโร ปุตฺตํ อิจฺฉนฺติ กุเล ชายมานํ.\csromanspacer{} กตมานิ ปญฺจ, ภโฏ วา โน ภริสฺสติ, กิจฺจํ วา โน กริสฺสติ.\csromanspacer{} กุลวํโส จิรํ ฐสฺสติ, ทายชฺชํ ปฏิปชฺชิสฺสติ.\csromanspacer{} อถ วา ปน เปตานํ กาลงฺกตานํ ทกฺขิณํ อนุปฺปทสฺสติ.\csromanspacer{} อิมานิ โข ภิกฺขเว ปญฺจ ฐานานิ สมฺปสฺสนฺตา มาตาปิตโร ปุตฺตํ อิจฺฉนฺติ กุเล ชายมานนฺติ.}

\csromangathameasure{ปญฺจ ฐานานิ สมฺปสฺสํ, ปุตฺตํ อิจฺฉนฺติ ปณฺฑิตา. \\ ภโฏ วา โน ภริสฺสติ, กิจฺจํ วา โน กริสฺสติ. \\ กุลวํโส จิรํ ติฏฺเฐ, ทายชฺชํ ปฏิปชฺชติ. \\ อถ วา ปน เปตานํ, ทกฺขิณํ อนุปฺปทสฺสติ. \\ ฐานาเนตานิ สมฺปสฺสํ, ปุตฺตํ อิจฺฉนฺติ ปณฺฑิตา. \\ ตสฺมา สนฺโต สปฺปุริสา, กตญฺญู กตเวทิโน. \\ ภรนฺติ มาตาปิตโร, ปุพฺเพ กตมนุสฺสรํ. \\ กโรนฺติ เตสํ กิจฺจานิ, ยถา ตํ ปุพฺพการินํ. \\ โอวาทการี ภฏโปสี, กุลวํสํ อหาปยํ. \\ สทฺโธ สีเลน สมฺปนฺโน, ปุตฺโต โหติ ปสํสิโย”ติ.}
\csromangathagroup{\csromangathastanza{ปญฺจ ฐานานิ สมฺปสฺสํ, ปุตฺตํ อิจฺฉนฺติ ปณฺฑิตา. \\ ภโฏ วา โน ภริสฺสติ, กิจฺจํ วา โน กริสฺสติ.}\csromangathastanza{กุลวํโส จิรํ ติฏฺเฐ, ทายชฺชํ ปฏิปชฺชติ. \\ อถ วา ปน เปตานํ, ทกฺขิณํ อนุปฺปทสฺสติ.}\csromangathastanza{ฐานาเนตานิ สมฺปสฺสํ, ปุตฺตํ อิจฺฉนฺติ ปณฺฑิตา. \\ ตสฺมา สนฺโต สปฺปุริสา, กตญฺญู กตเวทิโน.}\csromangathastanza{ภรนฺติ มาตาปิตโร, ปุพฺเพ กตมนุสฺสรํ. \\ กโรนฺติ เตสํ กิจฺจานิ, ยถา ตํ ปุพฺพการินํ.}\csromangathastanza{โอวาทการี ภฏโปสี, กุลวํสํ อหาปยํ. \\ สทฺโธ สีเลน สมฺปนฺโน, ปุตฺโต โหติ ปสํสิโย”ติ\footnote{อํ 2. 37 ปิฏฺเฐ.}.}}

\prose{อตฺเถว สุตฺตนฺโตติ, อามนฺตา.\csromanspacer{} เตน หิ อิโต ทินฺเนน ตตฺถ ยาเปนฺตีติ.}

\nitthitam{อิโต ทินฺนกถา นิฏฺฐิตา.}
\chapterheadpageread{7. สตฺตมวคฺค}
\vspace{\tipitakaparskip}
\csromancenter{\textbf{(69) 7.}\csromanspacer{}\textbf{ ปถวีกมฺมวิปาโกติกถา}}
\csromanmatikaanchor{mtk.94}
\csromantocmarkref{subsection}{(69) 7. ปถวีกมฺมวิปาโกติกถา}{mtk.94}
\bookmark[dest={mtk.94},level=2]{(69) 7. ปถวีกมฺมวิปาโกติกถา}
\proseitem{492}{\csromanfolio{262}\csromansymbolmark{*}ปถวี กมฺมวิปาโกติ, อามนฺตา.\csromanspacer{} สุขเวทนิยา, ทุกฺขเวทนิยา, อทุกฺขมสุขเวทนิยา, สุขาย เวทนาย สมฺปยุตฺตา, ทุกฺขาย เวทนาย สมฺปยุตฺตา, อทุกฺขมสุขาย เวทนาย สมฺปยุตฺตา, ผสฺเสน สมฺปยุตฺตา, เวทนาย สมฺปยุตฺตา, สญฺญาย สมฺปยุตฺตา, เจตนาย สมฺปยุตฺตา, จิตฺเตน สมฺปยุตฺตา, สารมฺมณา, อตฺถิ ตาย อาวฏฺฏนา\footnote{อาวชฺชนา (สฺยา, กํ)} อาโภโค สมนฺนาหาโร มนสิกาโร เจตนา ปตฺถนา ปณิธีติ, น เหวํ วตฺตพฺเพ ฯเปฯ.\csromanspacer{} นนุ น สุขเวทนิยา น ทุกฺขเวทนิยา น อทุกฺขมสุขเวทนิยา น สุขาย เวทนาย สมฺปยุตฺตา น ทุกฺขาย เวทนาย สมฺปยุตฺตา น อทุกฺขมสุขาย เวทนาย สมฺปยุตฺตา น ผสฺเสน สมฺปยุตฺตา น เวทนาย สมฺปยุตฺตา น สญฺญาย สมฺปยุตฺตา น เจตนาย สมฺปยุตฺตา น จิตฺเตน สมฺปยุตฺตา อนารมฺมณา นตฺถิ ตาย อาวฏฺฏนา อาโภโค สมนฺนาหาโร มนสิกาโร เจตนา ปตฺถนา ปณิธีติ, อามนฺตา.\csromanspacer{} หญฺจิ น สุขเวทนิยา น ทุกฺขเวทนิยา ฯเปฯ อนารมฺมณา นตฺถิ ตาย อาวฏฺฏนา ฯเปฯ ปณิธิ, โน จ วต เร วตฺตพฺเพ “ปถวี กมฺมวิปาโก”ติ.}

\prose{\csromansymbolmark{*}ผสฺโส กมฺมวิปาโก ผสฺโส สุขเวทนิโย ทุกฺขเวทนิโย อทุกฺขมสุขเวทนิโย สุขาย เวทนาย สมฺปยุตฺโต ทุกฺขาย ฯเปฯ อทุกฺขมสุขาย เวทนาย สมฺปยุตฺโต ผสฺเสน สมฺปยุตฺโต เวทนาย สมฺปยุตฺโต สญฺญาย สมฺปยุตฺโต เจตนาย สมฺปยุตฺโต จิตฺเตน สมฺปยุตฺโต สารมฺมโณ อตฺถิ ตสฺส อาวฏฺฏนา ฯเปฯ ปณิธีติ, อามนฺตา.\csromanspacer{} ปถวี กมฺมวิปาโก ปถวี สุขเวทนิยา ทุกฺขเวทนิยา อทุกฺขมสุขเวทนิยา สุขาย เวทนาย สมฺปยุตฺตา ทุกฺขาย ฯเปฯ อทุกฺขมสุขาย เวทนาย สมฺปยุตฺตา ผสฺเสน สมฺปยุตฺตา เวทนาย สมฺปยุตฺตา สญฺญาย สมฺปยุตฺตา เจตนาย สมฺปยุตฺตา จิตฺเตน สมฺปยุตฺตา สารมฺมณา อตฺถิ ตาย อาวฏฺฏนา อาโภโค สมนฺนาหาโร มนสิกาโร เจตนา ปตฺถนา ปณิธีติ, น เหวํ วตฺตพฺเพ ฯเปฯ.\csromanspacer{} ปถวี กมฺมวิปาโก ปถวี น สุขเวทนิยา น ทุกฺขเวทนิยา ฯเปฯ อนารมฺมณา นตฺถิ ตาย อาวฏฺฏนา ฯเปฯ ปณิธีติ, อามนฺตา.\csromanspacer{} ผสฺโส กมฺมวิปาโก ผสฺโส น สุขเวทนิโย น}

\vspace{\tipitakaparskip}
\csromancenter{สตฺตมวคฺค ทุกฺขเวทนิโย ฯเปฯ อนารมฺมโณ นตฺถิ ตสฺส อาวฏฺฏนา ฯเปฯ ปณิธีติ, น เหวํ วตฺตพฺเพ ฯเปฯ.}
\prose{\csromanfolio{263}\csromansymbolmark{*}ปถวี กมฺมวิปาโกติ, อามนฺตา.\csromanspacer{} ปถวี ปคฺคหนิคฺคหุปคา เฉทนเภทนุปคาติ, อามนฺตา.\csromanspacer{} กมฺมวิปาโก ปคฺคหนิคฺคหุปโค เฉทนเภทนุปโคติ, น เหวํ วตฺตพฺเพ ฯเปฯ.}

\prose{\csromansymbolmark{*}ลพฺภา ปถวี เกตุํ วิกฺเกตุํ อาฐเปตุํ โอจินิตุํ วิจินิตุนฺติ, อามนฺตา.\csromanspacer{} ลพฺภา กมฺมวิปาโก เกตุํ วิกฺเกตุํ อาฐเปตุํ โอจินิตุํ วิจินิตุนฺติ, น เหวํ วตฺตพฺเพ ฯเปฯ.}

\proseitem{493}{\csromansymbolmark{*}ปถวี ปเรสํ สาธารณาติ, อามนฺตา.\csromanspacer{} กมฺมวิปาโก ปเรสํ สาธารโณติ, น เหวํ วตฺตพฺเพ ฯเปฯ.}

\prose{กมฺมวิปาโก ปเรสํ สาธารโณติ, อามนฺตา.\csromanspacer{} นนุ วุตฺตํ ภควตา\csromandash{}}

\csromangathameasure{“อสาธารณมญฺเญสํ, อโจรหรโณ นิธิ. \\ กยิราถ มจฺโจ ปุญฺญานิ, สเจ สุจริตํ จเร”ติ. อตฺเถว สุตฺตนฺโตติ, อามนฺตา. เตน หิ น วตฺตพฺพํ “กมฺมวิปาโก ปเรสํ สาธารโณ”ติ.}
\csromangathagroup{\csromangathastanza{“อสาธารณมญฺเญสํ, อโจรหรโณ นิธิ. \\ กยิราถ มจฺโจ ปุญฺญานิ, สเจ สุจริตํ จเร”ติ\footnote{ขุ 1. 10 ขุทฺทกปาเฐ, โถกํ วิสทิสํ.}.\csromanspacer{} อตฺเถว สุตฺตนฺโตติ, อามนฺตา.\csromanspacer{} เตน หิ น วตฺตพฺพํ “กมฺมวิปาโก ปเรสํ สาธารโณ”ติ.}}

\prose{\csromansymbolmark{*}ปถวี กมฺมวิปาโกติ, อามนฺตา.\csromanspacer{} ปฐมํ ปถวี สณฺฐาติ ปจฺฉา สตฺตา อุปฺปชฺชนฺตีติ, อามนฺตา.\csromanspacer{} ปฐมํ วิปาโก อุปฺปชฺชติ ปจฺฉา วิปากปฏิลาภาย กมฺมํ กโรนฺตีติ, น เหวํ วตฺตพฺเพ ฯเปฯ.}

\prose{\csromansymbolmark{*}ปถวี สพฺพสตฺตานํ กมฺมวิปาโกติ, อามนฺตา.\csromanspacer{} สพฺเพ สตฺตา ปถวิํ ปริภุญฺชนฺตีติ, น เหวํ วตฺตพฺเพ ฯเปฯ.\csromanspacer{} สพฺเพ สตฺตา ปถวิํ ปริภุญฺชนฺตีติ, อามนฺตา.\csromanspacer{} อตฺถิ เกจิ ปถวิํ อปริภุญฺชิตฺวา ปรินิพฺพายนฺตีติ, อามนฺตา.\csromanspacer{} อตฺถิ เกจิ กมฺมวิปากํ อเขเปตฺวา ปรินิพฺพายนฺตีติ, น เหวํ วตฺตพฺเพ ฯเปฯ.}

\csromanmatikaanchor{mtk.95}
\csromantocmarkref{subsection}{(70) 8. ชรามรณํวิปาโกติกถา}{mtk.95}
\bookmark[dest={mtk.95},level=2]{(70) 8. ชรามรณํวิปาโกติกถา}
\prose{\csromansymbolmark{*}ปถวี จกฺกวตฺติสตฺตสฺส กมฺมวิปาโกติ, อามนฺตา.\csromanspacer{} อญฺเญ สตฺตา ปถวิํ ปริภุญฺชนฺตีติ, อามนฺตา.\csromanspacer{} จกฺกวตฺติสตฺตสฺส กมฺมวิปากํ อญฺเญ สตฺตา ปริภุญฺชนฺตีติ, น เหวํ วตฺตพฺเพ ฯเปฯ.\csromanspacer{} จกฺกวตฺติสตฺตสฺส กมฺมวิปากํ อญฺเญ สตฺตา ปริภุญฺชนฺตีติ, อามนฺตา.\csromanspacer{} จกฺกวตฺติสตฺตสฺส ผสฺสํ เวทนํ สญฺญํ \csromanfolio{264}เจตนํ จิตฺตํ สทฺธํ วีริยํ สติํ สมาธิํ ปญฺญํ อญฺเญ สตฺตา ปริภุญฺชนฺตีติ, น เหวํ วตฺตพฺเพ ฯเปฯ.}

\proseitem{494}{\csromansymbolmark{+}น วตฺตพฺพํ “ปถวี กมฺมวิปาโก”ติ, อามนฺตา.\csromanspacer{} นนุ อตฺถิ อิสฺสริยสํวตฺตนิยํ กมฺมํ อธิปจฺจสํวตฺตนิยํ กมฺมนฺติ, อามนฺตา.\csromanspacer{} หญฺจิ อตฺถิ อิสฺสริยสํวตฺตนิยํ กมฺมํ อธิปจฺจสํวตฺตนิยํ กมฺมํ, เตน วต เร วตฺตพฺเพ “ปถวี กมฺมวิปาโก”ติ.}

\nitthitam{ปถวี กมฺมวิปาโกติกถา นิฏฺฐิตา.}
\chapterhead{7. สตฺตมวคฺค}
\vspace{\tipitakaparskip}
\csromancenter{\textbf{(70) 8.}\csromanspacer{}\textbf{ ชรามรณํวิปาโกติกถา}}
\proseitem{495}{\csromansymbolmark{*}ชรามรณํ วิปาโกติ, อามนฺตา.\csromanspacer{} สุขเวทนิยํ ทุกฺขเวทนิยํ อทุกฺขมสุขเวทนิยํ, สุขาย เวทนาย สมฺปยุตฺตํ, ทุกฺขาย เวทนาย สมฺปยุตฺตํ, อทุกฺขมสุขาย เวทนาย สมฺปยุตฺตํ, ผสฺเสน สมฺปยุตฺตํ, เวทนาย สมฺปยุตฺตํ, สญฺญาย สมฺปยุตฺตํ, เจตนาย สมฺปยุตฺตํ, จิตฺเตน สมฺปยุตฺตํ, สารมฺมณํ, อตฺถิ ตสฺส อาวฏฺฏนา อาโภโค สมนฺนาหาโร มนสิกาโร เจตนา ปตฺถนา ปณิธีติ, น เหวํ วตฺตพฺเพ ฯเปฯ.\csromanspacer{} นนุ น สุขเวทนิยํ น ทุกฺขเวทนิยํ ฯเปฯ อนารมฺมณํ นตฺถิ ตสฺส อาวฏฺฏนา ฯเปฯ ปณิธีติ, อามนฺตา.\csromanspacer{} หญฺจิ น สุขเวทนิยํ น ทุกฺขเวทนิยํ ฯเปฯ อนารมฺมณํ, นตฺถิ ตสฺส อาวฏฺฏนา ฯเปฯ ปณิธิ, โน จ วต เร วตฺตพฺเพ “ชรามรณํ วิปาโก”ติ.}

\prose{\csromansymbolmark{*}ผสฺโส วิปาโก ผสฺโส สุขเวทนิโย ทุกฺขเวทนิโย ฯเปฯ สารมฺมโณ อตฺถิ ตสฺส อาวฏฺฏนา ฯเปฯ ปณิธีติ, อามนฺตา.\csromanspacer{} ชรามรณํ วิปาโก ชรามรณํ สุขเวทนิยํ ทุกฺขเวทนิยํ ฯเปฯ สารมฺมณํ อตฺถิ ตสฺส อาวฏฺฏนา ฯเปฯ ปณิธีติ, น เหวํ วตฺตพฺเพ ฯเปฯ.}

\prose{\csromansymbolmark{*}ชรามรณํ วิปาโก ชรามรณํ น สุขเวทนิยํ น ทุกฺขเวทนิยํ ฯเปฯ อนารมฺมณํ, นตฺถิ ตสฺส อาวฏฺฏนา ฯเปฯ ปณิธีติ, อามนฺตา.\csromanspacer{} ผสฺโส วิปาโก ผสฺโส น สุขเวทนิโย น ทุกฺขเวทนิโย ฯเปฯ อนารมฺมโณ, นตฺถิ ตสฺส อาวฏฺฏนา ฯเปฯ ปณิธีติ, น เหวํ วตฺตพฺเพ ฯเปฯ.}

\chapterheadpageread{7. สตฺตมวคฺค}
\proseitem{496}{\csromanfolio{265}\csromansymbolmark{*}อกุสลานํ ธมฺมานํ ชรามรณํ, อกุสลานํ ธมฺมานํ วิปาโกติ, อามนฺตา.\csromanspacer{} กุสลานํ ธมฺมานํ ชรามรณํ, กุสลานํ ธมฺมานํ วิปาโกติ, น เหวํ วตฺตพฺเพ ฯเปฯ.}

\prose{\csromansymbolmark{*}กุสลานํ ธมฺมานํ ชรามรณํ, น วตฺตพฺพํ “กุสลานํ ธมฺมานํ วิปาโก”ติ, อามนฺตา.\csromanspacer{} อกุสลานํ ธมฺมานํ ชรามรณํ, น วตฺตพฺพํ “อกุสลานํ ธมฺมานํ วิปาโก”ติ, น เหวํ วตฺตพฺเพ ฯเปฯ.}

\prose{\csromansymbolmark{*}กุสลานํ ธมฺมานํ ชรามรณํ, อกุสลานํ ธมฺมานํ วิปาโกติ, อามนฺตา.\csromanspacer{} อกุสลานํ ธมฺมานํ ชรามรณํ, กุสลานํ ธมฺมานํ วิปาโกติ, น เหวํ วตฺตพฺเพ ฯเปฯ.}

\prose{\csromansymbolmark{*}อกุสลานํ ธมฺมานํ ชรามรณํ, น วตฺตพฺพํ “กุสลานํ ธมฺมานํ วิปาโก”ติ, อามนฺตา.\csromanspacer{} กุสลานํ ธมฺมานํ ชรามรณํ, น วตฺตพฺพํ “อกุสลานํ ธมฺมานํ วิปาโก”ติ, น เหวํ วตฺตพฺเพ ฯเปฯ.}

\prose{\csromansymbolmark{*}กุสลานญฺจ อกุสลานญฺจ ธมฺมานํ ชรามรณํ, อกุสลานํ ธมฺมานํ วิปาโกติ, อามนฺตา.\csromanspacer{} กุสลานญฺจ อกุสลานญฺจ ธมฺมานํ ชรามรณํ, กุสลานํ ธมฺมานํ วิปาโกติ, น เหวํ วตฺตพฺเพ ฯเปฯ.}

\prose{\csromansymbolmark{*}กุสลานญฺจ อกุสลานญฺจ ธมฺมานํ ชรามรณํ, น วตฺตพฺพํ “กุสลานํ ธมฺมานํ วิปาโก”ติ, อามนฺตา.\csromanspacer{} กุสลานญฺจ อกุสลานญฺจ ธมฺมานํ ชรามรณํ, น วตฺตพฺพํ “อกุสลานํ ธมฺมานํ วิปาโก”ติ, น เหวํ วตฺตพฺเพ ฯเปฯ.}

\proseitem{497}{\csromansymbolmark{+}น วตฺตพฺพํ “ชรามรณํ วิปาโก”ติ, อามนฺตา.\csromanspacer{} นนุ อตฺถิ ทุพฺพณฺณสํวตฺตนิยํ กมฺมํ อปฺปายุกสํวตฺตนิยํ กมฺมนฺติ, อามนฺตา.\csromanspacer{} หญฺจิ อตฺถิ ทุพฺพณฺณสํวตฺตนิยํ กมฺมํ อปฺปายุกสํวตฺตนิยํ กมฺมํ, เตน วต เร วตฺตพฺเพ “ชรามรณํ วิปาโก”ติ.}

\nitthitam{ชรามรณํ วิปาโกติกถา นิฏฺฐิตา.}
\chapterheadpageread{7. สตฺตมวคฺค}
\vspace{\tipitakaparskip}
\csromancenter{\textbf{(71) 9.}\csromanspacer{}\textbf{ อริยธมฺมวิปากกถา}}
\csromanmatikaanchor{mtk.96}
\csromantocmarkref{subsection}{(71) 9. อริยธมฺมวิปากกถา}{mtk.96}
\bookmark[dest={mtk.96},level=2]{(71) 9. อริยธมฺมวิปากกถา}
\proseitem{498}{\csromanfolio{266}\csromansymbolmark{*}นตฺถิ อริยธมฺมวิปาโกติ, อามนฺตา.\csromanspacer{} นนุ มหปฺผลํ สามญฺญํ มหปฺผลํ พฺรหฺมญฺญนฺติ, อามนฺตา.\csromanspacer{} หญฺจิ มหปฺผลํ สามญฺญํ มหปฺผลํ พฺรหฺมญฺญํ, โน จ วต เร วตฺตพฺเพ “นตฺถิ อริยธมฺมวิปาโก”ติ.}

\prose{\csromansymbolmark{*}นตฺถิ อริยธมฺมวิปาโกติ, อามนฺตา.\csromanspacer{} นนุ อตฺถิ โสตาปตฺติผลนฺติ, อามนฺตา.\csromanspacer{} หญฺจิ อตฺถิ โสตาปตฺติผลํ, โน จ วต เร วตฺตพฺเพ “นตฺถิ อริยธมฺมวิปาโก”ติ, นนุ อตฺถิ สกทาคามิผลํ ฯเปฯ อนาคามิผลํ ฯเปฯ อรหตฺตผลนฺติ, อามนฺตา.\csromanspacer{} หญฺจิ อตฺถิ อรหตฺตผลํ, โน จ วต เร วตฺตพฺเพ “นตฺถิ อริยธมฺมวิปาโก”ติ.}

\prose{\csromansymbolmark{*}โสตาปตฺติผลํ น วิปาโกติ, อามนฺตา.\csromanspacer{} ทานผลํ น วิปาโกติ, น เหวํ วตฺตพฺเพ ฯเปฯ.\csromanspacer{} โสตาปตฺติผลํ น วิปาโกติ, อามนฺตา.\csromanspacer{} สีลผลํ ฯเปฯ ภาวนาผลํ น วิปาโกติ.\csromanspacer{} น เหวํ วตฺตพฺเพ ฯเปฯ.}

\prose{\csromansymbolmark{*}สกทาคามิผลํ ฯเปฯ.\csromanspacer{} อนาคามิผลํ ฯเปฯ.\csromanspacer{} อรหตฺตผลํ, น วิปาโกติ, อามนฺตา.\csromanspacer{} ทานผลํ น วิปาโกติ.\csromanspacer{} น เหวํ วตฺตพฺเพ ฯเปฯ.\csromanspacer{} อรหตฺตผลํ น วิปาโกติ, อามนฺตา.\csromanspacer{} สีลผลํ ฯเปฯ ภวนาผลํ น วิปาโกติ, น เหวํ วตฺตพฺเพ ฯเปฯ.\csromanspacer{} ทานผลํ วิปาโกติ, อามนฺตา.\csromanspacer{} โสตาปตฺติผลํ วิปาโกติ, น เหวํ วตฺตพฺเพ ฯเปฯ.}

\prose{\csromansymbolmark{*}ทานผลํ วิปาโกติ, อามนฺตา.\csromanspacer{} สกทาคามิผลํ ฯเปฯ.\csromanspacer{} อนาคามิผลํ ฯเปฯ.\csromanspacer{} อรหตฺตผลํ วิปาโกติ, น เหวํ วตฺตพฺเพ ฯเปฯ.\csromanspacer{} สีลผลํ ฯเปฯ.\csromanspacer{} ภาวนาผลํ วิปาโกติ, อามนฺตา.\csromanspacer{} โสตาปตฺติผลํ วิปาโกติ, น เหวํ วตฺตพฺเพ ฯเปฯ.\csromanspacer{} ภาวนาผลํ วิปาโกติ, อามนฺตา.\csromanspacer{} สกทาคามิผลํ ฯเปฯ.\csromanspacer{} อนาคามิผลํ ฯเปฯ.\csromanspacer{} อรหตฺตผลํ วิปาโกติ, น เหวํ วตฺตพฺเพ ฯเปฯ.}

\proseitem{499}{\csromansymbolmark{*}กามาวจรํ กุสลํ สวิปากนฺติ, อามนฺตา.\csromanspacer{} โลกุตฺตรํ กุสลํ สวิปากนฺติ, น เหวํ วตฺตพฺเพ ฯเปฯ.\csromanspacer{} รูปาวจรํ กุสลํ ฯเปฯ.\csromanspacer{} อรูปาวจรํ กุสลํ สวิปากนฺติ, อามนฺตา.\csromanspacer{} โลกุตฺตรํ กุสลํ สวิปากนฺติ, น เหวํ วตฺตพฺเพ ฯเปฯ.}

\chapterheadpageread{7. สตฺตมวคฺค}
\csromanmatikaanchor{mtk.97}
\csromantocmarkref{subsection}{(72) 10. วิปาโกวิปากธมฺมธมฺโมติกถา}{mtk.97}
\bookmark[dest={mtk.97},level=2]{(72) 10. วิปาโกวิปากธมฺมธมฺโมติกถา}
\prose{\csromanfolio{267}\csromansymbolmark{*}โลกุตฺตรํ กุสลํ อวิปากนฺติ, อามนฺตา.\csromanspacer{} กามาวจรํ กุสลํ อวิปากนฺติ, น เหวํ วตฺตพฺเพ ฯเปฯ.\csromanspacer{} โลกุตฺตรํ กุสลํ อวิปากนฺติ, อามนฺตา.\csromanspacer{} รูปาวจรํ ฯเปฯ.\csromanspacer{} อรูปาวจรํ กุสลํ อวิปากนฺติ, น เหวํ วตฺตพฺเพ ฯเปฯ.}

\proseitem{500}{\csromansymbolmark{+}กามาวจรํ กุสลํ สวิปากํ อาจยคามีติ, อามนฺตา.\csromanspacer{} โลกุตฺตรํ กุสลํ สวิปากํ อาจยคามีติ, น เหวํ วตฺตพฺเพ ฯเปฯ.\csromanspacer{} รูปาวจรํ ฯเปฯ.\csromanspacer{} อรูปาวจรํ กุสลํ สวิปากํ อาจยคามีติ, อามนฺตา.\csromanspacer{} โลกุตฺตรํ กุสลํ สวิปากํ อาจยคามีติ, น เหวํ วตฺตพฺเพ ฯเปฯ.}

\prose{\csromansymbolmark{+}โลกุตฺตรํ กุสลํ สวิปากํ อปจยคามีติ, อามนฺตา.\csromanspacer{} กามาวจรํ กุสลํ สวิปากํ อปจยคามีติ, น เหวํ วตฺตพฺเพ ฯเปฯ.\csromanspacer{} โลกุตฺตรํ กุสลํ สวิปากํ อปจยคามีติ, อามนฺตา.\csromanspacer{} รูปาวจรํ ฯเปฯ.\csromanspacer{} อรูปาวจรํ กุสลํ สวิปากํ อปจยคามีติ, น เหวํ วตฺตพฺเพ ฯเปฯ.}

\nitthitam{อริยธมฺมวิปากกถา นิฏฺฐิตา.}
\chapterhead{7. สตฺตมวคฺค}
\vspace{\tipitakaparskip}
\csromancenter{\textbf{(72) 10.}\csromanspacer{}\textbf{ วิปาโกวิปากธมฺมธมฺโมติกถา}}
\proseitem{501}{\csromansymbolmark{*}วิปาโก วิปากธมฺมธมฺโมติ, อามนฺตา.\csromanspacer{} ตสฺส วิปาโก วิปากธมฺมธมฺโมติ, น เหวํ วตฺตพฺเพ ฯเปฯ.\csromanspacer{} ตสฺส วิปาโก วิปากธมฺมธมฺโมติ, อามนฺตา.\csromanspacer{} ตสฺส ตสฺเสว นตฺถิ ทุกฺขสฺส อนฺตกิริยา นตฺถิ วฏฺฏุปจฺเฉโท นตฺถิ อนุปาทาปรินิพฺพานนฺติ, น เหวํ วตฺตพฺเพ ฯเปฯ.}

\prose{\csromansymbolmark{*}วิปาโก วิปากธมฺมธมฺโมติ, อามนฺตา.\csromanspacer{} วิปาโกติ วา วิปากธมฺมธมฺโมติ วา, วิปากธมฺมธมฺโมติ วา วิปาโกติ วา เอเสเส เอกฏฺเฐ สเม สมภาเค ตชฺชาเตติ, น เหวํ วตฺตพฺเพ ฯเปฯ.}

\prose{\csromansymbolmark{*}วิปาโก วิปากธมฺมธมฺโมติ, อามนฺตา.\csromanspacer{} วิปาโก จ วิปากธมฺมธมฺโม จ, วิปากธมฺมธมฺโม จ วิปาโก จ สหคตา สหชาตา สํสฏฺฐา สมฺปยุตฺตา เอกุปฺปาทา เอกนิโรธา เอกวตฺถุกา เอการมฺมณาติ, น เหวํ วตฺตพฺเพ ฯเปฯ.}

\csromanmatikaanchor{mtk.98}
\csromanmatikaanchor{mtk.99}
\csromantocmarknopage{chapter}{8. อฏฺฐมวคฺค}{mtk.98}
\bookmark[dest={mtk.98},level=0]{8. อฏฺฐมวคฺค}
\csromantocmarkref{subsection}{(73) 1. ฉคติกถา}{mtk.99}
\bookmark[dest={mtk.99},level=2]{(73) 1. ฉคติกถา}
\prose{\csromanfolio{268}\csromansymbolmark{*}วิปาโก วิปากธมฺมธมฺโมติ, อามนฺตา.\csromanspacer{} ตญฺเญว อกุสลํ โส อกุสลสฺส วิปาโก, ตญฺเญว กุสลํ โส กุสลสฺส วิปาโกติ, น เหวํ วตฺตพฺเพ ฯเปฯ.}

\prose{\csromansymbolmark{*}วิปาโก วิปากธมฺมธมฺโมติ, อามนฺตา.\csromanspacer{} เยเนว จิตฺเตน ปาณํ หนติ, เตเนว จิตฺเตน นิรเย ปจฺจติ, เยเนว จิตฺเตน ทานํ เทติ, เตเนว จิตฺเตน สคฺเค โมทตีติ, น เหวํ วตฺตพฺเพ ฯเปฯ.}

\proseitem{502}{\csromansymbolmark{+}น วตฺตพฺพํ “วิปาโก วิปากธมฺมธมฺโม”ติ, อามนฺตา.\csromanspacer{} นนุ วิปากา จตฺตาโร ขนฺธา อรูปิโน อญฺญมญฺญปจฺจยาติ, อามนฺตา.\csromanspacer{} หญฺจิ วิปากา จตฺตาโร ขนฺธา อรูปิโน อญฺญมญฺญปจฺจยา, เตน วต เร วตฺตพฺเพ “วิปาโก วิปากธมฺมธมฺโม”ติ.\csromanspacer{} วิปาโก วิปากธมฺมธมฺโมติกถา นิฏฺฐิตา.\csromanspacer{} สตฺตโม วคฺโค.}

\titlehead{\textbf{ตสฺสุทฺทานํ}}
\prose{สงฺคโห, สมฺปยุตฺโต, เจตสิโก ธมฺโม, เจตสิกํ ทานํ ปริโภคมยํ ปุญฺญํ วฑฺฒติ, อิโต ทินฺเนน ตตฺถ ยาเปนฺติ, ปถวี กมฺมวิปาโก, ชรามรณํ วิปาโก, นตฺถิ อริยธมฺมวิปาโก, วิปาโก วิปากธมฺมธมฺโมติ.}

\chapterhead{8. อฏฺฐมวคฺค}
\vspace{\tipitakaparskip}
\csromancenter{\textbf{(73) 1.}\csromanspacer{}\textbf{ ฉคติกถา}}
\proseitem{503}{\csromansymbolmark{*}ฉ คติโยติ, อามนฺตา.\csromanspacer{} นนุ ปญฺจ คติโย วุตฺตา ภควตา “นิรโย ติรจฺฉานโยนิ เปตฺติวิสโย มนุสฺสา เทวา”ติ\footnote{ม 1. 106 ปิฏฺเฐ.}, อามนฺตา.\csromanspacer{} หญฺจิ ปญฺจ คติโย วุตฺตา ภควตา นิรโย ติรจฺฉานโยนิ เปตฺติวิสโย มนุสฺสา เทวา, โน จ วต เร วตฺตพฺเพ “ฉ คติโย”ติ.}

\chapterheadpageread{8. อฏฺฐมวคฺค}
\csromanmatikaanchor{mtk.100}
\csromantocmarkref{subsection}{(74) 2. อนฺตราภวกถา}{mtk.100}
\bookmark[dest={mtk.100},level=2]{(74) 2. อนฺตราภวกถา}
\prose{\csromanfolio{269}\csromansymbolmark{*}ฉ คติโยติ, อามนฺตา.\csromanspacer{} นนุ กาลกญฺจิกา\footnote{กาลกญฺชิกา (สี, สฺยา, กํ), กาฬกญฺชกา (อิ)} อสุรา เปตานํ สมานวณฺณา สมานโภคา สมานาหารา สมานายุกา เปเตหิ สห อาวาหวิวาหํ คจฺฉนฺตีติ, อามนฺตา.\csromanspacer{} หญฺจิ กาลกญฺจิกา อสุรา เปตานํ สมานวณฺณา สมานโภคา สมานาหารา สมานายุกา เปเตหิ สห อาวาหวิวาหํ คจฺฉนฺติ, โน จ วต เร วตฺตพฺเพ “ฉ คติโย”ติ.}

\prose{\csromansymbolmark{*}ฉ คติโยติ, อามนฺตา.\csromanspacer{} นนุ เวปจิตฺติปริสา เทวานํ สมานวณฺณา สมานโภคา สมานาหารา สมานายุกา เทเวหิ สห อาวาหวิวาหํ คจฺฉนฺตีติ, อามนฺตา.\csromanspacer{} หญฺจิ เวปจิตฺติปริสา เทวานํ สมานวณฺณา สมานโภคา สมานาหารา สมานายุกา เทเวหิ สห อาวาหวิวาหํ คจฺฉนฺติ, โน จ วต เร วตฺตพฺเพ “ฉ คติโย”ติ.}

\prose{\csromansymbolmark{*}ฉ คติโยติ, อามนฺตา.\csromanspacer{} นนุ เวปจิตฺติปริสา ปุพฺพเทวาติ, อามนฺตา.\csromanspacer{} หญฺจิ เวปจิตฺติปริสา ปุพฺพเทวา, โน จ วต เร วตฺตพฺเพ “ฉ คติโย”ติ.}

\proseitem{504}{\csromansymbolmark{+}น วตฺตพฺพํ “ฉ คติโย”ติ, อามนฺตา.\csromanspacer{} นนุ อตฺถิ อสุรกาโยติ, อามนฺตา.\csromanspacer{} หญฺจิ อตฺถิ อสุรกาโย, เตน วต เร วตฺตพฺเพ “ฉ คติโย”ติ.}

\nitthitam{ฉคติกถา นิฏฺฐิตา.}
\chapterhead{8. อฏฺฐมวคฺค}
\vspace{\tipitakaparskip}
\csromancenter{\textbf{(74) 2.}\csromanspacer{}\textbf{ อนฺตราภวกถา}}
\proseitem{505}{\csromansymbolmark{*}อตฺถิ อนฺตราภโวติ, อามนฺตา.\csromanspacer{} กามภโวติ, น เหวํ วตฺตพฺเพ ฯเปฯ.\csromanspacer{} อตฺถิ อนฺตราภโวติ, อามนฺตา.\csromanspacer{} รูปภโวติ, น เหวํ วตฺตพฺเพ ฯเปฯ.\csromanspacer{} อตฺถิ อนฺตราภโวติ, อามนฺตา.\csromanspacer{} อรูปภโวติ, น เหวํ วตฺตพฺเพ ฯเปฯ.\csromanspacer{} อตฺถิ อนฺตราภโวติ, อามนฺตา.\csromanspacer{} กามภวสฺส จ รูปภวสฺส จ อนฺตเร อตฺถิ อนฺตราภโวติ, น เหวํ วตฺตพฺเพ ฯเปฯ.\csromanspacer{} อตฺถิ อนฺตราภโวติ, อามนฺตา.\csromanspacer{} รูปภวสฺส จ อรูปภวสฺส จ อนฺตเร อตฺถิ อนฺตราภโวติ, น เหวํ วตฺตพฺเพ ฯเปฯ.}

\prose{\csromanfolio{270}\csromansymbolmark{*}กามภวสฺส จ รูปภวสฺส จ อนฺตเร นตฺถิ อนฺตราภโวติ, อามนฺตา.\csromanspacer{} หญฺจิ กามภวสฺส จ รูปภวสฺส จ อนฺตเร นตฺถิ อนฺตราภโว, โน จ วต เร วตฺตพฺเพ “อตฺถิ อนฺตราภโว”ติ.\csromanspacer{} รูปภวสฺส จ อรูปภวสฺส จ อนฺตเร นตฺถิ อนฺตราภโวติ, อามนฺตา.\csromanspacer{} หญฺจิ รูปภวสฺส จ อรูปภวสฺส จ อนฺตเร นตฺถิ อนฺตราภโว, โน จ วต เร วตฺตพฺเพ “อตฺถิ อนฺตราภโว”ติ.}

\proseitem{506}{\csromansymbolmark{*}อตฺถิ อนฺตราภโวติ, อามนฺตา.\csromanspacer{} ปญฺจมี สา โยนิ, ฉฏฺฐมี สา คติ, อฏฺฐมี สา วิญฺญาณฏฺฐิติ, ทสโม โส สตฺตาวาโสติ, น เหวํ วตฺตพฺเพ ฯเปฯ.\csromanspacer{} อตฺถิ อนฺตราภโวติ, อามนฺตา.\csromanspacer{} อนฺตราภโว ภโว คติ สตฺตาวาโส สํสาโร โยนิ วิญฺญาณฏฺฐิติ อตฺตภาวปฏิลาโภติ, น เหวํ วตฺตพฺเพ ฯเปฯ.\csromanspacer{} อตฺถิ อนฺตราภวูปคํ กมฺมนฺติ, น เหวํ วตฺตพฺเพ ฯเปฯ.\csromanspacer{} อตฺถิ อนฺตราภวูปคา สตฺตาติ, น เหวํ วตฺตพฺเพ ฯเปฯ.\csromanspacer{} อนฺตราภเว สตฺตา ชายนฺติ ชียนฺติ มียนฺติ จวนฺติ อุปปชฺชนฺตีติ, น เหวํ วตฺตพฺเพ ฯเปฯ.\csromanspacer{} อนฺตราภเว อตฺถิ รูปํ เวทนา สญฺญา สงฺขารา วิญฺญาณนฺติ, น เหวํ วตฺตพฺเพ ฯเปฯ.\csromanspacer{} อนฺตราภโว ปญฺจโวการภโวติ, น เหวํ วตฺตพฺเพ ฯเปฯ.}

\proseitem{507}{\csromansymbolmark{*}อตฺถิ กามภโว, กามภโว ภโว คติ สตฺตาวาโส สํสาโร โยนิ วิญฺญาณฏฺฐิติ อตฺตภาวปฏิลาโภติ, อามนฺตา.\csromanspacer{} อตฺถิ อนฺตราภโว, อนฺตราภโว ภโว คติ สตฺตาวาโส สํสาโร โยนิ วิญฺญาณฏฺฐิติ อตฺตภาวปฏิลาโภติ, น เหวํ วตฺตพฺเพ ฯเปฯ.\csromanspacer{} อตฺถิ กามภวูปคํ กมฺมนฺติ, อามนฺตา.\csromanspacer{} อตฺถิ อนฺตราภวูปคํ กมฺมนฺติ, น เหวํ วตฺตพฺเพ ฯเปฯ.\csromanspacer{} อตฺถิ กามภวูปคา สตฺตาติ, อามนฺตา.\csromanspacer{} อตฺถิ อนฺตราภวูปคา สตฺตาติ, น เหวํ วตฺตพฺเพ ฯเปฯ.\csromanspacer{} กามภเว สตฺตา ชายนฺติ ชียนฺติ มียนฺติ จวนฺติ อุปปชฺชนฺตีติ, อามนฺตา.\csromanspacer{} อนฺตราภเว สตฺตา ชายนฺติ ชียนฺติ มียนฺติ จวนฺติ อุปปชฺชนฺตีติ, น เหวํ วตฺตพฺเพ ฯเปฯ.\csromanspacer{} กามภเว อตฺถิ รูปํ เวทนา สญฺญา สงฺขารา วิญฺญาณนฺติ, อามนฺตา.\csromanspacer{} อนฺตราภเว อตฺถิ รูปํ เวทนา สญฺญา สงฺขารา วิญฺญาณนฺติ, น เหวํ วตฺตพฺเพ ฯเปฯ.\csromanspacer{} กามภโว ปญฺจโวการภโวติ, อามนฺตา.\csromanspacer{} อนฺตราภโว ปญฺจโวการภโวติ, น เหวํ วตฺตพฺเพ ฯเปฯ.}

\prose{\csromansymbolmark{*}อตฺถิ รูปภโว, รูปภโว ภโว คติ สตฺตาวาโส สํสาโร โยนิ วิญฺญาณฏฺฐิติ อตฺตภาวปฏิลาโภติ, อามนฺตา.\csromanspacer{} อตฺถิ อนฺตราภโว, อนฺตราภโว ภโว คติ สตฺตาวาโส สํสาโร โยนิ}

\vspace{\tipitakaparskip}
\csromancenter{อฏฺฐมวคฺค วิญฺญาณฏฺฐิติ อตฺตภาวปฏิลาโภติ, น เหวํ วตฺตพฺเพ ฯเปฯ.\csromanspacer{} อตฺถิ รูปภวูปคํ กมฺมนฺติ, อามนฺตา.\csromanspacer{} อตฺถิ อนฺตราภวูปคํ กมฺมนฺติ, น เหวํ วตฺตพฺเพ ฯเปฯ.\csromanspacer{} อตฺถิ รูปภวูปคา สตฺตาติ, อามนฺตา.\csromanspacer{} อตฺถิ อนฺตราภวูปคา สตฺตาติ, น เหวํ วตฺตพฺเพ ฯเปฯ.\csromanspacer{} รูปภเว สตฺตา ชายนฺติ ชียนฺติ มียนฺติ จวนฺติ อุปปชฺชนฺตีติ, อามนฺตา.\csromanspacer{} อนฺตราภเว สตฺตา ชายนฺติ ชียนฺติ มียนฺติ จวนฺติ อุปปชฺชนฺตีติ, น เหวํ วตฺตพฺเพ ฯเปฯ.\csromanspacer{} รูปภเว อตฺถิ รูปํ เวทนา สญฺญา สงฺขารา วิญฺญาณนฺติ, อามนฺตา.\csromanspacer{} อนฺตราภเว อตฺถิ รูปํ เวทนา สญฺญา สงฺขารา วิญฺญาณนฺติ, น เหวํ วตฺตพฺเพ ฯเปฯ.\csromanspacer{} รูปภโว ปญฺจโวการภโวติ, อามนฺตา.\csromanspacer{} อนฺตราภโว ปญฺจโวการภโวติ, น เหวํ วตฺตพฺเพ ฯเปฯ.}
\prose{\csromanfolio{271}\csromansymbolmark{*}อตฺถิ อรูปภโว, อรูปภโว ภโว คติ สตฺตาวาโส สํสาโร โยนิ วิญฺญาณฏฺฐิติ อตฺตภาวปฏิลาโภติ, อามนฺตา.\csromanspacer{} อตฺถิ อนฺตราภโว, อนฺตราภโว ภโว คติ สตฺตาวาโส สํสาโร โยนิ วิญฺญาณฏฺฐิติ อตฺตภาวปฏิลาโภติ, น เหวํ วตฺตพฺเพ ฯเปฯ.\csromanspacer{} อตฺถิ อรูปภวูปคํ กมฺมนฺติ, อามนฺตา.\csromanspacer{} อตฺถิ อนฺตราภวูปคํ กมฺมนฺติ, น เหวํ วตฺตพฺเพ ฯเปฯ.\csromanspacer{} อตฺถิ อรูปภวูปคา สตฺตาติ, อามนฺตา.\csromanspacer{} อตฺถิ อนฺตราภวูปคา สตฺตาติ, น เหวํ วตฺตพฺเพ ฯเปฯ.\csromanspacer{} อรูปภเว สตฺตา ชายนฺติ ชียนฺติ มียนฺติ จวนฺติ อุปปชฺชนฺตีติ, อามนฺตา.\csromanspacer{} อนฺตราภเว สตฺตา ชายนฺติ ชียนฺติ มียนฺติ จวนฺติ อุปปชฺชนฺตีติ, น เหวํ วตฺตพฺเพ ฯเปฯ.\csromanspacer{} อรูปภเว อตฺถิ เวทนา สญฺญา สงฺขารา วิญฺญาณนฺติ, อามนฺตา.\csromanspacer{} อนฺตราภเว อตฺถิ เวทนา สญฺญา สงฺขารา วิญฺญาณนฺติ, น เหวํ วตฺตพฺเพ ฯเปฯ.\csromanspacer{} อรูปภโว จตุโวการภโวติ, อามนฺตา.\csromanspacer{} อนฺตราภโว จตุโวการภโวติ, น เหวํ วตฺตพฺเพ ฯเปฯ.}

\proseitem{508}{\csromansymbolmark{*}อตฺถิ อนฺตราภโวติ, อามนฺตา.\csromanspacer{} สพฺเพสญฺเญว สตฺตานํ อตฺถิ อนฺตราภโวติ, น เหวํ วตฺตพฺเพ ฯเปฯ.\csromanspacer{} สพฺเพสญฺเญว สตฺตานํ นตฺถิ อนฺตราภโวติ, อามนฺตา.\csromanspacer{} หญฺจิ สพฺเพสญฺเญว สตฺตานํ นตฺถิ อนฺตราภโว, โน จ วต เร วตฺตพฺเพ “อตฺถิ อนฺตราภโว”ติ.}

\prose{\csromansymbolmark{*}อตฺถิ อนฺตราภโวติ, อามนฺตา.\csromanspacer{} อานนฺตริยสฺส ปุคฺคลสฺส อตฺถิ อนฺตราภโวติ, น เหวํ วตฺตพฺเพ ฯเปฯ.\csromanspacer{} อานนฺตริยสฺส ปุคฺคลสฺส นตฺถิ อนฺตราภโวติ, อามนฺตา.\csromanspacer{} หญฺจิ อานนฺตริยสฺส ปุคฺคลสฺส นตฺถิ อนฺตราภโว, โน จ วต ร วตฺตพฺเพ “อตฺถิ อนฺตราภโว”ติ.}

\csromanmatikaanchor{mtk.101}
\csromantocmarkref{subsection}{(75) 3. กามคุณกถา}{mtk.101}
\bookmark[dest={mtk.101},level=2]{(75) 3. กามคุณกถา}
\prose{\csromanfolio{272}\csromansymbolmark{*}น อานนฺตริยสฺส ปุคฺคลสฺส อตฺถิ อนฺตราภโวติ, อามนฺตา.\csromanspacer{} อานนฺตริยสฺส ปุคฺคลสฺส อตฺถิ อนฺตราภโวติ, น เหวํ วตฺตพฺเพ ฯเปฯ.\csromanspacer{} อานนฺตริยสฺส ปุคฺคลสฺส นตฺถิ อนฺตราภโวติ, อามนฺตา.\csromanspacer{} น อานนฺตริยสฺส ปุคฺคลสฺส นตฺถิ อนฺตราภโวติ, น เหวํ วตฺตพฺเพ ฯเปฯ.\csromanspacer{} นิรยูปคสฺส ปุคฺคลสฺส ฯเปฯ.\csromanspacer{} อสญฺญสตฺตูปคสฺส ปุคฺคลสฺส ฯเปฯ.\csromanspacer{} อรูปูปคสฺส ปุคฺคลสฺส อตฺถิ อนฺตราภโวติ, น เหวํ วตฺตพฺเพ ฯเปฯ.\csromanspacer{} อรูปูปคสฺส ปุคฺคลสฺส นตฺถิ อนฺตราภโวติ, อามนฺตา.\csromanspacer{} หญฺจิ อรูปูปคสฺส ปุคฺคลสฺส นตฺถิ อนฺตราภโว, โน จ วต เร วตฺตพฺเพ “อตฺถิ อนฺตราภโว”ติ.}

\prose{\csromansymbolmark{*}น อรูปูปคสฺส ปุคฺคลสฺส อตฺถิ อนฺตราภโวติ, อามนฺตา.\csromanspacer{} อรูปูปคสฺส ปุคฺคลสฺส อตฺถิ อนฺตราภโวติ, น เหวํ วตฺตพฺเพ ฯเปฯ.\csromanspacer{} อรูปูปคสฺส ปุคฺคลสฺส นตฺถิ อนฺตราภโวติ, อามนฺตา.\csromanspacer{} น อรูปูปคสฺส ปุคฺคลสฺส นตฺถิ อนฺตราภโวติ, น เหวํ วตฺตพฺเพ ฯเปฯ.}

\proseitem{509}{\csromansymbolmark{+}น วตฺตพฺพํ “อตฺถิ อนฺตราภโว”ติ, อามนฺตา.\csromanspacer{} นนุ อนฺตราปรินิพฺพายี ปุคฺคโล อตฺถีติ, อามนฺตา.\csromanspacer{} หญฺจิ อนฺตราปรินิพฺพายี ปุคฺคโล อตฺถิ, เตน วต เร\footnote{โน จ วต เร (สี, ก), โน วต เร (สฺยา)} วตฺตพฺเพ “อตฺถิ อนฺตราภโว”ติ.}

\prose{\csromansymbolmark{*}อนฺตราปรินิพฺพายี ปุคฺคโล อตฺถีติ กตฺวา อตฺถิ อนฺตราภโวติ, อามนฺตา.\csromanspacer{} อุปหจฺจปรินิพฺพายี ปุคฺคโล อตฺถีติ กตฺวา อตฺถิ อุปหจฺจภโวติ, น เหวํ วตฺตพฺเพ ฯเปฯ.\csromanspacer{} อนฺตราปรินิพฺพายี ปุคฺคโล อตฺถีติ กตฺวา อตฺถิ อนฺตราภโวติ, อามนฺตา.\csromanspacer{} อสงฺขารปรินิพฺพายี ปุคฺคโล ฯเปฯ.\csromanspacer{} สสงฺขารปรินิพฺพายี ปุคฺคโล อตฺถีติ กตฺวา อตฺถิ สสงฺขารภโวติ, น เหวํ วตฺตพฺเพ ฯเปฯ.}

\nitthitam{อนฺตราภวกถา นิฏฺฐิตา.}
\chapterhead{8. อฏฺฐมวคฺค}
\vspace{\tipitakaparskip}
\csromancenter{\textbf{(75) 3.}\csromanspacer{}\textbf{ กามคุณกถา}}
\proseitem{510}{\csromansymbolmark{*}ปญฺเจว กามคุณา กามธาตูติ, อามนฺตา.\csromanspacer{} นนุ อตฺถิ ตปฺปฏิสญฺญุตฺโต ฉนฺโทติ, อามนฺตา.\csromanspacer{} หญฺจิ อตฺถิ ตปฺปฏิสญฺญุตฺโต ฉนฺโท, โน}

\vspace{\tipitakaparskip}
\csromancenter{อฏฺฐมวคฺค จ วต เร วตฺตพฺเพ “ปญฺเจว กามคุณา กามธาตู”ติ.\csromanspacer{} นนุ อตฺถิ ตปฺปฏิสญฺญุตฺโต ราโค ตปฺปฏิสญฺญุตฺโต ฉนฺโท ตปฺปฏิสญฺญุตฺโต ฉนฺทราโค ตปฺปฏิสญฺญุตฺโต สงฺกปฺโป ตปฺปฏิสญฺญุตฺโต ราโค ตปฺปฏิสญฺญุตฺโต สงฺกปฺปราโค ตปฺปฏิสญฺญุตฺตา ปีติ ตปฺปฏิสญฺญุตฺตํ โสมนสฺสํ ตปฺปฏิสญฺญุตฺตํ ปีติโสมนสฺสนฺติ, อามนฺตา.\csromanspacer{} หญฺจิ อตฺถิ ตปฺปฏิสญฺญุตฺตํ ปีติโสมนสฺสํ, โน จ วต เร วตฺตพฺเพ “ปญฺเจว กามคุณา กามธาตู”ติ.}
\prose{\csromanfolio{273}\csromansymbolmark{*}ปญฺเจว กามคุณา กามธาตูติ, อามนฺตา.\csromanspacer{} มนุสฺสานํ จกฺขุ น กามธาตูติ, น เหวํ วตฺตพฺเพ ฯเปฯ.\csromanspacer{} มนุสฺสานํ โสตํ ฯเปฯ.\csromanspacer{} มนุสฺสานํ ฆานํ ฯเปฯ.\csromanspacer{} มนุสฺสานํ ชิวฺหา ฯเปฯ.\csromanspacer{} มนุสฺสานํ กาโย ฯเปฯ.\csromanspacer{} มนุสฺสานํ มโน น กามธาตูติ, น เหวํ วตฺตพฺเพ ฯเปฯ. \csromansymbolmark{*}มนุสฺสานํ มโน น กามธาตูติ, อามนฺตา.\csromanspacer{} นนุ วุตฺตํ ภควตา\csromandash{}}

\csromangathameasure{ปญฺจ กามคุณา โลเก, มโนจฺฉฏฺฐา ปเวทิตา. \\ เอตฺถ ฉนฺทํ วิราเชตฺวา, เอวํ ทุกฺขา ปมุจฺจตีติ. อตฺเถว สุตฺตนฺโตติ, อามนฺตา. เตน หิ น วตฺตพฺพํ “มนุสฺสานํ มโน น กามธาตู”ติ.}
\csromangathagroup{\csromangathastanza{ปญฺจ กามคุณา โลเก, มโนจฺฉฏฺฐา ปเวทิตา. \\ เอตฺถ ฉนฺทํ วิราเชตฺวา, เอวํ ทุกฺขา ปมุจฺจตีติ\footnote{ขุ 1. 304; สํ 1. 16 ปิฏฺเฐสุ.}.\csromanspacer{} อตฺเถว สุตฺตนฺโตติ, อามนฺตา.\csromanspacer{} เตน หิ น วตฺตพฺพํ “มนุสฺสานํ มโน น กามธาตู”ติ.}}

\proseitem{511}{\csromansymbolmark{*}ปญฺเจว กามคุณา กามธาตูติ, อามนฺตา.\csromanspacer{} กามคุณา ภโว คติ สตฺตาวาโส สํสาโร โยนิ วิญฺญาณฏฺฐิติ อตฺตภาวปฏิลาโภติ, น เหวํ วตฺตพฺเพ ฯเปฯ.\csromanspacer{} อตฺถิ กามคุณูปคํ กมฺมนฺติ, น เหวํ วตฺตพฺเพ ฯเปฯ.\csromanspacer{} อตฺถิ กามคุณูปคา สตฺตาติ, น เหวํ วตฺตพฺเพ ฯเปฯ.\csromanspacer{} กามคุเณ สตฺตา ชายนฺติ ชียนฺติ มียนฺติ จวนฺติ อุปปชฺชนฺตีติ, น เหวํ วตฺตพฺเพ ฯเปฯ.\csromanspacer{} กามคุเณ อตฺถิ รูปํ เวทนา สญฺญา สงฺขารา วิญฺญาณนฺติ, น เหวํ วตฺตพฺเพ ฯเปฯ.\csromanspacer{} กามคุณา ปญฺจโวการภโวติ, น เหวํ วตฺตพฺเพ ฯเปฯ.\csromanspacer{} กามคุเณ สมฺมาสมฺพุทฺธา อุปฺปชฺชนฺติ, ปจฺเจกสมฺพุทฺธา อุปฺปชฺชนฺติ, สาวกยุคํ อุปฺปชฺชตีติ, น เหวํ วตฺตพฺเพ ฯเปฯ.}

\csromanmatikaanchor{mtk.102}
\csromantocmarkref{subsection}{(76) 4. กามกถา}{mtk.102}
\bookmark[dest={mtk.102},level=2]{(76) 4. กามกถา}
\prose{\csromansymbolmark{*}กามธาตุ ภโว คติ สตฺตาวาโส สํสาโร โยนิ วิญฺญาณฏฺฐิติ อตฺตภาวปฏิลาโภติ, อามนฺตา.\csromanspacer{} กามคุณา ภโว คติ สตฺตาวาโส สํสาโร โยนิ วิญฺญาณฏฺฐิติ อตฺตภาวปฏิลาโภติ, น เหวํ วตฺตพฺเพ ฯเปฯ.\csromanspacer{} อตฺถิ กามธาตูปคํ กมฺมนฺติ, อามนฺตา.\csromanspacer{} อตฺถิ กามคุณูปคํ \csromanfolio{274}กมฺมนฺติ, น เหวํ วตฺตพฺเพ ฯเปฯ.\csromanspacer{} อตฺถิ กามธาตูปคา สตฺตาติ, อามนฺตา.\csromanspacer{} อตฺถิ กามคุณูปคา สตฺตาติ, น เหวํ วตฺตพฺเพ ฯเปฯ.}

\prose{\csromansymbolmark{*}กามธาตุยา สตฺตา ชายนฺติ ชียนฺติ มียนฺติ จวนฺติ อุปปชฺชนฺตีติ, อามนฺตา.\csromanspacer{} กามคุเณ สตฺตา ชายนฺติ ชียนฺติ มียนฺติ จวนฺติ อุปปชฺชนฺตีติ, น เหวํ วตฺตพฺเพ ฯเปฯ.\csromanspacer{} กามธาตุยา อตฺถิ รูปํ เวทนา สญฺญา สงฺขารา วิญฺญาณนฺติ, อามนฺตา.\csromanspacer{} กามคุเณ อตฺถิ รูปํ เวทนา สญฺญา สงฺขารา วิญฺญาณนฺติ, น เหวํ วตฺตพฺเพ ฯเปฯ.\csromanspacer{} กามธาตุ ปญฺจโวการภโวติ, อามนฺตา.\csromanspacer{} กามคุณา ปญฺจโวการภโวติ, น เหวํ วตฺตพฺเพ ฯเปฯ.\csromanspacer{} กามธาตุยา สมฺมาสมฺพุทฺธา อุปฺปชฺชนฺติ, ปจฺเจกสมฺพุทฺธา อุปฺปชฺชนฺติ, สาวกยุคํ อุปฺปชฺชตีติ, อามนฺตา.\csromanspacer{} กามคุเณ สมฺมาสมฺพุทฺธา อุปฺปชฺชนฺติ, ปจฺเจกสมฺพุทฺธา อุปฺปชฺชนฺติ, สาวกยุคํ อุปฺปชฺชตีติ, น เหวํ วตฺตพฺเพ ฯเปฯ.}

\proseitem{512}{\csromansymbolmark{+}น วตฺตพฺพํ “ปญฺเจว กามคุณา กามธาตู”ติ, อามนฺตา.\csromanspacer{} นนุ วุตฺตํ ภควตา “ปญฺจิเม ภิกฺขเว กามคุณา.\csromanspacer{} กตเม ปญฺจ, จกฺขุวิญฺเญยฺยา รูปา อิฏฺฐา กนฺตา มนาปา ปิยรูปา กามูปสํหิตา รชนียา.\csromanspacer{} โสตวิญฺเญยฺยา สทฺทา ฯเปฯ.\csromanspacer{} ฆานวิญฺเญยฺยา คนฺธา ฯเปฯ.\csromanspacer{} ชิวฺหาวิญฺเญยฺยา รสา ฯเปฯ.\csromanspacer{} กายวิญฺเญยฺยา โผฏฺฐพฺพา อิฏฺฐา กนฺตา มนาปา ปิยรูปา กามูปสํหิตา รชนียา.\csromanspacer{} อิเม โข ภิกฺขเว ปญฺจ กามคุณา”ติ อตฺเถว สุตฺตนฺโตติ, อามนฺตา.\csromanspacer{} เตน หิ ปญฺเจว กามคุณา กามธาตูติ.}

\nitthitam{กามคุณกถา นิฏฺฐิตา.}
\chapterhead{8. อฏฺฐมวคฺค}
\vspace{\tipitakaparskip}
\csromancenter{\textbf{(76) 4.}\csromanspacer{}\textbf{ กามกถา}}
\proseitem{513}{\csromansymbolmark{*}ปญฺเจวายตนา กามาติ, อามนฺตา.\csromanspacer{} นนุ อตฺถิ ตปฺปฏิสํยุตฺโต ฉนฺโทติ, อามนฺตา.\csromanspacer{} หญฺจิ อตฺถิ ตปฺปฏิสํยุตฺโต ฉนฺโท, โน จ วต เร วตฺตพฺเพ “ปญฺเจวายตนา กามา”ติ, นนุ อตฺถิ ตปฺปฏิสํยุตฺโต ราโค ตปฺปฏิสํยุตฺโต ฉนฺโท ตปฺปฏิสํยุตฺโต ฉนฺทราโค ตปฺปฏิสํยุตฺโต สงฺกปฺโป ตปฺปฏิสํยุตฺโต ราโค ตปฺปฏิสํยุตฺโต สงฺกปฺปราโค ตปฺปฏิสํยุตฺตา ปีติ ตปฺปฏิสํยุตฺตํ โสมนสฺสํ ตปฺปฏิสํยุตฺตํ ปีติโสมนสฺสนฺติ,}

\vspace{\tipitakaparskip}
\csromancenter{อฏฺฐมวคฺค อามนฺตา.\csromanspacer{} หญฺจิ อตฺถิ ตปฺปฏิสํยุตฺตํ ปีติโสมนสฺสํ, โน จ วต เร วตฺตพฺเพ “ปญฺเจวายตนา กามา”ติ.}
\csromanmatikaanchor{mtk.103}
\csromantocmarkref{subsection}{(77) 5. รูปธาตุกถา}{mtk.103}
\bookmark[dest={mtk.103},level=2]{(77) 5. รูปธาตุกถา}
\proseitem{514}{\csromanfolio{275}\csromansymbolmark{+}น วตฺตพฺพํ “ปญฺเจวายตนา กามา”ติ, อามนฺตา.\csromanspacer{} นนุ วุตฺตํ ภควตา “ปญฺจิเม ภิกฺขเว กามคุณา.\csromanspacer{} กตเม ปญฺจ, จกฺขุวิญฺเญยฺยา รูปา ฯเปฯ.\csromanspacer{} กายวิญฺเญยฺยา โผฏฺฐพฺพา อิฏฺฐา กนฺตา มนาปา ปิยรูปา กามูปสํหิตา รชนียา.\csromanspacer{} อิเม โข ภิกฺขเว ปญฺจ กามคุณา”ติ อตฺเถว สุตฺตนฺโตติ, อามนฺตา.\csromanspacer{} เตน หิ ปญฺเจวายตนา กามาติ.}

\prose{\csromansymbolmark{*}ปญฺเจวายตนา กามาติ, อามนฺตา.\csromanspacer{} นนุ วุตฺตํ ภควตา “ปญฺจิเม ภิกฺขเว กามคุณา.\csromanspacer{} กตเม ปญฺจ, จกฺขุวิญฺเญยฺยา รูปา ฯเปฯ กายวิญฺเญยฺยา โผฏฺฐพฺพา อิฏฺฐา กนฺตา มนาปา ปิยรูปา กามูปสํหิตา รชนียา.\csromanspacer{} อิเม โข ภิกฺขเว ปญฺจ กามคุณา”. “อปิ จ ภิกฺขเว เนเต กามา, กามคุณา นาเมเต อริยสฺส วินเย วุจฺจนฺติ.}

\csromangathameasure{สงฺกปฺปราโค ปุริสสฺส กาโม, \\ น เต กามา ยานิ จิตฺรานิ โลเก. \\ สงฺกปฺปราโค ปุริสสฺส กาโม, \\ ติฏฺฐนฺติ จิตฺรานิ ตเถว โลเก. \\ อเถตฺถ ธีรา วินยนฺติ ฉนฺทนฺ”ติ. อตฺเถว สุตฺตนฺโตติ, \\ อามนฺตา. เตน หิ น วตฺตพฺพํ “ปญฺเจวายตนา กามา”ติ.}
\csromangathagroup{\csromangathastanza{สงฺกปฺปราโค ปุริสสฺส กาโม, \\ น เต กามา ยานิ จิตฺรานิ โลเก. \\ สงฺกปฺปราโค ปุริสสฺส กาโม, \\ ติฏฺฐนฺติ จิตฺรานิ ตเถว โลเก. \\ อเถตฺถ ธีรา วินยนฺติ ฉนฺทนฺ”ติ\footnote{อํ 2. 359 ปิฏฺเฐ นิพฺเพธิกสุตฺเต.}.\csromanspacer{} อตฺเถว สุตฺตนฺโตติ, \\ อามนฺตา.\csromanspacer{} เตน หิ น วตฺตพฺพํ “ปญฺเจวายตนา กามา”ติ.}}

\nitthitam{กามกถา นิฏฺฐิตา.}
\chapterhead{8. อฏฺฐมวคฺค}
\vspace{\tipitakaparskip}
\csromancenter{\textbf{(77) 5.}\csromanspacer{}\textbf{ รูปธาตุกถา}}
\csromanmatikaanchor{mtk.104}
\csromantocmarkref{subsection}{(78) 6. อรูปธาตุกถา}{mtk.104}
\bookmark[dest={mtk.104},level=2]{(78) 6. อรูปธาตุกถา}
\proseitem{515}{\csromansymbolmark{*}รูปิโน ธมฺมา รูปธาตูติ, อามนฺตา.\csromanspacer{} รูปํ ภโว คติ สตฺตาวาโส สํสาโร โยนิ วิญฺญาณฏฺฐิติ อตฺตภาวปฏิลาโภติ, น \csromanfolio{276}เหวํ วตฺตพฺเพ ฯเปฯ.\csromanspacer{} อตฺถิ รูปูปคํ กมฺมนฺติ, น เหวํ วตฺตพฺเพ ฯเปฯ.\csromanspacer{} อตฺถิ รูปูปคา สตฺตาติ, น เหวํ วตฺตพฺเพ ฯเปฯ.\csromanspacer{} รูเป สตฺตา ชายนฺติ ชียนฺติ มียนฺติ จวนฺติ อุปปชฺชนฺตีติ, น เหวํ วตฺตพฺเพ ฯเปฯ.\csromanspacer{} รูเป อตฺถิ รูปํ เวทนา สญฺญา สงฺขารา วิญฺญาณนฺติ, น เหวํ วตฺตพฺเพ ฯเปฯ.\csromanspacer{} รูปํ ปญฺจโวการภโวติ, น เหวํ วตฺตพฺเพ ฯเปฯ.}

\prose{\csromansymbolmark{*}รูปธาตุ ภโว คติ ฯเปฯ อตฺตภาวปฏิลาโภติ, อามนฺตา.\csromanspacer{} รูปํ ภโว คติ ฯเปฯ อตฺตภาวปฏิลาโภติ, น เหวํ วตฺตพฺเพ ฯเปฯ.\csromanspacer{} อตฺถิ รูปธาตูปคํ กมฺมนฺติ, อามนฺตา.\csromanspacer{} อตฺถิ รูปูปคํ กมฺมนฺติ, น เหวํ วตฺตพฺเพ ฯเปฯ.\csromanspacer{} อตฺถิ รูปธาตูปคา สตฺตาติ, อามนฺตา.\csromanspacer{} อตฺถิ รูปูปคา สตฺตาติ, น เหวํ วตฺตพฺเพ ฯเปฯ.}

\prose{\csromansymbolmark{*}รูปธาตุยา สตฺตา ชายนฺติ ชียนฺติ มียนฺติ จวนฺติ อุปปชฺชนฺตีติ, อามนฺตา.\csromanspacer{} รูเป สตฺตา ชายนฺติ ชียนฺติ มียนฺติ จวนฺติ อุปปชฺชนฺตีติ, น เหวํ วตฺตพฺเพ ฯเปฯ.\csromanspacer{} รูปธาตุยา อตฺถิ รูปํ เวทนา สญฺญา สงฺขารา วิญฺญาณนฺติ, อามนฺตา.\csromanspacer{} รูเป อตฺถิ รูปํ เวทนา สญฺญา สงฺขารา วิญฺญาณนฺติ, น เหวํ วตฺตพฺเพ ฯเปฯ.\csromanspacer{} รูปธาตุ ปญฺจโวการภโวติ, อามนฺตา.\csromanspacer{} รูปํ ปญฺจโวการภโวติ, น เหวํ วตฺตพฺเพ ฯเปฯ.}

\proseitem{516}{\csromansymbolmark{*}รูปิโน ธมฺมา รูปธาตุ, กามธาตุยา อตฺถิ รูปนฺติ, อามนฺตา.\csromanspacer{} สาว กามธาตุ, สา รูปธาตูติ, น เหวํ วตฺตพฺเพ ฯเปฯ.\csromanspacer{} สาว กามธาตุ, สา รูปธาตูติ, อามนฺตา.\csromanspacer{} กามภเวน สมนฺนาคโต ปุคฺคโล ทฺวีหิ ภเวหิ สมนฺนาคโต โหติ กามภเวน จ รูปภเวน จาติ, น เหวํ วตฺตพฺเพ.}

\nitthitam{รูปธาตุกถา นิฏฺฐิตา.}
\chapterhead{8. อฏฺฐมวคฺค}
\vspace{\tipitakaparskip}
\csromancenter{\textbf{(78) 6.}\csromanspacer{}\textbf{ อรูปธาตุกถา}}
\proseitem{517}{\csromansymbolmark{*}อรูปิโน ธมฺมา อรูปธาตูติ, อามนฺตา.\csromanspacer{} เวทนา ภโว คติ สตฺตาวาโส สํสาโร โยนิ วิญฺญาณฏฺฐิติ อตฺตภาวปฏิลาโภติ,}

\vspace{\tipitakaparskip}
\csromancenter{อฏฺฐมวคฺค น เหวํ วตฺตพฺเพ ฯเปฯ.\csromanspacer{} อตฺถิ เวทนูปคํ กมฺมนฺติ, น เหวํ วตฺตพฺเพ ฯเปฯ.\csromanspacer{} อตฺถิ เวทนูปคา สตฺตาติ, น เหวํ วตฺตพฺเพ ฯเปฯ.\csromanspacer{} เวทนาย สตฺตา ชายนฺติ ชียนฺติ มียนฺติ จวนฺติ อุปปชฺชนฺตีติ, น เหวํ วตฺตพฺเพ ฯเปฯ.\csromanspacer{} เวทนาย อตฺถิ เวทนา สญฺญา สงฺขารา วิญฺญาณนฺติ, น เหวํ วตฺตพฺเพ ฯเปฯ.\csromanspacer{} เวทนา จตุโวการภโวติ, น เหวํ วตฺตพฺเพ ฯเปฯ.}
\prose{\csromanfolio{277}\csromansymbolmark{*}อรูปธาตุ ภโว คติ ฯเปฯ อตฺตภาวปฏิลาโภติ, อามนฺตา.\csromanspacer{} เวทนา ภโว คติ ฯเปฯ อตฺตภาวปฏิลาโภติ, น เหวํ วตฺตพฺเพ ฯเปฯ.\csromanspacer{} อตฺถิ อรูปธาตูปคํ กมฺมนฺติ, อามนฺตา.\csromanspacer{} อตฺถิ เวทนูปคํ กมฺมนฺติ, น เหวํ วตฺตพฺเพ ฯเปฯ.\csromanspacer{} อตฺถิ อรูปธาตูปคา สตฺตาติ, อามนฺตา.\csromanspacer{} อตฺถิ เวทนูปคา สตฺตาติ, น เหวํ วตฺตพฺเพ ฯเปฯ.}

\prose{\csromansymbolmark{*}อรูปธาตุยา สตฺตา ชายนฺติ ชียนฺติ มียนฺติ จวนฺติ อุปปชฺชนฺตีติ, อามนฺตา.\csromanspacer{} เวทนาย สตฺตา ชายนฺติ ชียนฺติ มียนฺติ จวนฺติ อุปปชฺชนฺตีติ, น เหวํ วตฺตพฺเพ ฯเปฯ.\csromanspacer{} อรูปธาตุยา อตฺถิ เวทนา สญฺญา สงฺขารา วิญฺญาณนฺติ, อามนฺตา.\csromanspacer{} เวทนาย อตฺถิ เวทนา สญฺญา สงฺขารา วิญฺญาณนฺติ, น เหวํ วตฺตพฺเพ ฯเปฯ.\csromanspacer{} อรูปธาตุ จตุโวการภโวติ, อามนฺตา.\csromanspacer{} เวทนา จตุโวการภโวติ, น เหวํ วตฺตพฺเพ ฯเปฯ.}

\proseitem{518}{\csromansymbolmark{*}อรูปิโน ธมฺมา อรูปธาตุ, กามธาตุยา อตฺถิ เวทนา สญฺญา สงฺขารา วิญฺญาณนฺติ, อามนฺตา.\csromanspacer{} สาว กามธาตุ, สา อรูปธาตูติ, น เหวํ วตฺตพฺเพ ฯเปฯ.\csromanspacer{} สาว กามธาตุ, สา อรูปธาตูติ, อามนฺตา.\csromanspacer{} กามภเวน สมนฺนาคโต ปุคฺคโล ทฺวีหิ ภเวหิ สมนฺนาคโต โหติ กามภเวน จ อรูปภเวน จาติ, น เหวํ วตฺตพฺเพ ฯเปฯ.}

\prose{\csromansymbolmark{*}รูปิโน ธมฺมา รูปธาตุ, อรูปิโน ธมฺมา อรูปธาตุ, กามธาตุยา อตฺถิ รูปํ เวทนา สญฺญา สงฺขารา วิญฺญาณนฺติ, อามนฺตา.\csromanspacer{} สาว กามธาตุ, สา รูปธาตุ, สา อรูปธาตูติ, น เหวํ วตฺตพฺเพ ฯเปฯ.\csromanspacer{} สาว กามธาตุ, สา รูปธาตุ, สา อรูปธาตูติ, อามนฺตา.\csromanspacer{} กามภเวน สมนฺนาคโต ปุคฺคโล ตีหิ ภเวหิ สมนฺนาคโต โหติ กามภเวน จ รูปภเวน จ อรูปภเวน จาติ, น เหวํ วตฺตพฺเพ ฯเปฯ.}

\nitthitam{อรูปธาตุกถา นิฏฺฐิตา.}
\chapterheadpageread{8. อฏฺฐมวคฺค}
\vspace{\tipitakaparskip}
\csromancenter{\textbf{(79) 7.}\csromanspacer{}\textbf{ รูปธาตุยา-อายตนกถา}}
\csromanmatikaanchor{mtk.105}
\csromantocmarkref{subsection}{(79) 7. รูปธาตุยา อายตนกถา}{mtk.105}
\bookmark[dest={mtk.105},level=2]{(79) 7. รูปธาตุยา อายตนกถา}
\proseitem{519}{\csromanfolio{278}\csromansymbolmark{*}อตฺถิ สฬายตนิโก อตฺตภาโว รูปธาตุยาติ, อามนฺตา.\csromanspacer{} อตฺถิ ตตฺถ ฆานายตนนฺติ, อามนฺตา.\csromanspacer{} อตฺถิ ตตฺถ คนฺธายตนนฺติ, น เหวํ วตฺตพฺเพ ฯเปฯ.\csromanspacer{} อตฺถิ ตตฺถ ชิวฺหายตนนฺติ, อามนฺตา.\csromanspacer{} อตฺถิ ตตฺถ รสายตนนฺติ, น เหวํ วตฺตพฺเพ ฯเปฯ.\csromanspacer{} อตฺถิ ตตฺถ กายายตนนฺติ, อามนฺตา.\csromanspacer{} อตฺถิ ตตฺถ โผฏฺฐพฺพายตนนฺติ, น เหวํ วตฺตพฺเพ ฯเปฯ.}

\prose{\csromansymbolmark{*}นตฺถิ ตตฺถ คนฺธายตนนฺติ, อามนฺตา.\csromanspacer{} นตฺถิ ตตฺถ ฆานายตนนฺติ, น เหวํ วตฺตพฺเพ ฯเปฯ.\csromanspacer{} นตฺถิ ตตฺถ รสายตนนฺติ, อามนฺตา.\csromanspacer{} นตฺถิ ตตฺถ ชิวฺหายตนนฺติ, น เหวํ วตฺตพฺเพ ฯเปฯ.\csromanspacer{} นตฺถิ ตตฺถ โผฏฺฐพฺพายตนนฺติ, อามนฺตา.\csromanspacer{} นตฺถิ ตตฺถ กายายตนนฺติ, น เหวํ วตฺตพฺเพ ฯเปฯ.}

\proseitem{520}{\csromansymbolmark{*}อตฺถิ ตตฺถ จกฺขายตนํ อตฺถิ รูปายตนนฺติ, อามนฺตา.\csromanspacer{} อตฺถิ ตตฺถ ฆานายตนํ อตฺถิ คนฺธายตนนฺติ, น เหวํ วตฺตพฺเพ ฯเปฯ.\csromanspacer{} อตฺถิ ตตฺถ จกฺขายตนํ อตฺถิ รูปายตนนฺติ, อามนฺตา.\csromanspacer{} อตฺถิ ตตฺถ ชิวฺหายตนํ อตฺถิ รสายตนํ ฯเปฯ.\csromanspacer{} อตฺถิ ตตฺถ กายายตนํ อตฺถิ โผฏฺฐพฺพายตนนฺติ, น เหวํ วตฺตพฺเพ ฯเปฯ.\csromanspacer{} อตฺถิ ตตฺถ โสตายตนํ อตฺถิ สทฺทายตนํ ฯเปฯ.\csromanspacer{} อตฺถิ ตตฺถ มนายตนํ อตฺถิ ธมฺมายตนนฺติ, อามนฺตา.\csromanspacer{} อตฺถิ ตตฺถ ฆานายตนํ อตฺถิ คนฺธายตนนฺติ, น เหวํ วตฺตพฺเพ ฯเปฯ.\csromanspacer{} อตฺถิ ตตฺถ มนายตนํ อตฺถิ ธมฺมายตนนฺติ, อามนฺตา.\csromanspacer{} อตฺถิ ตตฺถ ชิวฺหายตนํ อตฺถิ รสายตนนฺติ ฯเปฯ.\csromanspacer{} อตฺถิ ตตฺถ กายายตนํ อตฺถิ โผฏฺฐพฺพายตนนฺติ, น เหวํ วตฺตพฺเพ ฯเปฯ.}

\prose{\csromansymbolmark{*}อตฺถิ ตตฺถ ฆานายตนํ นตฺถิ คนฺธายตนนฺติ, อามนฺตา.\csromanspacer{} อตฺถิ ตตฺถ จกฺขายตนํ นตฺถิ รูปายตนนฺติ, น เหวํ วตฺตพฺเพ ฯเปฯ.\csromanspacer{} อตฺถิ ตตฺถ ฆานายตนํ นตฺถิ คนฺธายตนนฺติ, อามนฺตา.\csromanspacer{} อตฺถิ ตตฺถ โสตายตนํ นตฺถิ สทฺทายตนํ ฯเปฯ.\csromanspacer{} อตฺถิ ตตฺถ มนายตนํ นตฺถิ ธมฺมายตนนฺติ, น เหวํ วตฺตพฺเพ ฯเปฯ.\csromanspacer{} อตฺถิ ตตฺถ ชิวฺหายตนํ นตฺถิ รสายตนํ ฯเปฯ.\csromanspacer{} อตฺถิ ตตฺถ กายายตนํ นตฺถิ โผฏฺฐพฺพายตนนฺติ, อามนฺตา.\csromanspacer{} อตฺถิ ตตฺถ จกฺขายตนํ นตฺถิ รูปายตนนฺติ, น เหวํ วตฺตพฺเพ ฯเปฯ.\csromanspacer{} อตฺถิ ตตฺถ กายายตนํ นตฺถิ โผฏฺฐพฺพายตนนฺติ, อามนฺตา.\csromanspacer{} อตฺถิ ตตฺถ โสตายตนํ นตฺถิ}

\vspace{\tipitakaparskip}
\csromancenter{อฏฺฐมวคฺค สทฺทายตนํ ฯเปฯ.\csromanspacer{} อตฺถิ ตตฺถ มนายตนํ นตฺถิ ธมฺมายตนนฺติ, น เหวํ วตฺตพฺเพ ฯเปฯ.}
\proseitem{521}{\csromanfolio{279}\csromansymbolmark{*}อตฺถิ ตตฺถ จกฺขายตนํ อตฺถิ รูปายตนํ เตน จกฺขุนา ตํ รูปํ ปสฺสตีติ, อามนฺตา.\csromanspacer{} อตฺถิ ตตฺถ ฆานายตนํ อตฺถิ คนฺธายตนํ เตน ฆาเนน ตํ คนฺธํ ฆายตีติ, น เหวํ วตฺตพฺเพ ฯเปฯ.\csromanspacer{} อตฺถิ ตตฺถ จกฺขายตนํ อตฺถิ รูปายตนํ เตน จกฺขุนา ตํ รูปํ ปสฺสตีติ, อามนฺตา.\csromanspacer{} อตฺถิ ตตฺถ ชิวฺหายตนํ อตฺถิ รสายตนํ ตาย ชิวฺหาย ตํ รสํ สายติ ฯเปฯ.\csromanspacer{} อตฺถิ ตตฺถ กายายตนํ อตฺถิ โผฏฺฐพฺพายตนํ เตน กาเยน ตํ โผฏฺฐพฺพํ ผุสตีติ, น เหวํ วตฺตพฺเพ ฯเปฯ.}

\prose{\csromansymbolmark{*}อตฺถิ ตตฺถ โสตายตนํ อตฺถิ สทฺทายตนํ ฯเปฯ.\csromanspacer{} อตฺถิ ตตฺถ มนายตนํ อตฺถิ ธมฺมายตนํ เตน มเนน ตํ ธมฺมํ วิชานาตีติ, อามนฺตา.\csromanspacer{} อตฺถิ ตตฺถ ฆานายตนํ อตฺถิ คนฺธายตนํ เตน ฆาเนน ตํ คนฺธํ ฆายตีติ, น เหวํ วตฺตพฺเพ ฯเปฯ.\csromanspacer{} อตฺถิ ตตฺถ มนายตนํ อตฺถิ ธมฺมายตนํ เตน มเนน ตํ ธมฺมํ วิชานาตีติ, อามนฺตา.\csromanspacer{} อตฺถิ ตตฺถ ชิวฺหายตนํ อตฺถิ รสายตนํ ฯเปฯ.\csromanspacer{} อตฺถิ ตตฺถ กายายตนํ อตฺถิ โผฏฺฐพฺพายตนํ เตน กาเยน ตํ โผฏฺฐพฺพํ ผุสตีติ, น เหวํ วตฺตพฺเพ ฯเปฯ.}

\prose{\csromansymbolmark{*}อตฺถิ ตตฺถ ฆานายตนํ อตฺถิ คนฺธายตนํ, น จ เตน ฆาเนน ตํ คนฺธํ ฆายตีติ, อามนฺตา.\csromanspacer{} อตฺถิ ตตฺถ จกฺขายตนํ อตฺถิ รูปายตนํ น จ เตน จกฺขุนา ตํ รูปํ ปสฺสตีติ, น เหวํ วตฺตพฺเพ ฯเปฯ.\csromanspacer{} อตฺถิ ตตฺถ ฆานายตนํ อตฺถิ คนฺธายตนํ น จ เตน ฆาเนน ตํ คนฺธํ ฆายตีติ, อามนฺตา.\csromanspacer{} อตฺถิ ตตฺถ โสตายตนํ อตฺถิ สทฺทายตนํ ฯเปฯ.\csromanspacer{} อตฺถิ ตตฺถ มนายตนํ อตฺถิ ธมฺมายตนํ น จ เตน มเนน ตํ ธมฺมํ วิชานาตีติ, น เหวํ วตฺตพฺเพ ฯเปฯ.}

\prose{\csromansymbolmark{*}อตฺถิ ตตฺถ ชิวฺหายตนํ อตฺถิ รสายตนํ ฯเปฯ.\csromanspacer{} อตฺถิ ตตฺถ กายายตนํ อตฺถิ โผฏฺฐพฺพายตนํ น จ เตน กาเยน ตํ โผฏฺฐพฺพํ ผุสตีติ, อามนฺตา.\csromanspacer{} อตฺถิ ตตฺถ โสตายตนํ, อตฺถิ สทฺทายตนํ ฯเปฯ.\csromanspacer{} อตฺถิ ตตฺถ มนายตนํ อตฺถิ ธมฺมายตนํ น จ เตน มเนน ตํ ธมฺมํ วิชานาตีติ, น เหวํ วตฺตพฺเพ ฯเปฯ.}

\csromanmatikaanchor{mtk.106}
\csromantocmarkref{subsection}{(80) 8. อรูเป รูปกถา}{mtk.106}
\bookmark[dest={mtk.106},level=2]{(80) 8. อรูเป รูปกถา}
\proseitem{522}{\csromanfolio{280}\csromansymbolmark{*}อตฺถิ ตตฺถ ฆานายตนํ อตฺถิ คนฺธายตนํ เตน ฆาเนน ตํ คนฺธํ ฆายตีติ, อามนฺตา.\csromanspacer{} อตฺถิ ตตฺถ มูลคนฺโธ สารคนฺโธ ตจคนฺโธ ปตฺตคนฺโธ ปุปฺผคนฺโธ ผลคนฺโธ อามคนฺโธ วิสฺสคนฺโธ สุคนฺโธ ทุคฺคนฺโธติ, น เหวํ วตฺตพฺเพ ฯเปฯ.}

\prose{\csromansymbolmark{*}อตฺถิ ตตฺถ ชิวฺหายตนํ อตฺถิ รสายตนํ ตาย ชิวฺหาย ตํ รสํ สายตีติ, อามนฺตา.\csromanspacer{} อตฺถิ ตตฺถ มูลรโส ขนฺธรโส ตจรโส ปตฺตรโส ปุปฺผรโส ผลรโส อมฺพิลํ มธุรํ ติตฺตกํ กฏุกํ โลณิยํ ขาริยํ ลมฺพิลํ กสาโว สาทุ อสาทูติ, น เหวํ วตฺตพฺเพ ฯเปฯ.}

\prose{\csromansymbolmark{*}อตฺถิ ตตฺถ กายายตนํ อตฺถิ โผฏฺฐพฺพายตนํ เตน กาเยน ตํ โผฏฺฐพฺพํ ผุสตีติ, อามนฺตา.\csromanspacer{} อตฺถิ ตตฺถ กกฺขฬํ มุทุกํ สณฺหํ ผรุสํ สุขสมฺผสฺสํ ทุกฺขสมฺผสฺสํ ครุกํ ลหุกนฺติ, น เหวํ วตฺตพฺเพ ฯเปฯ.}

\proseitem{523}{\csromansymbolmark{+}น วตฺตพฺพํ “สฬายตนิโก อตฺตภาโว รูปธาตุยา”ติ, อามนฺตา.\csromanspacer{} นนุ อตฺถิ ตตฺถ ฆานนิมิตฺตํ ชิวฺหานิมิตฺตํ กายนิมิตฺตนฺติ, อามนฺตา.\csromanspacer{} หญฺจิ อตฺถิ ตตฺถ ฆานนิมิตฺตํ ชิวฺหานิมิตฺตํ กายนิมิตฺตํ, เตน วต เร วตฺตพฺเพ “สฬายตนิโก อตฺตภาโว รูปธาตุยา”ติ.}

\nitthitam{รูปธาตุยา อายตนกถา นิฏฺฐิตา.}
\chapterhead{8. อฏฺฐมวคฺค}
\vspace{\tipitakaparskip}
\csromancenter{\textbf{(80) 8.}\csromanspacer{}\textbf{ อรูเปรูปกถา}}
\proseitem{524}{\csromansymbolmark{*}อตฺถิ รูปํ อรูเปสูติ, อามนฺตา.\csromanspacer{} รูปภโว รูปคติ รูปสตฺตาวาโส รูปสํสาโร รูปโยนิ รูปตฺตภาวปฏิลาโภติ, น เหวํ วตฺตพฺเพ ฯเปฯ.\csromanspacer{} นนุ อรูปภโว อรูปคติ อรูปสตฺตาวาโส อรูปสํสาโร อรูปโยนิ อรูปตฺตภาวปฏิลาโภติ, อามนฺตา.\csromanspacer{} หญฺจิ อรูปภโว ฯเปฯ อรูปตฺตภาวปฏิลาโภ, โน จ วต เร วตฺตพฺเพ “อตฺถิ รูปํ อรูเปสู”ติ.}

\prose{\csromansymbolmark{*}อตฺถิ รูปํ อรูเปสูติ, อามนฺตา.\csromanspacer{} ปญฺจโวการภโว คติ สตฺตาวาโส สํสาโร โยนิ วิญฺญาณฏฺฐิติ อตฺตภาวปฏิลาโภติ, น}

\vspace{\tipitakaparskip}
\csromancenter{อฏฺฐมวคฺค เหวํ วตฺตพฺเพ ฯเปฯ.\csromanspacer{} นนุ จตุโวการภโว ฯเปฯ อตฺตภาวปฏิลาโภติ, อามนฺตา.\csromanspacer{} หญฺจิ จตุโวการภโว คติ ฯเปฯ อตฺตภาวปฏิลาโภ, โน จ วต เร วตฺตพฺเพ “อตฺถิ รูปํ อรูเปสู”ติ.}
\proseitem{525}{\csromanfolio{281}\csromansymbolmark{*}อตฺถิ รูปํ รูปธาตุยา, โส จ รูปภโว รูปคติ รูปสตฺตาวาโส รูปสํสาโร รูปโยนิ รูปตฺตภาวปฏิลาโภติ, อามนฺตา.\csromanspacer{} อตฺถิ รูปํ อรูเปสุ, โส จ รูปภโว รูปคติ ฯเปฯ รูปตฺตภาวปฏิลาโภติ, น เหวํ วตฺตพฺเพ ฯเปฯ.\csromanspacer{} อตฺถิ รูปํ รูปธาตุยา, โส จ ปญฺจโวการภโว คติ ฯเปฯ อตฺตภาวปฏิลาโภติ, อามนฺตา.\csromanspacer{} อตฺถิ รูปํ อรูเปสุ, โส จ ปญฺจโวการภโว คติ ฯเปฯ อตฺตภาวปฏิลาโภติ, น เหวํ วตฺตพฺเพ ฯเปฯ.}

\prose{\csromansymbolmark{*}อตฺถิ รูปํ อรูเปสุ, โส จ อรูปภโว อรูปคติ อรูปสตฺตาวาโส อรูปสํสาโร อรูปโยนิ อรูปตฺตภาวปฏิลาโภติ, อามนฺตา.\csromanspacer{} อตฺถิ รูปํ รูปธาตุยา\footnote{อรูปธาตุยา (สพฺพตฺถ)}, โส จ อรูปภโว อรูปคติ ฯเปฯ อรูปตฺตภาวปฏิลาโภติ, น เหวํ วตฺตพฺเพ ฯเปฯ.\csromanspacer{} อตฺถิ รูปํ อรูเปสุ, โส จ จตุโวการภโว คติ ฯเปฯ อตฺตภาวปฏิลาโภติ, อามนฺตา.\csromanspacer{} อตฺถิ รูปํ รูปธาตุยา1, โส จ จตุโวการภโว คติ ฯเปฯ อตฺตภาวปฏิลาโภติ, น เหวํ วตฺตพฺเพ ฯเปฯ.}

\proseitem{526}{\csromansymbolmark{*}อตฺถิ รูปํ อรูเปสูติ, อามนฺตา.\csromanspacer{} นนุ รูปานํ นิสฺสรณํ อรูปํ\footnote{ขุ 1. 237; อํ 2. 214; ที 3. 200 ปิฏฺเฐสุ.} วุตฺตํ ภควตาติ, อามนฺตา.\csromanspacer{} หญฺจิ รูปานํ นิสฺสรณํ อรูปํ2 วุตฺตํ ภควตา, โน จ วต เร วตฺตพฺเพ “อตฺถิ รูปํ อรูเปสู”ติ.}

\prose{\csromansymbolmark{*}รูปานํ นิสฺสรณํ อรูปํ วุตฺตํ ภควตา “อตฺถิ รูปํ อรูเปสู”ติ, อามนฺตา.\csromanspacer{} กามานํ นิสฺสรณํ เนกฺขมฺมํ วุตฺตํ ภควตา “อตฺถิ เนกฺขมฺเมสุ กามา, อตฺถิ อนาสเวสุ อาสวา, อตฺถิ อปริยาปนฺเนสุ ปริยาปนฺนา”ติ, น เหวํ วตฺตพฺเพ ฯเปฯ.}

\nitthitam{อรูเป รูปกถา นิฏฺฐิตา.}
\chapterheadpageread{8. อฏฺฐมวคฺค}
\vspace{\tipitakaparskip}
\csromancenter{\textbf{(81) 9.}\csromanspacer{}\textbf{ รูปํกมฺมนฺติกถา}}
\csromanmatikaanchor{mtk.107}
\csromantocmarkref{subsection}{(81) 9. รูปํ กมฺมนฺติกถา}{mtk.107}
\bookmark[dest={mtk.107},level=2]{(81) 9. รูปํ กมฺมนฺติกถา}
\proseitem{527}{\csromanfolio{282}\csromansymbolmark{*}กุสเลน จิตฺเตน สมุฏฺฐิตํ กายกมฺมํ รูปํ กุสลนฺติ, อามนฺตา.\csromanspacer{} สฺèรมฺมณํ, อตฺถิ ตสฺส อาวฏฺฏนา อาโภโค สมนฺนาหาโร มนสิกาโร เจตนา ปตฺถนา ปณิธีติ, น เหวํ วตฺตพฺเพ ฯเปฯ.\csromanspacer{} นนุ อนารมฺมณํ, นตฺถิ ตสฺส อาวฏฺฏนา อาโภโค สมนฺนาหาโร มนสิกาโร เจตนา ปตฺถนา ปณิธีติ, อามนฺตา.\csromanspacer{} หญฺจิ อนารมฺมณํ, นตฺถิ ตสฺส อาวฏฺฏนา ฯเปฯ ปณิธีติ, โน จ วต เร วตฺตพฺเพ “กุสเลน จิตฺเตน สมุฏฺฐิตํ กายกมฺมํ รูปํ กุสลนฺ”ติ.}

\prose{\csromansymbolmark{*}กุสเลน จิตฺเตน สมุฏฺฐิโต ผสฺโส กุสโล สารมฺมโณ, อตฺถิ ตสฺส อาวฏฺฏนา ฯเปฯ ปณิธีติ, อามนฺตา.\csromanspacer{} กุสเลน จิตฺเตน สมุฏฺฐิตํ กายกมฺมํ รูปํ กุสลํ สารมฺมณํ, อตฺถิ ตสฺส อาวฏฺฏนา ฯเปฯ ปณิธีติ, น เหวํ วตฺตพฺเพ ฯเปฯ.\csromanspacer{} กุสเลน จิตฺเตน สมุฏฺฐิตา เวทนา ฯเปฯ สญฺญา.\csromanspacer{} เจตนา.\csromanspacer{} สทฺธา.\csromanspacer{} วีริยํ.\csromanspacer{} สติ.\csromanspacer{} สมาธิ ฯเปฯ.\csromanspacer{} ปญฺญา กุสลา สารมฺมณา, อตฺถิ ตาย อาวฏฺฏนา ฯเปฯ ปณิธีติ, อามนฺตา.\csromanspacer{} กุสเลน จิตฺเตน สมุฏฺฐิตํ กายกมฺมํ รูปํ กุสลํ สารมฺมณํ, อตฺถิ ตสฺส อาวฏฺฏนา ฯเปฯ ปณิธีติ, น เหวํ วตฺตพฺเพ ฯเปฯ.}

\prose{\csromansymbolmark{*}กุสเลน จิตฺเตน สมุฏฺฐิตํ กายกมฺมํ รูปํ กุสลํ อนารมฺมณํ, นตฺถิ ตสฺส อาวฏฺฏนา ฯเปฯ ปณิธีติ, อามนฺตา.\csromanspacer{} กุสเลน จิตฺเตน สมุฏฺฐิโต ผสฺโส กุสโล อนารมฺมโณ, นตฺถิ ตสฺส อาวฏฺฏนา ฯเปฯ ปณิธีติ, น เหวํ วตฺตพฺเพ ฯเปฯ.\csromanspacer{} กุสเลน จิตฺเตน สมุฏฺฐิตํ กายกมฺมํ รูปํ กุสลํ อนารมฺมณํ, นตฺถิ ตสฺส อาวฏฺฏนา ฯเปฯ ปณิธีติ, อามนฺตา.\csromanspacer{} กุสเลน จิตฺเตน สมุฏฺฐิตา เวทนา ฯเปฯ สญฺญา.\csromanspacer{} เจตนา.\csromanspacer{} สทฺธา.\csromanspacer{} วีริยํ.\csromanspacer{} สติ.\csromanspacer{} สมาธิ ฯเปฯ ปญฺญา กุสลา อนารมฺมณา, นตฺถิ ตาย อาวฏฺฏนา ฯเปฯ ปณิธีติ, น เหวํ วตฺตพฺเพ ฯเปฯ.}

\proseitem{528}{\csromansymbolmark{*}กุสเลน จิตฺเตน สมุฏฺฐิตํ กายกมฺมํ รูปํ กุสลนฺติ, อามนฺตา.\csromanspacer{} ยํ กิญฺจิ กุสเลน จิตฺเตน สมุฏฺฐิตํ รูปํ สพฺพํ ตํ กุสลนฺติ, น เหวํ วตฺตพฺเพ ฯเปฯ.}

\chapterheadpageread{8. อฏฺฐมวคฺค}
\prose{\csromanfolio{283}\csromansymbolmark{*}กุสเลน จิตฺเตน สมุฏฺฐิตํ กายกมฺมํ รูปํ กุสลนฺติ, อามนฺตา.\csromanspacer{} กุสเลน จิตฺเตน สมุฏฺฐิตํ รูปายตนํ กุสลนฺติ, น เหวํ วตฺตพฺเพ ฯเปฯ.\csromanspacer{} กุสเลน จิตฺเตน สมุฏฺฐิตํ กายกมฺมํ รูปํ กุสลนฺติ, อามนฺตา.\csromanspacer{} กุสเลน จิตฺเตน สมุฏฺฐิตํ สทฺทายตนํ ฯเปฯ คนฺธายตนํ.\csromanspacer{} รสายตนํ.\csromanspacer{} โผฏฺฐพฺพายตนํ ฯเปฯ ปถวีธาตุ.\csromanspacer{} อาโปธาตุ.\csromanspacer{} เตโชธาตุ ฯเปฯ วาโยธาตุ กุสลาติ, น เหวํ วตฺตพฺเพ ฯเปฯ.}

\prose{\csromansymbolmark{*}กุสเลน จิตฺเตน สมุฏฺฐิตํ รูปายตนํ อพฺยากตนฺติ, อามนฺตา.\csromanspacer{} กุสเลน จิตฺเตน สมุฏฺฐิตํ กายกมฺมํ รูปํ อพฺยากตนฺติ, น เหวํ วตฺตพฺเพ ฯเปฯ.\csromanspacer{} กุสเลน จิตฺเตน สมุฏฺฐิตํ สทฺทายตนํ ฯเปฯ คนฺธายตนํ.\csromanspacer{} รสายตนํ.\csromanspacer{} โผฏฺฐพฺพายตนํ ฯเปฯ ปถวีธาตุ.\csromanspacer{} อาโปธาตุ.\csromanspacer{} เตโชธาตุ ฯเปฯ วาโยธาตุ อพฺยากตาติ, อามนฺตา.\csromanspacer{} กุสเลน จิตฺเตน สมุฏฺฐิตํ กายกมฺมํ รูปํ อพฺยากตนฺติ, น เหวํ วตฺตพฺเพ ฯเปฯ.}

\proseitem{529}{\csromansymbolmark{*}กุสเลน จิตฺเตน สมุฏฺฐิตํ กายกมฺมํ รูปํ อนารมฺมณํ กุสลนฺติ, อามนฺตา.\csromanspacer{} กุสเลน จิตฺเตน สมุฏฺฐิตํ รูปายตนํ อนารมฺมณํ กุสลนฺติ, น เหวํ วตฺตพฺเพ ฯเปฯ.\csromanspacer{} กุสเลน จิตฺเตน สมุฏฺฐิตํ กายกมฺมํ รูปํ อนารมฺมณํ กุสลนฺติ, อามนฺตา.\csromanspacer{} กุสเลน จิตฺเตน สมุฏฺฐิตํ สทฺทายตนํ ฯเปฯ คนฺธายตนํ.\csromanspacer{} รสายตนํ.\csromanspacer{} โผฏฺฐพฺพายตนํ ฯเปฯ ปถวีธาตุ.\csromanspacer{} อาโปธาตุ.\csromanspacer{} เตโชธาตุ ฯเปฯ วาโยธาตุ อนารมฺมณา กุสลาติ, น เหวํ วตฺตพฺเพ ฯเปฯ.}

\prose{\csromansymbolmark{*}กุสเลน จิตฺเตน สมุฏฺฐิตํ รูปายตนํ อนารมฺมณํ อพฺยากตนฺติ, อามนฺตา.\csromanspacer{} กุสเลน จิตฺเตน สมุฏฺฐิตํ กายกมฺมํ รูปํ อนารมฺมณํ อพฺยากตนฺติ, น เหวํ วตฺตพฺเพ ฯเปฯ.\csromanspacer{} กุสเลน จิตฺเตน สมุฏฺฐิตํ สทฺทายตนํ ฯเปฯ คนฺธายตนํ.\csromanspacer{} รสายตนํ.\csromanspacer{} โผฏฺฐพฺพายตนํ ฯเปฯ ปถวีธาตุ.\csromanspacer{} อาโปธาตุ.\csromanspacer{} เตโชธาตุ ฯเปฯ วาโยธาตุ อนารมฺมณา อพฺยากตาติ, อามนฺตา.\csromanspacer{} กุสเลน จิตฺเตน สมุฏฺฐิตํ กายกมฺมํ รูปํ อนารมฺมณํ อพฺยากตนฺติ, น เหวํ วตฺตพฺเพ ฯเปฯ.}

\proseitem{530}{\csromansymbolmark{*}กุสเลน จิตฺเตน สมุฏฺฐิตํ กายกมฺมํ รูปํ ผสฺสวิปฺปยุตฺตํ กุสลนฺติ, อามนฺตา.\csromanspacer{} กุสเลน จิตฺเตน สมุฏฺฐิตํ รูปายตนํ ผสฺสวิปฺปยุตฺตํ กุสลนฺติ, น เหวํ วตฺตพฺเพ ฯเปฯ.\csromanspacer{} กุสเลน จิตฺเตน สมุฏฺฐิตํ กายกมฺมํ รูปํ \csromanfolio{284}ผสฺสวิปฺปยุตฺตํ กุสลนฺติ, อามนฺตา.\csromanspacer{} กุสเลน จิตฺเตน สมุฏฺฐิตํ สทฺทายตนํ ฯเปฯ คนฺธายตนํ.\csromanspacer{} รสายตนํ.\csromanspacer{} โผฏฺฐพฺพายตนํ ฯเปฯ ปถวีธาตุ.\csromanspacer{} อาโปธาตุ.\csromanspacer{} เตโชธาตุ ฯเปฯ วาโยธาตุ ผสฺสวิปฺปยุตฺตา กุสลาติ, น เหวํ วตฺตพฺเพ ฯเปฯ.}

\prose{\csromansymbolmark{*}กุสเลน จิตฺเตน สมุฏฺฐิตํ รูปายตนํ ผสฺสวิปฺปยุตฺตํ อพฺยากตนฺติ, อามนฺตา.\csromanspacer{} กุสเลน จิตฺเตน สมุฏฺฐิตํ กายกมฺมํ รูปํ ผสฺสวิปฺปยุตฺตํ อพฺยากตนฺติ, น เหวํ วตฺตพฺเพ ฯเปฯ.\csromanspacer{} กุสเลน จิตฺเตน สมุฏฺฐิตํ สทฺทายตนํ ฯเปฯ คนฺธายตนํ.\csromanspacer{} รสายตนํ.\csromanspacer{} โผฏฺฐพฺพายตนํ ฯเปฯ ปถวีธาตุ.\csromanspacer{} อาโปธาตุ.\csromanspacer{} เตโชธาตุ ฯเปฯ วาโยธาตุ ผสฺสวิปฺปยุตฺตา อพฺยากตาติ, อามนฺตา.\csromanspacer{} กุสเลน จิตฺเตน สมุฏฺฐิตํ กายกมฺมํ รูปํ ผสฺสวิปฺปยุตฺตํ อพฺยากตนฺติ, น เหวํ วตฺตพฺเพ ฯเปฯ.}

\proseitem{531}{\csromansymbolmark{*}กุสเลน จิตฺเตน สมุฏฺฐิตํ กายกมฺมํ รูปํ อนารมฺมณํ ผสฺสวิปฺปยุตฺตํ กุสลนฺติ, อามนฺตา.\csromanspacer{} กุสเลน จิตฺเตน สมุฏฺฐิตํ รูปายตนํ อนารมฺมณํ ผสฺสวิปฺปยุตฺตํ กุสลนฺติ, น เหวํ วตฺตพฺเพ ฯเปฯ.\csromanspacer{} กุสเลน จิตฺเตน สมุฏฺฐิตํ กายกมฺมํ รูปํ อนารมฺมณํ ผสฺสวิปฺปยุตฺตํ กุสลนฺติ, อามนฺตา.\csromanspacer{} กุสเลน จิตฺเตน สมุฏฺฐิตํ สทฺทายตนํ ฯเปฯ คนฺธายตนํ.\csromanspacer{} รสายตนํ โผฏฺฐพฺพายตนํ ฯเปฯ ปถวีธาตุ.\csromanspacer{} อาโปธาตุ.\csromanspacer{} เตโชธาตุ ฯเปฯ วาโยธาตุ อนารมฺมณา ผสฺสวิปฺปยุตฺตา กุสลาติ, น เหวํ วตฺตพฺเพ ฯเปฯ.}

\prose{\csromansymbolmark{*}กุสเลน จิตฺเตน สมุฏฺฐิตํ รูปายตนํ อนารมฺมณํ ผสฺสวิปฺปยุตฺตํ อพฺยากตนฺติ, อามนฺตา.\csromanspacer{} กุสเลน จิตฺเตน สมุฏฺฐิตํ กายกมฺมํ รูปํ อนารมฺมณํ ผสฺสวิปฺปยุตฺตํ อพฺยากตนฺติ, น เหวํ วตฺตพฺเพ ฯเปฯ.\csromanspacer{} กุสเลน จิตฺเตน สมุฏฺฐิตํ สทฺทายตนํ ฯเปฯ คนฺธายตนํ.\csromanspacer{} รสายตนํ.\csromanspacer{} โผฏฺฐพฺพายตนํ ฯเปฯ ปถวีธาตุ.\csromanspacer{} อาโปธาตุ.\csromanspacer{} เตโชธาตุ ฯเปฯ วาโยธาตุ อนารมฺมณา ผสฺสวิปฺปยุตฺตา อพฺยากตาติ, อามนฺตา.\csromanspacer{} กุสเลน จิตฺเตน สมุฏฺฐิตํ กายกมฺมํ รูปํ อนารมฺมณํ ผสฺสวิปฺปยุตฺตํ อพฺยากตนฺติ, น เหวํ วตฺตพฺเพ ฯเปฯ.}

\proseitem{532}{\csromansymbolmark{*}กุสเลน จิตฺเตน สมุฏฺฐิตํ วจีกมฺมํ รูปํ กุสลนฺติ, อามนฺตา.\csromanspacer{} สารมฺมณํ, อตฺถิ ตสฺส อาวฏฺฏนา ฯเปฯ ปณิธีติ, น เหวํ วตฺตพฺเพ ฯเปฯ.\csromanspacer{} นนุ}

\vspace{\tipitakaparskip}
\csromancenter{อฏฺฐมวคฺค อนารมฺมณํ, นตฺถิ ตสฺส อาวฏฺฏนา ฯเปฯ ปณิธีติ, อามนฺตา.\csromanspacer{} หญฺจิ อนารมฺมณํ, นตฺถิ ตสฺส อาวฏฺฏนา ฯเปฯ ปณิธิ, โน จ วต เร วตฺตพฺเพ “กุสเลน จิตฺเตน สมุฏฺฐิตํ วจีกมฺมํ รูปํ กุสลนฺ”ติ.}
\prose{\csromanfolio{285}\csromansymbolmark{*}กุสเลน จิตฺเตน สมุฏฺฐิโต ผสฺโส กุสโล สารมฺมโณ, อตฺถิ ตสฺส อาวฏฺฏนา ฯเปฯ ปณิธีติ, อามนฺตา.\csromanspacer{} กุสเลน จิตฺเตน สมุฏฺฐิตํ วจีกมฺมํ รูปํ กุสลํ สารมฺมณํ, อตฺถิ ตสฺส อาวฏฺฏนา ฯเปฯ ปณิธีติ, น เหวํ วตฺตพฺเพ ฯเปฯ.\csromanspacer{} กุสเลน จิตฺเตน สมุฏฺฐิตา เวทนา ฯเปฯ สญฺญา.\csromanspacer{} เจตนา.\csromanspacer{} สทฺธา.\csromanspacer{} วีริยํ.\csromanspacer{} สติ.\csromanspacer{} สมาธิ ฯเปฯ ปญฺญา กุสลา สารมฺมณา, อตฺถิ ตาย อาวฏฺฏนา ฯเปฯ ปณิธีติ, อามนฺตา.\csromanspacer{} กุสเลน จิตฺเตน สมุฏฺฐิตํ วจีกมฺมํ รูปํ กุสลํ สารมฺมณํ, อตฺถิ ตสฺส อาวฏฺฏนา ฯเปฯ ปณิธีติ, น เหวํ วตฺตพฺเพ ฯเปฯ.}

\prose{\csromansymbolmark{*}กุสเลน จิตฺเตน สมุฏฺฐิตํ วจีกมฺมํ รูปํ กุสลํ อนารมฺมณํ, นตฺถิ ตสฺส อาวฏฺฏนา ฯเปฯ ปณิธีติ, อามนฺตา.\csromanspacer{} กุสเลน จิตฺเตน สมุฏฺฐิโต ผสฺโส กุสโล อนารมฺมโณ, นตฺถิ ตสฺส อาวฏฺฏนา ฯเปฯ ปณิธีติ, น เหวํ วตฺตพฺเพ ฯเปฯ.\csromanspacer{} กุสเลน จิตฺเตน สมุฏฺฐิตํ วจีกมฺมํ รูปํ กุสลํ อนารมฺมณํ, นตฺถิ ตสฺส อาวฏฺฏนา ฯเปฯ ปณิธีติ, อามนฺตา.\csromanspacer{} กุสเลน จิตฺเตน สมุฏฺฐิตา เวทนา ฯเปฯ สญฺญา.\csromanspacer{} เจตนา.\csromanspacer{} สทฺธา.\csromanspacer{} วีริยํ.\csromanspacer{} สติ.\csromanspacer{} สมาธิ ฯเปฯ ปญฺญา กุสลา อนารมฺมณา, นตฺถิ ตาย อาวฏฺฏนา ฯเปฯ ปณิธีติ, น เหวํ วตฺตพฺเพ ฯเปฯ.}

\prose{\csromansymbolmark{*}กุสเลน จิตฺเตน สมุฏฺฐิตํ วจีกมฺมํ รูปํ กุสลนฺติ, อามนฺตา.\csromanspacer{} ยํ กิญฺจิ กุสเลน จิตฺเตน สมุฏฺฐิตํ รูปํ สพฺพํ ตํ กุสลนฺติ, น เหวํ วตฺตพฺเพ ฯเปฯ.\csromanspacer{} ยถา กายกมฺมํ ตถา วจีกมฺมนฺติ.}

\proseitem{533}{\csromansymbolmark{*}อกุสเลน จิตฺเตน สมุฏฺฐิตํ กายกมฺมํ รูปํ อกุสลนฺติ, อามนฺตา.\csromanspacer{} สารมฺมณํ, อตฺถิ ตสฺส อาวฏฺฏนา อาโภโค สมนฺนาหาโร มนสิกาโร เจตนา ปตฺถนา ปณิธีติ, น เหวํ วตฺตพฺเพ ฯเปฯ.\csromanspacer{} นนุ อนารมฺมณํ, นตฺถิ ตสฺส อาวฏฺฏนา อาโภโค สมนฺนาหาโร มนสิกาโร เจตนา ปตฺถนา ปณิธีติ, อามนฺตา.\csromanspacer{} หญฺจิ อนารมฺมณํ, นตฺถิ ตสฺส อาวฏฺฏนา อาโภโค สมนฺนาหาโร มนสิกาโร เจตนา ปตฺถนา ปณิธิ, โน จ วต เร วตฺตพฺเพ “อกุสเลน จิตฺเตน สมุฏฺฐิตํ กายกมฺมํ รูปํ อกุสลนฺ”ติ.}

\prose{\csromanfolio{286}\csromansymbolmark{*}อกุสเลน จิตฺเตน สมุฏฺฐิโต ผสฺโส อกุสโล สารมฺมโณ, อตฺถิ ตสฺส อาวฏฺฏนา ฯเปฯ ปณิธีติ, อามนฺตา.\csromanspacer{} อกุสเลน จิตฺเตน สมุฏฺฐิตํ กายกมฺมํ รูปํ อกุสลํ สารมฺมณํ, อตฺถิ ตสฺส อาวฏฺฏนา ฯเปฯ ปณิธีติ, น เหวํ วตฺตพฺเพ ฯเปฯ.\csromanspacer{} อกุสเลน จิตฺเตน สมุฏฺฐิตา เวทนา ฯเปฯ สญฺญา.\csromanspacer{} เจตนา.\csromanspacer{} ราโค.\csromanspacer{} โทโส.\csromanspacer{} โมโห.\csromanspacer{} มาโน.\csromanspacer{} ทิฏฺฐิ.\csromanspacer{} วิจิกิจฺฉา.\csromanspacer{} ถินํ.\csromanspacer{} อุทฺธจฺจํ.\csromanspacer{} อหิริกํ ฯเปฯ อโนตฺตปฺปํ อกุสลํ สารมฺมณํ, อตฺถิ ตสฺส อาวฏฺฏนา ฯเปฯ ปณิธีติ, อามนฺตา.\csromanspacer{} อกุสเลน จิตฺเตน สมุฏฺฐิตํ กายกมฺมํ รูปํ อกุสลํ สารมฺมณํ, อตฺถิ ตสฺส อาวฏฺฏนา ฯเปฯ ปณิธีติ, น เหวํ วตฺตพฺเพ ฯเปฯ.}

\prose{\csromansymbolmark{*}อกุสเลน จิตฺเตน สมุฏฺฐิตํ กายกมฺมํ รูปํ อกุสลํ อนารมฺมณํ, นตฺถิ ตสฺส อาวฏฺฏนา ฯเปฯ ปณิธีติ, อามนฺตา.\csromanspacer{} อกุสเลน จิตฺเตน สมุฏฺฐิโต ผสฺโส อกุสโล อนารมฺมโณ, นตฺถิ ตสฺส อาวฏฺฏนา ฯเปฯ ปณิธีติ, น เหวํ วตฺตพฺเพ ฯเปฯ.\csromanspacer{} อกุสเลน จิตฺเตน สมุฏฺฐิตํ กายกมฺมํ รูปํ อกุสลํ อนารมฺมณํ, นตฺถิ ตสฺส อาวฏฺฏนา ฯเปฯ ปณิธีติ, อามนฺตา.\csromanspacer{} อกุสเลน จิตฺเตน สมุฏฺฐิตา เวทนา ฯเปฯ สญฺญา.\csromanspacer{} เจตนา.\csromanspacer{} ราโค.\csromanspacer{} โทโส.\csromanspacer{} โมโห.\csromanspacer{} มาโน.\csromanspacer{} ทิฏฺฐิ.\csromanspacer{} วิจิกิจฺฉา.\csromanspacer{} ถินํ.\csromanspacer{} อุทฺธจฺจํ.\csromanspacer{} อหิริกํ ฯเปฯ อโนตฺตปฺปํ อกุสลํ อนารมฺมณํ, นตฺถิ ตสฺส อาวฏฺฏนา ฯเปฯ ปณิธีติ, น เหวํ วตฺตพฺเพ ฯเปฯ.}

\prose{\csromansymbolmark{*}อกุสเลน จิตฺเตน สมุฏฺฐิตํ กายกมฺมํ รูปํ อกุสลนฺติ, อามนฺตา.\csromanspacer{} ยํ กิญฺจิ อกุสเลน จิตฺเตน สมุฏฺฐิตํ รูปํ สพฺพํ ตํ อกุสลนฺติ, น เหวํ วตฺตพฺเพ ฯเปฯ.}

\proseitem{534}{\csromansymbolmark{*}อกุสเลน จิตฺเตน สมุฏฺฐิตํ วจีกมฺมํ รูปํ อกุสลนฺติ, อามนฺตา.\csromanspacer{} สารมฺมณํ, อตฺถิ ตสฺส อาวฏฺฏนา ฯเปฯ ปณิธีติ, น เหวํ วตฺตพฺเพ ฯเปฯ.\csromanspacer{} นนุ อนารมฺมณํ, นตฺถิ ตสฺส อาวฏฺฏนา ฯเปฯ ปณิธีติ, อามนฺตา.\csromanspacer{} หญฺจิ อนารมฺมณํ, นตฺถิ ตสฺส อาวฏฺฏนา ฯเปฯ ปณิธิ, โน จ วต เร วตฺตพฺเพ “อกุสเลน จิตฺเตน สมุฏฺฐิตํ วจีกมฺมํ รูปํ อกุสลนฺ”ติ.}

\prose{\csromansymbolmark{*}อกุสเลน จิตฺเตน สมุฏฺฐิโต ผสฺโส อกุสโล สารมฺมโณ, อตฺถิ ตสฺส อาวฏฺฏนา ฯเปฯ ปณิธีติ, อามนฺตา.\csromanspacer{} อกุสเลน จิตฺเตน สมุฏฺฐิตํ วจีกมฺมํ รูปํ อกุสลํ สารมฺมณํ, อตฺถิ ตสฺส อาวฏฺฏนา ฯเปฯ ปณิธีติ, น เหวํ วตฺตพฺเพ ฯเปฯ.\csromanspacer{} อกุสเลน จิตฺเตน สมุฏฺฐิตา เวทนา ฯเปฯ}

\vspace{\tipitakaparskip}
\csromancenter{อฏฺฐมวคฺค สญฺญา.\csromanspacer{} เจตนา.\csromanspacer{} ราโค.\csromanspacer{} โทโส.\csromanspacer{} โมโห.\csromanspacer{} มาโน.\csromanspacer{} ทิฏฺฐิ.\csromanspacer{} วิจิกิจฺฉา.\csromanspacer{} ถินํ.\csromanspacer{} อุทฺธจฺจํ.\csromanspacer{} อหิริกํ ฯเปฯ อโนตฺตปฺปํ อกุสลํ สารมฺมณํ, อตฺถิ ตสฺส อาวฏฺฏนา ฯเปฯ ปณิธีติ, อามนฺตา.\csromanspacer{} อกุสเลน จิตฺเตน สมุฏฺฐิตํ วจีกมฺมํ รูปํ อกุสลํ สารมฺมณํ, อตฺถิ ตสฺส อาวฏฺฏนา ฯเปฯ ปณิธีติ, น เหวํ วตฺตพฺเพ ฯเปฯ.}
\prose{\csromanfolio{287}\csromansymbolmark{*}อกุสเลน จิตฺเตน สมุฏฺฐิตํ วจีกมฺมํ รูปํ อกุสลํ อนารมฺมณํ, นตฺถิ ตสฺส อาวฏฺฏนา ฯเปฯ ปณิธีติ, อามนฺตา.\csromanspacer{} อกุสเลน จิตฺเตน สมุฏฺฐิโต ผสฺโส อกุสโล อนารมฺมโณ, นตฺถิ ตสฺส อาวฏฺฏนา ฯเปฯ ปณิธีติ, น เหวํ วตฺตพฺเพ ฯเปฯ.\csromanspacer{} อกุสเลน จิตฺเตน สมุฏฺฐิตํ วจีกมฺมํ รูปํ อกุสลํ อนารมฺมณํ, นตฺถิ ตสฺส อาวฏฺฏนา ฯเปฯ ปณิธีติ, อามนฺตา.\csromanspacer{} อกุสเลน จิตฺเตน สมุฏฺฐิตา เวทนา ฯเปฯ สญฺญา.\csromanspacer{} เจตนา.\csromanspacer{} ราโค.\csromanspacer{} โทโส.\csromanspacer{} โมโห.\csromanspacer{} มาโน.\csromanspacer{} ทิฏฺฐิ.\csromanspacer{} วิจิกิจฺฉา.\csromanspacer{} ถินํ.\csromanspacer{} อุทฺธจฺจํ.\csromanspacer{} อหิริกํ ฯเปฯ อโนตฺตปฺปํ อกุสลํ อนารมฺมณํ, นตฺถิ ตสฺส อาวฏฺฏนา ฯเปฯ ปณิธีติ, น เหวํ วตฺตพฺเพ ฯเปฯ.}

\prose{\csromansymbolmark{*}อกุสเลน จิตฺเตน สมุฏฺฐิตํ วจีกมฺมํ รูปํ อกุสลนฺติ, อามนฺตา.\csromanspacer{} ยํ กิญฺจิ อกุสเลน จิตฺเตน สมุฏฺฐิตํ รูปํ สพฺพํ ตํ อกุสลนฺติ, น เหวํ วตฺตพฺเพ ฯเปฯ.}

\proseitem{535}{\csromansymbolmark{*}อกุสเลน จิตฺเตน สมุฏฺฐิตํ วจีกมฺมํ รูปํ อกุสลนฺติ, อามนฺตา.\csromanspacer{} อกุสเลน จิตฺเตน สมุฏฺฐิตํ รูปายตนํ อกุสลนฺติ, น เหวํ วตฺตพฺเพ ฯเปฯ.\csromanspacer{} อกุสเลน จิตฺเตน สมุฏฺฐิตํ วจีกมฺมํ รูปํ อกุสลนฺติ, อามนฺตา.\csromanspacer{} อกุสเลน จิตฺเตน สมุฏฺฐิตํ สทฺทายตนํ ฯเปฯ คนฺธายตนํ.\csromanspacer{} รสายตนํ.\csromanspacer{} โผฏฺฐพฺพายตนํ ฯเปฯ ปถวีธาตุ.\csromanspacer{} อาโปธาตุ.\csromanspacer{} เตโชธาตุ ฯเปฯ วาโยธาตุ อสุจิ อสฺสุ โลหิตํ เสโท อกุสโลติ, น เหวํ วตฺตพฺเพ ฯเปฯ.}

\prose{\csromansymbolmark{*}อกุสเลน จิตฺเตน สมุฏฺฐิตํ รูปายตนํ อพฺยากตนฺติ, อามนฺตา.\csromanspacer{} อกุสเลน จิตฺเตน สมุฏฺฐิตํ วจีกมฺมํ รูปํ อพฺยากตนฺติ, น เหวํ วตฺตพฺเพ ฯเปฯ.\csromanspacer{} อกุสเลน จิตฺเตน สมุฏฺฐิตํ สทฺทายตนํ ฯเปฯ คนฺธายตนํ.\csromanspacer{} รสายตนํ.\csromanspacer{} โผฏฺฐพฺพายตนํ ฯเปฯ ปถวีธาตุ.\csromanspacer{} อาโปธาตุ.\csromanspacer{} เตโชธาตุ ฯเปฯ วาโยธาตุ อสุจิ อสฺสุ โลหิตํ เสโท อพฺยากโตติ, อามนฺตา.\csromanspacer{} อกุสเลน จิตฺเตน สมุฏฺฐิตํ วจีกมฺมํ รูปํ อพฺยากตนฺติ, น เหวํ วตฺตพฺเพ ฯเปฯ.}

\proseitem{536}{\csromanfolio{288}\csromansymbolmark{*}อกุสเลน จิตฺเตน สมุฏฺฐิตํ วจีกมฺมํ รูปํ อนารมฺมณํ อกุสลนฺติ, อามนฺตา.\csromanspacer{} อกุสเลน จิตฺเตน สมุฏฺฐิตํ รูปายตนํ อนารมฺมณํ อกุสลนฺติ, น เหวํ วตฺตพฺเพ ฯเปฯ.\csromanspacer{} อกุสเลน จิตฺเตน สมุฏฺฐิตํ วจีกมฺมํ รูปํ อนารมฺมณํ อกุสลนฺติ, อามนฺตา.\csromanspacer{} อกุสเลน จิตฺเตน สมุฏฺฐิตํ สทฺทายตนํ ฯเปฯ คนฺธายตนํ ฯเปฯ รสายตนํ ฯเปฯ โผฏฺฐพฺพายตนํ.\csromanspacer{} ปถวีธาตุ.\csromanspacer{} อาโปธาตุ.\csromanspacer{} เตโชธาตุ ฯเปฯ วาโยธาตุ อสุจิ อสฺสุ โลหิตํ เสโท อนารมฺมโณ อกุสโลติ, น เหวํ วตฺตพฺเพ ฯเปฯ.}

\prose{\csromansymbolmark{*}อกุสเลน จิตฺเตน สมุฏฺฐิตํ รูปายตนํ อนารมฺมณํ อพฺยากตนฺติ, อามนฺตา.\csromanspacer{} อกุสเลน จิตฺเตน สมุฏฺฐิตํ วจีกมฺมํ รูปํ อนารมฺมณํ อพฺยากตนฺติ, น เหวํ วตฺตพฺเพ ฯเปฯ.\csromanspacer{} อกุสเลน จิตฺเตน สมุฏฺฐิตํ สทฺทายตนํ ฯเปฯ คนฺธายตนํ.\csromanspacer{} รสายตนํ.\csromanspacer{} โผฏฺฐพฺพายตนํ ฯเปฯ ปถวีธาตุ.\csromanspacer{} อาโปธาตุ.\csromanspacer{} เตโชธาตุ ฯเปฯ วาโยธาตุ อสุจิ อสฺสุ โลหิตํ เสโท อนารมฺมโณ อพฺยากโตติ, อามนฺตา.\csromanspacer{} อกุสเลน จิตฺเตน สมุฏฺฐิตํ วจีกมฺมํ รูปํ อนารมฺมณํ อพฺยากตนฺติ, น เหวํ วตฺตพฺเพ ฯเปฯ.}

\proseitem{537}{\csromansymbolmark{*}อกุสเลน จิตฺเตน สมุฏฺฐิตํ วจีกมฺมํ รูปํ ผสฺสวิปฺปยุตฺตํ อกุสลนฺติ, อามนฺตา.\csromanspacer{} อกุสเลน จิตฺเตน สมุฏฺฐิตํ รูปายตนํ ผสฺสวิปฺปยุตฺตํ อกุสลนฺติ, น เหวํ วตฺตพฺเพ ฯเปฯ.\csromanspacer{} อกุสเลน จิตฺเตน สมุฏฺฐิตํ วจีกมฺมํ รูปํ ผสฺสวิปฺปยุตฺตํ อกุสลนฺติ, อามนฺตา.\csromanspacer{} อกุสเลน จิตฺเตน สมุฏฺฐิตํ สทฺทายตนํ ฯเปฯ คนฺธายตนํ.\csromanspacer{} รสายตนํ.\csromanspacer{} โผฏฺฐพฺพายตนํ.\csromanspacer{} ปถวีธาตุ.\csromanspacer{} อาโปธาตุ.\csromanspacer{} เตโชธาตุ ฯเปฯ วาโยธาตุ อสุจิ อสฺสุ โลหิตํ เสโท ผสฺสวิปฺปยุตฺโต อกุสโลติ, น เหวํ วตฺตพฺเพ ฯเปฯ.}

\prose{\csromansymbolmark{*}อกุสเลน จิตฺเตน สมุฏฺฐิตํ รูปายตนํ ผสฺสวิปฺปยุตฺตํ อพฺยากตนฺติ, อามนฺตา.\csromanspacer{} อกุสเลน จิตฺเตน สมุฏฺฐิตํ วจีกมฺมํ รูปํ อนารมฺมณํ ผสฺสวิปฺปยุตฺตํ อพฺยากตนฺติ, น เหวํ วตฺตพฺเพ ฯเปฯ.\csromanspacer{} อกุสเลน จิตฺเตน สมุฏฺฐิตํ สทฺทายตนํ ฯเปฯ คนฺธายตนํ.\csromanspacer{} รสายตนํ.\csromanspacer{} โผฏฺฐพฺพายตนํ.\csromanspacer{} ปถวีธาตุ.\csromanspacer{} อาโปธาตุ.\csromanspacer{} เตโชธาตุ ฯเปฯ วาโยธาตุ อสุจิ อสฺสุ โลหิตํ เสโท อนารมฺมโณ ผสฺสวิปฺปยุตฺโต อพฺยากโตติ, อามนฺตา.\csromanspacer{} อกุสเลน จิตฺเตน สมุฏฺฐิตํ วจีกมฺมํ รูปํ อนารมฺมณํ ผสฺสวิปฺปยุตฺตํ อพฺยากตนฺติ, น เหวํ วตฺตพฺเพ ฯเปฯ.}

\chapterheadpageread{8. อฏฺฐมวคฺค}
\prose{\csromanfolio{289}\csromansymbolmark{*}อกุสเลน จิตฺเตน สมุฏฺฐิตํ วจีกมฺมํ รูปํ อนารมฺมณํ ผสฺสวิปฺปยุตฺตํ อกุสลนฺติ, อามนฺตา.\csromanspacer{} อกุสเลน จิตฺเตน สมุฏฺฐิตํ รูปายตนํ อนารมฺมณํ ผสฺสวิปฺปยุตฺตํ อกุสลนฺติ, น เหวํ วตฺตพฺเพ ฯเปฯ.\csromanspacer{} อกุสเลน จิตฺเตน สมุฏฺฐิตํ วจีกมฺมํ รูปํ อนารมฺมณํ ผสฺสวิปฺปยุตฺตํ อกุสลนฺติ, อามนฺตา.\csromanspacer{} อกุสเลน จิตฺเตน สมุฏฺฐิตํ สทฺทายตนํ ฯเปฯ คนฺธายตนํ.\csromanspacer{} รสายตนํ.\csromanspacer{} โผฏฺฐพฺพายตนํ.\csromanspacer{} ปถวีธาตุ.\csromanspacer{} อาโปธาตุ.\csromanspacer{} เตโชธาตุ ฯเปฯ วาโยธาตุ อสุจิ อสฺสุ โลหิตํ เสโท อนารมฺมโณ ผสฺสวิปฺปยุตฺโต อกุสโลติ, น เหวํ วตฺตพฺเพ ฯเปฯ.}

\prose{\csromansymbolmark{*}อกุสเลน จิตฺเตน สมุฏฺฐิตํ รูปายตนํ อนารมฺมณํ ผสฺสวิปฺปยุตฺตํ อพฺยากตนฺติ, อามนฺตา.\csromanspacer{} อกุสเลน จิตฺเตน สมุฏฺฐิตํ วจีกมฺมํ รูปํ อนารมฺมณํ ผสฺสวิปฺปยุตฺตํ อพฺยากตนฺติ, น เหวํ วตฺตพฺเพ ฯเปฯ.\csromanspacer{} อกุสเลน จิตฺเตน สมุฏฺฐิตํ สทฺทายตนํ ฯเปฯ คนฺธายตนํ.\csromanspacer{} รสายตนํ.\csromanspacer{} โผฏฺฐพฺพายตนํ.\csromanspacer{} ปถวีธาตุ.\csromanspacer{} อาโปธาตุ.\csromanspacer{} เตโชธาตุ ฯเปฯ วาโยธาตุ อสุจิ อสฺสุ โลหิตํ เสโท อนารมฺมโณ ผสฺสวิปฺปยุตฺโต อพฺยากโตติ, อามนฺตา.\csromanspacer{} อกุสเลน จิตฺเตน สมุฏฺฐิตํ วจีกมฺมํ รูปํ อนารมฺมณํ ผสฺสวิปฺปยุตฺตํ อพฺยากตนฺติ, น เหวํ วตฺตพฺเพ ฯเปฯ.}

\proseitem{538}{\csromansymbolmark{+}น วตฺตพฺพํ “รูปํ กุสลมฺปิ อกุสลมฺปี”ติ, อามนฺตา.\csromanspacer{} นนุ กายกมฺมํ วจีกมฺมํ กุสลมฺปิ อกุสลมฺปีติ, อามนฺตา.\csromanspacer{} หญฺจิ กายกมฺมํ วจีกมฺมํ กุสลมฺปิ อกุสลมฺปิ, เตน วต เร วตฺตพฺเพ “รูปํ กุสลมฺปิ อกุสลมฺปี”ติ.}

\prose{\csromansymbolmark{*}รูปํ กุสลมฺปิ อกุสลมฺปีติ, อามนฺตา.\csromanspacer{} จกฺขายตนํ กุสลมฺปิ อกุสลมฺปีติ, น เหวํ วตฺตพฺเพ ฯเปฯ.\csromanspacer{} รูปํ กุสลมฺปิ อกุสลมฺปีติ, อามนฺตา.\csromanspacer{} โสตายตนํ ฯเปฯ ฆานายตนํ.\csromanspacer{} ชิวฺหายตนํ.\csromanspacer{} กายายตนํ.\csromanspacer{} รูปายตนํ.\csromanspacer{} สทฺทายตมํ.\csromanspacer{} คนฺธายตนํ.\csromanspacer{} รสายตนํ.\csromanspacer{} โผฏฺฐพฺพายตนํ.\csromanspacer{} ปถวีธาตุ.\csromanspacer{} อาโปธาตุ.\csromanspacer{} เตโชธาตุ ฯเปฯ วาโยธาตุ อสุจิ อสฺสุ โลหิตํ เสโท กุสโลปิ อกุสโลปีติ, น เหวํ วตฺตพฺเพ ฯเปฯ.}

\prose{\csromansymbolmark{*}กาโย รูปํ กายกมฺมํ รูปนฺติ, อามนฺตา.\csromanspacer{} มโน รูปํ มโนกมฺมํ รูปนฺติ, น เหวํ วตฺตพฺเพ ฯเปฯ.\csromanspacer{} มโน อรูปํ มโนกมฺมํ อรูปนฺติ, อามนฺตา.\csromanspacer{} กาโย อรูปํ กายกมฺมํ อรูปนฺติ, น เหวํ วตฺตพฺเพ ฯเปฯ.}

\prose{\csromanfolio{290}\csromansymbolmark{*}กาโย รูปนฺติ กายกมฺมํ รูปนฺติ, อามนฺตา.\csromanspacer{} จกฺขายตนํ รูปนฺติ จกฺขุวิญฺญาณํ รูปนฺติ, น เหวํ วตฺตพฺเพ ฯเปฯ.\csromanspacer{} กาโย รูปนฺติ กายกมฺมํ รูปนฺติ, อามนฺตา.\csromanspacer{} โสตายตนํ รูปนฺติ โสตวิญฺญาณํ รูปนฺติ, น เหวํ วตฺตพฺเพ ฯเปฯ.\csromanspacer{} กาโย รูปนฺติ กายกมฺมํ รูปนฺติ, อามนฺตา.\csromanspacer{} ฆานายตนํ รูปนฺติ ฆานวิญฺญาณํ รูปนฺติ, น เหวํ วตฺตพฺเพ ฯเปฯ.\csromanspacer{} กาโย รูปนฺติ กายกมฺมํ รูปนฺติ, อามนฺตา.\csromanspacer{} ชิวฺหายตนํ รูปนฺติ ชิวฺหาวิญฺญาณํ รูปนฺติ, น เหวํ วตฺตพฺเพ ฯเปฯ.\csromanspacer{} กาโย รูปนฺติ กายกมฺมํ รูปนฺติ, อามนฺตา.\csromanspacer{} กายายตนํ รูปนฺติ กายวิญฺญาณํ รูปนฺติ, น เหวํ วตฺตพฺเพ ฯเปฯ.}

\proseitem{539}{\csromansymbolmark{*}รูปํ กมฺมนฺติ, อามนฺตา.\csromanspacer{} นนุ วุตฺตํ ภควตา “เจตนาหํ ภิกฺขเว กมฺมํ วทามิ เจตยิตฺวา กมฺมํ กโรติ กาเยน วาจาย มนสา”ติ\footnote{อํ 2. 363 ปิฏฺเฐ นิพฺเพธิกสุตฺเต.} อตฺเถว สุตฺตนฺโตติ, อามนฺตา.\csromanspacer{} เตน หิ น วตฺตพฺพํ “รูปํ กมฺมนฺ”ติ.}

\prose{\csromansymbolmark{*}รูปํ กมฺมนฺติ, อามนฺตา.\csromanspacer{} นนุ วุตฺตํ ภควตา “กาเย วา อานนฺท\footnote{หานนฺท (สํ 1. 275 ปิฏฺเฐ.)} สติ กายสญฺเจตนาเหตุ อุปฺปชฺชติ อชฺฌตฺตํ สุขทุกฺขํ, วาจาย วา อานนฺท2 สติ วจีสญฺเจตนาเหตุ อุปฺปชฺชติ อชฺฌตฺตํ สุขทุกฺขํ, มเน วา อานนฺท2 สติ มโนสญฺเจตนาเหตุ อุปฺปชฺชติ อชฺฌตฺตํ สุขทุกฺขนฺ”ติ\footnote{สํ 1. 275 ปิฏฺเฐ, อํ 1. 476 ปิฏฺเฐปิ.} อตฺเถว สุตฺตนฺโตติ, อามนฺตา.\csromanspacer{} เตน หิ น วตฺตพฺพํ “รูปํ กมฺมนฺ”ติ.}

\prose{\csromansymbolmark{*}รูปํ กมฺมนฺติ, อามนฺตา.\csromanspacer{} นนุ วุตฺตํ ภควตา “ติวิธา ภิกฺขเว กายสญฺเจตนา อกุสลํ กายกมฺมํ ทุกฺขุทฺรยํ ทุกฺขวิปากํ, จตุพฺพิธา ภิกฺขเว วจีสญฺเจตนา อกุสลํ วจีกมฺมํ ทุกฺขุทฺรยํ ทุกฺขวิปากํ, ติวิธา ภิกฺขเว มโนสญฺเจตนา อกุสลํ มโนกมฺมํ ทุกฺขุทฺรยํ ทุกฺขวิปากํ, ติวิธา ภิกฺขเว กายสญฺเจตนา กุสลํ กายกมฺมํ สุขุทฺรยํ สุขวิปากํ, จตุพฺพิธา ภิกฺขเว วจีสญฺเจตนา กุสลํ วจีกมฺมํ สุขุทฺรยํ สุขวิปากํ, ติวิธา ภิกฺขเว มโนสญฺเจตนา กุสลํ มโนกมฺมํ สุขุทฺรยํ สุขวิปากนฺ”ติ\footnote{อํ 3. 497 โถกํ ปน วิสทิสํ.} อตฺเถว สุตฺตนฺโตติ, อามนฺตา.\csromanspacer{} เตน หิ น วตฺตพฺพํ “รูปํ กมฺมนฺ”ติ.}

\prose{\csromansymbolmark{*}รูปํ กมฺมนฺติ, อามนฺตา.\csromanspacer{} นนุ วุตฺตํ ภควตา “สจายํ อานนฺท สมิทฺธิ โมฆปุริโส ปาฏลิปุตฺตสฺส\footnote{โปตลิปุตฺตสฺส (ม 3. 251 ปิฏฺเฐ.)} ปริพฺพาชกสฺส เอวํ ปุฏฺโฐ เอวํ พฺยากเรยฺย}

\vspace{\tipitakaparskip}
\csromancenter{อฏฺฐมวคฺค ‘สญฺเจตนิยํ อาวุโส ปาฏลิปุตฺต กมฺมํ กตฺวา กาเยน วาจาย มนสา สุขเวทนิยํ สุขํ โส เวทยติ.\csromanspacer{} สญฺเจตนิยํ อาวุโส ปาฏลิปุตฺต กมฺมํ กตฺวา กาเยน วาจาย มนสา ทุกฺขเวทนิยํ ทุกฺขํ โส เวทยติ.\csromanspacer{} สญฺเจตนิยํ อาวุโส ปาฏลิปุตฺต กมฺมํ กตฺวา กาเยน วาจาย มนสา อทุกฺขมสุขเวทนิยํ อทุกฺขมสุขํ โส เวทยตี’ติ.\csromanspacer{} เอวํ พฺยากรมาโน โข อานนฺท สมิทฺธิ โมฆปุริโส ปาฏลิปุตฺตสฺส ปริพฺพาชกสฺส สมฺมา พฺยากรมาโน พฺยากเรยฺยา”ติ\footnote{ม 3. 251 ปิฏฺเฐ.} อตฺเถว สุตฺตนฺโตติ, อามนฺตา.\csromanspacer{} เตน หิ น วตฺตพฺพํ “รูปํ กมฺมนฺ”ติ.}
\nitthitam{รูปํ กมฺมนฺติกถา นิฏฺฐิตา.}
\chapterhead{8. อฏฺฐมวคฺค}
\vspace{\tipitakaparskip}
\csromancenter{\textbf{(82) 10.}\csromanspacer{}\textbf{ ชีวิตินฺทฺริยกถา}}
\csromanmatikaanchor{mtk.108}
\csromantocmarkref{subsection}{(82) 10. ชีวิตินฺทฺริยกถา}{mtk.108}
\bookmark[dest={mtk.108},level=2]{(82) 10. ชีวิตินฺทฺริยกถา}
\proseitem{540}{\csromanfolio{291}\csromansymbolmark{*}นตฺถิ รูปชีวิตินฺทฺริยนฺติ, อามนฺตา.\csromanspacer{} นตฺถิ รูปีนํ ธมฺมานํ อายุ ฐิติ ยปนา ยาปนา อิริยนา วตฺตนา ปาลนาติ, น เหวํ วตฺตพฺเพ ฯเปฯ.\csromanspacer{} อตฺถิ รูปีนํ ธมฺมานํ อายุ ฐิติ ยปนา ยาปนา อิริยนา วตฺตนา ปาลนาติ, อามนฺตา.\csromanspacer{} หญฺจิ อตฺถิ รูปีนํ ธมฺมานํ อายุ ฐิติ ยปนา ยาปนา อิริยนา วตฺตนา ปาลนา, โน จ วต เร วตฺตพฺเพ “นตฺถิ รูปชีวิตินฺทฺริยนฺ”ติ.}

\prose{\csromansymbolmark{*}อตฺถิ อรูปีนํ ธมฺมานํ อายุ ฐิติ ยปนา ยาปนา อิริยนา วตฺตนา ปาลนา, อตฺถิ อรูปชีวิตินฺทฺริยนฺติ, อามนฺตา.\csromanspacer{} อตฺถิ รูปีนํ ธมฺมานํ อายุ ฐิติ ยปนา ยาปนา อิริยนา วตฺตนา ปาลนา, อตฺถิ รูปชีวิตินฺทฺริยนฺติ, น เหวํ วตฺตพฺเพ ฯเปฯ.\csromanspacer{} อตฺถิ รูปีนํ ธมฺมานํ อายุ ฐิติ ยปนา ยาปนา อิริยนา วตฺตนา ปาลนา, นตฺถิ รูปชีวิตินฺทฺริยนฺติ, อามนฺตา.\csromanspacer{} อตฺถิ อรูปีนํ ธมฺมานํ อายุ ฐิติ ยปนา ยาปนา อิริยนา วตฺตนา ปาลนา, นตฺถิ อรูปชีวิตินฺทฺริยนฺติ, น เหวํ วตฺตพฺเพ ฯเปฯ.}

\prose{\csromansymbolmark{*}อรูปีนํ ธมฺมานํ อายุ อรูปชีวิตินฺทฺริยนฺติ, อามนฺตา.\csromanspacer{} รูปีนํ ธมฺมานํ อายุ รูปชีวิตินฺทฺริยนฺติ, น เหวํ วตฺตพฺเพ ฯเปฯ.\csromanspacer{} รูปีนํ ธมฺมานํ อายุ น วตฺตพฺพํ “รูปชีวิตินฺทฺริยนฺ”ติ, อามนฺตา.\csromanspacer{} อรูปีนํ ธมฺมานํ อายุ น วตฺตพฺพํ “อรูปชีวิตินฺทฺริยนฺ”ติ, น เหวํ วตฺตพฺเพ ฯเปฯ.}

\proseitem{541}{\csromanfolio{292}\csromansymbolmark{*}รูปีนํ ธมฺมานํ อายุ อรูปชีวิตินฺทฺริยนฺติ, อามนฺตา.\csromanspacer{} อรูปีนํ ธมฺมานํ อายุ รูปชีวิตินฺทฺริยนฺติ.\csromanspacer{} น เหวํ วตฺตพฺเพ ฯเปฯ.}

\prose{\csromansymbolmark{*}อรูปีนํ ธมฺมานํ อายุ น วตฺตพฺพํ “รูปชีวิตินฺทฺริยนฺ”ติ, อามนฺตา.\csromanspacer{} รูปีนํ ธมฺมานํ อายุ น วตฺตพฺพํ “อรูปชีวิตินฺทฺริยนฺ”ติ, น เหวํ วตฺตพฺเพ ฯเปฯ.}

\prose{\csromansymbolmark{*}รูปีนญฺจ อรูปีนญฺจ ธมฺมานํ อายุ อรูปชีวิตินฺทฺริยนฺติ, อามนฺตา.\csromanspacer{} รูปีนญฺจ อรูปีนญฺจ ธมฺมานํ อายุ รูปชีวิตินฺทฺริยนฺติ, น เหวํ วตฺตพฺเพ ฯเปฯ.}

\prose{\csromansymbolmark{*}รูปีนญฺจ อรูปีนญฺจ ธมฺมานํ อายุ น วตฺตพฺพํ “รูปชีวิตินฺทฺริยนฺ”ติ, อามนฺตา.\csromanspacer{} รูปีนญฺจ อรูปีนญฺจ ธมฺมานํ อายุ น วตฺตพฺพํ “อรูปชีวิตินฺทฺริยนฺ”ติ, น เหวํ วตฺตพฺเพ ฯเปฯ.\csromanspacer{} นตฺถิ รูปชีวิตินฺทฺริยนฺติ, อามนฺตา.\csromanspacer{} นิโรธํ สมาปนฺนสฺส นตฺถิ ชีวิตินฺทฺริยนฺติ, น เหวํ วตฺตพฺเพ ฯเปฯ.}

\proseitem{542}{\csromansymbolmark{*}นิโรธํ สมาปนฺนสฺส อตฺถิ ชีวิตินฺทฺริยนฺติ, อามนฺตา.\csromanspacer{} หญฺจิ นิโรธํ สมาปนฺนสฺส อตฺถิ ชีวิตินฺทฺริยํ, โน จ วต เร วตฺตพฺเพ “นตฺถิ รูปชีวิตินฺทฺริยนฺ”ติ.}

\prose{\csromansymbolmark{*}นิโรธํ สมาปนฺนสฺส อตฺถิ ชีวิตินฺทฺริยนฺติ, อามนฺตา.\csromanspacer{} กตมกฺขนฺธปริยาปนฺนนฺติ, สงฺขารกฺขนฺธปริยาปนฺนนฺติ.\csromanspacer{} นิโรธํ สมาปนฺนสฺส อตฺถิ สงฺขารกฺขนฺโธติ, น เหวํ วตฺตพฺเพ ฯเปฯ.}

\prose{\csromansymbolmark{*}นิโรธํ สมาปนฺนสฺส อตฺถิ สงฺขารกฺขนฺโธติ, อามนฺตา.\csromanspacer{} นิโรธํ สมาปนฺนสฺส อตฺถิ เวทนากฺขนฺโธ ฯเปฯ สญฺญากฺขนฺโธ ฯเปฯ วิญฺญาณกฺขนฺโธติ, น เหวํ วตฺตพฺเพ ฯเปฯ.}

\prose{\csromansymbolmark{*}นิโรธํ สมาปนฺนสฺส อตฺถิ เวทนากฺขนฺโธ ฯเปฯ สญฺญากฺขนฺโธ ฯเปฯ วิญฺญาณกฺขนฺโธติ, อามนฺตา.\csromanspacer{} น นิโรธํ สมาปนฺโนติ.\csromanspacer{} น เหวํ วตฺตพฺเพ ฯเปฯ.}

\proseitem{543}{\csromansymbolmark{*}นตฺถิ รูปชีวิตินฺทฺริยนฺติ, อามนฺตา.\csromanspacer{} อสญฺญสตฺตานํ นตฺถิ ชีวิตินฺทฺริยนฺติ, น เหวํ วตฺตพฺเพ ฯเปฯ.}

\prose{\csromansymbolmark{*}อสญฺญสตฺตานํ อตฺถิ ชีวิตินฺทฺริยนฺติ, อามนฺตา.\csromanspacer{} หญฺจิ อสญฺญสตฺตานํ อตฺถิ ชีวิตินฺทฺริยํ, โน จ วต เร วตฺตพฺเพ “นตฺถิ รูปชีวิตินฺทฺริยนฺ”ติ.\csromanspacer{} อสญฺญสตฺตานํ อตฺถิ ชีวิตินฺทฺริยนฺติ, อามนฺตา.\csromanspacer{} กตมกฺขนฺธปริยาปนฺนนฺติ, สงฺขารกฺขนฺธปริยาปนฺนนฺติ.\csromanspacer{} อสญฺญสตฺตานํ อตฺถิ สงฺขารกฺขนฺโธติ, น เหวํ วตฺตพฺเพ ฯเปฯ.}

\chapterheadpageread{8. อฏฺฐมวคฺค}
\csromanmatikaanchor{mtk.109}
\csromantocmarkref{subsection}{(83) 11. กมฺมเหตุกถา}{mtk.109}
\bookmark[dest={mtk.109},level=2]{(83) 11. กมฺมเหตุกถา}
\prose{\csromanfolio{293}\csromansymbolmark{*}อสญฺญสตฺตานํ อตฺถิ สงฺขารกฺขนฺโธติ, อามนฺตา.\csromanspacer{} อสญฺญสตฺตานํ อตฺถิ เวทนากฺขนฺโธ ฯเปฯ สญฺญากฺขนฺโธ ฯเปฯ วิญฺญาณกฺขนฺโธติ, น เหวํ วตฺตพฺเพ ฯเปฯ.}

\prose{\csromansymbolmark{*}อสญฺญสตฺตานํ อตฺถิ เวทนากฺขนฺโธ ฯเปฯ สญฺญากฺขนฺโธ ฯเปฯ วิญฺญาณกฺขนฺโธติ, อามนฺตา.\csromanspacer{} ปญฺจโวการภโวติ, น เหวํ วตฺตพฺเพ ฯเปฯ.}

\proseitem{544}{\csromansymbolmark{*}อุปปตฺเตสิเยน จิตฺเตน สมุฏฺฐิตํ ชีวิตินฺทฺริยํ อุปปตฺเตสิเย จิตฺเต ภิชฺชมาเน เอกเทสํ ภิชฺชตีติ, อามนฺตา.\csromanspacer{} อุปปตฺเตสิเยน จิตฺเตน สมุฏฺฐิโต ผสฺโส อุปปตฺเตสิเย จิตฺเต ภิชฺชมาเน เอกเทโส ภิชฺชตีติ, น เหวํ วตฺตพฺเพ ฯเปฯ.}

\prose{\csromansymbolmark{*}อุปปตฺเตสิเยน จิตฺเตน สมุฏฺฐิโต ผสฺโส อุปปตฺเตสิเย จิตฺเต ภิชฺชมาเน อนวเสโส ภิชฺชตีติ, อามนฺตา.\csromanspacer{} อุปปตฺเตสิเยน จิตฺเตน สมุฏฺฐิตํ ชีวิตินฺทฺริยํ อุปปตฺเตสิเย จิตฺเต ภิชฺชมาเน อนวเสสํ ภิชฺชตีติ, น เหวํ วตฺตพฺเพ ฯเปฯ.}

\proseitem{545}{\csromansymbolmark{+}เทฺว ชีวิตินฺทฺริยานีติ, อามนฺตา.\csromanspacer{} ทฺวีหิ ชีวิเตหิ ชีวติ, ทฺวีหิ มรเณหิ มียตีติ, อามนฺตา\footnote{น เหวํ วตฺตพฺเพ (สฺยา, กํ, อิ)}.}

\nitthitam{ชีวิตินฺทฺริยกถา นิฏฺฐิตา.}
\chapterhead{8. อฏฺฐมวคฺค}
\vspace{\tipitakaparskip}
\csromancenter{\textbf{(83) 11.}\csromanspacer{}\textbf{ กมฺมเหตุกถา}}
\proseitem{546}{\csromansymbolmark{*}กมฺมเหตุ อรหา อรหตฺตา ปริหายตีติ, อามนฺตา.\csromanspacer{} กมฺมเหตุ โสตาปนฺโน โสตาปตฺติผลา ปริหายตีติ, น เหวํ วตฺตพฺเพ ฯเปฯ.\csromanspacer{} กมฺมเหตุ อรหา อรหตฺตา ปริหายตีติ, อามนฺตา.\csromanspacer{} กมฺมเหตุ สกทาคามี ฯเปฯ อนาคามี อนาคามิผลา ปริหายตีติ, น เหวํ วตฺตพฺเพ ฯเปฯ.}

\csromanmatikaanchor{mtk.110}
\csromanmatikaanchor{mtk.111}
\csromantocmarknopage{chapter}{9. นวมวคฺค}{mtk.110}
\bookmark[dest={mtk.110},level=0]{9. นวมวคฺค}
\csromantocmarkref{subsection}{(84) 1. อานิสํสทสฺสาวีกถา}{mtk.111}
\bookmark[dest={mtk.111},level=2]{(84) 1. อานิสํสทสฺสาวีกถา}
\prose{\csromansymbolmark{*}กมฺมเหตุ โสตาปนฺโน โสตาปตฺติผลา น ปริหายตีติ, อามนฺตา.\csromanspacer{} กมฺมเหตุ อรหา อรหตฺตา น ปริหายตีติ, น เหวํ \csromanfolio{294}วตฺตพฺเพ ฯเปฯ.\csromanspacer{} กมฺมเหตุ สกทาคามี ฯเปฯ อนาคามี อนาคามิผลา น ปริหายตีติ, อามนฺตา.\csromanspacer{} กมฺมเหตุ อรหา อรหตฺตา น ปริหายตีติ, น เหวํ วตฺตพฺเพ ฯเปฯ.}

\prose{\csromansymbolmark{*}กมฺมเหตุ อรหา อรหตฺตา ปริหายตีติ, อามนฺตา.\csromanspacer{} ปาณาติปาตกมฺมสฺส เหตูติ, น เหวํ วตฺตพฺเพ ฯเปฯ.\csromanspacer{} อทินฺนาทานกมฺมสฺส เหตุ ฯเปฯ.\csromanspacer{} กาเมสุมิจฺฉาจารกมฺมสฺส เหตุ.\csromanspacer{} มุสาวาทกมฺมสฺส เหตุ.\csromanspacer{} ปิสุณวาจากมฺมสฺส เหตุ.\csromanspacer{} ผรุสวาจากมฺมสฺส เหตุ.\csromanspacer{} สมฺผปฺปลาปกมฺมสฺส เหตุ.\csromanspacer{} มาตุฆาตกมฺมสฺส\footnote{มาตุฆาตกกมฺมสฺส (สฺยา), มาตุฆาติกมฺมสฺส (ก) ตถา ปิตุฆาต- อรหนฺตฆาตกมฺมสฺสาติปเทสุ.} เหตุ.\csromanspacer{} ปิตุฆาตกมฺมสฺส เหตุ.\csromanspacer{} อรหนฺตฆาตกมฺมสฺส เหตุ.\csromanspacer{} รุหิรุปฺปาทกมฺมสฺส เหตุ ฯเปฯ.\csromanspacer{} สํฆเภทกมฺมสฺส เหตูติ, น เหวํ วตฺตพฺเพ ฯเปฯ.}

\prose{\csromansymbolmark{*}กตมสฺส กมฺมสฺส เหตูติ, หนฺท หิ อรหนฺตานํ อพฺภาจิกฺขตีติ.\csromanspacer{} อรหนฺตานํ อพฺภาจิกฺขนกมฺมสฺส เหตุ อรหา อรหตฺตา ปริหายตีติ, อามนฺตา.\csromanspacer{} เย เกจิ อรหนฺตานํ อพฺภาจิกฺขนฺติ, สพฺเพ เต อรหตฺตํ สจฺฉิกโรนฺตีติ, น เหวํ วตฺตพฺเพ ฯเปฯ.\csromanspacer{} กมฺมเหตุกถา นิฏฺฐิตา.\csromanspacer{} อฏฺฐโม วคฺโค.}

\titlehead{\textbf{ตสฺสุทฺทานํ}}
\prose{ฉ คติโย, อนฺตราภโว, ปญฺเจว กามคุณา กามธาตุ, ปญฺเจว อายตนา กามา, รูปิโน ธมฺมา รูปธาตุ, อรูปิโน ธมฺมา อรูปธาตุ, สฬายตนิโก อตฺตภาโว รูปธาตุยา, อตฺถิ รูปํ อรูเปสุ, รูปํ กมฺมํ, รูปํ ชีวิตํ, กมฺมเหตุกา ปริหายตีติ.}

\chapterhead{9. นวมวคฺค}
\vspace{\tipitakaparskip}
\csromancenter{\textbf{(84) 1.}\csromanspacer{}\textbf{ อานิสํสทสฺสาวีกถา}}
\proseitem{547}{\csromansymbolmark{*}อานิสํสทสฺสาวิสฺส สํโยชนานํ ปหานนฺติ, อามนฺตา.\csromanspacer{} นนุ สงฺขาเร อนิจฺจโต มนสิกโรโต สํโยชนา ปหียนฺตีติ, อามนฺตา.}

\vspace{\tipitakaparskip}
\csromancenter{นวมวคฺค หญฺจิ สงฺขาเร อนิจฺจโต มนสิกโรโต สํโยชนา ปหียนฺติ, โน จ วต เร วตฺตพฺเพ “อานิสํสทสฺสาวิสฺส สํโยชนานํ ปหานนฺ”ติ.}
\csromanmatikaanchor{mtk.112}
\csromantocmarkref{subsection}{(85) 2. อมตารมฺมณกถา}{mtk.112}
\bookmark[dest={mtk.112},level=2]{(85) 2. อมตารมฺมณกถา}
\prose{\csromanfolio{295}\csromansymbolmark{*}นนุ สงฺขาเร ทุกฺขโต ฯเปฯ โรคโต.\csromanspacer{} คณฺฑโต.\csromanspacer{} สลฺลโต.\csromanspacer{} อฆโต.\csromanspacer{} อาพาธโต.\csromanspacer{} ปรโต.\csromanspacer{} ปโลกโต. ¢ติโต.\csromanspacer{} อุปทฺทวโต.\csromanspacer{} ภยโต.\csromanspacer{} อุปสคฺคโต.\csromanspacer{} จลโต.\csromanspacer{} ปภงฺคุโต.\csromanspacer{} อทฺธุวโต.\csromanspacer{} อตาณโต.\csromanspacer{} อเลณโต.\csromanspacer{} อสรณโต.\csromanspacer{} อสรณีภูตโต.\csromanspacer{} ริตฺตโต.\csromanspacer{} ตุจฺฉโต.\csromanspacer{} สุญฺญโต.\csromanspacer{} อนตฺตโต.\csromanspacer{} อาทีนวโต ฯเปฯ วิปริณามธมฺมโต มนสิกโรโต สํโยชนา ปหียนฺตีติ, อามนฺตา.\csromanspacer{} หญฺจิ สงฺขาเร วิปริณามธมฺมโต มนสิกโรโต สํโยชนา ปหียนฺติ, โน จ วต เร วตฺตพฺเพ “อานิสํสทสฺสาวิสฺส สํโยชนานํ ปหานนฺ”ติ.}

\prose{\csromansymbolmark{*}สงฺขาเร จ อนิจฺจโต มนสิ กโรติ นิพฺพาเน จ อานิสํสทสฺสาวี โหตีติ, น เหวํ วตฺตพฺเพ ฯเปฯ.\csromanspacer{} สงฺขาเร จ อนิจฺจโต มนสิ กโรติ นิพฺพาเน จ อานิสํสทสฺสาวี โหตีติ, อามนฺตา.\csromanspacer{} ทฺวินฺนํ ผสฺสานํ ฯเปฯ.\csromanspacer{} ทฺวินฺนํ จิตฺตานํ สโมธานํ โหตีติ, น เหวํ วตฺตพฺเพ ฯเปฯ.\csromanspacer{} สงฺขาเร จ ทุกฺขโต ฯเปฯ วิปริณามธมฺมโต มนสิ กโรติ นิพฺพาเน จ อานิสํสทสฺสาวี โหตีติ, น เหวํ วตฺตพฺเพ ฯเปฯ.\csromanspacer{} สงฺขาเร จ วิปริณามธมฺมโต มนสิ กโรติ นิพฺพาเน จ อานิสํสทสฺสาวี โหตีติ, อามนฺตา.\csromanspacer{} ทฺวินฺนํ ผสฺสานํ ฯเปฯ.\csromanspacer{} ทฺวินฺนํ จิตฺตานํ สโมธานํ โหตีติ, น เหวํ วตฺตพฺเพ ฯเปฯ.}

\proseitem{548}{\csromansymbolmark{+}น วตฺตพฺพํ “อานิสํสทสฺสาวิสฺส สํโยชนานํ ปหานนฺ”ติ, อามนฺตา.\csromanspacer{} นนุ วุตฺตํ ภควตา “อิธ ภิกฺขเว ภิกฺขุ นิพฺพาเน สุขานุปสฺสี วิหรติ, สุขสญฺญี สุขปฏิสํเวที สตตํ สมิตํ อพฺโพกิณฺณํ เจตสา อธิมุจฺจมาโน ปญฺญาย ปริโยคาหมาโน”ติ\footnote{อํ 2. 406 ปิฏฺเฐ.} อตฺเถว สุตฺตนฺโตติ, อามนฺตา.\csromanspacer{} เตน หิ อานิสํสทสฺสาวิสฺส สํโยชนานํ ปหานนฺติ.}

\nitthitam{อานิสํสทสฺสาวีกถา นิฏฺฐิตา.}
\chapterhead{9. นวมวคฺค}
\vspace{\tipitakaparskip}
\csromancenter{\textbf{(85) 2.}\csromanspacer{}\textbf{ อมตารมฺมณกถา}}
\proseitem{549}{\csromansymbolmark{*}อมตารมฺมณํ สํโยชนนฺติ, อามนฺตา.\csromanspacer{} อมตํ สํโยชนิยํ คนฺถนิยํ โอฆนิยํ โยคนิยํ นีวรณิยํ ปรามฏฺฐํ อุปาทานิยํ สํกิเลสิยนฺติ, \csromanfolio{296}น เหวํ วตฺตพฺเพ ฯเปฯ.\csromanspacer{} นนุ อมตํ อสํโยชนิยํ อคนฺถนิยํ ฯเปฯ อสํกิเลสิยนฺติ, อามนฺตา.\csromanspacer{} หญฺจิ อมตํ อสํโยชนิยํ ฯเปฯ อสํกิเลสิยํ, โน จ วต เร วตฺตพฺเพ “อมตารมฺมณํ สํโยชนนฺ”ติ.}

\prose{\csromansymbolmark{*}อมตํ อารพฺภ ราโค อุปฺปชฺชตีติ, อามนฺตา.\csromanspacer{} อมตํ ราคฏฺฐานิยํ รชนิยํ กมนิยํ มทนิยํ พนฺธนิยํ มุจฺฉนิยนฺติ, น เหวํ วตฺตพฺเพ ฯเปฯ.\csromanspacer{} นนุ อมตํ น ราคฏฺฐานิยํ น รชนิยํ น กมนิยํ น มทนิยํ น พนฺธนิยํ น มุจฺฉนิยนฺติ, อามนฺตา.\csromanspacer{} หญฺจิ อมตํ น ราคฏฺฐานิยํ น รชนิยํ น กมนิยํ น มทนิยํ น พนฺธนิยํ น มุจฺฉนิยํ, โน จ วต เร วตฺตพฺเพ “อมตํ อารพฺภ ราโค อุปฺปชฺชตี”ติ.}

\prose{\csromansymbolmark{*}อมตํ อารพฺภ โทโส อุปฺปชฺชตีติ, อามนฺตา.\csromanspacer{} อมตํ โทสฏฺฐานิยํ โกปฏฺฐานิยํ ปฏิฆฏฺฐานิยนฺติ, น เหวํ วตฺตพฺเพ ฯเปฯ.\csromanspacer{} นนุ อมตํ น โทสฏฺฐานิยํ น โกปฏฺฐานิยํ น ปฏิฆฏฺฐานิยนฺติ, อามนฺตา.\csromanspacer{} หญฺจิ อมตํ น โทสฏฺฐานิยํ น โกปฏฺฐานิยํ น ปฏิฆฏฺฐานิยํ, โน จ วต เร วตฺตพฺเพ “อมตํ อารพฺภ โทโส อุปฺปชฺชตี”ติ.}

\prose{\csromansymbolmark{*}อมตํ อารพฺภ โมโห อุปฺปชฺชตีติ, อามนฺตา.\csromanspacer{} อมตํ โมหฏฺฐานิยํ อญฺญาณกรณํ อจกฺขุกรณํ ปญฺญานิโรธิยํ วิฆาตปกฺขิยํ อนิพฺพานสํวตฺตนิยนฺติ, น เหวํ วตฺตพฺเพ ฯเปฯ.\csromanspacer{} นนุ อมตํ น โมหฏฺฐานิยํ น อญฺญาณกรณํ น อจกฺขุกรณํ ปญฺญาพุทฺธิยํ อวิฆาตปกฺขิยํ นิพฺพานสํวตฺตนิยนฺติ, อามนฺตา.\csromanspacer{} หญฺจิ อมตํ น โมหฏฺฐานิยํ น อญฺญาณกรณํ ฯเปฯ นิพฺพานสํวตฺตนิยํ, โน จ วต เร วตฺตพฺเพ “อมตํ อารพฺภ โมโห อุปฺปชฺชตี”ติ.}

\proseitem{550}{\csromansymbolmark{*}รูปํ อารพฺภ สํโยชนา อุปฺปชฺชนฺติ, รูปํ สํโยชนิยํ คนฺถนิยํ ฯเปฯ สํกิเลสิยนฺติ, อามนฺตา.\csromanspacer{} อมตํ อารพฺภ สํโยชนา อุปฺปชฺชนฺติ, อมตํ สํโยชนิยํ ฯเปฯ สํกิเลสิยนฺติ, น เหวํ วตฺตพฺเพ ฯเปฯ.}

\prose{\csromansymbolmark{*}รูปํ อารพฺภ ราโค อุปฺปชฺชติ, รูปํ ราคฏฺฐานิยํ รชนิยํ กมนิยํ มทนิยํ พนฺธนิยํ มุจฺฉนิยนฺติ, อามนฺตา.\csromanspacer{} อมตํ อารพฺภ ราโค อุปฺปชฺชติ, อมตํ ราคฏฺฐานิยํ ฯเปฯ มุจฺฉนิยนฺติ, น เหวํ วตฺตพฺเพ ฯเปฯ.}

\prose{\csromansymbolmark{*}รูปํ อารพฺภ โทโส อุปฺปชฺชติ, รูปํ โทสฏฺฐานิยํ โกปฏฺฐานิยํ ปฏิฆฏฺฐานิยนฺติ, อามนฺตา.\csromanspacer{} อมตํ อารพฺภ โทโส อุปฺปชฺชติ, อมตํ โทสฏฺฐานิยํ โกปฏฺฐานิยํ ปฏิฆฏฺฐานิยนฺติ, น เหวํ วตฺตพฺเพ ฯเปฯ.}

\chapterheadpageread{9. นวมวคฺค}
\prose{\csromanfolio{297}\csromansymbolmark{*}รูปํ อารพฺภ โมโห อุปฺปชฺชติ, รูปํ โมหฏฺฐานิยํ อญฺญาณกรณํ ฯเปฯ อนิพฺพานสํวตฺตนิยนฺติ, อามนฺตา.\csromanspacer{} อมตํ อารพฺภ โมโห อุปฺปชฺชติ, อมตํ โมหฏฺฐานิยํ อญฺญาณกรณํ ฯเปฯ อนิพฺพานสํวตฺตนิยนฺติ, น เหวํ วตฺตพฺเพ ฯเปฯ.}

\prose{\csromansymbolmark{*}อมตํ อารพฺภ สํโยชนา อุปฺปชฺชนฺติ, อมตํ อสํโยชนิยํ อคนฺถนิยํ อโนฆนิยํ อโยคนิยํ อนีวรณิยํ อปรามฏฺฐํ อนุปาทานิยํ อสํกิเลสิยนฺติ, อามนฺตา.\csromanspacer{} รูปํ อารพฺภ สํโยชนา อุปฺปชฺชนฺติ, รูปํ อสํโยชนิยํ อคนฺถนิยํ ฯเปฯ อสํกิเลสิยนฺติ, น เหวํ วตฺตพฺเพ ฯเปฯ.}

\prose{\csromansymbolmark{*}อมตํ อารพฺภ ราโค อุปฺปชฺชติ, อมตํ น ราคฏฺฐานิยํ น รชนิยํ น กมนิยํ น มทนิยํ น พนฺธนิยํ น มุจฺฉนิยนฺติ, อามนฺตา.\csromanspacer{} รูปํ อารมฺภ ราโค อุปฺปชฺชติ, รูปํ น ราคฏฺฐานิยํ ฯเปฯ น มุจฺฉนิยนฺติ, น เหวํ วตฺตพฺเพ ฯเปฯ.}

\prose{\csromansymbolmark{*}อมตํ อารพฺภ โทโส อุปฺปชฺชติ, อมตํ น โทสฏฺฐานิยํ น โกปฏฺฐานิยํ น ปฏิฆฏฺฐานิยนฺติ, อามนฺตา.\csromanspacer{} รูปํ อารพฺภ โทโส อุปฺปชฺชติ, รูปํ น โทสฏฺฐานิยํ น โกปฏฺฐานิยํ น ปฏิฆฏฺฐานิยนฺติ, น เหวํ วตฺตพฺเพ ฯเปฯ.}

\prose{\csromansymbolmark{*}อมตํ อารพฺภ โมโห อุปฺปชฺชติ, อมตํ น โมหฏฺฐานิยํ น อญฺญาณกรณํ ฯเปฯ นิพฺพานสํวตฺตนิยนฺติ, อามนฺตา.\csromanspacer{} รูปํ อารพฺภ โมโห อุปฺปชฺชติ, รูปํ น โมหฏฺฐานิยํ น อญฺญาณกรณํ ฯเปฯ นิพฺพานสํวตฺตนิยนฺติ, น เหวํ วตฺตพฺเพ ฯเปฯ.}

\proseitem{551}{\csromansymbolmark{+}น วตฺตพฺพํ “อมตารมฺมณํ สํโยชนนฺ”ติ, อามนฺตา.\csromanspacer{} นนุ วุตฺตํ ภควตา “นิพฺพานํ นิพฺพานโต สญฺชานาติ, นิพฺพานํ นิพฺพานโต สญฺชานิตฺวา นิพฺพานํ มญฺญติ, นิพฺพานสฺมิํ มญฺญติ, นิพฺพานโต มญฺญติ, นิพฺพานํ เมติ มญฺญติ, นิพฺพานํ อภินนฺทตี”ติ\footnote{ม 1. 4 ปิฏฺเฐ.} อตฺเถว สุตฺตนฺโตติ, อามนฺตา.\csromanspacer{} เตน หิ อมตารมฺมณํ สํโยชนนฺติ.}

\nitthitam{อมตารมฺมณกถา นิฏฺฐิต.}
\chapterheadpageread{9. นวมวคฺค}
\vspace{\tipitakaparskip}
\csromancenter{\textbf{(86) 3.}\csromanspacer{}\textbf{ รูปํสารมฺมณนฺติกถา}}
\csromanmatikaanchor{mtk.113}
\csromantocmarkref{subsection}{(86) 3. รูปํสารมฺมณนฺติกถา}{mtk.113}
\bookmark[dest={mtk.113},level=2]{(86) 3. รูปํสารมฺมณนฺติกถา}
\proseitem{552}{\csromanfolio{298}\csromansymbolmark{*}รูปํ สารมฺมณนฺติ, อามนฺตา.\csromanspacer{} อตฺถิ ตสฺส อาวฏฺฏนา อาโภโค สมนฺนาหาโร มนสิกาโร เจตนา ปตฺถนา ปณิธีติ, น เหวํ วตฺตพฺเพ ฯเปฯ.\csromanspacer{} นนุ นตฺถิ ตสฺส อาวฏฺฏนา อาโภโค ฯเปฯ ปณิธีติ, อามนฺตา.\csromanspacer{} หญฺจิ นตฺถิ ตสฺส อาวฏฺฏนา อาโภโค ฯเปฯ ปณิธิ, โน จ วต เร วตฺตพฺเพ “รูปํ สารมฺมณนฺ”ติ.}

\prose{\csromansymbolmark{*}ผสฺโส สารมฺมโณ, อตฺถิ ตสฺส อาวฏฺฏนา อาโภโค ฯเปฯ ปณิธีติ, อามนฺตา.\csromanspacer{} รูปํ สารมฺมณํ อตฺถิ ตสฺส อาวฏฺฏนา อาโภโค ฯเปฯ ปณิธีติ, น เหวํ วตฺตพฺเพ ฯเปฯ.}

\prose{\csromansymbolmark{*}เวทนา ฯเปฯ.\csromanspacer{} สญฺญา เจตนา จิตฺตํ สทฺธา วีริยํ สติ สมาธิ ปญฺญา.\csromanspacer{} ราโค โทโส โมโห มาโน ทิฏฺฐิ วิจิกิจฺฉา ถินํ อุทฺธจฺจํ อหิริกํ ฯเปฯ.\csromanspacer{} อโนตฺตปฺปํ สารมฺมณํ, อตฺถิ ตสฺส อาวฏฺฏนา อาโภโค ฯเปฯ ปณิธีติ, อามนฺตา.\csromanspacer{} รูปํ สารมฺมณํ อตฺถิ ตสฺส อาวฏฺฏนา อาโภโค ฯเปฯ ปณิธีติ, น เหวํ วตฺตพฺเพ ฯเปฯ.}

\prose{\csromansymbolmark{*}รูปํ สารมฺมณํ, นตฺถิ ตสฺส อาวฏฺฏนา อาโภโค ฯเปฯ ปณิธีติ, อามนฺตา.\csromanspacer{} ผสฺโส สารมฺมโณ, นตฺถิ ตสฺส อาวฏฺฏนา อาโภโค ฯเปฯ ปณิธีติ, น เหวํ วตฺตพฺเพ ฯเปฯ.}

\prose{\csromansymbolmark{*}รูปํ สารมฺมณํ, นตฺถิ ตสฺส อาวฏฺฏนา อาโภโค ฯเปฯ ปณิธีติ, อามนฺตา.\csromanspacer{} เวทนา ฯเปฯ.\csromanspacer{} อโนตฺตปฺปํ สารมฺมณํ, นตฺถิ ตสฺส อาวฏฺฏนา อาโภโค ฯเปฯ ปณิธีติ, น เหวํ วตฺตพฺเพ ฯเปฯ.}

\proseitem{553}{\csromansymbolmark{+}น วตฺตพฺพํ รูปํ สารมฺมณนฺติ, อามนฺตา.\csromanspacer{} นนุ รูปํ สปฺปจฺจยนฺติ, อามนฺตา.\csromanspacer{} หญฺจิ รูปํ สปฺปจฺจยํ, เตน วต เร วตฺตพฺเพ “รูปํ สารมฺมณนฺ”ติ.}

\nitthitam{รูปํ สารมฺมณนฺติกถา นิฏฺฐิตา.}
\chapterheadpageread{9. นวมวคฺค \textbf{9. นวมวคฺค}}
\vspace{\tipitakaparskip}
\csromancenter{\textbf{(87) 4.}\csromanspacer{}\textbf{ อนุสยา-อนารมฺมณกถา}}
\csromanmatikaanchor{mtk.114}
\csromantocmarkref{subsection}{(87) 4. อนุสยา-อนารมฺมณกถา}{mtk.114}
\bookmark[dest={mtk.114},level=2]{(87) 4. อนุสยา-อนารมฺมณกถา}
\proseitem{554}{\csromanfolio{299}อนุสยา อนารมฺมณาติ, อามนฺตา.\csromanspacer{} รูปํ นิพฺพานํ จกฺขายตนํ ฯเปฯ.\csromanspacer{} โผฏฺฐพฺพายตนนฺติ, น เหวํ วตฺตพฺเพ ฯเปฯ.}

\prose{\csromansymbolmark{*}กามราคานุสโย อนารมฺมโณติ, อามนฺตา.\csromanspacer{} กามราโค กามราคปริยุฏฺฐานํ กามราคสํโยชนํ กาโมโฆ กามโยโค กามจฺฉนฺทนีวรณํ อนารมฺมณนฺติ, น เหวํ วตฺตพฺเพ ฯเปฯ.\csromanspacer{} กามราโค กามราคปริยุฏฺฐานํ กามราคสํโยชนํ กาโมโฆ กามโยโค กามจฺฉนฺทนีวรณํ สารมฺมณนฺติ, อามนฺตา.\csromanspacer{} กามราคานุสโย สารมฺมโณติ, น เหวํ วตฺตพฺเพ ฯเปฯ.}

\prose{\csromansymbolmark{*}กามราคานุสโย อนารมฺมโณติ, อามนฺตา.\csromanspacer{} กตมกฺขนฺธปริยาปนฺโนติ? สงฺขารกฺขนฺธปริยาปนฺโนติ.\csromanspacer{} สงฺขารกฺขนฺโธ อนารมฺมโณติ, น เหวํ วตฺตพฺเพ ฯเปฯ.\csromanspacer{} สงฺขารกฺขนฺโธ อนารมฺมโณติ, อามนฺตา.\csromanspacer{} เวทนากฺขนฺโธ สญฺญากฺขนฺโธ วิญฺญาณกฺขนฺโธ อนารมฺมโณติ, น เหวํ วตฺตพฺเพ ฯเปฯ.}

\prose{\csromansymbolmark{*}กามราคานุสโย สงฺขารกฺขนฺธปริยาปนฺโน อนารมฺมโณติ, อามนฺตา.\csromanspacer{} กามราโค สงฺขารกฺขนฺธปริยาปนฺโน อนารมฺมโณติ, น เหวํ วตฺตพฺเพ ฯเปฯ.\csromanspacer{} กามราโค สงฺขารกฺขนฺธปริยาปนฺโน สารมฺมโณติ, อามนฺตา.\csromanspacer{} กามราคานุสโย สงฺขารกฺขนฺธปริยาปนฺโน สารมฺมโณติ, น เหวํ วตฺตพฺเพ ฯเปฯ.}

\prose{\csromansymbolmark{*}กามราคานุสโย สงฺขารกฺขนฺธปริยาปนฺโน อนารมฺมโณ, กามราโค สงฺขารกฺขนฺธปริยาปนฺโน สารมฺมโณติ, อามนฺตา.\csromanspacer{} สงฺขารกฺขนฺโธ เอกเทโส สารมฺมโณ เอกเทโส อนารมฺมโณติ, น เหวํ วตฺตพฺเพ ฯเปฯ.\csromanspacer{} สงฺขารกฺขนฺโธ เอกเทโส สารมฺมโณ เอกเทโส อนารมฺมโณติ, อามนฺตา.\csromanspacer{} เวทนากฺขนฺโธ สญฺญากฺขนฺโธ วิญฺญาณกฺขนฺโธ เอกเทโส สารมฺมโณ เอกเทโส อนารมฺมโณติ, น เหวํ วตฺตพฺเพ ฯเปฯ.}

\csromanmatikaanchor{mtk.115}
\csromantocmarkref{subsection}{(88) 5. ญาณํอนารมฺมณนฺติกถา}{mtk.115}
\bookmark[dest={mtk.115},level=2]{(88) 5. ญาณํอนารมฺมณนฺติกถา}
\proseitem{555}{\csromansymbolmark{*}ปฏิฆานุสโย มานานุสโย ทิฏฺฐานุสโย วิจิกิจฺฉานุสโย ภวราคานุสโย อวิชฺชานุสโย อนารมฺมโณติ, อามนฺตา.\csromanspacer{} อวิชฺชา อวิชฺโชโฆ อวิชฺชาโยโค อวิชฺชานุสโย อวิชฺชาปริยุฏฺฐานํ \csromanfolio{300}อวิชฺชาสํโยชนํ อวิชฺชานีวรณํ อนารมฺมณนฺติ, น เหวํ วตฺตพฺเพ ฯเปฯ.\csromanspacer{} อวิชฺชา อวิชฺโชโฆ ฯเปฯ อวิชฺชานีวรณํ สารมฺมณนฺติ, อามนฺตา.\csromanspacer{} อวิชฺชานุสโย สารมฺมโณติ, น เหวํ วตฺตพฺเพ ฯเปฯ.}

\prose{\csromansymbolmark{*}อวิชฺชานุสโย อนารมฺมโณติ, อามนฺตา.\csromanspacer{} กตมกฺขนฺธปริยาปนฺโนติ? สงฺขารกฺขนฺธปริยาปนฺโนติ.\csromanspacer{} สงฺขารกฺขนฺโธ อนารมฺมโณติ, น เหวํ วตฺตพฺเพ ฯเปฯ.\csromanspacer{} สงฺขารกฺขนฺโธ อนารมฺมโณติ, อามนฺตา.\csromanspacer{} เวทนากฺขนฺโธ สญฺญากนฺโธ วิญฺญาณกฺขนฺโธ อนารมฺมโณติ, น เหวํ วตฺตพฺเพ ฯเปฯ.}

\prose{\csromansymbolmark{*}อวิชฺชานุสโย สงฺขารกฺขนฺธปริยาปนฺโน อนารมฺมโณติ, อามนฺตา.\csromanspacer{} อวิชฺชา สงฺขารกฺขนฺธปริยาปนฺนา อนารมฺมณาติ, น เหวํ วตฺตพฺเพ ฯเปฯ.\csromanspacer{} อวิชฺชา สงฺขารกฺขนฺธปริยาปนฺนา สารมฺมณาติ, อามนฺตา.\csromanspacer{} อวิชฺชานุสโย สงฺขารกฺขนฺธปริยาปนฺโน สารมฺมโณติ, น เหวํ วตฺตพฺเพ ฯเปฯ.}

\prose{\csromansymbolmark{*}อวิชฺชานุสโย สงฺขารกฺขนฺธปริยาปนฺโน อนารมฺมโณ, อวิชฺชา สงฺขารกฺขนฺธปริยาปนฺนา สารมฺมณาติ, อามนฺตา.\csromanspacer{} สงฺขารกฺขนฺโธ เอกเทโส สารมฺมโณ เอกเทโส อนารมฺมโณติ, น เหวํ วตฺตพฺเพ ฯเปฯ.\csromanspacer{} สงฺขารกฺขนฺโธ เอกเทโส สารมฺมโณ เอกเทโส อนารมฺมโณติ, อามนฺตา.\csromanspacer{} เวทนากฺขนฺโธ สญฺญากฺขนฺโธ วิญฺญาณกฺขนฺโธ เอกเทโส สารมฺมโณ เอกเทโส อนารมฺมโณติ, น เหวํ วตฺตพฺเพ ฯเปฯ.}

\proseitem{556}{\csromansymbolmark{+}น วตฺตพฺพํ “อนุสยา อนารมฺมณา”ติ, อามนฺตา.\csromanspacer{} ปุถุชฺชโน กุสลาพฺยากเต จิตฺเต วตฺตมาเน “สานุสโย”ติ วตฺตพฺโพติ, อามนฺตา.\csromanspacer{} อตฺถิ เตสํ อนุสยานํ อารมฺมณนฺติ, น เหวํ วตฺตพฺเพ ฯเปฯ.\csromanspacer{} เตน หิ อนุสยา อนารมฺมณาติ.\csromanspacer{} ปุถุชฺชโน กุสลาพฺยากเต จิตฺเต วตฺตมาเน “สราโค”ติ วตฺตพฺโพติ, อามนฺตา.\csromanspacer{} อตฺถิ ตสฺส ราคสฺส อารมฺมณนฺติ, น เหวํ วตฺตพฺเพ ฯเปฯ.\csromanspacer{} เตน หิ ราโค อนารมฺมโณติ.}

\nitthitam{อนุสยา อนารมฺมณาติกถา นิฏฺฐิตา.}
\chapterhead{9. นวมวคฺค}
\vspace{\tipitakaparskip}
\csromancenter{\textbf{(88) 5.}\csromanspacer{}\textbf{ ญาณํอนารมฺมณนฺติกถา}}
\proseitem{557}{\csromansymbolmark{*}ญาณํ อนารมฺมณนฺติ, อามนฺตา.\csromanspacer{} รูปํ นิพฺพานํ จกฺขายตนํ ฯเปฯ โผฏฺฐพฺพายตนนฺติ, น เหวํ วตฺตพฺเพ ฯเปฯ.\csromanspacer{} ญาณํ อนารมฺมณนฺติ, อามนฺตา.}

\vspace{\tipitakaparskip}
\csromancenter{นวมวคฺค ปญฺญา ปญฺญินฺทฺริยํ ปญฺญาพลํ สมฺมาทิฏฺฐิ ธมฺมวิจยสมฺโพชฺฌงฺโค อนารมฺมโณติ, น เหวํ วตฺตพฺเพ ฯเปฯ.\csromanspacer{} ปญฺญา ปญฺญินฺทฺริยํ ปญฺญาพลํ สมฺมาทิฏฺฐิ ธมฺมวิจยสมฺโพชฺฌงฺโค สารมฺมโณติ, อามนฺตา.\csromanspacer{} ญาณํ สารมฺมณนฺติ, น เหวํ วตฺตพฺเพ ฯเปฯ.}
\csromanmatikaanchor{mtk.116}
\csromantocmarkref{subsection}{(89) 6. อตีตานาคตารมฺมณกถา}{mtk.116}
\bookmark[dest={mtk.116},level=2]{(89) 6. อตีตานาคตารมฺมณกถา}
\prose{\csromanfolio{301}\csromansymbolmark{*}ญาณํ อนารมฺมณนฺติ, อามนฺตา.\csromanspacer{} กตมกฺขนฺธปริยาปนฺนนฺติ? สงฺขารกฺขนฺธปริยาปนฺนนฺติ.\csromanspacer{} สงฺขารกฺขนฺโธ อนารมฺมโณติ, น เหวํ วตฺตพฺเพ ฯเปฯ.\csromanspacer{} สงฺขารกฺขนฺโธ อนารมฺมโณติ, อามนฺตา.\csromanspacer{} เวทนากฺขนฺโธ สญฺญากฺขนฺโธ วิญฺญาณกฺขนฺโธ อนารมฺมโณติ, น เหวํ วตฺตพฺเพ ฯเปฯ.\csromanspacer{} ญาณํ สงฺขารกฺขนฺธปริยาปนฺนํ อนารมฺมณนฺติ, อามนฺตา.\csromanspacer{} ปญฺญา สงฺขารกฺขนฺธปริยาปนฺนา อนารมฺมณาติ, น เหวํ วตฺตพฺเพ ฯเปฯ.\csromanspacer{} ปญฺญา สงฺขารนฺธปริยาปนฺนา สารมฺมณาติ, อามนฺตา.\csromanspacer{} ญาณํ สงฺขารกฺขนฺธปริยาปนฺนํ สารมฺมณนฺติ, น เหวํ วตฺตพฺเพ ฯเปฯ. \csromansymbolmark{*}ญาณํ สงฺขารกฺขนฺธปริยาปนฺนํ อนารมฺมณํ, ปญฺญา สงฺขารกฺขนฺธปริยาปนฺนา สารมฺมณาติ, อามนฺตา.\csromanspacer{} สงฺขารกฺขนฺโธ เอกเทโส สารมฺมโณ เอกเทโส อนารมฺมโณติ, น เหวํ วตฺตพฺเพ ฯเปฯ.\csromanspacer{} สงฺขารกฺขนฺโธ เอกเทโส สารมฺมโณ เอกเทโส อนารมฺมโณติ, อามนฺตา.\csromanspacer{} เวทนากฺขนฺโธ สญฺญากฺขนฺโธ วิญฺญาณกฺขนฺโธ เอกเทโส สารมฺมโณ เอกเทโส อนารมฺมโณติ, น เหวํ วตฺตพฺเพ ฯเปฯ.}

\proseitem{558}{\csromansymbolmark{+}น วตฺตพฺพํ “ญาณํ อนารมฺมณนฺ”ติ, อามนฺตา.\csromanspacer{} อรหา จกฺขุวิญฺญาณสมงฺคี “ญาณี”ติ วตฺตพฺโพติ, อามนฺตา.\csromanspacer{} อตฺถิ ตสฺส ญาณสฺส อารมฺมณนฺติ, น เหวํ วตฺตพฺเพ ฯเปฯ.\csromanspacer{} เตน หิ ญาณํ อนารมฺมณนฺติ.\csromanspacer{} อรหา จกฺขุวิญฺญาณสมงฺคี “ปญฺญวา”ติ วตฺตพฺโพติ, อามนฺตา.\csromanspacer{} อตฺถิ ตาย ปญฺญาย อารมฺมณนฺติ, น เหวํ วตฺตพฺเพ ฯเปฯ.\csromanspacer{} เตน หิ ปญฺญา อนารมฺมณาติ.}

\nitthitam{ญานํ อนารมฺมณนฺติกถา นิฏฺฐิตา.}
\chapterhead{9. นวมวคฺค}
\vspace{\tipitakaparskip}
\csromancenter{\textbf{(89) 6.}\csromanspacer{}\textbf{ อตีตานาคตารมฺมณกถา}}
\proseitem{559}{\csromansymbolmark{*}อตีตารมฺมณํ จิตฺตํ อนารมฺมณนฺติ, อามนฺตา.\csromanspacer{} นนุ อตีตารมฺมณนฺติ, อามนฺตา.\csromanspacer{} หญฺจิ อตีตารมฺมณํ, โน จ วต เร วตฺตพฺเพ \csromanfolio{302}“อตีตารมฺมณํ จิตฺตํ อนารมฺมณนฺ”ติ.\csromanspacer{} อตีตารมฺมณํ จิตฺตํ อนารมฺมณนฺติ, มิจฺฉา.\csromanspacer{} หญฺจิ วา ปน อนารมฺมณํ, โน จ วต เร วตฺตพฺเพ “อตีตารมฺมณนฺ”ติ.\csromanspacer{} อนารมฺมณํ อตีตารมฺมณนฺติ, มิจฺฉา.}

\prose{\csromansymbolmark{*}อตีตารมฺมณํ จิตฺตํ อนารมฺมณนฺติ, อามนฺตา.\csromanspacer{} นนุ อตีตํ อารพฺภ อตฺถิ อาวฏฺฏนา ฯเปฯ ปณิธีติ, อามนฺตา.\csromanspacer{} หญฺจิ อตีตํ อารพฺภ อตฺถิ อาวฏฺฏนา ฯเปฯ ปณิธิ, โน จ วต เร วตฺตพฺเพ “อตีตารมฺมณํ จิตฺตํ อนารมฺมณนฺ”ติ.}

\proseitem{560}{\csromansymbolmark{*}อนาคตารมฺมณํ จิตฺตํ อนารมฺมณนฺติ, อามนฺตา.\csromanspacer{} นนุ อนาคตารมฺมณนฺติ, อามนฺตา.\csromanspacer{} หญฺจิ อนาคตารมฺมณํ, โน จ วต เร วตฺตพฺเพ “อนาคตารมฺมณํ จิตฺตํ อนารมฺมณนฺ”ติ.\csromanspacer{} อนาคตารมฺมณํ จิตฺตํ อนารมฺมณนฺติ, มิจฺฉา.\csromanspacer{} หญฺจิ วา ปน อนารมฺมณํ, โน จ วร เร วตฺตพฺเพ “อนาคตารมฺมณนฺ”ติ.\csromanspacer{} อนารมฺมณํ อนาคตารมฺมณนฺติ, มิจฺฉา.}

\prose{\csromansymbolmark{*}อนาคตารมฺมณํ จิตฺตํ อนารมฺมณนฺติ, อามนฺตา.\csromanspacer{} นนุ อนาคตํ อารพฺภ อตฺถิ อาวฏฺฏนา ฯเปฯ ปณิธีติ, อามนฺตา.\csromanspacer{} หญฺจิ อนาคตํ อารพฺภ อตฺถิ อาวฏฺฏนา ฯเปฯ ปณิธิ, โน จ วต เร วตฺตพฺเพ “อนาคตารมฺมณํ จิตฺตํ อนารมฺมณนฺ”ติ.}

\prose{\csromansymbolmark{*}ปจฺจุปฺปนฺนํ อารพฺภ อตฺถิ อาวฏฺฏนา ฯเปฯ ปณิธิ, ปจฺจุปฺปนฺนารมฺมณํ จิตฺตํ สารมฺมณนฺติ, อามนฺตา.\csromanspacer{} อตีตํ อารพฺภ อตฺถิ อาวฏฺฏนา ฯเปฯ ปณิธิ, อตีตารมฺมณํ จิตฺตํ สารมฺมณนฺติ, น เหวํ วตฺตพฺเพ ฯเปฯ.\csromanspacer{} ปจฺจุปฺปนฺนํ อารพฺภ อตฺถิ อาวฏฺฏนา ฯเปฯ ปณิธิ, ปจฺจุปฺปนฺนารมฺมณํ จิตฺตํ สารมฺมณนฺติ, อามนฺตา.\csromanspacer{} อนาคตํ อารพฺภ อตฺถิ อาวฏฺฏนา ฯเปฯ ปณิธิ, อนาคตารมฺมณํ จิตฺตํ สารมฺมณนฺติ, น เหวํ วตฺตพฺเพ ฯเปฯ.}

\prose{\csromansymbolmark{*}อตีตํ อารพฺภ อตฺถิ อาวฏฺฏนา ฯเปฯ ปณิธิ, อตีตารมฺมณํ จิตฺตํ อนารมฺมณนฺติ, อามนฺตา.\csromanspacer{} ปจฺจุปฺปนฺนํ อารพฺภ อตฺถิ อาวฏฺฏนา ฯเปฯ ปณิธิ, ปจฺจุปฺปนฺนารมฺมณํ จิตฺตํ อนารมฺมณนฺติ, น เหวํ วตฺตพฺเพ ฯเปฯ.}

\prose{\csromansymbolmark{*}อนาคตํ อารพฺภ อตฺถิ อาวฏฺฏนา ฯเปฯ ปณิธิ, อนาคตารมฺมณํ จิตฺตํ อนารมฺมณนฺติ, อามนฺตา.\csromanspacer{} ปจฺจุปฺปนฺนํ อารพฺภ อตฺถิ อาวฏฺฏนา ฯเปฯ ปณิธิ, ปจฺจุปฺปนฺนารมฺมณํ จิตฺตํ อนารมฺมณนฺติ, น เหวํ วตฺตพฺเพ ฯเปฯ.}

\chapterheadpageread{9. นวมวคฺค}
\csromanmatikaanchor{mtk.117}
\csromantocmarkref{subsection}{(90) 7. วิตกฺกานุปติตกถา}{mtk.117}
\bookmark[dest={mtk.117},level=2]{(90) 7. วิตกฺกานุปติตกถา}
\proseitem{561}{\csromanfolio{303}\csromansymbolmark{+}น วตฺตพฺพํ “อตีตานาคตารมฺมณํ จิตฺตํ อนารมฺมณนฺ”ติ, อามนฺตา.\csromanspacer{} นนุ อตีตานาคตํ นตฺถีติ, อามนฺตา.\csromanspacer{} หญฺจิ อตีตานาคตํ นตฺถิ, เตน วต เร วตฺตพฺเพ “อตีตานาคตารมฺมณํ จิตฺตํ อนารมฺมณนฺ”ติ ฯเปฯ.}

\nitthitam{อตีตานาคตารมฺมณกถา นิฏฺฐิตา.}
\chapterhead{9. นวมวคฺค}
\vspace{\tipitakaparskip}
\csromancenter{\textbf{(90) 7.}\csromanspacer{}\textbf{ วิตกฺกานุปติตกถา}}
\proseitem{562}{\csromansymbolmark{*}สพฺพํ จิตฺตํ วิตกฺกานุปติตนฺติ, อามนฺตา.\csromanspacer{} สพฺพํ จิตฺตํ วิจารานุปติตํ ปีตานุปติตํ สุขานุปติตํ ทุกฺขานุปติตํ โสมนสฺสานุปติตํ โทมนสฺสานุปติตํ อุเปกฺขานุปติตํ สทฺธานุปติตํ วีริยานุปติตํ สตานุปติตํ สมาธานุปติตํ ปญฺญานุปติตํ ราคานุปติตํ โทสานุปติตํ ฯเปฯ อโนตฺตปฺปานุปติตนฺติ, น เหวํ วตฺตพฺเพ ฯเปฯ.}

\prose{\csromansymbolmark{*}สพฺพํ จิตฺตํ วิตกฺกานุปติตนฺติ, อามนฺตา.\csromanspacer{} นนุ อตฺถิ อวิตกฺโก วิจารมตฺโต สมาธีติ, อามนฺตา.\csromanspacer{} หญฺจิ อตฺถิ อวิตกฺโก วิจารมตฺโต สมาธิ, โน จ วต เร วตฺตพฺเพ “สพฺพํ จิตฺตํ วิตกฺกานุปติตนฺ”ติ.}

\prose{\csromansymbolmark{*}สพฺพํ จิตฺตํ วิตกฺกานุปติตนฺติ, อามนฺตา.\csromanspacer{} นนุ อตฺถิ อวิตกฺโก อวิจาโร สมาธีติ, อามนฺตา.\csromanspacer{} หญฺจิ อตฺถิ อวิตกฺโก อวิจาโร สมาธิ, โน จ วต เร วตฺตพฺเพ “สพฺพํ จิตฺตํ วิตกฺกานุปติตนฺ”ติ.\csromanspacer{} สพฺพํ จิตฺตํ วิตกฺกานุปติตนฺติ, อามนฺตา.\csromanspacer{} นนุ ตโย สมาธี วุตฺตา ภควตา “สวิตกฺโก สวิจาโร สมาธิ, อวิตกฺโก วิจารมตฺโต สมาธิ, อวิตกฺโก อวิจาโร สมาธี”ติ\footnote{ที 3. 184, 229 ปิฏฺเฐสุ.}, อามนฺตา.\csromanspacer{} หญฺจิ ตโย สมาธี วุตฺตา ภควตา สวิตกฺโก สวิจาโร สมาธิ, อวิตกฺโก วิจารมตฺโต สมาธิ, อวิตกฺโก อวิจาโร สมาธิ, โน จ วต เร วตฺตพฺเพ “สพฺพํ จิตฺตํ วิตกฺกานุปติตนฺ”ติ.}

\nitthitam{วิตกฺกานุปติตกถา นิฏฺฐิตา.}
\chapterheadpageread{9. นวมวคฺค}
\vspace{\tipitakaparskip}
\csromancenter{\textbf{(91) 8.}\csromanspacer{}\textbf{ วิตกฺกวิปฺผารสทฺทกถา}}
\csromanmatikaanchor{mtk.118}
\csromanmatikaanchor{mtk.119}
\csromantocmarkref{subsection}{(91) 8. วิตกฺกวิปฺผารสทฺทกถา}{mtk.118}
\bookmark[dest={mtk.118},level=2]{(91) 8. วิตกฺกวิปฺผารสทฺทกถา}
\csromantocmarkref{subsection}{(92) 9. น ยถาจิตฺตสฺสวาจาติกถา}{mtk.119}
\bookmark[dest={mtk.119},level=2]{(92) 9. น ยถาจิตฺตสฺสวาจาติกถา}
\proseitem{563}{\csromanfolio{304}\csromansymbolmark{*}สพฺพโส วิตกฺกยโต วิจารยโต วิตกฺกวิปฺผาโร สทฺโทติ, อามนฺตา.\csromanspacer{} สพฺพโส ผุสยโต ผสฺสวิปฺผาโร สทฺโท.\csromanspacer{} สพฺพโส เวทยโต เวทนาวิปฺผาโร สทฺโท.\csromanspacer{} สพฺพโส สญฺชานโต สญฺญาวิปฺผาโร สทฺโท.\csromanspacer{} สพฺพโส เจตยโต เจตนาวิปฺผาโร สทฺโท.\csromanspacer{} สพฺพโส จินฺตยโต จิตฺตวิปฺผาโร สทฺโท.\csromanspacer{} สพฺพโส สรโต สติวิปฺผาโร สทฺโท.\csromanspacer{} สพฺพโส ปชานโต ปญฺญาวิปฺผาโร สทฺโทติ, น เหวํ วตฺตพฺเพ ฯเปฯ.}

\prose{\csromansymbolmark{*}สพฺพโส วิตกฺกยโต วิจารยโต วิตกฺกวิปฺผาโร สทฺโทติ, อามนฺตา.\csromanspacer{} วิตกฺกวิปฺผาโร สทฺโท โสตวิญฺเญยฺโย โสตสฺมิํ ปฏิหญฺญติ โสตสฺส อาปาถํ อาคจฺฉตีติ, น เหวํ วตฺตพฺเพ ฯเปฯ.\csromanspacer{} นนุ วิตกฺกวิปฺผาโร สทฺโท น โสตวิญฺเญยฺโย น โสตสฺมิํ ปฏิหญฺญติ น โสตสฺส อาปาถํ อาคจฺฉตีติ, อามนฺตา.\csromanspacer{} หญฺจิ วิตกฺกวิปฺผาโร สทฺโท น โสตวิญฺเญยฺโย น โสตสฺมิํ ปฏิหญฺญติ น โสตสฺส อาปาถํ อาคจฺฉติ, โน จ วต เร วตฺตพฺเพ “สพฺพโส วิตกฺกยโต วิจารยโต วิตกฺกวิปฺผาโร สทฺโท”ติ.}

\nitthitam{วิตกฺกวิปฺผารสทฺทกถา นิฏฺฐิตา.}
\chapterhead{9. นวมวคฺค}
\prose{\textbf{(92) 9.}\csromanspacer{}\textbf{ นยถาจิตฺตสฺสวาจาติกถา}}

\proseitem{564}{น ยถาจิตฺตสฺส วาจาติ, อามนฺตา.\csromanspacer{} อผสฺสกสฺส วาจา อเวทนกสฺส วาจา อสญฺญกสฺส วาจา อเจตนกสฺส วาจา อจิตฺตกสฺส วาจาติ, น เหวํ วตฺตพฺเพ ฯเปฯ.\csromanspacer{} นนุ สผสฺสกสฺส วาจา สเวทนกสฺส วาจา สสญฺญกสฺส วาจา สเจตนกสฺส วาจา สจิตฺตกสฺส วาจาติ, อามนฺตา.\csromanspacer{} หญฺจิ สผสฺสกสฺส วาจา ฯเปฯ สจิตฺตกสฺส วาจา, โน จ วต เร วตฺตพฺเพ “น ยถาจิตฺตสฺส วาจา”ติ.}

\chapterheadpageread{9. นวมวคฺค}
\prose{\csromanfolio{305}\csromansymbolmark{*}น ยถาจิตฺตสฺส วาจาติ, อามนฺตา.\csromanspacer{} อนาวฏฺเฏนฺตสฺส\footnote{อนาวฏฺฏนฺตสฺส (อิ), อนาวชฺชนฺตสฺส (สี, สฺยา, ก) 257 ปิฏฺเฐปิ ปสฺสิตพฺพํ.} วาจา ฯเปฯ.\csromanspacer{} อนาโภคสฺส วาจา ฯเปฯ.\csromanspacer{} อปฺปณิทหนฺตสฺส วาจาติ, น เหวํ วตฺตพฺเพ ฯเปฯ.\csromanspacer{} นนุ อาวฏฺเฏนฺตสฺส วาจา อาโภคสฺส วาจา ฯเปฯ ปณิทหนฺตสฺส วาจาติ, อามนฺตา.\csromanspacer{} หญฺจิ อาวฏฺเฏนฺตสฺส วาจา อาโภคสฺส วาจา ปณิทหนฺตสฺส วาจา, โน จ วต เร วตฺตพฺเพ “น ยถาจิตฺตสฺส วาจา”ติ.}

\prose{\csromansymbolmark{*}น ยถาจิตฺตสฺส วาจาติ, อามนฺตา.\csromanspacer{} นนุ วาจา จิตฺตสมุฏฺฐานา จิตฺเตน สหชาตา จิตฺเตน สห เอกุปฺปาทาติ, อามนฺตา.\csromanspacer{} หญฺจิ วาจา จิตฺตสมุฏฺฐานา จิตฺเตน สหชาตา จิตฺเตน สห เอกุปฺปาทา, โน จ วต เร วตฺตพฺเพ “น ยถาจิตฺตสฺส วาจา”ติ.}

\prose{\csromansymbolmark{*}น ยถาจิตฺตสฺส วาจาติ, อามนฺตา.\csromanspacer{} น ภณิตุกาโม ภณติ, น กเถตุกาโม กเถติ, น อาลปิตุกาโม อาลปติ, น โวหริตุกาโม โวหรตีติ, น เหวํ วตฺตพฺเพ ฯเปฯ.\csromanspacer{} นนุ ภณิตุกาโม ภณติ, กเถตุกาโม กเถติ, อาลปิตุกาโม อาลปติ, โวหริตุกาโม โวหรตีติ, อามนฺตา.\csromanspacer{} หญฺจิ ภณิตุกาโม ภณติ, กเถตุกาโม กเถติ, อาลปิตุกาโม อาลปติ, โวหริตุกาโม โวหรติ, โน จ วต เร วตฺตพฺเพ “น ยถาจิตฺตสฺส วาจา”ติ.}

\proseitem{565}{\csromansymbolmark{+}น วตฺตพฺพํ “น ยถาจิตฺตสฺส วาจา”ติ, อามนฺตา.\csromanspacer{} นนุ อตฺถิ โกจิ “อญฺญํ ภณิสฺสามี”ติ อญฺญํ ภณติ, “อญฺญํ กเถสฺสามี”ติ อญฺญํ กเถติ, “อญฺญํ อาลปิสฺสามี”ติ อญฺญํ อาลปติ, “อญฺญํ โวหริสฺสามี”ติ อญฺญํ โวหรตีติ, อามนฺตา.\csromanspacer{} หญฺจิ อตฺถิ โกจิ “อญฺญํ ภณิสฺสามี”ติ อญฺญํ ภณติ ฯเปฯ “อญฺญํ โวหริสฺสามี”ติ อญฺญํ โวหรติ, เตน วต เร วตฺตพฺเพ “น ยถาจิตฺตสฺส วาจา”ติ.}

\nitthitam{น ยถาจิตฺตสฺส วาจาติกถา นิฏฺฐิตา.}
\chapterheadpageread{9. นวมวคฺค}
\vspace{\tipitakaparskip}
\csromancenter{\textbf{(93) 10.}\csromanspacer{}\textbf{ นยถาจิตฺตสฺสกายกมฺมนฺติกถา}}
\csromanmatikaanchor{mtk.120}
\csromantocmarkref{subsection}{(93) 10. น ยถาจิตฺตสฺสกายกมฺมนฺติกถา}{mtk.120}
\bookmark[dest={mtk.120},level=2]{(93) 10. น ยถาจิตฺตสฺสกายกมฺมนฺติกถา}
\proseitem{566}{\csromanfolio{306}\csromansymbolmark{*}น ยถาจิตฺตสฺส กายกมฺมนฺติ, อามนฺตา.\csromanspacer{} อผสฺสกสฺส กายกมฺมํ ฯเปฯ.\csromanspacer{} อจิตฺตกสฺส กายกมฺมนฺติ, น เหวํ วตฺตพฺเพ ฯเปฯ.\csromanspacer{} นนุ สผสฺสกสฺส กายกมฺมํ ฯเปฯ สจิตฺตกสฺส กายกมฺมนฺติ, อามนฺตา.\csromanspacer{} หญฺจิ สผสฺสกสฺส กายกมฺมํ ฯเปฯ สจิตฺตกสฺส กายกมฺมํ, โน จ วต เร วตฺตพฺเพ “น ยถาจิตฺตสฺส กายกมฺมนฺ”ติ.}

\prose{\csromansymbolmark{*}น ยถาจิตฺตสฺส กายกมฺมนฺติ, อามนฺตา.\csromanspacer{} อนาวฏฺเฏนฺตสฺส กายกมฺมํ ฯเปฯ.\csromanspacer{} อปฺปณิทหนฺตสฺส กายกมฺมนฺติ, น เหวํ วตฺตพฺเพ ฯเปฯ.\csromanspacer{} นนุ อาวฏฺเฏนฺตสฺส กายกมฺมํ ฯเปฯ ปณิทหนฺตสฺส กายกมฺมนฺติ, อามนฺตา.\csromanspacer{} หญฺจิ อาวฏฺเฏนฺตสฺส กายกมฺมํ ฯเปฯ ปณิทหนฺตสฺส กายกมฺมํ, โน จ วต เร วตฺตพฺเพ “น ยถาจิตฺตสฺส กายกมฺมนฺ”ติ.}

\prose{\csromansymbolmark{*}น ยถาจิตฺตสฺส กายกมฺมนฺติ, อามนฺตา.\csromanspacer{} นนุ กายกมฺมํ จิตฺตสมุฏฺฐานํ จิตฺเตน สหชาตํ จิตฺเตน สห เอกุปฺปาทนฺติ, อามนฺตา.\csromanspacer{} หญฺจิ กายกมฺมํ จิตฺตสมุฏฺฐานํ จิตฺเตน สหชาตํ จิตฺเตน สห เอกุปฺปาทํ, โน จ วต เร วตฺตพฺเพ “น ยถาจิตฺตสฺส กายกมฺมนฺ”ติ.}

\prose{\csromansymbolmark{*}น ยถาจิตฺตสฺส กายกมฺมนฺติ, อามนฺตา.\csromanspacer{} น อภิกฺกมิตุกาโม อภิกฺกมติ, น ปฏิกฺกมิตุกาโม ปฏิกฺกมติ, น อาโลเกตุกาโม อาโลเกติ, น วิโลเกตุกาโม วิโลเกติ, น สมิญฺชิตุกาโม สมิญฺเชติ, น ปสาเรตุกาโม ปสาเรตีติ, น เหวํ วตฺตพฺเพ ฯเปฯ.\csromanspacer{} นนุ อภิกฺกมิตุกาโม อภิกฺกมติ, ปฏิกฺกมิตุกาโม ปฏิกฺกมติ, อาโลเกตุกาโม อาโลเกติ, วิโลเกตุกาโม วิโลเกติ, สมิญฺชิตุกาโม สมิญฺเชติ, ปสาเรตุกาโม ปสาเรตีติ, อามนฺตา.\csromanspacer{} หญฺจิ อภิกฺกมิตุกาโม อภิกฺกมติ ฯเปฯ ปสาเรตุกาโม ปสาเรติ, โน จ วต เร วตฺตพฺเพ “น ยถาจิตฺตสฺส กายกมฺมนฺ”ติ.}

\proseitem{567}{\csromansymbolmark{+}น วตฺตพฺพํ “น ยถาจิตฺตสฺส กายกมฺมนฺ”ติ, อามนฺตา.\csromanspacer{} นนุ อตฺถิ โกจิ “อญฺญตฺร คจฺฉิสฺสามี”ติ อญฺญตฺร คจฺฉติ ฯเปฯ “อญฺญํ ปสาเรสฺสามี”ติ อญฺญํ ปสาเรตีติ, อามนฺตา.\csromanspacer{} หญฺจิ อตฺถิ โกจิ}

\vspace{\tipitakaparskip}
\csromancenter{นวมวคฺค อญฺญตฺร คจฺฉิสฺสามี”ติ อญฺญตฺร คจฺฉติ ฯเปฯ “อญฺญํ ปสาเรสฺสามี”ติ อญฺญํ ปสาเรติ, เตน วต เร วตฺตพฺเพ “น ยถาจิตฺตสฺส กายกมฺมนฺ”ติ.}
\nitthitam{น ยถาจิตฺตสฺส กายกมฺมนฺติกถา นิฏฺฐิตา.}
\chapterhead{9. นวมวคฺค}
\vspace{\tipitakaparskip}
\csromancenter{\textbf{(94) 11.}\csromanspacer{}\textbf{ อตีตานาคตสมนฺนาคตกถา}}
\csromanmatikaanchor{mtk.121}
\csromantocmarkref{subsection}{(94) 11. อตีตานาคตสมนฺนาคตกถา}{mtk.121}
\bookmark[dest={mtk.121},level=2]{(94) 11. อตีตานาคตสมนฺนาคตกถา}
\proseitem{568}{\csromanfolio{307}\csromansymbolmark{*}อตีเตน สมนฺนาคโตติ, อามนฺตา.\csromanspacer{} นนุ อตีตํ นิรุทฺธํ วิคตํ วิปริณตํ อตฺถงฺคตํ อพฺภตฺถงฺคตนฺติ, อามนฺตา.\csromanspacer{} หญฺจิ อตีตํ นิรุทฺธํ วิคตํ วิปริณตํ อตฺถงฺคตํ อพฺภตฺถงฺคตํ, โน จ วต เร วตฺตพฺเพ “อตีเตน สมนฺนาคโต”ติ.}

\prose{\csromansymbolmark{*}อนาคเตน สมนฺนาคโตติ, อามนฺตา.\csromanspacer{} นนุ อนาคตํ อชาตํ อภูตํ อสญฺชาตํ อนิพฺพตฺตํ อนภินิพฺพตฺตํ อปาตุภูตนฺติ, อามนฺตา.\csromanspacer{} หญฺจิ อนาคตํ อชาตํ อภูตํ อสญฺชาตํ อนิพฺพตฺตํ อนภินิพฺพตฺตํ อปาตุภูตํ, โน จ วต เร วตฺตพฺเพ “อนาคเตน สมนฺนาคโต”ติ.}

\proseitem{569}{\csromansymbolmark{*}อตีเตน รูปกฺขนฺเธน สมนฺนาคโต, อนาคเตน รูปกฺขนฺเธน สมนฺนาคโต, ปจฺจุปฺปนฺเนน รูปกฺขนฺเธน สมนฺนาคโตติ, อามนฺตา.\csromanspacer{} ตีหิ รูปกฺขนฺเธหิ สมนฺนาคโตติ, น เหวํ วตฺตพฺเพ ฯเปฯ.\csromanspacer{} อตีเตหิ ปญฺจหิ ขนฺเธหิ สมนฺนาคโต, อนาคเตหิ ปญฺจหิ ขนฺเธหิ สมนฺนาคโต, ปจฺจุปฺปนฺเนหิ ปญฺจหิ ขนฺเธหิ สมนฺนาคโตติ, อามนฺตา.\csromanspacer{} ปนฺนรสหิ ขนฺเธหิ สมนฺนาคโตติ, น เหวํ วตฺตพฺเพ ฯเปฯ.}

\prose{\csromansymbolmark{*}อตีเตน จกฺขายตเนน สมนฺนาคโต, อนาคเตน จกฺขายตเนน สมนฺนาคโต, ปจฺจุปฺปนฺเนน จกฺขายตเนน สมนฺนาคโตติ, อามนฺตา.\csromanspacer{} ตีหิ จกฺขายตเนหิ สมนฺนาคโตติ, น เหวํ วตฺตพฺเพ ฯเปฯ.\csromanspacer{} อตีเตหิ ทฺวาทสหิ อายตเนหิ สมนฺนาคโต, อนาคเตหิ ทฺวาทสหิ อายตเนหิ สมนฺนาคโต, ปจฺจุปฺปนฺเนหิ ทฺวาทสหิ อายตเนหิ สมนฺนาคโตติ, อามนฺตา.\csromanspacer{} ฉตฺติํสายตเนหิ สมนฺนาคโตติ, น เหวํ วตฺตพฺเพ ฯเปฯ.}

\csromanmatikaanchor{mtk.122}
\csromanmatikaanchor{mtk.123}
\csromantocmarknopage{chapter}{10. ทสมวคฺค}{mtk.122}
\bookmark[dest={mtk.122},level=0]{10. ทสมวคฺค}
\csromantocmarkref{subsection}{(95) 1. นิโรธกถา}{mtk.123}
\bookmark[dest={mtk.123},level=2]{(95) 1. นิโรธกถา}
\prose{\csromansymbolmark{*}อตีตาย จกฺขุธาตุยา สมนฺนาคโต, อนาคตาย จกฺขุธาตุยา สมนฺนาคโต, ปจฺจุปฺปนฺนาย จกฺขุธาตุยา สมนฺนาคโตติ, \csromanfolio{308}อามนฺตา.\csromanspacer{} ตีหิ จกฺขุธาตูหิ สมนฺนาคโตติ, น เหวํ วตฺตพฺเพ ฯเปฯ.\csromanspacer{} อตีตาหิ อฏฺฐารสหิ ธาตูหิ สมนฺนาคโต, อนาคตาหิ อฏฺฐารสหิ ธาตูหิ สมนฺนาคโต, ปจฺจุปฺปนฺนาหิ อฏฺฐารสหิ ธาตูหิ สมนฺนาคโตติ, อามนฺตา.\csromanspacer{} จตุปญฺญาสธาตูหิ สมนฺนาคโตติ, น เหวํ วตฺตพฺเพ ฯเปฯ.}

\prose{\csromansymbolmark{*}อตีเตน จกฺขุนฺทฺริเยน สมนฺนาคโต, อนาคเตน จกฺขุนฺทฺริเยน สมนฺนาคโต, ปจฺจุปฺปนฺเนน จกฺขุนฺทฺริเยน สมนฺนาคโตติ, อามนฺตา.\csromanspacer{} ตีหิ จกฺขุนฺทฺริเยหิ สมนฺนาคโตติ, น เหวํ วตฺตพฺเพ ฯเปฯ.\csromanspacer{} อตีเตหิ พาวีสตินฺทฺริเยหิ สมนฺนาคโต, อนาคเตหิ พาวีสตินฺทฺริเยหิ สมนฺนาคโต, ปจฺจุปฺปนฺเนหิ พาวีสตินฺทฺริเยหิ สมนฺนาคโตติ, อามนฺตา.\csromanspacer{} ฉสฏฺฐินฺทฺริเยหิ สมนฺนาคโตติ, น เหวํ วตฺตพฺเพ ฯเปฯ.}

\proseitem{570}{\csromansymbolmark{+}น วตฺตพฺพํ “อตีตานาคเตหิ สมนฺนาคโต”ติ, อามนฺตา.\csromanspacer{} นนุ อตฺถิ อฏฺฐวิโมกฺขฌายี จตุนฺนํ ฌานานํ นิกามลาภี นวนฺนํ อนุปุพฺพวิหารสมาปตฺตีนํ ลาภีติ, อามนฺตา.\csromanspacer{} หญฺจิ อตฺถิ อฏฺฐวิโมกฺขฌายี จตุนฺนํ ฌานานํ นิกามลาภี นวนฺนํ อนุปุพฺพวิหารสมาปตฺตีนํ ลาภี, เตน วเต เร วตฺตพฺเพ “อตีตานาคเตหิ สมนฺนาคโต”ติ.\csromanspacer{} อตีตานาคตปจฺจุปฺปนฺนกถา นิฏฺฐิตา.\csromanspacer{} นวโม วคฺโค.}

\titlehead{\textbf{ตสฺสุทฺทานํ}}
\prose{อานิสํสทสฺสาวิสฺส สํโยชนานํ ปหานํ, อมตารมฺมณํ สํโยชนํ, รูปํ สารมฺมณํ, อนุสยา อนารมฺมณา, เอวเมวํ ญาณํ, อตีตานาคตารมฺมณํ จิตฺตํ, สพฺพํ จิตฺตํ วิตกฺกานุปติตํ, สพฺพโส วิตกฺกยโต วิจารยโต วิตกฺกวิปฺผาโร สทฺโท, น ยถาจิตฺตสฺส วาจา, ตเถว กายกมฺมํ, อตีตานาคเตหิ สมนฺนาคโตติ.}

\chapterhead{10. ทสมวคฺค}
\vspace{\tipitakaparskip}
\csromancenter{\textbf{(95) 1.}\csromanspacer{}\textbf{ นิโรธกถา}}
\proseitem{571}{\csromansymbolmark{*}อุปปตฺเตสิเย ปญฺจกฺขนฺเธ อนิรุทฺเธ กิริยา ปญฺจกฺขนฺธา อุปฺปชฺชนฺตีติ, อามนฺตา.\csromanspacer{} ทสนฺนํ ขนฺธานํ สโมธานํ โหติ.\csromanspacer{} ทส ขนฺธา สมฺมุขีภาวํ}

\vspace{\tipitakaparskip}
\csromancenter{ทสมวคฺค อาคจฺฉนฺตีติ, น เหวํ วตฺตพฺเพ ฯเปฯ.\csromanspacer{} ทสนฺนํ ขนฺธานํ สโมธานํ โหติ.\csromanspacer{} ทส ขนฺธา สมฺมุขีภาวํ อาคจฺฉนฺตีติ, อามนฺตา.\csromanspacer{} ทฺวินฺนํ ผสฺสานํ ฯเปฯ.\csromanspacer{} ทฺวินฺนํ จิตฺตานํ สโมธานํ โหตีติ, น เหวํ วตฺตพฺเพ ฯเปฯ.}
\csromanmatikaanchor{mtk.124}
\csromantocmarkref{subsection}{(96) 2. รูปํมคฺโคติกถา}{mtk.124}
\bookmark[dest={mtk.124},level=2]{(96) 2. รูปํมคฺโคติกถา}
\prose{\csromanfolio{309}\csromansymbolmark{*}อุปปตฺเตสิเย ปญฺจกฺขนฺเธ อนิรุทฺเธ กิริยา จตฺตาโร ขนฺธา อุปฺปชฺชนฺตีติ, อามนฺตา.\csromanspacer{} นวนฺนํ ขนฺธานํ สโมธานํ โหติ.\csromanspacer{} นว ขนฺธา สมฺมุขีภาวํ อาคจฺฉนฺตีติ, น เหวํ วตฺตพฺเพ ฯเปฯ.\csromanspacer{} นวนฺนํ ขนฺธานํ สโมธานํ โหติ.\csromanspacer{} นว ขนฺธา สมฺมุขีภาวํ อาคจฺฉนฺตีติ, อามนฺตา.\csromanspacer{} ทฺวินฺนํ ผสฺสานํ ฯเปฯ.\csromanspacer{} ทฺวินฺนํ จิตฺตานํ สโมธานํ โหตีติ, น เหวํ วตฺตพฺเพ ฯเปฯ.}

\prose{\csromansymbolmark{*}อุปปตฺเตสิเย ปญฺจกฺขนฺเธ อนิรุทฺเธ กิริยาญาณํ อุปฺปชฺชตีติ, อามนฺตา.\csromanspacer{} ฉนฺนํ ขนฺธานํ สโมธานํ โหติ.\csromanspacer{} ฉ ขนฺธา สมฺมุขีภาวํ อาคจฺฉนฺตีติ, น เหวํ วตฺตพฺเพ ฯเปฯ.\csromanspacer{} ฉนฺนํ ขนฺธานํ สโมธานํ โหติ.\csromanspacer{} ฉ ขนฺธา สมฺมุขีภาวํ อาคจฺฉนฺตีติ, อามนฺตา.\csromanspacer{} ทฺวินฺนํ ผสฺสานํ ฯเปฯ.\csromanspacer{} ทฺวินฺนํ จิตฺตานํ สโมธานํ โหตีติ, น เหวํ วตฺตพฺเพ ฯเปฯ.}

\proseitem{572}{\csromansymbolmark{+}อุปปตฺเตสิเย ปญฺจกฺขนฺเธ นิรุทฺเธ มคฺโค อุปฺปชฺชตีติ, อามนฺตา.\csromanspacer{} มโต มคฺคํ ภาเวติ กาลงฺกโต มคฺคํ ภาเวตีติ, น เหวํ วตฺตพฺเพ ฯเปฯ.}

\nitthitam{นิโรธกถา นิฏฺฐิตา.}
\chapterhead{10. ทสมวคฺค}
\vspace{\tipitakaparskip}
\csromancenter{\textbf{(96) 2.}\csromanspacer{}\textbf{ รูปํมคฺโคติกถา}}
\proseitem{573}{\csromansymbolmark{*}มคฺคสมงฺคิสฺส รูปํ มคฺโคติ, อามนฺตา.\csromanspacer{} สารมฺมณํ, อตฺถิ ตสฺส อาวฏฺฏนา อาโภโค สมนฺนาหาโร มนสิกาโร เจตนา ปตฺถนา ปณิธีติ, น เหวํ วตฺตพฺเพ ฯเปฯ.\csromanspacer{} นนุ อนารมฺมณํ, นตฺถิ ตสฺส อาวฏฺฏนา ฯเปฯ ปณิธีติ, อามนฺตา.\csromanspacer{} หญฺจิ อนารมฺมณํ, นตฺถิ ตสฺส อาวฏฺฏนา ฯเปฯ ปณิธิ, โน จ วต เร วตฺตพฺเพ “มคฺคสมงฺคิสฺส รูปํ มคฺโค”ติ.}

\prose{\csromansymbolmark{*}สมฺมาวาจา มคฺโคติ, อามนฺตา.\csromanspacer{} สารมฺมณา, อตฺถิ ตาย อาวฏฺฏนา ฯเปฯ ปณิธีติ, น เหวํ วตฺตพฺเพ ฯเปฯ.\csromanspacer{} นนุ อนารมฺมณา, นตฺถิ ตาย อาวฏฺฏนา ฯเปฯ ปณิธีติ, อามนฺตา.\csromanspacer{} หญฺจิ อนารมฺมณา, นตฺถิ ตาย อาวฏฺฏนา ฯเปฯ ปณิธิ, โน จ วต เร วตฺตพฺเพ “สมฺมาวาจา มคฺโค”ติ.}

\prose{\csromanfolio{310}\csromansymbolmark{*}สมฺมากมฺมนฺโต ฯเปฯ.\csromanspacer{} สมฺมา อาชีโว มคฺโคติ, อามนฺตา.\csromanspacer{} สารมฺมโณ, อตฺถิ ตสฺส อาวฏฺฏนา ฯเปฯ ปณิธีติ, น เหวํ วตฺตพฺเพ ฯเปฯ.\csromanspacer{} นนุ อนารมฺมโณ, นตฺถิ ตสฺส อาวฏฺฏนา ฯเปฯ ปณิธีติ, อามนฺตา.\csromanspacer{} หญฺจิ อนารมฺมโณ, นตฺถิ ตสฺส อาวฏฺฏนา ฯเปฯ ปณิธิ, โน จ วต เร วตฺตพฺเพ “สมฺมา-อาชีโว มคฺโค”ติ.}

\proseitem{574}{\csromansymbolmark{*}สมฺมาทิฏฺฐิ มคฺโค, สา จ สารมฺมณา, อตฺถิ ตาย อาวฏฺฏนา ฯเปฯ ปณิธีติ, อามนฺตา.\csromanspacer{} สมฺมาวาจา มคฺโค, สา จ สารมฺมณา, อตฺถิ ตาย อาวฏฺฏนา ฯเปฯ ปณิธีติ, น เหวํ วตฺตพฺเพ ฯเปฯ.\csromanspacer{} สมฺมาทิฏฺฐิ มคฺโค, สา จ สารมฺมณา, อตฺถิ ตาย อาวฏฺฏนา ฯเปฯ ปณิธีติ, อามนฺตา.\csromanspacer{} สมฺมากมฺมนฺโต ฯเปฯ.\csromanspacer{} สมฺมา อาชีโว มคฺโค, โส จ สารมฺมโณ, อตฺถิ ตสฺส อาวฏฺฏนา ฯเปฯ ปณิธีติ, น เหวํ วตฺตพฺเพ ฯเปฯ.\csromanspacer{} สมฺมาสงฺกปฺโป ฯเปฯ.\csromanspacer{} สมฺมาวายาโม ฯเปฯ.\csromanspacer{} สมฺมาสติ ฯเปฯ.}

\prose{\csromansymbolmark{*}สมฺมาสมาธิ มคฺโค, โส จ สารมฺมโณ, อตฺถิ ตสฺส อาวฏฺฏนา ฯเปฯ ปณิธีติ, อามนฺตา.\csromanspacer{} สมฺมาวาจา มคฺโค, สา จ สารมฺมณา, อตฺถิ ตาย อาวฏฺฏนา ฯเปฯ ปณิธีติ, น เหวํ วตฺตพฺเพ ฯเปฯ.\csromanspacer{} สมฺมาสมาธิ มคฺโค, โส จ สารมฺมโณ, อตฺถิ ตสฺส อาวฏฺฏนา ฯเปฯ ปณิธีติ, อามนฺตา.\csromanspacer{} สมฺมากมฺมนฺโต ฯเปฯ.\csromanspacer{} สมฺมา-อาชีโว มคฺโค, โส จ สารมฺมโณ, อตฺถิ ตสฺส อาวฏฺฏนา ฯเปฯ ปณิธีติ, น เหวํ วตฺตพฺเพ ฯเปฯ.}

\prose{\csromansymbolmark{*}สมฺมาวาจา มคฺโค, สา จ อนารมฺมณา, นตฺถิ ตาย อาวฏฺฏนา ฯเปฯ ปณิธีติ, อามนฺตา.\csromanspacer{} สมฺมาทิฏฺฐิ มคฺโค, สา จ อนารมฺมณา, นตฺถิ ตาย อาวฏฺฏนา ฯเปฯ ปณิธีติ, น เหวํ วตฺตพฺเพ ฯเปฯ.\csromanspacer{} สมฺมาวาจา มคฺโค, สา จ อนารมฺมณา, นตฺถิ ตาย อาวฏฺฏนา ฯเปฯ ปณิธีติ, อามนฺตา.\csromanspacer{} สมฺมาสงฺกปฺโป ฯเปฯ.\csromanspacer{} สมฺมาวายาโม สมฺมาสติ สมฺมาสมาธิ มคฺโค, โส จ อนารมฺมโณ, นตฺถิ ตสฺส อาวฏฺฏนา ฯเปฯ ปณิธีติ, น เหวํ วตฺตพฺเพ ฯเปฯ.}

\prose{\csromansymbolmark{*}สมฺมากมฺมนฺโต ฯเปฯ.\csromanspacer{} สมฺมา อาชีโว มคฺโค, โส จ อนารมฺมโณ, นตฺถิ ตสฺส อาวฏฺฏนา ฯเปฯ ปณิธีติ, อามนฺตา.\csromanspacer{} สมฺมาทิฏฺฐิ สมฺมาสงฺกปฺโป สมฺมาวายาโม สมฺมาสติ ฯเปฯ.\csromanspacer{} สมฺมาสมาธิ มคฺโค, โส จ อนารมฺมโณ, นตฺถิ ตสฺส อาวฏฺฏนา ฯเปฯ ปณิธีติ, น เหวํ วตฺตพฺเพ ฯเปฯ.}

\chapterheadpageread{10. ทสมวคฺค}
\csromanmatikaanchor{mtk.125}
\csromantocmarkref{subsection}{(97) 3. ปญฺจวิญฺญาณสมงฺคิสฺสมคฺคกถา}{mtk.125}
\bookmark[dest={mtk.125},level=2]{(97) 3. ปญฺจวิญฺญาณสมงฺคิสฺสมคฺคกถา}
\proseitem{575}{\csromanfolio{311}\csromansymbolmark{+}น วตฺตพฺพํ “มคฺคสมงฺคิสฺส รูปํ มคฺโค”ติ, อามนฺตา.\csromanspacer{} นนุ สมฺมาวาจา สมฺมากมฺมนฺโต สมฺมา-อาชีโว มคฺโคติ, อามนฺตา.\csromanspacer{} หญฺจิ สมฺมาวาจา สมฺมากมฺมนฺโต สมฺมา-อาชีโว มคฺโค, เตน วต เร วตฺตพฺเพ “มคฺคสมงฺคิสฺส รูปํ มคฺโค”ติ.}

\nitthitam{รูปํ มคฺโคติกถา นิฏฺฐิตา.}
\chapterhead{10. ทสมวคฺค}
\vspace{\tipitakaparskip}
\csromancenter{\textbf{(97) 3.}\csromanspacer{}\textbf{ ปญฺจวิญฺญาณสมงฺคิสฺสมคฺคกถา}}
\proseitem{576}{\csromansymbolmark{*}ปญฺจวิญฺญาณสมงฺคิสฺส อตฺถิ มคฺคภาวนาติ, อามนฺตา.\csromanspacer{} นนุ ปญฺจวิญฺญาณา อุปฺปนฺนวตฺถุกา อุปฺปนฺนารมฺมณาติ, อามนฺตา.\csromanspacer{} หญฺจิ ปญฺจวิญฺญาณา อุปฺปนฺนวตฺถุกา อุปฺปนฺนารมฺมณา, โน จ วต เร วตฺตพฺเพ “ปญฺจวิญฺญาณสมงฺคิสฺส อตฺถิ มคฺคภาวนา”ติ.}

\prose{\csromansymbolmark{*}นนุ ปญฺจวิญฺญาณา ปุเรชาตวตฺถุกา ปุเรชาตารมฺมณา อชฺฌตฺติกวตฺถุกา พาหิรารมฺมณา อสมฺภินฺนวตฺถุกา อสมฺภินฺนารมฺมณา นานาวตฺถุกา นานารมฺมณา น อญฺญมญฺญสฺส โคจรวิสยํ ปจฺจนุโภนฺติ, น อสมนฺนาหารา อุปฺปชฺชนฺติ, น อมนสิการา อุปฺปชฺชนฺติ, น อพฺโพกิณฺณา อุปฺปชฺชนฺติ, น อปุพฺพํ อจริมํ อุปฺปชฺชนฺติ, น อญฺญมญฺญสฺส สมนนฺตรา อุปฺปชฺชนฺติ, นนุ ปญฺจวิญฺญาณา อนาโภคาติ, อามนฺตา.\csromanspacer{} หญฺจิ ปญฺจวิญฺญาณา อนาโภคา, โน จ วต เร วตฺตพฺเพ “ปญฺจวิญฺญาณสมงฺคิสฺส อตฺถิ มคฺคภาวนา”ติ.}

\proseitem{577}{\csromansymbolmark{*}จกฺขุวิญฺญาณสมงฺคิสฺส อตฺถิ มคฺคภาวนาติ, อามนฺตา.\csromanspacer{} จกฺขุวิญฺญาณํ สุญฺญตํ อารพฺภ อุปฺปชฺชตีติ, น เหวํ วตฺตพฺเพ ฯเปฯ.\csromanspacer{} จกฺขุวิญฺญาณํ สุญฺญตํ อารพฺภ อุปฺปชฺชตีติ, อามนฺตา.\csromanspacer{} จกฺขุญฺจ ปฏิจฺจ สุญฺญตญฺจ อุปฺปชฺชติ จกฺขุวิญฺญาณนฺติ, น เหวํ วตฺตพฺเพ ฯเปฯ.\csromanspacer{} จกฺขุญฺจ ปฏิจฺจ สุญฺญตญฺจ อุปฺปชฺชติ จกฺขุวิญฺญาณนฺติ, อามนฺตา. “จกฺขุญฺจ ปฏิจฺจ สุญฺญตญฺจ อุปฺปชฺชติ จกฺขุวิญฺญาณนฺ”ติ อตฺเถว สุตฺตนฺโตติ, นตฺถิ.\csromanspacer{} จกฺขุญฺจ ปฏิจฺจ รูเป จ อุปฺปชฺชติ จกฺขุวิญฺญาณนฺติ อตฺเถว สุตฺตนฺโตติ, อามนฺตา.\csromanspacer{} หญฺจิ จกฺขุญฺจ ปฏิจฺจ รูเป จ อุปฺปชฺชติ จกฺขุวิญฺญาณนฺติ อตฺเถว สุตฺตนฺโต, โน จ วต เร วตฺตพฺเพ “จกฺขุญฺจ ปฏิจฺจ สุญฺญตญฺจ อุปฺปชฺชติ จกฺขุวิญฺญาณนฺ”ติ.}

\csromanmatikaanchor{mtk.126}
\csromantocmarkref{subsection}{(98) 4. ปญฺจวิญฺญาณากุสลาปิ-อกุสลาปีติกถา}{mtk.126}
\bookmark[dest={mtk.126},level=2]{(98) 4. ปญฺจวิญฺญาณากุสลาปิ-อกุสลาปีติกถา}
\proseitem{578}{\csromanfolio{312}\csromansymbolmark{*}จกฺขุวิญฺญาณสมงฺคิสฺส อตฺถิ มคฺคภาวนาติ, อามนฺตา.\csromanspacer{} จกฺขุวิญฺญาณํ อตีตานาคตํ อารพฺภ อุปฺปชฺชตีติ, น เหวํ วตฺตพฺเพ ฯเปฯ.\csromanspacer{} จกฺขุวิญฺญาณสมงฺคิสฺส อตฺถิ มคฺคภาวนาติ, อามนฺตา.\csromanspacer{} จกฺขุวิญฺญาณํ ผสฺสํ อารพฺภ.\csromanspacer{} เวทนํ อารพฺภ.\csromanspacer{} สญฺญํ อารพฺภ.\csromanspacer{} เจตนํ อารพฺภ.\csromanspacer{} จิตฺตํ อารพฺภ.\csromanspacer{} จกฺขุํ อารพฺภ ฯเปฯ.\csromanspacer{} กายํ อารพฺภ ฯเปฯ.\csromanspacer{} สทฺทํ อารพฺภ ฯเปฯ.\csromanspacer{} โผฏฺฐพฺพํ อารพฺภ อุปฺปชฺชตีติ, น เหวํ วตฺตพฺเพ ฯเปฯ.}

\prose{\csromansymbolmark{*}มโนวิญฺญาณสมงฺคิสฺส อตฺถิ มคฺคภาวนา, มโนวิญฺญาณํ สุญฺญตํ อารพฺภ อุปฺปชฺชตีติ, อามนฺตา.\csromanspacer{} จกฺขุวิญฺญาณสมงฺคิสฺส อตฺถิ มคฺคภาวนา, จกฺขุวิญฺญาณํ สุญฺญตํ อารพฺภ อุปฺปชฺชตีติ, น เหวํ วตฺตพฺเพ ฯเปฯ.\csromanspacer{} มโนวิญฺญาณสมงฺคิสฺส อตฺถิ มคฺคภาวนา, มโนวิญฺญาณํ อตีตานาคตํ อารพฺภ อุปฺปชฺชตีติ, อามนฺตา.\csromanspacer{} จกฺขุวิญฺญาณสมงฺคิสฺส อตฺถิ มคฺคภาวนา, จกฺขุวิญฺญาณํ อตีตานาคตํ อารพฺภ อุปฺปชฺชตีติ, น เหวํ วตฺตพฺเพ ฯเปฯ.\csromanspacer{} มโนวิญฺญาณสมงฺคิสฺส อตฺถิ มคฺคภาวนา, มโนวิญฺญาณํ ผสฺสํ อารพฺภ.\csromanspacer{} เวทนํ อารพฺภ ฯเปฯ.\csromanspacer{} โผฏฺฐพฺพํ อารพฺภ อุปฺปชฺชตีติ, อามนฺตา.\csromanspacer{} จกฺขุวิญฺญาณสมงฺคิสฺส อตฺถิ มคฺคภาวนา, จกฺขุวิญฺญาณํ ผสฺสํ อารพฺภ เวทนํ อารพฺภ ฯเปฯ.\csromanspacer{} โผฏฺฐพฺพํ อารพฺภ อุปฺปชฺชตีติ, น เหวํ วตฺตพฺเพ ฯเปฯ.}

\proseitem{579}{\csromansymbolmark{+}น วตฺตพฺพํ “ปญฺจวิญฺญาณสมงฺคิสฺส อตฺถิ มคฺคภาวนา”ติ, อามนฺตา.\csromanspacer{} นนุ วุตฺตํ ภควตา “อิธ ภิกฺขเว ภิกฺขุ จกฺขุนา รูปํ ทิสฺวา น นิมิตฺตคฺคาหี โหติ นานุพฺยญฺชนคฺคาหี ฯเปฯ โสเตน สทฺทํ สุตฺวา ฯเปฯ ฆาเนน คนฺธํ ฆายิตฺวา ฯเปฯ ชิวฺหาย รสํ สายิตฺวา ฯเปฯ กาเยน โผฏฺฐพฺพํ ผุสิตฺวา น นิมิตฺตคฺคาหี โหติ นานุพฺยญฺชนคฺคาหี”ติ\footnote{ม 1. 285; อํ 1. 348 ปิฏฺเฐ.} อตฺเถว สุตฺตนฺโตติ, อามนฺตา.\csromanspacer{} เตน หิ ปญฺจวิญฺญาณสมงฺคิสฺส อตฺถิ มคฺคภาวนาติ.}

\nitthitam{ปญฺจวิญฺญาณสมงฺคิสฺส มคฺคกถา นิฏฺฐิตา.}
\chapterhead{10. ทสมวคฺค}
\vspace{\tipitakaparskip}
\csromancenter{\textbf{(98) 4.}\csromanspacer{}\textbf{ ปญฺจวิญฺญาณากุสลาปิ-อกุสลาปีติกถา}}
\proseitem{580}{\csromansymbolmark{*}ปญฺจวิญฺญาณา กุสลาปิ อกุสลาปีติ, อามนฺตา.\csromanspacer{} นนุ ปญฺจวิญฺญาณา อุปฺปนฺนวตฺถุกา อุปฺปนฺนารมฺมณาติ, อามนฺตา.\csromanspacer{} หญฺจิ ปญฺจวิญฺญาณา}

\vspace{\tipitakaparskip}
\csromancenter{ทสมวคฺค อุปฺปนฺนวตฺถุกา อุปฺปนฺนารมฺมณา, โน จ วต เร วตฺตพฺเพ “ปญฺจวิญฺญาณา กุสลาปิ อกุสลาปี”ติ.\csromanspacer{} นนุ ปญฺจวิญฺญาณา ปุเรชาตวตฺถุกา ปุเรชาตารมฺมณา อชฺฌตฺติกวตฺถุกา พาหิรารมฺมณา.\csromanspacer{} อสมฺภินฺนวตฺถุกา อสมฺภินฺนารมฺมณา นานาวตฺถุกา นานารมฺมณา.\csromanspacer{} น อญฺญมญฺญสฺส โคจรวิสยํ ปจฺจนุโภนฺติ.\csromanspacer{} น อสมนฺนาหารา อุปฺปชฺชนฺติ น อมนสิการา อุปฺปชฺชนฺติ น อพฺโพกิณฺณา อุปฺปชฺชนฺติ.\csromanspacer{} น อปุพฺพํ อจริมํ อุปฺปชฺชนฺติ.\csromanspacer{} น อญฺญมญฺญสฺส สมนนฺตรา อุปฺปชฺชนฺติ.\csromanspacer{} นนุ ปญฺจวิญฺญาณา อนาโภคาติ, อามนฺตา.\csromanspacer{} หญฺจิ ปญฺจวิญฺญาณา อนาโภคา, โน จ วต เร วตฺตนฺเน “ปญฺจวิญฺญาณา กุสลาปิ อกุสลาปี”ติ.}
\proseitem{581}{\csromanfolio{313}\csromansymbolmark{*}จกฺขุวิญฺญาณํ กุสลนฺติ, อามนฺตา.\csromanspacer{} จกฺขุวิญฺญาณํ สุญฺญต อารพฺภ อุปฺปชฺชตีติ, น เหวํ วตฺตพฺเพ ฯเปฯ.\csromanspacer{} จกฺขุวิญฺญาณํ สุญฺญตํ อารพฺภ อุปฺปชฺชตีติ, อามนฺตา.\csromanspacer{} จกฺขุญฺจ ปฏิจฺจ สุญฺญตญฺจ อุปฺปชฺชติ จกฺขุวิญฺญาณนฺติ, น เหวํ วตฺตพฺเพ, จกฺขุญฺจ ปฏิจฺจ สุญฺญตญฺจ อุปฺปชฺชติ จกฺขุวิญฺญาณนฺติ, อามนฺตา. “จกฺขุญฺจ ปฏิจฺจ สุญฺญตญฺจ อุปฺปชฺชติ จกฺขุวิญฺญาณนฺ”ติ อตฺเถว สุตฺตนฺโตติ, นตฺถิ. “จกฺขุญฺจ ปฏิจฺจ รูเป จ อุปฺปชฺชติ จกฺขุวิญฺญาณนฺ”ติ อตฺเถว สุตฺตนฺโตติ, อามนฺตา.\csromanspacer{} หญฺจิ จกฺขุญฺจ ปฏิจฺจ รูเป จ อุปฺปชฺชติ จกฺขุวิญฺญาณนฺติ อตฺเถว สุตฺตนฺโต, โน จ วต เร วตฺตพฺเพ “จกฺขุญฺจ ปฏิจฺจ สุญฺญตญฺจ อุปฺปชฺชติ จกฺขุวิญฺญาณนฺ”ติ.}

\proseitem{582}{\csromansymbolmark{*}จกฺขุวิญฺญาณํ กุสลมฺปิ อกุสลมฺปีติ, อามนฺตา.\csromanspacer{} จกฺขุวิญฺญาณํ อตีตานาคตํ อารพฺภ อุปฺปชฺชตีติ, น เหวํ วตฺตพฺเพ ฯเปฯ.\csromanspacer{} จกฺขุวิญฺญาณํ กุสลมฺปิ อกุสลมฺปีติ, อามนฺตา.\csromanspacer{} จกฺขุวิญฺญาณํ ผสฺสํ อารพฺภ ฯเปฯ จิตฺตํ อารพฺภ.\csromanspacer{} จกฺขุํ อารพฺภ ฯเปฯ กาย อารพฺภ.\csromanspacer{} สทฺทํ อารพฺภ ฯเปฯ โผฏฺฐพฺพํ อารพฺภ อุปฺปชฺชตีติ.\csromanspacer{} น เหวํ วตฺตพฺเพ ฯเปฯ.}

\csromanmatikaanchor{mtk.127}
\csromantocmarkref{subsection}{(99) 5. สาโภคาติกถา}{mtk.127}
\bookmark[dest={mtk.127},level=2]{(99) 5. สาโภคาติกถา}
\prose{\csromansymbolmark{*}มโนวิญฺญาณํ กุสลมฺปิ อกุสลมฺปิ, มโนวิญฺญาณํ สุญฺญตํ อารพฺภ อุปฺปชฺชตีติ, อามนฺตา.\csromanspacer{} จกฺขุวิญฺญาณํ กุสลมฺปิ อกุสลมฺปิ, จกฺขุวิญฺญาณํ สุญฺญตํ อารพฺภ อุปฺปชฺชตีติ, น เหวํ วตฺตพฺเพ ฯเปฯ.\csromanspacer{} มโนวิญฺญาณํ กุสลมฺปิ อกุสลมฺปิ, มโนวิญฺญาณํ อตีตานาคตํ อารพฺภ อุปฺปชฺชตีติ, อามนฺตา.\csromanspacer{} จกฺขุวิญฺญาณํ กุสลมฺปิ อกุสลมฺปิ, จกฺขุวิญฺญาณํ อตีตานาคตํ อารพฺภ อุปฺปชฺชตีติ, น เหวํ วตฺตพฺเพ ฯเปฯ.\csromanspacer{} มโนวิญฺญาณํ กุสลมฺปิ อกุสลมฺปิ, มโนวิญฺญาณํ ผสฺสํ อารพฺภ ฯเปฯ จิตฺตํ อารพฺภ ฯเปฯ กายํ อารพฺภ.\csromanspacer{} สทฺทํ \csromanfolio{314}อารพฺภ ฯเปฯ โผฏฺฐพฺพํ อารพฺภ อุปฺปชฺชตีติ, อามนฺตา.\csromanspacer{} จกฺขุวิญฺญาณํ กุสลมฺปิ อกุสลมฺปิ, จกฺขุวิญฺญาณํ ผสฺสํ อารพฺภ ฯเปฯ โผฏฺฐพฺพํ อารพฺภ อุปฺปชฺชตีติ, น เหวํ วตฺตพฺเพ ฯเปฯ.}

\proseitem{583}{\csromansymbolmark{+}น วตฺตพฺพํ “ปญฺจวิญฺญาณา กุสลาปิ อกุสลาปี”ติ, อามนฺตา.\csromanspacer{} นนุ วุตฺตํ ภควตา “อิธ ภิกฺขเว ภิกฺขุ จกฺขุนา รูปํ ทิสฺวา นิมิตฺตคฺคาหี โหติ ฯเปฯ น นิมิตฺตคฺคาหี โหติ ฯเปฯ โสเตน สทฺทํ สุตฺวา ฯเปฯ กาเยน โผฏฺฐพฺพํ ผุสิตฺวา นิมิตฺตคฺคาหี โหติ ฯเปฯ น นิมิตฺตคฺคาหี โหตี”ติ อตฺเถว สุตฺตนฺโตติ, อามนฺตา.\csromanspacer{} เตน หิ ปญฺจวิญฺญาณา กุสลาปิ อกุสลาปีติ.}

\nitthitam{ปญฺจวิญฺญาณา กุสลาปิ อกุสลาปีติกถา นิฏฺฐิตา.}
\chapterhead{10. ทสมวคฺค}
\vspace{\tipitakaparskip}
\csromancenter{\textbf{(99) 5.}\csromanspacer{}\textbf{ สาโภคาติกถา}}
\proseitem{584}{\csromansymbolmark{*}ปญฺจวิญฺญาณา สาโภคาติ, อามนฺตา.\csromanspacer{} นนุ ปญฺจวิญฺญาณา อุปฺปนฺนวตฺถุกา อุปฺปนฺนารมฺมณาติ, อามนฺตา.\csromanspacer{} หญฺจิ ปญฺจวิญฺญาณา อุปฺปนฺนวตฺถุกา อุปฺปนฺนารมฺมณา, โน จ วต เร วตฺตพฺเพ “ปญฺจวิญฺญาณา สาโภคา”ติ.\csromanspacer{} นนุ ปญฺจวิญฺญาณา ปุเรชาตวตฺถุกา ปุเรชาตฺèรมฺมณา อชฺฌตฺติกวตฺถุกา พาหิรารมฺมณา อสมฺภินฺนวตฺถุกา อสมฺภินฺนารมฺมณา นานาวตฺถุกา นานารมฺมณา น อญฺญมญฺญสฺส โคจรวิสยํ ปจฺจนุโภนฺติ, น อสมนฺนาหารา อุปฺปชฺชนฺติ, น อมนสิการา อุปฺปชฺชนฺติ, น อพฺโพกิณฺณา อุปฺปชฺชนฺติ, น อปุพฺพํ อจริมํ อุปฺปชฺชนฺติ.\csromanspacer{} นนุ ปญฺจวิญฺญาณา น อญฺญมญฺญสฺส สมนนฺตรา อุปฺปชฺชนฺตีติ, อามนฺตา.\csromanspacer{} หญฺจิ ปญฺจวิญฺญาณา น อญฺญมญฺญสฺส สมนนฺตรา อุปฺปชฺชนฺติ, โน จ วต เร วตฺตพฺเพ “ปญฺจวิญฺญาณา สาโภคา”ติ.}

\proseitem{585}{\csromansymbolmark{*}จกฺขุวิญฺญาณํ สาโภคนฺติ, อามนฺตา.\csromanspacer{} จกฺขุวิญฺญาณํ สุญฺญตํ อารพฺภ อุปฺปชฺชตีติ, น เหวํ วตฺตพฺเพ ฯเปฯ.\csromanspacer{} จกฺขุวิญฺญาณํ สุญฺญตํ อารพฺภ อุปฺปชฺชตีติ, อามนฺตา.\csromanspacer{} จกฺขุญฺจ ปฏิจฺจ สุญฺญตญฺจ อุปฺปชฺชติ จกฺขุวิญฺญาณนฺติ, น เหวํ วตฺตพฺเพ ฯเปฯ.\csromanspacer{} จกฺขุญฺจ ปฏิจฺจ สุญฺญตญฺจ อุปฺปชฺชติ จกฺขุวิญฺญาณนฺติ, อามนฺตา.}

\vspace{\tipitakaparskip}
\csromancenter{ทสมวคฺค “จกฺขุญฺจ ปฏิจฺจ สุญฺญตญฺจ อุปฺปชฺชติ จกฺขุวิญฺญาณนฺ”ติ อตฺเถว สุตฺตนฺโตติ, นตฺถิ ฯเปฯ. “จกฺขุญฺจ ปฏิจฺจ รูเป จ อุปฺปชฺชติ จกฺขุวิญฺญาณนฺ”ติ อตฺเถว สุตฺตนฺโตติ, อามนฺตา.\csromanspacer{} หญฺจิ “จกฺขุญฺจ ปฏิจฺจ รูเป จ อุปฺปชฺชติ จกฺขุวิญฺญาณนฺ”ติ อตฺเถว สุตฺตนฺโต, โน จ วต เร วตฺตพฺเพ “จกฺขุญฺจ ปฏิจฺจ สุญฺญตญฺจ อุปฺปชฺชติ จกฺขุวิญฺญาณนฺ”ติ.}
\csromanmatikaanchor{mtk.128}
\csromantocmarkref{subsection}{(100) 6. ทฺวีหิสีเลหิกถา}{mtk.128}
\bookmark[dest={mtk.128},level=2]{(100) 6. ทฺวีหิสีเลหิกถา}
\prose{\csromanfolio{315}\csromansymbolmark{*}จกฺขุวิญฺญาณํ สาโภคนฺติ, อามนฺตา.\csromanspacer{} จกฺขุวิญฺญาณํ อตีตานาคตํ อารพฺภ อุปฺปชฺชตีติ, น เหวํ วตฺตพฺเพ ฯเปฯ.\csromanspacer{} จกฺขุวิญฺญาณํ สาโภคนฺติ, อามนฺตา.\csromanspacer{} จกฺขุวิญฺญาณํ ผสฺสํ อารพฺภ ฯเปฯ โผฏฺฐพฺพํ อารพฺภ อุปฺปชฺชตีติ, น เหวํ วตฺตพฺเพ ฯเปฯ.\csromanspacer{} มโนวิญฺญาณํ สาโภคํ มโนวิญฺญาณํ สุญฺญตํ อารพฺภ อุปฺปชฺชตีติ, อามนฺตา.\csromanspacer{} จกฺขุวิญฺญาณํ สาโภคํ จกฺขุวิญฺญาณํ สุญฺญตํ อารพฺภ อุปฺปชฺชตีติ, น เหวํ วตฺตพฺเพ ฯเปฯ.\csromanspacer{} มโนวิญฺญาณํ สาโภคํ มโนวิญฺญาณํ อตีตานาคตํ อารพฺภ อุปฺปชฺชตีติ, อามนฺตา.\csromanspacer{} จกฺขุวิญฺญาณํ สาโภคํ จกฺขุวิญฺญาณํ อตีตานาคตํ อารพฺภ อุปฺปชฺชตีติ, น เหวํ วตฺตพฺเพ ฯเปฯ.\csromanspacer{} มโนวิญฺญาณํ สาโภคํ มโนวิญฺญาณํ ผสฺสํ อารพฺภ ฯเปฯ โผฏฺฐพฺพํ อารพฺภ อุปฺปชฺชตีติ, อามนฺตา.\csromanspacer{} จกฺขุวิญฺญาณํ สาโภคํ จกฺขุวิญฺญาณํ ผสฺสํ อารพฺภ ฯเปฯ โผฏฺฐพฺพํ อารพฺภ อุปฺปชฺชตีติ, น เหวํ วตฺตพฺเพ ฯเปฯ.}

\proseitem{586}{\csromansymbolmark{+}น วตฺตพฺพํ “ปญฺจวิญฺญาณา สาโภคา”ติ, อามนฺตา.\csromanspacer{} นนุ วุตฺตํ ภควตา “อิธ ภิกฺขเว ภิกฺขุ จกฺขุนา รูปํ ทิสฺวา นิมิตฺตคฺคาหี โหติ ฯเปฯ น นิมิตฺตคฺคาหี โหติ ฯเปฯ กาเยน โผฏฺฐพฺพํ ผุสิตฺวา นิมิตฺตคฺคาหี โหติ ฯเปฯ น นิมิตฺตคฺคาหี โหตี”ติ อตฺเถว สุตฺตนฺโตติ, อามนฺตา.\csromanspacer{} เตน หิ ปญฺจวิญฺญาณา สาโภคาติ ฯเปฯ.}

\nitthitam{สาโภคาติกถา นิฏฺฐิตา.}
\chapterhead{10. ทสมวคฺค}
\vspace{\tipitakaparskip}
\csromancenter{\textbf{(100) 6.}\csromanspacer{}\textbf{ ทฺวีหิสีเลหิกถา}}
\proseitem{587}{\csromansymbolmark{*}มคฺคสมงฺคี ทฺวีหิ สีเลหิ สมนฺนาคโตติ, อามนฺตา.\csromanspacer{} มคฺคสมงฺคี ทฺวีหิ ผสฺเสหิ ทฺวีหิ เวทนาหิ ทฺวีหิ สญฺญาหิ ทฺวีหิ เจตนาหิ ทฺวีหิ \csromanfolio{316}จิตฺเตหิ ทฺวีหิ สทฺธาหิ ทฺวีหิ วีริเยหิ ทฺวีหิ สตีหิ ทฺวีหิ สมาธีหิ ทฺวีหิ ปญฺญาหิ สมนฺนาคโตติ, น เหวํ วตฺตพฺเพ ฯเปฯ.}

\prose{\csromansymbolmark{*}มคฺคสมงฺคี โลกิเยน สีเลน สมนฺนาคโตติ, อามนฺตา.\csromanspacer{} มคฺคสมงฺคี โลกิเยน ผสฺเสน โลกิยาย เวทนาย โลกิยาย ปญฺญาย โลกิยาย เจตนาย โลกิเยน จิตฺเตน โลกิยาย สทฺธาย โลกิเยน วีริเยน โลกิยาย สติยา โลกิเยน สมาธินา โลกิยาย ปญฺญาย สมนฺนาคโตติ, น เหวํ วตฺตพฺเพ ฯเปฯ.}

\prose{\csromansymbolmark{*}มคฺคสมงฺคี โลกิเยน จ โลกุตฺตเรน จ สีเลน สมนฺนาคโตติ, อามนฺตา.\csromanspacer{} มคฺคสมงฺคี โลกิเยน จ โลกุตฺตเรน จ ผสฺเสน สมนฺนาคโต ฯเปฯ โลกิยาย จ โลกุตฺตราย จ ปญฺญาย สมนฺนาคโตติ, น เหวํ วตฺตพฺเพ ฯเปฯ.\csromanspacer{} มคฺคสมงฺคี โลกิเยน สีเลน สมนฺนาคโตติ, อามนฺตา.\csromanspacer{} มคฺคสมงฺคี ปุถุชฺชโนติ, น เหวํ วตฺตพฺเพ ฯเปฯ.}

\proseitem{588}{\csromansymbolmark{*}มคฺคสมงฺคี โลกิยาย สมฺมาวาจาย สมนฺนาคโตติ, อามนฺตา.\csromanspacer{} มคฺคสมงฺคี โลกิยาย สมฺมาทิฏฺฐิยา สมนฺนาคโตติ, น เหวํ วตฺตพฺเพ ฯเปฯ.\csromanspacer{} มคฺคสมงฺคี โลกิยาย สมฺมาวาจาย สมนฺนาคโตติ, อามนฺตา.\csromanspacer{} มคฺคสมงฺคี โลกิเยน สมฺมาสงฺกปฺเปน ฯเปฯ โลกิเยน สมฺมาวายาเมน ฯเปฯ โลกิยาย สมฺมาสติยา ฯเปฯ โลกิเยน สมฺมาสมาธินา สมนฺนาคโตติ, น เหวํ วตฺตพฺเพ ฯเปฯ.\csromanspacer{} มคฺคสมงฺคี โลกิเยน สมฺมากมฺมนฺเตน ฯเปฯ โลกิเยน สมฺมา อาชีเวน สมนฺนาคโตติ, อามนฺตา.\csromanspacer{} มคฺคสมงฺคี โลกิยาย สมฺมาทิฏฺฐิยา ฯเปฯ โลกิเยน สมฺมาสมาธินา สมนฺนาคโตติ, น เหวํ วตฺตพฺเพ ฯเปฯ.}

\prose{\csromansymbolmark{*}มคฺคสมงฺคี โลกิยาย จ โลกุตฺตราย จ สมฺมาวาจาย สมนฺนาคโตติ, อามนฺตา.\csromanspacer{} มคฺคสมงฺคี โลกิยาย จ โลกุตฺตราย จ สมฺมาทิฏฺฐิยา สมนฺนาคโตติ, น เหวํ วตฺตพฺเพ ฯเปฯ.\csromanspacer{} มคฺคสมงฺคี โลกิยาย จ โลกุตฺตราย จ สมฺมาวาจาย สมนฺนาคโตติ, อามนฺตา.\csromanspacer{} มคฺคสมงฺคี โลกิเยน จ โลกุตฺตเรน จ สมฺมาสงฺกปฺเปน ฯเปฯ โลกิเยน จ โลกุตฺตเรน จ สมฺมาวายาเมน ฯเปฯ โลกิยาย จ โลกุตฺตราย จ สมฺมาสติยา ฯเปฯ โลกิเยน จ โลกุตฺตเรน จ สมฺมาสมาธินา สมนฺนาคโตติ, น เหวํ วตฺตพฺเพ ฯเปฯ.}

\chapterheadpageread{10. ทสมวคฺค}
\csromanmatikaanchor{mtk.129}
\csromantocmarkref{subsection}{(101) 7. สีลํอเจตสิกนฺติกถา}{mtk.129}
\bookmark[dest={mtk.129},level=2]{(101) 7. สีลํอเจตสิกนฺติกถา}
\prose{\csromanfolio{317}\csromansymbolmark{*}มคฺคสมงฺคี โลกิเยน จ โลกุตฺตเรน จ สมฺมากมฺมนฺเตน ฯเปฯ สมฺมา อาชีเวน สมนฺนาคโตติ, อามนฺตา.\csromanspacer{} มคฺคสมงฺคี โลกิยาย จ โลกุตฺตราย จ สมฺมาทิฏฺฐิยา สมนฺนาคโตติ, น เหวํ วตฺตพฺเพ ฯเปฯ.\csromanspacer{} มคฺคสมงฺคี โลกิเยน จ โลกุตฺตเรน จ สมฺมา-อาชีเวน สมนฺนาคโตติ, อามนฺตา.\csromanspacer{} มคฺคสมงฺคี โลกิเยน จ โลกุตฺตเรน จ สมฺมาสงฺกปฺเปน ฯเปฯ โลกิเยน จ โลกุตฺตเรน จ สมฺมาสมาธินา สมนฺนาคโตติ, น เหวํ วตฺตพฺเพ ฯเปฯ.}

\proseitem{589}{\csromansymbolmark{+}น วตฺตพฺพํ “มคฺคสมงฺคี ทฺวีหิ สีเลหิ สมนฺนาคโต”ติ, อามนฺตา.\csromanspacer{} โลกิเย สีเล นิรุทฺเธ มคฺโค อุปฺปชฺชตีติ, อามนฺตา.\csromanspacer{} ทุสฺสีโล ขณฺฑสีโล ฉินฺนสีโล มคฺคํ ภาเวตีติ, น เหวํ วตฺตพฺเพ ฯเปฯ.\csromanspacer{} เตน หิ มคฺคสมงฺคี ทฺวีหิ สีเลหิ สมนฺนาคโตติ.}

\nitthitam{ทฺวีหิสีเลหิกถา นิฏฺฐิตา.}
\chapterhead{10. ทสมวคฺค}
\vspace{\tipitakaparskip}
\csromancenter{\textbf{(101) 7.}\csromanspacer{}\textbf{ สีลํอเจตสิกนฺติกถา}}
\proseitem{590}{\csromansymbolmark{*}สีลํ อเจตสิกนฺติ, อามนฺตา.\csromanspacer{} รูปํ.\csromanspacer{} นิพฺพานํ.\csromanspacer{} จกฺขายตนํ ฯเปฯ.\csromanspacer{} กายายตนํ.\csromanspacer{} รูปายตนํ ฯเปฯ.\csromanspacer{} โผฏฺฐพฺพายตนนฺติ, น เหวํ วตฺตพฺเพ ฯเปฯ.\csromanspacer{} สีลํ อเจตสิกนฺติ, อามนฺตา.\csromanspacer{} ผสฺโส อเจตสิโกติ, น เหวํ วตฺตพฺเพ ฯเปฯ.\csromanspacer{} สีลํ อเจตสิกนฺติ, อามนฺตา.\csromanspacer{} เวทนา ฯเปฯ.\csromanspacer{} สญฺญา ฯเปฯ.\csromanspacer{} เจตนา ฯเปฯ.\csromanspacer{} สทฺธา ฯเปฯ.\csromanspacer{} วีริยํ ฯเปฯ.\csromanspacer{} สติ ฯเปฯ.\csromanspacer{} สมาธิ ฯเปฯ.\csromanspacer{} ปญฺญา อเจตสิกาติ, น เหวํ วตฺตพฺเพ ฯเปฯ.}

\prose{\csromansymbolmark{*}ผสฺโส เจตสิโกติ, อามนฺตา.\csromanspacer{} สีลํ เจตสิกนฺติ, น เหวํ วตฺตพฺเพ ฯเปฯ.\csromanspacer{} เวทนา ฯเปฯ.\csromanspacer{} สญฺญา ฯเปฯ.\csromanspacer{} เจตนา ฯเปฯ.\csromanspacer{} สทฺธา ฯเปฯ.\csromanspacer{} วีริยํ ฯเปฯ.\csromanspacer{} สติ ฯเปฯ.\csromanspacer{} สมาธิ ฯเปฯ.\csromanspacer{} ปญฺญา เจตสิกาติ, อามนฺตา.\csromanspacer{} สีลํ เจตสิกนฺติ, น เหวํ วตฺตพฺเพ ฯเปฯ.}

\proseitem{591}{\csromansymbolmark{*}สีลํ อเจตสิกนฺติ, อามนฺตา.\csromanspacer{} อนิฏฺฐผลนฺติ, น เหวํ วตฺตพฺเพ ฯเปฯ.\csromanspacer{} นนุ อิฏฺฐผลนฺติ, อามนฺตา.\csromanspacer{} หญฺจิ อิฏฺฐผลํ, โน จ วต เร วตฺตพฺเพ “สีลํ อเจตสิกนฺ”ติ.}

\prose{\csromanfolio{318}\csromansymbolmark{*}สทฺธา อิฏฺฐผลา, สทฺธา เจตสิกาติ, อามนฺตา.\csromanspacer{} สีลํ อิฏฺฐผลํ, สีลํ เจตสิกนฺติ, น เหวํ วตฺตพฺเพ ฯเปฯ.\csromanspacer{} วีริยํ ฯเปฯ.\csromanspacer{} สติ ฯเปฯ.\csromanspacer{} สมาธิ ฯเปฯ.\csromanspacer{} ปญฺญา อิฏฺฐผลา, ปญฺญา เจตสิกาติ, อามนฺตา.\csromanspacer{} สีลํ อิฏฺฐผลํ, สีลํ เจตสิกนฺติ, น เหวํ วตฺตพฺเพ ฯเปฯ.}

\prose{\csromansymbolmark{*}สีลํ อิฏฺฐผลํ, สีลํ อเจตสิกนฺติ, อามนฺตา.\csromanspacer{} สทฺธา อิฏฺฐผลา, สทฺธา อเจตสิกาติ, น เหวํ วตฺตพฺเพ ฯเปฯ.\csromanspacer{} สีลํ อิฏฺฐผลํ, สีลํ อเจตสิกนฺติ, อามนฺตา.\csromanspacer{} วีริยํ ฯเปฯ.\csromanspacer{} สติ ฯเปฯ.\csromanspacer{} สมาธิ ฯเปฯ.\csromanspacer{} ปญฺญา อิฏฺฐผลา, ปญฺญา อเจตสิกาติ, น เหวํ วตฺตพฺเพ ฯเปฯ.}

\proseitem{592}{\csromansymbolmark{*}สีลํ อเจตสิกนฺติ, อามนฺตา.\csromanspacer{} อผลํ อวิปากนฺติ, น เหวํ วตฺตพฺเพ ฯเปฯ.\csromanspacer{} นนุ สผลํ สวิปากนฺติ, อามนฺตา.\csromanspacer{} หญฺจิ สผลํ สวิปากํ, โน จ วต เร วตฺตพฺเพ “สีลํ อเจตสิกนฺ”ติ ฯเปฯ.}

\prose{\csromansymbolmark{*}จกฺขายตนํ อเจตสิกํ อวิปากนฺติ, อามนฺตา.\csromanspacer{} สีลํ อเจตสิกํ อวิปากนฺติ, น เหวํ วตฺตพฺเพ ฯเปฯ.\csromanspacer{} โสตายตนํ ฯเปฯ.\csromanspacer{} กายายตนํ ฯเปฯ.\csromanspacer{} รูปายตนํ ฯเปฯ.\csromanspacer{} โผฏฺฐพฺพายตนํ อเจตสิกํ อวิปากนฺติ, อามนฺตา.\csromanspacer{} สีลํ อเจตสิกํ อวิปากนฺติ, น เหวํ วตฺตพฺเพ ฯเปฯ.}

\prose{\csromansymbolmark{*}สีลํ อเจตสิกํ สวิปากนฺติ, อามนฺตา.\csromanspacer{} จกฺขายตนํ อเจตสิกํ สวิปากนฺติ, น เหวํ วตฺตพฺเพ ฯเปฯ.\csromanspacer{} สีลํ อเจตสิกํ สวิปากนฺติ, อามนฺตา.\csromanspacer{} โสตายตนํ ฯเปฯ.\csromanspacer{} กายายตนํ.\csromanspacer{} รูปายตนํ ฯเปฯ.\csromanspacer{} โผฏฺฐพฺพายตนํ อเจตสิกํ สวิปากนฺติ, น เหวํ วตฺตพฺเพ ฯเปฯ.}

\proseitem{593}{\csromansymbolmark{*}สมฺมาวาจา อเจตสิกาติ, อามนฺตา.\csromanspacer{} สมฺมาทิฏฺฐิ อเจตสิกาติ, น เหวํ วตฺตพฺเพ ฯเปฯ.\csromanspacer{} สมฺมาวาจา อเจตสิกาติ, อามนฺตา.\csromanspacer{} สมฺมาสงฺกปฺโป ฯเปฯ.\csromanspacer{} สมฺมาวายาโม ฯเปฯ.\csromanspacer{} สมฺมาสติ ฯเปฯ.\csromanspacer{} สมฺมาสมาธิ อเจตสิโกติ, น เหวํ วตฺตพฺเพ ฯเปฯ.\csromanspacer{} สมฺมากมฺมนฺโต ฯเปฯ.\csromanspacer{} สมฺมา อาชีโว อเจตสิโกติ, อามนฺตา.\csromanspacer{} สมฺมาทิฏฺฐิ อเจตสิกาติ, น เหวํ วตฺตพฺเพ ฯเปฯ.\csromanspacer{} สมฺมา อาชีโว อเจตสิโกติ, อามนฺตา.\csromanspacer{} สมฺมาสงฺกปฺโป.\csromanspacer{} สมฺมาวายาโม.\csromanspacer{} สมฺมาสติ.\csromanspacer{} สมฺมาสมาธิ อเจตสิโกติ, น เหวํ วตฺตพฺเพ ฯเปฯ.}

\prose{\csromansymbolmark{*}สมฺมาทิฏฺฐิ เจตสิกาติ, อามนฺตา.\csromanspacer{} สมฺมาวาจา เจตสิกาติ, น เหวํ วตฺตพฺเพ ฯเปฯ.\csromanspacer{} สมฺมาทิฏฺฐิ เจตสิกาติ, อามนฺตา.\csromanspacer{} สมฺมากมฺมนฺโต ฯเปฯ.}

\vspace{\tipitakaparskip}
\csromancenter{ทสมวคฺค สมฺมา อาชีโว เจตสิโกติ, น เหวํ วตฺตพฺเพ ฯเปฯ.\csromanspacer{} สมฺมาสงฺกปฺโป สมฺมาวายาโม ฯเปฯ.\csromanspacer{} สมฺมาสติ ฯเปฯ.\csromanspacer{} สมฺมาสมาธิ เจตสิโกติ, อามนฺตา.\csromanspacer{} สมฺมาวาจา เจตสิกาติ, น เหวํ วตฺตพฺเพ ฯเปฯ.\csromanspacer{} สมฺมาสมาธิ เจตสิโกติ, อามนฺตา.\csromanspacer{} สมฺมากมฺมนฺโต ฯเปฯ.\csromanspacer{} สมฺมา อาชีโว เจตสิโกติ, น เหวํ วตฺตพฺเพ ฯเปฯ.}
\csromanmatikaanchor{mtk.130}
\csromantocmarkref{subsection}{(102) 8. สีลํนจิตฺตานุปริวตฺตีติกถา}{mtk.130}
\bookmark[dest={mtk.130},level=2]{(102) 8. สีลํนจิตฺตานุปริวตฺตีติกถา}
\proseitem{594}{\csromanfolio{319}\csromansymbolmark{+}น วตฺตพฺพํ “สีลํ อเจตสิกนฺ”ติ, อามนฺตา.\csromanspacer{} สีเล อุปฺปชฺชิตฺวา นิรุทฺเธ ทุสฺสีโล โหตีติ, น เหวํ วตฺตพฺเพ.\csromanspacer{} เตน หิ สีลํ อเจตสิกนฺติ.}

\nitthitam{สีลํ อเจตสิกนฺติกถา นิฏฺฐิตา.}
\chapterhead{10. ทสมวคฺค}
\vspace{\tipitakaparskip}
\csromancenter{\textbf{(102) 8.}\csromanspacer{}\textbf{ สีลํนจิตฺตานุปริวตฺตีติกถา}}
\proseitem{595}{\csromansymbolmark{*}สีลํ น จิตฺตานุปริวตฺตีติ, อามนฺตา.\csromanspacer{} รูปํ.\csromanspacer{} นิพฺพานํ.\csromanspacer{} จกฺขายตนํ ฯเปฯ.\csromanspacer{} โผฏฺฐพฺพายตนนฺติ, น เหวํ วตฺตพฺเพ ฯเปฯ.\csromanspacer{} สีลํ น จิตฺตานุปริวตฺตีติ, อามนฺตา.\csromanspacer{} ผสฺโส น จิตฺตานุปริวตฺตีติ, น เหวํ วตฺตพฺเพ ฯเปฯ.\csromanspacer{} สีลํ น จิตฺตานุปริวตฺตีติ, อามนฺตา.\csromanspacer{} เวทนา ฯเปฯ.\csromanspacer{} สญฺญา.\csromanspacer{} เจตนา.\csromanspacer{} สทฺธา.\csromanspacer{} วีริยํ.\csromanspacer{} สติ.\csromanspacer{} สมาธิ ฯเปฯ.\csromanspacer{} ปญฺญา น จิตฺตานุปริวตฺตีติ, น เหวํ วตฺตพฺเพ ฯเปฯ.}

\prose{\csromansymbolmark{*}ผสฺโส จิตฺตานุปริวตฺตีติ, อามนฺตา.\csromanspacer{} สีลํ จิตฺตานุปริวตฺตีติ, น เหวํ วตฺตพฺเพ ฯเปฯ.\csromanspacer{} เวทนา ฯเปฯ.\csromanspacer{} สญฺญา.\csromanspacer{} เจตนา.\csromanspacer{} สทฺธา.\csromanspacer{} วีริยํ.\csromanspacer{} สติ.\csromanspacer{} สมาธิ ฯเปฯ.\csromanspacer{} ปญฺญา จิตฺตานุปริวตฺตีติ, อามนฺตา.\csromanspacer{} สีลํ จิตฺตานุปริวตฺตีติ, น เหวํ วตฺตพฺเพ ฯเปฯ.}

\proseitem{596}{\csromansymbolmark{*}สมฺมาวาจา น จิตฺตานุปริวตฺตีติ, อามนฺตา.\csromanspacer{} สมฺมาทิฏฺฐิ น จิตฺตานุปริวตฺตีติ, น เหวํ วตฺตพฺเพ ฯเปฯ.\csromanspacer{} สมฺมาวาจา น จิตฺตานุปริวตฺตีติ, อามนฺตา.\csromanspacer{} สมฺมาสงฺกปฺโป ฯเปฯ.\csromanspacer{} สมฺมาวายาโม ฯเปฯ.\csromanspacer{} สมฺมาสติ ฯเปฯ.\csromanspacer{} สมฺมาสมาธิ น จิตฺตานุปริวตฺตีติ, น เหวํ วตฺตพฺเพ ฯเปฯ.}

\csromanmatikaanchor{mtk.131}
\csromantocmarkref{subsection}{(103) 9. สมาทานเหตุกถา}{mtk.131}
\bookmark[dest={mtk.131},level=2]{(103) 9. สมาทานเหตุกถา}
\prose{\csromansymbolmark{*}สมฺมากมฺมนฺโต ฯเปฯ.\csromanspacer{} สมฺมา อาชีโว น จิตฺตานุปริวตฺตีติ, อามนฺตา.\csromanspacer{} สมฺมาทิฏฺฐิ น จิตฺตานุปริวตฺตีติ, น เหวํ วตฺตพฺเพ ฯเปฯ.\csromanspacer{} สมฺมา อาชีโว น \csromanfolio{320}จิตฺตานุปริวตฺตีติ, อามนฺตา.\csromanspacer{} สมฺมาสงฺกปฺโป ฯเปฯ.\csromanspacer{} สมฺมาวายาโม ฯเปฯ.\csromanspacer{} สมฺมาสติ ฯเปฯ.\csromanspacer{} สมฺมาสมาธิ น จิตฺตานุปริวตฺตีติ, น เหวํ วตฺตพฺเพ ฯเปฯ.}

\prose{\csromansymbolmark{*}สมฺมาทิฏฺฐิ จิตฺตานุปริวตฺตีติ, อามนฺตา.\csromanspacer{} สมฺมาวาจา จิตฺตานุปริวตฺตีติ, น เหวํ วตฺตพฺเพ ฯเปฯ.\csromanspacer{} สมฺมาทิฏฺฐิ จิตฺตานุปริวตฺตีติ, อามนฺตา.\csromanspacer{} สมฺมากมฺมนฺโต ฯเปฯ.\csromanspacer{} สมฺมา อาชีโว จิตฺตานุปริวตฺตีติ.\csromanspacer{} น เหวํ วตฺตพฺเพ ฯเปฯ. \csromansymbolmark{*}สมฺมาสงฺกปฺโป ฯเปฯ.\csromanspacer{} สมฺมาวายาโม ฯเปฯ.\csromanspacer{} สมฺมาสติ ฯเปฯ.\csromanspacer{} สมฺมาสมาธิ จิตฺตานุปริวตฺตีติ, อามนฺตา.\csromanspacer{} สมฺมาวาจา จิตฺตานุปริวตฺตีติ, น เหวํ วตฺตพฺเพ ฯเปฯ.\csromanspacer{} สมฺมาสมาธิ จิตฺตานุปริวตฺตีติ, อามนฺตา.\csromanspacer{} สมฺมากมฺมนฺโต ฯเปฯ.\csromanspacer{} สมฺมา อาชีโว จิตฺตานุปริวตฺตีติ.\csromanspacer{} น เหวํ วตฺตพฺเพ ฯเปฯ.}

\proseitem{597}{\csromansymbolmark{+}น วตฺตพฺพํ “สีลํ น จิตฺตานุปริวตฺตี”ติ, อามนฺตา.\csromanspacer{} สีเล อุปฺปชฺชิตฺวา นิรุทฺเธ ทุสฺสีโล โหตีติ, น เหวํ วตฺตพฺเพ.\csromanspacer{} เตน หิ สีลํ น จิตฺตานุปริวตฺตีติ.}

\nitthitam{สีลํ น จิตฺตานุปริวตฺตีติกถา นิฏฺฐิตา.}
\chapterhead{10. ทสมวคฺค}
\vspace{\tipitakaparskip}
\csromancenter{\textbf{(103) 9.}\csromanspacer{}\textbf{ สมาทานเหตุกถา}}
\proseitem{598}{\csromansymbolmark{*}สมาทานเหตุกํ สีลํ วฑฺฒตีติ, อามนฺตา.\csromanspacer{} สมาทานเหตุโก ผสฺโส วฑฺฒติ เวทนา วฑฺฒติ สญฺญา วฑฺฒติ เจตนา วฑฺฒติ จิตฺตํ วฑฺฒติ สทฺธา วฑฺฒติ วีริยํ วฑฺฒติ สติ วฑฺฒติ สมาธิ วฑฺฒติ ปญฺญา วฑฺฒตีติ, น เหวํ วตฺตพฺเพ ฯเปฯ.}

\prose{\csromansymbolmark{*}สมาทานเหตุกํ สีลํ วฑฺฒตีติ, อามนฺตา.\csromanspacer{} ลตา วิย วฑฺฒติ มาลุวา วิย วฑฺฒติ รุกฺโข วิย วฑฺฒติ ติณํ วิย วฑฺฒติ มุญฺชปุญฺโช วิย วฑฺฒตีติ, น เหวํ วตฺตพฺเพ ฯเปฯ.}

\proseitem{599}{\csromansymbolmark{*}สมาทานเหตุกํ สีลํ วฑฺฒตีติ, อามนฺตา.\csromanspacer{} สีลํ สมาทิยิตฺวา กามวิตกฺกํ วิตกฺเกนฺตสฺส พฺยาปาทวิตกฺกํ วิตกฺเกนฺตสฺส วิหิํ สาวิตกฺกํ วิตกฺเกนฺตสฺส สีลํ วฑฺฒตีติ, อามนฺตา.\csromanspacer{} ทฺวินฺนํ ผสฺสานํ ฯเปฯ.\csromanspacer{} ทฺวินฺนํ จิตฺตานํ}

\vspace{\tipitakaparskip}
\csromancenter{ทสมวคฺค สโมธานํ โหตีติ, น เหวํ วตฺตพฺเพ ฯเปฯ.\csromanspacer{} ทฺวินฺนํ ผสฺสานํ ฯเปฯ.\csromanspacer{} ทฺวินฺนํ จิตฺตานํ สโมธานํ โหตีติ, อามนฺตา.\csromanspacer{} กุสลากุสลา สาวชฺชานวชฺชา หีนปณีตา กณฺหสุกฺกสปฺปฏิภาคา ธมฺมา สมฺมุขีภาวํ อาคจฺฉนฺตีติ, น เหวํ วตฺตพฺเพ ฯเปฯ.}
\csromanmatikaanchor{mtk.132}
\csromantocmarkref{subsection}{(104) 10. วิญฺญตฺติสีลนฺติกถา}{mtk.132}
\bookmark[dest={mtk.132},level=2]{(104) 10. วิญฺญตฺติสีลนฺติกถา}
\prose{\csromanfolio{321}\csromansymbolmark{*}กุสลากุสลา สาวชฺชานวชฺชา หีนปณีตา กณฺหสุกฺกสปฺปฏิภาคา ธมฺมา สมฺมุขีภาวํ อาคจฺฉนฺตีติ, อามนฺตา.\csromanspacer{} นนุ วุตฺตํ ภควตา “จตฺตาริมานิ ภิกฺขเว สุวิทูรวิทูรานิ.\csromanspacer{} กตมานิ จตฺตาริ, นภญฺจ ภิกฺขเว ปถวี จ อิทํ ปฐมํ สุวิทูรวิทูรํ ฯเปฯ.\csromanspacer{} ตสฺมา สตํ ธมฺโม อสพฺภิ อารกา”ติ อตฺเถว สุตฺตนฺโตติ, อามนฺตา.\csromanspacer{} เตน หิ น วตฺตพฺพํ “กุสลากุสลา สาวชฺชานวชฺชา หีนปณีตา กณฺหสุกฺกสปฺปฏิภาคา ธมฺมา สมฺมุขีภาวํ อาคจฺฉนฺตี”ติ.}

\proseitem{600}{\csromansymbolmark{+}น วตฺตพฺพํ “สมาทานเหตุกํ สีลํ วฑฺฒตี”ติ, อามนฺตา.\csromanspacer{} นนุ วุตฺตํ ภควตา “อารามโรปา วนโรปา ฯเปฯ ธมฺมฏฺฐา สีลสมฺปนฺนา, เต ชนา สคฺคคามิโน”ติ อตฺเถว สุตฺตนฺโตติ, อามนฺตา.\csromanspacer{} เตน หิ สมาทานเหตุกํ สีลํ วฑฺฒตีติ.}

\nitthitam{สมาทานเหตุกถา นิฏฺฐิตา.}
\chapterhead{10. ทสมวคฺค}
\vspace{\tipitakaparskip}
\csromancenter{\textbf{(104) 10.}\csromanspacer{}\textbf{ วิญฺญตฺติสีลนฺติกถา}}
\proseitem{601}{\csromansymbolmark{*}วิญฺญตฺติ สีลนฺติ, อามนฺตา.\csromanspacer{} ปาณาติปาตา เวรมณีติ, น เหวํ วตฺตพฺเพ ฯเปฯ.\csromanspacer{} อทินฺนาทานา เวรมณีติ, น เหวํ วตฺตพฺเพ ฯเปฯ.\csromanspacer{} กาเมสุมิจฺฉาจารา เวรมณีติ, น เหวํ วตฺตพฺเพ ฯเปฯ.\csromanspacer{} มุสาวาทา เวรมณีติ, น เหวํ วตฺตพฺเพ ฯเปฯ.\csromanspacer{} สุราเมรยมชฺชปมาทฏฺฐานา เวรมณีติ, น เหวํ วตฺตพฺเพ ฯเปฯ.}

\prose{\csromansymbolmark{*}อภิวาทนํ สีลํ, ปจฺจุฏฺฐานํ สีลํ, อญฺชลิกมฺมํ สีลํ, สามีจิกมฺมํ สีลํ, อาสนาภิหาโร สีลํ, เสยฺยาภิหาโร สีลํ, ปาโททกาภิหาโร สีลํ, ปาทกถลิกาภิหาโร สีลํ, นฺหาเน ปิฏฺฐิปริกมฺมํ สีลนฺติ, อามนฺตา.\csromanspacer{} ปาณาติปาตา เวรมณีติ, น เหวํ วตฺตพฺเพ ฯเปฯ.\csromanspacer{} สุราเมรยมชฺชปมาทฏฺฐานา เวรมณีติ, น เหวํ วตฺตพฺเพ ฯเปฯ.}

\csromanmatikaanchor{mtk.133}
\csromantocmarkref{subsection}{(105) 11. อวิญฺญตฺติทุสฺสิลฺยนฺติกถา}{mtk.133}
\bookmark[dest={mtk.133},level=2]{(105) 11. อวิญฺญตฺติทุสฺสิลฺยนฺติกถา}
\proseitem{602}{\csromanfolio{322}\csromansymbolmark{+}น วตฺตพฺพํ “วิญฺญตฺติ สีลนฺ”ติ, อามนฺตา.\csromanspacer{} ทุสฺสิลฺยนฺติ, น เหวํ วตฺตพฺเพ ฯเปฯ.\csromanspacer{} เตน หิ วิญฺญตฺติ สีลนฺตี.}

\nitthitam{วิญฺญตฺติ สีลนฺติกถา นิฏฺฐิตา.}
\chapterhead{10. ทสมวคฺค}
\vspace{\tipitakaparskip}
\csromancenter{\textbf{(105) 11.}\csromanspacer{}\textbf{ อวิญฺญตฺติทุสฺสิลฺยนฺติกถา}}
\proseitem{603}{\csromansymbolmark{*}อวิญฺญตฺติ ทุสฺสิลฺยนฺติ, อามนฺตา.\csromanspacer{} ปาณาติปาโตติ, น เหวํ วตฺตพฺเพ ฯเปฯ.\csromanspacer{} อทินฺนาทานนฺติ, น เหวํ วตฺตพฺเพ ฯเปฯ.\csromanspacer{} กาเมสุมิจฺฉาจาโรติ, น เหวํ วตฺตพฺเพ ฯเปฯ.\csromanspacer{} มุสาวาโทติ, น เหวํ วตฺตพฺเพ ฯเปฯ.\csromanspacer{} สุราเมรยมชฺชปมาทฏฺฐานนฺติ, น เหวํ วตฺตพฺเพ ฯเปฯ.}

\prose{\csromansymbolmark{*}ปาปกมฺมํ สมาทิยิตฺวา ทานํ ททนฺตสฺส ปุญฺญญฺจ อปุญฺญญฺจ อุโภ วฑฺฒนฺตีติ, น เหวํ วตฺตพฺเพ ฯเปฯ ปุญฺญญฺจ อปุญฺญญฺจ อุโภ วฑฺฒนฺตีติ, อามนฺตา.\csromanspacer{} ทฺวินฺนํ ผสฺสานํ ฯเปฯ.\csromanspacer{} ทฺวินฺนํ จิตฺตานํ สโมธานํ โหตีติ, น เหวํ วตฺตพฺเพ ฯเปฯ.\csromanspacer{} ทฺวินฺนํ ผสฺสานํ ฯเปฯ.\csromanspacer{} ทฺวินฺนํ จิตฺตานํ สโมธานํ โหตีติ, อามนฺตา.\csromanspacer{} กุสลากุสลา สาวชฺชานวชฺชา หีนปณีตา กณฺหสุกฺกสปฺปฏิภาคา ธมฺมา สมฺมุขีภาวํ อาคจฺฉนฺตีติ, น เหวํ วตฺตพฺเพ ฯเปฯ.}

\prose{\csromansymbolmark{*}กุสลากุสลา สาวชฺชานวชฺชา หีนปณีตา กณฺหสุกฺกสปฺปฏิภาคา ธมฺมา สมฺมุขีภาวํ อาคจฺฉนฺตีติ, อามนฺตา.\csromanspacer{} นนุ วุตฺตํ ภควตา “จตฺตาริมานิ ภิกฺขเว สุวิทูรวิทูรานิ.\csromanspacer{} กตมานิ จตฺตาริ, นภญฺจ ภิกฺขเว ปถวี จ อิทํ ปฐมํ สุวิทูรวิทูรํ ฯเปฯ.\csromanspacer{} ตสฺมา สตํ ธมฺโม อสพฺภิ อารกา”ติ อตฺเถว สุตฺตนฺโตติ, อามนฺตา.\csromanspacer{} เตน หิ น วตฺตพฺพํ “กุสลากุสลา สาวชฺชานวชฺชา หีนปณีตา กณฺหสุกฺกสปฺปฏิภาคา ธมฺมา สมฺมุขีภาวํ อาคจฺฉนฺตี”ติ.}

\proseitem{604}{\csromansymbolmark{+}น วตฺตพฺพํ “อวิญฺญตฺติ ทุสฺสิลฺยนฺ”ติ, อามนฺตา.\csromanspacer{} นนุ ปาปกมฺมํ สมาทินฺโน อาสีติ, อามนฺตา.\csromanspacer{} หญฺจิ ปาปกมฺมํ สมาทินฺโน อาสิ, เตน วต เร วตฺตพฺเพ “อวิญฺญตฺติ ทุสฺสิลฺยนฺ”ติ.}

\nitthitam{อวิญฺญตฺติ ทุสฺสิลฺยนฺติกถา นิฏฺฐิตา.}
\vspace{\tipitakaparskip}
\csromancenter{ทสโม วคฺโค.}
\chapterheadpageread{11. เอกาทสมวคฺค \textbf{ตสฺสุทฺทานํ}}
\csromanmatikaanchor{mtk.134}
\csromanmatikaanchor{mtk.135}
\csromantocmarknopage{chapter}{11. เอกาทสมวคฺค}{mtk.134}
\bookmark[dest={mtk.134},level=0]{11. เอกาทสมวคฺค}
\csromantocmarkref{subsection}{(106-8) 1-3. ติสฺโสปิ-อนุสยกถา}{mtk.135}
\bookmark[dest={mtk.135},level=2]{(106-8) 1-3. ติสฺโสปิ-อนุสยกถา}
\prose{\csromanfolio{323}อุปปตฺเตสิเย ปญฺจกฺขนฺเธ อนิรุทฺเธ กิริยา ปญฺจกฺขนฺธา อุปฺปชฺชนฺติ, มคฺคสมงฺคิสฺส รูปํ มคฺโค, ปญฺจวิญฺญาณสมงฺคิสฺส อตฺถิ มคฺคภาวนา, ปญฺจวิญฺญาณา กุสลาปิ อกุสลาปิ, ปญฺจวิญฺญาณา สาโภคา, มคฺคสมงฺคี ทฺวีหิ สีเลหิ สมนฺนาคโต, สีลํ อเจตสิกํ, สีลํ น จิตฺตานุปริวตฺติ, สมาทานเหตุกํ สีลํ วฑฺฒตีติ, วิญฺญตฺติสีลํ, อวิญฺญตฺติ ทุสฺสิลฺยนฺติ.\csromanspacer{} ทุติโย ปณฺณาสโก.}

\titlehead{\textbf{ตสฺสุทฺทานํ}}
\csromangathameasure{นิยาม สงฺคห คตานิสํสตา จ นิโรโธติ.}
\csromangathagroup{\csromangathastanza{นิยาม สงฺคห คตานิสํสตา จ นิโรโธติ\footnote{นิยาม สงฺคหคตานิสํสตา อตฺถิ ปญฺจโวการภโว นิโรโธ มคฺคสมงฺคิสฺส ปญฺจโวการภโวติ (สี, สฺยา, ก), นิยาม สงฺคห คตานิสํสตา อตฺถิ ปญฺจโวการภโว มคฺคสมงฺคิสฺส โน ปญฺจโวการภโว (อิ)}.}}

\chapterhead{11. เอกาทสมวคฺค}
\vspace{\tipitakaparskip}
\csromancenter{\textbf{(106-108) 1-3.}\csromanspacer{}\textbf{ ติสฺโสปิ-อนุสยกถา}}
\proseitem{605}{\csromansymbolmark{*}อนุสยา อพฺยากตาติ, อามนฺตา.\csromanspacer{} วิปากาพฺยากตา กิริยาพฺยากตา รูปํ นิพฺพานํ จกฺขายตนํ ฯเปฯ โผฏฺฐพฺพายตนนฺติ, น เหวํ วตฺตพฺเพ ฯเปฯ.}

\prose{\csromansymbolmark{*}กามราคานุสโย อพฺยากโตติ, อามนฺตา.\csromanspacer{} กามราโค กามราคปริยุฏฺฐานํ กามราคสํโยชนํ กาโมโฆ กามโยโค กามจฺฉนฺทนีวรณํ อพฺยากตนฺติ, น เหวํ วตฺตพฺเพ ฯเปฯ.\csromanspacer{} กามราโค กามราคปริยุฏฺฐานํ ฯเปฯ กามจฺฉนฺทนีวรณํ อกุสลนฺติ, อามนฺตา.\csromanspacer{} กามราคานุสโย อกุสโลติ, น เหวํ วตฺตพฺเพ ฯเปฯ.}

\prose{\csromansymbolmark{*}ปฏิฆานุสโย อพฺยากโตติ, อามนฺตา.\csromanspacer{} ปฏิฆํ ปฏิฆปริยุฏฺฐานํ ปฏิฆสํโยชนํ อพฺยากตนฺติ, น เหวํ วตฺตพฺเพ ฯเปฯ.\csromanspacer{} ปฏิฆํ ปฏิฆปริยุฏฺฐานํ ปฏิฆสํโยชนํ อกุสลนฺติ, อามนฺตา.\csromanspacer{} ปฏิฆานุสโย อกุสโลติ, น เหวํ วตฺตพฺเพ ฯเปฯ.}

\prose{\csromanfolio{324}\csromansymbolmark{*}มานานุสโย อพฺยากโตติ, อามนฺตา.\csromanspacer{} มาโน มานปริยุฏฺฐานํ มานสํโยชนํ อพฺยากตนฺติ, น เหวํ วตฺตพฺเพ ฯเปฯ.\csromanspacer{} มาโน มานปริยุฏฺฐานํ มานสํโยชนํ อกุสลนฺติ, อามนฺตา.\csromanspacer{} มานานุสโย อกุสโลติ, น เหวํ วตฺตพฺเพ ฯเปฯ.}

\prose{\csromansymbolmark{*}ทิฏฺฐานุสโย อพฺยากโตติ, อามนฺตา.\csromanspacer{} ทิฏฺฐิ ทิฏฺโฐโฆ ทิฏฺฐิโยโค ทิฏฺฐิปริยุฏฺฐานํ ทิฏฺฐิสํโยชนํ อพฺยากตนฺติ, น เหวํ วตฺตพฺเพ ฯเปฯ.\csromanspacer{} ทิฏฺฐิ ทิฏฺโฐโฆ ทิฏฺฐิโยโค ทิฏฺฐิปริยุฏฺฐานํ ทิฏฺฐิสํโยชนํ อกุสลนฺติ, อามนฺตา.\csromanspacer{} ทิฏฺฐานุสโย อกุสโลติ, น เหวํ วตฺตพฺเพ ฯเปฯ.}

\prose{\csromansymbolmark{*}วิจิกิจฺฉานุสโย อพฺยากโตติ, อามนฺตา.\csromanspacer{} วิจิกิจฺฉา วิจิกิจฺฉาปริยุฏฺฐานํ วิจิกิจฺฉาสํโยชนํ วิจิกิจฺฉานีวรณํ อพฺยากตนฺติ, น เหวํ วตฺตพฺเพ ฯเปฯ.\csromanspacer{} วิจิกิจฺฉา วิจิกิจฺฉาปริยุฏฺฐานํ วิจิกิจฺฉาสํโยชนํ วิจิกิจฺฉานีวรณํ อกุสลนฺติ, อามนฺตา.\csromanspacer{} วิจิกิจฺฉานุสโย อกุสโลติ, น เหวํ วตฺตพฺเพ ฯเปฯ.}

\prose{\csromansymbolmark{*}ภวราคานุสโย อพฺยากโตติ, อามนฺตา.\csromanspacer{} ภวราโค ภวราคปริยุฏฺฐานํ ภวราคสํโยชนํ อพฺยากตนฺติ, น เหวํ วตฺตพฺเพ ฯเปฯ.\csromanspacer{} ภวราโค ภวราคปริยุฏฺฐานํ ภวราคสํโยชนํ อกุสลนฺติ, อามนฺตา.\csromanspacer{} ภวราคานุสโย อกุสโลติ, น เหวํ วตฺตพฺเพ ฯเปฯ.}

\prose{\csromansymbolmark{*}อวิชฺชานุสโย อพฺยากโตติ, อามนฺตา.\csromanspacer{} อวิชฺชา อวิชฺโชโฆ อวิชฺชาโยโค อวิชฺชาปริยุฏฺฐานํ อวิชฺชาสํโยชนํ อวิชฺชานีวรณํ อพฺยากตนฺติ, น เหวํ วตฺตพฺเพ ฯเปฯ.\csromanspacer{} อวิชฺชา อวิชฺโชโฆ อวิชฺชาโยโค อวิชฺชาปริยุฏฺฐานํ อวิชฺชาสํโยชนํ อวิชฺชานีวรณํ อกุสลนฺติ, อามนฺตา.\csromanspacer{} อวิชฺชานุสโย อกุสโลติ, น เหวํ วตฺตพฺเพ ฯเปฯ.}

\proseitem{606}{\csromansymbolmark{+}น วตฺตพฺพํ “อนุสยา อพฺยากตา”ติ, อามนฺตา.\csromanspacer{} ปุถุชฺชโน กุสลาพฺยากเต จิตฺเต วตฺตมาเน “สานุสโย”ติ วตฺตพฺโพติ, อามนฺตา.\csromanspacer{} กุสลากุสลา ธมฺมา สมฺมุขีภาวํ อาคจฺฉนฺตีติ, น เหวํ วตฺตพฺเพ ฯเปฯ.\csromanspacer{} เตน หิ อนุสยา อพฺยากตาติ.\csromanspacer{} ปุถุชฺชโน กุสลาพฺยากเต จิตฺเต วตฺตมาเน “สราโค”ติ วตฺตพฺโพติ, อามนฺตา.\csromanspacer{} กุสลากุสลา ธมฺมา สมฺมุขีภาวํ อาคจฺฉนฺตีติ, น เหวํ วตฺตพฺเพ ฯเปฯ.\csromanspacer{} เตน หิ ราโค อพฺยากโตติ.}

\chapterheadpageread{11. เอกาทสมวคฺค}
\proseitem{607}{\csromanfolio{325}\csromansymbolmark{*}อนุสยา อเหตุกาติ, อามนฺตา.\csromanspacer{} รูปํ นิพฺพานํ จกฺขายตนํ ฯเปฯ โผฏฺฐพฺพายตนนฺติ, น เหวํ วตฺตพฺเพ ฯเปฯ.}

\prose{\csromansymbolmark{*}กามราคานุสโย อเหตุโกติ, อามนฺตา.\csromanspacer{} กามราโค กามราคปริยุฏฺฐานํ กามราคสํโยชนํ กามจฺฉนฺทนีวรณํ อเหตุกนฺติ, น เหวํ วตฺตพฺเพ ฯเปฯ.\csromanspacer{} กามราโค กามราคปริยุฏฺฐานํ ฯเปฯ.\csromanspacer{} กามจฺฉนฺทนีวรณํ สเหตุกนฺติ, อามนฺตา.\csromanspacer{} กามราคานุสโย สเหตุโกติ, น เหวํ วตฺตพฺเพ ฯเปฯ.\csromanspacer{} ปฏิฆานุสโย ฯเปฯ.\csromanspacer{} มานานุสโย.\csromanspacer{} ทิฏฺฐานุสโย.\csromanspacer{} วิจิกิจฺฉานุสโย.\csromanspacer{} ภวราคานุสโย.\csromanspacer{} อวิชฺชานุสโย.\csromanspacer{} อเหตุโกติ, อามนฺตา.\csromanspacer{} อวิชฺชา อวิชฺโชโฆ อวิชฺชาโยโค อวิชฺชาปริยุฏฺฐานํ อวิชฺชาสํโยชนํ อวิชฺชานีวรณํ อเหตุกนฺติ, น เหวํ วตฺตพฺเพ ฯเปฯ.\csromanspacer{} อวิชฺชา อวิชฺโชโฆ ฯเปฯ.\csromanspacer{} อวิชฺชานีวรณํ สเหตุกนฺติ, อามนฺตา.\csromanspacer{} อวิชฺชานุสโย สเหตุโกติ, น เหวํ วตฺตพฺเพ ฯเปฯ.}

\proseitem{608}{\csromansymbolmark{+}น วตฺตพฺพํ “อนุสยา อเหตุกา”ติ, อามนฺตา.\csromanspacer{} ปุถุชฺชโน กุสลาพฺยากเต จิตฺเต วตฺตมาเน “สานุสโย”ติ วตฺตพฺโพติ, อามนฺตา.\csromanspacer{} อนุสยา เตน เหตุนา สเหตุกาติ, น เหวํ วตฺตพฺเพ ฯเปฯ.\csromanspacer{} เตน หิ อนุสยา อเหตุกาติ.\csromanspacer{} ปุถุชฺชโน กุสลาพฺยากเต จิตฺเต วตฺตมาเน “สราโค”ติ วตฺตพฺโพติ, อามนฺตา.\csromanspacer{} ราโค เตน เหตุนา สเหตุโกติ, น เหวํ วตฺตพฺเพ ฯเปฯ.\csromanspacer{} เตน หิ ราโค อเหตุโกติ.}

\proseitem{609}{\csromansymbolmark{*}อนุสยา จิตฺตวิปฺปยุตฺตาติ, อามนฺตา.\csromanspacer{} รูปํ นิพฺพานํ จกฺขายตนํ ฯเปฯ โผฏฺฐพฺพายตนนฺติ, น เหวํ วตฺตพฺเพ ฯเปฯ.}

\prose{\csromansymbolmark{*}กามราคานุสโย จิตฺตวิปฺปยุตฺโตติ, อามนฺตา.\csromanspacer{} กามราโค กามราคปริยุฏฺฐานํ กามราคสํโยชนํ กาโมโฆ กามโยโค กามจฺฉนฺทนีวรณํ จิตฺตวิปฺปยุตฺตนฺติ, น เหวํ วตฺตพฺเพ ฯเปฯ.\csromanspacer{} กามราโค กามราคปริยุฏฺฐานํ ฯเปฯ กามจฺฉนฺทนีวรณํ จิตฺตสมฺปยุตฺตนฺติ, อามนฺตา.\csromanspacer{} กามราคานุสโย จิตฺตสมฺปยุตฺโตติ, น เหวํ วตฺตพฺเพ ฯเปฯ.}

\proseitem{610}{\csromansymbolmark{*}กามราคานุสโย จิตฺตวิปฺปยุตฺโตติ, อามนฺตา.\csromanspacer{} กตมกฺขนฺธปริยาปนฺโนติ, สงฺขารกฺขนฺธปริยาปนฺโนติ.\csromanspacer{} สงฺขารกฺขนฺโธ จิตฺตวิปฺปยุตฺโตติ, \csromanfolio{326}น เหวํ วตฺตพฺเพ.\csromanspacer{} สงฺขารกฺขนฺโธ จิตฺตวิปฺปยุตฺโตติ, อามนฺตา.\csromanspacer{} เวทนากฺขนฺโธ สญฺญากฺขนฺโธ จิตฺตวิปฺปยุตฺโตติ, น เหวํ วตฺตพฺเพ ฯเปฯ.}

\prose{\csromansymbolmark{*}กามราคานุสโย สงฺขารกฺขนฺธปริยาปนฺโน จิตฺตวิปฺปยุตฺโตติ, อามนฺตา.\csromanspacer{} กามราโค สงฺขารกฺขนฺธปริยาปนฺโน จิตฺตวิปฺปยุตฺโตติ, น เหวํ วตฺตพฺเพ ฯเปฯ.\csromanspacer{} กามราโค สงฺขารกฺขนฺธปริยาปนฺโน จิตฺตสมฺปยุตฺโตติ, อามนฺตา.\csromanspacer{} กามราคานุสโย สงฺขารกฺขนฺธปริยาปนฺโน จิตฺตสมฺปยุตฺโตติ, น เหวํ วตฺตพฺเพ ฯเปฯ.}

\prose{\csromansymbolmark{*}กามราคานุสโย สงฺขารกฺขนฺธปริยาปนฺโน จิตฺตวิปฺปยุตฺโต, กามราโค สงฺขารกฺขนฺธปริยาปนฺโน จิตฺตสมฺปยุตฺโตติ, อามนฺตา.\csromanspacer{} สงฺขารกฺขนฺโธ เอกเทโส จิตฺตสมฺปยุตฺโต เอกเทโส จิตฺตวิปฺปยุตฺโตติ, น เหวํ วตฺตพฺเพ ฯเปฯ. \csromansymbolmark{*}สงฺขารกฺขนฺโธ เอกเทโส จิตฺตสมฺปยุตฺโต เอกเทโส จิตฺตวิปฺปยุตฺโตติ, อามนฺตา.\csromanspacer{} เวทนากฺขนฺโธ สญฺญากฺขนฺโธ เอกเทโส จิตฺตสมฺปยุตฺโต เอกเทโส จิตฺตวิปฺปยุตฺโตติ, น เหวํ วตฺตพฺเพ ฯเปฯ.}

\proseitem{611}{\csromansymbolmark{*}ปฏิฆานุสโย มานานุสโย ทิฏฺฐานุสโย วิจิกิจฺฉานุสโย ภวราคานุสโย อวิชฺชานุสโย จิตฺตวิปฺปยุตฺโตติ, อามนฺตา.\csromanspacer{} อวิชฺชา อวิชฺโชโฆ อวิชฺชาโยโค อวิชฺชานีวรณํ จิตฺตวิปฺปยุตฺตนฺติ, น เหวํ วตฺตพฺเพ ฯเปฯ.\csromanspacer{} อวิชฺชา อวิชฺโชโฆ อวิชฺชาโยโค อวิชฺชานีวรณํ จิตฺตสมฺปยุตฺตนฺติ, อามนฺตา.\csromanspacer{} อวิชฺชานุสโย จิตฺตสมฺปยุตฺโตติ, น เหวํ วตฺตพฺเพ ฯเปฯ.}

\proseitem{612}{\csromansymbolmark{*}อวิชฺชานุสโย จิตฺตวิปฺปยุตฺโตติ, อามนฺตา.\csromanspacer{} กตมกฺขนฺธปริยาปนฺโนติ, สงฺขารกฺขนฺธปริยาปนฺโนติ.\csromanspacer{} สงฺขารกฺขนฺโธ จิตฺตวิปฺปยุตฺโตติ, น เหวํ วตฺตพฺเพ.\csromanspacer{} สงฺขารกฺขนฺโธ จิตฺตวิปฺปยุตฺโตติ, อามนฺตา.\csromanspacer{} เวทนากฺขนฺโธ สญฺญากฺขนฺโธ จิตฺตวิปฺปยุตฺโตติ, น เหวํ วตฺตพฺเพ ฯเปฯ.}

\prose{\csromansymbolmark{*}อวิชฺชานุสโย สงฺขารกฺขนฺธปริยาปนฺโน จิตฺตวิปฺปยุตฺโตติ, อามนฺตา.\csromanspacer{} อวิชฺชา สงฺขารกฺขนฺธปริยาปนฺนา จิตฺตวิปฺปยุตฺตาติ, น เหวํ วตฺตพฺเพ ฯเปฯ.\csromanspacer{} อวิชฺชา สงฺขารกฺขนฺธปริยาปนฺนา จิตฺตสมฺปยุตฺตาติ, อามนฺตา.\csromanspacer{} อวิชฺชานุสโย สงฺขารกฺขนฺธปริยาปนฺโน จิตฺตสมฺปยุตฺโตติ, น เหวํ วตฺตพฺเพ ฯเปฯ.}

\chapterheadpageread{11. เอกาทสมวคฺค}
\csromanmatikaanchor{mtk.136}
\csromantocmarkref{subsubsection}{(109) 4. ญาณกถา}{mtk.136}
\bookmark[dest={mtk.136},level=3]{(109) 4. ญาณกถา}
\prose{\csromanfolio{327}\csromansymbolmark{*}อวิชฺชานุสโย สงฺขารกฺขนฺธปริยาปนฺโน จิตฺตวิปฺปยุตฺโต, อวิชฺชาสงฺขารกฺขนฺธปริยาปนฺนา จิตฺตสมฺปยุตฺตาติ, อามนฺตา.\csromanspacer{} สงฺขารกฺขนฺโธ เอกเทโส จิตฺตสมฺปยุตฺโต เอกเทโส จิตฺตวิปฺปยุตฺโตติ, น เหวํ วตฺตพฺเพ.}

\prose{\csromansymbolmark{*}สงฺขารกฺขนฺโธ เอกเทโส จิตฺตสมฺปยุตฺโต เอกเทโส จิตฺตวิปฺปยุตฺโตติ, อามนฺตา.\csromanspacer{} เวทนากฺขนฺโธ สญฺญากฺขนฺโธ เอกเทโส จิตฺตสมฺปยุตฺโต เอกเทโส จิตฺตวิปฺปยุตฺโตติ, น เหวํ วตฺตพฺเพ ฯเปฯ.}

\proseitem{613}{\csromansymbolmark{+}น วตฺตพฺพํ “อนุสยา จิตฺตวิปฺปยุตฺตา”ติ, อามนฺตา.\csromanspacer{} ปุถุชฺชโน กุสลาพฺยากเต จิตฺเต วตฺตมาเน “สานุสโย”ติ วตฺตพฺโพติ, อามนฺตา.\csromanspacer{} อนุสยา เตน จิตฺเตน สมฺปยุตฺตาติ, น เหวํ วตฺตพฺเพ.\csromanspacer{} เตน หิ อนุสยา จิตฺตวิปฺปยุตฺตาติ.\csromanspacer{} ปุถุชฺชโน กุสลาพฺยากเต จิตฺเต วตฺตมาเน “สราโค”ติ วตฺตพฺโพติ, อามนฺตา.\csromanspacer{} ราโค เตน จิตฺเตน สมฺปยุตฺโตติ, น เหวํ วตฺตพฺเพ.\csromanspacer{} เตน หิ ราโค จิตฺตวิปฺปยุตฺโตติ.}

\nitthitam{ติสฺโสปิ อนุสยกถา นิฏฺฐิตา.}
\chapterhead{11. เอกาทสมวคฺค}
\vspace{\tipitakaparskip}
\csromancenter{\textbf{(109) 4.}\csromanspacer{}\textbf{ ญาณกถา}}
\proseitem{614}{\csromansymbolmark{*}อญฺญาเณ วิคเต ญาณวิปฺปยุตฺเต จิตฺเต วตฺตมาเน น วตฺตพฺพํ “ญาณี”ติ, อามนฺตา.\csromanspacer{} ราเค วิคเต น วตฺตพฺพํ “วีตราโค”ติ, น เหวํ วตฺตพฺเพ ฯเปฯ.\csromanspacer{} อญฺญาเณ วิคเต ญาณวิปฺปยุตฺเต จิตฺเต วตฺตมาเน น วตฺตพฺพํ “ญาณี”ติ, อามนฺตา.\csromanspacer{} โทเส วิคเต โมเห วิคเต กิเลเส วิคเต น วตฺตพฺพํ “นิกฺกิเลโส”ติ, น เหวํ วตฺตพฺเพ ฯเปฯ.}

\prose{\csromansymbolmark{*}ราเค วิคเต วตฺตพฺพํ “วีตราโค”ติ, อามนฺตา.\csromanspacer{} อญฺญาเณ วิคเต ญาณวิปฺปยุตฺเต จิตฺเต วตฺตมาเน วตฺตพฺพํ “ญาณี”ติ, น เหวํ วตฺตพฺเพ ฯเปฯ.\csromanspacer{} โทเส วิคเต โมเห วิคเต กิเลเส วิคเต วตฺตพฺพํ “นิกฺกิเลโส”ติ, อามนฺตา.\csromanspacer{} อญฺญาเณ วิคเต ญาณวิปฺปยุตฺเต จิตฺเต วตฺตมาเน วตฺตพฺพํ “ญาณี”ติ, น เหวํ วตฺตพฺเพ ฯเปฯ.}

\csromanmatikaanchor{mtk.137}
\csromantocmarkref{subsubsection}{(110) 5. ญาณํจิตฺตวิปฺปยุตฺตนฺติกถา}{mtk.137}
\bookmark[dest={mtk.137},level=3]{(110) 5. ญาณํจิตฺตวิปฺปยุตฺตนฺติกถา}
\proseitem{615}{\csromanfolio{328}\csromansymbolmark{+}อญฺญาเณ วิคเต ญาณวิปฺปยุตฺเต จิตฺเต วตฺตมาเน วตฺตพฺพํ “ญาณี”ติ, อามนฺตา.\csromanspacer{} อตีเตน ญาเณน ญาณี นิรุทฺเธน วิคเตน ปฏิปสฺสทฺเธน ญาเณน ญาณีติ, น เหวํ วตฺตพฺเพ ฯเปฯ.}

\nitthitam{ญาณกถา นิฏฺฐิตา.}
\chapterhead{11. เอกาทสมวคฺค}
\vspace{\tipitakaparskip}
\csromancenter{\textbf{(110) 5.}\csromanspacer{}\textbf{ ญาณํจิตฺตวิปฺปยุตฺตนฺติกถา}}
\proseitem{616}{\csromansymbolmark{*}ญาณํ จิตฺตวิปฺปยุตฺตนฺติ, อามนฺตา.\csromanspacer{} รูปํ นิพฺพานํ จกฺขายตนํ ฯเปฯ โผฏฺฐพฺพายตนนฺติ, น เหวํ วตฺตพฺเพ ฯเปฯ.\csromanspacer{} ญาณํ จิตฺตวิปฺปยุตฺตนฺติ, อามนฺตา.\csromanspacer{} ปญฺญา ปญฺญินฺทฺริยํ ปญฺญาพลํ สมฺมาทิฏฺฐิ ธมฺมวิจยสมฺโพชฺฌงฺโค จิตฺตวิปฺปยุตฺโตติ, น เหวํ วตฺตพฺเพ ฯเปฯ.\csromanspacer{} ปญฺญา ปญฺญินฺทฺริยํ ปญฺญาพลํ สมฺมาทิฏฺฐิ ธมฺมวิจยสมฺโพชฺฌงฺโค จิตฺตสมฺปยุตฺโตติ, อามนฺตา.\csromanspacer{} ญาณํ จิตฺตสมฺปยุตฺตนฺติ, น เหวํ วตฺตพฺเพ ฯเปฯ.}

\prose{\csromansymbolmark{*}ญาณํ จิตฺตวิปฺปยุตฺตนฺติ, อามนฺตา.\csromanspacer{} กตมกฺขนฺธปริยาปนฺนนฺติ, สงฺขารกฺขนฺธปริยาปนฺนนฺติ.\csromanspacer{} สงฺขารกฺขนฺโธ จิตฺตวิปฺปยุตฺโตติ, น เหวํ วตฺตพฺเพ ฯเปฯ.\csromanspacer{} สงฺขารกฺขนฺโธ จิตฺตวิปฺปยุตฺโตติ, อามนฺตา.\csromanspacer{} เวทนากฺขนฺโธ สญฺญาขนฺโธ จิตฺตวิปฺปยุตฺโตติ, น เหวํ วตฺตพฺเพ ฯเปฯ.\csromanspacer{} ญาณํ สงฺขารกฺขนฺธปริยาปนฺนํ จิตฺตวิปฺปยุตฺตนฺติ, อามนฺตา.\csromanspacer{} ปญฺญา สงฺขารกฺขนฺธปริยาปนฺนา จิตฺตวิปฺปยุตฺตาติ, น เหวํ วตฺตพฺเพ ฯเปฯ.\csromanspacer{} ปญฺญา สงฺขารกฺขนฺธปริยาปนฺนา จิตฺตสมฺปยุตฺตาติ, อามนฺตา.\csromanspacer{} ญาณํ สงฺขารกฺขนฺธปริยาปนฺนํ จิตฺตสมฺปยุตฺตนฺติ, น เหวํ วตฺตพฺเพ ฯเปฯ.\csromanspacer{} ญาณํ สงฺขารกฺขนฺธปริยาปนฺนํ จิตฺตวิปฺปยุตฺตํ ปญฺญา สงฺขารกฺขนฺธปริยาปนฺนา จิตฺตสมฺปยุตฺตาติ, อามนฺตา.\csromanspacer{} สงฺขารกฺขนฺโธ เอกเทโส จิตฺตสมฺปยุตฺโต เอกเทโส จิตฺตวิปฺปยุตฺโตติ, น เหวํ วตฺตพฺเพ ฯเปฯ.\csromanspacer{} สงฺขารกฺขนฺโธ เอกเทโส จิตฺตสมฺปยุตฺโต เอกเทโส จิตฺตวิปฺปยุตฺโตติ, อามนฺตา.\csromanspacer{} เวทนากฺขนฺโธ สญฺญากฺขนฺโธ เอกเทโส จิตฺตสมฺปยุตฺโต เอกเทโส จิตฺตวิปฺปยุตฺโตติ, น เหวํ วตฺตพฺเพ ฯเปฯ.}

\proseitem{617}{\csromansymbolmark{+}น วตฺตพฺพํ “ญาณํ จิตฺตวิปฺปยุตฺตนฺ”ติ, อามนฺตา.\csromanspacer{} อรหา จกฺขุวิญฺญาณสมงฺคี “ญาณี”ติ วตฺตพฺโพติ, อามนฺตา.\csromanspacer{} ญาณํ เตน จิตฺเตน สมฺปยุตฺตนฺติ, น เหวํ วตฺตพฺเพ.\csromanspacer{} เตน หิ ญาณํ จิตฺตวิปฺปยุตฺตนฺติ.}

\chapterheadpageread{11. เอกาทสมวคฺค}
\csromanmatikaanchor{mtk.138}
\csromantocmarkref{subsubsection}{(111) 6. อิทํทุกฺขนฺติกถา}{mtk.138}
\bookmark[dest={mtk.138},level=3]{(111) 6. อิทํทุกฺขนฺติกถา}
\prose{\csromanfolio{329}อรหา จกฺขุวิญฺญาณสมงฺคี “ปญฺญวา”ติ วตฺตพฺโพติ\footnote{สกวาทีปุจฺฉา วิย ทิสฺสติ.}, อามนฺตา.\csromanspacer{} ปญฺญา เตน จิตฺเตน สมฺปยุตฺตาติ, น เหวํ วตฺตพฺเพ.\csromanspacer{} เตน หิ ปญฺญา จิตฺตวิปฺปยุตฺตาติ.}

\nitthitam{ญาณํ จิตฺตวิปฺปยุตฺตนฺติกถา นิฏฺฐิตา.}
\chapterhead{11. เอกาทสมวคฺค}
\vspace{\tipitakaparskip}
\csromancenter{\textbf{(111) 6.}\csromanspacer{}\textbf{ อิทํทุกฺขนฺติกถา}}
\proseitem{618}{\csromansymbolmark{*}“อิทํ ทุกฺขนฺ”ติ วาจํ ภาสโต “อิทํ ทุกฺขนฺ”ติ ญาณํ ปวตฺตตีติ, อามนฺตา. “อยํ ทุกฺขสมุทโย”ติ วาจํ ภาสโต “อยํ ทุกฺขสมุทโย”ติ ญาณํ ปวตฺตตีติ, น เหวํ วตฺตพฺเพ ฯเปฯ “อิทํ ทุกฺขนฺ”ติ วาจํ ภาสโต “อิทํ ทุกฺขนฺ”ติ ญาณํ ปวตฺตตีติ, อามนฺตา. “อยํ ทุกฺขนิโรโธ”ติ วาจํ ภาสโต “อยํ ทุกฺขนิโรโธ”ติ ญาณํ ปวตฺตตีติ, น เหวํ วตฺตพฺเพ ฯเปฯ “อิทํ ทุกฺขนฺ”ติ วาจํ ภาสโต “อิทํ ทุกฺขนฺ”ติ ญาณํ ปวตฺตตีติ, อามนฺตา. “อยํ มคฺโค”ติ วาจํ ภาสโต “อยํ มคฺโค”ติ ญาณํ ปวตฺตตีติ, น เหวํ วตฺตพฺเพ ฯเปฯ.}

\prose{\csromansymbolmark{*}“อยํ สมุทโย”ติ วาจํ ภาสโต น จ “อยํ สมุทโย”ติ ญาณํ ปวตฺตตีติ, อามนฺตา. “อิทํ ทุกฺขนฺ”ติ วาจํ ภาสโต น จ “อิทํ ทุกฺขนฺ”ติ ญาณํ ปวตฺตตีติ, น เหวํ วตฺตพฺเพ ฯเปฯ “อยํ นิโรโธ”ติ. “อยํ มคฺโค”ติ วาจํ ภาสโต น จ “อยํ มคฺโค”ติ ญาณํ ปวตฺตตีติ, อามนฺตา. “อิทํ ทุกฺขนฺ”ติ วาจํ ภาสโต น จ “อิทํ ทุกฺขนฺ”ติ ญาณํ ปวตฺตตีติ, น เหวํ วตฺตพฺเพ ฯเปฯ.}

\proseitem{619}{\csromansymbolmark{*}“อิทํ ทุกฺขนฺ”ติ วาจํ ภาสโต “อิทํ ทุกฺขนฺ”ติ ญาณํ ปวตฺตตีติ, อามนฺตา. “รูปํ อนิจฺจนฺ”ติ วาจํ ภาสโต “รูปํ อนิจฺจนฺ”ติ ญาณํ ปวตฺตตีติ, น เหวํ วตฺตพฺเพ ฯเปฯ “อิทํ ทุกฺขนฺ”ติ วาจํ ภาสโต “อิทํ ทุกฺขนฺ”ติ ญาณํ ปวตฺตตีติ, อามนฺตา. “เวทนา.\csromanspacer{} สญฺญา.\csromanspacer{} สงฺขารา.\csromanspacer{} วิญฺญาณํ อนิจฺจนฺ”ติ วาจํ ภาสโต “วิญฺญาณํ อนิจฺจนฺ”ติ ญาณํ ปวตฺตตีติ, น เหวํ วตฺตพฺเพ ฯเปฯ.}

\csromanmatikaanchor{mtk.139}
\csromantocmarkref{subsubsection}{(112) 7. อิทฺธิพลกถา}{mtk.139}
\bookmark[dest={mtk.139},level=3]{(112) 7. อิทฺธิพลกถา}
\prose{\csromanfolio{330}\csromansymbolmark{*}“อิทํ ทุกฺขนฺ”ติ วาจํ ภาสโต “อิทํ ทุกฺขนฺ”ติ ญาณํ ปวตฺตตีติ, อามนฺตา. “รูปํ อนตฺตา”ติ วาจํ ภาสโต “รูปํ อนตฺตา”ติ ญาณํ ปวตฺตตีติ, น เหวํ วตฺตพฺเพ ฯเปฯ “อิทํ ทุกฺขนฺ”ติ วาจํ ภาสโต “อิทํ ทุกฺขนฺ”ติ ญาณํ ปวตฺตตีติ, อามนฺตา. “เวทนา.\csromanspacer{} สญฺญา.\csromanspacer{} สงฺขารา.\csromanspacer{} วิญฺญาณํ อนตฺตา”ติ วาจํ ภาสโต “วิญฺญาณํ อนตฺตา”ติ ญาณํ ปวตฺตตีติ, น เหวํ วตฺตพฺเพ ฯเปฯ.}

\prose{\csromansymbolmark{*}“รูปํ อนิจฺจนฺ”ติ วาจํ ภาสโต น จ “รูปํ อนิจฺจนฺ”ติ ญาณํ ปวตฺตตีติ, อามนฺตา. “อิทํ ทุกฺขนฺ”ติ วาจํ ภาสโต “น จ อิทํ ทุกฺขนฺ”ติ ญาณํ ปวตฺตตีติ, น เหวํ วตฺตพฺเพ ฯเปฯ “เวทนา.\csromanspacer{} สญฺญา.\csromanspacer{} สงฺขารา.\csromanspacer{} วิญฺญาณํ อนิจฺจนฺ”ติ วาจํ ภาสโต น จ “วิญฺญาณํ อนิจฺจนฺ”ติ ญาณํ ปวตฺตตีติ, อามนฺตา. “อิทํ ทุกฺขนฺ”ติ วาจํ ภาสโต น จ “อิทํ ทุกฺขนฺ”ติ ญาณํ ปวตฺตตีติ, น เหวํ วตฺตพฺเพ ฯเปฯ.}

\prose{\csromansymbolmark{*}“รูปํ อนตฺตา”ติ วาจํ ภาสโต น จ “รูปํ อนตฺตา”ติ ญาณํ ปวตฺตตีติ, อามนฺตา. “อิทํ ทุกฺขนฺ”ติ วาจํ ภาสโต น จ “อิทํ ทุกฺขนฺ”ติ ญาณํ ปวตฺตตีติ, น เหวํ วตฺตพฺเพ ฯเปฯ “เวทนา.\csromanspacer{} สญฺญา.\csromanspacer{} สงฺขารา.\csromanspacer{} วิญฺญาณํ อนตฺตา”ติ วาจํ ภาสโต น จ “วิญฺญาณํ อนตฺตา”ติ ญาณํ ปวตฺตตีติ, อามนฺตา. “อิทํ ทุกฺขนฺ”ติ วาจํ ภาสโต น จ “ทุกฺขนฺ”ติ ญาณํ ปวตฺตตีติ, น เหวํ วตฺตพฺเพ ฯเปฯ.}

\proseitem{620}{\csromansymbolmark{*}“อิทํ ทุกฺขนฺ”ติ วาจํ ภาสโต “อิทํ ทุกฺขนฺ”ติ ญาณํ ปวตฺตตีติ, อามนฺตา. “อิ”ติ\footnote{¢ติ (สฺยา, อิ)} จ “ทนฺ”ติ จ “ทุ”ติ\footnote{ทูติ (สฺยา, อิ)} จ “ขนฺ”ติ จ ญาณํ ปวตฺตตีติ, น เหวํ วตฺตพฺเพ ฯเปฯ.}

\nitthitam{อิทํ ทุกฺขนฺติกถา นิฏฺฐิตา.}
\chapterhead{11. เอกาทสมวคฺค}
\vspace{\tipitakaparskip}
\csromancenter{\textbf{(112) 7.}\csromanspacer{}\textbf{ อิทฺธิพลกถา}}
\proseitem{621}{อิทฺธิพเลน สมนฺนาคโต กปฺปํ ติฏฺเฐยฺยาติ, อามนฺตา.\csromanspacer{} อิทฺธิมยิโก โส อายุ, อิทฺธิมยิกา สา คติ, อิทฺธิมยิโก โส อตฺตภาวปฺปฏิลาโภติ, น เหวํ วตฺตพฺเพ ฯเปฯ.}

\chapterheadpageread{11. เอกาทสมวคฺค}
\prose{\csromanfolio{331}\csromansymbolmark{*}อิทฺธิพเลน สมนฺนาคโต กปฺปํ ติฏฺเฐยฺยาติ, อามนฺตา.\csromanspacer{} อตีตํ กปฺปํ ติฏฺเฐยฺย, อนาคตํ กปฺปํ ติฏฺเฐยฺยาติ, น เหวํ วตฺตพฺเพ ฯเปฯ.\csromanspacer{} อิทฺธิพเลน สมนฺนาคโต กปฺปํ ติฏฺเฐยฺยาติ, อามนฺตา.\csromanspacer{} เทฺว กปฺเป ติฏฺเฐยฺย, ตโย กปฺเป ติฏฺเฐยฺย, จตฺตาโร กปฺเป ติฏฺเฐยฺยาติ, น เหวํ วตฺตพฺเพ ฯเปฯ.\csromanspacer{} อิทฺธิพเลน สมนฺนาคโต กปฺปํ ติฏฺเฐยฺยาติ, อามนฺตา.\csromanspacer{} สติ ชีวิเต ชีวิตาวเสเส ติฏฺเฐยฺย, อสติ ชีวิเต ชีวิตาวเสเส ติฏฺเฐยฺยาติ? สติ ชีวิเต ชีวิตาวเสเส ติฏฺเฐยฺยาติ.\csromanspacer{} หญฺจิ สติ ชีวิเต ชีวิตาวเสเส ติฏฺเฐยฺย, โน จ วต เร วตฺตพฺเพ “อิทฺธิพเลน สมนฺนาคโต กปฺปํ ติฏฺเฐยฺยา”ติ.\csromanspacer{} อสติ ชีวิเต ชีวิตาวเสเส ติฏฺเฐยฺยาติ.\csromanspacer{} มโต ติฏฺเฐยฺย กาลกโต ติฏฺเฐยฺยาติ, น เหวํ วตฺตพฺเพ ฯเปฯ.}

\proseitem{622}{\csromansymbolmark{*}อิทฺธิพเลน สมนฺนาคโต กปฺปํ ติฏฺเฐยฺยาติ, อามนฺตา. “อุปฺปนฺโน ผสฺโส มา นิรุชฺฌี”ติ ลพฺภา อิทฺธิยา ปคฺคเหตุนฺติ, น เหวํ วตฺตพฺเพ ฯเปฯ “อุปฺปนฺนา เวทนา ฯเปฯ.\csromanspacer{} อุปฺปนฺนา สญฺญา ฯเปฯ.\csromanspacer{} อุปฺปนฺนา เจตนา ฯเปฯ.\csromanspacer{} อุปฺปนฺนํ จิตฺตํ.\csromanspacer{} อุปฺปนฺนา สทฺธา ฯเปฯ.\csromanspacer{} อุปฺปนฺนํ วีริยํ.\csromanspacer{} อุปฺปนฺนา สติ.\csromanspacer{} อุปฺปนฺโน สมาธิ ฯเปฯ.\csromanspacer{} อุปฺปนฺนา ปญฺญา มา นิรุชฺฌี”ติ ลพฺภา อิทฺธิยา ปคฺคเหตุนฺติ, น เหวํ วตฺตพฺเพ ฯเปฯ.}

\prose{\csromansymbolmark{*}อิทฺธิพเลน สมนฺนาคโต กปฺปํ ติฏฺเฐยฺยาติ, อามนฺตา. “รูปํ นิจฺจํ โหตู”ติ ลพฺภา อิทฺธิยา ปคฺคเหตุนฺติ, น เหวํ วตฺตพฺเพ ฯเปฯ. “เวทนา ฯเปฯ.\csromanspacer{} สญฺญา.\csromanspacer{} สงฺขารา วิญฺญาณํ นิจฺจํ โหตู”ติ ลพฺภา อิทฺธิยา ปคฺคเหตุนฺติ, น เหวํ วตฺตพฺเพ ฯเปฯ.}

\prose{\csromansymbolmark{*}อิทฺธิพเลน สมนฺนาคโต กปฺปํ ติฏฺเฐยฺยาติ, อามนฺตา. “ชาติธมฺมา สตฺตา มา ชายิํสู”ติ ลพฺภา อิทฺธิยา ปคฺคเหตุนฺติ, น เหวํ วตฺตพฺเพ ฯเปฯ “ชราธมฺมา สตฺตา มา ชีริํสู”ติ ฯเปฯ “พฺยาธิธมฺมา สตฺตา มา พฺยาธิยิํสู”ติ ฯเปฯ “มรณธมฺมา สตฺตา มา มียิํสู”ติ ลพฺภา อิทฺธิยา ปคฺคเหตุนฺติ, น เหวํ วตฺตพฺเพ ฯเปฯ.}

\csromanmatikaanchor{mtk.140}
\csromantocmarkref{subsubsection}{(113) 8. สมาธิกถา}{mtk.140}
\bookmark[dest={mtk.140},level=3]{(113) 8. สมาธิกถา}
\proseitem{623}{\csromansymbolmark{+}น วตฺตพฺพํ “อิทฺธิพเลน สมนฺนาคโต กปฺปํ ติฏฺเฐยฺยา”ติ, อามนฺตา.\csromanspacer{} นนุ วุตฺตํ ภควตา “ยสฺส กสฺสจิ อานนฺท จตฺตาโร อิทฺธิปาทา ภาวิตา พหุลีกตา ยานีกตา วตฺถุกตา อนุฏฺฐิตา ปริจิตา สุสมารทฺธา, โส อากงฺขมาโน กปฺปํ วา ติฏฺเฐยฺย กปฺปาวเสสํ วา”ติ\footnote{ที 2. 86; สํ 3. 227 ปิฏฺเฐ; ขุ 1. 152 อุทาเนปิ.} \csromanfolio{332}อตฺเถว สุตฺตนฺโตติ, อามนฺตา.\csromanspacer{} เตน หิ อิทฺธิพเลน สมนฺนาคโต กปฺปํ ติฏฺเฐยฺยาติ.}

\proseitem{624}{\csromansymbolmark{*}อิทฺธิพเลน สมนฺนาคโต กปฺปํ ติฏฺเฐยฺยาติ, อามนฺตา.\csromanspacer{} นนุ วุตฺตํ ภควตา “จตุนฺนํ ภิกฺขเว ธมฺมานํ นตฺถิ โกจิ ปาฏิโภโค สมโณ วา พฺราหฺมโณ วา เทโว วา มาโร วา พฺรหฺมา วา โกจิ วา โลกสฺมิํ.\csromanspacer{} กตเมสํ จตุนฺนํ, ‘ชราธมฺโม มา ชีรี’ติ นตฺถิ โกจิ ปาฏิโภโค สมโณ วา พฺราหฺมโณ วา เทโว วา มาโร วา พฺรหฺมา วา โกจิ วา โลกสฺมิํ. ‘พฺยาธิธมฺโม มา พฺยาธิยี’ติ ฯเปฯ ‘มรณธมฺโม มา มียี’ติ ฯเปฯ.\csromanspacer{} ยานิ โข ปน ตานิ ปุพฺเพ กตานิ ปาปกานิ กมฺมานิ สํกิเสสิกานิ โปโนพฺภวิกานิ\footnote{โปโนภวิกานิ (สี, อิ)} สทรานิ\footnote{นิสฺสารานิ (สี, อิ, ก), ทุกฺขุทฺรยานิ (สฺยา)} ทุกฺขวิปากานิ อายติํ ชาติชรามรณิยานิ, ‘เตสํ วิปาโก มา นิพฺพตฺตี’ติ นตฺถิ โกจิ ปาฏิโภโค สมโณ วา พฺราหฺมโณ วา เทโว วา มาโร วา พฺรหฺมา วา โกจิ วา โลกสฺมิํ.\csromanspacer{} อิเมสํ โข ภิกฺขเว จตุนฺนํ ธมฺมานํ นตฺถิ โกจิ ปาฏิโภโค สมโณ วา พฺราหฺมโณ วา เทโว วา มาโร วา พฺรหฺมา วา โกจิ วา โลกสฺมินฺ”ติ\footnote{อํ 1. 491 ปิฏฺเฐ.} อตฺเถว สุตฺตนฺโตติ, อามนฺตา.\csromanspacer{} เตน หิ น วตฺตพฺพํ “อิทฺธิพเลน สมนฺนาคโต กปฺปํ ติฏฺเฐยฺยา”ติ.}

\nitthitam{อิทฺธิพลกถา นิฏฺฐิตา.}
\chapterhead{11. เอกาทสมวคฺค}
\vspace{\tipitakaparskip}
\csromancenter{\textbf{(113) 8.}\csromanspacer{}\textbf{ สมาธิกถา}}
\proseitem{625}{\csromansymbolmark{*}จิตฺตสนฺตติ สมาธีติ, อามนฺตา.\csromanspacer{} อตีตา จิตฺตสนฺตติ สมาธีติ, น เหวํ วตฺตพฺเพ ฯเปฯ.\csromanspacer{} จิตฺตสนฺตติ สมาธีติ, อามนฺตา.\csromanspacer{} อนาคตา จิตฺตสนฺตติ สมาธีติ, น เหวํ วตฺตพฺเพ ฯเปฯ.\csromanspacer{} จิตฺตสนฺตติ สมาธีติ, อามนฺตา.\csromanspacer{} นนุ อตีตํ นิรุทฺธํ อนาคตํ อชาตนฺติ, อามนฺตา.\csromanspacer{} หญฺจิ อตีตํ นิรุทฺธํ อนาคตํ อชาตํ, โน จ วต เร วตฺตพฺเพ “จิตฺตสนฺตติ สมาธี”ติ.}

\chapterheadpageread{11. เอกาทสมวคฺค}
\csromanmatikaanchor{mtk.141}
\csromantocmarkref{subsubsection}{(114) 9. ธมฺมฏฺฐิตตากถา}{mtk.141}
\bookmark[dest={mtk.141},level=3]{(114) 9. ธมฺมฏฺฐิตตากถา}
\proseitem{626}{\csromanfolio{333}\csromansymbolmark{+}เอกจิตฺตกฺขณิโก สมาธีติ, อามนฺตา.\csromanspacer{} จกฺขุวิญฺญาณสมงฺคี สมาปนฺโนติ, น เหวํ วตฺตพฺเพ ฯเปฯ.\csromanspacer{} โสตวิญฺญาณสมงฺคี ฯเปฯ.\csromanspacer{} ฆานวิญฺญาณสมงฺคี.\csromanspacer{} ชิวฺหาวิญฺญาณสมงฺคี.\csromanspacer{} กายวิญฺญาณสมงฺคี ฯเปฯ.\csromanspacer{} อกุสลจิตฺตสมงฺคี ฯเปฯ.\csromanspacer{} ราคสหคตจิตฺตสมงฺคี ฯเปฯ.\csromanspacer{} โทสสหคตจิตฺตสมงฺคี ฯเปฯ.\csromanspacer{} โมหสหคตจิตฺตสมงฺคี ฯเปฯ.\csromanspacer{} อโนตฺตปฺปสหคตจิตฺตสมงฺคี สมาปนฺโนติ, น เหวํ วตฺตพฺเพ ฯเปฯ.}

\prose{\csromansymbolmark{*}จิตฺตสนฺตติ สมาธีติ, อามนฺตา.\csromanspacer{} อกุสลจิตฺตสนฺตติ สมาธีติ, น เหวํ วตฺตพฺเพ ฯเปฯ.\csromanspacer{} ราคสหคตา ฯเปฯ.\csromanspacer{} โทสสหคตา ฯเปฯ.\csromanspacer{} โมหสหคตา ฯเปฯ.\csromanspacer{} อโนตฺตปฺปสหคตา จิตฺตสนฺตติ สมาธีติ, น เหวํ วตฺตพฺเพ ฯเปฯ.}

\prose{\csromansymbolmark{+}น วตฺตพฺพํ “จิตฺตสนฺตติ สมาธี”ติ, อามนฺตา.\csromanspacer{} นนุ วุตฺตํ ภควตา “อหํ โข อาวุโส นิคณฺฐา ปโหมิ อนิญฺชมาโน กาเยน อภาสมาโน วาจํ สตฺต รตฺตินฺทิวานิ\footnote{รตฺติทิวานิ (ก)} เอกนฺตสุขํ ปฏิสํเวที วิหริตุนฺ”ติ\footnote{ม 1. 132 ปิฏฺเฐ.} อตฺเถว สุตฺตนฺโตติ, อามนฺตา.\csromanspacer{} เตน หิ จิตฺตสนฺตติ สมาธีติ.}

\nitthitam{สมาธิกถา นิฏฺฐิตา.}
\chapterhead{11. เอกาทสมวคฺค}
\vspace{\tipitakaparskip}
\csromancenter{\textbf{(114) 9.}\csromanspacer{}\textbf{ ธมฺมฏฺฐิตตากถา}}
\proseitem{627}{\csromansymbolmark{*}ธมฺมฏฺฐิตตา ปรินิปฺผนฺนาติ, อามนฺตา.\csromanspacer{} ตาย ฐิตตา ปรินิปฺผนฺนาติ, น เหวํ วตฺตพฺเพ ฯเปฯ.\csromanspacer{} ตาย ฐิตตา ปรินิปฺผนฺนาติ, อามนฺตา.\csromanspacer{} ตาย ตาเยว นตฺถิ ทุกฺขสฺสนฺตกิริยา, นตฺถิ วฏฺฏุปจฺเฉโท, นตฺถิ อนุปาทาปรินิพฺพานนฺติ, น เหวํ วตฺตพฺเพ ฯเปฯ.}

\prose{\csromansymbolmark{*}รูปสฺส ฐิตตา ปรินิปฺผนฺนาติ, อามนฺตา.\csromanspacer{} ตาย ฐิตตา ปรินิปฺผนฺนาติ, น เหวํ วตฺตพฺเพ ฯเปฯ.\csromanspacer{} ตาย ฐิตตา ปรินิปฺผนฺนาติ, อามนฺตา.\csromanspacer{} ตาย ตาเยว นตฺถิ ทุกฺขสฺสนฺตกิริยา, นตฺถิ วฏฺฏุปจฺเฉโท, นตฺถิ อนุปาทาปรินิพฺพานนฺติ, น เหวํ วตฺตพฺเพ ฯเปฯ.}

\csromanmatikaanchor{mtk.142}
\csromantocmarkref{subsubsection}{(115) 10. อนิจฺจตากถา}{mtk.142}
\bookmark[dest={mtk.142},level=3]{(115) 10. อนิจฺจตากถา}
\prose{\csromanfolio{334}\csromansymbolmark{*}เวทนาย ฐิตตา ฯเปฯ.\csromanspacer{} สญฺญาย ฐิตตา ฯเปฯ.\csromanspacer{} สงฺขารานํ ฐิตตา ฯเปฯ.\csromanspacer{} วิญฺญาณสฺส ฐิตตา ปรินิปฺผนฺนาติ, อามนฺตา.\csromanspacer{} ตาย ฐิตตา ปรินิปฺผนฺนาติ, น เหวํ วตฺตพฺเพ ฯเปฯ.\csromanspacer{} ตาย ฐิตตา ปรินิปฺผนฺนาติ, อามนฺตา.\csromanspacer{} ตาย ตาเยว นตฺถิ ทุกฺขสฺสนฺตกิริยา, นตฺถิ วฏฺฏุปจฺเฉโท, นตฺถิ อนุปาทาปรินิพฺพานนฺติ, น เหวํ วตฺตพฺเพ ฯเปฯ.}

\nitthitam{ธมฺมฏฺฐิตตากถา นิฏฺฐิตา.}
\chapterhead{11. เอกาทสมวคฺค}
\vspace{\tipitakaparskip}
\csromancenter{\textbf{(115) 10.}\csromanspacer{}\textbf{ อนิจฺจตากถา}}
\proseitem{628}{\csromansymbolmark{*}อนิจฺจตา ปรินิปฺผนฺนาติ, อามนฺตา.\csromanspacer{} ตาย อนิจฺจตาย อนิจฺจตา ปรินิปฺผนฺนาติ, น เหวํ วตฺตพฺเพ ฯเปฯ.\csromanspacer{} ตาย อนิจฺจตาย อนิจฺจตา ปรินิปฺผนฺนาติ, อามนฺตา.\csromanspacer{} ตาย ตาเยว นตฺถิ ทุกฺขสฺสนฺตกิริยา, นตฺถิ วฏฺฏุปจฺเฉโท, นตฺถิ อนุปาทาปรินิพฺพานนฺติ, น เหวํ วตฺตพฺเพ.}

\prose{\csromansymbolmark{*}ชรา ปรินิปฺผนฺนาติ, อามนฺตา.\csromanspacer{} ตาย ชราย ชรา ปรินิปฺผนฺนาติ, น เหวํ วตฺตพฺเพ ฯเปฯ.\csromanspacer{} ตาย ชราย ชรา ปรินิปฺผนฺนาติ, อามนฺตา.\csromanspacer{} ตาย ตาเยว นตฺถิ ทุกฺขสฺสนฺตกิริยา, นตฺถิ วฏฺฏุปจฺเฉโท, นตฺถิ อนุปาทาปรินิพฺพานนฺติ, น เหวํ วตฺตพฺเพ ฯเปฯ.}

\prose{\csromansymbolmark{*}มรณํ ปรินิปฺผนฺนนฺติ, อามนฺตา.\csromanspacer{} ตสฺส มรณสฺส มรณํ ปรินิปฺผนฺนนฺติ, น เหวํ วตฺตพฺเพ.\csromanspacer{} ตสฺส มรณสฺส มรณํ ปรินิปฺผนฺนนฺติ, อามนฺตา.\csromanspacer{} ตสฺส ตสฺเสว นตฺถิ ทุกฺขสฺสนฺตกิริยา, นตฺถิ วฏฺฏุปจฺเฉโท, นตฺถิ อนุปาทาปรินิพฺพานนฺติ, น เหวํ วตฺตพฺเพ ฯเปฯ.}

\proseitem{629}{\csromansymbolmark{*}รูปํ ปรินิปฺผนฺนํ รูปสฺส อนิจฺจตา อตฺถีติ, อามนฺตา.\csromanspacer{} อนิจฺจตา ปรินิปฺผนฺนา อนิจฺจตาย อนิจฺจตา อตฺถีติ, น เหวํ วตฺตพฺเพ ฯเปฯ.\csromanspacer{} รูปํ ปรินิปฺผนฺนํ รูปสฺส ชรา อตฺถีติ, อามนฺตา.\csromanspacer{} ชรา ปรินิปฺผนฺนา ชราย ชรา อตฺถีติ, น เหวํ วตฺตพฺเพ ฯเปฯ.}

\prose{\csromansymbolmark{*}รูปํ ปรินิปฺผนฺนํ รูปสฺส เภโท อตฺถิ, อนฺตรธานํ อตฺถีติ, อามนฺตา.\csromanspacer{} มรณํ ปรินิปฺผนฺนํ มรณสฺส เภโท อตฺถิ, อนฺตรธานํ อตฺถีติ, น เหวํ วตฺตพฺเพ ฯเปฯ.}

\chapterheadpageread{12. ทฺวาทสมวคฺค}
\csromanmatikaanchor{mtk.143}
\csromanmatikaanchor{mtk.144}
\csromantocmarknopage{chapter}{12. ทฺวาทสมวคฺค}{mtk.143}
\bookmark[dest={mtk.143},level=0]{12. ทฺวาทสมวคฺค}
\csromantocmarkref{subsection}{(116) 1. สํวโรกมฺมนฺติกถา}{mtk.144}
\bookmark[dest={mtk.144},level=2]{(116) 1. สํวโรกมฺมนฺติกถา}
\prose{\csromanfolio{335}\csromansymbolmark{*}เวทนา ฯเปฯ.\csromanspacer{} สญฺญา.\csromanspacer{} สงฺขารา ฯเปฯ.\csromanspacer{} วิญฺญาณํ ปรินิปฺผนฺนํ วิญฺญาณสฺส อนิจฺจตา อตฺถีติ, อามนฺตา.\csromanspacer{} อนิจฺจตา ปรินิปฺผนฺนา อนิจฺจตาย อนิจฺจตา อตฺถีติ, น เหวํ วตฺตพฺเพ ฯเปฯ.\csromanspacer{} วิญฺญาณํ ปรินิปฺผนฺนํ วิญฺญาณสฺส ชรา อตฺถีติ, อามนฺตา.\csromanspacer{} ชรา ปรินิปฺผนฺนา ชราย ชรา อตฺถีติ, น เหวํ วตฺตพฺเพ ฯเปฯ.}

\prose{\csromansymbolmark{*}วิญฺญาณํ ปรินิปฺผนฺนํ วิญฺญาณสฺส เภโท อตฺถิ อนฺตรธานํ อตฺถีติ, อามนฺตา.\csromanspacer{} มรณํ ปรินิปฺผนฺนํ มรณสฺส เภโท อตฺถิ, อนฺตรธานํ อตฺถีติ, น เหวํ วตฺตพฺเพ ฯเปฯ.\csromanspacer{} อนิจฺจตากถา นิฏฺฐิตา.\csromanspacer{} เอกาทสโม วคฺโค.}

\titlehead{\textbf{ตสฺสุทฺทานํ}}
\prose{อนุสยา อพฺยากตา, อเหตุกา, จิตฺตวิปฺปยุตฺตา, อญฺญาเณ วิคเต ญาณี, ญาณํ จิตฺตวิปฺปยุตฺตํ, ยตฺถ สทฺเท\footnote{ยตฺถ สทฺโท (สี), ยถาสทฺทํ (?)} ญาณํ ปวตฺตติ, อิทฺธิพเลน สมนฺนาคโต กปฺปํ ติฏฺเฐยฺย, จิตฺตสนฺตติ สมาธิ, ธมฺมฏฺฐิตตา, อนิจฺจตาติ.}

\chapterhead{12. ทฺวาทสมวคฺค}
\vspace{\tipitakaparskip}
\csromancenter{\textbf{(116) 1.}\csromanspacer{}\textbf{ สํวโรกมฺมนฺติกถา}}
\proseitem{630}{\csromansymbolmark{*}สํวโร กมฺมนฺติ, อามนฺตา.\csromanspacer{} จกฺขุนฺทฺริยสํวโร จกฺขุกมฺมนฺติ, น เหวํ วตฺตพฺเพ ฯเปฯ.\csromanspacer{} โสตินฺทฺริยสํวโร ฯเปฯ.\csromanspacer{} ฆานินฺทฺริยสํวโร ฯเปฯ.\csromanspacer{} ชิวฺหินฺทฺริยสํวโร ฯเปฯ.\csromanspacer{} กายินฺทฺริยสํวโร กายกมฺมนฺติ, น เหวํ วตฺตพฺเพ ฯเปฯ.}

\prose{\csromansymbolmark{*}กายินฺทฺริยสํวโร กายกมฺมนฺติ, อามนฺตา.\csromanspacer{} จกฺขุนฺทฺริยสํวโร จกฺขุกมฺมนฺติ, น เหวํ วตฺตพฺเพ ฯเปฯ.\csromanspacer{} กายินฺทฺริยสํวโร กายกมฺมนฺติ, อามนฺตา.\csromanspacer{} โสตินฺทฺริยสํวโร ฯเปฯ.\csromanspacer{} ฆานินฺทฺริยสํวโร ฯเปฯ.\csromanspacer{} ชิวฺหินฺทฺริยสํวโร ชิวฺหากมฺมนฺติ, น เหวํ วตฺตพฺเพ ฯเปฯ.\csromanspacer{} มนินฺทฺริยสํวโร มโนกมฺมนฺติ, น เหวํ วตฺตพฺเพ ฯเปฯ.}

\csromanmatikaanchor{mtk.145}
\csromantocmarkref{subsection}{(117) 2. กมฺมกถา}{mtk.145}
\bookmark[dest={mtk.145},level=2]{(117) 2. กมฺมกถา}
\prose{\csromansymbolmark{*}มนินฺทฺริยสํวโร มโนกมฺมนฺติ, อามนฺตา.\csromanspacer{} จกฺขุนฺทฺริยสํวโร จกฺขุกมฺมนฺติ, น เหวํ วตฺตพฺเพ ฯเปฯ.\csromanspacer{} มนินฺทฺริยสํวโร มโนกมฺมนฺติ, อามนฺตา. \csromanfolio{336}โสตินฺทฺริยสํวโร ฯเปฯ.\csromanspacer{} ฆานินฺทฺริยสํวโร ฯเปฯ.\csromanspacer{} ชิวฺหินฺทฺริยสํวโร ฯเปฯ.\csromanspacer{} กายินฺทฺริยสํวโร กายกมฺมนฺติ, น เหวํ วตฺตพฺเพ ฯเปฯ.}

\proseitem{631}{\csromansymbolmark{*}อสํวโร กมฺมนฺติ, อามนฺตา.\csromanspacer{} จกฺขุนฺทฺริย-อสํวโร จกฺขุกมฺมนฺติ, น เหวํ วตฺตพฺเพ ฯเปฯ.\csromanspacer{} โสตินฺทฺริย-อสํวโร ฯเปฯ.\csromanspacer{} ฆานินฺทฺริย-อสํวโร ฯเปฯ.\csromanspacer{} ชิวฺหินฺทฺริย-อสํวโร ฯเปฯ.\csromanspacer{} กายินฺทฺริย- อสํวโร กายกมฺมนฺติ, น เหวํ วตฺตพฺเพ ฯเปฯ.}

\prose{\csromansymbolmark{*}กายินฺทฺริย-อสํวโร กายกมฺมนฺติ, อามนฺตา.\csromanspacer{} จกฺขุนฺทฺริย- อสํวโร จกฺขุกมฺมนฺติ, น เหวํ วตฺตพฺเพ ฯเปฯ.\csromanspacer{} กายินฺทฺริย-อสํวโร กายกมฺมนฺติ, อามนฺตา.\csromanspacer{} โสตินฺทฺริย-อสํวโร ฯเปฯ.\csromanspacer{} ฆานินฺทฺริย-อสํวโร ฯเปฯ.\csromanspacer{} ชิวฺหินฺทฺริย-อสํวโร ชิวฺหากมฺมนฺติ, น เหวํ วตฺตพฺเพ ฯเปฯ.\csromanspacer{} มนินฺทฺริย-อสํวโร มโนกมฺมนฺติ, น เหวํ วตฺตพฺเพ ฯเปฯ.}

\prose{\csromansymbolmark{*}มนินฺทฺริย-อสํวโร มโนกมฺมนฺติ, อามนฺตา.\csromanspacer{} จกฺขุนฺทฺริย- อสํวโร จกฺขุกมฺมนฺติ, น เหวํ วตฺตพฺเพ ฯเปฯ.\csromanspacer{} มนินฺทฺริย-อสํวโร มโนกมฺมนฺติ, อามนฺตา.\csromanspacer{} โสตินฺทฺริย-อสํวโร ฯเปฯ.\csromanspacer{} ฆานินฺทฺริย-อสํวโร ฯเปฯ.\csromanspacer{} ชิวฺหินฺทฺริย-อสํวโร ฯเปฯ.\csromanspacer{} กายินฺทฺริย-อสํวโร กายกมฺมนฺติ, น เหวํ วตฺตพฺเพ ฯเปฯ.}

\proseitem{632}{\csromansymbolmark{+}น วตฺตพฺพํ “สํวโรปิ อสํวโรปิ กมฺมนฺ”ติ, อามนฺตา.\csromanspacer{} นนุ วุตฺตํ ภควตา “อิธ ภิกฺขเว ภิกฺขุ จกฺขุนา รูปํ ทิสฺวา นิมิตฺตคฺคาหี โหตี ฯเปฯ น นิมิตฺตคฺคาหี โหติ.\csromanspacer{} โสเตน สทฺทํ สุตฺวา ฯเปฯ.\csromanspacer{} มนสา ธมฺมํ วิญฺญาย นิมิตฺตคฺคาหี โหติ ฯเปฯ น นิมิตฺตคฺคาหี โหตี”ติ อตฺเถว สุตฺตนฺโตติ, อามนฺตา.\csromanspacer{} เตน หิ สํวโรปิ อสํวโรปิ กมฺมนฺติ.}

\nitthitam{สํวโร กมฺมนฺติกถา นิฏฺฐิตา.}
\chapterhead{12. ทฺวาทสมวคฺค}
\vspace{\tipitakaparskip}
\csromancenter{\textbf{(117) 2.}\csromanspacer{}\textbf{ กมฺมกถา}}
\proseitem{633}{\csromansymbolmark{*}สพฺพํ กมฺมํ สวิปากนฺติ, อามนฺตา.\csromanspacer{} สพฺพา เจตนา สวิปากาติ, น เหวํ วตฺตพฺเพ ฯเปฯ.\csromanspacer{} สพฺพา เจตนา สวิปากาติ, อามนฺตา.\csromanspacer{} วิปากาพฺยากตา เจตนา สวิปากาติ, น เหวํ วตฺตพฺเพ ฯเปฯ.\csromanspacer{} สพฺพา เจตนา}

\vspace{\tipitakaparskip}
\csromancenter{ทฺวาทสมวคฺค สวิปากาติ, อามนฺตา.\csromanspacer{} กิริยาพฺยากตา เจตนา สวิปากาติ, น เหวํ วตฺตพฺเพ ฯเปฯ.}
\prose{\csromanfolio{337}\csromansymbolmark{*}สพฺพา เจตนา สวิปากาติ, อามนฺตา.\csromanspacer{} กามาวจรา วิปากาพฺยากตา เจตนา สวิปากาติ, น เหวํ วตฺตพฺเพ ฯเปฯ.\csromanspacer{} สพฺพา เจตนา สวิปากาติ, อามนฺตา.\csromanspacer{} รูปาวจรา อรูปาวจรา อปริยาปนฺนา วิปากาพฺยากตา เจตนา สวิปากาติ, น เหวํ วตฺตพฺเพ ฯเปฯ.}

\prose{\csromansymbolmark{*}สพฺพา เจตนา สวิปากาติ, อามนฺตา.\csromanspacer{} กามาวจรา กิริยาพฺยากตา เจตนา สวิปากาติ, น เหวํ วตฺตพฺเพ ฯเปฯ.\csromanspacer{} สพฺพา เจตนา สวิปากาติ, อามนฺตา.\csromanspacer{} รูปาวจรา อรูปาวจรา กิริยาพฺยากตา เจตนา สวิปากาติ, น เหวํ วตฺตพฺเพ ฯเปฯ.}

\proseitem{634}{\csromansymbolmark{*}วิปากาพฺยากตา เจตนา อวิปากาติ, อามนฺตา.\csromanspacer{} หญฺจิ วิปากาพฺยากตา เจตนา อวิปากา, โน จ วต เร วตฺตพฺเพ “สพฺพา เจตนา สวิปากา”ติ.}

\prose{\csromansymbolmark{*}กิริยาพฺยากตา เจตนา อวิปากาติ, อามนฺตา.\csromanspacer{} หญฺจิ กิริยาพฺยากตา เจตนา อวิปากา, โน จ วต เร วตฺตพฺเพ “สพฺพา เจตนา สวิปากา”ติ.}

\prose{\csromansymbolmark{*}กามาวจรา รูปาวจรา อรูปาวจรา อปริยาปนฺนา วิปากาพฺยากตา เจตนา อวิปากาติ, อามนฺตา.\csromanspacer{} หญฺจิ อปริยาปนฺนา วิปากาพฺยากตา เจตนา อวิปากา, โน จ วต เร วตฺตพฺเพ “สพฺพา เจตนา สวิปากา”ติ.}

\prose{\csromansymbolmark{*}กามาวจรา รูปาวจรา อรูปาวจรา กิริยาพฺยากตา เจตนา อวิปากาติ, อามนฺตา.\csromanspacer{} หญฺจิ อรูปาวจรา กิริยาพฺยากตา เจตนา อวิปากา, โน จ วต เร วตฺตพฺเพ “สพฺพา เจตนา สวิปากา”ติ.}

\proseitem{635}{\csromansymbolmark{+}น วตฺตพฺพํ “สพฺพํ กมฺมํ สวิปากนฺ”ติ, อามนฺตา.\csromanspacer{} นนุ วุตฺตํ ภควตา “นาหํ ภิกฺขเว สญฺเจตนิกานํ กมฺมานํ กตานํ อุปจิตานํ อปฺปฏิสํวิทิตฺวา พฺยนฺติภาวํ วทามิ, ตญฺจ โข ทิฏฺเฐว ธมฺเม อุปปชฺเช\footnote{อุปปชฺชํ (อํ 3. 497 ปิฏฺเฐ)} วา อปเร วา ปริยาเย”ติ\footnote{อํ 3. 497 ปิฏฺเฐ.} อตฺเถว สุตฺตนฺโตติ, อามนฺตา.\csromanspacer{} เตน หิ สพฺพํ กมฺมํ สวิปากนฺติ.}

\nitthitam{กมฺมกถา นิฏฺฐิตา.}
\chapterheadpageread{12. ทฺวาทสมวคฺค}
\vspace{\tipitakaparskip}
\csromancenter{\textbf{(118) 3.}\csromanspacer{}\textbf{ สทฺโทวิปาโกติกถา}}
\csromanmatikaanchor{mtk.146}
\csromantocmarkref{subsection}{(118) 3. สทฺโทวิปาโกติกถา}{mtk.146}
\bookmark[dest={mtk.146},level=2]{(118) 3. สทฺโทวิปาโกติกถา}
\proseitem{636}{\csromanfolio{338}\csromansymbolmark{*}สทฺโท วิปาโกติ, อามนฺตา.\csromanspacer{} สุขเวทนิโย ทุกฺขเวทนิโย อทุกฺขมสุขเวทนิโย สุขาย เวทนาย สมฺปยุตฺโต, ทุกฺขาย เวทนาย สมฺปยุตฺโต, อทุกฺขมสุขาย เวทนาย สมฺปยุตฺโต, ผสฺเสน สมฺปยุตฺโต, เวทนาย สมฺปยุตฺโต, สญฺญาย สมฺปยุตฺโต, เจตนาย สมฺปยุตฺโต, จิตฺเตน สมฺปยุตฺโต สารมฺมโณ, อตฺถิ ตสฺส อาวฏฺฏนา อาโภโค สมนฺนาหาโร มนสิกาโร เจตนา ปตฺถนา ปณิธีติ, น เหวํ วตฺตพฺเพ ฯเปฯ.\csromanspacer{} นนุ น สุขเวทนิโย น ทุกฺขเวทนิโย ฯเปฯ อนารมฺมโณ นตฺถิ ตสฺส อาวฏฺฏนา ฯเปฯ ปณิธีติ, อามนฺตา.\csromanspacer{} หญฺจิ น สุขเวทนิโย ฯเปฯ อนารมฺมโณ นตฺถิ ตสฺส อาวฏฺฏนา ฯเปฯ ปณิธิ, โน จ วต เร วตฺตพฺเพ “สทฺโท วิปาโก”ติ.}

\prose{ผสฺโส วิปาโก, ผสฺโส สุขเวทนิโย ฯเปฯ สารมฺมโณ, อตฺถิ ตสฺส อาวฏฺฏนา ฯเปฯ ปณิธีติ, อามนฺตา.\csromanspacer{} สทฺโท วิปาโก, สทฺโท สุขเวทนิโย ฯเปฯ สารมฺมโณ, อตฺถิ ตสฺส อาวฏฺฏนา ฯเปฯ ปณิธีติ, น เหวํ วตฺตพฺเพ ฯเปฯ.}

\prose{\csromansymbolmark{*}สทฺโท วิปาโก, สทฺโท น สุขเวทนิโย ฯเปฯ อนารมฺมโณ, นตฺถิ ตสฺส อาวฏฺฏนา ฯเปฯ ปณิธีติ, อามนฺตา.\csromanspacer{} ผสฺโส วิปาโก, ผสฺโส น สุขเวทนิโย น ทุกฺขเวทนิโย ฯเปฯ อนารมฺมโณ, นตฺถิ ตสฺส อาวฏฺฏนา ฯเปฯ ปณิธีติ, น เหวํ วตฺตพฺเพ ฯเปฯ.}

\proseitem{637}{\csromansymbolmark{+}น วตฺตพฺพํ “สทฺโท วิปาโก”ติ, อามนฺตา.\csromanspacer{} นนุ วุตฺตํ ภควตา “โส ตสฺส กมฺมสฺส กตตฺตา อุปจิตตฺตา อุสฺสนฺนตฺตา วิปุลตฺตา พฺรหฺมสฺสโร โหติ กรวิกภาณี”ติ\footnote{ที 3. 148 ปิฏฺเฐ.} อตฺเถว สุตฺตนฺโตติ, อามนฺตา.\csromanspacer{} เตน หิ สทฺโท วิปาโกติ.}

\nitthitam{สทฺโท วิปาโกติ กถา นิฏฺฐิตา.}
\chapterheadpageread{12. ทฺวาทสมวคฺค \textbf{12. ทฺวาทสมวคฺค}}
\vspace{\tipitakaparskip}
\csromancenter{\textbf{(119) 4.}\csromanspacer{}\textbf{ สฬายตนกถา}}
\csromanmatikaanchor{mtk.147}
\csromantocmarkref{subsection}{(119) 4. สฬายตนกถา}{mtk.147}
\bookmark[dest={mtk.147},level=2]{(119) 4. สฬายตนกถา}
\proseitem{638}{\csromanfolio{339}\csromansymbolmark{*}จกฺขายตนํ วิปาโกติ, อามนฺตา.\csromanspacer{} สุขเวทนิยํ, ทุกฺขเวทนิยํ ฯเปฯ สารมฺมณํ, อตฺถิ ตสฺส อาวฏฺฏนา ฯเปฯ ปณิธีติ, น เหวํ วตฺตพฺเพ ฯเปฯ.\csromanspacer{} นนุ น สุขเวทนิยํ ฯเปฯ อนารมฺมณํ นตฺถิ ตสฺส อาวฏฺฏนา ฯเปฯ ปณิธีติ, อามนฺตา.\csromanspacer{} หญฺจิ น สุขเวทนิยํ ฯเปฯ อนารมฺมณํ นตฺถิ ตสฺส อาวฏฺฏนา ฯเปฯ ปณิธิ, โน จ วต เร วตฺตพฺเพ “จกฺขายตนํ วิปาโก”ติ ฯเปฯ.}

\prose{\csromansymbolmark{*}ผสฺโส วิปาโก, ผสฺโส สุขเวทนิโย ฯเปฯ สารมฺมโณ, อตฺถิ ตสฺส อาวฏฺฏนา ฯเปฯ ปณิธีติ, อามนฺตา.\csromanspacer{} จกฺขายตนํ วิปาโก, จกฺขายตนํ สุขเวทนิยํ ฯเปฯ สารมฺมณํ, อตฺถิ ตสฺส อาวฏฺฏนา ฯเปฯ ปณิธีติ, น เหวํ วตฺตพฺเพ ฯเปฯ.}

\prose{\csromansymbolmark{*}จกฺขายตนํ วิปาโก, จกฺขายตนํ น สุขเวทนิยํ ฯเปฯ อนารมฺมณํ, นตฺถิ ตสฺส อาวฏฺฏนา ฯเปฯ ปณิธีติ, อามนฺตา.\csromanspacer{} ผสฺโส วิปาโก, ผสฺโส น สุขเวทนิโย ฯเปฯ อนารมฺมโณ, นตฺถิ ตสฺส อาวฏฺฏนา ฯเปฯ ปณิธีติ.\csromanspacer{} น เหวํ วตฺตพฺเพ ฯเปฯ.}

\proseitem{639}{\csromansymbolmark{*}โสตายตนํ ฯเปฯ.\csromanspacer{} ฆานายตนํ ฯเปฯ.\csromanspacer{} ชิวฺหายตนํ ฯเปฯ.\csromanspacer{} กายายตนํ วิปาโกติ, อามนฺตา.\csromanspacer{} สุขเวทนิยํ ฯเปฯ สารมฺมณํ, อตฺถิ ตสฺส อาวฏฺฏนา ฯเปฯ ปณิธีติ, น เหวํ วตฺตพฺเพ ฯเปฯ.\csromanspacer{} นนุ น สุขเวทนิยํ ฯเปฯ อนารมฺมณํ, นตฺถิ ตสฺส อาวฏฺฏนา ฯเปฯ ปณิธีติ, อามนฺตา.\csromanspacer{} หญฺจิ น สุขเวทนิยํ ฯเปฯ อนารมฺมณํ, นตฺถิ ตสฺส อาวฏฺฏนา ฯเปฯ ปณิธิ, โน จ วต เร วตฺตพฺเพ “กายายตนํ วิปาโก”ติ.}

\prose{\csromansymbolmark{*}ผสฺโส วิปาโก, ผสฺโส สุขเวทนิโย ฯเปฯ สารมฺมโณ, อตฺถิ ตสฺส อาวฏฺฏนา ฯเปฯ ปณิธีติ, อามนฺตา.\csromanspacer{} กายายตนํ วิปาโก, กายายตนํ สุขเวทนิยํ ฯเปฯ สารมฺมณํ, อตฺถิ ตสฺส อาวฏฺฏนา ฯเปฯ ปณิธีติ, น เหวํ วตฺตพฺเพ ฯเปฯ.\csromanspacer{} กายายตนํ วิปาโก, กายายตนํ น สุขเวทนิยํ ฯเปฯ อนารมฺมณํ, นตฺถิ ตสฺส อาวฏฺฏนา ฯเปฯ ปณิธีติ, อามนฺตา.\csromanspacer{} ผสฺโส วิปาโก, ผสฺโส น สุขเวทนิโย ฯเปฯ อนารมฺมโณ, นตฺถิ ตสฺส อาวฏฺฏนา ฯเปฯ ปณิธีติ, น เหวํ วตฺตพฺเพ ฯเปฯ.}

\csromanmatikaanchor{mtk.148}
\csromantocmarkref{subsection}{(120) 5. สตฺตกฺขตฺตุปรมกถา}{mtk.148}
\bookmark[dest={mtk.148},level=2]{(120) 5. สตฺตกฺขตฺตุปรมกถา}
\proseitem{640}{\csromanfolio{340}\csromansymbolmark{+}น วตฺตพฺพํ “สฬายตนํ วิปาโก”ติ, อามนฺตา.\csromanspacer{} นนุ สฬายตนํ กมฺมสฺส กตตฺตา อุปฺปนฺนนฺติ, อามนฺตา.\csromanspacer{} หญฺจิ สฬายตนํ กมฺมสฺส กตตฺตา อุปฺปนฺนํ, เตน วต เร วตฺตพฺเพ “สฬายตนํ วิปาโก”ติ.}

\nitthitam{สฬายตนกถา นิฏฺฐิตา.}
\chapterhead{12. ทฺวาทสมวคฺค}
\vspace{\tipitakaparskip}
\csromancenter{\textbf{(120) 5.}\csromanspacer{}\textbf{ สตฺตกฺขตฺตุปรมกถา}}
\proseitem{641}{\csromansymbolmark{*}สตฺตกฺขตฺตุปรโม ปุคฺคโล สตฺตกฺขตฺตุปรมตานิยโตติ, อามนฺตา.\csromanspacer{} มาตา ชีวิตา โวโรปิตา.\csromanspacer{} ปิตา ชีวิตา โวโรปิโต.\csromanspacer{} อรหา ชีวิตา โวโรปิโต.\csromanspacer{} ทุฏฺเฐน จิตฺเตน ตถาคตสฺส โลหิตํ อุปฺปาทิตํ.\csromanspacer{} สํโฆ ภินฺโนติ, น เหวํ วตฺตพฺเพ ฯเปฯ.}

\prose{\csromansymbolmark{*}สตฺตกฺขตฺตุปรโม ปุคฺคโล สตฺตกฺขตุปรมตานิยโตติ, อามนฺตา.\csromanspacer{} อภพฺโพ อนฺตรา ธมฺมํ อภิสเมตุนฺติ, น เหวํ วตฺตพฺเพ ฯเปฯ.\csromanspacer{} อภพฺโพ อนฺตรา ธมฺมํ อภิสเมตุนฺติ, อามนฺตา.\csromanspacer{} มาตา ชีวิตา โวโรปิตา.\csromanspacer{} ปิตา ชีวิตา โวโรปิโต.\csromanspacer{} อรหา ชีวิตา โวโรปิโต.\csromanspacer{} ทุฏฺเฐน จิตฺเตน ตถาคตสฺส โลหิตํ อุปฺปาทิตํ.\csromanspacer{} สํโฆ ภินฺโนติ, น เหวํ วตฺตพฺเพ ฯเปฯ.}

\proseitem{642}{\csromansymbolmark{*}สตฺตกฺขตฺตุปรโม ปุคฺคโล สตฺตกฺขตฺตุปรมตานิยโตติ, อามนฺตา.\csromanspacer{} อตฺถิ โส นิยโม เยน นิยเมน สตฺตกฺขตฺตุปรโม ปุคฺคโล สตฺตกฺขตฺตุปรมตานิยโตติ, น เหวํ วตฺตพฺเพ ฯเปฯ.\csromanspacer{} อตฺถิ เต สติปฏฺฐานา ฯเปฯ.\csromanspacer{} สมฺมปฺปธานา.\csromanspacer{} อิทฺธิปาทา, อินฺทฺริยา.\csromanspacer{} พลา.\csromanspacer{} โพชฺฌงฺคา เยหิ โพชฺฌงฺเคหิ สตฺตกฺขตฺตุปรโม ปุคฺคโล สตฺตกฺขตฺตุปรมตานิยโตติ, น เหวํ วตฺตพฺเพ ฯเปฯ.}

\proseitem{643}{\csromansymbolmark{*}นตฺถิ โส นิยโม เยน นิยเมน สตฺตกฺขตฺตุปรโม ปุคฺคโล สตฺตกฺขตฺตุปรมตานิยโตติ, อามนฺตา.\csromanspacer{} หญฺจิ นตฺถิ โส นิยโม เยน นิยเมน สตฺตกฺขตฺตุปรโม ปุคฺคโล สตฺตกฺขตฺตุปรมตานิยโต, โน จ วต เร วตฺตพฺเพ “สตฺตกฺขตฺตุปรโม ปุคฺคโล สตฺตกฺขตฺตุปรมตานิยโต”ติ.}

\chapterheadpageread{12. ทฺวาทสมวคฺค}
\csromanmatikaanchor{mtk.149}
\csromantocmarkref{subsection}{(121) 6. โกลํโกลกถา}{mtk.149}
\bookmark[dest={mtk.149},level=2]{(121) 6. โกลํโกลกถา}
\prose{\csromanfolio{341}\csromansymbolmark{*}นตฺถิ เต สติปฏฺฐานา ฯเปฯ โพชฺฌงฺคา เยหิ โพชฺฌงฺเคหิ สตฺตกฺขตฺตุปรโม ปุคฺคโล สตฺตกฺขตฺตุปรมตานิยโตติ, อามนฺตา.\csromanspacer{} หญฺจิ นตฺถิ เต โพชฺฌงฺคา เยหิ โพชฺฌงฺเคหิ สตฺตกฺขตฺตุปรโม ปุคฺคโล สตฺตกฺขตฺตุปรมตานิยโต, โน จ วต เร วตฺตพฺเพ “สตฺตกฺขตฺตุปรโม ปุคฺคโล สตฺตกฺขตฺตุปรมตานิยโต”ติ.}

\proseitem{644}{\csromansymbolmark{*}สตฺตกฺขตฺตุปรโม ปุคฺคโล สตฺตกฺขตฺตุปรมตานิยโตติ, อามนฺตา.\csromanspacer{} สกทาคามินิยเมนาติ, น เหวํ วตฺตพฺเพ ฯเปฯ.\csromanspacer{} อนาคามินิยเมนาติ, น เหวํ วตฺตพฺเพ ฯเปฯ.\csromanspacer{} อรหตฺตนิยเมนาติ, น เหวํ วตฺตพฺเพ ฯเปฯ.}

\prose{\csromansymbolmark{*}กตเมน นิยเมนาติ, โสตาปตฺตินิยเมนาติ.\csromanspacer{} สตฺตกฺขตฺตุปรโม ปุคฺคโล สตฺตกฺขตฺตุปรมตานิยโตติ, อามนฺตา.\csromanspacer{} เย เกจิ โสตาปตฺตินิยามํ โอกฺกมนฺติ, สพฺเพ เต สตฺตกฺขตฺตุปรมตานิยตาติ, น เหวํ วตฺตพฺเพ ฯเปฯ.}

\proseitem{645}{\csromansymbolmark{+}น วตฺตพฺพํ “สตฺตกฺขตฺตุปรโม ปุคฺคโล สตฺตกฺขตฺตุปรมตานิยโต”ติ, อามนฺตา.\csromanspacer{} นนุ โส สตฺตกฺขตฺตุปรโมติ, อามนฺตา.\csromanspacer{} หญฺจิ โส สตฺตกฺขตฺตุปรโม, เตน วต เร วตฺตพฺเพ “สตฺตกฺขตฺตุปรโม ปุคฺคโล สตฺตกฺขตฺตุปรมตานิยโต”ติ.}

\nitthitam{สตฺตกฺขตฺตุปรมกถา นิฏฺฐิตา.}
\chapterhead{12. ทฺวาทสมวคฺค}
\vspace{\tipitakaparskip}
\csromancenter{\textbf{(121) 6.}\csromanspacer{}\textbf{ โกลํโกลกถา}}
\proseitem{646}{\csromansymbolmark{+}น วตฺตพฺพํ “โกลํโกโล ปุคฺคโล โกลํโกลตานิยโต”ติ, อามนฺตา.\csromanspacer{} นนุ โส โกลํโกโลติ, อามนฺตา.\csromanspacer{} หญฺจิ โส โกลํโกโล, เตน วต เร วตฺตพฺเพ “โกลํโกโล ปุคฺคโล โกลํโกลตานิยโต”ติ.}

\nitthitam{โกลํโกลกถา นิฏฺฐิตา.}
\chapterheadpageread{12. ทฺวาทสมวคฺค}
\vspace{\tipitakaparskip}
\csromancenter{\textbf{(122) 7.}\csromanspacer{}\textbf{ เอกพีชีกถา}}
\csromanmatikaanchor{mtk.150}
\csromanmatikaanchor{mtk.151}
\csromantocmarkref{subsection}{(122) 7. เอกพีชีกถา}{mtk.150}
\bookmark[dest={mtk.150},level=2]{(122) 7. เอกพีชีกถา}
\csromantocmarkref{subsection}{(123) 8. ชีวิตาโวโรปนกถา}{mtk.151}
\bookmark[dest={mtk.151},level=2]{(123) 8. ชีวิตาโวโรปนกถา}
\proseitem{647}{\csromanfolio{342}\csromansymbolmark{+}น วตฺตพฺพํ “เอกพีชี ปุคฺคโล เอกพีชิตานิยโต”ติ, อามนฺตา.\csromanspacer{} นนุ โส เอกพีชีติ, อามนฺตา.\csromanspacer{} หญฺจิ โส เอกพีชี, เตน วต เร วตฺตพฺเพ “เอกพีชี ปุคฺคโล เอกพีชิตานิยโต”ติ.}

\nitthitam{เอกพีชีกถา นิฏฺฐิตา.}
\chapterhead{12. ทฺวาทสมวคฺค}
\vspace{\tipitakaparskip}
\csromancenter{\textbf{(123) 8.}\csromanspacer{}\textbf{ ชีวิตาโวโรปนกถา}}
\proseitem{648}{\csromansymbolmark{*}ทิฏฺฐิสมฺปนฺโน ปุคฺคโล สญฺจิจฺจ ปาณํ ชีวิตา โวโรเปยฺยาติ, อามนฺตา.\csromanspacer{} ทิฏฺฐิสมฺปนฺโน ปุคฺคโล สญฺจิจฺจ มาตรํ ชีวิตา โวโรเปยฺย ฯเปฯ ปิตรํ ชีวิตา โวโรเปยฺย ฯเปฯ อรหนฺตํ ชีวิตา โวโรเปยฺย ฯเปฯ ทุฏฺเฐน จิตฺเตน ตถาคตสฺส โลหิตํ อุปฺปาเทยฺย ฯเปฯ สํฆํ ภินฺเทยฺยาติ, น เหวํ วตฺตพฺเพ ฯเปฯ.}

\prose{\csromansymbolmark{*}ทิฏฺฐิสมฺปนฺโน ปุคฺคโล สญฺจิจฺจ ปาณํ ชีวิตา โวโรเปยฺยาติ, อามนฺตา.\csromanspacer{} ทิฏฺฐิสมฺปนฺโน ปุคฺคโล สตฺถริ อคารโวติ, น เหวํ วตฺตพฺเพ ฯเปฯ ธมฺเม ฯเปฯ สํเฆ ฯเปฯ สิกฺขาย อคารโวติ, น เหวํ วตฺตพฺเพ ฯเปฯ.}

\prose{\csromansymbolmark{*}นนุ ทิฏฺฐิสมฺปนฺโน ปุคฺคโล สตฺถริ สคารโวติ, อามนฺตา.\csromanspacer{} หญฺจิ ทิฏฺฐิสมฺปนฺโน ปุคฺคโล สตฺถริ สคารโว, โน จ วต เร วตฺตพฺเพ “ทิฏฺฐิสมฺปนฺโน ปุคฺคโล สญฺจิจฺจ ปาณํ ชีวิตา โวโรเปยฺยา”ติ.\csromanspacer{} นนุ ทิฏฺฐิสมฺปนฺโน ปุคฺคโล ธมฺเม ฯเปฯ สํเฆ ฯเปฯ สิกฺขาย สคารโวติ, อามนฺตา.\csromanspacer{} หญฺจิ ทิฏฺฐิสมฺปนฺโน ปุคฺคโล สิกฺขาย สคารโว, โน จ วต เร วตฺตพฺเพ “ทิฏฺฐิสมฺปนฺโน ปุคฺคโล สญฺจิจฺจ ปาณํ ชีวิตา โวโรเปยฺยา”ติ.}

\proseitem{649}{\csromansymbolmark{*}ทิฏฺฐิสมฺปนฺโน ปุคฺคโล สตฺถริ อคารโวติ, อามนฺตา.\csromanspacer{} ทิฏฺฐิสมฺปนฺโน ปุคฺคโล พุทฺธถูเป โอหเทยฺย โอมุตฺเตยฺย นิฏฺฐุเภยฺย พุทฺธถูเป อปพฺยามโต\footnote{อสพฺยากโต (สี, ก)} กเรยฺยาติ, น เหวํ วตฺตพฺเพ ฯเปฯ.}

\chapterheadpageread{12. ทฺวาทสมวคฺค}
\csromanmatikaanchor{mtk.152}
\csromantocmarkref{subsection}{(124) 9. ทุคฺคติกถา}{mtk.152}
\bookmark[dest={mtk.152},level=2]{(124) 9. ทุคฺคติกถา}
\prose{\csromanfolio{343}\csromansymbolmark{*}ทิฏฺฐิสมฺปนฺโน ปุคฺคโล สญฺจิจฺจ ปาณํ ชีวิตา โวโรเปยฺยาติ, อามนฺตา.\csromanspacer{} นนุ วุตฺตํ ภควตา “เสยฺยถาปิ ภิกฺขเว มหาสมุทฺโท ฐิตธมฺโม เวลํ นาติวตฺตติ.\csromanspacer{} เอวเมว โข ภิกฺขเว ยํ มยา สาวกานํ สิกฺขาปทํ ปญฺญตฺตํ, ตํ มม สาวกา ชีวิตเหตุปิ นาติกฺกมนฺตี”ติ\footnote{วิ 4. 421 ปิฏฺเฐ; อํ 3. 46 ปิฏฺเฐ; ขุ 1. 142 อุทาเน จ.} อตฺเถว สุตฺตนฺโตติ, อามนฺตา.\csromanspacer{} เตน หิ น วตฺตพฺพํ “ทิฏฺฐิสมฺปนฺโน ปุคฺคโล สญฺจิจฺจ ปาณํ ชีวิตา โวโรเปยฺยา”ติ.}

\nitthitam{ชีวิตา โวโรปนกถา นิฏฺฐิตา.}
\chapterhead{12. ทฺวาทสมวคฺค}
\vspace{\tipitakaparskip}
\csromancenter{\textbf{(124) 9.}\csromanspacer{}\textbf{ ทุคฺคติกถา}}
\proseitem{650}{\csromansymbolmark{*}ทิฏฺฐิสมฺปนฺนสฺส ปุคฺคลสฺส ปหีนา ทุคฺคตีติ, อามนฺตา.\csromanspacer{} ทิฏฺฐิสมฺปนฺโน ปุคฺคโล อาปายิเก รูเป รชฺเชยฺยาติ, อามนฺตา.\csromanspacer{} หญฺจิ ทิฏฺฐิสมฺปนฺโน ปุคฺคโล อาปายิเก รูเป รชฺเชยฺย, โน จ วต เร วตฺตพฺเพ “ทิฏฺฐิสมฺปนฺนสฺส ปุคฺคลสฺส ปหีนา ทุคฺคตี”ติ.}

\prose{\csromansymbolmark{*}ทิฏฺฐิสมฺปนฺนสฺส ปุคฺคลสฺส ปหีนา ทุคฺคตีติ, อามนฺตา.\csromanspacer{} ทิฏฺฐิสมฺปนฺโน ปุคฺคโล อาปายิเก สทฺเท ฯเปฯ คนฺเธ.\csromanspacer{} รเส.\csromanspacer{} โผฏฺฐพฺเพ ฯเปฯ อมนุสฺสิตฺถิยา ติรจฺฉานคติตฺถิยา นาคกญฺญาย เมถุนํ ธมฺมํ ปฏิเสเวยฺย.\csromanspacer{} อเชฬกํ ปฏิคฺคเณฺหยฺย.\csromanspacer{} กุกฺกุฏสูกรํ ปฏิคฺคเณฺหยฺย.\csromanspacer{} หตฺถิควสฺสวฬวํ ปฏิคฺคเณฺหยฺย.\csromanspacer{} ติตฺติรวฏฺฏกโมรกปิญฺชรํ\footnote{...กปิญฺชลํ (สฺยา, กํ, อิ)} ปฏิคฺคเณฺหยฺยาติ, อามนฺตา.\csromanspacer{} หญฺจิ ทิฏฺฐิสมฺปนฺโน ปุคฺคโล ติตฺติรวฏฺฏกโมรกปิญฺชรํ ปฏิคฺคเณฺหยฺย, โน จ วต เร วตฺตพฺเพ “ทิฏฺฐิสมฺปนฺนสฺส ปุคฺคลสฺส ปหีนา ทุคฺคตี”ติ.}

\csromanmatikaanchor{mtk.153}
\csromantocmarkref{subsection}{(125) 10. สตฺตมภวิกกถา}{mtk.153}
\bookmark[dest={mtk.153},level=2]{(125) 10. สตฺตมภวิกกถา}
\proseitem{651}{\csromansymbolmark{*}ทิฏฺฐิสมฺปนฺนสฺส ปุคฺคลสฺส ปหีนา ทุคฺคติ, ทิฏฺฐิสมฺปนฺโน ปุคฺคโล อาปายิเก รูเป รชฺเชยฺยาติ, อามนฺตา.\csromanspacer{} อรหโต ปหีนา ทุคฺคติ, อรหา อาปายิเก รูเป รชฺเชยฺยาติ, น เหวํ วตฺตพฺเพ ฯเปฯ.\csromanspacer{} ทิฏฺฐิสมฺปนฺนสฺส ปุคฺคลสฺส ปหีนา ทุคฺคติ, ทิฏฺฐิสมฺปนฺโน ปุคฺคโล อาปายิเก สทฺเท.\csromanspacer{} คนฺเธ.\csromanspacer{} รเส.\csromanspacer{} โผฏฺฐพฺเพ ฯเปฯ ติตฺติรวฏฺฏกโมรกปิญฺชรํ ปฏิคฺคเณฺหยฺยาติ, อามนฺตา. \csromanfolio{344}อรหโต ปหีนา ทุคฺคติ, อรหา ติตฺติรวฏฺฏกโมรกปิญฺชรํ ปฏิคฺคเณฺหยฺยาติ, น เหวํ วตฺตพฺเพ ฯเปฯ.}

\prose{\csromansymbolmark{*}อรหโต ปหีนา ทุคฺคติ, น จ อรหา อาปายิเก รูเป รชฺเชยฺยาติ, อามนฺตา.\csromanspacer{} ทิฏฺฐิสมฺปนฺนสฺส ปุคฺคลสฺส ปหีนา ทุคฺคติ, น จ ทิฏฺฐิสมฺปนฺโน ปุคฺคโล อาปายิเก รูเป รชฺเชยฺยาติ, น เหวํ วตฺตพฺเพ ฯเปฯ.\csromanspacer{} อรหโต ปหีนา ทุคฺคติ, น จ อรหา อาปายิเก สทฺเท ฯเปฯ คนฺเธ ฯเปฯ รเส ฯเปฯ โผฏฺฐพฺเพ ฯเปฯ อมนุสฺสิตฺถิยา ติรจฺฉานคติตฺถิยา นาคกญฺญาย เมถุนํ ธมฺมํ ปฏิเสเวยฺย.\csromanspacer{} อเชฬกํ ปฏิคฺคเณฺหยฺย.\csromanspacer{} กุกฺกุฏสูกรํ ปฏิคฺคเณฺหยฺย.\csromanspacer{} หตฺถิควสฺสวฬวํ ปฏิคฺคเณฺหยฺย ฯเปฯ ติตฺติรวฏฺฏกโมรกปิญฺชรํ ปฏิคฺคเณฺหยฺยาติ, อามนฺตา.\csromanspacer{} ทิฏฺฐิสมฺปนฺนสฺส ปุคฺคลสฺส ปหีนา ทุคฺคติ, น จ ทิฏฺฐิสมฺปนฺโน ปุคฺคโล ติตฺติรวฏฺฏกโมรกปิญฺชรํ ปฏิคฺคเณฺหยฺยาติ, น เหวํ วตฺตพฺเพ ฯเปฯ.}

\proseitem{652}{\csromansymbolfootnote{+}{อฏฺฐกถานุโลมํ ปรวาทีปุจฺฉาลกฺขณํ. ตถาปายํ ปุจฺฉา สกวาทิสฺส, ปุริมาโย จ อิมิสฺสํ ทุคฺคติกถายํ ปรวาทิสฺสาติ คเหตพฺพา วิย ทิสฺสนฺติ.}น วตฺตพฺพํ “ทิฏฺฐิสมฺปนฺนสฺส ปุคฺคลสฺส ปหีนา ทุคฺคตี”ติ, อามนฺตา.\csromanspacer{} ทิฏฺฐิสมฺปนฺโน ปุคฺคโล นิรยํ อุปปชฺเชยฺย ฯเปฯ ติรจฺฉานโยนิํ อุปปชฺเชยฺย.\csromanspacer{} เปตฺติวิสยํ อุปปชฺเชยฺยาติ, น เหวํ วตฺตพฺเพ.\csromanspacer{} เตน หิ ทิฏฺฐิสมฺปนฺนสฺส ปุคฺคลสฺส ปหีนา ทุคฺคตีติ.}

\nitthitam{ทุคฺคติกถา นิฏฺฐิตา.}
\chapterhead{12. ทฺวาทสมวคฺค}
\vspace{\tipitakaparskip}
\csromancenter{\textbf{(125) 10.}\csromanspacer{}\textbf{ สตฺตมภวิกกถา}}
\proseitem{635}{\csromansymbolmark{+}น วตฺตพฺพํ “สตฺตมภวิกสฺส ปุคฺคลสฺส ปหีนา ทุคฺคตี”ติ, อามนฺตา.\csromanspacer{} สตฺตมภวิโก ปุคฺคโล นิรยํ อุปปชฺเชยฺย.\csromanspacer{} ติรจฺฉานโยนิํ อุปปชฺเชยฺย.\csromanspacer{} เปตฺติวิสยํ อุปปชฺเชยฺยาติ, น เหวํ วตฺตพฺเพ.\csromanspacer{} เตน หิ สตฺตมภวิกสฺส ปุคฺคลสฺส ปหีนา ทุคฺคตีติ.}

\nitthitam{สตฺตมภวิกกถา นิฏฺฐิตา.}
\vspace{\tipitakaparskip}
\csromancenter{พารสโม วคฺโค.}
\chapterheadpageread{13. เตรสมวคฺค \textbf{ตสฺสุทฺทานํ}}
\csromanmatikaanchor{mtk.154}
\csromanmatikaanchor{mtk.155}
\csromantocmarknopage{chapter}{13. เตรสมวคฺค}{mtk.154}
\bookmark[dest={mtk.154},level=0]{13. เตรสมวคฺค}
\csromantocmarkref{subsection}{(126) 1. กปฺปฏฺฐกถา}{mtk.155}
\bookmark[dest={mtk.155},level=2]{(126) 1. กปฺปฏฺฐกถา}
\prose{\csromanfolio{345}สํวโร กมฺมํ ตเถว อสํวโร, สพฺพกมฺมํ สวิปากํ, สทฺโท วิปาโก, สฬายตนํ วิปาโก, สตฺตกฺขตฺตุปรโม ปุคฺคโล สตฺตกฺขตฺตุปรมตานิยโต, โกลํโกลปุคฺคโล โกลํโกลตานิยโต, เอกพีชี ปุคฺคโล เอกวีชิตานิยโต, ทิฏฺฐิสมฺปนฺโน ปุคฺคโล สญฺจิจฺจ ปาณํ ชีวิตา โวโรเปยฺย, ทิฏฺฐิสมฺปนฺนสฺส ปุคฺคลสฺส ปหีนา ทุคฺคติ, ตเถว สตฺตมภวิกสฺสาติ.}

\chapterhead{13. เตรสมวคฺค}
\vspace{\tipitakaparskip}
\csromancenter{\textbf{(126) 1.}\csromanspacer{}\textbf{ กปฺปฏฺฐกถา}}
\proseitem{654}{\csromansymbolmark{*}กปฺปฏฺโฐ กปฺปํ ติฏฺเฐยฺยาติ, อามนฺตา.\csromanspacer{} กปฺโป จ สณฺฐาติ พุทฺโธ จ โลเก อุปฺปชฺชตีติ, น เหวํ วตฺตพฺเพ ฯเปฯ.\csromanspacer{} กปฺปฏฺโฐ กปฺปํ ติฏฺเฐยฺยาติ, อามนฺตา.\csromanspacer{} กปฺโป จ สณฺฐาติ สํโฆ จ ภิชฺชตีติ, น เหวํ วตฺตพฺเพ ฯเปฯ.\csromanspacer{} กปฺปฏฺโฐ กปฺปํ ติฏฺเฐยฺยาติ, อามนฺตา.\csromanspacer{} กปฺโป จ สณฺฐาติ กปฺปฏฺโฐ จ กปฺปฏฺฐิยํ กมฺมํ กโรตีติ, น เหวํ วตฺตพฺเพ ฯเปฯ.\csromanspacer{} กปฺปฏฺโฐ กปฺปํ ติฏฺเฐยฺยาติ, อามนฺตา.\csromanspacer{} กปฺโป จ สณฺฐาติ กปฺปฏฺโฐ จ ปุคฺคโล กาลํ กโรตีติ, น เหวํ วตฺตพฺเพ ฯเปฯ.}

\proseitem{655}{\csromansymbolmark{*}กปฺปฏฺโฐ กปฺปํ ติฏฺเฐยฺยาติ, อามนฺตา.\csromanspacer{} อตีตํ กปฺปํ ติฏฺเฐยฺย.\csromanspacer{} อนาคตํ กปฺปํ ติฏฺเฐยฺยาติ, น เหวํ วตฺตพฺเพ ฯเปฯ.\csromanspacer{} กปฺปฏฺโฐ กปฺปํ ติฏฺเฐยฺยาติ, อามนฺตา.\csromanspacer{} เทฺว กปฺเป ติฏฺเฐยฺย.\csromanspacer{} ตโย กปฺเป ติฏฺเฐยฺย.\csromanspacer{} จตฺตาโร กปฺเป ติฏฺเฐยฺยาติ, น เหวํ วตฺตพฺเพ ฯเปฯ.}

\csromanmatikaanchor{mtk.156}
\csromantocmarkref{subsection}{(127) 2. กุสลปฏิลาภกถา}{mtk.156}
\bookmark[dest={mtk.156},level=2]{(127) 2. กุสลปฏิลาภกถา}
\proseitem{656}{\csromansymbolmark{*}กปฺปฏฺโฐ กปฺปํ ติฏฺเฐยฺยาติ, อามนฺตา.\csromanspacer{} กปฺปฏฺโฐ กปฺเป ฑยฺหนฺเต กตฺถ คจฺฉตีติ? อญฺญํ โลกธาตุํ คจฺฉตีติ.\csromanspacer{} มโต คจฺฉติ, เวหาสํ คจฺฉตีติ? มโต คจฺฉตีติ.\csromanspacer{} กปฺปฏฺฐิยํ กมฺมํ อปราปริยเวปกฺกนฺติ, น เหวํ วตฺตพฺเพ ฯเปฯ.1 เวหาสํ คจฺฉตีติ, อามนฺตา1.\csromanspacer{} กปฺปฏฺโฐ อิทฺธิมาติ, น เหวํ วตฺตพฺเพ ฯเปฯ.\csromanspacer{} กปฺปฏฺโฐ อิทฺธิมาติ, อามนฺตา. \csromanfolio{346}กปฺปฏฺเฐน ฉนฺทิทฺธิปาโท ภาวิโต วีริยิทฺธิปาโท ภาวิโต จิตฺติทฺธิปาโท ภาวิโต วีมํสิทฺธิปาโท ภาวิโตติ, น เหวํ วตฺตพฺเพ ฯเปฯ.}

\proseitem{657}{\csromansymbolmark{+}น วตฺตพฺพํ “กปฺปฏฺโฐ กปฺปํ ติฏฺเฐยฺยา”ติ, อามนฺตา.\csromanspacer{} นนุ วุตฺตํ ภควตา\csromandash{}}

\csromangathameasure{“อาปายิโก เนรยิโก, กปฺปฏฺโฐ สํฆเภทโก. \\ วคฺครโต อธมฺมฏฺโฐ, โยคกฺเขมา ปธํสติ. \\ สํฆํ สมคฺคํ เภตฺวาน, กปฺปํ นิรยมฺหิ ปจฺจตี”ติ.}
\csromangathagroup{\csromangathastanza{“อาปายิโก เนรยิโก, กปฺปฏฺโฐ สํฆเภทโก. \\ วคฺครโต อธมฺมฏฺโฐ, โยคกฺเขมา ปธํสติ. \\ สํฆํ สมคฺคํ เภตฺวาน\footnote{ภินฺทิตฺวา (สี, ก)}, กปฺปํ นิรยมฺหิ ปจฺจตี”ติ\footnote{วิ 4. 369 ปิฏฺเฐ; อํ 3. 315 ปิฏฺเฐ; ขุ 1. 203 ปิฏฺเฐปิ.}.}}

\prose{อตฺเถว สุตฺตนฺโตติ, อามนฺตา.\csromanspacer{} เตน หิ กปฺปฏฺโฐ กปฺปํ ติฏฺเฐยฺยาติ.}

\nitthitam{กปฺปฏฺฐกถา นิฏฺฐิตา.}
\chapterhead{13. เตรสมวคฺค}
\vspace{\tipitakaparskip}
\csromancenter{\textbf{(127) 2.}\csromanspacer{}\textbf{ กุสลปฏิลาภกถา}}
\proseitem{658}{\csromansymbolmark{*}กปฺปฏฺโฐ กุสลํ จิตฺตํ น ปฏิลเภยฺยาติ, อามนฺตา.\csromanspacer{} กปฺปฏฺโฐ ทานํ ทเทยฺยาติ, อามนฺตา.\csromanspacer{} หญฺจิ กปฺปฏฺโฐ ทานํ ทเทยฺย, โน จ วต เร วตฺตพฺเพ “กปฺปฏฺโฐ กุสลํ จิตฺตํ น ปฏิลเภยฺยา”ติ.}

\prose{\csromansymbolmark{*}กปฺปฏฺโฐ กุสลํ จิตฺตํ น ปฏิลเภยฺยาติ, อามนฺตา.\csromanspacer{} กปฺปฏฺโฐ จีวรํ ทเทยฺย ฯเปฯ ปิณฺฑปาตํ ทเทยฺย ฯเปฯ เสนาสนํ ทเทยฺย ฯเปฯ คิลานปจฺจยเภสชฺช ปริกฺขารํ ทเทยฺย.\csromanspacer{} ขาทนียํ ทเทยฺย.\csromanspacer{} โภชนียํ ทเทยฺย.\csromanspacer{} ปานียํ ทเทยฺย.\csromanspacer{} เจติยํ วนฺเทยฺย.\csromanspacer{} เจติเย มาลํ อาโรเปยฺย.\csromanspacer{} คนฺธํ อาโรเปยฺย.\csromanspacer{} วิเลปนํ อาโรเปยฺย ฯเปฯ เจติยํ อภิทกฺขิณํ\footnote{ปทกฺขิณํ (อิ)} กเรยฺยาติ, อามนฺตา.\csromanspacer{} หญฺจิ กปฺปฏฺโฐ เจติยํ อภิทกฺขิณํ กเรยฺย, โน จ วต เร วตฺตพฺเพ “กปฺปฏฺโฐ กุสลํ จิตฺตํ น ปฏิลเภยฺยา”ติ ฯเปฯ.}

\chapterheadpageread{13. เตรสมวคฺค}
\csromanmatikaanchor{mtk.157}
\csromantocmarkref{subsection}{(128) 3. อนนฺตราปยุตฺตกถา}{mtk.157}
\bookmark[dest={mtk.157},level=2]{(128) 3. อนนฺตราปยุตฺตกถา}
\proseitem{659}{\csromanfolio{347}\csromansymbolmark{+}กปฺปฏฺโฐ กุสลํ จิตฺตํ ปฏิลเภยฺยาติ, อามนฺตา.\csromanspacer{} ตโต วุฏฺฐานํ กุสลํ จิตฺตํ ปฏิลเภยฺยาติ, อามนฺตา.\csromanspacer{} รูปาวจรํ ฯเปฯ.\csromanspacer{} อรูปาวจรํ ฯเปฯ.\csromanspacer{} โลกุตฺตรํ กุสลํ จิตฺตํ ปฏิลเภยฺยาติ, น เหวํ วตฺตพฺเพ ฯเปฯ.}

\nitthitam{กุสลปฏิลาภกถา นิฏฺฐิตา.}
\chapterhead{13. เตรสมวคฺค}
\vspace{\tipitakaparskip}
\csromancenter{\textbf{(128) 3.}\csromanspacer{}\textbf{ อนนฺตราปยุตฺตกถา}}
\proseitem{660}{\csromansymbolmark{+}อนนฺตราปยุตฺโต ปุคฺคโล สมฺมตฺตนิยามํ โอกฺกเมยฺยาติ, อามนฺตา.\csromanspacer{} มิจฺฉตฺตนิยามญฺจ สมฺมตฺตนิยามญฺจ อุโภ โอกฺกเมยฺยาติ, น เหวํ วตฺตพฺเพ ฯเปฯ.}

\prose{\csromansymbolmark{+}อนนฺตราปยุตฺโต ปุคฺคโล สมฺมตฺตนิยามํ โอกฺกเมยฺยาติ, อามนฺตา.\csromanspacer{} นนุ ตํ กมฺมํ ปยุตฺตํ กุกฺกุจฺจํ อุปฺปาทิตํ วิปฺปฏิสาริยํ ชนิตนฺติ, อามนฺตา.\csromanspacer{} หญฺจิ ตํ กมฺมํ ปยุตฺตํ กุกฺกุจฺจํ อุปฺปาทิตํ วิปฺปฏิสาริยํ ชนิตํ, โน จ วต เร วตฺตพฺเพ “อนนฺตราปยุตฺโต ปุคฺคโล สมฺมตนิยามํ โอกฺกเมยฺยา”ติ.}

\proseitem{661}{\csromansymbolmark{*}อนนฺตราปยุตฺโต ปุคฺคโล อภพฺโพ สมฺมตฺตนิยามํ โอกฺกมิตุนฺติ, อามนฺตา.\csromanspacer{} มาตา ชีวิตา โวโรปิตา.\csromanspacer{} ปิตา ชีวิตา โวโรปิโต.\csromanspacer{} อรหา ชีวิตา โวโรปิโต.\csromanspacer{} ทุฏฺเฐน จิตฺเตน ตถาคตสฺส โลหิตํ อุปฺปาทิตํ.\csromanspacer{} สํโฆ ภินฺโนติ, น เหวํ วตฺตพฺเพ ฯเปฯ.}

\prose{\csromansymbolmark{*}อนนฺตราปยุตฺโต ปุคฺคโล ตํ กมฺมํ ปฏิสํหริตฺวา กุกฺกุจฺจํ ปฏิวิโนเทตฺวา วิปฺปฏิสาริยํ ปฏิวิเนตฺวา\footnote{ปฏิวิโนเทตฺวา (สี, อิ, ก)} อภพฺโพ สมฺมตฺตนิยามํ โอกฺกมิตุนฺติ, อามนฺตา.\csromanspacer{} มาตา ชีวิตา โวโรปิตา.\csromanspacer{} ปิตา ชีวิตา โวโรปิโต ฯเปฯ.\csromanspacer{} สํโฆ ภินฺโนติ, น เหวํ วตฺตพฺเพ ฯเปฯ.}

\csromanmatikaanchor{mtk.158}
\csromantocmarkref{subsection}{(129) 4. นิยตสฺสนิยามกถา}{mtk.158}
\bookmark[dest={mtk.158},level=2]{(129) 4. นิยตสฺสนิยามกถา}
\prose{\csromansymbolmark{*}อนนฺตราปยุตฺโต ปุคฺคโล ตํ กมฺมํ ปฏิสํหริตฺวา กุกฺกุจฺจํ ปฏิวิโนเทตฺวา วิปฺปฏิสาริยํ ปฏิวิเนตฺวา อภพฺโพ สมฺมตฺตนิยามํ โอกฺกมิตุนฺติ, อามนฺตา.\csromanspacer{} นนุ ตํ กมฺมํ ปฏิสํหตํ กุกฺกุจฺจํ ปฏิวิโนทิตํ วิปฺปฏิสาริยํ ปฏิวินีตนฺติ, อามนฺตา.\csromanspacer{} หญฺจิ ตํ กมฺมํ ปฏิสํหฏํ กุกฺกุจฺจํ ปฏิวิโนทิตํ \csromanfolio{348}วิปฺปฏิสาริยํ ปฏิวินีตํ, โน จ วต เร วตฺตพฺเพ “อนนฺตราปยุตฺโต ปุคฺคโล ตํ กมฺมํ ปฏิสํหริตฺวา กุกฺกุจฺจํ ปฏิวิโนเทตฺวา วิปฺปฏิสาริยํ ปฏิวิเนตฺวา อภพฺโพ สมฺมตฺตนิยามํ โอกฺกมิตุนฺ”ติ.}

\proseitem{662}{\csromansymbolmark{+}อนนฺตราปยุตฺโต ปุคฺคโล สมฺมตฺตนิยามํ โอกฺกเมยฺยาติ, อามนฺตา.\csromanspacer{} นนุ ตํ กมฺมํ ปยุตฺโต อาสีติ, อามนฺตา.\csromanspacer{} หญฺจิ ตํ กมฺมํ ปยุตฺโต อาสิ, โน จ วต เร วตฺตพฺเพ “อนนฺตราปยุตฺโต ปุคฺคโล สมฺมตฺตนิยามํ โอกฺกเมยฺยาติ.}

\nitthitam{อนนฺตราปยุตฺตกถา นิฏฺฐิตา.}
\chapterhead{13. เตรสมวคฺค}
\vspace{\tipitakaparskip}
\csromancenter{\textbf{(129) 4.}\csromanspacer{}\textbf{ นิยตสฺสนิยามกถา}}
\proseitem{663}{\csromansymbolmark{*}นิยโต นิยามํ โอกฺกมตีติ, อามนฺตา.\csromanspacer{} มิจฺฉตฺตนิยโต สมฺมตฺตนิยามํ โอกฺกมติ, สมฺมตฺตนิยโต มิจฺฉตฺตนิยามํ โอกฺกมตีติ, น เหวํ วตฺตพฺเพ ฯเปฯ.}

\prose{\csromansymbolmark{*}นิยโต นิยามํ โอกฺกมตีติ, อามนฺตา.\csromanspacer{} ปุพฺเพ มคฺคํ ภาเวตฺวา ปจฺฉา นิยามํ โอกฺกมตีติ, น เหวํ วตฺตพฺเพ ฯเปฯ.\csromanspacer{} ปุพฺเพ โสตาปตฺติมคฺคํ ภาเวตฺวา ปจฺฉา โสตาปตฺตินิยามํ โอกฺกมตีติ, น เหวํ วตฺตพฺเพ ฯเปฯ.\csromanspacer{} ปุพฺเพ สกทาคามิ ฯเปฯ อนาคามิ ฯเปฯ อรหตฺตมคฺคํ ภาเวตฺวา ปจฺฉา อรหตฺตนิยามํ โอกฺกมตีติ, น เหวํ วตฺตพฺเพ ฯเปฯ.}

\prose{\csromansymbolmark{*}ปุพฺเพ สติปฏฺฐานํ ฯเปฯ.\csromanspacer{} สมฺมปฺปธานํ.\csromanspacer{} อิทฺธิปาทํ.\csromanspacer{} อินฺทฺริยํ.\csromanspacer{} พลํ.\csromanspacer{} โพชฺฌงฺคํ ภาเวตฺวา ปจฺฉา นิยามํ โอกฺกมตีติ, น เหวํ วตฺตพฺเพ ฯเปฯ.}

\proseitem{664}{\csromansymbolmark{+}น วตฺตพฺพํ “นิยโต นิยามํ โอกฺกมตี’ติ, อามนฺตา.\csromanspacer{} ภพฺโพ โพธิสตฺโต ตาย ชาติยา ธมฺมํ นาภิสเมตุนฺติ, น เหวํ วตฺตพฺเพ.\csromanspacer{} เตน หิ นิยโต นิยามํ โอกฺกมตีติ.}

\nitthitam{นิยตสฺส นิยามกถา นิฏฺฐิตา.}
\chapterheadpageread{13. เตรสมวคฺค \textbf{13. เตรสมวคฺค}}
\vspace{\tipitakaparskip}
\csromancenter{\textbf{(130) 5.}\csromanspacer{}\textbf{ นิวุตกถา}}
\csromanmatikaanchor{mtk.159}
\csromantocmarkref{subsection}{(130) 5. นิวุตกถา}{mtk.159}
\bookmark[dest={mtk.159},level=2]{(130) 5. นิวุตกถา}
\proseitem{665}{\csromanfolio{349}\csromansymbolmark{*}นิวุโต นีวรณํ ชหตีติ, อามนฺตา.\csromanspacer{} รตฺโต ราคํ ชหติ.\csromanspacer{} ทุฏฺโฐ โทสํ ชหติ.\csromanspacer{} มูโฬฺห โมหํ ชหติ.\csromanspacer{} กิลิฏฺโฐ กิเลเส ชหตีติ, น เหวํ วตฺตพฺเพ ฯเปฯ.\csromanspacer{} ราเคน ราคํ ชหติ.\csromanspacer{} โทเสน โทสํ ชหติ.\csromanspacer{} โมเหน โมหํ ชหติ.\csromanspacer{} กิเลเสหิ กิเลเส ชหตีติ, น เหวํ วตฺตพฺเพ ฯเปฯ.}

\prose{\csromansymbolmark{*}ราโค จิตฺตสมฺปยุตฺโต, มคฺโค จิตฺตสมฺปยุตฺโตติ, อามนฺตา.\csromanspacer{} ทฺวินฺนํ ผสฺสานํ ฯเปฯ.\csromanspacer{} ทฺวินฺนํ จิตฺตานํ สโมธานํ โหตีติ, น เหวํ วตฺตพฺเพ ฯเปฯ.\csromanspacer{} ราโค อกุสโล, มคฺโค กุสโลติ, อามนฺตา.\csromanspacer{} กุสลากุสลา สาวชฺชานวชฺชา หีนปณีตา กณฺหสุกฺกสปฺปฏิภาคา ธมฺมา สมฺมุขีภาวํ อาคจฺฉนฺตีติ, น เหวํ วตฺตพฺเพ ฯเปฯ.}

\proseitem{666}{\csromansymbolmark{*}กุสลากุสลา สาวชฺชานวชฺชา หีนปณีตา กณฺหสุกฺกสปฺปฏิภาคา ธมฺมา สมฺมุขีภาวํ อาคจฺฉนฺตีติ, อามนฺตา.\csromanspacer{} นนุ วุตฺตํ ภควตา “จตฺตาริมานิ ภิกฺขเว สุวิทูรวิทูรานิ.\csromanspacer{} กตมานิ จตฺตาริ, นภญฺจ ภิกฺขเว ปถวี จ อิทํ ปฐมํ สุวิทูรวิทูรํ ฯเปฯ.\csromanspacer{} ตสฺมา สตํ ธมฺโม อสพฺภิ อารกา”ติ อตฺเถว สุตฺตนฺโตติ, อามนฺตา.\csromanspacer{} เตน หิ น วตฺตพฺพํ “กุสลากุสลา ฯเปฯ สมฺมุขีภาวํ อาคจฺฉนฺตี”ติ.}

\prose{\csromansymbolmark{*}นิวุโต นีวรณํ ชหตีติ, อามนฺตา.\csromanspacer{} นนุ วุตฺตํ ภควตา “โส เอวํ สมาหิเต จิตฺเต ปริสุทฺเธ ปริโยทาเต อนงฺคเณ วิคตูปกฺกิเลเส มุทุภูเต กมฺมนิเย ฐิเต อาเนญฺชปฺปตฺเต อาสวานํ ขยญาณาย จิตฺตํ อภินินฺนาเมตี”ติ อตฺเถว สุตฺตนฺโตติ, อามนฺตา.\csromanspacer{} เตน หิ น วตฺตพฺพํ “นิวุโต นีวรณํ ชหตี”ติ ฯเปฯ.}

\proseitem{667}{\csromansymbolmark{*}น วตฺตพฺพํ “นิวุโต นีวรณํ ชหตี”ติ, อามนฺตา.\csromanspacer{} นนุ วุตฺตํ ภควตา “ตสฺส เอวํ ชานโต เอวํ ปสฺสโต กามาสวาปิ จิตฺตํ วิมุจฺจติ ฯเปฯ.\csromanspacer{} อวิชฺชาสวาปิ จิตฺตํ วิมุจฺจตี”ติ อตฺเถว สุตฺตนฺโตติ, อามนฺตา.\csromanspacer{} เตน หิ นิวุโต นีวรณํ ชหตีติ.}

\nitthitam{นิวุตกถา นิฏฺฐิตา.}
\chapterheadpageread{13. เตรสมวคฺค}
\vspace{\tipitakaparskip}
\csromancenter{\textbf{(131) 6.}\csromanspacer{}\textbf{ สมฺมุขีภูตกถา}}
\csromanmatikaanchor{mtk.160}
\csromantocmarkref{subsection}{(131) 6. สมฺมุขีภูตกถา}{mtk.160}
\bookmark[dest={mtk.160},level=2]{(131) 6. สมฺมุขีภูตกถา}
\proseitem{668}{\csromanfolio{350}\csromansymbolmark{*}สมฺมุขีภูโต สํโยชนํ ชหตีติ, อามนฺตา.\csromanspacer{} รตฺโต ราคํ ชหติ.\csromanspacer{} ทุฏฺโฐ โทสํ ชหติ.\csromanspacer{} มูโฬฺห โมหํ ชหติ.\csromanspacer{} กิลิฏฺโฐ กิเลเส ชหตีติ, น เหวํ วตฺตพฺเพ ฯเปฯ.\csromanspacer{} ราเคน ราคํ ชหติ.\csromanspacer{} โทเสน โทสํ ชหติ.\csromanspacer{} โมเหน โมหํ ชหติ.\csromanspacer{} กิเลเสหิ กิเลเส ชหตีติ, น เหวํ วตฺตพฺเพ ฯเปฯ.}

\prose{\csromansymbolmark{*}ราโค จิตฺตสมฺปยุตฺโต, มคฺโค จิตฺตสมฺปยุตฺโตติ, อามนฺตา.\csromanspacer{} ทฺวินฺนํ ผสฺสานํ ฯเปฯ.\csromanspacer{} ทฺวินฺนํ จิตฺตานํ สโมธานํ โหตีติ, น เหวํ วตฺตพฺเพ ฯเปฯ.\csromanspacer{} ราโค อกุสโล, มคฺโค กุสโลติ, อามนฺตา.\csromanspacer{} กุสลากุสลา สาวชฺชานวชฺชา หีนปณีตา กณฺหสุกฺกสปฺปฏิภาคา ธมฺมา สมฺมุขีภาวํ อาคจฺฉนฺตีติ, น เหวํ วตฺตพฺเพ ฯเปฯ.}

\proseitem{669}{\csromansymbolmark{*}กุสลากุสลา ฯเปฯ สมฺมุขีภาวํ อาคจฺฉนฺตีติ, อามนฺตา.\csromanspacer{} นนุ วุตฺตํ ภควตา “จตฺตาริมานิ ภิกฺขเว สุวิทูรวิทูรานิ.\csromanspacer{} กตมานิ จตฺตาริ, นภญฺจ ภิกฺขเว ปถวี จ อิทํ ปฐมํ สุวิทูรวิทูรํ ฯเปฯ.\csromanspacer{} ตสฺมา สตํ ธมฺโม อสพฺภิ อารกา”ติ อตฺเถว สุตฺตนฺโตติ, อามนฺตา.\csromanspacer{} เตน หิ น วตฺตพฺพํ “กุสลากุสลา ฯเปฯ สมฺมุขีภาวํ อาคจฺฉนฺตี”ติ.}

\prose{\csromansymbolmark{*}สมฺมุขีภูโต สํโยชนํ ชหตีติ, อามนฺตา.\csromanspacer{} นนุ วุตฺตํ ภควตา “โส เอวํ สมาหิเต จิตฺเต ฯเปฯ อาสวานํ ขยญาณาย จิตฺตํ อภินินฺนาเมตี”ติ อตฺเถว สุตฺตนฺโตติ, อามนฺตา.\csromanspacer{} เตน หิ น วตฺตพฺพํ “สมฺมุขีภูโต สํโยชนํ ชหตี”ติ.}

\proseitem{670}{\csromansymbolmark{+}น วตฺตพฺพํ “สมฺมุขีภูโต สํโยชนํ ชหตี”ติ, อามนฺตา.\csromanspacer{} นนุ วุตฺตํ ภควตา “ตสฺส เอวํ ชานโต เอวํ ปสฺสโต กามาสวาปิ จิตฺตํ วิมุจฺจติ ฯเปฯ อวิชฺชาสวาปิ จิตฺตํ วิมุจฺจตี”ติ อตฺเถว สุตฺตนฺโตติ, อามนฺตา.\csromanspacer{} เตน หิ สมฺมุขีภูโต สํโยชนํ ชหตีติ.}

\nitthitam{สมฺมุขีภูตกถา นิฏฺฐิตา.}
\chapterheadpageread{13. เตรสมวคฺค \textbf{13. เตรสมวคฺค}}
\vspace{\tipitakaparskip}
\csromancenter{\textbf{(132) 7.}\csromanspacer{}\textbf{ สมาปนฺโน-อสฺสาเทติกถา}}
\csromanmatikaanchor{mtk.161}
\csromantocmarkref{subsection}{(132) 7. สมาปนฺโน-อสฺสาเทติกถา}{mtk.161}
\bookmark[dest={mtk.161},level=2]{(132) 7. สมาปนฺโน-อสฺสาเทติกถา}
\proseitem{671}{\csromanfolio{351}\csromansymbolmark{*}สมาปนฺโน อสฺสาเทติ, ฌานนิกนฺติ ฌานารมฺมณาติ, อามนฺตา.\csromanspacer{} ตํ ฌานํ ตสฺส ฌานสฺส อารมฺมณนฺติ, น เหวํ วตฺตพฺเพ ฯเปฯ.\csromanspacer{} ตํ ฌานํ ตสฺส ฌานสฺส อารมฺมณนฺติ, อามนฺตา.\csromanspacer{} เตน ผสฺเสน ตํ ผสฺสํ ผุสติ, ตาย เวทนาย ตํ เวทนํ เวเทติ, ตาย สญฺญาย ตํ สญฺญํ สญฺชานาติ, ตาย เจตนาย ตํ เจตนํ เจเตติ, เตน จิตฺเตน ตํ จิตฺตํ จินฺเตติ, เตน วิตกฺเกน ตํ วิตกฺกํ วิตกฺเกติ, เตน วิจาเรน ตํ วิจารํ วิจาเรติ, ตาย ปีติยา ตํ ปีติํ ปิยายติ, ตาย สติยา ตํ สติํ สรติ, ตาย ปญฺญาย ตํ ปญฺญํ ปชานาตีติ, น เหวํ วตฺตพฺเพ ฯเปฯ.}

\prose{\csromansymbolmark{*}ฌานนิกนฺติ จิตฺตสมฺปยุตฺตา, ฌานํ จิตฺตสมฺปยุตฺตนฺติ, อามนฺตา.\csromanspacer{} ทฺวินฺนํ ผสฺสานํ ฯเปฯ.\csromanspacer{} ทฺวินฺนํ จิตฺตานํ สโมธานํ โหตีติ, น เหวํ วตฺตพฺเพ ฯเปฯ.}

\prose{\csromansymbolmark{*}ฌานนิกนฺติ อกุสลํ, ฌานํ กุสลนฺติ, อามนฺตา.\csromanspacer{} กุสลากุสลา สาวชฺชานวชฺชา หีนปณีตา กณฺหสุกฺกสปฺปฏิภาคา ธมฺมา สมฺมุขีภาวํ อาคจฺฉนฺตีติ, น เหวํ วตฺตพฺเพ ฯเปฯ.}

\proseitem{672}{\csromansymbolmark{*}กุสลากุสลา สาวชฺชานวชฺชา หีนปณีตา กณฺหสุกฺกสปฺปฏิภาคา ธมฺมา สมฺมุขีภาวํ อาคจฺฉนฺตีติ, อามนฺตา.\csromanspacer{} นนุ วุตฺตํ ภควตา “จตฺตาริมานิ ภิกฺขเว สุวิทูรวิทูรานิ.\csromanspacer{} กตมานิ จตฺตาริ, นภญฺจ ภิกฺขเว ปถวี จ อิทํ ปฐมํ สุวิทูรวิทูรํ ฯเปฯ.\csromanspacer{} ตสฺมา สตํ ธมฺโม อสพฺภิ อารกา”ติ อตฺเถว สุตฺตนฺโตติ, อามนฺตา.\csromanspacer{} เตน หิ จ วตฺตพฺพํ “กุสลากุสลา สาวชฺชานวชฺชา หีนปณีตา กณฺหสุกฺกสปฺปฏิภาคา ธมฺมา สมฺมุขีภาวํ อาคจฺฉนฺตี”ติ.}

\csromanmatikaanchor{mtk.162}
\csromantocmarkref{subsection}{(133) 8. อสาตราคกถา}{mtk.162}
\bookmark[dest={mtk.162},level=2]{(133) 8. อสาตราคกถา}
\proseitem{673}{\csromansymbolmark{+}น วตฺตพฺพํ “สมาปนฺโน อสฺสาเทติ, ฌานนิกนฺติ ฌานารมฺมณา”ติ, อามนฺตา.\csromanspacer{} นนุ วุตฺตํ ภควตา “อิธ ภิกฺขเว ภิกฺขุ วิวิจฺเจว กาเมหิ วิวิจฺจ อกุสเลหิ ธมฺเมหิ ปฐมํ ฌานํ อุปสมฺปชฺช วิหรติ, โส ตํ อสฺสาเทติ ตํ นิกาเมติ เตน จ วิตฺติํ อาปชฺชติ.\csromanspacer{} วิตกฺกวิจารานํ วูปสมา ฯเปฯ ทุติยํ ฌานํ ฯเปฯ ตติยํ ฌานํ ฯเปฯ จตุตฺถํ ฌานํ อุปสมฺปชฺช \csromanfolio{352}วิหรติ, โส ตํ อสฺสาเทติ ตํ นิกาเมติ เตน จ วิตฺติํ อาปชฺชตี”ติ\footnote{อํ 1. 441 ปิฏฺเฐ.} อตฺเถว สุตฺตนฺโตติ, อามนฺตา.\csromanspacer{} เตน หิ สมาปนฺโน อสฺสาเทติ, ฌานนิกนฺติ ฌานารมฺมณาติ.}

\nitthitam{สมาปนฺโน อสฺสาเทติกถา นิฏฺฐิตา.}
\chapterhead{13. เตรสมวคฺค}
\vspace{\tipitakaparskip}
\csromancenter{\textbf{(133) 8.}\csromanspacer{}\textbf{ อสาตราคกถา}}
\proseitem{674}{\csromansymbolmark{*}อตฺถิ อสาตราโคติ, อามนฺตา.\csromanspacer{} ทุกฺขาภินนฺทิโน สตฺตา, อตฺถิ เกจิ ทุกฺขํ ปตฺเถนฺติ ปิเหนฺติ เอสนฺติ คเวสนฺติ ปริเยสนฺติ ทุกฺขํ อชฺโฌสาย ติฏฺฐนฺตีติ, น เหวํ วตฺตพฺเพ ฯเปฯ.\csromanspacer{} นนุ สุขาภินนฺทิโน สตฺตา, อตฺถิ เกจิ สุขํ ปตฺเถนฺติ ปิเหนฺติ เอสนฺติ คเวสนฺติ ปริเยสนฺติ สุขํ อชฺโฌสาย ติฏฺฐนฺตีติ, อามนฺตา.\csromanspacer{} หญฺจิ สุขาภินนฺทิโน สตฺตา, อตฺถิ เกจิ สุขํ ปตฺเถนฺติ ปิเหนฺติ เอสนฺติ คเวสนฺติ ปริเยสนฺติ สุขํ อชฺโฌสาย ติฏฺฐนฺติ, โน จ วต เร วตฺตพฺเพ “อตฺถิ อสาตราโค”ติ.}

\prose{\csromansymbolmark{*}อตฺถิ อสาตราโคติ, อามนฺตา.\csromanspacer{} ทุกฺขาย เวทนาย ราคานุสโย อนุเสติ, สุขาย เวทนาย ปฏิฆานุสโย อนุเสตีติ, น เหวํ วตฺตพฺเพ ฯเปฯ.\csromanspacer{} นนุ สุขาย เวทนาย ราคานุสโย อนุเสติ, ทุกฺขาย เวทนาย ปฏิฆานุสโย อนุเสตีติ, อามนฺตา.\csromanspacer{} หญฺจิ สุขาย เวทนาย ราคานุสโย อนุเสติ, ทุกฺขาย เวทนาย ปฏิฆานุสโย อนุเสติ, โน จ วต เร วตฺตพฺเพ “อตฺถิ อสาตราโค”ติ.}

\proseitem{675}{\csromansymbolmark{+}น วตฺตพฺพํ “อตฺถิ อสาตราโค”ติ, อามนฺตา.\csromanspacer{} นนุ วุตฺตํ ภควตา “โส เอวํ อนุโรธวิโรธํ สมาปนฺโน ยํ กิญฺจิ เวทนํ เวทยติ สุขํ วา ทุกฺข วา อทุกฺขมสุขํ วา, โส ตํ เวทนํ อภินนฺทติ อภิวทติ อชฺโฌสาย ติฏฺฐตี”ติ\footnote{ม 1. 333 ปิฏฺเฐ มหาตณฺหาสงฺขเย.} อตฺเถว สุตฺตนฺโตติ, อามนฺตา.\csromanspacer{} เตน หิ อตฺถิ อสาตราโคติ.}

\nitthitam{อสาตราคกถา นิฏฺฐิตา.}
\chapterheadpageread{13. เตรสมวคฺค \textbf{13. เตรสมวคฺค}}
\vspace{\tipitakaparskip}
\csromancenter{\textbf{(134) 9.}\csromanspacer{}\textbf{ ธมฺมตณฺหา-อพฺยากตาติกถา}}
\csromanmatikaanchor{mtk.163}
\csromantocmarkref{subsection}{(134) 9. ธมฺมตณฺหา-อพฺยากตาติกถา}{mtk.163}
\bookmark[dest={mtk.163},level=2]{(134) 9. ธมฺมตณฺหา-อพฺยากตาติกถา}
\proseitem{676}{\csromanfolio{353}\csromansymbolmark{*}ธมฺมตณฺหา อพฺยากตาติ, อามนฺตา.\csromanspacer{} วิปากาพฺยากตา กิริยาพฺยากตา รูปํ นิพฺพานํ จกฺขายตนํ ฯเปฯ โผฏฺฐพฺพายตนนฺติ, น เหวํ วตฺตพฺเพ ฯเปฯ.}

\prose{\csromansymbolmark{*}ธมฺมตณฺหา อพฺยากตาติ, อามนฺตา.\csromanspacer{} รูปตณฺหา อพฺยากตาติ, น เหวํ วตฺตพฺเพ ฯเปฯ.\csromanspacer{} ธมฺมตณฺหา อพฺยากตาติ, อามนฺตา.\csromanspacer{} สทฺทตณฺหา ฯเปฯ.\csromanspacer{} คนฺธตณฺหา ฯเปฯ.\csromanspacer{} รสตณฺหา ฯเปฯ.\csromanspacer{} โผฏฺฐพฺพตณฺหา อพฺยากตาติ, น เหวํ วตฺตพฺเพ ฯเปฯ.}

\prose{\csromansymbolmark{*}รูปตณฺหา อกุสลาติ, อามนฺตา.\csromanspacer{} ธมฺมตณฺหา อกุสลาติ, น เหวํ วตฺตพฺเพ ฯเปฯ.\csromanspacer{} สทฺทตณฺหา ฯเปฯ.\csromanspacer{} โผฏฺฐพฺพตณฺหา อกุสลาติ, อามนฺตา.\csromanspacer{} ธมฺมตณฺหา อกุสลาติ, น เหวํ วตฺตพฺเพ ฯเปฯ.}

\proseitem{677}{\csromansymbolmark{*}ธมฺมตณฺหา อพฺยากตาติ, อามนฺตา.\csromanspacer{} นนุ ตณฺหา อกุสลา วุตฺตา ภควตาติ, อามนฺตา.\csromanspacer{} หญฺจิ ตณฺหา อกุสลา วุตฺตา ภควตา, โน จ วต เร วตพฺเพ “ธมฺมตณฺหา อพฺยากตา”ติ.}

\prose{\csromansymbolmark{*}ธมฺมตณฺหา อพฺยากตาติ, อามนฺตา.\csromanspacer{} นนุ โลโภ อกุสโล วุตฺโต ภควตา, ธมฺมตณฺหา โลโภติ, อามนฺตา.\csromanspacer{} หญฺจิ โลโภ อกุสโล วุตฺโต ภควตา, ธมฺมตณฺหา โลโภ, โน จ วต เร วตฺตพฺเพ “ธมฺมตณฺหา อพฺยากตา”ติ.}

\proseitem{678}{\csromansymbolmark{*}ธมฺมตณฺหา โลโภ อพฺยากโตติ, อามนฺตา.\csromanspacer{} รูปตณฺหา โลโภ อพฺยากโตติ, น เหวํ วตฺตพฺเพ ฯเปฯ.\csromanspacer{} ธมฺมตณฺหา โลโภ อพฺยากโตติ, อามนฺตา.\csromanspacer{} สทฺทตณฺหา ฯเปฯ.\csromanspacer{} โผฏฺฐพฺพตณฺหา โลโภ อพฺยากโตติ, น เหวํ วตฺตพฺเพ ฯเปฯ.}

\prose{\csromansymbolmark{*}รูปตณฺหา โลโภ อกุสโลติ, อามนฺตา.\csromanspacer{} ธมฺมตณฺหา โลโภ อกุสโลติ, น เหวํ วตฺตพฺเพ ฯเปฯ.\csromanspacer{} สทฺทตณฺหา ฯเปฯ.\csromanspacer{} โผฏฺฐพฺพตณฺหา โลโภ อกุสโลติ, อามนฺตา.\csromanspacer{} ธมฺมตณฺหา โลโภ อกุสโลติ, น เหวํ วตฺตพฺเพ ฯเปฯ.}

\csromanmatikaanchor{mtk.164}
\csromantocmarkref{subsection}{(135) 10. ธมฺมตณฺหานทุกฺขสมุทโยติกถา ...}{mtk.164}
\bookmark[dest={mtk.164},level=2]{(135) 10. ธมฺมตณฺหานทุกฺขสมุทโยติกถา ...}
\proseitem{679}{\csromanfolio{354}\csromansymbolmark{*}ธมฺมตณฺหา อพฺยากตาติ, อามนฺตา.\csromanspacer{} นนุ วุตฺตํ ภควตา “ยายํ ตณฺหา โปโนพฺภวิกา นนฺทีราคสหคตา ตตฺรตตฺราภินนฺทินี.\csromanspacer{} เสยฺยถิทํ, กามตณฺหา ภวตณฺหา วิภวตณฺหา”ติ\footnote{วิ 3. 15 ปิฏฺเฐ; ที 2. 246 ปิฏฺฐาทีสุ จ.} อตฺเถว สุตฺตนฺโตติ, อามนฺตา.\csromanspacer{} เตน หิ น วตฺตพฺพํ “ธมฺมตณฺหา อพฺยากตา”ติ.}

\proseitem{680}{\csromansymbolmark{+}น วตฺตพฺพํ “ธมฺมตณฺหา อพฺยากตา”ติ, อามนฺตา.\csromanspacer{} นนุ สา ธมฺมตณฺหาติ, อามนฺตา.\csromanspacer{} หญฺจิ สา ธมฺมตณฺหา, เตน วต เร วตฺตพฺเพ “ธมฺมตณฺหา อพฺยากตา”ติ.}

\nitthitam{ธมฺมตณฺหา อพฺยากตาติกถา นิฏฺฐิตา.}
\chapterhead{13. เตรสมวคฺค}
\vspace{\tipitakaparskip}
\csromancenter{\textbf{(135) 10.}\csromanspacer{}\textbf{ ธมฺมตณฺหานทุกฺขสมุทโยติกถา}}
\proseitem{681}{\csromansymbolmark{*}ธมฺมตณฺหา น ทุกฺขสมุทโยติ, อามนฺตา.\csromanspacer{} รูปตณฺหา น ทุกฺขสมุทโยติ, น เหวํ วตฺตพฺเพ ฯเปฯ.\csromanspacer{} ธมฺมตณฺหา น ทุกฺขสมุทโยติ, อามนฺตา.\csromanspacer{} สทฺทตณฺหา ฯเปฯ.\csromanspacer{} คนฺธตณฺหา ฯเปฯ.\csromanspacer{} รสตณฺหา ฯเปฯ.\csromanspacer{} โผฏฺฐพฺพตณฺหา น ทุกฺขสมุทโยติ, น เหวํ วตฺตพฺเพ ฯเปฯ.}

\prose{\csromansymbolmark{*}รูปตณฺหา ทุกฺขสมุทโยติ, อามนฺตา.\csromanspacer{} ธมฺมตณฺหา ทุกฺขสมุทโยติ, น เหวํ วตฺตพฺเพ ฯเปฯ.\csromanspacer{} สทฺทตณฺหา ฯเปฯ.\csromanspacer{} คนฺธตณฺหา ฯเปฯ.\csromanspacer{} รสตณฺหา ฯเปฯ.\csromanspacer{} โผฏฺฐพฺพตณฺหา ทุกฺขสมุทโยติ, อามนฺตา.\csromanspacer{} ธมฺมตณฺหา ทุกฺขสมุทโยติ, น เหวํ วตฺตพฺเพ ฯเปฯ.}

\proseitem{682}{\csromansymbolmark{*}ธมฺมตณฺหา น ทุกฺขสมุทโยติ, อามนฺตา.\csromanspacer{} นนุ ตณฺหา ทุกฺขสมุทโย วุตฺโต ภควตาติ, อามนฺตา.\csromanspacer{} หญฺจิ ตณฺหาทุกฺขสมุทโย วุตฺโต ภควตา, โน จ วต เร วตฺตพฺเพ “ธมฺมตณฺหา น ทุกฺขสมุทโย”ติ.\csromanspacer{} ธมฺมตณฺหา น ทุกฺขสมุทโยติ, อามนฺตา.\csromanspacer{} นนุ โลโภ ทุกฺขสมุทโย วุตฺโต ภควตา, ธมฺมตณฺหา โลโภติ, อามนฺตา.\csromanspacer{} หญฺจิ โลโภ}

\vspace{\tipitakaparskip}
\csromancenter{เตรสมวคฺค ทุกฺขสมุทโย วุตฺโต ภควตา, ธมฺมตณฺหา โลโภ, โน จ วต เร วตฺตพฺเพ “ธมฺมตณฺหา น ทุกฺขสมุทโย”ติ.}
\proseitem{683}{\csromanfolio{355}\csromansymbolmark{*}ธมฺมตณฺหา โลโภ น ทุกฺขสมุทโยติ, อามนฺตา.\csromanspacer{} รูปตณฺหา โลโภ น ทุกฺขสมุทโยติ, น เหวํ วตฺตพฺเพ ฯเปฯ.\csromanspacer{} ธมฺมตณฺหา โลโภ น ทุกฺขสมุทโยติ, อามนฺตา.\csromanspacer{} สทฺทตณฺหา ฯเปฯ.\csromanspacer{} คนฺธตณฺหา ฯเปฯ.\csromanspacer{} รสตณฺหา ฯเปฯ.\csromanspacer{} โผฏฺฐพฺพตณฺหา โลโภ น ทุกฺขสมุทโยติ, น เหวํ วตฺตพฺเพ ฯเปฯ.}

\prose{\csromansymbolmark{*}รูปตณฺหา โลโภ ทุกฺขสมุทโยติ, อามนฺตา.\csromanspacer{} ธมฺมตณฺหา โลโภ ทุกฺขสมุทโยติ, น เหวํ วตฺตพฺเพ ฯเปฯ.\csromanspacer{} สทฺทตณฺหา ฯเปฯ.\csromanspacer{} โผฏฺฐพฺพตณฺหา โลโภ ทุกฺขสมุทโยติ, อามนฺตา.\csromanspacer{} ธมฺมตณฺหา โลโภ ทุกฺขสมุทโยติ, น เหวํ วตฺตพฺเพ ฯเปฯ.}

\proseitem{684}{\csromansymbolmark{*}ธมฺมตณฺหา น ทุกฺขสมุทโยติ, อามนฺตา.\csromanspacer{} นนุ วุตฺตํ ภควตา “ยายํ ตณฺหา โปโนพฺภวิกา นนฺทีราคสหคตา ตตฺรตตฺราภินนฺทินี.\csromanspacer{} เสยฺยถิทํ, กามตณฺหา ภวตณฺหา วิภวตณฺหา”ติ อตฺเถวสุตฺตนฺโตติ, อามนฺตา.\csromanspacer{} เตน หิ น วตฺตพฺพํ “ธมฺมตณฺหา น ทุกฺขสมุทโย”ติ.}

\proseitem{685}{\csromansymbolmark{+}น วตฺตพฺพํ “ธมฺมตณฺหา น ทุกฺขสมุทโย”ติ, อามนฺตา.\csromanspacer{} นนุ สา ธมฺมตณฺหาติ, อามนฺตา.\csromanspacer{} หญฺจิ สา ธมฺมตณฺหา, เตน วต เร วตฺตพฺเพ “ธมฺมตณฺหา น ทุกฺขสมุทโย”ติ.\csromanspacer{} ธมฺมตณฺหา น ทุกฺขสมุทโยติกถา นิฏฺฐิตา.\csromanspacer{} เตรสโม วคฺโค.}

\titlehead{\textbf{ตสฺสุทฺทานํ}}
\prose{กปฺปฏฺโฐ กปฺปํ ติฏฺเฐยฺย, กปฺปฏฺโฐ กุสลํ จิตฺตํ น ปฏิลเภยฺย, อนนฺตราปยุตฺโต ปุคฺคโล สมฺมตฺตนิยามํ โอกฺกเมยฺย, นิยโต นิยามํ โอกฺกมติ, นิวุโต นีวรณํ ชหติ, สมฺมุขีภูโต สํโยชนํ ชหติ, ฌานนิกนฺติ, อสาตราโค, ธมฺมตณฺหา อพฺยากตา, ธมฺมตณฺหา น ทุกฺขสมุทโยติ.}

\csromanmatikaanchor{mtk.165}
\csromantocmarknopage{chapter}{14. จุทฺทสมวคฺค}{mtk.165}
\bookmark[dest={mtk.165},level=0]{14. จุทฺทสมวคฺค}
\chapterheadpageread{14. จุทฺทสมวคฺค}
\vspace{\tipitakaparskip}
\csromancenter{\textbf{(136) 1.}\csromanspacer{}\textbf{ กุสลากุสลปฏิสนฺทหนกถา}}
\csromanmatikaanchor{mtk.166}
\csromantocmarkref{subsection}{(136) 1. กุสลากุสลปฏิสนฺทหนกถา}{mtk.166}
\bookmark[dest={mtk.166},level=2]{(136) 1. กุสลากุสลปฏิสนฺทหนกถา}
\proseitem{686}{\csromanfolio{356}อกุสลมูลํ ปฏิสนฺทหติ กุสลมูลนฺติ, อามนฺตา.\csromanspacer{} ยา อกุสลสฺส อุปฺปาทาย อาวฏฺฏนา ฯเปฯ ปณิธิ, สาว กุสลสฺส อุปฺปาทาย อาวฏฺฏนา ฯเปฯ ปณิธีติ, น เหวํ วตฺตพฺเพ ฯเปฯ.}

\prose{\csromansymbolmark{*}อกุสลมูลํ ปฏิสนฺทหติ กุสลมูลํ, น วตฺตพฺพํ “ยา อกุสลสฺส อุปฺปาทาย อาวฏฺฏนา ฯเปฯ ปณิธิ, สาว กุสลสฺส อุปฺปาทาย อาวฏฺฏนา ฯเปฯ ปณิธี”ติ, อามนฺตา.\csromanspacer{} กุสลํ อนาวฏฺเฏนฺตสฺส\footnote{อนาวฏฺฏนฺตสฺส (สี, อิ, ก), อนาวชฺชนฺตสฺส (สฺยา, กํ)} อุปฺปชฺชติ ฯเปฯ อปฺปณิทหนฺตสฺส อุปฺปชฺชตีติ, น เหวํ วตฺตพฺเพ ฯเปฯ.\csromanspacer{} นนุ กุสลํ อาวฏฺเฏนฺตสฺส อุปฺปชฺชติ ฯเปฯ ปณิทหนฺตสฺส อุปฺปชฺชตีติ, อามนฺตา.\csromanspacer{} หญฺจิ กุสลํ อาวฏฺเฏนฺตสฺส อุปฺปชฺชติ ฯเปฯ ปณิทหนฺตสฺส อุปฺปชฺชติ, โน จ วต เร วตฺตพฺเพ “อกุสลมูลํ ปฏิสนฺทหติ กุสลมูลนฺ”ติ.}

\proseitem{687}{\csromansymbolmark{*}อกุสลมูลํ ปฏิสนฺทหติ กุสลมูลนฺติ, อามนฺตา.\csromanspacer{} อกุสลมูลํ อโยนิโส มนสิกโรโต อุปฺปชฺชตีติ, อามนฺตา.\csromanspacer{} กุสลํ อโยนิโส มนสิกโรโต อุปฺปชฺชตีติ, น เหวํ วตฺตพฺเพ ฯเปฯ.\csromanspacer{} นนุ กุสลํ โยนิโส มนสิกโรโต อุปฺปชฺชตีติ, อามนฺตา.\csromanspacer{} หญฺจิ กุสลํ โยนิโส มนสิกโรโต อุปฺปชฺชติ, โน จ วต เร วตฺตพฺเพ “อกุสลมูลํ ปฏิสนฺทหติ กุสลมูลนฺ”ติ.}

\prose{\csromansymbolmark{*}อกุสลมูลํ ปฏิสนฺทหติ กุสลมูลนฺติ, อามนฺตา.\csromanspacer{} กามสญฺญาย อนนฺตรา เนกฺขมฺมสญฺญา อุปฺปชฺชติ, พฺยาปาทสญฺญาย อนนฺตรา อพฺยาปาทสญฺญา อุปฺปชฺชติ, วิหิํสาสญฺญาย อนนฺตรา อวิหิํสาสญฺญา อุปฺปชฺชติ, พฺยาปาทสฺส อนนฺตรา เมตฺตา อุปฺปชฺชติ, วิหิํสาย อนนฺตรา กรุณา อุปฺปชฺชติ, อรติยา อนนฺตรา มุทิตา อุปฺปชฺชติ, ปฏิฆสฺส อนนฺตรา อุเปกฺขา อุปฺปชฺชตีติ, น เหวํ วตฺตพฺเพ ฯเปฯ.}

\proseitem{688}{\csromansymbolmark{*}กุสลมูลํ ปฏิสนฺทหติ อกุสลมูลนฺติ, อามนฺตา.\csromanspacer{} ยา กุสลสฺส อุปฺปาทาย อาวฏฺฏนา ฯเปฯ ปณิธิ, สาว อกุสลสฺส อุปฺปาทาย อาวฏฺฏนา ฯเปฯ ปณิธีติ, น เหวํ วตฺตพฺเพ ฯเปฯ.}

\chapterheadpageread{14. จุทฺทสมวคฺค}
\prose{\csromanfolio{357}\csromansymbolmark{*}กุสลมูลํ ปฏิสนฺทหติ อกุสลมูลํ น วตฺตพฺพํ “ยา กุสลสฺส อุปฺปาทาย อาวฏฺฏนา ฯเปฯ ปณิธิ, สาว อกุสลสฺส อุปฺปาทาย อาวฏฺฏนา ฯเปฯ ปณิธี”ติ, อามนฺตา.\csromanspacer{} อกุสลํ อนาวฏฺเฏนฺตสฺส อุปฺปชฺชติ ฯเปฯ อปฺปณิทหนฺตสฺส อุปฺปชฺชตีติ, น เหวํ วตฺตพฺเพ ฯเปฯ.\csromanspacer{} นนุ อกุสลํ อาวฏฺเฏนฺตสฺส อุปฺปชฺชติ ฯเปฯ ปณิทหนฺตสฺส อุปฺปชฺชตีติ, อามนฺตา.\csromanspacer{} หญฺจิ อกุสลํ อาวฏฺเฏนฺตสฺส อุปฺปชฺชติ ฯเปฯ ปณิทหนฺตสฺส อุปฺปชฺชติ, โน จ วต เร วตฺตพฺเพ “กุสลมูลํ ปฏิสนฺทหติ อกุสลมูลนฺ”ติ.}

\proseitem{689}{\csromansymbolmark{*}กุสลมูลํ ปฏิสนฺทหติ อกุสลมูลนฺติ, อามนฺตา.\csromanspacer{} กุสลํ โยนิโส มนสิกโรโต อุปฺปชฺชตีติ, อามนฺตา.\csromanspacer{} อกุสลํ โยนิโส มนสิกโรโต อุปฺปชฺชตีติ, น เหวํ วตฺตพฺเพ ฯเปฯ.\csromanspacer{} นนุ อกุสลํ อโยนิโส มนสิกโรโต อุปฺปชฺชตีติ, อามนฺตา.\csromanspacer{} หญฺจิ อกุสลํ อโยนิโส มนสิกโรโต อุปฺปชฺชติ, โน จ วต เร วตฺตพฺเพ “กุสลมูลํ ปฏิสนฺทหติ อกุสลมูลนฺ”ติ.}

\prose{\csromansymbolmark{*}กุสลมูลํ ปฏิสนฺทหติ อกุสลมูลนฺติ, อามนฺตา.\csromanspacer{} เนกฺขมฺมสญฺญาย อนนฺตรา กามสญฺญา อุปฺปชฺชติ, อพฺยาปาทสญฺญาย อนนฺตรา พฺยาปาทสญฺญา อุปฺปชฺชติ, อวิหิํสาสญฺญาย อนนฺตรา วิหิํสาสญฺญา อุปฺปชฺชติ, เมตฺตาย อนนฺตรา พฺยาปาโท อุปฺปชฺชติ, กรุณาย อนนฺตรา วิหิํสา อุปฺปชฺชติ, มุทิตาย อนนฺตรา อรติ อุปฺปชฺชติ, อุเปกฺขาย อนนฺตรา ปฏิฆํ อุปฺปชฺชตีติ, น เหวํ วตฺตพฺเพ ฯเปฯ.}

\proseitem{690}{\csromansymbolmark{+}น วตฺตพฺพํ “อกุสลมูลํ ปฏิสนฺทหติ กุสลมูลํ, กุสลมูลํ ปฏิสนฺทหติ อกุสลมูลนฺ”ติ, อามนฺตา.\csromanspacer{} นนุ ยสฺมิญฺเญว วตฺถุสฺมิํ รชฺชติ, ตสฺมิญฺเญว วตฺถุสฺมิํ วิรชฺชติ.\csromanspacer{} ยสฺมิญฺเญว วตฺถุสฺมิํ วิรชฺชติ, ตสฺมิญฺเญว วตฺถุสฺมิํ รชฺชตีติ, อามนฺตา.\csromanspacer{} หญฺจิ ยสฺมิญฺเญว วตฺถุสฺมิํ รชฺชติ, ตสฺมิญฺเญว วตฺถุสฺมิํ วิรชฺชติ.\csromanspacer{} ยสฺมิญฺเญว วตฺถุสฺมิํ วิรชฺชติ, ตสฺมิญฺเญว วตฺถุสฺมิํ รชฺชติ, เตน วต เร วตฺตพฺเพ “อกุสลมูลํ ปฏิสนฺทหติ กุสลมูลํ, กุสลมูลํ ปฏิสนฺทหติ อกุสลมูลนฺ”ติ.}

\nitthitam{กุสลากุสลปฏิสนฺทหนกถา นิฏฺฐิตา.}
\chapterheadpageread{14. จุทฺทสมวคฺค}
\vspace{\tipitakaparskip}
\csromancenter{\textbf{(137) 2.}\csromanspacer{}\textbf{ สฬายตนุปฺปตฺติกถา}}
\csromanmatikaanchor{mtk.167}
\csromantocmarkref{subsection}{(137) 2. สฬายตนุปฺปตฺติกถา}{mtk.167}
\bookmark[dest={mtk.167},level=2]{(137) 2. สฬายตนุปฺปตฺติกถา}
\proseitem{691}{\csromanfolio{358}\csromansymbolmark{*}สฬายตนํ อปุพฺพํ อจริมํ มาตุกุจฺฉิสฺมิํ สณฺฐาตีติ, อามนฺตา.\csromanspacer{} สพฺพงฺคปจฺจงฺคี อหีนินฺทฺริโย มาตุกุจฺฉิสฺมิํ โอกฺกมตีติ, น เหวํ วตฺตพฺเพ ฯเปฯ.}

\prose{\csromansymbolmark{*}อุปปตฺเตสิเยน จิตฺเตน จกฺขายตนํ สณฺฐาตีติ, อามนฺตา.\csromanspacer{} อุปปตฺเตสิเยน จิตฺเตน หตฺถา สณฺฐนฺติ, ปาทา สณฺฐนฺติ, สีสํ สณฺฐาติ, กณฺโณ สณฺฐาติ, นาสิกา สณฺฐาติ, มุขํ สณฺฐาติ, ทนฺตา สณฺฐนฺตีติ, น เหวํ วตฺตพฺเพ ฯเปฯ.}

\prose{\csromansymbolmark{*}อุปปตฺเตสิเยน จิตฺเตน โสตายตนํ ฯเปฯ ฆานายตนํ ฯเปฯ ชิวฺหายตนํ สณฺฐาตีติ, อามนฺตา.\csromanspacer{} อุปปตฺเตสิเยน จิตฺเตน หตฺถา สณฺฐนฺติ, ปาทา สณฺฐนฺติ, สีสํ สณฺฐาติ, กณฺโณ สณฺฐาติ, นาสิกา สณฺฐาติ, มุขํ สณฺฐาติ, ทนฺตา สณฺฐนฺตีติ, น เหวํ วตฺตพฺเพ ฯเปฯ.}

\proseitem{692}{\csromansymbolmark{+}มาตุกุจฺฉิคตสฺส ปจฺฉา จกฺขายตนํ อุปฺปชฺชตีติ, อามนฺตา.\csromanspacer{} มาตุกุจฺฉิสฺมิํ จกฺขุปฏิลาภาย กมฺมํ กโรตีติ, น เหวํ วตฺตพฺเพ ฯเปฯ.\csromanspacer{} มาตุกุจฺฉิคตสฺส ปจฺฉา โสตายตนํ ฯเปฯ ฆานายตนํ ฯเปฯ ชิวฺหายตนํ อุปฺปชฺชตีติ, อามนฺตา.\csromanspacer{} มาตุกุจฺฉิสฺมิํ ชิวฺหาปฏิลาภาย กมฺมํ กโรตีติ, น เหวํ วตฺตพฺเพ ฯเปฯ.}

\prose{\csromansymbolmark{*}มาตุกุจฺฉิคตสฺสปจฺฉา เกสา, โลมา, นขา, ทนฺตา, อฏฺฐี อุปฺปชฺชนฺตีติ, อามนฺตา.\csromanspacer{} มาตุกุจฺฉิสฺมิํ อฏฺฐิปฏิลาภาย กมฺมํ กโรตีติ, น เหวํ วตฺตพฺเพ ฯเปฯ.}

\prose{\csromansymbolmark{*}น วตฺตพฺพํ “มาตุกุจฺฉิคตสฺส ปจฺฉา เกสา, โลมา นขา, ทนฺตา, อฏฺฐี อุปฺปชฺชนฺตี”ติ, อามนฺตา.\csromanspacer{} นนุ วุตฺตํ ภควตา\csromandash{}}

\csromangathameasure{“ปฐมํ กลลํ โหติ, กลลา โหติ อพฺพุทํ. \\ อพฺพุทา ชายเต เปสิ, เปสิ นิพฺพตฺตเต ฆโน. \\ ฆนา ปสาขา ชายนฺติ, เกสา โลมา นขาปิ จ.}
\csromangathagroup{\csromangathastanza{“ปฐมํ กลลํ โหติ, กลลา โหติ อพฺพุทํ. \\ อพฺพุทา ชายเต เปสิ\footnote{เปสี (สฺยา, กํ, อิ)}, เปสิ นิพฺพตฺตเต\footnote{นิพฺพตฺตตี (สี, สฺยา, สํ 1. 208 ปิฏฺเฐ), นิพฺพตฺตติ (อิ, ก)} ฆโน. \\ ฆนา ปสาขา ชายนฺติ, เกสา โลมา นขาปิ จ.}}

\chapterheadpageread{14. จุทฺทสมวคฺค}
\csromanmatikaanchor{mtk.168}
\csromantocmarkref{subsection}{(138) 3. อนนฺตรปจฺจยกถา}{mtk.168}
\bookmark[dest={mtk.168},level=2]{(138) 3. อนนฺตรปจฺจยกถา}
\csromangathameasure{ยญฺจสฺส ภุญฺชติ มาตา, อนฺนํ ปานญฺจ โภชนํ. \\ เตน โส ตตฺถ ยาเปติ, มาตุกุจฺฉิคโต นโร”ติ.}
\csromangathagroup{\csromangathastanza{\csromanfolio{359}ยญฺจสฺส ภุญฺชติ มาตา, อนฺนํ ปานญฺจ โภชนํ. \\ เตน โส ตตฺถ ยาเปติ, มาตุกุจฺฉิคโต นโร”ติ\footnote{สํ 1. 208 ปิฏฺเฐ.}.}}

\prose{อตฺเถว สุตฺตนฺโตติ, อามนฺตา.\csromanspacer{} เตน หิ มาตุกุจฺฉิคตสฺส ปจฺฉา เกสา, โลมา, นขา, ทนฺตา, อฏฺฐี อุปฺปชฺชนฺตีติ.}

\nitthitam{สฬายตนุปฺปตฺติกถา นิฏฺฐิตา.}
\chapterhead{14. จุทฺทสมวคฺค}
\vspace{\tipitakaparskip}
\csromancenter{\textbf{(138) 3.}\csromanspacer{}\textbf{ อนนฺตรปจฺจยกถา}}
\proseitem{693}{\csromansymbolmark{*}จกฺขุวิญฺญาณสฺส อนนฺตรา โสตวิญฺญาณํ อุปฺปชฺชตีติ, อามนฺตา.\csromanspacer{} ยา จกฺขุวิญฺญาณสฺส อุปฺปาทาย อาวฏฺฏนา ฯเปฯ ปณิธิ, สาว โสตวิญฺญาณสฺส อุปฺปาทาย อาวฏฺฏนา ฯเปฯ ปณิธีติ, น เหวํ วตฺตพฺเพ ฯเปฯ.}

\prose{\csromansymbolmark{*}จกฺขุวิญฺญาณสฺส อนนฺตรา โสตวิญฺญาณํ อุปฺปชฺชติ, น วตฺตพฺพํ “ยา จกฺขุวิญฺญาณสฺส อุปฺปาทาย อาวฏฺฏนา ฯเปฯ ปณิธิ, สาว โสตวิญฺญาณสฺส อุปฺปทาย อาวฏฺฏนา ฯเปฯ ปณิธี”ติ, อามนฺตา.\csromanspacer{} โสตวิญฺญาณํ อนาวฏฺเฏนฺตสฺส อุปฺปชฺชติ ฯเปฯ อปฺปณิทหนฺตสฺส อุปฺปชฺชตีติ, น เหวํ วตฺตพฺเพ ฯเปฯ.\csromanspacer{} นนุ โสตวิญฺญาณํ อาวฏฺเฏนฺตสฺส อุปฺปชฺชติ ฯเปฯ ปณิทหนฺตสฺส อุปฺปชฺชตีติ, อามนฺตา.\csromanspacer{} หญฺจิ โสตวิญฺญาณํ อาวฏฺเฏนฺตสฺส อุปฺปชฺชติ ฯเปฯ ปณิทหนฺตสฺส อุปฺปชฺชติ, โน จ วต เร วตฺตพฺเพ “จกฺขุวิญฺญาณสฺส อนนฺตรา โสตวิญฺญาณํ อุปฺปชฺชตี”ติ.}

\proseitem{694}{\csromansymbolmark{*}จกฺขุวิญฺญาณสฺส อนนฺตรา โสตวิญฺญาณํ อุปฺปชฺชตีติ, อามนฺตา.\csromanspacer{} จกฺขุวิญฺญาณํ รูปนิมิตฺตํ มนสิกโรโต อุปฺปชฺชตีติ, อามนฺตา.\csromanspacer{} โสตวิญฺญาณํ รูปนิมิตฺตํ มนสิกโรโต อุปฺปชฺชตีติ, น เหวํ วตฺตพฺเพ ฯเปฯ.}

\prose{\csromansymbolmark{*}จกฺขุวิญฺญาณํ รูปารมฺมณญฺเญว น อญฺญารมฺมณนฺติ, อามนฺตา.\csromanspacer{} โสตวิญฺญาณํ รูปารมฺมณญฺเญว น อญฺญารมฺมณนฺติ, น เหวํ วตฺตพฺเพ ฯเปฯ.}

\prose{\csromansymbolmark{*}จกฺขุญฺจ ปฏิจฺจ รูเป จ อุปฺปชฺชติ จกฺขุวิญฺญาณนฺติ, อามนฺตา.\csromanspacer{} จกฺขุญฺจ ปฏิจฺจ รูเป จ อุปฺปชฺชติ โสตวิญฺญาณนฺติ, น เหวํ วตฺตพฺเพ ฯเปฯ.}

\prose{\csromanfolio{360}\csromansymbolmark{*}จกฺขุญฺจ ปฏิจฺจ รูเป จ อุปฺปชฺชติ โสตวิญฺญาณนฺติ, อามนฺตา. “จกฺขุญฺจ ปฏิจฺจ รูเป จ อุปฺปชฺชติ โสตวิญฺญาณนฺ”ติ อตฺเถว สุตฺตนฺโตติ, นตฺถิ. “จกฺขุญฺจ ปฏิจฺจ รูเป จ อุปฺปชฺชติ จกฺขุวิญฺญาณนฺ”ติ อตฺเถว สุตฺตนฺโตติ, อามนฺตา.\csromanspacer{} หญฺจิ “จกฺขุญฺจ ปฏิจฺจ รูเป จ อุปฺปชฺชติ จกฺขุวิญฺญาณนฺ”ติ อตฺเถว สุตฺตนฺโต, โน จ วต เร วตฺตพฺเพ “จกฺขุญฺจ ปฏิจฺจ รูเป จ อุปฺปชฺชติ โสตวิญฺญาณนฺ”ติ.}

\prose{\csromansymbolmark{*}จกฺขุวิญฺญาณสฺส อนนฺตรา โสตวิญฺญาณํ อุปฺปชฺชตีติ, อามนฺตา.\csromanspacer{} ตญฺเญว จกฺขุวิญฺญาณํ ตํ โสตวิญฺญาณนฺติ, น เหวํ วตฺตพฺเพ ฯเปฯ.}

\proseitem{695}{\csromansymbolmark{*}โสตวิญฺญาณสฺส อนนฺตรา ฆานวิญฺญาณํ อุปฺปชฺชติ ฯเปฯ.\csromanspacer{} ฆานวิญฺญาณสฺส อนนฺตรา ชิวฺหาวิญฺญาณํ อุปฺปชฺชติ ฯเปฯ.\csromanspacer{} ชิวฺหาวิญฺญาณสฺส อนนฺตรา กายวิญฺญาณํ อุปฺปชฺชตีติ, ณมนฺตา.\csromanspacer{} ยา ชิวฺหาวิญฺญาณสฺส อุปฺปาทาย อาวฏฺฏนา ฯเปฯ ปณิธิ, สาว กายวิญฺญาณสฺส อุปฺปาทาย อาวฏฺฏนา ฯเปฯ ปณิธีติ, น เหวํ วตฺตพฺเพ ฯเปฯ.\csromanspacer{} ชิวฺหาวิญฺญาณสฺส อนนฺตรา กายวิญฺญาณํ อุปฺปชฺชติ, น วตฺตพฺพํ “ยา ชิวฺหาวิญฺญาณสฺส อุปฺปาทาย อาวฏฺฏนา ฯเปฯ ปณิธิ, สาว กายวิญฺญาณสฺส อุปฺปาทาย อาวฏฺฏนา ฯเปฯ ปณิธี”ติ, อามนฺตา.\csromanspacer{} กายวิญฺญาณํ อนาวฏฺเฏนฺตสฺส อุปฺปชฺชติ ฯเปฯ อปฺปณิทหนฺตสฺส อุปฺปชฺชตีติ, น เหวํ วตฺตพฺเพ ฯเปฯ.\csromanspacer{} นนุ กายวิญฺญาณํ อาวฏฺเฏนฺตสฺส อุปฺปชฺชติ ฯเปฯ ปณิทหนฺตสฺส อุปฺปชฺชตีติ, อามนฺตา.\csromanspacer{} หญฺจิ กายวิญฺญาณํ อาวฏฺเฏนฺตสฺส อุปฺปชฺชติ ฯเปฯ ปณิทหนฺตสฺส อุปฺปชฺชติ, โน จ วต เร วตฺตพฺเพ “ชิวฺหาวิญฺญาณสฺส อนนฺตรา กายวิญฺญาณํ อุปฺปชฺชตี”ติ.}

\proseitem{696}{\csromansymbolmark{*}ชิวฺหาวิญฺญาณสฺส อนนฺตรา กายวิญฺญาณํ อุปฺปชฺชตีติ, อามนฺตา.\csromanspacer{} ชิวฺหาวิญฺญาณํ รสนิมิตฺตํ มนสิกโรโต อุปฺปชฺชตีติ, อามนฺตา.\csromanspacer{} กายวิญฺญาณํ รสนิมิตฺตํ มนสิกโรโต อุปฺปชฺชตีติ, น เหวํ วตฺตพฺเพ ฯเปฯ.}

\prose{\csromansymbolmark{*}ชิวฺหาวิญฺญาณํ รสารมฺมณญฺเญว น อญฺญารมฺมณนฺติ, อามนฺตา.\csromanspacer{} กายวิญฺญาณํ รสารมฺมณญฺเญว น อญฺญารมฺมณนฺติ, น เหวํ วตฺตพฺเพ ฯเปฯ.}

\prose{\csromansymbolmark{*}ชิวฺหญฺจ ปฏิจฺจ รเส จ อุปฺปชฺชติ ชิวฺหาวิญฺญาณนฺติ, อามนฺตา.\csromanspacer{} ชิวฺหญฺจ ปฏิจฺจ รเส จ อุปฺปชฺชติ กายวิญฺญาณนฺติ, น เหวํ วตฺตพฺเพ ฯเปฯ.}

\prose{\csromansymbolmark{*}ชิวฺหญฺจ ปฏิจฺจ รเส จ อุปฺปชฺชติ กายวิญฺญาณนฺติ, อามนฺตา. “ชิวฺหญฺจ ปฏิจฺจ รเส จ อุปฺปชฺชติ กายวิญฺญาณนฺ”ติ อตฺเถว สุตฺตนฺโตติ, นตฺถิ.}

\vspace{\tipitakaparskip}
\csromancenter{จุทฺทสมวคฺค “ชิวฺหญฺจ ปฏิจฺจ รเส จ อุปฺปชฺชติ ชิวฺหาวิญฺญาณนฺ”ติ อตฺเถว สุตฺตนฺโตติ, อามนฺตา.\csromanspacer{} หญฺจิ “ชิวฺหญฺจ ปฏิจฺจ รเส จ อุปฺปชฺชติ ชิวฺหาวิญฺญาณนฺ”ติ อตฺเถว สุตฺตนฺโตติ, โน จ วต เร วตฺตพฺเพ “ชิวฺหญฺจ ปฏิจฺจ รเส จ อุปฺปชฺชติ กายวิญฺญาณนฺ”ติ.}
\csromanmatikaanchor{mtk.169}
\csromantocmarkref{subsection}{(139) 4. อริยรูปกถา}{mtk.169}
\bookmark[dest={mtk.169},level=2]{(139) 4. อริยรูปกถา}
\prose{\csromanfolio{361}\csromansymbolmark{*}ชิวฺหาวิญฺญาณสฺส อนนฺตรา กายวิญฺญาณํ อุปฺปชฺชตีติ, อามนฺตา.\csromanspacer{} ตญฺเญว ชิวฺหาวิญฺญาณํ ตํ กายวิญฺญาณนฺติ, น เหวํ วตฺตพฺเพ ฯเปฯ.}

\proseitem{697}{\csromansymbolmark{+}น วตฺตพฺพํ “ปญฺจวิญฺญาณา อญฺญมญฺญสฺส สมนนฺตรา อุปฺปชฺชนฺตี”ติ, อามนฺตา.\csromanspacer{} นนุ อตฺถิ โกจิ นจฺจติ, คายติ, วาเทติ, รูปญฺจ ปสฺสติ, สทฺทญฺจ สุณาติ, คนฺธญฺจ ฆายติ, รสญฺจ สายติ, โผฏฺฐพฺพญฺจ ผุสตีติ, อามนฺตา.\csromanspacer{} หญฺจิ อตฺถิ โกจิ นจฺจติ, คายติ, วาเทติ, รูปญฺจ ปสฺสติ, สทฺทญฺจ สุณาติ, คนฺธญฺจ ฆายติ, รสญฺจ สายติ, โผฏฺฐพฺพญฺจ ผุสติ, เตน วต เร วตฺตพฺเพ “ปญฺจวิญฺญาณา อญฺญมญฺญสฺส สมนนฺตรา อุปฺปชฺชนฺตี”ติ.}

\nitthitam{อนนฺตรปจฺจยกถา นิฏฺฐิตา.}
\chapterhead{14. จุทฺทสมวคฺค}
\vspace{\tipitakaparskip}
\csromancenter{\textbf{(139) 4.}\csromanspacer{}\textbf{ อริยรูปกถา}}
\proseitem{698}{\csromansymbolmark{*}อริยรูปํ มหาภูตานํ อุปาทายาติ, อามนฺตา.\csromanspacer{} อริยรูปํ กุสลนฺติ, อามนฺตา.\csromanspacer{} มหาภูตา กุสลาติ, น เหวํ วตฺตพฺเพ ฯเปฯ.\csromanspacer{} มหาภูตา อพฺยากตาติ, อามนฺตา.\csromanspacer{} อริยรูปํ อพฺยากตนฺติ, น เหวํ วตฺตพฺเพ ฯเปฯ.\csromanspacer{} อริยรูปํ มหาภูตานํ อุปาทายาติ, อามนฺตา.\csromanspacer{} อริยรูปํ อนาสวํ อสํโยชนิยํ อคนฺถนิยํ อโนฆนิยํ อโยคนิยํ อนีวรณิยํ อปรามฏฺฐํ อนุปาทานิยํ อสํกิเลสิยนฺติ, อามนฺตา.\csromanspacer{} มหาภูตา อนาสวา ฯเปฯ อสํกิเลสิยาติ, น เหวํ วตฺตพฺเพ ฯเปฯ.\csromanspacer{} มหาภูตา สาสวา สํโยชนิยา ฯเปฯ สํกิเลสิยาติ, อามนฺตา.\csromanspacer{} อริยรูปํ สาสวํ สํโยชนิยํ ฯเปฯ สํกิเลสิยนฺติ, น เหวํ วตฺตพฺเพ ฯเปฯ.}

\csromanmatikaanchor{mtk.170}
\csromantocmarkref{subsection}{(140) 5. อญฺโญ-อนุสโยติกถา}{mtk.170}
\bookmark[dest={mtk.170},level=2]{(140) 5. อญฺโญ-อนุสโยติกถา}
\proseitem{699}{\csromansymbolmark{+}น วตฺตพฺพํ “อริยรูปํ มหาภูตานํ อุปาทายา”ติ, อามนฺตา.\csromanspacer{} นนุ วุตฺตํ ภควตา “ยํ กิญฺจิ ภิกฺขเว รูปํ จตฺตาริ มหาภูตานิ จตุนฺนญฺจ \csromanfolio{362}มหาภูตานํ อุปาทายรูปนฺ”ติ\footnote{ม 1. 282 ปิฏฺเฐ; อํ 3. 543 ปิฏฺเฐ จ.} อตฺเถว สุตฺตนฺโตติ, อามนฺตา.\csromanspacer{} เตน หิ อริยรูปํ มหาภูตานํ อุปาทายาติ.}

\nitthitam{อริยรูปกถา นิฏฺฐิตา.}
\chapterhead{14. จุทฺทสมวคฺค}
\vspace{\tipitakaparskip}
\csromancenter{\textbf{(140) 5.}\csromanspacer{}\textbf{ อญฺโญ-อนุสโยติกถา}}
\proseitem{700}{\csromansymbolmark{*}อญฺโญ กามราคานุสโย อญฺญํ กามราคปริยุฏฺฐานนฺติ, อามนฺตา.\csromanspacer{} อญฺโญ กามราโค อญฺญํ กามราคปริยุฏฺฐานนฺติ, น เหวํ วตฺตพฺเพ ฯเปฯ.\csromanspacer{} เสฺวว กามราโค ตํ กามราคปริยุฏฺฐานนฺติ, อามนฺตา.\csromanspacer{} เสฺวว กามราคานุสโย ตํ กามราคปริยุฏฺฐานนฺติ, น เหวํ วตฺตพฺเพ ฯเปฯ.}

\prose{\csromansymbolmark{*}อญฺโญ ปฏิฆานุสโย อญฺญํ ปฏิฆปริยุฏฺฐานนฺติ, อามนฺตา.\csromanspacer{} อญฺญํ ปฏิฆํ อญฺญํ ปฏิฆปริยุฏฺฐานนฺติ, น เหวํ วตฺตพฺเพ ฯเปฯ.\csromanspacer{} ตญฺเญว ปฏิฆํ ตํ ปฏิฆปริยุฏฺฐานนฺติ, อามนฺตา.\csromanspacer{} เสฺวว ปฏิฆานุสโย ตํ ปฏิฆปริยุฏฺฐานนฺติ, น เหวํ วตฺตพฺเพ ฯเปฯ.}

\prose{\csromansymbolmark{*}อญฺโญ มานานุสโย อญฺญํ มานปริยุฏฺฐานนฺติ, อามนฺตา.\csromanspacer{} อญฺโญ มาโน อญฺญํ มานปริยุฏฺฐานนฺติ, น เหวํ วตฺตพฺเพ ฯเปฯ.\csromanspacer{} เสฺวว มาโน ตํ มานปริยุฏฺฐานนฺติ, อามนฺตา.\csromanspacer{} เสฺวว มานานุสโย ตํ มานปริยุฏฺฐานนฺติ, น เหวํ วตฺตพฺเพ ฯเปฯ.}

\prose{\csromansymbolmark{*}อญฺโญ ทิฏฺฐานุสโย อญฺญํ ทิฏฺฐิปริยุฏฺฐานนฺติ, อามนฺตา.\csromanspacer{} อญฺญา ทิฏฺฐิ อญฺญํ ทิฏฺฐิปริยุฏฺฐานนฺติ, น เหวํ วตฺตพฺเพ ฯเปฯ.\csromanspacer{} สาว ทิฏฺฐิ ตํ ทิฏฺฐิปริยุฏฺฐานนฺติ, อามนฺตา.\csromanspacer{} เสฺวว ทิฏฺฐานุสโย ตํ ทิฏฺฐิปริยุฏฺฐานนฺติ, น เหวํ วตฺตพฺเพ ฯเปฯ.}

\prose{\csromansymbolmark{*}อญฺโญ วิจิกิจฺฉานุสโย อญฺญํ วิจิกิจฺฉาปริยุฏฺฐานนฺติ, อามนฺตา.\csromanspacer{} อญฺญา วิจิกิจฺฉา อญฺญํ วิจิกิจฺฉาปริยุฏฺฐานนฺติ, น เหวํ วตฺตพฺเพ ฯเปฯ.\csromanspacer{} สาว วิจิกิจฺฉา ตํ วิจิกิจฺฉาปริยุฏฺฐานนฺติ, อามนฺตา.\csromanspacer{} เสฺวว วิจิกิจฺฉานุสโย ตํ วิจิกิจฺจาปริยุฏฺฐานนฺติ, น เหวํ วตฺตพฺเพ ฯเปฯ.}

\chapterheadpageread{14. จุทฺทสมวคฺค}
\csromanmatikaanchor{mtk.171}
\csromantocmarkref{subsection}{(141) 6. ปริยุฏฺฐานํจิตฺตวิปฺปยุตฺตนฺติกถา}{mtk.171}
\bookmark[dest={mtk.171},level=2]{(141) 6. ปริยุฏฺฐานํจิตฺตวิปฺปยุตฺตนฺติกถา}
\prose{\csromanfolio{363}\csromansymbolmark{*}อญฺโญ ภวราคานุสโย อญฺญํ ภวราคปริยุฏฺฐานนฺติ, อามนฺตา.\csromanspacer{} อญฺโญ ภวราโค อญฺญํ ภวราคปริยุฏฺฐานนฺติ, น เหวํ วตฺตพฺเพ ฯเปฯ.\csromanspacer{} เสฺวว ภวราโค ตํ ภวราคปริยุฏฺฐานนฺติ, อามนฺตา.\csromanspacer{} เสฺวว ภวราคานุสโย ตํ ภวราคปริยุฏฺฐานนฺติ, น เหวํ วตฺตพฺเพ ฯเปฯ.}

\prose{\csromansymbolmark{*}อญฺโญ อวิชฺชานุสโย อญฺญํ อวิชฺชาปริยุฏฺฐานนฺติ, อามนฺตา.\csromanspacer{} อญฺญา อวิชฺชา อญฺญํ อวิชฺชาปริยุฏฺฐานนฺติ, น เหวํ วตฺตพฺเพ ฯเปฯ.\csromanspacer{} สาว อวิชฺชา ตํ อวิชฺชาปริยุฏฺฐานนฺติ, อามนฺตา.\csromanspacer{} เสฺวว อวิชฺชานุสโย ตํ อวิชฺชาปริยุฏฺฐานนฺติ, น เหวํ วตฺตพฺเพ ฯเปฯ.}

\proseitem{701}{\csromansymbolmark{+}น วตฺตพฺพํ “อญฺโญ อนุสโย อญฺญํ ปริยุฏฺฐานนฺ”ติ, อามนฺตา.\csromanspacer{} ปุถุชฺชโน กุสลาพฺยากเต จิตฺเต วตฺตมาเน “สานุสโย”ติ วตฺตพฺโพติ, อามนฺตา. “ปริยุฏฺฐิโต”ติ วตฺตพฺโพติ, น เหวํ วตฺตพฺเพ.\csromanspacer{} เตน หิ อญฺโญ อนุสโย อญฺญํ ปริยุฏฺฐานนฺติ.\csromanspacer{} ปุถุชฺชโน กุสลาพฺยากเต จิตฺเต วตฺตมาเน “สราโค”ติ วตฺตพฺโพติ, อามนฺตา. “ปริยุฏฺฐิโต”ติ วตฺตพฺโพติ, น เหวํ วตฺตพฺเพ.\csromanspacer{} เตน หิ อญฺโญ ราโค อญฺญํ ปริยุฏฺฐานนฺติ.}

\nitthitam{อญฺโญ อนุสโยติกถา นิฏฺฐิตา.}
\chapterhead{14. จุทฺทสมวคฺค}
\vspace{\tipitakaparskip}
\csromancenter{\textbf{(141) 6.}\csromanspacer{}\textbf{ ปริยุฏฺฐานํจิตฺตวิปฺปยุตฺตนฺติกถา}}
\proseitem{702}{\csromansymbolmark{*}ปริยุฏฺฐานํ จิตฺตวิปฺปยุตฺตนฺติ, อามนฺตา.\csromanspacer{} รูปํ นิพฺพานํ จกฺขายตนํ ฯเปฯ โผฏฺฐพฺพายตนนฺติ, น เหวํ วตฺตพฺเพ ฯเปฯ.\csromanspacer{} ปริยุฏฺฐานํ จิตฺตวิปฺปยุตฺตนฺติ, อามนฺตา.\csromanspacer{} นตฺถิ สราคํ จิตฺตํ.\csromanspacer{} สโทสํ จิตฺตํ.\csromanspacer{} สโมหํ จิตฺตํ ฯเปฯ.\csromanspacer{} อกุสลํ จิตฺตํ สํกิลิฏฺฐํ จิตฺตนฺติ, น เหวํ วตฺตพฺเพ ฯเปฯ.\csromanspacer{} นนุ อตฺถิ สราคํ จิตฺตํ.\csromanspacer{} สโทสํ จิตฺตํ.\csromanspacer{} สโมหํ จิตฺตํ ฯเปฯ.\csromanspacer{} อกุสลํ จิตฺตํ สํกิลิฏฺฐํ จิตฺตนฺติ, อามนฺตา.\csromanspacer{} หญฺจิ อตฺถิ สราคํ จิตฺตํ.\csromanspacer{} สโทสํ จิตฺตํ.\csromanspacer{} สโมหํ จิตฺตํ ฯเปฯ.\csromanspacer{} อกุสลํ จิตฺตํ สํกิลิฏฺฐํ จิตฺตํ, โน จ วต เร วตฺตพฺเพ “ปริยุฏฺฐานํ จิตฺตวิปฺปยุตฺตนฺ”ติ.}

\nitthitam{ปริยุฏฺฐานํ จิตฺตวิปฺปยุตฺตนฺติกถา นิฏฺฐิตา.}
\chapterheadpageread{14. จุทฺทสมวคฺค}
\vspace{\tipitakaparskip}
\csromancenter{\textbf{(142) 7.}\csromanspacer{}\textbf{ ปริยาปนฺนกถา}}
\csromanmatikaanchor{mtk.172}
\csromantocmarkref{subsection}{(142) 7. ปริยาปนฺนกถา}{mtk.172}
\bookmark[dest={mtk.172},level=2]{(142) 7. ปริยาปนฺนกถา}
\proseitem{703}{\csromanfolio{364}\csromansymbolmark{*}รูปราโค รูปธาตุํ อนุเสติ รูปธาตุปริยาปนฺโนติ, อามนฺตา.\csromanspacer{} สมาปตฺเตสิโย อุปปตฺเตสิโย ทิฏฺฐธมฺมสุขวิหาโร, สมาปตฺเตสิเยน จิตฺเตน อุปปตฺเตสิเยน จิตฺเตน ทิฏฺฐธมฺมสุขวิหาเรน จิตฺเตน สหคโต สหชาโต สํสฏฺโฐ สมฺปยุตฺโต เอกุปฺปาโท เอกนิโรโธ เอกวตฺถุโก เอการมฺมโณติ, น เหวํ วตฺตพฺเพ ฯเปฯ.\csromanspacer{} นนุ น สมาปตฺเตสิโย น อุปปตฺเตสิโย น ทิฏฺฐธมฺมสุขวิหาโร, น สมาปตฺเตสิเยน จิตฺเตน น อุปปตฺเตสิเยน จิตฺเตน น ทิฏฺฐธมฺมสุขวิหาเรน จิตฺเตน สหคโต สหชาโต สํสฏฺโฐ สมฺปยุตฺโต เอกุปฺปาโท เอกนิโรโธ เอกวตฺถุโก เอการมฺมโณติ, อามนฺตา.\csromanspacer{} หญฺจิ น สมาปตฺเตสิโย น อุปปตฺเตสิโย น ทิฏฺฐธมฺมสุขวิหาโร ฯเปฯ เอการมฺมโณ, โน จ วต เร วตฺตพฺเพ “รูปราโค รูปธาตุํ อนุเสติ รูปธาตุปริยาปนฺโน”ติ ฯเปฯ.}

\prose{\csromansymbolmark{*}รูปราโค รูปธาตุํ อนุเสติ รูปธาตุปริยาปนฺโนติ, อามนฺตา.\csromanspacer{} สทฺทราโค สทฺทธาตุํ อนุเสติ สทฺทธาตุปริยาปนฺโนติ, น เหวํ วตฺตพฺเพ ฯเปฯ.\csromanspacer{} รูปราโค รูปธาตุํ อนุเสติ รูปธาตุปริยาปนฺโนติ, อามนฺตา.\csromanspacer{} คนฺธราโค ฯเปฯ.\csromanspacer{} รสราโค ฯเปฯ.\csromanspacer{} โผฏฺฐพฺพราโค โผฏฺฐพฺพธาตุํ อนุเสติ โผฏฺฐพฺพธาตุปริยาปนฺโนติ, น เหวํ วตฺตพฺเพ ฯเปฯ.}

\prose{\csromansymbolmark{*}สทฺทราโค สทฺทธาตุํ อนุเสติ, น วตฺตพฺพํ “สทฺทธาตุปริยาปนฺโน”ติ, อามนฺตา.\csromanspacer{} รูปราโค รูปธาตุํ อนุเสติ, น วตฺตพฺพํ “รูปธาตุปริยาปนฺโน”ติ, น เหวํ วตฺตพฺเพ ฯเปฯ.\csromanspacer{} คนฺธราโค ฯเปฯ.\csromanspacer{} รสราโค ฯเปฯ.\csromanspacer{} โผฏฺฐพฺพราโค โผฏฺฐพฺพธาตุํ อนุเสติ, น วตฺตพฺพํ “โผฏฺฐพฺพธาตุปริยาปนฺโน”ติ, อามนฺตา.\csromanspacer{} รูปราโค รูปธาตุํ อนุเสติ, น วตฺตพฺพํ “รูปธาตุปริยาปนฺโน”ติ, น เหวํ วตฺตพฺเพ ฯเปฯ.}

\proseitem{704}{\csromansymbolmark{*}อรูปราโค อรูปธาตุํ อนุเสติ, อรูปธาตุปริยาปนฺโนติ, อามนฺตา.\csromanspacer{} สมาปตฺเตสิโย อุปปตฺเตสิโย ทิฏฺฐธมฺมสุขวิหาโร, สมาปฺปตฺเตสิเยน จิตฺเตน อุปปตฺเตสิเยน จิตฺเตน ทิฏฺฐธมฺมสุขวิหาเรน จิตฺเตน สหคโต สหชาโต สํสฏฺโฐ สมฺปยุตฺโต เอกุปฺปาโท เอกนิโรโธ}

\vspace{\tipitakaparskip}
\csromancenter{จุทฺทสมวคฺค เอกวตฺถุโก เอการมฺมโณติ, น เหวํ วตฺตพฺเพ ฯเปฯ.\csromanspacer{} นนุ น สมาปตฺเตสิโย น อุปปตฺเตสิโย น ทิฏฺฐธมฺมสุขวิหาโร, น สมาปตฺเตสิเยน จิตฺเตน ฯเปฯ เอการมฺมโณติ, อามนฺตา.\csromanspacer{} หญฺจิ น สมาปตฺเตสิโย น อุปปตฺเตสิโย ฯเปฯ เอการมฺมโณ, โน จ วต เร วตฺตพฺเพ “อรูปราโค อรูปธาตุํ อนุเสติ อรูปธาตุปริยาปนฺโน”ติ.}
\csromanmatikaanchor{mtk.173}
\csromantocmarkref{subsection}{(143) 8. อพฺยากตกถา}{mtk.173}
\bookmark[dest={mtk.173},level=2]{(143) 8. อพฺยากตกถา}
\prose{\csromanfolio{365}\csromansymbolmark{*}อรูปราโค อรูปธาตุํ อนุเสติ อรูปธาตุปริยาปนฺโนติ, อามนฺตา.\csromanspacer{} สทฺทราโค สทฺทธาตุํ อนุเสติ สทฺทธาตุปริยาปนฺโนติ, น เหวํ วตฺตพฺเพ ฯเปฯ.\csromanspacer{} อรูปราโค อรูปธาตุํ อนุเสติ อรูปธาตุปริยาปนฺโนติ, อามนฺตา.\csromanspacer{} คนฺธราโค ฯเปฯ.\csromanspacer{} รสราโค ฯเปฯ.\csromanspacer{} โผฏฺฐพฺพราโค โผฏฺฐพฺพธาตุํ อนุเสติ โผฏฺฐพฺพธาตุปริยาปนฺโนติ, น เหวํ วตฺตพฺเพ ฯเปฯ. \csromansymbolmark{*}สทฺทราโค สทฺทธาตุํ อนุเสติ, น วตฺตพฺพํ “สทฺทธาตุปริยาปนฺโน”ติ, อามนฺตา.\csromanspacer{} อรูปราโค อรูปธาตุํ อนุเสติ, น วตฺตพฺพํ “อรูปธาตุปริยาปนฺโน”ติ, น เหวํ วตฺตพฺเพ ฯเปฯ.\csromanspacer{} คนฺธราโค ฯเปฯ.\csromanspacer{} รสราโค ฯเปฯ.\csromanspacer{} โผฏฺฐพฺพราโค โผฏฺฐพฺพธาตุํ อนุเสติ, น วตฺตพฺพํ “โผฏฺฐพฺพธาตุปริยาปนฺโน”ติ, อามนฺตา.\csromanspacer{} อรูปราโค อรูปธาตุํ อนุเสติ, น วตฺตพฺพํ “อรูปธาตุปริยาปนฺโน”ติ, น เหวํ วตฺตพฺเพ ฯเปฯ.}

\proseitem{705}{\csromansymbolmark{+}น วตฺตพฺพํ “รูปราโค รูปธาตุํ อนุเสติ รูปธาตุปริยาปนฺโน, อรูปราโค อรูปธาตุํ อนุเสติ อรูปธาตุปริยาปนฺโน”ติ, อามนฺตา.\csromanspacer{} นนุ กามราโค กามธาตุํ อนุเสติ กามธาตุปริยาปนฺโนติ, อามนฺตา.\csromanspacer{} หญฺจิ กามราโค กามธาตุํ อนุเสติ กามธาตุปริยาปนฺโน, เตน วต เร วตฺตพฺเพ “รูปราโค รูปธาตุํ อนุเสติ รูปธาตุปริยาปนฺโน, อรูปราโค อรูปธาตุํ อนุเสติ อรูปธาตุปริยาปนฺโน”ติ.}

\nitthitam{ปริยาปนฺนกถา นิฏฺฐิตา.}
\chapterhead{14. จุทฺทสมวคฺค}
\vspace{\tipitakaparskip}
\csromancenter{\textbf{(143) 8.}\csromanspacer{}\textbf{ อพฺยากตกถา}}
\proseitem{706}{\csromansymbolmark{*}ทิฏฺฐิคตํ อพฺยากตนฺติ, อามนฺตา.\csromanspacer{} วิปากาพฺยากตํ กิริยาพฺยากตํ รูปํ นิพฺพานํ จกฺขายตนํ ฯเปฯ โผฏฺฐพฺพายตนนฺติ, น เหวํ วตฺตพฺเพ ฯเปฯ. \csromanfolio{366}ทิฏฺฐิคตํ อพฺยากตนฺติ, อามนฺตา.\csromanspacer{} ทิฏฺฐิคตสมฺปยุตฺโต ผสฺโส อพฺยากโตติ, น เหวํ วตฺตพฺเพ ฯเปฯ.\csromanspacer{} ทิฏฺฐิคตํ อพฺยากตนฺติ, อามนฺตา.\csromanspacer{} ทิฏฺฐิคตสมฺปยุตฺตา เวทนา ฯเปฯ สญฺญา ฯเปฯ เจตนา ฯเปฯ จิตฺตํ อพฺยากตนฺติ, น เหวํ วตฺตพฺเพ ฯเปฯ.}

\prose{\csromansymbolmark{*}ทิฏฺฐิคตสมฺปยุตฺโต ผสฺโส อกุสโลติ, อามนฺตา.\csromanspacer{} ทิฏฺฐิคตํ อกุสลนฺติ, น เหวํ วตฺตพฺเพ ฯเปฯ.\csromanspacer{} ทิฏฺฐิคตสมฺปยุตฺตา เวทนา.\csromanspacer{} สญฺญา.\csromanspacer{} เจตนา.\csromanspacer{} จิตฺตํ อกุสลนฺติ, อามนฺตา.\csromanspacer{} ทิฏฺฐิคตํ อกุสลนฺติ, น เหวํ วตฺตพฺเพ ฯเปฯ.}

\proseitem{707}{\csromansymbolmark{*}ทิฏฺฐิคตํ อพฺยากตนฺติ, อามนฺตา.\csromanspacer{} อผลํ อวิปากนฺติ, น เหวํ วตฺตพฺเพ ฯเปฯ.\csromanspacer{} นนุ สผลํ สวิปากนฺติ, อามนฺตา.\csromanspacer{} หญฺจิ สผลํ สวิปากํ, โน จ วต เร วตฺตพฺเพ “ทิฏฺฐิคตํ อพฺยากตนฺ”ติ.}

\prose{\csromansymbolmark{*}ทิฏฺฐิคตํ อพฺยากตนฺติ, อามนฺตา.\csromanspacer{} นนุ มิจฺฉาทิฏฺฐิปรมานิ วชฺชานิ\footnote{อํ 1. 35 ปิฏฺเฐ.} วุตฺตานิ ภควตาติ, อามนฺตา.\csromanspacer{} หญฺจิ มิจฺฉาทิฏฺฐิปรมานิ วชฺชานิ วุตฺตานิ ภควตา, โน จ วต เร วตฺตพฺเพ “ทิฏฺฐิคตํ อพฺยากตนฺ”ติ.}

\prose{\csromansymbolmark{*}ทิฏฺฐิคตํ อพฺยากตนฺติ, อามนฺตา.\csromanspacer{} นนุ วุตฺตํ ภควตา มิจฺฉาทิฏฺฐิ โข วจฺฉ อกุสลา,\footnote{อกุสลํ (ม 2. 157 ปิฏฺเฐ) 3. กุสลํ (ม 2. 157 ปิฏฺเฐ)} สมฺมาทิฏฺฐิ กุสลา3 ติ\footnote{ม 2. 157 ปิฏฺเฐ.} อตฺเถว สุตฺตนฺโตติ, อามนฺตา.\csromanspacer{} เตน หิ น วตฺตพฺพํ “ทิฏฺฐิคตํ อพฺยากตนฺ”ติ.}

\prose{\csromansymbolmark{*}ทิฏฺฐิคตํ อพฺยากตนฺติ, อามนฺตา.\csromanspacer{} นนุ วุตฺตํ ภควตา “มิจฺฉาทิฏฺฐิสฺส โข อหํ ปุณฺณ ทฺวินฺนํ คตีนํ อญฺญตรํ คติํ วทามิ นิรยํ วา ติรจฺฉานโยนิํ วา”ติ\footnote{ม 2. 50 ปิฏฺเฐ.} อตฺเถว สุตฺตนฺโตติ, อามนฺตา.\csromanspacer{} เตน หิ จ วตฺตพฺพํ “ทิฏฺฐิคตํ อพฺยากตนฺ”ติ.}

\proseitem{708}{\csromansymbolmark{+}น วตฺตพฺพํ “ทิฏฺฐิคตํ อพฺยากตนฺ”ติ, อามนฺตา.\csromanspacer{} นนุ วุตฺตํ ภควตา\csromandash{}“สสฺสโต โลโก”ติ โข วจฺฉ อพฺยากตเมตํ, “อสสฺสโต โลโก”ติ โข วจฺฉ อพฺยากตเมตํ, “อนฺตวา โลโก”ติ โข วจฺฉ อพฺยากตเมตํ, “อนนฺตวา โลโก”ติ โข วจฺฉ ฯเปฯ “ตํ ชีวํ ตํ สรีรนฺ”ติ โข วจฺฉ ฯเปฯ “อญฺญํ ชีวํ อญฺญํ สรีรนฺ”ติ}

\vspace{\tipitakaparskip}
\csromancenter{จุทฺทสมวคฺค โข วจฺฉ ฯเปฯ “โหติ ตถาคโต ปรํ มรณา”ติ โข วจฺฉ ฯเปฯ “น โหติ ตถาคโต ปรํ มรณา”ติ โข วจฺฉ ฯเปฯ “โหติ จ น จ โหติ ตถาคโต ปรํ มรณา”ติ โข วจฺฉ ฯเปฯ “เนว โหติ น น โหติ ตถาคโต ปรํ มรณา”ติ โข วจฺฉ อพฺยากตเมตนฺติ\footnote{สํ 2. 559 ปิฏฺเฐ, โถกํ ปน วิสทิสํ.} อตฺเถว สุตฺตนฺโตติ, อามนฺตา.\csromanspacer{} เตน หิ ทิฏฺฐิคตํ อพฺยากตนฺติ.}
\csromanmatikaanchor{mtk.174}
\csromantocmarkref{subsection}{(144) 9. อปริยาปนฺนกถา}{mtk.174}
\bookmark[dest={mtk.174},level=2]{(144) 9. อปริยาปนฺนกถา}
\prose{\csromanfolio{367}\csromansymbolmark{*}ทิฏฺฐิคตํ อพฺยากตนฺติ, อามนฺตา.\csromanspacer{} นนุ วุตฺตํ ภควตา “มิจฺฉาทิฏฺฐิกสฺส ภิกฺขเว ปุริสปุคฺคลสฺส ยญฺเจว กายกมฺมํ ยถาทิฏฺฐิสมตฺตํ สมาทินฺนํ, ยญฺจ วจีกมฺมํ ฯเปฯ ยญฺจ มโนกมฺมํ ยา จ เจตนา ยา จ ปตฺถนา โย จ ปณิธิ เย จ สงฺขารา, สพฺเพ เต ธมฺมา อนิฏฺฐาย อกนฺตาย อมนาปาย อหิตาย ทุกฺขาย สํวตฺตนฺตี”ติ\footnote{อํ 1. 33 ปิฏฺเฐ.} อตฺเถว สุตฺตนฺโตติ, อามนฺตา.\csromanspacer{} เตน หิ น วตฺตพฺพํ “ทิฏฺฐิคตํ อพฺยากตนฺ”ติ.}

\nitthitam{อพฺยากตกถา นิฏฺฐิตา.}
\chapterhead{14. จุทฺทสมวคฺค}
\vspace{\tipitakaparskip}
\csromancenter{\textbf{(144) 9.}\csromanspacer{}\textbf{ อปริยาปนฺนกถา}}
\proseitem{709}{\csromansymbolmark{*}ทิฏฺฐิคตํ อปริยาปนฺนนฺติ, อามนฺตา.\csromanspacer{} มคฺโค ผลํ นิพฺพานํ โสตาปตฺติมคฺโค โสตาปตฺติผลํ, สกทาคามิมคฺโค สกทาคามิผลํ, อนาคามิมคฺโค อนาคามิผลํ, อรหตฺตมคฺโค อรหตฺตผลํ สติปฏฺฐานํ สมฺมปฺปธานํ อิทฺธิปาโท อินฺทฺริยํ พลํ โพชฺฌงฺโคติ, น เหวํ วตฺตพฺเพ ฯเปฯ.}

\proseitem{710}{\csromansymbolmark{+}น วตฺตพฺพํ “ทิฏฺฐิคตํ อปริยาปนฺนนฺ”ติ, อามนฺตา.\csromanspacer{} ปุถุชฺชโน “กาเมสุ วีตราโค”ติ วตฺตพฺโพติ, อามนฺตา. “วิคตทิฏฺฐิโย”ติ วตฺตพฺโพติ, น เหวํ วตฺตพฺเพ.\csromanspacer{} เตน หิ ทิฏฺฐิคตํ อปริยาปนฺนนฺติ.}

\nitthitam{อปริยาปนฺนกถา นิฏฺฐิตา.}
\vspace{\tipitakaparskip}
\csromancenter{จุทฺทสโม วคฺโค.}
\titlehead{\textbf{ตสฺสุทฺทานํ}}
\csromanmatikaanchor{mtk.175}
\csromanmatikaanchor{mtk.176}
\csromantocmarknopage{chapter}{15. ปนฺนารสมวคฺค}{mtk.175}
\bookmark[dest={mtk.175},level=0]{15. ปนฺนารสมวคฺค}
\csromantocmarkref{subsection}{(145) 1. ปจฺจยตากถา}{mtk.176}
\bookmark[dest={mtk.176},level=2]{(145) 1. ปจฺจยตากถา}
\prose{\csromanfolio{368}อกุสลมูลํ ปฏิสนฺทหติ กุสลมูลํ, กุสลมูลํ ปฏิสนฺทหติ อกุสลมูลํ, สฬายตนํ, ฉวิญฺญาณกายา, อริยรูปํ มหาภูตานํ อุปาทาย, เสฺวว อนุสโย ตํ ปริยุฏฺฐานํ, ปริยุฏฺฐานํ จิตฺตวิปฺปยุตฺตํ, ยถาธาตุ ตญฺเญว อนุเสติ, ทิฏฺฐิคตํ อพฺยากตํ, ทิฏฺฐิคตํ อปริยาปนฺนนฺติ.}

\chapterhead{15. ปนฺนรสมวคฺค}
\vspace{\tipitakaparskip}
\csromancenter{\textbf{(145) 1.}\csromanspacer{}\textbf{ ปจฺจยตากถา}}
\proseitem{711}{\csromansymbolmark{*}ปจฺจยตา ววตฺถิตาติ, อามนฺตา.\csromanspacer{} นนุ วีมํสา เหตุ, โส จ อธิปตีติ, อามนฺตา.\csromanspacer{} หญฺจิ วีมํสา เหตุ, โส จ อธิปติ, เตน วต เร วตฺตพฺเพ “เหตุปจฺจเยน ปจฺจโย, อธิปติปจฺจเยน ปจฺจโย”ติ.}

\prose{\csromansymbolmark{*}นนุ ฉนฺธาธิปติ สหชาตานํ ธมฺมานํ อธิปตีติ, อามนฺตา.\csromanspacer{} หญฺจิ ฉนฺทาธิปติ สหชาตานํ ธมฺมานํ อธิปติ, เตน วต เร วตฺตพฺเพ “อธิปติปจฺจเยน ปจฺจโย, สหชาตปจฺจเยน ปจฺจโย”ติ.}

\proseitem{712}{\csromansymbolmark{*}นนุ วีริยาธิปติ สหชาตานํ ธมฺมานํ อธิปตีติ, อามนฺตา.\csromanspacer{} หญฺจิ วีริยาธิปติ สหชาตานํ ธมฺมานํ อธิปติ, เตน วต เร วตฺตพฺเพ “อธิปติปจฺจเยน ปจฺจโย, สหชาตปจฺจเยน ปจฺจโย”ติ.}

\prose{\csromansymbolmark{*}นนุ วีริยาธิปติ สหชาตานํ ธมฺมานํ อธิปติ, ตญฺจ อินฺทฺริยนฺติ, อามนฺตา.\csromanspacer{} หญฺจิ วีริยาธิปติ สหชาตานํ ธมฺมานํ อธิปติ, ตญฺจ อินฺทฺริยํ, เตน วต เร วตฺตพฺเพ “อธิปติปจฺจเยน ปจฺจโย, อินฺทฺริยปจฺจเยน ปจฺจโย”ติ.}

\prose{\csromansymbolmark{*}นนุ วีริยาธิปติ สหชาตานํ ธมฺมานํ อธิปติ, ตญฺจ มคฺคงฺคนฺติ, อามนฺตา.\csromanspacer{} หญฺจิ วีริยาธิปติ สหชาตานํ ธมฺมานํ อธิปติ, ตญฺจ มคฺคงฺคํ, เตน วต เร วตฺตพฺเพ “อธิปติปจฺจเยน ปจฺจโย, มคฺคปจฺจเยน ปจฺจโย”ติ.}

\chapterheadpageread{15. ปนฺนรสมวคฺค}
\proseitem{713}{\csromanfolio{369}\csromansymbolmark{*}นนุ จิตฺตาธิปติ สหชาตานํ ธมฺมานํ อธิปตีติ, อามนฺตา.\csromanspacer{} หญฺจิ จิตฺตาธิปติ สหชาตานํ ธมฺมานํ อธิปติ, เตน วต เร วตฺตพฺเพ “อธิปติปจฺจเยน ปจฺจโย, สหชาตปจฺจเยน ปจฺจโย”ติ.}

\prose{\csromansymbolmark{*}นนุ จิตฺตาธิปติ สหชาตานํ ธมฺมานํ อธิปติ, โส จ อาหาโรติ, อามนฺตา.\csromanspacer{} หญฺจิ จิตฺตาธิปติ สหชาตานํ ธมฺมานํ อธิปติ, โส จ อาหาโร, เตน วต เร วตฺตพฺเพ “อธิปติปจฺจเยน ปจฺจโย, อาหารปจฺจเยน ปจฺจโย”ติ.}

\prose{\csromansymbolmark{*}นนุ จิตฺตาธิปติ สหชาตานํ ธมฺมานํ อธิปติ, ตญฺจ อินฺทฺริยนฺติ, อามนฺตา.\csromanspacer{} หญฺจิ จิตฺตาธิปติ สหชาตานํ ธมฺมานํ อธิปติ, ตญฺจ อินฺทฺริยํ, เตน วต เร วตฺตพฺเพ “อธิปติปจฺจเยน ปจฺจโย, อินฺทฺริยปจฺจเยน ปจฺจโย”ติ.}

\proseitem{714}{\csromansymbolmark{*}นนุ วีมํสาธิปติ สหชาตานํ ธมฺมานํ อธิปตีติ, อามนฺตา.\csromanspacer{} หญฺจิ วีมํสาธิปติ สหชาตานํ ธมฺมานํ อธิปติ, เตน วต เร วตฺตพฺเพ “อธิปติปจฺจเยน ปจฺจโย, สหชาตปจฺจเยน ปจฺจโย”ติ.}

\prose{\csromansymbolmark{*}นนุ วีมํสาธิปติ สหชาตานํ ธมฺมานํ อธิปติ, ตญฺจ อินฺทฺริยนฺติ, อามนฺตา.\csromanspacer{} หญฺจิ วีมํสาธิปติ สหชาตานํ ธมฺมานํ อธิปติ, ตญฺจ อินฺทฺริยํ, เตน วต เร วตฺตพฺเพ “อธิปติปจฺจเยน ปจฺจโย, อินฺทฺริยปจฺจเยน ปจฺจโย”ติ.}

\prose{\csromansymbolmark{*}นนุ วีมํสาธิปติ สหชาตานํ ธมฺมานํ อธิปติ, ตญฺจ มคฺคงฺคนฺติ, อามนฺตา.\csromanspacer{} หญฺจิ วีมํสาธิปติ สหชาตานํ ธมฺมานํ อธิปติ, ตญฺจ มคฺคงฺคํ, เตน วต เร วตฺตพฺเพ “อธิปติปจฺจเยน ปจฺจโย, มคฺคปจฺจเยน ปจฺจโย”ติ.}

\proseitem{715}{\csromansymbolmark{*}นนุ อริยํ ธมฺมํ ครุํ กตฺวา อุปฺปชฺชติ ปจฺจเวกฺขณา, ตญฺจารมฺมณนฺติ, อามนฺตา.\csromanspacer{} หญฺจิ อริยํ ธมฺมํ ครุํ กตฺวา อุปฺปชฺชติ ปจฺจเวกฺขณา, ตญฺจารมฺมณํ, เตน วต เร วตฺตพฺเพ “อธิปติปจฺจเยน ปจฺจโย, อารมฺมณปจฺจเยน ปจฺจโย”ติ.}

\csromanmatikaanchor{mtk.177}
\csromantocmarkref{subsection}{(146) 2. อญฺญมญฺญปจฺจยกถา}{mtk.177}
\bookmark[dest={mtk.177},level=2]{(146) 2. อญฺญมญฺญปจฺจยกถา}
\proseitem{716}{\csromansymbolmark{*}นนุ ปุริมา ปุริมา กุสลา ธมฺมา ปจฺฉิมานํ ปจฺฉิมานํ กุสลานํ ธมฺมานํ อนนฺตรปจฺจเยน ปจฺจโย, สา จ อาเสวนาติ, อามนฺตา.\csromanspacer{} หญฺจิ \csromanfolio{370}ปุริมา ปุริมา กุสลา ธมฺมา ปจฺฉิมานํ ปจฺฉิมานํ กุสลานํ ธมฺมานํ อนนฺตรปจฺจเยน ปจฺจโย, สา จ อาเสวนา, เตน วต เร วตฺตพฺเพ “อนนฺตรปจฺจเยน ปจฺจโย, อาเสวนปจฺจเยน ปจฺจโย”ติ.}

\prose{\csromansymbolmark{*}นนุ ปุริมา ปุริมา อกุสลา ธมฺมา ปจฺฉิมานํ ปจฺฉิมานํ อกุสลานํ ธมฺมานํ อนนฺตรปจฺจเยน ปจฺจโย, สา จ อาเสวนาติ, อามนฺตา.\csromanspacer{} หญฺจิ ปุริมา ปุริมา อกุสลา ธมฺมา ปจฺฉิมานํ ปจฺฉิมานํ อกุสลานํ ธมฺมานํ อนนฺตรปจฺจเยน ปจฺจโย, สา จ อาเสวนา, เตน วต เร วตฺตพฺเพ “อนนฺตรปจฺจเยน ปจฺจโย, อาเสวนปจฺจเยน ปจฺจโย”ติ.}

\prose{\csromansymbolmark{*}นนุ ปุริมา ปุริมา กิริยาพฺยากตา ธมฺมา ปจฺฉิมานํ ปจฺฉิมานํ กิริยาพฺยากตานํ ธมฺมานํ อนนฺตรปจฺจเยน ปจฺจโย, สา จ อาเสวนาติ, อามนฺตา.\csromanspacer{} หญฺจิ ปุริมา ปุริมา กิริยาพฺยากตา ธมฺมา ปจฺฉิมานํ ปจฺฉิมานํ กิริยาพฺยากตานํ ธมฺมานํ อนนฺตรปจฺจเยน ปจฺจโย, สา จ อาเสวนา, เตน วต เร วตฺตพฺเพ “อนนฺตรปจฺจเยน ปจฺจโย, อาเสวนปจฺจเยน ปจฺจโย”ติ.}

\proseitem{717}{\csromansymbolmark{*}น วตฺตพฺพํ “ปจฺจยตา ววตฺถิตา”ติ, อามนฺตา.\csromanspacer{} เหตุปจฺจเยน ปจฺจโย โหติ, อารมฺมณปจฺจเยน ปจฺจโย โหติ, อนนฺตรปจฺจเยน ปจฺจโย โหติ, สมนนฺตรปจฺจเยน ปจฺจโย โหตีติ, น เหวํ วตฺตพฺเพ.\csromanspacer{} เตน หิ “ปจฺจยตา ววตฺถิตา”ติ.}

\nitthitam{ปจฺจยตากถา นิฏฺฐิตา.}
\chapterhead{15. ปนฺนรสมวคฺค}
\vspace{\tipitakaparskip}
\csromancenter{\textbf{(146) 2.}\csromanspacer{}\textbf{ อญฺญมญฺญปจฺจยกถา}}
\proseitem{718}{\csromansymbolmark{*}อวิชฺชาปจฺจยาว สงฺขารา, น วตฺตพฺพํ “สงฺขารปจฺจยาปิ อวิชฺชา”ติ, อามนฺตา.\csromanspacer{} นนุ อวิชฺชา สงฺขาเรน สหชาตาติ, อามนฺตา.\csromanspacer{} หญฺจิ อวิชฺชา สงฺขาเรน สหชาตา, เตน วต เร วตฺตพฺเพ “อวิชฺชาปจฺจยาปิ สงฺขารา, สงฺขารปจฺจยาปิ อวิชฺชา”ติ.}

\prose{\csromansymbolmark{*}ตณฺหาปจฺจยาว อุปาทานํ, น วตฺตพฺพํ “อุปาทานปจฺจยาปิ ตณฺหา”ติ, อามนฺตา.\csromanspacer{} นนุ ตณฺหา อุปาทาเนน สหชาตาติ, อามนฺตา.\csromanspacer{} หญฺจิ ตณฺหา}

\vspace{\tipitakaparskip}
\csromancenter{ปนฺนรสมวคฺค อุปาทาเนน สหชาตา, เตน วต เร วตฺตพฺเพ “ตณฺหาปจฺจยาปิ อุปาทานํ, อุปาทานปจฺจยาปิ ตณฺหา”ติ.}
\csromanmatikaanchor{mtk.178}
\csromantocmarkref{subsection}{(147) 3. อทฺธากถา}{mtk.178}
\bookmark[dest={mtk.178},level=2]{(147) 3. อทฺธากถา}
\proseitem{719}{\csromanfolio{371}\csromansymbolmark{+}“ชรามรณปจฺจยา ภิกฺขเว ชาติ ชาติปจฺจยา ภโว”ติ อตฺเถว สุตฺตนฺโตติ, นตฺถิ.\csromanspacer{} เตน หิ อวิชฺชาปจฺจยาว สงฺขารา, น วตฺตพฺพํ “สงฺขารปจฺจยาปิ อวิชฺชา”ติ, ตณฺหาปจฺจยาว อุปาทานํ, น วตฺตพฺพํ “อุปาทานปจฺจยาปิ ตณฺหา”ติ.}

\prose{\csromansymbolmark{*}“วิญฺญาณปจฺจยา ภิกฺขเว นามรูปํ นามรูปปจฺจยาปิ วิญฺญาณนฺ”ติ\footnote{ที 2. 28 ปิฏฺเฐ, โถกํ ปน วิสทิสํ.} อตฺเถว สุตฺตนฺโตติ, อามนฺตา.\csromanspacer{} เตน หิ อวิชฺชาปจฺจยาปิ สงฺขารา, สงฺขารปจฺจยาปิ อวิชฺชา, ตณฺหาปจฺจยาปิ อุปาทานํ, อุปาทานปจฺจยาปิ ตณฺหาติ.}

\nitthitam{อญฺญมญฺญปจฺจยกถา นิฏฺฐิตา.}
\chapterhead{15. ปนฺนรสมวคฺค}
\vspace{\tipitakaparskip}
\csromancenter{\textbf{(147) 3.}\csromanspacer{}\textbf{ อทฺธากถา}}
\proseitem{720}{\csromansymbolmark{*}อทฺธา ปรินิปฺผนฺโนติ, อามนฺตา.\csromanspacer{} รูปนฺติ, น เหวํ วตฺตพฺเพ ฯเปฯ.\csromanspacer{} เวทนา ฯเปฯ.\csromanspacer{} สญฺญา ฯเปฯ.\csromanspacer{} สงฺขารา ฯเปฯ.\csromanspacer{} วิญฺญาณนฺติ, น เหวํ วตฺตพฺเพ ฯเปฯ.\csromanspacer{} อตีโต อทฺธา ปรินิปฺผนฺโนติ, อามนฺตา.\csromanspacer{} รูปนฺติ, น เหวํ วตฺตพฺเพ ฯเปฯ.\csromanspacer{} เวทนา ฯเปฯ.\csromanspacer{} สญฺญา ฯเปฯ.\csromanspacer{} สงฺขารา ฯเปฯ.\csromanspacer{} วิญฺญาณนฺติ, น เหวํ วตฺตพฺเพ ฯเปฯ.\csromanspacer{} อนาคโต อทฺธา ปรินิปฺผนฺโนติ, อามนฺตา.\csromanspacer{} รูปนฺติ, น เหวํ วตฺตพฺเพ ฯเปฯ.\csromanspacer{} เวทนา ฯเปฯ.\csromanspacer{} สญฺญา ฯเปฯ.\csromanspacer{} สงฺขารา ฯเปฯ.\csromanspacer{} วิญฺญาณนฺติ, น เหวํ วตฺตพฺเพ ฯเปฯ.\csromanspacer{} ปจฺจุปฺปนฺโน อทฺธา ปรินิปฺผนฺโนติ, อามนฺตา.\csromanspacer{} รูปนฺติ, น เหวํ วตฺตพฺเพ ฯเปฯ.\csromanspacer{} เวทนา ฯเปฯ.\csromanspacer{} สญฺญา ฯเปฯ.\csromanspacer{} สงฺขารา ฯเปฯ.\csromanspacer{} วิญฺญาณนฺติ, น เหวํ วตฺตพฺเพ ฯเปฯ.}

\prose{\csromansymbolmark{*}อตีตํ รูปํ เวทนา สญฺญา สงฺขารา วิญฺญาณํ อตีโต อทฺธาติ, อามนฺตา.\csromanspacer{} อตีตา ปญฺจทฺธาติ, น เหวํ วตฺตพฺเพ ฯเปฯ.\csromanspacer{} อนาคตํ รูปํ เวทนา สญฺญา สงฺขารา วิญฺญาณํ อนาคโต อทฺธาติ, อามนฺตา.\csromanspacer{} อนาคตา ปญฺจทฺธาติ, น เหวํ วตฺตพฺเพ ฯเปฯ.\csromanspacer{} ปจฺจุปฺปนฺนํ รูปํ เวทนา สญฺญา สงฺขารา วิญฺญาณํ ปจฺจุปฺปนฺโน อทฺธาติ, อามนฺตา.\csromanspacer{} ปจฺจุปฺปนฺนา ปญฺจทฺธาติ, น เหวํ วตฺตพฺเพ ฯเปฯ.}

\csromanmatikaanchor{mtk.179}
\csromantocmarkref{subsection}{(148) 4. ขณลยมุหุตฺตกถา}{mtk.179}
\bookmark[dest={mtk.179},level=2]{(148) 4. ขณลยมุหุตฺตกถา}
\prose{\csromanfolio{372}\csromansymbolmark{*}อตีตา ปญฺจกฺขนฺธา อตีโต อทฺธา, อนาคตา ปญฺจกฺขนฺธา อนาคโต อทฺธา, ปจฺจุปฺปนฺนา ปญฺจกฺขนฺธา ปจฺจุปฺปนฺโน อทฺธาติ, อามนฺตา.\csromanspacer{} ปนฺนรสทฺธาติ, น เหวํ วตฺตพฺเพ ฯเปฯ.}

\prose{\csromansymbolmark{*}อตีตานิ ทฺวาทสายตนานิ อตีโต อทฺธา, อนาคตานิ ทฺวาทสายตนานิ อนาคโต อทฺธา, ปจฺจุปฺปนฺนานิ ทฺวาทสายตนานิ ปจฺจุปฺปนฺโน อทฺธาติ, อามนฺตา.\csromanspacer{} ฉตฺติํส อทฺธาติ, น เหวํ วตฺตพฺเพ ฯเปฯ.}

\prose{\csromansymbolmark{*}อตีตา อฏฺฐารส ธาตุโย อตีโต อทฺธา, อนาคตา อฏฺฐารส ธาตุโย อนาคโต อทฺธา, ปจฺจุปฺปนฺนา อฏฺฐารส ธาตุโย ปจฺจุปฺปนฺโน อทฺธาติ, อามนฺตา.\csromanspacer{} จตุปญฺญาส อทฺธาติ, น เหวํ วตฺตพฺเพ ฯเปฯ.}

\prose{\csromansymbolmark{*}อตีตานิ พาวีสตินฺทฺริยานิ อตีโต อทฺธา, อนาคตานิ พาวีสิตินฺทฺริยานิ อนาคโต อทฺธา, ปจฺจุปฺปนฺนานิ พาวีสตินฺทฺริยานิ ปจฺจุปฺปนฺโน อทฺธาติ, อามนฺตา.\csromanspacer{} ฉสฏฺฐิ อทฺธาติ, น เหวํ วตฺตพฺเพ ฯเปฯ.}

\proseitem{721}{\csromansymbolmark{+}น วตฺตพฺพํ “อทฺธา ปรินิปฺผนฺโน”ติ, อามนฺตา.\csromanspacer{} นนุ วุตฺตํ ภควตา “ตีณิมานิ ภิกฺขเว กถาวตฺถูนิ.\csromanspacer{} กตมานิ ตีณิ, อตีตํ วา ภิกฺขเว อทฺธานํ อารพฺภ กถํ กเถยฺย ‘เอวํ อโหสิ อตีตมทฺธานนฺ’ติ, อนาคตํ วา ภิกฺขเว อทฺธานํ อารพฺภ กถํ กเถยฺย ‘เอวํ ภวิสฺสติ อนาคตมทฺธานนฺ’ติ, เอตรหิ วา ภิกฺขเว ปจฺจุปฺปนฺนํ อทฺธานํ อารพฺภ กถํ กเถยฺย ‘เอวํ โหติ เอตรหิ ปจฺจุปฺปนฺนนฺ’ติ.\csromanspacer{} อิมานิ โข ภิกฺขเว ตีณิ กถาวตฺถูนี”ติ\footnote{อํ 1. 197 ปิฏฺเฐ; ที 3. 184 ปิฏฺเฐ จ.} อตฺเถว สุตฺตนฺโตติ, อามนฺตา.\csromanspacer{} เตน หิ อทฺธา ปรินิปฺผนฺโนติ.}

\nitthitam{อทฺธากถา นิฏฺฐิตา.}
\chapterhead{15. ปนฺนรสมวคฺค}
\vspace{\tipitakaparskip}
\csromancenter{\textbf{(148) 4.}\csromanspacer{}\textbf{ ขณลยมุหุตฺตกถา}}
\proseitem{722}{\csromansymbolmark{*}ขโณ ปรินิปฺผนฺโน, ลโย ปรินิปฺผนฺโน, มุหุตฺตํ ปรินิปฺผนฺนนฺติ, อามนฺตา.\csromanspacer{} รูปนฺติ, น เหวํ วตฺตพฺเพ ฯเปฯ.\csromanspacer{} เวทนา.\csromanspacer{} สญฺญา.\csromanspacer{} สงฺขารา.\csromanspacer{} วิญฺญาณนฺติ, น เหวํ วตฺตพฺเพ ฯเปฯ.}

\chapterheadpageread{15. ปนฺนรสมวคฺค}
\csromanmatikaanchor{mtk.180}
\csromanmatikaanchor{mtk.181}
\csromantocmarkref{subsection}{(149) 5. อาสวกถา}{mtk.180}
\bookmark[dest={mtk.180},level=2]{(149) 5. อาสวกถา}
\csromantocmarkref{subsection}{(150) 6. ชรามรณกถา}{mtk.181}
\bookmark[dest={mtk.181},level=2]{(150) 6. ชรามรณกถา}
\proseitem{723}{\csromanfolio{373}\csromansymbolmark{*}น วตฺตพฺพํ “มุหุตฺตํ ปรินิปฺผนฺนนฺ”ติ, อามนฺตา.\csromanspacer{} นนุ วุตฺตํ ภควตา “ตีณิมานิ ภิกฺขเว กถาวตฺถูนิ.\csromanspacer{} กตมานิ ตีณิ, อตีตํ วา ภิกฺขเว อทฺธานํ อารพฺภ กถํ กเถยฺย ‘เอวํ อโหสิ อตีตมทฺธานนฺ’ติ.\csromanspacer{} อนาคตํ วา ภิกฺขเว อทฺธานํ อารพฺภ กถํ กเถยฺย ‘เอวํ ภวิสฺสติ อนาคตมทฺธานนฺ’ติ.\csromanspacer{} เอตรหิ วา ภิกฺขเว ปจฺจุปฺปนฺนํ อทฺธานํ อารพฺภ กถํ กเถยฺย ‘เอวํ โหติ เอตรหิ ปจฺจุปฺปนฺนนฺ’ติ.\csromanspacer{} อิมานิ โข ภิกฺขเว ตีณิ กถาวตฺถูนี”ติ อตฺเถว สุตฺตนฺโตติ, อามนฺตา.\csromanspacer{} เตน หิ มุหุตฺตํ ปรินิปฺผนฺนนฺติ.}

\nitthitam{ขณลยมุหุตฺตกถา นิฏฺฐิตา.}
\chapterhead{15. ปนฺนรสมวคฺค}
\vspace{\tipitakaparskip}
\csromancenter{\textbf{(149) 5.}\csromanspacer{}\textbf{ อาสวกถา}}
\proseitem{724}{\csromansymbolmark{*}จตฺตาโร อาสวา อนาสวาติ, อามนฺตา.\csromanspacer{} มคฺโค ผลํ นิพฺพานํ โสตาปตฺติมคฺโค โสตาปตฺติผลํ ฯเปฯ.\csromanspacer{} โพชฺฌงฺโคติ, น เหวํ วตฺตพฺเพ ฯเปฯ.}

\proseitem{725}{\csromansymbolmark{+}น วตฺตพฺพํ “จตฺตาโร อาสวา อนาสวา”ติ, อามนฺตา.\csromanspacer{} อตฺถญฺเญว อาสวา เยหิ อาสเวหิ เต อาสวา สาสวา โหนฺตีติ, น เหวํ วตฺตพฺเพ.\csromanspacer{} เตน หิ จตฺตาโร อาสวา อนาสวาติ.}

\nitthitam{อาสวกถา นิฏฺฐิตา.}
\chapterhead{15. ปนฺนรสมวคฺค}
\vspace{\tipitakaparskip}
\csromancenter{\textbf{(150) 6.}\csromanspacer{}\textbf{ ชรามรณกถา}}
\csromanmatikaanchor{mtk.182}
\csromanmatikaanchor{mtk.183}
\csromantocmarkref{subsection}{(151) 7. สญฺญาเวทยิตกถา}{mtk.182}
\bookmark[dest={mtk.182},level=2]{(151) 7. สญฺญาเวทยิตกถา}
\csromantocmarkref{subsection}{(152) 8. ทุติยสญฺญาเวทยิตกถา}{mtk.183}
\bookmark[dest={mtk.183},level=2]{(152) 8. ทุติยสญฺญาเวทยิตกถา}
\proseitem{726}{\csromansymbolmark{*}โลกุตฺตรานํ ธมฺมานํ ชรามรณํ โลกุตฺตรนฺติ, อามนฺตา.\csromanspacer{} มคฺโค ผลํ นิพฺพานํ โสตาปตฺติมคฺโค โสตาปตฺติผลํ ฯเปฯ.\csromanspacer{} โพชฺฌงฺโคติ, น เหวํ วตฺตพฺเพ ฯเปฯ.\csromanspacer{} โสตาปตฺติมคฺคสฺส ชรามรณํ โสตาปตฺติมคฺโคติ, น เหวํ วตฺตพฺเพ ฯเปฯ.\csromanspacer{} โสตาปตฺติมคฺคสฺส ชรามรณํ โสตาปตฺติมคฺโคติ, อามนฺตา.\csromanspacer{} โสตาปตฺติผลสฺส ชรามรณํ โสตาปตฺติผลนฺติ, น เหวํ \csromanfolio{374}วตฺตพฺเพ ฯเปฯ.\csromanspacer{} สกทาคามิมคฺคสฺส ฯเปฯ.\csromanspacer{} สกทาคามิผลสฺส ฯเปฯ.\csromanspacer{} อนาคามิมคฺคสฺส ฯเปฯ.\csromanspacer{} อนาคามิผลสฺส ฯเปฯ.\csromanspacer{} อรหตฺตมคฺคสฺส ชรามรณํ อรหตฺตมคฺโคติ, น เหวํ วตฺตพฺเพ ฯเปฯ.\csromanspacer{} อรหตฺตมคฺคสฺส ชรามรณํ อรหตฺตมคฺโคติ, อามนฺตา.\csromanspacer{} อรหตฺตผลสฺส ชรามรณํ อรหตฺตผลนฺติ, น เหวํ วตฺตพฺเพ ฯเปฯ.\csromanspacer{} สติปฏฺฐานานํ.\csromanspacer{} สมฺมปฺปธานานํ.\csromanspacer{} อิทฺธิปาทานํ.\csromanspacer{} อินฺทฺริยานํ.\csromanspacer{} พลานํ.\csromanspacer{} โพชฺฌงฺคานํ ชรามรณํ โพชฺฌงฺโคติ, น เหวํ วตฺตพฺเพ ฯเปฯ.}

\proseitem{727}{\csromansymbolmark{+}น วตฺตพฺพํ “โลกุตฺตรานํ ธมฺมานํ ชรามรณํ โลกุตฺตรนฺ”ติ, อามนฺตา.\csromanspacer{} โลกิยนฺติ, น เหวํ วตฺตพฺเพ.\csromanspacer{} เตน หิ โลกุตฺตรนฺติ.}

\nitthitam{ชรามรณกถา นิฏฺฐิตา.}
\chapterhead{15. ปนฺนรสมวคฺค}
\vspace{\tipitakaparskip}
\csromancenter{\textbf{(151) 7.}\csromanspacer{}\textbf{ สญฺญาเวทยิตกถา}}
\proseitem{728}{\csromansymbolmark{*}สญฺญาเวทยิตนิโรธสมาปตฺติ โลกุตฺตราติ, อามนฺตา.\csromanspacer{} มคฺโค ผลํ นิพฺพานํ โสตาปตฺติมคฺโค โสตาปตฺติผลํ ฯเปฯ.\csromanspacer{} โพชฺฌงฺโคติ, น เหวํ วตฺตพฺเพ ฯเปฯ.}

\proseitem{729}{\csromansymbolmark{+}น วตฺตพฺพํ “สญฺญาเวทยิตนิโรธสมาปตฺติ โลกุตฺตรา”ติ, อามนฺตา.\csromanspacer{} โลกิยาติ, น เหวํ วตฺตพฺเพ.\csromanspacer{} เตน หิ โลกุตฺตราติ.}

\nitthitam{สญฺญาเวทยิตกถา นิฏฺฐิตา.}
\chapterhead{15. ปนฺนรสมวคฺค}
\vspace{\tipitakaparskip}
\csromancenter{\textbf{(152) 8.}\csromanspacer{}\textbf{ ทุติยสญฺญาเวทยิตกถา}}
\proseitem{730}{\csromansymbolmark{*}สญฺญาเวทยิตนิโรธสมาปตฺติ โลกิยาติ, อามนฺตา.\csromanspacer{} รูปนฺติ, น เหวํ วตฺตพฺเพ ฯเปฯ.\csromanspacer{} เวทนา.\csromanspacer{} สญฺญา.\csromanspacer{} สงฺขารา.\csromanspacer{} วิญฺญาณนฺติ, น เหวํ วตฺตพฺเพ ฯเปฯ.\csromanspacer{} กามาวจราติ, น เหวํ วตฺตพฺเพ ฯเปฯ.\csromanspacer{} รูปาวจราติ, น เหวํ วตฺตพฺเพ ฯเปฯ.\csromanspacer{} อรูปาวจราติ, น เหวํ วตฺตพฺเพ ฯเปฯ.}

\chapterheadpageread{15. ปนฺนรสมวคฺค}
\csromanmatikaanchor{mtk.184}
\csromantocmarkref{subsection}{(153) 9. ตติยสญฺญาเวทยิตกถา}{mtk.184}
\bookmark[dest={mtk.184},level=2]{(153) 9. ตติยสญฺญาเวทยิตกถา}
\proseitem{731}{\csromanfolio{375}\csromansymbolmark{+}น วตฺตพฺพํ “สญฺญาเวทยิตนิโรธสมาปตฺติ โลกิยา”ติ, อามนฺตา.\csromanspacer{} โลกุตฺตราติ, น เหวํ วตฺตพฺเพ.\csromanspacer{} เตน หิ โลกิยาติ.}

\nitthitam{ทุติยสญฺญาเวทยิตกถา นิฏฺฐิตา.}
\chapterhead{15. ปนฺนรสมวคฺค}
\vspace{\tipitakaparskip}
\csromancenter{\textbf{(153) 9.}\csromanspacer{}\textbf{ ตติยสญฺญาเวทยิตกถา}}
\proseitem{732}{\csromansymbolmark{*}สญฺญาเวทยิตนิโรธํ สมาปนฺโน กาลํ กเรยฺยาติ, อามนฺตา.\csromanspacer{} อตฺถิ สญฺญาเวทยิตนิโรธํ สมาปนฺนสฺส มารณนฺติโย ผสฺโส, มารณนฺติยา เวทนา, มารณนฺติยา สญฺญา, มรณนฺติยา เจตนา, มารณนฺติยํ จิตฺตนฺติ, น เหวํ วตฺตพฺเพ ฯเปฯ.\csromanspacer{} นตฺถิ สญฺญาเวทยิตนิโรธํ สมาปนฺนสฺส มารณนฺติโย ผสฺโส, มารณนฺติยา เวทนา, มารณนฺติยา สญฺญา, มารณนฺติยา เจตนา, มารณนฺติยํ จิตฺตนฺติ, อามนฺตา.\csromanspacer{} หญฺจิ นตฺถิ สญฺญาเวทยิตนิโรธํ สมาปนฺนสฺส มารณนฺติโย ผสฺโส, มารณนฺติยา เวทนา, มารณนฺติยา สญฺญา, มารณนฺติยา เจตนา, มารณนฺติยํ จิตฺตํ, โน จ วต เร วตฺตพฺเพ “สญฺญาเวทยิตนิโรธํ สมาปนฺโน กาลํ กเรยฺยา”ติ.}

\prose{\csromansymbolmark{*}สญฺญาเวทยิตนิโรธํ สมาปนฺโน กาลํ กเรยฺยาติ, อามนฺตา.\csromanspacer{} อตฺถิ สญฺญาเวทยิตนิโรธํ สมาปนฺนสฺส ผสฺโส เวทนา สญฺญา เจตนา จิตฺตนฺติ, น เหวํ วตฺตพฺเพ ฯเปฯ.\csromanspacer{} นตฺถิ สญฺญาเวทยิตนิโรธํ สมาปนฺนสฺส ผสฺโส เวทนา สญฺญา เจตนา จิตฺตนฺติ, อามนฺตา.\csromanspacer{} อผสฺสกสฺส กาลํ กิริยา อเวทนกสฺส กาลํ กิริยา ฯเปฯ.\csromanspacer{} อจิตฺตกสฺส กาลํ กิริยาติ, น เหวํ วตฺตพฺเพ ฯเปฯ.\csromanspacer{} นนุ สผสฺสกสฺส กาลํ กิริยา ฯเปฯ.\csromanspacer{} สจิตฺตกสฺส กาลํ กิริยาติ, อามนฺตา.\csromanspacer{} หญฺจิ สผสฺสกสฺส กาลํ กิริยา ฯเปฯ.\csromanspacer{} สจิตฺตกสฺส กาลํ กิริยา, โน จ วต เร วตฺตพฺเพ “สญฺญาเวทยิตนิโรธํ สมาปนฺโน กาลํ กเรยฺยา”ติ.}

\csromanmatikaanchor{mtk.185}
\csromantocmarkref{subsection}{(154) 10. อสญฺญสตฺตุปิกกถา}{mtk.185}
\bookmark[dest={mtk.185},level=2]{(154) 10. อสญฺญสตฺตุปิกกถา}
\prose{\csromansymbolmark{*}สญฺญาเวทยิตนิโรธํ สมาปนฺโน กาลํ กเรยฺยาติ, อามนฺตา.\csromanspacer{} สญฺญาเวทยิตนิโรธํ สมาปนฺนสฺส กาเย วิสํ กเมยฺย, สตฺถํ กเมยฺย, อคฺคิ กเมยฺยาติ, น เหวํ วตฺตพฺเพ ฯเปฯ.\csromanspacer{} สญฺญาเวทยิตนิโรธํ สมาปนฺนสฺส กาเย วิสํ น กเมยฺย, สตฺถํ น กเมยฺย, อคฺคิ น กเมยฺยาติ, อามนฺตา. \csromanfolio{376}หญฺจิ สญฺญาเวทยิตนิโรธํ สมาปนฺนสฺส กาเย วิสํ น กเมยฺย, สตฺถํ น กเมยฺย, อคฺคิ น กเมยฺย, โน จ วต เร วตฺตพฺเพ “สญฺญาเวทยิตนิโรธํ สมาปนฺโน กาลํ กเรยฺยา”ติ.}

\prose{\csromansymbolmark{*}สญฺญาเวทยิตนิโรธํ สมาปนฺโน กาลํ กเรยฺยาติ, อามนฺตา.\csromanspacer{} สญฺญาเวทยิตนิโรธํ สมาปนฺนสฺส กาเย วิสํ กเมยฺย, สตฺถํ กเมยฺย, อคฺคิ กเมยฺยาติ, อามนฺตา.\csromanspacer{} น นิโรธํ สมาปนฺโนติ, น เหวํ วตฺตพฺเพ ฯเปฯ.}

\proseitem{733}{\csromansymbolmark{+}สญฺญาเวทยิตนิโรธํ สมาปนฺโน น กาลํ กเรยฺยาติ, อามนฺตา.\csromanspacer{} อตฺถิ โส นิยาโม เยน นิยาเมน นิยโต สญฺญาเวทยิตนิโรธํ สมาปนฺโน น กาลํ กเรยฺยาติ, นตฺถิ.\csromanspacer{} หญฺจิ นตฺถิ โส นิยาโม เยน นิยาเมน นิยโต สญฺญาเวทยิตนิโรธํ สมาปนฺโน น กาลํ กเรยฺย, โน จ วต เร วตฺตพฺเพ “สญฺญาเวทยิตนิโรธํ สมาปนฺโน น กาลํ กเรยฺยา”ติ.}

\proseitem{734}{\csromansymbolmark{*}จกฺขุวิญฺญาณสมงฺคี น กาลํ กเรยฺยาติ, อามนฺตา.\csromanspacer{} อตฺถิ โส นิยาโม เยน นิยาเมน นิยโต จกฺขุวิญฺญาณสมงฺคี น กาลํ กเรยฺยาติ, นตฺถิ.\csromanspacer{} หญฺจิ นตฺถิ โส นิยาโม เยน นิยาเมน นิยโต, จกฺขุวิญฺญาณสมงฺคี น กาลํ กเรยฺย, โน จ วต เร วตฺตพฺเพ “จกฺขุวิญฺญาณสมงฺคี น กาลํ กเรยฺยา”ติ.}

\nitthitam{ตติยสญฺญาเวทยิตกถา นิฏฺฐิตา.}
\chapterhead{15. ปนฺนรสมวคฺค}
\vspace{\tipitakaparskip}
\csromancenter{\textbf{(154) 10.}\csromanspacer{}\textbf{ อสญฺญสตฺตุปิกกถา}}
\proseitem{735}{\csromansymbolmark{*}สญฺญา เวทยิต นิโรธ สมาปตฺติ อสญฺญ สตฺถุ ปิกาติ, อามนฺตา.\csromanspacer{} อตฺถิ สญฺญาเวทยิตนิโรธํ สมาปนฺนสฺส อโลโภ กุสลมูลํ.\csromanspacer{} อโทโส กุสลมูลํ.\csromanspacer{} อโมโห กุสลมูลํ.\csromanspacer{} สทฺธา.\csromanspacer{} วีริยํ.\csromanspacer{} สติ.\csromanspacer{} สมาธิ.\csromanspacer{} ปญฺญาติ, น เหวํ วตฺตพฺเพ ฯเปฯ.\csromanspacer{} นตฺถิ สญฺญาเวทยิตนิโรธํ สมาปนฺนสฺส อโลโภ กุสลมูลํ.\csromanspacer{} อโทโส กุสลมูลํ ฯเปฯ.\csromanspacer{} ปญฺญาติ, อามนฺตา.\csromanspacer{} หญฺจิ นตฺถิ สญฺญาเวทยิตนิโรธํ สมาปนฺนสฺส อโลโภ}

\vspace{\tipitakaparskip}
\csromancenter{ปนฺนรสมวคฺค กุสลมูลํ.\csromanspacer{} อโทโส กุสลมูลํ.\csromanspacer{} อโมโห กุสลมูลํ.\csromanspacer{} สทฺธา.\csromanspacer{} วีริยํ.\csromanspacer{} สติ.\csromanspacer{} สมาธิ.\csromanspacer{} ปญฺญา, โน จ วต เร วตฺตพฺเพ “สญฺญาเวทยิตนิโรธสมาปตฺติ อสญฺญสตฺตุปิกา”ติ.}
\csromanmatikaanchor{mtk.186}
\csromantocmarkref{subsection}{(155) 11. กมฺมูปจยกถา}{mtk.186}
\bookmark[dest={mtk.186},level=2]{(155) 11. กมฺมูปจยกถา}
\prose{\csromanfolio{377}\csromansymbolmark{*}สญฺญาเวทยิตนิโรธสมาปตฺติ อสญฺญสตฺตุปิกาติ, อามนฺตา.\csromanspacer{} อตฺถิ สญฺญาเวทยิตนิโรธํ สมาปนฺนสฺส ผสฺโส เวทนา สญฺญา เจตนา จิตฺตนฺติ, น เหวํ วตฺตพฺเพ ฯเปฯ.\csromanspacer{} นตฺถิ สญฺญาเวทยิตนิโรธํ สมาปนฺนสฺส ผสฺโส เวทนา สญฺญา เจตนา จิตฺตนฺติ, อามนฺตา.\csromanspacer{} อผสฺสกสฺส มคฺคภาวนา ฯเปฯ.\csromanspacer{} อจิตฺตกสฺส มคฺคภาวนาติ, น เหวํ วตฺตพฺเพ ฯเปฯ.\csromanspacer{} นนุ สผสฺสกสฺส มคฺคภาวนา ฯเปฯ.\csromanspacer{} สจิตฺตกสฺส มคฺคภาวนาติ, อามนฺตา.\csromanspacer{} หญฺจิ สผสฺสกสฺส มคฺคภาวนา ฯเปฯ.\csromanspacer{} สจิตฺตกสฺส มคฺคภาวนา, โน จ วต เร วตฺตพฺเพ “สญฺญาเวทยิตนิโรธสมาปตฺติ อสญฺญสตฺตุปิกา”ติ.}

\prose{\csromansymbolmark{*}สญฺญาเวทยิตนิโรธสมาปตฺติ อสญฺญสตฺตุปิกาติ, อามนฺตา.\csromanspacer{} เย เกจิ สญฺญาเวทยิตนิโรธํ สมาปชฺชนฺติ, สพฺเพ เต อสญฺญสตฺตุปิกาติ, น เหวํ วตฺตพฺเพ ฯเปฯ.}

\proseitem{736}{\csromansymbolmark{+}น วตฺตพฺพํ “สญฺญาเวทยิตนิโรธสมาปตฺติ อสญฺญสตฺตุปิกา”ติ, อามนฺตา.\csromanspacer{} นนุ อิธาปิ อสญฺญี ตตฺราปิ อสญฺญีติ, อามนฺตา.\csromanspacer{} หญฺจิ อิธาปิ อสญฺญี ตตฺราปิ อสญฺญี, เตน วต เร วตฺตพฺเพ “สญฺญาเวทยิตนิโรธสมาปตฺติ อสญฺญสตฺตุปิกา”ติ.}

\nitthitam{อสญฺญสตฺตุปิกกถา นิฏฺฐิตา.}
\chapterhead{15. ปนฺนรสมวคฺค}
\vspace{\tipitakaparskip}
\csromancenter{\textbf{(155) 11.}\csromanspacer{}\textbf{ กมฺมูปจยกถา}}
\proseitem{737}{\csromansymbolmark{*}อญฺญํ กมฺมํ อญฺโญ กมฺมูปจโยติ, อามนฺตา.\csromanspacer{} อญฺโญ ผสฺโส อญฺโญ ผสฺสูปจโย.\csromanspacer{} อญฺญา เวทนา อญฺโญ เวทนูปจโย.\csromanspacer{} อญฺญา สญฺญา อญฺโญ สญฺญูปจโย.\csromanspacer{} อญฺญา เจตนา อญฺโญ เจตนูปจโย.\csromanspacer{} อญฺญํ จิตฺตํ อญฺโญ จิตฺตูปจโย.\csromanspacer{} อญฺญา สทฺธา อญฺโญ สทฺธูปจโย.\csromanspacer{} อญฺญํ วีริยํ อญฺโญ วีริยูปจโย.\csromanspacer{} อญฺญา สติ อญฺโญ \csromanfolio{378}สตูปจโย.\csromanspacer{} อญฺโญ สมาธิ อญฺโญ สมาธูปจโย.\csromanspacer{} อญฺญา ปญฺญา อญฺโญ ปญฺญูปจโย.\csromanspacer{} อญฺโญ ราโค อญฺโญ ราคูปจโย ฯเปฯ.\csromanspacer{} อญฺญํ อโนตฺตปฺปํ อญฺโญ อโนตฺตปฺปูปจโยติ, น เหวํ วตฺตพฺเพ ฯเปฯ.}

\proseitem{738}{\csromansymbolmark{*}อญฺญํ กมฺมํ อญฺโญ กมฺมูปจโยติ, อามนฺตา.\csromanspacer{} กมฺมูปจโย กมฺเมน สหชาโตติ, น เหวํ วตฺตพฺเพ ฯเปฯ.}

\prose{\csromansymbolmark{*}กมฺมูปจโย กมฺเมน สหชาโตติ, อามนฺตา.\csromanspacer{} กุสเลน กมฺเมน สหชาโต กมฺมูปจโย กุสโลติ, น เหวํ วตฺตพฺเพ ฯเปฯ.}

\prose{\csromansymbolmark{*}กุสเลน กมฺเมน สหชาโต กมฺมูปจโย กุสโลติ, อามนฺตา.\csromanspacer{} สุขาย เวทนาย สมฺปยุตฺเตน กมฺเมน สหชาโต กมฺมูปจโย สุขาย เวทนาย สมฺปยุตฺโตติ, น เหวํ วตฺตพฺเพ ฯเปฯ.\csromanspacer{} ทุกฺขาย เวทนาย ฯเปฯ.\csromanspacer{} อทุกฺขมสุขาย เวทนาย สมฺปยุตฺเตน กมฺเมน สหชาโต กมฺมูปจโย อทุกฺขมสุขาย เวทนาย สมฺปยุตฺโตติ, น เหวํ วตฺตพฺเพ ฯเปฯ.}

\proseitem{739}{\csromansymbolmark{*}กมฺมูปจโย กมฺเมน สหชาโตติ, อามนฺตา.\csromanspacer{} อกุสเลน กมฺเมน สหชาโต กมฺมูปจโย อกุสโลติ, น เหวํ วตฺตพฺเพ ฯเปฯ.}

\prose{\csromansymbolmark{*}อกุสเลน กมฺเมน สหชาโต กมฺมูปจโย อกุสโลติ, อามนฺตา.\csromanspacer{} สุขาย เวทนาย สมฺปยุตฺเตน กมฺเมน สหชาโต กมฺมูปจโย สุขาย เวทนาย สมฺปยุตฺโตติ, น เหวํ วตฺตพฺเพ ฯเปฯ.\csromanspacer{} ทุกฺขาย เวทนาย ฯเปฯ.\csromanspacer{} อทุกฺขมสุขาย เวทนาย สมฺปยุตฺเตน กมฺเมน สหชาโต กมฺมูปจโย อทุกฺขมสุขาย เวทนาย สมฺปยุตฺโตติ, น เหวํ วตฺตพฺเพ ฯเปฯ.}

\proseitem{740}{\csromansymbolmark{*}กมฺมํ จิตฺเตน สหชาตํ, กมฺมํ สารมฺมณนฺติ, อามนฺตา.\csromanspacer{} กมฺมูปจโย จิตฺเตน สหชาโต, กมฺมูปจโย สารมฺมโณติ, น เหวํ วตฺตพฺเพ ฯเปฯ.\csromanspacer{} กมฺมูปจโย จิตฺเตน สหชาโต, กมฺมูปจโย อนารมฺมโณติ, อามนฺตา.\csromanspacer{} กมฺมํ จิตฺเตน สหชาตํ, กมฺมํ อนารมฺมณนฺติ, น เหวํ วตฺตพฺเพ ฯเปฯ.}

\chapterheadpageread{15. ปนฺนรสมวคฺค}
\prose{\csromanfolio{379}\csromansymbolmark{*}กมฺมํ จิตฺเตน สหชาตํ, จิตฺตํ ภิชฺชมานํ กมฺมํ ภิชฺชตีติ, อามนฺตา.\csromanspacer{} กมฺมูปจโย จิตฺเตน สหชาโต, จิตฺตํ ภิชฺชมานํ กมฺมูปจโย ภิชฺชตีติ, น เหวํ วตฺตพฺเพ ฯเปฯ.}

\prose{\csromansymbolmark{*}กมฺมูปจโย จิตฺเตน สหชาโต, จิตฺตํ ภิชฺชมานํ กมฺมูปจโย น ภิชฺชตีติ, อามนฺตา.\csromanspacer{} กมฺมํ จิตฺเตน สหชาตํ, จิตฺตํ ภิชฺชมานํ กมฺมํ น ภิชฺชตีติ, น เหวํ วตฺตพฺเพ ฯเปฯ.}

\proseitem{741}{\csromansymbolmark{*}กมฺมมฺหิ กมฺมูปจโยติ, อามนฺตา.\csromanspacer{} ตญฺเญว กมฺมํ โส กมฺมูปจโยติ, น เหวํ วตฺตพฺเพ ฯเปฯ.}

\prose{\csromansymbolmark{*}กมฺมมฺหิ กมฺมูปจโย, กมฺมูปจยโต วิปาโก นิพฺพตฺตตีติ, อามนฺตา.\csromanspacer{} ตญฺเญว กมฺมํ, โส กมฺมูปจโย, โส กมฺมวิปาโกติ, น เหวํ วตฺตพฺเพ ฯเปฯ.}

\prose{\csromansymbolmark{*}กมฺมมฺหิ กมฺมูปจโย, กมฺมูปจยโต วิปาโก นิพฺพตฺตติ, วิปาโก สารมฺมโณติ, อามนฺตา.\csromanspacer{} กมฺมูปจโย สารมฺมโณติ, น เหวํ วตฺตพฺเพ ฯเปฯ.\csromanspacer{} กมฺมูปจโย อนารมฺมโณติ, อามนฺตา.\csromanspacer{} วิปาโก อนารมฺมโณติ, น เหวํ วตฺตพฺเพ ฯเปฯ.}

\csromanmatikaanchor{mtk.187}
\csromanmatikaanchor{mtk.188}
\csromantocmarknopage{chapter}{16. โสฬสมวคฺค}{mtk.187}
\bookmark[dest={mtk.187},level=0]{16. โสฬสมวคฺค}
\csromantocmarkref{subsection}{(156) 1. นิคฺคหกถา}{mtk.188}
\bookmark[dest={mtk.188},level=2]{(156) 1. นิคฺคหกถา}
\proseitem{742}{\csromansymbolmark{*}อญฺญํ กมฺมํ อญฺโญ กมฺมูปจโยติ, อามนฺตา.\csromanspacer{} นนุ วุตฺตํ ภควตา “อิธ ปุณฺณ เอกจฺโจ สพฺยาพชฺฌมฺปิ อพฺยาพชฺฌมฺปิ\footnote{สพฺยาปชฺฌมฺปิ อพฺยาปชฺฌมฺปิ (ก) ม 2. 52 ปิฏฺเฐ ปน ปสฺสิตพฺพํ.} กายสงฺขารํ อภิสงฺขโรติ, สพฺยาพชฺฌมฺปิ อพฺยาพชฺฌมฺปิ วจีสงฺขารํ ฯเปฯ มโนสงฺขารํ อภิสงฺขโรติ, โส สพฺยาพชฺฌมฺปิ อพฺยาพชฺฌมฺปิ กายสงฺขารํ อภิสงฺขริตฺวา สพฺยาพชฺฌมฺปิ อพฺยาพชฺฌมฺปิ วจีสงฺขารํ ฯเปฯ มโนสงฺขารํ อภิสงฺขริตฺวา สพฺยาพชฺฌมฺปิ อพฺยาพชฺฌมฺปิ โลกํ อุปปชฺชติ, ตเมนํ สพฺยาพชฺฌมฺปิ อพฺยาพชฺฌมฺปิ โลกํ อุปปนฺนํ สมานํ สพฺยาพชฺฌาปิ อพฺยาพชฺฌาปิ ผสฺสา ผุสนฺติ, โส สพฺยาพชฺเฌหิปิ อพฺยาพชฺเฌหิปิ ผสฺเสหิ ผุฏฺโฐ สมาโน สพฺยาพชฺฌมฺปิ อพฺยาพชฺฌมฺปิ เวทนํ เวเทติ โวกิณฺณสุขทุกฺขํ, เสยฺยถาปิ มนุสฺสา เอกจฺเจ จ เทวา เอกจฺเจ จ วินิปาติกา.\csromanspacer{} อิติ โข ปุณฺณ ภูตา ภูตสฺส อุปปตฺติ โหติ, ยํ กโรติ เตน อุปปชฺชติ อุปปนฺนเมตํ ผสฺสา ผุสนฺติ, เอวมฺปาหํ ปุณฺณ \csromanfolio{380}กมฺมทายาทา สตฺตาติ วทามี”ติ\footnote{ม 2. 52 ปิฏฺเฐ.} อตฺเถว สุตฺตนฺโตติ, อามนฺตา.\csromanspacer{} เตน หิ น วตฺตพฺพํ “อญฺญํ กมฺมํ อญฺโญ กมฺมูปจโย”ติ.\csromanspacer{} กมฺมูปจยกถา นิฏฺฐิตา.\csromanspacer{} ปนฺนรสโม วคฺโค.}

\titlehead{\textbf{ตสฺสุทฺทานํ}}
\prose{ปจฺจยตา ววตฺถิตา, ปฏิจฺจสมุปฺปาโท, อทฺธา, ขโณ ลโย มุหุตฺตํ, จตฺตาโร อาสวา อนาสวา, โลกุตฺตรานํ ธมฺมานํ ชรามรณํ โลกุตฺตรา, สญฺญาเวทยิตนิโรธสมาปตฺติ โลกุตฺตรา, สญฺญาเวทยิตนิโรธสมาปตฺติ โลกิยา, สญฺญาเวทยิตนิโรธํ สมาปนฺโน กาลงฺกเรยฺย, เสฺวว มคฺโค อสญฺญสตฺตุปปตฺติยา, อญฺญํ กมฺมํ อญฺโญ กมฺมูปจโยติ.\csromanspacer{} ตติโย ปณฺณาสโก.}

\titlehead{\textbf{ตสฺสุทฺทานํ}}
\csromangathameasure{อนุสยา สํวโร กปฺโป มูลํ จ ววตฺถิตาติ.}
\csromangathagroup{\csromangathastanza{อนุสยา สํวโร กปฺโป มูลํ จ ววตฺถิตาติ.}}

\chapterhead{16. โสฬสมวคฺค}
\vspace{\tipitakaparskip}
\csromancenter{\textbf{(156) 1.}\csromanspacer{}\textbf{ นิคฺคหกถา}}
\proseitem{743}{ปโร ปรสฺส จิตฺตํ นิคฺคณฺหาตีติ, อามนฺตา.\csromanspacer{} ปโร ปรสฺส จิตฺตํ มา รชฺชิ, มา ทุสฺสิ, มา มุยฺหิ, มา กิลิสฺสีติ นิคฺคณฺหาตีติ, น เหวํ วตฺตพฺเพ ฯเปฯ.\csromanspacer{} ปโร ปรสฺส จิตฺตํ นิคฺคณฺหาตีติ, อามนฺตา.\csromanspacer{} ปโร ปรสฺส อุปฺปนฺโน ผสฺโส มา นิรุชฺฌีติ นิคฺคณฺหาตีติ, น เหวํ วตฺตพฺเพ ฯเปฯ.\csromanspacer{} ปโร ปรสฺส อุปฺปนฺนา เวทนา ฯเปฯ อุปฺปนฺนา สญฺญา.\csromanspacer{} อุปฺปนฺนา เจตนา.\csromanspacer{} อุปฺปนฺนํ จิตฺตํ.\csromanspacer{} อุปฺปานฺนา สทฺธา.\csromanspacer{} อุปฺปนฺนํ วีริยํ.\csromanspacer{} อุปฺปนฺนา สติ.\csromanspacer{} อุปฺปนฺโน สมาธิ ฯเปฯ อุปฺปนฺนา ปญฺญา มา นิรุชฺฌีติ นิคฺคณฺหาตีติ, น เหวํ วตฺตพฺเพ ฯเปฯ.}

\chapterheadpageread{16. โสฬสมวคฺค}
\csromanmatikaanchor{mtk.189}
\csromantocmarkref{subsection}{(157) 2. ปคฺคหกถา}{mtk.189}
\bookmark[dest={mtk.189},level=2]{(157) 2. ปคฺคหกถา}
\prose{\csromanfolio{381}\csromansymbolmark{*}ปโร ปรสฺส จิตฺตํ นิคฺคณฺหาตีติ, อามนฺตา.\csromanspacer{} ปโร ปรสฺส อตฺถาย ราคํ ปชหติ.\csromanspacer{} โทสํ ปชหติ ฯเปฯ อโนตฺตปฺปํ ปชหตีติ, น เหวํ วตฺตพฺเพ ฯเปฯ.}

\prose{\csromansymbolmark{*}ปโร ปรสฺส จิตฺตํ นิคฺคณฺหาตีติ, อามนฺตา.\csromanspacer{} ปโร ปรสฺส อตฺถาย มคฺคํ ภาเวติ.\csromanspacer{} สติปฏฺฐานํ ภาเวติ ฯเปฯ.\csromanspacer{} โพชฺฌงฺคํ ภาเวตีติ, น เหวํ วตฺตพฺเพ ฯเปฯ.}

\prose{\csromansymbolmark{*}ปโร ปรสฺส จิตฺตํ นิคฺคณาตีติ, อามนฺตา.\csromanspacer{} ปโร ปรสฺส อตฺถาย ทุกฺขํ ปริชานาติ, สมุทยํ ปชหติ, นิโรธํ สจฺฉิกโรติ, มคฺคํ ภาเวตีติ, น เหวํ วตฺตพฺเพ ฯเปฯ.}

\prose{\csromansymbolmark{*}ปโร ปรสฺส จิตฺตํ นิคฺคณฺหาตีติ, อามนฺตา.\csromanspacer{} อญฺโญ อญฺญสฺส การโก, ปรํกตํ สุขํ ทุกฺขํ, อญฺโญ กโรติ อญฺโญ ปฏิสํเวเทตีติ, น เหวํ วตฺตพฺเพ ฯเปฯ.}

\prose{\csromansymbolmark{*}ปโร ปรสฺส จิตฺตํ นิคฺคณฺหาตีติ, อามนฺตา.\csromanspacer{} นนุ วุตฺตํ ภควตา\csromandash{}}

\csromangathameasure{“อตฺตนาว กตํ ปาปํ, อตฺตนา สํกิลิสฺสติ. \\ อตฺตนา อกตํ ปาปํ, อตฺตนาว วิสุชฺฌติ. \\ สุทฺธิ อสุทฺธิ ปจฺจตฺตํ, นาญฺโญ อญฺญํ วิโสธเย”ติ.}
\csromangathagroup{\csromangathastanza{“อตฺตนาว กตํ ปาปํ, อตฺตนา สํกิลิสฺสติ. \\ อตฺตนา อกตํ ปาปํ, อตฺตนาว วิสุชฺฌติ. \\ สุทฺธิ อสุทฺธิ ปจฺจตฺตํ, นาญฺโญ อญฺญํ วิโสธเย”ติ\footnote{ขุ 1. 38 ธมฺมปเท.}.}}

\prose{อตฺเถว สุตฺตนฺโตติ, อามนฺตา.\csromanspacer{} เตน หิ น วตฺตพฺพํ “ปโร ปรสฺส จิตฺตํ นิคฺคณฺหาตี”ติ.}

\proseitem{744}{\csromansymbolmark{+}น วตฺตพฺพํ “ปโร ปรสฺส จิตฺตํ นิคฺคณฺหาตี”ติ, อามนฺตา.\csromanspacer{} นนุ อตฺถิ พลปฺปตฺตา, อตฺถิ วสีภูตาติ, อามนฺตา.\csromanspacer{} หญฺจิ อตฺถิ พลปฺปตฺตา, อตฺถิ วสีภูตา, เตน วต เร วตฺตพฺเพ “ปโร ปรสฺส จิตฺตํ นิคฺคณฺหาตี”ติ.}

\nitthitam{นิคฺคหกถา นิฏฺฐิตา.}
\chapterhead{16. โสฬสมวคฺค}
\vspace{\tipitakaparskip}
\csromancenter{\textbf{(157) 2.}\csromanspacer{}\textbf{ ปคฺคหกถา}}
\proseitem{745}{\csromansymbolmark{*}ปโร ปรสฺส จิตฺตํ ปคฺคณฺหาตีติ, อามนฺตา.\csromanspacer{} ปโร ปรสฺส จิตฺตํ มา รชฺชิ, มา ทุสฺสิ, มา มุยฺหิ, มา กิลิสฺสีติ ปคฺคณฺหาตีติ, น เหวํ วตฺตพฺเพ ฯเปฯ. \csromanfolio{382}ปโร ปรสฺส จิตฺตํ ปคฺคณฺหาตีติ, อามนฺตา.\csromanspacer{} ปโร ปรสฺส อโลภํ กุสลมูลํ ชเนติ.\csromanspacer{} อโทสํ กุสลมูลํ ชเนติ, อโมหํ กุสลมูลํ ชเนติ.\csromanspacer{} สทฺธํ ชเนติ.\csromanspacer{} วีริยํ ชเนติ.\csromanspacer{} สติํ ชเนติ.\csromanspacer{} สมาธิํ ชเนติ.\csromanspacer{} ปญฺญํ ชเนตีติ, น เหวํ วตฺตพฺเพ ฯเปฯ.\csromanspacer{} ปโร ปรสฺส จิตฺตํ ปคฺคณฺหาตีติ, อามนฺตา.\csromanspacer{} ปโร ปรสฺส อุปฺปนฺโน ผสฺโส มา นิรุชฺฌีติ ปคฺคณฺหาตีติ, น เหวํ วตฺตพฺเพ ฯเปฯ.\csromanspacer{} ปโร ปรสฺส อุปฺปนฺนา เวทนา ฯเปฯ อุปฺปนฺนา ปญฺญา มา นิรุชฺฌีติ ปคฺคณฺหาตีติ, น เหวํ วตฺตพฺเพ ฯเปฯ. \csromansymbolmark{*}ปโร ปรสฺส จิตฺตํ ปคฺคณฺหาตีติ, อามนฺตา.\csromanspacer{} ปโร ปรสฺส อตฺถาย ราคํ ปชหติ.\csromanspacer{} โทสํ ปชหติ.\csromanspacer{} โมหํ ปชหติ ฯเปฯ อโนตฺตปฺปํ ปชหตีติ, น เหวํ วตฺตพฺเพ ฯเปฯ. \csromansymbolmark{*}ปโร ปรสฺส จิตฺตํ ปคฺคณฺหาตีติ, อามนฺตา.\csromanspacer{} ปโร ปรสฺส อตฺถาย มคฺคํ ภาเวติ.\csromanspacer{} สติปฏฺฐานํ ภาเวติ ฯเปฯ.\csromanspacer{} โพชฺฌงฺคํ ภาเวตีติ, น เหวํ วตฺตพฺเพ ฯเปฯ. \csromansymbolmark{*}ปโร ปรสฺส จิตฺตํ ปคฺคณฺหาตีติ, อามนฺตา.\csromanspacer{} ปโร ปรสฺส อตฺถาย ทุกฺขํ ปริชานาติ ฯเปฯ มคฺคํ ภาเวตีติ, น เหวํ วตฺตพฺเพ ฯเปฯ. \csromansymbolmark{*}ปโร ปรสฺส จิตฺตํ ปคฺคณฺหาตีติ, อามนฺตา.\csromanspacer{} อญฺโญ อญฺญสฺส การโก, ปรํกตํ สุขํ ทุกฺขํ, อญฺโญ กโรติ อญฺโญ ปฏิสํเวเทตีติ, น เหวํ วตฺตพฺเพ ฯเปฯ. \csromansymbolmark{*}ปโร ปรสฺส จิตฺตํ ปคฺคณฺหาตีติ, อามนฺตา.\csromanspacer{} นนุ วุตฺตํ ภควตา\csromandash{}}

\prose{“อตฺตนาว กตํ ปาปํ ฯเปฯ นาญฺโญ อญฺญํ วิโสธเย”ติ.}

\prose{อตฺเถว สุตฺตนฺโตติ, อามนฺตา.\csromanspacer{} เตน หิ น วตฺตพฺพํ “ปโร ปรสฺส จิตฺตํ ปคฺคณฺหาตี”ติ.}

\proseitem{746}{\csromansymbolmark{+}น วตฺตพฺพํ “ปโร ปรสฺส จิตฺตํ ปคฺคณฺหาตี”ติ, อามนฺตา.\csromanspacer{} นนุ อตฺถิ พลปฺปตฺตา, อตฺถิ วสีภูตาติ, อามนฺตา.\csromanspacer{} หญฺจิ อตฺถิ พลปฺปตฺตา, อตฺถิ วสีภูตา, เตน วต เร วตฺตพฺเพ “ปโร ปรสฺส จิตฺตํ ปคฺคณฺหาตี”ติ.}

\nitthitam{ปคฺคหกถา นิฏฺฐิตา.}
\chapterheadpageread{16. โสฬสมวคฺค \textbf{16. โสฬสมวคฺค}}
\vspace{\tipitakaparskip}
\csromancenter{\textbf{(158) 3.}\csromanspacer{}\textbf{ สุขานุปฺปทานกถา}}
\csromanmatikaanchor{mtk.190}
\csromanmatikaanchor{mtk.191}
\csromantocmarkref{subsection}{(158) 3. สุขานุปฺปทานกถา}{mtk.190}
\bookmark[dest={mtk.190},level=2]{(158) 3. สุขานุปฺปทานกถา}
\csromantocmarkref{subsection}{(159) 4. อธิคยฺหมนสิการกถา}{mtk.191}
\bookmark[dest={mtk.191},level=2]{(159) 4. อธิคยฺหมนสิการกถา}
\proseitem{747}{\csromanfolio{383}\csromansymbolmark{*}ปโร ปรสฺส สุขํ อนุปฺปเทตีติ, อามนฺตา.\csromanspacer{} ปโร ปรสฺส ทุกฺขํ อนุปฺปเทตีติ, น เหวํ วตฺตพฺเพ ฯเปฯ.\csromanspacer{} ปโร ปรสฺส ทุกฺขํ น อนุปฺปเทตีติ, อามนฺตา.\csromanspacer{} ปโร ปรสฺส สุขํ น อนุปฺปเทตีติ, น เหวํ วตฺตพฺเพ ฯเปฯ.\csromanspacer{} ปโร ปรสฺส สุขํ อนุปฺปเทตีติ, อามนฺตา.\csromanspacer{} ปโร ปรสฺส อตฺตโน สุขํ อนุปฺปเทติ, อญฺเญสํ สุขํ อนุปฺปเทติ ตสฺส สุขํ อนุปฺปเทตีติ, น เหวํ วตฺตพฺเพ ฯเปฯ.\csromanspacer{} ปโร ปรสฺส เนวตฺตโน, น อญฺเญสํ, น ตสฺส สุขํ อนุปฺปเทตีติ, อามนฺตา.\csromanspacer{} หญฺจิ ปโร ปรสฺส เนวตฺตโน, น อญฺเญสํ, น ตสฺส สุขํ อนุปฺปเทติ, โน จ วต เร วตฺตพฺเพ “ปโร ปรสฺส สุขํ อนุปฺปเทตี”ติ.}

\prose{\csromansymbolmark{*}ปโร ปรสฺส สุขํ อนุปฺปเทตีติ, อามนฺตา.\csromanspacer{} อญฺโญ อญฺญสฺส การโก, ปรํกตํ สุขํ ทุกฺขํ, อญฺโญ กโรติ อญฺโญ ปฏิสํเวเทตีติ, น เหวํ วตฺตพฺเพ ฯเปฯ.}

\proseitem{748}{\csromansymbolmark{+}น วตฺตพฺพํ “ปโร ปรสฺส สุขํ อนุปฺปเทตี”ติ, อามนฺตา.\csromanspacer{} นนุ อายสฺมา อุทายี เอตทโวจ “พหูนํ วต โน ภควา ทุกฺขธมฺมานํ อปหตฺตา, พหูนํ วต โน ภควา สุขธมฺมานํ อุปหตฺตา, พหูนํ วต โน ภควา อกุสลานํ ธมฺมานํ อปหตฺตา, พหูนํ วต โน ภควา กุสลานํ ธมฺมานํ อุปหตฺตา”ติ\footnote{ม 2. 111 ปิฏฺเฐ ลฏุกิโกปเม.} อตฺเถว สุตฺตนฺโตติ, อามนฺตา.\csromanspacer{} เตน หิ ปโร ปรสฺส สุขํ อนุปฺปเทตีติ.}

\nitthitam{สุขานุปฺปทานกถา นิฏฺฐิตา.}
\chapterhead{16. โสฬสมวคฺค}
\vspace{\tipitakaparskip}
\csromancenter{\textbf{(159) 4.}\csromanspacer{}\textbf{ อธิคยฺหมนสิการกถา}}
\proseitem{749}{\csromansymbolmark{*}อธิคยฺห มนสิ กโรตีติ, อามนฺตา.\csromanspacer{} เตน จิตฺเตน ตํ จิตฺตํ ปชานาตีติ, น เหวํ วตฺตพฺเพ ฯเปฯ.\csromanspacer{} เตน จิตฺเตน ตํ จิตฺตํ ปชานาตีติ, \csromanfolio{384}อามนฺตา.\csromanspacer{} เตน จิตฺเตน ตํ จิตฺตํ “จิตฺตนฺ”ติ ปชานาตีติ, น เหวํ วตฺตพฺเพ ฯเปฯ.\csromanspacer{} เตน จิตฺเตน ตํ จิตฺตํ “จิตฺตนฺ”ติ ปชานาตีติ, อามนฺตา.\csromanspacer{} ตํ จิตฺตํ ตสฺส จิตฺตสฺส อารมฺมณนฺติ, น เหวํ วตฺตพฺเพ ฯเปฯ.}

\prose{\csromansymbolmark{*}ตํ จิตฺตํ ตสฺส จิตฺตสฺส อารมฺมณนฺติ, อามนฺตา.\csromanspacer{} เตน ผสฺเสน ตํ ผสฺสํ ผุสติ, ตาย เวทนาย ฯเปฯ.\csromanspacer{} ตาย สญฺญาย.\csromanspacer{} ตาย เจตนาย.\csromanspacer{} เตน จิตฺเตน.\csromanspacer{} เตน วิตกฺเกน.\csromanspacer{} เตน วิจาเรน.\csromanspacer{} ตาย ปีติยา.\csromanspacer{} ตาย สติยา.\csromanspacer{} ตาย ปญฺญาย ตํ ปญฺญํ ปชานาตีติ, น เหวํ วตฺตพฺเพ ฯเปฯ.}

\proseitem{750}{\csromansymbolmark{*}อตีตํ “อตีตนฺ”ติ มนสิกโรนฺโต อนาคตํ “อนาคตนฺ”ติ มนสิ กโรตีติ, น เหวํ วตฺตพฺเพ ฯเปฯ.\csromanspacer{} อตีตํ “อตีตนฺ”ติ มนสิกโรนฺโต อนาคตํ “อนาคตนฺ”ติ มนสิ กโรตีติ, อามนฺตา.\csromanspacer{} ทฺวินฺนํ ผสฺสานํ ฯเปฯ.\csromanspacer{} ทฺวินฺนํ จิตฺตานํ สโมธานํ โหตีติ, น เหวํ วตฺตพฺเพ ฯเปฯ.}

\prose{\csromansymbolmark{*}อตีตํ “อตีตนฺ”ติ มนสิกโรนฺโต ปจฺจุปฺปนฺนํ “ปจฺจุปฺปนฺนนฺ”ติ มนสิ กโรตีติ, น เหวํ วตฺตพฺเพ ฯเปฯ.\csromanspacer{} อตีตํ “อตีตนฺ”ติ มนสิกโรนฺโต ปจฺจุปฺปนฺนํ “ปจฺจุปฺปนฺนนฺ”ติ มนสิ กโรตีติ, อามนฺตา.\csromanspacer{} ทฺวินฺนํ ผสฺสานํ ฯเปฯ.\csromanspacer{} ทฺวินฺนํ จิตฺตานํ สโมธานํ โหตีติ, น เหวํ วตฺตพฺเพ ฯเปฯ.}

\prose{\csromansymbolmark{*}อตีตํ “อตีตนฺ”ติ มนสิกโรนฺโต อนาคตํ “อนาคตนฺ”ติ มนสิ กโรติ, ปจฺจุปฺปนฺนํ “ปจฺจุปฺปนฺนนฺ”ติ มนสิ กโรตีติ, น เหวํ วตฺตพฺเพ ฯเปฯ.\csromanspacer{} อตีตํ “อตีตนฺ”ติ มนสิกโรนฺโต อนาคตํ “อนาคตนฺ”ติ มนสิ กโรติ, ปจฺจุปฺปนฺนํ “ปจฺจุปฺปนฺนนฺ”ติ มนสิ กโรตีติ, อามนฺตา.\csromanspacer{} ติณฺณํ ผสฺสานํ ฯเปฯ.\csromanspacer{} ติณฺณํ จิตฺตานํ สโมธานํ โหตีติ, น เหวํ วตฺตพฺเพ ฯเปฯ.}

\proseitem{751}{\csromansymbolmark{*}อนาคตํ “อนาคตนฺ”ติ มนสิกโรนฺโต อตีตํ “อตีตนฺ”ติ มนสิ กโรตีติ, น เหวํ วตฺตพฺเพ ฯเปฯ.\csromanspacer{} อนาคตํ “อนาคตนฺ”ติ มนสิกโรนฺโต, อตีตํ “อตีตนฺ”ติ มนสิ กโรตีติ, อามนฺตา.\csromanspacer{} ทฺวินฺนํ ผสฺสานํ ฯเปฯ.\csromanspacer{} ทฺวินฺนํ จิตฺตานํ สโมธานํ โหตีติ, น เหวํ วตฺตพฺเพ ฯเปฯ.}

\prose{\csromansymbolmark{*}อนาคตํ “อนาคตนฺ”ติ มนสิกโรนฺโต ปจฺจุปฺปนฺนํ “ปจฺจุปฺปนฺนนฺ”ติ มนสิ กโรตีติ, น เหวํ วตฺตพฺเพ ฯเปฯ.\csromanspacer{} อนาคตํ “อนาคตนฺ”ติ มนสิกโรนฺโต ปจฺจุปฺปนฺนํ “ปจฺจุปฺปนฺนนฺ”ติ มนสิ กโรตีติ, อามนฺตา.\csromanspacer{} ทฺวินฺนํ ผสฺสานํ ฯเปฯ.\csromanspacer{} ทฺวินฺนํ จิตฺตานํ สโมธานํ โหตีติ, น เหวํ วตฺตพฺเพ ฯเปฯ.}

\chapterheadpageread{16. โสฬสมวคฺค}
\prose{\csromanfolio{385}\csromansymbolmark{*}อนาคตํ “อนาคตนฺ”ติ มนสิกโรนฺโต อตีตํ “อตีตนฺ”ติ มนสิ กโรติ, ปจฺจุปฺปนฺนํ “ปจฺจุปฺปนฺนนฺ”ติ มนสิ กโรตีติ, น เหวํ วตฺตพฺเพ ฯเปฯ.\csromanspacer{} อนาคตํ “อนาคตนฺ”ติ มนสิกโรนฺโต อตีตํ “อตีตนฺ”ติ มนสิ กโรติ, ปจฺจุปฺปนฺนํ “ปจฺจุปฺปนฺนนฺ”ติ มนสิ กโรตีติ, อามนฺตา.\csromanspacer{} ติณฺณํ ผสฺสานํ ฯเปฯ.\csromanspacer{} ติณฺณํ จิตฺตานํ สโมธานํ โหตีติ, น เหวํ วตฺตพฺเพ ฯเปฯ.}

\proseitem{752}{\csromansymbolmark{*}ปจฺจุปฺปนฺนํ “ปจฺจุปฺปนฺนนฺ”ติ มนสิกโรนฺโต อตีตํ “อตีตนฺ”ติ มนสิ กโรตีติ.\csromanspacer{} น เหวํ วตฺตพฺเพ ฯเปฯ.\csromanspacer{} ปจฺจุปฺปนฺนํ “ปจฺจุปฺปนฺนนฺ”ติ มนสิกโรนฺโต อตีตํ “อตีตนฺ”ติ มนสิ กโรตีติ, อามนฺตา.\csromanspacer{} ทฺวินฺนํ ผสฺสานํ ฯเปฯ.\csromanspacer{} ทฺวินฺนํ จิตฺตานํ สโมธานํ โหตีติ, น เหวํ วตฺตพฺเพ ฯเปฯ.}

\prose{\csromansymbolmark{*}ปจฺจุปฺปนฺนํ “ปจฺจุปฺปนฺนนฺ”ติ มนสิกโรนฺโต อนาคตํ “อนาคตนฺ”ติ มนสิ กโรตีติ, น เหวํ วตฺตพฺเพ ฯเปฯ.\csromanspacer{} ปจฺจุปฺปนฺนํ “ปจฺจุปฺปนฺนนฺ”ติ มนสิกโรนฺโต อนาคตํ “อนาคตนฺ”ติ มนสิ กโรตีติ, อามนฺตา.\csromanspacer{} ทฺวินฺนํ ผสฺสานํ ฯเปฯ.\csromanspacer{} ทฺวินฺนํ จิตฺตานํ สโมธานํ โหตีติ, น เหวํ วตฺตพฺเพ ฯเปฯ.}

\prose{\csromansymbolmark{*}ปจฺจุปฺปนฺนํ “ปจฺจุปฺปนฺนนฺ”ติ มนสิกโรนฺโต อตีตํ “อตีตนฺ”ติ มนสิ กโรติ, อนาคตํ “อนาคตนฺ”ติ มนสิ กโรตีติ, น เหวํ วตฺตพฺเพ ฯเปฯ.\csromanspacer{} ปจฺจุปฺปนฺนํ “ปจฺจุปฺปนฺนนฺ”ติ มนสิกโรนฺโต อตีตํ “อตีตนฺ”ติ มนสิ กโรติ, อนาคตํ “อนาคตนฺ”ติ มนสิ กโรตีติ, อามนฺตา.\csromanspacer{} ติณฺณํ ผสฺสานํ ฯเปฯ.\csromanspacer{} ติณฺณํ จิตฺตานํ สโมธานํ โหตีติ, น เหวํ วตฺตพฺเพ ฯเปฯ.}

\proseitem{753}{\csromansymbolmark{+}น วตฺตพฺพํ “อธิคยฺห มนสิ กโรตี”ติ, อามนฺตา.\csromanspacer{} นนุ วุตฺตาํ ภควตา\csromandash{}}

\csromangathameasure{“สพฺเพ สงฺขารา อนิจฺจา”ติ, ยทา ปญฺญาย ปสฺสติ. \\ อถ นิพฺพินฺทติ ทุกฺเข, เอส มคฺโค วิสุทฺธิยา. \\ “สพฺเพ สงฺขารา ทุกฺขา”ติ, ยทา ปญฺญาย ปสฺสติ. \\ อถ นิพฺพินฺทติ ทุกฺเข, เอส มคฺโค วิสุทฺธิยา. \\ “สพฺเพ ธมฺมา อนตฺตา”ติ, ยทา ปญฺญาย ปสฺสติ. \\ อถ นิพฺพินฺทติ ทุกฺเข, เอส มคฺโค วิสุทฺธิยาติ.}
\csromangathagroup{\csromangathastanza{“สพฺเพ สงฺขารา อนิจฺจา”ติ, ยทา ปญฺญาย ปสฺสติ. \\ อถ นิพฺพินฺทติ ทุกฺเข, เอส มคฺโค วิสุทฺธิยา.}\csromangathastanza{“สพฺเพ สงฺขารา ทุกฺขา”ติ, ยทา ปญฺญาย ปสฺสติ. \\ อถ นิพฺพินฺทติ ทุกฺเข, เอส มคฺโค วิสุทฺธิยา.}\csromangathastanza{“สพฺเพ ธมฺมา อนตฺตา”ติ, ยทา ปญฺญาย ปสฺสติ. \\ อถ นิพฺพินฺทติ ทุกฺเข, เอส มคฺโค วิสุทฺธิยาติ\footnote{ขุ 1. 53 ธมฺมปเท.}.}}

\prose{อตฺเถว สุตฺตนฺโตติ, อามนฺตา.\csromanspacer{} เตน หิ อธิคยฺห มนสิ กโรตีติ.}

\nitthitam{อธิคยฺหมนสิการกถา นิฏฺฐิตา.}
\chapterheadpageread{16. โสฬสมวคฺค}
\vspace{\tipitakaparskip}
\csromancenter{\textbf{(160) 5.}\csromanspacer{}\textbf{ รูปํเหตูติกถา}}
\csromanmatikaanchor{mtk.192}
\csromantocmarkref{subsection}{(160) 5. รูปํเหตูติกถา}{mtk.192}
\bookmark[dest={mtk.192},level=2]{(160) 5. รูปํเหตูติกถา}
\proseitem{754}{\csromanfolio{386}\csromansymbolmark{*}รูปํ เหตูติ, อามนฺตา.\csromanspacer{} อโลโภ เหตูติ, น เหวํ วตฺตพฺเพ ฯเปฯ.\csromanspacer{} อโทโส เหตุ ฯเปฯ.\csromanspacer{} อโมโห เหตุ.\csromanspacer{} โลโภ เหตุ.\csromanspacer{} โทโส เหตุ.\csromanspacer{} โมโห เหตูติ, น เหวํ วตฺตพฺเพ ฯเปฯ.}

\prose{\csromansymbolmark{*}รูปํ เหตูติ, อามนฺตา.\csromanspacer{} สารมฺมณํ, อตฺถิ ตสฺส อาวฏฺฏนา ฯเปฯ ปณิธีติ, น เหวํ วตฺตพฺเพ ฯเปฯ.\csromanspacer{} นานุ อนารมฺมณํ, นตฺถิ ตสฺส อาวฏฺฏนา ฯเปฯ ปณิธีติ, อามนฺตา.\csromanspacer{} หญฺจิ อนารมฺมณํ, นตฺถิ ตสฺส อาวฏฺฏนา ฯเปฯ ปณิธิ, โน จ วต เร วตฺตพฺเพ “รูปํ เหตู”ติ.}

\proseitem{755}{\csromansymbolmark{*}อโลโภ เหตุ สารมฺมโณ, อตฺถิ ตสฺส อาวฏฺฏนา ฯเปฯ ปณิธีติ, อามนฺตา.\csromanspacer{} รูปํ เหตุ สารมฺมณํ, อตฺถิ ตสฺส อาวฏฺฏนา ฯเปฯ ปณิธีติ, น เหวํ วตฺตพฺเพ ฯเปฯ.\csromanspacer{} อโทโส เหตุ.\csromanspacer{} อโมโห เหตุ.\csromanspacer{} โลโภ เหตุ.\csromanspacer{} โทโส เหตุ.\csromanspacer{} โมโห เหตุ สารมฺมโณ, อตฺถิ ตสฺส อาวฏฺฏนา ฯเปฯ ปณิธีติ, อามนฺตา.\csromanspacer{} รูปํ เหตุ สารมฺมณํ, อตฺถิ ตสฺส อาวฏฺฏนา ฯเปฯ ปณิธีติ, น เหวํ วตฺตพฺเพ ฯเปฯ.}

\prose{\csromansymbolmark{*}รูปํ เหตุ อนารมฺมณํ, นตฺถิ ตสฺส อาวฏฺฏนา ฯเปฯ ปณิธีติ, อามนฺตา.\csromanspacer{} อโลโภ เหตุ อนารมฺมโณ, นตฺถิ ตสฺส อาวฏฺฏนา ฯเปฯ ปณิธีติ, น เหวํ วตฺตพฺเพ ฯเปฯ.\csromanspacer{} รูปํ เหตุ อนารมฺมณํ, นตฺถิ ตสฺส อาวฏฺฏนา ฯเปฯ ปณิธีติ, อามนฺตา.\csromanspacer{} อโทโส เหตุ.\csromanspacer{} อโมโห เหตุ.\csromanspacer{} โลโภ เหตุ.\csromanspacer{} โทโส เหตุ.\csromanspacer{} โมโห เหตุ อนารมฺมโณ, นตฺถิ ตสฺส อาวฏฺฏนา ฯเปฯ ปณิธีติ, น เหวํ วตฺตพฺเพ ฯเปฯ.}

\proseitem{756}{\csromansymbolmark{+}น วตฺตพฺพํ “รูปํ เหตู”ติ, อามนฺตา.\csromanspacer{} นนุ มหาภูตา อุปาทายรูปานํ\footnote{อุปาทารูปานํ (สี, อิ, ก)} อุปาทายเหตูติ, อามนฺตา.\csromanspacer{} หญฺจิ มหาภูตา อุปาทายรูปานํ อุปาทายเหตุ.\csromanspacer{} เตน วต เร วตฺตพฺเพ “รูปํ เหตู”ติ.}

\nitthitam{รูปํ เหตูติกถา นิฏฺฐิตา.}
\chapterheadpageread{16. โสฬสมวคฺค \textbf{16. โสฬสมวคฺค}}
\vspace{\tipitakaparskip}
\csromancenter{\textbf{(161) 6.}\csromanspacer{}\textbf{ รูปํสเหตุกนฺติกถา}}
\csromanmatikaanchor{mtk.193}
\csromantocmarkref{subsection}{(161) 6. รูปํสเหตุกนฺติกถา}{mtk.193}
\bookmark[dest={mtk.193},level=2]{(161) 6. รูปํสเหตุกนฺติกถา}
\proseitem{757}{\csromanfolio{387}\csromansymbolmark{*}รูปํ สเหตุกนฺติ, อามนฺตา.\csromanspacer{} อโลภเหตุนาติ, น เหวํ วตฺตพฺเพ ฯเปฯ.\csromanspacer{} อโทสเหตุนาติ ฯเปฯ.\csromanspacer{} อโมหเหตุนาติ ฯเปฯ.\csromanspacer{} โลภเหตุนา ฯเปฯ.\csromanspacer{} โทสเหตุนา ฯเปฯ.\csromanspacer{} โมหเหตุนาติ, น เหวํ วตฺตพฺเพ ฯเปฯ.}

\prose{\csromansymbolmark{*}รูปํ สเหตุกนฺติ, อามนฺตา.\csromanspacer{} สารมฺมณํ, อตฺถิ ตสฺส อาวฏฺฏนา ฯเปฯ ปณิธีติ, น เหวํ วตฺตพฺเพ ฯเปฯ.\csromanspacer{} นนุ อนารมฺมณํ, นตฺถิ ตสฺส อาวฏฺฏนา ฯเปฯ ปณิธีติ, อามนฺตา.\csromanspacer{} หญฺจิ อนารมฺมณํ, นตฺถิ ตสฺส อาวฏฺฏนา ฯเปฯ ปณิธิ, โน จ วต เร วตฺตพฺเพ “รูปํ สเหตุกนฺ”ติ.}

\proseitem{758}{\csromansymbolmark{*}อโลโภ สเหตุโก สารมฺมโณ, อตฺถิ ตสฺส อาวฏฺฏนา ฯเปฯ ปณิธีติ, อามนฺตา.\csromanspacer{} รูปํ สเหตุกํ สารมฺมณํ, อตฺถิ ตสฺส อาวฏฺฏนา ฯเปฯ ปณิธีติ, น เหวํ วตฺตพฺเพ ฯเปฯ.\csromanspacer{} อโทโส สเหตุโก ฯเปฯ.\csromanspacer{} อโมโห.\csromanspacer{} สทฺธา.\csromanspacer{} วีริยํ.\csromanspacer{} สติ.\csromanspacer{} สมาธิ.\csromanspacer{} ปญฺญา.\csromanspacer{} โลโภ.\csromanspacer{} โทโส.\csromanspacer{} โมโห.\csromanspacer{} มาโน.\csromanspacer{} ทิฏฺฐิ.\csromanspacer{} วิจิกิจฺฉา.\csromanspacer{} ถินํ.\csromanspacer{} อุทฺธจฺจํ.\csromanspacer{} อหิริกํ.\csromanspacer{} อโนตฺตปฺปํ สเหตุกํ สารมฺมณํ, อตฺถิ ตสฺส อาวฏฺฏนา ฯเปฯ ปณิธีติ, อามนฺตา.\csromanspacer{} รูปํ สเหตุกํ สารมฺมณํ, อตฺถิ ตสฺส อาวฏฺฏนา ฯเปฯ ปณิธีติ, น เหวํ วตฺตพฺเพ ฯเปฯ.}

\prose{\csromansymbolmark{*}รูปํ สเหตุกํ อนารมฺมณํ, นตฺถิ ตสฺส อาวฏฺฏนา ฯเปฯ ปณิธีติ, อามนฺตา.\csromanspacer{} อโลโภ สเหตุโก อนารมฺมโณ, นตฺถิ ตสฺส อาวฏฺฏนา ฯเปฯ ปณิธีติ, น เหวํ วตฺตพฺเพ ฯเปฯ.\csromanspacer{} รูปํ สเหตุกํ อนารมฺมณํ, นตฺถิ ตสฺส อาวฏฺฏนา ฯเปฯ ปณิธีติ, อามนฺตา.\csromanspacer{} อโทโส สเหตุโก ฯเปฯ.\csromanspacer{} อโนตฺตปฺปํ สเหตุกํ อนารมฺมณํ, นตฺถิ ตสฺส อาวฏฺฏนา ฯเปฯ ปณิธีติ, น เหวํ วตฺตพฺเพ ฯเปฯ.}

\proseitem{759}{\csromansymbolmark{+}น วตฺตพฺพํ “รูปํ สเหตุกนฺ”ติ, อามนฺตา.\csromanspacer{} นนุ รูปํ สปฺปจฺจยนฺติ, อามนฺตา.\csromanspacer{} หญฺจิ รูปํ สปฺปจฺจยํ.\csromanspacer{} เตน วต เร วตฺตพฺเพ “รูปํ สเหตุกนฺ”ติ.}

\nitthitam{รูปํ สเหตุกนฺติกถา นิฏฺฐิตา.}
\chapterheadpageread{16. โสฬสมวคฺค}
\vspace{\tipitakaparskip}
\csromancenter{\textbf{(162) 7.}\csromanspacer{}\textbf{ รูปํกุสลากุสลนฺติกถา}}
\csromanmatikaanchor{mtk.194}
\csromantocmarkref{subsection}{(162) 7. รูปํกุสลากุสลนฺติกถา}{mtk.194}
\bookmark[dest={mtk.194},level=2]{(162) 7. รูปํกุสลากุสลนฺติกถา}
\proseitem{760}{\csromanfolio{388}\csromansymbolmark{*}รูปํ กุสลนฺติ, อามนฺตา.\csromanspacer{} สารมฺมณํ, อตฺถิ ตสฺส อาวฏฺฏนา ฯเปฯ ปณิธีติ, น เหวํ วตฺตพฺเพ ฯเปฯ.\csromanspacer{} นนุ อนารมฺมณํ, นตฺถิ ตสฺส อาวฏฺฏนา ฯเปฯ ปณิธีติ, อามนฺตา.\csromanspacer{} หญฺจิ อนารมฺมณํ, นตฺถิ ตสฺส อาวฏฺฏนา ฯเปฯ ปณิธิ, โน จ วต เร วตฺตพฺเพ “รูปํ กุสลนฺ”ติ.}

\proseitem{761}{\csromansymbolmark{*}อโลโภ กุสโล สารมฺมโณ, อตฺถิ ตสฺส อาวฏฺฏนา ฯเปฯ ปณิธีติ, อามนฺตา.\csromanspacer{} รูปํ กุสลํ สารมฺมณํ, อตฺถิ ตสฺส อาวฏฺฏนา ฯเปฯ ปณิธีติ, น เหวํ วตฺตพฺเพ ฯเปฯ.}

\prose{\csromansymbolmark{*}อโทโส กุสโล ฯเปฯ.\csromanspacer{} อโมโห กุสโล ฯเปฯ.\csromanspacer{} สทฺธา.\csromanspacer{} วีริยํ.\csromanspacer{} สติ.\csromanspacer{} สมาธิ ฯเปฯ.\csromanspacer{} ปญฺญา กุสลา สารมฺมณา, อตฺถิ ตาย อาวฏฺฏนา ฯเปฯ ปณิธีติ, อามนฺตา.\csromanspacer{} รูปํ กุสลํ สารมฺมณํ, อตฺถิ ตสฺส อาวฏฺฏนา ฯเปฯ ปณิธีติ, น เหวํ วตฺตพฺเพ ฯเปฯ.}

\prose{\csromansymbolmark{*}รูปํ กุสลํ อนารมฺมณํ, นตฺถิ ตสฺส อาวฏฺฏนา ฯเปฯ ปณิธีติ, อามนฺตา.\csromanspacer{} อโลโภ กุสโล อนารมฺมโณ, นตฺถิ ตสฺส อาวฏฺฏนา ฯเปฯ ปณิธีติ, น เหวํ วตฺตพฺเพ ฯเปฯ.}

\prose{\csromansymbolmark{*}รูปํ กุสลํ อนารมฺมณํ, นตฺถิ ตสฺส อาวฏฺฏนา ฯเปฯ ปณิธีติ, อามนฺตา.\csromanspacer{} อโทโส กุสโล ฯเปฯ.\csromanspacer{} ปญฺญา กุสลา อนารมฺมณา, นตฺถิ ตาย อาวฏฺฏนา ฯเปฯ ปณิธีติ, น เหวํ วตฺตพฺเพ ฯเปฯ.}

\proseitem{762}{\csromansymbolmark{*}รูปํ อกุสลนฺติ, อามนฺตา.\csromanspacer{} สารมฺมณํ, อตฺถิ ตสฺส อาวฏฺฏนา ฯเปฯ ปณิธีติ, น เหวํ วตฺตพฺเพ ฯเปฯ.\csromanspacer{} นนุ อนารมฺมณํ, นตฺถิ ตสฺส อาวฏฺฏนา ฯเปฯ ปณิธีติ, อามนฺตา.\csromanspacer{} หญฺจิ อนารมฺมณํ, นตฺถิ ตสฺส อาวฏฺฏนา ฯเปฯ ปณิธิ, โน จ วต เร วตฺตพฺเพ “รูปํ อกุสลนฺ”ติ ฯเปฯ.}

\proseitem{763}{\csromansymbolmark{*}โลโภ อกุสโล สารมฺมโณ, อตฺถิ ตสฺส อาวฏฺฏนา ฯเปฯ ปณิธีติ, อามนฺตา.\csromanspacer{} รูปํ อกุสลํ สารมฺมณํ, อตฺถิ ตสฺส อาวฏฺฏนา ฯเปฯ ปณิธีติ, น เหวํ วตฺตพฺเพ ฯเปฯ.\csromanspacer{} โทโส.\csromanspacer{} โมโห.\csromanspacer{} มาโน ฯเปฯ.\csromanspacer{} อโนตฺตปฺปํ อกุสลํ สารมฺมณํ, อตฺถิ ตสฺส อาวฏฺฏนา ฯเปฯ}

\vspace{\tipitakaparskip}
\csromancenter{โสฬสมวคฺค ปณิธีติ, อามนฺตา.\csromanspacer{} รูปํ อกุสลํ สารมฺมณํ, อตฺถิ ตสฺส อาวฏฺฏนา ฯเปฯ ปณิธีติ, น เหวํ วตฺตพฺเพ ฯเปฯ.}
\csromanmatikaanchor{mtk.195}
\csromantocmarkref{subsection}{(163) 8. รูปํวิปาโกติกถา}{mtk.195}
\bookmark[dest={mtk.195},level=2]{(163) 8. รูปํวิปาโกติกถา}
\prose{\csromanfolio{389}\csromansymbolmark{*}รูปํ อกุสลํ อนารมฺมณํ, นตฺถิ ตสฺส อาวฏฺฏนา ฯเปฯ ปณิธีติ, อามนฺตา.\csromanspacer{} โลโภ อกุสโล อนารมฺมโณ, นตฺถิ ตสฺส อาวฏฺฏนา ฯเปฯ ปณิธีติ, น เหวํ วตฺตพฺเพ ฯเปฯ.\csromanspacer{} รูปํ อกุสลํ อนารมฺมณํ, นตฺถิ ตสฺส อาวฏฺฏนา ฯเปฯ ปณิธีติ, อามนฺตา.\csromanspacer{} โทโส.\csromanspacer{} โมโห ฯเปฯ.\csromanspacer{} อโนตฺตปฺปํ อกุสลํ อนารมฺมณํ, นตฺถิ ตสฺส อาวฏฺฏนา ฯเปฯ ปณิธีติ, น เหวํ วตฺตพฺเพ ฯเปฯ.}

\proseitem{746}{\csromansymbolmark{+}น วตฺตพฺพํ “รูปํ กุสลมฺปิ อกุสลมฺปี”ติ, อามนฺตา.\csromanspacer{} นนุ กายกมฺมํ วจีกมฺมํ กุสลมฺปิ อกุสลมฺปีติ, อามนฺตา.\csromanspacer{} หญฺจิ กายกมฺมํ วจีกมฺมํ กุสลมฺปิ อกุสลมฺปิ.\csromanspacer{} เตน วต เร วตฺตพฺเพ “รูปํ กุสลมฺปิ อกุสลมฺปี”ติ.}

\nitthitam{รูปํ กุสลากุสลนฺติกถา นิฏฺฐิตา.}
\chapterhead{16. โสฬสมวคฺค}
\vspace{\tipitakaparskip}
\csromancenter{\textbf{(163) 8.}\csromanspacer{}\textbf{ รูปํวิปาโกติกถา}}
\proseitem{765}{\csromansymbolmark{*}รูปํ วิปาโกติ, อามนฺตา.\csromanspacer{} รูปํ สุขเวทนิยํ ทุกฺขเวทนิยํ อทุกฺขมสุขเวทนิยํ, สุขาย เวทนาย สมฺปยุตฺตํ ทุกฺขาย เวทนาย สมฺปยุตฺตํ อทุกฺขมสุขาย เวทนาย สมฺปยุตฺตํ, ผสฺเสน สมฺปยุตฺตํ ฯเปฯ.\csromanspacer{} จิตฺเตน สมฺปยุตฺตํ สารมฺมณํ, อตฺถิ ตสฺส อาวฏฺฏนา ฯเปฯ ปณิธีติ, น เหวํ วตฺตพฺเพ ฯเปฯ.\csromanspacer{} นนุ น สุขเวทนิยํ.\csromanspacer{} น ทุกฺขเวทนิยํ ฯเปฯ อนารมฺมณํ, นตฺถิ ตสฺส อาวฏฺฏนา ฯเปฯ ปณิธีติ, อามนฺตา.\csromanspacer{} หญฺจิ น สุขเวทนิยํ.\csromanspacer{} น ทุกฺขเวทนิยํ ฯเปฯ อนารมฺมณํ, นตฺถิ ตสฺส อาวฏฺฏนา ฯเปฯ ปณิธิ, โน จ วต เร วตฺตพฺเพ “รูปํ วิปาโก”ติ.}

\proseitem{766}{\csromansymbolmark{*}ผสฺโส วิปาโก.\csromanspacer{} ผสฺโส สุขเวทนิโย.\csromanspacer{} ทุกฺขเวทนิโย ฯเปฯ สารมฺมโณ, อตฺถิ ตสฺส อาวฏฺฏนา ฯเปฯ ปณิธีติ, อามนฺตา.\csromanspacer{} รูปํ วิปาโก รูปํ สุขเวทนิยํ.\csromanspacer{} ทุกฺขเวทนิยํ ฯเปฯ สารมฺมณํ, อตฺถิ ตสฺส อาวฏฺฏนา ฯเปฯ ปณิธีติ, น เหวํ วตฺตพฺเพ ฯเปฯ.}

\csromanmatikaanchor{mtk.196}
\csromantocmarkref{subsection}{(164) 9. รูปํรูปาวจรารูปาวจรนฺติกถา}{mtk.196}
\bookmark[dest={mtk.196},level=2]{(164) 9. รูปํรูปาวจรารูปาวจรนฺติกถา}
\prose{\csromanfolio{390}\csromansymbolmark{*}รูปํ วิปาโก, รูปํ น สุขเวทนิยํ.\csromanspacer{} น ทุกฺขเวทนิยํ ฯเปฯ อนารมฺมณํ, นตฺถิ ตสฺส อาวฏฺฏนา ฯเปฯ ปณิธีติ, อามนฺตา.\csromanspacer{} ผสฺโส วิปาโก, ผสฺโส น สุขเวทนิโย.\csromanspacer{} น ทุกฺขเวทนิโย ฯเปฯ อนารมฺมโณ, นตฺถิ ตสฺส อาวฏฺฏนา ฯเปฯ ปณิธีติ, น เหวํ วตฺตพฺเพ ฯเปฯ.}

\proseitem{767}{\csromansymbolmark{+}น วตฺตพฺพํ “รูปํ วิปาโก”ติ, อามนฺตา.\csromanspacer{} นนุ กมฺมสฺส กตตฺตา อุปฺปนฺนา จิตฺตเจตสิกา ธมฺมา วิปาโกติ, อามนฺตา.\csromanspacer{} หญฺจิ กมฺมสฺส กตตฺตา อุปฺปนฺนา จิตฺตเจตสิกา ธมฺมา วิปาโก.\csromanspacer{} เตน วต เร วตฺตพฺเพ “กมฺมสฺส กตตฺตา อุปฺปนฺนํ รูปํ วิปาโก”ติ.}

\nitthitam{รูปํ วิปาโกติกถา นิฏฺฐิตา.}
\chapterhead{16. โสฬสมวคฺค}
\vspace{\tipitakaparskip}
\csromancenter{\textbf{(164) 9.}\csromanspacer{}\textbf{ รูปํรูปาวจรารูปาวจรนฺติกถา}}
\proseitem{768}{\csromansymbolmark{*}อตฺถิ รูปํ รูปาวจรนฺติ, อามนฺตา.\csromanspacer{} สมาปตฺเตสิยํ อุปปตฺเตสิยํ ทิฏฺฐธมฺมสุขวิหารํ, สมาปตฺเตสิเยน จิตฺเตน อุปปตฺเตสิเยน จิตฺเตน ทิฏฺฐธมฺมสุขวิหาเรน จิตฺเตน สหคตํ สหชาตํ สํสฏฺฐํ สมฺปยุตฺตํ เอกุปฺปาทํ เอกนิโรธํ เอกวตฺถุกํ เอการมฺมณนฺติ, น เหวํ วตฺตพฺเพ ฯเปฯ.\csromanspacer{} นนุ น สมาปตฺเตสิยํ น อุปปตฺเตสิยํ น ทิฏฺฐธมฺมสุขวิหารํ, น สมาปตฺเตสิเยน จิตฺเตน น อุปปตฺเตสิเยน จิตฺเตน น ทิฏฺฐธมฺมสุขวิหาเรน จิตฺเตน สหคตํ สหชาตํ สํสฏฺฐํ สมฺปยุตฺตํ เอกุปฺปาทํ เอกนิโรธํ เอกวตฺถุกํ เอการมฺมณนฺติ, อามนฺตา.\csromanspacer{} หญฺจิ น สมาปตฺเตสิยํ น อุปปตฺเตสิยํ น ทิฏฺฐธมฺมสุขวิหารํ, น สมาปตฺเตสิเยน จิตฺเตน ฯเปฯ เอการมฺมณํ, โน จ วต เร วตฺตพฺเพ อตฺถิ “รูปํ รูปาวจรนฺ”ติ.}

\proseitem{769}{\csromansymbolmark{*}อตฺถิ รูปํ อรูปาวจรนฺติ, อามนฺตา.\csromanspacer{} สมาปตฺเตสิยํ อุปปตฺเตสิยํ ทิฏฺฐธมฺมสุขวิหารํ, สมาปตฺเตสิเยน จิตฺเตน อุปปตฺเตสิเยน จิตฺเตน ทิฏฺฐธมฺมสุขวิหาเรน จิตฺเตน สหคตํ สหชาตํ สํสฏฺฐํ สมฺปยุตฺตํ เอกุปฺปาทํ เอกนิโรธํ เอกวตฺถุกํ เอการมฺมณนฺติ, น เหวํ วตฺตพฺเพ ฯเปฯ.\csromanspacer{} นนุ น สมาปตฺเตสิยํ น อุปปตฺเตสิยํ น ทิฏฺฐธมฺมสุขวิหารํ, น สมาปตฺเตสิเยน}

\vspace{\tipitakaparskip}
\csromancenter{โสฬสมวคฺค จิตฺเตน ฯเปฯ เอการมฺมณนฺติ, อามนฺตา.\csromanspacer{} หญฺจิ น สมาปตฺเตสิยํ น อุปปตฺเตสิยํ ฯเปฯ เอกวตฺถุกํ เอการมฺมณํ, โน จ วต เร วตฺตพฺเพ “อตฺถิ รูปํ อรูปาวจรนฺ”ติ.}
\csromanmatikaanchor{mtk.197}
\csromantocmarkref{subsection}{(165) 10. รูปารูปธาตุปริยาปนฺนกถา}{mtk.197}
\bookmark[dest={mtk.197},level=2]{(165) 10. รูปารูปธาตุปริยาปนฺนกถา}
\proseitem{770}{\csromanfolio{391}\csromansymbolmark{+}น วตฺตพฺพํ “อตฺถิ รูปํ รูปาวจรํ, อตฺถิ รูปํ อรูปาวจรนฺ”ติ, อามนฺตา.\csromanspacer{} นนุ กามาวจรกมฺมสฺส กตตฺตา รูปํ กามาวจรนฺติ, อามนฺตา.\csromanspacer{} หญฺจิ กามาวจรกมฺมสฺส กตตฺตา รูปํ กามาวจรํ.\csromanspacer{} เตน วต เร วตฺตพฺเพ “รูปาวจรกมฺมสฺส กตตฺตา รูปํ รูปาวจรํ, อรูปาวจรกมฺมสฺส กตตฺตา รูปํ อรูปาวจรนฺ”ติ.}

\nitthitam{รูปํ รูปาวจรารูปาวจรนฺติกถา นิฏฺฐิตา.}
\chapterhead{16. โสฬสมวคฺค}
\vspace{\tipitakaparskip}
\csromancenter{\textbf{(165) 10.}\csromanspacer{}\textbf{ รูปารูปธาตุปริยาปนฺนกถา}}
\proseitem{771}{\csromansymbolmark{*}รูปราโค รูปธาตุปริยาปนฺโนติ, อามนฺตา.\csromanspacer{} สมาปตฺเตสิโย อุปปตฺเตสิโย ทิฏฺฐธมฺมสุขวิหาโร, สมาปตฺเตสิเยน จิตฺเตน อุปปตฺเตสิเยน จิตฺเตน ทิฏฺฐธมฺมสุขวิหาเรน จิตฺเตน สหคโต สหชาโต สํสฏฺโฐ สมฺปยุตฺโต เอกุปฺปาโท เอกนิโรโธ เอกวตฺถุโก เอการมฺมโณติ, น เหวํ วตฺตพฺเพ ฯเปฯ.\csromanspacer{} นนุ น สมาปตฺเตสิโย น อุปปตฺเตสิโย น ทิฏฺฐธมฺมสุขวิหาโร, น สมาปตฺเตสิเยน จิตฺเตน ฯเปฯ เอกวตฺถุโก เอการมฺมโณติ, อามนฺตา.\csromanspacer{} หญฺจิ น สมาปตฺเตสิโย น อุปปตฺเตสิโย น ทิฏฺฐธมฺมสุขวิหาโร, น สมาปตฺเตสิเยน จิตฺเตน ฯเปฯ เอกวตฺถุโก เอการมฺมโณ, โน จ วต เร วตฺตพฺเพ “รูปราโค รูปธาตุปริยาปนฺโน”ติ.}

\proseitem{772}{\csromansymbolmark{*}รูปราโค รูปธาตุปริยาปนฺโนติ, อามนฺตา.\csromanspacer{} สทฺทราโค สทฺทธาตุปริยาปนฺโนติ, น เหวํ วตฺตพฺเพ ฯเปฯ.\csromanspacer{} รูปราโค รูปธาตุปริยาปนฺโนติ, อามนฺตา.\csromanspacer{} คนฺธราโค ฯเปฯ.\csromanspacer{} รสราโค ฯเปฯ.\csromanspacer{} โผฏฺฐพฺพราโค โผฏฺฐพฺพธาตุปริยาปนฺโนติ, น เหวํ วตฺตพฺเพ ฯเปฯ.}

\prose{\csromanfolio{392}\csromansymbolmark{*}สทฺทราโค น วตฺตพฺพํ “สทฺทธาตุปริยาปนฺโน”ติ, อามนฺตา.\csromanspacer{} รูปราโค น วตฺตพฺพํ “รูปธาตุปริยาปนฺโน”ติ, น เหวํ วตฺตพฺเพ ฯเปฯ.\csromanspacer{} คนฺธราโค ฯเปฯ.\csromanspacer{} รสราโค ฯเปฯ.\csromanspacer{} โผฏฺฐพฺพราโค น วตฺตพฺพํ “โผฏฺฐพฺพธาตุปริยาปนฺโน”ติ, อามนฺตา.\csromanspacer{} รูปราโค น วตฺตพฺพํ “รูปธาตุปริยาปนฺโน”ติ, น เหวํ วตฺตพฺเพ ฯเปฯ.}

\proseitem{773}{\csromansymbolmark{*}อรูปราโค อรูปธาตุปริยาปนฺโนติ, อามนฺตา.\csromanspacer{} อรูปราโค น วตฺตพฺพํ “อรูปธาตุปริยาปนฺโน”ติ, น เหวํ วตฺตพฺเพ ฯเปฯ.\csromanspacer{} อรูปราโค อรูปธาตุปริยาปนฺโนติ, อามนฺตา.\csromanspacer{} สมาปตฺเตสิโย อุปปตฺเตสิโย ทิฏฺฐธมฺมสุขวิหาโร, สมาปตฺเตสิเยน จิตฺเตน อุปปตฺเตสิเยน จิตฺเตน ทิฏฺฐธมฺมสุขวิหาเรน จิตฺเตน สหคโต สหชาโต สํสฏฺโฐ สมฺปยุตฺโต เอกุปฺปาโท เอกนิโรโธ เอกวตฺถุโก เอการมฺมโณติ, น เหวํ วตฺตพฺเพ ฯเปฯ.\csromanspacer{} นนุ น สมาปตฺเตสิโย น อุปปตฺเตสิโย น ทิฏฺฐธมฺมสุขวิหาโร, น สมาปตฺเตสิเยน จิตฺเตน ฯเปฯ เอกวตฺถุโก เอการมฺมโณติ, อามนฺตา.\csromanspacer{} หญฺจิ น สมาปตฺเตสิโย น อุปปตฺเตสิโย น ทิฏฺฐธมฺมสุขวิหาโร, น สมาปตฺเตสิเยน จิตฺเตน น อุปปตฺเตสิเยน จิตฺเตน น ทิฏฺฐธมฺมสุขวิหาเรน จิตฺเตน สหคโต สหชาโต สํสฏฺโฐ สมฺปยุตฺโต เอกุปฺปาโท เอกนิโรโธ เอกวตฺถุโก เอการมฺมโณ, โน จ วต เร วตฺตพฺเพ “อรูปราโค อรูปธาตุปริยาปนฺโน”ติ.}

\proseitem{774}{\csromansymbolmark{*}อรูปราโค อรูปธาตุปริยาปนฺโนติ, อามนฺตา.\csromanspacer{} สทฺทราโค สทฺทธาตุปริยาปนฺโนติ, น เหวํ วตฺตพฺเพ ฯเปฯ.\csromanspacer{} อรูปราโค อรูปธาตุปริยาปนฺโนติ, อามนฺตา.\csromanspacer{} คนฺธราโค ฯเปฯ.\csromanspacer{} รสราโค ฯเปฯ.\csromanspacer{} โผฏฺฐพฺพราโค โผฏฺฐพฺพธาตุปริยาปนฺโนติ, น เหวํ วตฺตพฺเพ ฯเปฯ.}

\prose{\csromansymbolmark{*}สทฺทราโค น วตฺตพฺพํ “สทฺทธาตุปริยาปนฺโน”ติ, อามนฺตา.\csromanspacer{} อรูปราโค น วตฺตพฺพํ “อรูปธาตุปริยาปนฺโน”ติ, น เหวํ วตฺตพฺเพ ฯเปฯ.\csromanspacer{} คนฺธราโค ฯเปฯ.\csromanspacer{} รสราโค ฯเปฯ.\csromanspacer{} โผฏฺฐพฺพราโค น วตฺตพฺพํ “โผฏฺฐพฺพธาตุปริยาปนฺโน”ติ, อามนฺตา.\csromanspacer{} อรูปราโค น วตฺตพฺพํ “อรูปธาตุปริยาปนฺโน”ติ, น เหวํ วตฺตพฺเพ ฯเปฯ.}

\proseitem{775}{\csromansymbolmark{+}น วตฺตพฺพํ “รูปราโค รูปธาตุปริยาปนฺโน, อรูปราโค อรูปธาตุปริยาปนฺโน”ติ, อามนฺตา.\csromanspacer{} นนุ กามราโค กามธาตุ-}

\vspace{\tipitakaparskip}
\csromancenter{สตฺตรสมวคฺค ปริยาปนฺโนติ, อามนฺตา.\csromanspacer{} หญฺจิ กามราโค กามธาตุปริยาปนฺโน, เตน วต เร วตฺตพฺเพ “รูปราโค รูปธาตุปริยาปนฺโน, อรูปราโค อรูปธาตุปริยาปนฺโน”ติ.}
\csromanmatikaanchor{mtk.198}
\csromanmatikaanchor{mtk.199}
\csromantocmarknopage{chapter}{17. สตฺตรสมวคฺค}{mtk.198}
\bookmark[dest={mtk.198},level=0]{17. สตฺตรสมวคฺค}
\csromantocmarkref{subsection}{(166) 1. อรหโตปุญฺญูปจยกถา}{mtk.199}
\bookmark[dest={mtk.199},level=2]{(166) 1. อรหโตปุญฺญูปจยกถา}
\prose{\csromanfolio{393}รูปราโค รูปธาตุปริยาปนฺโน อรูปราโค อรูปธาตุปริยาปนฺโนติกถา นิฏฺฐิตา.\csromanspacer{} โสฬสโม วคฺโค.}

\titlehead{\textbf{ตสฺสุทฺทานํ}}
\prose{จิตฺตนิคฺคโห, จิตฺตปคฺคโห, สุขานุปฺปทานํ, อธิคยฺห มนสิกาโร, รูปํ เหตุ, รูปํ สเหตุกํ, รูปํ กุสลมฺปิ อกุสลมฺปิ, รูปํ วิปาโก, อตฺถิ รูปํ รูปาวจรํ อตฺถิ รูปํ อรูปาวจรํ, สพฺเพ กิเลสา กามธาตุปริยาปนฺนาติ.}

\chapterhead{17. สตฺตรสมวคฺค}
\vspace{\tipitakaparskip}
\csromancenter{\textbf{(166) 1.}\csromanspacer{}\textbf{ อรหโตปุญฺญูปจยกถา}}
\proseitem{776}{\csromansymbolmark{*}อตฺถิ อรหโต ปุญฺญูปจโยติ, อามนฺตา.\csromanspacer{} อตฺถิ อรหโต อปุญฺญูปจโยติ, น เหวํ วตฺตพฺเพ ฯเปฯ.\csromanspacer{} นตฺถิ อรหโต อปุญฺญูปจโยติ, อามนฺตา.\csromanspacer{} นตฺถิ อรหโต ปุญฺญูปจโยติ, น เหวํ วตฺตพฺเพ ฯเปฯ.}

\proseitem{777}{\csromansymbolmark{*}อตฺถิ อรหโต ปุญฺญูปจโยติ, อามนฺตา.\csromanspacer{} อรหา ปุญฺญาภิสงฺขารํ อภิสงฺขโรติ.\csromanspacer{} อาเนญฺชาภิสงฺขารํ อภิสงฺขโรติ.\csromanspacer{} คติสํวตฺตนิยํ กมฺมํ กโรติ.\csromanspacer{} ภวสํวตฺตนิยํ กมฺมํ กโรติ.\csromanspacer{} อิสฺสริยสํวตฺตนิยํ กมฺมํ กโรติ.\csromanspacer{} อธิปจฺจสํวตฺตนิยํ\footnote{อาธิปจฺจสํวตฺตนิกํ (?)} กมฺมํ กโรติ.\csromanspacer{} มหาโภคสํวตฺตนิยํ กมฺมํ กโรติ.\csromanspacer{} มหาปริวารสํวตฺตนิยํ กมฺมํ กโรติ.\csromanspacer{} เทวโสภคฺยสํวตฺตนิยํ กมฺมํ กโรติ.\csromanspacer{} มนุสฺสโสภคฺยสํวตฺตนิยํ กมฺมํ กโรตีติ, น เหวํ วตฺตพฺเพ ฯเปฯ.}

\csromanmatikaanchor{mtk.200}
\csromantocmarkref{subsection}{(167) 2. นตฺถิ-อรหโต-อกาลมจฺจูติกถา}{mtk.200}
\bookmark[dest={mtk.200},level=2]{(167) 2. นตฺถิ-อรหโต-อกาลมจฺจูติกถา}
\proseitem{778}{\csromansymbolmark{*}อตฺถิ อรหโต ปุญฺญูปจโยติ, อามนฺตา.\csromanspacer{} อรหา อาจินาตีติ, น เหวํ วตฺตพฺเพ ฯเปฯ.\csromanspacer{} อรหา อปจินาตีติ, น เหวํ \csromanfolio{394}วตฺตพฺเพ ฯเปฯ.\csromanspacer{} อรหา ปชหตีติ ฯเปฯ.\csromanspacer{} อรหา อุปาทิยตีติ ฯเปฯ.\csromanspacer{} อรหา วิสิเนตีติ ฯเปฯ.\csromanspacer{} อรหา อุสฺสิเนตีติ ฯเปฯ.\csromanspacer{} อรหา วิธูเปตีติ ฯเปฯ.\csromanspacer{} อรหา สนฺธูเปตีติ, น เหวํ วตฺตพฺเพ ฯเปฯ.\csromanspacer{} นนุ อรหา เนวาจินาติ น อปจินาติ อปจินิตฺวา ฐิโตติ, อามนฺตา.\csromanspacer{} หญฺจิ อรหา เนวาจินาติ นาปจินาติ อปจินิตฺวา ฐิโต, โน จ วต เร วตฺตพฺเพ “อตฺถิ อรหโต ปุญฺญูปจโย”ติ.}

\prose{\csromansymbolmark{*}นนุ อรหา เนว ปชหติ น อุปาทิยติ ปชหิตฺวา ฐิโต, เนว วิสิเนติ น อุสฺสิเนติ วิสิเนตฺวา ฐิโต, เนว วิธูเปติ น สนฺธูเปติ วิธูเปตฺวา ฐิโตติ, อามนฺตา.\csromanspacer{} หญฺจิ อรหา เนว วิธูเปติ น สนฺธูเปติ วิธูเปตฺวา ฐิโต, โน จ วต เร วตฺตพฺเพ “อตฺถิ อรหโต ปุญฺญูปจโย”ติ.}

\proseitem{779}{\csromansymbolmark{+}นตฺถิ อรหโต ปุญฺญูปจโยติ, อามนฺตา.\csromanspacer{} อรหา ทานํ ทเทยฺยาติ, อามนฺตา.\csromanspacer{} หญฺจิ อรหา ทานํ ทเทยฺย, โน จ วต เร วตฺตพฺเพ “นตฺถิ อรหโต ปุญฺญูปจโย”ติ.}

\prose{\csromansymbolmark{+}อรหา จีวรํ ทเทยฺย ฯเปฯ ปิณฺฑปาตํ ทเทยฺย.\csromanspacer{} เสนาสนํ ทเทยฺย.\csromanspacer{} คิลานปจฺจยเภสชฺชปริกฺขารํ ทเทยฺย.\csromanspacer{} ขาทนียํ ทเทยฺย.\csromanspacer{} โภชนียํ ทเทยฺย.\csromanspacer{} ปานียํ ทเทยฺย.\csromanspacer{} เจติยํ วนฺเทยฺย.\csromanspacer{} เจติเย มาลํ อาโรเปยฺย.\csromanspacer{} คนฺธํ อาโรเปยฺย.\csromanspacer{} วิเลปนํ อาโรเปยฺย ฯเปฯ เจติยํ อภิทกฺขิณํ\footnote{ปทกฺขิณํ (?)} กเรยฺยาติ, อามนฺตา.\csromanspacer{} หญฺจิ อรหา เจติยํ อภิทกฺขิณํ กเรยฺย, โน จ วต เร วตฺตพฺเพ “นตฺถิ อรหโต ปุญฺญูปจโย”ติ.}

\nitthitam{อตฺถิ อรหโต ปุญฺญูปจโยติกถา นิฏฺฐิตา.}
\chapterhead{17. สตฺตรสมวคฺค}
\vspace{\tipitakaparskip}
\csromancenter{\textbf{(167) 2.}\csromanspacer{}\textbf{ นตฺถิ-อรหโต-อกาลมจฺจูติกถา}}
\proseitem{780}{\csromansymbolmark{*}นตฺถิ อรหโต อกาลมจฺจูติ, อามนฺตา.\csromanspacer{} นตฺถิ อรหนฺตฆาตโกติ, น เหวํ วตฺตพฺเพ ฯเปฯ.\csromanspacer{} อตฺถิ อรหนฺตฆาตโกติ, อามนฺตา.}

\vspace{\tipitakaparskip}
\csromancenter{สตฺตรสมวคฺค อตฺถิ อรหโต อกาลมจฺจูติ, น เหวํ วตฺตพฺเพ ฯเปฯ.\csromanspacer{} นตฺถิ อรหโต อกาลมจฺจูติ, อามนฺตา.\csromanspacer{} โย อรหนฺตํ ชีวิตา โวโรเปติ, สติ ชีวิเต ชีวิตาวเสเส ชีวิตา โวโรเปติ, อสติ ชีวิเต ชีวิตาวเสเส ชีวิตา โวโรเปตีติ? สติ ชีวิเต ชีวิตาวเสเส ชีวิตา โวโรเปตีติ.\csromanspacer{} หญฺจิ สติ ชีวิเต ชีวิตาวเสเส ชีวิตา โวโรเปติ, โน จ วต เร วตฺตพฺเพ “นตฺถิ อรหโต อกาลมจฺจู”ติ.\csromanspacer{} อสติ ชีวิเต ชีวิตาวเสเส ชีวิตา โวโรเปตีติ.\csromanspacer{} นตฺถิ อรหนฺตฆาตโกติ, น เหวํ วตฺตพฺเพ ฯเปฯ.}
\csromanmatikaanchor{mtk.201}
\csromantocmarkref{subsection}{(168) 3. สพฺพมิทํกมฺมโตติกถา}{mtk.201}
\bookmark[dest={mtk.201},level=2]{(168) 3. สพฺพมิทํกมฺมโตติกถา}
\proseitem{781}{\csromanfolio{395}\csromansymbolmark{*}นตฺถิ อรหโต อกาลมจฺจูติ, อามนฺตา.\csromanspacer{} อรหโต กาเย วิสํ น กเมยฺย, สตฺถํ น กเมยฺย, อคฺคิ น กเมยฺยาติ, น เหวํ วตฺตพฺเพ ฯเปฯ.\csromanspacer{} นนุ อรหโต กาเย วิสํ กเมยฺย, สตฺถํ กเมยฺย, อคฺคิ กเมยฺยาติ, อามนฺตา.\csromanspacer{} หญฺจิ อรหโต กาเย วิสํ กเมยฺย, สตฺถํ กเมยฺย, อคฺคิ กเมยฺย, โน จ วต เร วตฺตพฺเพ “นตฺถิ อรหโต อกาลมจฺจู”ติ.}

\prose{\csromansymbolmark{*}อรหโต กาเย วิสํ น กเมยฺย, สตฺถํ น กเมยฺย, อคฺคิ น กเมยฺยาติ, อามนฺตา.\csromanspacer{} นตฺถิ อรหนฺตฆาตโกติ, น เหวํ วตฺตพฺเพ ฯเปฯ.}

\proseitem{782}{\csromansymbolmark{+}อตฺถิ อรหโต อกาลมจฺจูติ, อามนฺตา.\csromanspacer{} นนุ วุตฺตํ ภควตา “นาหํ ภิกฺขเว สญฺเจตนิกานํ กมฺมานํ กตานํ อุปจิตานํ อปฺปฏิสํเวทิตฺวา พฺยนฺตีภาวํ วทามิ.\csromanspacer{} ตญฺจ โข ทิฏฺเฐว ธมฺเม อุปปชฺชํ วา อปเร วา ปริยาเย”ติ อตฺเถว สุตฺตนฺโตติ, อามนฺตา.\csromanspacer{} เตน หิ นตฺถิ อรหโต อกาลมจฺจูติ.}

\nitthitam{นตฺถิ อรหโต อกาลมจฺจูติกถา นิฏฺฐิตา.}
\chapterhead{17. สตฺตรสมวคฺค}
\vspace{\tipitakaparskip}
\csromancenter{\textbf{(168) 3.}\csromanspacer{}\textbf{ สพฺพมิทํกมฺมโตติกถา}}
\proseitem{783}{\csromansymbolmark{*}สพฺพมิทํ กมฺมโตติ, อามนฺตา, กมฺมมฺปิ กมฺมโตติ.\csromanspacer{} น เหวํ วตฺตพฺเพ ฯเปฯ.\csromanspacer{} สพฺพมิทํ กมฺมโตติ, อามนฺตา.\csromanspacer{} สพฺพมิทํ ปุพฺเพกตเหตูติ, น \csromanfolio{396}เหวํ วตฺตพฺเพ ฯเปฯ.\csromanspacer{} สพฺพมิทํ กมฺมโตติ, อามนฺตา.\csromanspacer{} สพฺพมิทํ กมฺมวิปากโตติ, น เหวํ วตฺตพฺเพ ฯเปฯ.}

\proseitem{784}{\csromansymbolmark{*}สพฺพมิทํ กมฺมวิปากโตติ, อามนฺตา.\csromanspacer{} กมฺมวิปาเกน ปาณํ หเนยฺยาติ, อามนฺตา.\csromanspacer{} ปาณาติปาโต สผโลติ, อามนฺตา.\csromanspacer{} กมฺมวิปาโก สผโลติ, น เหวํ วตฺตพฺเพ ฯเปฯ.\csromanspacer{} กมฺมวิปาโก อผโลติ, อามนฺตา.\csromanspacer{} ปาณาติปาโต อผโลติ, น เหวํ วตฺตพฺเพ ฯเปฯ.}

\prose{\csromansymbolmark{*}กมฺมวิปาเกน อทินฺนํ อาทิเยยฺย ฯเปฯ มุสา ภเณยฺย.\csromanspacer{} ปิสุณํ ภเณยฺย.\csromanspacer{} ผรุสํ ภเณยฺย.\csromanspacer{} สมฺผํ ปลเปยฺย.\csromanspacer{} สนฺธิํ ฉินฺเทยฺย.\csromanspacer{} นิลฺโลปํ หเรยฺย.\csromanspacer{} เอกาคาริกํ กเรยฺย.\csromanspacer{} ปริปนฺเถ ติฏฺเฐยฺย.\csromanspacer{} ปรทารํ คจฺเฉยฺย.\csromanspacer{} คามฆาตกํ กเรยฺย.\csromanspacer{} นิคมฆาตกํ กเรยฺย.\csromanspacer{} กมฺมวิปาเกน ทานํ ทเทยฺย.\csromanspacer{} จีวรํ ทเทยฺย.\csromanspacer{} ปิณฺฑปาตํ ทเทยฺย.\csromanspacer{} เสนาสนํ ทเทยฺย.\csromanspacer{} คิลานปจฺจยเภสชฺช ปริกฺขารํ ทเทยฺยาติ, อามนฺตา.\csromanspacer{} คิลานปจฺจยเภสชฺชปริกฺขาโร สผโลติ, อามนฺตา.\csromanspacer{} กมฺมวิปาโก สผโลติ, น เหวํ วตฺตพฺเพ ฯเปฯ.\csromanspacer{} กมฺมวิปาโก อผโลติ, อามนฺตา.\csromanspacer{} คิลานปจฺจยเภสชฺชปริกฺขาโร อผโลติ, น เหวํ วตฺตพฺเพ ฯเปฯ.}

\proseitem{785}{\csromansymbolmark{+}น วตฺตพฺพํ “สพฺพมิทํ กมฺมโต”ติ, อามนฺตา.\csromanspacer{} นนุ วุตฺตํ ภควตา\csromandash{}}

\csromangathameasure{“กมฺมุนา วตฺตตี โลโก, กมฺมุนา วตฺตตี1 ปชา. \\ กมฺมนิพนฺธนา สตฺตา, รถสฺสาณีว ยายโต. \\ กมฺเมน กิตฺติํ ลภเต ปสํสํ, \\ กมฺเมน ชานิํ จ วธญฺจ พนฺธํ. \\ ตํ กมฺมํ นานากรณํ วิทิตฺวา, \\ กสฺมา วเท ‘นตฺถิ กมฺมนฺ’ติ โลเก”ติ.}
\csromangathagroup{\csromangathastanza{“กมฺมุนา วตฺตตี\footnote{วตฺตติ (อิ, ก, ม 2. 412), วตฺตเต (?)} โลโก, กมฺมุนา วตฺตตี1 ปชา. \\ กมฺมนิพนฺธนา สตฺตา, รถสฺสาณีว ยายโต\footnote{วตฺตติ (อิ, ก, ม 2. 412), วตฺตเต (?)}.}\csromangathastanza{กมฺเมน กิตฺติํ ลภเต ปสํสํ, \\ กมฺเมน ชานิํ จ วธญฺจ พนฺธํ. \\ ตํ กมฺมํ นานากรณํ วิทิตฺวา, \\ กสฺมา วเท ‘นตฺถิ กมฺมนฺ’ติ โลเก”ติ.}}

\prose{อตฺเถว สุตฺตนฺโตติ, อามนฺตา.\csromanspacer{} เตน หิ สพฺพมิทํ กมฺมโตติ.}

\nitthitam{สพฺพมิทํ กมฺมโตติกถา นิฏฺฐิตา.}
\chapterheadpageread{17. สตฺตรสมวคฺค \textbf{17. สตฺตรสมวคฺค}}
\vspace{\tipitakaparskip}
\csromancenter{\textbf{(169) 4.}\csromanspacer{}\textbf{ อินฺทฺริยพทฺธกถา}}
\csromanmatikaanchor{mtk.202}
\csromantocmarkref{subsection}{(169) 4. อินฺทฺริยพทฺธกถา}{mtk.202}
\bookmark[dest={mtk.202},level=2]{(169) 4. อินฺทฺริยพทฺธกถา}
\proseitem{786}{\csromanfolio{397}\csromansymbolmark{*}อินฺทฺริยพทฺธญฺเญว ทุกฺขนฺติ, อามนฺตา.\csromanspacer{} อินฺทฺริยพทฺธญฺเญว อนิจฺจํ สงฺขตํ ปฏิจฺจสมุปฺปนฺนํ ขยธมฺมํ วยธมฺมํ วิราคธมฺมํ นิโรธธมฺมํ วิปริณามธมฺมนฺติ, น เหวํ วตฺตพฺเพ ฯเปฯ.\csromanspacer{} นนุ อนินฺทฺริยพทฺธํ อนิจฺจํ สงฺขตํ ปฏิจฺจสมุปฺปนฺนํ ขยธมฺมํ วยธมฺมํ วิราคธมฺมํ นิโรธธมฺมํ วิปริณามธมฺมนฺติ, อามนฺตา.\csromanspacer{} หญฺจิ อนินฺทฺริยพทฺธํ อนิจฺจํ สงฺขตํ ปฏิจฺจสมุปฺปนฺนํ ขยธมฺมํ วยธมฺมํ วิราคธมฺมํ นิโรธธมฺมํ วิปริณามธมฺมํ, โน จ วต เร วตฺตพฺเพ “อินฺทฺริยพทฺธญฺเญว ทุกฺขนฺ”ติ.}

\prose{\csromansymbolmark{*}อนินฺทฺริยพทฺธํ อนิจฺจํ สงฺขตํ ปฏิจฺจสมุปฺปนฺนํ ฯเปฯ วิปริณามธมฺมํ ตญฺจ น ทุกฺขนฺติ, อามนฺตา.\csromanspacer{} อินฺทฺริยพทฺธํ อนิจฺจํ สงฺขตํ ฯเปฯ วิปริณามธมฺมํ ตญฺจ น ทุกฺขนฺติ, น เหวํ วตฺตพฺเพ ฯเปฯ.}

\prose{\csromansymbolmark{*}อินฺทฺริยพทฺธํ อนิจฺจํ สงฺขตํ ฯเปฯ วิปริณามธมฺมํ ตญฺจ ทุกฺขนฺติ1, อามนฺตา.\csromanspacer{} อนินฺทฺริยพทฺธํ อนิจฺจํ สงฺขตํ ฯเปฯ วิปริณามธมฺมํ ตญฺจ ทุกฺขนฺติ2, น เหวํ วตฺตพฺเพ ฯเปฯ.}

\proseitem{787}{\csromansymbolmark{*}อินฺทฺริยพทฺธญฺเญว ทุกฺขนฺติ, อามนฺตา.\csromanspacer{} นนุ ‘ยทนิจฺจํ ตํ ทุกฺขํ’\footnote{สํ 2. 19 ปิฏฺเฐ.} วุตฺตํ ภควตา, “อนินฺทฺริยพทฺธํ อนิจฺจนฺ”ติ, อามนฺตา.\csromanspacer{} หญฺจิ ‘ยทนิจฺจํ ตํ ทุกฺขํ’3 วุตฺตํ ภควตา, อนินฺทฺริยพทฺธํ อนิจฺจํ, โน จ วต เร วตฺตพฺเพ “อินฺทฺริยพทฺธญฺเญว ทุกฺขนฺ”ติ.}

\proseitem{788}{\csromansymbolmark{+}น วตฺตพฺพํ “อินฺทฺริยพทฺธญฺเญว ทุกฺขนฺ”ติ, อามนฺตา.\csromanspacer{} ยถา อินฺทฺริยพทฺธสฺส ทุกฺขสฺส ปริญฺญาย ภควติ พฺรหฺมจริยํ วุสฺสติ, เอวเมวํ อนินฺทฺริยพทฺธสฺส ทุกฺขสฺส ปริญฺญาย ภควติ พฺรหฺมจริยํ วุสฺสตีติ, น เหวํ วตฺตพฺเพ ฯเปฯ.\csromanspacer{} ยถา อินฺทฺริยพทฺธํ ทุกฺขํ ปริญฺญาตํ น ปุน อุปฺปชฺชติ, เอวเมวํ อนินฺทฺริยพทฺธํ ทุกฺขํ ปริญฺญาตํ น ปุน อุปฺปชฺชตีติ, น เหวํ วตฺตพฺเพ.\csromanspacer{} เตน หิ อินฺทฺริยพทฺธญฺเญว ทุกฺขนฺติ.}

\nitthitam{อินฺทฺริยพทฺธกถา นิฏฺฐิตา.}
\chapterheadpageread{17. สตฺตรสมวคฺค}
\vspace{\tipitakaparskip}
\csromancenter{\textbf{(170) 5.}\csromanspacer{}\textbf{ ฐเปตฺวา-อริยมคฺคนฺติกถา}}
\csromanmatikaanchor{mtk.203}
\csromanmatikaanchor{mtk.204}
\csromantocmarkref{subsection}{(170) 5. ฐเปตฺวา-อริยมคฺคนฺติกถา}{mtk.203}
\bookmark[dest={mtk.203},level=2]{(170) 5. ฐเปตฺวา-อริยมคฺคนฺติกถา}
\csromantocmarkref{subsection}{(171) 6. นวตฺตพฺพํสํโฆทกฺขิณํปฏิคฺคณฺหาติกถา}{mtk.204}
\bookmark[dest={mtk.204},level=2]{(171) 6. นวตฺตพฺพํสํโฆทกฺขิณํปฏิคฺคณฺหาติกถา}
\proseitem{789}{\csromanfolio{398}\csromansymbolmark{*}ฐเปตฺวา อริยมคฺคํ อวเสสา สงฺขารา ทุกฺขาติ, อามนฺตา.\csromanspacer{} ทุกฺขสมุทโยปิ ทุกฺโขติ, น เหวํ วตฺตพฺเพ ฯเปฯ.\csromanspacer{} ทุกฺขสมุทโยปิ ทุกฺโขติ, อามนฺตา.\csromanspacer{} ตีเณว อริยสจฺจานีติ, น เหวํ วตฺตพฺเพ ฯเปฯ.\csromanspacer{} ตีเณว อริยสจฺจานีติ, อามนฺตา.\csromanspacer{} นนุ จตฺตาริ อริยสจฺจานิ วุตฺตานิ ภควตา “ทุกฺขํ, ทุกฺขสมุทโย, ทุกฺขนิโรโธ, ทุกฺขนิโรธคามินี ปฏิปทา”ติ, อามนฺตา.\csromanspacer{} หญฺจิ จตฺตาริ อริยสจฺจานิ วุตฺตานิ ภควตา ทุกฺขํ, ทุกฺขสมุทโย, ทุกฺขนิโรโธ, ทุกฺขนิโรธคามินี ปฏิปทา.\csromanspacer{} โน จ วต เร วตฺตพฺเพ “ตีเณว อริยสจฺจานี”ติ.}

\prose{\csromansymbolmark{*}ทุกฺขสมุทโยปิ ทุกฺโขติ, อามนฺตา.\csromanspacer{} เกนฏฺเฐนาติ อนิจฺจฏฺเฐน.\csromanspacer{} อริยมคฺโค อนิจฺโจติ, อามนฺตา.\csromanspacer{} อริยมคฺโค ทุกฺโขติ, น เหวํ วตฺตพฺเพ ฯเปฯ.}

\prose{\csromansymbolmark{*}อริยมคฺโค อนิจฺโจ โส จ น ทุกฺโขติ, อมนฺตา.\csromanspacer{} ทุกฺขสมุทโย อนิจฺโจ โส จ น ทุกฺโขติ, น เหวํ วตฺตพฺเพ ฯเปฯ.\csromanspacer{} ทุกฺขสมุทโย อนิจฺโจ โส จ ทุกฺโขติ, อามนฺตา.\csromanspacer{} อริยมคฺโค อนิจฺโจ โส จ ทุกฺโขติ, น เหวํ วตฺตพฺเพ ฯเปฯ.}

\proseitem{790}{\csromansymbolmark{+}น วตฺตพฺพํ “ฐเปตฺวา อริยมคฺคํ อวเสสา สงฺขารา ทุกฺขา”ติ, อามนฺตา.\csromanspacer{} นนุ สา ทุกฺขนิโรธคามินี ปฏิปทาติ, อามนฺตา.\csromanspacer{} หญฺจิ สา ทุกฺขนิโรธคามินี ปฏิปทา, เตน วต เร วตฺตพฺเพ “ฐเปตฺวา อริยมคฺคํ อวเสสา สงฺขารา ทุกฺขา”ติ.}

\nitthitam{ฐเปตฺวา อริยมคฺคนฺติกถา นิฏฺฐิตา.}
\chapterhead{17. สตฺตรสมวคฺค}
\prose{\textbf{(171) 6.}\csromanspacer{}\textbf{ นวตฺตพฺพํสํโฆทกฺขิณํปฏิคฺคณฺหาติกถา}}

\proseitem{791}{\csromansymbolmark{*}น วตฺตพฺพํ “สํโฆ ทกฺขิณํ ปฏิคฺคณฺหาตี”ติ, อามนฺตา.\csromanspacer{} นนุ สํโฆ อาหุเนยฺโย ปาหุเนยฺโย ทกฺขิเณยฺโย อญฺชลิกรณีโย อนุตฺตรํ ปุญฺญกฺเขตฺตํ โลกสฺสาติ, อามนฺตา.\csromanspacer{} หญฺจิ สํโฆ อาหุเนยฺโย}

\vspace{\tipitakaparskip}
\csromancenter{สตฺตรสมวคฺค ปาหุเนยฺโย ทกฺขิเณยฺโย อญฺชลิกรณีโย อนุตฺตรํ ปุญฺญกฺเขตฺตํ โลกสฺส.\csromanspacer{} เตน วต เร วตฺตพฺเพ “สํโฆ ทกฺขิณํ ปฏิคฺคณฺหาติ”ติ.}
\csromanmatikaanchor{mtk.205}
\csromantocmarkref{subsection}{(172) 7. นวตฺตพฺพํสํโฆทกฺขิณํ ฯเปฯ กถา ...}{mtk.205}
\bookmark[dest={mtk.205},level=2]{(172) 7. นวตฺตพฺพํสํโฆทกฺขิณํ ฯเปฯ กถา ...}
\prose{\csromanfolio{399}\csromansymbolmark{*}น วตฺตพฺพํ “สํโฆ ทกฺขิณํ ปฏิคฺคณฺหาตี”ติ, อามนฺตา.\csromanspacer{} นนุ จตฺตาโร ปุริสยุคา อฏฺฐปุริสปุคฺคลา ทกฺขิเณยฺยา วุตฺตา ภควตาติ, อามนฺตา.\csromanspacer{} หญฺจิ จตฺตาโร ปุริสยุคา อฏฺฐ ปุริสปุคฺคลา ทกฺขิเณยฺยา วุตฺตา ภควตา.\csromanspacer{} เตน วต เร วตฺตพฺเพ “สํโฆ ทกฺขิณํ ปฏิคฺคณฺหาตี”ติ.}

\prose{\csromansymbolmark{*}น วตฺตพฺพํ “สํโฆ ทกฺขิณํ ปฏิคฺคณฺหาตี”ติ, อามนฺตา.\csromanspacer{} นนุ อตฺถิ เกจิ สํฆสฺส ทานํ เทนฺตีติ, อามนฺตา.\csromanspacer{} หญฺจิ อตฺถิ เกจิ สํฆสฺส ทานํ เทนฺติ.\csromanspacer{} เตน วต เร วตฺตพฺเพ “สํโฆ ทกฺขิณํ ปฏิคฺคณฺหาตี”ติ.\csromanspacer{} นนุ อตฺถิ เกจิ สํฆสฺส จีวรํ เทนฺติ ฯเปฯ ปิณฺฑปาตํ เทนฺติ.\csromanspacer{} เสนาสนํ เทนฺติ.\csromanspacer{} คิลานปจฺจยเภสชฺชปริกฺขารํ เทนฺติ.\csromanspacer{} ขาทนียํ เทนฺติ.\csromanspacer{} โภชนียํ เทนฺติ ฯเปฯ ปานียํ เทนฺตีติ, อามนฺตา.\csromanspacer{} หญฺจิ อตฺถิ เกจิ สํฆสฺส ปานียํ เทนฺติ.\csromanspacer{} เตน วต เร วตฺตพฺเพ “สํโฆ ทกฺขิณํ ปฏิคฺคณฺหาตี”ติ.}

\prose{\csromansymbolmark{*}น วตฺตพฺพํ “สํโฆ ทกฺขิณํ ปฏิคฺคณฺหาตี”ติ, อามนฺตา.\csromanspacer{} นนุ วุตฺตํ ภควตา\csromandash{}}

\csromangathameasure{“อาหุติํ ชาตเวโทว, มหาเมฆํว เมทนี. \\ สํโฆ สมาธิสมฺปนฺโน, ปฏิคฺคณฺหาติ ทกฺขิณนฺ”ติ.}
\csromangathagroup{\csromangathastanza{“อาหุติํ ชาตเวโทว, มหาเมฆํว เมทนี. \\ สํโฆ สมาธิสมฺปนฺโน, ปฏิคฺคณฺหาติ ทกฺขิณนฺ”ติ.}}

\prose{อตฺเถว สุตฺตนฺโตติ, อามนฺตา.\csromanspacer{} เตน หิ สํโฆ ทกฺขิณํ ปฏิคฺคณฺหาตีติ.}

\proseitem{792}{\csromansymbolmark{+}สํโฆ ทกฺขิณํ ปฏิคฺคณฺหาตีติ, อามนฺตา.\csromanspacer{} มคฺโค ปฏิคฺคณฺหาติ ผลํ ปฏิคฺคณฺหาตีติ, น เหวํ วตฺตพฺเพ ฯเปฯ.}

\nitthitam{น วตฺตพฺพํ สํโฆ ทกฺขิณํ ปฏิคฺคณฺหาตีติกถา นิฏฺฐิตา.}
\chapterhead{17. สตฺตรสมวคฺค}
\vspace{\tipitakaparskip}
\csromancenter{\textbf{(172) 7.}\csromanspacer{}\textbf{ นวตฺตพฺพํสํโฆทกฺขิณํ ฯเปฯ กถา}}
\csromanmatikaanchor{mtk.206}
\csromantocmarkref{subsection}{(173) 8. นวตฺตพฺพํสํโฆภุญฺชตีติกถา}{mtk.206}
\bookmark[dest={mtk.206},level=2]{(173) 8. นวตฺตพฺพํสํโฆภุญฺชตีติกถา}
\proseitem{793}{\csromansymbolmark{*}น วตฺตพฺพํ “สํโฆ ทกฺขิณํ วิโสเธตี”ติ, อามนฺตา.\csromanspacer{} นนุ สํโฆ อาหุเนยฺโย ปาหุเนยฺโย ทกฺขิเณยฺโย อญฺชลิกรณีโย \csromanfolio{400}อนุตฺตรํ ปุญฺญกฺเขตฺตํ โลกสฺสาติ, อามนฺตา.\csromanspacer{} หญฺจิ สํโฆ อาหุเนยฺโย ฯเปฯ อนุตฺตรํ ปุญฺญกฺเขตฺตํ โลกสฺส.\csromanspacer{} เตน วต เร วตฺตพฺเพ “สํโฆ ทกฺขิณํ วิโสเธตี”ติ.}

\prose{\csromansymbolmark{*}น วตฺตพฺพํ “สํโฆ ทกฺขิณํ วิโสเธตี”ติ, อามนฺตา.\csromanspacer{} นนุ จตฺตาโร ปุริสยุคา อฏฺฐ ปุริสปุคฺคลา ทกฺขิเณยฺยา วุตฺตา ภควตาติ, อามนฺตา.\csromanspacer{} หญฺจิ จตฺตาโร ปุริสยุคา อฏฺฐ ปุริสปุคฺคลา ทกฺขิเณยฺยา วุตฺตา ภควตา.\csromanspacer{} เตน วต เร วตฺตพฺเพ “สํโฆ ทกฺขิณํ วิโสเธตี”ติ.}

\prose{\csromansymbolmark{*}น วตฺตพฺพํ “สํโฆ ทกฺขิณํ วิโสเธตี”ติ, อามนฺตา.\csromanspacer{} นนุ อตฺถิ เกจิ สํฆสฺส ทานํ ทตฺวา ทกฺขิณํ อาราเธนฺตีติ, อามนฺตา.\csromanspacer{} หญฺจิ อตฺถิ เกจิ สํฆสฺส ทานํ ทตฺวา ทกฺขิณํ อาราเธนฺติ.\csromanspacer{} เตน วต เร วตฺตพฺเพ “สํโฆ ทกฺขิณํ วิโสเธตี”ติ.}

\prose{\csromansymbolmark{*}นนุ อตฺถิ เกจิ สํฆสฺส จีวรํ ทตฺวา ฯเปฯ ปิณฺฑปาตํ ทตฺวา ฯเปฯ เสนาสนํ ทตฺวา คิลานปจฺจยเภสชฺชปริกฺขารํ ทตฺวา, ขาทนียํ ทตฺวา, โภชนียํ ทตฺวา ฯเปฯ ปานียํ ทตฺวา ทกฺขิณํ อาราเธนฺตีติ, อามนฺตา.\csromanspacer{} หญฺจิ อตฺถิ เกจิ สํฆสฺส ปานียํ ทตฺวา ทกฺขิณํ อาราเธนฺติ.\csromanspacer{} เตน วต เร วตฺตพฺเพ “สํโฆ ทกฺขิณํ วิโสเธตี”ติ.}

\proseitem{794}{\csromansymbolmark{+}สํโฆ ทกฺขิณํ วิโสเธตีติ, อามนฺตา.\csromanspacer{} มคฺโค วิโสเธติ ผลํ วิโสเธตีติ, น เหวํ วตฺตพฺเพ ฯเปฯ.}

\nitthitam{น วตฺตพฺพํ สํโฆ ทกฺขิณํ วิโสเธตีติกถา นิฏฺฐิตา.}
\chapterhead{17. สตฺตรสมวคฺค}
\vspace{\tipitakaparskip}
\csromancenter{\textbf{(173) 8.}\csromanspacer{}\textbf{ นวตฺตพฺพํสํโฆภุญฺชตีติกถา}}
\proseitem{795}{\csromansymbolmark{*}น วตฺตพฺพํ “สํโฆ ภุญฺชติ ปิวติ ขาทติ สายตี”ติ, อามนฺตา.\csromanspacer{} นนุ อตฺถิ เกจิ สํฆภตฺตานิ กโรนฺติ อุทฺเทสภตฺตานิ กโรนฺติ ยาคุปานานิ กโรนฺตีติ, อามนฺตา.\csromanspacer{} หญฺจิ อตฺถิ เกจิ สํฆภตฺตานิ กโรนฺติ อุทฺเทสภตฺตานิ กโรนฺติ ยาคุปานานิ กโรนฺติ.\csromanspacer{} เตน วต เร วตฺตพฺเพ “สํโฆ ภุญฺชติ ปิวติ ขาทติ สายตี”ติ.}

\chapterheadpageread{17. สตฺตรสมวคฺค}
\csromanmatikaanchor{mtk.207}
\csromantocmarkref{subsection}{(174) 9. นวตฺตพฺพํ ฯเปฯ มหปฺผลนฺติกถา}{mtk.207}
\bookmark[dest={mtk.207},level=2]{(174) 9. นวตฺตพฺพํ ฯเปฯ มหปฺผลนฺติกถา}
\prose{\csromanfolio{401}\csromansymbolmark{*}น วตฺตพฺพํ “สํโฆ ภุญฺชติ ปิวติ ขาทติ สายตี”ติ, อามนฺตา.\csromanspacer{} นนุ วุตฺตํ ภควตา “คณโภชนํ ปรมฺปรโภชนํ อติริตฺตโภชนํ อนติริตฺตโภชนนฺ”ติ, อามนฺตา.\csromanspacer{} หญฺจิ วุตฺตํ ภควตา คณโภชนํ ปรมฺปรโภชนํ อติริตฺตโภชนํ อนติริตฺตโภชนํ, เตน วต เร วตฺตพฺเพ “สํโฆ ภุญฺชติ ปิวติ ขาทติ สายตี”ติ.}

\prose{\csromansymbolmark{*}น วตฺตพฺพํ “สํโฆ ภุญฺชติ ปิวติ ขาทติ สายตี”ติ, อามนฺตา.\csromanspacer{} นนุ อฏฺฐ ปานานิ วุตฺตานิ ภควตา “อมฺพปานํ ชมฺพุปานํ โจจปานํ โมจปานํ มธุกปานํ\footnote{มธุปานํ (สี, สฺยา, กํ, อิ) วิ 3. 344 ปิฏฺเฐ ปน ปสฺสิตพฺพํ.} มุทฺทิกปานํ สาลุกปานํ ผารุสกปานนฺ”ติ, อามนฺตา.\csromanspacer{} หญฺจิ อฏฺฐ ปานานิ วุตฺตานิ ภควตา อมฺพปานํ ชมฺพุปานํ โจจปานํ โมจปานํ มธุกปานํ มุทฺทิกปานํ สาลุกปานํ ผารุสกปานํ, เตน วต เร วตฺตพฺเพ “สํโฆ ภุญฺชติ ปิวติ ขาทติ สายตี”ติ.}

\proseitem{796}{\csromansymbolmark{+}สํโฆ ภุญฺชติ ปิวติ ขาทติ สายตีติ, อามนฺตา.\csromanspacer{} มคฺโค ภุญฺชติ ปิวติ ขาทติ สายติ, ผลํ ภุญฺชติ ปิวติ ขาทติ สายตีติ, น เหวํ วตฺตพฺเพ ฯเปฯ.}

\nitthitam{น วตฺตพฺพํ สํโฆ ภุญฺชตีติกถา นิฏฺฐิตา.}
\chapterhead{17. สตฺตรสมวคฺค}
\vspace{\tipitakaparskip}
\csromancenter{\textbf{(174) 9.}\csromanspacer{}\textbf{ นวตฺตพฺพํ ฯเปฯ มหปฺผลนฺติกถา}}
\proseitem{797}{\csromansymbolmark{*}น วตฺตพฺพํ “สํฆสฺส ทินฺนํ มหปฺผลนฺ”ติ, อามนฺตา.\csromanspacer{} นนุ สํโฆ อาหุเนยฺโย ปาหุเนยฺโย ทกฺขิเณยฺโย อญฺชลิกรณีโย อนุตฺตรํ ปุญฺญกฺเขตฺตํ โลกสฺสาติ, อามนฺตา.\csromanspacer{} หญฺจิ สํโฆ อาหุเนยฺโย ฯเปฯ อนุตฺตรํ ปุญฺญกฺเขตฺตํ โลกสฺส, เตน วต เร วตฺตพฺเพ “สํฆสฺส ทินฺนํ มหปฺผลนฺ”ติ.}

\prose{\csromansymbolmark{*}น วตฺตพฺพํ “สํฆสฺส ทินฺนํ มหปฺผลนฺ”ติ, อามนฺตา.\csromanspacer{} นนุ จตฺตาโร ปุริสยุคา อฏฺฐ ปุริสปุคฺคลา ทกฺขิเณยฺยา วุตฺตา ภควตาติ, อามนฺตา.\csromanspacer{} หญฺจิ จตฺตาโร ปุริสยุคา อฏฺฐ ปุริสปุคฺคลา ทกฺขิเณยฺยา วุตฺตา ภควตา, เตน วต เร วตฺตพฺเพ “สํฆสฺส ทินฺนํ มหปฺผลนฺ”ติ.}

\proseitem{798}{\csromanfolio{402}\csromansymbolmark{*}น วตฺตพฺพํ “สํฆสฺส ทินฺนํ มหปฺผลนฺ”ติ, อามนฺตา.\csromanspacer{} นนุ วุตฺตํ ภควตา “สํเฆ โคตมิ เทหิ, สํเฆ เต ทินฺเน อหญฺเจว ปูชิโต ภวิสฺสามิ สํโฆ จา”ติ\footnote{ม 3. 296 ปิฏฺเฐ.} อตฺเถว สุตฺตนฺโตติ, อามนฺตา.\csromanspacer{} เตน หิ “สํฆสฺส ทินฺนํ มหปฺผลนฺ”ติ.}

\prose{\csromansymbolmark{*}น วตฺตพฺพํ “สํฆสฺส ทินฺนํ มหปฺผลนฺ”ติ, อามนฺตา.\csromanspacer{} นนุ สกฺโก เทวานมินฺโท ภควนฺตํ เอตทโวจ\csromandash{}}

\csromangathameasure{“ยชมานานํ มนุสฺสานํ, ปุญฺญเปกฺขาน ปาณินํ. \\ กโรตํ โอปธิกํ ปุญฺญํ, กตฺถ ทินฺนํ มหปฺผลนฺ”ติ. \\ “จตฺตาโร จ ปฏิปนฺนา, จตฺตาโร จ ผเล ฐิตา. \\ เอส สํโฆ อุชุภูโต, ปญฺญาสีลสมาหิโต. \\ ยชมานานํ มนุสฺสานํ, ปุญฺญเปกฺขาน ปาณินํ. \\ กโรตํ โอปธิกํ ปุญฺญํ, สํเฆ ทินฺนํ มหปฺผลนฺ”ติ. \\ เอโส หิ สํโฆ วิปุโล มหคฺคโต, \\ เอสปฺปเมยฺโย อุทธีว สาคโร. \\ เอเต หิ เสฏฺฐา นรวีรสาวกา, \\ ปภงฺกรา ธมฺมมุทีรยนฺติ. \\ เตสํ สุทินฺนํ สุหุตํ สุยิฏฺฐํ, \\ เย สํฆมุทฺทิสฺส ททนฺติ ทานํ. \\ สา ทกฺขิณา สํฆคตา ปติฏฺฐิตา, \\ มหปฺผลา โลกวิทูน วณฺณิตา. \\ เอตาทิสํ ยญฺญมนุสฺสรนฺตา, \\ เย เวทชาตา วิจรนฺติ โลเก. \\ วิเนยฺย มจฺเฉรมลํ สมูลํ, \\ อนินฺทิตา สคฺคมุเปนฺติ ฐานนฺ”ติ.}
\csromangathagroup{\csromangathastanza{“ยชมานานํ มนุสฺสานํ, ปุญฺญเปกฺขาน ปาณินํ. \\ กโรตํ โอปธิกํ ปุญฺญํ, กตฺถ ทินฺนํ มหปฺผลนฺ”ติ.}\csromangathastanza{“จตฺตาโร จ ปฏิปนฺนา, จตฺตาโร จ ผเล ฐิตา. \\ เอส สํโฆ อุชุภูโต, ปญฺญาสีลสมาหิโต.}\csromangathastanza{ยชมานานํ มนุสฺสานํ, ปุญฺญเปกฺขาน ปาณินํ. \\ กโรตํ โอปธิกํ ปุญฺญํ, สํเฆ ทินฺนํ มหปฺผลนฺ”ติ\footnote{สํ 1. 235 ปิฏฺเฐ; ขุ 2. 50, 62 ปิฏฺเฐ จ.}.}\csromangathastanza{เอโส หิ สํโฆ วิปุโล มหคฺคโต, \\ เอสปฺปเมยฺโย อุทธีว สาคโร. \\ เอเต หิ เสฏฺฐา นรวีรสาวกา\footnote{นรสีหสาวกา (ก)}, \\ ปภงฺกรา ธมฺมมุทีรยนฺติ.}\csromangathastanza{เตสํ สุทินฺนํ สุหุตํ สุยิฏฺฐํ, \\ เย สํฆมุทฺทิสฺส ททนฺติ ทานํ. \\ สา ทกฺขิณา สํฆคตา ปติฏฺฐิตา, \\ มหปฺผลา โลกวิทูน วณฺณิตา.}\csromangathastanza{เอตาทิสํ ยญฺญมนุสฺสรนฺตา, \\ เย เวทชาตา วิจรนฺติ\footnote{วิหรนฺติ (สิ, ก)} โลเก. \\ วิเนยฺย มจฺเฉรมลํ สมูลํ, \\ อนินฺทิตา สคฺคมุเปนฺติ ฐานนฺ”ติ\footnote{วิหรนฺติ (สิ, ก)}.}}

\prose{อตฺเถว สุตฺตนฺโตติ, อามนฺตา.\csromanspacer{} เตน หิ สํฆสฺส ทินฺนํ มหปฺผลนฺติ.}

\nitthitam{น วตฺตพฺพํ สํฆสฺส ทินฺนํ มหปฺผลนฺติกถา นิฏฺฐิตา.}
\chapterheadpageread{17. สตฺตรสมวคฺค \textbf{17. สตฺตรสมวคฺค}}
\csromanmatikaanchor{mtk.208}
\csromanmatikaanchor{mtk.209}
\csromantocmarkref{subsection}{(175) 10. นวตฺตพฺพํพุทฺธสฺสทินฺนํ มหปฺผลนฺติกถา ...}{mtk.208}
\bookmark[dest={mtk.208},level=2]{(175) 10. นวตฺตพฺพํพุทฺธสฺสทินฺนํ มหปฺผลนฺติกถา ...}
\csromantocmarkref{subsection}{(176) 11. ทกฺขิณาวิสุทฺธิกถา}{mtk.209}
\bookmark[dest={mtk.209},level=2]{(176) 11. ทกฺขิณาวิสุทฺธิกถา}
\prose{\csromanfolio{403}\textbf{(175) 10.}\csromanspacer{}\textbf{ นวตฺตพฺพํพุทฺธสฺสทินฺนํมหปฺผลนฺติกถา}}

\proseitem{779}{\csromansymbolmark{*}น วตฺตพฺพํ “พุทฺธสฺส ทินฺนํ มหปฺผลนฺ”ติ, อามนฺตา.\csromanspacer{} นนุ ภควา ทฺวิปทานํ อคฺโค ทฺวิปทานํ เสฏฺโฐ ทฺวิปทานํ ปโมกฺโข ทฺวิปทานํ อุตฺตโม ทฺวิปทานํ ปวโร อสโม อสมสโม อปฺปฏิสโม อปฺปฏิภาโค อปฺปฏิปุคฺคโลติ, อามนฺตา.\csromanspacer{} หญฺจิ ภควา ทฺวิปทานํ อคฺโค ทฺวิปทานํ เสฏฺโฐ ทฺวิปทานํ ปโมกฺโข ทฺวิปทานํ อุตฺตโม ทฺวิปทานํ ปวโร อสโม อสมสโม อปฺปฏิสโม อปฺปฏิภาโค อปฺปฏิปุคฺคโล, เตน วต เร วตฺตพฺเพ “พุทฺธสฺส ทินฺนํ มหปฺผลนฺ”ติ.}

\prose{\csromansymbolmark{*}น วตฺตพฺพํ “พุทฺธสฺส ทินฺนํ มหปฺผลนฺ”ติ, อามนฺตา.\csromanspacer{} อตฺถิ โกจิ พุทฺเธน สมสโม สีเลน สมาธินา ปญฺญายาติ, นตฺถิ.\csromanspacer{} หญฺจิ นตฺถิ โกจิ พุทฺเธน สมสโม สีเลน สมาธินา ปญฺญาย, เตน วต เร วตฺตพฺเพ “พุทฺธสฺส ทินฺนํ มหปฺผลนฺ”ติ.}

\prose{\csromansymbolmark{*}น วตฺตพฺพํ “พุทฺธสฺส ทินฺนํ มหปฺผลนฺ”ติ, อามนฺตา.\csromanspacer{} นนุ วุตฺตํ ภควตา\csromandash{}}

\csromangathameasure{“นยิมสฺมิํ วา โลเก ปรสฺมิํ วา ปน, \\ พุทฺเธน เสฏฺโฐ จ สโม จ วิชฺชติ. \\ ยมาหุเนยฺยานํ อคฺคตํ คโต, \\ ปุญฺญตฺถิกานํ วิปุลปฺผเลสินนฺ”ติ.}
\csromangathagroup{\csromangathastanza{“นยิมสฺมิํ วา โลเก ปรสฺมิํ วา ปน, \\ พุทฺเธน เสฏฺโฐ จ สโม จ วิชฺชติ. \\ ยมาหุเนยฺยานํ อคฺคตํ คโต, \\ ปุญฺญตฺถิกานํ วิปุลปฺผเลสินนฺ”ติ\footnote{ขุ 2. 95 วิมาเน.}.}}

\prose{อตฺเถว สุตฺตนฺโตติ, อามนฺตา.\csromanspacer{} เตน หิ พุทฺธสฺส ทินฺนํ มหปฺผลนฺติ.}

\nitthitam{น วตฺตพฺพํ พุทฺธสฺส ทินฺนํ มหปฺผลนฺติกถา นิฏฺฐิตา.}
\chapterhead{17. สตฺตรสมวคฺค}
\vspace{\tipitakaparskip}
\csromancenter{\textbf{(176) 11.}\csromanspacer{}\textbf{ ทกฺขิณาวิสุทฺธิกถา}}
\proseitem{800}{\csromansymbolmark{*}ทายกโตว ทานํ วิสุชฺฌติ โน ปฏิคฺคาหกโตติ, อามนฺตา.\csromanspacer{} นนุ อตฺถิ เกจิ ปฏิคฺคาหกา อาหุเนยฺยา ปาหุเนยฺยา ทกฺขิเณยฺยา อญฺชลิกรณียา อนุตฺตรํ ปุญฺญกฺเขตฺตํ โลกสฺสาติ, \csromanfolio{404}อามนฺตา.\csromanspacer{} หญฺจิ อตฺถิ เกจิ ปฏิคฺคาหกา อาหุเนยฺยา ปาหุเนยฺยา ทกฺขิเณยฺยา อญฺชลิกรณียา อนุตฺตรํ ปุญฺญกฺเขตฺตํ โลกสฺส, โน จ วต เร วตฺตพฺเพ “ทายกโตว ทานํ วิสุชฺฌติ โน ปฏิคฺคาหกโต”ติ.}

\prose{\csromansymbolmark{*}ทายกโตว ทานํ วิสุชฺฌติ โน ปฏิคฺคาหกโตติ, อามนฺตา.\csromanspacer{} นนุ จตฺตาโร ปุริสยุคา อฏฺฐ ปุริสปุคฺคลา ทกฺขิเณยฺยา วุตฺตา ภควตาติ, อามนฺตา.\csromanspacer{} หญฺจิ จตฺตาโร ปุริสยุคา อฏฺฐ ปุริสปุคฺคลา ทกฺขิเณยฺยา วุตฺตา ภควตา, โน จ วต เร วตฺตพฺเพ “ทายกโตว ทานํ วิสุชฺฌติ โน ปฏิคฺคาหกโต”ติ.}

\prose{\csromansymbolmark{*}ทายกโตว ทานํ วิสุชฺฌติ โน ปฏิคฺคาหกโตติ, อามนฺตา.\csromanspacer{} นนุ อตฺถิ เกจิ โสตาปนฺเน ทานํ ทตฺวา ทกฺขิณํ อาราเธนฺตีติ, อามนฺตา.\csromanspacer{} หญฺจิ อตฺถิ เกจิ โสตาปนฺเน ทานํ ทตฺวา ทกฺขิณํ อาราเธนฺติ, โน จ วต เร วตฺตพฺเพ “ทายกโตว ทานํ วิสุชฺฌติ โน ปฏิคฺคาหกโต”ติ.}

\prose{\csromansymbolmark{*}นนุ อตฺถิ เกจิ สกทาคามิสฺส ฯเปฯ อนาคามิสฺส ฯเปฯ อรหโต ทานํ ทตฺวา ทกฺขิณํ อาราเธนฺตีติ, อามนฺตา.\csromanspacer{} หญฺจิ อตฺถิ เกจิ อรหโต ทานํ ทตฺวา ทกฺขิณํ อาราเธนฺติ, โน จ วต เร วตฺตพฺเพ “ทายกโตว ทานํ วิสุชฺฌติ โน ปฏิคฺคาหกโต”ติ.}

\proseitem{801}{\csromansymbolmark{+}ปฏิคฺคาหกโต ทานํ วิสุชฺฌตีติ, อามนฺตา.\csromanspacer{} อญฺโญ อญฺญสฺส การโก, ปรกตํ สุขทุกฺขํ, อญฺโญ กโรติ อญฺโญ ปฏิสํเวเทตีติ, น เหวํ วตฺตพฺเพ ฯเปฯ.}

\prose{\csromansymbolmark{*}ทายกโตว ทานํ วิสุชฺฌติ โน ปฏิคฺคาหกโตติ, อามนฺตา.\csromanspacer{} นนุ วุตฺตํ ภควตา “จตสฺโส โข อิมา อานนฺท ทกฺขิณา วิสุทฺธิโย.\csromanspacer{} กตมา จตสฺโส, อตฺถานนฺท ทกฺขิณา ทายกโต วิสุชฺฌติ โน ปฏิคฺคาหกโต, อตฺถานนฺท ทกฺขิณา ปฏิคฺคาหกโต วิสุชฺฌติ โน ทายกโต, อตฺถานนฺท ทกฺขิณา ทายกโต เจว วิสุชฺฌติ ปฏิคฺคาหกโต จ, อตฺถานนฺท ทกฺขิณา เนว ทายกโต วิสุชฺฌติ โน ปฏิคฺคาหกโต, อิมา โข อานนฺท จตสฺโส ทกฺขิณา วิสุทฺธิโย”ติ\footnote{ม 3. 299 ปิฏฺเฐ.}}

\vspace{\tipitakaparskip}
\csromancenter{อฏฺฐารสมวคฺค อตฺเถว สุตฺตนฺโตติ, อามนฺตา.\csromanspacer{} เตน หิ น วตฺตพฺพํ “ทายกโตว ทานํ วิสุชฺฌติ โน ปฏิคฺคาหกโต”ติ.\csromanspacer{} ทกฺขิณาวิสุทฺธิกถา นิฏฺฐิตา.\csromanspacer{} สตฺตรสโม วคฺโค.}
\titlehead{\textbf{ตสฺสุทฺทานํ}}
\csromanmatikaanchor{mtk.210}
\csromanmatikaanchor{mtk.211}
\csromantocmarknopage{chapter}{18. อฏฺฐารสมวคฺค}{mtk.210}
\bookmark[dest={mtk.210},level=0]{18. อฏฺฐารสมวคฺค}
\csromantocmarkref{subsection}{(177) 1. มนุสฺสโลกกถา}{mtk.211}
\bookmark[dest={mtk.211},level=2]{(177) 1. มนุสฺสโลกกถา}
\prose{\csromanfolio{405}อตฺถิ อรหโต ปุญฺญูปจโย, นตฺถิ อรหโต อกาลมจฺจุ, สพฺพมิทํ กมฺมโต, อินฺทฺริยพทฺธญฺเญว ทุกฺขํ, ฐเปตฺวา อริยมคฺคํ อวเสสา สงฺขารา ทุกฺขา, สํโฆ ทกฺขิณํ ปฏิคฺคณฺหาติ, สํโฆ ทกฺขิณํ วิโสเธติ, สํโฆ ภุญฺชติ ปิวติ ขาทติ สายติ, สํฆสฺส ทินฺนํ มหปฺผลํ, อตฺถิ ทานํ วิสุทฺธิยาติ\footnote{วิสุชฺฌตีติ (สฺยา)}.}

\chapterhead{18. อฏฺฐารสมวคฺค}
\vspace{\tipitakaparskip}
\csromancenter{\textbf{(177) 1.}\csromanspacer{}\textbf{ มนุสฺสโลกกถา}}
\proseitem{802}{\csromansymbolmark{*}น วตฺตพฺพํ “พุทฺโธ ภควา มนุสฺสโลเก อฏฺฐาสี”ติ, อามนฺตา.\csromanspacer{} นนุ อตฺถิ พุทฺธวุตฺถานิ เจติยานิ อารามวิหารคามนิคมนครานิ รฏฺฐานิ ชนปทานีติ, อามนฺตา.\csromanspacer{} หญฺจิ อตฺถิ พุทฺธวุตฺถานิ เจติยานิ อารามวิหารคามนิคมนครานิ รฏฺฐานิ ชนปทานิ, เตน วต เร วตฺตพฺเพ “พุทฺโธ ภควา มนุสฺสโลเก อฏฺฐาสี”ติ.}

\prose{\csromansymbolmark{*}น วตฺตพฺพํ “พุทฺโธ ภควา มนุสฺสโลเก อฏฺฐาสี”ติ, อามนฺตา.\csromanspacer{} นนุ ภควา ลุมฺพินิยา ชาโต, โพธิยา มูเล อภิสมฺพุทฺโธ, พาราณสิยํ ภควตา ธมฺมจกฺกํ ปวตฺติตํ, จาปาเล เจติเย อายุสงฺขาโร โอสฺสฏฺโฐ, กุสินารายํ ภควา ปรินิพฺพุโตติ, อามนฺตา.\csromanspacer{} หญฺจิ ภควา ลุมฺพินิยา ชาโต ฯเปฯ กุสินารายํ ภควา ปรินิพฺพุโต, เตน วต เร วตฺตพฺเพ “พุทฺโธ ภควา มนุสฺสโลเก อฏฺฐาสี”ติ.}

\csromanmatikaanchor{mtk.212}
\csromantocmarkref{subsection}{(178) 2. ธมฺมเทสนากถา}{mtk.212}
\bookmark[dest={mtk.212},level=2]{(178) 2. ธมฺมเทสนากถา}
\prose{\csromansymbolmark{*}น วตฺตพฺพํ “พุทฺโธ ภควา มนุสฺสโลเก อฏฺฐาสี”ติ, อามนฺตา.\csromanspacer{} นนุ วุตฺตํ ภควตา “เอกมิทาหํ ภิกฺขเว สมยํ อุกฺกฏฺฐายํ วิหรามิ สุภควเน \csromanfolio{406}สาลราชมูเล”ติ\footnote{ที 2. 42 ปิฏฺเฐ.}. “เอกมิทาหํ ภิกฺขเว สมยํ อุรุเวลายํ วิหรามิ อชปาลนิโคฺรเธ ปฐมาภิสมฺพุทฺโธ”ติ\footnote{อํ 1. 328 ปิฏฺเฐ.}. “เอกมิทาหํ ภิกฺขเว สมยํ ราชคเห วิหรามิ เวฬุวเน กลนฺทกนิวาเป”ติ\footnote{ที 2. 97 ปิฏฺเฐ.}. “เอกมิทาหํ ภิกฺขเว สมยํ สาวตฺถิยํ วิหรามิ เชตวเน อนาถปิณฺฑิกสฺส อาราเม”ติ. “เอกมิทาหํ ภิกฺขเว สมยํ เวสาลิยํ วิหรามิ มหาวเน กูฏาคารสาลายนฺ”ติ\footnote{ที 3. 7 ปิฏฺเฐ.} อตฺเถว สุตฺตนฺโตติ, อามนฺตา.\csromanspacer{} เตน หิ พุทฺโธ ภควา มนุสฺสโลเก อฏฺฐาสีติ.}

\proseitem{803}{\csromansymbolmark{+}พุทฺโธ ภควา มนุสฺสโลเก อฏฺฐาสีติ, อามนฺตา.\csromanspacer{} นนุ ภควา โลเก ชาโต, โลเก สํวฑฺโฒ, โลกํ อภิภุยฺย วิหรติ อนุปลิตฺโต โลเกนาติ\footnote{อํ 1. 347 ปิฏฺเฐ.}, อามนฺตา.\csromanspacer{} หญฺจิ ภควา โลเก ชาโต, โลเก สํวฑฺโฒ, โลกํ อภิภุยฺย วิหรติ อนุปลิตฺโต โลเกน, โน จ วต เร วตฺตพฺเพ “พุทฺโธ ภควา มนุสฺสโลเก อฏฺฐาสี”ติ.}

\nitthitam{มนุสฺสโลกกถา นิฏฺฐิตา.}
\chapterhead{18. อฏฺฐารสมวคฺค}
\vspace{\tipitakaparskip}
\csromancenter{\textbf{(178) 2.}\csromanspacer{}\textbf{ ธมฺมเทสนากถา}}
\proseitem{804}{\csromansymbolmark{*}น วตฺตพฺพํ “พุทฺเธน ภควตา ธมฺโม เทสิโต”ติ, อามนฺตา.\csromanspacer{} เกน เทสิโตติ, อภินิมฺมิเตน เทสิโตติ.\csromanspacer{} อภินิมฺมิโต ชิโน สตฺถา สมฺมาสมฺพุทฺโธ สพฺพญฺญู สพฺพทสฺสาวี ธมฺมสฺสามี ธมฺมปฺปฏิสรโณติ, น เหวํ วตฺตพฺเพ ฯเปฯ.}

\prose{\csromansymbolmark{*}น วตฺตพฺพํ “พุทฺเธน ภควตา ธมฺโม เทสิโต”ติ, อามนฺตา.\csromanspacer{} เกน เทสิโตติ, อายสฺมตา อานนฺเทน เทสิโตติ.\csromanspacer{} อายสฺมา อานนฺโท ชิโน สตฺถา สมฺมาสมฺพุทฺโธ สพฺพญฺญู สพฺพทสฺสาวี ธมฺมสฺสามี ธมฺมปฺปฏิสรโณติ, น เหวํ วตฺตพฺเพ ฯเปฯ.}

\proseitem{805}{\csromansymbolmark{*}น วตฺตพฺพํ “พุทฺเธน ภควตา ธมฺโม เทสิโต”ติ, อามนฺตา.\csromanspacer{} นนุ วุตฺตํ ภควตา “สํขิตฺเตนปิ โข อหํ สาริปุตฺต ธมฺมํ เทเสยฺยํ,}

\vspace{\tipitakaparskip}
\csromancenter{อฏฺฐารสมวคฺค วิตฺถาเรนปิ โข อหํ สาริปุตฺต ธมฺมํ เทเสยฺยํ, สํขิตฺตวิตฺถาเรนปิ โข อหํ สาริปุตฺต ธมฺมํ เทเสยฺยํ, อญฺญาตาโร จ ทุลฺลภา”ติ\footnote{อํ 1. 132 ปิฏฺเฐ สาริปุตฺตสุตฺเต.} อตฺเถว สุตฺตนฺโตติ, อามนฺตา.\csromanspacer{} เตน หิ พุทฺเธน ภควตา ธมฺโม เทสิโตติ.}
\csromanmatikaanchor{mtk.213}
\csromantocmarkref{subsection}{(179) 3. กรุณากถา}{mtk.213}
\bookmark[dest={mtk.213},level=2]{(179) 3. กรุณากถา}
\proseitem{806}{\csromanfolio{407}\csromansymbolmark{*}น วตฺตพฺพํ “พุทฺเธน ภควตา ธมฺโม เทสิโต”ติ, อามนฺตา.\csromanspacer{} นนุ วุตฺตํ ภควตา “อภิญฺญายาหํ ภิกฺขเว ธมฺมํ เทเสมิ โน อนภิญฺญาย, สนิทานาหํ ภิกฺขเว ธมฺมํ เทเสมิ โน อนิทานํ, สปฺปาฏิหาริยาหํ ภิกฺขเว ธมฺมํ เทเสมิ โน อปฺปาฏิหาริยํ.\csromanspacer{} ตสฺส\footnote{ยญฺจสฺส (สี, อิ, ก)} มยฺหํ ภิกฺขเว อภิญฺญาย ธมฺมํ เทสยโต โน อนภิญฺญาย สนิทานํ ธมฺมํ เทสยโต โน อนิทานํ สปฺปาฏิหาริยํ ธมฺมํ เทสยโต โน อปฺปาฏิหาริยํ กรณีโย โอวาโท กรณียา อนุสาสนี, อลญฺจ ปน โว ภิกฺขเว ตุฏฺฐิยา อลํ อตฺตมนตาย อลํ โสมนสฺสาย ‘สมฺมาสมฺพุทฺโธ ภควา สฺวากฺขาโต ธมฺโม สุปฺปฏิปนฺโน สํโฆ’ติ.\csromanspacer{} อิมสฺมิํ จ ปน เวยฺยากรณสฺมิํ ภญฺญมาเน ทสสหสฺสิโลกธาตุ อกมฺปิตฺถา”ติ\footnote{อํ 1. 280 ปิฏฺเฐ.} อตฺเถว สุตฺตนฺโตติ, อามนฺตา.\csromanspacer{} เตน หิ พุทฺเธน ภควตา ธมฺโม เทสิโตติ.}

\nitthitam{ธมฺมเทสนากถา นิฏฺฐิตา.}
\chapterhead{18. อฏฺฐารสมวคฺค}
\vspace{\tipitakaparskip}
\csromancenter{\textbf{(179) 3.}\csromanspacer{}\textbf{ กรุณากถา}}
\proseitem{807}{\csromansymbolmark{*}นตฺถิ พุทฺธสฺส ภควโต กรุณาติ, อามนฺตา.\csromanspacer{} นตฺถิ พุทฺธสฺส ภควโต เมตฺตาติ, น เหวํ วตฺตพฺเพ ฯเปฯ.\csromanspacer{} นตฺถิ พุทฺธสฺส ภควโต กรุณาติ, อามนฺตา.\csromanspacer{} นตฺถิ พุทฺธสฺส ภควโต มุทิตา ฯเปฯ อุเปกฺขาติ น เหวํ วตฺตพฺเพ ฯเปฯ.}

\prose{\csromansymbolmark{*}อตฺถิ พุทฺธสฺส ภควโต เมตฺตาติ, อามนฺตา.\csromanspacer{} อตฺถิ พุทฺธสฺส ภควโต กรุณาติ, น เหวํ วตฺตพฺเพ ฯเปฯ.\csromanspacer{} อตฺถิ พุทฺธสฺส ภควโต มุทิตา ฯเปฯ อุเปกฺขาติ, อามนฺตา.\csromanspacer{} อตฺถิ พุทฺธสฺส ภควโต กรุณาติ, น เหวํ วตฺตพฺเพ ฯเปฯ.}

\csromanmatikaanchor{mtk.214}
\csromantocmarkref{subsection}{(180) 4. คนฺธชาติกถา}{mtk.214}
\bookmark[dest={mtk.214},level=2]{(180) 4. คนฺธชาติกถา}
\prose{\csromanfolio{408}\csromansymbolmark{*}นตฺถิ พุทฺธสฺส ภควโต กรุณาติ, อามนฺตา.\csromanspacer{} ภควา อการุณิโกติ, น เหวํ วตฺตพฺเพ ฯเปฯ.\csromanspacer{} นนุ ภควา การุณิโก โลกหิโต โลกานุกมฺปโก โลกตฺถจโรติ, อามนฺตา.\csromanspacer{} หญฺจิ ภควา การุณิโก โลกหิโต โลกานุกมฺปโก โลกตฺถจโร, โน จ วต เร วตฺตพฺเพ “นตฺถิ พุทฺธสฺส ภควโต กรุณา”ติ ฯเปฯ.}

\prose{\csromansymbolmark{*}นตฺถิ พุทฺธสฺส ภควโต กรุณาติ, อามนฺตา.\csromanspacer{} นนุ ภควา มหากรุณาสมาปตฺติํ สมาปชฺชีติ, อามนฺตา.\csromanspacer{} หญฺจิ ภควา มหากรุณาสมาปตฺติํ สมาปชฺชิ, โน จ วต เร วตฺตพฺเพ “นตฺถิ พุทฺธสฺส ภควโต กรุณา”ติ.}

\proseitem{808}{\csromansymbolmark{+}อตฺถิ พุทฺธสฺส ภควโต กรุณาติ, อามนฺตา.\csromanspacer{} ภควา สราโคติ, น เหวํ วตฺตพฺเพ.\csromanspacer{} เตน หิ นตฺถิ พุทฺธสฺส ภควโต กรุณาติ.}

\nitthitam{กรุณากถา นิฏฺฐิตา.}
\chapterhead{18. อฏฺฐารสมวคฺค}
\vspace{\tipitakaparskip}
\csromancenter{\textbf{(180) 4.}\csromanspacer{}\textbf{ คนฺธชาติกถา}}
\proseitem{809}{\csromansymbolmark{*}พุทฺธสฺส ภควโต อุจฺจารปสฺสาโว อติวิย อญฺเญ คนฺธชาเต อธิคฺคณฺหาตีติ, อามนฺตา.\csromanspacer{} ภควา คนฺธโภชีติ, น เหวํ วตฺตพฺเพ ฯเปฯ.\csromanspacer{} นนุ ภควา โอทนกุมฺมาสํ ภุญฺชตีติ, อามนฺตา.\csromanspacer{} หญฺจิ ภควา โอทนกุมฺมาสํ ภุญฺชติ, โน จ วต เร วตฺตพฺเพ “พุทฺธสฺส ภควโต อุจฺจารปสฺสาโว อติวิย อญฺเญ คนฺธชาเต อธิคฺคณฺหาตี”ติ.}

\prose{\csromansymbolmark{*}พุทฺธสฺส ภควโต อุจฺจารปสฺสาโว อติวิย อญฺเญ คนฺธชาเต อธิคฺคณฺหาตีติ, อามนฺตา.\csromanspacer{} อตฺถิ เกจิ พุทฺธสฺส ภควโต อุจฺจารปสฺสาวํ นฺหายนฺติ วิลิมฺปนฺติ อุจฺฉาเทนฺติ\footnote{อุจฺจาเรนฺติ (สี, อิ, ก)} เปฬาย ปฏิสาเมนฺติ กรณฺฑาย นิกฺขิปนฺติ อาปเณ ปสาเรนฺติ, เตน จ คนฺเธน คนฺธกรณียํ กโรนฺตีติ, น เหวํ วตฺตพฺเพ ฯเปฯ.}

\nitthitam{คนฺธชาติกถา นิฏฺฐิตา.}
\chapterheadpageread{18. อฏฺฐารสมวคฺค \textbf{18. อฏฺฐารสมวคฺค}}
\vspace{\tipitakaparskip}
\csromancenter{\textbf{(181) 5.}\csromanspacer{}\textbf{ เอกมคฺคกถา}}
\csromanmatikaanchor{mtk.215}
\csromantocmarkref{subsection}{(181) 5. เอกมคฺคกถา}{mtk.215}
\bookmark[dest={mtk.215},level=2]{(181) 5. เอกมคฺคกถา}
\proseitem{810}{\csromanfolio{409}\csromansymbolmark{*}เอเกน อริยมคฺเคน จตฺตาริ สามญฺญผลานิ สจฺฉิกโรตีติ อามนฺตา.\csromanspacer{} จตุนฺนํ ผสฺสานํ ฯเปฯ.\csromanspacer{} จตุนฺนํ สญฺญานํ สโมธานํ โหตีติ, น เหวํ วตฺตพฺเพ ฯเปฯ.\csromanspacer{} เอเกน อริยมคฺเคน จตฺตาริ สามญฺญผลานิ สจฺฉิกโรตีติ อามนฺตา.\csromanspacer{} โสตาปตฺติมคฺเคนาติ, น เหวํ วตฺตพฺเพ ฯเปฯ.\csromanspacer{} สกทาคามิ ฯเปฯ.\csromanspacer{} อนาคามิมคฺเคนาติ, น เหวํ วตฺตพฺเพ ฯเปฯ.}

\prose{\csromansymbolmark{*}กตเมน มคฺเคนาติ, อรหตฺตมคฺเคนาติ.\csromanspacer{} อรหตฺตมคฺเคน สกฺกายทิฏฺฐิํ วิจิกิจฺฉํ สีลพฺพตปรามาสํ ชหตีติ, น เหวํ วตฺตพฺเพ ฯเปฯ.\csromanspacer{} อรหตฺตมคฺเคน สกฺกายทิฏฺฐิํ วิจิกิจฺฉํ สีลพฺพตปรามาสํ ชหตีติ, อามนฺตา.\csromanspacer{} นนุ ติณฺณํ สํโยชนานํ ปหานํ โสตาปตฺติผลํ วุตฺตํ ภควตาติ, อามนฺตา.\csromanspacer{} หญฺจิ ติณฺณํ สํโยชนานํ ปหานํ โสตาปตฺติผลํ วุตฺตํ ภควตา, โน จ วต เร วตฺตพฺเพ “อรหตฺตมคฺเคน สกฺกายทิฏฺฐิํ วิจิกิจฺฉํ สีลพฺพตปรามาสํ ชหตี”ติ.}

\prose{\csromansymbolmark{*}อรหตฺตมคฺเคน โอฬาริกํ กามราคํ โอฬาริกํ พฺยาปาทํ ชหตีติ, น เหวํ วตฺตพฺเพ ฯเปฯ.\csromanspacer{} อรหตฺตมคฺเคน โอฬาริกํ กามราคํ โอฬาริกํ พฺยาปาทํ ชหตีติ, อามนฺตา.\csromanspacer{} นนุ กามราคพฺยาปาทานํ ตนุภาวํ สกทาคามิผลํ วุตฺตํ ภควตาติ, อามนฺตา.\csromanspacer{} หญฺจิ กามราคพฺยาปาทานํ ตนุภาวํ สกทาคามิผลํ วุตฺตํ ภควตา, โน จ วต เร วตฺตพฺเพ “อรหตฺตมคฺเคน โอฬาริกํ กามราคํ โอฬาริกํ พฺยาปาทํ ชหตี”ติ.}

\prose{\csromansymbolmark{*}อรหตฺตมคฺเคน อณุสหคตํ กามราคํ อณุสหคตํ พฺยาปาทํ ชหตีติ, น เหวํ วตฺตพฺเพ ฯเปฯ.\csromanspacer{} อรหตฺตมคฺเคน อณุสหคตํ กามราคํ อณุสหคตํ พฺยาปาทํ ชหตีติ, อามนฺตา.\csromanspacer{} นนุ กามราคพฺยาปาทานํ อนวเสสปฺปหานํ อนาคามิผลํ วุตฺตํ ภควตาติ, อามนฺตา.\csromanspacer{} หญฺจิ กามราคพฺยาปาทานํ อนวเสสปฺปหานํ อนาคามิผลํ วุตฺตํ ภควตา, โน จ วต เร วตฺตพฺเพ “อรหตฺตมคฺเคน อณุสหคตํ กามราคํ อณุสหคตํ พฺยาปาทํ ชหตี”ติ.}

\csromanmatikaanchor{mtk.216}
\csromantocmarkref{subsection}{(182) 6. ฌานสงฺกนฺติกถา}{mtk.216}
\bookmark[dest={mtk.216},level=2]{(182) 6. ฌานสงฺกนฺติกถา}
\proseitem{811}{\csromansymbolmark{+}น วตฺตพฺพํ “เอเกน อริยมคฺเคน จตฺตาริ สามญฺญผลานิ สจฺฉิกโรตี”ติ, อามนฺตา.\csromanspacer{} ภควตา โสตาปตฺติมคฺโค ภาวิโตติ, \csromanfolio{410}อามนฺตา.\csromanspacer{} ภควา โสตาปนฺโนติ, น เหวํ วตฺตพฺเพ ฯเปฯ.\csromanspacer{} ภควตา สกทาคามิ ฯเปฯ อนาคามิมคฺโค ภาวิโตติ, อามนฺตา.\csromanspacer{} ภควา อนาคามีติ, น เหวํ วตฺตพฺเพ ฯเปฯ.}

\proseitem{812}{\csromansymbolmark{*}ภควา เอเกน อริยมคฺเคน จตฺตาริ สามญฺญผลานิ สจฺฉิกโรติ, สาวกา จตูหิ อริยมคฺเคหิ จตฺตาริ สามญฺญผลานิ สจฺฉิกโรนฺตีติ, อามนฺตา.\csromanspacer{} สาวกา พุทฺธสฺส ภควโต อทิฏฺฐํ ทกฺขนฺติ อนธิคตํ อธิคจฺฉนฺติ อสจฺฉิกตํ สจฺฉิกโรนฺตีติ, น เหวํ วตฺตพฺเพ ฯเปฯ.}

\nitthitam{เอกมคฺคกถา นิฏฺฐิตา.}
\chapterhead{18. อฏฺฐารสมวคฺค}
\vspace{\tipitakaparskip}
\csromancenter{\textbf{(182) 6.}\csromanspacer{}\textbf{ ฌานสงฺกนฺติกถา}}
\proseitem{813}{\csromansymbolmark{*}ฌานา ฌานํ สงฺกมตีหิ, อามนฺตา.\csromanspacer{} ปฐมา ฌานา ตติยํ ฌานํ สงฺกมตีติ, น เหวํ วตฺตพฺเพ ฯเปฯ.\csromanspacer{} ฌานา ฌานํ สงฺกมตีติ, อามนฺตา.\csromanspacer{} ทุติยา ฌานา จตุตฺถํ ฌานํ สงฺกมตีติ, น เหวํ วตฺตพฺเพ ฯเปฯ.}

\prose{\csromansymbolmark{*}ปฐมา ฌานา ทุติยํ ฌานํ สงฺกมตีติ, อามนฺตา.\csromanspacer{} ยา ปฐมสฺส ฌานสฺส อุปฺปาทาย อาวฏฺฏนา ฯเปฯ ปณิธิ, สาว ทุติยสฺส ฌานสฺส อุปฺปาทาย อาวฏฺฏนา ฯเปฯ ปณิธีติ, น เหวํ วตฺตพฺเพ ฯเปฯ.}

\prose{\csromansymbolmark{*}ปฐมา ฌานา ทุติยํ ฌานํ สงฺกมติ, น วตฺตพฺพํ “ยา ปฐมสฺส ฌานสฺส อุปฺปาทาย อาวฏฺฏนา ฯเปฯ ปณิธิ, สาว ทุติยสฺส ฌานสฺส อุปฺปาทาย อาวฏฺฏนา ฯเปฯ ปณิธี”ติ, อามนฺตา.\csromanspacer{} ทุติยํ ฌานํ อนาวฏฺเฏนฺตสฺส อุปฺปชฺชติ ฯเปฯ อปฺปณิทหนฺตสฺส อุปฺปชฺชตีติ, น เหวํ วตฺตพฺเพ ฯเปฯ.\csromanspacer{} นนุ ทุติยํ ฌานํ อาวฏฺเฏนฺตสฺส อุปฺปชฺชติ ฯเปฯ ปณิทหนฺตสฺส อุปฺปชฺชตีติ, อามนฺตา.\csromanspacer{} หญฺจิ ทุติยํ ฌานํ อาวฏฺเฏนฺตสฺส อุปฺปชฺชติ ฯเปฯ ปณิทหนฺตสฺส อุปฺปชฺชติ, โน จ วต เร วตฺตพฺเพ “ปฐมา ฌานา ทุติยํ ฌานํ สงฺกมตี”ติ ฯเปฯ.}

\prose{\csromansymbolmark{*}ปฐมา ฌานา ทุติยํ ฌานํ สงฺกมตีติ, อามนฺตา.\csromanspacer{} ปฐมํ ฌานํ กาเม อาทีนวโต มนสิกโรโต อุปฺปชฺชตีติ, อามนฺตา.\csromanspacer{} ทุติยํ ฌานํ กาเม อาทีนวโต มนสิกโรโต อุปฺปชฺชตีติ, น เหวํ วตฺตพฺเพ ฯเปฯ.}

\chapterheadpageread{18. อฏฺฐารสมวคฺค}
\prose{\csromanfolio{411}\csromansymbolmark{*}ปฐมํ ฌานํ สวิตกฺกํ สวิจารนฺติ, อามนฺตา.\csromanspacer{} ทุติยํ ฌานํ สวิตกฺกํ สวิจารนฺติ, น เหวํ วตฺตพฺเพ ฯเปฯ.\csromanspacer{} ปฐมา ฌานา ทุติยํ ฌานํ สงฺกมตีติ, อามนฺตา.\csromanspacer{} ตญฺเญว ปฐมํ ฌานํ ตํ ทุติยํ ฌานนฺติ, น เหวํ วตฺตพฺเพ ฯเปฯ.}

\proseitem{814}{\csromansymbolmark{*}ทุติยา ฌานา ตติยํ ฌานํ สงฺกมตีติ, อามนฺตา.\csromanspacer{} ยา ทุติยสฺส ฌานสฺส อุปฺปาทาย อาวฏฺฏนา ฯเปฯ ปณิธิ, สาว ตติยสฺส ฌานสฺส อุปฺปาทาย อาวฏฺฏนา ฯเปฯ ปณิธีติ, น เหวํ วตฺตพฺเพ ฯเปฯ.}

\prose{\csromansymbolmark{*}ทุติยา ฌานา ตติยํ ฌานํ สงฺกมติ, น วตฺตพฺพํ “ยา ทุติยสฺส ฌานสฺส อุปฺปาทาย อาวฏฺฏนา ฯเปฯ ปณิธิ, สาว ตติยสฺส ฌานสฺส อุปฺปาทาย อาวฏฺฏนา ฯเปฯ ปณิธี”ติ, อามนฺตา.\csromanspacer{} ตติยํ ฌานํ อนาวฏฺเฏนฺตสฺส อุปฺปชฺชติ ฯเปฯ อปฺปณิทหนฺตสฺส อุปฺปชฺชตีติ, น เหวํ วตฺตพฺเพ ฯเปฯ.\csromanspacer{} นนุ ตติยํ ฌานํ อาวฏฺเฏนฺตสฺส อุปฺปชฺชติ ฯเปฯ ปณิทหนฺตสฺส อุปฺปชฺชตีติ, อามนฺตา.\csromanspacer{} หญฺจิ ตติยํ ฌานํ อาวฏฺเฏนฺตสฺส อุปฺปชฺชติ ฯเปฯ ปณิทหนฺตสฺส อุปฺปชฺชติ, โน จ วต เร วตฺตพฺเพ “ทุติยา ฌานา ตติยํ ฌานํ สงฺกมตี”ติ.}

\prose{\csromansymbolmark{*}ทุติยา ฌานา ตติยํ ฌานํ สงฺกมตีติ, อามนฺตา.\csromanspacer{} ทุติยํ ฌานํ วิตกฺกวิจาเร อาทีนวโต มนสิกโรโต อุปฺปชฺชตีติ, อามนฺตา.\csromanspacer{} ตติยํ ฌานํ วิตกฺกวิจาเร อาทีนวโต มนสิกโรโต อุปฺปชฺชตีติ, น เหวํ วตฺตพฺเพ ฯเปฯ.}

\prose{\csromansymbolmark{*}ทุติยํ ฌานํ สปฺปีติกนฺติ, อามนฺตา.\csromanspacer{} ตติยํ ฌานํ สปฺปีติกนฺติ, น เหวํ วตฺตพฺเพ ฯเปฯ.\csromanspacer{} ทุติยา ฌานา ตติยํ ฌานํ สงฺกมตีติ, อามนฺตา.\csromanspacer{} ตญฺเญว ทุติยํ ฌานํ ตํ ตติยํ ฌานนฺติ, น เหวํ วตฺตพฺเพ ฯเปฯ.}

\proseitem{815}{\csromansymbolmark{*}ตติยา ฌานา จตุตฺถํ ฌานํ สงฺกมตีติ, อามนฺตา.\csromanspacer{} ยา ตติยสฺส ฌานสฺส อุปฺปาทาย อาวฏฺฏนา ฯเปฯ ปณิธิ, สาว จตุตฺถสฺส ฌานสฺส อุปฺปาทาย อาวฏฺฏนา ฯเปฯ ปณิธีติ, น เหวํ วตฺตพฺเพ ฯเปฯ.}

\csromanmatikaanchor{mtk.217}
\csromantocmarkref{subsection}{(183) 7. ฌานนฺตริกกถา}{mtk.217}
\bookmark[dest={mtk.217},level=2]{(183) 7. ฌานนฺตริกกถา}
\prose{ตติยา ฌานา จตุตฺถํ ฌานํ สงฺกมติ, น วตฺตพฺพํ “ยา ตติยสฺส ฌานสฺส อุปฺปาทาย อาวฏฺฏนา ฯเปฯ ปณิธิ, สาว จตุตฺถสฺส ฌานสฺส อุปฺปาทาย อาวฏฺฏนา ฯเปฯ ปณิธี”ติ, อามนฺตา.\csromanspacer{} จตุตฺถํ ฌานํ อนาวฏฺเฏนฺตสฺส อุปฺปชฺชติ ฯเปฯ อปฺปณิทหนฺตสฺส อุปฺปชฺชตีติ, น เหวํ วตฺตพฺเพ ฯเปฯ.\csromanspacer{} นนุ จตุตฺถํ ฌานํ อาวฏฺเฏนฺตสฺส อุปฺปชฺชติ ฯเปฯ ปณิทหนฺตสฺส อุปฺปชฺชตีติ, อามนฺตา. \csromanfolio{412}หญฺจิ จตุตฺถํ ฌานํ อาวฏฺเฏนฺตสฺส อุปฺปชฺชติ ฯเปฯ ปณิทหนฺตสฺส อุปฺปชฺชติ, โน จ วต เร วตฺตพฺเพ “ตติยา ฌานา จตุตฺถํ ฌานํ สงฺกมตี”ติ.}

\prose{\csromansymbolmark{*}ตติยา ฌานา จตุตฺถํ ฌานํ สงฺกมตีติ, อามนฺตา.\csromanspacer{} ตติยํ ฌานํ ปีติํ อาทีนวโต มนสิกโรโต อุปฺปชฺชตีติ, อามนฺตา.\csromanspacer{} จตุตฺถํ ฌานํ ปีติํ อาทีนวโต มนสิกโรโต อุปฺปชฺชตีติ, น เหวํ วตฺตพฺเพ ฯเปฯ.}

\prose{\csromansymbolmark{*}ตติยํ ฌานํ สุขสหคตนฺติ, อามนฺตา.\csromanspacer{} จตุตฺถํ ฌานํ สุขสหคตนฺติ, น เหวํ วตฺตพฺเพ ฯเปฯ.\csromanspacer{} ตติยา ฌานา จตุตฺถํ ฌานํ สงฺกมตีติ, อามนฺตา.\csromanspacer{} ตญฺเญว ตติยํ ฌานํ ตํ จตุตฺถํ ฌานนฺติ.\csromanspacer{} น เหวํ วตฺตพฺเพ ฯเปฯ.}

\proseitem{816}{\csromansymbolmark{+}น วตฺตพฺพํ “ฌานา ฌานํ สงฺกมตี”ติ, อามนฺตา.\csromanspacer{} นนุ วุตฺตํ ภควตา “อิธ ภิกฺขเว ภิกฺขุ วิวิจฺเจว กาเมหิ ฯเปฯ.\csromanspacer{} จตุตฺถํ ฌานํ อุปสมฺปชฺช วิหรตี”ติ\footnote{อํ 1. 55, 469 ปิฏฺเฐสุ.} อตฺเถว สุตฺตนฺโตติ, อามนฺตา.\csromanspacer{} เตน หิ ฌานา ฌานํ สงฺกมตีติ.}

\nitthitam{ฌานสงฺกนฺติกถา นิฏฺฐิตา.}
\chapterhead{18. อฏฺฐารสมวคฺค}
\vspace{\tipitakaparskip}
\csromancenter{\textbf{(183) 7.}\csromanspacer{}\textbf{ ฌานนฺตริกกถา}}
\proseitem{817}{\csromansymbolmark{*}อตฺถิ ฌานนฺตริกาติ, อามนฺตา.\csromanspacer{} อตฺถิ ผสฺสนฺตริกา ฯเปฯ.\csromanspacer{} อตฺถิ สญฺญนฺตริกาติ, น เหวํ วตฺตพฺเพ ฯเปฯ.}

\prose{\csromansymbolmark{*}อตฺถิ ฌานนฺตริกาติ, อามนฺตา.\csromanspacer{} ทุติยสฺส จ ฌานสฺส ตติยสฺส จ ฌานสฺส อนฺตเร อตฺถิ ฌานนฺตริกาติ, น เหวํ วตฺตพฺเพ ฯเปฯ.}

\prose{\csromansymbolmark{*}อตฺถิ ฌานนฺตริกาติ, อามนฺตา.\csromanspacer{} ตติยสฺส จ ฌานสฺส จตุตฺถสฺส จ ฌานสฺส อนฺตเร อตฺถิ ฌานนฺตริกาติ, น เหวํ วตฺตพฺเพ ฯเปฯ.}

\prose{\csromansymbolmark{*}ทุติยสฺส จ ฌานสฺส ตติยสฺส จ ฌานสฺส อนฺตเร นตฺถิ ฌานนฺตริกาติ, อามนฺตา.\csromanspacer{} หญฺจิ ทุติยสฺส จ ฌานสฺส ตติยสฺส จ ฌานสฺส อนฺตเร นตฺถิ ฌานนฺตริกา, โน จ วต เร วตฺตพฺเพ “อตฺถิ ฌานนฺตริกา”ติ.}

\chapterheadpageread{18. อฏฺฐารสมวคฺค}
\prose{\csromanfolio{413}\csromansymbolmark{*}ตติยสฺส จ ฌานสฺส จตุตฺถสฺส จ ฌานสฺส อนฺตเร นตฺถิ ฌานนฺตริกาติ, อามนฺตา.\csromanspacer{} หญฺจิ ตติยสฺส จ ฌานสฺส จตุตฺถสฺส จ ฌานสฺส อนฺตเร นตฺถิ ฌานนฺตริกา, โน จ วต เร วตฺตพฺเพ “อตฺถิ ฌานนฺตริกา”ติ.}

\proseitem{818}{\csromansymbolmark{*}ปฐมสฺส จ ฌานสฺส ทุติยสฺส จ ฌานสฺส อนฺตเร อตฺถิ ฌานนฺตริกาติ, อามนฺตา.\csromanspacer{} ทุติยสฺส จ ฌานสฺส ตติยสฺส จ ฌานสฺส อนฺตเร อตฺถิ ฌานนฺตริกาติ, น เหวํ วตฺตพฺเพ ฯเปฯ.}

\prose{\csromansymbolmark{*}ปฐมสฺส จ ฌานสฺส ทุติยสฺส จ ฌานสฺส อนฺตเร อตฺถิ ฌานนฺตริกาติ, อามนฺตา.\csromanspacer{} ตติยสฺส จ ฌานสฺส จตุตฺถสฺส จ ฌานสฺส อนฺตเร อตฺถิ ฌานนฺตริกาติ, น เหวํ วตฺตพฺเพ ฯเปฯ.}

\prose{\csromansymbolmark{*}ทุติยสฺส จ ฌานสฺส ตติยสฺส จ ฌานสฺส อนฺตเร นตฺถิ ฌานนฺตริกาติ, อามนฺตา.\csromanspacer{} ปฐมสฺส จ ฌานสฺส ทุติยสฺส จ ฌานสฺส อนฺตเร นตฺถิ ฌานนฺตริกาติ, น เหวํ วตฺตพฺเพ ฯเปฯ.}

\prose{\csromansymbolmark{*}ตติยสฺส จ ฌานสฺส จตุตฺถสฺส จ ฌานสฺส อนฺตเร นตฺถิ ฌานนฺตริกาติ, อามนฺตา.\csromanspacer{} ปฐมสฺส จ ฌานสฺส ทุติยสฺส จ ฌานสฺส อนฺตเร นตฺถิ ฌานนฺตริกาติ, น เหวํ วตฺตพฺเพ ฯเปฯ.}

\proseitem{819}{อวิตกฺโก วิจารมตฺโต สมาธิ ฌานนฺตริกาติ, อมนฺตา.\csromanspacer{} สวิตกฺโก สวิจาโร สมาธิ ฌานนฺตริกาติ, น เหวํ วตฺตพฺเพ ฯเปฯ.}

\prose{\csromansymbolmark{*}อวิตกฺโก วิจารมตฺโต สมาธิ ฌานนฺตริกาติ, อามนฺตา.\csromanspacer{} อวิตกฺโก อวิจาโร สมาธิ ฌานนฺตริกาติ, น เหวํ วตฺตพฺเพ ฯเปฯ.}

\prose{\csromansymbolmark{*}สวิตกฺโก สวิจาโร สมาธิ น ฌานนฺตริกาติ, อามนฺตา.\csromanspacer{} อวิตกฺโก วิจารมตฺโต สมาธิ น ฌานนฺตริกาติ, เน เหวํ วตฺตพฺเพ ฯเปฯ.}

\prose{\csromansymbolmark{*}อวิตกฺโก อวิจาโร สมาธิ น ฌานนฺตริกาติ, อามนฺตา.\csromanspacer{} อวิตกฺโก วิจารมตฺโต สมาธิ น ฌานนฺตริกาติ, น เหวํ วตฺตพฺเพ ฯเปฯ.}

\csromanmatikaanchor{mtk.218}
\csromantocmarkref{subsection}{(184) 8. สทฺทํสุณาตีติกถา}{mtk.218}
\bookmark[dest={mtk.218},level=2]{(184) 8. สทฺทํสุณาตีติกถา}
\proseitem{820}{\csromansymbolmark{*}ทฺวินฺนํ ฌานานํ ปฏุปฺปนฺนานมนฺตเร อวิตกฺโก วิจารมตฺโต สมาธีติ, อามนฺตา.\csromanspacer{} นนุ อวิตกฺเก วิจารมตฺเต สมาธิมฺหิ วตฺตมาเน ปฐมํ ฌานํ นิรุทฺธํ ทุติยํ ฌานํ ปฏุปฺปนฺนนฺติ, อามนฺตา.\csromanspacer{} หญฺจิ อวิตกฺเก \csromanfolio{414}วิจารมตฺเต สมาธิมฺหิ วตฺตมาเน ปฐมํ ฌานํ นิรุทฺธํ ทุติยํ ฌานํ ปฏุปฺปนฺนํ, โน จ วต เร วตฺตพฺเพ “ทฺวินฺนํ ฌานานํ ปฏุปฺปนฺนานมนฺตเร อวิตกฺโก วิจารมตฺโต สมาธิ ฌานนฺตริกา”ติ.}

\proseitem{821}{\csromansymbolmark{+}อวิตกฺโก วิจารมตฺโต สมาธิ น ฌานนฺตริกาติ, อามนฺตา.\csromanspacer{} อวิตกฺโก วิจารมตฺโต สมาธิ ปฐมํ ฌานํ ฯเปฯ ทุติยํ ฌานํ ฯเปฯ ตติยํ ฌานํ ฯเปฯ จตุตฺถํ ฌานนฺติ, น เหวํ วตฺตพฺเพ ฯเปฯ.\csromanspacer{} เตน หิ อวิตกฺโก วิจารมตฺโต สมาธิ ฌานนฺตริกาติ.}

\proseitem{822}{\csromansymbolmark{*}อวิตกฺโก วิจารมตฺโต สมาธิ ฌานนฺตริกาติ, อามนฺตา.\csromanspacer{} นนุ ตโย สมาธี วุตฺตา ภควตา “สวิตกฺโก สวิจาโร สมาธิ, อวิตกฺโก วิจารมตฺโต สมาธิ, อวิตกฺโก อวิจาโร สมาธี”ติ\footnote{ที 3. 184, 229 ปิฏฺเฐสุ.}, อามนฺตา.\csromanspacer{} หญฺจิ ตโย สมาธี วุตฺตา ภควตา “สวิตกฺโก ฯเปฯ อวิจาโร สมาธิ”, โน จ วต เร วตฺตพฺเพ “อวิตกฺโก วิจารมตฺโต สมาธิ ฌานนฺตริกา”ติ.}

\nitthitam{ฌานนฺตริกกถา นิฏฺฐิตา.}
\chapterhead{18. อฏฺฐารสมวคฺค}
\vspace{\tipitakaparskip}
\csromancenter{\textbf{(184) 8.}\csromanspacer{}\textbf{ สทฺทํสุณาตีติกถา}}
\proseitem{823}{\csromansymbolmark{*}สมาปนฺโน สทฺทํ สุณาตีติ, อามนฺตา.\csromanspacer{} สมาปนฺโน จกฺขุนา รูปํ ปสฺสติ ฯเปฯ โสเตน ฯเปฯ ฆาเนน ฯเปฯ ชิวฺหาย ฯเปฯ กาเยน โผฏฺฐพฺพํ ผุสตีติ, น เหวํ วตฺตพฺเพ ฯเปฯ.}

\prose{\csromansymbolmark{*}สมาปนฺโน สทฺทํ สุณาตีติ, อามนฺตา.\csromanspacer{} โสตวิญฺญาณสมงฺคี สมาปนฺโนติ, น เหวํ วตฺตพฺเพ.\csromanspacer{} นนุ สมาธิ มโนวิญฺญาณสมงฺคีสฺสาติ, อามนฺตา.\csromanspacer{} หญฺจิ สมาธิ มโนวิญฺญาณสมงฺคิสฺส, โน จ วต เร วตฺตพฺเพ “สมาปนฺโน สทฺทํ สุณาตี”ติ.}

\prose{\csromansymbolmark{*}สมาธิ มโนวิญฺญาณสมงฺคิสฺส, โสตวิญฺญาณสมงฺคี สทฺทํ สุณาตีติ, อามนฺตา.\csromanspacer{} หญฺจิ สมาธิ มโนวิญฺญาณสมงฺคีสฺส, โสตวิญฺญาณสมงฺคี สทฺทํ สุณาติ, โน จ วต เร วตฺตพฺเพ “สมาปนฺโน สทฺทํ}

\vspace{\tipitakaparskip}
\csromancenter{อฏฺฐารสมวคฺค สุณาตี”ติ.\csromanspacer{} สมาธิ มโนวิญฺญาณสมงฺคิสฺส, โสตวิญฺญาณสมงฺคี สทฺทํ สุณาตีติ, อามนฺตา.\csromanspacer{} ทฺวินฺนํ ผสฺสานํ ฯเปฯ.\csromanspacer{} ทฺวินฺนํ จิตฺตานํ สโมธานํ โหตีติ, น เหวํ วตฺตพฺเพ ฯเปฯ.}
\csromanmatikaanchor{mtk.219}
\csromantocmarkref{subsection}{(185) 9. จกฺขุนารูปํปสฺสตีติกถา}{mtk.219}
\bookmark[dest={mtk.219},level=2]{(185) 9. จกฺขุนารูปํปสฺสตีติกถา}
\proseitem{824}{\csromanfolio{415}\csromansymbolmark{+}น วตฺตพฺพํ “สมาปนฺโน สทฺทํ สุณาตี”ติ, อามนฺตา.\csromanspacer{} นนุ ปฐมสฺส ฌานสฺส สทฺโท กณฺฑโก วุตฺโต ภควตาติ, อามนฺตา.\csromanspacer{} หญฺจิ ปฐมสฺส ฌานสฺส สทฺโท กณฺฑโก วุตฺโต ภควตา, เตน วต เร วตฺตพฺเพ “สมาปนฺโน สทฺทํ สุณาตี”ติ.}

\proseitem{825}{\csromansymbolmark{*}ปฐมสฺส ฌานสฺส สทฺโท กณฺฑโก วุตฺโต ภควตาติ สมาปนฺโน สทฺทํ สุณาตีติ, อามนฺตา.\csromanspacer{} ทุติยสฺส ฌานสฺส วิตกฺโก วิจาโร กณฺฑโก วุตฺโต ภควตา อตฺถิ ตสฺส วิตกฺกวิจาราติ, น เหวํ วตฺตพฺเพ ฯเปฯ.}

\prose{\csromansymbolmark{*}ปฐมสฺส ฌานสฺส สทฺโท กณฺฑโก วุตฺโต ภควตาติ สมาปนฺโน สทฺทํ สุณาตีติ, อามนฺตา.\csromanspacer{} ตติยสฺส ฌานสฺส ปีติ กณฺฑโก ฯเปฯ.\csromanspacer{} จตุตฺถสฺส ฌานสฺส อสฺสาสปสฺสาโส กณฺฑโก.\csromanspacer{} อากาสานญฺจายตนํ สมาปนฺนสฺส รูปสญฺญา กณฺฑโก.\csromanspacer{} วิญฺญาณญฺจายตนํ สมาปนฺนสฺส อากาสานญฺจายตนสญฺญา กณฺฑโก.\csromanspacer{} อากิญฺจญฺญายตนํ สมาปนฺนสฺส วิญฺญาณญฺจายตนสญฺญา กณฺฑโก.\csromanspacer{} เนวสญฺญานาสญฺญายตนํ สมาปนฺนสฺส อากิญฺจญฺญายตนสญฺญา กณฺฑโก.\csromanspacer{} สญฺญาเวทยิตนิโรธํ สมาปนฺนสฺส สญฺญา จ เวทนา จ กณฺฑโก วุตฺโต ภควตา อตฺถิ ตสฺส สญฺญา จ เวทนา จาติ, น เหวํ วตฺตพฺเพ ฯเปฯ.}

\nitthitam{สทฺทํ สุณาตีติกถา นิฏฺฐิตา.}
\chapterhead{18. อฏฺฐารสมวคฺค}
\vspace{\tipitakaparskip}
\csromancenter{\textbf{(185) 9.}\csromanspacer{}\textbf{ จกฺขุนารูปํปสฺสตีติกถา}}
\proseitem{826}{\csromansymbolmark{*}จกฺขุนา รูปํ ปสฺสตีติ, อามนฺตา.\csromanspacer{} รูเปน รูปํ ปสฺสตีติ, น เหวํ วตฺตพฺเพ ฯเปฯ.\csromanspacer{} รูเปน รูปํ ปสฺสตีติ, อามนฺตา.\csromanspacer{} รูเปน รูปํ ปฏิวิชานาตีติ, น เหวํ วตฺตพฺเพ ฯเปฯ.\csromanspacer{} รูเปน รูปํ ปฏิวิชานาตีติ, อามนฺตา.\csromanspacer{} รูปํ \csromanfolio{416}มโนวิญฺญาณนฺติ, น เหวํ วตฺตพฺเพ ฯเปฯ.\csromanspacer{} จกฺขุนา รูปํ ปสฺสตีติ, อามนฺตา.\csromanspacer{} อตฺถิ จกฺขุสฺส อาวฏฺฏนา ฯเปฯ ปณิธีติ, น เหวํ วตฺตพฺเพ ฯเปฯ.\csromanspacer{} นนุ นตฺถิ จกฺขุสฺส อาวฏฺฏนา ฯเปฯ ปณิธีติ, อามนฺตา.\csromanspacer{} หญฺจิ นตฺถิ จกฺขุสฺส อาวฏฺฏนา ฯเปฯ ปณิธิ, โน จ วต เร วตฺตพฺเพ “จกฺขุนา รูปํ ปสฺสตี”ติ.}

\prose{\csromansymbolmark{*}โสเตน สทฺทํ สุณาตีติ ฯเปฯ.\csromanspacer{} ฆาเนน คนฺธํ ฆายตีติ ฯเปฯ.\csromanspacer{} ชิวฺหาย รสํ สายตีติ ฯเปฯ.\csromanspacer{} กาเยน โผฏฺฐพฺพํ ผุสตีติ, อามนฺตา.\csromanspacer{} รูเปน รูปํ ผุสตีติ, น เหวํ วตฺตพฺเพ ฯเปฯ.}

\prose{\csromansymbolmark{*}รูเปน รูปํ ผุสตีติ, อามนฺตา.\csromanspacer{} รูเปน รูปํ ปฏิวิชานาตีติ, น เหวํ วตฺตพฺเพ ฯเปฯ.\csromanspacer{} รูเปน รูปํ ปฏิวิชานาตีติ, อามนฺตา.\csromanspacer{} รูปํ มโนวิญฺญาณนฺติ, น เหวํ วตฺตพฺเพ ฯเปฯ.\csromanspacer{} กาเยน โผฏฺฐพฺพํ ผุสตีติ, อามนฺตา.\csromanspacer{} อตฺถิ กายสฺส อาวฏฺฏนา ฯเปฯ ปณิธีติ, น เหวํ วตฺตพฺเพ ฯเปฯ.\csromanspacer{} นนุ นตฺถิ กายสฺส อาวฏฺฏนา ฯเปฯ ปณิธีติ, อามนฺตา.\csromanspacer{} หญฺจิ นตฺถิ กายสฺส อาวฏฺฏนา ฯเปฯ ปณิธิ, โน จ วต เร วตฺตพฺเพ “กาเยน โผฏฺฐพฺพํ ผุสตี”ติ ฯเปฯ.}

\proseitem{827}{\csromansymbolmark{+}น วตฺตพฺพํ “จกฺขุนา รูปํ ปสฺสตี”ติ ฯเปฯ “กาเยน โผฏฺฐพฺพํ ผุสตี”ติ, อามนฺตา.\csromanspacer{} นนุ วุตฺตํ ภควตา “อิธ ภิกฺขเว ภิกฺขุ จกฺขุนา รูปํ ปสฺสติ ฯเปฯ กาเยน โผฏฺฐพฺพํ ผุสตี”ติ\footnote{ม 1. 285; อํ 1. 348 ปิฏฺเฐสุ (อฏฺฐกถา ปสฺสิตพฺพา.)} อตฺเถว สุตฺตนฺโตติ, อามนฺตา.\csromanspacer{} เตน หิ จกฺขุนา รูปํ ปสฺสติ ฯเปฯ กาเยน โผฏฺฐพฺพํ ผุสตีติ.\csromanspacer{} จกฺขุนา รูปํ ปสฺสตีติกถา นิฏฺฐิตา.\csromanspacer{} อฏฺฐารสโม วคฺโค.}

\titlehead{\textbf{ตสฺสุทฺทานํ}}
\prose{พุทฺโธ ภควา มนุสฺสโลเก อฏฺฐาสิ, พุทฺเธน ภควตา ธมฺโม เทสิโต, นตฺถิ พุทฺธสฺส ภควโต กรุณา, พุทฺธสฺส ภควโต อุจฺจารปสฺสาโว อติวิย อญฺเญ คนฺธชาเต อธิคฺคณฺหาติ, เอเกน อริยมคฺเคน จตฺตาริ สามญฺญผลานิ สจฺฉิกโรติ, ฌานา ฌานํ สงฺกมติ, อตฺถิ ฌานนฺตริกา, สมาปนฺโน สทฺทํ สุณาติ, จกฺขุนา รูปํ ปสฺสติ กาเยน โผฏฺฐพฺพํ ผุสติ.}

\csromanmatikaanchor{mtk.220}
\csromantocmarknopage{chapter}{19. เอกูนวีสติมวคฺค}{mtk.220}
\bookmark[dest={mtk.220},level=0]{19. เอกูนวีสติมวคฺค}
\chapterheadpageread{19. เอกูนวีสติมวคฺค \textbf{19. เอกูนวีสติมวคฺค}}
\vspace{\tipitakaparskip}
\csromancenter{\textbf{(186) 1.}\csromanspacer{}\textbf{ กิเลสชหนกถา}}
\csromanmatikaanchor{mtk.221}
\csromantocmarkref{subsection}{(186) 1. กิเลสชหนกถา}{mtk.221}
\bookmark[dest={mtk.221},level=2]{(186) 1. กิเลสชหนกถา}
\proseitem{828}{\csromanfolio{417}\csromansymbolmark{*}อตีเต กิเลเส ชหตีติ, อามนฺตา.\csromanspacer{} นิรุทฺธํ นิโรเธติ วิคตํ วิคเมติ ขีณํ เขเปติ อตฺถงฺคตํ อตฺถงฺคเมติ อพฺภตฺถงฺคตํ อพฺภตฺถงฺคเมตีติ, น เหวํ วตฺตพฺเพ ฯเปฯ.\csromanspacer{} อตีเต กิเลเส ชหตีติ, อามนฺตา.\csromanspacer{} นนุ อตีตํ นิรุทฺธนฺติ, อามนฺตา.\csromanspacer{} หญฺจิ อตีตํ นิรุทฺธํ, โน จ วต เร วตฺตพฺเพ “อตีเต กิเลเส ชหตี”ติ.\csromanspacer{} อตีเต กิเลเส ชหตีติ, อามนฺตา.\csromanspacer{} นนุ อตีตํ นตฺถีติ, อามนฺตา.\csromanspacer{} หญฺจิ อตีตํ นตฺถิ, โน จ วต เร วตฺตพฺเพ “อตีเต กิเลเส ชหตี”ติ.}

\proseitem{829}{\csromansymbolmark{*}อนาคเต กิเลเส ชหตีติ, อามนฺตา.\csromanspacer{} อชาตํ อชเนติ อสญฺชาตํ อสญฺชเนติ อนิพฺพตฺตํ อนิพฺพตฺเตติ อปาตุภูตํ อปาตุภาเวตีติ, น เหวํ วตฺตพฺเพ ฯเปฯ.\csromanspacer{} อนาคเต กิเลเส ชหตีติ, อามนฺตา.\csromanspacer{} นนุ อนาคตํ อชาตนฺติ, อามนฺตา.\csromanspacer{} หญฺจิ อนาคตํ อชาตํ, โน จ วต เร วตฺตพฺเพ “อนาคเต กิเลเส ชหตี”ติ.\csromanspacer{} อนาคเต กิเลเส ชหตีติ, อามนฺตา.\csromanspacer{} นนุ อนาคตํ นตฺถีติ, อามนฺตา.\csromanspacer{} หญฺจิ อนาคตํ นตฺถิ, โน จ วต เร วตฺตพฺเพ “อนาคเต กิเลเส ชหตี”ติ.}

\proseitem{830}{\csromansymbolmark{*}ปจฺจุปฺปนฺเน กิเลเส ชหตีติ, อามนฺตา.\csromanspacer{} รตฺโต ราคํ ชาหติ, ทุฏฺโฐ โทสํ ชหติ, มูโฬฺห โมหํ ชหติ, กิลิฏฺโฐ โกเลเส ชหตีติ, น เหวํ วตฺตพฺเพ ฯเปฯ.\csromanspacer{} ราเคน รูคํ ชหติ, โทเสน โทสํ ชหติ, โมเหน โมหํ ชหติ, กิเลเสหิ กิเลเส ชหตีติ, น เหวํ วตฺตพฺเพ ฯเปฯ.}

\prose{\csromansymbolmark{*}ราโค จิตฺตสมฺปยุตฺโต, มคฺโค จิตฺตสมฺปยุตฺโตติ, อามนฺตา.\csromanspacer{} ทฺวินฺนํ จิตฺตานํ สโมธานํ โหตีติ, น เหวํ วตฺตพฺเพ ฯเปฯ.}

\prose{\csromansymbolmark{*}ราโค อกุสโล, มคฺโค กุสโลติ, อามนฺตา.\csromanspacer{} กุสลากุสลา สาวชฺชานวชฺชา หีนปณีตา กณฺหสุกฺกสปฺปฏิภาคา ธมฺมา สมฺมุขีภาวํ อาคจฺฉนฺตีติ, น เหวํ วตฺตพฺเพ ฯเปฯ.}

\csromanmatikaanchor{mtk.222}
\csromantocmarkref{subsection}{(187) 2. สุญฺญตากถา}{mtk.222}
\bookmark[dest={mtk.222},level=2]{(187) 2. สุญฺญตากถา}
\prose{\csromanfolio{418}\csromansymbolmark{*}กุสลากุสลา สาวชฺชานวชฺชา หีนปณีตา กณฺหสุกฺกสปฺปฏิภาคา ธมฺมา สมฺมุขีภาวํ อาคจฺฉนฺตีติ, อามนฺตา.\csromanspacer{} นนุ วุตฺตํ ภควตา “จตฺตาริมานิ ภิกฺขเว สุวิทูรวิทูรานิ.\csromanspacer{} กตมานิ จตฺตาริ, นภญฺจ ภิกฺขเว ปถวี จ อิทํ ปฐมํ สุวิทูรวิทูรํ ฯเปฯ.\csromanspacer{} ตสฺมา สตํ ธมฺโม อสพฺภิ อารกา”ติ อตฺเถว สุตฺตนฺโตติ, อามนฺตา.\csromanspacer{} เตน หิ น วตฺตพฺพํ “กุสลากุสลา ฯเปฯ สมฺมุขีภาวํ อาคจฺฉนฺตี”ติ.}

\proseitem{831}{\csromansymbolmark{+}น วตฺตพฺพํ “อตีเต กิเลเส ชหติ อนาคเต กิเลเส ชหติ ปจฺจุปฺปนฺเน กิเลเส ชหตี”ติ, อามนฺตา.\csromanspacer{} นตฺถิ กิเลเส ชหตี”ติ, น เหวํ วตฺตพฺเพ ฯเปฯ.\csromanspacer{} เตน หิ อตีเต กิเลเส ชหติ, อนาคเต กิเลเส ชหติ, ปจฺจุปฺปนฺเน กิเลเส ชหตีติ.}

\nitthitam{กิเลสชหนกถา นิฏฺฐิตา.}
\chapterhead{19. เอกูนวีสติมวคฺค}
\vspace{\tipitakaparskip}
\csromancenter{\textbf{(187) 2.}\csromanspacer{}\textbf{ สุญฺญตากถา}}
\proseitem{832}{\csromansymbolmark{*}สุญฺญตา สงฺขารกฺขนฺธปริยาปนฺนาติ, อามนฺตา.\csromanspacer{} อนิมิตฺตํ สงฺขารกฺขนฺธปริยาปนฺนนฺติ, น เหวํ วตฺตพฺเพ ฯเปฯ.\csromanspacer{} สุญฺญตา สงฺขารกฺขนฺธปริยาปนฺนาติ, อามนฺตา.\csromanspacer{} อปฺปณิหิโต สงฺขารกฺขนฺธปริยาปนฺโนติ, น เหวํ วตฺตพฺเพ ฯเปฯ.\csromanspacer{} อนิมิตฺตํ น วตฺตพฺพํ “สงฺขารกฺขนฺธปริยาปนฺนนฺ”ติ, อามนฺตา.\csromanspacer{} สุญฺญตา น วตฺตพฺพา “สงฺขารกฺขนฺธปริยาปนฺนา”ติ, น เหวํ วตฺตพฺเพ ฯเปฯ.\csromanspacer{} อปฺปณิหิโต น วตฺตพฺโพ “สงฺขารกฺขนฺธปริยาปนฺโน”ติ, อามนฺตา.\csromanspacer{} สุญฺญตา น วตฺตพฺพา “สงฺขารกฺขนฺธปริยาปนฺนา”ติ, น เหวํ วตฺตพฺเพ ฯเปฯ.}

\prose{\csromansymbolmark{*}สุญฺญตา “สงฺขารกฺขนฺธปริยาปนฺนา”ติ, อามนฺตา.\csromanspacer{} สงฺขารกฺขนฺโธ น อนิจฺโจ น สงฺขโต น ปฏิจฺจสมุปฺปนฺโน น ขยธมฺโม น วยธมฺโม น วิราคธมฺโม น นิโรธธมฺโม น วิปริณามธมฺโมติ, น เหวํ วตฺตพฺเพ ฯเปฯ.\csromanspacer{} นนุ สงฺขารกฺขนฺโธ อนิจฺโจ สงฺขโต ปฏิจฺจสมุปฺปนฺโน ขยธมฺโม วยธมฺโม วิราคธมฺโม นิโรธธมฺโม วิปริณามธมฺโมติ, อามนฺตา.\csromanspacer{} หญฺจิ สงฺขารกฺขนฺโธ อนิจฺโจ ฯเปฯ วิปริณามธมฺโม, โน จ วต เร วตฺตพฺเพ “สุญฺญตา สงฺขารกฺขนฺธปริยาปนฺนา”ติ.}

\chapterheadpageread{19. เอกูนวีสติมวคฺค}
\csromanmatikaanchor{mtk.223}
\csromantocmarkref{subsection}{(188) 3. สามญฺญผลกถา}{mtk.223}
\bookmark[dest={mtk.223},level=2]{(188) 3. สามญฺญผลกถา}
\proseitem{833}{\csromanfolio{419}\csromansymbolmark{*}รูปกฺขนฺธสฺส สุญฺญตา สงฺขารกฺขนฺธปริยาปนฺนาติ, อามนฺตา.\csromanspacer{} สงฺขารกฺขนฺธสฺส สุญฺญตา รูปกฺขนฺธปริยาปนฺนาติ, น เหวํ วตฺตพฺเพ ฯเปฯ.\csromanspacer{} เวทนากฺขนฺธสฺส ฯเปฯ.\csromanspacer{} สญฺญากฺขนฺธสฺส.\csromanspacer{} วิญฺญาณกฺขนฺธสฺส สุญฺญตา สงฺขารกฺขนฺธปริยาปนฺนาติ, อามนฺตา.\csromanspacer{} สงฺขารกฺขนฺธสฺส สุญฺญตา วิญฺญณกฺขนฺธปริยาปนฺนาติ, น เหวํ วตฺตพฺเพ ฯเปฯ.}

\prose{\csromansymbolmark{*}สงฺขารกฺขนฺธสฺส สุญฺญตา น วตฺตพฺพา “รูปกฺขนฺธปริยาปนฺนา”ติ, อามนฺตา.\csromanspacer{} รูปกฺขนฺธสฺส สุญฺญตา น วตฺตพฺพา “สงฺขารกฺขนฺธปริยาปนฺนา”ติ, น เหวํ วตฺตพฺเพ ฯเปฯ.\csromanspacer{} สงฺขารกฺขนฺธสฺส สุญฺญตา น วตฺตพฺพา “เวทนากฺขนฺธปริยาปนฺนา ฯเปฯ.\csromanspacer{} สญฺญากฺขนฺธปริยาปนฺนา.\csromanspacer{} วิญฺญาณกฺขนฺธปริยาปนฺนา”ติ, อามนฺตา.\csromanspacer{} วิญฺญาณกฺขนฺธสฺส สุญฺญตา น วตฺตพฺพา “สงฺขารกฺขนฺธปริยาปนฺนา”ติ, น เหวํ วตฺตพฺเพ ฯเปฯ.}

\proseitem{834}{\csromansymbolmark{+}น วตฺตพฺพํ “สุญฺญตา สงฺขารกฺขนฺธปริยาปนฺนา”ติ, อามนฺตา.\csromanspacer{} นนุ วุตฺตํ ภควตา “สุญฺญมิทํ ภิกฺขเว สงฺขารา อตฺเตน วา อตฺตนิเยน วา”ติ\footnote{ม 3. 50 ปิฏฺเฐ อาเนญฺชสปฺปาเย.} อตฺเถว สุตฺตนฺโตติ, อามนฺตา.\csromanspacer{} เตน หิ สุญฺญตา สงฺขารกฺขนฺธปริยาปนฺนาติ.}

\nitthitam{สุญฺญตากถา นิฏฺฐิตา.}
\chapterhead{19. เอกูนวีสติมวคฺค}
\vspace{\tipitakaparskip}
\csromancenter{\textbf{(188) 3.}\csromanspacer{}\textbf{ สามญฺญผลกถา}}
\proseitem{835}{\csromansymbolmark{*}สามญฺญผลํ อสงฺขตนฺติ, อามนฺตา.\csromanspacer{} นิพฺพานํ ตาณํ เลณํ สรณํ ปรายณํ อจฺจุตํ อมตนฺติ, น เหวํ วตฺตพฺเพ ฯเปฯ.\csromanspacer{} สามญฺญผลํ อสงฺขตํ, นิพฺพานํ อสงฺขตนฺติ, อามนฺตา.\csromanspacer{} เทฺว อสงฺขตานีติ, น เหวํ วตฺตพฺเพ ฯเปฯ.\csromanspacer{} เทฺว อสงฺขตานีติ, อามนฺตา.\csromanspacer{} เทฺว ตาณานิ ฯเปฯ อนฺตริกา วาติ, น เหวํ วตฺตพฺเพ ฯเปฯ.}

\proseitem{836}{\csromansymbolmark{*}สามญฺญผลํ อสงฺขตนฺติ, อามนฺตา.\csromanspacer{} สามญฺญํ อสงฺขตนฺติ, น เหวํ วตฺตพฺเพ ฯเปฯ.\csromanspacer{} สามญฺญํ สงฺขตนฺติ, อามนฺตา.\csromanspacer{} สามญฺญผลํ สงฺขตนฺติ, น เหวํ วตฺตพฺเพ ฯเปฯ.}

\csromanmatikaanchor{mtk.224}
\csromantocmarkref{subsection}{(189) 4. ปตฺติกถา}{mtk.224}
\bookmark[dest={mtk.224},level=2]{(189) 4. ปตฺติกถา}
\prose{\csromanfolio{420}\csromansymbolmark{*}โสตาปตฺติผลํ อสงฺขตนฺติ, อามนฺตา.\csromanspacer{} โสตาปตฺติมคฺโค อสงฺขโตติ, น เหวํ วตฺตพฺเพ ฯเปฯ.\csromanspacer{} โสตาปตฺติมคฺโค สงฺขโตติ, อามนฺตา.\csromanspacer{} โสตาปตฺติผลํ สงฺขตนฺติ, น เหวํ วตฺตพฺเพ ฯเปฯ.}

\prose{\csromansymbolmark{*}สกทาคามิผลํ ฯเปฯ.\csromanspacer{} อนาคามิผลํ ฯเปฯ.\csromanspacer{} อรหตฺตผลํ อสงฺขตนฺติ, อามนฺตา.\csromanspacer{} อรหตฺตมคฺโค อสงฺขโตติ, น เหวํ วตฺตพฺเพ ฯเปฯ.\csromanspacer{} อรหตฺตมคฺโค สงฺขโตติ, อามนฺตา.\csromanspacer{} อรหตฺตผลํ สงฺขตนฺติ, น เหวํ วตฺตพฺเพ ฯเปฯ.}

\prose{\csromansymbolmark{*}โสตาปตฺติผลํ อสงฺขตํ สกทาคามิผลํ ฯเปฯ อนาคามิผลํ ฯเปฯ อรหตฺตผลํ อสงฺขตํ, นิพฺพานํ อสงฺขตนฺติ, อามนฺตา.\csromanspacer{} ปญฺจ อสงฺขตานีติ, น เหวํ วตฺตพฺเพ ฯเปฯ.\csromanspacer{} ปญฺจ อสงฺขตานีติ, อามนฺตา.\csromanspacer{} ปญฺจ ตาณานิ ฯเปฯ อนฺตริกา วาติ, น เหวํ วตฺตพฺเพ ฯเปฯ.}

\nitthitam{สามญฺญผลกถา นิฏฺฐิตา.}
\chapterhead{19. เอกูนวีสติมวคฺค}
\vspace{\tipitakaparskip}
\csromancenter{\textbf{(189) 4.}\csromanspacer{}\textbf{ ปตฺติกถา}}
\proseitem{837}{\csromansymbolmark{*}ปตฺติ อสงฺขตาติ, อามนฺตา.\csromanspacer{} นิพฺพานํ ตาณํ เลณํ สรณํ ปรายณํ อจฺจุตํ อมตนฺติ, น เหวํ วตฺตพฺเพ ฯเปฯ.\csromanspacer{} ปตฺติ อสงฺขตา, นิพฺพานํ อสงฺขตนฺติ, อามนฺตา.\csromanspacer{} เทฺว อสงฺขตานีติ, น เหวํ วตฺตพฺเพ ฯเปฯ.\csromanspacer{} เทฺว อสงฺขตานีติ, อามนฺตา.\csromanspacer{} เทฺว ตาณานิ ฯเปฯ อนฺตริกา วาติ, น เหวํ วตฺตพฺเพ ฯเปฯ.}

\proseitem{838}{\csromansymbolmark{*}จีวรสฺส ปตฺติ อสงฺขตาติ, อามนฺตา.\csromanspacer{} นิพฺพานํ ฯเปฯ อมตนฺติ, น เหวํ วตฺตพฺเพ ฯเปฯ.\csromanspacer{} จีวรสฺส ปตฺติ อสงฺขตา, นิพฺพานํ อสงฺขตนฺติ, อามนฺตา.\csromanspacer{} เทฺว อสงฺขตานีติ, น เหวํ วตฺตพฺเพ ฯเปฯ.\csromanspacer{} เทฺว อสงฺขตานีติ, อามนฺตา.\csromanspacer{} เทฺว ตาณานิ ฯเปฯ อนฺตริกา วาติ, น เหวํ วตฺตพฺเพ ฯเปฯ.\csromanspacer{} ปิณฺฑปาตสฺส ฯเปฯ.\csromanspacer{} เสนาสนสฺส ฯเปฯ.\csromanspacer{} คิลานปจฺจยเภสชฺชปริกฺขารสฺส ปตฺติ อสงฺขตาติ, อามนฺตา.\csromanspacer{} นิพฺพานํ ฯเปฯ อจฺจุตํ อมตนฺติ, น เหวํ วตฺตพฺเพ ฯเปฯ.\csromanspacer{} คิลานปจฺจยเภสชฺชปริกฺขารสฺส ปตฺติ อสงฺขตา, นิพฺพานํ อสงฺขตนฺติ, อามนฺตา.\csromanspacer{} เทฺว อสงฺขตานีติ, น เหวํ วตฺตพฺเพ ฯเปฯ.\csromanspacer{} เทฺว อสงฺขตานีติ, อามนฺตา.\csromanspacer{} เทฺว ตาณานิ ฯเปฯ อนฺตริกา วาติ, น เหวํ วตฺตพฺเพ ฯเปฯ.}

\chapterheadpageread{19. เอกูนวีสติมวคฺค}
\csromanmatikaanchor{mtk.225}
\csromantocmarkref{subsection}{(190) 5. ตถตากถา}{mtk.225}
\bookmark[dest={mtk.225},level=2]{(190) 5. ตถตากถา}
\prose{\csromanfolio{421}\csromansymbolmark{*}จีวรสฺส ปตฺติ อสงฺขตา ฯเปฯ.\csromanspacer{} ปิณฺฑปาตสฺส ฯเปฯ.\csromanspacer{} เสนาสนสฺส ฯเปฯ.\csromanspacer{} คิลานปจฺจยเภสชฺชปริกฺขารสฺส ปตฺติ อสงฺขตา, นิพฺพานํ อสงฺขตนฺติ, อามนฺตา.\csromanspacer{} ปญฺจ อสงฺขตานีติ, น เหวํ วตฺตพฺเพ ฯเปฯ.\csromanspacer{} ปญฺจ อสงฺขตานีติ, อามนฺตา.\csromanspacer{} ปญฺจ ตาณานิ ฯเปฯ อนฺตริกา วาติ, น เหวํ วตฺตพฺเพ ฯเปฯ.}

\proseitem{839}{\csromansymbolmark{*}ปฐมสฺส ฌานสฺส ปตฺติ อสงฺขตา. (เอวํ สพฺพํ วิตฺถาเรตพฺพํ.) ทุติยสฺส ฌานสฺส.\csromanspacer{} ตติยสฺส ฌานสฺส.\csromanspacer{} จตุตฺถสฺส ฌานสฺส.\csromanspacer{} อากาสานญฺจายตนสฺส.\csromanspacer{} วิญฺญาณญฺจายตนสฺส.\csromanspacer{} อากิญฺจญฺญายตนสฺส.\csromanspacer{} เนวสญฺญานาสญฺญายตนสฺส.\csromanspacer{} โสตาปตฺติมคฺคสฺส.\csromanspacer{} โสตาปตฺติผลสฺส.\csromanspacer{} สกทาคามิมคฺคสฺส.\csromanspacer{} สกทาคามิผลสฺส.\csromanspacer{} อนาคามิมคฺคสฺส.\csromanspacer{} อนาคามิผลสฺส.\csromanspacer{} อรหตฺตมคฺคสฺส.\csromanspacer{} อรหตฺตผลสฺส ปตฺติ อสงฺขตาติ, อามนฺตา.\csromanspacer{} นิพฺพานํ ฯเปฯ อจฺจุตํ อมตนฺติ, น เหวํ วตฺตพฺเพ ฯเปฯ.\csromanspacer{} อรหตฺตผลสฺส ปตฺติ อสงฺขตา, นิพฺพานํ อสงฺขตนฺติ, อามนฺตา.\csromanspacer{} เทฺว อสงฺขตานีติ, น เหวํ วตฺตพฺเพ ฯเปฯ.\csromanspacer{} เทฺว อสงฺขตานีติ, อามนฺตา.\csromanspacer{} เทฺว ตาณานิ ฯเปฯ อนฺตริกา วาติ, น เหวํ วตฺตพฺเพ ฯเปฯ.}

\prose{\csromansymbolmark{*}โสตาปตฺติมคฺคสฺส ปตฺติ อสงฺขตา.\csromanspacer{} โสตาปตฺติผลสฺส ปตฺติ อสงฺขตา ฯเปฯ.\csromanspacer{} อรหตฺตมคฺคสฺส ปตฺติ อสงฺขตา.\csromanspacer{} อรหตฺตผลสฺส ปตฺติ อสงฺขตา นิพฺพานํ อสงฺขตนฺติ, อามนฺตา.\csromanspacer{} นว อสงฺขตานีติ, น เหวํ วตฺตพฺเพ ฯเปฯ.\csromanspacer{} นว อสงฺขตานีติ, อามนฺตา.\csromanspacer{} นว ตาณานิ ฯเปฯ อนฺตริกา วาติ, น เหวํ วตฺตพฺเพ ฯเปฯ.}

\proseitem{840}{\csromansymbolmark{+}น วตฺตพฺพํ “ปตฺติ อสงฺขตา”ติ, อามนฺตา.\csromanspacer{} ปตฺติ รูปํ.\csromanspacer{} เวทนา.\csromanspacer{} สญฺญา.\csromanspacer{} สงฺขารา.\csromanspacer{} วิญฺญาณนฺติ, น เหวํ วตฺตพฺเพ.\csromanspacer{} เตน หิ ปตฺติ อสงฺขตาติ.}

\nitthitam{ปตฺติกถา นิฏฺฐิตา.}
\chapterhead{19. เอกูนวีสติมวคฺค}
\vspace{\tipitakaparskip}
\csromancenter{\textbf{(190) 5.}\csromanspacer{}\textbf{ ตถตากถา}}
\csromanmatikaanchor{mtk.226}
\csromantocmarkref{subsection}{(191) 6. กุสลกถา}{mtk.226}
\bookmark[dest={mtk.226},level=2]{(191) 6. กุสลกถา}
\proseitem{841}{\csromansymbolmark{*}สพฺพธมฺมานํ ตถตา อสงฺขตาติ, อามนฺตา.\csromanspacer{} นิพฺพานํ ฯเปฯ อมตนฺติ, น เหวํ วตฺตพฺเพ ฯเปฯ.\csromanspacer{} สพฺพธมฺมานํ ตถตา อสงฺขตา, นิพฺพานํ \csromanfolio{422}อสงฺขตนฺติ, อามนฺตา.\csromanspacer{} เทฺว อสงฺขตานีติ, น เหวํ วตฺตพฺเพ ฯเปฯ.\csromanspacer{} เทฺว อสงฺขตานีติ, อามนฺตา.\csromanspacer{} เทฺว ตาณานิ ฯเปฯ อนฺตริกา วาติ, น เหวํ วตฺตพฺเพ ฯเปฯ.}

\proseitem{842}{\csromansymbolmark{*}รูปสฺส รูปตา, นนุ รูปตา อสงฺขตาติ, อามนฺตา.\csromanspacer{} นิพฺพานํ ฯเปฯ อมตนฺติ, น เหวํ วตฺตพฺเพ ฯเปฯ.\csromanspacer{} รูปสฺส รูปตา, นนุ รูปตา อสงฺขตา นิพฺพานํ อสงฺขตนฺติ, อามนฺตา.\csromanspacer{} เทฺว อสงฺขตานีติ, น เหวํ วตฺตพฺเพ ฯเปฯ.\csromanspacer{} เทฺว อสงฺขตานีติ, อามนฺตา.\csromanspacer{} เทฺว ตาณานิ ฯเปฯ อนฺตริกา วาติ, น เหวํ วตฺตพฺเพ ฯเปฯ.}

\prose{\csromansymbolmark{*}เวทนาย เวทนตา, นนุ เวทนตา ฯเปฯ.\csromanspacer{} สญฺญาย สญฺญตา, นนุ สญฺญตา ฯเปฯ.\csromanspacer{} สงฺขารานํ สงฺขารตา, นนุ สงฺขารตา ฯเปฯ.\csromanspacer{} วิญฺญาณสฺส วิญฺญาณตา, นนุ วิญฺญาณตา อสงฺขตาติ, อามนฺตา.\csromanspacer{} นิพฺพานํ ฯเปฯ อมตนฺติ, น เหวํ วตฺตพฺเพ ฯเปฯ.}

\prose{\csromansymbolmark{*}รูปสฺส รูปตา, นนุ รูปตา ฯเปฯ.\csromanspacer{} วิญฺญาณสฺส วิญฺญาณตา, นนุ วิญฺญาณตา อสงฺขตา, นิพฺพานํ อสงฺขตนฺติ, อามนฺตา.\csromanspacer{} ฉ อสงฺขตานีติ, น เหวํ วตฺตพฺเพ ฯเปฯ.\csromanspacer{} ฉ อสงฺขตานีติ, อามนฺตา.\csromanspacer{} ฉ ตาณานิ ฯเปฯ อนฺตริกา วาติ, น เหวํ วตฺตพฺเพ ฯเปฯ.}

\proseitem{843}{\csromansymbolmark{+}น วตฺตพฺพํ “สพฺพธมฺมานํ ตถตา อสงฺขตา”ติ, อามนฺตา.\csromanspacer{} สพฺพธมฺมานํ ตถตา รูปํ.\csromanspacer{} เวทนา.\csromanspacer{} สญฺญา.\csromanspacer{} สงฺขารา.\csromanspacer{} วิญฺญาณนฺติ, น เหวํ วตฺตพฺเพ.\csromanspacer{} เตน หิ สพฺพธมฺมานํ ตถตา อสงฺขตาติ.}

\nitthitam{ตถตากถา นิฏฺฐิตา.}
\chapterhead{19. เอกูนวีสติมวคฺค}
\vspace{\tipitakaparskip}
\csromancenter{\textbf{(191) 6.}\csromanspacer{}\textbf{ กุสลกถา}}
\proseitem{844}{\csromansymbolmark{*}นิพฺพานธาตุ กุสลาติ, อามนฺตา.\csromanspacer{} สารมฺมณา, อตฺถิ ตาย อาวฏฺฏนา ฯเปฯ ปณิธีติ, น เหวํ วตฺตพฺเพ ฯเปฯ.\csromanspacer{} นนุ อนารมฺมณา, นตฺถิ ตาย อาวฏฺฏนา ฯเปฯ ปณิธีติ, อามนฺตา.\csromanspacer{} หญฺจิ อนารมฺมณา, นตฺถิ}

\vspace{\tipitakaparskip}
\csromancenter{เอกูนวีสติมวคฺค ตาย อาวฏฺฏนา ฯเปฯ ปณิธิ, โน จ วต เร วตฺตพฺเพ “นิพฺพานธาตุ กุสลา”ติ.}
\csromanmatikaanchor{mtk.227}
\csromantocmarkref{subsection}{(192) 7. อจฺจนฺตนิยามกถา}{mtk.227}
\bookmark[dest={mtk.227},level=2]{(192) 7. อจฺจนฺตนิยามกถา}
\proseitem{845}{\csromanfolio{423}\csromansymbolmark{*}อโลโภ กุสโล สารมฺมโณ, อตฺถิ ตสฺส อาวฏฺฏนา ฯเปฯ ปณิธีติ, อามนฺตา.\csromanspacer{} นิพฺพานธาตุ กุสลา สารมฺมณา, อตฺถิ ตาย อาวฏฺฏนา ฯเปฯ ปณิธีติ, น เหวํ วตฺตพฺเพ ฯเปฯ.\csromanspacer{} อโทโส กุสโล ฯเปฯ.\csromanspacer{} อโมโห กุสโล ฯเปฯ.\csromanspacer{} สทฺธา.\csromanspacer{} วีริยํ.\csromanspacer{} สติ.\csromanspacer{} สมาธิ ฯเปฯ.\csromanspacer{} ปญฺญา กุสลา สารมฺมณา, อตฺถิ ตาย อาวฏฺฏนา ฯเปฯ ปณิธีติ, อามนฺตา.\csromanspacer{} นิพฺพานธาตุ กุสลา สารมฺมณา, อตฺถิ ตาย อาวฏฺฏนา ฯเปฯ ปณิธีติ, น เหวํ วตฺตพฺเพ ฯเปฯ.}

\prose{\csromansymbolmark{*}นิพฺพานธาตุ กุสลา อนารมฺมณา, นตฺถิ ตาย อาวฏฺฏนา ฯเปฯ ปณิธีติ, อามนฺตา.\csromanspacer{} อโลโภ กุสโล อนารมฺมโณ, นตฺถิ ตสฺส อาวฏฺฏนา ฯเปฯ ปณิธีติ, น เหวํ วตฺตพฺเพ ฯเปฯ.\csromanspacer{} นิพฺพานธาตุ กุสลา อนารมฺมณา, นตฺถิ ตาย อาวฏฺฏนา ฯเปฯ ปณิธีติ, อามนฺตา.\csromanspacer{} อโทโส กุสโล ฯเปฯ.\csromanspacer{} อโมโห กุสโล ฯเปฯ.\csromanspacer{} สทฺธา.\csromanspacer{} วีริยํ.\csromanspacer{} สติ.\csromanspacer{} สมาธิ ฯเปฯ.\csromanspacer{} ปญฺญา กุสลา อนารมฺมณา, นตฺถิ ตาย อาวฏฺฏนา ฯเปฯ ปณิธีติ, น เหวํ วตฺตพฺเพ ฯเปฯ.}

\proseitem{846}{\csromansymbolmark{+}น วตฺตพฺพํ “นิพฺพานธาตุ กุสลา”ติ, อามนฺตา.\csromanspacer{} นนุ นิพฺพานธาตุ อนวชฺชาติ, อามนฺตา.\csromanspacer{} หญฺจิ นิพฺพานธาตุ อนวชฺชา, เตน วต เร วตฺตพฺเพ “นิพฺพานธาตุ กุสลา”ติ.}

\nitthitam{กุสลกถา นิฏฺฐิตา.}
\chapterhead{19. เอกูนวีสติมวคฺค}
\vspace{\tipitakaparskip}
\csromancenter{\textbf{(192) 7.}\csromanspacer{}\textbf{ อจฺจนฺตนิยามกถา}}
\proseitem{847}{\csromansymbolmark{*}อตฺถิ ปุถุชฺชนสฺส อจฺจนฺตนิยามตาติ, อามนฺตา.\csromanspacer{} มาตุฆาตโก อจฺจนฺตนิยโต.\csromanspacer{} ปิตุฆาตโก ฯเปฯ.\csromanspacer{} อรหนฺตฆาตโก ฯเปฯ.\csromanspacer{} รุหิรุปฺปาทโก ฯเปฯ.\csromanspacer{} สํฆเภทโก อจฺจนฺตนิยโตติ, น เหวํ วตฺตพฺเพ ฯเปฯ.}

\proseitem{848}{\csromanfolio{424}\csromansymbolmark{*}อตฺถิ ปุถุชฺชนสฺส อจฺจนฺตนิยามตาติ, อามนฺตา.\csromanspacer{} อจฺจนฺตนิยตสฺส ปุคฺคลสฺส วิจิกิจฺฉา อุปฺปชฺเชยฺยาติ, อามนฺตา.\csromanspacer{} หญฺจิ อจฺจนฺตนิยตสฺส ปุคฺคลสฺส วิจิกิจฺฉา อุปฺปชฺเชยฺย, โน จ วต เร วตฺตพฺเพ “อตฺถิ ปุถุชฺชนสฺส อจฺจนฺตนิยามตา”ติ.}

\prose{\csromansymbolmark{*}อจฺจนฺตนิยตสฺส ปุคฺคลสฺส วิจิกิจฺฉา นุปฺปชฺเชยฺยาติ, อามนฺตา.\csromanspacer{} ปหีนาติ, น เหวํ วตฺตพฺเพ ฯเปฯ.\csromanspacer{} ปหีนาติ, อามนฺตา.\csromanspacer{} โสตาปตฺติมคฺเคนาติ, น เหวํ วตฺตพฺเพ ฯเปฯ.\csromanspacer{} สกทาคามิมคฺเคน ฯเปฯ.\csromanspacer{} อนาคามิมคฺเคน ฯเปฯ.\csromanspacer{} อรหตฺตมคฺเคนาติ.\csromanspacer{} น เหวํ วตฺตพฺเพ ฯเปฯ.}

\prose{\csromansymbolmark{*}กตเมน มคฺเคนาติ, อกุสเลน มคฺเคนาติ.\csromanspacer{} อกุสโล มคฺโค นิยฺยานิโก ขยคามี โพธคามี อนาสโว ฯเปฯ อสํกิเลสิโกติ, น เหวํ วตฺตพฺเพ ฯเปฯ.\csromanspacer{} นนุ อกุสโล มคฺโค อนิยฺยานิโก ฯเปฯ สํกิเลสิโกติ, อามนฺตา.\csromanspacer{} หญฺจิ อกุสโล มคฺโค อนิยฺยานิโก ฯเปฯ สํกิเลสิโก, โน จ วต เร วตฺตพฺเพ “อจฺจนฺตนิยตสฺส ปุคฺคลสฺส วิจิกิจฺฉา อกุสเลน มคฺเคน ปหีนา”ติ.}

\proseitem{849}{\csromansymbolmark{*}สสฺสตทิฏฺฐิยา นิยตสฺส ปุคฺคลสฺส อุจฺเฉททิฏฺฐิ อุปฺปชฺเชยฺยาติ, อามนฺตา.\csromanspacer{} หญฺจิ สสฺสตทิฏฺฐิยา นิยตสฺส ปุคฺคลสฺส อุจฺเฉททิฏฺฐิ อุปฺปชฺเชยฺย, โน จ วต เร วตฺตพฺเพ “อตฺถิ ปุถุชฺชนสฺส อจฺจนฺตนิยามตา”ติ.}

\prose{\csromansymbolmark{*}สสฺสตทิฏฺฐิยา นิยตสฺส ปุคฺคลสฺส อุจฺเฉททิฏฺฐิ นุปฺปชฺเชยฺยาติ, อามนฺตา.\csromanspacer{} ปหีนาติ, น เหวํ วตฺตพฺเพ ฯเปฯ.\csromanspacer{} ปหีนาติ, อามนฺตา.\csromanspacer{} โสตาปตฺติมคฺเคนาติ, น เหวํ วตฺตพฺเพ ฯเปฯ.\csromanspacer{} สกทาคามิมคฺเคน ฯเปฯ.\csromanspacer{} อนาคามิมคฺเคน ฯเปฯ.\csromanspacer{} อรหตฺตมคฺเคนาติ, น เหวํ วตฺตพฺเพ ฯเปฯ.}

\prose{\csromansymbolmark{*}กตเมน มคฺเคนาติ, อกุสเลน มคฺเคน.\csromanspacer{} อกุสโล มคฺโค ฯเปฯ.\csromanspacer{} โน จ วต เร วตฺตพฺเพ “สสฺสตทิฏฺฐิยา นิยตสฺส ปุคฺคลสฺส อุจฺเฉททิฏฺฐิ อกุสเลน มคฺเคน ปหีนา”ติ.}

\proseitem{850}{\csromansymbolmark{*}อุจฺเฉททิฏฺฐิยา นิยตสฺส ปุคฺคลสฺส สสฺสตทิฏฺฐิ อุปฺปชฺเชยฺยาติ, อามนฺตา.\csromanspacer{} หญฺจิ อุจฺเฉททิฏฺฐิยา นิยตสฺส ปุคฺคลสฺส สสฺสตทิฏฺฐิ อุปฺปชฺเชยฺย, โน จ วต เร วตฺตพฺเพ “อตฺถิ ปุถุชฺชนสฺส อจฺจนฺตนิยามตา”ติ.}

\prose{\csromansymbolmark{*}อุจฺเฉททิฏฺฐิยา นิยตสฺส ปุคฺคลสฺส สสฺสตทิฏฺฐิ นุปฺปชฺเชยฺยาติ, อามนฺตา.\csromanspacer{} ปหีนาติ, น เหวํ วตฺตพฺเพ ฯเปฯ.\csromanspacer{} ปหีนาติ, อามนฺตา.}

\chapterheadpageread{19. เอกูนวีสติมวคฺค}
\prose{\csromanfolio{425}โสตาปตฺติมคฺเคนาติ, น เหวํ วตฺตพฺเพ ฯเปฯ.\csromanspacer{} สกทาคามิมคฺเคน ฯเปฯ.\csromanspacer{} อนาคามิมคฺเคน.\csromanspacer{} อรหตฺตมคฺเคนาติ, น เหวํ วตฺตพฺเพ ฯเปฯ.}

\prose{\csromansymbolmark{*}กตเมน มคฺเคนาติ, อกุสเลน มคฺเคนาติ.\csromanspacer{} อกุสโล มคฺโค ฯเปฯ.\csromanspacer{} โน จ วต เร วตฺตพฺเพ “อุจฺเฉททิฏฺฐิยา นิยตสฺส ปุคฺคลสฺส สสฺสตทิฏฺฐิ อกุสเลน มคฺเคน ปหีนา”ติ.}

\proseitem{851}{\csromansymbolmark{+}น วตฺตพฺพํ “อตฺถิ ปุถุชฺชนสฺส อจฺจนฺตนิยามตา”ติ, อามนฺตา.\csromanspacer{} นนุ วุตฺตํ ภควตา “อิธ ภิกฺขเว เอกจฺโจ ปุคฺคโล สมนฺนาคโต โหติ เอกนฺตกาฬเกหิ อกุสเลหิ ธมฺเมหิ, โส\footnote{เอวํ โข ภิกฺขเว ปุคฺคโล (อํ 2. 403 ปิฏฺฐิ.)} สกิํ นิมุคฺโค นิมุคฺโคว โหตี”ติ\footnote{อํ 2. 403 ปิฏฺเฐ.} อตฺเถว สุตฺตนฺโตติ, อามนฺตา.\csromanspacer{} เตน หิ อตฺถิ ปุถุชฺชนสฺส อจฺจนฺตนิยามตาติ.}

\proseitem{852}{\csromansymbolmark{*}วุตฺตํ ภควตา “อิธ ภิกฺขเว เอกจฺโจ ปุคฺคโล สมนฺนาคโต โหติ เอกนฺตกาฬเกหิ อกุสเลหิ ธมฺเมหิ, โส สกิํ นิมุคฺโค นิมุคฺโคว โหตี”ติ กตฺวา เตน จ การเณน อตฺถิ ปุถุชฺชนสฺส อจฺจนฺตนิยามตาติ, อามนฺตา.\csromanspacer{} วุตฺตํ ภควตา “อิธ ภิกฺขเว เอกจฺโจ ปุคฺคโล อุมฺมุชฺชิตฺวา นิมุชฺชตี”ติ3 อตฺเถว สุตฺตนฺโตติ, อามนฺตา.\csromanspacer{} สพฺพกาลํ อุมฺมุชฺชิตฺวา นิมุชฺชตีติ, น เหวํ วตฺตพฺเพ ฯเปฯ.}

\prose{\csromansymbolmark{*}วุตฺตํ ภควตา “อิธ ภิกฺขเว เอกจฺโจ ปุคฺคโล สมนฺนาคโต โหติ เอกนฺตกาฬเกหิ อกุสเลหิ ธมฺเมหิ, โส สกิํ นิมุคฺโค นิมุคฺโคว โหตี”ติ กตฺวา เตน จ การเณน อตฺถิ ปุถุชฺชนสฺส อจฺจนฺตนิยามตาติ, อามนฺตา.\csromanspacer{} วุตฺตํ ภควตา “อิธ ภิกฺขเว เอกจฺโจ ปุคฺคโล อุมฺมุชฺชิตฺวา ฐิโต โหติ, อุมฺมุชฺชิตฺวา วิปสฺสติ วิโลเกติ, อุมฺมุชฺชิตฺวา ปตรติ, อุมฺมุชฺชิตฺวา ปติคาธปฺปตฺโต โหตี”ติ4 อตฺเถว สุตฺตนฺโตติ, อามนฺตา.\csromanspacer{} สพฺพกาลํ อุมฺมุชฺชิตฺวา ปติคาธปฺปตฺโต โหตีติ, น เหวํ วตฺตพฺเพ ฯเปฯ.}

\nitthitam{อจฺจนฺตนิยามกถา นิฏฺฐิตา.}
\chapterheadpageread{19. เอกูนวีสติมวคฺค}
\vspace{\tipitakaparskip}
\csromancenter{\textbf{(193) 8.}\csromanspacer{}\textbf{ อินฺทฺริยกถา}}
\csromanmatikaanchor{mtk.228}
\csromantocmarkref{subsection}{(193) 8. อินฺทฺริยกถา}{mtk.228}
\bookmark[dest={mtk.228},level=2]{(193) 8. อินฺทฺริยกถา}
\proseitem{853}{\csromanfolio{426}\csromansymbolmark{*}นตฺถิ โลกิยํ สทฺธินฺทฺริยนฺติ, อามนฺตา.\csromanspacer{} นตฺถิ โลกิยา สทฺธาติ, น เหวํ วตฺตพฺเพ ฯเปฯ.\csromanspacer{} นตฺถิ โลกิยํ วีริยินฺทฺริยํ ฯเปฯ สตินฺทฺริยํ ฯเปฯ สมาธินฺทฺริยํ ฯเปฯ ปญฺญินฺทฺริยนฺติ, อามนฺตา.\csromanspacer{} นตฺถิ โลกิยา ปญฺญาติ, น เหวํ วตฺตพฺเพ ฯเปฯ.}

\prose{\csromansymbolmark{*}อตฺถิ โลกิยา สทฺธาติ, อามนฺตา.\csromanspacer{} อตฺถิ โลกิยํ สทฺธินฺทฺริยนฺติ, น เหวํ วตฺตพฺเพ ฯเปฯ.\csromanspacer{} อตฺถิ โลกิยํ วีริยํ ฯเปฯ สติ ฯเปฯ สมาธิ ฯเปฯ ปญฺญาติ, อามนฺตา.\csromanspacer{} อตฺถิ โลกิยํ ปญฺญินฺทฺริยนฺติ, น เหวํ วตฺตพฺเพ ฯเปฯ.}

\prose{\csromansymbolmark{*}อตฺถิ โลกิโย มโน, อตฺถิ โลกิยํ มนินฺทฺริยนฺติ, อามนฺตา.\csromanspacer{} อตฺถิ โลกิยา สทฺธา, อตฺถิ โลกิยํ สทฺธินฺทฺริยนฺติ, น เหวํ วตฺตพฺเพ ฯเปฯ.\csromanspacer{} อตฺถิ โลกิโย มโน, อตฺถิ โลกิยํ มนินฺทฺริยนฺติ, อามนฺตา.\csromanspacer{} อตฺถิ โลกิยา ปญฺญา, อตฺถิ โลกิยํ ปญฺญินฺทฺริยนฺติ, น เหวํ วตฺตพฺเพ ฯเปฯ.}

\prose{\csromansymbolmark{*}อตฺถิ โลกิยํ โสมนสฺสํ, อตฺถิ โลกิยํ โสมนสฺสินฺทฺริยํ ฯเปฯ.\csromanspacer{} อตฺถิ โลกิยํ ชีวิตํ, อตฺถิ โลกิยํ ชีวิตินฺทฺริยนฺติ, อามนฺตา.\csromanspacer{} อตฺถิ โลกิยา สทฺธา, อตฺถิ โลกิยํ สทฺธินฺทฺริยนฺติ, น เหวํ วตฺตพฺเพ ฯเปฯ.\csromanspacer{} อตฺถิ โลกิยํ ชีวิตํ, อตฺถิ โลกิยํ ชีวิตินฺทฺริยนฺติ, อามนฺตา.\csromanspacer{} อตฺถิ โลกิยา ปญฺญา, อตฺถิ โลกิยํ ปญฺญินฺทฺริยนฺติ, น เหวํ วตฺตพฺเพ ฯเปฯ.}

\proseitem{854}{\csromansymbolmark{*}อตฺถิ โลกิยา สทฺธา, นตฺถิ โลกิยํ สทฺธินฺทฺริยนฺติ, อามนฺตา.\csromanspacer{} อตฺถิ โลกิโย มโน, นตฺถิ โลกิยํ มนินฺทฺริยนฺติ, น เหวํ วตฺตพฺเพ ฯเปฯ.\csromanspacer{} อตฺถิ โลกิยา ปญฺญา, นตฺถิ โลกิยํ ปญฺญินฺทฺริยนฺติ, อามนฺตา.\csromanspacer{} อตฺถิ โลกิโย มโน, นตฺถิ โลกิยํ มนินฺทฺริยนฺติ, น เหวํ วตฺตพฺเพ ฯเปฯ.}

\prose{\csromansymbolmark{*}อตฺถิ โลกิยา สทฺธา, นตฺถิ โลกิยํ สทฺธินฺทฺริยนฺติ, อามนฺตา.\csromanspacer{} อตฺถิ โลกิยํ โสมนสฺสํ, นตฺถิ โลกิยํ โสมนสฺสินฺทฺริยนฺติ, อตฺถิ โลกิยํ ชีวิตํ นตฺถิ โลกิยํ ชีวิตินฺทฺริยนฺติ, น เหวํ วตฺตพฺเพ ฯเปฯ.\csromanspacer{} อตฺถิ โลกิยา ปญฺญา, นตฺถิ โลกิยํ ปญฺญินฺทฺริยนฺติ, อามนฺตา.\csromanspacer{} อตฺถิ โลกิโย มโน, นตฺถิ โลกิยํ มนินฺทฺริยนฺติ ฯเปฯ.\csromanspacer{} อตฺถิ โลกิยํ ชีวิตํ, นตฺถิ โลกิยํ ชีวิตินฺทฺริยนฺติ, น เหวํ วตฺตพฺเพ ฯเปฯ.}

\chapterheadpageread{20. วีสติมวคฺค}
\csromanmatikaanchor{mtk.229}
\csromanmatikaanchor{mtk.230}
\csromantocmarknopage{chapter}{20. วีสติมวคฺค}{mtk.229}
\bookmark[dest={mtk.229},level=0]{20. วีสติมวคฺค}
\csromantocmarkref{subsection}{(194) 1. อสญฺจิจฺจกถา}{mtk.230}
\bookmark[dest={mtk.230},level=2]{(194) 1. อสญฺจิจฺจกถา}
\proseitem{855}{\csromanfolio{427}\csromansymbolmark{*}อตฺถิ โลกุตฺตรา สทฺธา, อตฺถิ โลกุตฺตรํ สทฺธินฺทฺริยนฺติ, อามนฺตา.\csromanspacer{} อตฺถิ โลกิยา สทฺธา, อตฺถิ โลกิยํ สทฺธินฺทฺริยนฺติ, น เหวํ วตฺตพฺเพ ฯเปฯ.\csromanspacer{} อตฺถิ โลกุตฺตรํ วีริยํ ฯเปฯ.\csromanspacer{} อตฺถิ โลกุตฺตรา ปญฺญา, อตฺถิ โลกุตฺตรํ ปญฺญินฺทฺริยนฺติ, อามนฺตา.\csromanspacer{} อตฺถิ โลกิยา ปญฺญา, อตฺถิ โลกิยํ ปญฺญินฺทฺริยนฺติ, น เหวํ วตฺตพฺเพ ฯเปฯ.}

\prose{\csromansymbolmark{*}อตฺถิ โลกิยา สทฺธา, นตฺถิ โลกิยํ สทฺธินฺทฺริยนฺติ, อามนฺตา.\csromanspacer{} อตฺถิ โลกุตฺตรา สทฺธา, นตฺถิ โลกุตฺตรํ สทฺธินฺทฺริยนฺติ, น เหวํ วตฺตพฺเพ ฯเปฯ.\csromanspacer{} อตฺถิ โลกิยํ วีริยํ ฯเปฯ.\csromanspacer{} อตฺถิ โลกิยา ปญฺญา, นตฺถิ โลกิยํ ปญฺญินฺทฺริยนฺติ, อามนฺตา.\csromanspacer{} อตฺถิ โลกุตฺตรา ปญฺญา, นตฺถิ โลกุตฺตรํ ปญฺญินฺทฺริยนฺติ, น เหวํ วตฺตพฺเพ ฯเปฯ.}

\proseitem{856}{\csromansymbolmark{*}นตฺถิ โลกิยานิ ปญฺจินฺทฺริยานีติ, อามนฺตา.\csromanspacer{} นนุ วุตฺตํ ภควตา “อทฺทสํ โข อหํ ภิกฺขเว พุทฺธจกฺขุนา โลกํ โวโลเกนฺโต สตฺเต อปฺปรชกฺเข มหารชกฺเข ติกฺขินฺทฺริเย มุทินฺทฺริเย สฺวากาเร สุวิญฺญาปเย อปฺเปกจฺเจ ปรโลกวชฺชภยทสฺสาวิโน วิหรนฺเต”ติ\footnote{ม 1. 225 ปิฏฺเฐ. ตตฺถ “ทฺวากาเร ทุวิญฺญาปเย” อิจฺจาทีนิปิ ปทานิ ทิสฺสนฺติ.} อตฺเถว สุตฺตนฺโตติ, อามนฺตา.\csromanspacer{} เตน หิ อตฺถิ โลกิยานิ ปญฺจินฺทฺริยานีติ.\csromanspacer{} อินฺทฺริยกถา นิฏฺฐิตา.\csromanspacer{} เอกูนวีสติโม วคฺโค.}

\titlehead{\textbf{ตสฺสุทฺทานํ}}
\prose{อตีเต กิเลเส ชหติ, อนาคเต กิเลเส ชหติ, ปจฺจุปฺปนฺเน กิเลเส ชหติ, สุญฺญตา สงฺขารกฺขนฺธปริยาปนฺนา, สามญฺญผลํ อสงฺขตํ, ปตฺติ อสงฺขตา, สพฺพธมฺมานํ ตถตา อสงฺขตา, นิพฺพานธาตุ กุสลา, อตฺถิ ปุถุชฺชนสฺส อจฺจนฺตนิยามตา, นตฺถิ โลกิยานิ ปญฺจินฺทฺริยานีติ.}

\chapterhead{20. วีสติมวคฺค}
\vspace{\tipitakaparskip}
\csromancenter{\textbf{(194) 1.}\csromanspacer{}\textbf{ อสญฺจิจฺจกถา}}
\proseitem{857}{\csromansymbolmark{*}อสญฺจิจฺจ มาตรํ ชีวิตา โวโรเปตฺวา อานนฺตริโก โหตีติ, อามนฺตา.\csromanspacer{} อสญฺจิจฺจ ปาณํ หนฺตฺวา ปาณาติปาตี โหตีติ, น \csromanfolio{428}เหวํ วตฺตพฺเพ ฯเปฯ.\csromanspacer{} อสญฺจิจฺจ มาตรํ ชีวิตา โวโรเปตฺวา อานนฺตริโก โหตีติ, อามนฺตา.\csromanspacer{} อสญฺจิจฺจ อทินฺนํ อาทิยิตฺวา ฯเปฯ.\csromanspacer{} มุสา ภณิตฺวา มุสาวาที โหตีติ, น เหวํ วตฺตพฺเพ ฯเปฯ.}

\prose{\csromansymbolmark{*}อสญฺจิจฺจ ปาณํ หนฺตฺวา ปาณาติปาตี น โหตีติ, อามนฺตา.\csromanspacer{} อสญฺจิจฺจ มาตรํ ชีวิตา โวโรเปตฺวา อานนฺตริโก น โหตีติ, น เหวํ วตฺตพฺเพ ฯเปฯ.\csromanspacer{} อสญฺจิจฺจ อทินฺนํ อาทิยิตฺวา ฯเปฯ.\csromanspacer{} มุสา ภณิตฺวา มุสาวาที น โหตีติ, อามนฺตา.\csromanspacer{} อสญฺจิจฺจ มาตรํ ชีวิตา โวโรเปตฺวา อานนฺตริโก น โหตีติ, น เหวํ วตฺตพฺเพ ฯเปฯ.}

\proseitem{858}{\csromansymbolmark{*}อสญฺจิจฺจ มาตรํ ชีวิตา โวโรเปตฺวา อานนฺตริโก โหตีติ, อามนฺตา. “อสญฺจิจฺจ มาตรํ ชีวิตา โวโรเปตฺวา อานนฺตริโก โหตี”ติ อตฺเถว สุตฺตนฺโตติ, นตฺถิ. “สญฺจิจฺจ มาตรํ ชีวิตา โวโรเปตฺวา อานนฺตริโก โหตี”ติ อตฺเถว สุตฺตนฺโตติ, อามนฺตา.\csromanspacer{} หญฺจิ “สญฺจิจฺจ มาตรํ ชีวิตา โวโรเปตฺวา อานนฺตริโก โหตี”ติ อตฺเถว สุตฺตนฺโต, โน จ วต เร วตฺตพฺเพ “อสญฺจิจฺจ มาตรํ ชีวิตา โวโรเปตฺวา อานนฺตริโก โหตี”ติ.}

\proseitem{859}{\csromansymbolmark{+}น วตฺตพฺพํ “มาตุฆาตโก อานนฺตริโก”ติ, อามนฺตา.\csromanspacer{} นนุ มาตา ชีวิตา โวโรปิตาติ, อามนฺตา.\csromanspacer{} หญฺจิ มาตา ชีวิตา โวโรปิตา, เตน วต เร วตฺตพฺเพ “มาตุฆาตโก อานนฺตริโก”ติ.}

\prose{\csromansymbolmark{+}น วตฺตพฺพํ “ปิตุฆาตโก อานนฺตริโก”ติ, อามนฺตา.\csromanspacer{} นนุ ปิตา ชีวิตา โวโรปิโตติ, อามนฺตา.\csromanspacer{} หญฺจิ ปิตา ชีวิตา โวโรปิโต, เตน วต เร วตฺตพฺเพ “ปิตุฆาตโก อานนฺตริโก”ติ.}

\prose{\csromansymbolmark{+}น วตฺตพฺพํ “อรหนฺตฆาตโก อานนฺตริโก”ติ, อามนฺตา.\csromanspacer{} นนุ อรหา ชีวิตา โวโรปิโตติ, อามนฺตา.\csromanspacer{} หญฺจิ อรหา ชีวิตา โวโรปิโต, เตน วต เร วตฺตพฺเพ “อรหนฺตฆาตโก อานนฺตริโก”ติ.}

\prose{\csromansymbolmark{+}น วตฺตพฺพํ “รุหิรุปฺปาทโก อานนฺตริโก”ติ, อามนฺตา.\csromanspacer{} นนุ ตถาคตสฺส โลหิตํ อุปฺปาทิตนฺติ, อามนฺตา.\csromanspacer{} หญฺจิ ตถาคตสฺส โลหิตํ อุปฺปาทิตํ, เตน วต เร วตฺตพฺเพ “รุหิรุปฺปาทโก อานนฺตริโก”ติ.}

\chapterheadpageread{20. วีสติมวคฺค}
\csromanmatikaanchor{mtk.231}
\csromantocmarkref{subsection}{(195) 2. ญาณกถา}{mtk.231}
\bookmark[dest={mtk.231},level=2]{(195) 2. ญาณกถา}
\proseitem{860}{\csromanfolio{429}\csromansymbolmark{*}สํฆเภทโก อานนฺตริโกติ, อามนฺตา.\csromanspacer{} สพฺเพ สํฆเภทกา อานนฺตริกาติ, น เหวํ วตฺตพฺเพ ฯเปฯ.\csromanspacer{} สพฺเพ สํฆเภทกา อานนฺตริกาติ, อามนฺตา.\csromanspacer{} ธมฺมสญฺญี สํฆเภทโก อานนฺตริโกติ, น เหวํ วตฺตพฺเพ ฯเปฯ.}

\proseitem{861}{\csromansymbolmark{*}ธมฺมสญฺญี สํฆเภทโก อานนฺตริโกติ, อามนฺตา.\csromanspacer{} นนุ วุตฺตํ ภควตา “อตฺถุปาลิ สํฆเภทโก อาปายิโก เนรยิโก กปฺปฏฺโฐ อเตกิจฺโฉ.\csromanspacer{} อตฺถุปาลิ สํฆเภทโก น อาปายิโก น เนรยิโก น กปฺปฏฺโฐ น อเตกิจฺโฉ”ติ อตฺเถว สุตฺตนฺโตติ, อามนฺตา.\csromanspacer{} เตน หิ น วตฺตพฺพํ “ธมฺมสญฺญี สํฆเภทโก อานนฺตริโก”ติ.}

\proseitem{862}{\csromansymbolmark{+}น วตฺตพฺพํ “ธมฺมสญฺญี สํฆเภทโก อานนฺตริโก”ติ, อามนฺตา.\csromanspacer{} นนุ วุตฺตํ ภควตา\csromandash{}}

\csromangathameasure{“อาปายิโก เนรยิโก, กปฺปฏฺโฐ สํฆเภทโก. \\ วคฺครโต อธมฺมฏฺโฐ, โยคกฺเขมา ปธํสติ. \\ สํฆํ สมคฺคํ เภตฺวาน, กปฺปํ นิรยมฺหิ ปจฺจตี”ติ.}
\csromangathagroup{\csromangathastanza{“อาปายิโก เนรยิโก, กปฺปฏฺโฐ สํฆเภทโก. \\ วคฺครโต อธมฺมฏฺโฐ, โยคกฺเขมา ปธํสติ. \\ สํฆํ สมคฺคํ เภตฺวาน, กปฺปํ นิรยมฺหิ ปจฺจตี”ติ\footnote{วิ 4. 369; อํ 3. 315; ขุ 1. 203 ปิฏฺเฐสุ.}.}}

\prose{อตฺเถว สุตฺตนฺโตติ, อามนฺตา.\csromanspacer{} เตน หิ สํฆเภทโก อานนฺตริโกติ.}

\nitthitam{อสญฺจิจฺจกถา นิฏฺฐิตา.}
\chapterhead{20. วีสติมวคฺค}
\vspace{\tipitakaparskip}
\csromancenter{\textbf{(195) 2.}\csromanspacer{}\textbf{ ญาณกถา}}
\proseitem{863}{\csromansymbolmark{*}นตฺถิ ปุถุชฺชนสฺส ญาณนฺติ, อามนฺตา.\csromanspacer{} นตฺถิ ปุถุชฺชนสฺส ปญฺญา ปชานนา วิจโย ปวิจโย ธมฺมวิจโย สลฺลกฺขณา อุปลกฺขณา ปจฺจุปลกฺขณาติ, น เหวํ วตฺตพฺเพ ฯเปฯ.\csromanspacer{} นนุ อตฺถิ ปุถุชฺชนสฺส ปญฺญา ปชานนา วิจโย ฯเปฯ ปจฺจุปลกฺขณาติ, อามนฺตา.\csromanspacer{} หญฺจิ อตฺถิ ปุถุชฺชนสฺส ปญฺญา ปชานนา วิจโย ฯเปฯ ปจฺจุปลกฺขณา, โน จ วต เร วตฺตพฺเพ “นตฺถิ ปุถุชฺชนสฺส ญาณนฺ”ติ.}

\csromanmatikaanchor{mtk.232}
\csromantocmarkref{subsection}{(196) 3. นิรยปาลกถา}{mtk.232}
\bookmark[dest={mtk.232},level=2]{(196) 3. นิรยปาลกถา}
\proseitem{864}{\csromanfolio{430}\csromansymbolmark{*}นตฺถิ ปุถุชฺชนสฺส ญาณนฺติ, อามนฺตา.\csromanspacer{} ปุถุชฺชโน ปฐมํ ฌานํ สมาปชฺเชยฺยาติ, อามนฺตา.\csromanspacer{} หญฺจิ ปุถุชฺชโน ปฐมํ ฌานํ สมาปชฺเชยฺย, โน จ วต เร วตฺตพฺเพ “นตฺถิ ปุถุชฺชนสฺส ญาณนฺ”ติ.}

\prose{\csromansymbolmark{*}ปุถุชฺชโน ทุติยํ ฌานํ ฯเปฯ ตติยํ ฌานํ ฯเปฯ จตุตฺถํ ฌานํ ฯเปฯ อากาสานญฺจายตนํ สมาปชฺเชยฺย.\csromanspacer{} วิญฺญาณญฺจายตนํ.\csromanspacer{} อากิญฺจญฺญายตนํ.\csromanspacer{} เนวสญฺญานาสญฺญายตนํ สมาปชฺเชยฺย.\csromanspacer{} ปุถุชฺชโน ทานํ ทเทยฺย ฯเปฯ จีวรํ ทเทยฺย.\csromanspacer{} ปิณฺฑปาตํ ทเทยฺย.\csromanspacer{} เสนาสนํ ทเทยฺย.\csromanspacer{} คิลานปจฺจยเภสชฺชปริกฺขารํ ทเทยฺยาติ, อามนฺตา.\csromanspacer{} หญฺจิ ปุถุชฺชโน คิลานปจฺจยเภสชฺชปริกฺขารํ ทเทยฺย, โน จ วต เร วตฺตพฺเพ “นตฺถิ ปุถุชฺชนสฺส ญาณนฺ”ติ.}

\proseitem{865}{\csromansymbolmark{+}อตฺถิ ปุถุชฺชนสฺส ญาณนฺติ, อามนฺตา.\csromanspacer{} ปุถุชฺชโน เตน ญาเณน ทุกฺขํ ปริชานาติ, สมุทยํ ปชหติ, นิโรธํ สจฺฉิกโรติ, มคฺคํ ภาเวตีติ, น เหวํ วตฺตพฺเพ ฯเปฯ.}

\nitthitam{ญาณกถา นิฏฺฐิตา.}
\chapterhead{20. วีสติมวคฺค}
\vspace{\tipitakaparskip}
\csromancenter{\textbf{(196) 3.}\csromanspacer{}\textbf{ นิรยปาลกถา}}
\proseitem{866}{\csromansymbolmark{*}นตฺถิ นิรเยสุ นิรยปาลาติ, อามนฺตา.\csromanspacer{} นตฺถิ นิรเยสุ กมฺมการณาติ, น เหวํ วตฺตพฺเพ ฯเปฯ.\csromanspacer{} อตฺถิ นิรเยสุ กมฺมการณาติ, อามนฺตา.\csromanspacer{} อตฺถิ นิรเยสุ นิรยปาลาติ, น เหวํ วตฺตพฺเพ ฯเปฯ.}

\prose{\csromansymbolmark{*}อตฺถิ มนุสฺเสสุ กมฺมการณา, อตฺถิ จ การณิกาติ, อามนฺตา.\csromanspacer{} อตฺถิ นิรเยสุ กมฺมการณา, อตฺถิ จ การณิกาติ, น เหวํ วตฺตพฺเพ ฯเปฯ.\csromanspacer{} อตฺถิ นิรเยสุ กมฺมการณา, นตฺถิ จ การณิกาติ, อามนฺตา.\csromanspacer{} อตฺถิ มนุสฺเสสุ กมฺมการณา, นตฺถิ จ การณิกาติ, น เหวํ วตฺตพฺเพ ฯเปฯ.}

\proseitem{867}{\csromansymbolmark{+}อตฺถิ นิรเยสุ นิรยปาลาติ, อามนฺตา.\csromanspacer{} นนุ วุตฺตํ ภควตา\csromandash{}}

\chapterheadpageread{20. วีสติมวคฺค}
\csromangathameasure{“น เวสฺสภู โนปิ จ เปตฺติราชา, \\ โสโม ยโม เวสฺสวโณ จ ราชา. \\ สกานิ กมฺมานิ หนนฺติ ตตฺถ, \\ อิโต ปณุนฺนํ ปรโลกปตฺตนฺ”ติ.}
\csromangathagroup{\csromangathastanza{\csromanfolio{431}“น เวสฺสภู โนปิ จ เปตฺติราชา, \\ โสโม ยโม เวสฺสวโณ จ ราชา. \\ สกานิ กมฺมานิ หนนฺติ ตตฺถ, \\ อิโต ปณุนฺนํ\footnote{ปณุณฺณํ (สฺยา, ก), ปนุนฺนํ (สี)} ปรโลกปตฺตนฺ”ติ.}}

\prose{อตฺเถว สุตฺตนฺโตติ, อามนฺตา.\csromanspacer{} เตน หิ นตฺถิ นิรเยสุ นิรยปาลาติ.}

\proseitem{868}{\csromansymbolmark{*}นตฺถิ นิรเยสุ นิรยปาลาติ, อามนฺตา.\csromanspacer{} นนุ วุตฺตํ ภควตา “ตเมนํ ภิกฺขเว นิรยปาลา ปญฺจวิธพนฺธนํ นาม การณํ กาเรนฺติ\footnote{กโรนฺติ (ม 3. 204; อํ 1. 139 ปิฏฺเฐ.)}, ตตฺตํ อโยขีลํ หตฺเถ คเมนฺติ, ตตฺตํ อโยขีลํ ทุติเย หตฺเถ คเมนฺติ, ตตฺตํ อโยขีลํ ปาเท คเมนฺติ, ตตฺตํ อโยขีลํ ทุติเย ปาเท คเมนฺติ, ตตฺตํ อโยขีลํ มชฺเฌ-อุรสฺมิํ คเมนฺติ, โส ตตฺถ ทุกฺขา ติพฺพา กฏุกา เวทนา เวเทติ, น จ ตาว กาลํ กโรติ, ยาว น ตํ ปาปกมฺมํ พฺยนฺตีโหตี”ติ\footnote{ม 3. 204 ปิฏฺเฐ; อํ 1. 139 ปิฏฺเฐ จ.} อตฺเถว สุตฺตนฺโตติ, อามนฺตา.\csromanspacer{} เตน หิ อตฺถิ นิรเยสุ นิรยปาลาติ.}

\csromanmatikaanchor{mtk.233}
\csromantocmarkref{subsection}{(197) 4. ติรจฺฉานกถา}{mtk.233}
\bookmark[dest={mtk.233},level=2]{(197) 4. ติรจฺฉานกถา}
\prose{\csromansymbolmark{*}นตฺถิ นิรเยสุ นิรยปาลาติ, อามนฺตา.\csromanspacer{} นนุ วุตฺตํ ภควตา “ตเมนํ ภิกฺขเว นิรยปาลา สํเวเสตฺวา กุฐารีหิ\footnote{กุธารีหิ (สพฺพตฺถ)} ตจฺฉนฺติ ฯเปฯ.\csromanspacer{} ตเมนํ ภิกฺขเว นิรยปาลา อุทฺธํปาทํ อโธสิรํ ฐเปตฺวา วาสีหิ ตจฺฉนฺติ ฯเปฯ.\csromanspacer{} ตเมนํ ภิกฺขเว นิรยปาลา รเถ โยเชตฺวา อาทิตฺตาย ปถวิยา สมฺปชฺชลิตาย สโชติภูตาย\footnote{สญฺโชติภูตาย (สฺยา)} สาเรนฺติปิ ปจฺจาสาเรนฺติปิ ฯเปฯ.\csromanspacer{} ตเมนํ ภิกฺขเว นิรยปาลา มหนฺตํ องฺคารปพฺพตํ อาทิตฺตํ สมฺปชฺชลิตํ สโชติภูตํ อาโรเปนฺติปิ โอโรเปนฺติปิ ฯเปฯ.\csromanspacer{} ตเมนํ ภิกฺขเว นิรยปาลา อุทฺธํปาทํ อโธสิรํ คเหตฺวา ตตฺตาย โลหกุมฺภิยา ปกฺขิปนฺติ อาทิตฺตาย สมฺปชฺชลิตาย สโชติภูตาย.\csromanspacer{} โส ตตฺถ เผณุทฺเทหกํ ปจฺจติ, โส ตตฺถ เผณุทฺเทหกํ ปจฺจมาโน สกิมฺปิ อุทฺธํ คจฺฉติ, สกิมฺปิ อโธ คจฺฉติ, สกิมฺปิ ติริยํ คจฺฉติ.\csromanspacer{} โส ตตฺถ ทุกฺขา ติพฺพา กฏุกา เวทนา เวทยติ, น จ ตาว กาลํ กโรติ, ยาว \csromanfolio{432}น ตํ ปาปกมฺมํ พฺยนฺตีโหติ.\csromanspacer{} ตเมนํ ภิกฺขเว นิรยปาลา มหานิรเย ปกฺขิปนฺติ.\csromanspacer{} โส โข ปน ภิกฺขเว มหานิรโย.}

\csromangathameasure{จตุกฺกณฺโณ จตุทฺวาโร, วิภตฺโต ภาคโส มิโต. \\ อโยปาการปริยนฺโต, อยสา ปฏิกุชฺชิโต. \\ ตสฺส อโยมยา ภูมิ, ชลิตา เตชสา ยุตา. \\ สมนฺตา โยชนสตํ, ผริตฺวา ติฏฺฐติ สพฺพาทา”ติ.}
\csromangathagroup{\csromangathastanza{จตุกฺกณฺโณ จตุทฺวาโร, วิภตฺโต ภาคโส มิโต. \\ อโยปาการปริยนฺโต, อยสา ปฏิกุชฺชิโต.}\csromangathastanza{ตสฺส อโยมยา ภูมิ, ชลิตา เตชสา ยุตา. \\ สมนฺตา โยชนสตํ, ผริตฺวา ติฏฺฐติ สพฺพาทา”ติ\footnote{ม 3. 204; อํ 1. 140; ขุ 2. 135 ปิฏฺเฐสุ.}.}}

\prose{อตฺเถว สุตฺตนฺโตติ, อามนฺตา.\csromanspacer{} เตน หิ อตฺถิ นิรเยสุ นิรยปาลาติ.}

\nitthitam{นิรยปาลกถา นิฏฺฐิตา.}
\chapterhead{20. วีสติมวคฺค}
\vspace{\tipitakaparskip}
\csromancenter{\textbf{(197) 4.}\csromanspacer{}\textbf{ ติรจฺฉานกถา}}
\proseitem{869}{\csromansymbolmark{*}อตฺถิ เทเวสุ ติรจฺฉานคตาติ, อามนฺตา.\csromanspacer{} อตฺถิ ติรจฺฉานคเตสุ เทวาติ, น เหวํ วตฺตพฺเพ ฯเปฯ.\csromanspacer{} อตฺถิ เทเวสุ ติรจฺฉานคตาติ, อามนฺตา.\csromanspacer{} เทวโลโก ติรจฺฉานโยนีติ, น เหวํ วตฺตพฺเพ ฯเปฯ.\csromanspacer{} อตฺถิ เทเวสุ ติรจฺฉานคตาติ, อามนฺตา.\csromanspacer{} อตฺถิ ตตฺถ กีฏา ปฏงฺคา มกสา มกฺขิกา อหี วิจฺฉิกา สตปที คณฺฑุปฺปาทาติ\footnote{คณฺฑุปาทาติ (สฺยา, กํ, อิ)}.\csromanspacer{} น เหวํ วตฺตพฺเพ ฯเปฯ.}

\proseitem{870}{\csromansymbolmark{+}นตฺถิ เทเวสุ ติรจฺฉานคตาติ, อามนฺตา.\csromanspacer{} นนุ อตฺถิ ตตฺถ เอราวโณ นาม หตฺถินาโค, สหสฺสยุตฺตํ ทิพฺพํ ยานนฺติ, อามนฺตา.\csromanspacer{} หญฺจิ อตฺถิ ตตฺถ เอราวโณ นาม หตฺถินาโค, สหสฺสยุตฺตํ ทิพฺพํ ยานํ, เตน วต เร วตฺตพฺเพ “อตฺถิ เทเวสุ ติรจฺฉานคตา”ติ.}

\proseitem{871}{\csromansymbolmark{*}อตฺถิ เทเวสุ ติรจฺฉานคตาติ, อามนฺตา.\csromanspacer{} อตฺถิ ตตฺถ หตฺถิพนฺธา อสฺสพนฺธา\footnote{หตฺถิภณฺฑา อสฺสภณฺฑา (?)} ยาวสิกา การณิกา ภตฺตการกาติ, น เหวํ วตฺตพฺเพ.\csromanspacer{} เตน หิ นตฺถิ เทเวสุ ติรจฺฉานคตาติ.}

\nitthitam{ติรจฺฉานกถา นิฏฺฐิตา.}
\chapterheadpageread{20. วีสติมวคฺค \textbf{20. วีสติมวคฺค}}
\vspace{\tipitakaparskip}
\csromancenter{\textbf{(198) 5.}\csromanspacer{}\textbf{ มคฺคกถา}}
\csromanmatikaanchor{mtk.234}
\csromantocmarkref{subsection}{(198) 5. มคฺคกถา}{mtk.234}
\bookmark[dest={mtk.234},level=2]{(198) 5. มคฺคกถา}
\proseitem{872}{\csromanfolio{433}\csromansymbolmark{*}ปญฺจงฺคิโก มคฺโคติ, อามนฺตา.\csromanspacer{} นนุ อฏฺฐงฺคิโก มคฺโค วุตฺโต ภควตา.\csromanspacer{} เสยฺยถิทํ.\csromanspacer{} สมฺมาทิฏฺฐิ ฯเปฯ สมฺมาสมาธีติ, อามนฺตา.\csromanspacer{} หญฺจิ อฏฺฐงฺคิโก มคฺโค วุตฺโต ภควตา.\csromanspacer{} เสยฺยถิทํ, สมฺมาทิฏฺฐิ ฯเปฯ สมฺมาสมาธิ, โน จ วต เร วตฺตพฺเพ “ปญฺจงฺคิโก มคฺโค”ติ.}

\prose{\csromansymbolmark{*}ปญฺจงฺคิโก มคฺโคติ, อามนฺตา.\csromanspacer{} นนุ วุตฺตํ ภควตา\csromandash{}}

\csromangathameasure{“มคฺคานํ อฏฺฐงฺคิโก เสฏฺโฐ, สจฺจานํ จตุโร ปทา. \\ วิราโค เสฏฺโฐ ธมฺมานํ, ทฺวิปทานญฺจ จกฺขุมา”ติ.}
\csromangathagroup{\csromangathastanza{“มคฺคานํ อฏฺฐงฺคิโก เสฏฺโฐ, สจฺจานํ จตุโร ปทา. \\ วิราโค เสฏฺโฐ ธมฺมานํ, ทฺวิปทานญฺจ จกฺขุมา”ติ\footnote{ขุ 1. 52 ธมฺมปเท.}.}}

\prose{อตฺเถว สุตฺตนฺโตติ, อามนฺตา.\csromanspacer{} เตน หิ อฏฺฐงฺคิโก มคฺโคติ.}

\proseitem{873}{\csromansymbolmark{*}สมฺมาวาจา มคฺคงฺคํ, สา จ น มคฺโคติ, อามนฺตา.\csromanspacer{} สมฺมาทิฏฺฐิ มคฺคงฺคํ, สา จ น มคฺโคติ, น เหวํ วตฺตพฺเพ ฯเปฯ.\csromanspacer{} สมฺมาวาจา มคฺคงฺคํ, สา จ น มคฺโคติ, อามนฺตา.\csromanspacer{} สมฺมาสงฺกปฺโป ฯเปฯ.\csromanspacer{} สมฺมาวายาโม ฯเปฯ.\csromanspacer{} สมฺมาสติ ฯเปฯ.\csromanspacer{} สมฺมาสมาธิ มคฺคงฺคํ, โส จ น มคฺโคติ, น เหวํ วตฺตพฺเพ ฯเปฯ.}

\prose{\csromansymbolmark{*}สมฺมากมฺมนฺโต ฯเปฯ.\csromanspacer{} สมฺมา อาชีโว มคฺคงฺคํ, โส จ น มคฺโคติ, อามนฺตา.\csromanspacer{} สมฺมาทิฏฺฐิ ฯเปฯ.\csromanspacer{} สมฺมาสมาธิ มคฺคงฺคํ, โส จ น มคฺโคติ, น เหวํ วตฺตพฺเพ ฯเปฯ.}

\prose{\csromansymbolmark{*}สมฺมาทิฏฺฐิ มคฺคงฺคํ, สา จ มคฺโคติ, อามนฺตา.\csromanspacer{} สมฺมาวาจา มคฺคงฺคํ, สา จ มคฺโคติ, น เหวํ วตฺตพฺเพ ฯเปฯ.\csromanspacer{} สมฺมาทิฏฺฐิ มคฺคงฺคํ, สา จ มคฺโคติ, อามนฺตา.\csromanspacer{} สมฺมากมฺมนฺโต ฯเปฯ.\csromanspacer{} สมฺมา-อาชีโว มคฺคงฺคํ, โส จ มคฺโคติ, น เหวํ วตฺตพฺเพ ฯเปฯ.}

\prose{\csromansymbolmark{*}สมฺมาสงฺกปฺโป ฯเปฯ.\csromanspacer{} สมฺมาวายาโม ฯเปฯ.\csromanspacer{} สมฺมาสติ ฯเปฯ.\csromanspacer{} สมฺมาสมาธิ มคฺคงฺคํ, โส จ มคฺโคติ, อามนฺตา.\csromanspacer{} สมฺมาวาจา ฯเปฯ.\csromanspacer{} สมฺมากมฺมนฺโต ฯเปฯ.\csromanspacer{} สมฺมา อาชีโว มคฺคงฺคํ, โส จ มคฺโคติ, น เหวํ วตฺตพฺเพ ฯเปฯ.}

\csromanmatikaanchor{mtk.235}
\csromantocmarkref{subsection}{(199) 6. ญาณกถา}{mtk.235}
\bookmark[dest={mtk.235},level=2]{(199) 6. ญาณกถา}
\proseitem{874}{\csromansymbolmark{+}อริโย อฏฺฐงฺคิโก มคฺโคติ, อามนฺตา.\csromanspacer{} นนุ วุตฺตํ ภควตา “ปุพฺเพว โข ปนสฺส กายกมฺมํ วจีกมฺมํ อาชีโว สุปริสุทฺโธ โหติ. \csromanfolio{434}เอวมสฺสายํ อริโย อฏฺฐงฺคิโก มคฺโค ภาวนาปาริปูริํ คจฺฉตี”ติ\footnote{ม 3. 337 ปิฏฺเฐ.} อตฺเถว สุตฺตนฺโตติ, อามนฺตา.\csromanspacer{} เตน หิ ปญฺจงฺคิโก มคฺโคติ.}

\proseitem{875}{\csromansymbolmark{*}ปญฺจงฺคิโก มคฺโคติ, อามนฺตา.\csromanspacer{} นนุ วุตฺตํ ภควตา “ยสฺมิํ โข สุภทฺท ธมฺมวินเย อริโย อฏฺฐงฺคิโก มคฺโค น อุปลพฺภติ, สมโณปิ ตตฺถ น อุปลพฺภติ, ทุติโยปิ ตตฺถ สมโณ น อุปลพฺภติ, ตติโยปิ ตตฺถ สมโณ น อุปลพฺภติ, จตุตฺโถปิ ตตฺถ สมโณ น อุปลพฺภติ.\csromanspacer{} ยสฺมิํ จ โข สุภทฺท ธมฺมวินเย อริโย อฏฺฐงฺคิโก มคฺโค อุปลพฺภติ, สมโณปิ ตตฺถ อุปลพฺภติ, ทุติโยปิ ฯเปฯ ตติโยปิ ฯเปฯ จตุตฺโถปิ ตตฺถ สมโณ อุปลพฺภติ.\csromanspacer{} อิมสฺมิํ โข สุภทฺท ธมฺมวินเย อริโย อฏฺฐงฺคิโก มคฺโค อุปลพฺภติ, อิเธว สุภทฺท สมโณ, อิธ ทุติโย สมโณ, อิธ ตติโย สมโณ, อิธ จตุตฺโถ สมโณ, สุญฺญา ปรปฺปวาทา สมเณภิ อญฺเญหี”ติ\footnote{ที 2. 125 ปิฏฺเฐ.} อตฺเถว สุตฺตนฺโตติ, อามนฺตา.\csromanspacer{} เตน หิ อฏฺฐงฺคิโก มคฺโคติ.}

\nitthitam{มคฺคกถา นิฏฺฐิตา.}
\chapterhead{20. วีสติมวคฺค}
\vspace{\tipitakaparskip}
\csromancenter{\textbf{(199) 6.}\csromanspacer{}\textbf{ ญาณกถา}}
\proseitem{876}{\csromansymbolmark{*}ทฺวาทสวตฺถุกํ ญาณํ โลกุตฺตรนฺติ, อามนฺตา.\csromanspacer{} ทฺวาทส โลกุตฺตรญาณานีติ, น เหวํ วตฺตพฺเพ ฯเปฯ.\csromanspacer{} ทฺวาทส โลกุตฺตรญาณานีติ, อามนฺตา.\csromanspacer{} ทฺวาทส โสตาปตฺติมคฺคาติ, น เหวํ วตฺตพฺเพ ฯเปฯ.\csromanspacer{} ทฺวาทส โสตาปตฺติมคฺคาติ, อามนฺตา.\csromanspacer{} ทฺวาทส โสตาปตฺติผลานีติ, น เหวํ วตฺตพฺเพ ฯเปฯ.\csromanspacer{} ทฺวาทส สกทาคามิมคฺคา ฯเปฯ.\csromanspacer{} อนาคามิมคฺคา ฯเปฯ.\csromanspacer{} อรหตฺตมคฺคาติ, น เหวํ วตฺตพฺเพ ฯเปฯ.\csromanspacer{} ทฺวาทส อรหตฺตมคฺคาติ, อามนฺตา.\csromanspacer{} ทฺวาทส อรหตฺตผลานีติ, น เหวํ วตฺตพฺเพ ฯเปฯ.}

\proseitem{877}{\csromansymbolmark{+}น วตฺตพฺพํ “ทฺวาทสวตฺถุกํ ญาณํ โลกุตฺตรนฺ”ติ, อามนฺตา.\csromanspacer{} นนุ วุตฺตํ ภควตา “อิทํ ทุกฺขํ อริยสจฺจนฺ”ติ เม ภิกฺขเว ปุพฺเพ อนนุสฺสุเตสุ}

\vspace{\tipitakaparskip}
\csromancenter{วีสติมวคฺค ธมฺเมสุ จกฺขุํ อุทปาทิ ญาณํ อุทปาทิ ปญฺญา อุทปาทิ วิชฺชา อุทปาทิ อาโลโก อุทปาทิ, “ตํ โข ปนิทํ ทุกฺขํ อริยสจฺจํ ปริญฺเญยฺยนฺ”ติ เม ภิกฺขเว ฯเปฯ ปริญฺญาตนฺ”ติ เม ภิกฺขเว ฯเปฯ “อิทํ ทุกฺขสมุทยํ\footnote{ทุกฺขสมุทโย (สฺยา, กํ, อิ)} อริยสจฺจนฺ”ติ เม ภิกฺขเว ฯเปฯ “ตํ โข ปนิทํ ทุกฺขสมุทยํ อริยสจฺจํ ปหาตพฺพนฺ”ติ เม ภิกฺขเว ฯเปฯ ปหีนนฺ”ติ เม ภิกฺขเว ฯเปฯ “อิทํ ทุกฺขนิโรธํ\footnote{ทุกฺขนิโรโธ (สฺยา, กํ, อิ)} อริยสจฺจนฺ”ติ เม ภิกฺขเว ฯเปฯ “ตํ โข ปนิทํ ทุกฺขนิโรธํ อริยสจฺจํ สจฺฉิกาตพฺพนฺ”ติ เม ภิกฺขเว ฯเปฯ สจฺฉิกตนฺ”ติ เม ภิกฺขเว ฯเปฯ “อิทํ ทุกฺขนิโรธคามินี ปฏิปทา อริยสจฺจนฺ”ติ เม ภิกฺขเว ฯเปฯ “ตํ โข ปนิทํ ทุกฺขนิโรธคามินี ปฏิปทา อริยสจฺจํ ภาเวตพฺพนฺ”ติ เม ภิกฺขเว ฯเปฯ ภาวิตนฺ”ติ เม ภิกฺขเว ปุพฺเพ อนนุสฺสุเตสุ ธมฺเมสุ จกฺขุํ อุทปาทิ ฯเปฯ อาโลโก อุทปาที”ติ\footnote{วิ 3. 15; สํ 3. 369; ขุ 9. 330 ปิฏฺเฐสุ.} อตฺเถว สุตฺตนฺโตติ, อามนฺตา.\csromanspacer{} เตน หิ ทฺวาทสวตฺถุกํ ญาณํ โลกุตฺตรนฺติ.\csromanspacer{} ญาณกถา นิฏฺฐิตา.\csromanspacer{} วีสติโม วคฺโค.}
\titlehead{\textbf{ตสฺสุทฺทานํ}}
\prose{\csromanfolio{435}มาตุฆาตโก อานนฺตริโก, ปิตุฆาตโก อานนฺตริโก, อรหนฺตฆาตโก อานนฺตริโก, รุหิรุปฺปาทโก อานนฺตริโก, สํฆเภทโก อานนฺตริโก, นตฺถิ ปุถุชฺชนสฺส ญาณํ, นตฺถิ นิรเยสุ นิรยปาลา, อตฺถิ เทเวสุ ติรจฺฉานคตา, ปญฺจงฺคิโก มคฺโค, ทฺวาทสวตฺถุกํ ญาณํ โลกุตฺตรนฺติ.\csromanspacer{} จตุตฺโถ ปณฺณาสโก.}

\titlehead{\textbf{ตสฺสุทฺทานํ}}
\csromangathameasure{นิคฺคโห, ปุญฺญสญฺจโย, อฏฺฐาสิ, อตีเตน จ, มาตุฆาตโก.}
\csromangathagroup{\csromangathastanza{นิคฺคโห, ปุญฺญสญฺจโย, อฏฺฐาสิ, อตีเตน จ, มาตุฆาตโก.}}

\csromanmatikaanchor{mtk.236}
\csromantocmarknopage{chapter}{21. เอกวีสติมวคฺค}{mtk.236}
\bookmark[dest={mtk.236},level=0]{21. เอกวีสติมวคฺค}
\chapterheadpageread{21. เอกวีสติมวคฺค}
\vspace{\tipitakaparskip}
\csromancenter{\textbf{(200) 1.}\csromanspacer{}\textbf{ สาสนกถา}}
\csromanmatikaanchor{mtk.237}
\csromanmatikaanchor{mtk.238}
\csromantocmarkref{subsection}{(200) 1. สาสนกถา}{mtk.237}
\bookmark[dest={mtk.237},level=2]{(200) 1. สาสนกถา}
\csromantocmarkref{subsection}{(201) 2. อวิวิตฺตกถา}{mtk.238}
\bookmark[dest={mtk.238},level=2]{(201) 2. อวิวิตฺตกถา}
\proseitem{878}{\csromanfolio{436}\csromansymbolmark{*}สาสนํ นวํ กตนฺติ, อามนฺตา.\csromanspacer{} สติปฏฺฐานา นวํ กตาติ, น เหวํ วตฺตพฺเพ ฯเปฯ.\csromanspacer{} สาสนํ นวํ กตนฺติ, อามนฺตา.\csromanspacer{} สมฺมปฺปธานา ฯเปฯ.\csromanspacer{} อิทฺธิปาทา ฯเปฯ.\csromanspacer{} อินฺทฺริยา ฯเปฯ.\csromanspacer{} พลา ฯเปฯ.\csromanspacer{} โพชฺฌงฺคา นวํ กตาติ, น เหวํ วตฺตพฺเพ ฯเปฯ.\csromanspacer{} ปุพฺเพ อกุสลํ ปจฺฉา กุสลํ กตนฺติ, น เหวํ วตฺตพฺเพ ฯเปฯ.\csromanspacer{} ปุพฺเพ สาสวํ ฯเปฯ สํโยชนิยํ.\csromanspacer{} คนฺถนิยํ.\csromanspacer{} โอฆนิยํ.\csromanspacer{} โยคนิยํ.\csromanspacer{} นีวรณิยํ.\csromanspacer{} ปรามฏฺฐํ.\csromanspacer{} อุปาทานิยํ ฯเปฯ สํกิเลสิยํ ปจฺฉา อสํกิเลสิยํ กตนฺติ, น เหวํ วตฺตพฺเพ ฯเปฯ.}

\prose{\csromansymbolmark{*}อตฺถิ โกจิ ตถาคตสฺส สาสนํ นวํ กโรตีติ, อามนฺตา.\csromanspacer{} อตฺถิ โกจิ สติปฏฺฐาเน นวํ กโรตีติ, น เหวํ วตฺตพฺเพ ฯเปฯ.\csromanspacer{} อตฺถิ โกจิ สมฺมปฺปธาเน ฯเปฯ อิทฺธิปาเท ฯเปฯ อินฺทฺริเย ฯเปฯ พเล ฯเปฯ โพชฺฌงฺเค นวํ กโรตีติ, น เหวํ วตฺตพฺเพ ฯเปฯ.\csromanspacer{} อตฺถิ โกจิ ปุพฺเพ อกุสลํ ปจฺฉา กุสลํ กโรตีติ, น เหวํ วตฺตพฺเพ ฯเปฯ.\csromanspacer{} อตฺถิ โกจิ ปุพฺเพ สาสวํ ฯเปฯ สํกิเลสิยํ ปจฺฉา อสํกิเลสิยํ กโรตีติ, น เหวํ วตฺตพฺเพ ฯเปฯ.}

\prose{\csromansymbolmark{*}ลพฺภา ตถาคตสฺส สาสนํ ปุน นวํ กาตุนฺติ, อามนฺตา.\csromanspacer{} ลพฺภา สติปฏฺฐานา ปุน นวํ กาตุนฺติ, น เหวํ วตฺตพฺเพ ฯเปฯ.\csromanspacer{} ลพฺภา สมฺมปฺปธานา ฯเปฯ อิทฺธิปาทา ฯเปฯ อินฺทฺริยา ฯเปฯ พลา ฯเปฯ โพชฺฌงฺคา ปุน นวํ กาตุนฺติ, น เหวํ วตฺตพฺเพ ฯเปฯ.\csromanspacer{} ลพฺภา ปุพฺเพ อกุสลํ ปจฺฉา กุสลํ กาตุนฺติ, น เหวํ วตฺตพฺเพ ฯเปฯ.\csromanspacer{} ลพฺภา ปุพฺเพ สาสวํ ฯเปฯ สํกิเลสิยํ ปจฺฉา อสํกิเลสิยํ กาตุนฺติ, น เหวํ วตฺตพฺเพ ฯเปฯ.}

\nitthitam{สาสนกถา นิฏฺฐิตา.}
\chapterhead{21. เอกวีสติมวคฺค}
\vspace{\tipitakaparskip}
\csromancenter{\textbf{(201) 2.}\csromanspacer{}\textbf{ อวิวิตฺตกถา}}
\proseitem{879}{\csromansymbolmark{*}ปุถุชฺชโน เตธาตุเกหิ ธมฺเมหิ อวิวิตฺโตติ, อามนฺตา.\csromanspacer{} ปุถุชฺชโน เตธาตุเกหิ ผสฺเสหิ ฯเปฯ เตธาตุกาหิ เวทนาหิ.}

\vspace{\tipitakaparskip}
\csromancenter{เอกวีสติมวคฺค สญฺญาหิ.\csromanspacer{} เจตนาหิ.\csromanspacer{} จิตฺเตหิ.\csromanspacer{} สทฺธาหิ.\csromanspacer{} วีริเยหิ.\csromanspacer{} สตีหิ.\csromanspacer{} สมาธีหิ ฯเปฯ เตธาตุกาหิ ปญฺญาหิ อวิวิตฺโตติ, น เหวํ วตฺตพฺเพ ฯเปฯ.}
\csromanmatikaanchor{mtk.239}
\csromantocmarkref{subsection}{(202) 3. สํโยชนกถา}{mtk.239}
\bookmark[dest={mtk.239},level=2]{(202) 3. สํโยชนกถา}
\prose{\csromanfolio{437}\csromansymbolmark{*}ปุถุชฺชโน เตธาตุเกหิ กมฺเมหิ อวิวิตฺโตติ, อามนฺตา.\csromanspacer{} ยสฺมิํ ขเณ ปุถุชฺชโน จีวรํ เทติ, ตสฺมิํ ขเณ ปฐมํ ฌานํ อุปสมฺปชฺช วิหรติ ฯเปฯ อากาสานญฺจายตนํ อุปสมฺปชฺช วิหรตีติ, น เหวํ วตฺตพฺเพ ฯเปฯ.\csromanspacer{} ยสมิํ ขเณ ปุถุชฺชโน ปิณฺฑปาตํ เทติ ฯเปฯ เสนาสนํ เทติ ฯเปฯ คิลานปจฺจยเภสชฺชปริกฺขารํ เทติ, ตสฺมิํ ขเณ จตุตฺถํ ฌานํ อุปสมฺปชฺช วิหรติ.\csromanspacer{} เนวสญฺญานาสญฺญายตนํ อุปสมฺปชฺช วิหรตีติ, น เหวํ วตฺตพฺเพ ฯเปฯ.}

\proseitem{880}{\csromansymbolmark{+}น วตฺตพฺพํ “ปุถุชฺชโน เตธาตุเกหิ กมฺเมหิ อวิวิตฺโต”ติ, อามนฺตา.\csromanspacer{} ปุถุชฺชนสฺส รูปธาตุ-อรูปธาตูปคํ กมฺมํ ปริญฺญาตนฺติ, น เหวํ วตฺตพฺเพ.\csromanspacer{} เตน หิ ปุถุชฺชโน เตธาตุเกหิ กมฺเมหิ อวิวิตฺโตติ ฯเปฯ.}

\nitthitam{อวิวิตฺตกถา นิฏฺฐิตา.}
\chapterhead{21. เอกวีสติมวคฺค}
\vspace{\tipitakaparskip}
\csromancenter{\textbf{(202) 3.}\csromanspacer{}\textbf{ สํโยชนกถา}}
\proseitem{881}{\csromansymbolmark{*}อตฺถิ กิญฺจิ สํโยชนํ อปฺปหาย อรหตฺตปฺปตฺตีติ, อามนฺตา.\csromanspacer{} อตฺถิ กิญฺจิ สกฺกายทิฏฺฐิํ อปฺปหาย ฯเปฯ วิจิกิจฺฉํ อปฺปหาย ฯเปฯ สีลพฺพตปรามาสํ อปฺปหาย.\csromanspacer{} ราคํ อปฺปหาย.\csromanspacer{} โทสํ อปฺปหาย.\csromanspacer{} โมหํ อปฺปหาย.\csromanspacer{} อโนตฺตปฺปํ อปฺปหาย อรหตฺตปฺปตฺตีติ, น เหวํ วตฺตพฺเพ ฯเปฯ.}

\prose{\csromansymbolmark{*}อตฺถิ กิญฺจิ สํโยชนํ อปฺปหาย อรหตฺตปฺปตฺตีติ, อามนฺตา.\csromanspacer{} อรหา สราโค สโทโส สโมโห สมาโน สมกฺโข สปฬาโส ส-อุปายาโส สกิเลโสติ, น เหวํ วตฺตพฺเพ ฯเปฯ.\csromanspacer{} นนุ อรหา นิราโค นิทฺโทโส นิมฺโมโห นิมฺมาโน นิมฺมกฺโข นิปฺปฬาโส นิรุปายาโส นิกฺกิเลโสติ, อามนฺตา.\csromanspacer{} หญฺจิ อรหา นิราโค ฯเปฯ นิกฺกิเลโส, โน จ วต เร วตฺตพฺเพ “อตฺถิ กิญฺจิ สํโยชนํ อปฺปหาย อรหตฺตปฺปตฺตี”ติ.}

\csromanmatikaanchor{mtk.240}
\csromantocmarkref{subsection}{(203) 4. อิทฺธิกถา}{mtk.240}
\bookmark[dest={mtk.240},level=2]{(203) 4. อิทฺธิกถา}
\proseitem{882}{\csromanfolio{438}\csromansymbolmark{+}น วตฺตพฺพํ “กิญฺจิ สํโยชนํ อปฺปหาย อรหตฺตปฺปตฺตี”ติ, อามนฺตา.\csromanspacer{} อรหา สพฺพํ พุทฺธวิสยํ ชานาตีติ, น เหวํ วตฺตพฺเพ.\csromanspacer{} เตน หิ อตฺถิ กิญฺจิ สํโยชนํ อปฺปหาย อรหตฺตปฺปตฺตีติ.}

\nitthitam{สํโยชนกถา นิฏฺฐิตา.}
\chapterhead{21. เอกวีสติมวคฺค}
\vspace{\tipitakaparskip}
\csromancenter{\textbf{(203) 4.}\csromanspacer{}\textbf{ อิทฺธิกถา}}
\proseitem{883}{\csromansymbolmark{*}อตฺถิ อธิปฺปาย-อิทฺธิ พุทฺธานํ วา สาวกานํ วาติ, อามนฺตา. “นิจฺจปณฺณา รุกฺขา โหนฺตู”ติ อตฺถิ อธิปฺปาย-อิทฺธิ พุทฺธานํ วา สาวกานํ วาติ, น เหวํ วตฺตพฺเพ ฯเปฯ “นิจฺจปุปฺผา รุกฺขา โหนฺตุ ฯเปฯ.\csromanspacer{} นิจฺจผลิกา รุกฺขา โหนฺตุ.\csromanspacer{} นิจฺจํ ชุณฺหํ โหตุ.\csromanspacer{} นิจฺจํ เขมํ โหตุ.\csromanspacer{} นิจฺจํ สุภิกฺขํ โหตุ.\csromanspacer{} นิจฺจํ สุวตฺถิ\footnote{สุวุฏฺฐิกํ (สฺยา)} โหตู”ติ อตฺถิ อธิปฺปาย-อิทฺธิ พุทฺธานํ วา สาวกานํ วาติ, น เหวํ วตฺตพฺเพ ฯเปฯ.}

\prose{\csromansymbolmark{*}อตฺถิ อธิปฺปาย-อิทฺธิ พุทฺธานํ วา สาวกานํ วาติ, อามนฺตา.\csromanspacer{} อุปฺปนฺโน ผสฺโส “มา นิรุชฺฌี”ติ อตฺถิ อธิปฺปาย-อิทฺธิ พุทฺธานํ วา สาวกานํ วาติ, น เหวํ วตฺตพฺเพ ฯเปฯ.}

\prose{\csromansymbolmark{*}อุปฺปนฺนา เวทนา ฯเปฯ สญฺญา.\csromanspacer{} เจตนา.\csromanspacer{} จิตฺตํ.\csromanspacer{} สทฺธา.\csromanspacer{} วีริยํ.\csromanspacer{} สติ.\csromanspacer{} สมาธิ ฯเปฯ.\csromanspacer{} ปญฺญา “มา นิรุชฺฌี”ติ อตฺถิ อธิปฺปาย-อิทฺธิ พุทฺธานํ วา สาวกานํ วาติ, น เหวํ วตฺตพฺเพ ฯเปฯ.}

\prose{\csromansymbolmark{*}อตฺถิ อธิปฺปาย-อิทฺธิ พุทฺธานํ วา สาวกานํ วาติ, อามนฺตา. “รูปํ นิจฺจํ โหตู”ติ อตฺถิ อธิปฺปาย-อิทฺธิ เวทนา ฯเปฯ.\csromanspacer{} สญฺญา ฯเปฯ.\csromanspacer{} สงฺขารา ฯเปฯ. “วิญฺญาณํ นิจฺจํ โหตู”ติ อตฺถิ อธิปฺปาย-อิทฺธิ พุทฺธานํ วา สาวกานํ วาติ, น เหวํ วตฺตพฺเพ ฯเปฯ.}

\prose{\csromansymbolmark{*}อตฺถิ อธิปฺปาย-อิทฺธิ พุทฺธานํ วา สาวกานํ วาติ, อามนฺตา.\csromanspacer{} ชาติธมฺมา สตฺตา “มา ชายิํสู”ติ อตฺถิ ฯเปฯ.\csromanspacer{} ชราธมฺมา สตฺตา “มา ชีริํสู”ติ ฯเปฯ.\csromanspacer{} พฺยาธิธมฺมา สตฺตา “มา พฺยาธิยิํสู”ติ ฯเปฯ.\csromanspacer{} มรณธมฺมา สตฺตา “มา มียิํสู”ติ อตฺถิ อธิปฺปาย-อิทฺธิ พุทฺธานํ วา สาวกานํ วาติ, น เหวํ วตฺตพฺเพ ฯเปฯ.}

\chapterheadpageread{21. เอกวีสติมวคฺค}
\csromanmatikaanchor{mtk.241}
\csromanmatikaanchor{mtk.242}
\csromantocmarkref{subsection}{(204) 5. พุทฺธกถา}{mtk.241}
\bookmark[dest={mtk.241},level=2]{(204) 5. พุทฺธกถา}
\csromantocmarkref{subsection}{(205) 6. สพฺพทิสากถา}{mtk.242}
\bookmark[dest={mtk.242},level=2]{(205) 6. สพฺพทิสากถา}
\proseitem{884}{\csromanfolio{439}\csromansymbolmark{+}น วตฺตพฺพํ “อตฺถิ อธิปฺปาย-อิทฺธิ พุทฺธานํ วา สาวกานํ วา”ติ, อามนฺตา.\csromanspacer{} นนุ อายสฺมา ปิลินฺทวจฺโฉ รญฺโญ มาคธสฺส เสนิยสฺส พิมฺพิสารสฺส ปาสาทํ “สุวณฺณนฺ”เตฺวว อธิมุจฺจิ, สุวณฺโณ จ ปน อาสีติ\footnote{วิ 1. 362; วิ 3. 302 ปิฏฺเฐสุ.}, อามนฺตา.\csromanspacer{} หญฺจิ อายสฺมา ปิลินฺทวจฺโฉ รญฺโญ มาคธสฺส เสนิยสฺส พิมฺพิสารสฺส ปาสาทํ สุวณฺณนฺเตฺวว อธิมุจฺจิ, สุวณฺโณ จ ปน อาสิ, เตน วต เร วตฺตพฺเพ “อตฺถิ อธิปฺปาย-อิทฺธิ พุทฺธานํ วา สาวกานํ วา”ติ.}

\nitthitam{อิทฺธิกถา นิฏฺฐิตา.}
\chapterhead{21. เอกวีสติมวคฺค}
\vspace{\tipitakaparskip}
\csromancenter{\textbf{(204) 5.}\csromanspacer{}\textbf{ พุทฺธกถา}}
\proseitem{885}{\csromansymbolmark{*}อตฺถิ พุทฺธานํ พุทฺเธหิ หีนาติเรกตาติ, อามนฺตา.\csromanspacer{} สติปฏฺฐานโตติ, น เหวํ วตฺตพฺเพ ฯเปฯ.\csromanspacer{} สมฺมปฺปธานโต ฯเปฯ.\csromanspacer{} อิทฺธิปาทโต.\csromanspacer{} อินฺทฺริยโต.\csromanspacer{} พลโต.\csromanspacer{} โพชฺฌงฺคโต.\csromanspacer{} วสิภาวโต ฯเปฯ.\csromanspacer{} สพฺพญฺญุตญาณทสฺสนโตติ, น เหวํ วตฺตพฺเพ ฯเปฯ.}

\nitthitam{พุทฺธกถา นิฏฺฐิตา.}
\chapterhead{21. เอกวีสติมวคฺค}
\vspace{\tipitakaparskip}
\csromancenter{\textbf{(205) 6.}\csromanspacer{}\textbf{ สพฺพทิสากถา}}
\proseitem{886}{\csromansymbolmark{*}สพฺพา ทิสา พุทฺธา ติฏฺฐนฺตีติ, อามนฺตา.\csromanspacer{} ปุรตฺถิมาย ทิสาย พุทฺโธ ติฏฺฐตีติ, น เหวํ วตฺตพฺเพ ฯเปฯ.\csromanspacer{} ปุรตฺถิมาย ทิสาย พุทฺโธ ติฏฺฐตีติ, อามนฺตา.\csromanspacer{} กินฺนาโม โส ภควา, กิํชจฺโจ, กิํโคตฺโต, กินฺนามา ตสฺส ภควโต มาตาปิตโร, กินฺนามํ ตสฺส ภควโต สาวกยุคํ, โกนาโม ตสฺส ภควโต อุปฏฺฐาโก, กีทิสํ จีวรํ ธาเรติ, กีทิสํ ปตฺตํ ธาเรติ, กตรสฺมิํ คาเม วา นิคเม วา นคเร วา รฏฺเฐ วา ชนปเท วาติ, น เหวํ วตฺตพฺเพ ฯเปฯ.}

\csromanmatikaanchor{mtk.243}
\csromantocmarkref{subsection}{(206) 7. ธมฺมกถา}{mtk.243}
\bookmark[dest={mtk.243},level=2]{(206) 7. ธมฺมกถา}
\prose{\csromansymbolmark{*}ทกฺขิณาย ทิสาย ฯเปฯ.\csromanspacer{} ปจฺฉิมาย ทิสาย ฯเปฯ.\csromanspacer{} อุตฺตราย ทิสาย ฯเปฯ.\csromanspacer{} เหฏฺฐิมาย ทิสาย พุทฺโธ ติฏฺฐตีติ, น เหวํ วตฺตพฺเพ ฯเปฯ.\csromanspacer{} เหฏฺฐิมาย \csromanfolio{440}ทิสาย พุทฺโธ ติฏฺฐตีติ, อามนฺตา.\csromanspacer{} กินฺนาโม โส ภควา ฯเปฯ.\csromanspacer{} ชนปเท วาติ, น เหวํ วตพฺเพ ฯเปฯ.}

\prose{\csromansymbolmark{*}อุปริมาย ทิสาย พุทฺโธ ติฏฺฐตีติ, น เหวํ วตฺตพฺเพ ฯเปฯ.\csromanspacer{} อุปริมาย ทิสาย พุทฺโธ ติฏฺฐตีติ, อามนฺตา.\csromanspacer{} จาตุมหาราชิเก ติฏฺฐติ ฯเปฯ.\csromanspacer{} ตาวติํเส ติฏฺฐติ ฯเปฯ.\csromanspacer{} ยาเม ติฏฺฐติ ฯเปฯ.\csromanspacer{} ตุสิเต ติฏฺฐติ ฯเปฯ.\csromanspacer{} นิมฺมานรติยา ติฏฺฐติ ฯเปฯ.\csromanspacer{} ปรนิมฺมิตวสวตฺติยา ติฏฺฐติ ฯเปฯ.\csromanspacer{} พฺรหฺมโลเก ติฏฺฐตีติ, น เหวํ วตฺตพฺเพ ฯเปฯ.}

\nitthitam{สพฺพทิสากถา นิฏฺฐิตา.}
\chapterhead{21. เอกวีสติมวคฺค}
\vspace{\tipitakaparskip}
\csromancenter{\textbf{(206) 7.}\csromanspacer{}\textbf{ ธมฺมกถา}}
\proseitem{887}{\csromansymbolmark{*}สพฺเพ ธมฺมา นิยตาติ, อามนฺตา.\csromanspacer{} มิจฺฉตฺตนิยตาติ, น เหวํ วตฺตพฺเพ ฯเปฯ.\csromanspacer{} สมฺมตฺตนิยตาติ, น เหวํ วตฺตพฺเพ ฯเปฯ.\csromanspacer{} นตฺถิ อนิยโต ราสีติ, น เหวํ วตฺตพฺเพ ฯเปฯ.\csromanspacer{} นนุ อตฺถิ อนิยโต ราสีติ, อามนฺตา.\csromanspacer{} หญฺจิ อตฺถิ อนิยโต ราสิ, โน จ วต เร วตฺตพฺเพ “สพฺเพ ธมฺมา นิยตา”ติ.}

\prose{\csromansymbolmark{*}สพฺเพ ธมฺมา นิยตาติ, อามนฺตา.\csromanspacer{} นนุ ตโย ราสี วุตฺตา ภควตา\csromandash{} มิจฺฉตฺตนิยโต ราสิ, สมฺมตฺตนิยโต ราสิ, อนิยโต ราสีติ\footnote{อภิ 1. 3, 211 ปิฏฺเฐสุ.}, อามนฺตา.\csromanspacer{} หญฺจิ ตโย ราสี วุตฺตา ภควตา\csromandash{}มิจฺฉตฺตนิยโต ราสิ, สมฺมตฺตนิยโต ราสิ, อนิยโต ราสิ, โน จ วต เร วตฺตพฺเพ “สพฺเพ ธมฺมา นิยตา”ติ.}

\prose{\csromansymbolmark{*}รูปํ รูปฏฺเฐน นิยตนฺติ, อามนฺตา.\csromanspacer{} มิจฺฉตฺตนิยตนฺติ, น เหวํ วตฺตพฺเพ ฯเปฯ.\csromanspacer{} สมฺมตฺตนิยตนฺติ, น เหวํ วตฺตพฺเพ ฯเปฯ.\csromanspacer{} เวทนา ฯเปฯ.\csromanspacer{} สญฺญา ฯเปฯ.\csromanspacer{} สงฺขารา ฯเปฯ.\csromanspacer{} วิญฺญาณํ วิญฺญาณฏฺเฐน นิยตนฺติ, อามนฺตา.\csromanspacer{} มิจฺฉตฺตนิยตนฺติ, น เหวํ วตฺตพฺเพ ฯเปฯ.\csromanspacer{} สมฺมตฺตนิยตนฺติ, น เหวํ วตฺตพฺเพ ฯเปฯ.}

\proseitem{888}{\csromansymbolmark{+}น วตฺตพฺพํ “รูปํ รูปฏฺเฐน นิยตํ ฯเปฯ เวทนา ฯเปฯ สญฺญา ฯเปฯ สงฺขารา ฯเปฯ วิญฺญาณํ วิญฺญาณฏฺเฐน นิยตนฺ”ติ, อามนฺตา.\csromanspacer{} รูปํ เวทนา}

\vspace{\tipitakaparskip}
\csromancenter{เอกวีสติมวคฺค โหติ ฯเปฯ สญฺญา โหติ.\csromanspacer{} สงฺขารา โหนฺติ.\csromanspacer{} วิญฺญาณํ โหติ.\csromanspacer{} เวทนา ฯเปฯ.\csromanspacer{} สญฺญา ฯเปฯ.\csromanspacer{} สงฺขารา ฯเปฯ.\csromanspacer{} วิญฺญาณํ รูปํ โหติ ฯเปฯ เวทนา โหติ.\csromanspacer{} สญฺญา โหติ.\csromanspacer{} สงฺขารา โหนฺตีติ, น เหวํ วตฺตพฺเพ.\csromanspacer{} เตน หิ รูปํ รูปฏฺเฐน นิยตํ, เวทนา ฯเปฯ สญฺญา ฯเปฯ สงฺขารา ฯเปฯ วิญฺญาณํ วิญฺญาณฏฺเฐน นิยตนฺติ.}
\nitthitam{ธมฺมกถา นิฏฺฐิตา.}
\chapterhead{21. เอกวีสติมวคฺค}
\vspace{\tipitakaparskip}
\csromancenter{\textbf{(207) 8.}\csromanspacer{}\textbf{ กมฺมกถา}}
\csromanmatikaanchor{mtk.244}
\csromantocmarkref{subsection}{(207) 8. กมฺมกถา}{mtk.244}
\bookmark[dest={mtk.244},level=2]{(207) 8. กมฺมกถา}
\proseitem{889}{\csromanfolio{441}\csromansymbolmark{*}สพฺเพ กมฺมา นิยตาติ, อามนฺตา.\csromanspacer{} มิจฺฉตฺตนิยตาติ, น เหวํ วตฺตพฺเพ ฯเปฯ.\csromanspacer{} สมฺมตฺตนิยตาติ, น เหวํ วตฺตพฺเพ ฯเปฯ.\csromanspacer{} นตฺถิ อนิยโต ราสีติ, น เหวํ วตฺตพฺเพ ฯเปฯ.\csromanspacer{} นนุ อตฺถิ อนิยโต ราสีติ, อามนฺตา.\csromanspacer{} หญฺจิ อตฺถิ อนิยโต ราสิ, โน จ วต เร วตฺตพฺเพ “สพฺเพ กมฺมา นิยตา”ติ.}

\prose{\csromansymbolmark{*}สพฺเพกมฺมา นิยตาติ, อามนฺตา.\csromanspacer{} นนุ ตโย ราสี วุตฺตา ภควตา\csromandash{} มิจฺฉตฺตนิยโต ราสิ, สมฺมตฺตนิยโต ราสิ, อนิยโต ราสีติ, อามนฺตา.\csromanspacer{} หญฺจิ ตโย ราสี วุตฺตา ภควตา\csromandash{}มิจฺฉตฺตนิยโต ราสิ, สมฺมตฺตนิยโต ราสิ, อนิยโต ราสิ, โน จ วต เร วตฺตพฺเพ “สพฺเพ กมฺมา นิยตา”ติ.}

\proseitem{890}{\csromansymbolmark{*}ทิฏฺฐธมฺมเวทนียํ กมฺมํ ทิฏฺฐธมฺมเวทนียฏฺเฐน นิยตนฺติ, อามนฺตา.\csromanspacer{} มิจฺฉตฺตนิยตนฺติ, น เหวํ วตฺตพฺเพ ฯเปฯ.\csromanspacer{} สมฺมตฺตนิยตนฺติ, น เหวํ วตฺตพฺเพ ฯเปฯ.\csromanspacer{} อุปปชฺชเวทนียํ กมฺมํ ฯเปฯ.\csromanspacer{} อปราปริยเวทนียํ กมฺมํ อปราปริยเวทนียฏฺเฐน นิยตนฺติ, อามนฺตา.\csromanspacer{} มิจฺฉตฺตนิยตนฺติ, น เหวํ วตฺตพฺเพ ฯเปฯ.\csromanspacer{} สมฺมตฺตนิยตนฺติ, น เหวํ วตฺตพฺเพ ฯเปฯ.}

\csromanmatikaanchor{mtk.245}
\csromanmatikaanchor{mtk.246}
\csromantocmarknopage{chapter}{22. พาวีสติมวคฺค}{mtk.245}
\bookmark[dest={mtk.245},level=0]{22. พาวีสติมวคฺค}
\csromantocmarkref{subsection}{(208) 1. ปรินิพฺพานกถา}{mtk.246}
\bookmark[dest={mtk.246},level=2]{(208) 1. ปรินิพฺพานกถา}
\proseitem{891}{\csromansymbolmark{+}น วตฺตพฺพํ “ทิฏฺฐธมฺมเวทนียํ กมฺมํ ทิฏฺฐธมฺมเวทนียฏฺเฐน นิยตํ.\csromanspacer{} อุปปชฺชเวทนียํ กมฺมํ ฯเปฯ.\csromanspacer{} อปราปริยเวทนียํ กมฺมํ อปราปริยเวทนียฏฺเฐน นิยตนฺ”ติ, อามนฺตา.\csromanspacer{} ทิฏฺฐธมฺมเวทนียํ กมฺมํ อุปปชฺชเวทนียํ โหติ, อปราปริยเวทนียํ โหติ ฯเปฯ.\csromanspacer{} อุปปชฺชเวทนียํ กมฺมํ ทิฏฺฐธมฺมเวทนียํ โหติ, \csromanfolio{442}อปราปริยเวทนียํ โหติ ฯเปฯ.\csromanspacer{} อปราปริยเวทนียํ กมฺมํ ทิฏฺฐธมฺมเวทนียํ โหติ, อุปปชฺชเวทนียํ โหตีติ, น เหวํ วตฺตพฺเพ.\csromanspacer{} เตน หิ ทิฏฺฐธมฺมเวทนียํ กมฺมํ ทิฏฺฐธมฺมเวทนียฏฺเฐน นิยตํ.\csromanspacer{} อุปปชฺชเวทนียํ กมฺมํ ฯเปฯ.\csromanspacer{} อปราปริยเวทนียํ กมฺมํ อปราปริยเวทนียฏฺเฐน นิยตนฺติ.\csromanspacer{} กมฺมกถา นิฏฺฐิตา.\csromanspacer{} เอกวีสติโม วคฺโค.}

\titlehead{\textbf{ตสฺสุทฺทานํ}}
\prose{สาสนํ นวํ กตํ, อตฺถิ โกจิ ตถาคตสฺส สาสนํ นวํ กโรติ, ลพฺภา ตถาคตสฺส สาสนํ ปุน นวํ กาตุํ, ปุถุชฺชโน เตธาตุเกหิ กมฺเมหิ อวิวิตฺโต, อตฺถิ กิญฺจิ สํโยชนํ อปฺปหาย อรหตฺตปฺปตฺติ, อตฺถิ อธิปฺปายิทฺธิ พุทฺธานํ วา สาวกานํ วา, อตฺถิ พุทฺธานํ พุทฺเธหิ หีนาติเรกตา, สพฺพา ทิสา พุทฺธา ติฏฺฐนฺติ, สพฺเพ ธมฺมา นิยตา, สพฺเพ กมฺมา นิยตาติ.}

\chapterhead{22. พาวีสติมวคฺค}
\vspace{\tipitakaparskip}
\csromancenter{\textbf{(208) 1.}\csromanspacer{}\textbf{ ปรินิพฺพานกถา}}
\proseitem{892}{\csromansymbolmark{*}อตฺถิ กิญฺจิ สํโยชนํ อปฺปหาย ปรินิพฺพานนฺติ, อามนฺตา.\csromanspacer{} อตฺถิ กิญฺจิ สกฺกายทิฏฺฐิํ อปฺปหาย ฯเปฯ.\csromanspacer{} อโนตฺตปฺปํ อปฺปหาย ปรินิพฺพานนฺติ, น เหวํ วตฺตพฺเพ ฯเปฯ.}

\prose{\csromansymbolmark{*}อตฺถิ กิญฺจิ สํโยชนํ อปฺปหาย ปรินิพฺพานนฺติ, อามนฺตา.\csromanspacer{} อรหา สราโค ฯเปฯ สกิเลโสติ, น เหวํ วตฺตพฺเพ ฯเปฯ.\csromanspacer{} นนุ อรหา นิราโค ฯเปฯ นิกฺกิเลโสติ, อามนฺตา.\csromanspacer{} หญฺจิ อรหา นิราโค ฯเปฯ นิกฺกิเลโส, โน จ วต เร วตฺตพฺเพ “อตฺถิ กิญฺจิ สํโยชนํ อปฺปหาย ปรินิพฺพานนฺ”ติ.}

\proseitem{893}{\csromansymbolmark{+}น วตฺตพฺพํ “อตฺถิ กิญฺจิ สํโยชนํ อปฺปหาย ปรินิพฺพานนฺติ, อามนฺตา.\csromanspacer{} อรหา สพฺพํ พุทฺธวิสยํ ชานาตีติ, น เหวํ วตฺตพฺเพ.\csromanspacer{} เตน หิ อตฺถิ กิญฺจิ สํโยชนํ อปฺปหาย ปรินิพฺพานนฺติ.}

\nitthitam{ปรินิพฺพานกถา นิฏฺฐิตา.}
\chapterheadpageread{22. พาวีสติมวคฺค \textbf{22. พาวีสติมวคฺค}}
\vspace{\tipitakaparskip}
\csromancenter{\textbf{(209) 2.}\csromanspacer{}\textbf{ กุสลจิตฺตกถา}}
\csromanmatikaanchor{mtk.247}
\csromantocmarkref{subsection}{(209) 2. กุสลจิตฺตกถา}{mtk.247}
\bookmark[dest={mtk.247},level=2]{(209) 2. กุสลจิตฺตกถา}
\proseitem{894}{\csromanfolio{443}\csromansymbolmark{*}อรหา กุสลจิตฺโต ปรินิพฺพายตีติ, อามนฺตา.\csromanspacer{} อรหา ปุญฺญาภิสงฺขารํ อภิสงฺขโรนฺโต.\csromanspacer{} อาเนญฺชาภิสงฺขารํ อภิสงฺขโรนฺโต.\csromanspacer{} คติสํวตฺตนิยํ กมฺมํ กโรนฺโต.\csromanspacer{} ภวสํวตฺตนิยํ กมฺมํ กโรนฺโต.\csromanspacer{} อิสฺสริยสํวตฺตนิยํ กมฺมํ กโรนฺโต.\csromanspacer{} อธิปจฺจสํวตฺตนิยํ กมฺมํ กโรนฺโต.\csromanspacer{} มหาโภคสํวตฺตนิยํ กมฺมํ กโรนฺโต.\csromanspacer{} มหาปริวารสํวตฺตนิยํ กมฺมํ กโรนฺโต.\csromanspacer{} เทวโสภคฺยสํวตฺตนิยํ กมฺมํ กโรนฺโต.\csromanspacer{} มนุสฺสโสภคฺยสํวตฺตนิยํ กมฺมํ กโรนฺโต ปรินิพฺพายตีติ, น เหวํ วตฺตพฺเพ ฯเปฯ.}

\prose{\csromansymbolmark{*}อรหา กุสลจิตฺโต ปรินิพฺพายตีติ, อามนฺตา.\csromanspacer{} อรหา อาจินนฺโต อปจินนฺโต.\csromanspacer{} ปชหนฺโต อุปาทิยนฺโต.\csromanspacer{} วิสิเนนฺโต อุสฺสิเนนฺโต.\csromanspacer{} วิธูเปนฺโต สนฺธูเปนฺโต ปรินิพฺพายตีติ, น เหวํ วตฺตพฺเพ ฯเปฯ.\csromanspacer{} นนุ อรหา เนวาจินาติ น อปจินาติ อปจินิตฺวา ฐิโตติ, อามนฺตา.\csromanspacer{} หญฺจิ อรหา เนวาจินาติ, น อปจินาติ, อปจินิตฺวา ฐิโต, โน จ วต เร วตฺตพฺเพ “อรหา กุสลจิตฺโต ปรินิพฺพายตี”ติ.\csromanspacer{} นนุ อรหา เนว ปชหติ น อุปาทิยติ, ปชหิตฺวา ฐิโต.\csromanspacer{} เนว วิสิเนติ น อุสฺสิเนติ, วิสิเนตฺวา ฐิโต.\csromanspacer{} นนุ อรหา เนว วิธูเปติ น สนฺธูเปติ, วิธูเปตฺวา ฐิโตติ, อามนฺตา.\csromanspacer{} หญฺจิ อรหา เนว วิธูเปติ น สนฺธูเปติ, วิธูเปตฺวา ฐิโต, โน จ วต เร วตฺตพฺเพ “อรหา กุสลจิตฺโต ปรินิพฺพายตี”ติ.}

\proseitem{895}{\csromansymbolmark{+}น วตฺตพฺพํ “อรหา กุสลจิตฺโต ปรินิพฺพายตี”ติ, อามนฺตา.\csromanspacer{} นนุ อรหา อุปฏฺฐิตสฺสติ สโต สมฺปชาโน ปรินิพฺพายตีติ, อามนฺตา.\csromanspacer{} หญฺจิ อรหา อุปฏฺฐิตสฺสติ สโต สมฺปชาโน ปรินิพฺพายติ, เตน วต เร วตฺตพฺเพ “อรหา กุสลจิตฺโต ปรินิพฺพายตี”ติ.}

\nitthitam{กุสลจิตฺตกถา นิฏฺฐิตา.}
\chapterheadpageread{22. พาวีสติมวคฺค}
\vspace{\tipitakaparskip}
\csromancenter{\textbf{(210) 3.}\csromanspacer{}\textbf{ อาเนญฺชกถา}}
\csromanmatikaanchor{mtk.248}
\csromanmatikaanchor{mtk.249}
\csromantocmarkref{subsection}{(210) 3. อาเนญฺชกถา}{mtk.248}
\bookmark[dest={mtk.248},level=2]{(210) 3. อาเนญฺชกถา}
\csromantocmarkref{subsection}{(211) 4. ธมฺมาภิสมยกถา}{mtk.249}
\bookmark[dest={mtk.249},level=2]{(211) 4. ธมฺมาภิสมยกถา}
\proseitem{896}{\csromanfolio{444}\csromansymbolmark{*}อรหา อาเนญฺเช ฐิโต ปรินิพฺพายตีติ, อามนฺตา.\csromanspacer{} นนุ อรหา ปกติจิตฺเต ฐิโต ปรินิพฺพายตีติ, อามนฺตา.\csromanspacer{} หญฺจิ อรหา ปกติจิตฺเต ฐิโต ปรินิพฺพายติ, โน จ วต เร วตฺตพฺเพ “อรหา อาเนญฺเช ฐิโต ปรินิพฺพายตี”ติ.}

\prose{\csromansymbolmark{*}อรหา อาเนญฺเช ฐิโต ปรินิพฺพายตีติ, อามนฺตา.\csromanspacer{} อรหา กิริยมเย จิตฺเต ฐิโต ปรินิพฺพายตีติ, น เหวํ วตฺตพฺเพ ฯเปฯ.\csromanspacer{} นนุ อรหา วิปากจิตฺเต ฐิโต ปรินิพฺพายตีติ, อามนฺตา.\csromanspacer{} หญฺจิ อรหา วิปากจิตฺเต ฐิโต ปรินิพฺพายติ, โน จ วต เร วตฺตพฺเพ “อรหา อาเนญฺเช ฐิโต ปรินิพฺพายตี”ติ.}

\prose{\csromansymbolmark{*}อรหา อาเนญฺเช ฐิโต ปรินิพฺพายตีติ, อามนฺตา.\csromanspacer{} อรหา กิริยาพฺยากเต จิตฺเต ฐิโต ปรินิพฺพายตีติ, น เหวํ วตฺตพฺเพ ฯเปฯ.\csromanspacer{} นนุ อรหา วิปากาพฺยากเต จิตฺเต ฐิโต ปรินิพฺพายตีติ, อามนฺตา.\csromanspacer{} หญฺจิ อรหา วิปากาพฺยากเต จิตฺเต ฐิโต ปรินิพฺพายติ, โน จ วต เร วตฺตพฺเพ “อรหา อาเนญฺเช ฐิโต ปรินิพฺพายตี”ติ.}

\prose{\csromansymbolmark{*}อรหา อาเนญฺเช ฐิโต ปรินิพฺพายตีติ, อามนฺตา.\csromanspacer{} นนุ ภควา จตุตฺถชฺฌานา วุฏฺฐหิตฺวา สมนนฺตรา ปรินิพฺพุโตติ\footnote{ที 2. 128 ปิฏฺเฐ.}, อามนฺตา.\csromanspacer{} หญฺจิ ภควา จตุตฺถชฺฌานา วุฏฺฐหิตฺวา สมนนฺตรา ปรินิพฺพุโต, โน จ วต เร วตฺตพฺเพ “อรหา อาเนญฺเช ฐิโต ปรินิพฺพายตี”ติ.}

\nitthitam{อาเนญฺชกถา นิฏฺฐิตา.}
\chapterhead{22. พาวีสติมวคฺค}
\vspace{\tipitakaparskip}
\csromancenter{\textbf{(211) 4.}\csromanspacer{}\textbf{ ธมฺมาภิสมยกถา}}
\proseitem{897}{\csromansymbolmark{*}อตฺถิ คพฺภเสยฺยาย ธมฺมาภิสมโยติ, อามนฺตา.\csromanspacer{} อตฺถิ คพฺภเสยฺยาย ธมฺมเทสนา, ธมฺมสฺสวนํ, ธมฺมสากจฺฉา, ปริปุจฺฉา, สีลสมาทานํ, อินฺทฺริเยสุ คุตฺตทฺวารตา, โภชเน มตฺตญฺญุตา, ปุพฺพรตฺตาปรรตฺตํ}

\vspace{\tipitakaparskip}
\csromancenter{พาวีสติมวคฺค ชาคริยานุโยโคติ, น เหวํ วตฺตพฺเพ ฯเปฯ.\csromanspacer{} นตฺถิ คพฺภเสยฺยาย ธมฺมเทสนา, ธมฺมสฺสวนํ ฯเปฯ ปุพฺพรตฺตาปรรตฺตํ ชาคริยานุโยโคติ, อามนฺตา.\csromanspacer{} หญฺจิ นตฺถิ คพฺภเสยฺยาย ธมฺมเทสนา, ธมฺมสฺสวนํ ฯเปฯ.\csromanspacer{} ปุพฺพรตฺตาปรรตฺตํ ชาคริยานุโยโค, โน จ วต เร วตฺตพฺเพ “อตฺถิ คพฺภเสยฺยาย ธมฺมาภิสฺqมโย”ติ.}
\csromanmatikaanchor{mtk.250}
\csromantocmarkref{subsection}{(212-4) 5-7. ติสฺโสปิกถา}{mtk.250}
\bookmark[dest={mtk.250},level=2]{(212-4) 5-7. ติสฺโสปิกถา}
\prose{\csromanfolio{445}\csromansymbolmark{*}อตฺถิ คพฺภเสยฺยาย ธมฺมาภิสมโยติ, อามนฺตา.\csromanspacer{} นนุ เทฺว ปจฺจยา สมฺมาทิฏฺฐิยา อุปฺปาตาย ปรโต จ โฆโส, โยนิโส จ มนสิกาโรติ, อามนฺตา.\csromanspacer{} หญฺจิ เทฺว ปจฺจยา สมฺมาทิฏฺฐิยา อุปฺปาทาย ปรโต จ โฆโส, โยนิโส จ มนสิกาโร, โน จ วต เร วตฺตพฺเพ “อตฺถิ คพฺภเสยฺยาย ธมฺมาภิสมโย”ติ.}

\prose{\csromansymbolmark{*}อตฺถิ คพฺภเสยฺยาย ธมฺมาภิสมโยติ, อามนฺตา.\csromanspacer{} สุตฺตสฺส ปมตฺตสฺส มุฏฺฐสฺสติสฺส อสมฺปชานสฺส ธมฺมาภิสมโยติ, น เหวํ วตฺตพฺเพ ฯเปฯ.}

\nitthitam{ธมฺมาภิสมยกถา นิฏฺฐิตา.}
\chapterhead{22. พาวีสติมวคฺค}
\vspace{\tipitakaparskip}
\csromancenter{\textbf{(212-4) 5-7.}\csromanspacer{}\textbf{ ติสฺโสปิกถา}}
\proseitem{898}{\csromansymbolmark{*}อตฺถิ คพฺภเสยฺยาย อรหตฺตปฺปตฺตีติ, อามนฺตา.\csromanspacer{} สุตฺตสฺส ปมตฺตสฺส มุฏฺฐสฺสติสฺส อสมฺปชานสฺส อรหตฺตปฺปตฺตีติ, น เหวํ วตฺตพฺเพ ฯเปฯ.}

\proseitem{899}{\csromansymbolmark{*}อตฺถิ สุปินคตสฺส ธมฺมาภิสมโยติ, อามนฺตา.\csromanspacer{} สุตฺตสฺส ปมตฺตสฺส มุฏฺฐสฺสติสฺส อสมฺปชานสฺส ธมฺมาภิสมโยติ, น เหวํ วตฺตพฺเพ ฯเปฯ.}

\proseitem{900}{\csromansymbolmark{*}อตฺถิ สุปินคตสฺส อรหตฺตปฺปตฺตีติ, อามนฺตา.\csromanspacer{} สุตฺตสฺส ปมตฺตสฺส มุฏฺฐสฺสติสฺส อสมฺปชานสฺส อรหตฺตปฺปตฺตีติ, น เหวํ วตฺตพฺเพ ฯเปฯ.}

\nitthitam{ติสฺโสปิกถา นิฏฺฐิตา.}
\chapterheadpageread{22. พาวีสติมวคฺค}
\vspace{\tipitakaparskip}
\csromancenter{\textbf{(215) 8.}\csromanspacer{}\textbf{ อพฺยากตกถา}}
\csromanmatikaanchor{mtk.251}
\csromanmatikaanchor{mtk.252}
\csromantocmarkref{subsubsection}{(215) 8. อพฺยากตกถา}{mtk.251}
\bookmark[dest={mtk.251},level=3]{(215) 8. อพฺยากตกถา}
\csromantocmarkref{subsubsection}{(216) 9. อาเสวนปจฺจยกถา}{mtk.252}
\bookmark[dest={mtk.252},level=3]{(216) 9. อาเสวนปจฺจยกถา}
\proseitem{901}{\csromanfolio{446}\csromansymbolmark{*}สพฺพํ สุปินคตสฺส จิตฺตํ อพฺยากตนฺติ, อามนฺตา.\csromanspacer{} สุปินนฺเตน ปาณํ หเนยฺยาติ, อามนฺตา.\csromanspacer{} หญฺจิ สุปินนฺเตน ปาณํ หเนยฺย, โน จ วต เร วตฺตพฺเพ “สพฺพํ สุปินคตสฺส จิตฺตํ อพฺยากตนฺ”ติ.}

\prose{\csromansymbolmark{*}สุปินนฺเตน อทินฺนํ อาทิเยยฺย ฯเปฯ มุสา ภเณยฺย.\csromanspacer{} ปิสุณํ ภเณยฺย, ผรุสํ ภเณยฺย.\csromanspacer{} สมฺผํ ปลเปยฺย.\csromanspacer{} สนฺธิํ ฉินฺเทยฺย.\csromanspacer{} นิลฺโลปํ หเรยฺย.\csromanspacer{} เอกาคาริกํ กเรยฺย.\csromanspacer{} ปริปนฺเถ ติฏฺเฐยฺย.\csromanspacer{} ปรทารํ คจฺเฉยฺย.\csromanspacer{} คามฆาตกํ กเรยฺย.\csromanspacer{} นิคมฆาตกํ กเรยฺย.\csromanspacer{} สุปินนฺเตน เมถุนํ ธมฺมํ ปฏิเสเวยฺย.\csromanspacer{} สุปินคตสฺส อสุจิ มุจฺเจยฺย.\csromanspacer{} สุปินนฺเตน ทานํ ทเทยฺย.\csromanspacer{} จีวรํ ทเทยฺย.\csromanspacer{} ปิณฺฑปาตํ ทเทยฺย.\csromanspacer{} เสนาสนํ ทเทยฺย.\csromanspacer{} คิลานปจฺจยเภสชฺชปริกฺขารํ ทเทยฺย.\csromanspacer{} ขาทนียํ ทเทยฺย.\csromanspacer{} โภชนียํ ทเทยฺย.\csromanspacer{} ปานียํ ทเทยฺย.\csromanspacer{} เจติยํ วนฺเทยฺย, เจติเย มาลํ อาโรเปยฺย.\csromanspacer{} คนฺธํ อาโรเปยฺย.\csromanspacer{} วิเลปนํ อาโรเปยฺย ฯเปฯ เจติยํ อภิทกฺขิณํ กเรยฺยาติ, อามนฺตา.\csromanspacer{} หญฺจิ สุปินนฺเตน เจติยํ อภิทกฺขิณํ กเรยฺย, โน จ วต เร วตฺตพฺเพ “สพฺพํ สุปินคตสฺส จิตฺตํ อพฺยากตนฺ”ติ.}

\proseitem{902}{\csromansymbolmark{+}น วตฺตพฺพํ “สพฺพํ สุปินคตสฺส จิตฺตํ อพฺยากตนฺ”ติ, อามนฺตา.\csromanspacer{} นนุ สุปินคตสฺส จิตฺตํ อพฺโพหาริยํ วุตฺตํ ภควตาติ, อามนฺตา.\csromanspacer{} หญฺจิ สุปินคตสฺส จิตฺตํ อพฺโพหาริยํ วุตฺตํ ภควตา, เตน วต เร วตฺตพฺเพ “สพฺพํ สุปินคตสฺส จิตฺตํ อพฺยากตนฺ”ติ.}

\nitthitam{อพฺยากตกถา นิฏฺฐิตา.}
\chapterhead{22. พาวีสติมวคฺค}
\vspace{\tipitakaparskip}
\csromancenter{\textbf{(216) 9.}\csromanspacer{}\textbf{ อาเสวนปจฺจยกถา}}
\proseitem{903}{\csromansymbolmark{*}นตฺถิ กาจิ อาเสวนปจฺจยตาติ, อามนฺตา.\csromanspacer{} นนุ วุตฺตํ ภควตา “ปาณาติปาโต ภิกฺขเว อาเสวิโต ภาวิโต}

\vspace{\tipitakaparskip}
\csromancenter{พาวีสติมวคฺค พหุลีกโต นิรยสํวตฺตนิโก ติรจฺฉานโยนิสํวตฺตนิโก เปตฺติวิสยสํวตฺตนิโก, โย สพฺพลหุโส ปาณาติปาตสฺส วิปาโก มนุสฺสภูตสฺส อปฺปายุกสํวตฺตนิโก โหตี”ติ\footnote{อํ 3. 77 ปิฏฺเฐ.} อตฺเถว สุตฺตนฺโตติ, อามนฺตา.\csromanspacer{} เตน หิ อตฺถิ กาจิ อาเสวนปจฺจยตาติ.}
\prose{\csromanfolio{447}\csromansymbolmark{*}นตฺถิ กาจิ อาเสวนปจฺจยตาติ, อามนฺตา.\csromanspacer{} นนุ วุตฺตํ ภควตา “อทินฺนาทานํ ภิกฺขเว อาเสวิตํ ภาวิตํ พหุลีกตํ นิรยสํวตฺตนิกํ ติรจฺฉานโยนิสํวตฺตนิกํ เปตฺติวิสยสํวตฺตนิกํ, โย สพฺพลหุโส อทินฺนาทานสฺส วิปาโก มนุสฺสภูตสฺส โภคพฺยสนสํวตฺตนิโก โหติ ฯเปฯ โย สพฺพลหุโส กาเมสุมิจฺฉาจารสฺส วิปาโก มนุสฺสภูตสฺส สปตฺตเวรสํวตฺตนิโก โหติ ฯเปฯ โย สพฺพลหุโส มุสาวาทสฺส วิปาโก มนุสฺสภูตสฺส อพฺภูตพฺภกฺขานสํวตฺตนิโก โหติ ฯเปฯ โย สพฺพลหุโส ปิสุณาย วาจาย วิปาโก มนุสฺสภูตสฺส มิตฺเตหิ เภทนสํวตฺตนิโก โหติ ฯเปฯ โย สพฺพลหุโส ผรุสาย วาจาย วิปาโก มนุสฺสภูตสฺส อมนาปสทฺทสํวตฺตนิโก โหติ ฯเปฯ โย สพฺพลหุโส สมฺผปฺปลาปสฺส วิปาโก มนุสฺสภูตสฺส อนาเทยฺยวาจาสํวตฺตนิโก โหติ.\csromanspacer{} สุราเมรยปานํ ภิกฺขเว อาเสวิตํ ฯเปฯ โย สพฺพลหุโส สุราเมรยปานสฺส วิปาโก มนุสฺสภูตสฺส อุมฺมตฺตกสํวตฺตนิโก โหตี”ติ\footnote{อํ 3. 77 ปิฏฺเฐ.} อตฺเถว สุตฺตนฺโตติ, อามนฺตา.\csromanspacer{} เตน หิ อตฺถิ กาจิ อาเสวนปจฺจยตาติ.}

\proseitem{904}{\csromansymbolmark{*}นตฺถิ กาจิ อาเสวนปจฺจยตาติ, อามนฺตา.\csromanspacer{} นนุ วุตฺตํ ภควตา “มิจฺฉาทิฏฺฐิ ภิกฺขเว อาเสวิตา ภาวิตา พหุลีกตา นิรยสํวตฺตนิกา ติรจฺฉานโยนิสํวตฺตนิกา เปตฺติวิสยสํวตฺตนิกา”ติ อตฺเถว สุตฺตนฺโตติ, อามนฺตา.\csromanspacer{} เตน หิ อตฺถิ กาจิ อาเสวนปจฺจยตาติ.}

\prose{\csromansymbolmark{*}นตฺถิ กาจิ อาเสวนปจฺจยตาติ, อามนฺตา.\csromanspacer{} นนุ วุตฺตํ ภควตา “มิจฺฉาสงฺกปฺโป ฯเปฯ มิจฺฉาสมาธิ ภิกฺขเว อาเสวิโต ภาวิโต ฯเปฯ เปตฺติวิสยสํวตฺตนิโก”ติ อตฺเถว สุตฺตนฺโตติ, อามนฺตา.\csromanspacer{} เตน หิ อตฺถิ กาจิ อาเสวนปจฺจยตาติ.}

\csromanmatikaanchor{mtk.253}
\csromantocmarkref{subsubsection}{(217) 10. ขณิกกถา}{mtk.253}
\bookmark[dest={mtk.253},level=3]{(217) 10. ขณิกกถา}
\proseitem{905}{\csromanfolio{448}\csromansymbolmark{*}นตฺถิ กาจิ อาเสวนปจฺจยตาติ, อามนฺตา.\csromanspacer{} นนุ วุตฺตํ ภควตา “สมฺมาทิฏฺฐิ ภิกฺขเว อาเสวิตา ภาวิตา พหุลีกตา อมโตคธา โหติ อมตปรายนา อมตปริโยสานา”ติ อตฺเถว สุตฺตนฺโตติ, อามนฺตา.\csromanspacer{} เตน หิ อตฺถิ กาจิ อาเสวนปจฺจยตาติ.}

\prose{\csromansymbolmark{*}นตฺถิ กาจิ อาเสวนปจฺจยตาติ, อามนฺตา.\csromanspacer{} นนุ วุตฺตํ ภควตา “สมฺมาสงฺกปฺโป ภิกฺขเว อาเสวิโต ภาวิโต พหุลีกโต ฯเปฯ.\csromanspacer{} สมฺมาสมาธิ ภิกฺขเว อาเสวิโต ภาวิโต พหุลีกโต อมโตคโธ โหติ อมตปรายโน อมตปริโยสาโน”ติ อตฺเถว สุตฺตนฺโตติ, อามนฺตา.\csromanspacer{} เตน หิ อตฺถิ กาจิ อาเสวนปจฺจยตาติ.}

\nitthitam{อาเสวนปจฺจยกถา นิฏฺฐิตา.}
\chapterhead{22. พาวีสติมวคฺค}
\vspace{\tipitakaparskip}
\csromancenter{\textbf{(217) 10.}\csromanspacer{}\textbf{ ขณิกกถา}}
\proseitem{906}{\csromansymbolmark{*}เอกจิตฺตกฺขณิกา สพฺเพ ธมฺมาติ, อามนฺตา.\csromanspacer{} จิตฺเต มหาปถวี สณฺฐาติ.\csromanspacer{} มหาสมุทฺโท สณฺฐาติ.\csromanspacer{} สิเนรุปพฺพตราชา สณฺฐาติ.\csromanspacer{} อาโป สณฺฐาติ.\csromanspacer{} เตโช สณฺฐาติ.\csromanspacer{} วาโย สณฺฐาติ.\csromanspacer{} ติณกฏฺฐวนปฺปตโย สณฺฐหนฺตีติ, น เหวํ วตฺตพฺเพ ฯเปฯ.}

\prose{\csromansymbolmark{*}เอกจิตฺตกฺขณิกา สพฺเพ ธมฺมาติ, อามนฺตา.\csromanspacer{} จกฺขายตนํ จกฺขุวิญฺญาเณน สหชาตนฺติ, น เหวํ วตฺตพฺเพ ฯเปฯ.\csromanspacer{} จกฺขายตนํ จกฺขุวิญฺญาเณน สหชาตนฺติ, อามนฺตา.\csromanspacer{} นนุ อายสฺมา สาริปุตฺโต เอตทโวจ “อชฺฌตฺติกญฺเจว อาวุโส จกฺขุํ อปริภินฺนํ โหติ, พาหิรา จ รูปา น อาปาถํ อาคจฺฉนฺติ, โน จ ตชฺโช สมนฺนาหาโร โหติ, เนว ตาว ตชฺชสฺส วิญฺญาณภาคสฺส\footnote{วิญฺญาณภาวสฺส (พหูสุ)} ปาตุภาโว โหติ.\csromanspacer{} อชฺฌตฺติกญฺเจว อาวุโส จกฺขุํ อปริภินฺนํ โหติ, พาหิรา จ รูปา อาปาถํ อาคจฺฉนฺติ, โน จ ตชฺโช สมนฺนาหาโร โหติ, เนว ตาว ตชฺชสฺส วิญฺญาณภาคสฺส ปาตุภาโว โหติ.\csromanspacer{} ยโต จ โข อาวุโส อชฺฌตฺติกญฺเจว จกฺขุํ อปริภินฺนํ}

\vspace{\tipitakaparskip}
\csromancenter{พาวีสติมวคฺค โหติ, พาหิรา จ รูปา อาปาถํ อาคจฺฉนฺติ, ตชฺโช จ สมนฺนาหาโร โหติ, เอวํ ตชฺชสฺส วิญฺญาณภาคสฺส ปาตุภาโว โหตี”ติ\footnote{ม 1. 248 ปิฏฺเฐ.} อตฺเถว สุตฺตนฺโตติ, อามนฺตา.\csromanspacer{} เตน หิ น วตฺตพฺพํ “จกฺขายตนํ จกฺขุวิญฺญาเณน สหชาตนฺ”ติ.}
\prose{\csromanfolio{449}\csromansymbolmark{*}โสตายตนํ ฯเปฯ.\csromanspacer{} ฆานายตนํ ฯเปฯ.\csromanspacer{} ชิวฺหายตนํ ฯเปฯ.\csromanspacer{} กายายตนํ กายวิญฺญาเณน สหชาตนฺติ, น เหวํ วตฺตพฺเพ ฯเปฯ.\csromanspacer{} กายายตนํ กายวิญฺญาเณน สหชาตนฺติ, อามนฺตา.\csromanspacer{} นนุ อายสฺมา สาริปุตฺโต เอตทโวจ “อชฺฌตฺติโก เจว อาวุโส กาโย อปริภินฺโน โหติ, พาหิรา จ โผฏฺฐพฺพา น อาปาถํ อาคจฺฉนฺติ, โน จ ฯเปฯ.\csromanspacer{} อชฺฌตฺติโก เจว อาวุโส กาโย อปริภินฺโน โหติ, พาหิรา จ โผฏฺฐพฺพา อาปาถํ อาคจฺฉนฺติ, โน จ ฯเปฯ.\csromanspacer{} ยโต จ โข อาวุโส อชฺฌตฺติโก เจว กาโย อปริภินฺโน โหติ, พาหิรา จ โผฏฺฐพฺพา อาปาถํ อาคจฺฉนฺติ, ตชฺโช จ สมนฺนาหาโร โหติ, เอวํ ตชฺชสฺส วิญฺญาณภาคสฺส ปาตุภาโว โหตี”ติ\footnote{ม 1. 249 ปิฏฺเฐ.} อตฺเถว สุตฺตนฺโตติ, อามนฺตา.\csromanspacer{} เตน หิ น วตฺตพฺพํ “กายายตนํ กายวิญฺญาเณน สหชาตนฺ”ติ.}

\proseitem{907}{\csromansymbolmark{+}น วตฺตพฺพํ “เอกจิตฺตกฺขณิกา สพฺเพ ธมฺมา”ติ, อามนฺตา.\csromanspacer{} สพฺเพ ธมฺมา นิจฺจา ธุวา สสฺสตา อวิปริณามธมฺมาติ, น เหวํ วตฺตพฺเพ.\csromanspacer{} เตน หิ เอกจิตฺตกฺขณิกา สพฺเพ ธมฺมาติ.\csromanspacer{} ขณิกกถา นิฏฺฐิตา.\csromanspacer{} พาวีสติโม วคฺโค.}

\titlehead{\textbf{ตสฺสุทฺทานํ}}
\prose{อตฺถิ กิญฺจิ สํโยชนํ อปฺปหาย ปรินิพฺพานํ, อรหา กุสลจิตฺโต ปรินิพฺพายติ, อรหา อาเนญฺเช ฐิโต ปรินิพฺพายติ, อตฺถิ คพฺภเสยฺยาย ธมฺมาภิสมโย, อตฺถิ คพฺภเสยฺยาย อรหตฺตปฺปตฺติ, อตฺถิ สุปินคตสฺส ธมฺมาภิสมโย, อตฺถิ สุปินคตสฺส อรหตฺตปฺปตฺติ, สพฺพํ สุปินคตสฺส จิตฺตํ อพฺยากตํ, นตฺถิ กาจิ อาเสวนปจฺจยตา, เอกจิตฺตกฺขณิกา สพฺเพ ธมฺมาติ.}

\csromanmatikaanchor{mtk.254}
\csromantocmarknopage{chapter}{23. เตวีสติมวคฺค}{mtk.254}
\bookmark[dest={mtk.254},level=0]{23. เตวีสติมวคฺค}
\chapterheadpageread{23. เตวีสติมวคฺค}
\vspace{\tipitakaparskip}
\csromancenter{\textbf{(218) 1.}\csromanspacer{}\textbf{ เอกาธิปฺปายกถา}}
\csromanmatikaanchor{mtk.255}
\csromanmatikaanchor{mtk.256}
\csromanmatikaanchor{mtk.257}
\csromantocmarkref{subsection}{(218) 1. เอกาธิปฺปายกถา}{mtk.255}
\bookmark[dest={mtk.255},level=2]{(218) 1. เอกาธิปฺปายกถา}
\csromantocmarkref{subsection}{(219) 2. อรหนฺตวณฺณกถา}{mtk.256}
\bookmark[dest={mtk.256},level=2]{(219) 2. อรหนฺตวณฺณกถา}
\csromantocmarkref{subsection}{(220-4) 3-7. อิสฺสริยกามการิกากถา}{mtk.257}
\bookmark[dest={mtk.257},level=2]{(220-4) 3-7. อิสฺสริยกามการิกากถา}
\proseitem{908}{\csromanfolio{450}\csromansymbolmark{*}เอกาธิปฺปาเยน เมถุโน ธมฺโม ปฏิเสวิตพฺโพติ, อามนฺตา.\csromanspacer{} เอกาธิปฺปาเยน อสฺสมเณน โหตพฺพํ อภิกฺขุนา โหตพฺพํ ฉินฺนมูเลน โหตพฺพํ ปาราชิเกน โหตพฺพนฺติ, น เหวํ วตฺตพฺเพ ฯเปฯ.\csromanspacer{} เอกาธิปฺปาเยน เมถุโน ธมฺโม ปฏิเสวิตพฺโพติ, อามนฺตา.\csromanspacer{} เอกาธิปฺปาเยน ปาโณ หนฺตพฺโพ.\csromanspacer{} อทินฺนํ อาทิยิตพฺพํ.\csromanspacer{} มุสา ภณิตพฺพา.\csromanspacer{} ปิสุณํ ภณิตพฺพํ.\csromanspacer{} ผรุสํ ภณิตพฺพํ.\csromanspacer{} สมฺผํ ปลปิตพฺพํ.\csromanspacer{} สนฺธิ เฉทิตพฺโพ.\csromanspacer{} นิลฺโลปํ หาตพฺพํ.\csromanspacer{} เอกาคาริกํ กาตพฺพํ.\csromanspacer{} ปริปนฺเถ ฐาตพฺพํ.\csromanspacer{} ปรทาโร คนฺตพฺโพ.\csromanspacer{} คามฆาตโก กาตพฺโพ.\csromanspacer{} นิคมฆาตโก กาตพฺโพติ, น เหวํ วตฺตพฺเพ ฯเปฯ.}

\nitthitam{เอกาธิปฺปายกถา นิฏฺฐิตา.}
\chapterhead{23. เตวีสติมวคฺค}
\vspace{\tipitakaparskip}
\csromancenter{\textbf{(219) 2.}\csromanspacer{}\textbf{ อรหนฺตวณฺณกถา}}
\proseitem{909}{\csromansymbolmark{*}อรหนฺตานํ วณฺเณน อมนุสฺสา เมถุนํ ธมฺมํ ปฏิเสวนฺตีติ, อามนฺตา.\csromanspacer{} อรหนฺตานํ วณฺเณน อมนุสฺสา ปาณํ หนนฺติ ฯเปฯ อทินฺนํ อาทิยนฺติ.\csromanspacer{} มุสา ภณนฺติ.\csromanspacer{} ปิสุณํ ภณนฺติ.\csromanspacer{} ผรุสํ ภณนฺติ.\csromanspacer{} สมฺผํ ปลปนฺติ.\csromanspacer{} สนฺธิํ ฉินฺทนฺติ.\csromanspacer{} นิลฺโลปํ หรนฺติ.\csromanspacer{} เอกาคาริกํ กโรนฺติ.\csromanspacer{} ปริปนฺเถ ติฏฺฐนฺติ.\csromanspacer{} ปรทารํ คจฺฉนฺติ.\csromanspacer{} คามฆาตกํ กโรนฺติ ฯเปฯ นิคมฆาตกํ กโรนฺตีติ, น เหวํ วตฺตพฺเพ ฯเปฯ.}

\nitthitam{อรหนฺตวณฺณกถา นิฏฺฐิตา.}
\chapterhead{23. เตวีสติมวคฺค}
\vspace{\tipitakaparskip}
\csromancenter{\textbf{(220-4) 3-7.}\csromanspacer{}\textbf{ อิสฺสริยกามการิกากถา}}
\proseitem{910}{\csromansymbolmark{*}โพธิสตฺโต อิสฺสริยกามการิกาเหตุ วินิปาตํ คจฺฉตีติ, อามนฺตา.\csromanspacer{} โพธิสตฺโต อิสฺสริยกามการิกาเหตุ นิรยํ}

\vspace{\tipitakaparskip}
\csromancenter{เตวีสติมวคฺค คจฺฉติ.\csromanspacer{} สญฺชีวํ คจฺฉติ.\csromanspacer{} กาลสุตฺตํ คจฺฉติ.\csromanspacer{} ตาปนํ คจฺฉติ.\csromanspacer{} มหาตาปนํ\footnote{ปตาปนํ (สี, สฺยา, กํ, อิ)} คจฺฉติ.\csromanspacer{} สํฆาตกํ คจฺฉติ.\csromanspacer{} โรรุวํ คจฺฉติ ฯเปฯ อวีจิํ คจฺฉตีติ, น เหวํ วตฺตพฺเพ ฯเปฯ.}
\prose{\csromanfolio{451}\csromansymbolmark{*}โพธิสตฺโต อิสฺสริยกามการิกาเหตุ วินิปาตํ คจฺฉตีติ, อามนฺตา. “โพธิสตฺโต อิสฺสริยกามการิกาเหตุ วินิปาตํ คจฺฉตี”ติ อตฺเถว สุตฺตนฺโตติ, นตฺถิ.\csromanspacer{} หญฺจิ “โพธิสตฺโต อิสฺสริยกามการิกาเหตุ วินิปาตํ คจฺฉตี”ติ นตฺเถว สุตฺตนฺโต, โน จ วต เร วตฺตพฺเพ “โพธิสตฺโต อิสฺสริยกามการิกาเหตุ วินิปาตํ คจฺฉตี”ติ.}

\proseitem{911}{\csromansymbolmark{*}โพธิสตฺโต อิสฺสริยกามการิกาเหตุ คพฺภเสยฺยํ โอกฺกมตีติ, อามนฺตา.\csromanspacer{} โพธิสตฺโต อิสฺสริยกามการิกาเหตุ นิรยํ อุปปชฺเชยฺย.\csromanspacer{} ติรจฺฉานโยนิํ อุปปชฺเชยฺยาติ, น เหวํ วตฺตพฺเพ ฯเปฯ.}

\prose{\csromansymbolmark{*}โพธิสตฺโต อิสฺสริยกามการิกาเหตุ คพฺภเสยฺยํ โอกฺกมตีติ, อามนฺตา.\csromanspacer{} โพธิสตฺโต อิทฺธิมาติ, น เหวํ วตฺตพฺเพ ฯเปฯ.\csromanspacer{} โพธิสตฺโต อิทฺธิมาติ, อามนฺตา.\csromanspacer{} โพธิสตฺเตน ฉนฺทิทฺธิปาโท ภาวิโต ฯเปฯ วีริยิทฺธิปาโท ฯเปฯ จิตฺติทฺธิปาโท ฯเปฯ วีมํสิทฺธิปาโท ภาวิโตติ, น เหวํ วตฺตพฺเพ ฯเปฯ.}

\prose{\csromansymbolmark{*}โพธิสตฺโต อิสฺสริยกามการิกาเหตุ คพฺภเสยฺยํ โอกฺกมตีติ, อามนฺตา. “โพธิสตฺโต อิสฺสริยกามการิกาเหตุ คพฺภเสยฺยํ โอกฺกมตี”ติ อตฺเถว สุตฺตนฺโตติ, นตฺถิ.\csromanspacer{} หญฺจิ “โพธิสตฺโต อิสฺสริยกามการิกาเหตุ คพฺภเสยฺยํ โอกฺกมตี”ติ นตฺเถว สุตฺตนฺโต, โน จ วต เร วตฺตพฺเพ “โพธิสตฺโต อิสฺสริยกามการิกาเหตุ คพฺภเสยฺยํ โอกฺกมตี”ติ.}

\csromanmatikaanchor{mtk.258}
\csromantocmarkref{subsubsection}{(225) 8. ปติรูปกถา}{mtk.258}
\bookmark[dest={mtk.258},level=3]{(225) 8. ปติรูปกถา}
\proseitem{912}{\csromansymbolmark{*}โพธิสตฺโต อิสฺสริยกามการิกาเหตุ ทุกฺกรการิกํ อกาสีติ, อามนฺตา.\csromanspacer{} โพธิสตฺโต อิสฺสริยกามการิกาเหตุ “สสฺสโต โลโก”ติ ปจฺจาคจฺฉิ. “อสสฺสโต โลโก”ติ ฯเปฯ “อนฺตวา โลโก”ติ ฯเปฯ “อนนฺตวา โลโก”ติ, “ตํ ชีวํ ตํ \csromanfolio{452}สรีรนฺ”ติ, “อญฺญํ ชีวํ อญฺญํ สรีรนฺ”ติ, “โหติ ตถาคโต ปรํ มรณา”ติ. “น โหติ ตถาคโต ปรํ มรณา”ติ, “โหติ จ น จ โหติ ตถาคโต ปรํ มรณา”ติ ฯเปฯ “เนว โหติ น น โหติ ตถาคโต ปรํ มรณา”ติ ปจฺจาคจฺฉีติ, น เหวํ วตฺตพฺเพ ฯเปฯ.}

\prose{\csromansymbolmark{*}โพธิสตฺโต อิสฺสริยกามการิกาเหตุ ทุกฺกรการิกํ อกาสีติ, อามนฺตา. “โพธิสตฺโต อิสฺสริยกามการิกาเหตุ ทุกฺกรการิกํ อกาสี”ติ อตฺเถว สุตฺตนฺโตติ, นตฺถิ.\csromanspacer{} หญฺจิ “โพธิสตฺโต อิสฺสริยกามการิกาเหตุ ทุกฺกรการิกํ อกาสี”ติ นตฺเถว สุตฺตนฺโต, โน จ วต เร วตฺตพฺเพ “โพธิสตฺโต อิสฺสริยกามการิกาเหตุ ทุกฺกรการิกํ อกาสี”ติ.}

\proseitem{913}{\csromansymbolmark{*}โพธิสตฺโต อิสฺสริยกามการิกาเหตุ อปรนฺตปํ\footnote{อมรํ ตปํ (สํ 1. 104 ปิฏฺเฐ.)} อกาสิ, อญฺญํ สตฺถารํ อุทฺทิสีติ, อามนฺตา.\csromanspacer{} โพธิสตฺโต อิสฺสริยกามการิกาเหตุ “สสฺสโต โลโก”ติ ปจฺจาคจฺฉิ ฯเปฯ “เนว โหติ น น โหติ ตถาคโต ปรํ มรณา”ติ ปจฺจาคจฺฉีติ, น เหวํ วตฺตพฺเพ ฯเปฯ.}

\proseitem{914}{\csromansymbolmark{*}โพธิสตฺโต อิสฺสริยกามการิกาเหตุ อญฺญํ สตฺถารํ อุทฺทิสีติ, อามนฺตา. “โพธิสตฺโต อิสฺสริยกามการิกาเหตุ อญฺญํ สตฺถารํ อุทฺทิสี”ติ อตฺเถว สุตฺตนฺโตติ, นตฺถิ.\csromanspacer{} หญฺจิ “โพธิสตฺโต อิสฺสริยกามการิกาเหตุ อญฺญํ สตฺถารํ อุทฺทิสี”ติ นตฺเถว สุตฺตนฺโต, โน จ วต เร วตฺตพฺเพ “โพธิสตฺโต อิสฺสริยกามการิกาเหตุ อญฺญํ สตฺถารํ อุทฺทิสี”ติ.}

\nitthitam{อิสฺสริยกามการิกากถา นิฏฺฐิตา.}
\chapterhead{23. เตวีสติมวคฺค}
\vspace{\tipitakaparskip}
\csromancenter{\textbf{(225) 8.}\csromanspacer{}\textbf{ ปติรูปกถา}}
\proseitem{915}{\csromansymbolmark{*}อตฺถิ น ราโค ราคปติรูปโกติ, อามนฺตา.\csromanspacer{} อตฺถิ น ผสฺโส ผสฺสปติรูปโก.\csromanspacer{} อตฺถิ น เวทนา เวทนาปติรูปิกา.\csromanspacer{} อตฺถิ น สญฺญา}

\vspace{\tipitakaparskip}
\csromancenter{เตวีสติมวคฺค สญฺญาปติรูปิกา.\csromanspacer{} อตฺถิ น เจตนา เจตนาปติรูปิกา.\csromanspacer{} อตฺถิ น จิตฺตํ จิตฺตปติรูปกํ.\csromanspacer{} อตฺถิ น สทฺธา สทฺธาปติรูปิกา.\csromanspacer{} อตฺถิ น วีริยํ วีริยปติรูปกํ.\csromanspacer{} อตฺถิ น สติ สติปติรูปิกา.\csromanspacer{} อตฺถิ น สมาธิ สมาธิปติรูปโก.\csromanspacer{} อตฺถิ น ปญฺญา ปญฺญาปติรูปิกาติ, น เหวํ วตฺตพฺเพ ฯเปฯ.}
\csromanmatikaanchor{mtk.259}
\csromantocmarkref{subsubsection}{(226) 9. อปรินิปฺผนฺนกถา}{mtk.259}
\bookmark[dest={mtk.259},level=3]{(226) 9. อปรินิปฺผนฺนกถา}
\proseitem{916}{\csromanfolio{453}\csromansymbolmark{*}อตฺถิ น โทโส โทสปติรูปโก.\csromanspacer{} อตฺถิ น โมโห โมหปติรูปโก.\csromanspacer{} อตฺถิ น กิเลโส กิเลสปติรูปโกติ, อามนฺตา.\csromanspacer{} อตฺถิ น ผสฺโส ผสฺสปติรูปโก ฯเปฯ.\csromanspacer{} อตฺถิ น ปญฺญา ปญฺญาปติรูปิกาติ, น เหวํ วตฺตพฺเพ ฯเปฯ.}

\nitthitam{ปติรูปกถา นิฏฺฐิตา.}
\chapterhead{23. เตวีสติมวคฺค}
\vspace{\tipitakaparskip}
\csromancenter{\textbf{(226) 9.}\csromanspacer{}\textbf{ อปรินิปฺผนฺนกถา}}
\proseitem{917}{\csromansymbolmark{*}รูปํ อปรินิปฺผนฺนนฺติ, อามนฺตา.\csromanspacer{} รูปํ น อนิจฺจํ น สงฺขตํ น ปฏิจฺจสมุปฺปนฺนํ น ขยธมฺมํ น วยธมฺมํ น วิราคธมฺมํ น นิโรธธมฺมํ น วิปริณามธมฺมนฺติ, น เหวํ วตฺตพฺเพ ฯเปฯ.\csromanspacer{} นนุ รูปํ อนิจฺจํ สงฺขตํ ปฏิจฺจสมุปฺปนฺนํ ขยธมฺมํ วยธมฺมํ วิราคธมฺมํ นิโรธธมฺมํ วิปริณามธมฺมนฺติ, อามนฺตา.\csromanspacer{} หญฺจิ รูปํ อนิจฺจํ สงฺขตํ ฯเปฯ วิปริณามธมฺมํ, โน จ วต เร วตฺตพฺเพ “รูปํ อปรินิปฺผนฺนนฺ”ติ.}

\prose{\csromansymbolmark{*}ทุกฺขญฺเญว ปรินิปฺผนฺนนฺติ, อามนฺตา.\csromanspacer{} นนุ ‘ยทนิจฺจํ ตํ ทุกฺขํ’\footnote{สํ 2. 19 ปิฏฺเฐ.} วุตฺตํ ภควตา รูปํ อนิจฺจนฺติ, อามนฺตา.\csromanspacer{} หญฺจิ ‘ยทนิจฺจํ ตํ ทุกฺขํ’ วุตฺตํ ภควตา รูปํ อนิจฺจํ, โน จ วต เร วตฺตพฺเพ “ทุกฺขญฺเญว ปรินิปฺผนฺนนฺ”ติ ฯเปฯ.}

\proseitem{918}{\csromansymbolmark{*}เวทนา ฯเปฯ.\csromanspacer{} สญฺญา.\csromanspacer{} สงฺขารา.\csromanspacer{} วิญฺญาณํ.\csromanspacer{} จกฺขายตนํ ฯเปฯ.\csromanspacer{} ธมฺมายตนํ.\csromanspacer{} จกฺขุธาตุ ฯเปฯ.\csromanspacer{} ธมฺมธาตุ.\csromanspacer{} จกฺขุนฺทฺริยํ ฯเปฯ.\csromanspacer{} อญฺญาตาวินฺทฺริยํ อปรินิปฺผนฺนนฺติ, อามนฺตา.\csromanspacer{} อญฺญาตาวินฺทฺริยํ น อนิจฺจํ ฯเปฯ.\csromanspacer{} น วิปริณามธมฺมนฺติ, น เหวํ วตฺตพฺเพ ฯเปฯ.\csromanspacer{} นนุ อญฺญาตาวินฺทฺริยํ อนิจฺจํ สงฺขตํ ฯเปฯ วิปริณามธมฺมนฺติ, อามนฺตา.\csromanspacer{} หญฺจิ อญฺญาตาวินฺทฺริยํ อนิจฺจํ สงฺขตํ ปฏิจฺจสมุปฺปนฺนํ ขยธมฺมํ วยธมฺมํ วิราคธมฺมํ นิโรธธมฺมํ วิปริณามธมฺมํ, โน จ วต เร วตฺตพฺเพ “อญฺญาตาวินฺทฺริยํ อปรินิปฺผนฺนนฺ”ติ.}

\prose{\csromanfolio{454}\csromansymbolmark{*}ทุกฺขญฺเญว ปรินิปฺผนฺนนฺติ, อามนฺตา.\csromanspacer{} นนุ ‘ยทนิจฺจํ ตํ ทุกฺขํ’ วุตฺตํ ภควตา อญฺญาตาวินฺทฺริยํ อนิจฺจนฺติ, อามนฺตา.\csromanspacer{} หญฺจิ ‘ยทนิจฺจํ ตํ ทุกฺขํ’ วุตฺตํ ภควตา อญฺญาตาวินฺทฺริยํ อนิจฺจํ, โน จ วต เร วตฺตพฺเพ “ทุกฺขญฺเญว ปรินิปฺผนฺนนฺ”ติ.\csromanspacer{} อปรินิปฺผนฺนกถา นิฏฺฐิตา.\csromanspacer{} เตวีสติโม วคฺโค.}

\titlehead{\textbf{ตสฺสุทฺทานํ}}
\prose{เอกาธิปฺปาเยน เมถุโน ธมฺโม ปฏิเสวิตพฺโพ, อรหนฺตานํ วณฺเณน อมนุสฺสา เมถุนํ ธมฺมํ ปฏิเสวนฺติ, โพธิสตฺโต อิสฺสริยกามการิกาเหตุ วินิปาตํ คจฺฉติ, คพฺภเสยฺยํ โอกฺกมติ, ทุกฺกรการิกํ อกาสิ, อปรนฺตปํ อกาสิ, อญฺญํ สตฺถารํ อุทฺทิสิ, อตฺถิ น ราโค ราคปติรูปโก, อตฺถิ น โทโส โทสปติรูปโก, อตฺถิ น โมโห โมหปติรูปโก, อตฺถิ น กิเลโส กิเลสปติรูปโก, รูปํ อปรินิปฺผนฺนํ อญฺญาตาวินฺทฺริยํ อปรินิปฺผนฺนนฺติ.\csromanspacer{} ขุทฺทโก อฑฺฒปณฺณาสโก.}

\titlehead{\textbf{ตสฺสุทฺทานํ}}
\csromangathameasure{นวํ, นิพฺพุติ, เอกาธิปฺปาโยติ.}
\csromangathagroup{\csromangathastanza{นวํ, นิพฺพุติ, เอกาธิปฺปาโยติ.}}

\csromanheader{2}{\textbf{ปณฺณาสกุทฺทานํ}}
\csromangathameasure{มหานิยาโม, อนุสยา, นิคฺคโห ขุทฺทกปญฺจโม. \\ ปรปฺปวาทมทฺทนา, สุตฺตมูลสมาหิตา. \\ อุชฺโชตนา สตฺถุสมเย, กถาวตฺถุปกรเณติ.}
\csromangathagroup{\csromangathastanza{มหานิยาโม, อนุสยา, นิคฺคโห ขุทฺทกปญฺจโม. \\ ปรปฺปวาทมทฺทนา, สุตฺตมูลสมาหิตา. \\ อุชฺโชตนา สตฺถุสมเย, กถาวตฺถุปกรเณติ.}}

\csromanheader{2}{\textbf{ปญฺจตฺติํสภาณวารํ}}
\nitthitam{\textbf{กถาวตฺถุปกรณํ นิฏฺฐิตํ.}}
\orphannote{ที 1. 79 ปิฏฺฐาทีสุ.}

\orphannote{สกวาทิวจนํ (อฏฺฐกถา ปสฺสิตพฺพา)}

\orphannote{เวหาสํ คจฺฉตีติ (?) อยํ หิ ปรวาทิสฺส ปฏิญฺญาเยว, 331 ปิฏฺเฐ นวมปนฺติยํ วิย.}

\orphannote{อิทํ ปุจฺฉาทฺวเยน ปุริมปุจฺฉาทฺวยสฺส ปุรโต ภวิตพฺพํ.}

\orphannote{อํ 2. 403 ปิฏฺเฐ.}

